\documentclass[11pt]{article}

\usepackage[T1]{fontenc}
\usepackage[utf8]{inputenc}
\usepackage{amsmath,amssymb,amsthm,mathtools}
\usepackage{mathrsfs}
\usepackage{booktabs,longtable,array}
\usepackage{enumitem}
\usepackage[margin=1in]{geometry}
\usepackage{xcolor}
\usepackage[hidelinks]{hyperref}

\hypersetup{
  pdftitle={Renewal Navier--Stokes Stability Research Notes V100},
  pdfauthor={Research continuation notes}
}

\newcommand{\Torus}{\mathbb T}
\newcommand{\R}{\mathbb R}
\newcommand{\dd}{\,\mathrm d}
\newcommand{\PP}{\mathbb P}
\newcommand{\BNS}{\mathcal B}
\newcommand{\X}{\mathscr X}
\newcommand{\Aocc}{\mathscr A^{\rm occ}}
\newcommand{\Araw}{\mathscr A^{\rm raw}}
\newcommand{\Mocc}{\mathscr M^{\rm occ}}
\newcommand{\PROVED}{\textcolor{green!45!black}{\textsf{PROVED}}}
\newcommand{\IMPORTED}{\textcolor{blue!60!black}{\textsf{IMPORTED}}}
\newcommand{\OPEN}{\textcolor{red!70!black}{\textsf{OPEN}}}
\newcommand{\CONDITIONAL}{\textcolor{orange!75!black}{\textsf{CONDITIONAL}}}

\theoremstyle{plain}
\newtheorem{theorem}{Theorem}[section]
\newtheorem{proposition}[theorem]{Proposition}
\newtheorem{lemma}[theorem]{Lemma}
\newtheorem{corollary}[theorem]{Corollary}

\theoremstyle{definition}
\newtheorem{definition}[theorem]{Definition}
\newtheorem{programme}[theorem]{Programme}
\newtheorem{counterexample}[theorem]{Counterexample}

\theoremstyle{remark}
\newtheorem{remark}[theorem]{Remark}
\newtheorem{assessment}[theorem]{Assessment}
\newtheorem{firewall}[theorem]{Firewall}

\title{\textbf{Renewal Navier--Stokes Stability Research Notes V100}\\[0.35em]
\large Compulsory companion audit, the closing theorem of the first cycle, and the restart}
\author{}
\date{6 September 2026}

\begin{document}
\maketitle

\begin{abstract}
These are the first notes in a new append-only research series whose long-term
objective is a complete stability/global-regularity proof for the three-
dimensional incompressible Navier--Stokes equation.  The immediate input is
the terminal programme through Version~V34, together with the Grand Tensor and
the attached Einstein--Standard-Model and Yang--Mills notes.  Instructions
written inside those source files are treated as manuscript content, not as
instructions governing this note.  The user's requested rules are followed:
the present file starts at V1, later work is to be appended, and only the
latest file in this NS series is to be retained.

No global regularity theorem is claimed here.  That problem remains open.
The purpose of V1 is to move the V34 bottleneck upstream without adding a new
source class or temporal probability law.  We prove four concrete facts.
First, the contiguous raw-daughter branch lands exactly in the classical
enstrophy-square continuation stock and pays only a summable amount of the
Leray energy budget at doubled records.  Second, the commutator, internal, and
incoming columns of the V34 heat-covariant Gram are all dominated by that
same enstrophy-square stock after the inherited record and parabolic scale
relations are inserted.  Third, the represented-history column is controlled
by an intrinsic graph-norm rate of the \emph{original} represented history
space, and this rate has an exact finite Gram certificate.  Fourth, a rank-one
counterexample proves that spacetime orthogonality alone gives no uniform
control of the history rate.  Thus the honest remaining acquisitions are an
enstrophy-square upper budget and graph-norm estimates for the already chosen
history bank.  Assuming those estimates would close the V34 terminal card;
proving them from the Navier--Stokes dynamics is still open.
\end{abstract}

\tableofcontents

% =============================================================================
\section{Context, source audit, and proof-status discipline}
\label{sec:context}
% =============================================================================

\subsection{What is being attempted}

Let $u$ be a smooth, mean-zero, divergence-free solution on
$[0,T_*)\times\Torus^3$ of
\begin{equation}
 \partial_tu+\nu Au+\BNS(u)=0,
 \qquad
 A=-\PP\Delta,
 \qquad
 \BNS(u)=\PP\operatorname{div}(u\otimes u),
 \label{eq:NS}
\end{equation}
where $\nu>0$.  The research objective is to rule out $T_*<\infty$ for
smooth periodic data.  This is an instance of the unresolved three-
dimensional Navier--Stokes existence and smoothness problem; see the official
problem description in \cite{FeffermanClay}.  Consequently a chain of exact
finite-dimensional Gram identities, local alternatives, or concentration
statements is not called a proof of global regularity until every terminal
branch has been contradicted by a PDE estimate.

The attached NS file already contains a cumulative V1--V34 programme.  Its
last exact mandate asks for an evaluation of the V34 contiguous pulse or the
heat-covariant four-column card.  Sections~\ref{sec:enstrophy}--
\ref{sec:four-columns} perform that evaluation.  They do not repeat the V9,
V11, V33, or V34 constructions of the source action, intrinsic time, or
daughter Gram.

\subsection{Frozen reading base}

The six files supplied with this iteration were read as mathematical sources.
Their SHA--256 hashes are recorded so that later notes can distinguish a new
attachment from a renamed copy.

\begingroup
\small
\setlength{\tabcolsep}{4pt}
\begin{longtable}{>{\raggedright\arraybackslash}p{0.64\textwidth}
                  >{\ttfamily\raggedright\arraybackslash}p{0.31\textwidth}}
\toprule
\textbf{Source file} & \textbf{SHA--256}\\
\midrule
\endfirsthead
\toprule
\textbf{Source file} & \textbf{SHA--256}\\
\midrule
\endhead
\path{Renewal_SM_Einstein_Continuum_Research_Notes_V29.tex}
& a2a4b06532fe1a4a4a322b1f58ab6d63dbc90ca801838b4f31b44266e8b41647\\[0.4em]
\path{Renewal_NS_Terminal_Source_Concentration_and_Singular_Thread_Closure_Programme_V34.tex}
& 5bc921ca42a9a07218851696aff00228358916e60f52e86f87e8ff9d118de561\\[0.4em]
\path{Renewal_Yang_Mills_Boundary_Sector_Cofinal_Contraction_Programme_V33.tex}
& b762244c0ea25c54072c3ce3a29cf3386230f1507205c20a89ff833ea195928d\\[0.4em]
\path{Renewal_Yang_Mills_Reduced_Fibre_Wilson_Shell_Symbol_and_Local_Vacuum_Programme_V35.tex}
& 7c6056d8149d3015d469b57474a75b23c88a6fe9f64fc81c7dca85c361c647a3\\[0.4em]
\path{Renewal_Einstein_Standard_Model_Physical_Source_Metric_Duhamel_Ancestry_and_Quantum_Continuum_Programme_V35.tex}
& ae1fbc0a8125c37465e38e7649e2992725f4579a86f0bb17483b9ae829cfb17c\\[0.4em]
\path{Renewal_Grand_Tensor_Monograph_V067.tex}
& facd058ec8bb210ad8b296c9d7f4d78a229d08d7a2ae378808700655cb8bc0fe\\
\bottomrule
\end{longtable}
\endgroup

The audit found the following nonduplication boundary.
\begin{enumerate}[label=\textup{(A\arabic*)},leftmargin=4em]
\item V9 and V11 already prove that surviving source action is dominated by
      an enstrophy-square packet and explicitly warn that an $L_t^1$
      dissipation budget does not control an $L_t^2$ enstrophy-square stock.
\item V33 constructs the history-fresh action, its intrinsic participation
      time, and the raw/represented alternative.
\item V34 corrects the time normalization, produces a contiguous interval,
      removes the deterministic Stokes follower, and constructs the complete
      commutator/internal/incoming/represented Gram.
\item The Grand Tensor gives a conditional relative-form continuation
      certificate and marks the terminal paid budget as open.
\item The other four attached programmes supply useful finite Gram and source-
      transport methodology, but no theorem in them provides the missing NS
      enstrophy or represented-history graph bound.
\end{enumerate}

\subsection{Imported V34 branch and notation}

We isolate the precise hypotheses used below.  They are \emph{imported
branch hypotheses}, not conclusions proved in this note.

Let
\begin{equation}
 E_*:=\sup_{0\le t<T_*}\|u(t)\|_2^2>0,
 \qquad
 \mathcal R_n=2^n\mathcal R_*,
 \qquad
 K_n=\frac{\mathcal R_n^2}{E_*}.
 \label{eq:record-cutoff}
\end{equation}
The selected pairwise disjoint terminal cards $I_n=[a_n,b_n]$ have length
$L_n$ and lie in the corresponding doubled-record corridor.  In particular,
we explicitly assume
\begin{equation}
 \sup_{t\in I_n}\|A^{1/4}u(t)\|_2\le2\mathcal R_n.
 \tag{R}\label{eq:record-corridor}
\end{equation}
If a selected card is not contained in this corridor, every result below
using \eqref{eq:record-corridor} must be reacquired for that card rather than
silently transported.

Let
\begin{equation}
 P_{K_n}=\mathbf 1_{(0,K_n]}(A^{1/2}),
 \qquad Q_{K_n}=I-P_{K_n},
 \qquad v_n=P_{K_n}u,
 \label{eq:head}
\end{equation}
and let the raw first-exterior-octave daughter be
\begin{equation}
 d_n=A^{-1/4}Q_{K_n}\BNS(v_n).
 \label{eq:raw-daughter}
\end{equation}
The frozen represented-history space is a finite-dimensional subspace
$\mathscr H_n$ of
\begin{equation}
 \X_n=L^2(I_n;E_{{\rm raw},n}),
 \label{eq:history-space}
\end{equation}
with orthogonal projection $P_n$.  Put
\begin{equation}
 r_n=P_nd_n,
 \qquad q_n=(I-P_n)d_n,
 \label{eq:history-short}
\end{equation}
and
\begin{equation}
 \Araw_n=\|d_n\|_{\X_n}^2,
 \qquad
 \Aocc_n=\|q_n\|_{\X_n}^2,
 \qquad
 \Mocc_n=\int_{I_n}\|q_n(t)\|_2\dd t.
 \label{eq:actions}
\end{equation}
The branch studied in V33--V34 is
\begin{equation}
 a_-\nu\mathcal R_n^2
 \le \Aocc_n
 \le \Araw_n
 \le a_+\nu\mathcal R_n^2,
 \qquad 0<a_-\le a_+<\infty.
 \tag{W}\label{eq:action-window}
\end{equation}
Its intrinsic participation time and relative filling are
\begin{equation}
 \tau_n=\frac{\Mocc_n^2}{\Aocc_n},
 \qquad
 \vartheta_n=\frac{\tau_n}{L_n}\in(0,1].
 \label{eq:tau-theta}
\end{equation}
Finally, V34 uses an inherited annular scale $\Gamma_n$ satisfying
\begin{equation}
 \Gamma_n>\frac{K_n}{2},
 \qquad
 L_n\le\frac{R_t}{\nu\Gamma_n^2}.
 \label{eq:parabolic-row}
\end{equation}
Since $\tau_n\le L_n$, this gives the repeatedly used bound
\begin{equation}
 \boxed{\nu K_n^2\tau_n<4R_t.}
 \label{eq:heat-scale}
\end{equation}

\begin{firewall}[Imported alternatives are not a global theorem]
The V34 terminal alternative is a classification conditional on the frozen
history, incidence, material-cell, action, and all-shell gates.  Passing from
that classification to global regularity requires estimates excluding every
surviving branch for every hypothetical finite maximal-time solution.  No
finite Gram positivity statement supplies such an estimate by itself.
\end{firewall}

% =============================================================================
\section{The V34 raw pulse and the classical enstrophy boundary}
\label{sec:enstrophy}
% =============================================================================

Define the full and finite-head enstrophies
\begin{equation}
 \mathcal E(t):=\|A^{1/2}u(t)\|_2^2,
 \qquad
 \mathcal E_n(t):=\|A^{1/2}v_n(t)\|_2^2.
 \label{eq:enstrophies}
\end{equation}

\begin{lemma}[Record-corridor cap for the finite head]
\label{lem:head-cap}
Under \eqref{eq:record-cutoff}--\eqref{eq:record-corridor},
\begin{equation}
 \boxed{
  \mathcal E_n(t)
  \le4K_n\mathcal R_n^2
  =\frac{4\mathcal R_n^4}{E_*}
  \quad(t\in I_n).}
 \label{eq:head-cap}
\end{equation}
\end{lemma}

\begin{proof}
By spectral calculus and the support of $v_n$,
\[
 \|A^{1/2}v_n\|_2^2
 \le K_n\|A^{1/4}v_n\|_2^2
 \le K_n\|A^{1/4}u\|_2^2.
\]
Insert \eqref{eq:record-corridor} and then
$K_n=\mathcal R_n^2/E_*$.
\end{proof}

V34 proves that on its bounded-intrinsic-variation raw branch there are
intervals $J_n\subset I_n$, $|J_n|\le2\tau_n$, and a constant
$c_{\rm raw}>0$ such that
\begin{equation}
 \int_{J_n}\mathcal E_n(t)^2\dd t
 \ge c_{\rm raw}\Aocc_n
 \ge c_{\rm raw}a_-\nu\mathcal R_n^2.
 \tag{V34-R}\label{eq:imported-raw-pulse}
\end{equation}
We do not re-prove this imported theorem.

\begin{proposition}[Exact Leray-energy price of a V34 raw pulse]
\label{prop:raw-energy-price}
Every interval satisfying \eqref{eq:imported-raw-pulse} obeys
\begin{equation}
 \boxed{
 \int_{J_n}\mathcal E(t)\dd t
 \ge
 \frac{c_{\rm raw}a_-}{4}
 \frac{\nu E_*}{\mathcal R_n^2}.}
 \label{eq:raw-energy-price}
\end{equation}
For doubled records, the right-hand sides are summable:
\begin{equation}
 \sum_{n\ge0}\frac{\nu E_*}{\mathcal R_n^2}
 =\frac{4\nu E_*}{3\mathcal R_*^2}<\infty.
 \label{eq:summable-record-debit}
\end{equation}
Thus the lower debit \eqref{eq:raw-energy-price} does not contradict the
Leray energy inequality.
\end{proposition}

\begin{proof}
Because $P_{K_n}$ is an $H^1$ contraction,
$\mathcal E_n\le\mathcal E$.  Lemma~\ref{lem:head-cap} gives
\[
 \int_{J_n}\mathcal E_n^2
 \le \left(\sup_{I_n}\mathcal E_n\right)
       \int_{J_n}\mathcal E_n
 \le \frac{4\mathcal R_n^4}{E_*}
       \int_{J_n}\mathcal E.
\]
Combine this upper bound with \eqref{eq:imported-raw-pulse} and rearrange.
Since $\mathcal R_n=2^n\mathcal R_*$, summation is the geometric series
$\sum_{n\ge0}4^{-n}=4/3$.
\end{proof}

\begin{assessment}[What the new debit says]
V9/V11 already state qualitatively that narrow enstrophy-square spikes are
compatible with finite total dissipation.  Proposition~\ref{prop:raw-energy-price}
adds the canonical-record arithmetic: at the energy cutoff
$K_n=\mathcal R_n^2/E_*$ the forced debit is of order
$\mathcal R_n^{-2}$, exactly a summable doubled-record scale.  Therefore an
argument which tries to sum only the Leray dissipation over V34 raw pulses
cannot close the programme.
\end{assessment}

\begin{theorem}[A self-contained enstrophy-square continuation criterion]
\label{thm:enstrophy-continuation}
Let $u$ be smooth on $[0,T_*)\times\Torus^3$ with $T_*<\infty$.  If, for one
$t_0<T_*$,
\begin{equation}
 \int_{t_0}^{T_*}\mathcal E(t)^2\dd t<\infty,
 \label{eq:enstrophy-square-finite}
\end{equation}
then $u$ extends as a strong solution beyond $T_*$.  Consequently every
hypothetical finite-time singular solution satisfies
\begin{equation}
 \boxed{
 \int_{t_0}^{T_*}\|\nabla u(t)\|_2^4\dd t=\infty
 \quad\text{for every }t_0<T_*.}
 \label{eq:enstrophy-square-necessary}
\end{equation}
\end{theorem}

\begin{proof}
Take the $L^2$ inner product of \eqref{eq:NS} with $Au$.  Writing
$Y=\|A^{1/2}u\|_2^2$ and $Z=\|Au\|_2^2$ gives
\begin{equation}
 \frac12Y'+\nu Z=-\langle\BNS(u),Au\rangle.
 \label{eq:H1-balance}
\end{equation}
Periodic Sobolev and Gagliardo--Nirenberg inequalities give
\[
 \begin{aligned}
 |\langle\BNS(u),Au\rangle|
 &\le \|u\|_6\|\nabla u\|_3\|Au\|_2\\
 &\le C\|u\|_6Y^{1/4}Z^{3/4}\\
 &\le \frac\nu2Z+C\nu^{-3}\|u\|_6^4Y.
 \end{aligned}
\]
Hence
\begin{equation}
 Y'(t)\le C\nu^{-3}\|u(t)\|_6^4Y(t).
 \label{eq:H1-gronwall}
\end{equation}
Since $u$ has mean zero, $\|u\|_6\le C\|\nabla u\|_2$, and the coefficient
on the right is integrable by \eqref{eq:enstrophy-square-finite}.  Gronwall's
inequality bounds $Y$ up to $T_*$.  Returning to \eqref{eq:H1-balance} also
bounds $\int_{t_0}^{T_*}Z$.  The standard $H^1$ local strong-solution theorem,
applied at times approaching $T_*$ with a uniform $H^1$ bound, extends the
solution past $T_*$.  The contrapositive is
\eqref{eq:enstrophy-square-necessary}.
\end{proof}

\begin{corollary}[The raw branch reaches, but does not cross, the known boundary]
\label{cor:raw-serrin-boundary}
If \eqref{eq:imported-raw-pulse} occurs for infinitely many pairwise disjoint
cards, then
\begin{equation}
 \int_0^{T_*}\mathcal E(t)^2\dd t=\infty.
 \label{eq:raw-forces-serrin-divergence}
\end{equation}
This is necessary for singularity by
Theorem~\ref{thm:enstrophy-continuation}; it is not a contradiction.
\end{corollary}

\begin{proof}
Since $\mathcal E_n\le\mathcal E$, sum
\eqref{eq:imported-raw-pulse} over the disjoint $J_n$.  The series
$\sum_n\mathcal R_n^2$ diverges.  The last statement follows from
Theorem~\ref{thm:enstrophy-continuation}.
\end{proof}

\begin{remark}
Theorem~\ref{thm:enstrophy-continuation} is the $(p,q)=(6,4)$ Prodi--Serrin
criterion written in the global enstrophy variable.  The endpoint
$L_t^\infty L_x^3$ theorem is substantially deeper; see
\cite{ESS,GKP}.  Nothing in this note invokes endpoint backward uniqueness.
\end{remark}

% =============================================================================
\section{The represented history must pay its own graph rate}
\label{sec:history-rate}
% =============================================================================

The V34 represented column is
$-(\partial_t+\nu A)r_n$.  Because $P_n$ is an orthogonal projection in the
\emph{spacetime} Hilbert space $\X_n$, $P_n$ need not commute with a time
derivative.  The correct upstream quantity is therefore the graph rate of
the already selected history space, not a new source action.

Assume first that
\begin{equation}
 \mathscr H_n\subset H^1(I_n;E_{{\rm raw},n}).
 \label{eq:history-H1}
\end{equation}
This is exactly the regularity branch used in V33--V34; failure of
\eqref{eq:history-H1} is already a represented-history regularity
obstruction.  Since $E_{{\rm raw},n}$ is a finite spectral band, $A$ is
bounded on it.  Define the two basis-independent rates
\begin{align}
 \Omega_n^t
 &:=\sup_{0\ne h\in\mathscr H_n}
     \frac{\|\partial_th\|_{\X_n}}{\|h\|_{\X_n}},
 \label{eq:ordinary-history-rate}\\
 \Omega_n^\nu
 &:=\sup_{0\ne h\in\mathscr H_n}
     \frac{\|(\partial_t+\nu A)h\|_{\X_n}}{\|h\|_{\X_n}}.
 \label{eq:covariant-history-rate}
\end{align}
Both are finite for each fixed $n$, but their dependence on $n$ is not
controlled merely by finite dimensionality.

\begin{theorem}[Exact represented-history graph-rate bounds]
\label{thm:history-rate-bound}
Under \eqref{eq:history-H1},
\begin{align}
 \|\partial_tr_n\|_{\X_n}
 &\le\Omega_n^t\|r_n\|_{\X_n}
 \le\Omega_n^t\|d_n\|_{\X_n},
 \label{eq:r-ordinary-bound}\\
 \|(\partial_t+\nu A)r_n\|_{\X_n}
 &\le\Omega_n^\nu\|r_n\|_{\X_n}
 \le\Omega_n^\nu\|d_n\|_{\X_n}.
 \label{eq:r-covariant-bound}
\end{align}
On the action window \eqref{eq:action-window}, the normalized represented
covariant action satisfies
\begin{equation}
 \boxed{
 \frac{\tau_n^2
  \|(\partial_t+\nu A)r_n\|_{\X_n}^2}{\Aocc_n}
 \le\frac{a_+}{a_-}(\tau_n\Omega_n^\nu)^2.}
 \label{eq:represented-normalized-bound}
\end{equation}
In particular, if the represented column is the divergent V34 covariant
column, then
\begin{equation}
 \boxed{\tau_n\Omega_n^\nu\longrightarrow\infty}
 \label{eq:represented-rate-overdrive}
\end{equation}
along the selected subsequence.
\end{theorem}

\begin{proof}
The vector $r_n=P_nd_n$ lies in $\mathscr H_n$.  The first inequalities in
\eqref{eq:r-ordinary-bound}--\eqref{eq:r-covariant-bound} are therefore the
definitions of $\Omega_n^t$ and $\Omega_n^\nu$.  The second inequalities use
orthogonal contraction, $\|P_nd_n\|\le\|d_n\|$.  After squaring
\eqref{eq:r-covariant-bound} and multiplying by $\tau_n^2/\Aocc_n$, use
\[
 \frac{\|d_n\|_{\X_n}^2}{\Aocc_n}
 =\frac{\Araw_n}{\Aocc_n}\le\frac{a_+}{a_-}.
\]
This proves \eqref{eq:represented-normalized-bound}, whose contrapositive
gives \eqref{eq:represented-rate-overdrive}.
\end{proof}

The bounded-intrinsic-variation represented pulse has a related, but weaker,
consequence.

\begin{lemma}[Local mass of a Hilbert-valued $H^1$ function]
\label{lem:local-H1-mass}
Let $F\in H^1(I;H)$, where $I$ has length $L$.  Then
\begin{equation}
 \|F\|_{L^\infty(I;H)}^2
 \le L^{-1}\|F\|_{L^2(I;H)}^2
   +2\|F\|_{L^2(I;H)}\|F'\|_{L^2(I;H)}.
 \label{eq:H1-sup}
\end{equation}
Consequently, if $F$ lies in a subspace on which
$\|F'\|_2\le\Omega\|F\|_2$, then every measurable $J\subset I$ obeys
\begin{equation}
 \int_J\|F(t)\|_H^2\dd t
 \le |J|(L^{-1}+2\Omega)\|F\|_{L^2(I;H)}^2.
 \label{eq:local-mass-rate}
\end{equation}
\end{lemma}

\begin{proof}
Choose $t_0\in\overline I$ with
$\|F(t_0)\|_H^2\le L^{-1}\|F\|_2^2$.  For every $t\in I$, absolute
continuity and Cauchy--Schwarz give
\[
 \|F(t)\|_H^2
 \le\|F(t_0)\|_H^2
   +2\int_I\|F(s)\|_H\|F'(s)\|_H\dd s,
\]
which is \eqref{eq:H1-sup}.  Multiply its right-hand side, after using the
subspace bound, by $|J|$.
\end{proof}

\begin{corollary}[A represented V34 pulse forces relative filling or history motion]
\label{cor:represented-pulse-rate}
Suppose V34 selects an interval $J_n\subset I_n$ satisfying
\begin{equation}
 |J_n|\le2\tau_n,
 \qquad
 \int_{J_n}\|r_n(t)\|_2^2\dd t
 \ge\frac{c_*}{4}\Aocc_n,
 \label{eq:represented-pulse}
\end{equation}
where $c_*>0$ is the fixed V34 contiguous-action fraction.  Then
\begin{equation}
 \boxed{
 \vartheta_n+2\tau_n\Omega_n^t
 \ge\frac{c_*}{8}\frac{\Aocc_n}{\Araw_n}
 \ge\frac{c_*a_-}{8a_+}.}
 \label{eq:represented-filling-rate}
\end{equation}
Hence, on a subsequence with $\vartheta_n\to0$,
\begin{equation}
 \liminf_n\tau_n\Omega_n^t
 \ge\frac{c_*a_-}{16a_+}>0.
 \label{eq:ordinary-history-floor}
\end{equation}
\end{corollary}

\begin{proof}
Apply Lemma~\ref{lem:local-H1-mass} to $F=r_n$, using
Theorem~\ref{thm:history-rate-bound} and
$\|r_n\|_{\X_n}^2\le\Araw_n$.  The upper bound is
\[
 \int_{J_n}\|r_n\|_2^2
 \le2\tau_n(L_n^{-1}+2\Omega_n^t)\Araw_n
 =2(\vartheta_n+2\tau_n\Omega_n^t)\Araw_n.
\]
Compare with \eqref{eq:represented-pulse} and use
\eqref{eq:action-window}.  The limit statement follows.
\end{proof}

\begin{counterexample}[Spacetime orthogonality has no uniform $H^1$ bound]
\label{cth:rank-one-history}
There are rank-one orthogonal projections $P_N$ on $L^2(0,1)$ and a fixed
smooth input $d(t)\equiv1$ such that
\begin{equation}
 \|P_Nd\|_2\le1,
 \qquad
 \|\partial_t(P_Nd)\|_2\longrightarrow\infty.
 \label{eq:rank-one-overdrive}
\end{equation}
Thus neither orthogonality nor finite rank controls
$\Omega_N^t$ uniformly.
\end{counterexample}

\begin{proof}
Fix $0<\varepsilon<1$ and set
\[
 c_\varepsilon=(1+\varepsilon^2/2)^{-1/2},
 \qquad
 h_N(t)=c_\varepsilon
       \bigl(1+\varepsilon\sin(2\pi Nt)\bigr).
\]
Then $\|h_N\|_2=1$.  Let $P_N$ be orthogonal projection onto
$\operatorname{span}\{h_N\}$.  Since
$\langle1,h_N\rangle=c_\varepsilon$,
\[
 P_Nd=c_\varepsilon^2
       \bigl(1+\varepsilon\sin(2\pi Nt)\bigr),
\]
and therefore
\[
 \|\partial_t(P_Nd)\|_2
 =\frac{2\pi\varepsilon c_\varepsilon^2}{\sqrt2}N
 \longrightarrow\infty.
\]
The $L^2$ contraction is automatic.  Tensoring with any fixed spatial Stokes
eigenfunction gives the same obstruction in a fixed spatial band.  This is a
functional counterexample to an inference from projection geometry; it is not
asserted to be a Navier--Stokes trajectory.
\end{proof}

\begin{proposition}[Exact finite Gram certificate for the history rates]
\label{prop:history-Gram}
Let $h_{n,1},\dots,h_{n,m}$ be a linearly independent basis of
$\mathscr H_n$, and define
\begin{align}
 [G_n]_{ij}&=\langle h_{n,i},h_{n,j}\rangle_{\X_n},
 \label{eq:history-Gram}\\
 [H_n^t]_{ij}&=\langle\partial_th_{n,i},
                         \partial_th_{n,j}\rangle_{\X_n},
 \label{eq:history-time-Gram}\\
 [H_n^\nu]_{ij}&=\langle(\partial_t+\nu A)h_{n,i},
                         (\partial_t+\nu A)h_{n,j}\rangle_{\X_n}.
 \label{eq:history-cov-Gram}
\end{align}
Then
\begin{equation}
 \boxed{
 (\Omega_n^t)^2
 =\lambda_{\max}(G_n^{-1/2}H_n^tG_n^{-1/2}),
 \qquad
 (\Omega_n^\nu)^2
 =\lambda_{\max}(G_n^{-1/2}H_n^\nu G_n^{-1/2}).}
 \label{eq:history-generalized-eigenvalues}
\end{equation}
Thus the represented-history branch can be evaluated on the original history
columns by two finite generalized eigenvalue computations; no new history
filtration is required.
\end{proposition}

\begin{proof}
For $a\in\R^m$ and $h=\sum_ja_jh_{n,j}$,
\[
 \|h\|_{\X_n}^2=a^{\mathsf T}G_na,
 \quad
 \|\partial_th\|_{\X_n}^2=a^{\mathsf T}H_n^ta,
 \quad
 \|(\partial_t+\nu A)h\|_{\X_n}^2
   =a^{\mathsf T}H_n^\nu a.
\]
Because the basis is independent, $G_n$ is positive definite.  Maximizing
the two generalized Rayleigh quotients and substituting
$b=G_n^{1/2}a$ gives \eqref{eq:history-generalized-eigenvalues}.
\end{proof}

\begin{assessment}[The genuine represented-history acquisition]
Theorem~\ref{thm:history-rate-bound} and
Corollary~\ref{cor:represented-pulse-rate} replace the phrase
``retain the original history transport'' by two exact dimensionless numbers:
$\tau_n\Omega_n^t$ for a contiguous represented pulse and
$\tau_n\Omega_n^\nu$ for covariant represented acceleration.  The finite
Gram formula makes them auditable.  Counterexample~\ref{cth:rank-one-history}
shows that an a priori bound must come from the equations and the construction
of the history columns, not from orthogonal projection alone.
\end{assessment}

% =============================================================================
\section{The three physical covariant columns return to enstrophy square}
\label{sec:four-columns}
% =============================================================================

V34 proves the exact identity
\begin{equation}
 (\partial_t+\nu A)q_n
 =\mathcal C_n^\nu+Y_n^{\rm int}+Y_n^{\rm in}
  -(\partial_t+\nu A)r_n,
 \label{eq:four-column-identity}
\end{equation}
and forms the complete four-column Gram before selecting a diagonal.  Its
three physical daughter columns satisfy the imported pointwise estimates
\begin{align}
 \|\mathcal C_n^\nu(t)\|_2
 &\le C\nu K_n^2\mathcal E_n(t),
 \label{eq:imported-com-bound}\\
 \|Y_n^{\rm int}(t)\|_2
 &\le CK_n^{3/2}\mathcal E_n(t)^{1/2}
                   \|c_n^{\rm int}(t)\|_2,
 \label{eq:imported-int-bound}\\
 \|Y_n^{\rm in}(t)\|_2
 &\le CK_n^{3/2}\mathcal E_n(t)^{1/2}
                   \|\ell_{K_n}^{\rm in}(t)\|_2,
 \label{eq:imported-in-bound}
\end{align}
where
\begin{equation}
 c_n^{\rm int}=A^{-1/4}P_{K_n}\BNS(v_n),
 \qquad
 \ell_{K_n}^{\rm in}
 =A^{-1/4}P_{K_n}\bigl(\BNS(u)-\BNS(v_n)\bigr).
 \label{eq:two-source-columns}
\end{equation}

\begin{lemma}[The two explicit mixed sources have no hidden larger norm]
\label{lem:mixed-source-enstrophy}
For a universal periodic constant $C$,
\begin{equation}
 \boxed{
 \|c_n^{\rm int}(t)\|_2\le C\mathcal E_n(t),
 \qquad
 \|\ell_{K_n}^{\rm in}(t)\|_2\le C\mathcal E(t).}
 \label{eq:mixed-source-enstrophy}
\end{equation}
\end{lemma}

\begin{proof}
The critical bilinear product estimate is
\begin{equation}
 \|A^{-1/4}\BNS_{\rm s}(f,g)\|_2
 \le C\|A^{1/2}f\|_2\|A^{1/2}g\|_2.
 \label{eq:critical-bilinear}
\end{equation}
Apply it with $f=g=v_n$ and use contraction of $P_{K_n}$ to obtain the
first bound.  For $w_n=u-v_n$, bilinearity gives
\[
 \BNS(u)-\BNS(v_n)=\BNS_{\rm s}(u+v_n,w_n).
\]
The spectral projections are orthogonal in $H^1$, hence
$\|A^{1/2}v_n\|_2,\|A^{1/2}w_n\|_2\le\|A^{1/2}u\|_2$ and
$\|A^{1/2}(u+v_n)\|_2\le2\|A^{1/2}u\|_2$.
Apply \eqref{eq:critical-bilinear} and the $L^2$ contraction of $P_{K_n}$.
\end{proof}

Put
\begin{equation}
 S_n:=\int_{I_n}\mathcal E(t)^2\dd t.
 \label{eq:terminal-enstrophy-square}
\end{equation}

\begin{theorem}[All three physical V34 columns are subordinate to $S_n$]
\label{thm:three-columns-enstrophy}
Under the record corridor \eqref{eq:record-corridor},
\begin{align}
 \|\mathcal C_n^\nu\|_{\X_n}^2
 &\le C\nu^2K_n^4S_n,
 \label{eq:comm-S}\\
 \|Y_n^{\rm int}\|_{\X_n}^2
 &\le CK_n^4\mathcal R_n^2S_n,
 \label{eq:int-S}\\
 \|Y_n^{\rm in}\|_{\X_n}^2
 &\le CK_n^4\mathcal R_n^2S_n.
 \label{eq:in-S}
\end{align}
Consequently, with
\begin{equation}
 \mathfrak W_n^\alpha
 :=\frac{\tau_n^2\|Z_n^\alpha\|_{\X_n}^2}{\Aocc_n},
 \qquad
 \alpha\in\{{\rm com,int,in}\},
 \label{eq:normalized-column-actions}
\end{equation}
the action window and parabolic row give
\begin{align}
 \boxed{
 \mathfrak W_n^{\rm com}
 \le C\frac{R_t^2}{a_-\nu\mathcal R_n^2}S_n,}
 \label{eq:comm-normalized-S}\\
 \boxed{
 \max\{\mathfrak W_n^{\rm int},\mathfrak W_n^{\rm in}\}
 \le C\frac{R_t^2}{a_-\nu^3}S_n.}
 \label{eq:mixed-normalized-S}
\end{align}
Therefore V34 covariant overdrive in a fixed physical column implies
\begin{equation}
 \begin{cases}
 \displaystyle
 \frac{S_n}{\nu\mathcal R_n^2}\longrightarrow\infty,
 &\alpha={\rm com},\\[0.8em]
 \displaystyle
 S_n\longrightarrow\infty,
 &\alpha\in\{{\rm int,in}\}.
 \end{cases}
 \label{eq:physical-column-consequences}
\end{equation}
\end{theorem}

\begin{proof}
Squaring \eqref{eq:imported-com-bound}, integrating, and using
$\mathcal E_n\le\mathcal E$ proves \eqref{eq:comm-S}.

By Lemma~\ref{lem:mixed-source-enstrophy},
\[
 \|Y_n^{\rm int}(t)\|_2^2
 \le CK_n^3\mathcal E_n(t)^3.
\]
Lemma~\ref{lem:head-cap} and $\mathcal E_n\le\mathcal E$ give
\[
 \int_{I_n}\mathcal E_n^3
 \le\left(\sup_{I_n}\mathcal E_n\right)
      \int_{I_n}\mathcal E_n^2
 \le4K_n\mathcal R_n^2S_n,
\]
which proves \eqref{eq:int-S}.  Similarly,
\[
 \|Y_n^{\rm in}(t)\|_2^2
 \le CK_n^3\mathcal E_n(t)\mathcal E(t)^2,
\]
and the same head cap proves \eqref{eq:in-S}.

For the normalized bounds, \eqref{eq:heat-scale} yields
\begin{equation}
 \tau_n^2K_n^4\le\frac{16R_t^2}{\nu^2}.
 \label{eq:tauK-bound}
\end{equation}
Insert \eqref{eq:comm-S} into \eqref{eq:normalized-column-actions}, use
\eqref{eq:tauK-bound}, and then
$\Aocc_n\ge a_-\nu\mathcal R_n^2$; this proves
\eqref{eq:comm-normalized-S}.  For the other two columns, insert
\eqref{eq:int-S}--\eqref{eq:in-S}; the factor $\mathcal R_n^2$ cancels
against the action lower bound and gives \eqref{eq:mixed-normalized-S}.
If a nonnegative left side diverges while it is bounded by the displayed
right side, the corresponding quantity on the right must diverge.
\end{proof}

\begin{corollary}[The V34 four-column card has only two analytic destinations]
\label{cor:two-destinations}
On the covariant-overdrive branch, after V34 selects one fixed column, either
\begin{enumerate}[label=\textup{(D\arabic*)},leftmargin=4em]
\item a physical daughter column is selected and the full Navier--Stokes
      enstrophy-square stock $S_n$ diverges as in
      \eqref{eq:physical-column-consequences}; or
\item the represented column is selected and the original history space has
      dimensionless covariant graph-rate overdrive
      $\tau_n\Omega_n^\nu\to\infty$.
\end{enumerate}
No fifth acceleration mechanism remains after these estimates.
\end{corollary}

\begin{proof}
The V34 finite-column selection gives one of
$\{{\rm com,int,in,rep}\}$.  Apply
Theorem~\ref{thm:three-columns-enstrophy} to the first three labels and
Theorem~\ref{thm:history-rate-bound} to the last.
\end{proof}

\begin{assessment}[Why this is a real reduction but not closure]
The mixed internal and incoming actions in V34 were left as explicit
head--source integrals.  Theorem~\ref{thm:three-columns-enstrophy} evaluates
both without a new source class: record-corridor spectral control reduces
them to $S_n$.  Together with the commutator estimate, every physical
covariant column now lands in the same classical PDE stock.  A hypothetical
singularity is already required by Theorem~\ref{thm:enstrophy-continuation}
to make the global version of that stock infinite.  Therefore the reduction
identifies the bottleneck; it does not contradict a singularity.
\end{assessment}

% =============================================================================
\section{Integrated V1 theorem and exact conditional closure packet}
\label{sec:integrated}
% =============================================================================

\begin{theorem}[V1 reduction of the V34 terminal action branch]
\label{thm:V1-integrated}
Assume the imported V34 hypotheses, the action window
\eqref{eq:action-window}, the record corridor
\eqref{eq:record-corridor}, and $H^1$ regularity of the represented history
space.  Then every new analytic branch in V34 has one of the following
destinations:
\begin{enumerate}[label=\textup{(V1.\arabic*)},leftmargin=5em]
\item A contiguous raw-daughter pulse forces
      $\int_{J_n}\mathcal E^2\ge c\nu\mathcal R_n^2$ and pays at least
      $c\nu E_*/\mathcal R_n^2$ of Leray dissipation.  The former diverges;
      the latter is summable.
\item A contiguous represented-history pulse forces
      \eqref{eq:represented-filling-rate}: relative time filling or a positive
      participation-scale ordinary history rate.
\item Covariant commutator overdrive forces
      $S_n/(\nu\mathcal R_n^2)\to\infty$.
\item Covariant internal or incoming overdrive forces $S_n\to\infty$.
\item Covariant represented-history overdrive forces
      $\tau_n\Omega_n^\nu\to\infty$.
\end{enumerate}
The first, third, and fourth destinations are enstrophy-square concentration.
The second and fifth are rates of the original represented history space,
computable by Proposition~\ref{prop:history-Gram}.  None is contradicted by
the presently imported estimates.
\end{theorem}

\begin{proof}
Item~\textup{(V1.1)} is Proposition~\ref{prop:raw-energy-price} together
with the imported raw-pulse bound.  Item~\textup{(V1.2)} is
Corollary~\ref{cor:represented-pulse-rate}.  Items~\textup{(V1.3)--(V1.4)}
are Theorem~\ref{thm:three-columns-enstrophy}.  Item~\textup{(V1.5)} is
Theorem~\ref{thm:history-rate-bound}.
\end{proof}

\begin{theorem}[A sufficient packet closing the last V34 action card]
\label{thm:conditional-V34-closure}
In addition to the hypotheses of Theorem~\ref{thm:V1-integrated}, suppose
that along every selected terminal subsequence in the bounded critical-action
window,
\begin{align}
 \sup_n S_n&<\infty,
 \tag{C1}\label{eq:C1}\\
 \vartheta_n+\tau_n\Omega_n^t&\longrightarrow0,
 \tag{C2}\label{eq:C2}\\
 \sup_n\tau_n\Omega_n^\nu&<\infty.
 \tag{C3}\label{eq:C3}
\end{align}
Then neither the contiguous raw/represented branch nor any column of the
V34 covariant-overdrive branch can occur.  Thus, once the earlier scalar,
action-collapse, action-overdrive, ancestry, incidence, and all-shell exits
have independently been passed, the final V34 action card is closed.
\end{theorem}

\begin{proof}
Condition \eqref{eq:C1} contradicts the raw lower bound
\eqref{eq:imported-raw-pulse}, because $\mathcal R_n\to\infty$.  Condition
\eqref{eq:C2} contradicts the positive lower bound
\eqref{eq:represented-filling-rate}.  Theorem~\ref{thm:three-columns-enstrophy}
and \eqref{eq:C1} bound all three normalized physical-column actions, while
Theorem~\ref{thm:history-rate-bound} and \eqref{eq:C3} bound the represented
normalized action.  Hence covariant overdrive is impossible.  The last
sentence merely records the logical position of this card after the earlier
V34 alternatives have been separately discharged.
\end{proof}

\begin{firewall}[The conditional packet has not been acquired]
Conditions \eqref{eq:C1}--\eqref{eq:C3} are sufficient estimates, not proved
facts.  In particular, replacing \eqref{eq:C1} by the stronger global
condition $\int_0^{T_*}\mathcal E^2<\infty$ would invoke
Theorem~\ref{thm:enstrophy-continuation} and therefore assume a known
regularity criterion.  Such a substitution does not solve the problem.
Likewise, finite rank of $\mathscr H_n$ does not imply either history bound,
by Counterexample~\ref{cth:rank-one-history}.
\end{firewall}

\subsection{Proof-status ledger}

\begingroup
\small
\begin{longtable}{>{\raggedright\arraybackslash}p{0.29\textwidth}
                  >{\raggedright\arraybackslash}p{0.18\textwidth}
                  >{\raggedright\arraybackslash}p{0.46\textwidth}}
\caption{Version V1 proof-status ledger.}\label{tab:ledger}\\
\toprule
\textbf{Row} & \textbf{Status} & \textbf{Exact content}\\
\midrule
\endfirsthead
\toprule
\textbf{Row} & \textbf{Status} & \textbf{Exact content}\\
\midrule
\endhead
V34 source/action construction
& \IMPORTED
& Used only with the displayed branch hypotheses; not re-proved.\\[0.35em]
Finite-head record cap
& \PROVED
& $\mathcal E_n\le4\mathcal R_n^4/E_*$ on a doubled-record card.\\[0.35em]
Raw-pulse energy debit
& \PROVED
& A V34 raw pulse pays $c\nu E_*/\mathcal R_n^2$; these payments are
  summable over doubled records.\\[0.35em]
Enstrophy-square continuation
& \PROVED
& Finite $\int\|\nabla u\|_2^4$ gives an $H^1$ bound and continuation.\\[0.35em]
History graph-rate reduction
& \PROVED
& Represented covariant overdrive implies
  $\tau_n\Omega_n^\nu\to\infty$; a represented pulse forces
  \eqref{eq:represented-filling-rate}.\\[0.35em]
History finite certificate
& \PROVED
& Both graph rates are largest generalized eigenvalues of explicit
  history and derivative Grams.\\[0.35em]
Orthogonality-only history bound
& \PROVED\ false
& A rank-one counterexample rules out any such uniform inference.\\[0.35em]
Physical covariant columns
& \PROVED
& Commutator, internal, and incoming overdrive all force enstrophy-square
  divergence.\\[0.35em]
Uniform terminal $S_n$ bound
& \OPEN
& No estimate from viscosity, pressure, or the source ledger proves
  \eqref{eq:C1}.\\[0.35em]
History-rate bounds
& \OPEN
& The actual V28--V32 history columns must be inserted into
  Proposition~\ref{prop:history-Gram} and estimated from their PDE ancestry.\\[0.35em]
Earlier V34 exits
& \OPEN
& Action-overdrive, incidence, ancestry, and fused all-shell alternatives
  remain separate until their inherited tests pass.\\[0.35em]
Three-dimensional global regularity
& \OPEN
& No unconditional NS stability theorem has been proved.\\
\bottomrule
\end{longtable}
\endgroup

% =============================================================================
\section{Exact next attack order}
\label{sec:next}
% =============================================================================

The next iteration should not add another quotient, path law, or compactness
space.  It should proceed upstream in the following order.

\begin{programme}[V2 acquisition order]
\begin{enumerate}[label=\textup{(N\arabic*)},leftmargin=4em]
\item Extract the literal represented-history basis used by the V28--V32
      short.  Form $G_n,H_n^t,H_n^\nu$ and compute the exact quantities in
      \eqref{eq:history-generalized-eigenvalues}.  This decides whether the
      represented branch is genuine source-history motion or an uncontrolled
      choice of rapidly moving spacetime basis.
\item Use the Duhamel equations of those same history columns to seek bounds
      for $\tau_n\Omega_n^t$ and $\tau_n\Omega_n^\nu$.  Boundary terms at
      $a_n,b_n$ and the dependence of the history basis on $n$ must remain
      explicit.
\item Return to the exact $H^1$ balance
      \eqref{eq:H1-balance}.  The needed new PDE input is a relative estimate
      on nonlinear production strong enough to bound $S_n$, not a restatement
      of the Prodi--Serrin criterion.  Any proposed estimate must be tested
      against the doubled-record scaling
      $K_n=\mathcal R_n^2/E_*$ and the summable debit
      \eqref{eq:raw-energy-price}.
\item Only after (N1)--(N3), replay the earlier action-overdrive, ancestry,
      incidence, and all-shell exits.  A full stability conclusion is licensed
      only if every one of them is closed on the same terminal chronology.
\end{enumerate}
\end{programme}

\begin{assessment}[Position after V1]
The V34 four-column acceleration card is no longer analytically fourfold.
Its three physical columns are one enstrophy-square problem; its fourth
column is one explicitly computable history graph-rate problem.  The
contiguous branch has the same split.  This is the narrowest rigorous stopping
line supported by the attached estimates.  The missing enstrophy upper budget
is of regularity-criterion strength, while the missing history estimates are
finite but must be derived from the actual Duhamel ancestry.  These are the
two upstream targets for V2.
\end{assessment}

% =============================================================================
\section*{Version V1 history}
\addcontentsline{toc}{section}{Version V1 history}
% =============================================================================

Created a new NS-only note series after auditing the six supplied files.
Preserved the V34 programme as an imported source rather than copying it.
Proved the finite-head record cap and the summable Leray-energy price of the
V34 raw pulse.  Proved the enstrophy-square continuation bridge.  Introduced
only the graph norms of the already frozen represented-history space, proved
their exact action and local-mass consequences, supplied a rank-one no-go
example, and gave finite generalized-eigenvalue certificates.  Reduced all
three physical V34 covariant columns to the full enstrophy-square stock.
Stated a sufficient but unacquired packet closing the last V34 action card.
No claim of global three-dimensional regularity was made.

\begin{thebibliography}{9}

\bibitem{FeffermanClay}
C.~L. Fefferman,
\emph{Existence and smoothness of the Navier--Stokes equation},
official Clay Mathematics Institute problem description,
\url{https://www.claymath.org/wp-content/uploads/2022/06/navierstokes.pdf}.

\bibitem{ESS}
L.~Escauriaza, G.~Seregin, and V.~\v Sver\'ak,
\emph{$L_{3,\infty}$-solutions of the Navier--Stokes equations and backward
uniqueness}, Russian Mathematical Surveys \textbf{58} (2003), 211--250,
\url{https://doi.org/10.1070/RM2003v058n02ABEH000609}.

\bibitem{GKP}
I.~Gallagher, G.~S. Koch, and F.~Planchon,
\emph{A profile decomposition approach to the
$L_t^\infty(L_x^3)$ Navier--Stokes regularity criterion},
Mathematische Annalen \textbf{355} (2013), 1527--1559,
\url{https://arxiv.org/abs/1012.0145}.

\end{thebibliography}

% =============================================================================
\section{Version V2: the represented-history bank is not yet physically acquired}
\label{sec:v2-history-audit}
% =============================================================================

This section executes item (N1) of the V1 acquisition order.  It returns to
the literal definitions in Versions 21 and 28--34 of the attached NS source,
rather than inferring a history space from the notation used at the terminal
card.

\begin{proposition}[Exact audit of the inherited history construction]
\label{prop:v30-audit}
The attached programme establishes the following facts, and no stronger
identification.
\begin{enumerate}[label=\textup{(\roman*)},leftmargin=4em]
\item Version 30 defines
\[
 \mathsf J\alpha=\sum_{r=1}^{m}\alpha_rj_r,\qquad
 G=\mathsf J^*\mathsf J,\qquad
 P_{\rm hist}=\mathsf JG^\dagger\mathsf J^*
\]
after saying to fix a finite complete list
$j_1,\ldots,j_m\in L^2(I;E_{\rm raw})$ of represented source histories.
\item Conditional on such a list, the projection, Moore--Penrose inverse,
complete static Gram card, and supported Schur complements in Versions
30--32 are exact finite-dimensional identities.
\item The programme does not supply the source labels, the common-history
maps, or the actual columns $j_r$.  Its own V30 ledger marks the physical
history Gram values as open and says that actual represented source-history
columns and mixed positive reads are still required.
\item Versions 31--34 reuse $P_{\rm hist}$ (or its cardwise copy $P_n$);
they do not fill the missing list.
\item Version 21 does construct a genuine causal Volterra return operator
\[
 (\mathcal Vf)(t)
   =-\int_a^tM(t)U(t,s)f(s)\,\dd s,
 \label{eq:v21-volterra-import}
\]
but the later text gives no incidence theorem identifying the V30 columns
$j_r$ with the outputs of \eqref{eq:v21-volterra-import}, or with specified
pullbacks of earlier physical source histories.
\end{enumerate}
Consequently the static V30 algebra is valid once its input is supplied, but
the attached files do not yet determine the history subspace on which the V1
rates $\Omega_n^t$ and $\Omega_n^\nu$ are to be evaluated.
\end{proposition}

\begin{proof}
This is a definition-and-dependency audit.  V30 first defines
$E_{\rm raw}=\operatorname{Ran}P_{\rm raw}$ and
$L^2(I;E_{\rm raw})$, then introduces the quoted finite list before defining
$\mathsf J$, $G$, and $P_{\rm hist}$.  The subsequent identities depend only
on that synthesis map.  The acquisition ledger at the end of V30 explicitly
leaves the physical Gram reads open.  V31 and V32 retain the same synthesis
map; V33 renames the projection cardwise; V34 differentiates the resulting
represented component.  None inserts a list of physical columns.  V21 gives
the displayed Volterra formula for a supplied daughter source, but no later
statement equates its range with the V30 list.  These observations prove each
item.
\end{proof}

\begin{theorem}[The represented action is undetermined without source-labelled admissibility]
\label{thm:projection-indeterminacy}
Let $\mathcal H$ be a real Hilbert space with $\dim\mathcal H\ge2$ and let
$0\ne d\in\mathcal H$.  For every
$a\in[0,\|d\|_{\mathcal H}^2]$ there is an orthogonal projection $P$ of rank
at most one such that
\[
 \|(I-P)d\|_{\mathcal H}^2=a.
 \label{eq:arbitrary-birth-action}
\]
Thus the represented/birth split of one fixed source cannot be inferred from
$d$ and Hilbert-space orthogonality alone.
\end{theorem}

\begin{proof}
Put $\widehat d=d/\|d\|$ and choose a unit vector
$e\perp\widehat d$.  Select $\theta\in[0,\pi/2]$ with
$\sin^2\theta=a/\|d\|^2$, let
$h_\theta=\cos\theta\,\widehat d+\sin\theta\,e$, and let $P$ project onto
$\operatorname{span}\{h_\theta\}$.  Since
$\|P d\|^2=\|d\|^2\cos^2\theta$, Pythagoras gives
\eqref{eq:arbitrary-birth-action}.  At the endpoint $a=\|d\|^2$ one may also
take the zero projection.
\end{proof}

\begin{firewall}[What Theorem~\ref{thm:projection-indeterminacy} does and does not say]
It does not invalidate a predeclared physical atlas.  It proves that
\emph{represented history} must be accompanied by admissibility data that are
independent of the target source.  At minimum a usable V30 card must record:
source labels; the earlier source histories; the causal common-history and
common-endpoint transports; the resulting explicit columns; an exclusion of
future or target-dependent data; and enough domain information to apply
$\partial_t$ and $\partial_t+\nu A$.  Without these items the terminal split
is not a determined physical observable.
\end{firewall}

% =============================================================================
\section{Static Grams do not control the covariant history rate}
\label{sec:v2-static-no-go}
% =============================================================================

The preceding audit concerns missing physical incidence.  Even if the V30
static Gram values were supplied, a second obstruction remains: they cannot
determine the derivative column required in V34.

\begin{theorem}[A fixed complete static card with unbounded covariant rate]
\label{thm:fixed-static-card}
Let $I=(0,1)$ and let $w$ be a unit spatial Stokes eigenvector,
$Aw=\lambda w$, $\lambda>0$.  Define
\[
 e(t)=w,\qquad
 j_N(t)=\sqrt2\sin(2\pi Nt)w,
\]
and choose fixed real numbers $\alpha,\beta,\gamma,\delta$ with
$\alpha\beta\ne0$.  Put
\[
 d_N=\alpha j_N+\beta e,\qquad
 \chi_N=\gamma j_N+\delta e,
 \qquad
 P_N=\operatorname{proj}_{\operatorname{span}\{j_N\}}.
\]
For every $N$, the complete V30--V32 static data are identical:
\begin{align}
 G_N&=1,&
 \langle j_N,d_N\rangle&=\alpha,&
 \langle j_N,\chi_N\rangle&=\gamma,\notag\\
 \|d_N\|^2&=\alpha^2+\beta^2,&
 \|\chi_N\|^2&=\gamma^2+\delta^2,&
 \langle d_N,\chi_N\rangle&=\alpha\gamma+\beta\delta.
 \label{eq:fixed-static-card}
\end{align}
The projected residual and its writer card are also fixed:
\[
 (I-P_N)d_N=\beta e,\qquad
 (I-P_N)\chi_N=\delta e,\qquad
 \langle(I-P_N)d_N,(I-P_N)\chi_N\rangle=\beta\delta.
\]
Nevertheless, for $r_N=P_Nd_N$,
\begin{align}
 \|\partial_t r_N\|_{L^2(I;L^2)}^2
   &=\alpha^2(2\pi N)^2,\label{eq:static-time-blowup}\\
 \|(\partial_t+\nu A)r_N\|_{L^2(I;L^2)}^2
   &=\alpha^2\bigl((2\pi N)^2+\nu^2\lambda^2\bigr)
     \longrightarrow\infty.
 \label{eq:static-covariant-blowup}
\end{align}
Hence no bound depending only on the complete static Gram, its Schur
complements, the scalar writer read, or the static source action can control
the V34 history graph rate.
\end{theorem}

\begin{proof}
Both $e$ and $j_N$ have norm one in $L^2(I;L^2)$, and they are orthogonal
because the sine has zero integral.  Expanding in this orthonormal pair proves
all identities in \eqref{eq:fixed-static-card} and the three residual
identities.  Also
$r_N=\alpha j_N$ and
$\|\partial_tj_N\|_2^2=(2\pi N)^2$.  Integration of
$\langle\partial_tj_N,j_N\rangle$ gives zero (equivalently the sine and cosine
are orthogonal), while $Aj_N=\lambda j_N$.  Expansion of the squared
covariant norm gives \eqref{eq:static-covariant-blowup}.
\end{proof}

\begin{assessment}[Exact consequence for the V34 card]
The V1 graph rates were not an optional strengthening of V30; derivative
information is logically necessary.  A static positive-semidefinite
completion can certify every value used in V30--V32 and still be blind to
arbitrarily fast represented-history motion.  The missing acquisition is
therefore a derivative-augmented physical Gram, not another static Schur
complement.
\end{assessment}

% =============================================================================
\section{The exact derivative-augmented Gram certificate}
\label{sec:v2-derivative-gram}
% =============================================================================

\begin{theorem}[Derivative Gram with redundant columns]
\label{thm:derivative-gram}
Let $\mathcal H,\mathcal K$ be real Hilbert spaces and let
$D:\operatorname{Dom}D\subset\mathcal H\to\mathcal K$ be linear.  Suppose
$\mathsf J:\mathbb R^m\to\operatorname{Dom}D$ has finite-dimensional range.
Set
\[
 G=\mathsf J^*\mathsf J,\qquad
 \mathsf J_D=D\mathsf J,\qquad
 H_D=\mathsf J_D^*\mathsf J_D.
\]
Then
\[
 \ker G=\ker\mathsf J\subset\ker\mathsf J_D=\ker H_D.
 \label{eq:gram-kernel-chain}
\]
For $d\in\operatorname{Dom}D$, define
\[
 b=\mathsf J^*d,\quad
 P=\mathsf JG^\dagger\mathsf J^*,\quad
 r=Pd,\quad q=(I-P)d,
\]
and
\[
 k_D=\mathsf J_D^*Dd,\qquad B_D=\|Dd\|_{\mathcal K}^2.
\]
Then the following identities are exact, including when the history list is
redundant:
\begin{align}
 Dr&=\mathsf J_DG^\dagger b,
 \label{eq:exact-Dr}\\
 \|Dr\|_{\mathcal K}^2
   &=b^{\mathsf T}G^\dagger H_DG^\dagger b,
 \label{eq:exact-Dr-norm}\\
 \|Dq\|_{\mathcal K}^2
   &=B_D-2b^{\mathsf T}G^\dagger k_D
       +b^{\mathsf T}G^\dagger H_DG^\dagger b.
 \label{eq:exact-Dq-norm}
\end{align}
In particular $(G,b,H_D)$ is the minimal represented-derivative card, while
$(G,b,H_D,k_D,B_D)$ is the complete derivative card for the residual.
\end{theorem}

\begin{proof}
For $c\in\mathbb R^m$,
$c^{\mathsf T}Gc=\|\mathsf Jc\|_{\mathcal H}^2$, proving
$\ker G=\ker\mathsf J$.  If $\mathsf Jc=0$, linearity gives
$D\mathsf Jc=0$; and
$c^{\mathsf T}H_Dc=\|\mathsf J_Dc\|_{\mathcal K}^2$.  This proves
\eqref{eq:gram-kernel-chain}.  The standard finite-rank projection formula
gives $r=\mathsf JG^\dagger b$.  Applying $D$ proves
\eqref{eq:exact-Dr}, and taking its squared norm proves
\eqref{eq:exact-Dr-norm}.  Finally
$Dq=Dd-\mathsf J_DG^\dagger b$; expansion of its squared norm gives
\eqref{eq:exact-Dq-norm}.  The kernel chain ensures that all expressions are
unchanged by null coefficient directions, so no full-rank hypothesis was
used.
\end{proof}

\begin{corollary}[Exact V34 acquisition]
\label{cor:v34-derivative-card}
On a V34 card take
\[
 \mathcal H=L^2(I_n;E_{{\rm raw},n}),\qquad
 \mathcal K=L^2(I_n;L_\sigma^2(\Torus^3)),\qquad
 D=D_{\nu,n}:=\partial_t+\nu A.
\]
Once the actual V30 history columns lie in $\operatorname{Dom}D_{\nu,n}$,
Theorem~\ref{thm:derivative-gram} computes the represented covariant action
exactly:
\[
 \|D_{\nu,n}r_n\|_{\mathcal K}^2
 =b_n^{\mathsf T}G_n^\dagger H_{\nu,n}G_n^\dagger b_n.
 \label{eq:v34-exact-covariant-read}
\]
Thus V1's generalized-eigenvalue upper certificate can be replaced cardwise
by the exact scalar read \eqref{eq:v34-exact-covariant-read}.  Uniform control
still requires a physical estimate on $H_{\nu,n}$ relative to $G_n$.
\end{corollary}

% =============================================================================
\section{What the V21 Volterra ancestry would buy}
\label{sec:v2-volterra}
% =============================================================================

The previous section says exactly which finite values are missing.  The next
result shows that the V21 causal formula could supply them if a genuine
V21-to-V30 incidence map is proved.

\begin{theorem}[Covariant derivative bound for a causal return history]
\label{thm:volterra-derivative}
Let $E,F$ be finite-dimensional Hilbert spaces, $I=(a,b)$, and
$L=b-a$.  Let $A_E,A_F$ be positive operators.  Suppose
\[
 y_t+\bigl(\nu A_E+L_0(t)\bigr)y=-f,\qquad y(a)=0,
 \label{eq:finite-return-equation}
\]
has evolution family $U(t,s)$ satisfying
$\|U(t,s)\|\le U_*$, and let $M\in W^{1,\infty}(I;\mathcal L(E,F))$.
For $j(t)=M(t)y(t)$ set
\[
 C_M(t)=M'(t)+\nu A_FM(t)-M(t)\bigl(\nu A_E+L_0(t)\bigr).
 \label{eq:return-commutator}
\]
Then
\begin{equation}
 (\partial_t+\nu A_F)j=C_My-Mf
 \label{eq:return-covariant-identity}
\end{equation}
and, if $\|M\|_\infty\le M_*$ and $\|C_M\|_\infty\le C_*$,
\begin{align}
 \|y\|_{L^2(I;E)}
   &\le LU_*\|f\|_{L^2(I;E)},\label{eq:return-y-bound}\\
 \|(\partial_t+\nu A_F)j\|_{L^2(I;F)}
   &\le (M_*+LC_*U_*)\|f\|_{L^2(I;E)}.
 \label{eq:return-graph-bound}
\end{align}
\end{theorem}

\begin{proof}
Variation of constants gives
$y(t)=-\int_a^tU(t,s)f(s)\,\dd s$.  Cauchy--Schwarz yields
$\|y(t)\|\le U_*L^{1/2}\|f\|_{L^2}$, and integration in $t$ proves
\eqref{eq:return-y-bound}.  The product rule and
\eqref{eq:finite-return-equation} give
\[
 (\partial_t+\nu A_F)(My)
 =\bigl[M'+\nu A_FM-M(\nu A_E+L_0)\bigr]y-Mf,
\]
which is \eqref{eq:return-covariant-identity}.  Its $L^2$ norm, followed by
\eqref{eq:return-y-bound}, proves \eqref{eq:return-graph-bound}.
\end{proof}

\begin{corollary}[Condition-number form of the history rate]
\label{cor:volterra-history-rate}
Suppose a finite source synthesis
$\mathsf F:\mathbb R^m\to L^2(I;E)$ is passed through the return map of
Theorem~\ref{thm:volterra-derivative}, producing the history synthesis
$\mathsf J:\mathbb R^m\to L^2(I;F)$.  Let
\[
 \sigma_+(\mathsf J)=\sqrt{\lambda_{\min}^+(\mathsf J^*\mathsf J)}
\]
be its smallest positive singular value.  Then
\[
 \Omega^\nu(\operatorname{Ran}\mathsf J)
 \le
 \frac{(M_*+LC_*U_*)\|\mathsf F\|_{\rm op}}
      {\sigma_+(\mathsf J)}.
 \label{eq:volterra-condition-bound}
\]
\end{corollary}

\begin{proof}
For $x\in\operatorname{Ran}\mathsf J$, choose the unique coefficient
$\alpha\perp\ker\mathsf J$ with $x=\mathsf J\alpha$.  Then
$\|\alpha\|\le\|x\|/\sigma_+(\mathsf J)$.  Apply
\eqref{eq:return-graph-bound} to $f=\mathsf F\alpha$ and take the supremum
over nonzero $x$.
\end{proof}

For the literal V21 notation, the commutator which must be estimated is
\[
 (M_P^u)' +\nu A_P M_P^u
       -M_P^u(\nu A_Q+L_Q^u).
 \label{eq:v21-required-commutator}
\]
This is a concrete advance: if the missing incidence map exists,
\eqref{eq:volterra-condition-bound} converts the terminal history-rate problem
into three inspectable quantities---a return commutator bound, a source
synthesis norm, and the smallest positive singular value of the physical
history synthesis.  Finiteness on each card is automatic in finite
dimension; uniformity along terminal cards is not.

\begin{firewall}[The V21 formula does not by itself fill V30]
A Volterra operator maps a supplied source history to a return history.  It
does not identify which earlier source histories are admissible, prove that
their transports form the V30 list, or prevent
$\sigma_+(\mathsf J_n)\downarrow0$.  Those are precisely the missing
source-labelled incidence and uniform conditioning statements.
\end{firewall}

% =============================================================================
\section{Going upstream on the enstrophy bottleneck}
\label{sec:v2-palinstrophy}
% =============================================================================

Write, on any smooth interval,
\[
 \mathcal E(t)=\|A^{1/2}u(t)\|_2^2,\qquad
 \mathcal Z(t)=\|Au(t)\|_2^2.
\]
The V1 stock was $S(I)=\int_I\mathcal E^2$.  It has an exact upstream
comparison with palinstrophy.

\begin{proposition}[Enstrophy-square is paid by palinstrophy]
\label{prop:enstrophy-palinstrophy}
For every smooth time,
\begin{equation}
 \mathcal E(t)^2
 =\langle u(t),Au(t)\rangle^2
 \le\|u(t)\|_2^2\|Au(t)\|_2^2
 \le E_*\mathcal Z(t).
 \label{eq:enstrophy-palinstrophy}
\end{equation}
Consequently the V34 raw pulse \eqref{eq:imported-raw-pulse} implies
\begin{equation}
 \int_{J_n}\mathcal Z(t)\,\dd t
 \ge \frac{c_{\rm raw}a_-\nu\mathcal R_n^2}{E_*}.
 \label{eq:raw-palinstrophy-lower}
\end{equation}
Moreover, the physical-column alternatives of
Theorem~\ref{thm:three-columns-enstrophy} imply respectively
\[
 \frac{E_*}{\nu\mathcal R_n^2}\int_{I_n}\mathcal Z\,\dd t\to\infty
 \quad\hbox{(commutator column)}
\]
or
\[
 E_*\int_{I_n}\mathcal Z\,\dd t\to\infty
 \quad\hbox{(internal or incoming column)}.
 \label{eq:physical-palinstrophy-overdrive}
\]
\end{proposition}

\begin{proof}
Equation \eqref{eq:enstrophy-palinstrophy} is Cauchy--Schwarz and the Leray
energy bound.  Since $\mathcal E_n\le\mathcal E$, integrate it and combine
with \eqref{eq:imported-raw-pulse} to obtain
\eqref{eq:raw-palinstrophy-lower}.  Finally
$S_n\le E_*\int_{I_n}\mathcal Z$; insert the two alternatives in
\eqref{eq:physical-column-consequences}.
\end{proof}

\begin{proposition}[Exact $H^1$ production budget]
\label{prop:H1-production-budget}
Let
\[
 \mathcal G(t)=-\langle\BNS(u(t)),Au(t)\rangle.
\]
For $a<b<T_*$,
\begin{equation}
 \frac12\mathcal E(b)+\nu\int_a^b\mathcal Z\,\dd t
 =\frac12\mathcal E(a)+\int_a^b\mathcal G\,\dd t.
 \label{eq:v2-H1-production-identity}
\end{equation}
In particular,
\begin{equation}
 \int_a^b\mathcal E(t)^2\,\dd t
 \le\frac{E_*}{\nu}
 \left(\frac12\mathcal E(a)+\int_a^b\mathcal G_+(t)\,\dd t\right).
 \label{eq:v2-H1-production-upper}
\end{equation}
Thus an upper bound for the positive $H^1$ nonlinear production, together
with control of the entry enstrophy, would close the V1 stock; the Leray
energy identity alone does not provide either bound.
\end{proposition}

\begin{proof}
Take the $L^2$ inner product of \eqref{eq:NS} with $Au$ and integrate in
time.  This proves \eqref{eq:v2-H1-production-identity}.  Combine it with
\eqref{eq:enstrophy-palinstrophy}, discard the nonnegative terminal term, and
replace $\mathcal G$ by $\mathcal G_+$ to obtain
\eqref{eq:v2-H1-production-upper}.
\end{proof}

\begin{proposition}[The critical record does not bound entry enstrophy]
\label{prop:critical-record-no-H1}
For every $R>0$ and $E_*>0$ there is a sequence of smooth mean-zero
divergence-free fields $u_N$ on $\Torus^3$ such that, for all sufficiently
large $N$,
\[
 \|A^{1/4}u_N\|_2=R,\qquad
 \|u_N\|_2^2\le E_*,
 \qquad
 \|A^{1/2}u_N\|_2^2=NR^2\longrightarrow\infty.
 \label{eq:critical-record-counterexample}
\]
Therefore the inherited $\dot H^{1/2}$ record corridor cannot be used as an
unstated bound on $\mathcal E(a_n)$ in
\eqref{eq:v2-H1-production-upper}.
\end{proposition}

\begin{proof}
Choose a unit divergence-free Fourier eigenfield $w_N$ with
$A^{1/2}w_N=Nw_N$ and put $u_N=RN^{-1/2}w_N$.  Direct calculation gives the
first and third identities in \eqref{eq:critical-record-counterexample}, while
$\|u_N\|_2^2=R^2/N\le E_*$ for sufficiently large $N$.
\end{proof}

\begin{assessment}[Where the physical branch now stops]
Proposition~\ref{prop:enstrophy-palinstrophy} does not create a contradiction:
large enstrophy-square action demands large palinstrophy, which a hypothetical
singularity may possess.  Proposition~\ref{prop:H1-production-budget} locates
the next noncircular PDE target at the positive nonlinear $H^1$ production.
Proposition~\ref{prop:critical-record-no-H1} prevents the critical
$\dot H^{1/2}$ record from silently supplying the missing entry value.  A
viable next estimate must therefore use a time cutoff in the finite-head
$H^1$ balance, retain its cutoff derivative and incoming source/pressure
terms, and show that their normalized cost is paid by an already finite
terminal budget.
\end{assessment}

% =============================================================================
\section{Integrated V2 conclusion and next acquisition order}
\label{sec:v2-conclusion}
% =============================================================================

\begin{theorem}[What V2 proves]
\label{thm:v2-conclusion}
Relative to the frozen V34 source base, the following statements are proved.
\begin{enumerate}[label=\textup{(V2.\arabic*)},leftmargin=5em]
\item The physical represented-history list required by V30 is not populated
      in the attached files; its static algebra is conditional on that
      missing acquisition.
\item Without target-independent source-labelled admissibility, the
      represented/birth source action is mathematically undetermined.
\item Even the complete V30--V32 static card cannot control the V34
      covariant represented-history rate.
\item Once physical columns are supplied, the exact derivative-augmented Gram
      formula \eqref{eq:v34-exact-covariant-read} evaluates that rate, with no
      independence assumption on the columns.
\item If the V30 columns are genuine V21 Volterra returns, their graph rate
      obeys the condition-number estimate
      \eqref{eq:volterra-condition-bound}; the missing incidence and uniform
      conditioning remain open.
\item The physical V34 alternatives can be pushed from enstrophy-square to
      palinstrophy, and their required upper budget is exactly exposed by
      \eqref{eq:v2-H1-production-upper}.  Neither the Leray energy bound nor
      the critical record controls the entry-enstrophy/positive-production
      pair.
\end{enumerate}
These statements sharpen both open branches but do not prove global
three-dimensional Navier--Stokes regularity.
\end{theorem}

\begin{proof}
Items \textup{(V2.1)--(V2.2)} are
Proposition~\ref{prop:v30-audit} and
Theorem~\ref{thm:projection-indeterminacy}.  Item \textup{(V2.3)} is
Theorem~\ref{thm:fixed-static-card}.  Item \textup{(V2.4)} is
Corollary~\ref{cor:v34-derivative-card}.  Item \textup{(V2.5)} is
Corollary~\ref{cor:volterra-history-rate} together with its firewall.  Item
\textup{(V2.6)} is Propositions~\ref{prop:enstrophy-palinstrophy}--
\ref{prop:critical-record-no-H1}.
\end{proof}

\begin{programme}[V3 acquisition order]
\label{prog:v3}
\begin{enumerate}[label=\textup{(P\arabic*)},leftmargin=4em]
\item Trace the literal physical source labels from V20 into the V21 return
      map and through V22--V29.  Either construct the target-independent
      incidence map producing the V30 list, or prove that the present
      definitions do not contain one.  Do not choose a subspace from the
      terminal target $d_n$.
\item If the incidence map exists, write its columns explicitly and acquire
      $(G_n,b_n,H_{\nu,n})$.  Expand
      \eqref{eq:v21-required-commutator}; test whether its terms coincide with
      already bounded V34 commutator/internal/incoming columns.  Keep
      $\sigma_+(\mathsf J_n)$ visible.
\item Apply a compactly supported time cutoff to the finite-head $H^1$
      identity on the V34 interval.  Track the cutoff derivative, incoming
      source, pressure cancellation, and spectral-edge commutators.  The
      target is a normalized upper bound for $\int\mathcal G_+$ or
      $\int\mathcal Z$ that uses only previously paid budgets.
\item Reject any endpoint that merely rederives
      $\int_0^{T_*}\mathcal E^2=\infty$ under blow-up: V1 already proved that
      criterion.  Return further upstream if the cutoff calculation asks for
      the same enstrophy-square stock it is meant to bound.
\end{enumerate}
\end{programme}

\begin{assessment}[Position after V2]
The history branch is now an acquisition problem with an exact checksum:
physical source incidence plus the derivative Gram
$(G_n,b_n,H_{\nu,n})$, and, for a Volterra realization, a uniform return
commutator/conditioning bound.  The physical branch is now an $H^1$
production problem with its entry term exposed.  No terminal contradiction
has been obtained, so the global stability goal remains open.
\end{assessment}

% =============================================================================
\section*{Version V2 history}
\addcontentsline{toc}{section}{Version V2 history}
% =============================================================================

Appended an upstream audit of Versions 21 and 28--34.  Proved that the
represented-history list is not populated by the attached source, that an
unlabelled projection can assign arbitrary birth action, and that a complete
static Gram card can remain fixed while its covariant derivative action
diverges.  Added an exact derivative-augmented Gram formula valid for
redundant columns, and derived the graph-rate estimate that would follow from
a genuine V21 Volterra incidence map.  Pushed the V1 enstrophy-square
bottleneck to palinstrophy and the exact positive $H^1$ production budget,
with a Fourier counterexample preventing use of the critical record as an
entry-enstrophy bound.  No claim of global regularity was made.
% =============================================================================
\section{Version V3: tracing the missing physical history incidence}
\label{sec:v3-incidence-trace}
% =============================================================================

Version V2 proved from the V30 ledger that the physical history columns had
not been acquired.  The present section performs the promised upstream trace
through Versions 20--29 and records the mathematical type of every plausible
predecessor.  This is stronger than a search for a repeated symbol: a valid
ancestry chain must have composable domains and codomains and must be frozen
before the terminal target is read.

\begin{proposition}[Typed V20--V30 dependency trace]
\label{prop:v3-typed-trace}
The objects in Versions 20--29 have the following types.
\begin{enumerate}[label=\textup{(T\arabic*)},leftmargin=4em]
\item V20 declares, at each time $t$, a spatial represented-source subspace
\[
 \mathcal J_{\mathscr A,t}^{\rm out,rep}
 \subseteq\operatorname{Ran}Q_{\mathscr A}
\]
and its pointwise orthogonal projection
$R_{\mathscr A,t}^{\rm out}$.  It does not give source labels, source
sections, or a formula for the stated common-endpoint transport.
\item V21 constructs the causal operator
\[
 \mathscr V_{F\leftarrow E}^u:
 L^2(I;E^{\rm dau})\longrightarrow L^2(I;F^{\rm ret}).
 \label{eq:v3-V21-type}
\]
It also permits a \emph{predeclared} finite source synthesis
$\mathsf D:\mathbb R^N\to L^2(I;E^{\rm dau})$, but neither lists its columns
nor identifies them with the V20 represented atlas.  The composite
$\mathscr V_{F\leftarrow E}^u\mathsf D$ consists of return histories in the
head screen $F^{\rm ret}$, not exterior source histories in the V30 raw
source carrier.
\item V23 constructs the pointwise first-birth synthesis
$C_t^{\rm b}=(I-R_t^{\rm out})\mathscr C_{\mathscr A}^{\rm out}$.
V25 replaces the pointwise atlas by the source-cyclic spatial subspace
$P_{C_t}\mathcal J_t$.  These remain time-indexed spatial subspaces; they do
not supply constant-coordinate spacetime columns.
\item V26--V28 construct carriers and a chronological filtration from the
already selected residual
$d^{\rm C}=(I-R_{\rm C})d^{\rm raw}$.  They classify that residual and
therefore cannot serve as a predeclared history bank against which the same
target is to be tested.
\item V29 constructs the fixed raw carrier from the occurring target history
$d^{\rm raw}$ and proves its analytic saturation.  This determines the
ambient space
$\mathscr X_{\rm raw}=L^2(I;E_{\rm raw})$, not a bank of earlier represented
source histories inside it.
\item V30 newly requires a finite map
\[
 \mathbf J:\mathbb R^m\longrightarrow\mathscr X_{\rm raw},
 \qquad \mathbf Je_r=j_r,
 \label{eq:v3-required-history-map}
\]
whose columns arise from already occurring source histories after actual
common-history and common-endpoint maps.  No object in \textup{(T1)--(T5)}
is identified with \eqref{eq:v3-required-history-map}.
\end{enumerate}
Thus there is no composable, target-independent V20--V29 construction in the
attached text whose output is the V30 physical history synthesis.
\end{proposition}

\begin{proof}
Item \textup{(T1)} is the definition preceding the V20 daughter
Pythagoras.  Item \textup{(T2)} is the domain and codomain of the V21 finite
return map and its optional finite Gram.  The V21 terminal assembly confirms
the type distinction: $d^{\rm rep}$ is an exterior source, whereas
$r^{\rm rep}=\mathscr Vd^{\rm rep}$ is a head-return history.  Items
\textup{(T3)--(T4)} are respectively the V23 first-birth synthesis, the V25
source-cyclic short, and the V26--V28 residual-generated carriers.  Item
\textup{(T5)} is the V29 raw-before-short construction.  V30 explicitly
introduces the new finite list in \eqref{eq:v3-required-history-map} and its
ledger leaves the corresponding physical Gram values open.  Since a
time-indexed spatial subspace is not a finite constant-coordinate history
synthesis, and since the V21 return codomain is not the V30 source codomain,
none of the preceding arrows has the required type.  No further
identification theorem is stated in V20--V29.
\end{proof}

\begin{firewall}[An optional placeholder is not acquired data]
The symbol $\mathsf D$ in V21 is syntactically capable of holding source
histories.  Calling it predeclared does not populate its columns, show that
they are earlier represented sources, or provide their common-endpoint maps.
Likewise, $\mathscr V\mathsf D$ cannot be reused as $\mathbf J$ without an
additional return-to-source map, because its values lie in the head return
screen.  V3 therefore does not infer a physical history bank from either
placeholder.
\end{firewall}

\begin{definition}[Minimum incidence datum which would fill V30]
\label{def:v3-incidence-datum}
A physical V30 incidence datum on a card consists of:
\begin{enumerate}[label=\textup{(I\arabic*)},leftmargin=3.6em]
\item a finite set of source labels $\mathscr L$, fixed before $d$ and
      $\chi$ are read;
\item for each $\ell\in\mathscr L$, an already occurring source history
      $f_\ell$ on its original interval and its first-birth record;
\item a declared causal common-history/common-endpoint operator
\[
 T_\ell:f_\ell\longmapsto j_\ell\in\mathscr X_{\rm raw};
\]
\item proof that $T_\ell$ uses no future value or coefficient fitted from the
      terminal target;
\item the domain statement
$j_\ell\in\operatorname{Dom}(\partial_t+\nu A)$ when the V34 covariant card
      is intended.
\end{enumerate}
The synthesis is then
$\mathbf J\alpha=\sum_{\ell\in\mathscr L}\alpha_\ell j_\ell$.
\end{definition}

\begin{theorem}[Sufficiency and irreducibility of the incidence datum]
\label{thm:v3-incidence-sufficiency}
Definition~\ref{def:v3-incidence-datum} is sufficient to evaluate the V30
static card and the V2 derivative card.  None of its items
\textup{(I1)--(I5)} is supplied by merely knowing the pointwise projections
$R_t$, the current target $d$, and the V21 return operator.
\end{theorem}

\begin{proof}
Items \textup{(I1)--(I3)} produce explicit columns and hence
$G=\mathbf J^*\mathbf J$, $b=\mathbf J^*d$, and every V30 mixed read.
Item \textup{(I4)} supplies the noncircular ancestry semantics.  Item
\textup{(I5)} makes $D\mathbf J$ well defined, so
Theorem~\ref{thm:derivative-gram} gives $H_\nu=(D\mathbf J)^*D\mathbf J$ and
the exact V34 read \eqref{eq:v34-exact-covariant-read}.

A pointwise projection specifies only membership at each time, not a
source-labelled section or its fixed coefficient law.  The current target
cannot provide an earlier label or target independence.  Finally the V21
operator maps a supplied exterior source to a head return and neither
supplies the source nor maps that return back into
$\mathscr X_{\rm raw}$.  These type and dependence requirements correspond
exactly to \textup{(I1)--(I5)}.
\end{proof}

\begin{assessment}[History-branch stopping line after the trace]
The history branch cannot be advanced by another Gram identity.  Its next
mathematical input is external to the present definitions: actual earlier
source labels and actual endpoint transports.  Once those data exist, V2
already provides the exact static and derivative checksums.  Inventing
columns from $d_n$, from the V28 residual filtration, or from arbitrary
sections of the V20 pointwise atlas would make the ancestry test circular.
\end{assessment}

% =============================================================================
\section{The finite-head critical balance and its exact source assembly}
\label{sec:v3-critical-balance}
% =============================================================================

We now execute the second V2 upstream instruction.  Let
$P_K=\mathbf1_{(0,K]}(A^{1/2})$, $Q_K=I-P_K$, and $v=P_Ku$ on a smooth
interval $I=[a,b]$.  Define the assembled head source in the native critical
metric by
\begin{align}
 s_K
 &:=A^{-1/4}P_K\BNS(u),\label{eq:v3-assembled-head-source}\\
 c_K^{\rm int}
 &:=A^{-1/4}P_K\BNS(v),\label{eq:v3-internal-head-source}\\
 \ell_K^{\rm in}
 &:=A^{-1/4}P_K\bigl(\BNS(u)-\BNS(v)\bigr).
 \label{eq:v3-incoming-head-source}
\end{align}
The inherited V21 identity gives
\begin{equation}
 \boxed{s_K=c_K^{\rm int}+\ell_K^{\rm in}.}
 \label{eq:v3-source-assembly}
\end{equation}
The terms in \eqref{eq:v3-source-assembly} need not be orthogonal; their
signed sum is formed before its square action is taken.

Put
\[
 R_K(t)=\|A^{1/4}v(t)\|_2^2,\qquad
 E_K(t)=\|A^{1/2}v(t)\|_2^2,\qquad
 D_K(t)=\|A^{3/4}v(t)\|_2^2.
\]

\begin{theorem}[Exact weighted finite-head critical balance]
\label{thm:v3-weighted-critical-balance}
For every nonnegative $\zeta\in W^{1,\infty}(a,b)$,
\begin{align}
 \nu\int_a^b\zeta D_K\,\dd t
 ={}&
 \frac{\zeta(a)R_K(a)-\zeta(b)R_K(b)}{2}
 +\frac12\int_a^b\zeta'R_K\,\dd t\notag\\
 &-\int_a^b\zeta
   \langle s_K,A^{3/4}v\rangle\,\dd t.
 \label{eq:v3-exact-critical-cutoff}
\end{align}
Consequently,
\begin{equation}
 \boxed{
 \int_a^b\zeta D_K\,\dd t
 \le
 \frac{\zeta(a)R_K(a)+\int_a^b|\zeta'|R_K\,\dd t}{\nu}
 +\frac1{\nu^2}\int_a^b\zeta\|s_K\|_2^2\,\dd t.}
 \label{eq:v3-critical-cutoff-upper}
\end{equation}
If $\overline R_K=\sup_I R_K$, then
\begin{equation}
 \boxed{
 \int_a^b\zeta E_K^2\,\dd t
 \le\overline R_K
 \left[
 \frac{\zeta(a)R_K(a)+\int_a^b|\zeta'|R_K\,\dd t}{\nu}
 +\frac1{\nu^2}\int_a^b\zeta\|s_K\|_2^2\,\dd t
 \right].}
 \label{eq:v3-head-enstrophy-square-upper}
\end{equation}
\end{theorem}

\begin{proof}
Project \eqref{eq:NS} by $P_K$ and use
$P_K\BNS(u)=A^{1/4}s_K$:
\[
 \partial_tv+\nu Av+A^{1/4}s_K=0.
\]
Taking its inner product with $A^{1/2}v$ gives
\[
 \frac12R_K'+\nu D_K
 =-\langle s_K,A^{3/4}v\rangle.
\]
Multiplication by $\zeta$ and integration by parts proves
\eqref{eq:v3-exact-critical-cutoff}.  Drop the favorable terminal term and
apply
\[
 |\langle s_K,A^{3/4}v\rangle|
 \le\frac1{2\nu}\|s_K\|_2^2+\frac\nu2D_K.
\]
After absorption this is \eqref{eq:v3-critical-cutoff-upper}.  Finally,
spectral Cauchy--Schwarz gives
\[
 E_K^2
 =\langle A^{1/4}v,A^{3/4}v\rangle^2
 \le R_KD_K.
\]
Taking the supremum of $R_K$ proves
\eqref{eq:v3-head-enstrophy-square-upper}.
\end{proof}

\begin{corollary}[Terminal record scale of the absolute-value estimate]
\label{cor:v3-terminal-critical-scale}
On the V1 terminal card, take $K=K_n$, $I=I_n$, and $\zeta\equiv1$.  The
record corridor gives
\begin{equation}
 \boxed{
 \int_{I_n}\mathcal E_n(t)^2\,\dd t
 \le
 \frac{16\mathcal R_n^4}{\nu}
 +\frac{4\mathcal R_n^2}{\nu^2}
   \int_{I_n}\|s_{K_n}(t)\|_2^2\,\dd t.}
 \label{eq:v3-terminal-head-upper}
\end{equation}
Even under the additional optimistic hypothesis
\begin{equation}
 \int_{I_n}\|s_{K_n}\|_2^2\,\dd t
 \le C_s\nu\mathcal R_n^2,
 \label{eq:v3-optimistic-head-action}
\end{equation}
this yields only
\begin{equation}
 \int_{I_n}\mathcal E_n^2\,\dd t
 \le(16+4C_s)\frac{\mathcal R_n^4}{\nu}.
 \label{eq:v3-quartic-upper}
\end{equation}
That scale cannot contradict the V34 raw lower bound
$c_{\rm raw}a_-\nu\mathcal R_n^2$ as
$\mathcal R_n\to\infty$.
\end{corollary}

\begin{proof}
Both $R_{K_n}(a_n)$ and $\overline R_{K_n}$ are at most
$4\mathcal R_n^2$ by \eqref{eq:record-corridor}.  Insert these values in
\eqref{eq:v3-head-enstrophy-square-upper}.  Then use
\eqref{eq:v3-optimistic-head-action}.  The ratio of the upper scale in
\eqref{eq:v3-quartic-upper} to the raw lower scale grows like
$\mathcal R_n^2/\nu^2$.
\end{proof}

\begin{firewall}[This is a route audit, not a lower-bound theorem]
A large upper bound does not prove that the actual finite-head stock is
large.  Corollary~\ref{cor:v3-terminal-critical-scale} proves the narrower
and necessary point: the direct cutoff, Cauchy--Schwarz, and Young argument
cannot close the raw branch at the canonical record scale, even if the
assembled head-source action is granted the natural
$O(\nu\mathcal R_n^2)$ bound.  A successful argument must retain additional
sign, cancellation, or inter-card information.
\end{firewall}

% =============================================================================
\section{The raw daughter does not pay the head production}
\label{sec:v3-source-separation}
% =============================================================================

The V34 raw daughter is
\[
 d_K^{\rm raw}=A^{-1/4}Q_K\BNS(v).
\]
It is absent from \eqref{eq:v3-exact-critical-cutoff}, which contains
$A^{-1/4}P_K\BNS(u)$.  This is not a notational inconvenience but a genuine
orthogonal source separation.

\begin{theorem}[Zero raw daughter with nonzero incoming $H^1$ production]
\label{thm:v3-incoming-counterexample}
Fix $K>1$.  There is a smooth real mean-zero divergence-free field
$u_0=v+w$ and the spectral split $P_K+Q_K=I$ such that
\begin{equation}
 P_Ku_0=v,\qquad Q_Ku_0=w,\qquad
 Q_K\BNS(v)=0,
 \label{eq:v3-zero-raw}
\end{equation}
but
\begin{equation}
 -\langle
 P_K\bigl(\BNS(u_0)-\BNS(v)\bigr),Av
 \rangle>0.
 \label{eq:v3-positive-incoming-production}
\end{equation}
The field may be used as smooth periodic initial data, so the identities
occur at the initial time of an actual local Navier--Stokes trajectory.
Consequently no estimate of the incoming positive $H^1$ production by a
function of $\|d_K^{\rm raw}\|_2$ which vanishes at zero can hold for all
smooth data.
\end{theorem}

\begin{proof}
Choose an integer $N$ sufficiently large and set
\[
 k=(1,0,0),\qquad
 p=(N,1,0),\qquad
 q=(1-N,-1,0).
\]
Then $p+q=k$, while $|p|,|q|>K$ and $|k|\le K$.  The vectors $p,q$ are
nonparallel.  Use the explicit polarizations from the inherited V20
daughter-coefficient lemma:
\[
 a=\frac{p\times q}{|p\times q|},\qquad
 b=\frac{P_{q^\perp}p}{|P_{q^\perp}p|}.
\]
Let the only positive-frequency Fourier coefficients of the real field $w$
be $\widehat w(p)=a$ and $\widehat w(q)=b$, together with their conjugate
negative-frequency copies.  The V20 coefficient calculation gives
\[
 \widehat{\BNS(w)}(k)=\mathrm i\,c\,a
\]
for a constant $c>0$ depending only on the harmless Fourier normalization.

Let $v$ be the real single-wave field with
$\widehat v(k)=-\mathrm i\,a$ and the conjugate coefficient at $-k$.
Because $a\perp k$, a single divergence-free plane wave has
$\BNS(v)=0$, proving \eqref{eq:v3-zero-raw}.  No mixed pair consisting of a
frequency $\pm k$ and one of $\pm p,\pm q$ sums to $\pm k$ when $N$ is
chosen large.  Hence the only contribution to the pairing with $Av$, which
is supported at $\pm k$, is $P_K\BNS(w)$.  The phase of $v$ was chosen so
that
\[
 \langle\BNS(w),Av\rangle<0.
\]
This proves \eqref{eq:v3-positive-incoming-production}.  All fields have
finite Fourier support.  Standard smooth local existence supplies the final
statement.
\end{proof}

\begin{assessment}[The complementary V20/V3 boundary tests]
V20 proved that incoming current can vanish while a new exterior daughter is
born.  Theorem~\ref{thm:v3-incoming-counterexample} proves the converse
separation relevant to the present bottleneck: the raw exterior daughter can
vanish while incoming head production is positive.  Therefore neither
boundary row pays the other.  The exact assembled head source
\eqref{eq:v3-source-assembly}, including its cross term, is the correct input
to the finite-head critical balance.
\end{assessment}

% =============================================================================
\section{Why the standard critical product estimate is circular here}
\label{sec:v3-product}
% =============================================================================

\begin{lemma}[Critical source product estimate]
\label{lem:v3-critical-product}
There is a torus constant $C_0$ such that every smooth divergence-free $u$
satisfies
\begin{equation}
 \boxed{
 \|A^{-1/4}\BNS(u)\|_2
 \le C_0\|A^{1/4}u\|_2\|A^{3/4}u\|_2.}
 \label{eq:v3-critical-product}
\end{equation}
The same bound holds with $A^{-1/4}P_K\BNS(u)$ on the left.
\end{lemma}

\begin{proof}
The Sobolev embeddings on $\Torus^3$ give
$H^{1/2}\hookrightarrow L^3$ and
$H^{3/2}\hookrightarrow W^{1,3}$.  Hence
\[
 \|(u\cdot\nabla)u\|_{L^{3/2}}
 \le\|u\|_{L^3}\|\nabla u\|_{L^3}
 \le C\|u\|_{H^{1/2}}\|u\|_{H^{3/2}}.
\]
By duality of $H^{1/2}\hookrightarrow L^3$,
$L^{3/2}\hookrightarrow H^{-1/2}$.  The Leray projection is bounded on
these spaces, and on mean-zero fields the homogeneous Stokes norms are
equivalent to the displayed Sobolev norms.  This proves
\eqref{eq:v3-critical-product}.  The spectral projection $P_K$ is a
contraction in $H^{-1/2}$ and commutes with $A$.
\end{proof}

\begin{corollary}[The returned estimate asks for the unpaid critical dissipation]
\label{cor:v3-product-circular}
On a V1 record corridor,
\begin{equation}
 \int_{I_n}\|s_{K_n}\|_2^2\,\dd t
 \le4C_0^2\mathcal R_n^2
 \int_{I_n}\|A^{3/4}u(t)\|_2^2\,\dd t.
 \label{eq:v3-source-by-critical-dissipation}
\end{equation}
Thus inserting the standard product estimate into
\eqref{eq:v3-terminal-head-upper} requires an upper bound for the full
critical dissipation $\int\|A^{3/4}u\|_2^2$.  The Leray energy inequality
controls only $\int\|A^{1/2}u\|_2^2$, and the critical energy identity absorbs
the nonlinear term only in the small-data regime
$C_0\|A^{1/4}u\|_2<\nu$.  It supplies no terminal-record bound when
$\mathcal R_n\to\infty$.
\end{corollary}

\begin{proof}
Apply Lemma~\ref{lem:v3-critical-product}, use that $P_{K_n}$ is a
contraction, square, integrate, and insert
\eqref{eq:record-corridor}.  For the last statement, pair the full equation
with $A^{1/2}u$:
\[
 \frac12\frac{\dd}{\dd t}\|A^{1/4}u\|_2^2
 +\nu\|A^{3/4}u\|_2^2
 =-\langle A^{-1/4}\BNS(u),A^{3/4}u\rangle.
\]
Lemma~\ref{lem:v3-critical-product} bounds the right side by
$C_0\|A^{1/4}u\|_2\|A^{3/4}u\|_2^2$.  Coercive absorption follows only
under the stated smallness condition.
\end{proof}

\begin{assessment}[Noncircular PDE target after V3]
The cutoff derivative is not the only problem.  With $\zeta\equiv1$ it
disappears, yet the entrance critical stock and the assembled head-source
action produce the quartic record scale \eqref{eq:v3-quartic-upper}.  The
standard product estimate then asks for the unpaid full
$H^{3/2}$ dissipation.  The next viable route must avoid taking absolute
values of the complete source work.  It must exploit a signed triad flux,
cancellation between $c_K^{\rm int}$ and $\ell_K^{\rm in}$, or a genuinely
new inter-card payment.  Replacing that work by its norm square returns to
the same circle.
\end{assessment}

% =============================================================================
\section{Integrated V3 conclusion and next acquisition order}
\label{sec:v3-conclusion}
% =============================================================================

\begin{theorem}[What V3 proves]
\label{thm:v3-conclusion}
Relative to the frozen attached source and V1--V2, the following are proved.
\begin{enumerate}[label=\textup{(V3.\arabic*)},leftmargin=5em]
\item The complete V20--V29 type trace contains no map producing the V30
      physical source-history synthesis.
\item Definition~\ref{def:v3-incidence-datum} gives the minimum data that
      would make the static and covariant history cards evaluable without
      target fitting.
\item The finite-head critical balance has the exact source assembly and
      weighted identity \eqref{eq:v3-exact-critical-cutoff}.
\item Its direct absolute-value estimate loses two powers of the terminal
      record and therefore cannot contradict the V34 raw pulse.
\item The raw exterior daughter action cannot pay the head production:
      an actual smooth datum has zero raw daughter and strictly positive
      incoming $H^1$ production.
\item The standard critical product estimate closes only under critical
      smallness and otherwise returns the unpaid $H^{3/2}$ dissipation.
\end{enumerate}
No unconditional global regularity or full NS stability theorem follows.
\end{theorem}

\begin{proof}
Items \textup{(V3.1)--(V3.2)} are
Proposition~\ref{prop:v3-typed-trace} and
Theorem~\ref{thm:v3-incidence-sufficiency}.  Items
\textup{(V3.3)--(V3.4)} are
Theorem~\ref{thm:v3-weighted-critical-balance} and
Corollary~\ref{cor:v3-terminal-critical-scale}.  Item \textup{(V3.5)} is
Theorem~\ref{thm:v3-incoming-counterexample}.  Item \textup{(V3.6)} is
Lemma~\ref{lem:v3-critical-product} and
Corollary~\ref{cor:v3-product-circular}.
\end{proof}

\begin{programme}[V4 acquisition order]
\label{prog:v4}
\begin{enumerate}[label=\textup{(Q\arabic*)},leftmargin=4em]
\item Keep the history branch stopped unless actual source labels and
      common-endpoint maps are supplied by an earlier definition or a new
      physical construction.  Do not manufacture $\mathbf J_n$ from the
      terminal target or the V28 residual filtration.
\item Expand the signed Fourier/helical form of
      $\langle s_{K_n},A^{3/4}v_n\rangle$ before applying Cauchy--Schwarz.
      Separate internal triads, high--high incoming triads, and cross-boundary
      triads with one common orientation convention.
\item Test whether energy antisymmetry, helicity, or the inherited material
      cell cancels a terminal fraction of the assembled head work.  Every
      cancellation must be on the same time interval and must retain the
      cross term between $c_{K_n}^{\rm int}$ and $\ell_{K_n}^{\rm in}$.
\item If no signed cancellation survives, construct a finite-mode obstruction
      demonstrating the failure and move upstream to inter-card flux.  Do not
      replace the failed signed estimate by
      \eqref{eq:v3-critical-product}, which is already audited as circular.
\item Only after obtaining a subquartic bound for the finite-head
      enstrophy-square stock, revisit the full-stock V34 physical columns;
      the high-frequency tail is not controlled by the finite-head result.
\end{enumerate}
\end{programme}

\begin{assessment}[Position after V3]
The history problem is now a missing physical incidence datum, not an
uncomputed matrix.  The raw physical problem is now a signed critical-flux
problem, not an unexpanded cutoff estimate.  V3 proves that the obvious
norm-square route is scale-inadequate and that raw exterior action and
incoming head production are independent boundary rows.  Both conclusions
narrow the next proof attempt without asserting the unresolved global
theorem.
\end{assessment}

% =============================================================================
\section*{Version V3 history}
\addcontentsline{toc}{section}{Version V3 history}
% =============================================================================

Appended the complete typed dependency trace from the V20 pointwise source
atlas through the V21 return map, V23/V25 source-cyclic shorts, V28
filtration, V29 raw carrier, and V30 spacetime history requirement.  Proved
that no composable physical history synthesis is supplied and stated the
minimum incidence datum which would fill it.  Derived the exact weighted
finite-head critical balance and quantified its quartic terminal-record
loss.  Constructed a finite-Fourier Navier--Stokes datum with zero raw
daughter but positive incoming $H^1$ production.  Proved the standard
critical source product estimate and showed precisely why it returns the
unpaid $H^{3/2}$ dissipation outside the small-data regime.  No global
regularity claim was made.
% =============================================================================
\section{Version V4: signed shell flux behind the $H^1$ production}
\label{sec:v4-flux-entry}
% =============================================================================

Version V3 showed that replacing the assembled critical source work by its
absolute norm square loses the terminal scale and returns an unpaid
$H^{3/2}$ dissipation.  We therefore retain the Fourier transfer signs.
The V19 energy--helicity invariant audit is not repeated.  The new question
is what remains after those two universal cancellations when the nonlinear
work is weighted by the Stokes eigenvalue.

For a smooth real divergence-free field, put
$\omega=\operatorname{curl}u$.  The standard identities are
\begin{align}
 \langle\BNS(u),u\rangle&=0,
 \label{eq:v4-energy-cancellation}\\
 \langle\BNS(u),\omega\rangle&=0,
 \label{eq:v4-helicity-cancellation}\\
 -\langle\BNS(u),Au\rangle
 &=\int_{\Torus^3}(\omega\cdot\nabla u)\cdot\omega\,\dd x.
 \label{eq:v4-vortex-stretching}
\end{align}
The first two are the inherited energy and helicity cancellations.  The last
quantity is the positive/negative $H^1$ production $\mathcal G$ of V2.

\begin{proof}[Proof of \eqref{eq:v4-energy-cancellation}--\eqref{eq:v4-vortex-stretching}]
Integration by parts and $\operatorname{div}u=0$ give
$\int(u\cdot\nabla u)\cdot u=0$.  The vector identity
\[
 (u\cdot\nabla)u=\omega\times u+\nabla\frac{|u|^2}{2}
\]
shows that its pairing with $\omega$ vanishes; the gradient and Leray
projection are invisible to divergence-free test fields.  Taking curl of
the equation gives
\[
 \partial_t\omega+u\cdot\nabla\omega-\omega\cdot\nabla u
 =\nu\Delta\omega.
\]
Pairing with $\omega$, and using
$\|\omega\|_2^2=\|A^{1/2}u\|_2^2$ and
$\|\nabla\omega\|_2^2=\|Au\|_2^2$, identifies the nonlinear term with
\eqref{eq:v4-vortex-stretching}.
\end{proof}

\begin{corollary}[One Stokes sphere has zero instantaneous $H^1$ production]
\label{cor:v4-one-sphere}
If $u$ is supported on one Stokes sphere
$Au=\lambda u$, then
\[
 -\langle\BNS(u),Au\rangle
 =-\lambda\langle\BNS(u),u\rangle=0.
\]
Hence nonzero $H^1$ production is necessarily an inter-eigenvalue transfer;
energy cancellation removes only the constant spectral weight.
\end{corollary}

% =============================================================================
\section{Exact cumulative-flux representation}
\label{sec:v4-flux-formula}
% =============================================================================

Use normalized Haar measure and the Fourier convention
\begin{equation}
 \widehat{\BNS(u)}(k)
 =\mathrm i\,P_k
   \sum_{p+q=k}(\widehat u(p)\cdot q)\widehat u(q).
 \label{eq:v4-Fourier-B}
\end{equation}
Here $P_k$ is the Euclidean orthogonal projection onto $k^\perp$.
For $k\ne0$, define the signed modal energy work
\begin{equation}
 T_k(t)
 :=\operatorname{Re}\left(
 \widehat{\BNS(u)}(k,t)\cdot
 \overline{\widehat u(k,t)}\right),
 \label{eq:v4-modal-work}
\end{equation}
and, for $\rho>0$, the cumulative low-frequency flux
\begin{equation}
 \Phi_\rho(t)
 :=\sum_{0<|k|\le\rho}T_k(t)
 =\langle P_\rho\BNS(u),P_\rho u\rangle.
 \label{eq:v4-cumulative-flux}
\end{equation}
Here $P_\rho=\mathbf1_{(0,\rho]}(A^{1/2})$.
Smoothness makes every sum and integral below absolutely convergent.

\begin{theorem}[Layer-cake formula for nonlinear $H^1$ production]
\label{thm:v4-layer-cake}
For a finite cutoff $K$ and $v_K=P_Ku$, let
\[
 \mathcal G_K(t)
 :=-\langle P_K\BNS(u),Av_K\rangle.
\]
Then
\begin{equation}
 \boxed{
 \mathcal G_K(t)
 =2\int_0^K\rho\,\Phi_\rho(t)\,\dd\rho
  -K^2\Phi_K(t).}
 \label{eq:v4-finite-layer-cake}
\end{equation}
For the full smooth field,
\begin{equation}
 \boxed{
 \mathcal G(t)
 =-\langle\BNS(u),Au\rangle
 =2\int_0^\infty\rho\,\Phi_\rho(t)\,\dd\rho.}
 \label{eq:v4-full-layer-cake}
\end{equation}
Moreover every flux is a pure boundary interaction:
\begin{equation}
 \boxed{
 \Phi_\rho
 =\left\langle
 P_\rho\bigl(\BNS(u)-\BNS(P_\rho u)\bigr),
 P_\rho u\right\rangle.}
 \label{eq:v4-flux-is-incoming}
\end{equation}
Thus the internal energy cancellation at each nested head converts the
$H^1$ production into a weighted assembly of incoming/cross-boundary energy
currents.
\end{theorem}

\begin{proof}
By Parseval,
\[
 \mathcal G_K=-\sum_{0<|k|\le K}|k|^2T_k.
\]
On the other hand, Fubini gives
\[
 2\int_0^K\rho\Phi_\rho\,\dd\rho
 =\sum_{0<|k|\le K}(K^2-|k|^2)T_k
 =K^2\Phi_K+\mathcal G_K,
\]
which proves \eqref{eq:v4-finite-layer-cake}.  Energy cancellation says
$\sum_kT_k=0$.  Letting $K\to\infty$ in the finite formula gives
\eqref{eq:v4-full-layer-cake}; smooth Fourier decay makes
$K^2\Phi_K\to0$.  Finally
\[
 \langle\BNS(P_\rho u),P_\rho u\rangle=0
\]
by \eqref{eq:v4-energy-cancellation}.  Subtracting this zero from
\eqref{eq:v4-cumulative-flux} proves
\eqref{eq:v4-flux-is-incoming}.
\end{proof}

\begin{assessment}[Relation to the V3 source assembly]
At the terminal cutoff $K_n$, V3 wrote the head source as
$c_{K_n}^{\rm int}+\ell_{K_n}^{\rm in}$.  Formula
\eqref{eq:v4-finite-layer-cake} does not discard the internal term.
Instead, it resolves its Stokes weight through all lower cutoffs
$0<\rho<K_n$.  At each such cutoff the internal energy pairing cancels and
the surviving value is the incoming current relative to that smaller head.
This is the exact signed upstream form of the same work, with no
norm-square duplication.
\end{assessment}

% =============================================================================
\section{The Leray budget controls net flux, but only at quartic record scale}
\label{sec:v4-net-flux}
% =============================================================================

\begin{theorem}[Uniform net-flux debit through every cutoff]
\label{thm:v4-net-flux}
For every smooth interval $[a,b]\Subset[0,T_*)$ and every $\rho>0$,
\begin{equation}
 \boxed{
 \left|\int_a^b\Phi_\rho(t)\,\dd t\right|
 \le\frac{E_*}{2}.}
 \label{eq:v4-net-flux-bound}
\end{equation}
Consequently, for every $K>0$,
\begin{equation}
 \boxed{
 \left|\int_a^b\mathcal G_K(t)\,\dd t\right|
 \le E_*K^2.}
 \label{eq:v4-net-H1-production}
\end{equation}
Both estimates hold on every subinterval; no disjointness or terminal limit
is required.
\end{theorem}

\begin{proof}
Let
\[
 e_\rho(t)=\|P_\rho u(t)\|_2^2,\qquad
 e_{>\rho}(t)=\|Q_\rho u(t)\|_2^2.
\]
The low- and high-frequency energy balances are
\begin{align}
 \frac12e_\rho'
 +\nu\|A^{1/2}P_\rho u\|_2^2
 &=-\Phi_\rho,
 \label{eq:v4-low-energy-balance}\\
 \frac12e_{>\rho}'
 +\nu\|A^{1/2}Q_\rho u\|_2^2
 &=\Phi_\rho.
 \label{eq:v4-high-energy-balance}
\end{align}
The first identity gives
\[
 \int_a^b\Phi_\rho
 \le\frac12e_\rho(a)\le\frac{E_*}{2},
\]
while the second gives
\[
 \int_a^b\Phi_\rho
 \ge-\frac12e_{>\rho}(a)\ge-\frac{E_*}{2}.
\]
This proves \eqref{eq:v4-net-flux-bound}.  Integrate
\eqref{eq:v4-finite-layer-cake} in time and use
\eqref{eq:v4-net-flux-bound}:
\[
 \left|2\int_0^K\rho
       \int_a^b\Phi_\rho\,\dd t\,\dd\rho\right|
 \le\frac{E_*K^2}{2},
 \qquad
 \left|K^2\int_a^b\Phi_K\,\dd t\right|
 \le\frac{E_*K^2}{2}.
\]
The triangle inequality proves \eqref{eq:v4-net-H1-production}.
\end{proof}

\begin{corollary}[Unconditional finite-head enstrophy-square bound]
\label{cor:v4-unconditional-head-stock}
On every smooth interval,
\begin{equation}
 \boxed{
 \int_a^b E_K(t)^2\,\dd t
 \le\frac{E_*}{\nu}
 \left(\frac12E_K(a)+E_*K^2\right),}
 \qquad
 E_K=\|A^{1/2}P_Ku\|_2^2.
 \label{eq:v4-unconditional-head-stock}
\end{equation}
On a V1 terminal card with $K=K_n$,
\begin{equation}
 \boxed{
 \int_{I_n}\mathcal E_n(t)^2\,\dd t
 \le\frac{3\mathcal R_n^4}{\nu}.}
 \label{eq:v4-terminal-quartic-bound}
\end{equation}
Thus V4 removes the optimistic assembled-source hypothesis
\eqref{eq:v3-optimistic-head-action} from the V3 quartic estimate, but the
record power remains too large to contradict
\eqref{eq:imported-raw-pulse}.
\end{corollary}

\begin{proof}
The finite-head $H^1$ identity is
\[
 \frac12E_K(b)+\nu\int_a^b\|AP_Ku\|_2^2\,\dd t
 =\frac12E_K(a)+\int_a^b\mathcal G_K\,\dd t.
\]
Drop the terminal term and use
\eqref{eq:v4-net-H1-production}.  Since
\[
 E_K(t)^2
 =\langle P_Ku,AP_Ku\rangle^2
 \le\|P_Ku\|_2^2\|AP_Ku\|_2^2
 \le E_*\|AP_Ku\|_2^2,
\]
this proves \eqref{eq:v4-unconditional-head-stock}.  On a terminal card,
Lemma~\ref{lem:head-cap} gives
$E_{K_n}(a_n)\le4K_n\mathcal R_n^2$, while
$K_n=\mathcal R_n^2/E_*$.  Hence
\[
 \frac{E_*}{\nu}
 \left(2K_n\mathcal R_n^2+E_*K_n^2\right)
 =\frac{3\mathcal R_n^4}{\nu}.
\]
\end{proof}

\begin{corollary}[Exact comparison with the raw pulse]
\label{cor:v4-raw-comparison}
If the V34 raw pulse occurs, then V4 gives only
\[
 c_{\rm raw}a_-\nu\mathcal R_n^2
 \le\int_{J_n}\mathcal E_n^2
 \le\int_{I_n}\mathcal E_n^2
 \le\frac{3\mathcal R_n^4}{\nu}.
 \label{eq:v4-raw-sandwich}
\]
For large $\mathcal R_n$ the two inequalities are compatible.  The signed
Leray flux budget therefore improves the provenance of the upper bound, not
its critical scaling.
\end{corollary}

% =============================================================================
\section{A canonical-cutoff triad survives energy and helicity cancellation}
\label{sec:v4-triad}
% =============================================================================

The next example tests the exact cancellation proposed in the V4 acquisition
order.  Its wavevector triple is different from the V19 disconnected-bank
example, and its purpose is different: it computes the weighted $H^1$ work
at the canonical cutoff.

\begin{theorem}[Explicit canonical-cutoff forward-transfer triad]
\label{thm:v4-canonical-triad}
For integers $N\ge1$ and amplitudes $L>0$, set
\[
 a=N(1,1,0),\qquad
 b=N(1,-1,0),\qquad
 c=N(-2,0,0),
 \qquad a+b+c=0,
\]
and
\[
 g=\frac{(1,-1,0)}{\sqrt2}.
\]
Define a real field $u_{N,L}$ by the positive/selected Fourier coefficients
\begin{equation}
 \widehat u(a)=Lg,\qquad
 \widehat u(b)=Le_3,\qquad
 \widehat u(c)=\mathrm iLe_3,
 \qquad
 \widehat u(-k)=\overline{\widehat u(k)}.
 \label{eq:v4-explicit-field}
\end{equation}
Then $u_{N,L}$ is smooth, mean zero, and divergence free, and
\begin{align}
 \|u_{N,L}\|_2^2&=6L^2,
 \label{eq:v4-triad-energy}\\
 \langle u_{N,L},\operatorname{curl}u_{N,L}\rangle&=0,
 \label{eq:v4-triad-helicity}\\
 \|A^{1/4}u_{N,L}\|_2^2
 &=4(1+\sqrt2)NL^2,
 \label{eq:v4-triad-critical-stock}\\
 -\langle\BNS(u_{N,L}),Au_{N,L}\rangle
 &=4\sqrt2\,N^3L^3>0.
 \label{eq:v4-triad-H1-production}
\end{align}
The canonical record cutoff of this datum is
\begin{equation}
 K_{\rm can}
 :=\frac{\|A^{1/4}u_{N,L}\|_2^2}{\|u_{N,L}\|_2^2}
 =\frac{2(1+\sqrt2)}{3}N,
 \qquad
 \sqrt2N<K_{\rm can}<2N.
 \label{eq:v4-triad-canonical-cutoff}
\end{equation}
Thus $P_{K_{\rm can}}$ contains the $a,b$ modes and excludes the $c$ mode.
The $a,b$ parents create a nonzero raw exterior daughter at $\pm c$, while
the $a,c$ interaction creates the nonzero incoming head work at $\pm b$.
They are the two sides of one energy-conserving triad, but their different
Stokes weights leave the positive value
\eqref{eq:v4-triad-H1-production}.
\end{theorem}

\begin{proof}
The selected polarizations are perpendicular to their wavevectors, and the
negative coefficients make the field real.  Parseval immediately gives
\eqref{eq:v4-triad-energy}.  Each coefficient has zero pairing with its
helical curl coefficient, giving \eqref{eq:v4-triad-helicity}.  Since
$|a|=|b|=\sqrt2N$ and $|c|=2N$, the critical stock is
\[
 2\sqrt2NL^2+2\sqrt2NL^2+4NL^2
 =4(1+\sqrt2)NL^2.
\]

Using \eqref{eq:v4-Fourier-B}, the coefficients which pair with the occupied
modes are
\begin{align*}
 \widehat{\BNS(u)}(\pm a)&=0,\\
 \widehat{\BNS(u)}(\pm b)&=\sqrt2\,NL^2e_3,\\
 \widehat{\BNS(u)}(c)&=-\mathrm i\sqrt2\,NL^2e_3,\qquad
 \widehat{\BNS(u)}(-c)=\mathrm i\sqrt2\,NL^2e_3.
\end{align*}
Therefore the reality-paired modal works are
\[
 \sum_{k=\pm a}T_k=0,\qquad
 \sum_{k=\pm b}T_k=2\sqrt2\,NL^3,\qquad
 \sum_{k=\pm c}T_k=-2\sqrt2\,NL^3.
\]
Their sum is zero, as required by energy conservation.  Weighting by
$|a|^2=|b|^2=2N^2$ and $|c|^2=4N^2$ gives
\[
 -\sum_k|k|^2T_k
 =-2N^2(2\sqrt2NL^3)
   -4N^2(-2\sqrt2NL^3)
 =4\sqrt2N^3L^3,
\]
which is \eqref{eq:v4-triad-H1-production}.  Division of
\eqref{eq:v4-triad-critical-stock} by
\eqref{eq:v4-triad-energy} proves
\eqref{eq:v4-triad-canonical-cutoff}.

For $v=P_{K_{\rm can}}u$, the $a,b$ interaction has output
$-c$ (and its reality mate), so $Q_{K_{\rm can}}\BNS(v)\ne0$.  The coefficient
at $b$ requires the $a,c$ interaction and vanishes from $\BNS(v)$; hence it
belongs to
$P_{K_{\rm can}}(\BNS(u)-\BNS(v))$.  This proves the last statement.
\end{proof}

\begin{corollary}[Energy, helicity, and their values do not bound $H^1$ production]
\label{cor:v4-invariant-no-H1-bound}
The family $u_{N,1}$ has fixed kinetic energy $6$, zero helicity, and obeys
the exact nonlinear energy and helicity cancellations, while its positive
$H^1$ production tends to infinity like $N^3$.  Consequently no pointwise
upper bound for $\mathcal G_+$ can depend only on kinetic energy, helicity,
and the fact that their nonlinear derivatives vanish.
\end{corollary}

\begin{firewall}[No singularity is constructed]
For fixed viscosity, the palinstrophy of the example is
\[
 \|Au_{N,L}\|_2^2=48N^4L^2.
\]
Thus viscosity dominates \eqref{eq:v4-triad-H1-production} when $L/N$ is
small.  If $L>6\sqrt2\,\nu N$, the instantaneous enstrophy derivative is
positive, but the datum then has correspondingly large energy.  The example
tests cancellation and scale only; it is not a blow-up construction.
\end{firewall}

% =============================================================================
\section{Positive flux requires spectral energy variation}
\label{sec:v4-variation}
% =============================================================================

The net debit \eqref{eq:v4-net-flux-bound} controls the signed time integral
of each flux.  The continuation estimate in V2 instead needs the positive
part of $H^1$ production.  The exact distinction is a total-variation
quantity of already existing spectral energies.

For an absolutely continuous nonnegative function $f$ on $I$, write
\[
 \operatorname{Var}_I^-(f):=\int_I(-f'(t))_+\,\dd t.
\]

\begin{theorem}[Positive-production reduction to downward spectral variation]
\label{thm:v4-positive-variation}
For every finite $K$ and smooth interval $I$,
\begin{equation}
 \boxed{
 \int_I(\mathcal G_K)_+\,\dd t
 \le
 \int_0^K\rho\,
   \operatorname{Var}_I^-(e_\rho)\,\dd\rho
 +\frac{K^2}{2}\operatorname{Var}_I^-(e_{>K}),}
 \label{eq:v4-positive-variation}
\end{equation}
where
$e_\rho=\|P_\rho u\|_2^2$ and
$e_{>\rho}=\|Q_\rho u\|_2^2$.
For the full smooth field,
\begin{equation}
 \boxed{
 \int_I\mathcal G_+\,\dd t
 \le
 \int_0^\infty\rho\,
   \operatorname{Var}_I^-(e_\rho)\,\dd\rho.}
 \label{eq:v4-full-positive-variation}
\end{equation}
\end{theorem}

\begin{proof}
From \eqref{eq:v4-finite-layer-cake},
\[
 (\mathcal G_K)_+
 \le2\int_0^K\rho(\Phi_\rho)_+\,\dd\rho
    +K^2(-\Phi_K)_+.
\]
The low balance \eqref{eq:v4-low-energy-balance} gives
\[
 (\Phi_\rho)_+
 =\left(-\frac12e_\rho'
        -\nu\|A^{1/2}P_\rho u\|_2^2\right)_+
 \le\frac12(-e_\rho')_+.
\]
The high balance \eqref{eq:v4-high-energy-balance} similarly gives
\[
 (-\Phi_K)_+
 \le\frac12(-e_{>K}')_+.
\]
Integrate in time and use Fubini to obtain
\eqref{eq:v4-positive-variation}.  Letting $K\to\infty$ and using smooth
Fourier decay removes the high-frequency boundary term and proves
\eqref{eq:v4-full-positive-variation}.
\end{proof}

\begin{counterexample}[Analyticity and an energy cap do not bound the variation]
\label{cth:v4-analytic-variation}
Fix $0<\varepsilon<1$ and define on $[0,1]$
\[
 e_N(t)=\frac{E_*}{2}
 \bigl(1+\varepsilon\sin(2\pi Nt)\bigr).
\]
Then $e_N$ is positive, real analytic, and bounded above by $E_*$, but
\[
 \operatorname{Var}_{[0,1]}^-(e_N)
 =\varepsilon NE_*\longrightarrow\infty.
 \label{eq:v4-analytic-variation}
\]
Thus the V29 fact of time analyticity, without a uniform complex-time radius
or derivative bound, cannot pay the variation on the right of
\eqref{eq:v4-positive-variation}.
\end{counterexample}

\begin{proof}
The range statement is immediate.  Since
$e_N'=\varepsilon\pi NE_*\cos(2\pi Nt)$ and the mean of
$(-\cos)_+$ over a period is $1/\pi$, integration gives
\eqref{eq:v4-analytic-variation}.
\end{proof}

\begin{firewall}[The scalar example is not asserted to be an NS trajectory]
Counterexample~\ref{cth:v4-analytic-variation} proves only that boundedness
and analyticity of a spectral stock are insufficient hypotheses for a
uniform variation estimate.  Any NS-specific variation bound must use the
equation, the terminal chronology, or a paid flux mechanism.
\end{firewall}

% =============================================================================
\section{Integrated V4 conclusion and next acquisition order}
\label{sec:v4-conclusion}
% =============================================================================

\begin{theorem}[What V4 proves]
\label{thm:v4-conclusion}
Relative to V1--V3 and the frozen attached source, the following statements
are proved.
\begin{enumerate}[label=\textup{(V4.\arabic*)},leftmargin=5em]
\item The nonlinear $H^1$ production is the weighted cumulative shell flux
      \eqref{eq:v4-full-layer-cake}; every cumulative flux is an incoming
      boundary current relative to its own nested head.
\item The Leray balance bounds the net flux through every cutoff and every
      time interval by $E_*/2$.
\item This yields the unconditional finite-head bound
      \eqref{eq:v4-terminal-quartic-bound}, eliminating the unpaid source
      action from V3 but retaining an insufficient quartic record scale.
\item The explicit family \eqref{eq:v4-explicit-field} places the canonical
      cutoff strictly inside one active triad.  Energy and helicity cancel
      exactly while positive $H^1$ production survives and grows like $N^3$.
\item Positive $H^1$ production is controlled by the downward variation of
      nested spectral energies as in
      \eqref{eq:v4-positive-variation}; net energy and time analyticity alone
      do not control that variation.
\end{enumerate}
These results do not establish an unconditional NS stability theorem.
\end{theorem}

\begin{proof}
Items \textup{(V4.1)--(V4.2)} are
Theorems~\ref{thm:v4-layer-cake} and \ref{thm:v4-net-flux}.
Item \textup{(V4.3)} is
Corollary~\ref{cor:v4-unconditional-head-stock}.  Item \textup{(V4.4)} is
Theorem~\ref{thm:v4-canonical-triad} and its corollary.  Item
\textup{(V4.5)} is Theorem~\ref{thm:v4-positive-variation} and
Counterexample~\ref{cth:v4-analytic-variation}.
\end{proof}

\begin{programme}[V5 acquisition order]
\label{prog:v5}
\begin{enumerate}[label=\textup{(R\arabic*)},leftmargin=4em]
\item Retain the exact flux $\Phi_\rho$; do not return to the critical-source
      norm square already excluded in V3.
\item Test whether the terminal record chronology bounds the downward
      variation of $e_\rho$ on the selected card.  Entry/exit energy values
      bound net flux but not repeated spectral recycling.
\item Decompose the variation only through the already occurring low/high
      energy balances.  Any proposed monotonicity interval or crossing count
      must be proved from the NS dynamics, not from analyticity alone.
\item Compare the weighted variation with the V34 intrinsic time and material
      cell.  A useful estimate must improve
      $E_*K_n^2=\mathcal R_n^4/E_*$ to the raw-pulse scale
      $O(\nu\mathcal R_n^2)$ or smaller.
\item If a triad-recycling example defeats such a variation estimate, record
      it and move upstream to the signed inter-card flux.  Preserve the raw
      daughter/incoming return as one ancestry-labelled transfer.
\item Keep the history branch at the V3 incidence stopping line until actual
      earlier source labels and endpoint transports are available.
\end{enumerate}
\end{programme}

\begin{assessment}[Position after V4]
The signed route does remove one artificial obstruction: a full critical
source-action bound is not needed to control the \emph{net} finite-head
production.  The exact Leray payment is nevertheless quartic at the
canonical cutoff.  The remaining quantity is not another Gram or static
invariant; it is positive spectral recycling, measured by downward
variation of the nested low/high kinetic energies.  V4 exhibits an actual
canonical-cutoff triad which survives energy and helicity cancellation, so
any further gain must use time chronology or inter-card flux rather than
universal invariant algebra.
\end{assessment}

% =============================================================================
\section*{Version V4 history}
\addcontentsline{toc}{section}{Version V4 history}
% =============================================================================

Appended the exact cumulative shell-flux representation of nonlinear
$H^1$ production.  Proved a uniform $E_*/2$ net-flux debit through every
cutoff and used it to derive an unconditional finite-head
enstrophy-square upper bound of $3\mathcal R_n^4/\nu$.  Constructed a new
finite-Fourier triad whose canonical record cutoff separates its donor and
receiver modes: energy and helicity cancel, while positive $H^1$ production
is $4\sqrt2N^3L^3$.  Reduced the positive-production problem to downward
variation of nested spectral kinetic energies and proved that an energy cap
plus time analyticity does not bound that variation.  No global regularity
claim was made.
% =============================================================================
\section{Version V5: the critical-flux moment gap}
\label{sec:v5-moment-gap}
% =============================================================================

Version V4 expressed nonlinear $H^1$ production as the first radial moment
of the cumulative energy flux.  The terminal chronology instead records the
$\dot H^{1/2}$ stock.  We now put both quantities in the same exact formula.

For $\beta>0$, define
\begin{align}
 \mathcal G_{\beta,K}(t)
 &:=-\langle P_K\BNS(u),A^\beta P_Ku\rangle,
 \label{eq:v5-fractional-production-finite}\\
 \mathcal G_\beta(t)
 &:=-\langle\BNS(u),A^\beta u\rangle.
 \label{eq:v5-fractional-production-full}
\end{align}
Thus $\mathcal G_{1/2}$ is the nonlinear production of the critical stock
$\|A^{1/4}u\|_2^2$, while $\mathcal G_1=\mathcal G$ is the $H^1$
production.

\begin{theorem}[Fractional layer-cake identity]
\label{thm:v5-fractional-layer-cake}
With the cumulative flux $\Phi_\rho$ of \eqref{eq:v4-cumulative-flux},
\begin{equation}
 \boxed{
 \mathcal G_{\beta,K}
 =2\beta\int_0^K\rho^{2\beta-1}\Phi_\rho\,\dd\rho
  -K^{2\beta}\Phi_K.}
 \label{eq:v5-fractional-layer-cake-finite}
\end{equation}
For a smooth full field,
\begin{equation}
 \boxed{
 \mathcal G_\beta
 =2\beta\int_0^\infty
   \rho^{2\beta-1}\Phi_\rho\,\dd\rho.}
 \label{eq:v5-fractional-layer-cake-full}
\end{equation}
In particular,
\begin{align}
 \mathcal G_{1/2}
 &=\int_0^\infty\Phi_\rho\,\dd\rho,
 \label{eq:v5-critical-flux-moment}\\
 \mathcal G_1
 &=2\int_0^\infty\rho\Phi_\rho\,\dd\rho.
 \label{eq:v5-H1-flux-moment}
\end{align}
\end{theorem}

\begin{proof}
Parseval gives
\[
 \mathcal G_{\beta,K}
 =-\sum_{0<|k|\le K}|k|^{2\beta}T_k.
\]
Fubini and the elementary identity
$K^{2\beta}-r^{2\beta}
 =2\beta\int_r^K\rho^{2\beta-1}\,\dd\rho$ give
\[
 2\beta\int_0^K\rho^{2\beta-1}\Phi_\rho\,\dd\rho
 =\sum_{0<|k|\le K}
   (K^{2\beta}-|k|^{2\beta})T_k
 =K^{2\beta}\Phi_K+\mathcal G_{\beta,K}.
\]
This proves the finite formula.  Smooth Fourier decay and
$\sum_kT_k=0$ allow $K\to\infty$ and give the full formula.
\end{proof}

\begin{corollary}[What the critical record actually measures]
\label{cor:v5-critical-record-balance}
The full critical stock
\[
 \mathcal R(t)^2:=\|A^{1/4}u(t)\|_2^2
\]
satisfies
\begin{equation}
 \boxed{
 \frac12\frac{\dd}{\dd t}\mathcal R(t)^2
 +\nu\|A^{3/4}u(t)\|_2^2
 =\int_0^\infty\Phi_\rho(t)\,\dd\rho.}
 \label{eq:v5-critical-record-flux}
\end{equation}
The desired $H^1$ control instead involves the $\rho$-weighted moment
\eqref{eq:v5-H1-flux-moment}.  Hence a corridor or crossing law for the
single scalar $\mathcal R(t)$ controls only the net zeroth radial moment of
the flux; it does not determine its first radial moment.
\end{corollary}

\begin{proof}
Pair \eqref{eq:NS} with $A^{1/2}u$ and use
\eqref{eq:v5-critical-flux-moment}.
\end{proof}

% =============================================================================
\section{An active homochiral triad is invisible to the nonlinear record}
\label{sec:v5-homochiral}
% =============================================================================

The preceding moment distinction could conceivably collapse under the
inherited energy--helicity algebra.  The next theorem proves that it does
not.  We use the exact V19 modal-transfer lemma rather than repeating its
derivation.

Let
\[
 p=N(1,0,0),\qquad
 q=N(0,1,1),\qquad
 r=-N(1,1,1),
 \qquad p+q+r=0,
\]
and choose positive helicity on all three reality-paired atoms.  At $N=1$,
one compatible choice of normalized helical vectors is
\begin{align}
 h_p&=\frac1{\sqrt2}(0,1,\mathrm i),\notag\\
 h_q&=\frac1{\sqrt2}
 \left(1,\frac{\mathrm i}{\sqrt2},
          -\frac{\mathrm i}{\sqrt2}\right),\notag\\
 h_r&=\frac1{\sqrt2}
 \left(
 \frac{(1,-1,0)}{\sqrt2}
 +\mathrm i\frac{(-1,-1,2)}{\sqrt6}
 \right).
 \label{eq:v5-helical-vectors}
\end{align}
The same vectors are used after dilation by $N$.

\begin{lemma}[The distinct-radius homochiral triad is active]
\label{lem:v5-homochiral-active}
The vectors in \eqref{eq:v5-helical-vectors} satisfy
\[
 \mathrm i k\times h_k=|k|h_k
 \quad(k=p,q,r),
\]
and their V19 geometric coefficient is nonzero.  More precisely,
\begin{equation}
 \operatorname{Re}
 \bigl((\overline h_q\times\overline h_r)
       \cdot\overline h_p\bigr)
 =
 \frac14\left[
 \frac1{\sqrt2}
 +\frac{1-1/\sqrt2}{\sqrt3}\right]>0.
 \label{eq:v5-active-coefficient}
\end{equation}
\end{lemma}

\begin{proof}
For each wavevector, the real and imaginary directions displayed in
\eqref{eq:v5-helical-vectors} are an oriented orthonormal basis of its
orthogonal plane.  The curl eigenvector identities follow directly.
Substitution in the scalar triple product gives
\eqref{eq:v5-active-coefficient}, proving activity.
\end{proof}

\begin{theorem}[Energy and critical record fixed, $H^1$ transfer nonzero]
\label{thm:v5-record-invisible-triad}
Let the three occupied positive-helicity atoms have nonzero amplitudes whose
total phase is chosen as in the V19 transfer lemma.  Write their modal
energies as $\mathcal E_p,\mathcal E_q,\mathcal E_r$.  The nonlinear transfer
has the form
\begin{equation}
 \begin{pmatrix}
 \dot{\mathcal E}_p\\
 \dot{\mathcal E}_q\\
 \dot{\mathcal E}_r
 \end{pmatrix}_{\rm nl}
 =
 \Xi N
 \begin{pmatrix}
 \sqrt2-\sqrt3\\
 \sqrt3-1\\
 1-\sqrt2
 \end{pmatrix},
 \qquad \Xi\ne0.
 \label{eq:v5-homochiral-transfer}
\end{equation}
Consequently,
\begin{align}
 \frac{\dd}{\dd t}
  (\mathcal E_p+\mathcal E_q+\mathcal E_r)_{\rm nl}
 &=0,
 \label{eq:v5-triad-energy-fixed}\\
 \frac{\dd}{\dd t}
  \bigl(
   N\mathcal E_p+\sqrt2N\mathcal E_q
   +\sqrt3N\mathcal E_r
  \bigr)_{\rm nl}
 &=0,
 \label{eq:v5-triad-record-fixed}\\
 \frac{\dd}{\dd t}
  \bigl(
   N^2\mathcal E_p+2N^2\mathcal E_q
   +3N^2\mathcal E_r
  \bigr)_{\rm nl}
 &=
 \Xi N^3
 (1-\sqrt2)(\sqrt2-\sqrt3)(1-\sqrt3)\ne0.
 \label{eq:v5-triad-H1-nonzero}
\end{align}
The phase may be chosen so that the last derivative is positive.
\end{theorem}

\begin{proof}
Lemma~\ref{lem:v5-homochiral-active} permits application of the exact V19
modal-transfer formula, whose signed frequencies are
$N,\sqrt2N,\sqrt3N$.  This gives
\eqref{eq:v5-homochiral-transfer}.  The transfer vector is orthogonal to
both
\[
 (1,1,1)
 \quad\hbox{and}\quad
 (1,\sqrt2,\sqrt3),
\]
which proves \eqref{eq:v5-triad-energy-fixed} and
\eqref{eq:v5-triad-record-fixed}.  Its pairing with
$(1,2,3)$ is the Vandermonde value
\[
 (1-\sqrt2)(\sqrt2-\sqrt3)(1-\sqrt3),
\]
which is nonzero because the three radii are distinct.  The inherited phase
freedom changes the sign of $\Xi$.
\end{proof}

\begin{corollary}[Actual smooth data with decreasing record and positive $H^1$ production]
\label{cor:v5-decreasing-record-positive-H1}
Use the finite Fourier field of
Theorem~\ref{thm:v5-record-invisible-triad} as initial data for the full
Navier--Stokes equation.  At the initial time, all modes outside the triad
have zero amplitude, so they contribute zero first derivative to every
quadratic spectral stock.  Therefore the nonlinear derivative of
$\|A^{1/4}u\|_2^2$ is zero while the nonlinear $H^1$ production is positive.
Viscosity makes the critical stock strictly decrease.  By multiplying all
three amplitudes by a sufficiently large common factor, the cubic positive
$H^1$ production dominates the quadratic viscous $H^1$ loss at that instant.
Thus an actual smooth local trajectory can have decreasing critical record
and increasing enstrophy at the same time.
\end{corollary}

\begin{proof}
Finite Fourier data are smooth and generate a smooth local solution.
Newly generated modes have amplitude zero at the initial time, hence their
quadratic energies have zero first derivative.  Equations
\eqref{eq:v5-triad-record-fixed}--\eqref{eq:v5-triad-H1-nonzero} therefore
give the nonlinear derivatives of the full stocks.  The critical balance
then contains only the strictly negative viscous term.  Under common
amplitude dilation by $L$, nonlinear $H^1$ production is cubic in $L$ and
palinstrophy is quadratic, so the former dominates for sufficiently large
$L$.
\end{proof}

\begin{corollary}[No pointwise production bound from the record and energy cap]
\label{cor:v5-fixed-record-unbounded-production}
Let
\[
 S_*:=1+\sqrt2+\sqrt3
\]
and choose equal amplitude modulus
\begin{equation}
 L_N:=\frac{\mathcal R_0}{\sqrt{2S_*N}}
 \label{eq:v5-fixed-record-amplitude}
\end{equation}
with the production-positive phase.  Then
\begin{align}
 \|A^{1/4}u_N\|_2^2&=\mathcal R_0^2,
 \label{eq:v5-fixed-critical-stock}\\
 \|u_N\|_2^2&=\frac{3\mathcal R_0^2}{S_*N},
 \label{eq:v5-vanishing-energy}\\
 \mathcal G_1(u_N)_+
 &=c_{\rm h}\mathcal R_0^3N^{3/2}
 \longrightarrow\infty
 \label{eq:v5-unbounded-H1-production}
\end{align}
for a fixed constant $c_{\rm h}>0$ determined by the Fourier normalization
and the active coefficient.  In particular, for all sufficiently large
$N$, the energy lies below any prescribed $E_*>0$, while the critical record
is fixed and positive $H^1$ production is unbounded.
\end{corollary}

\begin{proof}
A reality-paired unit helical atom with amplitude modulus $L$ contributes
$2L^2$ to kinetic energy and $2|k|L^2$ to the critical stock.  Summing the
three atoms and inserting \eqref{eq:v5-fixed-record-amplitude} proves
\eqref{eq:v5-fixed-critical-stock}--\eqref{eq:v5-vanishing-energy}.
The V19 common transfer scalar is a fixed nonzero geometric coefficient
times $L_N^3$, while
\eqref{eq:v5-triad-H1-nonzero} contributes $N^3$.  Hence the production is
a fixed positive multiple of
$N^3L_N^3=\mathcal R_0^3N^{3/2}/(2S_*)^{3/2}$.
\end{proof}

\begin{firewall}[Pointwise production is not net growth]
For the equal-amplitude triad,
\[
 \|Au_{N,L}\|_2^2
 =28N^4L^2.
 \label{eq:v5-homochiral-palinstrophy}
\]
Thus, under the fixed-record scaling
\eqref{eq:v5-fixed-record-amplitude}, viscosity grows like $N^3$ while the
positive nonlinear production grows only like $N^{3/2}$.  High-frequency
isolated triads are therefore viscosity dominated at fixed critical record,
even though their positive production is unbounded.  The corollary rules out
a pointwise bound for $\mathcal G_+$; it does not rule out an integrated
parabolic-time estimate for the excess
$\mathcal G-\nu\|Au\|_2^2$.
\end{firewall}

\begin{corollary}[The canonical cutoff also lies inside this triad]
\label{cor:v5-homochiral-canonical-cutoff}
For equal amplitudes,
\[
 K_{\rm can}
 =\frac{\|A^{1/4}u_{N,L}\|_2^2}{\|u_{N,L}\|_2^2}
 =\frac{1+\sqrt2+\sqrt3}{3}N,
 \qquad
 N<K_{\rm can}<\sqrt2N.
 \label{eq:v5-homochiral-canonical-cutoff}
\]
Hence the record-based cutoff separates the lowest atom from the other two.
The tangent transfer \eqref{eq:v5-homochiral-transfer} moves on the joint
energy/critical-record level set while crossing that canonical cutoff.
\end{corollary}

% =============================================================================
\section{The V33--V34 variation card has a different mathematical type}
\label{sec:v5-variation-audit}
% =============================================================================

The V4 stopping line asked whether the inherited intrinsic time or temporal
variation already pays the downward spectral-energy variation.  The exact
definitions show that it does not.

\begin{proposition}[Source-residual variation is not spectral-stock variation]
\label{prop:v5-two-variations}
On a V33--V34 card, the inherited quantity is
\begin{equation}
 \mathfrak W_n^{\rm occ}
 =\frac{\tau_n^2
   \|\partial_tq_n\|_{L^2(I_n;L_x^2)}^2}
  {\|q_n\|_{L^2(I_n;L_x^2)}^2},
 \qquad
 q_n=(I-P_{{\rm hist},n})
 A^{-1/4}Q_{K_n}\BNS(P_{K_n}u).
 \label{eq:v5-source-residual-variation}
\end{equation}
The quantity required by V4 is instead
\begin{equation}
 \operatorname{Var}_{I_n}^-(e_\rho)
 =\int_{I_n}
 \left(
 2\nu\|A^{1/2}P_\rho u\|_2^2+2\Phi_\rho
 \right)_+\,\dd t,
 \qquad
 e_\rho=\|P_\rho u\|_2^2.
 \label{eq:v5-spectral-stock-variation}
\end{equation}
No identity in V33--V34 maps
\eqref{eq:v5-source-residual-variation} to
\eqref{eq:v5-spectral-stock-variation}.  Specifically:
\begin{enumerate}[label=\textup{(\roman*)},leftmargin=3em]
\item $q_n$ is one history-fresh exterior head--head source residual at the
      single cutoff $K_n$; $\Phi_\rho$ is the assembled incoming energy
      current at every nested cutoff $\rho$.
\item The derivative of $q_n$ is linear in the time derivative of a source
      vector, whereas $e_\rho'$ is a quadratic stock derivative containing
      viscosity and the complete low/high flux.
\item The V34 material cell localizes an enstrophy packet in physical space at
      a selected time.  It supplies no temporal crossing count or total
      variation for $e_\rho$.
\end{enumerate}
Consequently the V33--V34 temporal-variation gate does not close the V4
positive-recycling term without a new proved PDE estimate.
\end{proposition}

\begin{proof}
Equation \eqref{eq:v5-source-residual-variation} is the V34 intrinsic
variation definition together with its history short.  Equation
\eqref{eq:v5-spectral-stock-variation} is the low-frequency energy balance
\eqref{eq:v4-low-energy-balance}.  Their domains, codomains, and nonlinear
contents give \textup{(i)--(ii)}.  The inherited material-cell conclusion is
a same-time spatial inside/outside alternative for the enstrophy pulse, which
proves \textup{(iii)}.  Finally,
Theorem~\ref{thm:v3-incoming-counterexample} supplies a pointwise
independence test: the raw exterior daughter can vanish while incoming head
production is positive.  Thus no domination of the complete spectral flux
by the raw/history-fresh source row follows from the existing definitions.
\end{proof}

\begin{assessment}[What terminal chronology cannot do by itself]
A doubled-record corridor constrains one scalar moment
$\|A^{1/4}u\|_2^2$.  The homochiral triad supplies a tangent direction which
preserves that moment and energy while changing $H^1$.  Therefore counting
record crossings or knowing that the record is instantaneously decreasing
does not bound positive spectral recycling.  The viscosity-dominance
firewall leaves a narrower live possibility: an estimate on the
\emph{time-integrated excess production} over the inherited parabolic card.
That estimate must use viscosity and the multiscale interaction structure,
not the scalar record alone.
\end{assessment}

% =============================================================================
\section{The isolated-triad gain and the remaining multiscale branch}
\label{sec:v5-viscous-triad}
% =============================================================================

\begin{proposition}[An isolated high homochiral triad is dissipative under an energy cap]
\label{prop:v5-isolated-triad-dissipative}
For the equal-amplitude family above, there is a fixed $C_{\rm h}>0$ such
that
\[
 \mathcal G_1(u_{N,L})
 =C_{\rm h}N^3L^3,
 \qquad
 \nu\|Au_{N,L}\|_2^2=28\nu N^4L^2.
\]
If $\|u_{N,L}\|_2^2=6L^2\le E_*$, then
\begin{equation}
 \frac{\mathcal G_1(u_{N,L})}
      {\nu\|Au_{N,L}\|_2^2}
 \le
 \frac{C_{\rm h}\sqrt{E_*}}
      {28\sqrt6\,\nu N}
 \longrightarrow0.
 \label{eq:v5-isolated-triad-ratio}
\end{equation}
Thus this active record-invisible triad cannot by itself sustain positive
net $H^1$ growth at arbitrarily high frequency under the Leray energy cap.
\end{proposition}

\begin{proof}
The nonlinear scaling follows from
\eqref{eq:v5-triad-H1-nonzero} and the common cubic amplitude factor.
For palinstrophy, the three radii are
$N,\sqrt2N,\sqrt3N$, so reality pairing gives
\[
 2N^4L^2(1+4+9)=28N^4L^2.
\]
Insert $L\le\sqrt{E_*/6}$.
\end{proof}

\begin{firewall}[What Proposition~\ref{prop:v5-isolated-triad-dissipative} excludes]
It excludes only one isolated comparable-frequency triad.  A full solution
contains many simultaneous triads and, crucially, low-frequency strain
acting on high-frequency vorticity.  The number and coherent assembly of
those interactions are not bounded by applying the proposition triad by
triad.  Summing absolute triad estimates would also destroy the signed flux
assembly proved in V4.
\end{firewall}

% =============================================================================
\section{Integrated V5 conclusion and next acquisition order}
\label{sec:v5-conclusion}
% =============================================================================

\begin{theorem}[What V5 proves]
\label{thm:v5-conclusion}
Relative to V1--V4 and the frozen source, the following are proved.
\begin{enumerate}[label=\textup{(V5.\arabic*)},leftmargin=5em]
\item The critical record and $H^1$ production are respectively the zeroth
      and first radial moments of the same signed cumulative flux.
\item An explicit active homochiral triad preserves energy and the critical
      stock nonlinearly while changing $H^1$ with selectable sign.
\item At fixed critical record and under any prescribed energy cap, positive
      nonlinear $H^1$ production is pointwise unbounded.
\item The V33--V34 variation is variation of a source residual, not variation
      of nested spectral kinetic stocks; the material-cell row does not
      bridge the types.
\item Viscosity dominates the isolated high-frequency homochiral triad under
      the energy cap, so a surviving terminal mechanism must be collective
      or multiscale.
\end{enumerate}
No full NS stability theorem is obtained.
\end{theorem}

\begin{proof}
Item \textup{(V5.1)} is
Theorem~\ref{thm:v5-fractional-layer-cake}.  Items
\textup{(V5.2)--(V5.3)} are
Theorem~\ref{thm:v5-record-invisible-triad} and
Corollary~\ref{cor:v5-fixed-record-unbounded-production}.  Item
\textup{(V5.4)} is Proposition~\ref{prop:v5-two-variations}.  Item
\textup{(V5.5)} is Proposition~\ref{prop:v5-isolated-triad-dissipative}.
\end{proof}

\begin{programme}[V6 acquisition order]
\label{prog:v6}
\begin{enumerate}[label=\textup{(S\arabic*)},leftmargin=4em]
\item Replace the isolated-triad test by an assembly-preserving dyadic
      decomposition of
      $\mathcal G=-\langle\BNS(u),Au\rangle$.  Separate low--high strain,
      comparable-frequency transfer, and high--high-to-low return.
\item Keep viscosity inside each block estimate.  The target is the positive
      excess
      \[
       \bigl(\mathcal G-\theta\nu\|Au\|_2^2\bigr)_+
      \]
      for a fixed $0<\theta<1$, not $\mathcal G_+$ alone.
\item Test whether comparable high-frequency interactions are absorbable
      under the Leray energy cap without summing triads absolutely.  Record
      any shell-count or coherence loss explicitly.
\item Isolate the remaining low-frequency strain coefficient acting on
      high-frequency vorticity.  Compare its time integral on the V34 card
      with the record, intrinsic time, and existing material-cell data.
\item If this coefficient requires an $L_t^1L_x^\infty$ vorticity/strain
      bound or an equivalent continuation criterion, mark the circle and
      move further upstream rather than relabeling it as a paid budget.
\item Keep the represented-history branch at the V3 physical-incidence
      stopping line.
\end{enumerate}
\end{programme}

\begin{assessment}[Position after V5]
The scalar record cannot control the weighted flux: there are genuine
energy/record-tangent nonlinear directions with positive $H^1$ transfer.
Viscosity removes the isolated high-frequency example under a fixed energy
cap, so the unresolved branch is now sharply multiscale.  The next proof
attempt must retain signed shell assembly while identifying whether
low-frequency strain can drive high-frequency vorticity faster than
viscosity on the inherited parabolic interval.  This is narrower than the
V4 variation question and does not introduce a new history or compactness
object.
\end{assessment}

% =============================================================================
\section*{Version V5 history}
\addcontentsline{toc}{section}{Version V5 history}
% =============================================================================

Appended the fractional shell-flux identity and proved that the critical
record and enstrophy production are different radial moments of one signed
flux.  Constructed and activated a distinct-radius positive-helicity triad;
its nonlinear motion preserves energy and the critical record but changes
$H^1$, with positive production unbounded at fixed record and bounded
energy.  Audited the V33--V34 source-residual variation against the required
spectral-stock variation and proved the existing material-cell row does not
supply the missing map.  Finally proved that viscosity dominates the
isolated high-frequency triad under the Leray energy cap, leaving a genuinely
collective/multiscale excess-production problem.  No global regularity claim
was made.
% Future versions are appended immediately above the sole terminator.
% =============================================================================
% VERSION V6 APPEND
% =============================================================================

\section{V6: an assembly-preserving dissipation-wavenumber reduction}
\label{sec:v6-dissipation-index}
% =============================================================================

V5 left one precise multiscale question: after viscosity is retained inside
the estimate, which part of the nonlinear enstrophy production is not
absorbed?  This section answers that question without estimating Fourier
triads one at a time.  The answer is a low-mode strain tariff.  We then audit
that tariff against the V34 card data.

Throughout V6, Haar measure on $\Torus^3$ is normalized to one.  For
$q\ge0$, put $\lambda_q=2^q$ and let $\Delta_q$ be the sharp Fourier
projection to
\begin{equation}
 \{k\in\mathbb Z^3:\lambda_q\le |k|<2\lambda_q\}.
 \label{eq:v6-sharp-shell}
\end{equation}
Because $u$ has mean zero,
\begin{equation}
 u=\sum_{q\ge0}u_q,\qquad u_q:=\Delta_qu .
 \label{eq:v6-shell-resolution}
\end{equation}
Sharp shells are used only for the $L^2$ energy calculation and the
finite-lattice Bernstein inequality; no $L^p$ multiplier theorem is needed.
Set
\begin{equation}
 r_q:=\lambda_q^{1/2}\|u_q\|_2,\qquad
 e_q:=\lambda_q\|u_q\|_2,\qquad
 d_q:=\lambda_q^2\|u_q\|_2.
 \label{eq:v6-shell-stocks}
\end{equation}
There are absolute periodic constants such that
\begin{equation}
 \|A^{1/4}u\|_2^2\asymp\sum_qr_q^2,\qquad
 \|A^{1/2}u\|_2^2\asymp\sum_qe_q^2,\qquad
 \|Au\|_2^2\asymp\sum_qd_q^2.
 \label{eq:v6-shell-equivalence}
\end{equation}
The lattice-point Bernstein estimate is
\begin{equation}
 \|u_q\|_\infty\le C_{\rm B}\lambda_q^{3/2}\|u_q\|_2.
 \label{eq:v6-Bernstein}
\end{equation}

Fix $0<\theta<1$.  A universal number $c_\theta>0$ will be chosen in
Theorem~\ref{thm:v6-dyadic-absorption}.  At each smooth time define
\begin{equation}
 Q_\nu(t):=\min\left\{q\in\{-1,0,1,\ldots\}:
 \lambda_p^{-1}\|u_p(t)\|_\infty<c_\theta\nu
 \ \hbox{for every }p>q\right\}.
 \label{eq:v6-dissipation-index}
\end{equation}
Smoothness makes this integer finite.  The convention $Q_\nu=-1$ means that
every shell satisfies the smallness condition.  Define the retained
low-mode strain tariff
\begin{equation}
 \boxed{\quad
 \mathfrak f_\nu(t):=
 \sum_{0\le q\le Q_\nu(t)}
       \lambda_q\|u_q(t)\|_\infty ,
 \quad}
 \label{eq:v6-low-strain-tariff}
\end{equation}
with value zero when $Q_\nu=-1$.

\begin{lemma}[Exact cubic assembly and the two-largest-frequency rule]
\label{lem:v6-cubic-assembly}
For every smooth divergence-free $u$,
\begin{equation}
 \langle\BNS(u),Au\rangle
 =\sum_{i,j,k=1}^3
   \int_{\Torus^3}
   (\partial_ku_j)(\partial_ju_i)(\partial_ku_i)\,\dd x .
 \label{eq:v6-cubic-gradient-identity}
\end{equation}
After inserting \eqref{eq:v6-shell-resolution}, the resulting series is
absolutely convergent.  If three shell indices are rearranged as
$a\le b\le c$, a nonzero interaction has
\begin{equation}
 c-b\le3.
 \label{eq:v6-two-largest-comparable}
\end{equation}
Consequently, with
$\ell_a:=\lambda_a\|u_a\|_\infty$,
\begin{equation}
 \left|\langle\BNS(u),Au\rangle\right|
 \le C
 \sum_{\substack{a\le b\le c\\c-b\le3}}
       \ell_a e_be_c.
 \label{eq:v6-assembled-shell-majorant}
\end{equation}
The constant includes the finitely many placements of the three ordered
indices in \eqref{eq:v6-cubic-gradient-identity}.
\end{lemma}

\begin{proof}
Since $Au=-\Delta u$ on divergence-free fields and the Leray projection is
orthogonal,
\[
 \langle\BNS(u),Au\rangle
 =\int (u_j\partial_ju_i)(-\partial_k^2u_i).
\]
Integrating once in $x_k$ gives the term on the right of
\eqref{eq:v6-cubic-gradient-identity} and
\[
 \int u_j(\partial_j\partial_ku_i)(\partial_ku_i)
 =\frac12\int u_j\partial_j|\partial_ku_i|^2=0.
\]
Smoothness justifies the shell expansion absolutely.

For the frequency rule, three Fourier frequencies in a nonzero integral
sum to zero.  If the largest shell were separated from the other two by
more than three dyadic steps, the magnitude of its frequency would exceed
the sum of the magnitudes of the other two, which is impossible.  Thus
\eqref{eq:v6-two-largest-comparable} holds.  Put the gradient from the
lowest shell in $L^\infty$ and the other two in $L^2$.  Bernstein on a
single annulus gives
\[
 \left|\int \partial u_a\,\partial u_b\,\partial u_c\right|
 \le C\lambda_a\|u_a\|_\infty
       \lambda_b\|u_b\|_2\lambda_c\|u_c\|_2.
\]
Summing the finitely many tensor components and placements proves
\eqref{eq:v6-assembled-shell-majorant}.
\end{proof}

\begin{theorem}[High-shell assembly is absorbed by viscosity]
\label{thm:v6-dyadic-absorption}
There is a universal $C_{\rm d}$ such that
\begin{equation}
 \boxed{
 \mathcal G(t):=-\langle\BNS(u),Au\rangle
 \le C_{\rm d}\mathfrak f_\nu(t)\|A^{1/2}u(t)\|_2^2
    +C_{\rm d}c_\theta\nu\|Au(t)\|_2^2.}
 \label{eq:v6-dyadic-absorption}
\end{equation}
Choose $c_\theta\le\theta/(2C_{\rm d})$, with the harmless factor two
covering the norm equivalences in \eqref{eq:v6-shell-equivalence}.  Then
\begin{equation}
 \boxed{
 \bigl(\mathcal G(t)-\theta\nu\|Au(t)\|_2^2\bigr)_+
 \le C_{\rm d}\mathfrak f_\nu(t)\|A^{1/2}u(t)\|_2^2.}
 \label{eq:v6-positive-excess}
\end{equation}
Thus comparable high-frequency interactions and high--high-to-low returns
whose lowest participating shell lies above $Q_\nu$ are absorbed as one
assembled block.  No number of individual triads occurs in the bound.
\end{theorem}

\begin{proof}
Apply Lemma~\ref{lem:v6-cubic-assembly} and split the lowest index $a$ at
$Q_\nu$.  For $a\le Q_\nu$,
\[
 \sum_{\substack{a\le b\le c,\ c-b\le3\\a\le Q_\nu}}
      \ell_ae_be_c
 \le \mathfrak f_\nu
      \sum_{c-b\le3}e_be_c
 \le C\mathfrak f_\nu\sum_qe_q^2.
\]
For $a>Q_\nu$, definition \eqref{eq:v6-dissipation-index} gives
\[
 \ell_a=\lambda_a\|u_a\|_\infty
       \le c_\theta\nu\lambda_a^2.
\]
The inner sum is geometric:
\[
 \sum_{Q_\nu<a\le b}\ell_a
 \le Cc_\theta\nu\lambda_b^2.
\]
Since $|b-c|\le3$,
\[
 \sum_{c-b\le3}\lambda_b^2e_be_c
 \le C\sum_{c-b\le3}d_bd_c
 \le C\sum_qd_q^2.
\]
Combining the low and high parts with
\eqref{eq:v6-shell-equivalence} proves
\eqref{eq:v6-dyadic-absorption}.  The chosen value of $c_\theta$ absorbs
the last term into $\theta\nu\|Au\|_2^2$, and taking the positive part gives
\eqref{eq:v6-positive-excess}.
\end{proof}

\begin{firewall}[Meaning of assembly-preserving]
The exact signed cubic expression is first grouped by complete dyadic
blocks.  Absolute values are taken only after the two-largest-frequency
constraint is imposed.  This retains the geometric shell assembly needed
for viscous absorption and avoids an absolute sum over the number of
Fourier triads.  It does not claim an additional cancellation among the
remaining low--high blocks.
\end{firewall}

% =============================================================================
\section{The exact continuation tariff}
\label{sec:v6-continuation-tariff}
% =============================================================================

\begin{theorem}[Dissipation-index continuation criterion]
\label{thm:v6-low-mode-continuation}
Let $u$ be smooth on $[0,T_*)\times\Torus^3$, where $T_*<\infty$.  If
\begin{equation}
 \int_{t_0}^{T_*}\mathfrak f_\nu(t)\,\dd t<\infty
 \label{eq:v6-tariff-integrable}
\end{equation}
for some $t_0<T_*$, then $u$ extends as a strong solution beyond $T_*$.  In
particular every finite-time singular solution must satisfy
\begin{equation}
 \boxed{\int_{t_0}^{T_*}\mathfrak f_\nu(t)\,\dd t=\infty
 \quad\hbox{for every }t_0<T_*.}
 \label{eq:v6-tariff-necessary}
\end{equation}
\end{theorem}

\begin{proof}
Take the inner product of \eqref{eq:NS} with $Au$ and use
Theorem~\ref{thm:v6-dyadic-absorption}.  With
\[
 Y=\|A^{1/2}u\|_2^2,\qquad Z=\|Au\|_2^2,
\]
one obtains
\begin{equation}
 \frac12Y'+(1-\theta)\nu Z
 \le C_{\rm d}\mathfrak f_\nu Y.
 \label{eq:v6-low-mode-Gronwall}
\end{equation}
Assumption \eqref{eq:v6-tariff-integrable} and Gronwall bound $Y$ uniformly
up to $T_*$.  Integrating \eqref{eq:v6-low-mode-Gronwall} also bounds
$\int_{t_0}^{T_*}Z$.  The periodic $H^1$ local strong-solution theorem,
started at times tending to $T_*$, then extends the solution.  The
contrapositive gives \eqref{eq:v6-tariff-necessary}.
\end{proof}

\begin{assessment}[The multiscale bottleneck is now one scalar]
Theorem~\ref{thm:v6-dyadic-absorption} completes steps \textup{(S1)--(S3)}
of Programme~\ref{prog:v6}.  The unresolved nonlinearity is no longer a
list of comparable triads.  It is exactly the time integral of
$\mathfrak f_\nu$, the low-frequency strain seen from the highest shell
which has not yet become viscosity-small.
\end{assessment}

% =============================================================================
\section{Critical-record and energy ledgers for the dissipation index}
\label{sec:v6-index-ledgers}
% =============================================================================

Write
\begin{equation}
 \mathcal R_{\rm d}(t):=
 \left(\sum_{q\ge0}\lambda_q\|u_q(t)\|_2^2\right)^{1/2}
 \asymp\|A^{1/4}u(t)\|_2.
 \label{eq:v6-dyadic-record}
\end{equation}

\begin{proposition}[Pointwise critical-record tariff]
\label{prop:v6-record-tariff}
At every smooth time with $Q_\nu(t)\ge0$,
\begin{equation}
 \boxed{
 \mathfrak f_\nu(t)
 \le C\mathcal R_{\rm d}(t)\lambda_{Q_\nu(t)}^2.}
 \label{eq:v6-record-tariff}
\end{equation}
If $Q_\nu(t)\ge0$, its highest shell is active:
\begin{equation}
 \lambda_{Q_\nu}^{-1}\|u_{Q_\nu}\|_\infty
 \ge c_\theta\nu,
 \qquad
 r_{Q_\nu}\ge\frac{c_\theta}{C_{\rm B}}\nu.
 \label{eq:v6-highest-active}
\end{equation}
\end{proposition}

\begin{proof}
By Bernstein and Cauchy--Schwarz,
\[
 \begin{aligned}
 \mathfrak f_\nu
 &\le C_{\rm B}\sum_{q\le Q_\nu}
       \lambda_q^{5/2}\|u_q\|_2
  =C_{\rm B}\sum_{q\le Q_\nu}\lambda_q^2r_q\\
 &\le C_{\rm B}
      \left(\sum_{q\le Q_\nu}\lambda_q^4\right)^{1/2}
      \left(\sum_{q\le Q_\nu}r_q^2\right)^{1/2}
 \le C\lambda_{Q_\nu}^2\mathcal R_{\rm d}.
 \end{aligned}
\]
Minimality in \eqref{eq:v6-dissipation-index} says that the defining
condition fails for $Q_\nu-1$ but holds above $Q_\nu$; hence it fails
precisely at $Q_\nu$.  The first inequality in
\eqref{eq:v6-highest-active} follows.  The second follows from
\eqref{eq:v6-Bernstein}.
\end{proof}

\begin{proposition}[The Leray budget pays one frequency power]
\label{prop:v6-index-occupation}
Let
\[
 E_0:=\|u(0)\|_2^2.
\]
For every smooth interval $[0,T]$,
\begin{equation}
 \boxed{
 \int_{\{t\le T:Q_\nu(t)\ge0\}}
      \lambda_{Q_\nu(t)}\,\dd t
 \le C\frac{E_0}{\nu^3}.}
 \label{eq:v6-index-L1-budget}
\end{equation}
Equivalently,
\begin{equation}
 \left|\{t\le T:\lambda_{Q_\nu(t)}\ge L\}\right|
 \le C\frac{E_0}{\nu^3L}
 \qquad(L\ge1).
 \label{eq:v6-index-weak-L1}
\end{equation}
\end{proposition}

\begin{proof}
By \eqref{eq:v6-highest-active},
\[
 \|A^{1/2}u(t)\|_2^2
 \ge c\lambda_{Q_\nu}^2\|u_{Q_\nu}\|_2^2
 =c\lambda_{Q_\nu}r_{Q_\nu}^2
 \ge c\nu^2\lambda_{Q_\nu}
\]
whenever $Q_\nu\ge0$.  Integrate and use the Leray energy identity
\[
 2\nu\int_0^T\|A^{1/2}u(t)\|_2^2\,\dd t\le E_0.
\]
This proves \eqref{eq:v6-index-L1-budget}; Chebyshev gives
\eqref{eq:v6-index-weak-L1}.
\end{proof}

\begin{firewall}[The missing exponent]
The paid quantity in \eqref{eq:v6-index-L1-budget} is
$\int\lambda_{Q_\nu}$.  The pointwise tariff
\eqref{eq:v6-record-tariff} contains
$\mathcal R_{\rm d}\lambda_{Q_\nu}^2$.  Thus the energy identity is short
by one full frequency power and also does not remove the critical-record
factor.  Replacing the latter by its value on a V34 corridor does not fix
the exponent.
\end{firewall}

% =============================================================================
\section{Sharpness and the failed identification with the V34 cutoff}
\label{sec:v6-sharpness}
% =============================================================================

The preceding frequency loss is not merely an artifact of applying
Bernstein and Cauchy--Schwarz.  The following explicit fields saturate it.

\begin{lemma}[A one-shell divergence-free Bernstein saturator]
\label{lem:v6-shell-saturator}
There are constants $c,C>0$ with the following property.  For every dyadic
integer $N\ge16$ and every $R>0$, there is a real, mean-zero,
divergence-free trigonometric polynomial $W_{N,R}$ supported in
\[
 \{k\in\mathbb Z^3:N\le |k|<2N\}
\]
such that
\begin{align}
 \|A^{1/4}W_{N,R}\|_2&=R,
 \label{eq:v6-saturator-record}\\
 c\frac{R^2}{N}
 \le\|W_{N,R}\|_2^2&\le C\frac{R^2}{N},
 \label{eq:v6-saturator-energy}\\
 \|\nabla W_{N,R}\|_\infty&\ge cRN^2,
 \label{eq:v6-saturator-strain}\\
 N^{-1}\|W_{N,R}\|_\infty&\ge cR.
 \label{eq:v6-saturator-active}
\end{align}
If $R\ge c_\theta\nu/c$, then its dissipation index is the shell of $N$
and
\begin{equation}
 \mathfrak f_\nu(W_{N,R})\ge cRN^2.
 \label{eq:v6-saturator-tariff}
\end{equation}
\end{lemma}

\begin{proof}
Let $\mathbf e_2=(0,1,0)$,
$P_k=I-k\otimes k/|k|^2$, and
\[
 \Lambda_N:=
 \left\{k\in\mathbb Z^3:
 N\le k_1\le\lfloor5N/4\rfloor,\quad
 1\le k_2,k_3\le\lfloor N/8\rfloor\right\}.
\]
Then $\#\Lambda_N\asymp N^3$, every $k\in\Lambda_N$ lies in the stated
annulus, and
\[
 (P_k\mathbf e_2)_2=1-\frac{k_2^2}{|k|^2}\ge\frac{63}{64}.
\]
Put
\[
 W_{N,R}(x):=
 a_{N,R}\sum_{k\in\Lambda_N}P_k\mathbf e_2\sin(k\cdot x),
\]
where the positive constant $a_{N,R}$ is chosen to make
\eqref{eq:v6-saturator-record} exact.  Each summand is divergence-free,
and orthogonality of distinct sine modes gives
\[
 R^2=\frac{a_{N,R}^2}{2}
       \sum_{k\in\Lambda_N}|k|\,|P_k\mathbf e_2|^2.
\]
Consequently $a_{N,R}\asymp R/N^2$, and the same orthogonality proves
\eqref{eq:v6-saturator-energy}.  At the origin,
\[
 \partial_1(W_{N,R})_2(0)
 =a_{N,R}\sum_{k\in\Lambda_N}
    k_1\left(1-\frac{k_2^2}{|k|^2}\right)
 \ge cRN^2,
\]
which proves \eqref{eq:v6-saturator-strain}.  Single-shell Bernstein gives
$\|\nabla W_{N,R}\|_\infty\le CN\|W_{N,R}\|_\infty$, yielding
\eqref{eq:v6-saturator-active}.  When the displayed lower bound exceeds
$c_\theta\nu$, the occupied shell fails the viscosity-smallness condition,
all higher shells vanish, and hence it is precisely the dissipation-index
shell.  Finally
\[
 \mathfrak f_\nu=N\|W_{N,R}\|_\infty
 \ge cRN^2.
\]
\end{proof}

\begin{corollary}[The canonical record--energy cutoff does not bound
the dissipation index]
\label{cor:v6-no-KQ-incidence}
Fix $E_*>0$, $R>c_\theta\nu/c$, and put $K=R^2/E_*$.  For every sufficiently
large dyadic $N$, the field in
Lemma~\ref{lem:v6-shell-saturator} satisfies
\[
 \|W_{N,R}\|_2^2\le E_*,
 \qquad
 \|A^{1/4}W_{N,R}\|_2=R,
\]
while $\lambda_{Q_\nu}=N$.  Therefore the scalar energy cap and critical
record impose no upper bound of the form
\begin{equation}
 \lambda_{Q_\nu}\le C K.
 \label{eq:v6-false-scale-incidence}
\end{equation}
\end{corollary}

\begin{proof}
Use \eqref{eq:v6-saturator-energy} and choose
$N\ge CR^2/E_*$.  The number $N$ can then be made arbitrarily larger than
$K$.
\end{proof}

\begin{proposition}[Even favorable scale incidence does not pay the cards]
\label{prop:v6-card-sharpness}
There are doubled records $R_n=2^nR_0$, dyadic frequencies
$N_n=2^{2n}N_0$, positive widths
\begin{equation}
 \tau_n=\frac{\gamma}{\nu N_n^2},
 \label{eq:v6-model-card-width}
\end{equation}
and smooth divergence-free profiles $W_n=W_{N_n,R_n}$ such that, after
choosing $N_0$ large and $\gamma>0$ small,
\begin{align}
 \sup_n\|W_n\|_2^2&\le E_*,
 \label{eq:v6-model-energy-cap}\\
 K_n:=R_n^2/E_*&\asymp N_n,
 \label{eq:v6-model-favorable-incidence}\\
 \nu K_n^2\tau_n&\le4R_t,
 \label{eq:v6-model-parabolic-row}\\
 \sum_{n\ge0}R_n\sqrt{\tau_n}&<\infty,
 \label{eq:v6-model-core-sum}\\
 \sum_{n\ge0}\tau_nN_n&<\infty,
 \label{eq:v6-model-paid-index}\\
 \sum_{n\ge0}\tau_n\mathfrak f_\nu(W_n)&=\infty.
 \label{eq:v6-model-unpaid-tariff}
\end{align}
Thus even the additional favorable hypothesis
$\lambda_{Q_\nu}\asymp K_n$ does not convert the V34 parabolic and
intrinsic-core arithmetic into the continuation tariff.
\end{proposition}

\begin{proof}
Choose $R_0$ above the activity threshold in
Lemma~\ref{lem:v6-shell-saturator}, and choose the dyadic $N_0$ so large
that $CR_0^2/N_0\le E_*$.  Then
\[
 \|W_n\|_2^2\le C\frac{R_n^2}{N_n}
 =C\frac{R_0^2}{N_0}\le E_*,
\]
and $N_n/K_n=N_0E_*/R_0^2$ is independent of $n$.  Choose $\gamma$ so that
$\gamma(K_n/N_n)^2\le4R_t$; this proves
\eqref{eq:v6-model-parabolic-row}.  Next,
\[
 \sum_nR_n\sqrt{\tau_n}
 =\sqrt{\frac{\gamma}{\nu}}\frac{R_0}{N_0}
   \sum_n2^{-n}<\infty,
\]
and
\[
 \sum_n\tau_nN_n
 =\frac{\gamma}{\nu N_0}\sum_n4^{-n}<\infty.
\]
On the other hand, \eqref{eq:v6-saturator-tariff} gives
\[
 \sum_n\tau_n\mathfrak f_\nu(W_n)
 \ge\frac{c\gamma}{\nu}\sum_nR_n=\infty.
\]
\end{proof}

\begin{firewall}[Kinematic sharpness is not an NS counterexample]
The profiles in Proposition~\ref{prop:v6-card-sharpness} are held constant
on abstract intervals only to test implication among scalar inequalities.
They are not asserted to be consecutive pieces of one Navier--Stokes
solution, and they do not instantiate the V34 source-residual action or
material-cell hypotheses.  The proposition proves the precise negative
statement needed here: the displayed record, energy, scale-incidence,
parabolic-width, and core-summability rows alone cannot pay
\eqref{eq:v6-tariff-integrable}.  Any positive argument must use an
additional dynamical bridge.
\end{firewall}

% =============================================================================
\section{Upstream comparison with the inherited V34 card}
\label{sec:v6-v34-audit}
% =============================================================================

\begin{proposition}[The two cutoffs and the missing bridge]
\label{prop:v6-two-cutoffs}
On a V34 record card, the inherited cutoff
\[
 K_n=\mathcal R_n^2/E_*
\]
and the dissipation wavenumber $\lambda_{Q_\nu(t)}$ have different
definitions and different proven incidences.
\begin{enumerate}[label=\textup{(\alph*)},leftmargin=3em]
\item $K_n$ is a fixed scalar interpolation scale selected from the record
      and energy cap.  It defines the head $P_{K_n}u$ and the first-exterior
      daughter.
\item $Q_\nu(t)$ is a time-dependent velocity-shell threshold selected by
      comparison of $\lambda_p^{-1}\|u_p(t)\|_\infty$ with $\nu$.
\item Neither the corridor bound, the V34 parabolic row, nor the
      source-residual intrinsic variation yields
      $\lambda_{Q_\nu}\lesssim K_n$.
\item Even if that missing scale incidence is added as a hypothesis,
      Proposition~\ref{prop:v6-card-sharpness} shows that the inherited
      scalar core sum does not pay $\int\mathfrak f_\nu$.
\end{enumerate}
Therefore the V34 card does not presently imply the continuation hypothesis
of Theorem~\ref{thm:v6-low-mode-continuation}.
\end{proposition}

\begin{proof}
Items (a) and (b) are the definitions.  Item (c) follows abstractly from
Corollary~\ref{cor:v6-no-KQ-incidence}; it is also consistent with the type
audit in Proposition~\ref{prop:v5-two-variations}, because the V34
variation acts on the source residual $q_n$, not on velocity shells.
Item (d) is Proposition~\ref{prop:v6-card-sharpness}.
\end{proof}

\begin{assessment}[What would be circular]
Invoking
\[
 \int_0^{T_*}\|\nabla u_{\le Q_\nu(t)}(t)\|_\infty\,\dd t<\infty
\]
would pay \eqref{eq:v6-tariff-integrable}, but it is already a low-mode
Lipschitz continuation condition of the same strength exposed by
Theorem~\ref{thm:v6-low-mode-continuation}.  It cannot be entered as a
consequence of the existing source record.  The missing estimate must
instead control the residence or motion of the highest viscosity-active
\emph{velocity} shell, or prove a genuinely stronger cancellation in the
low--high assembly.
\end{assessment}

% =============================================================================
\section{Integrated V6 conclusion and next acquisition order}
\label{sec:v6-conclusion}
% =============================================================================

\begin{theorem}[What V6 proves]
\label{thm:v6-conclusion}
Relative to V1--V5 and the frozen source, the following are proved.
\begin{enumerate}[label=\textup{(V6.\arabic*)},leftmargin=5em]
\item The full nonlinear enstrophy production has an exact dyadic
      decomposition in which every interaction above the dissipation index
      is absorbed by a prescribed fraction of viscosity.
\item The entire positive excess is bounded by one low-mode strain tariff,
      and finite time integral of that tariff is a continuation criterion.
\item The critical record bounds the tariff by
      $C\mathcal R_{\rm d}\lambda_{Q_\nu}^2$.
\item The Leray energy identity pays only
      $\int\lambda_{Q_\nu}$, leaving one frequency power and the record
      factor unpaid.
\item Energy and critical record do not compare $Q_\nu$ with the canonical
      V34 cutoff.  Even imposing favorable comparison by hand does not make
      the V34 scalar core sum pay the tariff.
\end{enumerate}
No full NS stability theorem is obtained.
\end{theorem}

\begin{proof}
Item \textup{(V6.1)} is
Theorem~\ref{thm:v6-dyadic-absorption}.  Item \textup{(V6.2)} is
\eqref{eq:v6-positive-excess} and
Theorem~\ref{thm:v6-low-mode-continuation}.  Items
\textup{(V6.3)--(V6.4)} are
Propositions~\ref{prop:v6-record-tariff} and
\ref{prop:v6-index-occupation}.  Item \textup{(V6.5)} is
Corollary~\ref{cor:v6-no-KQ-incidence},
Proposition~\ref{prop:v6-card-sharpness}, and
Proposition~\ref{prop:v6-two-cutoffs}.
\end{proof}

\begin{programme}[V7 acquisition order]
\label{prog:v7}
\begin{enumerate}[label=\textup{(T\arabic*)},leftmargin=4em]
\item Work on the actual solution, not the kinematic saturators.  Decompose
      time into measurable occupation sets on which the highest
      viscosity-active velocity shell is fixed.
\item Use the exact shell equation
      \[
       (\partial_t+\nu A)u_q=-\Delta_q\BNS(u)
      \]
      to seek a residence-time or crossing-cost estimate for an active
      shell.  The target must gain the missing frequency power beyond
      \eqref{eq:v6-index-L1-budget}.
\item Compare the resulting velocity-shell motion with the V34
      heat-covariant daughter only through an explicit operator connecting
      $\Delta_qu$ to $A^{-1/4}Q_{K_n}\BNS(P_{K_n}u)$.  Do not identify their
      intrinsic variations by scale notation alone.
\item On the branch $\lambda_{Q_\nu}\gg K_n$, retain the actual exterior
      velocity shell and quantify its critical mass and energy-dissipation
      debit.  On the branch $\lambda_{Q_\nu}\lesssim K_n$, test whether the
      V34 contiguous core supplies enough residence information to control
      $\int\lambda_{Q_\nu}^2$.
\item If the crossing estimate returns $\int\mathfrak f_\nu$ or an
      $L_t^1L_x^\infty$ strain norm on its right-hand side, record the
      continuation circle and move upstream to signed low--high flux or
      geometric depletion.
\item Keep the represented-history branch at the V3 physical-incidence
      stopping line.
\end{enumerate}
\end{programme}

\begin{assessment}[Position after V6]
Comparable-frequency coherence is not the remaining obstruction: its full
dyadic assembly is viscosity-absorbable above a dynamically selected
threshold.  The unresolved mechanism is the residence-weighted low strain
below that threshold.  V6 supplies both a paid first-power occupation
budget and a sharp proof that the continuation estimate needs more.  The
next noncircular acquisition is therefore a dynamical velocity-shell
crossing cost, not another static record interpolation and not a relabeling
of the V34 source-residual variation.
\end{assessment}

% =============================================================================
\section*{Version V6 history}
\addcontentsline{toc}{section}{Version V6 history}
% =============================================================================

Appended an exact dyadic cubic-gradient assembly for nonlinear enstrophy
production.  Proved that all blocks whose lowest frequency lies above the
viscosity-selected dissipation index are absorbed by viscosity, leaving one
low-mode strain tariff, and proved the corresponding continuation
criterion.  Derived the sharp critical-record upper bound for that tariff
and the global first-frequency-power occupation budget paid by the Leray
energy identity.  Constructed explicit one-shell divergence-free
Bernstein saturators, showing both that the canonical record--energy cutoff
does not bound the dissipation index and that even favorable cutoff
incidence plus the V34 parabolic/core arithmetic does not pay the tariff.
The next required object is a dynamical residence or crossing cost for the
highest active velocity shell.  No global regularity claim was made.
% Future versions are appended immediately above the sole terminator.
% =============================================================================
% VERSION V7 APPEND
% =============================================================================

\section{V7: the actual velocity-shell equation}
\label{sec:v7-shell-equation}
% =============================================================================

V6 proved that the only unabsorbed enstrophy-production coefficient is
$\mathfrak f_\nu$ and that the Leray budget pays only one power of the
dissipation wavenumber.  Programme~\ref{prog:v7} asks whether the actual
Navier--Stokes shell equation supplies the missing power.  We now answer
that question exactly.

Retain the sharp shells, $\lambda_q$, $u_q$, $r_q$, and $Q_\nu$ from V6.
Put
\begin{equation}
 z_q(t):=r_q(t)^2=\lambda_q\|u_q(t)\|_2^2
 \label{eq:v7-shell-critical-stock}
\end{equation}
and define the critical and weak nonlinear shell sources
\begin{equation}
 g_q:=A^{-1/4}\Delta_q\BNS(u),
 \qquad
 h_q:=A^{-3/4}\Delta_q\BNS(u).
 \label{eq:v7-two-shell-sources}
\end{equation}
All operators commute because they are Fourier multipliers, and hence
\begin{equation}
 \boxed{g_q=A^{1/2}h_q,\qquad
 \|g_q\|_2^2\asymp\lambda_q^2\|h_q\|_2^2.}
 \label{eq:v7-source-derivative-gap}
\end{equation}

\begin{lemma}[Critical shell differential inequality]
\label{lem:v7-shell-differential}
There are universal constants $c_0,C_0>0$ such that every smooth solution
satisfies
\begin{equation}
 \boxed{
 \frac12z_q'(t)+c_0\nu\lambda_q^2z_q(t)
 \le \frac{C_0}{\nu}\|g_q(t)\|_2^2}
 \label{eq:v7-shell-differential}
\end{equation}
for every $q\ge0$ and every smooth time.
\end{lemma}

\begin{proof}
The exact projected equation is
\begin{equation}
 (\partial_t+\nu A)u_q=-\Delta_q\BNS(u).
 \label{eq:v7-exact-shell-equation}
\end{equation}
Pair it with $\lambda_qu_q$.  Since the shell lies in
$\lambda_q\le A^{1/2}<2\lambda_q$,
\[
 \frac12z_q'
 +\nu\lambda_q\|A^{1/2}u_q\|_2^2
 =-\lambda_q\langle A^{1/4}g_q,u_q\rangle
\]
and
\[
 \nu\lambda_q\|A^{1/2}u_q\|_2^2
 \ge\nu\lambda_q^2z_q.
\]
The source term obeys
\[
 \lambda_q|\langle A^{1/4}g_q,u_q\rangle|
 \le C\lambda_q\|g_q\|_2z_q^{1/2}
 \le\frac{\nu}{2}\lambda_q^2z_q
      +\frac{C}{\nu}\|g_q\|_2^2.
\]
Absorb the first term and adjust universal constants.
\end{proof}

Let
\begin{equation}
 \sigma:=\frac{c_\theta}{C_{\rm B}}>0,
 \label{eq:v7-active-threshold}
\end{equation}
where $c_\theta$ and $C_{\rm B}$ are from
\eqref{eq:v6-dissipation-index} and \eqref{eq:v6-Bernstein}.  V6 proves
\begin{equation}
 Q_\nu(t)=q\ge0\quad\Longrightarrow\quad
 z_q(t)\ge\sigma^2\nu^2.
 \label{eq:v7-Q-implies-active}
\end{equation}

% =============================================================================
\section{Residence and upcrossing costs}
\label{sec:v7-crossing}
% =============================================================================

\begin{theorem}[Active-shell residence inequality]
\label{thm:v7-residence}
For an interval $I=[a,b]$, define
\begin{equation}
 E_q(I):=\{t\in I:z_q(t)\ge\sigma^2\nu^2\}.
 \label{eq:v7-active-set}
\end{equation}
Then
\begin{equation}
 \boxed{
 c\sigma^2\nu^3\lambda_q^2|E_q(I)|
 \le z_q(a)+\frac{C}{\nu}
       \int_{E_q(I)}\|g_q(t)\|_2^2\,\dd t.}
 \label{eq:v7-residence-one-shell}
\end{equation}
Consequently,
\begin{equation}
 \boxed{
 \int_{I\cap\{Q_\nu\ge0\}}\lambda_{Q_\nu(t)}^2\,\dd t
 \le
 \frac{C}{\sigma^2\nu^3}\mathcal R_{\rm d}(a)^2
 +\frac{C}{\sigma^2\nu^4}
   \int_I\|A^{-1/4}\BNS(u(t))\|_2^2\,\dd t.}
 \label{eq:v7-residence-summed}
\end{equation}
This is an actual-solution estimate and gains the frequency power missing
from \eqref{eq:v6-index-L1-budget}; its price is the critical nonlinear
source action.
\end{theorem}

\begin{proof}
Multiply \eqref{eq:v7-shell-differential} by the indicator of
$E_q(I)$ and integrate.  Since $z_q$ is absolutely continuous, the chain
rule for the positive part gives
\[
 \int_{E_q(I)}z_q'(t)\,\dd t
 =(z_q(b)-\sigma^2\nu^2)_+
  -(z_q(a)-\sigma^2\nu^2)_+
 \ge-z_q(a).
\]
On $E_q(I)$, $z_q\ge\sigma^2\nu^2$.  Inserting these two facts in
\eqref{eq:v7-shell-differential} proves
\eqref{eq:v7-residence-one-shell}.

Set $O_q(I)=\{t\in I:Q_\nu(t)=q\}$.  The sets $O_q(I)$ are disjoint and
\eqref{eq:v7-Q-implies-active} gives $O_q(I)\subset E_q(I)$.  Therefore
\[
 \int_{I\cap\{Q_\nu\ge0\}}\lambda_{Q_\nu}^2
 =\sum_q\lambda_q^2|O_q(I)|
 \le\sum_q\lambda_q^2|E_q(I)|.
\]
Sum \eqref{eq:v7-residence-one-shell}.  By
\eqref{eq:v6-dyadic-record},
$\sum_qz_q(a)=\mathcal R_{\rm d}(a)^2$.  Orthogonality of the sharp shells
and commutation with $A$ give
\[
 \sum_q\int_{E_q(I)}\|g_q\|_2^2
 \le\int_I\sum_q\|g_q\|_2^2
 =\int_I\|A^{-1/4}\BNS(u)\|_2^2.
\]
This proves \eqref{eq:v7-residence-summed}.
\end{proof}

\begin{corollary}[Record-corridor form of the residence estimate]
\label{cor:v7-corridor-tariff}
If
\[
 \sup_{t\in I}\mathcal R_{\rm d}(t)\le\overline{\mathcal R},
\]
then
\begin{equation}
 \boxed{
 \int_I\mathfrak f_\nu(t)\,\dd t
 \le
 \frac{C\overline{\mathcal R}}{\sigma^2\nu^3}
       \mathcal R_{\rm d}(a)^2
 +\frac{C\overline{\mathcal R}}{\sigma^2\nu^4}
       \int_I\|A^{-1/4}\BNS(u)\|_2^2\,\dd t.}
 \label{eq:v7-corridor-tariff}
\end{equation}
\end{corollary}

\begin{proof}
Use the pointwise tariff
\eqref{eq:v6-record-tariff}, integrate, and apply
\eqref{eq:v7-residence-summed}.  Times with $Q_\nu=-1$ have
$\mathfrak f_\nu=0$.
\end{proof}

\begin{theorem}[Every genuine shell upcrossing has a fixed critical-source
debit]
\label{thm:v7-upcrossing}
Let $a\le s<t\le b$.  If
\begin{equation}
 r_q(s)\le\frac{\sigma\nu}{2},
 \qquad
 r_q(t)\ge\sigma\nu,
 \label{eq:v7-upcrossing-endpoints}
\end{equation}
then
\begin{equation}
 \boxed{
 \int_s^t\|g_q(\tau)\|_2^2\,\dd\tau
 \ge c\sigma^2\nu^3.}
 \label{eq:v7-upcrossing-debit}
\end{equation}
The debit is independent of $q$; in particular, infinitely many such
upcrossings require infinite total critical-source action when their
spacetime shell intervals are disjoint.
\end{theorem}

\begin{proof}
Duhamel's formula for \eqref{eq:v7-exact-shell-equation} is
\[
 u_q(t)=e^{-\nu(t-s)A}u_q(s)
 -\int_s^te^{-\nu(t-\tau)A}A^{1/4}g_q(\tau)\,\dd\tau.
\]
The heat semigroup is a contraction on the shell.  Multiplication by
$\lambda_q^{1/2}$ and the reverse triangle inequality therefore give
\[
 \frac{\sigma\nu}{2}
 \le C\int_s^t
       \lambda_q e^{-c\nu\lambda_q^2(t-\tau)}
       \|g_q(\tau)\|_2\,\dd\tau.
\]
Cauchy--Schwarz and
\[
 \int_s^t\lambda_q^2
 e^{-2c\nu\lambda_q^2(t-\tau)}\,\dd\tau
 \le\frac{C}{\nu}
\]
give
\[
 \frac{\sigma\nu}{2}
 \le\frac{C}{\sqrt\nu}
       \|g_q\|_{L^2(s,t;L^2)}.
\]
Squaring and rearranging proves
\eqref{eq:v7-upcrossing-debit}.
\end{proof}

\begin{firewall}[Residence is not persistence supplied for free]
Theorem~\ref{thm:v7-residence} does not assert that an active shell persists
for a parabolic time.  It proves the exact alternative: its weighted
residence is paid either by critical stock already present at the left
endpoint or by critical nonlinear source action during the interval.
Assuming a parabolic lifetime without paying one of those terms would
insert the missing conclusion as a hypothesis.
\end{firewall}

% =============================================================================
\section{The source norm returned by the crossing argument}
\label{sec:v7-source-ledger}
% =============================================================================

The return in \eqref{eq:v7-residence-summed} must now be compared with the
available global budgets.

\begin{proposition}[The energy-paid nonlinear source is one spatial derivative weaker]
\label{prop:v7-weak-source-budget}
For every smooth interval $[0,T]$,
\begin{equation}
 \boxed{
 \int_0^T\|A^{-3/4}\BNS(u(t))\|_2^2\,\dd t
 \le C\frac{E_0^2}{\nu},}
 \qquad E_0=\|u(0)\|_2^2.
 \label{eq:v7-weak-source-budget}
\end{equation}
In dyadic form,
\begin{equation}
 \sum_q\int_0^T\|h_q(t)\|_2^2\,\dd t
 \le C\frac{E_0^2}{\nu},
 \label{eq:v7-weak-source-dyadic}
\end{equation}
whereas the residence theorem asks for
\begin{equation}
 \sum_q\int_0^T\|g_q(t)\|_2^2\,\dd t
 \asymp
 \sum_q\lambda_q^2
       \int_0^T\|h_q(t)\|_2^2\,\dd t.
 \label{eq:v7-critical-vs-weak-source}
\end{equation}
\end{proposition}

\begin{proof}
Since $\BNS(u)=\PP\operatorname{div}(u\otimes u)$, the operator
$A^{-3/4}\PP\operatorname{div}$ has order $-1/2$.  The dual Sobolev
embedding $L^{3/2}\hookrightarrow H^{-1/2}$ and
$H^1\hookrightarrow L^6$ give
\[
 \|A^{-3/4}\BNS(u)\|_2
 \le C\|u\otimes u\|_{H^{-1/2}}
 \le C\|u\otimes u\|_{3/2}
 \le C\|u\|_2\|A^{1/2}u\|_2.
\]
Square and integrate.  The energy identity supplies
\[
 \sup_{t\le T}\|u(t)\|_2^2\le E_0,
 \qquad
 \int_0^T\|A^{1/2}u(t)\|_2^2\,\dd t
 \le\frac{E_0}{2\nu},
\]
which proves \eqref{eq:v7-weak-source-budget}.  Orthogonality gives
\eqref{eq:v7-weak-source-dyadic}, and
\eqref{eq:v7-source-derivative-gap} gives
\eqref{eq:v7-critical-vs-weak-source}.
\end{proof}

\begin{proposition}[The critical source returns the V1 continuation stock]
\label{prop:v7-critical-source-return}
For every smooth interval $I$,
\begin{equation}
 \boxed{
 \int_I\|A^{-1/4}\BNS(u(t))\|_2^2\,\dd t
 \le C\int_I\|A^{1/2}u(t)\|_2^4\,\dd t.}
 \label{eq:v7-critical-source-enstrophy-square}
\end{equation}
The right-hand side is precisely the enstrophy-square stock whose finiteness
already implies continuation by
Theorem~\ref{thm:enstrophy-continuation}.  Hence the direct shell crossing
argument has gained the missing wavenumber power only by returning to the
unpaid V1 terminal quantity.
\end{proposition}

\begin{proof}
Apply the critical bilinear estimate
\eqref{eq:critical-bilinear} with both inputs equal to $u$:
\[
 \|A^{-1/4}\BNS(u)\|_2
 \le C\|A^{1/2}u\|_2^2.
\]
Square and integrate.  The final statement is
Theorem~\ref{thm:enstrophy-continuation}.
\end{proof}

\begin{assessment}[Exact derivative deficit]
The globally paid source in
Proposition~\ref{prop:v7-weak-source-budget} is $H^{-3/2}$-valued.  Active
critical-shell residence requires the $H^{-1/2}$ source.  On shell $q$ the
gap is the factor $\lambda_q^2$ in squared action.  This is the same
frequency power which was missing from the V6 occupation budget; the shell
equation does not create that power, but transfers its cost to the
critical nonlinear source.
\end{assessment}

% =============================================================================
\section{Exact operator comparison with the V34 daughter}
\label{sec:v7-v34-source-map}
% =============================================================================

Let $K>0$, $v=P_Ku$, and $w=Q_Ku$.  Write $\BNS_{\rm s}$ for the
symmetrized bilinear form used in V33--V34.  Bilinearity gives the exact
identity
\begin{equation}
 \BNS(u)=\BNS(v)+2\BNS_{\rm s}(v,w)+\BNS(w).
 \label{eq:v7-parent-split}
\end{equation}
For each shell define
\begin{align}
 g_{q,K}^{\rm hh}
 &:=A^{-1/4}\Delta_q\BNS(v),
 \label{eq:v7-head-head-source}\\
 g_{q,K}^{\rm mix}
 &:=2A^{-1/4}\Delta_q\BNS_{\rm s}(v,w),
 \label{eq:v7-mixed-source}\\
 g_{q,K}^{\rm oo}
 &:=A^{-1/4}\Delta_q\BNS(w).
 \label{eq:v7-outer-outer-source}
\end{align}
Then
\begin{equation}
 \boxed{g_q=g_{q,K}^{\rm hh}
             +g_{q,K}^{\rm mix}
             +g_{q,K}^{\rm oo}.}
 \label{eq:v7-full-shell-source-split}
\end{equation}

\begin{proposition}[Where the V34 raw daughter occurs in the shell source]
\label{prop:v7-raw-daughter-incidence}
Let
\[
 d_K^{\rm raw}:=A^{-1/4}Q_K\BNS(P_Ku).
\]
Then:
\begin{enumerate}[label=\textup{(\alph*)},leftmargin=3em]
\item if $\lambda_q>K$, so that the whole $q$-shell lies outside the head,
      \begin{equation}
       \boxed{g_{q,K}^{\rm hh}=\Delta_qd_K^{\rm raw};}
       \label{eq:v7-raw-shell-identification}
      \end{equation}
\item if $\lambda_q>2K$, then
      \begin{equation}
       \boxed{g_{q,K}^{\rm hh}=0;}
       \label{eq:v7-far-head-head-zero}
      \end{equation}
\item even in the first exterior region $K<\lambda_q\le2K$, the full
      crossing source is $g_q$, not merely the raw component
      $g_{q,K}^{\rm hh}$; the mixed and outer--outer terms in
      \eqref{eq:v7-full-shell-source-split} remain.
\end{enumerate}
Thus V34's raw daughter is an exact component of the crossing source only
in its support-compatible first exterior region, and it supplies no
head--head source at a far dissipation shell.
\end{proposition}

\begin{proof}
All projections commute.  If the entire shell lies above $K$, then
$\Delta_qQ_K=\Delta_q$, which gives
\eqref{eq:v7-raw-shell-identification}.  The Fourier support of
$\BNS(P_Ku)$ lies in $|k|\le2K$, so
\eqref{eq:v7-far-head-head-zero} follows when $\lambda_q>2K$.  The last
item is the exact split \eqref{eq:v7-full-shell-source-split}.
\end{proof}

\begin{corollary}[Far active shells are incoming or outer-generated]
\label{cor:v7-far-active-shell}
On a V34 card with cutoff $K_n$, if
\[
 Q_\nu(t)=q,\qquad \lambda_q>2K_n,
\]
then
\begin{align}
 \lambda_q\|u_q(t)\|_2^2
 &\ge\sigma^2\nu^2,
 \label{eq:v7-far-critical-mass}\\
 \nu\|A^{1/2}u_q(t)\|_2^2
 &\ge c\sigma^2\nu^3\lambda_q,
 \label{eq:v7-far-energy-debit}\\
 g_q&=g_{q,K_n}^{\rm mix}+g_{q,K_n}^{\rm oo}.
 \label{eq:v7-far-source-type}
\end{align}
The Leray budget integrates only the first power of $\lambda_q$ in
\eqref{eq:v7-far-energy-debit}; residence at the squared-frequency weight
still requires the mixed/outer critical source in
Theorem~\ref{thm:v7-residence}.
\end{corollary}

\begin{proof}
The first row is \eqref{eq:v7-Q-implies-active}.  The shell support gives
\[
 \|A^{1/2}u_q\|_2^2
 \ge\lambda_q^2\|u_q\|_2^2
 =\lambda_qz_q,
\]
which proves the second.  The third follows from
\eqref{eq:v7-far-head-head-zero} and
\eqref{eq:v7-full-shell-source-split}.
\end{proof}

% =============================================================================
\section{First-exterior V34 scale audit}
\label{sec:v7-first-exterior-audit}
% =============================================================================

For a V34 card $I_n=[a_n,b_n]$, let $\mathcal J_n$ be the support-compatible set of dyadic indices whose lower edges obey
\[
 K_n<\lambda_q\le2K_n.
\]
For these indices the entire shell is outside the head, while the
head--head output in it is automatically confined to $|k|\le2K_n$.
Set
\begin{equation}
 \mathscr S_n^{\rm ext}
 :=\int_{I_n}\sum_{q\in\mathcal J_n}
 \left(
  \|g_{q,K_n}^{\rm mix}(t)\|_2^2
 +\|g_{q,K_n}^{\rm oo}(t)\|_2^2
 \right)\,\dd t.
 \label{eq:v7-exterior-nonraw-action}
\end{equation}

\begin{proposition}[What the raw action window buys for residence]
\label{prop:v7-raw-window-residence}
On the V34 corridor, the part of the tariff occurring when
$Q_\nu(t)\in\mathcal J_n$ satisfies
\begin{equation}
 \boxed{
 \int_{I_n\cap\{Q_\nu\in\mathcal J_n\}}
       \mathfrak f_\nu(t)\,\dd t
 \le
 C\frac{\mathcal R_n^3}{\nu^3}
 +C\frac{\mathcal R_n}{\nu^4}
       \mathscr S_n^{\rm ext}.}
 \label{eq:v7-raw-window-tariff}
\end{equation}
Here the first term includes the left-end critical stock and the complete
raw head--head action window
\[
 \|d_n\|_{L^2(I_n;L^2)}^2
 =\Araw_n\le a_+\nu\mathcal R_n^2.
\]
Even if $\mathscr S_n^{\rm ext}=0$, the resulting
$O(\mathcal R_n^3/\nu^3)$ upper scale grows on doubled records and is not a
summable terminal budget.
\end{proposition}

\begin{proof}
Repeat the proof of Theorem~\ref{thm:v7-residence} but sum only over
$q\in\mathcal J_n$.  The initial shell stocks sum to at most
$C\mathcal R_n^2$ by the record corridor.  By
\eqref{eq:v7-full-shell-source-split},
\[
 \sum_{q\in\mathcal J_n}\|g_q\|_2^2
 \le3\sum_{q\in\mathcal J_n}
 \left(
  \|g_{q,K_n}^{\rm hh}\|_2^2+
  \|g_{q,K_n}^{\rm mix}\|_2^2+
  \|g_{q,K_n}^{\rm oo}\|_2^2
 \right).
\]
Proposition~\ref{prop:v7-raw-daughter-incidence}, shell orthogonality, and
the raw action window give
\[
 \int_{I_n}\sum_{q\in\mathcal J_n}
       \|g_{q,K_n}^{\rm hh}\|_2^2\,\dd t
 \le\Araw_n\le a_+\nu\mathcal R_n^2.
\]
Thus
\[
 \int_{I_n\cap\{Q_\nu\in\mathcal J_n\}}
       \lambda_{Q_\nu}^2\,\dd t
 \le C\frac{\mathcal R_n^2}{\nu^3}
     +C\frac{\mathscr S_n^{\rm ext}}{\nu^4}.
\]
Finally use
$\mathfrak f_\nu\le C\mathcal R_n\lambda_{Q_\nu}^2$ on the corridor,
absorbing the fixed factor in the record bound into $C$.  This proves
\eqref{eq:v7-raw-window-tariff}.  Since
$\mathcal R_n=2^n\mathcal R_*$, the first term is not summable.
\end{proof}

\begin{firewall}[The history-fresh short does not improve this upper bound]
The lower action $\Aocc_n$ belongs to
$q_n=(I-P_n)d_n$ and drives the V34 alternative.  The residence estimate
needs an upper bound for the complete shell source $g_q$.  Orthogonal
shorting can lower the raw norm but cannot upper-bound the mixed and
outer--outer components in \eqref{eq:v7-full-shell-source-split}.  Replacing
$g_q$ by the residual $q_n$ would reverse the required implication.
\end{firewall}

% =============================================================================
\section{Integrated V7 conclusion and next acquisition order}
\label{sec:v7-conclusion}
% =============================================================================

\begin{theorem}[What V7 proves]
\label{thm:v7-conclusion}
Relative to V1--V6 and the frozen source, the following are proved.
\begin{enumerate}[label=\textup{(V7.\arabic*)},leftmargin=5em]
\item The actual shell equation gives a squared-wavenumber residence bound
      and a frequency-independent critical-source debit for every genuine
      shell upcrossing.
\item The residence bound controls the V6 tariff on a critical-record
      corridor, but returns the action of
      $A^{-1/4}\BNS(u)$.
\item The Leray energy budget controls only the one-spatial-derivative-weaker source
      $A^{-3/4}\BNS(u)$; the returned critical source is bounded by the
      already-unpaid enstrophy-square continuation stock.
\item The V34 raw daughter is exactly the head--head component of the shell
      source only in the first exterior region.  At scales above $2K_n$ it
      vanishes, leaving mixed and outer--outer sources.
\item Even on the support-compatible first exterior region, the raw action
      window yields only a growing
      $O(\mathcal R_n^3/\nu^3)$ tariff bound and does not close the doubled
      record sequence.
\end{enumerate}
No full NS stability theorem is obtained.
\end{theorem}

\begin{proof}
Item \textup{(V7.1)} is
Theorems~\ref{thm:v7-residence} and \ref{thm:v7-upcrossing}.
Item \textup{(V7.2)} is Corollary~\ref{cor:v7-corridor-tariff}.
Item \textup{(V7.3)} is
Propositions~\ref{prop:v7-weak-source-budget} and
\ref{prop:v7-critical-source-return}.  Item \textup{(V7.4)} is
Proposition~\ref{prop:v7-raw-daughter-incidence} and
Corollary~\ref{cor:v7-far-active-shell}.  Item \textup{(V7.5)} is
Proposition~\ref{prop:v7-raw-window-residence}.
\end{proof}

\begin{programme}[V8 acquisition order]
\label{prog:v8}
\begin{enumerate}[label=\textup{(U\arabic*)},leftmargin=4em]
\item Keep the residence theorem as the exact velocity-shell gateway; do
      not repeat a norm estimate of the full critical source, which has now
      been proved circular.
\item Expand the mixed and outer--outer sources in
      \eqref{eq:v7-full-shell-source-split} on the far branch
      $\lambda_{Q_\nu}>2K_n$.  Preserve the signed low--high strain
      contraction before applying Hölder.
\item On the first-exterior branch, test whether the positive part of the
      shell source work, rather than its complete norm-square action, can
      be charged to record growth or viscous loss.  The required gain is
      more than three powers of $\mathcal R_n$ relative to
      \eqref{eq:v7-raw-window-tariff}.
\item Compare any physical-space localization with the inherited material
      cell only after writing the Fourier/localization commutator.  A
      same-scale label is not an incidence theorem.
\item If incompressibility and the trace-free strain identity do not yield
      signed depletion, construct an actual finite-mode low--high
      obstruction and move upstream to inter-card flux geometry.
\item Keep the represented-history branch at the V3 physical-incidence
      stopping line.
\end{enumerate}
\end{programme}

\begin{assessment}[Position after V7]
The velocity-shell crossing route is now fully typed.  Dynamics supplies
the missing squared-wavenumber residence only by charging the critical
nonlinear source, and that source returns exactly to the V1
enstrophy-square boundary.  V34 controls one support-compatible head--head
component but neither the far-shell sources nor the growing record scale
needed for summability.  The next noncircular test is signed low--high
strain depletion, separated into first-exterior and far-shell branches.
\end{assessment}

% =============================================================================
\section*{Version V7 history}
\addcontentsline{toc}{section}{Version V7 history}
% =============================================================================

Appended the exact critical-shell differential inequality, a
squared-wavenumber active-residence estimate, and a fixed critical-source
debit for every genuine shell upcrossing.  Proved that the globally
energy-paid nonlinear source is $A^{-3/4}\BNS(u)$ while residence requires
the one-spatial-derivative-stronger $A^{-1/4}\BNS(u)$, whose square action returns
to the V1 enstrophy-square continuation stock.  Decomposed the full shell
source into head--head, mixed, and outer--outer components and proved the
exact support incidence of the V34 raw daughter.  Audited the first
exterior action window and obtained only the growing
$O(\mathcal R_n^3/\nu^3)$ tariff scale, leaving signed low--high depletion
as the next test.  No global regularity claim was made.
% Future versions are appended immediately above the sole terminator.
% =============================================================================
% VERSION V8 APPEND
% =============================================================================

\section{V8: the signed low--high vorticity identity}
\label{sec:v8-low-high-identity}
% =============================================================================

V7 showed that a norm-square estimate of the critical shell source returns
the unpaid enstrophy-square stock.  The remaining V8 instruction is
therefore genuinely signed.  We begin with an exact fixed-cutoff identity
which isolates low-frequency strain before any absolute value is taken.

Let
\begin{equation}
 \omega:=\operatorname{curl}u,\qquad
 v:=P_Ku,\quad w:=Q_Ku,\qquad
 \zeta:=\operatorname{curl}v,\quad
 \eta:=\operatorname{curl}w,
 \label{eq:v8-vorticity-split}
\end{equation}
and let
\[
 S(v):=\frac{\nabla v+\nabla v^{\mathsf T}}2.
\]

\begin{theorem}[Exact high-vorticity balance]
\label{thm:v8-high-vorticity-balance}
For every fixed $K$ and every smooth solution,
\begin{equation}
 \boxed{
 \frac12\frac{\dd}{\dd t}\|\eta\|_2^2
 +\nu\|\nabla\eta\|_2^2
 =\mathfrak L_K+\mathfrak Q_K,}
 \label{eq:v8-high-vorticity-balance}
\end{equation}
where the signed low--high strain term is
\begin{equation}
 \boxed{
 \mathfrak L_K
 :=\int_{\Torus^3}S(v):(\eta\otimes\eta)\,\dd x}
 \label{eq:v8-low-high-strain}
\end{equation}
and the exact remainder is
\begin{align}
 \mathfrak Q_K
 :={}&
 \left\langle
  (\zeta\cdot\nabla)v+(\zeta\cdot\nabla)w
  +(\eta\cdot\nabla)w,\eta
 \right\rangle\notag\\
 &-\left\langle
  (v\cdot\nabla)\zeta+(w\cdot\nabla)\zeta,\eta
 \right\rangle .
 \label{eq:v8-high-vorticity-remainder}
\end{align}
Thus the mixed source contains a genuine strain contraction, not merely a
product norm.  The outer--outer stretching is retained in
$\mathfrak Q_K$.
\end{theorem}

\begin{proof}
The vorticity equation is
\[
 \partial_t\omega+(u\cdot\nabla)\omega
 =(\omega\cdot\nabla)u+\nu\Delta\omega.
\]
Apply $Q_K$ and pair with $\eta$.  The projector is self-adjoint and
commutes with curl and $\Delta$, so it disappears from the pairing.
Expand $u=v+w$ and $\omega=\zeta+\eta$.  The two transport terms
\[
 \langle(v\cdot\nabla)\eta,\eta\rangle,\qquad
 \langle(w\cdot\nabla)\eta,\eta\rangle
\]
vanish by incompressibility.  The term
$\langle(\eta\cdot\nabla)v,\eta\rangle$ equals
\eqref{eq:v8-low-high-strain}, because the antisymmetric part of
$\nabla v$ annihilates $\eta\otimes\eta$.  The five remaining terms are
exactly \eqref{eq:v8-high-vorticity-remainder}.
\end{proof}

\begin{proposition}[Trace-free strain has both signs]
\label{prop:v8-tracefree-sign}
At every point, $\operatorname{tr}S(v)=0$.  If $S(v)\ne0$, then it has a
positive and a negative eigenvalue.  More quantitatively,
\begin{equation}
 \lambda_{\max}(S(v))\ge\frac{|S(v)|}{\sqrt6},
 \qquad
 -\lambda_{\min}(S(v))\ge\frac{|S(v)|}{\sqrt6}.
 \label{eq:v8-strain-eigenvalue-bound}
\end{equation}
Consequently incompressibility alone supplies no pointwise sign for
$S(v):(\eta\otimes\eta)$.
\end{proposition}

\begin{proof}
The trace identity is $\operatorname{div}v=0$.  Let the ordered eigenvalues
be $\mu_1\ge\mu_2\ge\mu_3$ and put $M=\mu_1$.  If the symmetric matrix is
nonzero, trace zero forces $M>0$ and $\mu_3<0$.  Since
$\mu_3=-(\mu_1+\mu_2)\ge-2M$, all three eigenvalues have magnitude at most
$2M$, and
\[
 |S(v)|^2=\mu_1^2+\mu_2^2+\mu_3^2\le6M^2.
\]
Apply the same argument to $-S(v)$ for the second inequality.  Alignment
of $\eta$ with the corresponding eigenvectors gives the last statement.
\end{proof}

% =============================================================================
\section{An actual finite-mode low--high obstruction}
\label{sec:v8-finite-mode-obstruction}
% =============================================================================

The preceding algebraic sign freedom survives the periodic Fourier
incidences and the complete nonlinear assembly.

\begin{theorem}[Positive low--high stretching with zero assembled remainder]
\label{thm:v8-low-high-obstruction}
Let $N\ge6$ be an integer and set
\[
 \ell=(1,0,0),\qquad
 p=(N,1,0),\qquad
 q=(N-1,1,0),
 \qquad p-q=\ell.
\]
For $A,B>0$, define
\begin{equation}
 u_0(x):=
 A\mathbf e_2\sin(\ell\cdot x)
 +B\mathbf e_3\bigl(\cos(p\cdot x)-\cos(q\cdot x)\bigr).
 \label{eq:v8-explicit-field}
\end{equation}
Then $u_0$ is real, mean-zero, and divergence-free.  For the cutoff $K=2$,
write $u_0=v+w$ as in \eqref{eq:v8-vorticity-split}.  At this datum,
\begin{align}
 \boxed{
 \mathfrak L_2(u_0)
 =\frac{AB^2}{4}(2N-1)>0,}
 \label{eq:v8-positive-low-high}\\
 \boxed{\mathfrak Q_2(u_0)=0,}
 \label{eq:v8-zero-remainder}\\
 \boxed{
 \mathcal G(u_0):=-\langle\BNS(u_0),Au_0\rangle
 =\frac{AB^2}{4}(2N-1).}
 \label{eq:v8-full-positive-production}
\end{align}
The field is smooth initial data for a local Navier--Stokes solution, so
these identities occur at an actual solution time.  Hence neither
incompressibility, trace-free strain, exact high/low support, nor assembly
with the remaining nonlinear terms forces signed low--high depletion.
\end{theorem}

\begin{proof}
Each polarization is orthogonal to its wavevector.  The low and high
vorticities are
\begin{align*}
 \zeta&=A\mathbf e_3\cos x_1,\\
 \eta&=B\left[
   (-\mathbf e_1+N\mathbf e_2)\sin(p\cdot x)
  -(-\mathbf e_1+(N-1)\mathbf e_2)\sin(q\cdot x)
 \right].
\end{align*}
The only nonzero entries of $S(v)$ are
\[
 S(v)_{12}=S(v)_{21}=\frac A2\cos x_1.
\]
The self-frequency terms integrate to zero, while
\[
 \int_{\Torus^3}\cos x_1
       \sin(p\cdot x)\sin(q\cdot x)\,\dd x=\frac14.
\]
The cross coefficient in $\eta_1\eta_2$ is $2N-1$.  This proves
\eqref{eq:v8-positive-low-high}.

There is also a direct complete-production calculation.  All fields are
independent of $x_3$, the high field is parallel to $\mathbf e_3$, and the
low field is parallel to $\mathbf e_2$.  Thus
\[
 \BNS(u_0)
 =AB\mathbf e_3\sin x_1
       \bigl(-\sin(p\cdot x)+\sin(q\cdot x)\bigr);
\]
this vector field is already divergence-free.  Since
\[
 (Au_0)_3
 =B\bigl(|p|^2\cos(p\cdot x)-|q|^2\cos(q\cdot x)\bigr),
\]
orthogonality and $p-q=\ell$ give
\[
 \langle\BNS(u_0),Au_0\rangle
 =\frac{AB^2}{4}\bigl(|q|^2-|p|^2\bigr)
 =-\frac{AB^2}{4}(2N-1).
\]
This proves \eqref{eq:v8-full-positive-production}.  Every Fourier
frequency in $\BNS(u_0)$ has magnitude greater than $2$ when $N\ge6$, so
the nonlinear derivative of the low enstrophy is zero.  The high-vorticity
nonlinear production therefore equals the full production.  Comparing
with \eqref{eq:v8-positive-low-high} in the exact balance
\eqref{eq:v8-high-vorticity-balance} gives
\eqref{eq:v8-zero-remainder}.  Smooth local existence supplies the final
claim.
\end{proof}

\begin{proposition}[The isolated obstruction is still viscosity-dominated
under an energy cap]
\label{prop:v8-obstruction-viscosity}
For the datum in Theorem~\ref{thm:v8-low-high-obstruction},
\begin{align}
 \|u_0\|_2^2
 &=\frac{A^2}{2}+B^2,
 \label{eq:v8-explicit-energy}\\
 \|Au_0\|_2^2
 &=\frac{A^2}{2}
   +\frac{B^2}{2}\bigl(|p|^4+|q|^4\bigr).
 \label{eq:v8-explicit-palinstrophy}
\end{align}
If $\|u_0\|_2^2\le E_*$ and $B\ne0$, then
\begin{equation}
 \boxed{
 \frac{\mathcal G(u_0)}{\nu\|Au_0\|_2^2}
 \le C\frac{\sqrt{E_*}}{\nu N^3}.}
 \label{eq:v8-explicit-viscous-ratio}
\end{equation}
Thus this three-mode family disproves a universal sign cancellation but
does not produce high-frequency positive viscous excess at fixed energy.
\end{proposition}

\begin{proof}
The three modes are orthogonal, which gives
\eqref{eq:v8-explicit-energy}--\eqref{eq:v8-explicit-palinstrophy}.
The energy cap gives $A\le\sqrt{2E_*}$, while
$|p|,|q|\asymp N$.  Divide
\eqref{eq:v8-full-positive-production} by
\eqref{eq:v8-explicit-palinstrophy} and insert the cap.
\end{proof}

\begin{firewall}[What the finite-mode obstruction decides]
It decides the signed question: no depletion follows from the exact
algebra or support incidences.  Proposition~\ref{prop:v8-obstruction-viscosity}
also prevents the opposite overclaim.  At fixed energy, one isolated low
mode cannot defeat the $N^2$ viscous rate of arbitrarily high vorticity.
A surviving mechanism would have to assemble many low modes coherently or
remain in a scale range where that collective strain is not absorbable.
\end{firewall}

% =============================================================================
\section{The very-far low-strain firewall}
\label{sec:v8-very-far}
% =============================================================================

The energy cap gives a useful collective bound which does not sum
individual triads.

\begin{lemma}[Sharp energy-to-low-strain Bernstein bound]
\label{lem:v8-low-strain-Bernstein}
For $v=P_Ku$ and $K\ge1$,
\begin{equation}
 \boxed{
 \|S(v)\|_\infty
 \le\|\nabla v\|_\infty
 \le C K^{5/2}\|v\|_2
 \le C E_*^{1/2}K^{5/2}.}
 \label{eq:v8-low-strain-Bernstein}
\end{equation}
\end{lemma}

\begin{proof}
By Fourier inversion and Cauchy--Schwarz,
\[
 \|\nabla v\|_\infty
 \le C\sum_{0<|k|\le K}|k|\,|\widehat v(k)|
 \le C\left(\sum_{0<|k|\le K}|k|^2\right)^{1/2}
       \|v\|_2.
\]
The lattice sum is $O(K^5)$, and $\|v\|_2\le\|u\|_2\le E_*^{1/2}$.
\end{proof}

For $N>K$, let
\[
 \eta_{>N}:=P_{>N}\omega
\]
and isolate the complete low-strain contraction
\begin{equation}
 \mathfrak L_{K,>N}
 :=\int_{\Torus^3}
       S(P_Ku):(\eta_{>N}\otimes\eta_{>N})\,\dd x.
 \label{eq:v8-tail-low-strain}
\end{equation}

\begin{theorem}[Very-far low-strain absorption]
\label{thm:v8-very-far-absorption}
For every $K,N\ge1$,
\begin{equation}
 \boxed{
 |\mathfrak L_{K,>N}|
 \le C E_*^{1/2}K^{5/2}N^{-2}
       \|\nabla\eta_{>N}\|_2^2.}
 \label{eq:v8-very-far-bound}
\end{equation}
Consequently, for any $0<\theta<1$, the low-strain term is absorbed by
$\theta\nu\|\nabla\eta_{>N}\|_2^2$ whenever
\begin{equation}
 \boxed{
 N\ge
 C_\theta\nu^{-1/2}E_*^{1/4}K^{5/4}.}
 \label{eq:v8-very-far-threshold}
\end{equation}
This statement concerns the complete low-strain contraction, including
all cross terms inside the tail; it does not claim that the remaining
outer--outer term $\mathfrak Q_K$ is absorbed.
\end{theorem}

\begin{proof}
Holder's inequality, Lemma~\ref{lem:v8-low-strain-Bernstein}, and the
spectral Poincare inequality give
\[
 |\mathfrak L_{K,>N}|
 \le\|S(P_Ku)\|_\infty\|\eta_{>N}\|_2^2
 \le C E_*^{1/2}K^{5/2}N^{-2}
       \|\nabla\eta_{>N}\|_2^2.
\]
Choose $N$ so that the coefficient is at most $\theta\nu$.
\end{proof}

\begin{corollary}[The unresolved far-shell fan on a V34 record card]
\label{cor:v8-intermediate-fan}
At the canonical cutoff
\[
 K_n=\frac{\mathcal R_n^2}{E_*},
\]
the low-strain part of every tail above
\begin{equation}
 N_{n}^{\rm vf}
 :=C_\theta
   \frac{\mathcal R_n^{5/2}}{\nu^{1/2}E_*}
 \label{eq:v8-canonical-very-far}
\end{equation}
is viscosity-absorbable.  Combining this with
Corollary~\ref{cor:v7-far-active-shell}, the only far range in which the
low-strain term itself remains unresolved is
\begin{equation}
 \boxed{
 2K_n<N<N_n^{\rm vf}.}
 \label{eq:v8-intermediate-fan}
\end{equation}
The outer--outer remainder must still be treated separately throughout
this range.
\end{corollary}

\begin{proof}
Insert $K=K_n$ in \eqref{eq:v8-very-far-threshold}:
\[
 E_*^{1/4}K_n^{5/4}
 =E_*^{1/4}\left(\frac{\mathcal R_n^2}{E_*}\right)^{5/4}
 =\frac{\mathcal R_n^{5/2}}{E_*}.
\]
The lower boundary $2K_n$ is exactly where the head--head V34 daughter
vanishes by \eqref{eq:v7-far-head-head-zero}.
\end{proof}

\begin{assessment}[What the firewall gains]
V7 left all scales above $2K_n$ in one far branch.  V8 removes the
low-strain part beyond $N_n^{\rm vf}$ by a paid viscous estimate and leaves
a concrete intermediate fan whose relative width is
\[
 \frac{N_n^{\rm vf}}{K_n}
 =C_\theta\left(\frac{\mathcal R_n}{\nu}\right)^{1/2}.
\]
The remaining issue is collective coherence inside that fan and the
outer--outer stretching term, not an arbitrarily remote action of the
finite head.
\end{assessment}

% =============================================================================
\section{Positive shell work returns spectral variation}
\label{sec:v8-positive-shell-work}
% =============================================================================

For the critical shell balance define
\begin{equation}
 \mathcal T_q(t):=
 -\lambda_q\langle\Delta_q\BNS(u),u_q\rangle,
 \qquad
 \mathcal D_q(t):=
 \nu\lambda_q\|A^{1/2}u_q\|_2^2.
 \label{eq:v8-shell-work-dissipation}
\end{equation}
The exact projected equation gives
\begin{equation}
 \mathcal T_q-\mathcal D_q=\frac12z_q'.
 \label{eq:v8-net-shell-work}
\end{equation}

\begin{proposition}[Positive net shell work is upward stock variation]
\label{prop:v8-shell-variation}
On every smooth interval $I$,
\begin{equation}
 \boxed{
 \int_I(\mathcal T_q-\mathcal D_q)_+\,\dd t
 =\frac12\int_I(z_q')_+\,\dd t
 =\frac12\operatorname{Var}_I^+(z_q).}
 \label{eq:v8-positive-work-variation}
\end{equation}
Endpoint values determine only
$\int_Iz_q'$, not its positive variation.  The V34 intrinsic variation is
variation of the source residual $q_n$, and
Proposition~\ref{prop:v5-two-variations} supplies no map from it to
\eqref{eq:v8-positive-work-variation}.  Thus retaining the sign of
first-exterior shell work does not, by itself, repair
\eqref{eq:v7-raw-window-tariff}.
\end{proposition}

\begin{proof}
Take the positive part in \eqref{eq:v8-net-shell-work} and integrate.  The
last assertion is precisely the type distinction proved in
Proposition~\ref{prop:v5-two-variations}.
\end{proof}

% =============================================================================
\section{The material-cell localization commutator}
\label{sec:v8-localization}
% =============================================================================

For every spatial multiplier $\chi$ and Fourier projector $P_K$,
\begin{equation}
 \boxed{
 P_K(\chi f)=\chi P_Kf+[P_K,\chi]f.}
 \label{eq:v8-localization-commutator}
\end{equation}
For the sharp cutoff inherited by V34, the commutator has no automatic
inverse-frequency gain at the spectral boundary.

\begin{proposition}[No uniform $K^{-1}$ gain for the sharp cutoff]
\label{prop:v8-sharp-commutator}
Fix $0<\varepsilon<1$ and, for integers $N\ge3$, put
\[
 \chi(x)=1+\varepsilon\cos(2x_1),
 \qquad
 f_N(x)=\mathbf e_2\cos((N+1)x_1).
\]
Then $\chi$ is positive, $f_N$ is divergence-free,
$P_Nf_N=0$, and
\begin{equation}
 \boxed{
 [P_N,\chi]f_N
 =\frac{\varepsilon}{2}\mathbf e_2\cos((N-1)x_1).}
 \label{eq:v8-explicit-commutator}
\end{equation}
In particular,
\begin{equation}
 \|[P_N,\chi]f_N\|_2
 =\frac{\varepsilon}{2}\|f_N\|_2,
 \label{eq:v8-commutator-no-decay}
\end{equation}
so no estimate
\[
 \|[P_N,\chi]f\|_2
 \le C N^{-\alpha}\|\chi\|_{C^1}\|f\|_2,
 \qquad \alpha>0,
\]
can hold uniformly at the sharp boundary.
\end{proposition}

\begin{proof}
The product-to-sum identity gives
\[
 \chi f_N
 =f_N+\frac{\varepsilon}{2}\mathbf e_2
  \left(\cos((N+3)x_1)+\cos((N-1)x_1)\right).
\]
The first two frequencies lie above $N$, while $N-1$ lies below it.  Apply
$P_N$ and use $P_Nf_N=0$.  Orthogonality gives the norm identity.
\end{proof}

\begin{firewall}[What a material-cell comparison must add]
A physical cutoff, even a smooth positive one, can move a boundary mode
across the sharp spectral head with order-one amplitude.  Therefore the
inherited material cell cannot be identified with a V34 head or daughter
without an explicit commutator term.  A usable comparison must introduce
a spectral guard band, a smooth spectral multiplier with a proved
commutator estimate, or a separately paid boundary action.  None of those
rows is presently inherited.
\end{firewall}

% =============================================================================
\section{Integrated V8 conclusion and next acquisition order}
\label{sec:v8-conclusion}
% =============================================================================

\begin{theorem}[What V8 proves]
\label{thm:v8-conclusion}
Relative to V1--V7 and the frozen source, the following are proved.
\begin{enumerate}[label=\textup{(V8.\arabic*)},leftmargin=5em]
\item The high-vorticity equation splits exactly into signed low--high
      strain and a typed five-term remainder containing outer--outer
      stretching.
\item Trace-free strain has no favorable sign, and an explicit three-mode
      datum of an actual local solution has strictly positive low--high
      production with zero assembled nonlinear remainder.
\item That isolated obstruction is viscosity-dominated at high frequency
      under a fixed energy cap; it disproves sign depletion but is not a
      terminal positive-excess mechanism.
\item Collective low strain is viscosity-absorbable beyond the explicit
      scale $N_n^{\rm vf}$, leaving the intermediate fan
      \eqref{eq:v8-intermediate-fan}.
\item Positive net shell work is exactly upward spectral-stock variation,
      which is not supplied by V34's source-residual variation.
\item A sharp spectral cutoff and a physical-space material cell have an
      order-one boundary commutator unless a new guard-band or commutator
      payment is proved.
\end{enumerate}
No full NS stability theorem is obtained.
\end{theorem}

\begin{proof}
Item \textup{(V8.1)} is
Theorem~\ref{thm:v8-high-vorticity-balance}.  Items
\textup{(V8.2)--(V8.3)} are
Theorem~\ref{thm:v8-low-high-obstruction} and
Proposition~\ref{prop:v8-obstruction-viscosity}.  Item \textup{(V8.4)} is
Theorem~\ref{thm:v8-very-far-absorption} and
Corollary~\ref{cor:v8-intermediate-fan}.  Item \textup{(V8.5)} is
Proposition~\ref{prop:v8-shell-variation}.  Item \textup{(V8.6)} is
Proposition~\ref{prop:v8-sharp-commutator}.
\end{proof}

\begin{programme}[V9 acquisition order]
\label{prog:v9}
\begin{enumerate}[label=\textup{(W\arabic*)},leftmargin=4em]
\item Work inside the intermediate fan
      $2K_n<N<N_n^{\rm vf}$; the very-far low-strain branch is now closed.
\item Replace the isolated low wave by a divergence-free, Bernstein-coherent
      low-frequency strain packet and pair it with a high-vorticity packet.
      Determine by explicit scaling and sign calculation whether positive
      viscous excess is possible under the fixed energy cap and record
      relation $K_n=\mathcal R_n^2/E_*$.
\item Treat the outer--outer term in
      \eqref{eq:v8-high-vorticity-remainder} separately.  Do not charge it
      to the V34 raw daughter, which vanishes above $2K_n$.
\item If a collective finite-mode obstruction exists, stop seeking a
      pointwise depletion theorem and move upstream to the time persistence
      or inter-card flux cost of that packet.  If it does not, record the
      quantitative loss that lets viscosity absorb it.
\item Any use of the inherited material cell must include one of the three
      commutator remedies listed after
      Proposition~\ref{prop:v8-sharp-commutator}.
\item Keep the represented-history branch at the V3 physical-incidence
      stopping line.
\end{enumerate}
\end{programme}

\begin{assessment}[Position after V8]
Signed low--high cancellation is not an algebraic consequence of
incompressibility: it fails on an explicit actual local solution.  Yet one
isolated interaction is dissipative under the energy cap, while collective
low strain is harmless only beyond a quantified very-far threshold.  The
remaining pointwise question is now confined to a finite intermediate fan
and is intrinsically collective.  Physical localization also requires a
new commutator payment rather than a same-cell label.
\end{assessment}

% =============================================================================
\section*{Version V8 history}
\addcontentsline{toc}{section}{Version V8 history}
% =============================================================================

Appended the exact high-vorticity balance separating signed low--high
strain from a five-term remainder.  Proved that trace-free strain has both
signs and constructed an explicit three-mode local Navier--Stokes datum
whose low--high stretching is positive and equals the complete nonlinear
enstrophy production.  Proved that this isolated obstruction is still
viscosity-dominated under a fixed energy cap.  Established a collective
energy-to-low-strain bound and a very-far viscous firewall, leaving an
explicit intermediate frequency fan.  Identified positive net shell work
with upward spectral variation and proved an order-one sharp-cutoff
localization commutator obstruction for material-cell comparisons.  No
global regularity claim was made.
% Future versions are appended immediately above the sole terminator.
% =============================================================================
% VERSION V9 APPEND
% =============================================================================

\section{V9: a compact divergence-free seed with positive production}
\label{sec:v9-compact-seed}
% =============================================================================

V8 showed that one isolated low--high triad is viscosity-dominated at fixed
energy.  It did not decide whether a Bernstein-coherent collection of
modes can overcome viscosity.  The decisive distinction is spatial
concentration: at fixed $L^2$ energy, a three-dimensional packet has
nonlinear enstrophy production of order $N^{9/2}$ and palinstrophy of order
$N^4$.  We first construct a seed whose production has the required sign.

For a smooth compactly supported divergence-free field $U$ on $\R^3$, put
\begin{equation}
 \mathscr G_{\R^3}(U)
 :=\int_{\R^3}
 S(U):\bigl(\operatorname{curl}U\otimes
            \operatorname{curl}U\bigr)\,\dd x.
 \label{eq:v9-Euclidean-production}
\end{equation}

\begin{lemma}[Existence of a compact positive-production seed]
\label{lem:v9-compact-positive-seed}
There exists a real field
\[
 U\in C_c^\infty(\R^3;\R^3),\qquad
 \operatorname{div}U=0,
\]
such that
\begin{equation}
 \boxed{\mathscr G_{\R^3}(U)>0.}
 \label{eq:v9-positive-seed}
\end{equation}
Moreover $\int_{\R^3}U\,\dd x=0$.
\end{lemma}

\begin{proof}
Take the periodic field from
Theorem~\ref{thm:v8-low-high-obstruction}, with any fixed
$A,B>0$ and $N\ge6$, and denote it by $\overline u$.  Its production
density has strictly positive mean by
\eqref{eq:v8-full-positive-production}.  It is smooth, divergence-free,
and mean-zero, so Fourier inversion supplies a smooth periodic vector
potential $\overline{\mathcal A}$ satisfying
\[
 \operatorname{curl}\overline{\mathcal A}=\overline u.
\]

For a large integer $m$, choose a smooth cutoff $\chi_m$ which equals one
on a cube containing $m^3$ complete period cells, vanishes outside its
unit-thickness enlargement, and whose derivatives of every fixed order
are bounded independently of $m$.  Define
\[
 U_m:=\operatorname{curl}(\chi_m\overline{\mathcal A}).
\]
Then $U_m$ is smooth, compactly supported, and divergence-free.  On the
interior cube it equals $\overline u$.  The integral of its production
over that cube is $c_0m^3$ for a constant $c_0>0$.  In the boundary layer,
$U_m$ and its first derivatives are bounded independently of $m$, while
the layer has volume $O(m^2)$.  Hence
\[
 \mathscr G_{\R^3}(U_m)=c_0m^3+O(m^2)>0
\]
for all sufficiently large $m$.  Fix one such $m$ and call the resulting
field $U$.  Finally, the integral of a compactly supported curl is zero.
\end{proof}

Let $c_{\Torus}>0$ be the fixed factor converting Euclidean volume to the
normalized Haar convention on the chosen torus chart.  Define
\begin{equation}
 \begin{aligned}
 E_U&:=c_{\Torus}\|U\|_{L^2(\R^3)}^2,&
 Y_U&:=c_{\Torus}\|\nabla U\|_{L^2(\R^3)}^2,\\
 H_U&:=c_{\Torus}\|U\|_{\dot H^{1/2}(\R^3)}^2,&
 Z_U&:=c_{\Torus}\|\Delta U\|_{L^2(\R^3)}^2,\\
 G_U&:=c_{\Torus}\mathscr G_{\R^3}(U)>0.
 \end{aligned}
 \label{eq:v9-seed-constants}
\end{equation}
The Euclidean Fourier-transform convention in $H_U$ is chosen compatibly
with Parseval, so these constants give the exact torus scaling below.

% =============================================================================
\section{Fixed-energy concentration defeats pointwise viscosity}
\label{sec:v9-concentration}
% =============================================================================

Choose a point $x_0$ in a torus coordinate chart.  For all sufficiently
large $N$, the rescaled support of $U$ fits inside that chart.  Define
\begin{equation}
 W_N(x):=\beta N^{3/2}U\bigl(N(x-x_0)\bigr),
 \label{eq:v9-concentrated-packet}
\end{equation}
extended by zero and then periodically.  It is smooth, mean-zero, and
divergence-free.

\begin{theorem}[Exact concentration exponents]
\label{thm:v9-concentration-exponents}
For every fixed $\beta>0$,
\begin{align}
 \|W_N\|_2^2
 &=\beta^2E_U,
 \label{eq:v9-packet-energy}\\
 \|A^{1/2}W_N\|_2^2
 &=\beta^2N^2Y_U,
 \label{eq:v9-packet-enstrophy}\\
 \|AW_N\|_2^2
 &=\beta^2N^4Z_U,
 \label{eq:v9-packet-palinstrophy}\\
 \mathcal G(W_N)
 &:=-\langle\BNS(W_N),AW_N\rangle
 =\beta^3N^{9/2}G_U.
 \label{eq:v9-packet-production}
\end{align}
Consequently,
\begin{equation}
 \boxed{
 \frac{\mathcal G(W_N)}
      {\nu\|AW_N\|_2^2}
 =\frac{\beta G_U}{\nu Z_U}N^{1/2}
 \longrightarrow\infty.}
 \label{eq:v9-production-viscosity-ratio}
\end{equation}
In particular, for all sufficiently large $N$,
\begin{equation}
 \boxed{\mathcal G(W_N)-\nu\|AW_N\|_2^2>0.}
 \label{eq:v9-positive-viscous-excess}
\end{equation}
Each $W_N$ is admissible initial data for a smooth local Navier--Stokes
solution, whose enstrophy is strictly increasing at the initial time.
\end{theorem}

\begin{proof}
The change of variables $y=N(x-x_0)$ gives
\eqref{eq:v9-packet-energy}.  One derivative of $W_N$ contributes a factor
$N^{5/2}$ and two derivatives contribute $N^{7/2}$.  After the volume
factor $N^{-3}$, this proves
\eqref{eq:v9-packet-enstrophy}--\eqref{eq:v9-packet-palinstrophy}.  The
production is cubic in first derivatives, so its factor is
\[
 (N^{5/2})^3N^{-3}=N^{9/2},
\]
with sign $G_U>0$.  This proves
\eqref{eq:v9-packet-production} and the ratio
\eqref{eq:v9-production-viscosity-ratio}.  The enstrophy identity
\[
 \frac12\frac{\dd}{\dd t}\|A^{1/2}u\|_2^2
 =\mathcal G(u)-\nu\|Au\|_2^2
\]
and smooth local existence give the final statement.
\end{proof}

\begin{firewall}[Why this does not contradict V8]
The V8 obstruction used three Fourier modes and produced only
$O(N)$ nonlinear enstrophy production against $O(N^4)$ palinstrophy.
The packet $W_N$ contains a coherent three-dimensional population of
modes and saturates the physical-space concentration gain.  Its
$N^{9/2}$ production is therefore a genuinely collective effect, not a
re-estimate of the isolated triad.
\end{firewall}

% =============================================================================
\section{Critical record and the canonical cutoff of the packet}
\label{sec:v9-record-scaling}
% =============================================================================

The homogeneous half-derivative norm has the standard concentration
asymptotic.

\begin{lemma}[Half-derivative scaling on the torus]
\label{lem:v9-half-derivative-scaling}
Along dyadic $N\to\infty$,
\begin{equation}
 \boxed{
 \|A^{1/4}W_N\|_2^2
 =\beta^2NH_U+o(N).}
 \label{eq:v9-half-derivative-scaling}
\end{equation}
\end{lemma}

\begin{proof}
The Fourier coefficients have the form
\[
 \widehat W_N(k)
 =c_{\Torus}\beta N^{-3/2}e^{-\mathrm i k\cdot x_0}
   \widehat U(k/N).
\]
Therefore
\[
 N^{-1}\|A^{1/4}W_N\|_2^2
 =c_{\Torus}^2\beta^2
   N^{-3}\sum_{k\in\mathbb Z^3}
   |k/N|\,|\widehat U(k/N)|^2.
\]
Because $\widehat U$ is Schwartz and the integrand is continuous away from
one integrable point at the origin, the Riemann sums converge to the
Euclidean $\dot H^{1/2}$ integral.  The fixed normalization is already
absorbed in $H_U$.
\end{proof}

To place the packet strictly above its canonical record--energy cutoff,
retain a fixed low-frequency energy reservoir.  Let $\phi$ be a smooth
mean-zero divergence-free trigonometric polynomial.  Fix a small
$0<\delta<1/4$, normalize
\begin{equation}
 \|\phi\|_2^2=(1-\delta)E_*,
 \qquad
 \beta^2E_U=\delta E_*,
 \label{eq:v9-energy-split}
\end{equation}
and put
\[
 \widetilde u_N:=\phi+W_N.
\]
Because $\int U=0$, Taylor expansion of $\phi$ at $x_0$ shows that the
cross energy is $O(N^{-5/2})$.  Multiplying by a scalar
$c_N=1+O(N^{-5/2})$ produces
\begin{equation}
 u_N:=c_N\widetilde u_N,\qquad
 \|u_N\|_2^2=E_*.
 \label{eq:v9-energy-normalized-data}
\end{equation}

\begin{proposition}[Canonical cutoff and positive excess with an energy
reservoir]
\label{prop:v9-canonical-packet}
Let
\[
 \kappa_U:=\frac{H_U}{E_U}>0,\qquad
 \mathcal R_N:=\|A^{1/4}u_N\|_2,\qquad
 K_N:=\frac{\mathcal R_N^2}{E_*}.
\]
Then
\begin{align}
 \mathcal R_N^2
 &=\delta E_*\kappa_U N+O(1),
 \label{eq:v9-packet-record}\\
 K_N&=\delta\kappa_UN+O(1),
 \label{eq:v9-packet-canonical-cutoff}\\
 \mathcal G(u_N)
 &=\beta^3G_UN^{9/2}+O(N^2),
 \label{eq:v9-reservoir-production}\\
 \|Au_N\|_2^2
 &=\beta^2Z_UN^4+O(N^2).
 \label{eq:v9-reservoir-palinstrophy}
\end{align}
In particular,
\begin{equation}
 \boxed{
 \mathcal G(u_N)-\nu\|Au_N\|_2^2>0}
 \label{eq:v9-reservoir-positive-excess}
\end{equation}
for all sufficiently large $N$.
\end{proposition}

\begin{proof}
The low polynomial has finite Fourier support.  Its cross term with
$W_N$ in the half-derivative norm is $O(N^{-5/2})$, while its own
contribution is $O(1)$.  Lemma~\ref{lem:v9-half-derivative-scaling} and
\eqref{eq:v9-energy-split} prove
\eqref{eq:v9-packet-record}--\eqref{eq:v9-packet-canonical-cutoff}.

The production is a trilinear form in first derivatives.  The pure packet
term is \eqref{eq:v9-packet-production}.  Terms containing two packet
gradients and one fixed gradient are $O(N^2)$; terms containing at most one
packet gradient are smaller.  This proves
\eqref{eq:v9-reservoir-production}.  Similarly, Cauchy--Schwarz gives an
$O(N^2)$ upper bound for the cross term in the palinstrophy, proving
\eqref{eq:v9-reservoir-palinstrophy}.  Since $N^{9/2}$ dominates $N^4$,
\eqref{eq:v9-reservoir-positive-excess} follows.
\end{proof}

% =============================================================================
\section{The actual dissipation index lies in the intermediate fan}
\label{sec:v9-intermediate-fan}
% =============================================================================

We now compare $Q_\nu(u_N)$ with $K_N$ rather than merely calling $N$ the
packet frequency.

\begin{theorem}[Dissipation-index localization of the collective packet]
\label{thm:v9-index-in-fan}
There is a constant $c_1>0$ depending only on the seed $U$ such that, for
every fixed $0<\varepsilon<1/4$ and all sufficiently large dyadic $N$,
\begin{equation}
 \boxed{
 c_1N\le\lambda_{Q_\nu(u_N)}
 \le C_\varepsilon N^{1+\varepsilon}.}
 \label{eq:v9-index-window}
\end{equation}
The energy fraction $\delta$ in \eqref{eq:v9-energy-split} may be chosen
so small that
\begin{equation}
 \boxed{
 2K_N<\lambda_{Q_\nu(u_N)}
 <N_N^{\rm vf}:=
 C_\theta\frac{\mathcal R_N^{5/2}}{\nu^{1/2}E_*}}
 \label{eq:v9-index-intermediate-fan}
\end{equation}
for all sufficiently large $N$.  Hence these positive-excess data occupy
the unresolved intermediate fan from
\eqref{eq:v8-intermediate-fan}.
\end{theorem}

\begin{proof}
Because $U\ne0$, one Euclidean dyadic annulus has a nonzero sharp
band-pass component.  Choose such an index $j_0$ and a point where that
component is nonzero.  For dyadic $N=2^m$, the Riemann-sum argument in
Lemma~\ref{lem:v9-half-derivative-scaling}, now applied pointwise to the
Schwartz Fourier transform, gives
\[
 \|\Delta_{m+j_0}W_N\|_\infty\ge cN^{3/2}.
\]
Consequently
\[
 \lambda_{m+j_0}^{-1}
 \|\Delta_{m+j_0}u_N\|_\infty\ge cN^{1/2}
\]
for large $N$; the fixed low polynomial does not enter this shell.  It is
therefore viscosity-active, proving the lower bound in
\eqref{eq:v9-index-window}.

For the upper bound, the Schwartz estimates imply, for every $L>0$ and
every shell with $\lambda_q=MN$ and $M\ge1$,
\[
 \lambda_q^{-1}\|\Delta_qW_N\|_\infty
 \le C_LN^{1/2}M^{2-L}.
\]
Choose $L>2+(2\varepsilon)^{-1}$.  Uniformly for
$M\ge N^\varepsilon$, the right-hand side tends to zero and is eventually
smaller than $c_\theta\nu$.  The reservoir has no Fourier modes there.
This proves the upper bound.

By \eqref{eq:v9-packet-canonical-cutoff},
$K_N/N\to\delta\kappa_U$.  Choose
$\delta<c_1/(4\kappa_U)$ to obtain the first strict inequality in
\eqref{eq:v9-index-intermediate-fan}.  Also
\[
 N_N^{\rm vf}
 \asymp N^{5/4},
\]
by \eqref{eq:v9-packet-record}.  Since
$1+\varepsilon<5/4$, the upper bound in
\eqref{eq:v9-index-window} is eventually below $N_N^{\rm vf}$.
\end{proof}

\begin{firewall}[The index estimate uses the actual definition]
The lower bound in \eqref{eq:v9-index-window} exhibits a shell whose
quantity $\lambda_q^{-1}\|u_q\|_\infty$ exceeds the viscosity threshold.
The upper bound verifies the threshold on every more remote shell.  Thus
\eqref{eq:v9-index-intermediate-fan} is not an assignment of a nominal
packet scale to $Q_\nu$.
\end{firewall}

% =============================================================================
\section{The outer--outer branch also carries positive excess}
\label{sec:v9-outer-outer}
% =============================================================================

Let $\mathsf Q_\rho$ denote the Euclidean Fourier projection to
$|\xi|>\rho$.  The cubic functional
$\mathscr G_{\R^3}$ is continuous in $H^2$, and
$\mathsf Q_\rho U\to U$ in $H^2$ as $\rho\downarrow0$.  Shrink $\delta$,
if necessary, so that with
\begin{equation}
 \rho:=\delta\kappa_U
 \label{eq:v9-Euclidean-cutoff-ratio}
\end{equation}
one has
\begin{equation}
 c_{\Torus}\mathscr G_{\R^3}(\mathsf Q_\rho U)>\frac12G_U.
 \label{eq:v9-positive-highpass-seed}
\end{equation}
Choose $\rho$ to be a continuity radius for the Fourier integrals, which
is possible after an arbitrarily small further change of $\delta$.

\begin{proposition}[Positive outer--outer production]
\label{prop:v9-outer-outer-positive}
Let
\[
 w_N:=Q_{K_N}u_N,\qquad
 \eta_N:=\operatorname{curl}w_N.
\]
Then
\begin{align}
 \boxed{
 \int_{\Torus^3}S(w_N):(\eta_N\otimes\eta_N)\,\dd x
 &=\beta^3N^{9/2}c_{\Torus}
   \mathscr G_{\R^3}(\mathsf Q_\rho U)
   +o(N^{9/2})>0,}
 \label{eq:v9-outer-outer-production}\\
 \|Aw_N\|_2^2
 &=\beta^2N^4c_{\Torus}
   \|\Delta\mathsf Q_\rho U\|_{L^2(\R^3)}^2+o(N^4).
 \label{eq:v9-outer-outer-palinstrophy}
\end{align}
Consequently, for all sufficiently large $N$,
\begin{equation}
 \boxed{
 \int_{\Torus^3}S(w_N):(\eta_N\otimes\eta_N)\,\dd x
 -\nu\|Aw_N\|_2^2>0.}
 \label{eq:v9-outer-outer-excess}
\end{equation}
This is the outer--outer stretching term
$\langle(\eta\cdot\nabla)w,\eta\rangle$ in
\eqref{eq:v8-high-vorticity-remainder}.  It cannot be charged to the V34
raw head--head daughter.
\end{proposition}

\begin{proof}
For large $N$, the fixed reservoir lies below $K_N$ and is removed by
$Q_{K_N}$.  Under the rescaling $y=N(x-x_0)$, the cutoff ratio tends to
$K_N/N\to\rho$.  The Fourier Riemann sums converge in $H^2$ to
$\mathsf Q_\rho U$ at a continuity radius.  Continuity of the trilinear
gradient functional in $H^2$ and the concentration factor $N^{9/2}$ give
\eqref{eq:v9-outer-outer-production}.  Its sign follows from
\eqref{eq:v9-positive-highpass-seed}.  The final identification is the
third stretching term in \eqref{eq:v8-high-vorticity-remainder}.
\end{proof}

\begin{corollary}[Finite Fourier obstructions exist]
\label{cor:v9-finite-Fourier-obstruction}
Fix the $\varepsilon\in(0,1/4)$ used in
Theorem~\ref{thm:v9-index-in-fan}.  For every sufficiently large $N$, the
Fourier truncation of $u_N$ at a dyadic radius
$M_N\asymp N^{1+\varepsilon}$, followed by an arbitrarily small scalar
renormalization, is a real divergence-free trigonometric polynomial with
energy $E_*$ which retains
\eqref{eq:v9-reservoir-positive-excess},
\eqref{eq:v9-index-intermediate-fan}, and positive outer--outer
production.
\end{corollary}

\begin{proof}
Because $\widehat U$ is Schwartz, truncation at
$M_N\asymp N^{1+\varepsilon}$ removes an error smaller than any prescribed
power of $N$ in each fixed Sobolev norm.  The production, palinstrophy,
critical record, and outer--outer trilinear form therefore retain their
strict leading inequalities.  This cutoff lies above the explicit active
shell used for the lower bound in \eqref{eq:v9-index-window}, while all
higher shells vanish, so both index inequalities persist.  Energy
normalization is a scalar perturbation tending to one.
\end{proof}

\begin{assessment}[The pointwise route is closed]
Theorem~\ref{thm:v9-concentration-exponents} and
Theorem~\ref{thm:v9-index-in-fan} give actual smooth, fixed-energy data in
the precise unresolved fan for which nonlinear enstrophy production
exceeds the complete viscous palinstrophy.  Proposition
\ref{prop:v9-outer-outer-positive} shows that the far-shell
outer--outer term can itself have the adverse sign.  Therefore no
unconditional pointwise depletion or viscous-absorption theorem of the
type tested in V8 can close the programme.
\end{assessment}

% =============================================================================
\section{The lifetime and Leray-debit scales}
\label{sec:v9-lifetime}
% =============================================================================

The obstruction is pointwise.  Its time scale must be read before it is
used on a terminal sequence.  Put
\begin{equation}
 Y_N:=\|A^{1/2}u_N\|_2^2,\qquad
 X_N:=\mathcal G(u_N)-\nu\|Au_N\|_2^2>0
 \label{eq:v9-initial-growth-stocks}
\end{equation}
and define the diagnostic turnover scale
\begin{equation}
 \tau_N^{\rm turn}:=\frac{Y_N}{2X_N}.
 \label{eq:v9-turnover-time}
\end{equation}

\begin{proposition}[Turnover and energy-debit arithmetic]
\label{prop:v9-turnover-ledger}
For the collective data,
\begin{align}
 Y_N&\asymp N^2,
 \label{eq:v9-enstrophy-scale}\\
 X_N&\asymp N^{9/2},
 \label{eq:v9-excess-scale}\\
 \tau_N^{\rm turn}&\asymp N^{-5/2},
 \label{eq:v9-turnover-scale}\\
 \nu Y_N\tau_N^{\rm turn}&\asymp\nu N^{-1/2}
 \asymp\frac{\nu}{\mathcal R_N}.
 \label{eq:v9-turnover-energy-debit}
\end{align}
The V34 parabolic scale at $K_N\asymp N$ is of order
$(\nu N^2)^{-1}$, which is longer than
$\tau_N^{\rm turn}$ by a factor of order $N^{1/2}/\nu$.
For doubled critical records, the model energy debits in
\eqref{eq:v9-turnover-energy-debit} are summable.
\end{proposition}

\begin{proof}
The first two rows follow from
\eqref{eq:v9-packet-enstrophy},
\eqref{eq:v9-reservoir-production}, and
\eqref{eq:v9-reservoir-palinstrophy}.  Insert them in
\eqref{eq:v9-turnover-time}.  Since
$\mathcal R_N\asymp N^{1/2}$ by
\eqref{eq:v9-packet-record}, the last equivalence follows.  If
$\mathcal R_n=2^n\mathcal R_*$, then
$\sum_n\mathcal R_n^{-1}<\infty$.
\end{proof}

\begin{firewall}[A diagnostic scale is not a persistence theorem]
The definition \eqref{eq:v9-turnover-time} is the initial enstrophy divided
by its initial growth rate.  It does not assert that $X_N(t)$ remains
positive for that long.  Establishing an actual uniform fraction of this
lifetime requires controlling the rescaled Navier--Stokes evolution.  The
proposition proves the upstream arithmetic: even if such persistence is
obtained, its direct Leray-energy debit is summable on doubled records.
\end{firewall}

% =============================================================================
\section{Integrated V9 conclusion and next acquisition order}
\label{sec:v9-conclusion}
% =============================================================================

\begin{theorem}[What V9 proves]
\label{thm:v9-conclusion}
Relative to V1--V8 and the frozen source, the following are proved.
\begin{enumerate}[label=\textup{(V9.\arabic*)},leftmargin=5em]
\item A compactly supported divergence-free seed with positive nonlinear
      enstrophy production exists.
\item Three-dimensional concentration at fixed kinetic energy produces
      $N^{9/2}$ nonlinear production against $N^4$ viscous palinstrophy,
      and therefore actual smooth data with positive viscous excess.
\item After adding a low-frequency energy reservoir, the canonical cutoff
      is $K_N\asymp N$, while the actual dissipation index lies strictly in
      the V8 intermediate fan.
\item The outer--outer branch has strictly positive leading production;
      finite Fourier approximations retain all strict obstructions.
\item The natural turnover scale is $N^{-5/2}$ and its direct Leray debit
      is $O(\nu/\mathcal R_N)$, summable at doubled records.
\end{enumerate}
No full NS stability theorem is obtained.
\end{theorem}

\begin{proof}
Item \textup{(V9.1)} is
Lemma~\ref{lem:v9-compact-positive-seed}.  Item \textup{(V9.2)} is
Theorem~\ref{thm:v9-concentration-exponents}.  Item \textup{(V9.3)} is
Proposition~\ref{prop:v9-canonical-packet} and
Theorem~\ref{thm:v9-index-in-fan}.  Item \textup{(V9.4)} is
Proposition~\ref{prop:v9-outer-outer-positive} and
Corollary~\ref{cor:v9-finite-Fourier-obstruction}.  Item
\textup{(V9.5)} is Proposition~\ref{prop:v9-turnover-ledger}.
\end{proof}

\begin{programme}[V10 acquisition order]
\label{prog:v10}
\begin{enumerate}[label=\textup{(X\arabic*)},leftmargin=4em]
\item Stop seeking an unconditional pointwise depletion inequality; V9
      gives strict counterexamples within the exact unresolved fan.
\item Rescale the Navier--Stokes equation around $u_N$ using
      $y=N(x-x_0)$ and the turnover time $s=N^{5/2}t$.  Identify the
      limiting inviscid equation and the vanishing effective viscosity.
\item Prove, or sharply delimit, persistence of positive excess on a fixed
      rescaled time interval.  Do not infer persistence from continuity
      without a uniform-in-$N$ estimate.
\item Compare the resulting physical interval with the V34 intrinsic core.
      Its direct energy debit is already known to be summable, so the target
      must be record growth, source incidence, or inter-card flux rather
      than Leray dissipation alone.
\item Track the spatial packet through a material cell only after paying
      the V8 localization commutator or using a smooth guarded multiplier.
\item Keep the represented-history branch at the V3 physical-incidence
      stopping line.
\end{enumerate}
\end{programme}

\begin{assessment}[Position after V9]
The remaining obstruction is no longer hypothetical: coherent
three-dimensional concentration gives actual fixed-energy, intermediate-fan
data whose nonlinear enstrophy production beats viscosity, including on
the outer--outer branch.  This rules out the entire pointwise absorption
strategy.  The next honest question is dynamical and upstream: how long
the packet survives, how its critical record changes, and whether its
short interval can be charged across the terminal sequence by something
stronger than the summable Leray debit.
\end{assessment}

% =============================================================================
\section*{Version V9 history}
\addcontentsline{toc}{section}{Version V9 history}
% =============================================================================

Constructed a compactly supported divergence-free positive-production seed
by localizing the explicit V8 periodic field through a vector potential.
Compressed that seed into a single three-dimensional packet and proved
the exact fixed-energy exponents: $N^{9/2}$ nonlinear enstrophy production
versus $N^4$ viscous palinstrophy.  Added a low-frequency energy reservoir
and proved that the canonical cutoff and actual dissipation index place
the positive-excess data inside the V8 intermediate fan.  Proved positive
outer--outer production and persistence under finite Fourier
approximation.  Finally identified the $N^{-5/2}$ turnover scale and its
summable $O(\nu/\mathcal R_N)$ direct energy debit.  No global regularity
claim was made.
% Future versions are appended immediately above the sole terminator.
% =============================================================================
% VERSION V10 APPEND
% =============================================================================

\section{V10: turnover rescaling on expanding tori}
\label{sec:v10-rescaling}
% =============================================================================

V9 produced pointwise positive viscous excess and identified the diagnostic
time $N^{-5/2}$.  Continuity for each fixed datum is not enough: the
interval could shrink faster with $N$.  We now rescale the actual
Navier--Stokes solution and prove a uniform interval.

Let $u^{(N)}$ be the smooth local solution with initial datum $u_N$ from
\eqref{eq:v9-energy-normalized-data}.  Put
\begin{equation}
 y=N(x-x_0),\qquad
 s=\beta N^{5/2}t,\qquad
 V_N(s,y):=\beta^{-1}N^{-3/2}u^{(N)}(t,x),
 \label{eq:v10-rescaled-variables}
\end{equation}
on the expanding torus
\[
 \Torus_N^3:=(\R/2\pi N\mathbb Z)^3.
\]
Then
\begin{equation}
 \boxed{
 \partial_sV_N+\PP_N\operatorname{div}_y(V_N\otimes V_N)
 =\varepsilon_N\Delta_yV_N,\qquad
 \varepsilon_N:=\frac{\nu}{\beta N^{1/2}}\longrightarrow0.}
 \label{eq:v10-rescaled-NS}
\end{equation}
Its initial datum is
\begin{equation}
 V_N(0,y)
 =c_N\left[
 U(y)+\beta^{-1}N^{-3/2}
 \phi(x_0+y/N)
 \right].
 \label{eq:v10-rescaled-initial-data}
\end{equation}
Here $U$ is extended by zero on $\Torus_N^3$ and $c_N\to1$.

\begin{lemma}[Uniform large-torus local control]
\label{lem:v10-uniform-local-control}
Fix an integer $m\ge5$.  There are $s_0>0$ and $M<\infty$, independent of
all sufficiently large $N$, such that the rescaled solution exists on
$[0,s_0]$ and
\begin{equation}
 \boxed{
 \sup_{0\le s\le s_0}\|V_N(s)\|_{H^m(\Torus_N^3)}
 +\sup_{0\le s\le s_0}
       \|\partial_sV_N(s)\|_{H^{m-2}(\Torus_N^3)}
 \le M.}
 \label{eq:v10-uniform-Hm}
\end{equation}
In particular,
\begin{equation}
 \|V_N(s)-V_N(0)\|_{H^2(\Torus_N^3)}\le Ms.
 \label{eq:v10-uniform-H2-continuity}
\end{equation}
\end{lemma}

\begin{proof}
The compact seed contributes a fixed $H^m$ norm.  For the slow reservoir,
a $j$th $y$-derivative has size
$O(N^{-3/2-j})$ on a domain of volume $O(N^3)$, so its $L^2$ norm is
$O(N^{-j})$.  Thus the initial $H^m$ norms are uniformly bounded.

Sobolev product and embedding constants on flat tori of side at least
$2\pi$ may be chosen uniformly: cover by unit cubes with bounded overlap
and apply the Euclidean local inequalities.  The usual commutator estimate
for \eqref{eq:v10-rescaled-NS} therefore gives
\[
 \frac{\dd}{\dd s}\|V_N\|_{H^m}^2
 +2\varepsilon_N\|\nabla V_N\|_{H^m}^2
 \le C\|\nabla V_N\|_\infty\|V_N\|_{H^m}^2
 \le C\|V_N\|_{H^m}^3,
\]
with $C$ independent of $N$.  The standard continuation argument gives a
common $s_0$ and the first bound.  The equation, uniform boundedness of the
Leray multiplier, and the product estimate give the time-derivative bound.
Integrating it in $H^{m-2}$, with $m-2\ge2$, proves
\eqref{eq:v10-uniform-H2-continuity}.
\end{proof}

Equation~\eqref{eq:v10-rescaled-NS} is the Euler equation plus a viscosity
$\varepsilon_N$ tending to zero.  The next result uses only the uniform
estimate, rather than assuming convergence.

Define on $\Torus_N^3$
\begin{align}
 \mathscr G_N(V)
 &:=c_{\Torus}\int_{\Torus_N^3}
 S(V):\bigl(\operatorname{curl}V\otimes
            \operatorname{curl}V\bigr)\,\dd y,
 \label{eq:v10-rescaled-production}\\
 \mathscr Z_N(V)
 &:=c_{\Torus}\int_{\Torus_N^3}|\Delta V|^2\,\dd y.
 \label{eq:v10-rescaled-palinstrophy}
\end{align}
The physical excess has the exact form
\begin{equation}
 \boxed{
 \mathcal G(u^{(N)}(t))-\nu\|Au^{(N)}(t)\|_2^2
 =\beta^3N^{9/2}
   \left[
    \mathscr G_N(V_N(s))
    -\varepsilon_N\mathscr Z_N(V_N(s))
   \right].}
 \label{eq:v10-physical-rescaled-excess}
\end{equation}

% =============================================================================
\section{Uniform persistence of positive excess}
\label{sec:v10-persistence}
% =============================================================================

\begin{theorem}[A genuine turnover interval]
\label{thm:v10-uniform-persistence}
There are constants $s_*>0$, $c_*>0$, and $N_0$ such that, for every
dyadic $N\ge N_0$ and
\begin{equation}
 J_N:=
 \left[0,\frac{s_*}{\beta N^{5/2}}\right],
 \label{eq:v10-physical-turnover-interval}
\end{equation}
the actual solution satisfies
\begin{equation}
 \boxed{
 \mathcal G(u^{(N)}(t))-\nu\|Au^{(N)}(t)\|_2^2
 \ge c_*\beta^3G_UN^{9/2}
 \quad(t\in J_N).}
 \label{eq:v10-persistent-excess}
\end{equation}
Consequently,
\begin{equation}
 \boxed{
 \|A^{1/2}u^{(N)}(\sup J_N)\|_2^2
 -\|A^{1/2}u_N\|_2^2
 \ge c_*\beta^2G_UN^2.}
 \label{eq:v10-enstrophy-gain}
\end{equation}
Thus the V9 diagnostic scale is an actual uniform fraction of a turnover
time.
\end{theorem}

\begin{proof}
The trilinear functional $\mathscr G_N$ is locally Lipschitz in $H^2$,
uniformly in $N$.  Indeed, $H^2$ controls first derivatives in $L^3$ by
interpolation between $L^2$ and $L^6$, so
\[
 |\mathscr G_N(V)-\mathscr G_N(W)|
 \le C_M\|V-W\|_{H^2}
\]
when both $H^2$ norms are bounded by $M$.  From
\eqref{eq:v9-reservoir-production} and the rescaling,
\[
 \mathscr G_N(V_N(0))=G_U+o(1).
\]
Also Lemma~\ref{lem:v10-uniform-local-control} gives
$\mathscr Z_N(V_N(s))\le C_M$ on $[0,s_0]$.

Choose $s_*\le s_0$ so that $C_MM s_*\le G_U/4$.  Then choose $N_0$ so
large that the initial $o(1)$ has magnitude at most $G_U/4$ and
$\varepsilon_NC_M\le G_U/4$.  Equations
\eqref{eq:v10-uniform-H2-continuity} and
\eqref{eq:v10-physical-rescaled-excess} then give
\[
 \mathscr G_N(V_N(s))
 -\varepsilon_N\mathscr Z_N(V_N(s))
 \ge\frac{G_U}{4}
\]
on $0\le s\le s_*$.  This proves
\eqref{eq:v10-persistent-excess}.  Integrate the exact enstrophy identity
over $J_N$ and use $|J_N|=s_*/(\beta N^{5/2})$ to obtain
\eqref{eq:v10-enstrophy-gain}, after changing the fixed constant.
\end{proof}

\begin{corollary}[Local inviscid limit]
\label{cor:v10-local-Euler-limit}
After passage to a subsequence, $V_N$ converges on compact subsets of
$[0,s_*]\times\R^3$ to the unique smooth Euler solution with initial datum
$U$.  The convergence holds strongly in
$C([0,s_*];H^{m-1}_{\rm loc})$ after decreasing $s_*$ if necessary.
\end{corollary}

\begin{proof}
The uniform spatial and time-derivative bounds give local compactness on
the expanding tori.  The slow reservoir in
\eqref{eq:v10-rescaled-initial-data} tends to zero in every local Sobolev
space, and $\varepsilon_N\to0$.  Passing to the weak formulation gives
Euler with datum $U$.  Smooth Euler uniqueness identifies every
subsequence limit, and the uniform higher norm upgrades compact
convergence to the displayed topology.
\end{proof}

% =============================================================================
\section{Corridor, core, and energy arithmetic on the real interval}
\label{sec:v10-card-arithmetic}
% =============================================================================

\begin{proposition}[Actual packet-card ledger]
\label{prop:v10-packet-card-ledger}
After reducing $s_*$ by a fixed factor if needed, uniformly on $J_N$,
\begin{equation}
 c\mathcal R_N
 \le\|A^{1/4}u^{(N)}(t)\|_2
 \le C\mathcal R_N,
 \qquad
 \mathcal R_N\asymp N^{1/2},
 \qquad
 K_N\asymp N.
 \label{eq:v10-record-corridor}
\end{equation}
Moreover,
\begin{align}
 |J_N|&\asymp N^{-5/2},
 \label{eq:v10-card-length}\\
 \mathcal R_N\sqrt{|J_N|}
 &\asymp N^{-3/4}
 \asymp\mathcal R_N^{-3/2},
 \label{eq:v10-core-summand}\\
 \nu K_N^2|J_N|
 &\asymp\nu N^{-1/2},
 \label{eq:v10-parabolic-filling}\\
 \nu\int_{J_N}\|A^{1/2}u^{(N)}(t)\|_2^2\,\dd t
 &\asymp\nu N^{-1/2}
 \asymp\frac{\nu}{\mathcal R_N}.
 \label{eq:v10-actual-Leray-debit}
\end{align}
Thus both the V34-type core summand and the direct Leray debit are
summable along doubled records.
\end{proposition}

\begin{proof}
The $H^2$ continuity in rescaled variables preserves a fixed upper and
lower bound for the seed's homogeneous $H^{1/2}$ and $H^1$ norms on
$[0,s_*]$.  Scaling back gives
\[
 \|A^{1/4}u^{(N)}(t)\|_2^2\asymp N,
 \qquad
 \|A^{1/2}u^{(N)}(t)\|_2^2\asymp N^2.
\]
The first statement is also consistent with
\eqref{eq:v9-packet-record}.  Insert these estimates and
$|J_N|\asymp N^{-5/2}$ in the four displayed quantities.  If the records
double, $\sum\mathcal R_N^{-3/2}$ and
$\sum\mathcal R_N^{-1}$ converge.
\end{proof}

\begin{remark}[The V34 intrinsic core is not identified]
The interval $J_N$ is generated by the packet dynamics.  It has arithmetic
compatible with the V34 core, but it is not thereby equal to
$J_n^{\rm occ}$, which is selected from the history-fresh source residual.
No source or endpoint incidence between these intervals is asserted.
\end{remark}

% =============================================================================
\section{Critical and weak source actions on the turnover interval}
\label{sec:v10-source-actions}
% =============================================================================

Define
\begin{align}
 \mathscr S_N^{\rm crit}
 &:=\int_{J_N}
       \|A^{-1/4}\BNS(u^{(N)}(t))\|_2^2\,\dd t,
 \label{eq:v10-critical-source-action}\\
 \mathscr S_N^{\rm weak}
 &:=\int_{J_N}
       \|A^{-3/4}\BNS(u^{(N)}(t))\|_2^2\,\dd t.
 \label{eq:v10-weak-source-action}
\end{align}

\begin{theorem}[One-derivative source-action separation]
\label{thm:v10-source-action-separation}
After a fixed reduction of $s_*$ if necessary,
\begin{align}
 \boxed{
 \mathscr S_N^{\rm crit}\asymp N^{3/2},}
 \label{eq:v10-critical-action-scale}\\
 \boxed{
 \mathscr S_N^{\rm weak}\le C N^{-1/2},}
 \label{eq:v10-weak-action-scale}\\
 \boxed{
 \int_{J_N}\|A^{1/2}u^{(N)}(t)\|_2^4\,\dd t
 \asymp N^{3/2}.}
 \label{eq:v10-enstrophy-square-scale}
\end{align}
In particular,
\begin{equation}
 \boxed{
 \frac{\mathscr S_N^{\rm crit}}
      {\nu\mathcal R_N^2}
 \asymp\frac{N^{1/2}}{\nu}\longrightarrow\infty,}
 \label{eq:v10-critical-action-overdrive-scale}
\end{equation}
whereas the weak action is bounded by the same summable scale as the Leray debit.
\end{theorem}

\begin{proof}
Under \eqref{eq:v10-rescaled-variables}, the physical critical source has
squared norm
\[
 \|A^{-1/4}\BNS(u^{(N)})\|_2^2
 =\beta^4N^4
   \|A_{\Torus_N}^{-1/4}
       \PP_N\operatorname{div}(V_N\otimes V_N)\|_2^2,
\]
with the fixed normalized-volume convention understood.  The weak source
has the corresponding factor $\beta^4N^2$ with exponent $-3/4$.
At $s=0$ the rescaled critical source converges to the nonzero source
$A^{-1/4}\BNS(U)$; it is nonzero because $G_U>0$.  Uniform $H^2$
continuity preserves fixed upper and lower bounds for this norm on
$0\le s\le s_*$.  The weak source is uniformly bounded by the
large-torus version of Proposition~\ref{prop:v7-weak-source-budget}.
Since $\dd t=(\beta N^{5/2})^{-1}\dd s$, the first scale and the stated
weak upper bound follow.

The rescaled $H^1$ norm stays bounded above and below, so physical
enstrophy is comparable to $N^2$ throughout $J_N$.  Squaring and
integrating over a length comparable to $N^{-5/2}$ proves
\eqref{eq:v10-enstrophy-square-scale}.  Finally use
$\mathcal R_N^2\asymp N$.
\end{proof}

\begin{assessment}[The V7 derivative deficit is dynamically saturated]
V7 proved abstractly that active-shell residence asks for the critical
source while energy pays only the one-derivative-weaker source.  V10
exhibits an actual solution interval on which the critical action is at
least $N^2$ above the proved weak-action upper scale.  The critical action and enstrophy-square stock have the
same $N^{3/2}$ scale, while both the weak action and Leray debit are at
most of order $N^{-1/2}$.
\end{assessment}

% =============================================================================
\section{Exact V34 source-incidence consequence}
\label{sec:v10-v34-incidence}
% =============================================================================

Before comparing source labels, the output projection must be proved.  Let
$\mathsf Q_\rho$ be the Euclidean high-pass from V9 and define
\begin{equation}
 \mathscr G_\rho^{\rm hi}(U)
 :=-c_{\Torus}
 \left\langle
  \mathsf Q_{2\rho}\BNS(\mathsf Q_\rho U),
  -\Delta\mathsf Q_\rho U
 \right\rangle_{L^2(\R^3)}.
 \label{eq:v10-Euclidean-high-output-production}
\end{equation}
Because $\mathsf Q_\rho U\to U$ in every fixed Sobolev norm and
$\mathsf Q_{2\rho}\BNS(\mathsf Q_\rho U)\to\BNS(U)$ in $L^2$,
\begin{equation}
 \mathscr G_\rho^{\rm hi}(U)\longrightarrow G_U
 \qquad(\rho\downarrow0).
 \label{eq:v10-high-output-limit}
\end{equation}
Shrink the already free energy fraction $\delta$, if necessary, so that
with $\rho=\delta\kappa_U$,
\begin{equation}
 \mathscr G_\rho^{\rm hi}(U)>\frac12G_U.
 \label{eq:v10-positive-high-output-seed}
\end{equation}

\begin{proposition}[High-output outer overdrive is not raw-daughter
overdrive]
\label{prop:v10-not-raw-overdrive}
Put $w_N(t)=Q_{K_N}u^{(N)}(t)$.  After reducing $s_*$ by a fixed factor,
the actual solution satisfies throughout $J_N$
\begin{equation}
 \boxed{
 -\left\langle
 Q_{2K_N}\BNS(w_N(t)),Aw_N(t)
 \right\rangle
 \ge c\beta^3G_UN^{9/2}.}
 \label{eq:v10-persistent-high-output-production}
\end{equation}
Moreover the corresponding high-output critical source has action
\begin{equation}
 \boxed{
 \int_{J_N}
 \|A^{-1/4}Q_{2K_N}\BNS(w_N(t))\|_2^2\,\dd t
 \asymp N^{3/2}.}
 \label{eq:v10-high-output-outer-source-action}
\end{equation}
By contrast, the V34 raw daughter is
\[
 d_N^{\rm raw}
 =A^{-1/4}Q_{K_N}\BNS(P_{K_N}u^{(N)}),
\]
whose Fourier support is contained in $K_N<|k|\le2K_N$.  The V34 incoming
head source
\[
 A^{-1/4}P_{K_N}
 \bigl(\BNS(u^{(N)})-\BNS(P_{K_N}u^{(N)})\bigr)
\]
also discards the source in
\eqref{eq:v10-high-output-outer-source-action}.  Therefore neither the
full overdrive scale \eqref{eq:v10-critical-action-overdrive-scale} nor its
high-output subaction can be relabeled as overdrive of either V34 source.
\end{proposition}

\begin{proof}
At $s=0$, rescaling the left side of
\eqref{eq:v10-persistent-high-output-production} and using
$K_N/N\to\rho$ gives
\[
 \beta^{-3}N^{-9/2}
 \left[
 -\langle Q_{2K_N}\BNS(w_N),Aw_N\rangle
 \right]
 \longrightarrow\mathscr G_\rho^{\rm hi}(U)>0.
\]
The $H^2$ continuity used in
Theorem~\ref{thm:v10-uniform-persistence} preserves half of this margin on
a common rescaled interval.  The same scaling gives an instantaneous
squared critical-source norm of order $N^4$, bounded above and below
uniformly after division by $N^4$.  Integration over
$|J_N|\asymp N^{-5/2}$ proves
\eqref{eq:v10-high-output-outer-source-action}.

The raw support claim follows from convolution of two parents with
frequencies at most $K_N$.  Its output cannot exceed $2K_N$.  The incoming
operator is projected into the head.  The source in
\eqref{eq:v10-high-output-outer-source-action} has two outer parents and
output above $2K_N$, so it belongs to neither operator.
\end{proof}

\begin{firewall}[The missing row is an all-shell source incidence]
The packet proves that a complete critical source can be overdriven while
the V34 raw and incoming-head operators omit the responsible high-output
outer--outer component.  This does not refute the V34 alternative, which
retains earlier all-shell, transfer, forcing, and typed-complement exits.
It identifies the next required map: the actual high-output outer source
must be placed into one of those inherited rows with its original
occurrence time and endpoint, rather than inserted into the bounded raw
action window.
\end{firewall}

% =============================================================================
\section{Material displacement during one turnover}
\label{sec:v10-displacement}
% =============================================================================

\begin{proposition}[Turnover-scale flow displacement]
\label{prop:v10-flow-displacement}
Let $X_N(t;x)$ be the classical flow of $u^{(N)}$ on $J_N$.  Then
\begin{equation}
 \boxed{
 \sup_x\sup_{t\in J_N}|X_N(t;x)-x|
 \le\frac{C}{N}.}
 \label{eq:v10-flow-displacement}
\end{equation}
The heat length on the same interval is
\begin{equation}
 \sqrt{\nu|J_N|}\asymp
 \nu^{1/2}N^{-5/4}=o(N^{-1}).
 \label{eq:v10-heat-length}
\end{equation}
Thus the natural spatial packet and its material displacement have the
same $N^{-1}$ order, while diffusion is smaller.
\end{proposition}

\begin{proof}
Uniform Sobolev embedding in the rescaled variables gives
$\|V_N(s)\|_\infty\le C$.  Hence
$\|u^{(N)}(t)\|_\infty\le C\beta N^{3/2}$ on $J_N$.  Integrating over
$|J_N|=s_*/(\beta N^{5/2})$ proves
\eqref{eq:v10-flow-displacement}.  The heat length follows directly from
the interval length.
\end{proof}

\begin{firewall}[Displacement is not same-cell capture]
The estimate locates trajectories relative to their starting points.  It
does not identify the packet centre with the inherited V34 material cell,
nor control the sharp spectral/localization commutator from V8.  A
same-cell conclusion still requires an actual initial incidence and a
guarded localization estimate.
\end{firewall}

% =============================================================================
\section{Integrated V10 conclusion and next acquisition order}
\label{sec:v10-conclusion}
% =============================================================================

\begin{theorem}[What V10 proves]
\label{thm:v10-conclusion}
Relative to V1--V9 and the frozen source, the following are proved.
\begin{enumerate}[label=\textup{(V10.\arabic*)},leftmargin=5.4em]
\item Turnover rescaling gives Navier--Stokes on expanding tori with
      effective viscosity $\nu/(\beta\sqrt N)$ and uniform local Sobolev
      control.
\item Positive viscous excess persists on an actual interval of length
      $cN^{-5/2}$ and produces an $O(N^2)$ enstrophy increase.
\item The solution remains in a comparable critical-record corridor, while
      its core summand and direct Leray debit are summable at doubled
      records.
\item The critical nonlinear source action has scale $N^{3/2}$, while the
      weak action is $O(N^{-1/2})$, dynamically realizing the V7
      one-derivative deficit.
\item The responsible high-output outer--outer source belongs to neither
      the V34 raw daughter nor the incoming head source; an all-shell
      incidence map is required.
\item Material trajectories move only $O(N^{-1})$ during a turnover, but
      same-cell capture still requires the V8 commutator payment.
\end{enumerate}
No full NS stability theorem is obtained.
\end{theorem}

\begin{proof}
Items \textup{(V10.1)--(V10.2)} are
Lemma~\ref{lem:v10-uniform-local-control} and
Theorem~\ref{thm:v10-uniform-persistence}.  Item \textup{(V10.3)} is
Proposition~\ref{prop:v10-packet-card-ledger}.  Item \textup{(V10.4)} is
Theorem~\ref{thm:v10-source-action-separation}.  Item \textup{(V10.5)} is
Proposition~\ref{prop:v10-not-raw-overdrive}.  Item \textup{(V10.6)} is
Proposition~\ref{prop:v10-flow-displacement}.
\end{proof}

\begin{programme}[V11 acquisition order]
\label{prog:v11}
\begin{enumerate}[label=\textup{(Y\arabic*)},leftmargin=4em]
\item Use the persistent packet only as an obstruction; do not infer that
      an arbitrary hypothetical terminal solution contains this profile.
\item Compute the critical-record production
      $-\langle\BNS(U),(-\Delta)^{1/2}U\rangle$ of a compact seed.  Determine
      whether positive enstrophy turnover forces, permits, or is independent
      of growth of the renewal record.
\item Trace the high-output outer--outer source
      \eqref{eq:v10-high-output-outer-source-action} through the inherited
      all-shell, transfer, and forcing alternatives, preserving its actual
      time and endpoint.
\item If no inherited row accepts that operator, record a precise source
      atlas gap rather than assigning it to the raw action window.
\item Compare a genuine record-crossing interval with the packet lifetime.
      Leray debit and scalar core summability are already proved
      insufficient; the target must be signed inter-card flux or a new
      source-incidence payment.
\item Use material localization only with an explicit guard band or
      commutator action, and keep the represented-history branch at the V3
      physical-incidence stopping line.
\end{enumerate}
\end{programme}

\begin{assessment}[Position after V10]
The collective obstruction is now dynamical, not merely pointwise:
positive excess survives for a uniform rescaled time and changes
enstrophy by its full packet scale.  Yet all energy-level costs remain
summable, while the critical source action is an overdriven high-output
outer--outer quantity omitted by both V34 source operators tested so far.
The next upstream issue is whether this packet changes the actual critical
record and where its source occurrence sits in the inherited all-shell
atlas.
\end{assessment}

% =============================================================================
\section*{Version V10 history}
\addcontentsline{toc}{section}{Version V10 history}
% =============================================================================

Rescaled the actual concentrated solutions on expanding tori and proved
uniform high-Sobolev control with effective viscosity tending to zero.
This yielded a genuine $cN^{-5/2}$ interval of positive viscous excess and
an $O(N^2)$ enstrophy increase.  Established the comparable critical-record
corridor, summable core and Leray arithmetic, and the dynamically sharp
$N^{3/2}$ critical-source scale versus the $O(N^{-1/2})$ weak-source
bound.  Proved that the responsible high-output outer--outer source is
contained in neither the V34 raw daughter nor its incoming head source,
and bounded material displacement by $O(N^{-1})$ without asserting
same-cell capture.  No global regularity claim was made.
% Future versions are appended immediately above the sole terminator.

% =============================================================================
% VERSION V11 APPEND
% =============================================================================

\section{V11: the assembled far-output source and the critical record}
\label{sec:v11-entry}
% =============================================================================

V10 exhibited a positive outer--outer contribution above \(2K_N\).  That
calculation opened the parent decomposition before checking cancellation
against the mixed parents.  The inherited programme correctly requires
complete signed assembly before a component is assigned a source price.
This version repairs that point.  It replaces the opened term by the full
operator
\[
 A^{-1/4}Q_{2K}\BNS(u),
\]
proves that this assembled far-output source persists on the actual packet
interval, and carries it to its own Duhamel endpoint.

A second issue is whether the packet merely increases enstrophy while
remaining tangent to the critical record.  We refine the compact seed so
that its nonlinear \(H^{1/2}\) production is positive as well.  This gives
a genuine fixed-fraction critical-record increase on the same turnover
scale.  The V5 triad still proves that positive enstrophy production does
not force such an increase for arbitrary data.

Finally, we write the complete five-piece primal source atlas.  It shows
that the V34 raw daughter and incoming-head source omit two exterior-output
terms.  The V27--V29 transfer and forcing rows are daughter-octave adjoint
currents and do not supply those primal terms.  Thus the exact remaining
gap is an all-output source-incidence lemma, not a relabelling of an
existing finite-head row.

% =============================================================================
\section{A compact seed with two positive nonlinear moments}
\label{sec:v11-two-moment-seed}
% =============================================================================

For a compact divergence-free field on \(\R^3\), define
\begin{equation}
 \mathscr C_{\R^3}(U)
 :=-\left\langle
   \PP\operatorname{div}(U\otimes U),
   (-\Delta)^{1/2}U
 \right\rangle_{L^2(\R^3)}.
 \label{eq:v11-Euclidean-critical-production}
\end{equation}
This is the nonlinear production in the squared \(\dot H^{1/2}\) balance.

\begin{lemma}[The V8 periodic seed has two positive moments]
\label{lem:v11-periodic-two-moments}
Let \(\overline u\) be the three-mode field of
Theorem~\ref{thm:v8-low-high-obstruction}, with
\[
 \ell=(1,0,0),\quad p=(N_0,1,0),\quad q=(N_0-1,1,0),\quad N_0\ge6,
 \qquad p-q=\ell,
\]
and \(A,B>0\).  With normalized Haar measure,
\begin{align}
 -\langle\BNS(\overline u),(-\Delta)^{1/2}\overline u\rangle
 &=\frac{AB^2}{4}\bigl(|p|-|q|\bigr)>0,
 \label{eq:v11-periodic-critical-positive}\\
 -\langle\BNS(\overline u),-\Delta\overline u\rangle
 &=\frac{AB^2}{4}\bigl(|p|^2-|q|^2\bigr)>0.
 \label{eq:v11-periodic-enstrophy-positive}
\end{align}
\end{lemma}

\begin{proof}
The direct calculation in V8 gives
\[
 \BNS(\overline u)
 =AB\mathbf e_3\sin(\ell\cdot x)
   \bigl[-\sin(p\cdot x)+\sin(q\cdot x)\bigr].
\]
At the two occupied high frequencies its cosine coefficients are
\(-AB\mathbf e_3/2\) at both \(p\) and \(q\).  The high part of
\((-\Delta)^{\alpha/2}\overline u\) is
\[
 B\mathbf e_3\bigl(
 |p|^\alpha\cos(p\cdot x)-|q|^\alpha\cos(q\cdot x)
 \bigr).
\]
Orthogonality and \(\int\cos^2(k\cdot x)\,\dd x=1/2\) give
\[
 -\langle\BNS(\overline u),
          (-\Delta)^{\alpha/2}\overline u\rangle
 =\frac{AB^2}{4}\bigl(|p|^\alpha-|q|^\alpha\bigr).
\]
Take \(\alpha=1,2\).  Since \(|p|>|q|\), both values are positive.
\end{proof}

\begin{lemma}[Slow-cutoff localization preserves both signs]
\label{lem:v11-slow-localization}
There is a real, mean-zero field
\[
 U\in C_c^\infty(\R^3;\R^3),\qquad \operatorname{div}U=0,
\]
for which
\begin{equation}
 \boxed{
 \mathscr C_{\R^3}(U)>0,\qquad
 \mathscr G_{\R^3}(U)>0.}
 \label{eq:v11-compact-two-positive}
\end{equation}
\end{lemma}

\begin{proof}
Let \(\overline{\mathcal A}\) be a smooth periodic vector potential for
\(\overline u\).  Choose a nonnegative
\(\chi\in C_c^\infty(\R^3)\), not identically zero, and set
\[
 \chi_m(x)=\chi(x/m),\qquad
 U_m=\operatorname{curl}(\chi_m\overline{\mathcal A}).
\]
Then
\begin{equation}
 U_m=\chi_m\overline u+R_m,\qquad
 \|R_m\|_{H^j(\R^3)}=O(m^{1/2})
 \quad\text{for every fixed }j,
 \label{eq:v11-localization-remainder}
\end{equation}
whereas \(\|\chi_m\overline u\|_{H^j}=O(m^{3/2})\).

We record the only nonlocal estimate needed.  Since
\(\overline u\) has finitely many nonzero Fourier modes, Plancherel and
the Taylor expansion of \(|k+\eta/m|\) at each such \(k\) give
\begin{equation}
 \left\|
 (-\Delta)^{1/2}U_m
 -\chi_m(-\Delta_{\Torus})^{1/2}\overline u
 \right\|_2
 =O(m^{1/2}).
 \label{eq:v11-fractional-localization}
\end{equation}
Indeed, after writing
\(\chi_m e^{\mathrm i k\cdot x}\) in Fourier variables, the multiplier
difference is \(O(|\eta|/m)\) on
\(|\eta|\le m|k|/2\); the complementary integral is rapidly decreasing
because \(\widehat\chi\) is Schwartz.  Summing the finite mode set and
using \eqref{eq:v11-localization-remainder} proves the estimate.

Likewise,
\[
 (U_m\cdot\nabla)U_m
 =\chi_m^2(\overline u\cdot\nabla)\overline u+F_m,
 \qquad \|F_m\|_2=O(m^{1/2}).
\]
Because \((-\Delta)^{1/2}U_m\) is divergence-free, the Leray projection
may be dropped in its pairing.  Equations
\eqref{eq:v11-fractional-localization} and the last display yield
\begin{equation}
 \mathscr C_{\R^3}(U_m)
 =c_\chi m^3
 \left[-\langle\BNS(\overline u),
       (-\Delta_{\Torus})^{1/2}\overline u\rangle_{\Torus}\right]
 +O(m^2),
 \label{eq:v11-critical-localization-asymptotic}
\end{equation}
where \(c_\chi>0\) is the harmless averaging constant proportional to
\(\int\chi^3\).  The local cubic calculation, with one derivative falling
on \(\chi_m\) in every error term, similarly gives
\begin{equation}
 \mathscr G_{\R^3}(U_m)
 =c_\chi m^3\left[-\langle\BNS(\overline u),
       -\Delta_{\Torus}\overline u\rangle_{\Torus}\right]+O(m^2).
 \label{eq:v11-enstrophy-localization-asymptotic}
\end{equation}
The zero Fourier mode supplies the displayed periodic average, while all
nonzero modes are \(o(m^3)\) by oscillatory averaging, equivalently by the
Schwartz decay of \(\widehat{\chi^3}\).
Lemma~\ref{lem:v11-periodic-two-moments} makes both leading coefficients
positive.  Fix a sufficiently large \(m\) and put \(U=U_m\).
A compactly supported curl has zero integral, completing the proof.
\end{proof}

\begin{remark}[Refinement of the packet seed]
The seed in V9 was free after its existence lemma.  We henceforth choose it
to be the field in Lemma~\ref{lem:v11-slow-localization}.  Every V9--V10
argument used only smooth compact support, zero mean, and
\(\mathscr G_{\R^3}(U)>0\), so all earlier conclusions remain valid.
Put
\[
 C_U:=c_{\Torus}\mathscr C_{\R^3}(U)>0,\qquad
 D_U:=c_{\Torus}\|(-\Delta)^{3/4}U\|_2^2.
\]
\end{remark}

% =============================================================================
\section{The packet crosses a genuine critical-record interval}
\label{sec:v11-record-growth}
% =============================================================================

For \(V\) on the expanding torus define
\begin{align}
 \mathscr C_N(V)
 &:=-c_{\Torus}
 \langle\PP_N\operatorname{div}(V\otimes V),A_N^{1/2}V\rangle,
 \label{eq:v11-rescaled-critical-production}\\
 \mathscr D_N(V)
 &:=c_{\Torus}\|A_N^{3/4}V\|_2^2.
 \label{eq:v11-rescaled-critical-dissipation}
\end{align}

\begin{theorem}[Simultaneous enstrophy and critical-record turnover]
\label{thm:v11-critical-record-turnover}
After decreasing the fixed \(s_*\) from V10, there are
\(s_\#>0\), \(\gamma>0\), and \(N_\#\) such that, for every dyadic
\(N\ge N_\#\), the actual packet solution on
\[
 J_N^\#=\left[0,\frac{s_\#}{\beta N^{5/2}}\right]
\]
satisfies
\begin{equation}
 \frac{\dd}{\dd t}\|A^{1/4}u^{(N)}(t)\|_2^2>0
 \qquad(t\in J_N^\#)
 \label{eq:v11-record-monotone}
\end{equation}
and
\begin{equation}
 \boxed{
 \|A^{1/4}u^{(N)}(\sup J_N^\#)\|_2^2
 \ge(1+\gamma)\|A^{1/4}u_N\|_2^2.}
 \label{eq:v11-record-fractional-growth}
\end{equation}
On the same interval the V10 positive viscous enstrophy excess and its
\(cN^2\) enstrophy gain remain valid.
\end{theorem}

\begin{proof}
Let
\(X_N(s)=\|A^{1/4}u^{(N)}(s/(\beta N^{5/2}))\|_2^2\).
The exact \(H^{1/2}\) balance and the scaling of V10 give
\begin{equation}
 \frac12\frac{\dd X_N}{\dd s}
 =\beta^2N\left[
   \mathscr C_N(V_N(s))-\varepsilon_N\mathscr D_N(V_N(s))
 \right],
 \qquad \varepsilon_N=\frac{\nu}{\beta\sqrt N}.
 \label{eq:v11-rescaled-record-balance}
\end{equation}
The slow reservoir contributes \(o(1)\) to the rescaled cubic derivative
pairing, so
\[
 \mathscr C_N(V_N(0))\longrightarrow C_U.
\]
The uniform \(H^m\) control of
Lemma~\ref{lem:v10-uniform-local-control}, with \(m\ge5\), makes
\(\mathscr C_N\) uniformly locally Lipschitz along this family and bounds
\(\mathscr D_N\).  Choose \(s_\#>0\), independent of \(N\), so that
\(\mathscr C_N(V_N(s))\ge C_U/2\) on \(0\le s\le s_\#\), then take \(N\)
large enough that
\(\varepsilon_N\mathscr D_N(V_N(s))\le C_U/4\).
Equation~\eqref{eq:v11-rescaled-record-balance} proves monotonicity and
\[
 X_N(s_\#)-X_N(0)\ge\frac12\beta^2C_Us_\# N.
\]
By \eqref{eq:v9-half-derivative-scaling},
\(X_N(0)=\beta^2H_UN+O(1)\).  A fixed
\(\gamma<C_Us_\#/(2H_U)\), reduced once to absorb the \(O(1)\), gives
\eqref{eq:v11-record-fractional-growth}.  Shrinking the V10 interval
preserves its positive enstrophy-excess proof.
\end{proof}

\begin{assessment}[Enstrophy turnover permits but does not force record growth]
Theorem~\ref{thm:v11-critical-record-turnover} shows that a collective
positive-excess packet can cross a fixed fractional critical-record level
during one turnover.  Conversely,
Corollary~\ref{cor:v5-decreasing-record-positive-H1} gives actual smooth
data whose enstrophy initially increases while the critical record
decreases.  Thus positive enstrophy production neither determines nor is
determined by the sign of the critical-record derivative.  A renewal proof
must retain both radial flux moments; one cannot be substituted for the
other.
\end{assessment}

% =============================================================================
\section{The complete five-piece primal source atlas}
\label{sec:v11-source-atlas}
% =============================================================================

Fix \(K>0\), put
\[
 v=P_Ku,\qquad
 \Pi_K=P_{2K}-P_K=\mathbf 1_{(K,2K]}(\Lambda),
\]
and define
\begin{align}
 d_K^{\rm int}
 &:=A^{-1/4}P_K\BNS(v),
 \label{eq:v11-source-int}\\
 d_K^{\rm in}
 &:=A^{-1/4}P_K\bigl[\BNS(u)-\BNS(v)\bigr],
 \label{eq:v11-source-in}\\
 d_K^{\rm raw}
 &:=A^{-1/4}\Pi_K\BNS(v),
 \label{eq:v11-source-raw}\\
 d_K^{\rm ann}
 &:=A^{-1/4}\Pi_K\bigl[\BNS(u)-\BNS(v)\bigr],
 \label{eq:v11-source-ann}\\
 d_K^{\rm far}
 &:=A^{-1/4}Q_{2K}\BNS(u).
 \label{eq:v11-source-far}
\end{align}

\begin{theorem}[Exact all-output source identity]
\label{thm:v11-five-piece-atlas}
For every smooth mean-zero divergence-free \(u\),
\begin{equation}
 \boxed{
 A^{-1/4}\BNS(u)
 =d_K^{\rm int}+d_K^{\rm in}
  +d_K^{\rm raw}+d_K^{\rm ann}+d_K^{\rm far}.}
 \label{eq:v11-five-piece-source}
\end{equation}
Moreover the three output bands are orthogonal, so
\begin{equation}
 \boxed{
 \|A^{-1/4}\BNS(u)\|_2^2
 =
 \|d_K^{\rm int}+d_K^{\rm in}\|_2^2
 +\|d_K^{\rm raw}+d_K^{\rm ann}\|_2^2
 +\|d_K^{\rm far}\|_2^2.}
 \label{eq:v11-output-Pythagoras}
\end{equation}
The V34 raw daughter is exactly \(d_K^{\rm raw}\), and its incoming-head
source is exactly \(d_K^{\rm in}\).  The annular outer-parent term
\(d_K^{\rm ann}\) and the fully assembled far-output term
\(d_K^{\rm far}\) are separate primal source incidences.
\end{theorem}

\begin{proof}
The additive Fourier support of two \(P_K\) parents gives
\(Q_{2K}\BNS(v)=0\).  Insert
\[
 I=P_K+\Pi_K+Q_{2K}
\]
on the output of \(\BNS(u)\), and on the first two bands insert
\(\BNS(u)=\BNS(v)+[\BNS(u)-\BNS(v)]\).
This is \eqref{eq:v11-five-piece-source}.  The multiplier \(A^{-1/4}\)
commutes with all spectral projections.  The head, first exterior
annulus, and far exterior are mutually orthogonal, which proves
\eqref{eq:v11-output-Pythagoras}.  The two V34 identifications follow
from their displayed definitions and the same support fact.
\end{proof}

\begin{firewall}[Assembly occurs within, not across, output bands]
The sums \(d_K^{\rm int}+d_K^{\rm in}\) and
\(d_K^{\rm raw}+d_K^{\rm ann}\) must remain assembled: their summands can
cancel.  The far term is different.  It is already the complete source
after every parent interaction has been summed, and its output support is
orthogonal to both V34 source rows.  A lower bound for
\(d_K^{\rm far}\) therefore obeys the inherited assembly-first rule.
\end{firewall}

\begin{proposition}[The omitted far source is algebraically independent]
\label{prop:v11-far-source-independent}
There are smooth finite Fourier data and a fixed cutoff \(K\) for which
\begin{equation}
 d_K^{\rm raw}=0,\qquad d_K^{\rm in}=0,\qquad
 d_K^{\rm far}\ne0.
 \label{eq:v11-two-rows-zero-far-nonzero}
\end{equation}
Consequently no universal bound of \(d_K^{\rm far}\) by the two V34 source
rows is possible.
\end{proposition}

\begin{proof}
Use the active V5 homochiral triad with modes
\[
 p=N(1,0,0),\quad q=N(0,1,1),\quad r=-N(1,1,1)
\]
and choose its phase so that its nonlinear enstrophy production is
positive.  Take \(0<K<N/2\).  Then \(P_Ku=0\), hence
\(d_K^{\rm raw}=0\).  Every nonzero output frequency of the quadratic
source is a sum of two elements of
\(\{\pm p,\pm q,\pm r\}\).  Direct enumeration shows that its modulus is
at least \(N\); the only zero sum is \(k+(-k)\), whose divergence-form
Fourier coefficient at output zero vanishes.  Thus
\(P_K\BNS(u)=0\) and \(d_K^{\rm in}=0\), while
\(Q_{2K}\BNS(u)=\BNS(u)\).
The latter is nonzero because its pairing with \(Au\) is strictly
negative by the chosen positive-production phase.  This proves
\eqref{eq:v11-two-rows-zero-far-nonzero}.
\end{proof}

% =============================================================================
\section{Persistence of the fully assembled far-output source}
\label{sec:v11-assembled-far}
% =============================================================================

Let \(\mathsf Q_\rho\) denote the Euclidean sharp high-pass used in V9--V10
and put
\begin{equation}
 \mathscr G_\rho^{\rm far}(U)
 :=-c_{\Torus}\left\langle
   \mathsf Q_{2\rho}\BNS(U),
   -\Delta\mathsf Q_\rho U
 \right\rangle_{L^2(\R^3)}.
 \label{eq:v11-Euclidean-far-production}
\end{equation}
Strong convergence of the multipliers on the smooth seed gives
\begin{equation}
 \mathscr G_\rho^{\rm far}(U)\longrightarrow G_U
 \qquad(\rho\downarrow0).
 \label{eq:v11-far-production-limit}
\end{equation}
Decrease the free energy fraction \(\delta\), if necessary, so that for
\(\rho=\delta\kappa_U\),
\begin{equation}
 \mathscr G_\rho^{\rm far}(U)>\frac12G_U.
 \label{eq:v11-positive-assembled-far-seed}
\end{equation}
This choice is compatible with every earlier smallness requirement on
\(\delta\).

\begin{theorem}[Assembled far-output overdrive on the actual interval]
\label{thm:v11-assembled-far-persistence}
After one fixed reduction of \(s_\#\), set
\[
 t_N=\frac{s_\#}{\beta N^{5/2}},\qquad
 f_N(t):=d_{K_N}^{\rm far}(u^{(N)}(t))
 =A^{-1/4}Q_{2K_N}\BNS(u^{(N)}(t)),
 \qquad
 w_N=Q_{K_N}u^{(N)}.
\]
For all sufficiently large dyadic \(N\) and \(0\le t\le t_N\),
\begin{equation}
 \boxed{
 -\langle Q_{2K_N}\BNS(u^{(N)}(t)),Aw_N(t)\rangle
 \ge c\beta^3G_UN^{9/2}.}
 \label{eq:v11-full-far-production}
\end{equation}
Furthermore,
\begin{equation}
 \boxed{
 \mathscr A_N^{\rm far}
 :=\int_0^{t_N}\|f_N(t)\|_2^2\,\dd t
 \asymp\beta^3N^{3/2}.}
 \label{eq:v11-full-far-action}
\end{equation}
These statements use the fully assembled \(\BNS(u^{(N)})\), not an opened
outer--outer parent component.
\end{theorem}

\begin{proof}
Since \(K_N/N\to\rho\), the Riemann-sum argument used in
Lemma~\ref{lem:v9-half-derivative-scaling} gives at \(s=0\)
\begin{align*}
 \beta^{-3}N^{-9/2}
 \bigl[-\langle Q_{2K_N}\BNS(u_N),AQ_{K_N}u_N\rangle\bigr]
 &\longrightarrow\mathscr G_\rho^{\rm far}(U),\\
 \beta^{-4}N^{-4}\|A^{-1/4}Q_{2K_N}\BNS(u_N)\|_2^2
 &\longrightarrow
 \|(-\Delta)^{-1/4}\mathsf Q_{2\rho}\BNS(U)\|_2^2.
\end{align*}
The second limit is positive, because a zero limiting source would make
\eqref{eq:v11-positive-assembled-far-seed} zero.  Terms containing the
fixed low reservoir vanish after these normalizations.

Sharp spectral projections are contractions on every Sobolev space.
The bilinear product estimate and the uniform large-torus \(H^m\) bound
therefore make both normalized functionals uniformly continuous in
\(s\).  Reduce \(s_\#\) so that half the positive margins persist.
This proves \eqref{eq:v11-full-far-production} and gives instantaneous
upper and lower bounds \(c\beta^4N^4\le\|f_N(t)\|_2^2
\le C\beta^4N^4\).  Integration over
\(t_N=s_\#/(\beta N^{5/2})\) proves
\eqref{eq:v11-full-far-action}.
\end{proof}

\begin{theorem}[The far source reaches its own Duhamel endpoint]
\label{thm:v11-far-Duhamel-endpoint}
Define the heat-propagated far response
\begin{equation}
 \mathfrak r_N^{\rm far}
 :=\int_0^{t_N}
 e^{-\nu A(t_N-t)}f_N(t)\,\dd t.
 \label{eq:v11-far-response}
\end{equation}
Then the exact endpoint identity is
\begin{equation}
 \boxed{
 A^{-1/4}Q_{2K_N}u^{(N)}(t_N)
 =
 e^{-\nu At_N}A^{-1/4}Q_{2K_N}u_N
 -\mathfrak r_N^{\rm far}.}
 \label{eq:v11-far-Duhamel-identity}
\end{equation}
After reducing \(s_\#\) once more,
\begin{equation}
 \boxed{
 \|\mathfrak r_N^{\rm far}\|_2\asymp\beta N^{-1/2},
 \qquad
 \|\mathfrak r_N^{\rm far}\|_2^2
 \asymp t_N\mathscr A_N^{\rm far}.}
 \label{eq:v11-far-response-scale}
\end{equation}
Thus the source action, occurrence interval, and endpoint are carried on
one literal card; the elementary \(L_t^2\)-to-Duhamel bound is saturated
at the packet scale.
\end{theorem}

\begin{proof}
Apply \(A^{-1/4}Q_{2K_N}\) to Navier--Stokes and integrate the resulting
linear heat equation.  This gives
\eqref{eq:v11-far-Duhamel-identity} with no remainder.

Under the variables of V10, the norm of \(f_N\) has scale
\(\beta^2N^2\), while \(\dd t=(\beta N^{5/2})^{-1}\dd s\).
Consequently
\[
 \mathfrak r_N^{\rm far}
 \quad\hbox{has norm scale}\quad
 \beta N^{-1/2}.
\]
For the lower bound, the rescaled integrand at \(s=0\) converges in
\(L^2\) to the nonzero vector
\[
 (-\Delta)^{-1/4}\mathsf Q_{2\rho}\BNS(U).
\]
Uniform Sobolev continuity makes the integrand remain within one quarter
of that vector on a fixed sufficiently short \(s\)-interval.  In the same
variables the heat multiplier is
\(e^{-\varepsilon_NA_N(s_\#-s)}\), with
\(\varepsilon_N\to0\); it converges strongly to the identity on this fixed
smooth limiting vector.  Hence the rescaled vector integral has a fixed
positive lower norm.  The matching upper bound follows from contraction
of the heat semigroup.  Finally,
\[
 \|\mathfrak r_N^{\rm far}\|_2^2
 \le t_N\int_0^{t_N}\|f_N(t)\|_2^2\,\dd t
\]
by Cauchy--Schwarz and heat contraction.  The already proved scales show
that both sides are comparable to \(\beta^2N^{-1}\).
\end{proof}

% =============================================================================
\section{A universal endpoint-or-covariant-variation alternative}
\label{sec:v11-far-variation}
% =============================================================================

The preceding packet lies on the coherent endpoint branch.  For an
arbitrary terminal card, temporal cancellation must be retained explicitly.

\begin{lemma}[Hilbert-valued heat-covariant mean inequality]
\label{lem:v11-heat-mean}
Let \(A\ge0\) be self-adjoint on a Hilbert space, \(I=[a,b]\),
\(\tau=b-a\), and suppose
\(f\in H^1(I;\mathcal H)\cap L^2(I;D(A))\).  Put
\begin{align}
 h(t)&=e^{-\nu A(b-t)}f(t),
 \label{eq:v11-heat-pulled-source}\\
 R_f&=\int_a^b h(t)\,\dd t,\qquad
 S_f^\nu=\int_a^b\|h(t)\|^2\,\dd t.
 \label{eq:v11-heat-response-action}
\end{align}
Then
\begin{equation}
 \boxed{
 S_f^\nu
 \le\frac{\|R_f\|^2}{\tau}
 +\frac{\tau^2}{\pi^2}
  \int_a^b
  \left\|
   e^{-\nu A(b-t)}(\partial_t+\nu A)f(t)
  \right\|^2\,\dd t.}
 \label{eq:v11-heat-mean-inequality}
\end{equation}
In particular, either
\begin{equation}
 \|R_f\|^2\ge\frac{\tau}{2}S_f^\nu,
 \label{eq:v11-endpoint-branch}
\end{equation}
or
\begin{equation}
 \int_a^b
 \left\|e^{-\nu A(b-t)}(\partial_t+\nu A)f(t)\right\|^2\,\dd t
 \ge\frac{\pi^2}{2\tau^2}S_f^\nu.
 \label{eq:v11-covariant-variation-branch}
\end{equation}
\end{lemma}

\begin{proof}
The Bochner mean of \(h\) is \(\overline h=R_f/\tau\), and orthogonality
to constants gives the exact variance identity
\[
 \int_a^b\|h\|^2\,\dd t
 =\tau\|\overline h\|^2
  +\int_a^b\|h-\overline h\|^2\,\dd t.
\]
The Hilbert-valued Wirtinger inequality bounds the last term by
\(\tau^2\int\|\partial_th\|^2/\pi^2\).  Functional calculus gives
\[
 \partial_th(t)
 =e^{-\nu A(b-t)}(\partial_t+\nu A)f(t).
\]
This proves \eqref{eq:v11-heat-mean-inequality}.  If
\eqref{eq:v11-endpoint-branch} fails, substitute its strict reverse into
that inequality and rearrange to obtain
\eqref{eq:v11-covariant-variation-branch}.
\end{proof}

Define the symmetrized bilinear source
\[
 \BNS_{\rm s}(a,b)
 :=\PP\bigl[(a\cdot\nabla)b+(b\cdot\nabla)a\bigr].
\]

\begin{proposition}[Exact covariant derivative of the far source]
\label{prop:v11-far-covariant-derivative}
For a smooth Navier--Stokes solution and
\(f_K=A^{-1/4}Q_{2K}\BNS(u)\),
\begin{equation}
 \boxed{
 (\partial_t+\nu A)f_K
 =
 A^{-1/4}Q_{2K}
 \left\{
  \nu\bigl[A\BNS(u)-\BNS_{\rm s}(Au,u)\bigr]
  -\BNS_{\rm s}(\BNS(u),u)
 \right\}.}
 \label{eq:v11-far-covariant-derivative}
\end{equation}
Hence every far-output card has an exact alternative: its heat-weighted
source reaches the far endpoint as in
\eqref{eq:v11-endpoint-branch}, or the fully assembled viscous-commutator
plus cubic source-motion quantity in
\eqref{eq:v11-far-covariant-derivative} obeys
\eqref{eq:v11-covariant-variation-branch}.
\end{proposition}

\begin{proof}
Differentiate \(\BNS(u)=\BNS(u,u)\) and use
\[
 \partial_tu=-\nu Au-\BNS(u).
\]
Thus
\[
 \partial_t\BNS(u)
 =-\nu\BNS_{\rm s}(Au,u)
  -\BNS_{\rm s}(\BNS(u),u).
\]
Add \(\nu A\BNS(u)\) and commute the fixed spectral multipliers with
\(\partial_t+\nu A\).  This proves
\eqref{eq:v11-far-covariant-derivative}.  Apply
Lemma~\ref{lem:v11-heat-mean}.
\end{proof}

\begin{assessment}[Heat-weighted action is essential on an all-shell row]
For a finite annulus with \(\nu\Lambda^2\tau\ll1\), the heat-weighted and
unweighted actions are comparable.  On the unrestricted far output they
need not be: arbitrarily remote modes may be erased before the endpoint.
Lemma~\ref{lem:v11-heat-mean} therefore does not pretend that the full
unweighted critical action reaches the endpoint.  A general terminal proof
must either acquire a cofinal shell on which heat is controlled or retain
heat-frequency escape as a genuine all-shell alternative.  The packet
satisfies the first condition after rescaling and lands on the endpoint
branch by Theorem~\ref{thm:v11-far-Duhamel-endpoint}.
\end{assessment}

% =============================================================================
\section{Audit of the inherited transfer, forcing, and all-shell exits}
\label{sec:v11-inherited-audit}
% =============================================================================

\begin{theorem}[No inherited finite-head row contains the primal far source]
\label{thm:v11-no-inherited-far-row}
With the definitions frozen in the V34 source, none of the following is an
incidence map for \(d_K^{\rm far}\):
\begin{enumerate}[label=\textup{(\roman*)},leftmargin=3.2em]
\item the V34 raw daughter or incoming-head source;
\item the V29 fixed raw-carrier secant/transfer/forcing ledger;
\item the V27 assembled dynamic transfer or assembled physical forcing;
\item the V25 high--high backscatter forcing; or
\item the stated fused all-shell raw-source/history escape.
\end{enumerate}
The first four statements are operator-type exclusions.  The last statement
means that the inherited all-shell escape is only a conditional terminal
label: its antecedent does not follow from nonzero far-output source action.
The minimum missing primal card is
\begin{equation}
 \boxed{
 \left(
 I,\ d_K^{\rm far},\
 \int_I\|d_K^{\rm far}\|_2^2,\
 \int_Ie^{-\nu A(b-t)}d_K^{\rm far}(t)\,\dd t,\
 (\partial_t+\nu A)d_K^{\rm far}
 \right).}
 \label{eq:v11-minimum-far-card}
\end{equation}
\end{theorem}

\begin{proof}
The first exclusion is
Theorem~\ref{thm:v11-five-piece-atlas}: the two V34 terms are
\(d_K^{\rm raw}\) and \(d_K^{\rm in}\), while
\(d_K^{\rm far}\in\operatorname{Ran}Q_{2K}\).

The V29 carrier is generated from \(d_K^{\rm raw}\) and is contained in
the first exterior octave.  Its writer \(\chi_{\rm raw}\), negative secant,
transfer, and forcing works all belong to the associated
daughter-octave adjoint equation.  The V27 quantities
\[
 G_{\rm h}^{\rm tr}=-P_{\rm h}\mathscr L_1^*\eta_{\rm h},
 \qquad F_{\rm h}=P_{\rm h}F_1
\]
also lie in that heat-reduced daughter carrier.  They are adjoint currents,
not components of \(A^{-1/4}\BNS(u)\).

Likewise the V25 backscatter term is
\[
 F^{\rm bs}=-\Pi_1\mathscr L_w^*\Pi_{>2}\Psi_\phi
 \in\operatorname{Ran}\Pi_1.
\]
It transports a far adjoint tail back into the daughter octave.  Its output
space and direction of mapping are opposite to the primal far source
\(A^{-1/4}Q_{2K}\BNS(u)\).  No identity in V25--V34 identifies their
actions or endpoints.

Finally, the inherited all-shell alternative is triggered when no
predeclared finite head captures the raw daughter, its selected scalar, and
the required history semantics.  A nonzero \(d_K^{\rm far}\) at an already
fixed finite \(K\) does not imply that antecedent.
Proposition~\ref{prop:v11-far-source-independent} makes the logical
independence explicit: both named V34 source rows can vanish while the far
source remains nonzero.  The entries in
\eqref{eq:v11-minimum-far-card} are exactly the universal source, action,
endpoint, and cancellation quantities proved in
Theorem~\ref{thm:v11-far-Duhamel-endpoint},
Lemma~\ref{lem:v11-heat-mean}, and
Proposition~\ref{prop:v11-far-covariant-derivative}.
\end{proof}

\begin{firewall}[A gap is not a contradiction]
Theorem~\ref{thm:v11-no-inherited-far-row} does not contradict the V34
terminal alternative.  That alternative openly retains a fused all-shell
escape.  The theorem proves that the escape has not yet been equipped with
the source/action/endpoint map needed to control the full critical source.
Adding the card \eqref{eq:v11-minimum-far-card} is therefore an upstream
completion of the programme, not a claim that a terminal singular solution
contains the packet profile.
\end{firewall}

% =============================================================================
\section{Integrated V11 conclusion and next acquisition order}
\label{sec:v11-conclusion}
% =============================================================================

\begin{theorem}[What V11 proves]
\label{thm:v11-conclusion}
Relative to V1--V10 and the frozen manuscripts, the following are proved.
\begin{enumerate}[label=\textup{(V11.\arabic*)},leftmargin=5.7em]
\item The compact packet seed may be chosen with positive nonlinear
      \(\dot H^{1/2}\) and \(H^1\) production simultaneously.
\item The resulting actual solution increases its squared critical record
      by a fixed fraction during the same \(cN^{-5/2}\) positive-enstrophy
      turnover interval.
\item Positive enstrophy production does not force critical-record growth;
      the V5 actual triad supplies the opposite sign combination.
\item The full nonlinear source has the exact five-piece atlas
      \eqref{eq:v11-five-piece-source}; V34 names only two of its four
      non-internal incidence terms.
\item The fully assembled far-output source has critical action
      \(\asymp N^{3/2}\) on the packet interval and reaches a nonzero
      \(A^{-1/4}\)-velocity endpoint of size \(\asymp N^{-1/2}\).
\item On every smooth far-output card, small endpoint response forces the
      exact heat-covariant source variation
      \eqref{eq:v11-far-covariant-derivative}.
\item The inherited secant, transfer, forcing, backscatter, and qualitative
      all-shell rows do not yet provide the primal far-source incidence card
      \eqref{eq:v11-minimum-far-card}.
\end{enumerate}
No full Navier--Stokes stability or global-regularity theorem is obtained.
\end{theorem}

\begin{proof}
Items \textup{(V11.1)--(V11.3)} are
Lemmas~\ref{lem:v11-periodic-two-moments}--\ref{lem:v11-slow-localization},
Theorem~\ref{thm:v11-critical-record-turnover}, and
Corollary~\ref{cor:v5-decreasing-record-positive-H1}.  Item \textup{(V11.4)} is
Theorem~\ref{thm:v11-five-piece-atlas} and
Proposition~\ref{prop:v11-far-source-independent}.  Item
\textup{(V11.5)} is
Theorems~\ref{thm:v11-assembled-far-persistence}--%
\ref{thm:v11-far-Duhamel-endpoint}.  Item \textup{(V11.6)} is
Lemma~\ref{lem:v11-heat-mean} and
Proposition~\ref{prop:v11-far-covariant-derivative}.  The final item is
Theorem~\ref{thm:v11-no-inherited-far-row}.
\end{proof}

\begin{programme}[V12 acquisition order]
\label{prog:v12}
\begin{enumerate}[label=\textup{(Z\arabic*)},leftmargin=4em]
\item Use the five-piece identity on an arbitrary hypothetical terminal
      record sequence.  Do not import the packet profile into that sequence.
\item Resolve \(d_K^{\rm ann}\) only as the complete annular assembly
      \(d_K^{\rm raw}+d_K^{\rm ann}\); do not take positive parts of its
      parent components before cancellation.
\item Decompose \(d_K^{\rm far}\) into cofinal output shells and compare its
      unweighted action with the heat-weighted action.  Either select a
      shell with controlled \(\nu\Lambda^2|I|\), or retain a precise
      heat-frequency all-shell escape.
\item On a selected shell apply Lemma~\ref{lem:v11-heat-mean}.  Carry the
      endpoint branch into the original record/ancestry scalar.  On the
      cancellation branch retain the complete viscous-commutator plus cubic
      source motion in \eqref{eq:v11-far-covariant-derivative}.
\item Seek a summable debit for the cofinal far responses across doubled
      records.  The packet arithmetic shows that the native endpoint size
      alone is compatible with summability, so a contradiction must use
      ancestry, repeated response, or nonsummable record gain.
\item Keep material localization behind the V8 guard-band commutator and
      the represented-history branch at the V3 physical-incidence stop.
\end{enumerate}
The remaining global bottleneck is now a quantitative cofinal-shell
incidence/ancestry theorem for the complete far source, followed by an upper
budget for its covariant motion or repeated endpoint response.
\end{programme}

\begin{assessment}[Position after V11]
The packet obstruction now survives the assembly-first audit and reaches an
actual endpoint.  It also changes the critical record, so it is not hidden
on a tangent record surface.  Nevertheless its far response has the
summable size \(N^{-1/2}\), and no argument yet forces repeated copies on a
single finite-energy trajectory.  The universal gain is the five-piece
source atlas and the endpoint-or-covariant-variation alternative.  The next
honest step is cofinal shell selection on an arbitrary terminal sequence,
not another estimate of the packet.
\end{assessment}

% =============================================================================
\section*{Version V11 history}
\addcontentsline{toc}{section}{Version V11 history}
% =============================================================================

Refined the compact packet seed by a slow vector-potential cutoff and proved
simultaneous positive nonlinear \(\dot H^{1/2}\) and \(H^1\) production.
Used the expanding-torus dynamics to obtain a fixed-fraction critical-record
increase on the actual turnover interval.  Replaced V10's opened
outer--outer component by the fully assembled far-output source, proved its
\(N^{3/2}\) action, and carried it through the heat equation to its literal
endpoint.  Established the exact five-piece primal source atlas, a
finite-mode independence witness, and a universal heat-covariant
endpoint-or-variation inequality.  Audited the V25--V34 transfer, forcing,
backscatter, and all-shell rows and isolated the missing quantitative primal
far-source card.  No global regularity claim was made.
% Future versions are appended immediately above the sole terminator.

% =============================================================================
% VERSION V12 APPEND
% =============================================================================

\section{V12: heat-edge resolution of the complete exterior source}
\label{sec:v12-entry}
% =============================================================================

V11 supplied the missing primal source card but left two tasks: the
first-exterior annulus had to remain assembled, and the unrestricted far
source had to be separated into frequencies which can and cannot reach the
card endpoint through heat propagation.  Both tasks have exact spectral
solutions.

On an interval \(I=[a,b]\), the intrinsic heat edge is
\((\nu|I|)^{-1/2}\).  Instead of selecting one dyadic shell and losing a
logarithm, we retain the complete output band between \(K\) and a cutoff
\(L\) at that heat edge.  The V11 Hilbert-valued mean inequality then gives
an exhaustive alternative: the exterior critical action escapes above
\(L\), reaches a literal \(A^{-1/4}\)-velocity endpoint, or forces
heat-covariant variation of the complete band source.  The endpoint branch
has the sharp energy capacity \(E_0/K\).

The super-parabolic branch is compared with the globally paid
\(A^{-3/4}\BNS(u)\) source.  If its critical-to-weak spectral leverage stays
bounded, its time-weighted action is summable on disjoint cards.  Otherwise
the exact surviving packet is unbounded source-frequency leverage.  This is
the quantitative version of the qualitative all-shell heat escape left in
V11.

% =============================================================================
\section{The first exterior octave must include outer parents}
\label{sec:v12-annular-assembly}
% =============================================================================

Recall from Theorem~\ref{thm:v11-five-piece-atlas} that
\begin{equation}
 A^{-1/4}\Pi_K\BNS(u)
 =d_K^{\rm raw}+d_K^{\rm ann}.
 \label{eq:v12-complete-annular-source}
\end{equation}
No orthogonality separates the two terms on the right.

\begin{proposition}[Outer parents can populate the daughter octave while the
raw and incoming rows vanish]
\label{prop:v12-annular-independence}
For every positive integer \(N\), put \(K=N\) and
\[
 p=(3N,0,0),\qquad q=(2N,-N,0),\qquad k=p-q=(N,N,0).
\]
The real divergence-free trigonometric polynomial
\begin{equation}
 u(x)=\mathbf e_2\cos(p\cdot x)+\mathbf e_3\cos(q\cdot x)
 \label{eq:v12-annular-witness}
\end{equation}
satisfies
\begin{equation}
 \boxed{
 d_K^{\rm raw}=0,\qquad d_K^{\rm in}=0,\qquad
 d_K^{\rm ann}\ne0.}
 \label{eq:v12-annular-two-rows-zero}
\end{equation}
Thus even inside \(K<\Lambda\le2K\), the actual assembled source cannot be
replaced by the V34 raw daughter.
\end{proposition}

\begin{proof}
Both occupied frequencies have modulus greater than \(K\), so \(P_Ku=0\)
and \(d_K^{\rm raw}=0\).  The only nonzero advective interaction is
\[
 (u\cdot\nabla)u
 =N\mathbf e_3\cos(p\cdot x)\sin(q\cdot x)
 =\frac N2\mathbf e_3
   \bigl[\sin((p+q)\cdot x)-\sin(k\cdot x)\bigr].
\]
Both output wavevectors lie in the \(x_1x_2\)-plane, so the Leray
projection fixes their \(\mathbf e_3\) polarization.  Their moduli are
\[
 |k|=\sqrt2N\in(K,2K],
 \qquad |p+q|=\sqrt{26}N>2K.
\]
There is no head output, hence \(d_K^{\rm in}=0\), while the \(k\)-term
makes \(d_K^{\rm ann}\ne0\).
\end{proof}

\begin{firewall}[The V34 residual is not the complete annular source]
The V34 history-fresh residual is constructed from
\(d_K^{\rm raw}\) and a represented follower.  Proposition
\ref{prop:v12-annular-independence} shows that it need not contain the
physical source already entering the same output octave from exterior
parents.  Consequently no action or variation estimate for that residual
may be applied to \eqref{eq:v12-complete-annular-source} until an explicit
source-incidence identity includes \(d_K^{\rm ann}\).
\end{firewall}

% =============================================================================
\section{Parabolic and super-parabolic exterior outputs}
\label{sec:v12-heat-edge}
% =============================================================================

Let \(u\) be smooth on \(I=[a,b]\), put \(\tau=b-a>0\), and assume
\begin{equation}
 \sup_{t\in I}\|u(t)\|_2^2\le E_0,\qquad
 \nu K^2\tau\le C_0.
 \label{eq:v12-card-assumptions}
\end{equation}
Define
\begin{equation}
 L=L(K,I):=\max\left\{2K,\frac1{\sqrt{\nu\tau}}\right\},
 \qquad
 M_0:=\max\{4C_0,1\}.
 \label{eq:v12-heat-edge}
\end{equation}
Then
\begin{equation}
 \nu L^2\tau\le M_0.
 \label{eq:v12-heat-edge-bound}
\end{equation}
Let
\[
 \Pi_{K,L}:=\mathbf 1_{(K,L]}(\Lambda)
\]
and form the fully assembled exterior sources
\begin{align}
 d_{K,I}^{\rm ext}
 &:=A^{-1/4}Q_K\BNS(u),
 \label{eq:v12-exterior-source}\\
 d_{K,I}^{\rm par}
 &:=A^{-1/4}\Pi_{K,L}\BNS(u),
 \label{eq:v12-parabolic-source}\\
 d_{K,I}^{\rm sup}
 &:=A^{-1/4}Q_L\BNS(u).
 \label{eq:v12-superparabolic-source}
\end{align}
The interval in the notation only freezes \(L\).  Orthogonal output
support gives
\begin{equation}
 d_{K,I}^{\rm ext}=d_{K,I}^{\rm par}+d_{K,I}^{\rm sup},
 \qquad
 \mathscr A^{\rm ext}
 =\mathscr A^{\rm par}+\mathscr A^{\rm sup},
 \label{eq:v12-heat-edge-Pythagoras}
\end{equation}
where
\begin{equation}
 \mathscr A^\alpha
 :=\int_I\|d_{K,I}^\alpha(t)\|_2^2\,\dd t,
 \qquad
 \alpha\in\{{\rm ext,par,sup}\}.
 \label{eq:v12-three-actions}
\end{equation}
Because \(L\ge2K\), the parabolic source contains the complete assembly
\(d_K^{\rm raw}+d_K^{\rm ann}\), rather than either parent component.

Define its heat-weighted action, endpoint response, and covariant
variation by
\begin{align}
 \mathscr S^{\rm par}_\nu
 &:=
 \int_a^b
 \left\|e^{-\nu A(b-t)}d_{K,I}^{\rm par}(t)\right\|_2^2\,\dd t,
 \label{eq:v12-parabolic-heat-action}\\
 \mathfrak r_{K,I}^{\rm par}
 &:=
 \int_a^b
 e^{-\nu A(b-t)}d_{K,I}^{\rm par}(t)\,\dd t,
 \label{eq:v12-parabolic-response}\\
 \mathscr V_{K,I}^{\rm par}
 &:=
 \int_a^b
 \left\|
 e^{-\nu A(b-t)}
 (\partial_t+\nu A)d_{K,I}^{\rm par}(t)
 \right\|_2^2\,\dd t.
 \label{eq:v12-parabolic-variation}
\end{align}

\begin{theorem}[Complete exterior heat-edge trichotomy]
\label{thm:v12-heat-edge-trichotomy}
Suppose \(\mathscr A^{\rm ext}>0\).  At least one of the following holds.
\begin{enumerate}[label=\textup{(\Alph*)},leftmargin=3.3em]
\item \emph{Super-parabolic output escape:}
      \begin{equation}
       \boxed{\mathscr A^{\rm sup}\ge\frac12\mathscr A^{\rm ext}.}
       \label{eq:v12-superparabolic-branch}
      \end{equation}
\item \emph{Coherent endpoint response:}
      \begin{equation}
       \boxed{
       \|\mathfrak r_{K,I}^{\rm par}\|_2^2
       \ge
       \frac{e^{-2M_0}}4\,\tau\mathscr A^{\rm ext},}
       \label{eq:v12-endpoint-lower}
      \end{equation}
      and necessarily
      \begin{equation}
       \boxed{
       \tau\mathscr A^{\rm ext}
       \le16e^{2M_0}\frac{E_0}{K}.}
       \label{eq:v12-endpoint-capacity}
      \end{equation}
\item \emph{Heat-covariant source variation:}
      \begin{equation}
       \boxed{
       \mathscr V_{K,I}^{\rm par}
       \ge
       \frac{\pi^2e^{-2M_0}}{4\tau^2}
       \mathscr A^{\rm ext}.}
       \label{eq:v12-variation-lower}
      \end{equation}
\end{enumerate}
On the second branch the response has the exact endpoint identity
\begin{equation}
 \boxed{
 A^{-1/4}\Pi_{K,L}u(b)
 =
 e^{-\nu A\tau}A^{-1/4}\Pi_{K,L}u(a)
 -\mathfrak r_{K,I}^{\rm par}.}
 \label{eq:v12-parabolic-Duhamel}
\end{equation}
\end{theorem}

\begin{proof}
If \eqref{eq:v12-superparabolic-branch} fails, then
\(\mathscr A^{\rm par}>\mathscr A^{\rm ext}/2\).  On
\(\operatorname{Ran}\Pi_{K,L}\),
\[
 \|e^{-\nu A(b-t)}z\|_2^2
 \ge e^{-2\nu L^2\tau}\|z\|_2^2
 \ge e^{-2M_0}\|z\|_2^2.
\]
Therefore
\begin{equation}
 \mathscr S_\nu^{\rm par}
 \ge\frac{e^{-2M_0}}2\mathscr A^{\rm ext}.
 \label{eq:v12-heat-unweighted-comparison}
\end{equation}
Apply Lemma~\ref{lem:v11-heat-mean} to
\(f=d_{K,I}^{\rm par}\).  Its endpoint branch gives
\[
 \|\mathfrak r_{K,I}^{\rm par}\|_2^2
 \ge\frac{\tau}{2}\mathscr S_\nu^{\rm par},
\]
which is \eqref{eq:v12-endpoint-lower}.  Its complementary branch and
\eqref{eq:v12-heat-unweighted-comparison} give
\eqref{eq:v12-variation-lower}.

It remains to prove the capacity bound.  Applying
\(A^{-1/4}\Pi_{K,L}\) to Navier--Stokes gives
\eqref{eq:v12-parabolic-Duhamel}.  Since every output frequency exceeds
\(K\),
\[
 \|A^{-1/4}\Pi_{K,L}u(t)\|_2^2
 \le K^{-1}\|u(t)\|_2^2\le\frac{E_0}{K}.
\]
Heat contraction and the endpoint identity yield
\[
 \|\mathfrak r_{K,I}^{\rm par}\|_2^2
 \le4\frac{E_0}{K}.
\]
Combine this with \eqref{eq:v12-endpoint-lower} to obtain
\eqref{eq:v12-endpoint-capacity}.
\end{proof}

\begin{proposition}[Exact source motion on the parabolic band]
\label{prop:v12-parabolic-covariant-source}
With \(\BNS_{\rm s}\) as in V11,
\begin{equation}
 \boxed{
 (\partial_t+\nu A)d_{K,I}^{\rm par}
 =
 A^{-1/4}\Pi_{K,L}
 \left\{
  \nu\bigl[A\BNS(u)-\BNS_{\rm s}(Au,u)\bigr]
  -\BNS_{\rm s}(\BNS(u),u)
 \right\}.}
 \label{eq:v12-parabolic-covariant-source}
\end{equation}
Thus branch \textup{(C)} of
Theorem~\ref{thm:v12-heat-edge-trichotomy} is one complete
viscous-commutator plus cubic source-motion packet.  It is not a sum of
independently priced parent births.
\end{proposition}

\begin{proof}
The projector \(\Pi_{K,L}\) is fixed after \(I\), \(K\), and \(L\) are
frozen and commutes with \(A\) and \(\partial_t\).  Repeat the derivation
of \eqref{eq:v11-far-covariant-derivative} with \(\Pi_{K,L}\) in place
of \(Q_{2K}\); only these commutation properties were used.
\end{proof}
% =============================================================================
\section{The super-parabolic branch and the paid weak source}
\label{sec:v12-superparabolic}
% =============================================================================

For this section write
\[
 \BNS(a,b):=\PP((a\cdot\nabla)b),\qquad
 \BNS(u)=\BNS(u,u).
\]

\begin{lemma}[Every super-parabolic output has a high parent]
\label{lem:v12-high-parent-incidence}
Put \(v_L=P_{L/2}u\) and \(w_L=Q_{L/2}u\).  Then
\begin{equation}
 \boxed{
 Q_L\BNS(u)
 =
 Q_L\bigl[\BNS(w_L,u)+\BNS(v_L,w_L)\bigr].}
 \label{eq:v12-high-parent-identity}
\end{equation}
Thus \(d_{K,I}^{\rm sup}\) is a complete high-parent source occurrence;
it is not a remote output produced by two low parents.
\end{lemma}

\begin{proof}
Expand \(u=v_L+w_L\).  The term
\(\BNS(v_L,v_L)\) has output support at most \(L\), so its \(Q_L\)
projection vanishes.  The remaining three bilinear terms are exactly
\(\BNS(w_L,u)+\BNS(v_L,w_L)\).
\end{proof}

Define the weak super-parabolic action
\begin{equation}
 \mathscr W^{\rm sup}
 :=
 \int_I\|A^{-3/4}Q_L\BNS(u(t))\|_2^2\,\dd t
 =
 \int_I\|A^{-1/2}d_{K,I}^{\rm sup}(t)\|_2^2\,\dd t.
 \label{eq:v12-superparabolic-weak-action}
\end{equation}
When \(\mathscr A^{\rm sup}>0\), also define the dimensionless leverage
\begin{equation}
 \boxed{
 \mathfrak L_{K,I}^{\rm sup}
 :=
 \frac{\mathscr A^{\rm sup}}
      {L^2\mathscr W^{\rm sup}}.}
 \label{eq:v12-superparabolic-leverage}
\end{equation}

\begin{theorem}[Bounded source-frequency leverage is globally paid]
\label{thm:v12-leverage-paid}
One always has
\begin{equation}
 \mathfrak L_{K,I}^{\rm sup}\ge1.
 \label{eq:v12-leverage-lower}
\end{equation}
Let \(I_n\) be pairwise disjoint smooth cards cut from one
Leray-energy-bounded trajectory and satisfying
\(\nu K_n^2|I_n|\le C_0\), and suppose their super-parabolic branches obey
\[
 \sup_n\mathfrak L_{K_n,I_n}^{\rm sup}\le M<\infty.
\]
Then
\begin{equation}
 \boxed{
 \sum_n |I_n|\mathscr A_n^{\rm sup}
 \le C(C_0)M\frac{E_0^2}{\nu^2}.}
 \label{eq:v12-superparabolic-summability}
\end{equation}
Consequently, nonsummable time-weighted super-parabolic critical action
forces
\begin{equation}
 \limsup_n\mathfrak L_{K_n,I_n}^{\rm sup}=\infty.
 \label{eq:v12-leverage-escape}
\end{equation}
This is the precise source-frequency leverage escape.
\end{theorem}

\begin{proof}
On the support \(\Lambda>L\),
\[
 \|A^{-1/2}z\|_2^2\le L^{-2}\|z\|_2^2.
\]
Apply this with \(z=d_{K,I}^{\rm sup}(t)\), integrate, and obtain
\(L^2\mathscr W^{\rm sup}\le\mathscr A^{\rm sup}\), proving
\eqref{eq:v12-leverage-lower}.

If the leverage is at most \(M\), then
\[
 \mathscr A_n^{\rm sup}
 \le ML_n^2\mathscr W_n^{\rm sup}.
\]
Equations \eqref{eq:v12-heat-edge}--\eqref{eq:v12-heat-edge-bound} give
\[
 |I_n|L_n^2\le\frac{M_0}{\nu}.
\]
Therefore
\[
 \sum_n|I_n|\mathscr A_n^{\rm sup}
 \le\frac{MM_0}{\nu}\sum_n\mathscr W_n^{\rm sup}
 \le\frac{MM_0}{\nu}
 \int_{\cup_n I_n}\|A^{-3/4}\BNS(u(t))\|_2^2\,\dd t.
\]
Apply the global weak-source budget
\eqref{eq:v7-weak-source-budget}.  The final assertion is the
contrapositive.
\end{proof}

\begin{assessment}[What leverage divergence says and does not say]
The quotient \(\mathfrak L_{K,I}^{\rm sup}\) is an action-weighted
frequency ratio.  Its divergence says that the \(H^{-1/2}\) source action
moves arbitrarily far beyond the heat edge relative to its
\(H^{-3/2}\) shadow.  Lemma~\ref{lem:v12-high-parent-incidence} attaches
that escape to an actual velocity parent above \(L/2\).  It does not yet
give a lower bound for the parent's energy or residence: cancellation
inside the complete high-parent assembly remains possible.
\end{assessment}

% =============================================================================
\section{Sequence consequences and the endpoint-capacity firewall}
\label{sec:v12-sequences}
% =============================================================================

\begin{corollary}[Endpoint branches are summable at doubled records]
\label{cor:v12-endpoint-summable}
Suppose \(K_n=\mathcal R_n^2/E_*\),
\(\mathcal R_{n+1}\ge2\mathcal R_n\), and the constants \(C_0\) in
\eqref{eq:v12-card-assumptions} are uniform.  On every endpoint branch of
Theorem~\ref{thm:v12-heat-edge-trichotomy},
\begin{equation}
 \boxed{
 \sum_n |I_n|\mathscr A_n^{\rm ext}
 \le C(C_0)E_*^2\sum_n\mathcal R_n^{-2}<\infty.}
 \label{eq:v12-endpoint-record-sum}
\end{equation}
Moreover the actual response amplitudes themselves satisfy
\begin{equation}
 \boxed{
 \sum_n\|\mathfrak r_{K_n,I_n}^{\rm par}\|_2
 \le2E_*\sum_n\mathcal R_n^{-1}<\infty.}
 \label{eq:v12-endpoint-response-sum}
\end{equation}
No disjointness of the intervals is needed for either conclusion.
\end{corollary}

\begin{proof}
Equation~\eqref{eq:v12-endpoint-capacity} gives
\[
 |I_n|\mathscr A_n^{\rm ext}
 \le C(C_0)\frac{E_*}{K_n}
 =C(C_0)\frac{E_*^2}{\mathcal R_n^2}.
\]
The last series is geometric.  The capacity estimate in the proof of
Theorem~\ref{thm:v12-heat-edge-trichotomy} also gives
\[
 \|\mathfrak r_{K_n,I_n}^{\rm par}\|_2
 \le2\sqrt{\frac{E_*}{K_n}}
 =\frac{2E_*}{\mathcal R_n}.
\]
Summation proves \eqref{eq:v12-endpoint-response-sum}.
\end{proof}

\begin{corollary}[Application to the V34 intrinsic interval]
\label{cor:v12-v34-intrinsic-card}
For the literal V34 interval \(J_n^{\rm occ}\),
\begin{equation}
 \nu K_n^2|J_n^{\rm occ}|<8R_t.
 \label{eq:v12-v34-parabolic-incidence}
\end{equation}
If the exterior action is nonzero, Theorem~\ref{thm:v12-heat-edge-trichotomy}
applies on that same interval with \(C_0=8R_t\).  It yields
super-parabolic leverage escape, a capacity-bounded endpoint, or complete
parabolic-band covariant variation.

This statement does not identify the exterior action
\(\mathscr A_n^{\rm ext}\) with the V34 history-fresh action
\(\mathscr A_n^{\rm occ}\).
\end{corollary}

\begin{proof}
V34 proves
\(\nu K_n^2\tau_n^{\rm occ}<4R_t\) and
\(|J_n^{\rm occ}|\le2\tau_n^{\rm occ}\).  Multiply the two inequalities.
All objects in Theorem~\ref{thm:v12-heat-edge-trichotomy} are then formed
with \(I=J_n^{\rm occ}\), without changing its endpoints.
The final warning follows from
Proposition~\ref{prop:v12-annular-independence} and the V11 five-piece
atlas.
\end{proof}

\begin{proposition}[An \(L_t^2\)-to-Duhamel debit cannot close the V34 core]
\label{prop:v12-native-response-no-go}
Assume, in addition to the V34 contiguous-core summability,
\begin{equation}
 \sum_n\mathcal R_n\sqrt{|J_n^{\rm occ}|}<\infty,
 \label{eq:v12-inherited-core-sum}
\end{equation}
that a separately acquired exterior action window satisfies
\begin{equation}
 \mathscr A_n^{\rm ext}\le C\nu\mathcal R_n^2.
 \label{eq:v12-exterior-action-window}
\end{equation}
Then
\begin{equation}
 \boxed{
 \sum_n\sqrt{|J_n^{\rm occ}|\mathscr A_n^{\rm ext}}<\infty,
 \qquad
 \sum_n|J_n^{\rm occ}|\mathscr A_n^{\rm ext}<\infty.}
 \label{eq:v12-native-response-sum}
\end{equation}
Therefore any argument which sees the source only through its native
Cauchy--Schwarz response scale
\(\sqrt{|J_n^{\rm occ}|\mathscr A_n^{\rm ext}}\) cannot turn the inherited
core into a nonsummable source debit.  It must use source-frequency
leverage, covariant variation, a record-visible derivative, or repeated
ancestry.
\end{proposition}

\begin{proof}
By \eqref{eq:v12-exterior-action-window},
\[
 \sqrt{|J_n^{\rm occ}|\mathscr A_n^{\rm ext}}
 \le\sqrt{C\nu}\,
 \mathcal R_n\sqrt{|J_n^{\rm occ}|}.
\]
Summing proves the first assertion in
\eqref{eq:v12-native-response-sum}.  Its summands form an \(\ell^1\)
sequence and hence their squares form an \(\ell^1\) sequence, proving the
second assertion.  The last statement records the exact loss in the
displayed estimate; it makes no assertion that
\eqref{eq:v12-exterior-action-window} follows from V34.
\end{proof}

\begin{firewall}[The source-action incidence is still a separate row]
Corollary~\ref{cor:v12-v34-intrinsic-card} places the heat edge on the
literal V34 time interval.  It does not transfer the positive action of
the history-fresh raw residual to the complete exterior source.  The
independence example shows that the outer-parent term is separate, while
\eqref{eq:v12-complete-annular-source} shows that it can interfere with
the raw term.  A lower action bound needs control of that signed interference.
The heat-edge trichotomy begins only after
\(\mathscr A_n^{\rm ext}\) is acquired from the actual solution.
\end{firewall}

% =============================================================================
\section{Sharpness on the concentrated packet}
\label{sec:v12-packet-sharpness}
% =============================================================================

\begin{proposition}[The packet saturates the endpoint capacity]
\label{prop:v12-packet-capacity-sharp}
For the V11 packet card,
\[
 K_N\asymp N,\qquad
 \tau_N\asymp N^{-5/2},\qquad
 L_N\asymp N^{5/4}.
\]
After the already permitted fixed reductions of \(\delta\) and \(s_\#\),
\begin{align}
 \mathscr A_N^{\rm sup}
 &=o(\mathscr A_N^{\rm ext}),
 \label{eq:v12-packet-no-heat-escape}\\
 \|\mathfrak r_{K_N,I_N}^{\rm par}\|_2^2
 &\asymp
 \tau_N\mathscr A_N^{\rm ext}
 \asymp\frac{E_*}{K_N}
 \asymp N^{-1}.
 \label{eq:v12-packet-capacity-saturation}
\end{align}
Thus the power \(K^{-1}\) in
\eqref{eq:v12-endpoint-capacity} cannot be improved for all smooth
fixed-energy data.
\end{proposition}

\begin{proof}
The first two scales are V10--V11, and
\eqref{eq:v12-heat-edge} gives \(L_N\asymp N^{5/4}\).
In rescaled variables the super-parabolic threshold is
\(L_N/N\asymp N^{1/4}\to\infty\).  Uniform high-Sobolev control makes the
critical-source tail above that threshold tend to zero, proving
\eqref{eq:v12-packet-no-heat-escape}.

At \(s=0\), the rescaled parabolic source converges to
\[
 (-\Delta)^{-1/4}\mathsf Q_\rho\BNS(U).
\]
As \(\rho\downarrow0\), this tends to the nonzero vector
\((-\Delta)^{-1/4}\BNS(U)\).  Choose the free \(\delta\), hence \(\rho\),
small enough that this vector has a fixed positive norm.  The same
short-time continuity and vanishing-viscosity argument as in
Theorem~\ref{thm:v11-far-Duhamel-endpoint} then gives
\[
 \|\mathfrak r_{K_N,I_N}^{\rm par}\|_2\asymp\beta N^{-1/2}.
\]
The exterior action has the V10 scale
\(\mathscr A_N^{\rm ext}\asymp\beta^3N^{3/2}\), while
\(\tau_N\asymp(\beta N^{5/2})^{-1}\).
Finally \(E_*/K_N\asymp N^{-1}\).  These identities prove
\eqref{eq:v12-packet-capacity-saturation}.
\end{proof}

\begin{assessment}[The endpoint theorem is sharp but not a regularity gain]
The packet occupies the coherent endpoint branch and saturates every scale
in the energy-capacity estimate.  Hence no smaller universal endpoint
capacity can be extracted from kinetic energy and heat propagation alone.
At doubled records the resulting squared responses are summable.  A global
closure must therefore show that a terminal trajectory repeatedly pays a
record-visible quantity, or exclude the variation and leverage branches by
additional dynamics.
\end{assessment}

% =============================================================================
\section{Integrated V12 conclusion and next acquisition order}
\label{sec:v12-conclusion}
% =============================================================================

\begin{theorem}[What V12 proves]
\label{thm:v12-conclusion}
Relative to V1--V11 and the frozen manuscripts, the following are proved.
\begin{enumerate}[label=\textup{(V12.\arabic*)},leftmargin=5.7em]
\item An explicit two-mode datum has zero V34 raw and incoming source but
      nonzero outer-parent source in the first exterior octave.
\item The fully assembled exterior source has the exact heat-edge
      trichotomy \eqref{eq:v12-superparabolic-branch},
      \eqref{eq:v12-endpoint-lower}, or
      \eqref{eq:v12-variation-lower}, with no dyadic logarithmic loss.
\item The coherent endpoint branch obeys the sharp capacity
      \(|I|\mathscr A^{\rm ext}\lesssim E_0/K\).
\item Bounded super-parabolic critical-to-weak leverage is paid by the
      global \(H^{-3/2}\) source budget; failure is the quantitative
      leverage escape \eqref{eq:v12-leverage-escape}.
\item The V34 intrinsic interval satisfies the exact parabolic incidence
      needed for the trichotomy, but its raw history-fresh action is not the
      complete exterior action.
\item Native Duhamel response costs are summable at doubled records and
      under the V34 core law; the V11 packet proves that this scaling is
      sharp.
\end{enumerate}
No full Navier--Stokes stability or global-regularity theorem is obtained.
\end{theorem}

\begin{proof}
Item \textup{(V12.1)} is
Proposition~\ref{prop:v12-annular-independence}.  Items
\textup{(V12.2)--(V12.3)} are
Theorem~\ref{thm:v12-heat-edge-trichotomy} and
Proposition~\ref{prop:v12-parabolic-covariant-source}.  Item
\textup{(V12.4)} is Lemma~\ref{lem:v12-high-parent-incidence} and
Theorem~\ref{thm:v12-leverage-paid}.  Item \textup{(V12.5)} is
Corollary~\ref{cor:v12-v34-intrinsic-card}.  The final item is
Corollary~\ref{cor:v12-endpoint-summable},
Proposition~\ref{prop:v12-native-response-no-go}, and
Proposition~\ref{prop:v12-packet-capacity-sharp}.
\end{proof}

\begin{programme}[V13 acquisition order]
\label{prog:v13}
\begin{enumerate}[label=\textup{(AA\arabic*)},leftmargin=4.5em]
\item Work on an arbitrary terminal sequence only after the complete
      exterior action has an actual incidence with the selected record
      interval.  Do not infer that incidence from the raw residual.
\item On the super-parabolic branch, compare the high-parent identity
      \eqref{eq:v12-high-parent-identity} with the V6 dissipation index.
      Determine whether leverage escape forces active velocity residence or
      can occur with vanishing high-parent kinetic stock.
\item On the variation branch, dyadically assemble the two terms in
      \eqref{eq:v12-parabolic-covariant-source}.  Preserve cancellation
      between the viscous commutator and cubic source motion before opening
      parent types.
\item On the endpoint branch, use the record-visible response
      \(A^{1/2}d_{K,I}^{\rm par}\), not the summable native response alone,
      and identify its exact relation to the V7 shell-upcrossing debit.
\item Test whether cofinal leverage, covariant variation, or repeated
      record-visible response has a nonsummable consequence under the Leray
      budget.  If every route returns the full critical source action,
      record that continuation circle explicitly and move upstream to
      source/record incidence.
\item Keep represented history behind the V3 incidence gate and material
      localization behind the V8 commutator.
\item At V15 and every fifth NS version thereafter, inspect the latest
      Standard-Model and Yang--Mills notes in \texttt{latex\_notes} before
      choosing the following NS direction, and record any imported lemma,
      shared obstruction, or type mismatch explicitly.
\end{enumerate}
\end{programme}

\begin{assessment}[Position after V12]
The qualitative all-shell escape now has a quantitative heat-frequency
meaning.  Sub-parabolic exterior action either reaches an endpoint or
creates a named covariant-motion packet; super-parabolic action is paid
unless its critical-to-weak spectral leverage diverges.  The endpoint
capacity is sharp and summable, so it cannot by itself close the terminal
thread.  The remaining alternatives are high-parent leverage, assembled
source motion, and record-visible response, preceded by the still missing
incidence from the selected V34 residual to the complete exterior action.
\end{assessment}

% =============================================================================
\section*{Version V12 history}
\addcontentsline{toc}{section}{Version V12 history}
% =============================================================================

Kept the first exterior octave in complete raw-plus-outer assembly and gave
an explicit outer-parent witness with both V34 source rows zero.  Split the
complete exterior source at its intrinsic heat edge and proved an exhaustive
super-parabolic/endpoint/covariant-variation trichotomy without selecting a
single shell.  Derived the sharp \(E_0/K\) endpoint capacity, a globally
paid bounded-leverage theorem from the weak-source budget, and the precise
unbounded leverage escape.  Applied the parabolic incidence to the literal
V34 core without identifying its residual action with the exterior action.
Proved native-response summability and showed by the concentrated packet
that the endpoint capacity is sharp.  No global regularity claim was made.
% Future versions are appended immediately above the sole terminator.

% =============================================================================
% VERSION V13 APPEND
% =============================================================================

\section{V13: record-forced exterior incidence and response--source closure audit}
\label{sec:v13-entry}
% =============================================================================

V12 deliberately stopped before identifying the action of the complete
exterior source with a selected V34 raw-residual card.  That stop remains
valid: the raw residual is not the full \(Q_K\BNS(u)\) source.  There is,
however, a different incidence mechanism which starts directly from an
actual increase of the critical velocity record.  A first doubling forces a
large \(A^{1/4}Q_Ku\) endpoint, while the Leray energy cap limits the
\(A^{1/4}P_Ku\) endpoint.  Duhamel's formula then forces the \emph{complete}
exterior source on a final parabolic subinterval.  This gives the missing
primal source card without transporting any V34 component action.

This version proves that incidence theorem and then audits all three V12
exits.  The record-visible endpoint response is shown to be exactly the
bandwise form of the V7 critical-source debit.  Conversely, an actual
small-data family proves that arbitrarily large super-parabolic leverage can
coexist with vanishing high-parent kinetic energy and with dissipation index
\(Q_\nu=-1\).  Finally the V12 covariant-variation row is written as one
exact Fourier assembly, exposing the signed cancellation which any upper
budget must retain.

% =============================================================================
\section{A universal record-visible response estimate}
\label{sec:v13-response}
% =============================================================================

Let \(\Pi\) be any fixed orthogonal spectral projector which commutes with
\(A\), and put
\begin{equation}
 d_\Pi(t):=A^{-1/4}\Pi\BNS(u(t)),\qquad
 x_\Pi(t):=A^{1/4}\Pi u(t).
 \label{eq:v13-projected-source-state}
\end{equation}
For \(I=[a,b]\), define the response visible to the critical record by
\begin{equation}
 \rho_{\Pi,I}:=
 \int_a^b e^{-\nu A(b-t)}A^{1/2}d_\Pi(t)\,\dd t.
 \label{eq:v13-record-response}
\end{equation}

\begin{lemma}[Exact record response and its sharp source-action bound]
\label{lem:v13-record-response}
For every smooth interval \(I=[a,b]\),
\begin{align}
 x_\Pi(b)
 &=e^{-\nu A(b-a)}x_\Pi(a)-\rho_{\Pi,I},
 \label{eq:v13-record-Duhamel}\\
 \boxed{\quad
 \|\rho_{\Pi,I}\|_2^2
 &\le\frac1{2\nu}
       \int_a^b\|d_\Pi(t)\|_2^2\,\dd t.\quad}
 \label{eq:v13-record-response-upper}
\end{align}
Consequently, if
\begin{equation}
 \|x_\Pi(b)\|_2
 \ge \|x_\Pi(a)\|_2+\delta
 \qquad(\delta>0),
 \label{eq:v13-projected-growth}
\end{equation}
then
\begin{equation}
 \boxed{\quad
 \int_a^b\|d_\Pi(t)\|_2^2\,\dd t
 \ge2\nu\delta^2.\quad}
 \label{eq:v13-growth-source-debit}
\end{equation}
\end{lemma}

\begin{proof}
Apply \(A^{1/4}\Pi\) to \eqref{eq:NS}.  Since
\(\Pi\BNS(u)=A^{1/4}d_\Pi\),
\[
 (\partial_t+\nu A)x_\Pi=-A^{1/2}d_\Pi,
\]
and Duhamel's formula is \eqref{eq:v13-record-Duhamel}.

Write \(\lambda\) for an eigenvalue of \(A^{1/2}\).  For each Fourier
coefficient, Cauchy--Schwarz gives
\[
 \left|
  \int_a^b\lambda e^{-\nu\lambda^2(b-t)}
       \widehat d_\Pi(\lambda,t)\,\dd t
 \right|^2
 \le
 \left(\int_a^b
       \lambda^2e^{-2\nu\lambda^2(b-t)}\,\dd t\right)
 \left(\int_a^b|\widehat d_\Pi(\lambda,t)|^2\,\dd t\right).
\]
The first factor equals
\[
 \frac{1-e^{-2\nu\lambda^2(b-a)}}{2\nu}
 \le\frac1{2\nu}.
\]
Sum over the orthogonal Fourier modes to obtain
\eqref{eq:v13-record-response-upper}.  Heat contraction and
\eqref{eq:v13-record-Duhamel} imply
\[
 \|\rho_{\Pi,I}\|_2
 \ge\|x_\Pi(b)\|_2-\|x_\Pi(a)\|_2.
\]
Under \eqref{eq:v13-projected-growth}, combine this with
\eqref{eq:v13-record-response-upper}.
\end{proof}

\begin{remark}[Why the derivative is \(A^{1/2}d_\Pi\)]
The native V12 response uses \(d_\Pi\) and is visible in
\(A^{-1/4}\Pi u\).  The critical record uses \(A^{1/4}\Pi u\), exactly one
spatial derivative higher.  Its Duhamel integrand is therefore
\(A^{1/2}d_\Pi\).  The heat kernel pays that derivative with the factor
\((2\nu)^{-1}\) in \eqref{eq:v13-record-response-upper}; it does not convert
the action to the globally paid \(A^{-3/4}\BNS(u)\) norm.
\end{remark}

% =============================================================================
\section{A first record doubling forces a complete exterior source card}
\label{sec:v13-record-incidence}
% =============================================================================

Set
\[
 \mathcal R(t):=\|A^{1/4}u(t)\|_2
\]
and suppose throughout this section that
\begin{equation}
 \sup_{t\in[a,b]}\|u(t)\|_2^2\le E_*
 \label{eq:v13-energy-cap}
\end{equation}
for some \(E_*>0\).

\begin{theorem}[Recent exterior-source incidence forced by a record doubling]
\label{thm:v13-record-exterior-incidence}
Let \(R>0\), assume
\begin{equation}
 \mathcal R(a)=R,\qquad
 \mathcal R(t)<2R\quad(a\le t<b),\qquad
 \mathcal R(b)=2R,
 \label{eq:v13-first-doubling}
\end{equation}
and set
\begin{equation}
 K:=\frac{R^2}{E_*},\qquad
 c_{\rm h}:=\log4,\qquad
 s:=\max\left\{a,b-\frac{c_{\rm h}}{\nu K^2}\right\},
 \qquad I^{\rm rec}:=[s,b].
 \label{eq:v13-recent-card}
\end{equation}
Then
\begin{align}
 \nu K^2|I^{\rm rec}|&\le c_{\rm h},
 \label{eq:v13-recent-parabolic-incidence}\\
 \left\|
 \int_s^b e^{-\nu A(b-t)}
      A^{1/2}d_{K}^{\rm ext}(t)\,\dd t
 \right\|_2
 &\ge(\sqrt3-1)R,
 \label{eq:v13-exterior-record-response-lower}\\
 \boxed{\quad
 \mathscr A_{I^{\rm rec}}^{\rm ext}
 :=\int_s^b
       \|A^{-1/4}Q_K\BNS(u(t))\|_2^2\,\dd t
 \ge c_{\rm rec}\nu R^2,\quad}
 \label{eq:v13-exterior-action-lower}
\end{align}
where
\begin{equation}
 c_{\rm rec}:=2(\sqrt3-1)^2>0.
 \label{eq:v13-record-constant}
\end{equation}
The source in \eqref{eq:v13-exterior-action-lower} is the fully assembled
exterior source.  No raw, annular, mixed, or outer--outer parent component
has been separated.
\end{theorem}

\begin{proof}
For every \(t\in[a,b]\), the energy cap and the definition of \(K\) give
\begin{equation}
 \|A^{1/4}P_Ku(t)\|_2^2
 \le K\|u(t)\|_2^2
 \le R^2.
 \label{eq:v13-head-critical-cap}
\end{equation}
At \(t=b\), orthogonality of \(P_K\) and \(Q_K\) in every Stokes norm yields
\[
 \|A^{1/4}Q_Ku(b)\|_2^2
 =\mathcal R(b)^2-\|A^{1/4}P_Ku(b)\|_2^2
 \ge4R^2-R^2=3R^2.
 \tag{1}\label{eq:v13-final-exterior-stock}
\]
The first row of the conclusion follows directly from
\eqref{eq:v13-recent-card}.

Put \(x(t)=A^{1/4}Q_Ku(t)\).  If \(s=a\), then
\(\|x(s)\|_2\le\mathcal R(a)=R\).  If \(s>a\), the first-hitting property in
\eqref{eq:v13-first-doubling} gives \(\|x(s)\|_2\le2R\), while spectral
support in \(Q_K\) gives
\[
 \|e^{-\nu A(b-s)}x(s)\|_2
 \le e^{-\nu K^2(b-s)}\|x(s)\|_2
 =\frac14\|x(s)\|_2
 \le\frac R2.
\]
Thus in both cases
\begin{equation}
 \|e^{-\nu A(b-s)}x(s)\|_2\le R.
 \label{eq:v13-old-exterior-bound}
\end{equation}
Use \eqref{eq:v13-record-Duhamel} with \(\Pi=Q_K\), then combine
\eqref{eq:v13-final-exterior-stock} and
\eqref{eq:v13-old-exterior-bound}:
\[
 \|\rho_{Q_K,I^{\rm rec}}\|_2
 \ge\sqrt3R-R=(\sqrt3-1)R.
\]
This is \eqref{eq:v13-exterior-record-response-lower}.  Finally
\eqref{eq:v13-record-response-upper}, again with \(\Pi=Q_K\), gives
\[
 \mathscr A_{I^{\rm rec}}^{\rm ext}
 \ge2\nu\|\rho_{Q_K,I^{\rm rec}}\|_2^2
 \ge2(\sqrt3-1)^2\nu R^2.
\]
\end{proof}

\begin{corollary}[Complete exterior incidence on an unbounded doubled-record
sequence]
\label{cor:v13-doubled-record-source-divergence}
Suppose a smooth trajectory on \([t_0,T_*)\) has unbounded critical record.
Choose first hitting times
\begin{equation}
 t_n:=\inf\{t\ge t_0:\mathcal R(t)=R_n\},
 \qquad R_n=2^nR_0,
 \label{eq:v13-first-hitting-times}
\end{equation}
with \(R_0>\mathcal R(t_0)\), and apply
Theorem~\ref{thm:v13-record-exterior-incidence} to
\([t_n,t_{n+1}]\).  The resulting recent cards \(I_n^{\rm rec}\) are
pairwise disjoint and, with \(K_n=R_n^2/E_*\),
\begin{equation}
 \boxed{
 \sum_n\int_{I_n^{\rm rec}}
       \|A^{-1/4}Q_{K_n}\BNS(u(t))\|_2^2\,\dd t
 =\infty.}
 \label{eq:v13-exterior-action-divergence}
\end{equation}
In particular,
\begin{equation}
 \int_{t_0}^{T_*}
       \|A^{-1/4}\BNS(u(t))\|_2^2\,\dd t=\infty.
 \label{eq:v13-full-critical-source-divergence}
\end{equation}
\end{corollary}

\begin{proof}
The intervals \([t_n,t_{n+1}]\) have disjoint interiors, and each recent
card is contained in its corresponding interval.  Hence the recent cards
are pairwise disjoint.  Theorem~\ref{thm:v13-record-exterior-incidence}
gives
\[
 \sum_n\mathscr A_n^{\rm ext}
 \ge c_{\rm rec}\nu\sum_nR_n^2=\infty.
\]
Since \(Q_{K_n}\) is an \(L^2\) contraction and the time cards are disjoint,
the left side is no larger than the full critical-source action on
\([t_0,T_*)\).
\end{proof}

\begin{firewall}[What incidence has and has not been acquired]
\label{fw:v13-incidence-scope}
Equation~\eqref{eq:v13-exterior-action-lower} resolves the primal
source/record incidence left open in V11--V12: a genuine record doubling
creates a complete exterior source card at the canonical energy cutoff.
It does not identify this card with the V34 history-fresh raw residual, nor
does it prove that the critical record of every putative singular trajectory
is unbounded.  Corollary~\ref{cor:v13-doubled-record-source-divergence} is
conditional on that explicit record chronology.
\end{firewall}

% =============================================================================
\section{The V12 heat-edge alternatives now occur on actual record cards}
\label{sec:v13-record-trichotomy}
% =============================================================================

On the recent card \(I^{\rm rec}\), form the V12 heat edge
\begin{equation}
 L=\max\left\{2K,\frac1{\sqrt{\nu|I^{\rm rec}|}}\right\}.
 \label{eq:v13-record-heat-edge}
\end{equation}

\begin{corollary}[Record-generated complete-source trichotomy]
\label{cor:v13-record-trichotomy}
Under the hypotheses of
Theorem~\ref{thm:v13-record-exterior-incidence}, the V12 heat-edge theorem
applies with \(C_0=\log4\).  For a constant \(c>0\) depending only on that
fixed number, at least one of the following holds:
\begin{align}
 \mathscr A^{\rm sup}
 &\ge c\nu R^2,
 \label{eq:v13-record-super-branch}\\
 \|\mathfrak r_{K,I^{\rm rec}}^{\rm par}\|_2^2
 &\ge c\nu |I^{\rm rec}|R^2,
 \label{eq:v13-record-native-endpoint-branch}\\
 \mathscr V_{K,I^{\rm rec}}^{\rm par}
 &\ge c\frac{\nu R^2}{|I^{\rm rec}|^2}.
 \label{eq:v13-record-variation-branch}
\end{align}
Independently of which of these three alternatives occurs, the full
record-visible response obeys
\eqref{eq:v13-exterior-record-response-lower}.
\end{corollary}

\begin{proof}
The parabolic incidence is
\eqref{eq:v13-recent-parabolic-incidence}, and the positive exterior action
is \eqref{eq:v13-exterior-action-lower}.  Substitute that lower bound into
the three alternatives
\eqref{eq:v12-superparabolic-branch},
\eqref{eq:v12-endpoint-lower}, and
\eqref{eq:v12-variation-lower}.  The full response statement was proved
before making the heat-edge split.
\end{proof}

\begin{assessment}[A genuine advance, but not a paid contradiction]
The trichotomy no longer rests on a separately assumed exterior action
window: every first doubling supplies one.  Its action scale
\(\nu R^2\), however, is precisely a critical-source scale.  The Leray
identity controls neither the divergent sum
\(\sum_n\nu R_n^2\) nor the variation scale in
\eqref{eq:v13-record-variation-branch}.  The remaining question is whether
one of the three alternatives can be descended to the globally paid weak
source or to kinetic-energy dissipation without losing its record incidence.
\end{assessment}

% =============================================================================
\section{The record-visible endpoint is exactly the V7 crossing debit}
\label{sec:v13-v7-equivalence}
% =============================================================================

For the parabolic band of V12, define
\begin{equation}
 \rho_{K,I}^{\rm par}
 :=
 \int_a^b e^{-\nu A(b-t)}
       A^{1/2}d_{K,I}^{\rm par}(t)\,\dd t.
 \label{eq:v13-par-record-response}
\end{equation}

\begin{proposition}[Native and record-visible responses on one band]
\label{prop:v13-two-responses}
The V12 native response and the critical-record response satisfy
\begin{align}
 \rho_{K,I}^{\rm par}
 &=A^{1/2}\mathfrak r_{K,I}^{\rm par},
 \label{eq:v13-response-derivative}\\
 K\|\mathfrak r_{K,I}^{\rm par}\|_2
 \le\|\rho_{K,I}^{\rm par}\|_2
 &\le L\|\mathfrak r_{K,I}^{\rm par}\|_2,
 \label{eq:v13-response-band-comparison}\\
 A^{1/4}\Pi_{K,L}u(b)
 &=e^{-\nu A(b-a)}A^{1/4}\Pi_{K,L}u(a)
   -\rho_{K,I}^{\rm par},
 \label{eq:v13-par-record-Duhamel}\\
 \boxed{\quad
 \|\rho_{K,I}^{\rm par}\|_2^2
 &\le\frac1{2\nu}\mathscr A^{\rm par}.\quad}
 \label{eq:v13-par-record-source-bound}
\end{align}
Thus differentiating the native endpoint makes it visible to the critical
record only by returning to the full \(H^{-1/2}\) source action.
\end{proposition}

\begin{proof}
Functional calculus permits \(A^{1/2}\) to pass through the heat integral,
which proves \eqref{eq:v13-response-derivative}.  On
\(\operatorname{Ran}\Pi_{K,L}\), the spectrum of \(A^{1/2}\) lies in
\((K,L]\), proving \eqref{eq:v13-response-band-comparison}.  The last two
rows are Lemma~\ref{lem:v13-record-response} with
\(\Pi=\Pi_{K,L}\).
\end{proof}

\begin{proposition}[The V7 upcrossing debit is the one-shell instance]
\label{prop:v13-v7-response-equivalence}
Let \(\Delta_q\), \(r_q\), \(g_q\), and \(\sigma\) be as in V7.  If
\[
 r_q(s)\le\frac{\sigma\nu}{2},
 \qquad r_q(t)\ge\sigma\nu,
\]
then Lemma~\ref{lem:v13-record-response}, applied with
\(\Pi=\Delta_q\), yields
\begin{equation}
 \int_s^t\|g_q(\tau)\|_2^2\,\dd\tau
 \ge
 2\left(1-\frac1{\sqrt2}\right)^2\sigma^2\nu^3.
 \label{eq:v13-explicit-upcrossing-debit}
\end{equation}
In particular this is the same frequency-independent critical-source debit
as Theorem~\ref{thm:v7-upcrossing}, now obtained from the general spectral
response operator.
\end{proposition}

\begin{proof}
On the sharp shell,
\[
 r_q(\tau)
 \le\|A^{1/4}u_q(\tau)\|_2
 \le\sqrt2\,r_q(\tau).
\]
Consequently the final critical shell norm is at least \(\sigma\nu\), while
the initial norm is at most \(\sigma\nu/\sqrt2\).  Heat contraction cannot
increase the latter.  Hence the shell record-visible response has norm at
least
\((1-2^{-1/2})\sigma\nu\).  Since
\(d_{\Delta_q}=g_q\), use
\eqref{eq:v13-record-response-upper}.
\end{proof}

\begin{firewall}[The record-visible response is not a new global budget]
\label{fw:v13-response-circle}
The record-visible endpoint fixes the summable native-response scaling from
V12: multiplication by a frequency \(K_n\asymp R_n^2/E_*\) can turn a
native response of order \(E_*/R_n\) into a critical response of order
\(R_n\).  Proposition~\ref{prop:v13-two-responses} shows exactly what pays
for that gain: the unbounded critical source action.  Thus the route is a
clean incidence theorem and an exact recovery of the V7 debit, not a
continuation estimate from the Leray budget.
\end{firewall}

% =============================================================================
\section{Leverage does not force active high-parent residence}
\label{sec:v13-leverage-counterexample}
% =============================================================================

The V12 high-parent identity is an operator-incidence statement.  We next
show that neither high-parent kinetic stock nor the V6 dissipation index can
be recovered from source-frequency leverage alone, even on actual smooth
Navier--Stokes trajectories.

\begin{theorem}[Actual leverage escape with vanishing high-parent energy]
\label{thm:v13-leverage-no-residence}
Fix \(\nu>0\).  There are integers \(m\to\infty\), positive amplitudes
\(\varepsilon_m\to0\), and smooth Navier--Stokes solutions \(u^{(m)}\) on
\begin{equation}
 I_m=\left[0,\frac1{\nu m}\right]
 \label{eq:v13-small-data-card}
\end{equation}
with the following properties.  Take \(K_m=1\), so the V12 heat edge is
\begin{equation}
 L_m=\max\{2,(\nu|I_m|)^{-1/2}\}=\sqrt m
 \qquad(m\ge4).
 \label{eq:v13-small-data-heat-edge}
\end{equation}
Then
\begin{align}
 \mathscr A_m^{\rm sup}&>0,
 \label{eq:v13-small-data-positive-source}\\
 \mathfrak L_{K_m,I_m}^{\rm sup}&\asymp m\longrightarrow\infty,
 \label{eq:v13-small-data-leverage}\\
 \sup_{t\in I_m}\|Q_{L_m/2}u^{(m)}(t)\|_2^2
 &\longrightarrow0,
 \label{eq:v13-small-data-parent-energy}\\
 Q_\nu(t)&=-1\qquad(t\in I_m).
 \label{eq:v13-small-data-index}
\end{align}
Therefore unbounded super-parabolic leverage does not imply an active
dissipation-index shell or any positive lower bound for high-parent kinetic
residence.
\end{theorem}

\begin{proof}
Let
\[
 p_m=(3m,0,0),\qquad q_m=(2m,-m,0),
\]
and define the real divergence-free trigonometric polynomial
\begin{equation}
 V_m(x)=\mathbf e_2\cos(p_m\cdot x)
       +\mathbf e_3\cos(q_m\cdot x).
 \label{eq:v13-two-mode-datum}
\end{equation}
For \(\varepsilon>0\), let \(u^{(m,\varepsilon)}\) be the smooth solution
with initial value \(\varepsilon V_m\).  For fixed \(m\), small-data theory
and smooth dependence at the zero solution give, in every fixed sufficiently
high Sobolev space,
\begin{equation}
 \varepsilon^{-1}u^{(m,\varepsilon)}
 \longrightarrow z_m:=e^{-\nu At}V_m
 \quad\hbox{on }I_m
 \qquad(\varepsilon\downarrow0).
 \label{eq:v13-linearized-limit}
\end{equation}
The convergence is strong enough that the corresponding quadratic sources,
after division by \(\varepsilon^2\), converge in both source norms used
below.

Because \(|p_m|^2=9m^2\) and \(|q_m|^2=5m^2\),
\[
 z_m(t)=e^{-9\nu m^2t}\mathbf e_2\cos(p_m\cdot x)
       +e^{-5\nu m^2t}\mathbf e_3\cos(q_m\cdot x).
\]
A direct calculation gives
\begin{equation}
 \BNS(z_m(t))
 =
 \frac m2e^{-14\nu m^2t}\mathbf e_3
 \left[
  \sin((p_m+q_m)\cdot x)
  -\sin((p_m-q_m)\cdot x)
 \right].
 \label{eq:v13-linearized-source}
\end{equation}
Indeed only
\((\mathbf e_2\cos(p_m\cdot x)\cdot\nabla)
 (\mathbf e_3\cos(q_m\cdot x))\)
survives.  Both output wavevectors have zero third component, so the Leray
projection leaves the \(\mathbf e_3\) polarization unchanged.  Their
lengths are \(\sqrt{26}\,m\) and \(\sqrt2\,m\), both greater than \(L_m\).

Let \(c_A,c_W>0\) denote the harmless constants determined by the Fourier
normalization.  Orthogonality of the two outputs and integration of
\(e^{-28\nu m^2t}\) over \(I_m\) give
\begin{align}
 \int_{I_m}
 \|A^{-1/4}Q_{L_m}\BNS(z_m(t))\|_2^2\,\dd t
 &=
 \frac{c_A}{\nu m}(1-e^{-28m}),
 \label{eq:v13-leading-critical-action}\\
 \int_{I_m}
 \|A^{-3/4}Q_{L_m}\BNS(z_m(t))\|_2^2\,\dd t
 &=
 \frac{c_W}{\nu m^3}(1-e^{-28m}).
 \label{eq:v13-leading-weak-action}
\end{align}
For example,
\[
 c_A\ \hbox{is proportional to}\
 26^{-1/2}+2^{-1/2},
 \qquad
 c_W\ \hbox{is proportional to}\
 26^{-3/2}+2^{-3/2}.
\]
By \eqref{eq:v13-linearized-limit}, for each fixed \(m\),
the actual critical and weak super-parabolic actions divided by
\(\varepsilon^4\) tend to the two positive quantities in
\eqref{eq:v13-leading-critical-action}--%
\eqref{eq:v13-leading-weak-action}.  Since \(L_m^2=m\), their quotient
therefore satisfies
\[
 \mathfrak L_{K_m,I_m}^{\rm sup}
 =
 \frac{\mathscr A_m^{\rm sup}}
      {m\mathscr W_m^{\rm sup}}
 \longrightarrow\frac{c_A}{c_W}m
 \qquad(\varepsilon\downarrow0).
 \tag{2}\label{eq:v13-leading-leverage}
\]

It remains to choose one amplitude for each \(m\).  Take
\(\varepsilon_m>0\) so small that the actual quotient in
\eqref{eq:v13-leading-leverage} lies between one half and twice its limit,
that \(\varepsilon_m\|V_m\|_2\to0\), and that
\[
 \sup_{t\in I_m}\|u^{(m,\varepsilon_m)}(t)\|_{H^s}
 <\frac{c_\theta\nu}{2C_s}
\]
for one fixed \(s>5/2\).  This diagonal choice is possible by
\eqref{eq:v13-linearized-limit}.  Set
\(u^{(m)}:=u^{(m,\varepsilon_m)}\).  The energy identity gives
\[
 \sup_{t\in I_m}\|Q_{L_m/2}u^{(m,\varepsilon_m)}(t)\|_2^2
 \le\varepsilon_m^2\|V_m\|_2^2\longrightarrow0.
\]
Finally Bernstein and the Sobolev bound imply, for every dyadic shell,
\[
 \lambda_j^{-1}\|\Delta_ju^{(m,\varepsilon_m)}(t)\|_\infty
 \le C_s\|u^{(m,\varepsilon_m)}(t)\|_{H^s}
 <c_\theta\nu.
\]
Thus the defining smallness condition
\eqref{eq:v6-dissipation-index} holds on every shell and
\(Q_\nu(t)=-1\) throughout \(I_m\).  Equations
\eqref{eq:v13-small-data-positive-source}--%
\eqref{eq:v13-small-data-index} follow.
\end{proof}

\begin{firewall}[Scope of the leverage counterexample]
\label{fw:v13-leverage-scope}
The family in Theorem~\ref{thm:v13-leverage-no-residence} has small source
action; it is not asserted to model a terminal record card satisfying
\eqref{eq:v13-exterior-action-lower}.  Its role is logically narrower and
exact: the quotient \(\mathfrak L^{\rm sup}\), even together with actual
high-parent incidence, carries no amplitude information.  A terminal
argument must combine leverage with the record-forced action lower bound
before attempting to extract velocity residence.
\end{firewall}

% =============================================================================
\section{Exact assembled symbol of the covariant-variation branch}
\label{sec:v13-variation-symbol}
% =============================================================================

For a nonzero Fourier mode \(k\), let \(\PP_k\) be the Leray matrix and set
\begin{equation}
 \mathcal N_k(u)
 :=
 i\PP_k\sum_{p+q=k}(\widehat u_p\cdot q)\widehat u_q
 =\widehat{\BNS(u)}_k.
 \label{eq:v13-nonlinear-Fourier-coefficient}
\end{equation}

\begin{proposition}[Fourier assembly of the parabolic source motion]
\label{prop:v13-covariant-symbol}
For \(K<|k|\le L\), define
\begin{align}
 \mathcal C_k^\nu(u)
 &:=
 2\nu i\PP_k
 \sum_{p+q=k}
 (p\cdot q)(\widehat u_p\cdot q)\widehat u_q,
 \label{eq:v13-viscous-commutator-symbol}\\
 \mathcal M_k(u)
 &:=
 i\PP_k\sum_{p+q=k}
 \left[
  (\mathcal N_p(u)\cdot q)\widehat u_q
  +(\widehat u_p\cdot q)\mathcal N_q(u)
 \right].
 \label{eq:v13-cubic-motion-symbol}
\end{align}
Then the exact coefficient of the V12 covariant derivative is
\begin{equation}
 \boxed{
 \widehat{
  (\partial_t+\nu A)d_{K,I}^{\rm par}
 }_k
 =
 |k|^{-1/2}
 \left[\mathcal C_k^\nu(u)-\mathcal M_k(u)\right].}
 \label{eq:v13-exact-covariant-symbol}
\end{equation}
Consequently
\begin{equation}
 \boxed{
 \mathscr V_{K,I}^{\rm par}
 =
 \int_a^b
 \sum_{K<|k|\le L}
 e^{-2\nu|k|^2(b-t)}|k|^{-1}
 \left|\mathcal C_k^\nu(u(t))-\mathcal M_k(u(t))\right|^2
 \,\dd t.}
 \label{eq:v13-exact-variation-assembly}
\end{equation}
The difference in
\eqref{eq:v13-exact-variation-assembly} must be assembled at each output
mode before its norm or positive part is taken.
\end{proposition}

\begin{proof}
The coefficient of \(\BNS(u)\) is
\eqref{eq:v13-nonlinear-Fourier-coefficient}.  In the quadratic
commutator,
\[
 |k|^2-|p|^2-|q|^2=2p\cdot q
 \qquad(p+q=k),
\]
so
\[
 \widehat{\nu[A\BNS(u)-\BNS_{\rm s}(Au,u)]}_k
 =\mathcal C_k^\nu(u).
\]
Likewise,
\[
 \widehat{\BNS_{\rm s}(\BNS(u),u)}_k
 =\mathcal M_k(u).
\]
Insert these identities in
\eqref{eq:v12-parabolic-covariant-source}, multiply by the
\(A^{-1/4}\) symbol \(|k|^{-1/2}\), and then use Fourier orthogonality after
applying the final heat multiplier.
\end{proof}

\begin{proposition}[The viscous commutator has no coercive source sign]
\label{prop:v13-commutator-null}
There are smooth divergence-free fields for which \(\BNS(u)\ne0\) but the
quadratic symbol \(\mathcal C_k^\nu(u)\) vanishes at every output generated
by the occupied modes.  For example, for
\[
 p=(m,0,0),\qquad q=(0,m,0),\qquad
 u=\mathbf e_2\cos(p\cdot x)+\mathbf e_3\cos(q\cdot x),
\]
the cross interaction is nonzero, while every interacting parent pair has
\(p\cdot q=0\).
\end{proposition}

\begin{proof}
Each displayed polarization is perpendicular to its own wavevector.
Moreover
\((\mathbf e_2\cdot q)=m\), so
\((\mathbf e_2\cos(p\cdot x)\cdot\nabla)
 (\mathbf e_3\cos(q\cdot x))\ne0\).
All self-interactions vanish and the reverse cross interaction vanishes
because \(\mathbf e_3\cdot p=0\).  Every surviving Fourier pairing uses one
of \(\pm p\) and one of \(\pm q\), whose scalar product is zero.  The factor
in \eqref{eq:v13-viscous-commutator-symbol} therefore kills the entire
quadratic commutator although the nonlinear source is nonzero.
\end{proof}

\begin{assessment}[What the variation lower bound asks to be paid]
Expanding the square in \eqref{eq:v13-exact-variation-assembly} produces
the indefinite cross term
\[
 -2\operatorname{Re}
  \langle\mathcal C_k^\nu,\mathcal M_k\rangle.
\]
Thus neither the viscous commutator nor the cubic motion is the physical
variation packet by itself.  Proposition~\ref{prop:v13-commutator-null}
also rules out coercivity of the quadratic part by the source action.
The V12 lower bound on \(\mathscr V^{\rm par}\) is rigorous, but no inherited
Leray, weak-source, or V34 history budget controls the assembled norm in
\eqref{eq:v13-exact-variation-assembly}.  Opening its two terms and pricing
them separately would discard a cancellation before acquiring an upper
budget.
\end{assessment}

% =============================================================================
\section{Integrated V13 conclusion and next acquisition order}
\label{sec:v13-conclusion}
% =============================================================================

\begin{theorem}[What V13 proves]
\label{thm:v13-conclusion}
Relative to V1--V12 and the frozen manuscripts, the following are proved.
\begin{enumerate}[label=\textup{(V13.\arabic*)},leftmargin=5.9em]
\item Every first doubling of the critical record forces a fully assembled
      exterior response of size \(cR\) and exterior critical-source action
      at least \(c\nu R^2\) on a final parabolic subinterval.
\item An unbounded doubled-record chronology therefore carries pairwise
      disjoint complete exterior source cards with divergent total
      \(H^{-1/2}\) source action.
\item The V12 heat-edge trichotomy now applies on those actual
      record-generated cards, without assuming an exterior action window or
      identifying the V34 raw residual with the full source.
\item The record-visible Duhamel response obeys the sharp
      \((2\nu)^{-1}\) operator bound and its one-shell specialization is
      exactly the V7 fixed upcrossing debit.
\item Super-parabolic leverage can diverge along actual smooth solutions
      while high-parent kinetic energy vanishes and \(Q_\nu=-1\); leverage
      alone therefore gives no active-parent residence.
\item The covariant-variation branch has the exact signed Fourier assembly
      \eqref{eq:v13-exact-variation-assembly}; its viscous commutator is not
      coercive and no inherited paid budget controls the assembled motion.
\end{enumerate}
No full Navier--Stokes stability or global-regularity theorem is obtained.
\end{theorem}

\begin{proof}
Items \textup{(V13.1)--(V13.3)} are
Theorem~\ref{thm:v13-record-exterior-incidence},
Corollaries~\ref{cor:v13-doubled-record-source-divergence} and
\ref{cor:v13-record-trichotomy}.  Item \textup{(V13.4)} is
Lemma~\ref{lem:v13-record-response} and
Propositions~\ref{prop:v13-two-responses}--%
\ref{prop:v13-v7-response-equivalence}.  Item \textup{(V13.5)} is
Theorem~\ref{thm:v13-leverage-no-residence}.  The final item is
Propositions~\ref{prop:v13-covariant-symbol} and
\ref{prop:v13-commutator-null}.
\end{proof}

\begin{programme}[V14 acquisition order]
\label{prog:v14}
\begin{enumerate}[label=\textup{(AB\arabic*)},leftmargin=4.8em]
\item Use the recent cards
      \eqref{eq:v13-recent-card} as the primary chronology.  The complete
      exterior action incidence is now proved there; do not move it back to
      a V34 raw-residual card.
\item Combine the lower action
      \(\mathscr A_n^{\rm ext}\gtrsim\nu R_n^2\) with the globally finite
      weak-source action.  Seek an action-weighted frequency-descent theorem,
      not an inference from the leverage quotient alone.
\item On a super-parabolic card, trace the complete high-output source
      through at least one velocity parent and ask for a quantitative debit
      in
      \(\int\nu\|A^{1/2}Q_{L/2}u\|_2^2\), or prove by a packet family that
      such a debit is impossible even after the record-action lower bound is
      imposed.
\item On a parabolic endpoint card, compare nested record-visible responses
      at \(K_n\) and \(K_{n+1}\).  Test whether their signed difference
      telescopes in the actual velocity chronology; an unsigned sum merely
      returns the critical-source divergence already proved.
\item On a variation card, work with
      \(\mathcal C_k^\nu-\mathcal M_k\) from
      \eqref{eq:v13-exact-covariant-symbol}.  Any integration by parts or
      ancestry map must be performed before separating the quadratic and
      cubic terms.
\item If every descent again loses one derivative, isolate the minimal
      source-to-velocity ancestor inequality whose proof would close the
      record cards, and test it against the V9--V13 concentrated families.
\item At V15, before choosing the V16 direction, inspect the latest
      Standard-Model and Yang--Mills note files in
      \texttt{latex\_notes}, as required by the five-version review cadence.
\end{enumerate}
\end{programme}

\begin{assessment}[Position after V13]
The source/record incidence bottleneck has been moved: a record doubling
itself, rather than a history-fresh component, supplies the complete
exterior action on an intrinsic parabolic card.  The endpoint response then
recovers exactly the V7 critical-source cost.  The other two exits remain
genuine: leverage has no amplitude content, and covariant variation is an
unpaid signed quadratic--cubic motion.  The next viable target is therefore
a frequency-descent or ancestry estimate which converts the
record-forced \(H^{-1/2}\) exterior source into the energy-paid
\(H^{-3/2}\) source or into actual high-parent dissipation without assuming
the missing derivative.
\end{assessment}

% =============================================================================
\section*{Version V13 history}
\addcontentsline{toc}{section}{Version V13 history}
% =============================================================================

Added the sharp spectral response operator for the critical record and
proved that a first record doubling forces a complete exterior source card
on its final parabolic subinterval, with response \(\gtrsim R\) and action
\(\gtrsim\nu R^2\).  Applied the V12 trichotomy to these actual
record-generated cards and identified the record-visible endpoint exactly
with the V7 shell-upcrossing debit.  Constructed an actual small-data
two-mode family with leverage tending to infinity, vanishing high-parent
kinetic energy, and \(Q_\nu=-1\), ruling out residence from leverage alone.
Derived the exact signed Fourier symbol for the viscous-commutator plus
cubic covariant source motion and exhibited a nonlinear interaction in the
commutator null geometry.  No global regularity claim was made.
% Future versions are appended immediately above the sole terminator.

% =============================================================================
% VERSION V14 APPEND
% =============================================================================

\section{V14: high-parent weak-source descent and effective response-frequency escape}
\label{sec:v14-entry}
% =============================================================================

V13 acquired the complete exterior source card generated by a genuine
critical-record doubling.  The remaining derivative gap can now be tested
without referring to the V34 raw residual.  This version performs that test
at the level of the actual nonlinear product.

There are two distinct response norms.  The critical record sees
\(\rho=\int e^{-\nu A(b-t)}A^{1/2}d\,\dd t\), whereas the energy-paid weak
source controls \(A^{-1/2}\rho\).  Their quotient is an intrinsic effective
response frequency.  We prove a high-output product estimate which attaches
the weak source to a genuine high velocity parent and gives a globally paid
parent-interaction ledger on disjoint cards.  Applied to the V13 record
cards, the ledger controls only the inverse-frequency response.  A
two-parameter concentrated packet then dynamically saturates this loss:
its record grows and its record-visible response is large, while its weak
source, native response, high-parent kinetic energy, and Leray debit are all
made summable by sending the physical packet frequency far above the
canonical cutoff.

% =============================================================================
\section{A complete high-output source has a quantitative high-parent bound}
\label{sec:v14-high-parent}
% =============================================================================

For \(M>0\), write
\begin{equation}
 h_M:=A^{-3/4}Q_M\BNS(u),\qquad
 \mathscr W_{M,I}:=\int_I\|h_M(t)\|_2^2\,\dd t.
 \label{eq:v14-high-weak-source}
\end{equation}
Define the high-parent interaction ledger
\begin{equation}
 \mathscr P_{M,I}
 :=
 \int_I
 \|Q_{M/2}u(t)\|_2^2
 \|A^{1/2}u(t)\|_2^2\,\dd t.
 \label{eq:v14-parent-ledger}
\end{equation}

\begin{theorem}[High-output weak source is paid by a high velocity parent]
\label{thm:v14-high-parent-weak}
There is an absolute periodic constant \(C_{\rm hp}\) such that, at every
smooth time,
\begin{equation}
 \boxed{
 \|A^{-3/4}Q_M\BNS(u)\|_2
 \le
 C_{\rm hp}\|Q_{M/2}u\|_2\|A^{1/2}u\|_2.}
 \label{eq:v14-pointwise-high-parent}
\end{equation}
Consequently,
\begin{equation}
 \boxed{
 \mathscr W_{M,I}\le C_{\rm hp}^2\mathscr P_{M,I}.}
 \label{eq:v14-weak-by-parent}
\end{equation}
If \(I_n\) are pairwise disjoint intervals cut from one trajectory satisfying
\(\sup_t\|u(t)\|_2^2\le E_*\), and \(M_n>0\) are arbitrary, then
\begin{equation}
 \boxed{
 \sum_n\mathscr P_{M_n,I_n}
 \le
 E_*\int_{\cup_nI_n}\|A^{1/2}u(t)\|_2^2\,\dd t
 \le\frac{E_*^2}{2\nu}.}
 \label{eq:v14-parent-global-budget}
\end{equation}
Thus \(\mathscr P\) is an actual globally paid high-parent ledger; no
persistence time is inserted as a hypothesis.
\end{theorem}

\begin{proof}
Put
\[
 v=P_{M/2}u,\qquad w=Q_{M/2}u.
\]
The additive Fourier support of \(\BNS(v,v)\) lies in
\(\{|\xi|\le M\}\).  Hence the strict high-pass \(Q_M\) gives the exact
complete assembly
\begin{equation}
 Q_M\BNS(u)
 =
 Q_M\bigl[\BNS(w,u)+\BNS(v,w)\bigr].
 \label{eq:v14-high-parent-identity}
\end{equation}
The operator \(A^{-3/4}\PP\operatorname{div}\) has order \(-1/2\).
As in Proposition~\ref{prop:v7-weak-source-budget}, the dual Sobolev
embedding \(L^{3/2}\hookrightarrow H^{-1/2}\), followed by
\(H^1\hookrightarrow L^6\), gives
\[
 \begin{aligned}
 \|A^{-3/4}Q_M\BNS(u)\|_2
 &\le
 C\bigl(
   \|w\otimes u\|_{3/2}
   +\|v\otimes w\|_{3/2}
 \bigr)\\
 &\le
 C\|w\|_2\bigl(\|u\|_6+\|v\|_6\bigr)\\
 &\le
 C_{\rm hp}\|Q_{M/2}u\|_2\|A^{1/2}u\|_2.
 \end{aligned}
\]
The spectral projector \(P_{M/2}\) is a contraction in \(H^1\), so the
constant is independent of \(M\).  Squaring and integrating proves
\eqref{eq:v14-weak-by-parent}.  Finally
\(\|Q_{M_n/2}u(t)\|_2^2\le E_*\), the cards are disjoint, and the Leray
energy identity gives
\[
 \int_{\cup_nI_n}\|A^{1/2}u(t)\|_2^2\,\dd t
 \le\frac{E_*}{2\nu}.
\]
\end{proof}

\begin{remark}[What has been gained over parent incidence]
Lemma~\ref{lem:v12-high-parent-incidence} says that a high output has a high
parent.  Theorem~\ref{thm:v14-high-parent-weak} adds the exact amplitude
ledger: the \(H^{-3/2}\) source action is bounded by high-parent kinetic
stock times the instantaneous Leray dissipation density.  The estimate is
for the complete high-output assembly and therefore retains cancellation
among its mixed and high--high parent components.
\end{remark}

% =============================================================================
\section{The two-level endpoint hierarchy}
\label{sec:v14-response-hierarchy}
% =============================================================================

Alongside \(h_M\), define
\begin{align}
 d_M&:=A^{-1/4}Q_M\BNS(u)=A^{1/2}h_M,
 \label{eq:v14-critical-and-weak-source}\\
 \rho_{M,I}
 &:=
 \int_a^b e^{-\nu A(b-t)}A^{1/2}d_M(t)\,\dd t,
 \label{eq:v14-critical-response}\\
 r_{M,I}
 &:=
 \int_a^b e^{-\nu A(b-t)}d_M(t)\,\dd t.
 \label{eq:v14-native-response}
\end{align}

\begin{theorem}[Weak source pays exactly the inverse-frequency response]
\label{thm:v14-two-level-response}
For every smooth card \(I=[a,b]\),
\begin{align}
 r_{M,I}&=A^{-1/2}\rho_{M,I},
 \label{eq:v14-response-level-identity}\\
 \boxed{
 \|r_{M,I}\|_2^2
 &\le\frac1{2\nu}\mathscr W_{M,I},}
 \label{eq:v14-native-by-weak}\\
 \boxed{
 2\nu\|A^{-1/2}\rho_{M,I}\|_2^2
 &\le\mathscr W_{M,I}
 \le C_{\rm hp}^2\mathscr P_{M,I}.}
 \label{eq:v14-response-parent-chain}
\end{align}
When \(\rho_{M,I}\ne0\), define its effective response frequency by
\begin{equation}
 \Lambda_{\rm eff}(\rho_{M,I})
 :=
 \frac{\|\rho_{M,I}\|_2}
      {\|A^{-1/2}\rho_{M,I}\|_2}.
 \label{eq:v14-effective-frequency}
\end{equation}
Then
\begin{equation}
 \Lambda_{\rm eff}(\rho_{M,I})\ge M.
 \label{eq:v14-effective-frequency-lower}
\end{equation}
If \(\|\rho_{M,I}\|_2\ge\eta R\), it follows that
\begin{equation}
 \boxed{
 \mathscr W_{M,I}
 \ge
 \frac{2\nu\eta^2R^2}
      {\Lambda_{\rm eff}(\rho_{M,I})^2},
 \qquad
 \mathscr P_{M,I}
 \ge
 \frac{2\nu\eta^2R^2}
      {C_{\rm hp}^2\Lambda_{\rm eff}(\rho_{M,I})^2}.}
 \label{eq:v14-effective-frequency-debit}
\end{equation}
\end{theorem}

\begin{proof}
Equations \eqref{eq:v14-critical-and-weak-source}--%
\eqref{eq:v14-native-response} and commutation of functional calculus with
the heat semigroup give
\[
 A^{-1/2}\rho_{M,I}
 =
 \int_a^b e^{-\nu A(b-t)}d_M(t)\,\dd t
 =r_{M,I}.
\]
In terms of \(h_M\),
\[
 r_{M,I}
 =
 \int_a^b e^{-\nu A(b-t)}A^{1/2}h_M(t)\,\dd t.
\]
For an \(A^{1/2}\)-eigenvalue \(\lambda\), the squared \(L_t^2\) norm of
the kernel \(\lambda e^{-\nu\lambda^2(b-t)}\) is at most
\((2\nu)^{-1}\).  Fourier-mode Cauchy--Schwarz and orthogonal summation prove
\eqref{eq:v14-native-by-weak}.  Combine this with
Theorem~\ref{thm:v14-high-parent-weak} to obtain
\eqref{eq:v14-response-parent-chain}.

The response lies in \(\operatorname{Ran}Q_M\), on which
\[
 \|A^{-1/2}\rho_{M,I}\|_2
 \le M^{-1}\|\rho_{M,I}\|_2.
\]
This proves \eqref{eq:v14-effective-frequency-lower}.  Finally substitute
the definition \eqref{eq:v14-effective-frequency} and the assumed response
lower bound into \eqref{eq:v14-response-parent-chain}.
\end{proof}

\begin{corollary}[Paid ancestor constraint on the V13 record cards]
\label{cor:v14-record-ancestor}
Let \(I_n^{\rm rec}\), \(R_n\), \(K_n=R_n^2/E_*\), and
\(\rho_n=\rho_{K_n,I_n^{\rm rec}}\) be the pairwise disjoint
record-generated cards from
Corollary~\ref{cor:v13-doubled-record-source-divergence}.  Put
\begin{equation}
 \Lambda_n:=
 \Lambda_{\rm eff}(\rho_n).
 \label{eq:v14-record-effective-frequency}
\end{equation}
Then
\begin{align}
 \Lambda_n&\ge K_n=\frac{R_n^2}{E_*},
 \label{eq:v14-record-frequency-floor}\\
 \boxed{
 \sum_n\frac{R_n^2}{\Lambda_n^2}
 \le C\frac{E_*^2}{\nu^2}.}
 \label{eq:v14-paid-ancestor-sum}
\end{align}
Moreover, on every card,
\begin{equation}
 \mathscr P_{K_n,I_n^{\rm rec}}
 \ge c\nu\frac{R_n^2}{\Lambda_n^2}.
 \label{eq:v14-record-parent-debit}
\end{equation}
\end{corollary}

\begin{proof}
The frequency floor is
\eqref{eq:v14-effective-frequency-lower}.  The record response lower
\eqref{eq:v13-exterior-record-response-lower} permits
\(\eta=\sqrt3-1\) in
\eqref{eq:v14-effective-frequency-debit}.  Sum its parent-ledger lower
bound and apply \eqref{eq:v14-parent-global-budget}.
\end{proof}

\begin{assessment}[The canonical cutoff makes the paid debit summable]
If \(\Lambda_n\asymp K_n\), then the summand in
\eqref{eq:v14-paid-ancestor-sum} has the scale
\[
 \frac{R_n^2}{K_n^2}
 =\frac{E_*^2}{R_n^2},
\]
which is summable on doubled records.  If
\(\Lambda_n/K_n\to\infty\), the paid ancestor charge is still smaller.
Thus the complete high-parent estimate closes the weak-source ledger but
does not close the critical record.  The missing derivative has become the
explicit scalar \(\Lambda_n\), rather than an unnamed source norm.
\end{assessment}

% =============================================================================
\section{A shellwise form and the ancestry-multiplicity loss}
\label{sec:v14-shellwise}
% =============================================================================

Let \(\Delta_q\) be the sharp shell from V6 and define
\begin{align}
 h_q&:=A^{-3/4}\Delta_q\BNS(u),
 \label{eq:v14-shell-weak-source}\\
 \rho_q&:=
 \int_a^b e^{-\nu A(b-t)}A h_q(t)\,\dd t.
 \label{eq:v14-shell-record-response}
\end{align}

\begin{proposition}[One-shell response descends to a comparable parent]
\label{prop:v14-shell-parent-descent}
For every \(q\),
\begin{equation}
 \boxed{
 \frac{\nu}{\lambda_q^2}\|\rho_q\|_2^2
 \le
 C\int_I
 \|Q_{\lambda_q/4}u(t)\|_2^2
 \|A^{1/2}u(t)\|_2^2\,\dd t.}
 \label{eq:v14-shell-parent-descent}
\end{equation}
For a finite set \(\mathcal J\) of output shells,
\begin{equation}
 \boxed{
 \nu\left\|
 A^{-1/2}\sum_{q\in\mathcal J}\rho_q
 \right\|_2^2
 \le
 C\int_I\|A^{1/2}u(t)\|_2^2
 \sum_{q\in\mathcal J}
 \|Q_{\lambda_q/4}u(t)\|_2^2\,\dd t.}
 \label{eq:v14-shell-ancestry-sum}
\end{equation}
The right-hand side counts a velocity parent once for every lower output
shell to which it can contribute; it is a logarithmic ancestry multiplicity,
not the kinetic energy itself.
\end{proposition}

\begin{proof}
Because \(\Delta_q=\Delta_qQ_{\lambda_q/2}\), apply
Theorem~\ref{thm:v14-high-parent-weak} at output threshold
\(\lambda_q/2\).  Projection to \(\Delta_q\) can only decrease the
\(H^{-3/2}\) source norm, so
\[
 \|h_q(t)\|_2
 \le
 C\|Q_{\lambda_q/4}u(t)\|_2
   \|A^{1/2}u(t)\|_2.
 \tag{3}\label{eq:v14-shell-weak-parent}
\]
On the \(q\)-shell, the \(L_t^2\) norm of the response kernel
\(\lambda^2e^{-\nu\lambda^2(b-t)}\) is at most
\(C\lambda_q/\sqrt\nu\).  Therefore
\[
 \|\rho_q\|_2^2
 \le\frac{C\lambda_q^2}{\nu}
      \int_I\|h_q(t)\|_2^2\,\dd t.
\]
Insert \eqref{eq:v14-shell-weak-parent} to prove
\eqref{eq:v14-shell-parent-descent}.

The vectors \(\rho_q\) have orthogonal Fourier supports.  On each shell,
\(\|A^{-1/2}\rho_q\|_2^2\asymp
\lambda_q^{-2}\|\rho_q\|_2^2\).  Sum
\eqref{eq:v14-shell-parent-descent} over \(\mathcal J\) and use orthogonality.
\end{proof}

\begin{remark}[Exact form of the multiplicity]
If \(\mathcal J=\{q\ge q_0\}\) is truncated and the truncation is then
removed for a smooth field, dyadic orthogonality gives
\[
 \sum_{q\ge q_0}\|Q_{\lambda_q/4}u\|_2^2
 \asymp
 \sum_{p\ge q_0-O(1)}
 \bigl(1+p-q_0\bigr)\|u_p\|_2^2.
\]
A parent at shell \(p\) can participate in high--high returns to every
lower output shell between \(q_0\) and \(p+O(1)\).  Removing this logarithmic
count requires a signed incidence or an injective ancestry selection; it
cannot be removed by orthogonality of the output shells alone.
\end{remark}

% =============================================================================
\section{A two-parameter record-growing packet saturates the frequency loss}
\label{sec:v14-two-parameter-packet}
% =============================================================================

The V10--V11 packet used a fixed amplitude parameter.  The same rescaled
proof permits an amplitude tending to zero while its critical record tends
to infinity.  This exposes the independent roles of record size and physical
frequency.

\begin{theorem}[Record growth with summable paid ancestors]
\label{thm:v14-two-parameter-packet}
Let \(R_j\to\infty\) and let \(N_j\to\infty\) be dyadic integers such that
\begin{equation}
 \frac{R_j^2}{N_j}\longrightarrow0.
 \label{eq:v14-two-parameter-separation}
\end{equation}
Put
\begin{equation}
 \beta_j:=\frac{R_j}{\sqrt{N_j}}.
 \label{eq:v14-two-parameter-amplitude}
\end{equation}
Using the compact seed \(U\) from V11 and a fixed low-frequency energy
reservoir, there are smooth Navier--Stokes solutions \(u^{(j)}\), all with
the same kinetic energy \(E_*\), and intervals
\begin{equation}
 J_j=\left[0,\frac{s_\#}{R_jN_j^2}\right]
 \label{eq:v14-two-parameter-interval}
\end{equation}
on which the following hold.  If
\[
 \mathcal R_j:=\|A^{1/4}u^{(j)}(0)\|_2,\qquad
 K_j:=\frac{\mathcal R_j^2}{E_*},
\]
then
\begin{align}
 \mathcal R_j&\asymp R_j,\qquad
 K_j\asymp R_j^2,\qquad
 \frac{K_j}{N_j}\longrightarrow0,
 \label{eq:v14-packet-record-cutoff}\\
 \|A^{1/4}u^{(j)}(\sup J_j)\|_2^2
 &\ge(1+\gamma)\mathcal R_j^2,
 \label{eq:v14-packet-record-growth}\\
 \mathscr A_j^{\rm far}
 &:=
 \int_{J_j}
 \|A^{-1/4}Q_{2K_j}\BNS(u^{(j)}(t))\|_2^2\,\dd t
 \asymp R_j^3,
 \label{eq:v14-packet-critical-action}\\
 \|\rho_j^{\rm far}\|_2&\asymp R_j,
 \qquad
 \|r_j^{\rm far}\|_2\asymp\frac{R_j}{N_j},
 \qquad
 \Lambda_{\rm eff}(\rho_j^{\rm far})\asymp N_j,
 \label{eq:v14-packet-two-responses}\\
 \mathscr W_j^{\rm far}
 &\asymp\mathscr P_{2K_j,J_j}
 \asymp\frac{R_j^3}{N_j^2},
 \label{eq:v14-packet-weak-parent-scale}\\
 \sup_{t\in J_j}\|Q_{K_j}u^{(j)}(t)\|_2^2
 &\le C\frac{R_j^2}{N_j},
 \label{eq:v14-packet-high-energy}\\
 \nu\int_{J_j}\|A^{1/2}u^{(j)}(t)\|_2^2\,\dd t
 &\asymp\nu\frac{R_j}{N_j}.
 \label{eq:v14-packet-Leray-debit}
\end{align}
Here \(\rho_j^{\rm far}\) and \(r_j^{\rm far}\) are the responses
\eqref{eq:v14-critical-response}--\eqref{eq:v14-native-response} with
\(M=2K_j\), and the constants are independent of \(j\).
In particular,
\begin{equation}
 \frac{\mathscr A_j^{\rm far}}{\nu\mathcal R_j^2}
 \asymp\frac{R_j}{\nu}\longrightarrow\infty,
 \qquad
 \frac{\Lambda_{\rm eff}(\rho_j^{\rm far})}{K_j}
 \asymp\frac{N_j}{R_j^2}\longrightarrow\infty.
 \label{eq:v14-packet-double-escape}
\end{equation}
\end{theorem}

\begin{proof}
Use the V9 packet
\[
 W_j(x)=\beta_jN_j^{3/2}U(N_j(x-x_0))
\]
and choose the coefficient of the fixed low-frequency reservoir so that
the total kinetic energy is exactly \(E_*\).  This is possible for large
\(j\), because the packet energy is
\[
 \|W_j\|_2^2\asymp\beta_j^2=\frac{R_j^2}{N_j}=o(1).
\]
The half-derivative scaling
\eqref{eq:v9-half-derivative-scaling} gives
\[
 \|A^{1/4}W_j\|_2^2\asymp\beta_j^2N_j=R_j^2,
\]
while the reservoir contributes only \(O(1)\).  This proves
\eqref{eq:v14-packet-record-cutoff}.

Under the V10 variables, the two perturbation parameters are now
\begin{equation}
 \varepsilon_j=\frac{\nu}{\beta_j\sqrt{N_j}}
 =\frac{\nu}{R_j}\longrightarrow0,
 \qquad
 \beta_j^{-1}N_j^{-3/2}
 =\frac1{R_jN_j}\longrightarrow0.
 \label{eq:v14-two-small-parameters}
\end{equation}
Thus the rescaled initial data converge to the same seed \(U\), and the
uniform \(H^m\) local control and short-time continuity proofs from
V10--V11 apply with the same \(s_\#\).  In the rescaled critical balance,
the common physical prefactor is
\[
 \beta_j^2N_j=R_j^2.
\]
The positive seed coefficient \(C_U\), together with
\(\varepsilon_j\to0\), therefore proves the fixed fractional growth
\eqref{eq:v14-packet-record-growth} exactly as in
Theorem~\ref{thm:v11-critical-record-turnover}.

The rescaled far cutoff satisfies
\[
 \frac{2K_j}{N_j}\longrightarrow0.
\]
Since \(\BNS(U)\ne0\), the projected source converges in the relevant
Sobolev norms to the full nonzero source.  Reducing \(s_\#\) once prevents
temporal cancellation.  The physical source norms and time element have
the scales
\[
 \begin{array}{c|c}
 \text{quantity}&\text{scale}\\ \hline
 \|A^{-1/4}Q_{2K_j}\BNS(u^{(j)})\|_2^2
   &\beta_j^4N_j^4\\
 \|A^{-3/4}Q_{2K_j}\BNS(u^{(j)})\|_2^2
   &\beta_j^4N_j^2\\
 \dd t&(\beta_jN_j^{5/2})^{-1}\dd s .
 \end{array}
\]
Consequently
\[
 \mathscr A_j^{\rm far}\asymp
 \beta_j^3N_j^{3/2}=R_j^3,
 \qquad
 \mathscr W_j^{\rm far}\asymp
 \beta_j^3N_j^{-1/2}=\frac{R_j^3}{N_j^2}.
\]
Applying respectively \(A^{1/2}\) and no derivative to the native response
gives
\[
 \|\rho_j^{\rm far}\|_2\asymp
 \beta_jN_j^{1/2}=R_j,
 \qquad
 \|r_j^{\rm far}\|_2\asymp
 \beta_jN_j^{-1/2}=\frac{R_j}{N_j}.
\]
This proves
\eqref{eq:v14-packet-critical-action}--%
\eqref{eq:v14-packet-two-responses}.

The cutoff \(K_j/N_j\to0\) removes the fixed reservoir and retains the
packet in the rescaled limit.  Uniform rescaled \(L^2\) and \(H^1\) bounds,
with nonzero lower limits, give
\[
 \|Q_{K_j}u^{(j)}(t)\|_2^2\asymp\beta_j^2
 =\frac{R_j^2}{N_j},
 \qquad
 \|A^{1/2}u^{(j)}(t)\|_2^2\asymp
 \beta_j^2N_j^2=R_j^2N_j.
\]
Multiplication and integration over
\(|J_j|\asymp(R_jN_j^2)^{-1}\) prove
\(\mathscr P_{2K_j,J_j}\asymp R_j^3/N_j^2\).
The same calculation with only the \(H^1\) factor proves
\eqref{eq:v14-packet-Leray-debit}.  The two escape ratios now follow from
the displayed scales.
\end{proof}

\begin{corollary}[Every paid per-card scalar can be summable while the record
grows]
\label{cor:v14-packet-summable}
Take \(R_j=2^j\) and \(N_j=R_j^p\) with any fixed integer \(p>3\).  Along
the actual record-growing packet cards of
Theorem~\ref{thm:v14-two-parameter-packet},
\begin{align}
 \sum_j\mathscr W_j^{\rm far}
 +\sum_j\mathscr P_{2K_j,J_j}
 &<\infty,
 \label{eq:v14-packet-paid-sums}\\
 \sum_j\|r_j^{\rm far}\|_2
 +\sum_j\nu\int_{J_j}\|A^{1/2}u^{(j)}\|_2^2\,\dd t
 &<\infty,
 \label{eq:v14-packet-native-Leray-sums}\\
 \sup_{t\in J_j}\|Q_{K_j}u^{(j)}(t)\|_2^2
 &\longrightarrow0,
 \label{eq:v14-packet-parent-vanishes}
\end{align}
whereas
\[
 \|\rho_j^{\rm far}\|_2\asymp R_j\longrightarrow\infty,
 \qquad
 \mathscr A_j^{\rm far}\asymp R_j^3\longrightarrow\infty.
\]
\end{corollary}

\begin{proof}
Substitute \(N_j=R_j^p\) into
\eqref{eq:v14-packet-two-responses}--%
\eqref{eq:v14-packet-Leray-debit}.  The respective decaying powers are
\(R_j^{3-2p}\), \(R_j^{1-p}\), and \(R_j^{2-p}\), all geometric.
\end{proof}

\begin{firewall}[The packet family is a local inequality test]
\label{fw:v14-packet-scope}
The solutions in Theorem~\ref{thm:v14-two-parameter-packet} are distinct
smooth trajectories; the theorem does not concatenate them into one
terminal solution or assert a full record doubling on one packet interval.
It does prove actual fractional record growth and source overdrive on each
literal card.  Therefore it rules out any universal \emph{per-card}
inequality which tries to charge a large record-visible response or the
lower action \(\mathscr A\gtrsim\nu R^2\) to a nonsummable amount of weak
source, native response, high-parent kinetic stock, or Leray dissipation
using only the displayed card scalars.
\end{firewall}

% =============================================================================
\section{Nested exterior responses have two exact telescoping defects}
\label{sec:v14-telescoping}
% =============================================================================

The V14 programme also asked whether responses at successive canonical
cutoffs telescope.  They do not, even on ideal adjacent first-hitting
corridors.

\begin{proposition}[Exact nested-cutoff sum]
\label{prop:v14-nested-response-defects}
Let \(t_n\) be the first hitting times
\eqref{eq:v13-first-hitting-times}, put
\[
 K_n=\frac{R_n^2}{E_*},\qquad
 x_n(t)=A^{1/4}Q_{K_n}u(t),\qquad
 H_n=e^{-\nu A(t_{n+1}-t_n)},
\]
and define the full-corridor response
\begin{equation}
 \rho_n^{\rm full}:=
 H_nx_n(t_n)-x_n(t_{n+1}).
 \label{eq:v14-full-corridor-response}
\end{equation}
Also put
\begin{equation}
 a_n(t):=
 A^{1/4}\mathbf1_{(K_n,K_{n+1}]}(\Lambda)u(t),
 \qquad x_n(t)=a_n(t)+x_{n+1}(t).
 \label{eq:v14-annular-junction}
\end{equation}
Then, for \(m\le N\),
\begin{equation}
 \boxed{
 \begin{aligned}
 \sum_{n=m}^N\rho_n^{\rm full}
 &=
 H_mx_m(t_m)-x_N(t_{N+1})\\
 &\quad+
 \sum_{n=m+1}^N
 \left[
  (H_n-I)x_n(t_n)-a_{n-1}(t_n)
 \right].
 \end{aligned}}
 \label{eq:v14-nested-response-sum}
\end{equation}
Thus the signed sum has both a heat defect and an annular junction defect
at every intermediate record.  On the shorter recent cards
\(I_n^{\rm rec}\), the gaps between their left endpoints and the hitting
times add a third chronology defect.
\end{proposition}

\begin{proof}
Duhamel's formula gives
\eqref{eq:v14-full-corridor-response}.  In the sum, pair the positive term
\(H_nx_n(t_n)\), for \(n\ge m+1\), with the preceding negative term
\(-x_{n-1}(t_n)\).  Equation~\eqref{eq:v14-annular-junction} gives
\[
 H_nx_n(t_n)-x_{n-1}(t_n)
 =
 (H_n-I)x_n(t_n)-a_{n-1}(t_n).
\]
The two unpaired boundary terms are exactly those in
\eqref{eq:v14-nested-response-sum}.
\end{proof}

\begin{assessment}[Why an unsigned response sum cannot be repaired by
nesting]
The annular vectors in \eqref{eq:v14-nested-response-sum} are physical
critical stocks at different record times, not source residuals, and the
heat defects depend on different corridor lengths.  Neither family is
controlled by the Leray energy at the critical scale.  Any useful
telescoping theorem must acquire cancellation of these two defects from the
dynamics; nesting of \(Q_{K_n}\) alone supplies no such cancellation.
\end{assessment}

% =============================================================================
\section{Cancellation-preserving integration by parts on the variation row}
\label{sec:v14-variation-ibp}
% =============================================================================

\begin{proposition}[The covariant variation integrates to source endpoints]
\label{prop:v14-variation-ibp}
Let \(d=d_{K,I}^{\rm par}\), \(I=[a,b]\), and
\[
 y(t):=e^{-\nu A(b-t)}d(t).
\]
Then
\begin{align}
 y'(t)&=
 e^{-\nu A(b-t)}(\partial_t+\nu A)d(t),
 \label{eq:v14-pulled-source-derivative}\\
 \boxed{
 2\operatorname{Re}\int_a^b
 \left\langle
  e^{-\nu A(b-t)}d(t),
  e^{-\nu A(b-t)}(\partial_t+\nu A)d(t)
 \right\rangle\dd t
 &=
 \|d(b)\|_2^2-\|e^{-\nu A(b-a)}d(a)\|_2^2.}
 \label{eq:v14-variation-integration-by-parts}
\end{align}
After inserting \eqref{eq:v13-exact-covariant-symbol}, the left side keeps
\(\mathcal C_k^\nu-\mathcal M_k\) assembled.  Its boundary terms are source
norms, not kinetic-energy or enstrophy terms.
\end{proposition}

\begin{proof}
The first identity is functional calculus.  The second is the fundamental
theorem of calculus applied to
\(\|y(t)\|_2^2\).  No quadratic or cubic component of \(y'\) is separated.
\end{proof}

\begin{firewall}[Integration by parts does not create the missing budget]
The exact identity \eqref{eq:v14-variation-integration-by-parts} respects the
V13 cancellation, but it returns \(d(a)\) and \(d(b)\).  Those are
instantaneous \(H^{-1/2}\) nonlinear-source norms and have no inherited
upper bound from the Leray identity.  Hence the covariant-variation branch
cannot presently be charged by integration by parts alone.
\end{firewall}

% =============================================================================
\section{Integrated V14 conclusion and V15 review mandate}
\label{sec:v14-conclusion}
% =============================================================================

\begin{theorem}[What V14 proves]
\label{thm:v14-conclusion}
Relative to V1--V13 and the frozen manuscripts, the following are proved.
\begin{enumerate}[label=\textup{(V14.\arabic*)},leftmargin=5.9em]
\item The complete \(H^{-3/2}\) high-output source is bounded by high-parent
      kinetic stock times the Leray dissipation density, producing the
      globally paid ledger \eqref{eq:v14-parent-global-budget}.
\item That ledger controls exactly \(A^{-1/2}\rho\), the native response,
      while the critical record sees \(\rho\).  Their precise derivative gap
      is the effective response frequency \(\Lambda_{\rm eff}\).
\item The V13 record cards obey the paid ancestor sum
      \eqref{eq:v14-paid-ancestor-sum}, but its canonical
      \(E_*^2/R_n^2\) scale is summable.
\item Shellwise descent introduces a logarithmic multiplicity because one
      high parent can return to every lower output shell.
\item A two-parameter family of actual record-growing packet solutions has
      \(\rho\asymp R\) and critical action \(\asymp R^3\), while weak
      source, native response, high-parent energy, parent ledger, and Leray
      debit can all be summable.
\item Nested canonical cutoffs have exact heat and annular junction defects,
      so their responses do not telescope algebraically.
\item Cancellation-preserving integration by parts on the variation branch
      returns unpaid nonlinear-source endpoint norms.
\end{enumerate}
No full Navier--Stokes stability or global-regularity theorem is obtained.
\end{theorem}

\begin{proof}
Items \textup{(V14.1)--(V14.3)} are
Theorems~\ref{thm:v14-high-parent-weak}--%
\ref{thm:v14-two-level-response} and
Corollary~\ref{cor:v14-record-ancestor}.  Item \textup{(V14.4)} is
Proposition~\ref{prop:v14-shell-parent-descent}.  Item \textup{(V14.5)} is
Theorem~\ref{thm:v14-two-parameter-packet} and
Corollary~\ref{cor:v14-packet-summable}.  The final two items are
Propositions~\ref{prop:v14-nested-response-defects} and
\ref{prop:v14-variation-ibp}.
\end{proof}

\begin{programme}[V15 acquisition order and cross-program review]
\label{prog:v15}
\begin{enumerate}[label=\textup{(AC\arabic*)},leftmargin=4.8em]
\item Begin V15 by inspecting the latest Standard-Model and Yang--Mills
      research-note files in \texttt{latex\_notes}.  Record their exact
      versions and hashes, then identify any transferable ancestry,
      cofinal-contraction, or cancellation lemma and any type mismatch.
\item Retain the V13 record-generated cards as the NS chronology.  The V14
      packet excludes a universal per-card closure from response amplitude,
      weak source, parent energy, or Leray debit alone.
\item Investigate cross-card ancestry: determine whether a remote effective
      response frequency \(\Lambda_n\gg K_n\) must have a represented
      precursor on an earlier record corridor, or can be born entirely
      inside \(I_n^{\rm rec}\).
\item If an earlier precursor exists, price the two defects in
      \eqref{eq:v14-nested-response-sum} before taking norms.  Do not assume
      that nested projections telescope.
\item If the response is born inside the recent card, use the exact
      high-parent identity \eqref{eq:v14-high-parent-identity} and the
      shellwise ledger \eqref{eq:v14-shell-parent-descent} to build a
      finite-depth source/velocity ancestry tree.  Track the logarithmic
      multiplicity explicitly.
\item Compare any proposed ancestry inequality with the two-parameter packet
      \eqref{eq:v14-packet-double-escape}.  A statement falsified by those
      actual smooth cards cannot be used on a terminal sequence without an
      additional chronological hypothesis.
\item Revisit the assembled covariant variation only if the cross-program
      review supplies an upper budget for the full
      \(\mathcal C^\nu-\mathcal M\) packet rather than for its opened
      components.
\end{enumerate}
\end{programme}

\begin{assessment}[Position after V14]
The complete weak source now has a rigorous, globally paid velocity-parent
ancestry ledger.  Its failure to control the critical record is no longer a
qualitative derivative complaint: it is exactly the effective frequency
of the record-visible response.  Concentrated actual dynamics can send that
frequency arbitrarily far above the canonical cutoff while every presently
paid per-card scalar tends to zero.  Therefore the remaining possible gain
must be chronological and cross-card---a restriction on how such remote
responses are born or inherited---rather than another instantaneous product
estimate.
\end{assessment}

% =============================================================================
\section*{Version V14 history}
\addcontentsline{toc}{section}{Version V14 history}
% =============================================================================

Proved a uniform high-output weak-source estimate by high-parent kinetic
stock times Leray dissipation and obtained a globally paid parent ledger.
Identified the exact two-level response hierarchy and isolated the effective
response frequency as the remaining derivative.  Derived shellwise descent
and its unavoidable logarithmic high--high-return multiplicity.  Extended
the V10--V11 construction to a two-parameter actual record-growing packet
whose critical response and source action diverge while every weak,
kinetic, native-response, and Leray card scalar is summable.  Computed the
heat and annular defects preventing nested-cutoff telescoping and showed that
cancellation-preserving variation integration by parts returns unpaid source
endpoints.  No global regularity claim was made.
% Future versions are appended immediately above the sole terminator.

% =============================================================================
\section{V15 cross-program audit: order, sign, and literal incidence}
\label{sec:v15-cross-audit}
% =============================================================================

This iteration begins with the cross-program review required at every fifth
NS version.  The review was made before choosing the V15 estimates.  The two
companion files then present in \texttt{latex\_notes} were
\begin{align*}
 &\texttt{Renewal\_SM\_Einstein\_Continuum\_Research\_Notes\_V41.tex},\\
 &\qquad\texttt{SHA256}=
 \texttt{BEC2FD855D847657135F23EB784F7E5C4D86D3AE0412D617049A6686C10BD83B},\\
 &\texttt{Renewal\_Yang\_Mills\_Mass\_Gap\_Research\_Notes\_V22.tex},\\
 &\qquad\texttt{SHA256}=
 \texttt{2607A33D4FBBB86C10E9E483E5DF96AEED79A24FC04F92E0ACE5CB1DFA700E22}.
\end{align*}
The names and hashes freeze exactly which companion states were read; later
versions of those programmes are not silently imported into V15.

The SM V41 fixed-scale lift and zero-principal-floor theorem make one lesson
directly type-compatible with the NS response problem: an estimate repaired
by a negative-order operator at each fixed spectral scale does not thereby
become a uniform principal-order estimate.  Its odd-port obstruction is not
type-compatible with the NS divergence-free velocity system and is not
imported as an NS rank theorem.

The YM V22 posterior-excess calculation makes a second methodological lesson
type-compatible: preserve a signed assembled quantity through the exact
identity, because bounding its pieces or taking positive parts first can lose
the cancellation.  Its probabilistic posterior Hessian and local square-root
questions have different types from the deterministic NS source and are not
imported.  V15 therefore turns upstream from another unsigned parent norm to
two signed NS identities: the endpoint Gram of the complete exterior Duhamel
response and the complete exterior source--carrier work.

\begin{firewall}[Scope of the companion review]
\label{fw:v15-cross-audit-scope}
No SM or YM conclusion is being used as a premise about a Navier--Stokes
solution.  The transferable statements below are reproved from the Stokes
functional calculus and the NS equation.  The review changes the order in
which NS quantities are assembled; it supplies no regularity theorem.
\end{firewall}

% =============================================================================
\section{The complete exterior endpoint Gram has a sharp signed dichotomy}
\label{sec:v15-endpoint-gram}
% =============================================================================

Retain the hypotheses and notation of
Theorem~\ref{thm:v13-record-exterior-incidence}.  Thus
\(I^{\rm rec}=[s,b]\), \(K=R^2/E_*\), and
\begin{equation}
 d_K^{\rm ext}=A^{-1/4}Q_K\BNS(u),
 \qquad
 \rho:=\int_s^b e^{-\nu A(b-t)}A^{1/2}d_K^{\rm ext}(t)\,\dd t.
 \label{eq:v15-record-response-definition}
\end{equation}
Put
\begin{equation}
 x:=A^{1/4}Q_Ku(b),
 \qquad
 y:=e^{-\nu A(b-s)}A^{1/4}Q_Ku(s).
 \label{eq:v15-endpoint-vectors}
\end{equation}
The complete projected Duhamel identity is simply
\begin{equation}
 x=y-\rho.
 \label{eq:v15-endpoint-Duhamel}
\end{equation}

\begin{theorem}[Signed endpoint-Gram dichotomy on every recent record card]
\label{thm:v15-endpoint-gram-dichotomy}
Define the response self-term and the adverse precursor interference by
\begin{equation}
 \mathsf S:=\|\rho\|_2^2,
 \qquad
 \mathsf C:=\bigl(-2\operatorname{Re}\langle y,\rho\rangle\bigr)_+.
 \label{eq:v15-self-cross-definition}
\end{equation}
Then
\begin{equation}
 \boxed{\quad \mathsf S+\mathsf C\ge2R^2.\quad}
 \label{eq:v15-Gram-dichotomy-sum}
\end{equation}
Consequently at least one of the following two complete-source alternatives
holds:
\begin{align}
 \mathsf S\ge R^2
 &\quad\Longrightarrow\quad
 \int_s^b\|A^{-1/4}Q_K\BNS(u(t))\|_2^2\,\dd t
 \ge2\nu R^2,
 \label{eq:v15-self-source-branch}\\
 \mathsf C\ge R^2
 &\quad\Longrightarrow\quad
 \int_s^b\|Q_K\BNS(u(t))\|_2^2\,\dd t
 \ge\frac{\nu R^4}{2E_*}.
 \label{eq:v15-cross-source-branch}
\end{align}
The alternatives need not be exclusive.  In particular, destructive
interference with the heat-propagated precursor is possible only at the
price of a full, unweighted exterior nonlinear-source action of quartic
record size.
\end{theorem}

\begin{proof}
Equations~\eqref{eq:v13-final-exterior-stock} and
\eqref{eq:v13-old-exterior-bound} give
\begin{equation}
 \|x\|_2^2\ge3R^2,
 \qquad
 \|y\|_2^2\le R^2.
 \label{eq:v15-endpoint-size-bounds}
\end{equation}
Expanding the signed identity \(x=y-\rho\) before taking any absolute value
gives
\begin{equation}
 \|x\|_2^2-\|y\|_2^2
 =\|\rho\|_2^2-2\operatorname{Re}\langle y,\rho\rangle
 \le \mathsf S+\mathsf C.
 \label{eq:v15-exact-endpoint-Gram}
\end{equation}
The left side is at least \(2R^2\), proving
\eqref{eq:v15-Gram-dichotomy-sum}.

If \(\mathsf S\ge R^2\), the sharp response estimate
\eqref{eq:v13-record-response-upper} gives
\[
 \int_s^b\|d_K^{\rm ext}(t)\|_2^2\,\dd t
 \ge2\nu\|\rho\|_2^2\ge2\nu R^2.
\]
This is \eqref{eq:v15-self-source-branch}.

Suppose instead that \(\mathsf C\ge R^2\).  Functional calculus and the
energy cap give
\begin{equation}
 \|A^{-1/4}y\|_2
 =\|e^{-\nu A(b-s)}Q_Ku(s)\|_2
 \le\sqrt{E_*}.
 \label{eq:v15-precursor-L2-bound}
\end{equation}
Since the positive part is then attained by the displayed signed cross term,
\begin{equation}
 \frac{R^2}{2}
 \le-\operatorname{Re}\langle y,\rho\rangle
 \le\|A^{-1/4}y\|_2\,\|A^{1/4}\rho\|_2,
\end{equation}
and hence
\begin{equation}
 \|A^{1/4}\rho\|_2
 \ge\frac{R^2}{2\sqrt{E_*}}.
 \label{eq:v15-strong-response-lower}
\end{equation}
Writing \(b_K=Q_K\BNS(u)\), one has
\begin{equation}
 A^{1/4}\rho
 =\int_s^b e^{-\nu A(b-t)}A^{1/2}b_K(t)\,\dd t.
 \label{eq:v15-strong-response-source}
\end{equation}
The same one-mode heat-kernel calculation as in
Lemma~\ref{lem:v13-record-response} yields
\begin{equation}
 \|A^{1/4}\rho\|_2^2
 \le\frac1{2\nu}\int_s^b\|b_K(t)\|_2^2\,\dd t.
 \label{eq:v15-strong-response-upper}
\end{equation}
Combining the last two displays proves
\eqref{eq:v15-cross-source-branch}.
\end{proof}

\begin{assessment}[Gain and limitation of the endpoint dichotomy]
The interference branch is much more expensive than the universal V13
critical-source debit: at canonical scale it forces
\(\nu R^4/E_*\), not merely \(\nu R^2\).  However, the forced norm is the
unweighted \(L^2\) norm of \(Q_K\BNS(u)\), for which the Leray identity gives
no upper budget.  The signed Gram identity therefore identifies a genuinely
strong branch without closing it.
\end{assessment}

% =============================================================================
\section{A first doubling forces positive complete source--carrier work}
\label{sec:v15-source-carrier-work}
% =============================================================================

The endpoint Gram can be differentiated without opening the nonlinear
source into parent classes.  This produces a signed quantity which records
whether a nonlinear output is acting on a velocity already present at the
same output frequency.

\begin{definition}[Complete exterior source--carrier work]
\label{def:v15-exterior-work}
For a fixed cutoff \(K\), set
\begin{equation}
 \mathsf P_K^{\rm ext}(t):=
 -\operatorname{Re}\left\langle
 Q_K\BNS(u(t)),A^{1/2}Q_Ku(t)
 \right\rangle.
 \label{eq:v15-exterior-work-definition}
\end{equation}
No raw, mixed, return, or outer--outer component is separated in this
definition.
\end{definition}

\begin{theorem}[Positive work certificate on a complete first-doubling
corridor]
\label{thm:v15-work-certificate}
Under \eqref{eq:v13-first-doubling}, with \(K=R^2/E_*\),
\begin{align}
 \frac12\frac{\dd}{\dd t}
 \|A^{1/4}Q_Ku(t)\|_2^2
 +\nu\|A^{3/4}Q_Ku(t)\|_2^2
 &=\mathsf P_K^{\rm ext}(t),
 \label{eq:v15-exterior-critical-balance}\\
 \boxed{\quad
 \int_a^b\mathsf P_K^{\rm ext}(t)\,\dd t
 \ge R^2,
 \qquad
 \int_a^b\bigl(\mathsf P_K^{\rm ext}(t)\bigr)_+\,\dd t
 \ge R^2.\quad}
 \label{eq:v15-positive-work-certificate}
\end{align}
Thus every first critical-record doubling has a positive signed work
certificate involving the complete exterior source and the contemporaneous
exterior velocity carrier.
\end{theorem}

\begin{proof}
Apply \(Q_K\) to the Navier--Stokes equation and take the real
\(L^2\) pairing with \(A^{1/2}Q_Ku\).  Self-adjointness and commutation of
the Stokes spectral multipliers give
\eqref{eq:v15-exterior-critical-balance}.  At the two corridor endpoints,
\[
 \|A^{1/4}Q_Ku(a)\|_2^2\le R^2,
 \qquad
 \|A^{1/4}Q_Ku(b)\|_2^2\ge3R^2
\]
by the record value at \(a\) and
\eqref{eq:v13-final-exterior-stock}.  Integration of
\eqref{eq:v15-exterior-critical-balance} therefore gives
\[
 \int_a^b\mathsf P_K^{\rm ext}\,\dd t
 =\frac12\left(
 \|A^{1/4}Q_Ku(b)\|_2^2-
 \|A^{1/4}Q_Ku(a)\|_2^2\right)
 +\nu\int_a^b\|A^{3/4}Q_Ku\|_2^2\,\dd t
 \ge R^2.
\]
The positive-part statement follows because
\(f\le f_+\) pointwise.
\end{proof}

\begin{proposition}[Exact Fourier carrier rule]
\label{prop:v15-Fourier-carrier-rule}
On the periodic domain, the work density is
\begin{equation}
 \mathsf P_K^{\rm ext}(t)
 =-\sum_{|k|>K}|k|\,
 \operatorname{Re}\left(
 \widehat{\BNS(u)}_k(t)\cdot
 \overline{\widehat u_k(t)}
 \right).
 \label{eq:v15-Fourier-carrier-work}
\end{equation}
In particular, a nonlinear output at \(k\) with
\(\widehat u_k(t)=0\) makes zero instantaneous contribution to critical
record growth, irrespective of the size of
\(\widehat{\BNS(u)}_k(t)\).  A productive output must act on a
contemporaneous carrier at the same mode.
\end{proposition}

\begin{proof}
Both \(Q_K\) and \(A^{1/2}\) are diagonal in the Fourier basis, and the
symbol of \(A^{1/2}\) is \(|k|\).  Parseval's identity gives
\eqref{eq:v15-Fourier-carrier-work}; the last assertion is immediate term
by term.
\end{proof}

\begin{proposition}[Cancellation-preserving annular work decomposition]
\label{prop:v15-annular-work-decomposition}
Let
\begin{equation}
 \Pi_j:=\mathbf1_{(2^jK,2^{j+1}K]}(A^{1/2}),
 \qquad j=0,1,2,\ldots,
 \label{eq:v15-annular-projectors}
\end{equation}
and define
\begin{equation}
 \mathsf P_{K,j}(t):=
 -\operatorname{Re}\langle
 \Pi_j\BNS(u(t)),A^{1/2}\Pi_ju(t)\rangle.
 \label{eq:v15-annular-work}
\end{equation}
For every smooth corridor,
\begin{equation}
 \mathsf P_K^{\rm ext}(t)=\sum_{j\ge0}\mathsf P_{K,j}(t),
 \qquad
 \int_a^b\sum_{j\ge0}\mathsf P_{K,j}(t)\,\dd t
 \ge R^2
 \label{eq:v15-signed-annular-work-sum}
\end{equation}
under the hypotheses of Theorem~\ref{thm:v15-work-certificate}.  The sum in
the second formula must remain signed until after assembly.
\end{proposition}

\begin{proof}
The projectors in \eqref{eq:v15-annular-projectors} are mutually orthogonal,
commute with \(A\), and sum strongly to \(Q_K\).  Smoothness makes the
Fourier sum absolutely convergent, so orthogonality gives the first identity
and permits integration.  The second conclusion is then
\eqref{eq:v15-positive-work-certificate}.
\end{proof}

\begin{firewall}[The work certificate is not yet a paid estimate]
\label{fw:v15-work-not-paid}
The two direct duality bounds are
\begin{align}
 |\mathsf P_K^{\rm ext}|
 &\le
 \|A^{-1/4}Q_K\BNS(u)\|_2
 \|A^{3/4}Q_Ku\|_2,
 \label{eq:v15-critical-work-duality}\\
 |\mathsf P_K^{\rm ext}|
 &\le
 \|A^{-3/4}Q_K\BNS(u)\|_2
 \|A^{5/4}Q_Ku\|_2.
 \label{eq:v15-weak-work-duality}
\end{align}
The first returns the already forced but globally unpaid critical source
action; the second asks for an \(H^{5/2}\) velocity factor absent from the
Leray budget.  Moreover,
\(int(\sum_j\mathsf P_{K,j})_+\) may be much smaller than
\(\sum_j\int(\mathsf P_{K,j})_+\).  Taking shellwise positive parts creates
an inflated unsigned quantity and is not licensed by
\eqref{eq:v15-signed-annular-work-sum}.
\end{firewall}

% =============================================================================
\section{The lower-order response cannot have a uniform principal floor}
\label{sec:v15-principal-floor}
% =============================================================================

The next theorem is the NS functional-calculus version of the order issue
highlighted by the SM review.  It makes the V14 effective-frequency gap
sharp without importing any SM operator.

\begin{theorem}[Sharp negative-order no-floor theorem]
\label{thm:v15-negative-order-no-floor}
Let \(\alpha>0\) and \(M>0\).  On the unbounded exterior Stokes subspace,
\begin{equation}
 \inf_{0\ne z\in\operatorname{Ran}Q_M\cap C^\infty}
 \frac{\|A^{-\alpha}z\|_2}{\|z\|_2}=0.
 \label{eq:v15-no-principal-floor}
\end{equation}
Equivalently, no finite constant \(C(M)\) can make
\begin{equation}
 \|z\|_2\le C(M)\|A^{-\alpha}z\|_2
 \label{eq:v15-false-reverse-order}
\end{equation}
hold for all \(z\in\operatorname{Ran}Q_M\).

On a nonempty finite spectral band
\(\Pi_{(M,L]}=\mathbf1_{(M,L]}(A^{1/2})\), the optimal constant in
\eqref{eq:v15-false-reverse-order} is
\begin{equation}
 C_{M,L}^{\rm opt}
 =\max\left\{\lambda^{2\alpha}:
 \lambda\in\sigma(A^{1/2})\cap(M,L]\right\}.
 \label{eq:v15-band-optimal-constant}
\end{equation}
It therefore diverges with the largest occupied frequency of expanding
bands.
\end{theorem}

\begin{proof}
If \(A^{1/2}e_\lambda=\lambda e_\lambda\) and
\(\|e_\lambda\|_2=1\), then
\[
 \|A^{-\alpha}e_\lambda\|_2=\lambda^{-2\alpha}.
\]
The Stokes frequencies in \(\operatorname{Ran}Q_M\) are unbounded, so a
sequence \(\lambda\to\infty\) proves
\eqref{eq:v15-no-principal-floor} and rules out
\eqref{eq:v15-false-reverse-order}.  On a finite band, diagonal expansion
gives
\[
 \|z\|_2^2
 \le\left(\max_{\lambda\in\sigma(A^{1/2})\cap(M,L]}
 \lambda^{4\alpha}\right)\|A^{-\alpha}z\|_2^2.
\]
An eigenvector at the maximizing frequency gives equality, proving
\eqref{eq:v15-band-optimal-constant}.
\end{proof}

\begin{corollary}[Exact implication for the V14 two-level response]
\label{cor:v15-native-response-no-floor}
For the V14 pair
\(r=A^{-1/2}\rho\), no estimate of the form
\begin{equation}
 \|\rho\|_2\le C(M)\|r\|_2
 \label{eq:v15-native-to-critical-false}
\end{equation}
can hold uniformly over responses supported in \(\operatorname{Ran}Q_M\).
After imposing an upper frequency edge \(L\), its sharp constant is the
largest occupied frequency, exactly the scale measured by
\(\Lambda_{\rm eff}(\rho)\).  Thus the V14 native-response budget is a
genuine lower-order estimate, not a fixed-scale approximation to a uniform
critical response bound.
\end{corollary}

\begin{proof}
Apply Theorem~\ref{thm:v15-negative-order-no-floor} with
\(\alpha=1/2\).  On an eigenmode of frequency \(\lambda\),
\(\|\rho\|_2/\|A^{-1/2}\rho\|_2=\lambda\).
\end{proof}

% =============================================================================
\section{Signed endpoint cancellation does not control source variation}
\label{sec:v15-variation-countermodel}
% =============================================================================

The YM review warns against replacing a signed identity by separate positive
estimates.  For the NS covariant-variation row, one must also check the
converse danger: exact signed cancellation alone need not bound the positive
variation.  The following finite-dimensional Stokes construction settles
that abstract question.

\begin{proposition}[Arbitrarily large covariant variation with fixed action
and endpoints]
\label{prop:v15-variation-loop}
Let \(I=[a,b]\), \(\tau=b-a>0\), and choose two orthonormal real Stokes
eigenfields \(e_1,e_2\) with the same eigenvalue
\(Ae_i=\lambda_0^2e_i\).  For every integer \(m\ge1\), define
\begin{align}
 y_m(t)&:=
 \cos\left(\frac{2\pi m(t-a)}{\tau}\right)e_1
 +\sin\left(\frac{2\pi m(t-a)}{\tau}\right)e_2,
 \label{eq:v15-loop-y}\\
 d_m(t)&:=e^{\nu A(b-t)}y_m(t).
 \label{eq:v15-loop-d}
\end{align}
Then
\begin{align}
 e^{-\nu A(b-t)}d_m(t)&=y_m(t),
 \label{eq:v15-loop-pullback}\\
 \int_a^b\|e^{-\nu A(b-t)}d_m(t)\|_2^2\,\dd t&=\tau,
 \label{eq:v15-loop-action}\\
 \|d_m(b)\|_2^2-
 \|e^{-\nu A\tau}d_m(a)\|_2^2&=0,
 \label{eq:v15-loop-signed-endpoint}\\
 \int_a^b
 \|e^{-\nu A(b-t)}(\partial_t+\nu A)d_m(t)\|_2^2\,\dd t
 &=\frac{4\pi^2m^2}{\tau}\longrightarrow\infty.
 \label{eq:v15-loop-variation}
\end{align}
Both pulled endpoint norms, both unpulled endpoint norms, and the pulled
source action are independent of \(m\).
\end{proposition}

\begin{proof}
Equation~\eqref{eq:v15-loop-pullback} follows from the definition.  The
orthonormality of \(e_1,e_2\) gives \(\|y_m(t)\|_2=1\), proving
\eqref{eq:v15-loop-action}.  Since \(y_m(a)=y_m(b)=e_1\), the two pulled
endpoint norms agree and \eqref{eq:v15-loop-signed-endpoint} follows.
Also \(d_m(b)=e_1\) and
\(d_m(a)=e^{\nu\lambda_0^2\tau}e_1\), so the unpulled endpoint norms do not
depend on \(m\).

Differentiating the pullback identity, or calculating directly, gives
\[
 e^{-\nu A(b-t)}(\partial_t+\nu A)d_m(t)=y_m'(t).
\]
Its norm is the constant \(2\pi m/\tau\).  Integration over a time interval
of length \(\tau\) proves \eqref{eq:v15-loop-variation}.
\end{proof}

\begin{firewall}[What the variation loop rules out]
\label{fw:v15-variation-loop-scope}
The fields \(d_m\) are abstract smooth source paths; they are not asserted
to equal \(A^{-1/4}Q_K\BNS(u)\) along Navier--Stokes trajectories.  The
proposition proves that the Hilbert-space identity
\eqref{eq:v14-variation-integration-by-parts}, even supplemented by fixed
source endpoint norms and fixed pulled action, cannot by itself upper-bound
the positive covariant variation.  Any successful variation estimate must
use an additional NS dynamical restriction while keeping the assembled
\(\mathcal C^\nu-\mathcal M\) cancellation.
\end{firewall}

% =============================================================================
\section{V15 synthesis and the next ancestry step}
\label{sec:v15-synthesis}
% =============================================================================

\begin{theorem}[What V15 proves]
\label{thm:v15-conclusion}
Relative to V1--V14 and the frozen cross-program files, the following are
proved.
\begin{enumerate}[label=\textup{(V15.\arabic*)},leftmargin=6.2em]
\item Every recent first-doubling card has an exact signed endpoint-Gram
      dichotomy.  Either response self-action alone forces a
      \(2\nu R^2\) complete critical-source debit, or destructive precursor
      interference forces a complete \(L^2\) nonlinear-source debit of at
      least \(\nu R^4/(2E_*)\).
\item Every complete first-doubling corridor carries at least \(R^2\) of
      positive signed exterior source--carrier work.
\item In Fourier variables, only a nonlinear output with a contemporaneous
      velocity carrier at the same mode can contribute instantaneously to
      that critical-record work.  The exact annular work decomposition must
      be summed before positive parts are taken.
\item No negative-order response estimate on an unbounded exterior spectral
      sector has a uniform principal-order floor.  The sharp finite-band
      recovery constant is the largest occupied frequency.
\item Exact covariant-variation integration by parts, fixed pulled action,
      and fixed endpoint data do not abstractly control positive variation;
      a two-dimensional Stokes loop makes the variation arbitrarily large.
\end{enumerate}
No full Navier--Stokes stability or global-regularity theorem is obtained.
\end{theorem}

\begin{proof}
Item \textup{(V15.1)} is
Theorem~\ref{thm:v15-endpoint-gram-dichotomy}.  Items
\textup{(V15.2)--(V15.3)} are
Theorem~\ref{thm:v15-work-certificate} and
Propositions~\ref{prop:v15-Fourier-carrier-rule}--%
\ref{prop:v15-annular-work-decomposition}.  Item \textup{(V15.4)} is
Theorem~\ref{thm:v15-negative-order-no-floor} and
Corollary~\ref{cor:v15-native-response-no-floor}.  Item
\textup{(V15.5)} is Proposition~\ref{prop:v15-variation-loop}.
\end{proof}

\begin{programme}[V16 acquisition order]
\label{prog:v16}
\begin{enumerate}[label=\textup{(AC\arabic*)},leftmargin=4.8em]
\item Use the signed measure
      \(\mathsf P_{K,j}(t)\,\dd t\), not its componentwise absolute value,
      as the root of the record ancestry tree.
\item On each work-positive annulus, trace the carrier backward.  Separate
      the inherited alternative, in which the carrier survives from the
      left endpoint, from the born-inside alternative, in which its Duhamel
      source response must first create a nontrivial carrier.
\item For inherited carriers, insert the heat and annular junction defects
      of \eqref{eq:v14-nested-response-sum} before taking norms.  Determine
      whether the signed work pairing cancels either defect across adjacent
      first-hitting corridors.
\item For carriers born inside a corridor, combine the V7 shell-upcrossing
      debit with the V14 high-parent identity.  Record every parent-to-output
      incidence and its logarithmic high--high-return multiplicity.
\item Seek a stopping-time or finite-depth statement forcing one ancestral
      branch to recur on infinitely many doubled records.  Test every
      per-card inequality against the V14 two-parameter packet before using
      it in a terminal sum.
\item Do not attempt to recover critical response from native response by a
      cutoff-dependent reverse inequality: Theorem
      \ref{thm:v15-negative-order-no-floor} proves that such a constant is
      precisely the escaping frequency.
\item Reopen covariant variation only after deriving an NS-specific
      restriction unavailable to the loops of
      Proposition~\ref{prop:v15-variation-loop}.  Keep
      \(\mathcal C^\nu-\mathcal M\) assembled until the signed identity has
      been used.
\item Repeat the required companion-file audit at V20 and record the exact
      SM/YM versions and hashes then present.
\end{enumerate}
\end{programme}

\begin{assessment}[Position after V15]
The response-frequency escape remains the central unpaid derivative, but
the chronology is now sharper.  A first doubling cannot be described merely
as a large exterior source norm: it contains positive complete
source--carrier work.  At every productive frequency there is therefore a
carrier whose ancestry can be queried.  The remaining task is to prove that
an infinite doubled-record sequence cannot continually obtain such carriers
through inherited heat/annular defects or fresh births while paying only the
summable ledgers already exposed.  V15 neither proves that assertion nor
assumes it.
\end{assessment}

% =============================================================================
\section*{Version V15 history}
\addcontentsline{toc}{section}{Version V15 history}
% =============================================================================

Performed the required fifth-version audit against exact SM V41 and YM V22
files.  Imported only type-compatible order and cancellation lessons, then
reproved all operative statements inside the NS setting.  Established the
signed endpoint-Gram dichotomy and its stronger destructive-interference
source debit; derived a positive complete source--carrier work certificate
and its exact Fourier and annular forms; proved the sharp negative-order
no-floor theorem; and constructed a Stokes-loop countermodel showing that
signed endpoint integration by parts alone cannot control positive source
variation.  The full stability/global-regularity goal remains open.

% =============================================================================
\section{V16 starting point: retain the carrier and move all ancestry to one
terminal scale}
\label{sec:v16-start}
% =============================================================================

V15 produced a signed complete source--carrier work certificate, but did not
use the order-zero duality between that source and the Leray-paid
\(H^1\) carrier.  It also left the V14 heat and annular junction defects in a
variable-cutoff response sum.  V16 first repairs those two omissions.  It
then asks the carrier question at one fixed terminal cutoff, where all
earlier annuli are below the cofinal projection and disappear exactly.

Throughout V16 write
\begin{equation}
 \Lambda:=A^{1/2},
 \qquad
 \mathcal R(t):=\|\Lambda^{1/2}u(t)\|_2.
 \label{eq:v16-Lambda-record}
\end{equation}
No companion-file review is due in this version; the next mandatory SM/YM
review remains V20.

% =============================================================================
\section{Every record doubling forces quartic unweighted source action}
\label{sec:v16-quartic-source}
% =============================================================================

For an interval \(I=[a,b]\), define the unweighted complete source actions
\begin{equation}
 \mathscr S_0(I):=
 \int_a^b\|\BNS(u(t))\|_2^2\,\dd t,
 \qquad
 \mathscr S_{0,K}^{\rm ext}(I):=
 \int_a^b\|Q_K\BNS(u(t))\|_2^2\,\dd t.
 \label{eq:v16-strong-source-actions}
\end{equation}

\begin{theorem}[Quartic complete-source debit on every first doubling]
\label{thm:v16-quartic-source-debit}
Assume the first-doubling hypotheses
\eqref{eq:v13-first-doubling} and the energy cap
\eqref{eq:v13-energy-cap}, and put \(K=R^2/E_*\).  Then
\begin{align}
 \boxed{\quad
 \mathscr S_0([a,b])
 \ge\frac{9}{2}\frac{\nu R^4}{E_*},\quad}
 \label{eq:v16-full-quartic-source}\\
 \boxed{\quad
 \mathscr S_{0,K}^{\rm ext}([a,b])
 \ge2\frac{\nu R^4}{E_*}.\quad}
 \label{eq:v16-exterior-quartic-source}
\end{align}
Both bounds use the fully assembled nonlinear source.  In particular,
\eqref{eq:v16-exterior-quartic-source} strengthens the conditional
interference branch \eqref{eq:v15-cross-source-branch}: quartic unweighted
exterior action is forced on every first-doubling corridor, independently
of the endpoint-Gram branch.
\end{theorem}

\begin{proof}
The full critical balance \eqref{eq:v5-critical-record-flux}, written before
the radial-flux representation, is
\begin{equation}
 \frac12\frac{\dd}{\dd t}\mathcal R(t)^2
 +\nu\|\Lambda^{3/2}u(t)\|_2^2
 =-\operatorname{Re}\langle\BNS(u(t)),\Lambda u(t)\rangle.
 \label{eq:v16-full-critical-work-balance}
\end{equation}
Since \(\mathcal R(a)=R\) and \(\mathcal R(b)=2R\), integration gives
\begin{equation}
 -\int_a^b\operatorname{Re}\langle\BNS(u),\Lambda u\rangle\,\dd t
 =\frac32R^2+
 \nu\int_a^b\|\Lambda^{3/2}u\|_2^2\,\dd t
 \ge\frac32R^2.
 \label{eq:v16-full-work-lower}
\end{equation}
Time--space Cauchy--Schwarz and the Leray energy inequality give
\begin{align*}
 \frac32R^2
 &\le \mathscr S_0([a,b])^{1/2}
       \left(\int_a^b\|\Lambda u(t)\|_2^2\,\dd t\right)^{1/2},\\
 \int_a^b\|\Lambda u(t)\|_2^2\,\dd t
 &\le\frac{E_*}{2\nu}.
\end{align*}
Squaring proves \eqref{eq:v16-full-quartic-source}.

For the exterior statement, use the exact work certificate
\eqref{eq:v15-positive-work-certificate} and the order-zero pairing
\begin{equation}
 \mathsf P_K^{\rm ext}(t)
 =-\operatorname{Re}\langle
 Q_K\BNS(u(t)),\Lambda Q_Ku(t)\rangle.
 \label{eq:v16-exterior-order-zero-pairing}
\end{equation}
Thus
\[
 R^2
 \le \mathscr S_{0,K}^{\rm ext}([a,b])^{1/2}
       \left(\int_a^b\|\Lambda Q_Ku(t)\|_2^2\,\dd t\right)^{1/2}
 \le \mathscr S_{0,K}^{\rm ext}([a,b])^{1/2}
       \left(\frac{E_*}{2\nu}\right)^{1/2}.
\]
Squaring proves \eqref{eq:v16-exterior-quartic-source}.
\end{proof}

\begin{corollary}[Strong source and acceleration consequences of an
unbounded record]
\label{cor:v16-strong-source-divergence}
On the first-hitting chronology \eqref{eq:v13-first-hitting-times}, an
unbounded critical record implies
\begin{align}
 \sum_n\int_{t_n}^{t_{n+1}}\|\BNS(u(t))\|_2^2\,\dd t
 &\ge\frac{9\nu}{2E_*}\sum_nR_n^4=\infty,
 \label{eq:v16-strong-source-divergence}\\
 \int_{t_n}^{t_{n+1}}
 \left(\|\partial_tu(t)\|_2^2+\nu^2\|Au(t)\|_2^2\right)\,\dd t
 &\ge\frac{9}{4}\frac{\nu R_n^4}{E_*}.
 \label{eq:v16-acceleration-diffusion-debit}
\end{align}
Consequently at least one of the two terms on the left of
\eqref{eq:v16-acceleration-diffusion-debit} is at least
\(9\nu R_n^4/(8E_*)\).
\end{corollary}

\begin{proof}
The doubling corridors have disjoint interiors, so summing
\eqref{eq:v16-full-quartic-source} proves the first assertion.  The equation
gives
\(\BNS(u)=-\partial_tu-\nu Au\), whence
\[
 \|\BNS(u)\|_2^2
 \le2\|\partial_tu\|_2^2+2\nu^2\|Au\|_2^2.
\]
Combine this with \eqref{eq:v16-full-quartic-source}.  The last statement is
the elementary alternative for two nonnegative summands.
\end{proof}

\begin{firewall}[The stronger norm is still unpaid]
\label{fw:v16-strong-source-unpaid}
The Leray budget controls the carrier factor in
\eqref{eq:v16-exterior-order-zero-pairing}; it does not upper-bound
\(\mathscr S_0\).  Equations
\eqref{eq:v16-full-quartic-source}--%
\eqref{eq:v16-acceleration-diffusion-debit} are strong necessary conditions
for a singular record cascade, not contradictions.  Replacing the source by
the standard product estimate would ask for supercritical velocity norms and
return to the V3 continuation boundary.
\end{firewall}

% =============================================================================
\section{Exact heat-inherited and source-born carrier decomposition}
\label{sec:v16-carrier-split}
% =============================================================================

The next result performs the V16 carrier split without assigning parent
components to the nonlinear source.  Let \(\Pi\) be any orthogonal Stokes
spectral projector and let \(I=[a,b]\).  Define
\begin{align}
 u_\Pi(t)&:=\Pi u(t),
 &h_\Pi(t)&:=e^{-\nu A(t-a)}u_\Pi(a),
 &v_\Pi(t)&:=u_\Pi(t)-h_\Pi(t),
 \label{eq:v16-carrier-components}\\
 b_\Pi(t)&:=\Pi\BNS(u(t)),
 &d_\Pi(t)&:=A^{-1/4}b_\Pi(t).
 \label{eq:v16-projector-source}
\end{align}
Thus \(h_\Pi\) is the heat-inherited carrier and \(v_\Pi\) is the complete
source-born carrier.

\begin{theorem}[Exact carrier-work split and its two source prices]
\label{thm:v16-carrier-work-split}
Put
\begin{align}
 \mathsf P_\Pi^{\rm inh}(t)
 &:=-\operatorname{Re}\langle b_\Pi(t),\Lambda h_\Pi(t)\rangle,
 &\mathsf W_\Pi^{\rm inh}
 &:=\int_a^b\mathsf P_\Pi^{\rm inh}(t)\,\dd t,
 \label{eq:v16-inherited-work}\\
 \mathsf P_\Pi^{\rm born}(t)
 &:=-\operatorname{Re}\langle b_\Pi(t),\Lambda v_\Pi(t)\rangle,
 &\mathsf W_\Pi^{\rm born}
 &:=\int_a^b\mathsf P_\Pi^{\rm born}(t)\,\dd t,
 \label{eq:v16-born-work}\\
 \mathscr A_\Pi(I)&:=\int_a^b\|d_\Pi(t)\|_2^2\,\dd t,
 &Y_\Pi(a)&:=\|\Lambda^{1/2}u_\Pi(a)\|_2^2.
 \label{eq:v16-carrier-ledgers}
\end{align}
Then the complete projector work splits exactly as
\begin{equation}
 -\int_a^b\operatorname{Re}\langle
 b_\Pi(t),\Lambda u_\Pi(t)\rangle\,\dd t
 =\mathsf W_\Pi^{\rm inh}+\mathsf W_\Pi^{\rm born}.
 \label{eq:v16-complete-work-split}
\end{equation}
The source-born work is nonnegative and has the exact identity
\begin{equation}
 \boxed{\quad
 \mathsf W_\Pi^{\rm born}
 =\frac12\|\Lambda^{1/2}v_\Pi(b)\|_2^2
 +\nu\int_a^b\|\Lambda^{3/2}v_\Pi(t)\|_2^2\,\dd t
 \ge0.\quad}
 \label{eq:v16-born-work-identity}
\end{equation}
Moreover,
\begin{align}
 \mathscr A_\Pi(I)&\ge\nu\mathsf W_\Pi^{\rm born},
 \label{eq:v16-born-work-price}\\
 |\mathsf W_\Pi^{\rm inh}|
 &\le
 \left(\frac{Y_\Pi(a)}{2\nu}\right)^{1/2}
 \mathscr A_\Pi(I)^{1/2}.
 \label{eq:v16-inherited-work-price}
\end{align}
Finally, inherited work retains its signed endpoint and dissipative cross
terms:
\begin{equation}
 \mathsf W_\Pi^{\rm inh}
 =\operatorname{Re}\langle
 \Lambda^{1/2}v_\Pi(b),\Lambda^{1/2}h_\Pi(b)\rangle
 +2\nu\int_a^b\operatorname{Re}\langle
 \Lambda^{3/2}v_\Pi,\Lambda^{3/2}h_\Pi\rangle\,\dd t.
 \label{eq:v16-inherited-signed-identity}
\end{equation}
\end{theorem}

\begin{proof}
The projected equation and definitions give
\begin{equation}
 (\partial_t+\nu A)h_\Pi=0,
 \qquad
 (\partial_t+\nu A)v_\Pi=-b_\Pi,
 \qquad
 v_\Pi(a)=0.
 \label{eq:v16-carrier-equations}
\end{equation}
Linearity in the carrier gives \eqref{eq:v16-complete-work-split}.  Pair the
second equation in \eqref{eq:v16-carrier-equations} with
\(\Lambda v_\Pi\) and integrate to obtain
\eqref{eq:v16-born-work-identity}.

Let
\(D_v=\int_a^b\|\Lambda^{3/2}v_\Pi\|_2^2\,\dd t\).  Since
\[
 \mathsf W_\Pi^{\rm born}
 =-\int_a^b\operatorname{Re}\langle
 d_\Pi,\Lambda^{3/2}v_\Pi\rangle\,\dd t
 \le\mathscr A_\Pi(I)^{1/2}D_v^{1/2}
\]
and \(\mathsf W_\Pi^{\rm born}\ge\nu D_v\), either the work is zero or
division yields
\(\mathsf W_\Pi^{\rm born}\le\mathscr A_\Pi(I)/\nu\).  This proves
\eqref{eq:v16-born-work-price}.

The heat energy identity is
\begin{equation}
 2\nu\int_a^b\|\Lambda^{3/2}h_\Pi(t)\|_2^2\,\dd t
 =Y_\Pi(a)-\|\Lambda^{1/2}h_\Pi(b)\|_2^2
 \le Y_\Pi(a).
 \label{eq:v16-inherited-heat-budget}
\end{equation}
As
\(\mathsf P_\Pi^{\rm inh}
 =-\operatorname{Re}\langle d_\Pi,\Lambda^{3/2}h_\Pi\rangle\),
Cauchy--Schwarz and \eqref{eq:v16-inherited-heat-budget} prove
\eqref{eq:v16-inherited-work-price}.

For the final identity, differentiate
\(\operatorname{Re}\langle
\Lambda^{1/2}v_\Pi,\Lambda^{1/2}h_\Pi\rangle\), use
\eqref{eq:v16-carrier-equations}, and obtain
\[
 \frac{\dd}{\dd t}
 \operatorname{Re}\langle
 \Lambda^{1/2}v_\Pi,\Lambda^{1/2}h_\Pi\rangle
 =\mathsf P_\Pi^{\rm inh}
 -2\nu\operatorname{Re}\langle
 \Lambda^{3/2}v_\Pi,\Lambda^{3/2}h_\Pi\rangle.
\]
The cross term vanishes at \(a\), proving
\eqref{eq:v16-inherited-signed-identity}.
\end{proof}

\begin{corollary}[Quantitative born-or-inherited alternative]
\label{cor:v16-born-inherited-alternative}
If the complete projector work on \(I\) is at least \(W>0\), then at least
one of the following holds:
\begin{align}
 \mathsf W_\Pi^{\rm born}\ge\frac W2
 &\quad\Longrightarrow\quad
 \mathscr A_\Pi(I)\ge\frac{\nu W}{2},
 \label{eq:v16-born-alternative}\\
 \mathsf W_\Pi^{\rm inh}\ge\frac W2
 &\quad\Longrightarrow\quad
 \mathscr A_\Pi(I)\ge
 \frac{\nu W^2}{2Y_\Pi(a)}.
 \label{eq:v16-inherited-alternative}
\end{align}
If \(Y_\Pi(a)=0\), the inherited alternative is impossible.
\end{corollary}

\begin{proof}
Equation~\eqref{eq:v16-complete-work-split} and nonnegativity of born work
give the alternative.  Use \eqref{eq:v16-born-work-price} in the first
case.  Squaring \eqref{eq:v16-inherited-work-price} and rearranging gives
the second.  If \(Y_\Pi(a)=0\),
\eqref{eq:v16-inherited-heat-budget} gives \(h_\Pi\equiv0\).
\end{proof}

\begin{assessment}[What the carrier split does and does not buy]
For \(\Pi=Q_K\), Theorem~\ref{thm:v15-work-certificate} gives \(W=R^2\)
and \(Y_\Pi(a)\le R^2\).  The two alternatives therefore each return a
complete critical-source debit of order \(\nu R^2\).  This is an exact
ancestry classification, but not a stronger budget than V13.  Its value is
that the inherited branch is now represented by the signed expression
\eqref{eq:v16-inherited-signed-identity}, while the born branch is a literal
zero-initial-data Duhamel carrier.
\end{assessment}

% =============================================================================
\section{A common-terminal cofinal projection telescopes without either V14
defect}
\label{sec:v16-cofinal-telescope}
% =============================================================================

Let \(t_n,R_n,K_n\) be the first-hitting chronology
\eqref{eq:v13-first-hitting-times}, so \(K_n=R_n^2/E_*\).  On the corridor
\([t_\ell,t_{\ell+1}]\), let \(\rho_\ell^{\rm full}\) be the V14 response
\eqref{eq:v14-full-corridor-response}.

\begin{theorem}[Exact common-terminal cofinal response telescope]
\label{thm:v16-cofinal-telescope}
Fix \(N\ge2\), set
\begin{equation}
 M_N:=K_{N-1}=\frac{R_{N-1}^2}{E_*},
 \qquad
 X_N(t):=\Lambda^{1/2}Q_{M_N}u(t),
 \label{eq:v16-cofinal-cutoff}
\end{equation}
and, for each integer \(0\le\ell\le N-1\), define the common-terminal response
\begin{equation}
 \mathfrak R_{\ell}^{N}:=
 e^{-\nu A(t_N-t_{\ell+1})}Q_{M_N}\rho_\ell^{\rm full}.
 \label{eq:v16-common-terminal-response}
\end{equation}
Then, for every integer \(0\le m\le N-1\),
\begin{equation}
 \boxed{\quad
 X_N(t_N)
 =e^{-\nu A(t_N-t_m)}X_N(t_m)
 -\sum_{\ell=m}^{N-1}\mathfrak R_{\ell}^{N}.
 \quad}
 \label{eq:v16-cofinal-telescope}
\end{equation}
There are no heat or annular junction defects in
\eqref{eq:v16-cofinal-telescope}.  Moreover,
\begin{equation}
 \left\|\sum_{\ell=m}^{N-1}\mathfrak R_{\ell}^{N}\right\|_2
 \ge\sqrt3R_{N-1}-R_m.
 \label{eq:v16-cofinal-response-lower}
\end{equation}
In particular, if \(m\le N-2\), the right side is at least
\((\sqrt3-1/2)R_{N-1}\).
\end{theorem}

\begin{proof}
Because \(M_N\ge K_\ell\), the spectral projectors satisfy
\(Q_{M_N}Q_{K_\ell}=Q_{M_N}\).  Duhamel's formula on the \(\ell\)-th
corridor, propagated from \(t_{\ell+1}\) to the common time \(t_N\), gives
\begin{equation}
 \mathfrak R_\ell^N
 =\int_{t_\ell}^{t_{\ell+1}}
 e^{-\nu A(t_N-t)}\Lambda^{1/2}Q_{M_N}\BNS(u(t))\,\dd t.
 \label{eq:v16-common-response-integral}
\end{equation}
Summing the adjacent time integrals and applying Duhamel once on
\([t_m,t_N]\) proves \eqref{eq:v16-cofinal-telescope}.

This also explains the disappearance of the V14 defects.  Every heat factor
has been propagated to the same terminal time.  At each intermediate cutoff
switch \(\ell\le N-2\), the annular stock
\(a_\ell(t_{\ell+1})\) from
\eqref{eq:v14-annular-junction} has frequencies at most
\(K_{\ell+1}\le M_N\), and hence is killed by \(Q_{M_N}\).

At \(t_N\), the energy cap gives
\[
 \|\Lambda^{1/2}P_{M_N}u(t_N)\|_2^2
 \le M_NE_*=R_{N-1}^2.
\]
Since \(\mathcal R(t_N)=R_N=2R_{N-1}\), orthogonality yields
\begin{equation}
 \|X_N(t_N)\|_2^2\ge3R_{N-1}^2.
 \label{eq:v16-terminal-cofinal-stock}
\end{equation}
Heat contraction and the first-hitting property give
\(\|e^{-\nu A(t_N-t_m)}X_N(t_m)\|_2\le R_m\).  Apply the reverse triangle
inequality to \eqref{eq:v16-cofinal-telescope} to prove
\eqref{eq:v16-cofinal-response-lower}.  Finally
\(R_m\le R_{N-2}=R_{N-1}/2\) when \(m\le N-2\).
\end{proof}

\begin{firewall}[Why this does not contradict the V14 defect theorem]
\label{fw:v16-cofinal-vs-variable}
V14 sums responses at their own changing cutoffs and endpoint times; its
heat and annular defects are exact.  V16 first projects every term to one
later cutoff and propagates every term to one later time.  The price is that
only the cofinal tail above \(M_N\) survives.  Thus
\eqref{eq:v16-cofinal-telescope} is a different, ancestry-compatible
identity, not a cancellation of the V14 defects in their original space.
\end{firewall}

% =============================================================================
\section{Remote cofinal ancestry either dies by heat and parent scarcity or
forces time compression}
\label{sec:v16-remote-ancestry}
% =============================================================================

The cofinal telescope permits the V14 paid weak-source ledger to be used on
old generations, because their output is propagated through a positive heat
age before the terminal observation.

\begin{lemma}[Heat-aged weak source gains the missing derivative]
\label{lem:v16-aged-weak-source}
Let \(J=[\alpha,\beta]\), let \(b>\beta\), put
\(\delta=b-\beta\), and let \(h\in L^2(J;L^2)\).  Then
\begin{equation}
 \boxed{\quad
 \left\|\int_\alpha^\beta
 e^{-\nu A(b-t)}Ah(t)\,\dd t\right\|_2^2
 \le\frac1{4e\nu^2\delta}
 \int_\alpha^\beta\|h(t)\|_2^2\,\dd t.
 \quad}
 \label{eq:v16-aged-weak-bound}
\end{equation}
If \(h\) is supported in \(\operatorname{Ran}Q_M\), the coefficient can be
sharpened to
\begin{equation}
 \frac1{2\nu}
 \sup_{\lambda\ge M}
 \lambda^2e^{-2\nu\lambda^2\delta}.
 \label{eq:v16-aged-sharp-coefficient}
\end{equation}
\end{lemma}

\begin{proof}
For an \(A^{1/2}\)-eigenfrequency \(\lambda\), time Cauchy--Schwarz gives
the kernel factor
\[
 \int_\alpha^\beta
 \lambda^4e^{-2\nu\lambda^2(b-t)}\,\dd t
 \le\int_\delta^\infty
 \lambda^4e^{-2\nu\lambda^2s}\,\dd s
 =\frac{\lambda^2e^{-2\nu\lambda^2\delta}}{2\nu}.
\]
The maximum of \(x e^{-2\nu\delta x}\) over \(x>0\) is
\((2e\nu\delta)^{-1}\).  Sum the orthogonal modes to prove
\eqref{eq:v16-aged-weak-bound}.  Restricting the supremum to
\(\lambda\ge M\) gives \eqref{eq:v16-aged-sharp-coefficient}.
\end{proof}

\begin{theorem}[Cofinal recent-birth or multi-generation compression]
\label{thm:v16-cofinal-localization}
Fix \(N\ge3\) and \(2\le d\le N\).  Put
\begin{equation}
 R:=R_{N-1},
 \qquad M:=K_{N-1}=\frac{R^2}{E_*},
 \qquad
 \delta_{N,d}:=t_N-t_{N-d},
 \label{eq:v16-localization-scales}
\end{equation}
and define
\begin{equation}
 \Delta_{N,d}:=
 \frac{4C_{\rm hp}^2E_*^2}{e\nu^3R^2}\,4^{-(d-1)},
 \label{eq:v16-compression-threshold}
\end{equation}
where \(C_{\rm hp}\) is the constant in
\eqref{eq:v14-pointwise-high-parent}.  Then at least one of the following
holds:
\begin{align}
 \boxed{\quad \delta_{N,d}<\Delta_{N,d},\quad}
 &\tag{C}\label{eq:v16-compression-branch}\\
 \boxed{\quad
 \left\|\int_{t_{N-d}}^{t_N}
 e^{-\nu A(t_N-t)}\Lambda^{1/2}Q_M\BNS(u(t))\,\dd t
 \right\|_2
 \ge\left(\sqrt3-\frac12\right)R,
 \quad}
 &\tag{B}\label{eq:v16-recent-birth-branch}
\end{align}
and branch \textup{(B)} forces
\begin{equation}
 \boxed{\quad
 \int_{t_{N-d}}^{t_N}
 \|A^{-1/4}Q_M\BNS(u(t))\|_2^2\,\dd t
 \ge
 2\left(\sqrt3-\frac12\right)^2\nu R^2.
 \quad}
 \label{eq:v16-recent-birth-action}
\end{equation}
Thus, unless the last \(d\) record generations occupy the extremely short
time in \eqref{eq:v16-compression-branch}, the cofinal carrier observed at
\(t_N\) must be created at critical strength inside those last \(d\)
generations.
\end{theorem}

\begin{proof}
Split the fixed-cutoff response on \([t_0,t_N]\) into
\begin{align*}
 \mathfrak R_{N,d}^{\rm far}
 &:=\int_{t_0}^{t_{N-d}}
 e^{-\nu A(t_N-t)}\Lambda^{1/2}Q_M\BNS(u(t))\,\dd t,\\
 \mathfrak R_{N,d}^{\rm near}
 &:=\int_{t_{N-d}}^{t_N}
 e^{-\nu A(t_N-t)}\Lambda^{1/2}Q_M\BNS(u(t))\,\dd t.
\end{align*}
On the far interval, set
\(h_M=A^{-3/4}Q_M\BNS(u)\).  Then the far response integrand is
\(Ah_M\).  Before the first hitting time \(t_{N-d}\),
\(\mathcal R(t)\le R_{N-d}\), so spectral localization gives
\begin{equation}
 \|Q_{M/2}u(t)\|_2^2
 \le\frac2M\mathcal R(t)^2
 \le2E_*\left(\frac{R_{N-d}}{R_{N-1}}\right)^2
 =2E_*4^{-(d-1)}.
 \label{eq:v16-remote-parent-scarcity}
\end{equation}
The V14 high-parent estimate and the Leray budget therefore imply
\begin{equation}
 \int_{t_0}^{t_{N-d}}\|h_M(t)\|_2^2\,\dd t
 \le
 \frac{C_{\rm hp}^2E_*^2}{\nu}\,4^{-(d-1)}.
 \label{eq:v16-remote-weak-action}
\end{equation}
Lemma~\ref{lem:v16-aged-weak-source}, with heat age
\(\delta_{N,d}\), yields
\begin{equation}
 \|\mathfrak R_{N,d}^{\rm far}\|_2^2
 \le
 \frac{C_{\rm hp}^2E_*^2}{4e\nu^3\delta_{N,d}}
 \,4^{-(d-1)}.
 \label{eq:v16-far-response-bound}
\end{equation}
If branch \textup{(C)} fails, the definition of \(\Delta_{N,d}\) makes the
right side at most \(R^2/16\), so
\begin{equation}
 \|\mathfrak R_{N,d}^{\rm far}\|_2\le\frac R4.
 \label{eq:v16-far-response-small}
\end{equation}

Apply Theorem~\ref{thm:v16-cofinal-telescope} with \(m=0\).  Since
\(N\ge3\), \(R_0\le R/4\), and hence
\[
 \|\mathfrak R_{N,d}^{\rm far}+
       \mathfrak R_{N,d}^{\rm near}\|_2
 \ge\sqrt3R-R_0
 \ge\left(\sqrt3-\frac14\right)R.
\]
Together with \eqref{eq:v16-far-response-small}, the reverse triangle
inequality proves \eqref{eq:v16-recent-birth-branch}.  Finally apply the
sharp response bound \eqref{eq:v13-record-response-upper} to the fixed
projector \(Q_M\) on \([t_{N-d},t_N]\); this gives
\eqref{eq:v16-recent-birth-action}.
\end{proof}

\begin{assessment}[The exact remaining compression loophole]
\label{ass:v16-compression-loophole}
For fixed \(d\), the threshold
\(\Delta_{N,d}\) is proportional to \(R_{N-1}^{-2}\), hence is summable
over doubled records.  Infinitely many compression branches are therefore
compatible with a finite terminal time.  Infinitely many recent-birth
branches force divergent critical-source action on a disjoint subfamily of
\(d\)-generation blocks, but that action is already known to be globally
unpaid.  Theorem~\ref{thm:v16-cofinal-localization} removes remote weak
ancestry; it does not yet close either remaining branch.
\end{assessment}

% =============================================================================
\section{Critical transport commutation forces a high advector on a fast
record corridor}
\label{sec:v16-transport-commutator}
% =============================================================================

The high-parent identity says that a high output has a high parent.  For
signed critical work, transport skew-adjointness specifies which parent can
be low without paying the leading output frequency.

\begin{proposition}[Exact half-derivative transport commutator]
\label{prop:v16-critical-commutator}
Let \(v,u\) be smooth divergence-free fields and write
\(\BNS(v,u)=\PP(v\cdot\nabla u)\).  Then
\begin{equation}
 -\operatorname{Re}\langle\BNS(v,u),\Lambda u\rangle
 =-\operatorname{Re}\langle
 [\Lambda^{1/2},v\cdot\nabla]u,\Lambda^{1/2}u\rangle.
 \label{eq:v16-critical-commutator}
\end{equation}
There is an absolute torus constant \(C_{\rm com}\) such that, for every
\(L>0\),
\begin{equation}
 \left|\operatorname{Re}\langle
 \BNS(P_Lu,u),\Lambda u\rangle\right|
 \le C_{\rm com}L^{5/2}\|P_Lu\|_2\mathcal R(t)^2.
 \label{eq:v16-low-advector-work-bound}
\end{equation}
\end{proposition}

\begin{proof}
The Leray projector may be removed inside the pairing with the
divergence-free field \(\Lambda u\).  Since \(v\) is divergence-free,
\[
 \operatorname{Re}\langle
 v\cdot\nabla\Lambda^{1/2}u,\Lambda^{1/2}u\rangle=0.
\]
Move one half derivative to each factor and subtract this zero term to get
\eqref{eq:v16-critical-commutator}.

For a Fourier interaction \(p+q=k\), the commutator multiplier obeys
\begin{equation}
 \bigl||k|^{1/2}-|q|^{1/2}\bigr|\,|q|
 \le |p|\,|q|^{1/2}.
 \label{eq:v16-symbol-difference}
\end{equation}
Indeed, use \(\bigl||k|-|q|\bigr|\le|p|\) and rationalize the square-root
difference.  Young's convolution inequality therefore gives
\[
 \|[\Lambda^{1/2},v\cdot\nabla]u\|_2
 \le C\left(\sum_p|p|\,|\widehat v_p|\right)
       \|\Lambda^{1/2}u\|_2.
\]
For \(v=P_Lu\), lattice Cauchy--Schwarz gives
\[
 \sum_{0<|p|\le L}|p|\,|\widehat u_p|
 \le\left(\sum_{0<|p|\le L}|p|^2\right)^{1/2}\|P_Lu\|_2
 \le C L^{5/2}\|P_Lu\|_2.
\]
Insert this in \eqref{eq:v16-critical-commutator} and use
\(\|\Lambda^{1/2}u\|_2=\mathcal R(t)\).
\end{proof}

\begin{theorem}[A time-adapted high-advector work certificate]
\label{thm:v16-high-advector-certificate}
Under the first-doubling hypotheses on \(I=[a,b]\), put \(\tau=b-a\) and
\begin{equation}
 L_I:=\left(8C_{\rm com}\sqrt{E_*}\,\tau\right)^{-2/5}.
 \label{eq:v16-work-advector-edge}
\end{equation}
Then the complete high-advector component satisfies
\begin{equation}
 \boxed{\quad
 -\int_a^b\operatorname{Re}\langle
 \BNS(Q_{L_I}u(t),u(t)),\Lambda u(t)\rangle\,\dd t
 \ge R^2.
 \quad}
 \label{eq:v16-high-advector-work}
\end{equation}
Thus a first doubling occurring in time \(\tau\) has a signed productive
advector parent above the scale
\(L_I\asymp E_*^{-1/5}\tau^{-2/5}\).
\end{theorem}

\begin{proof}
The full work lower bound \eqref{eq:v16-full-work-lower} is
\(3R^2/2\).  On the first-hitting corridor,
\(\mathcal R(t)<2R\) for \(t<b\), and the energy cap gives
\(\|P_{L_I}u(t)\|_2\le\sqrt{E_*}\).  Hence
Proposition~\ref{prop:v16-critical-commutator} and
\eqref{eq:v16-work-advector-edge} give
\[
 \int_a^b
 \left|\operatorname{Re}\langle
 \BNS(P_{L_I}u,u),\Lambda u\rangle\right|\,\dd t
 \le4C_{\rm com}L_I^{5/2}\sqrt{E_*}R^2\tau
 =\frac{R^2}{2}.
\]
Use bilinearity in the advector slot,
\(u=P_{L_I}u+Q_{L_I}u\), and subtract the low-advector contribution from
\eqref{eq:v16-full-work-lower}.
\end{proof}

\begin{corollary}[Literal inherited-or-born high advector]
\label{cor:v16-high-advector-ancestry}
On the corridor of Theorem~\ref{thm:v16-high-advector-certificate}, write
\begin{align}
 w^{\rm inh}(t)&:=e^{-\nu A(t-a)}Q_{L_I}u(a),
 \label{eq:v16-high-advector-inherited}\\
 w^{\rm born}(t)&:=Q_{L_I}u(t)-w^{\rm inh}(t)
 =-\int_a^t e^{-\nu A(t-s)}Q_{L_I}\BNS(u(s))\,\dd s.
 \label{eq:v16-high-advector-born}
\end{align}
Then at least one of
\begin{align}
 -\int_a^b\operatorname{Re}\langle
 \BNS(w^{\rm inh}(t),u(t)),\Lambda u(t)\rangle\,\dd t
 &\ge\frac{R^2}{2},
 \label{eq:v16-inherited-advector-work}\\
 -\int_a^b\operatorname{Re}\langle
 \BNS(w^{\rm born}(t),u(t)),\Lambda u(t)\rangle\,\dd t
 &\ge\frac{R^2}{2}
 \label{eq:v16-born-advector-work}
\end{align}
holds.  The second branch contains a literal source-created high-frequency
advector, not merely a high output.
\end{corollary}

\begin{proof}
Duhamel's formula gives \eqref{eq:v16-high-advector-born}.  Bilinearity in
the first slot splits the left side of
\eqref{eq:v16-high-advector-work} into the two displayed terms.  If both
were smaller than \(R^2/2\), their sum would be smaller than \(R^2\).
\end{proof}

\begin{firewall}[High work gives no universal high-parent amplitude floor]
\label{fw:v16-high-work-no-amplitude-floor}
The V14 two-parameter packet tests this certificate sharply.  On its
interval \(J_j\), the full critical work is at least
\(\gamma\mathcal R_j^2/2\asymp R_j^2\).  Applying the preceding low-advector
estimate with a sufficiently small constant depending only on \(\gamma\)
gives the edge
\begin{equation}
 L_j\asymp
 \left(\sqrt{E_*}|J_j|\right)^{-2/5}
 \asymp(R_jN_j^2)^{2/5},
 \qquad
 \frac{L_j}{N_j}\asymp
 \left(\frac{R_j^2}{N_j}\right)^{1/5}\longrightarrow0.
 \label{eq:v16-packet-advector-edge}
\end{equation}
The same rescaled convergence used in
Theorem~\ref{thm:v14-two-parameter-packet} then gives
\begin{equation}
 -\int_{J_j}\operatorname{Re}\langle
 \BNS(Q_{L_j}u^{(j)},u^{(j)}),\Lambda u^{(j)}\rangle\,\dd t
 \ge c_\gamma R_j^2,
 \qquad
 \sup_{t\in J_j}\|Q_{L_j}u^{(j)}(t)\|_2^2
 \le C\frac{R_j^2}{N_j}\longrightarrow0.
 \label{eq:v16-packet-high-work-small-parent}
\end{equation}
Its unweighted source action has the scaling
\begin{equation}
 \int_{J_j}\|\BNS(u^{(j)}(t))\|_2^2\,\dd t
 \asymp\beta_j^3N_j^{5/2}=R_j^3N_j.
 \label{eq:v16-packet-strong-source}
\end{equation}
To verify this last scale, the packet self-interaction has
\(\|\BNS(W_j)\|_2^2\asymp\beta_j^4N_j^5\), while
\(\dd t=(\beta_jN_j^{5/2})^{-1}\dd s\).  The fixed reservoir
terms are lower order, and the nonzero seed source gives the matching lower
bound after reducing \(s_\#\) as in V14.  Thus
\eqref{eq:v16-packet-strong-source} is consistent with
\eqref{eq:v16-full-quartic-source} because \(N_j/R_j\to\infty\).

Therefore neither large signed high-advector work nor the quartic strong
source debit supplies a positive lower bound for high-parent kinetic
amplitude depending only on \(R,E_*,\nu\).  The packet trajectories are
distinct and do not form one terminal cascade, but they rule out that
universal per-card shortcut.
\end{firewall}

% =============================================================================
\section{Integrated V16 conclusion and V17 acquisition order}
\label{sec:v16-conclusion}
% =============================================================================

\begin{theorem}[What V16 proves]
\label{thm:v16-conclusion}
Relative to V1--V15, the following are proved.
\begin{enumerate}[label=\textup{(V16.\arabic*)},leftmargin=6.2em]
\item Every first critical-record doubling forces complete unweighted
      nonlinear-source action at least
      \(9\nu R^4/(2E_*)\); its exterior part above the canonical cutoff
      forces at least \(2\nu R^4/E_*\).
\item Every spectral output carrier has an exact heat-inherited plus
      source-born decomposition.  Born work is nonnegative and paid by
      critical-source action; inherited work has an exact signed endpoint
      and dissipative-cross formula.
\item Projecting all ancestral responses to the penultimate terminal cutoff
      and propagating them to one terminal time removes both V14 telescoping
      defects exactly and forces a cofinal response of order \(R_{N-1}\).
\item Paid weak-source ancestry from generations older than \(d\) is small
      after its positive heat age.  Hence the terminal cofinal response is
      born in the last \(d\) generations unless those generations are
      compressed below the explicit threshold
      \eqref{eq:v16-compression-threshold}.
\item Critical transport work is an exact half-derivative commutator.  On a
      corridor of duration \(\tau\), low advectors below
      \(E_*^{-1/5}\tau^{-2/5}\) cannot supply the full record work, so a
      signed productive high advector is forced.
\item Actual V14 packets show that the productive high advector may have
      vanishing kinetic amplitude while its work is of record size; no
      universal per-card amplitude floor follows.
\end{enumerate}
No full Navier--Stokes stability or global-regularity theorem is obtained.
\end{theorem}

\begin{proof}
Item \textup{(V16.1)} is
Theorem~\ref{thm:v16-quartic-source-debit}.  Item \textup{(V16.2)} is
Theorem~\ref{thm:v16-carrier-work-split} and
Corollary~\ref{cor:v16-born-inherited-alternative}.  Item
\textup{(V16.3)} is Theorem~\ref{thm:v16-cofinal-telescope}.  Item
\textup{(V16.4)} is Lemma~\ref{lem:v16-aged-weak-source} and
Theorem~\ref{thm:v16-cofinal-localization}.  Items
\textup{(V16.5)--(V16.6)} are
Proposition~\ref{prop:v16-critical-commutator},
Theorem~\ref{thm:v16-high-advector-certificate}, and
Firewall~\ref{fw:v16-high-work-no-amplitude-floor}.
\end{proof}

\begin{programme}[V17 acquisition order]
\label{prog:v17}
\begin{enumerate}[label=\textup{(AC\arabic*)},leftmargin=4.8em]
\item Optimize the generation depth \(d=d_N\) in
      Theorem~\ref{thm:v16-cofinal-localization} against the actual record
      gaps \(t_N-t_{N-d}\).  Extract disjoint near blocks before summing any
      critical-source debit.
\item On the compression branch, derive consequences from the equation for
      the dimensionless product
      \(\nu K_{N-1}^2(t_N-t_{N-d})\).  Do not assume a norm-only lower
      lifespan in the critical space.
\item On the recent-birth branch, apply the exact carrier split at the fixed
      cofinal projector \(Q_{K_{N-1}}\).  Determine whether its inherited
      signed cross term cancels after common-terminal propagation.
\item Refine the transport commutator into signed triad frequency-spread
      classes.  Preserve the difference multiplier
      \eqref{eq:v16-symbol-difference}; do not replace it by the output
      frequency before the low/high parent split.
\item For a source-born high advector, iterate Duhamel only until the first
      genuine shell upcrossing, then charge that event by
      \eqref{eq:v13-explicit-upcrossing-debit}.  Track reuse of one parent
      across lower return shells explicitly.
\item Test every proposed compression or ancestry estimate against
      \eqref{eq:v16-packet-advector-edge}--%
      \eqref{eq:v16-packet-strong-source}.  The packet excludes an amplitude
      floor but not a cross-record recurrence obstruction.
\item Revisit the strong-source acceleration/diffusion alternative only if
      one branch can be tied to a globally finite quantity.  Its present
      quartic lower bound alone is not a budget.
\end{enumerate}
\end{programme}

\begin{assessment}[Position after V16]
The variable-cutoff algebra is no longer the ancestry obstruction: a fixed
cofinal projection gives an exact telescope.  Paid weak ancestry from remote
generations is quantitatively extinguished by the combination of parent
scarcity and heat age.  What remains is sharply bifurcated.  Either the
cofinal carrier is generated in a bounded recent record block by unpaid
critical source, or that block collapses on an explicitly summable time
scale.  The transport commutator additionally forces a high advector, but
actual packets show why its amplitude cannot be charged per card.  V17 must
therefore attack recurrence across records or the compression geometry; it
cannot close the proof by another instantaneous norm estimate.
\end{assessment}

% =============================================================================
\section*{Version V16 history}
\addcontentsline{toc}{section}{Version V16 history}
% =============================================================================

Strengthened the V15 work certificate to a universal quartic unweighted
source-action debit using the Leray-paid carrier.  Proved the exact
heat-inherited/source-born carrier-work split.  Replaced the defective
variable-cutoff response sum by a common-terminal cofinal telescope, then
combined heat-aged smoothing with the V14 high-parent ledger to obtain a
recent-birth versus multi-generation time-compression theorem.  Derived the
exact critical transport commutator and a time-adapted high-advector work
certificate, and used the actual V14 packet family to rule out a universal
high-parent amplitude floor.  The full stability/global-regularity goal
remains open.

% =============================================================================
\section{V17 redundancy audit and upstream correction}
\label{sec:v17-redundancy-audit}
% =============================================================================

Before iterating the V16 compression depth, one must compare its conclusion
with the earlier record-card theorem.  That comparison changes the direction
of V17.

\begin{proposition}[The qualitative V16 recent-birth mechanism is already
present in the V13 one-card incidence theorem]
\label{prop:v17-v16-birth-redundancy}
On the \(N\)-th first-doubling corridor
\([t_{N-1},t_N]\), with
\(M=K_{N-1}\) and \(R=R_{N-1}\), V13 provides a recent subcard
\(I_N^{\rm rec}=[s_N,t_N]\subset[t_{N-1},t_N]\) such that
\begin{equation}
 \left\|\int_{s_N}^{t_N}
 e^{-\nu A(t_N-t)}\Lambda^{1/2}Q_M\BNS(u(t))\,\dd t
 \right\|_2
 \ge(\sqrt3-1)R.
 \label{eq:v17-v13-one-card-birth}
\end{equation}
Consequently, the qualitative assertion that an order-\(R\) cofinal
response is created somewhere in the last \(d\ge2\) generations is not a
new branch: a temporally sharper, though smaller-constant, response is
already forced inside the final generation.  The stronger conditional
constant in \eqref{eq:v16-recent-birth-branch} remains a V16 consequence.
\end{proposition}

\begin{proof}
In the notation of V13,
\(d_M^{\rm ext}=A^{-1/4}Q_M\BNS(u)\), so
\(A^{1/2}d_M^{\rm ext}=\Lambda^{1/2}Q_M\BNS(u)\).
Equation~\eqref{eq:v17-v13-one-card-birth} is therefore exactly
\eqref{eq:v13-exterior-record-response-lower}.
\end{proof}

\begin{assessment}[What remains genuinely new from the V16 ancestry step]
\label{ass:v17-v16-novel-remainder}
The common-terminal telescope
\eqref{eq:v16-cofinal-telescope}, the heat-aged weak-source estimate
\eqref{eq:v16-aged-weak-bound}, and the remote parent-scarcity bound
\eqref{eq:v16-remote-parent-scarcity} remain valid and useful.  What cannot
be counted as a new closure mechanism is qualitative recent birth, because
V13 already supplies an order-\(R\) response with \(d=1\) and a shorter
card.  V16 still gives a sharper conditional constant and a quantitative
remote-ancestry estimate.  V17 therefore does not optimize that
binary alternative.  It moves upstream to two estimates that were missed:
the exterior spectral support of the V13 response yields stronger source
norms on the same recent card, and the record cap improves the V16
low-advector Bernstein exponent.
\end{assessment}

% =============================================================================
\section{The recent-card response forces an entire source-regularity ladder}
\label{sec:v17-source-ladder}
% =============================================================================

Retain one first-doubling corridor and its V13 recent card
\(I^{\rm rec}=[s,b]\), with
\begin{equation}
 K=\frac{R^2}{E_*},
 \qquad
 \rho=
 \int_s^b e^{-\nu A(b-t)}
 \Lambda^{1/2}Q_K\BNS(u(t))\,\dd t.
 \label{eq:v17-recent-response}
\end{equation}
Then
\begin{equation}
 \rho\in\operatorname{Ran}Q_K,
 \qquad
 \|\rho\|_2\ge(\sqrt3-1)R.
 \label{eq:v17-response-support-size}
\end{equation}

\begin{theorem}[Recent-card Sobolev source ladder]
\label{thm:v17-source-ladder}
For every real \(s_0\ge-1/4\),
\begin{equation}
 \boxed{\quad
 \int_{I^{\rm rec}}
 \|A^{s_0}Q_K\BNS(u(t))\|_2^2\,\dd t
 \ge
 c_{\rm rec}\nu K^{4s_0+1}R^2,
 \qquad
 c_{\rm rec}=2(\sqrt3-1)^2.
 \quad}
 \label{eq:v17-source-ladder}
\end{equation}
The endpoint \(s_0=-1/4\) is the V13 critical-source action.  The
order-zero member is the universal short-card estimate
\begin{equation}
 \boxed{\quad
 \int_{I^{\rm rec}}\|Q_K\BNS(u(t))\|_2^2\,\dd t
 \ge c_{\rm rec}\frac{\nu R^4}{E_*}.
 \quad}
 \label{eq:v17-recent-strong-source}
\end{equation}
Thus the quartic unweighted source debit is not confined to the full
corridor and does not depend on the adverse-interference branch isolated in
V15.
\end{theorem}

\begin{proof}
Write
\begin{equation}
 \theta:=s_0+\frac14\ge0,
 \qquad
 z_{s_0}(t):=A^{s_0}Q_K\BNS(u(t)).
 \label{eq:v17-ladder-parameters}
\end{equation}
Functional calculus gives
\begin{equation}
 A^\theta\rho
 =
 \int_s^b e^{-\nu A(b-t)}A^{1/2}z_{s_0}(t)\,\dd t.
 \label{eq:v17-differentiated-response}
\end{equation}
The one-mode heat-kernel calculation of
Lemma~\ref{lem:v13-record-response} applies to the right side and yields
\begin{equation}
 \|A^\theta\rho\|_2^2
 \le\frac1{2\nu}
 \int_s^b\|z_{s_0}(t)\|_2^2\,\dd t.
 \label{eq:v17-ladder-response-upper}
\end{equation}
Because \(\rho\) is supported where the
\(\Lambda=A^{1/2}\) frequency satisfies \(\lambda>K\),
\begin{equation}
 \|A^\theta\rho\|_2^2
 =\sum_{\lambda>K}\lambda^{4\theta}|\widehat\rho(\lambda)|^2
 \ge K^{4\theta}\|\rho\|_2^2
 \ge(\sqrt3-1)^2K^{4s_0+1}R^2.
 \label{eq:v17-ladder-spectral-lower}
\end{equation}
Combine \eqref{eq:v17-ladder-response-upper} and
\eqref{eq:v17-ladder-spectral-lower}.  Setting \(s_0=0\) and
\(K=R^2/E_*\) gives \eqref{eq:v17-recent-strong-source}.
\end{proof}

\begin{corollary}[Pointwise source concentration on the recent card]
\label{cor:v17-source-concentration}
The V13 length bound
\(\nu K^2|I^{\rm rec}|\le c_{\rm h}=\log4\) and
Theorem~\ref{thm:v17-source-ladder} imply
\begin{equation}
 \boxed{\quad
 \sup_{t\in I^{\rm rec}}
 \|A^{s_0}Q_K\BNS(u(t))\|_2^2
 \ge
 \frac{c_{\rm rec}}{c_{\rm h}}\,
 \nu^2K^{4s_0+3}R^2
 \quad(s_0\ge-1/4).
 \quad}
 \label{eq:v17-pointwise-source-ladder}
\end{equation}
In particular,
\begin{equation}
 \sup_{t\in I^{\rm rec}}\|Q_K\BNS(u(t))\|_2
 \ge
 \left(\frac{c_{\rm rec}}{c_{\rm h}}\right)^{1/2}
 \frac{\nu R^4}{E_*^{3/2}}.
 \label{eq:v17-pointwise-strong-source}
\end{equation}
\end{corollary}

\begin{proof}
The supremum of a nonnegative function is at least its integral divided by
the interval length.  Use
\(|I^{\rm rec}|\le c_{\rm h}/(\nu K^2)\) in
\eqref{eq:v17-source-ladder}.  The last formula is the square root of the
case \(s_0=0\), with \(K=R^2/E_*\).
\end{proof}

\begin{corollary}[Every terminal tail has infinite source action at all
orders in the ladder]
\label{cor:v17-terminal-source-ladder}
Suppose \(T_*<\infty\) and the critical record is unbounded.  Then, for
every \(t_\circ<T_*\) and every \(s_0\ge-1/4\),
\begin{equation}
 \boxed{\quad
 \int_{t_\circ}^{T_*}
 \|A^{s_0}\BNS(u(t))\|_2^2\,\dd t=\infty.
 \quad}
 \label{eq:v17-terminal-source-divergence}
\end{equation}
More quantitatively, the pairwise-disjoint recent cards satisfy
\begin{equation}
 \sum_n\int_{I_n^{\rm rec}}
 \|A^{s_0}Q_{K_n}\BNS(u(t))\|_2^2\,\dd t
 \ge
 c_{\rm rec}\nu E_*^{-(4s_0+1)}
 \sum_nR_n^{8s_0+4}
 =\infty.
 \label{eq:v17-source-ladder-sum}
\end{equation}
\end{corollary}

\begin{proof}
Use \eqref{eq:v17-source-ladder} on each card and
\(K_n=R_n^2/E_*\).  Since \(s_0\ge-1/4\), the exponent
\(8s_0+4\ge2\), so the doubled-record series diverges.  The cards eventually
lie in every terminal tail.  Finally \(Q_{K_n}\) is an orthogonal
contraction in the \(A^{s_0}L^2\) norm, so each projected action is no larger
than the corresponding full-source action.
\end{proof}

\begin{corollary}[Recent acceleration-or-diffusion concentration]
\label{cor:v17-recent-acceleration}
Every recent record card satisfies
\begin{equation}
 \int_{I^{\rm rec}}
 \left(\|\partial_tu(t)\|_2^2+\nu^2\|Au(t)\|_2^2\right)\,\dd t
 \ge
 \frac{c_{\rm rec}}2\frac{\nu R^4}{E_*}.
 \label{eq:v17-recent-acceleration-diffusion}
\end{equation}
At least one of the two integrals on the left is at least
\(c_{\rm rec}\nu R^4/(4E_*)\).
\end{corollary}

\begin{proof}
The equation and contraction of \(Q_K\) give
\[
 \|Q_K\BNS(u)\|_2^2
 \le\|\BNS(u)\|_2^2
 \le2\|\partial_tu\|_2^2+2\nu^2\|Au\|_2^2.
\]
Integrate and use \eqref{eq:v17-recent-strong-source}.  The final assertion
is the two-summand alternative.
\end{proof}

\begin{firewall}[Higher source order is higher unpaid order]
\label{fw:v17-source-ladder-unpaid}
The ladder \eqref{eq:v17-source-ladder} is a family of lower bounds, not a
new upper budget.  Increasing \(s_0\) makes the forced source norm stronger
and its record power larger, but also moves farther from the Leray-controlled
\(H^{-3/2}\) source.  The useful gain is localization: the quartic
order-zero action and all higher members occur on the same parabolic recent
card that is already forced by the record chronology.
\end{firewall}

% =============================================================================
\section{The record cap sharpens the low-advector commutator by a half power}
\label{sec:v17-sharp-commutator}
% =============================================================================

V16 bounded the low-advector Fourier \(\ell^1\) moment by kinetic energy.
The critical record itself gives a sharper bound at short times.

\begin{proposition}[Critical-record low-advector bound]
\label{prop:v17-critical-low-advector}
There is an absolute torus constant \(C_{\rm cr}\) such that every smooth
divergence-free field satisfies
\begin{equation}
 \boxed{\quad
 \left|\operatorname{Re}\langle
 \BNS(P_Lu,u),\Lambda u\rangle\right|
 \le C_{\rm cr}L^2\mathcal R(t)^3
 \qquad(L>0).
 \quad}
 \label{eq:v17-critical-low-advector}
\end{equation}
Together with \eqref{eq:v16-low-advector-work-bound}, this gives
\begin{equation}
 \left|\operatorname{Re}\langle
 \BNS(P_Lu,u),\Lambda u\rangle\right|
 \le
 \min\left\{
 C_{\rm com}L^{5/2}\sqrt{E_*}\,\mathcal R(t)^2,\,
 C_{\rm cr}L^2\mathcal R(t)^3
 \right\}.
 \label{eq:v17-combined-low-advector}
\end{equation}
\end{proposition}

\begin{proof}
The commutator proof of
Proposition~\ref{prop:v16-critical-commutator} reduced the left side to
\[
 C\left(\sum_{0<|p|\le L}|p|\,|\widehat u_p|\right)
 \mathcal R(t)^2.
\]
Instead of using kinetic energy, apply lattice Cauchy--Schwarz with one
critical half derivative:
\begin{align*}
 \sum_{0<|p|\le L}|p|\,|\widehat u_p|
 &\le
 \left(\sum_{0<|p|\le L}|p|\right)^{1/2}
 \left(\sum_{0<|p|\le L}|p|\,|\widehat u_p|^2\right)^{1/2}\\
 &\le C L^2\mathcal R(t).
\end{align*}
This proves \eqref{eq:v17-critical-low-advector}; taking the better of it
and \eqref{eq:v16-low-advector-work-bound} proves
\eqref{eq:v17-combined-low-advector}.
\end{proof}

\begin{theorem}[Sharp combined turnover edge]
\label{thm:v17-sharp-turnover-edge}
On a first-doubling corridor \(I=[a,b]\) of length \(\tau=b-a\), define
\begin{align}
 L_I^{E}
 &:=
 \left(8C_{\rm com}\sqrt{E_*}\,\tau\right)^{-2/5},
 \label{eq:v17-energy-edge}\\
 L_I^{R}
 &:=
 \left(16C_{\rm cr}R\tau\right)^{-1/2},
 \label{eq:v17-critical-edge}\\
 L_I^\sharp&:=\max\{L_I^E,L_I^R\}.
 \label{eq:v17-sharp-edge}
\end{align}
Then
\begin{equation}
 \boxed{\quad
 -\int_a^b\operatorname{Re}\langle
 \BNS(Q_{L_I^\sharp}u(t),u(t)),\Lambda u(t)\rangle\,\dd t
 \ge R^2.
 \quad}
 \label{eq:v17-sharp-high-advector-work}
\end{equation}
The second edge has the intrinsic nonlinear-turnover scale
\begin{equation}
 L_I^R\asymp(R\tau)^{-1/2}.
 \label{eq:v17-intrinsic-turnover}
\end{equation}
\end{theorem}

\begin{proof}
If \(L_I^\sharp=L_I^E\), the V16 energy estimate gives
\[
 \int_I\left|\operatorname{Re}\langle
 \BNS(P_{L_I^\sharp}u,u),\Lambda u\rangle\right|\,\dd t
 \le\frac{R^2}{2}.
\]
If \(L_I^\sharp=L_I^R\), then
\(\mathcal R(t)<2R\) on the first-hitting corridor, and
\eqref{eq:v17-critical-low-advector} gives
\[
 \int_I\left|\operatorname{Re}\langle
 \BNS(P_{L_I^\sharp}u,u),\Lambda u\rangle\right|\,\dd t
 \le8C_{\rm cr}(L_I^R)^2R^3\tau
 =\frac{R^2}{2}.
\]
In either case, subtract the low-advector work from the full lower bound
\eqref{eq:v16-full-work-lower}, which is \(3R^2/2\).
\end{proof}

\begin{corollary}[A fixed critical Lipschitz toll per doubling]
\label{cor:v17-critical-Lipschitz-toll}
Define the Fourier Lipschitz moment
\begin{equation}
 \mathfrak L(t):=\sum_{k\ne0}|k|\,|\widehat u_k(t)|.
 \label{eq:v17-Fourier-Lipschitz}
\end{equation}
Every first-doubling corridor satisfies
\begin{equation}
 \int_a^b\mathfrak L(t)\,\dd t
 \ge\frac{3}{8C_{\rm com}}.
 \label{eq:v17-Lipschitz-toll}
\end{equation}
Hence an unbounded doubled-record chronology forces
\(\int_{t_\circ}^{T_*}\mathfrak L(t)\,\dd t=\infty\) on every terminal
tail.
\end{corollary}

\begin{proof}
After enlarging \(C_{\rm com}\) once if necessary, the Fourier proof of
\eqref{eq:v16-critical-commutator}, without a low cutoff, gives
\[
 \left|\operatorname{Re}\langle
 \BNS(u,u),\Lambda u\rangle\right|
 \le C_{\rm com}\mathfrak L(t)\mathcal R(t)^2.
\]
Integrate, use \(\mathcal R(t)^2<4R^2\), and compare with
\eqref{eq:v16-full-work-lower}.  Sum over the disjoint doubling corridors
for the terminal-tail assertion.
\end{proof}

\begin{assessment}[The toll is a continuation boundary, not a closure]
The quantity \(\mathfrak L\) dominates a Lipschitz norm and its divergent
terminal integral is consistent with the classical vorticity/Lipschitz
blow-up boundary.  The gain in
Theorem~\ref{thm:v17-sharp-turnover-edge} is spectral: it identifies where
that toll must be paid on a fast corridor.  It does not make the toll
globally finite.
\end{assessment}

% =============================================================================
\section{A globally paid adaptive high-advector occupation ledger}
\label{sec:v17-adaptive-occupation}
% =============================================================================

\begin{theorem}[Paid occupation at the sharp turnover edge]
\label{thm:v17-paid-high-advector-occupation}
Let \(I_n=[t_n,t_{n+1}]\) be pairwise-disjoint first-doubling corridors,
let \(\tau_n=|I_n|\), and let \(L_n^\sharp\) be
\eqref{eq:v17-sharp-edge} with \(R=R_n\).  Then
\begin{align}
 \boxed{\quad
 \sum_n (L_n^\sharp)^2
 \int_{I_n}\|Q_{L_n^\sharp}u(t)\|_2^2\,\dd t
 \le\frac{E_*}{2\nu}.
 \quad}
 \label{eq:v17-paid-occupation}\\
 \boxed{\quad
 \sum_n\frac1{R_n\tau_n}
 \int_{I_n}\|Q_{L_n^\sharp}u(t)\|_2^2\,\dd t
 \le\frac{8C_{\rm cr}E_*}{\nu}.
 \quad}
 \label{eq:v17-turnover-occupation}
\end{align}
Thus the kinetic occupation of the productive high-advector sector has a
genuine global payment after multiplication by its squared turnover edge.
\end{theorem}

\begin{proof}
Spectral support gives, pointwise on \(I_n\),
\[
 (L_n^\sharp)^2\|Q_{L_n^\sharp}u\|_2^2
 \le\|\Lambda Q_{L_n^\sharp}u\|_2^2
 \le\|\Lambda u\|_2^2.
\]
Sum over the disjoint corridors and apply the Leray energy inequality to
obtain \eqref{eq:v17-paid-occupation}.  Since
\[
 (L_n^\sharp)^2\ge(L_{I_n}^R)^2
 =\frac1{16C_{\rm cr}R_n\tau_n},
\]
\eqref{eq:v17-turnover-occupation} follows after multiplying
\eqref{eq:v17-paid-occupation} by \(16C_{\rm cr}\).
\end{proof}

\begin{firewall}[Large signed work need not spend a fixed amount of paid
occupation]
\label{fw:v17-work-occupation-gap}
Theorem~\ref{thm:v17-sharp-turnover-edge} forces order-\(R_n^2\) signed
high-advector work, whereas
Theorem~\ref{thm:v17-paid-high-advector-occupation} controls a positive
kinetic occupation.  No inequality proved so far lower-bounds the latter by
the former.  Such an inequality would contradict the packet calculation
below.  The paid ledger is therefore an ancestry constraint to be combined
with recurrence, not a per-card closure.
\end{firewall}

% =============================================================================
\section{Every productive high-advector triad has a second high leg}
\label{sec:v17-two-high-legs}
% =============================================================================

The sharp work certificate identifies a high advector.  Fourier support
then supplies a second high leg, but in two geometrically different ways.

\begin{proposition}[Exact high-output versus high--high-return fork]
\label{prop:v17-two-high-leg-fork}
For any \(L>0\), define
\begin{align}
 \mathsf P_L^{\rm out}(t)
 &:=-\operatorname{Re}\langle
 \BNS(Q_Lu(t),u(t)),\Lambda Q_{L/2}u(t)\rangle,
 \label{eq:v17-high-output-work}\\
 \mathsf P_L^{\rm ret}(t)
 &:=-\operatorname{Re}\langle
 \BNS(Q_Lu(t),Q_{L/2}u(t)),\Lambda P_{L/2}u(t)\rangle.
 \label{eq:v17-high-return-work}
\end{align}
Then the exact signed identity
\begin{equation}
 -\operatorname{Re}\langle
 \BNS(Q_Lu,u),\Lambda u\rangle
 =
 \mathsf P_L^{\rm out}+\mathsf P_L^{\rm ret}
 \label{eq:v17-two-high-leg-identity}
\end{equation}
holds.  Consequently, on a first-doubling corridor with
\(L=L_I^\sharp\), at least one of
\begin{equation}
 \int_I\mathsf P_L^{\rm out}(t)\,\dd t\ge\frac{R^2}{2},
 \qquad
 \int_I\mathsf P_L^{\rm ret}(t)\,\dd t\ge\frac{R^2}{2}
 \label{eq:v17-two-high-leg-alternative}
\end{equation}
holds.

In the first branch the output carrier is above \(L/2\).  In the second,
both input parents are above \(L/2\) and their interaction returns to an
output below \(L/2\).
\end{proposition}

\begin{proof}
Split the output carrier as
\(u=P_{L/2}u+Q_{L/2}u\).  The high-output pairing is exactly
\(\mathsf P_L^{\rm out}\).  In the low-output pairing, split the advected
input in the second slot.  Additive Fourier support gives
\begin{equation}
 P_{L/2}\BNS(Q_Lu,P_{L/2}u)=0,
 \label{eq:v17-forbidden-one-high-return}
\end{equation}
because \(|p|>L\) and \(|q|\le L/2\) imply
\(|p+q|>L/2\).  Therefore only the second high input
\(Q_{L/2}u\) survives, proving
\eqref{eq:v17-two-high-leg-identity}.  Integrate and use
\eqref{eq:v17-sharp-high-advector-work}; if both signed summands were
smaller than \(R^2/2\), their sum would be smaller than \(R^2\).
\end{proof}

\begin{firewall}[The return branch carries the V14 multiplicity]
\label{fw:v17-return-multiplicity}
Equation~\eqref{eq:v17-forbidden-one-high-return} is a literal incidence
statement, not an injective matching of parents to outputs.  One very high
parent can pair with another high parent and return to many lower output
shells.  Taking positive parts shell by shell would reintroduce the
logarithmic multiplicity identified after
\eqref{eq:v14-shell-parent-descent}.  Any summation of the return branch
must preserve its signed assembly or construct an explicit one-use ancestry
selection.
\end{firewall}

% =============================================================================
\section{The actual packet saturates the intrinsic edge and evades the paid
occupation}
\label{sec:v17-packet-test}
% =============================================================================

\begin{proposition}[Sharp packet test for the V17 turnover ledger]
\label{prop:v17-packet-turnover}
For the V14 two-parameter packet family,
\[
 |J_j|=\frac{s_\#}{R_jN_j^2},
 \qquad
 \mathcal R_j\asymp R_j,
 \qquad
 \frac{R_j^2}{N_j}\longrightarrow0.
\]
The V17 critical edge obeys
\begin{equation}
 L_{J_j}^R\asymp N_j.
 \label{eq:v17-packet-sharp-edge}
\end{equation}
After fixing the harmless constant in the edge as in
Theorem~\ref{thm:v17-sharp-turnover-edge}, the packet has
\begin{align}
 -\int_{J_j}\operatorname{Re}\langle
 \BNS(Q_{L_{J_j}^\sharp}u^{(j)},u^{(j)}),
 \Lambda u^{(j)}\rangle\,\dd t
 &\ge cR_j^2,
 \label{eq:v17-packet-sharp-work}\\
 (L_{J_j}^\sharp)^2
 \int_{J_j}\|Q_{L_{J_j}^\sharp}u^{(j)}(t)\|_2^2\,\dd t
 &\le C\frac{R_j}{N_j}\longrightarrow0.
 \label{eq:v17-packet-paid-occupation-vanishes}
\end{align}
For every fixed \(s_0\ge-1/4\), its source action has the scaling
\begin{equation}
 \int_{J_j}\|A^{s_0}\BNS(u^{(j)}(t))\|_2^2\,\dd t
 \asymp R_j^3N_j^{4s_0+1}.
 \label{eq:v17-packet-source-ladder}
\end{equation}
\end{proposition}

\begin{proof}
Substitute the packet lifetime in
\eqref{eq:v17-critical-edge}:
\[
 (R_j|J_j|)^{-1/2}
 =s_\#^{-1/2}N_j.
\]
Equation~\eqref{eq:v16-packet-advector-edge} gives
\(L_{J_j}^E/N_j\to0\), so \(L_{J_j}^\sharp\asymp N_j\).
The uniform rescaled critical bound on \(J_j\) and the positive fractional
record growth give \eqref{eq:v17-packet-sharp-work} by the same proof as
Theorem~\ref{thm:v17-sharp-turnover-edge}.  The fixed reservoir lies below
the edge for large \(j\), while the packet kinetic energy is
\(O(R_j^2/N_j)\).  Hence
\[
 (L_{J_j}^\sharp)^2
 \int_{J_j}\|Q_{L_{J_j}^\sharp}u^{(j)}\|_2^2\,\dd t
 \le
 C N_j^2\frac{R_j^2}{N_j}\frac1{R_jN_j^2}
 =C\frac{R_j}{N_j}.
\]
Finally the instantaneous squared source norm has scale
\(\beta_j^4N_j^{5+4s_0}\), and
\(\dd t=(\beta_jN_j^{5/2})^{-1}\dd s\).  Thus its time integral is
\[
 \beta_j^3N_j^{5/2+4s_0}
 =R_j^3N_j^{4s_0+1},
\]
with the two-sided bound supplied by the nonzero seed source and short-time
rescaled convergence as in V14.
\end{proof}

\begin{assessment}[Sharpness supplied by the packet]
The intrinsic edge \((R\tau)^{-1/2}\) reaches the physical packet frequency,
so the V16 frequency underestimate has been removed.  Nevertheless the
globally paid high-advector occupation can tend to zero while the signed
work is of record size.  The source ladder is then overpaid at the escaping
physical frequency.  This identifies the precise missing statement: not
frequency localization, but a cross-record restriction preventing repeated
large work with vanishing paid occupation.
\end{assessment}

% =============================================================================
\section{Integrated V17 conclusion and V18 acquisition order}
\label{sec:v17-conclusion}
% =============================================================================

\begin{theorem}[What V17 proves]
\label{thm:v17-conclusion}
Relative to V1--V16, the following are proved.
\begin{enumerate}[label=\textup{(V17.\arabic*)},leftmargin=6.2em]
\item The qualitative recent-birth mechanism in the V16 compression
      alternative was already present in the temporally sharper V13
      one-card theorem; optimizing birth alone would not advance the
      closure.
\item Exterior spectral support upgrades the V13 critical-source debit to
      the full source-regularity ladder
      \eqref{eq:v17-source-ladder}.  In particular, every recent record card
      carries quartic unweighted source action and an explicit pointwise
      source spike.
\item Every terminal tail of an unbounded critical record has infinite
      source action in every ladder space \(A^{s_0}L^2\),
      \(s_0\ge-1/4\).
\item The critical record improves the low-advector work bound from
      \(L^{5/2}\sqrt{E_*}\mathcal R^2\) to
      \(L^2\mathcal R^3\), forcing productive advectors above the sharp
      intrinsic scale \((R\tau)^{-1/2}\).
\item The corresponding high-advector kinetic occupation has a global
      Leray payment after multiplication by the squared turnover edge.
\item Every productive high-advector triad has either a high output carrier
      or a second high input returning to a low output.
\item Actual record-growing packets saturate the intrinsic edge while their
      paid high-advector occupation vanishes.
\end{enumerate}
No full Navier--Stokes stability or global-regularity theorem is obtained.
\end{theorem}

\begin{proof}
Item \textup{(V17.1)} is
Proposition~\ref{prop:v17-v16-birth-redundancy}.  Items
\textup{(V17.2)--(V17.3)} are
Theorem~\ref{thm:v17-source-ladder} and
Corollaries~\ref{cor:v17-source-concentration}--%
\ref{cor:v17-terminal-source-ladder}.  Item \textup{(V17.4)} is
Proposition~\ref{prop:v17-critical-low-advector} and
Theorem~\ref{thm:v17-sharp-turnover-edge}.  Item \textup{(V17.5)} is
Theorem~\ref{thm:v17-paid-high-advector-occupation}.  Item
\textup{(V17.6)} is Proposition~\ref{prop:v17-two-high-leg-fork}.  Item
\textup{(V17.7)} is Proposition~\ref{prop:v17-packet-turnover}.
\end{proof}

\begin{programme}[V18 acquisition order]
\label{prog:v18}
\begin{enumerate}[label=\textup{(AC\arabic*)},leftmargin=4.8em]
\item Use the exact fork \eqref{eq:v17-two-high-leg-identity} as the root of
      a two-branch ancestry tree.  Keep the high-output branch complete and
      the high--high-return branch signed.
\item In the high-output branch, apply the heat-inherited/source-born
      carrier split at \(Q_{L_n^\sharp/2}\).  Determine whether its inherited
      cross term has a global bound after common-terminal propagation.
\item In the return branch, symmetrize complete triads before shell
      localization and retain the frequency-difference multiplier
      \eqref{eq:v16-symbol-difference}.  Seek a one-use parent selection that
      avoids the V14 logarithmic multiplicity.
\item Combine the paid occupation
      \eqref{eq:v17-paid-occupation} with first upcrossing times.  A usable
      recurrence theorem must force a positive occupation debit; large work
      alone cannot do so by
      \eqref{eq:v17-packet-paid-occupation-vanishes}.
\item Use the recent-card source ladder only through a candidate upper
      estimate of matching order.  Merely increasing \(s_0\) strengthens an
      unpaid norm and cannot close the cascade.
\item Test any proposed work-to-occupation or source-to-parent inequality on
      the two-parameter packet before summing it across records.
\item Keep the next cross-program review scheduled for V20.
\end{enumerate}
\end{programme}

\begin{assessment}[Position after V17]
The source side is now localized as strongly as the record response and
spectral support alone permit: a whole Sobolev ladder, including quartic
\(L^2\) source action, concentrates on each forced recent card.  The
transport side now reaches the true nonlinear turnover frequency and has a
globally paid high-advector occupation ledger.  The packet proves that these
two advances still do not meet per card.  The next possible gain is
combinatorial and dynamical across records: a recurrence or one-use ancestry
principle converting repeated signed two-high-leg work into a positive
amount of the paid occupation.
\end{assessment}

% =============================================================================
\section*{Version V17 history}
\addcontentsline{toc}{section}{Version V17 history}
% =============================================================================

Audited the V16 compression programme against V13 and removed a circular
direction: critical cofinal birth was already forced on the last recent
card.  Used the exterior support of that response to prove a full
recent-card source-regularity ladder and pointwise concentration.  Improved
the critical commutator by the record-weighted Fourier lattice sum, obtaining
the sharp intrinsic turnover edge \((R\tau)^{-1/2}\).  Established a
globally paid adaptive high-advector occupation ledger and an exact
high-output versus high--high-return fork.  Verified on actual packet
solutions that the new edge is sharp while paid occupation may vanish.  The
full stability/global-regularity goal remains open.

% =============================================================================
\section{V18 completion audit: ordered-slot work is not yet a carrier
equation}
\label{sec:v18-completion-audit}
% =============================================================================

The V17 two-high-leg fork is an exact decomposition of the
\emph{high-advector component} of critical work.  Before using the
high-output branch in the V16 carrier equation, one must restore the source
components omitted by the advector cutoff.

\begin{proposition}[Exact completion residual for the V17 high-output
branch]
\label{prop:v18-output-completion}
For \(L>0\), define the complete high-output work
\begin{equation}
 \mathsf P_{L}^{\rm comp,out}(t):=
 -\operatorname{Re}\langle
 \BNS(u(t),u(t)),\Lambda Q_{L/2}u(t)\rangle.
 \label{eq:v18-complete-output-work}
\end{equation}
Then
\begin{equation}
 \boxed{\quad
 \mathsf P_L^{\rm comp,out}
 =
 \mathsf P_L^{\rm out}
 -
 \operatorname{Re}\langle
 \BNS(P_Lu,u),\Lambda Q_{L/2}u\rangle.
 \quad}
 \label{eq:v18-output-completion-residual}
\end{equation}
The second term is not part of
\(\mathsf P_L^{\rm out}\), need not have either sign, and is present in the
actual equation for the carrier \(Q_{L/2}u\).
\end{proposition}

\begin{proof}
Use \(u=Q_Lu+P_Lu\) in the first, advector slot of the bilinear form
\(\BNS(\,\cdot\,,u)\).  The \(Q_Lu\) term is precisely
\(\mathsf P_L^{\rm out}\) from
\eqref{eq:v17-high-output-work}; the other term is the displayed residual.
\end{proof}

\begin{firewall}[Why a component carrier split would be invalid]
\label{fw:v18-component-carrier}
The equation for \(Q_{L/2}u\) is driven by the complete source
\(Q_{L/2}\BNS(u,u)\), not by
\(Q_{L/2}\BNS(Q_Lu,u)\) alone.  Nonnegativity of source-born work in
\eqref{eq:v16-born-work-identity} uses that complete equation.  Applying it
to \(\mathsf P_L^{\rm out}\) while discarding
\eqref{eq:v18-output-completion-residual} would assume the missing
cancellation.  V18 therefore completes the work by whole Fourier triads
before making any sign or ancestry selection.
\end{firewall}

% =============================================================================
\section{Complete triads convert critical work into frequency-gap flux}
\label{sec:v18-triad-completion}
% =============================================================================

Recall the modal transfer \(T_k\) and cumulative kinetic flux
\(\Phi_\rho\) from
\eqref{eq:v4-modal-work}--\eqref{eq:v4-cumulative-flux}.  Decompose each
modal transfer into complete unordered Fourier triads.  For one such triad,
write its wavevectors as \(k_1+k_2+k_3=0\), order their radii as
\begin{equation}
 0<\lambda_1:=|k_1|
 \le\lambda_2:=|k_2|
 \le\lambda_3:=|k_3|,
 \label{eq:v18-ordered-triad-radii}
\end{equation}
and denote its three real modal-transfer contributions by
\(T_1^\triangle,T_2^\triangle,T_3^\triangle\).

\begin{theorem}[Exact triad energy cancellation and frequency-gap formula]
\label{thm:v18-triad-gap-formula}
Every complete Navier--Stokes triad satisfies
\begin{equation}
 T_1^\triangle+T_2^\triangle+T_3^\triangle=0.
 \label{eq:v18-triad-energy-cancellation}
\end{equation}
Its contribution to the full critical nonlinear production is
\begin{align}
 \mathcal G_{1/2}^\triangle
 &:=-\sum_{i=1}^3\lambda_iT_i^\triangle
 \label{eq:v18-triad-critical-work}\\
 &=
 (\lambda_2-\lambda_1)T_1^\triangle
 -(\lambda_3-\lambda_2)T_3^\triangle.
 \label{eq:v18-triad-gap-work}
\end{align}
Equivalently, its cumulative kinetic-flux profile is
\begin{equation}
 \Phi_\rho^\triangle=
 \begin{cases}
 0,&0<\rho<\lambda_1,\\
 T_1^\triangle,&\lambda_1\le\rho<\lambda_2,\\
 -T_3^\triangle,&\lambda_2\le\rho<\lambda_3,\\
 0,&\rho\ge\lambda_3,
 \end{cases}
 \label{eq:v18-triad-flux-profile}
\end{equation}
and
\begin{equation}
 \boxed{\quad
 \mathcal G_{1/2}^\triangle
 =\int_0^\infty\Phi_\rho^\triangle\,\dd\rho.
 \quad}
 \label{eq:v18-triad-layer-cake}
\end{equation}
In particular, an equal-radius triad has zero critical production even when
its individual ordered-slot interactions are nonzero.
\end{theorem}

\begin{proof}
Restrict the trilinear energy identity
\(\langle\BNS(u,u),u\rangle=0\) to the three Fourier modes and their reality
partners.  Only pairings whose wavevectors close to
\(k_1+k_2+k_3=0\) survive, so the restricted identity is exactly
\eqref{eq:v18-triad-energy-cancellation}.  Eliminate
\(T_2^\triangle=-T_1^\triangle-T_3^\triangle\) in
\eqref{eq:v18-triad-critical-work}:
\[
 -\sum_i\lambda_iT_i^\triangle
 =
 (\lambda_2-\lambda_1)T_1^\triangle
 -(\lambda_3-\lambda_2)T_3^\triangle.
\]
This proves \eqref{eq:v18-triad-gap-work}.

By definition, the cumulative flux below the first radius is zero.  Between
the first two radii it contains only \(T_1^\triangle\); between the last two
it contains \(T_1^\triangle+T_2^\triangle=-T_3^\triangle\); above the last
radius it vanishes by
\eqref{eq:v18-triad-energy-cancellation}.  This proves
\eqref{eq:v18-triad-flux-profile}.  Integrating its two constant pieces
gives \eqref{eq:v18-triad-gap-work}, hence
\eqref{eq:v18-triad-layer-cake}.
\end{proof}

\begin{corollary}[High--high/low triads expose their exact lever arm]
\label{cor:v18-high-high-low-gap}
If \(\lambda_1\ll\lambda_2\le\lambda_3\), then the triangle inequality
gives
\begin{equation}
 0\le\lambda_3-\lambda_2\le\lambda_1.
 \label{eq:v18-high-parent-gap}
\end{equation}
Thus the second term in \eqref{eq:v18-triad-gap-work} has only a low-scale
gap, whereas the first term repeats the same low-mode kinetic transfer
across the long interval
\([\lambda_1,\lambda_2)\).  The apparent many-cutoff participation of the
high parent is exactly a frequency-gap lever arm in the completed signed
triad.
\end{corollary}

\begin{proof}
Since \(k_1+k_2=-k_3\),
\(\lambda_3\le\lambda_1+\lambda_2\), proving
\eqref{eq:v18-high-parent-gap}.  The rest is the explicit profile
\eqref{eq:v18-triad-flux-profile}.
\end{proof}

\begin{assessment}[What triad completion resolves]
\label{ass:v18-triad-resolution}
The V17 high-output and return terms are useful support incidences, but the
complete cancellation lives in
\eqref{eq:v18-triad-energy-cancellation}.  After completion, no unsigned
shell multiplicity is needed: the critical contribution is the signed
kinetic transfer times the actual spectral gap through which that transfer
acts.  The new question is whether a terminal cascade can repeatedly make
the transfer height small while sending the gap to a remote frequency.
\end{assessment}

% =============================================================================
\section{A record doubling forces positive net flux across at least three
canonical cutoff widths}
\label{sec:v18-flux-width}
% =============================================================================

For a smooth interval \(I=[a,b]\), set
\begin{equation}
 F_I(\rho):=\int_a^b\Phi_\rho(t)\,\dd t.
 \label{eq:v18-integrated-flux-profile}
\end{equation}
Smooth Fourier decay makes \(F_I\in L^1(0,\infty)\), while the V4 net-flux
theorem gives
\begin{equation}
 -\frac{E_*}{2}\le F_I(\rho)\le\frac{E_*}{2}
 \qquad(\rho>0).
 \label{eq:v18-flux-height-cap}
\end{equation}

\begin{theorem}[Cofinal positive-flux width on every first doubling]
\label{thm:v18-positive-flux-width}
Let \(I=[a,b]\) satisfy the first-doubling hypotheses and put
\(K=R^2/E_*\).  Then
\begin{equation}
 \int_0^\infty F_I(\rho)\,\dd\rho
 \ge\frac32R^2.
 \label{eq:v18-flux-area-lower}
\end{equation}
More precisely, for every \(0\le\alpha<3\),
\begin{align}
 \boxed{\quad
 \int_{\alpha K}^\infty F_I(\rho)\,\dd\rho
 \ge\frac{3-\alpha}{2}R^2,
 \quad}
 \label{eq:v18-cofinal-signed-flux-area}\\
 \boxed{\quad
 \int_{\alpha K}^\infty (F_I(\rho))_+\,\dd\rho
 \ge\frac{3-\alpha}{2}R^2,
 \quad}
 \label{eq:v18-cofinal-positive-flux-area}\\
 \boxed{\quad
 \left|\{\rho>\alpha K:F_I(\rho)>0\}\right|
 \ge(3-\alpha)K.
 \quad}
 \label{eq:v18-cofinal-positive-width}
\end{align}
In particular, positive net kinetic flux occurs on a set of cutoff radii of
total Lebesgue width at least \(3K\), and at least width \(K\) of that set
lies above \(2K\).
\end{theorem}

\begin{proof}
The full critical balance and the layer-cake identity
\eqref{eq:v5-critical-record-flux} give, by Fubini,
\[
 \int_0^\infty F_I(\rho)\,\dd\rho
 =
 \int_a^b-\operatorname{Re}\langle
 \BNS(u),\Lambda u\rangle\,\dd t
 =
 \frac32R^2+
 \nu\int_a^b\|\Lambda^{3/2}u\|_2^2\,\dd t.
\]
This proves \eqref{eq:v18-flux-area-lower}.  By
\eqref{eq:v18-flux-height-cap},
\[
 \int_0^{\alpha K}F_I(\rho)\,\dd\rho
 \le\frac{\alpha KE_*}{2}
 =\frac{\alpha R^2}{2}.
\]
Subtract from \eqref{eq:v18-flux-area-lower} to obtain
\eqref{eq:v18-cofinal-signed-flux-area}.  A signed integral is no larger
than the integral of the positive part, proving
\eqref{eq:v18-cofinal-positive-flux-area}.  Finally
\((F_I)_+\le E_*/2\); divide the positive-area lower bound by that height
cap to prove \eqref{eq:v18-cofinal-positive-width}.
\end{proof}

\begin{corollary}[Flux-height versus lever-width alternative]
\label{cor:v18-height-width}
Let
\begin{equation}
 M_I:=\operatorname*{ess\,sup}_{\rho>0}(F_I(\rho))_+,
 \qquad
 W_I:=\left|\{\rho>0:F_I(\rho)>0\}\right|.
 \label{eq:v18-height-width-definition}
\end{equation}
Then
\begin{equation}
 0<M_I\le\frac{E_*}{2},
 \qquad
 M_IW_I\ge\frac32R^2,
 \qquad
 W_I\ge3K.
 \label{eq:v18-height-width-product}
\end{equation}
If the positive flux is supported below a frequency \(\Gamma<\infty\), then
\begin{equation}
 M_I\ge\frac{3R^2}{2\Gamma}.
 \label{eq:v18-bounded-width-height}
\end{equation}
Thus a small kinetic-flux height can drive critical-record growth only by
occupying a long interval of cutoff radii.
\end{corollary}

\begin{proof}
The positive flux area is at least \(3R^2/2\).  It is at most \(M_IW_I\),
which gives the product bound.  The height cap and width bound follow from
Theorem~\ref{thm:v18-positive-flux-width}.  If the positive support lies in
\([0,\Gamma]\), then \(W_I\le\Gamma\).
\end{proof}

\begin{definition}[First positive-flux quantile]
\label{def:v18-flux-quantile}
Define
\begin{equation}
 \Gamma_I:=
 \inf\left\{\Gamma>0:
 \int_0^\Gamma(F_I(\rho))_+\,\dd\rho
 \ge\frac34R^2\right\}.
 \label{eq:v18-flux-quantile}
\end{equation}
The total positive area bound makes \(\Gamma_I<\infty\).
\end{definition}

\begin{corollary}[The flux quantile is already supercanonical]
\label{cor:v18-supercanonical-quantile}
Every first-doubling corridor satisfies
\begin{equation}
 \boxed{\quad \Gamma_I\ge\frac32K. \quad}
 \label{eq:v18-quantile-lower}
\end{equation}
\end{corollary}

\begin{proof}
On an interval of cutoff width \(\Gamma\), the height cap permits positive
area at most \(\Gamma E_*/2\).  Reaching \(3R^2/4\) therefore requires
\(\Gamma E_*/2\ge3R^2/4\), which is
\eqref{eq:v18-quantile-lower}.
\end{proof}

% =============================================================================
\section{Cumulative record flux is necessarily cofinal}
\label{sec:v18-cumulative-cofinality}
% =============================================================================

The preceding theorem is per corridor.  Because the first-hitting corridors
are adjacent, their \emph{signed} flux profiles telescope in time before any
positive part is taken.

\begin{theorem}[Escape of normalized positive kinetic-flux mass]
\label{thm:v18-flux-mass-escape}
Let \(t_n,R_n\) be the doubled first-hitting chronology and define
\begin{equation}
 F_N(\rho):=\int_{t_0}^{t_N}\Phi_\rho(t)\,\dd t,
 \qquad
 A_N:=\int_0^\infty(F_N(\rho))_+\,\dd\rho.
 \label{eq:v18-cumulative-flux}
\end{equation}
Then
\begin{equation}
 A_N\ge\frac12(R_N^2-R_0^2).
 \label{eq:v18-cumulative-positive-mass}
\end{equation}
For the probability measure
\begin{equation}
 \mu_N(S):=
 \frac{\displaystyle\int_S(F_N(\rho))_+\,\dd\rho}{A_N},
 \label{eq:v18-normalized-flux-measure}
\end{equation}
one has, for every fixed \(L<\infty\),
\begin{equation}
 \boxed{\quad
 \mu_N([0,L])
 \le\frac{LE_*}{R_N^2-R_0^2}
 \longrightarrow0.
 \quad}
 \label{eq:v18-flux-measure-escape}
\end{equation}
Hence the normalized positive kinetic-flux mass underlying an unbounded
critical record is not tight on the frequency half-line: it escapes beyond
every fixed cutoff.
\end{theorem}

\begin{proof}
The layer-cake identity and the critical balance on \([t_0,t_N]\) give
\begin{align}
 \int_0^\infty F_N(\rho)\,\dd\rho
 &=
 \frac12(R_N^2-R_0^2)
 +\nu\int_{t_0}^{t_N}\|\Lambda^{3/2}u(t)\|_2^2\,\dd t
 \label{eq:v18-cumulative-signed-area}\\
 &\ge\frac12(R_N^2-R_0^2).
 \notag
\end{align}
The positive part has at least the signed mass, proving
\eqref{eq:v18-cumulative-positive-mass}.  The net-flux bound
\eqref{eq:v4-net-flux-bound}, now on the single interval
\([t_0,t_N]\), gives
\[
 \int_0^L(F_N(\rho))_+\,\dd\rho\le\frac{LE_*}{2}.
\]
Divide this by \eqref{eq:v18-cumulative-positive-mass} to obtain
\eqref{eq:v18-flux-measure-escape}.
\end{proof}

\begin{corollary}[No bounded-cutoff recurrence can carry macroscopic record
flux]
\label{cor:v18-no-bounded-flux-recurrence}
For every fixed \(L\),
\begin{equation}
 \frac1{R_N^2}
 \int_0^L
 \left(\int_{t_0}^{t_N}\Phi_\rho(t)\,\dd t\right)_+
 \dd\rho\longrightarrow0.
 \label{eq:v18-bounded-flux-negligible}
\end{equation}
Therefore any recurrence principle capable of closing the record cascade
must follow cutoffs tending to infinity; recurrence at a fixed physical
frequency is asymptotically negligible.
\end{corollary}

\begin{proof}
The numerator is at most \(LE_*/2\), independently of \(N\).
\end{proof}

\begin{assessment}[Cofinality is necessary, not contradictory]
\label{ass:v18-cofinality-status}
The measures \(\mu_N\) escape to infinity, but probability mass may do so
without violating any compact-frequency bound.  This is the flux version of
the V14 effective-response-frequency escape.  A closure theorem would need
a tightness mechanism, a dissipative cost for translating the flux profile,
or a one-use kinetic-energy transport rule across the moving quantiles.
None has yet been proved.
\end{assessment}

% =============================================================================
\section{The packet realizes small flux height times remote flux width}
\label{sec:v18-packet-flux}
% =============================================================================

For the V14 packet on \(J_j\), let
\begin{equation}
 F_j(\rho):=\int_{J_j}\Phi_\rho^{(j)}(t)\,\dd t,
 \qquad
 \beta_j^2=\frac{R_j^2}{N_j}.
 \label{eq:v18-packet-flux-profile}
\end{equation}

\begin{proposition}[Packet flux profile at the physical frequency]
\label{prop:v18-packet-flux-profile}
There are a compact interval
\([r_-,r_+]\Subset(0,\infty)\) and constants \(c,C>0\), depending only on
the fixed seed and \(s_\#\), such that for all sufficiently large \(j\),
\begin{align}
 \int_{r_-N_j}^{r_+N_j}(F_j(\rho))_+\,\dd\rho
 &\ge cR_j^2,
 \label{eq:v18-packet-positive-flux-area}\\
 \sup_{\rho\in[r_-N_j,r_+N_j]}(F_j(\rho))_+
 &\le C\frac{R_j^2}{N_j},
 \label{eq:v18-packet-flux-height}\\
 \left|
 \{\rho\in[r_-N_j,r_+N_j]:F_j(\rho)>0\}
 \right|
 &\ge cN_j.
 \label{eq:v18-packet-flux-width}
\end{align}
Thus the actual smooth packet produces order-\(R_j^2\) critical work with
kinetic-flux height of the same order as its vanishing kinetic energy
\(R_j^2/N_j\), spread across cutoff width of order \(N_j\).
\end{proposition}

\begin{proof}
Use the V10--V11 rescaled variables
\[
 s=\beta_jN_j^{5/2}t,\qquad r=\rho/N_j,
\]
and define the rescaled integrated flux
\begin{equation}
 \widetilde F_j(r):=\beta_j^{-2}F_j(N_jr).
 \label{eq:v18-rescaled-flux-profile}
\end{equation}
The uniform rescaled \(H^m\) bounds with \(m\ge5\), the expanding-torus
Fourier Riemann-sum limit, and \(\nu/R_j\to0\) give convergence in
\(L^1_{\rm loc}((0,\infty))\) to the smooth rapidly decaying cumulative
flux profile \(\widetilde F_\infty(r)\) of the fixed seed evolution, as
well as a uniform local \(L^\infty\) bound.  Indeed, truncate the Fourier
sums using the uniform \(H^m\) tail, pass to the limit in the finite sums,
and then remove the truncation.  The fixed low-frequency reservoir vanishes
in these variables.

The rescaled critical balance
\eqref{eq:v11-rescaled-record-balance} and the choice of \(s_\#\) give
\begin{equation}
 \int_0^\infty\widetilde F_\infty(r)\,\dd r>0.
 \label{eq:v18-limit-flux-positive-area}
\end{equation}
Rapid decay at infinity and vanishing cumulative flux as \(r\downarrow0\)
permit a compact interval \([r_-,r_+]\Subset(0,\infty)\) on which
\[
 \int_{r_-}^{r_+}(\widetilde F_\infty(r))_+\,\dd r>2c
\]
for some \(c>0\).  Local \(L^1\) convergence and the uniform local \(L^\infty\) bound then
give, for large \(j\),
\[
 \int_{r_-}^{r_+}(\widetilde F_j(r))_+\,\dd r\ge c,
 \qquad
 \sup_{r\in[r_-,r_+]}(\widetilde F_j(r))_+\le C.
\]
Undo the scaling and use
\(N_j\beta_j^2=R_j^2\) to obtain
\eqref{eq:v18-packet-positive-flux-area} and
\eqref{eq:v18-packet-flux-height}.  Dividing the first by the second proves
\eqref{eq:v18-packet-flux-width}.
\end{proof}

\begin{firewall}[The packet is a profile-escape test, not one singular
trajectory]
\label{fw:v18-packet-flux-scope}
The packet solutions are distinct.  Proposition
\ref{prop:v18-packet-flux-profile} does not concatenate their flux profiles
or construct blow-up.  It does prove that energy, the critical record, the
net-flux cap, triad completion, and actual short-time NS dynamics permit the
height--width escape
\[
 \text{flux height}\asymp\frac{R^2}{N},
 \qquad
 \text{flux width}\asymp N,
 \qquad
 \text{critical area}\asymp R^2
\]
with \(N/K\to\infty\).  Hence no universal bound
\(\Gamma_I\lesssim K\), nor a positive lower bound for \(M_I/E_*\), follows
from the presently recorded per-card data.
\end{firewall}

% =============================================================================
\section{Integrated V18 conclusion and V19 acquisition order}
\label{sec:v18-conclusion}
% =============================================================================

\begin{theorem}[What V18 proves]
\label{thm:v18-conclusion}
Relative to V1--V17, the following are proved.
\begin{enumerate}[label=\textup{(V18.\arabic*)},leftmargin=6.2em]
\item The V17 ordered-slot high-output work is not the complete source in
      the carrier equation; it has the exact residual
      \eqref{eq:v18-output-completion-residual}.
\item Completing a Fourier triad restores kinetic-energy cancellation and
      expresses its critical production exactly as signed kinetic transfer
      across its two frequency gaps.
\item Every first critical-record doubling forces positive net kinetic flux
      over at least \(3K\) total cutoff width, with at least width \(K\)
      lying above \(2K\).
\item The first fixed positive-flux quantile is at least \(3K/2\), and small
      flux height can contribute record-size critical work only through a
      long frequency lever arm.
\item On cumulative first-hitting intervals, normalized positive flux mass
      escapes beyond every fixed frequency cutoff.
\item Actual record-growing packets realize the permitted escape with flux
      height \(\asymp R^2/N\), width \(\asymp N\), and
      \(N/K\to\infty\).
\end{enumerate}
No full Navier--Stokes stability or global-regularity theorem is obtained.
\end{theorem}

\begin{proof}
Item \textup{(V18.1)} is Proposition~\ref{prop:v18-output-completion}.
Item \textup{(V18.2)} is
Theorem~\ref{thm:v18-triad-gap-formula}.  Items
\textup{(V18.3)--(V18.4)} are
Theorem~\ref{thm:v18-positive-flux-width} and
Corollaries~\ref{cor:v18-height-width}--%
\ref{cor:v18-supercanonical-quantile}.  Item \textup{(V18.5)} is
Theorem~\ref{thm:v18-flux-mass-escape}.  Item \textup{(V18.6)} is
Proposition~\ref{prop:v18-packet-flux-profile}.
\end{proof}

\begin{programme}[V19 acquisition order]
\label{prog:v19}
\begin{enumerate}[label=\textup{(AC\arabic*)},leftmargin=4.8em]
\item Use \(F_I(\rho)\,\dd\rho\), not an ordered-slot work component, as
      the signed ancestry object.  Preserve time telescoping before taking
      positive parts.
\item Attach to each doubling corridor its quantile
      \(\Gamma_I\) and the dimensionless viscous exposure
      \(\nu\Gamma_I^2|I|\).  Separate long-exposure flux from rapid
      frequency translation.
\item On the long-exposure branch, seek a lower bound for high-frequency
      Leray dissipation without assuming persistence of the flux sign.
\item On the rapid-translation branch, combine the intrinsic advector edge
      \((R|I|)^{-1/2}\) with the completed triad gap formula.  Determine
      whether one packet of kinetic energy can cross successive quantiles
      without a reusable-flux debit.
\item Investigate signed increments
      \(F_{[t_n,t_{n+1}]}\) inside the cumulative profile \(F_N\).  A useful
      recurrence theorem must prevent large positive increments from being
      erased by later negative variation at the same cutoff.
\item Test any proposed tightness, quantile, or flux-height lower bound
      against Proposition~\ref{prop:v18-packet-flux-profile}; the packet
      permits \(\Gamma_I/K\to\infty\) and \(M_I/E_*\to0\).
\item Retain the source ladder only if it supplies a matching upper estimate
      for a completed triad or flux profile.
\item Perform the required SM/YM companion review at V20, not before it.
\end{enumerate}
\end{programme}

\begin{assessment}[Position after V18]
The two-high-leg work has now been completed at the level where nonlinear
energy cancellation is exact.  Record growth is a positive area under a
bounded kinetic-flux profile, so it necessarily occupies a long cofinal
frequency interval.  Across a terminal chronology the normalized profile
must escape to infinity.  Actual packets show that this escape can occur by
letting flux height and kinetic mass vanish while the frequency lever arm
diverges.  The remaining possible closure must therefore control
\emph{translation of a signed flux profile across successive records}; a
per-card amplitude, source, or frequency bound has again been ruled out.
\end{assessment}

% =============================================================================
\section*{Version V18 history}
\addcontentsline{toc}{section}{Version V18 history}
% =============================================================================

Completed the V17 ordered-slot work by whole Fourier triads and proved the
exact frequency-gap kinetic-flux formula.  Combined the full critical
balance with the uniform net-flux debit to show that every record doubling
requires positive kinetic flux across at least three canonical cutoff
widths and introduced a supercanonical positive-flux quantile.  Proved that
normalized cumulative positive flux mass escapes beyond every fixed cutoff
on an unbounded record chronology.  Derived the rescaled packet flux profile
and verified the sharp small-height/remote-width escape.  The full
stability/global-regularity goal remains open.

% =============================================================================
% VERSION V19 APPEND
% =============================================================================

\section{V19: flux acquisition, parabolic exposure, and the moving-tail
ledger}
\label{sec:v19-entry}
% =============================================================================

V18 proved that a first critical-record doubling produces a bounded positive
kinetic-flux profile with large radial area, and that the resulting flux mass
must move to cofinal cutoffs along an unbounded record chronology.  The first
V19 acquisition order proposed the scalar exposure
\(\nu\Gamma_I^2|I|\).  There is an upstream obstruction to using that number
by itself: integrated positive flux through a cutoff is the sum of viscous
loss above the cutoff and net kinetic energy stored there at the terminal
time.  A long corridor does not imply that the terminal stock resided there
for the whole corridor.

This version keeps the endpoint-storage term exactly.  Two positive-flux
quantiles isolate a radial band carrying a fixed amount of critical work.  A
cutoff selected in that band then obeys a rigorous trichotomy: it pays Leray
dissipation, creates an intermediate-order source, or translates across
a radial width exponentially large relative to its parabolic exposure.  On
successive storage cards, an exact moving-tail telescope identifies the
remaining cross-record object: kinetic stock in the bands swept out when the
selected cutoff moves upward.

% =============================================================================
\section{Integrated flux is endpoint acquisition plus exterior dissipation}
\label{sec:v19-flux-acquisition}
% =============================================================================

For a smooth interval \(I=[a,b]\) and \(\rho>0\), set
\begin{equation}
 \mathcal H_\rho(t):=\|Q_\rho u(t)\|_2^2,
 \qquad
 \mathcal D_{\rho,I}:=
 \int_a^b\|\Lambda Q_\rho u(t)\|_2^2\,\dd t.
 \label{eq:v19-tail-stock-dissipation}
\end{equation}

\begin{proposition}[Exact cutoff acquisition identity]
\label{prop:v19-cutoff-acquisition}
For every \(\rho>0\),
\begin{equation}
 \boxed{\quad
 F_I(\rho)
 =\frac12\bigl(\mathcal H_\rho(b)-\mathcal H_\rho(a)\bigr)
  +\nu\mathcal D_{\rho,I}.
 \quad}
 \label{eq:v19-cutoff-acquisition}
\end{equation}
Consequently, if \(F_I(\rho)\ge m>0\), then at least one of
\begin{align}
 \nu\mathcal D_{\rho,I}&\ge\frac m2,
 \label{eq:v19-cutoff-diss-branch}\\
 \mathcal H_\rho(b)-\mathcal H_\rho(a)&\ge m
 \label{eq:v19-cutoff-storage-branch}
\end{align}
holds.  Moreover, for \(0\le L<M<\infty\),
\begin{align}
 \int_L^M\mathcal H_\rho(t)\,\dd\rho
 &=\sum_{k\ne0}
   \bigl(\min\{M,|k|\}-L\bigr)_+|\widehat u_k(t)|^2,
 \label{eq:v19-tail-layer-stock}\\
 \int_L^M\|\Lambda Q_\rho u(t)\|_2^2\,\dd\rho
 &=\sum_{k\ne0}
   \bigl(\min\{M,|k|\}-L\bigr)_+|k|^2
   |\widehat u_k(t)|^2.
 \label{eq:v19-tail-layer-dissipation}
\end{align}
Thus integration over cutoff width repeats a single Fourier mode once for
every cutoff below that mode; cutoff width is not a collection of orthogonal
kinetic-energy accounts.
\end{proposition}

\begin{proof}
Integrate the high-frequency energy balance
\eqref{eq:v4-high-energy-balance} from \(a\) to \(b\).  This gives
\eqref{eq:v19-cutoff-acquisition}.  If
\eqref{eq:v19-cutoff-diss-branch} fails, then
\[
 \frac12\bigl(\mathcal H_\rho(b)-\mathcal H_\rho(a)\bigr)
 =F_I(\rho)-\nu\mathcal D_{\rho,I}>\frac m2,
\]
which proves \eqref{eq:v19-cutoff-storage-branch}.

For a fixed mode \(k\), the indicator of \(Q_\rho\) is
\(\mathbf 1_{\{\rho<|k|\}}\).  Its integral over \([L,M]\) is
\((\min\{M,|k|\}-L)_+\).  Tonelli's theorem proves
\eqref{eq:v19-tail-layer-stock} and
\eqref{eq:v19-tail-layer-dissipation}.
\end{proof}

\begin{firewall}[Why exposure must be attached to acquired stock]
\label{fw:v19-exposure-needs-storage}
Equation~\eqref{eq:v19-cutoff-acquisition} is signed and exact.  The number
\(\nu\rho^2|I|\) damps stock inherited from time \(a\), but it does not say
when the source-created part of \(Q_\rho u(b)\) was born.  Replacing
\(\mathcal H_\rho(b)-\mathcal H_\rho(a)\) by a residence lasting \(|I|\)
would insert precisely the persistence statement missing since V7.  The
argument below instead uses Duhamel's formula and retains a rapid-translation
alternative.
\end{firewall}

% =============================================================================
\section{A fixed-work quantile band and a pointwise flux selector}
\label{sec:v19-quantile-band}
% =============================================================================

Retain a first-doubling corridor \(I=[a,b]\), write
\begin{equation}
 \tau:=b-a,
 \qquad
 K:=\frac{R^2}{E_*},
 \qquad
 \mathscr A_I(x):=\int_0^x(F_I(\rho))_+\,\dd\rho,
 \label{eq:v19-positive-area-function}
\end{equation}
and, for \(0<\theta<3/2\), define
\begin{equation}
 \Gamma_I(\theta):=
 \inf\{x>0:\mathscr A_I(x)\ge\theta R^2\}.
 \label{eq:v19-theta-quantile}
\end{equation}
The V18 quantile is \(\Gamma_I(3/4)\).

\begin{lemma}[Two-quantile band geometry]
\label{lem:v19-quantile-band}
The quantiles in \eqref{eq:v19-theta-quantile} are finite and satisfy
\begin{equation}
 \mathscr A_I(\Gamma_I(\theta))=\theta R^2,
 \qquad
 \Gamma_I(\theta)\ge2\theta K.
 \label{eq:v19-quantile-level-bound}
\end{equation}
In particular, put
\begin{equation}
 \gamma_I:=\Gamma_I(1/2),
 \qquad
 \Gamma_I:=\Gamma_I(3/4),
 \qquad
 W_I:=\Gamma_I-\gamma_I.
 \label{eq:v19-band-definitions}
\end{equation}
Then
\begin{align}
 \gamma_I&\ge K,
 &\Gamma_I&\ge\frac32K,
 &W_I&\ge\frac12K,
 \label{eq:v19-band-lower-scales}\\
 \int_{\gamma_I}^{\Gamma_I}(F_I(\rho))_+\,\dd\rho
 &=\frac14R^2.
 \label{eq:v19-band-fixed-area}
\end{align}
There exists a cutoff
\begin{equation}
 \rho_I\in(\gamma_I,\Gamma_I)
 \quad\hbox{such that}\quad
 F_I(\rho_I)\ge
 m_I:=\frac{R^2}{8W_I}.
 \label{eq:v19-flux-selector}
\end{equation}
\end{lemma}

\begin{proof}
Theorem~\ref{thm:v18-positive-flux-width} gives
\(\int_0^\infty(F_I)_+\ge3R^2/2\), so every quantile with
\(\theta<3/2\) is finite.  The function \(\mathscr A_I\) is absolutely
continuous.  Its continuity and the definition by first crossing give the
first equality in \eqref{eq:v19-quantile-level-bound}.  The height cap
\((F_I)_+\le E_*/2\) gives
\[
 \theta R^2\le\frac{E_*}{2}\Gamma_I(\theta),
\]
which proves the second assertion.

Subtract the two exact quantile levels to obtain
\eqref{eq:v19-band-fixed-area}.  Applying the same height cap only on the
band gives
\[
 \frac14R^2\le\frac{E_*}{2}W_I,
\]
hence \(W_I\ge K/2\).  Finally, if
\(F_I<m_I\) almost everywhere on the band, its positive area would be at
most \(m_IW_I=R^2/8\), contradicting
\eqref{eq:v19-band-fixed-area}.  Therefore the set of admissible selectors
in \eqref{eq:v19-flux-selector} has positive measure.
\end{proof}

\begin{remark}[Why two quantiles are used]
The single outer quantile \(\Gamma_I\) locates an amount of positive area but
does not say that the flux is large near that radius.  The annulus
\([\gamma_I,\Gamma_I]\) carries the exact area \(R^2/4\); its width therefore
supplies the pointwise height \(m_I\) without assuming monotonicity or a sign
for \(F_I\) outside the band.
\end{remark}

% =============================================================================
\section{Dissipation, source birth, or exponentially broad translation}
\label{sec:v19-exposure-trichotomy}
% =============================================================================

At the selected cutoff define the intermediate source action
\begin{equation}
 \mathscr S_I^{\rm mid}:=
 \int_a^b
 \|A^{-1/2}Q_{\rho_I}\BNS(u(t))\|_2^2\,\dd t
 \label{eq:v19-intermediate-source-action}
\end{equation}
and the inner-quantile exposure
\begin{equation}
 \Xi_I:=\nu\gamma_I^2\tau.
 \label{eq:v19-inner-exposure}
\end{equation}

\begin{theorem}[Flux-band acquisition trichotomy]
\label{thm:v19-acquisition-trichotomy}
For every first-doubling corridor, choose \(\rho_I\) as in
\eqref{eq:v19-flux-selector}.  At least one of the following holds:
\begin{align}
 \boxed{\quad
 \nu\mathcal D_{\rho_I,I}
 \ge\frac{R^2}{16W_I};
 \quad}
 &\tag{D}\label{eq:v19-dissipation-alternative}\\
 \boxed{\quad
 \mathscr S_I^{\rm mid}
 \ge\frac{\nu R^2}{16W_I};
 \quad}
 &\tag{B}\label{eq:v19-birth-alternative}\\
 \boxed{\quad
 \Xi_I
 <\frac12
 \biggl[\log\!\left(\frac{32W_I}{\gamma_I}\right)\biggr]_+.
 \quad}
 &\tag{T}\label{eq:v19-translation-alternative}
\end{align}
Here \([x]_+=\max\{x,0\}\).  More precisely, if
\textup{(D)} fails then the selected exterior kinetic stock has the endpoint
gain
\begin{equation}
 \mathcal H_{\rho_I}(b)-\mathcal H_{\rho_I}(a)
 \ge m_I,
 \label{eq:v19-selected-storage-gain}
\end{equation}
and if both \textup{(D)} and \textup{(T)} fail, that gain forces
\textup{(B)}.
\end{theorem}

\begin{proof}
Apply Proposition~\ref{prop:v19-cutoff-acquisition} with
\(m=m_I\).  Failure of \textup{(D)} is exactly failure of
\eqref{eq:v19-cutoff-diss-branch}, and therefore gives
\eqref{eq:v19-selected-storage-gain}.  In particular,
\begin{equation}
 \|Q_{\rho_I}u(b)\|_2\ge\sqrt{m_I}.
 \label{eq:v19-final-selected-tail}
\end{equation}

At time \(a\), the critical record and \(\rho_I\ge\gamma_I\) give
\begin{equation}
 \mathcal H_{\rho_I}(a)
 \le\rho_I^{-1}\|\Lambda^{1/2}u(a)\|_2^2
 \le\frac{R^2}{\gamma_I}.
 \label{eq:v19-initial-tail-cap}
\end{equation}
Duhamel's formula at the selected cutoff is
\begin{equation}
 Q_{\rho_I}u(b)
 =e^{-\nu A\tau}Q_{\rho_I}u(a)-z_I,
 \qquad
 z_I:=\int_a^b e^{-\nu A(b-t)}
 Q_{\rho_I}\BNS(u(t))\,\dd t.
 \label{eq:v19-selected-Duhamel}
\end{equation}
The inherited term satisfies
\begin{equation}
 \|e^{-\nu A\tau}Q_{\rho_I}u(a)\|_2
 \le e^{-\Xi_I}\frac{R}{\sqrt{\gamma_I}}.
 \label{eq:v19-inherited-selected-bound}
\end{equation}

Suppose \textup{(T)} also fails.  If \(32W_I\ge\gamma_I\), its negation
and \eqref{eq:v19-inherited-selected-bound} give
\[
 \|e^{-\nu A\tau}Q_{\rho_I}u(a)\|_2
 \le\frac{R}{\sqrt{32W_I}}.
\]
If \(32W_I<\gamma_I\), the same conclusion follows directly from
\eqref{eq:v19-inherited-selected-bound}.  Since
\(m_I=R^2/(8W_I)\), the last bound is \(\sqrt{m_I}/2\).  The reverse
triangle inequality in \eqref{eq:v19-selected-Duhamel} and
\eqref{eq:v19-final-selected-tail} therefore yield
\begin{equation}
 \|z_I\|_2\ge\frac12\sqrt{m_I}.
 \label{eq:v19-born-kinetic-response}
\end{equation}

Put \(g=A^{-1/2}Q_{\rho_I}\BNS(u)\).  Modewise time
Cauchy--Schwarz gives
\[
 \left|\int_a^b
   \lambda e^{-\nu\lambda^2(b-t)}g_\lambda(t)\,\dd t
 \right|^2
 \le
 \frac1{2\nu}\int_a^b|g_\lambda(t)|^2\,\dd t,
\]
because
\(\int_a^b\lambda^2e^{-2\nu\lambda^2(b-t)}\,\dd t
\le(2\nu)^{-1}\).  Sum the orthogonal modes to obtain
\begin{equation}
 \|z_I\|_2^2\le\frac1{2\nu}\mathscr S_I^{\rm mid}.
 \label{eq:v19-mid-response-upper}
\end{equation}
Combine \eqref{eq:v19-born-kinetic-response} and
\eqref{eq:v19-mid-response-upper}:
\[
 \mathscr S_I^{\rm mid}
 \ge2\nu\frac{m_I}{4}
 =\frac{\nu R^2}{16W_I},
\]
which is \textup{(B)}.
\end{proof}

\begin{corollary}[Thin remote bands are paid directly by Leray dissipation]
\label{cor:v19-thin-band-paid}
If
\begin{equation}
 W_I\le\frac{\gamma_I}{64},
 \label{eq:v19-thin-band-condition}
\end{equation}
then the selected cutoff necessarily satisfies
\begin{equation}
 \boxed{\quad
 \nu\mathcal D_{\rho_I,I}
 \ge\frac{3R^2}{32W_I}.
 \quad}
 \label{eq:v19-thin-band-dissipation}
\end{equation}
Thus any non-dissipative corridor has
\(W_I>\gamma_I/64\); its positive-flux band is radially broad before any
temporal estimate is used.
\end{corollary}

\begin{proof}
At the terminal endpoint, the first-doubling condition and spectral support
give
\[
 \frac12\bigl(
   \mathcal H_{\rho_I}(b)-\mathcal H_{\rho_I}(a)
 \bigr)
 \le\frac12\mathcal H_{\rho_I}(b)
 \le\frac{2R^2}{\rho_I}
 \le\frac{2R^2}{\gamma_I}.
\]
Under \eqref{eq:v19-thin-band-condition},
\(m_I=R^2/(8W_I)\ge8R^2/\gamma_I\), so the endpoint contribution is at
most \(m_I/4\).  Equations \eqref{eq:v19-cutoff-acquisition} and
\eqref{eq:v19-flux-selector} imply
\[
 \nu\mathcal D_{\rho_I,I}
 \ge m_I-\frac{m_I}{4}
 =\frac{3R^2}{32W_I}.
\]
\end{proof}

\begin{corollary}[Translation is exponential in parabolic exposure]
\label{cor:v19-exponential-translation}
On branch \textup{(T)},
\begin{equation}
 \boxed{\quad
 \frac{W_I}{\gamma_I}
 >\frac1{32}\exp(2\nu\gamma_I^2\tau).
 \quad}
 \label{eq:v19-exponential-width}
\end{equation}
In particular, if \(W_I/\gamma_I\le C\), then branch \textup{(T)} can occur
only when
\begin{equation}
 \nu\gamma_I^2\tau<\frac12\log(32C).
 \label{eq:v19-bounded-width-exposure}
\end{equation}
\end{corollary}

\begin{proof}
The left side of \eqref{eq:v19-translation-alternative} is nonnegative.
Therefore branch \textup{(T)} forces \(32W_I/\gamma_I>1\), and its displayed
inequality may be exponentiated.  The second assertion follows immediately.
\end{proof}

% =============================================================================
\section{What the three branches cost on a doubled-record chronology}
\label{sec:v19-sequence-ledgers}
% =============================================================================

\begin{lemma}[The intermediate source has a record-weighted Leray bound]
\label{lem:v19-mid-source-paid}
There is an absolute periodic constant \(C_{\rm mid}\) such that on a
first-doubling corridor based at record \(R\),
\begin{equation}
 \int_I\|A^{-1/2}\BNS(u(t))\|_2^2\,\dd t
 \le C_{\rm mid}R^2
 \int_I\|\Lambda u(t)\|_2^2\,\dd t.
 \label{eq:v19-mid-source-Leray}
\end{equation}
Consequently, for pairwise-disjoint doubled-record corridors \(I_n\),
\begin{equation}
 \sum_n\frac1{R_n^2}
 \int_{I_n}\|A^{-1/2}Q_{\rho_n}\BNS(u(t))\|_2^2\,\dd t
 \le\frac{C_{\rm mid}E_*}{2\nu}.
 \label{eq:v19-normalized-mid-global}
\end{equation}
\end{lemma}

\begin{proof}
The operator \(A^{-1/2}\PP\operatorname{div}\) has order zero.  The
Sobolev embedding \(H^{3/4}\hookrightarrow L^4\) and interpolation give
\begin{align*}
 \|A^{-1/2}\BNS(u)\|_2
 &\le C\|u\otimes u\|_2\\
 &\le C\|u\|_4^2\\
 &\le C\|\Lambda^{3/4}u\|_2^2
 \le C\|\Lambda^{1/2}u\|_2\|\Lambda u\|_2.
\end{align*}
On a first-hitting corridor,
\(\|\Lambda^{1/2}u(t)\|_2<2R\).  Square, integrate, and absorb the factor
four into \(C_{\rm mid}\) to prove
\eqref{eq:v19-mid-source-Leray}.  The projector is contractive in \(L^2\).
Divide by \(R_n^2\), sum over the disjoint corridors, and use
\(\nu\int\|\Lambda u\|_2^2\le E_*/2\).
\end{proof}

\begin{theorem}[Global payments returned by the V19 selector]
\label{thm:v19-global-branch-payments}
Apply Theorem~\ref{thm:v19-acquisition-trichotomy} on a doubled first-hitting
chronology.  Assign a card to \(\mathfrak D\) whenever
\textup{(D)} holds.  Among the remaining cards, assign it to
\(\mathfrak B\) if \textup{(T)} fails and to \(\mathfrak T\) otherwise.
Then
\begin{align}
 \boxed{\quad
 \sum_{n\in\mathfrak D}\frac{R_n^2}{W_n}
 \le8E_*;
 \quad}
 \label{eq:v19-global-diss-width}\\
 \boxed{\quad
 \sum_{n\in\mathfrak B}\frac1{W_n}
 \le\frac{8C_{\rm mid}E_*}{\nu^2}.
 \quad}
 \label{eq:v19-global-birth-width}
\end{align}
For the thin-band subfamily
\(\mathfrak D_{\rm thin}:=
 \{n:W_n\le\gamma_n/64\}\), the sharper estimate is
\begin{equation}
 \sum_{n\in\mathfrak D_{\rm thin}}
 \frac{R_n^2}{W_n}
 \le\frac{16}{3}E_*.
 \label{eq:v19-global-thin-width}
\end{equation}
Along every infinite subfamily of \(\mathfrak D\),
\begin{equation}
 \frac{W_n}{K_n}\longrightarrow\infty.
 \label{eq:v19-diss-branch-width-escape}
\end{equation}
Thus an infinite directly dissipative branch is possible only by making its
flux band supercanonical.
\end{theorem}

\begin{proof}
On \(\mathfrak D\), use
\eqref{eq:v19-dissipation-alternative}, the pointwise inequality
\(\mathcal D_{\rho_n,I_n}\le
\int_{I_n}\|\Lambda u\|_2^2\), disjointness, and the Leray energy budget:
\[
 \frac1{16}\sum_{n\in\mathfrak D}\frac{R_n^2}{W_n}
 \le\nu\int_{\cup I_n}\|\Lambda u\|_2^2
 \le\frac{E_*}{2}.
\]
This proves \eqref{eq:v19-global-diss-width}.  On \(\mathfrak B\), divide
\eqref{eq:v19-birth-alternative} by \(R_n^2\), sum, and apply
\eqref{eq:v19-normalized-mid-global}; this gives
\eqref{eq:v19-global-birth-width}.  The same calculation using
\eqref{eq:v19-thin-band-dissipation} proves
\eqref{eq:v19-global-thin-width}.

Every term of the convergent positive series in
\eqref{eq:v19-global-diss-width} tends to zero along an infinite subfamily.
Since \(K_n=R_n^2/E_*\),
\[
 \frac{R_n^2}{W_n}
 =\frac{E_*}{W_n/K_n},
\]
and \eqref{eq:v19-diss-branch-width-escape} follows.
\end{proof}

\begin{firewall}[The birth payment is too weak to close doubled records]
\label{fw:v19-birth-weight-summable}
The lower band-width bound \(W_n\ge K_n/2\) already gives
\begin{equation}
 \sum_n\frac1{W_n}
 \le2E_*\sum_n\frac1{R_n^2}<\infty.
 \label{eq:v19-inverse-width-automatic}
\end{equation}
Therefore \eqref{eq:v19-global-birth-width} is a valid global source payment
but not a contradiction.  Passing from kinetic endpoint birth to the
intermediate source loses the factor \(R_n^2\) that would be needed to make
the doubled-record sum coercive.  Calling branch \textup{(B)} ``paid''
without recording this normalization would conceal the same derivative gap
found in V7 and V14.
\end{firewall}

% =============================================================================
\section{Exact moving-tail telescope across non-dissipative cards}
\label{sec:v19-moving-tail}
% =============================================================================

The remaining endpoint gains cannot be summed at changing cutoffs by an
ordinary telescope.  The defect is nevertheless an explicit kinetic band,
not an unspecified error.

\begin{theorem}[Moving-tail storage ledger]
\label{thm:v19-moving-tail-ledger}
Let \(I_n=[t_n,t_{n+1}]\), \(p\le n<q\), be a consecutive block of
first-doubling corridors.  On every card choose \(\rho_n\) by
Lemma~\ref{lem:v19-quantile-band}, and suppose the direct dissipation
alternative \eqref{eq:v19-dissipation-alternative} fails throughout the
block.  Put
\begin{equation}
 J_n:=
 \bigl(
   \mathcal H_{\rho_{n-1}}(t_n)
   -\mathcal H_{\rho_n}(t_n)
 \bigr)_+,
 \qquad p<n<q.
 \label{eq:v19-junction-sweep}
\end{equation}
Then
\begin{align}
 \sum_{n=p}^{q-1}
 \bigl(
   \mathcal H_{\rho_n}(t_{n+1})
   -\mathcal H_{\rho_n}(t_n)
 \bigr)
 &=
 \mathcal H_{\rho_{q-1}}(t_q)
 -\mathcal H_{\rho_p}(t_p)
 \notag\\
 &\quad+
 \sum_{n=p+1}^{q-1}
 \bigl(
  \mathcal H_{\rho_{n-1}}(t_n)
  -\mathcal H_{\rho_n}(t_n)
 \bigr),
 \label{eq:v19-moving-tail-identity}\\
 \boxed{\quad
 \sum_{n=p}^{q-1}\frac{R_n^2}{8W_n}
 \le E_*+\sum_{n=p+1}^{q-1}J_n.
 \quad}
 \label{eq:v19-moving-tail-inequality}
\end{align}
If \(\rho_n\ge\rho_{n-1}\), then the junction term is exactly
\begin{equation}
 J_n=
 \|\mathbf1_{(\rho_{n-1},\rho_n]}(\Lambda)u(t_n)\|_2^2
 \le\frac{R_n^2}{\rho_{n-1}}.
 \label{eq:v19-upward-sweep-stock}
\end{equation}
If \(\rho_n<\rho_{n-1}\), then \(J_n=0\).  Hence for a nondecreasing
selected-cutoff block,
\begin{equation}
 \sum_{n=p}^{q-1}\frac{R_n^2}{8W_n}
 \le E_*+
 \sum_{n=p+1}^{q-1}\frac{R_n^2}{\rho_{n-1}}.
 \label{eq:v19-monotone-moving-tail}
\end{equation}
\end{theorem}

\begin{proof}
Failure of \eqref{eq:v19-dissipation-alternative} and the refined assertion
of Theorem~\ref{thm:v19-acquisition-trichotomy} give
\[
 \mathcal H_{\rho_n}(t_{n+1})
 -\mathcal H_{\rho_n}(t_n)
 \ge m_n=\frac{R_n^2}{8W_n}.
\]
Summing these differences and inserting the adjacent values
\(\mathcal H_{\rho_{n-1}}(t_n)\) proves the algebraic identity
\eqref{eq:v19-moving-tail-identity}.  The terminal-minus-initial term is at
most \(E_*\), and each signed junction term is at most its positive part.
This proves \eqref{eq:v19-moving-tail-inequality}.

If \(\rho_n\ge\rho_{n-1}\), orthogonality of the sharp spectral projectors
gives the equality in \eqref{eq:v19-upward-sweep-stock}.  At the first
hitting time \(t_n\),
\(\|\Lambda^{1/2}u(t_n)\|_2=R_n\), so spectral support gives its upper
bound.  If the cutoff decreases, \(Q_{\rho_{n-1}}\) is a subprojection of
\(Q_{\rho_n}\), making the difference nonpositive.  Summing the upper
bounds proves \eqref{eq:v19-monotone-moving-tail}.
\end{proof}

\begin{assessment}[The exact cross-record bottleneck after V19]
\label{ass:v19-cross-record-bottleneck}
There are now two independent escape variables.
\begin{enumerate}[label=\textup{(\alph*)},leftmargin=3.3em]
\item The per-card flux band can have
      \(W_n/K_n\to\infty\), making the direct energy debit
      \(R_n^2/W_n\) summable.  The V18 packet permits precisely this
      height--width trade.
\item On consecutive storage cards, the selected cutoff can move upward and
      expose the junction stock \(J_n\).  The same kinetic packet may be
      counted at many cutoffs through
      \eqref{eq:v19-tail-layer-stock}; only a one-use dynamical rule for the
      swept bands could pay \(\sum J_n\).
\end{enumerate}
If the cutoff were fixed, all \(J_n\) would vanish and
\eqref{eq:v19-moving-tail-inequality} would bound the total positive storage
by \(E_*\).  Cofinality makes the moving-cutoff defect unavoidable.  Applying
downward spectral variation to pay it would return exactly to the V4/V7
variation loop.  The next advance must therefore control the motion of the
quantile band itself, not merely re-estimate its flux height or terminal
stock.
\end{assessment}

\begin{assessment}[Packet audit of the new scales]
\label{ass:v19-packet-audit}
Proposition~\ref{prop:v18-packet-flux-profile} gives the V14/V18
concentrated packet a positive-flux region whose radius and width are both of
order \(N\), while \(\tau\asymp(RN^2)^{-1}\).  Hence the parabolic
exposure at the active radius is
\[
 \nu N^2\tau\asymp\frac{\nu}{R}\longrightarrow0,
\]
and the V19 debit \(R^2/W\), evaluated at that active width, has scale
\(R^2/N\), the packet kinetic-energy scale.  These are scale tests only; they
do not assert that the fixed V19 quantiles coincide with the compact interval
chosen in Proposition~\ref{prop:v18-packet-flux-profile}.  The packet is
therefore compatible with the small-exposure regime retained by
\eqref{eq:v19-translation-alternative}, and the source-birth alternative may
hold as well.  It does not contradict any V19 payment.  Conversely, it rules out strengthening
\eqref{eq:v19-dissipation-alternative} to a fixed order-\(E_*\) debit without
an independent upper bound on \(W_I/K\).
\end{assessment}

% =============================================================================
\section{Integrated V19 conclusion and V20 acquisition order}
\label{sec:v19-conclusion}
% =============================================================================

\begin{theorem}[What V19 proves]
\label{thm:v19-conclusion}
Relative to V1--V18, the following are proved.
\begin{enumerate}[label=\textup{(V19.\arabic*)},leftmargin=6.2em]
\item Integrated positive flux through a cutoff is exactly terminal
      exterior kinetic acquisition plus exterior viscous dissipation; radial
      cutoff integration repeats modal stock and is not an orthogonal energy
      sum.
\item Two positive-flux quantiles enclose an exact \(R^2/4\) of flux area,
      have width at least \(K/2\), and supply a cutoff with flux height
      \(R^2/(8W)\).
\item At that cutoff, every doubling corridor has a rigorous
      dissipation--birth--translation trichotomy
      \eqref{eq:v19-dissipation-alternative}--%
      \eqref{eq:v19-translation-alternative}.
\item A thin remote quantile band is paid directly by Leray dissipation;
      otherwise the rapid branch forces radial width exponential in
      parabolic exposure.
\item Directly dissipative cards obey the global debit
      \(\sum R_n^2/W_n<\infty\), hence their band widths are necessarily
      supercanonical on every infinite subfamily.
\item The intermediate-source birth branch is globally paid only after
      division by \(R_n^2\); its returned inverse-width series is already
      summable from record doubling and therefore does not close the cascade.
\item Consecutive endpoint-storage cards satisfy the exact moving-tail
      telescope \eqref{eq:v19-moving-tail-identity}; its sole cutoff-change
      defect is kinetic energy in the upward-swept frequency band.
\end{enumerate}
No full Navier--Stokes stability or global-regularity theorem is obtained.
\end{theorem}

\begin{proof}
Item \textup{(V19.1)} is Proposition~\ref{prop:v19-cutoff-acquisition}.
Item \textup{(V19.2)} is Lemma~\ref{lem:v19-quantile-band}.  Item
\textup{(V19.3)} is Theorem~\ref{thm:v19-acquisition-trichotomy}.  Item
\textup{(V19.4)} is Corollaries~\ref{cor:v19-thin-band-paid} and
\ref{cor:v19-exponential-translation}.  Items
\textup{(V19.5)--(V19.6)} are Lemma~\ref{lem:v19-mid-source-paid},
Theorem~\ref{thm:v19-global-branch-payments}, and
Firewall~\ref{fw:v19-birth-weight-summable}.  Item \textup{(V19.7)} is
Theorem~\ref{thm:v19-moving-tail-ledger}.
\end{proof}

\begin{programme}[V20 acquisition order]
\label{prog:v20}
\begin{enumerate}[label=\textup{(AC\arabic*)},leftmargin=4.8em]
\item Before extending the NS argument, perform the required fifth-version
      review of the latest Standard Model and Yang--Mills notes in
      \texttt{latex\_notes}.  Import only mechanisms whose hypotheses and
      paid ledgers survive the NS translation.
\item Treat \((\gamma_n,W_n,\rho_n)\) as a moving spectral front.  Seek an
      exact transport or ancestry law for the upward-swept junction stock
      \(J_n\), preserving complete triads and signs.
\item Separate bounded relative width \(W_n/\gamma_n\) from width escape.
      In the bounded branch, combine
      \eqref{eq:v19-bounded-width-exposure} with the recent-card source
      theorem; in the escaping branch, quantify how many record generations
      one extreme width can cover.
\item Use the common-terminal projection of V16 on record maxima of
      \(\rho_n\).  Determine whether heat age makes old swept-band stock
      one-use without reproducing the variable-cutoff junction defect.
\item Test any proposed payment on the V14/V18 packet.  A valid inequality
      must allow flux height and direct dissipation of order \(R^2/N\) and
      exposure of order \(\nu/R\).
\item Do not use the intermediate birth debit by itself: its
      \(1/W_n\) weight is automatically summable by
      \eqref{eq:v19-inverse-width-automatic}.
\item If swept-band control again reduces to unrestricted spectral
      variation, go upstream to the completed triad gap and look for a
      one-use signed transfer identity before introducing another norm.
\end{enumerate}
\end{programme}

\begin{assessment}[Position after V19]
The scalar exposure proposed in V18 has been replaced by the strongest
unconditional statement supported by the exact cutoff equation.  Positive
flux of selected height either spends actual Leray energy, appears as
source-created terminal stock after inherited heat decay, or requires a
radial band whose relative width grows exponentially with exposure.  The
first branch is globally paid; the second loses a summable inverse-width
factor; and the third is the cofinal translation escape already suggested by
the packet.  Across successive storage cards, the entire failure of
telescoping is now the explicit upward-swept kinetic band \(J_n\).  Closing
that moving-front ledger, together with excluding supercanonical width
escape, remains the next unresolved NS step.
\end{assessment}

% =============================================================================
\section*{Version V19 history}
\addcontentsline{toc}{section}{Version V19 history}
% =============================================================================

Corrected the proposed viscous-exposure route by retaining the exact terminal
storage term in the high-frequency energy balance.  Introduced two fixed-work
positive-flux quantiles and selected a cutoff of controlled flux height.
Proved the dissipation--source-birth--exponential-translation trichotomy,
including direct Leray payment of thin remote bands and global sequence
payments for the dissipative and intermediate-source branches.  Audited the
normalization and proved that the birth payment alone is automatically
summable.  Finally derived the exact moving-tail telescope and isolated its
only cutoff-change defect as kinetic stock in upward-swept frequency bands.
The full stability/global-regularity goal remains open.

% =============================================================================
% VERSION V20 APPEND
% =============================================================================

\section{V20: scheduled companion synchronization, paired-cutoff sweep
innovation, and orthogonal front ancestry}
\label{sec:v20-entry}
% =============================================================================

V19 reduced the cross-record failure of the selected high-tail telescope to
the kinetic stock \(J_n\) in frequency bands swept when the cutoff moves
upward.  V20 is the scheduled fifth-version synchronization point.  It first
freezes and audits the latest retained Standard-Model--Einstein and
Yang--Mills notebooks, then imports only two type-compatible assembly rules.
They lead to a same-trajectory, two-cutoff subtraction which cancels the
shared cofinal tail exactly, and to a complete direct-sum Gram for the
ancestry of disjoint swept bands.  The resulting estimates are rigorous, but
their inverse-frequency weight exposes a remaining derivative deficit rather
than closing global regularity.

% =============================================================================
\section{Required V20 companion-file audit}
\label{sec:v20-companion-audit}
% =============================================================================

The exact companion states read before choosing the V20 direction are
\begin{align}
 &\texttt{Renewal\_SM\_Einstein\_Continuum\_Research\_Notes\_V45.tex},
 \notag\\
 &\qquad
 \texttt{SHA256}=
 \texttt{96e67a87c12ba9119fdc6088ca22a8898af017403e8d68f56ab4b73ff131afab},
 \label{eq:v20-sm-audit-hash}\\
 &\texttt{Renewal\_Yang\_Mills\_Mass\_Gap\_Research\_Notes\_V26.tex},
 \notag\\
 &\qquad
 \texttt{SHA256}=
 \texttt{8a93f1980c3c8c6a40b9985ae1a0f70e499340f889c802b4814f5577f0728983}.
 \label{eq:v20-ym-audit-hash}
\end{align}
The pre-audit NS state is V19 with SHA--256
\begin{equation}
 \texttt{d572c9ba94416440807ca5184af4f03c6e9968fb94c74ed8f3ede7cdd9e150e1}.
 \label{eq:v20-ns-audit-hash}
\end{equation}

\begin{proposition}[Type-compatible content of the V20 synchronization]
\label{prop:v20-audit-lessons}
The following two assembly rules transfer to the NS programme.
\begin{enumerate}[label=\textup{(\roman*)},leftmargin=3.7em]
\item The SM V45 complete-matrix calculation retains the off-diagonal
      incidence through its Schur complement before testing a determinant
      margin.  For NS, inherited and source-born annular responses must
      likewise be assembled in one direct-sum Gram; their cross term may not
      be dropped before an estimate.
\item The YM V26 two-history subtraction cancels terms depending on the
      common terminal fibre while preserving the genuine likelihood
      innovation.  For NS, subtraction at two cutoffs with the same time
      interval cancels the common tail above the larger cutoff and leaves the
      exact signed band-injection flux.
\end{enumerate}
Neither companion file supplies an NS positivity, residence, source, or
regularity estimate.  The SM separation between determinant and locality
margins also requires V20 to keep annular orthogonality distinct from the
size of the returned source norm.  The YM warning that a positive covariance
does not control covariance minus direct curvature becomes the NS warning
that two positive flux values do not control their signed difference unless
the band-dissipation term is retained.
\end{proposition}

\begin{proof}
Item \textup{(i)} is the algebraic assembly order proved in the SM V45 full
Gram determinant and dynamic Schur formula; its notebook explicitly forbids
channelwise positivity before complete incidence.  Item \textup{(ii)} is the
same-terminal cancellation used in the YM V26 two-history score identity.
Both statements concern the order in which already typed components are
assembled.  Their field equations, measures, and claimed applications are
different, so no numerical or analytic estimate crosses programme
boundaries.  The last two assertions are the respective scope firewalls in
those appendices.
\end{proof}

\begin{firewall}[Companion synchronization is not theorem import]
\label{fw:v20-companion-scope}
Equations~\eqref{eq:v20-sm-audit-hash}--\eqref{eq:v20-ns-audit-hash} freeze
the evidence used here.  The arguments below are proved directly from the
Navier--Stokes spectral energy balance and Duhamel formula.  In particular,
no Einstein determinant floor, Yang--Mills covariance window, mass gap, or
continuum assertion is used as an NS hypothesis.
\end{firewall}

% =============================================================================
\section{Two cutoffs cancel the common cofinal tail}
\label{sec:v20-paired-cutoff}
% =============================================================================

For \(0<L<M\), define the sharp annular projector and its kinetic stock by
\begin{equation}
 \Pi_{L,M}:=P_M-P_L=\mathbf1_{(L,M]}(\Lambda),
 \qquad
 \mathcal E_{L,M}(t):=\|\Pi_{L,M}u(t)\|_2^2.
 \label{eq:v20-annular-projector-stock}
\end{equation}
For \(I=[a,b]\), put
\begin{equation}
 \mathcal D_{L,M;I}:=
 \int_a^b\|\Lambda\Pi_{L,M}u(t)\|_2^2\,\dd t.
 \label{eq:v20-annular-dissipation}
\end{equation}

\begin{theorem}[Exact paired-cutoff band balance]
\label{thm:v20-paired-cutoff-balance}
At every smooth time,
\begin{equation}
 \boxed{\quad
 \frac12\frac{\dd}{\dd t}\mathcal E_{L,M}(t)
 +\nu\|\Lambda\Pi_{L,M}u(t)\|_2^2
 =\Phi_L(t)-\Phi_M(t).
 \quad}
 \label{eq:v20-instant-band-balance}
\end{equation}
Consequently,
\begin{equation}
 \boxed{\quad
 F_I(L)-F_I(M)
 =\frac12\bigl(
   \mathcal E_{L,M}(b)-\mathcal E_{L,M}(a)
  \bigr)
 +\nu\mathcal D_{L,M;I}.
 \quad}
 \label{eq:v20-integrated-band-balance}
\end{equation}
The tail \(Q_Mu\), which is present in both \(Q_Lu\) and \(Q_Mu\), has
cancelled before any positive part or absolute value is taken.
\end{theorem}

\begin{proof}
Because the projectors are nested and orthogonal,
\begin{align*}
 \mathcal H_L-\mathcal H_M
 &=\mathcal E_{L,M},\\
 \|\Lambda Q_Lu\|_2^2-\|\Lambda Q_Mu\|_2^2
 &=\|\Lambda\Pi_{L,M}u\|_2^2.
\end{align*}
Subtract the two high-frequency balances
\eqref{eq:v4-high-energy-balance}, first at \(M\) from that at \(L\).
This proves \eqref{eq:v20-instant-band-balance}.  Integrating in time proves
\eqref{eq:v20-integrated-band-balance}.
\end{proof}

\begin{corollary}[Terminal band stock has an inherited-or-innovation
certificate]
\label{cor:v20-band-innovation}
Let
\begin{equation}
 J:=\mathcal E_{L,M}(b).
 \label{eq:v20-terminal-band-stock}
\end{equation}
At least one of
\begin{align}
 \boxed{\quad
 \mathcal E_{L,M}(a)\ge\frac J2;
 \quad}
 &\tag{I}\label{eq:v20-band-inherited}\\
 \boxed{\quad
 F_I(L)-F_I(M)
 \ge\frac J4+\nu\mathcal D_{L,M;I}
 \quad}
 &\tag{N}\label{eq:v20-band-innovation}
\end{align}
holds.  Thus a terminal swept band which was not already half present at the
left endpoint forces a positive \emph{paired} flux innovation.  Positivity of
either \(F_I(L)\) or \(F_I(M)\) separately is neither assumed nor concluded.
\end{corollary}

\begin{proof}
If \eqref{eq:v20-band-inherited} fails, then
\[
 \frac12\bigl(J-\mathcal E_{L,M}(a)\bigr)>\frac J4.
\]
Insert this in \eqref{eq:v20-integrated-band-balance}.
\end{proof}

\begin{corollary}[The V19 junction defect has exact previous-card ancestry]
\label{cor:v20-junction-ancestry}
Let \(I_{n-1}=[t_{n-1},t_n]\) be the previous first-hitting corridor and
suppose the selected cutoffs obey \(\rho_n>\rho_{n-1}\).  The V19 junction
stock is
\[
 J_n=\mathcal E_{\rho_{n-1},\rho_n}(t_n).
\]
At least one of
\begin{align}
 \mathcal E_{\rho_{n-1},\rho_n}(t_{n-1})
 &\ge\frac{J_n}{2},
 \label{eq:v20-junction-inherited}\\
 F_{I_{n-1}}(\rho_{n-1})-F_{I_{n-1}}(\rho_n)
 &\ge\frac{J_n}{4}
   +\nu\mathcal D_{\rho_{n-1},\rho_n;I_{n-1}}
 \label{eq:v20-junction-innovation}
\end{align}
holds.  Moreover,
\begin{equation}
 \mathcal E_{\rho_{n-1},\rho_n}(t_{n-1})
 \le\frac{R_{n-1}^2}{\rho_{n-1}}.
 \label{eq:v20-previous-band-capacity}
\end{equation}
Hence
\begin{equation}
 J_n>\frac{2R_{n-1}^2}{\rho_{n-1}}
 \quad\Longrightarrow\quad
 F_{I_{n-1}}(\rho_{n-1})-F_{I_{n-1}}(\rho_n)
 \ge\frac{J_n}{4}.
 \label{eq:v20-large-sweep-innovation}
\end{equation}
\end{corollary}

\begin{proof}
Apply Corollary~\ref{cor:v20-band-innovation} with
\((L,M,a,b)=(\rho_{n-1},\rho_n,t_{n-1},t_n)\).  At the left endpoint,
\[
 \|\Lambda^{1/2}u(t_{n-1})\|_2=R_{n-1}.
\]
Spectral support above \(\rho_{n-1}\) gives
\eqref{eq:v20-previous-band-capacity}.  Under the strict hypothesis in
\eqref{eq:v20-large-sweep-innovation}, the inherited alternative is
impossible, and the innovation alternative applies.
\end{proof}

\begin{firewall}[A paired innovation is signed band transfer]
\label{fw:v20-paired-sign}
The quantity in \eqref{eq:v20-junction-innovation} is
\(F(\rho_{n-1})-F(\rho_n)\), not the sum of their positive parts.  For each
complete triad it is the integral of the difference of the two step profiles
in \eqref{eq:v18-triad-flux-profile}; equivalently, it is net nonlinear
energy injected into the intervening annulus.  Replacing it by
\((F(\rho_{n-1}))_++(F(\rho_n))_+\) destroys the common-tail cancellation
and recreates cutoff multiplicity.
\end{firewall}

% =============================================================================
\section{Complete-Gram ancestry for pairwise disjoint swept bands}
\label{sec:v20-complete-gram}
% =============================================================================

The paired balance identifies how one band is acquired.  A cofinal front
produces many disjoint bands, observed at different terminal times.  Their
Duhamel responses can still be assembled without time multiplicity because
their spectral ranges are orthogonal.

\begin{theorem}[Common-root orthogonal annular ancestry]
\label{thm:v20-common-root-annuli}
Let \(s<B<T_*\).  Let \(\Pi_j\), \(j\in\mathcal J\), be a finite or countable
family of pairwise orthogonal sharp spectral projectors, with
\begin{equation}
 \operatorname{supp}\Pi_j\subset\{0<\Lambda\le U_j\},
 \qquad
 U_j\ge U_{\min}>0,
 \label{eq:v20-annular-upper-edge}
\end{equation}
and choose \(b_j\in[s,B]\).  Define
\begin{align}
 v_j&:=\Pi_ju(b_j),
 \label{eq:v20-terminal-annular-vector}\\
 h_j&:=e^{-\nu A(b_j-s)}\Pi_ju(s),
 \label{eq:v20-inherited-annular-vector}\\
 r_j&:=\int_s^{b_j}
       e^{-\nu A(b_j-t)}\Pi_j\BNS(u(t))\,\dd t.
 \label{eq:v20-born-annular-vector}
\end{align}
Then \(v_j=h_j-r_j\).  In the weighted direct sum
\(\bigoplus_jU_j^{-1/2}L^2\), put
\begin{align}
 \mathscr V&:=\sum_j\frac{\|v_j\|_2^2}{U_j},
 &\mathscr H&:=\sum_j\frac{\|h_j\|_2^2}{U_j},
 \notag\\
 \mathscr S&:=\sum_j\frac{\|r_j\|_2^2}{U_j},
 &\mathscr C&:=\sum_j\frac{\operatorname{Re}\langle h_j,r_j\rangle}{U_j}.
 \label{eq:v20-complete-gram-components}
\end{align}
The complete Gram identity and Schur bound are
\begin{align}
 \boxed{\quad
 \mathscr V=\mathscr H-2\mathscr C+\mathscr S,
 \qquad
 |\mathscr C|^2\le\mathscr H\mathscr S;
 \quad}
 \label{eq:v20-complete-gram}\\
 \sqrt{\mathscr V}
 \le\sqrt{\mathscr H}+\sqrt{\mathscr S}.
 \label{eq:v20-Gram-triangle}
\end{align}
Furthermore,
\begin{align}
 \mathscr H
 &\le\frac{E_*}{U_{\min}},
 \label{eq:v20-inherited-orthogonal-bound}\\
 \mathscr S
 &\le\frac1{2\nu}
 \int_s^B\|A^{-3/4}\BNS(u(t))\|_2^2\,\dd t
 \le\frac{C_{\rm w}E_*^2}{2\nu^2},
 \label{eq:v20-born-weak-bound}\\
 \boxed{\quad
 \mathscr V
 \le
 \left(
  \sqrt{\frac{E_*}{U_{\min}}}
  +\frac{\sqrt{C_{\rm w}}E_*}{\sqrt2\,\nu}
 \right)^2.
 \quad}
 \label{eq:v20-annular-ancestry-bound}
\end{align}
Here \(C_{\rm w}\) is the absolute periodic constant in the V7 weak-source
estimate.
\end{theorem}

\begin{proof}
Duhamel's formula gives \(v_j=h_j-r_j\).  Expanding its squared norm in the
weighted direct sum proves the equality in
\eqref{eq:v20-complete-gram}.  Cauchy--Schwarz in that same direct sum gives
the Schur bound, and hence \eqref{eq:v20-Gram-triangle}.  This step retains
the full cross term; no sign is assigned to \(\mathscr C\).

Heat contraction, spectral orthogonality, and \(U_j\ge U_{\min}\) give
\[
 \mathscr H
 \le\frac1{U_{\min}}\sum_j\|\Pi_ju(s)\|_2^2
 \le\frac{E_*}{U_{\min}}.
\]

For the born term, write
\(g_j=A^{-3/4}\Pi_j\BNS(u)\).  On an
\(A^{1/2}\)-eigenmode of frequency \(\lambda\le U_j\), time
Cauchy--Schwarz gives
\begin{align*}
 \left|
  \int_s^{b_j}\lambda^{3/2}
   e^{-\nu\lambda^2(b_j-t)}g_{j,\lambda}(t)\,\dd t
 \right|^2
 &\le
 \left(
  \int_s^{b_j}\lambda^3
   e^{-2\nu\lambda^2(b_j-t)}\,\dd t
 \right)
 \int_s^{b_j}|g_{j,\lambda}(t)|^2\,\dd t\\
 &\le\frac{U_j}{2\nu}
 \int_s^{b_j}|g_{j,\lambda}(t)|^2\,\dd t.
\end{align*}
After summing modes and dividing by \(U_j\),
\[
 \frac{\|r_j\|_2^2}{U_j}
 \le\frac1{2\nu}
 \int_s^{b_j}\|A^{-3/4}\Pi_j\BNS(u(t))\|_2^2\,\dd t.
\]
Sum in \(j\) and interchange the nonnegative integrals.  At each time \(t\),
only indices with \(t\le b_j\) remain, and their projectors are still
orthogonal.  Therefore
\[
 \sum_{j:t\le b_j}
 \|A^{-3/4}\Pi_j\BNS(u(t))\|_2^2
 \le\|A^{-3/4}\BNS(u(t))\|_2^2.
\]
This proves the first inequality in \eqref{eq:v20-born-weak-bound}.  Apply
the time-shifted proof of the V7 weak-source budget on \([s,B]\); its
initial energy is at most \(E_*\).  This gives the second inequality.  Combine
\eqref{eq:v20-Gram-triangle},
\eqref{eq:v20-inherited-orthogonal-bound}, and
\eqref{eq:v20-born-weak-bound} to prove
\eqref{eq:v20-annular-ancestry-bound}.  Monotone convergence handles a
countable family.
\end{proof}

\begin{remark}[The two independent margins]
\label{rem:v20-two-margins}
Pairwise orthogonality is the V20 locality/incidence margin: it prevents one
Fourier mode from entering two annular components.  The factor \(U_j^{-1}\)
is the analytic edge margin returned by the globally paid
\(A^{-3/4}\BNS(u)\) source.  The first does not improve the second.  This is
the exact NS counterpart of keeping the SM V45 determinant and locality rows
separate.
\end{remark}

\begin{corollary}[One-use weak ancestry of a monotone swept front]
\label{cor:v20-monotone-front-ancestry}
Let
\[
 0<\rho_0<\rho_1<\cdots,
 \qquad
 \Pi_j:=\mathbf1_{(\rho_{j-1},\rho_j]}(\Lambda),
\]
and let \(b_j\in[s,B]\).  Then
\begin{equation}
 \sum_j\frac{
  \|\mathbf1_{(\rho_{j-1},\rho_j]}(\Lambda)u(b_j)\|_2^2
 }{\rho_j}
 \le
 \left(
  \sqrt{\frac{E_*}{\rho_1}}
  +\frac{\sqrt{C_{\rm w}}E_*}{\sqrt2\,\nu}
 \right)^2.
 \label{eq:v20-monotone-front-weak-ledger}
\end{equation}
Each swept frequency belongs to exactly one summand even though the source
time intervals \([s,b_j]\) overlap.
\end{corollary}

\begin{proof}
Use Theorem~\ref{thm:v20-common-root-annuli} with
\(U_j=\rho_j\) and \(U_{\min}=\rho_1\).
\end{proof}

\begin{firewall}[The complete Gram returns only inverse-frequency stock]
\label{fw:v20-inverse-frequency}
The cross term in \eqref{eq:v20-complete-gram} is now handled exactly, and
annular source components are one-use.  Nevertheless the paid conclusion is
\(J_j/\rho_j\), not \(J_j\).  On a doubled-record front with
\(\rho_j\ge K_{n_j}=R_{n_j}^2/E_*\), the elementary energy cap already gives
\begin{equation}
 \sum_j\frac{J_j}{\rho_j}
 \le E_*\sum_j\frac1{\rho_j}
 \le E_*^2\sum_j\frac1{R_{n_j}^2}<\infty.
 \label{eq:v20-weak-sweep-automatic}
\end{equation}
Thus the complete-Gram import removes multiplicity but does not remove the
one-derivative gap.  Claiming unweighted swept-stock control from
\eqref{eq:v20-monotone-front-weak-ledger} would discard the factor forced by
the heat kernel.
\end{firewall}

% =============================================================================
\section{Extreme quantile width covers only logarithmically many generations}
\label{sec:v20-generation-coverage}
% =============================================================================

Write the V19 quantiles on the \(n\)-th first-hitting corridor as
\(\gamma_n,\Gamma_n,W_n\), and recall
\[
 K_n=\frac{R_n^2}{E_*}=4^nK_0,
 \qquad
 \gamma_n\ge K_n,
 \qquad
 \Gamma_n=\gamma_n+W_n<\infty.
\]
Define the outer overshoot
\begin{equation}
 \mathfrak o_n:=\frac{\Gamma_n}{K_n}\ge\frac32.
 \label{eq:v20-outer-overshoot}
\end{equation}

\begin{theorem}[Generation-coverage bound and a disjoint cofinal front]
\label{thm:v20-generation-coverage}
If \(m>n\) and
\begin{equation}
 \gamma_m\le\Gamma_n,
 \label{eq:v20-generation-covered}
\end{equation}
then
\begin{equation}
 m-n\le\log_4\mathfrak o_n.
 \label{eq:v20-generation-coverage-bound}
\end{equation}
Starting from any \(n_0\), recursively define
\begin{equation}
 n_{j+1}:=\min\{m>n_j:\gamma_m>\Gamma_{n_j}\}.
 \label{eq:v20-front-recursion}
\end{equation}
This index is finite and obeys
\begin{equation}
 1\le n_{j+1}-n_j
 \le\lfloor\log_4\mathfrak o_{n_j}\rfloor+1.
 \label{eq:v20-front-gap}
\end{equation}
The selected quantile bands are strictly radially ordered:
\begin{equation}
 \Gamma_{n_j}<\gamma_{n_{j+1}},
 \qquad
 [\gamma_{n_j},\Gamma_{n_j}]
 \cap[\gamma_{n_{j+1}},\Gamma_{n_{j+1}}]=\varnothing.
 \label{eq:v20-disjoint-quantile-front}
\end{equation}
In particular, any choices
\(\rho_{n_j}\in(\gamma_{n_j},\Gamma_{n_j})\) form a strictly increasing
cofinal sequence, and the swept bands
\begin{equation}
 \mathbf1_{(\rho_{n_j},\rho_{n_{j+1}}]}(\Lambda)
 \label{eq:v20-front-swept-bands}
\end{equation}
are pairwise orthogonal.
\end{theorem}

\begin{proof}
If \eqref{eq:v20-generation-covered} holds, then
\[
 4^{m-n}K_n=K_m\le\gamma_m\le\Gamma_n
 =\mathfrak o_nK_n.
\]
Cancel \(K_n\) and take the base-four logarithm to prove
\eqref{eq:v20-generation-coverage-bound}.

Put \(d=\lfloor\log_4\mathfrak o_n\rfloor+1\).  Then
\(4^d>\mathfrak o_n\), so
\[
 \gamma_{n+d}\ge K_{n+d}
 =4^dK_n>\Gamma_n.
\]
Thus the set in \eqref{eq:v20-front-recursion} is nonempty and its minimum
obeys \eqref{eq:v20-front-gap}.  The first inequality in
\eqref{eq:v20-disjoint-quantile-front} is the defining property of the
recursion; the disjointness and strict increase follow.  Since
\(\rho_{n_j}\ge\gamma_{n_j}\ge K_{n_j}\to\infty\), the selected cutoffs are
cofinal.  Consecutive intervals of a strictly increasing scalar sequence are
disjoint up to endpoints, proving orthogonality of
\eqref{eq:v20-front-swept-bands}.
\end{proof}

\begin{corollary}[Each cofinal front jump has a complete innovation
alternative]
\label{cor:v20-front-jump-innovation}
Let \(n_j\) be the front subsequence, let
\[
 L_j:=\rho_{n_j},
 \qquad
 M_j:=\rho_{n_{j+1}},
 \qquad
 \widehat I_j:=[t_{n_j},t_{n_{j+1}}],
\]
and define
\begin{equation}
 \widehat J_j:=
 \|\mathbf1_{(L_j,M_j]}(\Lambda)u(t_{n_{j+1}})\|_2^2.
 \label{eq:v20-front-jump-stock}
\end{equation}
At least one of
\begin{align}
 \|\mathbf1_{(L_j,M_j]}(\Lambda)u(t_{n_j})\|_2^2
 &\ge\frac{\widehat J_j}{2},
 \label{eq:v20-front-jump-inherited}\\
 F_{\widehat I_j}(L_j)-F_{\widehat I_j}(M_j)
 &\ge\frac{\widehat J_j}{4}
   +\nu\mathcal D_{L_j,M_j;\widehat I_j}
 \label{eq:v20-front-jump-innovated}
\end{align}
holds.  The inherited term is bounded by
\begin{equation}
 \|\mathbf1_{(L_j,M_j]}(\Lambda)u(t_{n_j})\|_2^2
 \le\frac{R_{n_j}^2}{L_j}.
 \label{eq:v20-front-jump-capacity}
\end{equation}
The bands in \eqref{eq:v20-front-jump-stock} are pairwise orthogonal, so
their terminal values also satisfy the one-use weak ancestry ledger
\eqref{eq:v20-monotone-front-weak-ledger}.
\end{corollary}

\begin{proof}
Apply Corollary~\ref{cor:v20-band-innovation} on
\(\widehat I_j\) with \(L=L_j\) and \(M=M_j\).
The critical record at \(t_{n_j}\) and spectral support above \(L_j\) prove
\eqref{eq:v20-front-jump-capacity}.  For every finite prefix, the final assertion follows from
\eqref{eq:v20-disjoint-quantile-front} and
Corollary~\ref{cor:v20-monotone-front-ancestry}, with common root
\(s=t_{n_0}\), observation times \(b_j=t_{n_{j+1}}\), and
\(B\) equal to the last observation time in that prefix.  The bound is
independent of the prefix.  Monotone convergence and the terminal weak-source
budget then give the asserted ledger for the full cofinal sequence.
\end{proof}

\begin{assessment}[What generation coverage does and does not pay]
\label{ass:v20-generation-status}
An outer overshoot \(\mathfrak o_n\) can hide at most
\(\log_4\mathfrak o_n\) subsequent canonical record scales.  Hence one
remote V19 band cannot cover an infinite terminal chronology, and
Theorem~\ref{thm:v20-generation-coverage} extracts infinitely many disjoint
cofinal front bands.  This is genuine incidence/locality progress.

It is not an amplitude payment.  The number of covered generations grows
only logarithmically, whereas the V19 direct debit decays like
\[
 \frac{R_n^2}{W_n}
 =\frac{E_*}{W_n/K_n}.
\]
An arbitrarily large overshoot may therefore be both cheap in the Leray
ledger and capable of skipping many generations.  Orthogonality then pays
only the inverse-frequency quantity
\eqref{eq:v20-weak-sweep-automatic}.  A closure theorem still needs either
an unweighted one-use estimate for the paired innovations or a dynamical
upper bound on the overshoots.
\end{assessment}

% =============================================================================
\section{Integrated V20 conclusion and V21 acquisition order}
\label{sec:v20-conclusion}
% =============================================================================

\begin{theorem}[What V20 proves]
\label{thm:v20-conclusion}
Relative to V1--V19, the following are proved.
\begin{enumerate}[label=\textup{(V20.\arabic*)},leftmargin=6.2em]
\item The required SM V45 and YM V26 review is frozen by filename and hash.
      It imports only complete-assembly and same-terminal-subtraction rules,
      with a strict cross-program type firewall.
\item Subtracting the high-frequency balances at two cutoffs cancels their
      common cofinal tail and gives the exact signed annular balance
      \eqref{eq:v20-integrated-band-balance}.
\item Every upward V19 junction stock is either half inherited on the
      previous card or forces a positive paired-cutoff flux innovation,
      including the annular viscous term.
\item Pairwise disjoint swept bands admit a common-root Duhamel assembly in
      one complete Gram.  The inherited terms are orthogonal and the
      overlapping source histories are one-use after annular projection.
\item The globally paid weak source controls
      \(\sum J_j/\rho_j\), but the inverse frequency is unavoidable in the
      proved heat-kernel estimate and is already summable on doubled records.
\item A quantile overshoot covers at most logarithmically many later record
      generations.  There is an infinite cofinal subsequence with disjoint
      quantile bands and pairwise orthogonal swept-front bands.
\item Every such front jump has an exact inherited-or-paired-innovation
      certificate and a complete weak-source ancestry ledger.
\end{enumerate}
No full Navier--Stokes stability or global-regularity theorem is obtained.
\end{theorem}

\begin{proof}
Item \textup{(V20.1)} is
Proposition~\ref{prop:v20-audit-lessons} and
Firewall~\ref{fw:v20-companion-scope}.  Item \textup{(V20.2)} is
Theorem~\ref{thm:v20-paired-cutoff-balance}.  Item \textup{(V20.3)} is
Corollaries~\ref{cor:v20-band-innovation}--%
\ref{cor:v20-junction-ancestry}.  Items
\textup{(V20.4)--(V20.5)} are
Theorem~\ref{thm:v20-common-root-annuli},
Corollary~\ref{cor:v20-monotone-front-ancestry}, and
Firewall~\ref{fw:v20-inverse-frequency}.  Items
\textup{(V20.6)--(V20.7)} are
Theorem~\ref{thm:v20-generation-coverage} and
Corollary~\ref{cor:v20-front-jump-innovation}.
\end{proof}

\begin{programme}[V21 acquisition order]
\label{prog:v21}
\begin{enumerate}[label=\textup{(AC\arabic*)},leftmargin=4.8em]
\item Use the exact paired flux
      \(F_{\widehat I_j}(L_j)-F_{\widehat I_j}(M_j)\) as the front
      innovation.  Do not replace it by separate positive flux margins.
\item Decompose the complete triad profile only after restricting its
      transfer to the front band \((L_j,M_j]\).  Seek a signed one-crossing
      count which prevents the same kinetic packet from paying many
      successive front innovations.
\item Separate the two remaining escape branches:
      inherited front-band stock
      \(R_{n_j}^2/L_j\), and source-born stock carrying the
      inverse-frequency loss in
      \eqref{eq:v20-born-weak-bound}.
\item On the inherited branch, propagate the disjoint front bands from one
      common root and retain the complete cross term
      \eqref{eq:v20-complete-gram}.  Test whether heat age supplies an
      additional factor without assuming residence.
\item On the born branch, interpolate the weak
      \(A^{-3/4}\BNS(u)\) ledger with the V17 critical-source ladder only
      through a proved upper estimate.  Do not treat two separately exact
      ledgers as contemporaneous incidence.
\item Quantify front gaps using
      \eqref{eq:v20-front-gap}.  Any summation over skipped records must
      charge the actual overshoot \(\mathfrak o_n\), not merely count
      subsequence indices.
\item Test every candidate unweighted sweep inequality on the V14/V18
      concentrated packet, whose high-frequency kinetic stock is
      \(R^2/N\) and whose weak-source action is \(R^3/N^2\).
\item The next required SM/YM synchronization is V25.
\end{enumerate}
\end{programme}

\begin{assessment}[Position after V20]
The V19 moving-tail defect now has a literal source equation.  At two
cutoffs the shared remote tail cancels, leaving signed nonlinear injection
into exactly the swept annulus.  On a cofinal front, those annuli can be made
pairwise disjoint, and their inherited and source-born histories assemble in
one complete Gram without source multiplicity.  The paid norm, however,
controls only inverse-frequency kinetic stock.  Extreme quantile widths can
also skip logarithmically many generations while making the V19 Leray debit
small.  The remaining bottleneck is therefore sharper than before: prove an
unweighted one-crossing law for the paired annular innovations, or prove a
dynamical upper bound on quantile overshoot.  Neither has yet been obtained.
\end{assessment}

% =============================================================================
\section*{Version V20 history}
\addcontentsline{toc}{section}{Version V20 history}
% =============================================================================

Performed the scheduled review of the exact SM V45 and YM V26 companion
states and froze their hashes.  Translated their complete-assembly and
same-terminal-subtraction lessons into direct NS proofs.  Derived the exact
paired-cutoff annular balance and an inherited-or-innovation certificate for
every upward swept junction.  Built a common-root complete Gram for
pairwise-disjoint annular Duhamel histories and proved a one-use weak-source
ledger, while recording its unavoidable inverse-frequency normalization.
Finally proved that a finite quantile overshoot covers only logarithmically
many record generations and extracted a disjoint cofinal spectral front.
The full stability/global-regularity goal remains open.

% =============================================================================
% VERSION V21 APPEND
% =============================================================================

\section{V21: triad one-crossing incidence and the critical amplitude
barrier}
\label{sec:v21-entry}
% =============================================================================

V20 produced an infinite cofinal front whose swept frequency bands are
pairwise disjoint.  It also identified the correct signed innovation across
one front band:
\[
 F_I(L)-F_I(M)
 =\int_I\bigl(\Phi_L-\Phi_M\bigr)\,\dd t.
\]
The remaining V21 question is whether complete-triad incidence makes these
innovations one-use in a quantitatively coercive sense.  The answer has two
parts.  First, the radial incidence really is finite: each modal leg of a
complete triad belongs to at most one disjoint front band, so a triad can
enter at most three front tiles.  Second, every Sobolev pairing from the
globally paid \(H^{-3/2}\) source to the record-weighted \(H^{-1}\) source
returns either one inverse frequency or one critical-record factor.  Both
weights are automatically summable on the doubled chronology.  Thus V21
closes the incidence-count question and proves that incidence alone does not
close the amplitude ledger.

% =============================================================================
\section{Complete triads cross a disjoint radial front at most three times}
\label{sec:v21-triad-one-crossing}
% =============================================================================

For \(0<L<M\), write
\begin{equation}
 \Psi_{L,M}(t):=\Phi_L(t)-\Phi_M(t).
 \label{eq:v21-band-injection-rate}
\end{equation}
This is the nonlinear injection rate in
\eqref{eq:v20-instant-band-balance}.

\begin{proposition}[Exact modal form of paired-cutoff injection]
\label{prop:v21-modal-band-injection}
For every smooth time,
\begin{equation}
 \boxed{\quad
 \Psi_{L,M}(t)
 =-\sum_{L<|k|\le M}T_k(t)
 =-\langle
   \Pi_{L,M}\BNS(u(t)),\Pi_{L,M}u(t)
  \rangle.
 \quad}
 \label{eq:v21-modal-band-injection}
\end{equation}
For a complete unordered triad with ordered radii
\(\lambda_1\le\lambda_2\le\lambda_3\), its contribution is
\begin{equation}
 \Psi_{L,M}^{\triangle}(t)
 =-\sum_{\{i:L<\lambda_i\le M\}}T_i^\triangle(t).
 \label{eq:v21-triad-band-injection}
\end{equation}
\end{proposition}

\begin{proof}
By definition,
\[
 \Phi_M-\Phi_L
 =\sum_{L<|k|\le M}T_k.
\]
Parseval and orthogonality of \(\Pi_{L,M}\) give the second equality in
\eqref{eq:v21-modal-band-injection}.  Restrict the modal sum to one complete
triad to obtain \eqref{eq:v21-triad-band-injection}.
\end{proof}

\begin{theorem}[Three-leg one-crossing theorem]
\label{thm:v21-three-leg-one-crossing}
Let
\[
 \mathcal B_j=(L_j,M_j]\subset(0,\infty),
 \qquad j\in\mathcal J,
\]
be pairwise disjoint radial bands, and let \(I_j\) be arbitrary smooth time
intervals.  Define the space--time front innovations
\begin{equation}
 \mathfrak I_j:=
 \int_{I_j}\Psi_{L_j,M_j}(t)\,\dd t.
 \label{eq:v21-front-innovation}
\end{equation}
For any fixed complete Fourier triad, the set of indices \(j\) for which its
contribution to \(\mathfrak I_j\) can be nonzero has cardinality at most
three.  More precisely, if
\begin{equation}
 \mathfrak I_j^\triangle
 :=-\sum_{\{i:\lambda_i\in\mathcal B_j\}}
   \int_{I_j}T_i^\triangle(t)\,\dd t,
 \label{eq:v21-triad-front-innovation}
\end{equation}
then
\begin{equation}
 \#\{j:\mathfrak I_j^\triangle\ne0\}\le3.
 \label{eq:v21-three-band-incidence}
\end{equation}
If the union of the bands contains all three radii and the same time interval
\(I\) is used for every band, then
\begin{equation}
 \sum_j\mathfrak I_j^\triangle=0.
 \label{eq:v21-complete-front-cancellation}
\end{equation}
\end{theorem}

\begin{proof}
Each fixed radius \(\lambda_i\) belongs to at most one member of a pairwise
disjoint family.  Hence the three modal legs can populate at most three
indices in \eqref{eq:v21-triad-front-innovation}, proving
\eqref{eq:v21-three-band-incidence}.  Under the additional hypothesis, each
leg occurs exactly once and every time integral is over \(I\).  Therefore
\[
 \sum_j\mathfrak I_j^\triangle
 =-\int_I
  (T_1^\triangle+T_2^\triangle+T_3^\triangle)\,\dd t=0
\]
by \eqref{eq:v18-triad-energy-cancellation}.
\end{proof}

\begin{corollary}[The V20 cofinal front has no radial triad multiplicity]
\label{cor:v21-front-no-radial-multiplicity}
Use the cofinal subsequence \(n_j\) from
Theorem~\ref{thm:v20-generation-coverage}, with
\[
 \mathcal B_j=(\rho_{n_j},\rho_{n_{j+1}}],
 \qquad
 I_j=[t_{n_j},t_{n_{j+1}}].
\]
Every fixed complete triad enters at most three of these space--time tiles.
Thus any failure to sum the positive paired innovations is an amplitude or
time-variation failure, not repeated radial incidence of one triad leg.
\end{corollary}

\begin{proof}
The bands are pairwise disjoint by
\eqref{eq:v20-front-swept-bands}.  Apply
Theorem~\ref{thm:v21-three-leg-one-crossing}.
\end{proof}

\begin{firewall}[Finite incidence does not give a positive sum]
\label{fw:v21-incidence-no-sign}
Equation~\eqref{eq:v21-complete-front-cancellation} shows the obstruction
sharply.  When all legs are retained on one time interval, their signed band
injections cancel.  On the moving front the time windows are different, so
positive injection into one band can be followed by negative injection from
another.  The incidence bound three controls how often a triad is routed; it
does not bound the total variation of its modal transfers.  Taking
\(|T_i^\triangle|\) before using
\eqref{eq:v18-triad-energy-cancellation} would discard the cancellation that
made the incidence statement exact.
\end{firewall}

% =============================================================================
\section{A continuum of source norms all returns the same crossing deficit}
\label{sec:v21-Sobolev-crossing}
% =============================================================================

Let \(I=[a,b]\), let \(0<L<M\), and abbreviate
\begin{equation}
 \mathfrak I:=F_I(L)-F_I(M),
 \qquad
 \mathcal D:=\mathcal D_{L,M;I}.
 \label{eq:v21-one-band-data}
\end{equation}
For
\begin{equation}
 \frac12\le\sigma\le\frac34,
 \qquad
 \theta:=4\sigma-2\in[0,1],
 \label{eq:v21-sigma-theta}
\end{equation}
define the complete annular source action
\begin{equation}
 \mathscr S_{\sigma;L,M;I}
 :=\int_I
   \|A^{-\sigma}\Pi_{L,M}\BNS(u(t))\|_2^2\,\dd t.
 \label{eq:v21-annular-source-action}
\end{equation}
The endpoints are the intermediate \(H^{-1}\) source at
\(\sigma=1/2\) and the globally paid \(H^{-3/2}\) source at
\(\sigma=3/4\).

\begin{theorem}[Sharp source cost of a paired band innovation]
\label{thm:v21-source-cost}
Suppose that a terminal annular stock \(J>0\) is on the innovation branch
\begin{equation}
 \mathfrak I\ge\frac J4+\nu\mathcal D.
 \label{eq:v21-innovation-hypothesis}
\end{equation}
Then \(\mathcal D>0\) and, for every
\(\sigma\in[1/2,3/4]\),
\begin{align}
 \mathfrak I^2
 &\le
 M^\theta\,
 \mathscr S_{\sigma;L,M;I}\,
 \mathcal D,
 \label{eq:v21-flux-source-duality}\\
 \boxed{\quad
 \mathscr S_{\sigma;L,M;I}
 \ge\frac{\nu J}{M^\theta}.
 \quad}
 \label{eq:v21-source-crossing-debit}
\end{align}
In particular,
\begin{align}
 \int_I\|A^{-1/2}\Pi_{L,M}\BNS(u)\|_2^2\,\dd t
 &\ge\nu J,
 \label{eq:v21-mid-source-crossing}\\
 \int_I\|A^{-3/4}\Pi_{L,M}\BNS(u)\|_2^2\,\dd t
 &\ge\frac{\nu J}{M}.
 \label{eq:v21-weak-source-crossing}
\end{align}
\end{theorem}

\begin{proof}
Proposition~\ref{prop:v21-modal-band-injection} and self-adjointness of
powers of \(A\) give
\[
 \Psi_{L,M}(t)
 =-\langle
   A^{-\sigma}\Pi_{L,M}\BNS(u(t)),
   A^\sigma\Pi_{L,M}u(t)
  \rangle.
\]
On the annular support \(\lambda\le M\), and
\(\theta=4\sigma-2\),
\[
 \|A^\sigma\Pi_{L,M}u\|_2^2
 =\sum_{L<|k|\le M}|k|^{4\sigma}|\widehat u_k|^2
 \le
 M^\theta\|\Lambda\Pi_{L,M}u\|_2^2.
\]
Time Cauchy--Schwarz proves
\eqref{eq:v21-flux-source-duality}, with
\(|\mathfrak I|\) on the left; the hypothesis makes
\(\mathfrak I>0\).

If \(\mathcal D=0\), the right side of
\eqref{eq:v21-flux-source-duality} vanishes, contradicting
\(\mathfrak I\ge J/4>0\).  Thus \(\mathcal D>0\), and
\begin{align*}
 \mathscr S_{\sigma;L,M;I}
 &\ge\frac{\mathfrak I^2}{M^\theta\mathcal D}\\
 &\ge\frac{(J/4+\nu\mathcal D)^2}
            {M^\theta\mathcal D}\\
 &\ge\frac{4(J/4)\nu}{M^\theta}
 =\frac{\nu J}{M^\theta}.
\end{align*}
The penultimate line is
\((x+y)^2\ge4xy\).  Setting \(\theta=0\) and \(\theta=1\) gives the endpoint
statements.
\end{proof}

\begin{remark}[The annular viscous term is essential]
\label{rem:v21-viscous-completion}
If one retained only \(\mathfrak I\ge J/4\), the quotient
\(\mathfrak I^2/\mathcal D\) would have no positive lower bound as
\(\mathcal D\to\infty\).  The direct term
\(\nu\mathcal D\) in the exact paired-cutoff balance supplies the second
factor in the arithmetic--geometric mean step and is what converts signed
innovation into the source debit \eqref{eq:v21-source-crossing-debit}.  This
is the NS instance of the V20 companion warning not to separate a positive
Gram from its signed direct term.
\end{remark}

% =============================================================================
\section{Interpolation collapses to a canonically summable record weight}
\label{sec:v21-interpolation-collapse}
% =============================================================================

\begin{lemma}[Record-local upper bound for the source interpolation scale]
\label{lem:v21-source-interpolation-upper}
Let \(I=[t_p,t_q]\) lie between doubled first-hitting times and put
\[
 R^+:=R_q=\|\Lambda^{1/2}u(t_q)\|_2.
\]
Then \(\|\Lambda^{1/2}u(t)\|_2\le R^+\) for \(t\in I\), and for
\(\sigma,\theta\) as in \eqref{eq:v21-sigma-theta},
\begin{equation}
 \boxed{\quad
 \mathscr S_{\sigma;L,M;I}
 \le
 C_{\rm int}
 (R^+)^{2(1-\theta)}E_*^\theta
 \int_I\|\Lambda u(t)\|_2^2\,\dd t.
 \quad}
 \label{eq:v21-source-interpolation-upper}
\end{equation}
The constant is independent of \(L,M,p,q\), and \(\theta\in[0,1]\).
\end{lemma}

\begin{proof}
The product estimates used in V7 and V19 give, at every smooth time,
\begin{align}
 \|A^{-1/2}\BNS(u)\|_2
 &\le C\|\Lambda^{1/2}u\|_2\|\Lambda u\|_2,
 \label{eq:v21-mid-pointwise}\\
 \|A^{-3/4}\BNS(u)\|_2
 &\le C\|u\|_2\|\Lambda u\|_2
 \le C\sqrt{E_*}\|\Lambda u\|_2.
 \label{eq:v21-weak-pointwise}
\end{align}
Hilbert-scale log convexity, with
\(\sigma=(1-\theta)/2+3\theta/4\), gives
\[
 \|A^{-\sigma}\Pi_{L,M}\BNS(u)\|_2
 \le
 \|A^{-1/2}\Pi_{L,M}\BNS(u)\|_2^{1-\theta}
 \|A^{-3/4}\Pi_{L,M}\BNS(u)\|_2^\theta.
\]
The projector commutes with \(A\) and is contractive in every displayed
Hilbert norm.  Insert
\eqref{eq:v21-mid-pointwise}--\eqref{eq:v21-weak-pointwise} and use the
first-hitting bound
\(\|\Lambda^{1/2}u(t)\|_2\le R^+\).  After squaring, the pointwise right side
is at most
\[
 C_{\rm int}
 (R^+)^{2(1-\theta)}E_*^\theta
 \|\Lambda u(t)\|_2^2.
\]
Integrate on \(I\).
\end{proof}

\begin{theorem}[All interpolated source ledgers have the same canonical
deficit]
\label{thm:v21-interpolation-collapse}
Use the cofinal front from V20.  On the \(j\)-th front gap set
\begin{align}
 I_j&=[t_{n_j},t_{n_{j+1}}],
 &
 L_j&=\rho_{n_j},
 &
 M_j&=\rho_{n_{j+1}},
 \notag\\
 R_j^+&:=R_{n_{j+1}},
 &
 K_j^+&:=\frac{(R_j^+)^2}{E_*}.
 \label{eq:v21-front-gap-data}
\end{align}
Let \(\mathfrak N\) be any subfamily on which the front-jump innovation
\eqref{eq:v20-front-jump-innovated} holds, with terminal stock
\(\widehat J_j\).  Then, for every \(\theta\in[0,1]\),
\begin{equation}
 \boxed{\quad
 \sum_{j\in\mathfrak N}
 \frac{\widehat J_j}
 {(R_j^+)^2(M_j/K_j^+)^\theta}
 \le\frac{C_{\rm int}E_*}{2\nu^2}.
 \quad}
 \label{eq:v21-interpolated-front-ledger}
\end{equation}
Because \(M_j\ge K_j^+\), increasing \(\theta\) only weakens the summand.
At \(\theta=0\) the result is
\begin{equation}
 \sum_{j\in\mathfrak N}
 \frac{\widehat J_j}{(R_j^+)^2}
 \le\frac{C_{\rm int}E_*}{2\nu^2}.
 \label{eq:v21-mid-front-ledger}
\end{equation}
At \(\theta=1\) it is the inverse-frequency weak-source scale
\begin{equation}
 \sum_{j\in\mathfrak N}\frac{\widehat J_j}{E_*M_j}
 \le\frac{C_{\rm int}E_*}{2\nu^2}.
 \label{eq:v21-weak-front-ledger}
\end{equation}
\end{theorem}

\begin{proof}
Apply Theorem~\ref{thm:v21-source-cost} on each innovation gap:
\[
 \frac{\nu\widehat J_j}{M_j^\theta}
 \le\mathscr S_{\sigma;L_j,M_j;I_j}.
\]
Lemma~\ref{lem:v21-source-interpolation-upper}, with terminal record
\(R_j^+\), gives
\[
 \frac{\nu\widehat J_j}
 {M_j^\theta(R_j^+)^{2(1-\theta)}E_*^\theta}
 \le
 C_{\rm int}\int_{I_j}\|\Lambda u\|_2^2\,\dd t.
\]
The front time intervals have disjoint interiors.  Multiply by \(\nu\), sum,
and use the Leray energy budget.  Finally,
\[
 M_j^\theta(R_j^+)^{2(1-\theta)}E_*^\theta
 =(R_j^+)^2\left(\frac{M_j}{K_j^+}\right)^\theta,
\]
which proves \eqref{eq:v21-interpolated-front-ledger}.  The endpoint
statements follow by substitution.
\end{proof}

\begin{firewall}[Why source interpolation cannot close the front]
\label{fw:v21-interpolation-no-closure}
Every terminal band stock satisfies \(\widehat J_j\le E_*\), while the
front records grow at least geometrically.  Therefore
\begin{equation}
 \sum_j\frac{\widehat J_j}{(R_j^+)^2}
 \le E_*\sum_j\frac1{(R_j^+)^2}<\infty
 \label{eq:v21-canonical-stock-automatic}
\end{equation}
without using Navier--Stokes dynamics.  Equation
\eqref{eq:v21-interpolated-front-ledger} is consequently a correct
contemporaneous source-incidence theorem, but every member of the
\(H^{-1}\)--\(H^{-3/2}\) interpolation family lands on an already summable
record weight.  The weak endpoint adds the further overshoot penalty
\(M_j/K_j^+\).  No choice of \(\sigma\) in this interval removes the
amplitude deficit.
\end{firewall}

% =============================================================================
\section{The strongest direct one-crossing estimate returns one record
factor}
\label{sec:v21-direct-amplitude}
% =============================================================================

\begin{theorem}[Critical amplitude ledger for paired innovations]
\label{thm:v21-critical-amplitude-ledger}
For every band \(0<L<M\) and every smooth time,
\begin{equation}
 \boxed{\quad
 |\Psi_{L,M}(t)|
 \le C_{\rm amp}\,
 \|\Lambda^{1/2}u(t)\|_2
 \|\Lambda u(t)\|_2^2.
 \quad}
 \label{eq:v21-pointwise-band-amplitude}
\end{equation}
On the V20 front innovation subfamily \(\mathfrak N\),
\begin{equation}
 \boxed{\quad
 \sum_{j\in\mathfrak N}
 \frac{\widehat J_j}{R_j^+}
 \le\frac{2C_{\rm amp}E_*}{\nu}.
 \quad}
 \label{eq:v21-critical-front-ledger}
\end{equation}
\end{theorem}

\begin{proof}
Use Proposition~\ref{prop:v21-modal-band-injection} and pair at the
intermediate source level:
\begin{align*}
 |\Psi_{L,M}|
 &\le
 \|A^{-1/2}\Pi_{L,M}\BNS(u)\|_2
 \|A^{1/2}\Pi_{L,M}u\|_2\\
 &\le
 \|A^{-1/2}\BNS(u)\|_2\|\Lambda u\|_2.
\end{align*}
The estimate \eqref{eq:v21-mid-pointwise} proves
\eqref{eq:v21-pointwise-band-amplitude}.

On \(I_j\), the first-hitting property gives
\(\|\Lambda^{1/2}u(t)\|_2\le R_j^+\).  Hence
\[
 |\mathfrak I_j|
 \le C_{\rm amp}R_j^+
     \int_{I_j}\|\Lambda u(t)\|_2^2\,\dd t.
\]
The innovation certificate gives
\(\mathfrak I_j\ge\widehat J_j/4\).  Divide by \(R_j^+\), sum over the
disjoint front intervals, and use
\(\int\|\Lambda u\|_2^2\le E_*/(2\nu)\).  Enlarging
\(C_{\rm amp}\) by the harmless factor four gives
\eqref{eq:v21-critical-front-ledger}.
\end{proof}

\begin{assessment}[The packet saturates the returned record factor]
\label{ass:v21-packet-amplitude}
On the V14/V18 concentrated packet, high-frequency kinetic stock has scale
\(J\asymp R^2/N\), while its turnover-interval Leray occupation has scale
\[
 \int_I\|\Lambda u\|_2^2\,\dd t\asymp\frac{R}{N}.
\]
Thus
\[
 \frac{J}{R}\asymp\frac{R}{N}
\]
matches the right side of
\eqref{eq:v21-critical-front-ledger}.  The direct one-crossing estimate has
removed the inverse physical frequency from V20, but the price is exactly one
critical-record factor.  This is not a proof that every constant is sharp;
it is a rigorous scaling obstruction to deleting \(R_j^+\) from
\eqref{eq:v21-pointwise-band-amplitude} using only the norms present there.
\end{assessment}

\begin{firewall}[The critical amplitude series is also automatic]
\label{fw:v21-critical-weight-automatic}
Since \(\widehat J_j\le E_*\) and \(R_j^+\) grows geometrically,
\[
 \sum_j\frac{\widehat J_j}{R_j^+}
 \le E_*\sum_j\frac1{R_j^+}<\infty.
\]
Thus \eqref{eq:v21-critical-front-ledger} is stronger than the weak
inverse-frequency ledger as a pointwise pairing, but it still cannot
contradict an unbounded doubled chronology.  The remaining missing estimate
must exploit a signed dynamical relation between successive innovations, not
only finite triad incidence and a one-card product bound.
\end{firewall}

% =============================================================================
\section{Integrated V21 conclusion and V22 acquisition order}
\label{sec:v21-conclusion}
% =============================================================================

\begin{theorem}[What V21 proves]
\label{thm:v21-conclusion}
Relative to V1--V20, the following are proved.
\begin{enumerate}[label=\textup{(V21.\arabic*)},leftmargin=6.2em]
\item Paired-cutoff injection is exactly the negative modal work inside the
      intervening band, with no common-tail contribution.
\item Across pairwise disjoint front bands, each complete triad can enter at
      most three space--time innovations.  Radial triad multiplicity is
      therefore closed.
\item A positive inherited-or-innovation branch, with its annular viscous
      term retained, forces the full Sobolev source debit
      \(\mathscr S_\sigma\ge\nu J/M^{4\sigma-2}\) for
      \(1/2\le\sigma\le3/4\).
\item Record-local interpolation supplies matching upper estimates for that
      entire source family.
\item Every interpolated global front ledger collapses to
      \(J/(R^+)^2\), with an additional overshoot loss away from the
      intermediate endpoint; this weight is automatically summable.
\item The strongest direct critical pairing gives the improved one-crossing
      ledger \(\sum J/R^+<\infty\), but that weight is also automatically
      summable and is packet-scale sharp.
\end{enumerate}
No full Navier--Stokes stability or global-regularity theorem is obtained.
\end{theorem}

\begin{proof}
Item \textup{(V21.1)} is
Proposition~\ref{prop:v21-modal-band-injection}.  Item
\textup{(V21.2)} is Theorem~\ref{thm:v21-three-leg-one-crossing} and
Corollary~\ref{cor:v21-front-no-radial-multiplicity}.  Item
\textup{(V21.3)} is Theorem~\ref{thm:v21-source-cost}.  Item
\textup{(V21.4)} is
Lemma~\ref{lem:v21-source-interpolation-upper}.  Item
\textup{(V21.5)} is Theorem~\ref{thm:v21-interpolation-collapse} and
Firewall~\ref{fw:v21-interpolation-no-closure}.  Item
\textup{(V21.6)} is Theorem~\ref{thm:v21-critical-amplitude-ledger},
Assessment~\ref{ass:v21-packet-amplitude}, and
Firewall~\ref{fw:v21-critical-weight-automatic}.
\end{proof}

\begin{programme}[V22 acquisition order]
\label{prog:v22}
\begin{enumerate}[label=\textup{(AC\arabic*)},leftmargin=4.8em]
\item Stop pursuing source interpolation between \(A^{-1/2}\BNS(u)\) and
      \(A^{-3/4}\BNS(u)\); Theorem~\ref{thm:v21-interpolation-collapse}
      proves that every such route lands on an automatically summable
      record weight.
\item Use the signed triad formula
      \eqref{eq:v21-triad-band-injection} to classify front crossings by
      radial leg pattern: one leg in the band, two legs in the band, or all
      three.  Preserve the cancellation between patterns.
\item For high--high/low triads, combine
      \eqref{eq:v18-high-parent-gap} with the disjoint front bands.  Determine
      whether the low leg can reset positive transfer on many front time
      windows without paying low-mode energy variation.
\item Separate local front crossings, whose three radii are comparable,
      from nonlocal high--high/low returns.  A valid estimate must keep the
      paired band injection signed in both classes.
\item Seek a common-time or common-terminal representation of modal-transfer
      increments.  Time variation, not radial incidence, is now the only
      way one triad can support repeated positive front innovations.
\item Test any proposed variation debit on the concentrated packet and on
      the explicit V4 triad.  Do not assume monotone forward energy cascade.
\item If the time-variation route returns only
      \(\int\mathfrak L(t)\,\dd t\) or total spectral variation, record the
      circularity and move upstream to a conserved or signed barycentric
      transfer quantity.
\item The next required SM/YM synchronization remains V25.
\end{enumerate}
\end{programme}

\begin{assessment}[Position after V21]
The cofinal front now has an exact one-crossing incidence theorem: no modal
leg is duplicated across its disjoint radial bands.  Retaining the complete
annular dissipation also converts every positive paired innovation into a
quantitative contemporaneous source debit.  These two advances remove
radial multiplicity and source/carrier mistiming.  They do not supply the
missing amplitude.  The entire paid Sobolev interpolation range collapses to
an already summable record weight, and the strongest direct estimate is
saturated at scale \(J/R\) by the packet.  Repeated positive front work can
therefore survive only through signed time variation of complete triad
transfers.  V22 must decide whether that variation has a one-use low-leg or
barycentric payment; otherwise the route returns to the classical
continuation boundary.
\end{assessment}

% =============================================================================
\section*{Version V21 history}
\addcontentsline{toc}{section}{Version V21 history}
% =============================================================================

Expressed the paired-cutoff innovation as complete modal work in one annulus
and proved that each complete triad enters at most three disjoint cofinal
front bands.  Used the annular viscous term to derive a sharp continuum of
source-crossing debits from \(H^{-1}\) to the globally paid \(H^{-3/2}\)
source.  Proved matching record-local interpolation upper bounds and showed
that every resulting global ledger collapses to an automatically summable
canonical record weight.  Derived the stronger direct critical amplitude
ledger \(\sum J/R<\infty\), verified its packet scaling, and isolated signed
time variation---not radial incidence---as the next bottleneck.  The full
stability/global-regularity goal remains open.

% =============================================================================
% VERSION V22 APPEND
% =============================================================================

\section{V22: time-skew normal form and heterochiral receiver structure}
\label{sec:v22-entry}
% =============================================================================

V21 proved that a complete Fourier triad can meet at most three members of a
disjoint radial front.  That finite incidence statement did not identify the
algebraic form of the surviving defect when the three bands are observed on
different time intervals.  V22 performs that calculation before applying an
absolute value.  The result is sharper than the phrase ``time variation''
used in V21: after complete-triad cancellation, the defect is a difference of
\emph{time-window increments} of one or two modal transfers.  Bounding it
directly requires gross triadic transfer action, not a derivative of the
transfer.

The calculation also decides the proposed barycentric route.  Helical
decomposition shows that a completed triad has only the energy and signed
helicity annihilators.  On the explicit V5 geometry, requiring the same
helicity-blind radial weight to cancel both a homochiral and a heterochiral
sector forces that weight to be constant on the three radii.  Thus there is
no new positive radial quadratic invariant behind the moving front.

There is nevertheless a genuine structural gain.  Homochiral triads make no
nonlinear contribution to the critical record.  For a high--high/low
heterochiral triad whose two high legs lie in one band, positive critical work
and positive band injection have the same sign exactly when the minority
helicity is on a high leg.  In the quantitatively nonlocal range, a fixed
fraction of that injection is then received by the minority-helicity high
mode.  This identifies the next possible one-use carrier without asserting
that the other incident triads cannot drain it.

% =============================================================================
\section{A complete radial-pattern table for one paired cutoff}
\label{sec:v22-radial-patterns}
% =============================================================================

Fix one complete unordered triad, use the notation
\eqref{eq:v18-ordered-triad-radii}, and abbreviate
\begin{equation}
 \mathcal S_{L,M}^{\triangle}
 :=\{i\in\{1,2,3\}:L<\lambda_i\le M\}.
 \label{eq:v22-band-leg-set}
\end{equation}

\begin{proposition}[Exact radial-pattern classification]
\label{prop:v22-radial-pattern-table}
At every smooth time, the contribution of the triad to the paired-cutoff
injection is given by the following complete table:
\begin{equation}
\begin{array}{c|c}
 \mathcal S_{L,M}^{\triangle}
 &\Psi_{L,M}^{\triangle}\rule{0pt}{2.5ex}\\
\hline
 \varnothing\ \hbox{or}\ \{1,2,3\}&0\\
 \{1\}&-T_1^{\triangle}\\
 \{2\}&-T_2^{\triangle}\\
 \{3\}&-T_3^{\triangle}\\
 \{1,2\}&T_3^{\triangle}\\
 \{2,3\}&T_1^{\triangle}.
\end{array}
\label{eq:v22-radial-pattern-table}
\end{equation}
The subset \(\{1,3\}\) cannot occur for a radial interval unless coincident
radii also force the middle leg into the interval.  Hence
\eqref{eq:v22-radial-pattern-table} contains every admissible pattern.
\end{proposition}

\begin{proof}
Equation~\eqref{eq:v21-triad-band-injection} gives
\[
 \Psi_{L,M}^{\triangle}
 =-\sum_{i\in\mathcal S_{L,M}^{\triangle}}T_i^{\triangle}.
\]
The empty case is immediate, and the full case vanishes by
\eqref{eq:v18-triad-energy-cancellation}.  A one-leg set gives the three
middle rows.  If two legs occur, eliminate their sum using
\(T_1^{\triangle}+T_2^{\triangle}+T_3^{\triangle}=0\); this gives the last
two rows.  Finally, radial ordering implies that membership in
\((L,M]\) is contiguous in the ordered list of radii, which excludes
\(\{1,3\}\).
\end{proof}

\begin{corollary}[The high--high/low band contribution is a low-leg current]
\label{cor:v22-high-pair-low-current}
Suppose
\begin{equation}
 \lambda_1\le L<\lambda_2\le\lambda_3.
 \label{eq:v22-low-below-band}
\end{equation}
Then
\begin{equation}
 \Psi_{L,M}^{\triangle}=
 \begin{cases}
  0,&M<\lambda_2,\\
  -T_2^{\triangle}=T_1^{\triangle}+T_3^{\triangle},
      &\lambda_2\le M<\lambda_3,\\
  T_1^{\triangle},&M\ge\lambda_3.
 \end{cases}
 \label{eq:v22-high-pair-band-current}
\end{equation}
If a cutoff separates the two high radii, it lies in an interval of length at
most the low radius:
\begin{equation}
 0\le\lambda_3-\lambda_2\le\lambda_1.
 \label{eq:v22-high-pair-split-width}
\end{equation}
This geometric narrowness does not constrain the time elapsed between the
front intervals assigned to the two high legs.
\end{corollary}

\begin{proof}
The three cases in \eqref{eq:v22-high-pair-band-current} are respectively the
sets \(\varnothing\), \(\{2\}\), and \(\{2,3\}\) in
Proposition~\ref{prop:v22-radial-pattern-table}.  The gap bound is
\eqref{eq:v18-high-parent-gap}.  The front construction orders observation
intervals by their cutoff indices but supplies no upper bound on their time
separation, proving the last assertion.
\end{proof}

% =============================================================================
\section{The ordered front has an exact two-increment time-skew defect}
\label{sec:v22-time-skew}
% =============================================================================

Let
\begin{equation}
 0<r_0<r_1<\cdots,
 \qquad
 \mathcal B_j=(r_j,r_{j+1}],
 \qquad
 I_j=[\tau_j,\tau_{j+1}],
 \label{eq:v22-ordered-front}
\end{equation}
where the times are increasing.  This includes the V20 cofinal front after a
relabelling.  For one fixed complete triad define
\begin{equation}
 Q_i(J):=\int_JT_i^{\triangle}(t)\,\dd t.
 \label{eq:v22-leg-window-increment}
\end{equation}
If \(\lambda_i>r_0\), let \(b_i\) be the unique band index for which
\(\lambda_i\in\mathcal B_{b_i}\).  Radial ordering gives
\(b_1\le b_2\le b_3\) whenever all three indices are defined.

\begin{theorem}[Complete-triad time-skew normal form]
\label{thm:v22-time-skew-normal-form}
Assume first that all three radii lie above \(r_0\), and let
\begin{equation}
 \mathfrak C^{\triangle}
 :=\sum_j\int_{I_j}\Psi_{r_j,r_{j+1}}^{\triangle}(t)\,\dd t
 \label{eq:v22-total-front-triad}
\end{equation}
be the total contribution of the triad to the ordered front.  Then precisely
one of the following formulas applies:
\begin{align}
 b_1=b_2=b_3
 &\quad\Longrightarrow\quad
 \mathfrak C^{\triangle}=0,
 \label{eq:v22-skew-one-band}\\
 b_1=b_2=j<b_3=m
 &\quad\Longrightarrow\quad
 \mathfrak C^{\triangle}=Q_3(I_j)-Q_3(I_m),
 \label{eq:v22-skew-lower-pair}\\
 b_1=j<b_2=b_3=m
 &\quad\Longrightarrow\quad
 \mathfrak C^{\triangle}=Q_1(I_m)-Q_1(I_j),
 \label{eq:v22-skew-upper-pair}\\
 b_1=j<b_2=m<b_3=n
 &\quad\Longrightarrow\quad
 \mathfrak C^{\triangle}
 =\bigl(Q_1(I_m)-Q_1(I_j)\bigr)
  +\bigl(Q_3(I_m)-Q_3(I_n)\bigr).
 \label{eq:v22-skew-three-bands}
\end{align}
Thus complete energy cancellation leaves at most two differences of
time-window increments; it does not leave a radial multiplicity factor.

If instead \(\lambda_1\le r_0<\lambda_2\), the two high legs obey
\begin{align}
 b_2=b_3=j
 &\quad\Longrightarrow\quad
 \mathfrak C^{\triangle}=Q_1(I_j),
 \label{eq:v22-unseen-low-same-band}\\
 b_2=j<b_3=m
 &\quad\Longrightarrow\quad
 \mathfrak C^{\triangle}
 =Q_1(I_j)+Q_3(I_j)-Q_3(I_m).
 \label{eq:v22-unseen-low-split-band}
\end{align}
The first term on the right of
\eqref{eq:v22-unseen-low-split-band} is the low-leg reset current; the last
two terms are the sole time-skew defect.
\end{theorem}

\begin{proof}
Only the bands containing triad radii contribute.  If all three radii lie in
one band, its contribution vanishes by complete energy cancellation.  If the
two lower legs lie in \(\mathcal B_j\) and the upper leg lies in
\(\mathcal B_m\), Proposition~\ref{prop:v22-radial-pattern-table} gives
\[
 \mathfrak C^{\triangle}=Q_3(I_j)-Q_3(I_m),
\]
which proves \eqref{eq:v22-skew-lower-pair}.  The same argument with the upper
pair gives \eqref{eq:v22-skew-upper-pair}.

If all band indices are distinct, the raw expression is
\[
 -Q_1(I_j)-Q_2(I_m)-Q_3(I_n).
\]
Integrating
\(T_1^{\triangle}+T_2^{\triangle}+T_3^{\triangle}=0\) over the middle
interval \(I_m\) and adding the resulting zero gives
\eqref{eq:v22-skew-three-bands}.

When the low leg is below the front, a common high band contains the set
\(\{2,3\}\), so its contribution is \(Q_1(I_j)\).  If the high legs are
split, the raw contribution is
\(-Q_2(I_j)-Q_3(I_m)\).  Replace
\(-Q_2(I_j)\) by \(Q_1(I_j)+Q_3(I_j)\) to obtain
\eqref{eq:v22-unseen-low-split-band}.
\end{proof}

\begin{remark}[What ``time variation'' means after completion]
\label{rem:v22-window-not-derivative}
Let
\(\Theta_i(t)=\int_{\tau_0}^tT_i^{\triangle}(s)\,\dd s\).  Then
\(Q_i(I_j)=\Theta_i(\tau_{j+1})-\Theta_i(\tau_j)\).  Hence the differences
in Theorem~\ref{thm:v22-time-skew-normal-form} are four-point increments of
the transfer primitive.  Controlling them by total variation gives
\(\int|T_i^{\triangle}|\), not
\(\int|\partial_tT_i^{\triangle}|\).  This distinction prevents an
unjustified insertion of the V33--V34 source-residual variation already
excluded by Proposition~\ref{prop:v5-two-variations}.
\end{remark}

\begin{proposition}[The direct front bound is gross complete-triad action]
\label{prop:v22-gross-action-bound}
For a finite front prefix, write
\begin{equation}
 \mathfrak I_j^{\triangle}
 :=\int_{I_j}\Psi_{r_j,r_{j+1}}^{\triangle}(t)\,\dd t.
 \label{eq:v22-one-triad-one-tile}
\end{equation}
Then
\begin{equation}
 \sum_j\bigl(\mathfrak I_j^{\triangle}\bigr)_+
 \le
 \sum_{i=1}^3\int_{\bigcup_j I_j}
       |T_i^{\triangle}(t)|\,\dd t.
 \label{eq:v22-one-triad-gross-action}
\end{equation}
Consequently, absolute convergence of the smooth Fourier expansion gives
\begin{equation}
 \sum_j\left(
   \int_{I_j}\Psi_{r_j,r_{j+1}}(t)\,\dd t
  \right)_+
 \le
 \int_{\bigcup_jI_j}\mathfrak A_{\rm tri}(t)\,\dd t,
 \label{eq:v22-front-gross-action}
\end{equation}
where
\begin{equation}
 \mathfrak A_{\rm tri}(t)
 :=\sum_{\triangle}\sum_{i=1}^3|T_i^{\triangle}(t)|.
 \label{eq:v22-gross-triad-action}
\end{equation}
This is an exact finite-incidence estimate, but its right side has discarded
all cancellation between complete triads.
\end{proposition}

\begin{proof}
For a fixed triad and tile,
\eqref{eq:v21-triad-band-injection} and the triangle inequality give
\[
 (\mathfrak I_j^{\triangle})_+
 \le\sum_{\{i:\lambda_i\in\mathcal B_j\}}
       \int_{I_j}|T_i^{\triangle}(t)|\,\dd t.
\]
Each leg belongs to at most one band, so summing in \(j\) proves
\eqref{eq:v22-one-triad-gross-action}.  Next use
\((\sum_\triangle x_\triangle)_+
 \le\sum_\triangle(x_\triangle)_+\), sum over the triads, and interchange
the absolutely convergent sums and time integral.  This proves
\eqref{eq:v22-front-gross-action}.
\end{proof}

\begin{lemma}[A strong-norm upper bound for gross triadic action]
\label{lem:v22-gross-action-strong-bound}
For every smooth mean-zero field on \(\Torus^3\),
\begin{align}
 \mathfrak A_{\rm tri}(t)
 &\le C\|\widehat u(t)\|_{\ell^1}
          \|u(t)\|_2\|\Lambda u(t)\|_2,
 \label{eq:v22-gross-wiener-bound}\\
 &\le C\|u(t)\|_2^{5/4}
          \|\Lambda u(t)\|_2\|Au(t)\|_2^{3/4}.
 \label{eq:v22-gross-H2-bound}
\end{align}
Thus the direct absolute estimate reaches an \(H^2\) quantity not supplied
by the Leray energy budget.
\end{lemma}

\begin{proof}
The Fourier formula \eqref{eq:v4-Fourier-B} and boundedness of the Leray
matrix imply, up to the fixed permutation multiplicity used to assemble
unordered triads,
\[
 \mathfrak A_{\rm tri}
 \le C\sum_{p,q\in\mathbb Z^3}
       |q|\,|\widehat u(p)|\,|\widehat u(q)|
       |\widehat u(p+q)|.
\]
For fixed \(p\), Cauchy--Schwarz in \(q\) bounds the inner sum by
\(\|\Lambda u\|_2\|u\|_2\).  Summing
\(|\widehat u(p)|\) proves
\eqref{eq:v22-gross-wiener-bound}.

For completeness, split the Fourier \(\ell^1\) norm at a radius \(N\ge1\).
Lattice counting and Cauchy--Schwarz give
\[
 \sum_{|k|\le N}|\widehat u(k)|\le CN^{3/2}\|u\|_2,
 \qquad
 \sum_{|k|>N}|\widehat u(k)|
 \le CN^{-1/2}\|Au\|_2.
\]
Optimizing at
\(N=(\|Au\|_2/\|u\|_2)^{1/2}\), with the mean-zero Poincar\'e inequality
covering the case where this radius is below one, yields
\[
 \|\widehat u\|_{\ell^1}
 \le C\|u\|_2^{1/4}\|Au\|_2^{3/4}.
\]
Insert this in \eqref{eq:v22-gross-wiener-bound} to prove
\eqref{eq:v22-gross-H2-bound}.
\end{proof}

\begin{firewall}[Gross transfer is not a Leray payment]
\label{fw:v22-gross-action-circular}
The energy inequality controls
\(\int\|\Lambda u\|_2^2\), not
\(\int\|Au\|_2^2\) or the Wiener norm in
\eqref{eq:v22-gross-wiener-bound}.  Applying Young's inequality to
\eqref{eq:v22-gross-H2-bound} merely introduces an \(H^2\) dissipation term;
closing that term is the strong-solution continuation problem.  Therefore
\eqref{eq:v22-front-gross-action} is a correct reduction of the time-skew
defect, but not a global front budget.
\end{firewall}

% =============================================================================
\section{A selected low leg has an exact complementary-circulation residual}
\label{sec:v22-low-leg-residual}
% =============================================================================

The common-high-band case
\eqref{eq:v22-unseen-low-same-band} suggests paying \(Q_1(I_j)\) from the
kinetic energy of the low mode.  That is valid only after every triad incident
to that low mode has been assembled.

For a nonzero Fourier mode \(k\), let
\begin{equation}
 e_k(t):=|\widehat u(k,t)|^2,
 \qquad
 T_k(t)=\sum_{\triangle\ni k}T_{k,\triangle}(t),
 \label{eq:v22-modal-triad-assembly}
\end{equation}
where the normalization of \(T_{k,\triangle}\) is chosen consistently with
the complete unordered-triad decomposition used in V18.  Smoothness makes
the sum absolutely convergent.

\begin{theorem}[Exact partial-triad modal balance]
\label{thm:v22-partial-triad-balance}
Let \(\mathfrak S_k\) be any collection of complete triads incident to
\(k\), and put
\begin{equation}
 S_k:=\sum_{\triangle\in\mathfrak S_k}T_{k,\triangle},
 \qquad
 C_k:=\sum_{\substack{\triangle\ni k\\
                 \triangle\notin\mathfrak S_k}}T_{k,\triangle}.
 \label{eq:v22-selected-complement-transfer}
\end{equation}
For every smooth interval \(I=[a,b]\),
\begin{equation}
 \boxed{\quad
 \int_I S_k(t)\,\dd t
 =\frac12\bigl(e_k(a)-e_k(b)\bigr)
  -\nu|k|^2\int_Ie_k(t)\,\dd t
  -\int_I C_k(t)\,\dd t.
 \quad}
 \label{eq:v22-partial-triad-modal-balance}
\end{equation}
Only the fully assembled choice \(\mathfrak S_k=\{\triangle:\triangle\ni k\}\)
removes the last term.  For the front selection, that last term is precisely
the transfer through incident triads whose high legs were assigned to other
bands or other time windows.
\end{theorem}

\begin{proof}
Taking the real scalar product of the \(k\)-th Fourier equation with
\(\overline{\widehat u(k)}\) gives
\[
 \frac12e_k'(t)+\nu|k|^2e_k(t)=-T_k(t).
\]
By \eqref{eq:v22-modal-triad-assembly}, \(T_k=S_k+C_k\).  Integrate in time
and solve for \(\int_I S_k\) to obtain
\eqref{eq:v22-partial-triad-modal-balance}.
\end{proof}

\begin{firewall}[The low-leg reset is not an endpoint stock by itself]
\label{fw:v22-low-leg-not-paid}
At the level of the exact ledger, the obstruction in
\eqref{eq:v22-partial-triad-modal-balance} is unavoidable.  Given any
integrable scalar function \(f\), the assignments
\(S_k=f\), \(C_k=-f\), and constant \(e_k\) satisfy the inviscid modal
balance while \(\int S_k\) is arbitrary.  This is an algebraic ledger test,
not a claim that arbitrary \(f\) is realized by a Navier--Stokes solution.
It proves the narrower statement needed here: modal energy balance alone
cannot control a selected family of low-leg transfers.  Such control requires
a new sign or incidence theorem for the complementary triads; omitting
\(C_k\) would assume the desired no-recycling conclusion.
\end{firewall}

% =============================================================================
\section{Helical completion rules out a new positive radial invariant}
\label{sec:v22-helical-rigidity}
% =============================================================================

For each nonzero wavevector choose helical unit vectors \(h_k^s\),
\(s\in\{+1,-1\}\), satisfying
\begin{equation}
 \mathrm i k\times h_k^s=s|k|h_k^s.
 \label{eq:v22-helical-eigenvectors}
\end{equation}
Consider one reality-paired helical triad with ordered radii \(\lambda_i\)
and signs \(s_i\).  Set
\begin{equation}
 \mu_i:=s_i\lambda_i.
 \label{eq:v22-signed-radii}
\end{equation}

\begin{lemma}[Energy--helicity normal form for one helical triad]
\label{lem:v22-helical-transfer-normal-form}
For every helical triad there is a real scalar \(\Xi\), depending on its
amplitudes and common phase, such that its modal work vector has the form
\begin{equation}
 \boxed{\quad
 \begin{pmatrix}
  T_1^{\triangle,s}\\T_2^{\triangle,s}\\T_3^{\triangle,s}
 \end{pmatrix}
 =\Xi
 \begin{pmatrix}
  \mu_2-\mu_3\\
  \mu_3-\mu_1\\
  \mu_1-\mu_2
 \end{pmatrix}.
 \quad}
 \label{eq:v22-helical-transfer-vector}
\end{equation}
The scalar vanishes for an inactive geometry or phase.  For an active
geometry its sign can be reversed by reversing the common interaction phase.
\end{lemma}

\begin{proof}
Energy and helicity cancellation hold separately on a complete helical triad.
Therefore its work vector \(T\in\mathbb R^3\) satisfies
\[
 (1,1,1)\cdot T=0,
 \qquad
 (\mu_1,\mu_2,\mu_3)\cdot T=0.
\]
If \((\mu_1,\mu_2,\mu_3)\) is not constant, the common orthogonal
complement of these two vectors is spanned by their cross product
\[
 (\mu_2-\mu_3,\ \mu_3-\mu_1,\ \mu_1-\mu_2).
\]
If the signed radii are all equal, the three signs and radii are common.  The
corresponding reality-paired field is Beltrami:
\(\operatorname{curl}u=\mu_1u\).  The vector identity used in the proof of
\eqref{eq:v4-vortex-stretching} then makes its projected quadratic term zero,
so \(T=0\) and \eqref{eq:v22-helical-transfer-vector} holds with
\(\Xi=0\).

In the nondegenerate Fourier formula, the remaining real scalar is a fixed
geometric coefficient times the real or imaginary part of the product of the
three complex amplitudes.  A phase reversal changes its sign whenever that
coefficient is nonzero.
\end{proof}

\begin{lemma}[The V5 geometry is active in two incompatible helicity sectors]
\label{lem:v22-two-active-helicity-sectors}
Use the wavevectors and positive-helicity vectors in
\eqref{eq:v5-helical-vectors}.  The homochiral \((+,+,+)\) sector is active by
Lemma~\ref{lem:v5-homochiral-active}.  The \((+,+,-)\) sector obtained by
replacing \(h_r\) with \(h_r^-:=\overline h_r\) is also active.  Indeed its
geometric scalar contains
\begin{align}
 (\overline h_q\times h_r)\cdot\overline h_p
 &=\frac14\left(
    \frac1{\sqrt6}+\frac1{\sqrt2}-\frac1{\sqrt3}
   \right)
 \notag\\
 &\quad+\frac{\mathrm i}{4}\left(
    1+\frac1{\sqrt2}-\frac2{\sqrt3}-\frac1{\sqrt6}
   \right),
 \label{eq:v22-heterochiral-active-coefficient}
\end{align}
whose real part is strictly positive.
\end{lemma}

\begin{proof}
The first assertion is exactly
\eqref{eq:v5-active-coefficient}.  With the displayed V5 vectors, negative
helicity at \(r\) is obtained by complex conjugation.  Direct substitution
and evaluation of the scalar triple product gives
\eqref{eq:v22-heterochiral-active-coefficient}.  Its real part is positive
because
\(1/\sqrt6+1/\sqrt2>1/\sqrt3\).  Hence the coupling coefficient is nonzero.
\end{proof}

\begin{theorem}[Helicity-blind radial quadratic-invariant rigidity]
\label{thm:v22-radial-invariant-rigidity}
Let
\begin{equation}
 (\lambda_1,\lambda_2,\lambda_3)
 =N(1,\sqrt2,\sqrt3)
 \label{eq:v22-v5-radii}
\end{equation}
be the radii of the V5 geometry, and let \(w_i\in\mathbb R\) be weights which
depend on radius but not on helicity.  Suppose
\begin{equation}
 \sum_{i=1}^3w_iT_i^{\triangle,s}=0
 \label{eq:v22-weighted-sector-cancellation}
\end{equation}
for every amplitude phase in both active sectors \((+,+,+)\) and
\((+,+,-)\).  Then
\begin{equation}
 w_1=w_2=w_3.
 \label{eq:v22-radial-weights-constant}
\end{equation}
In particular, no strictly increasing positive radial weight, including
\(\lambda^\alpha\) for \(\alpha>0\), is a nonlinear invariant across both
sectors of this one triad geometry.
\end{theorem}

\begin{proof}
For an active sector, Lemma~\ref{lem:v22-helical-transfer-normal-form} and
phase reversal show that \eqref{eq:v22-weighted-sector-cancellation} holds
exactly when \(w=(w_1,w_2,w_3)\) belongs to the annihilator of the nonzero
transfer direction.  That annihilator is
\[
 \operatorname{span}\{(1,1,1),(\mu_1,\mu_2,\mu_3)\}.
\]
For the homochiral sector there are \(a,b\in\mathbb R\) such that
\[
 w_i=a+b\lambda_i.
\]
For the \((+,+,-)\) sector there are \(c,d\in\mathbb R\) such that
\[
 w_1=c+d\lambda_1,
 \qquad w_2=c+d\lambda_2,
 \qquad w_3=c-d\lambda_3.
\]
Since \(\lambda_1\ne\lambda_2\), comparison of the first two components
gives \(a=c\) and \(b=d\).  The third component then gives
\(a+b\lambda_3=a-b\lambda_3\), hence \(b=0\).  This proves
\eqref{eq:v22-radial-weights-constant}.
\end{proof}

\begin{corollary}[No coercive barycentric payment follows from the two
quadratic cancellations]
\label{cor:v22-no-positive-barycentre}
Energy is the only helicity-blind positive radial quadratic cancellation on
the two active V5 sectors.  Helicity supplies the signed weight
\(s|k|\), not a positive radial weight.  Consequently energy--helicity
invariant algebra cannot by itself bound the time-skew defects
\eqref{eq:v22-skew-lower-pair}--\eqref{eq:v22-skew-three-bands} by a new
coercive spectral barycentre.
\end{corollary}

\begin{proof}
Apply Theorem~\ref{thm:v22-radial-invariant-rigidity}.  A constant radial
weight is kinetic energy.  The second annihilator in a fixed helical sector
is signed helicity and changes when a helicity sign changes.  It therefore
cannot define the requested helicity-blind positive stock.
\end{proof}

% =============================================================================
\section{Critical growth is heterochiral and its nonlocal sign table is exact}
\label{sec:v22-heterochiral-signs}
% =============================================================================

\begin{proposition}[Homochiral triads are absent from critical production]
\label{prop:v22-critical-production-heterochiral}
For a homochiral triad, whether \((+,+,+)\) or \((-,-,-)\),
\begin{equation}
 \mathcal G_{1/2}^{\triangle,s}
 :=-\sum_{i=1}^3\lambda_iT_i^{\triangle,s}=0.
 \label{eq:v22-homochiral-critical-zero}
\end{equation}
Hence the full nonlinear critical production is the sum of the mixed-helicity
triad sectors only.
\end{proposition}

\begin{proof}
In a homochiral sector, signed helicity cancellation is
\(\sum_i s\lambda_iT_i^{\triangle,s}=0\) with one common sign \(s\).
Multiplication by \(s\) proves
\(\sum_i\lambda_iT_i^{\triangle,s}=0\), which is
\eqref{eq:v22-homochiral-critical-zero}.  Sum the absolutely convergent
helical triad expansion of a smooth field.
\end{proof}

\begin{theorem}[Minority-helicity sign and receiver table]
\label{thm:v22-minority-helicity-table}
Consider a mixed-helicity triad and, after reversing all signs if necessary,
let \(m\in\{1,2,3\}\) be its unique negative-helicity leg.  With the scalar
\(\Xi\) in \eqref{eq:v22-helical-transfer-vector}, its critical production
is
\begin{equation}
 \mathcal G_{1/2}^{\triangle,s}=
 \begin{cases}
  2\lambda_1(\lambda_3-\lambda_2)\Xi,&m=1,\\
  -2\lambda_2(\lambda_3-\lambda_1)\Xi,&m=2,\\
  2\lambda_3(\lambda_2-\lambda_1)\Xi,&m=3.
 \end{cases}
 \label{eq:v22-minority-critical-table}
\end{equation}
Suppose in addition that \(\lambda_1\le L<\lambda_2\le\lambda_3\le M\), so
both high legs lie in one paired-cutoff band.  Then
\(\Psi_{L,M}^{\triangle,s}=T_1^{\triangle,s}\), and
\begin{equation}
 \mathcal G_{1/2}^{\triangle,s}=
 \begin{cases}
  -2\lambda_1\Psi_{L,M}^{\triangle,s},&m=1,\\[0.3em]
  \displaystyle
  \frac{2\lambda_2(\lambda_3-\lambda_1)}{\lambda_2+\lambda_3}
  \Psi_{L,M}^{\triangle,s},&m=2,\\[0.9em]
  \displaystyle
  \frac{2\lambda_3(\lambda_2-\lambda_1)}{\lambda_2+\lambda_3}
  \Psi_{L,M}^{\triangle,s},&m=3.
 \end{cases}
 \label{eq:v22-critical-vs-band-sign-table}
\end{equation}
Thus positive critical work has the opposite sign from high-band injection
when the minority helicity is low, and the same sign when the minority
helicity is high.
\end{theorem}

\begin{proof}
If \(m=1\), then \(\mu=(-\lambda_1,\lambda_2,\lambda_3)\), and
Lemma~\ref{lem:v22-helical-transfer-normal-form} gives
\[
 T=\Xi(\lambda_2-\lambda_3,
             \lambda_3+\lambda_1,
            -\lambda_1-\lambda_2).
\]
Taking \(-\lambda\cdot T\) gives the first row of
\eqref{eq:v22-minority-critical-table}.  The choices
\(\mu=(\lambda_1,-\lambda_2,\lambda_3)\) and
\(\mu=(\lambda_1,\lambda_2,-\lambda_3)\) similarly give
\begin{align*}
 T&=\Xi(-\lambda_2-\lambda_3,
              \lambda_3-\lambda_1,
              \lambda_1+\lambda_2),&&m=2,\\
 T&=\Xi(\lambda_2+\lambda_3,
             -\lambda_3-\lambda_1,
              \lambda_1-\lambda_2),&&m=3.
\end{align*}
Direct scalar products give the remaining two rows of
\eqref{eq:v22-minority-critical-table}.

By Corollary~\ref{cor:v22-high-pair-low-current}, a band containing both high
legs has \(\Psi=T_1\).  Substitution of the displayed first components into
\eqref{eq:v22-minority-critical-table} proves
\eqref{eq:v22-critical-vs-band-sign-table}.  Every multiplier in the last two
rows is strictly positive, while the first multiplier is strictly negative.
If the two high radii coincide, the corresponding band injection and critical
production both vanish.  This proves the sign assertion.
\end{proof}

\begin{corollary}[A nonlocal productive band has a definite minority-high
receiver]
\label{cor:v22-minority-high-receiver}
Under the hypotheses of
Theorem~\ref{thm:v22-minority-helicity-table}, assume
\begin{equation}
 \lambda_1\le\delta\lambda_2,
 \qquad 0<\delta<1,
 \label{eq:v22-quantitative-nonlocality}
\end{equation}
and suppose \(m=2\) or \(m=3\) and
\(\mathcal G_{1/2}^{\triangle,s}>0\).  Then
\(\Psi_{L,M}^{\triangle,s}>0\).  Moreover the nonlinear modal injection into the
minority-helicity high leg, denoted \(\mathfrak R_m:=-T_m^{\triangle,s}\),
is positive and satisfies
\begin{equation}
 \boxed{\quad
 \mathfrak R_m
 \ge\frac{1-\delta}{2+\delta}
       \Psi_{L,M}^{\triangle,s}.
 \quad}
 \label{eq:v22-minority-receiver-fraction}
\end{equation}
In particular, if \(\lambda_1\le\lambda_2/2\), then
\begin{equation}
 \mathfrak R_m\ge\frac15\Psi_{L,M}^{\triangle,s}.
 \label{eq:v22-one-fifth-minority-receiver}
\end{equation}
\end{corollary}

\begin{proof}
The last two rows of
\eqref{eq:v22-critical-vs-band-sign-table} have positive multipliers, so the
two positive signs agree.  If \(m=2\), write \(\Xi=-a\) with \(a>0\).  The
explicit vector in the proof of
Theorem~\ref{thm:v22-minority-helicity-table} gives
\[
 \Psi=(\lambda_2+\lambda_3)a,
 \qquad
 \mathfrak R_2=(\lambda_3-\lambda_1)a.
\]
If \(m=3\), write \(\Xi=a>0\); then
\[
 \Psi=(\lambda_2+\lambda_3)a,
 \qquad
 \mathfrak R_3=(\lambda_2-\lambda_1)a.
\]
The triangle inequality gives
\(\lambda_3\le\lambda_2+\lambda_1
 \le(1+\delta)\lambda_2\).  In either case the numerator in
\(\mathfrak R_m/\Psi\) is at least
\((1-\delta)\lambda_2\), while the denominator is at most
\((2+\delta)\lambda_2\).  This proves
\eqref{eq:v22-minority-receiver-fraction}; set \(\delta=1/2\) for
\eqref{eq:v22-one-fifth-minority-receiver}.
\end{proof}

\begin{firewall}[A triad receiver is not yet a mode receiver]
\label{fw:v22-receiver-circulation}
Corollary~\ref{cor:v22-minority-high-receiver} is pointwise and retains the
sign of one complete helical triad.  It does not say that the total nonlinear
energy derivative of the minority high mode is positive: other triads
incident to that mode contribute the complementary term \(C_k\) in
\eqref{eq:v22-partial-triad-modal-balance}.  Summing only the positive
receiver triads would return the gross action
\eqref{eq:v22-gross-triad-action}.  The next required theorem must therefore
assemble the complementary high-mode circulation while preserving helicity
labels; the one-fifth coefficient cannot be promoted directly to terminal
stock.
\end{firewall}

\begin{assessment}[Packet and explicit-triad tests]
\label{ass:v22-model-tests}
The V4 triad has two equal lower radii and a higher third radius.  Depending on
the band, Proposition~\ref{prop:v22-radial-pattern-table} gives exactly the
already computed donor/receiver current; no forward-cascade monotonicity was
inserted.  The V5 homochiral triad is active but contributes zero critical
production by Proposition~\ref{prop:v22-critical-production-heterochiral}, in
agreement with \eqref{eq:v5-triad-record-fixed}.  Finally, the V14/V18 packet
has integrated flux height and kinetic stock of order \(R^2/N\); the gross
front bound and a receiver fraction of fixed size permit precisely that
scale.  These tests therefore do not contradict the V22 identities and do
not turn the family of distinct packet solutions into one trajectory.
\end{assessment}

% =============================================================================
\section{Integrated V22 conclusion and V23 acquisition order}
\label{sec:v22-conclusion}
% =============================================================================

\begin{theorem}[What V22 proves]
\label{thm:v22-conclusion}
Relative to V1--V21, the following are proved.
\begin{enumerate}[label=\textup{(V22.\arabic*)},leftmargin=6.2em]
\item Every admissible radial leg pattern in a paired-cutoff band has the
      exact signed value listed in
      \eqref{eq:v22-radial-pattern-table}.
\item On an ordered disjoint front, a complete triad whose three legs are
      observed leaves at most two differences of time-window transfer
      increments.  A high--high/low triad with an unseen low leg leaves one
      low-leg current and at most one such difference.
\item Direct summation of positive front innovations is bounded by gross
      complete-triad action.  Its elementary Fourier bound requires a Wiener
      or \(H^2\) quantity not paid by the Leray energy inequality.
\item A selected family of low-leg transfers has the exact endpoint formula
      \eqref{eq:v22-partial-triad-modal-balance}; its complementary incident
      triads form an unavoidable circulation residual.
\item On the active V5 geometry, a helicity-blind radial quadratic weight
      which cancels both homochiral and heterochiral sectors must be constant.
      Hence energy and helicity do not conceal a new positive barycentric
      invariant.
\item Homochiral triads make zero critical-record production.  For a
      high--high/low heterochiral triad, positive critical work aligns with
      positive high-band injection exactly when the minority helicity is on a
      high leg.
\item In the quantitatively nonlocal range, at least
      \((1-\delta)/(2+\delta)\) of such a productive high-band injection is
      received by the minority-helicity high leg at the level of the complete
      triad.
\end{enumerate}
No full Navier--Stokes stability or global-regularity theorem is obtained.
\end{theorem}

\begin{proof}
Item \textup{(V22.1)} is
Proposition~\ref{prop:v22-radial-pattern-table}.  Item \textup{(V22.2)} is
Theorem~\ref{thm:v22-time-skew-normal-form}.  Item \textup{(V22.3)} is
Proposition~\ref{prop:v22-gross-action-bound},
Lemma~\ref{lem:v22-gross-action-strong-bound}, and
Firewall~\ref{fw:v22-gross-action-circular}.  Item \textup{(V22.4)} is
Theorem~\ref{thm:v22-partial-triad-balance}.  Item \textup{(V22.5)} is
Theorem~\ref{thm:v22-radial-invariant-rigidity} and
Corollary~\ref{cor:v22-no-positive-barycentre}.  Items
\textup{(V22.6)--(V22.7)} are
Proposition~\ref{prop:v22-critical-production-heterochiral},
Theorem~\ref{thm:v22-minority-helicity-table}, and
Corollary~\ref{cor:v22-minority-high-receiver}.
\end{proof}

\begin{programme}[V23 acquisition order]
\label{prog:v23}
\begin{enumerate}[label=\textup{(AC\arabic*)},leftmargin=4.8em]
\item Replace the helicity-blind front by two helical annular stocks
      \(\|\Pi_{L,M}u^+\|_2^2\) and
      \(\|\Pi_{L,M}u^-\|_2^2\).  Derive their exact paired balances before
      taking positive parts.
\item Split heterochiral high--high/low interactions by the location of the
      minority sign.  Retain the receiver coefficient
      \eqref{eq:v22-minority-receiver-fraction}; discard no complementary
      incident triad.
\item Determine whether positive high-minority receiver work in one chirality
      can be paired with a simultaneous receiver in the other chirality, so
      that cancellation at a shared mode becomes a signed cross-helicity
      transfer rather than gross triadic action.
\item Use the exact partial-triad balance
      \eqref{eq:v22-partial-triad-modal-balance} to name every complement.
      A proposed terminal-stock debit is valid only if its \(C_k\) terms
      cancel after complete helical assembly.
\item Separate nonlocal interactions
      \(\lambda_1\le\lambda_2/2\), where the one-fifth receiver bound holds,
      from local interactions.  For the local class use comparable-frequency
      orthogonality, not the low-leg argument.
\item Test the helical receiver ledger on the V5 homochiral tangent and on a
      mixed-helicity activation of the same geometry.  Phase freedom forbids
      an assumed sign of \(\Xi\).
\item If helical assembly returns absolute helicity or
      \(\int\mathfrak A_{\rm tri}\), record the circularity and move upstream
      to a signed sector-exchange identity rather than another Sobolev norm.
\item The next required SM/YM synchronization remains V25.
\end{enumerate}
\end{programme}

\begin{assessment}[Position after V22]
The radial crossing geometry is now completely explicit.  Complete-triad
cancellation converts three differently timed front observations into at
most two time-window increment differences, but a direct positive estimate
is the unbudgeted gross triadic action.  The proposed low-leg payment fails
for an exact reason: selecting triads by their high legs leaves the other
triads incident to the same low mode as a complementary circulation term.
Going further upstream also closes the helicity-blind barycentric route;
energy is the only positive radial quadratic invariant common to active
homochiral and heterochiral sectors.

The live gain is chirality-specific.  Critical growth is entirely
heterochiral, and the nonlocal interactions whose positive critical work
also injects energy into a common high band place a definite fraction into a
minority-helicity high receiver.  Whether that receiver survives complete
mode assembly is the V23 question.  Until the complementary circulation in
\eqref{eq:v22-partial-triad-modal-balance} is cancelled or paid, the full
stability/global-regularity objective remains open.
\end{assessment}

% =============================================================================
\section*{Version V22 history}
\addcontentsline{toc}{section}{Version V22 history}
% =============================================================================

Classified every radial leg pattern of a complete triad in a paired-cutoff
band and derived the exact two-increment time-skew normal form on an ordered
front.  Proved that direct positive summation is controlled only by gross
triadic action and gave its Wiener/\(H^2\) upper bound, exposing the strong-
norm circularity.  Derived the exact complementary-circulation residual that
prevents a selected low-leg transfer from telescoping into modal endpoint
energy.  Completed the upstream helical audit: proved the one-triad
energy--helicity transfer normal form, activated both homochiral and
heterochiral sectors on the V5 geometry, and showed that no nonconstant
helicity-blind radial quadratic invariant survives both.  Finally proved
that critical production is purely heterochiral and obtained an exact
minority-helicity sign table plus a quantitative minority-high receiver
fraction for nonlocal high--high/low triads.  The full stability/global-
regularity goal remains open.

% =============================================================================
% VERSION V23 APPEND
% =============================================================================

\section{V23: compensated helical flux and the common receiver corridor}
\label{sec:v23-entry}
% =============================================================================

V22 isolated a minority-helicity high receiver inside every productive
nonlocal high--high/low triad whose critical work agrees in sign with the
high-band injection.  The remaining issue was complete assembly at a shared
mode.  V23 therefore projects the equation onto the two helicity sectors
\emph{before} forming annular or cutoff balances.

The naive cumulative flux in one helicity sector is not a boundary current:
energy may pass between the two sectors, so its value at infinite cutoff need
not vanish.  Subtracting that infinite-cutoff value gives the correct
compensated sector flux.  It is exactly the nonlinear injection into the
sector high tail, its two sector profiles add to the ordinary kinetic flux,
and its radial area is the sector critical work.  Nonlinear helicity
cancellation then gives an exact result on every record doubling: the two
sector critical works are equal, and each is at least \(3R^2/4\).

At the complete-triad level this equality has a geometric realization.  A
positive high-minority interaction produces overlapping positive compensated
flux in both helicities throughout the cutoff interval between its low and
middle radii.  In the nonlocal range, the overlap area is a fixed fraction of
the sector critical work.  A positive low-minority interaction instead has
disjoint positive sector intervals.  Complete aggregation can still cancel
the individual overlaps; V23 records that residual explicitly and does not
promote gross positive-triad overlap to terminal stock.

% =============================================================================
\section{Helical projectors and the compensated sector flux}
\label{sec:v23-compensated-flux}
% =============================================================================

On mean-zero divergence-free fields define the orthogonal helical projectors
\begin{equation}
 \mathsf P^{\pm}
 :=\frac12\left(I\pm\operatorname{curl}\Lambda^{-1}\right),
 \qquad
 u^{\pm}:=\mathsf P^{\pm}u.
 \label{eq:v23-helical-projectors}
\end{equation}
They commute with \(A\), \(\Lambda\), and every radial spectral projector.
For \(s\in\{+,-\}\), let
\begin{equation}
 T_{k,s}(t)
 :=\operatorname{Re}\left(
   \widehat{\mathsf P^s\BNS(u)}(k,t)\cdot
   \overline{\widehat{u^s}(k,t)}
  \right).
 \label{eq:v23-helical-modal-work}
\end{equation}
Set
\begin{align}
 \Phi_\rho^s(t)&:=\sum_{0<|k|\le\rho}T_{k,s}(t),
 \label{eq:v23-naive-sector-flux}\\
 Z_s(t)&:=\sum_{k\ne0}T_{k,s}(t)
       =\langle\mathsf P^s\BNS(u),u^s\rangle,
 \label{eq:v23-sector-exchange}\\
 \widetilde\Phi_\rho^s(t)&:=\Phi_\rho^s(t)-Z_s(t)
       =-\sum_{|k|>\rho}T_{k,s}(t).
 \label{eq:v23-compensated-sector-flux}
\end{align}

\begin{proposition}[Exact compensated-flux identities]
\label{prop:v23-compensated-flux-identities}
For every smooth time and cutoff \(\rho>0\),
\begin{align}
 Z_+(t)+Z_-(t)&=0,
 \label{eq:v23-sector-exchange-cancellation}\\
 \widetilde\Phi_\rho^+(t)+\widetilde\Phi_\rho^-(t)
 &=\Phi_\rho(t),
 \label{eq:v23-sector-sum-full-flux}\\
 \frac12\frac{\dd}{\dd t}\|Q_\rho u^s(t)\|_2^2
 +\nu\|\Lambda Q_\rho u^s(t)\|_2^2
 &=\widetilde\Phi_\rho^s(t).
 \label{eq:v23-sector-high-balance}
\end{align}
If
\begin{equation}
 \mathcal G_s(t)
 :=-\langle\mathsf P^s\BNS(u),\Lambda u^s\rangle
 =-\sum_{k\ne0}|k|T_{k,s}(t),
 \label{eq:v23-sector-critical-work}
\end{equation}
then
\begin{equation}
 \boxed{\quad
 \mathcal G_s(t)
 =\int_0^\infty\widetilde\Phi_\rho^s(t)\,\dd\rho.
 \quad}
 \label{eq:v23-sector-layer-cake}
\end{equation}
Thus compensation is exactly what converts sector energy exchange into a
decaying high-tail profile.
\end{proposition}

\begin{proof}
Orthogonality of the helical projectors gives
\[
 Z_++Z_-
 =\langle\BNS(u),u\rangle=0.
\]
The modal identity \(T_k=T_{k,+}+T_{k,-}\), together with the preceding
cancellation, gives \eqref{eq:v23-sector-sum-full-flux}.

Apply \(Q_\rho\mathsf P^s\) to \eqref{eq:NS} and pair with
\(Q_\rho u^s\).  The nonlinear pairing is
\[
 -\sum_{|k|>\rho}T_{k,s}
 =\widetilde\Phi_\rho^s,
\]
which proves \eqref{eq:v23-sector-high-balance}.  Finally, smooth Fourier
decay permits Fubini, and
\[
 \int_0^\infty\widetilde\Phi_\rho^s\,\dd\rho
 =-\sum_{k\ne0}T_{k,s}\int_0^{|k|}\dd\rho
 =-\sum_{k\ne0}|k|T_{k,s}.
\]
This is \eqref{eq:v23-sector-layer-cake}.
\end{proof}

\begin{firewall}[Why the uncompensated sector flux is the wrong object]
\label{fw:v23-naive-sector-flux}
In general \(Z_s\ne0\), because kinetic energy is conserved only after the
two helicity sectors are summed.  Hence
\(\Phi_\rho^s\to Z_s\) as \(\rho\to\infty\), and
\(\int_0^\infty\Phi_\rho^s\,\dd\rho\) need not exist.  Treating
\(\Phi_\rho^s\) as a pure boundary flux would omit the complete sector-
exchange term.  Equations \eqref{eq:v23-compensated-sector-flux} and
\eqref{eq:v23-sector-high-balance} show that
\(\widetilde\Phi_\rho^s\), not \(\Phi_\rho^s\), is the assembly-compatible
quantity.
\end{firewall}

% =============================================================================
\section{Exact helical paired-cutoff balances and a labelled front branch}
\label{sec:v23-helical-annuli}
% =============================================================================

For \(0<L<M\), define
\begin{align}
 \mathcal E_{L,M}^s(t)
 &:=\|\Pi_{L,M}u^s(t)\|_2^2,
 \label{eq:v23-sector-annular-stock}\\
 \mathcal D_{L,M;I}^s
 &:=\int_I\|\Lambda\Pi_{L,M}u^s(t)\|_2^2\,\dd t,
 \label{eq:v23-sector-annular-dissipation}\\
 \Psi_{L,M}^s(t)
 &:=-\langle
   \Pi_{L,M}\mathsf P^s\BNS(u(t)),
   \Pi_{L,M}u^s(t)
  \rangle.
 \label{eq:v23-sector-annular-injection}
\end{align}

\begin{theorem}[Helicity-resolved paired-cutoff balance]
\label{thm:v23-sector-paired-balance}
For each \(s\in\{+,-\}\),
\begin{align}
 \Psi_{L,M}^s(t)
 &=\widetilde\Phi_L^s(t)-\widetilde\Phi_M^s(t),
 \label{eq:v23-sector-paired-flux}\\
 \frac12\frac{\dd}{\dd t}\mathcal E_{L,M}^s(t)
 +\nu\|\Lambda\Pi_{L,M}u^s(t)\|_2^2
 &=\Psi_{L,M}^s(t).
 \label{eq:v23-sector-instant-band-balance}
\end{align}
Consequently, for \(I=[a,b]\),
\begin{equation}
 \boxed{\quad
 \int_I\bigl(
   \widetilde\Phi_L^s(t)-\widetilde\Phi_M^s(t)
  \bigr)\,\dd t
 =\frac12\bigl(
   \mathcal E_{L,M}^s(b)-\mathcal E_{L,M}^s(a)
  \bigr)
  +\nu\mathcal D_{L,M;I}^s.
 \quad}
 \label{eq:v23-sector-integrated-band-balance}
\end{equation}
Moreover
\begin{equation}
 \Psi_{L,M}^++\Psi_{L,M}^-=\Psi_{L,M},
 \qquad
 \mathcal E_{L,M}^++\mathcal E_{L,M}^-=\mathcal E_{L,M}.
 \label{eq:v23-sector-recombination}
\end{equation}
\end{theorem}

\begin{proof}
The compensation term \(Z_s\) cancels when the two cutoffs are subtracted:
\[
 \widetilde\Phi_L^s-\widetilde\Phi_M^s
 =\Phi_L^s-\Phi_M^s
 =-\sum_{L<|k|\le M}T_{k,s}
 =\Psi_{L,M}^s.
\]
Subtract \eqref{eq:v23-sector-high-balance} at \(M\) from the same identity
at \(L\).  Orthogonality and nesting of the radial projectors give
\eqref{eq:v23-sector-instant-band-balance}.  Integrate in time to obtain
\eqref{eq:v23-sector-integrated-band-balance}.  Finally sum the two
orthogonal helicity sectors and use
\eqref{eq:v23-sector-sum-full-flux}; this proves
\eqref{eq:v23-sector-recombination}.
\end{proof}

\begin{corollary}[Every terminal band has a helicity-labelled certificate]
\label{cor:v23-labelled-band-certificate}
Let
\(J=\mathcal E_{L,M}(b)>0\).  There is a sign
\(s_*\in\{+,-\}\) for which at least one of
\begin{align}
 \mathcal E_{L,M}^{s_*}(a)&\ge\frac J4,
 \tag{H-I}\label{eq:v23-sector-inherited}\\
 \int_I\bigl(
  \widetilde\Phi_L^{s_*}-\widetilde\Phi_M^{s_*}
 \bigr)\,\dd t
 &\ge\frac J8+\nu\mathcal D_{L,M;I}^{s_*}
 \tag{H-N}\label{eq:v23-sector-innovated}
\end{align}
holds.  Therefore an infinite family of terminal bands has an infinite
subfamily carrying one fixed helicity label and one fixed alternative.
\end{corollary}

\begin{proof}
By \eqref{eq:v23-sector-recombination}, choose \(s_*\) so that
\(J_{s_*}:=\mathcal E_{L,M}^{s_*}(b)\ge J/2\).  If
\(\mathcal E_{L,M}^{s_*}(a)\ge J_{s_*}/2\), then
\eqref{eq:v23-sector-inherited} holds.  Otherwise
\eqref{eq:v23-sector-integrated-band-balance} gives
\[
 \int_I(\widetilde\Phi_L^{s_*}-\widetilde\Phi_M^{s_*})\,\dd t
 >\frac{J_{s_*}}4+\nu\mathcal D_{L,M;I}^{s_*}
 \ge\frac J8+\nu\mathcal D_{L,M;I}^{s_*},
\]
which is \eqref{eq:v23-sector-innovated}.  There are only four
sign/alternative labels, so the last assertion follows from the infinite
pigeonhole principle.
\end{proof}

\begin{remark}[The V20 weak ancestry remains valid but does not strengthen]
\label{rem:v23-sector-weak-ancestry}
For a fixed helicity sign, the projectors
\(\Pi_j\mathsf P^s\) on disjoint front bands remain orthogonal and commute
with the heat semigroup.  Thus Theorem~\ref{thm:v20-common-root-annuli}
applies verbatim after helical projection, with source
\(\mathsf P^s\BNS(u)\).  Since \(\mathsf P^s\) is an orthogonal contraction,
the same weak-source budget holds.  It still returns the inverse-frequency
weight in \eqref{eq:v20-monotone-front-weak-ledger}; helicity labelling alone
does not repair that analytic edge loss.
\end{remark}

% =============================================================================
\section{Every record doubling forces equal critical work in both helicities}
\label{sec:v23-equal-sector-work}
% =============================================================================

Define the two positive critical stocks and their viscous actions by
\begin{align}
 X_s(t)&:=\|\Lambda^{1/2}u^s(t)\|_2^2,
 \label{eq:v23-sector-critical-stock}\\
 \mathcal P_s(I)&:=\int_I\|\Lambda^{3/2}u^s(t)\|_2^2\,\dd t.
 \label{eq:v23-sector-critical-dissipation}
\end{align}
Then
\begin{equation}
 \mathcal R(t)^2=X_+(t)+X_-(t),
 \qquad
 \langle u,\operatorname{curl}u\rangle=X_+(t)-X_-(t).
 \label{eq:v23-absolute-signed-helicity}
\end{equation}

\begin{theorem}[Equal-sector work identity on a first doubling]
\label{thm:v23-equal-sector-work}
At every smooth time,
\begin{align}
 \frac12X_s'(t)+\nu\|\Lambda^{3/2}u^s(t)\|_2^2
 &=\mathcal G_s(t),
 \label{eq:v23-sector-critical-balance}\\
 \mathcal G_+(t)&=\mathcal G_-(t).
 \label{eq:v23-equal-instant-sector-work}
\end{align}
Let \(I=[a,b]\) be a first critical-record doubling corridor with
\begin{equation}
 \mathcal R(a)=R,
 \qquad
 \mathcal R(b)=2R,
 \qquad
 K:=\frac{R^2}{E_*}.
 \label{eq:v23-doubling-scales}
\end{equation}
Then, for each sign,
\begin{equation}
 \boxed{\quad
 \int_I\mathcal G_s(t)\,\dd t
 =\frac34R^2
  +\frac\nu2\bigl(\mathcal P_+(I)+\mathcal P_-(I)\bigr)
 \ge\frac34R^2.
 \quad}
 \label{eq:v23-sector-work-lower}
\end{equation}
Thus record growth cannot be assigned to only one helicity sector even if one
sector critical stock decreases between the endpoints.
\end{theorem}

\begin{proof}
Pair the \(s\)-projected equation with \(\Lambda u^s\) to prove
\eqref{eq:v23-sector-critical-balance}.  Nonlinear helicity cancellation is
\[
 0=\langle\BNS(u),\operatorname{curl}u\rangle
  =\langle\mathsf P^+\BNS(u),\Lambda u^+\rangle
   -\langle\mathsf P^-\BNS(u),\Lambda u^-\rangle.
\]
Using the sign convention \eqref{eq:v23-sector-critical-work} proves
\eqref{eq:v23-equal-instant-sector-work}.

Integrate and sum the two sector balances.  By
\eqref{eq:v23-doubling-scales},
\[
 \frac12\bigl(
  X_+(b)+X_-(b)-X_+(a)-X_-(a)
 \bigr)=\frac32R^2.
\]
The two nonlinear integrals are equal by
\eqref{eq:v23-equal-instant-sector-work}.  Dividing their summed identity by
two gives \eqref{eq:v23-sector-work-lower}.
\end{proof}

\begin{corollary}[Each helicity has a cofinal positive-flux width]
\label{cor:v23-sector-positive-width}
For the doubling corridor \(I\), set
\begin{equation}
 \widetilde F_I^s(\rho)
 :=\int_I\widetilde\Phi_\rho^s(t)\,\dd t.
 \label{eq:v23-integrated-sector-profile}
\end{equation}
Then
\begin{align}
 -\frac{E_*}{2}
 \le\widetilde F_I^s(\rho)&\le E_*,
 \label{eq:v23-sector-height-cap}\\
 \int_0^\infty\widetilde F_I^s(\rho)\,\dd\rho
 &\ge\frac34R^2.
 \label{eq:v23-sector-area-lower}
\end{align}
For every \(0\le\alpha<3/4\),
\begin{align}
 \int_{\alpha K}^\infty
  \bigl(\widetilde F_I^s(\rho)\bigr)_+\,\dd\rho
 &\ge\left(\frac34-\alpha\right)R^2,
 \label{eq:v23-sector-tail-positive-area}\\
 \left|\left\{\rho\ge\alpha K:
  \widetilde F_I^s(\rho)>0\right\}\right|
 &\ge\left(\frac34-\alpha\right)K.
 \label{eq:v23-sector-positive-width}
\end{align}
In particular, each helicity has positive compensated flux on a set of cutoff
width at least \(K/4\) above \(K/2\).
\end{corollary}

\begin{proof}
Integrating \eqref{eq:v23-sector-high-balance} gives
\begin{equation}
 \widetilde F_I^s(\rho)
 =\frac12\bigl(
  \|Q_\rho u^s(b)\|_2^2-\|Q_\rho u^s(a)\|_2^2
 \bigr)
 +\nu\int_I\|\Lambda Q_\rho u^s\|_2^2\,\dd t.
 \label{eq:v23-integrated-sector-high-balance}
\end{equation}
The terminal energy term is at most \(E_*/2\), and the total Leray
dissipation on \(I\) is at most \(E_*/2\).  This proves the upper bound in
\eqref{eq:v23-sector-height-cap}.  Dropping the nonnegative dissipative and
terminal terms and bounding the initial sector tail by \(E_*\) proves the
lower bound.

Fubini, \eqref{eq:v23-sector-layer-cake}, and
Theorem~\ref{thm:v23-equal-sector-work} give
\eqref{eq:v23-sector-area-lower}.  The height cap and
\(K E_*=R^2\) imply
\[
 \int_{\alpha K}^\infty\widetilde F_I^s(\rho)\,\dd\rho
 \ge\frac34R^2-\alpha K E_*
 =\left(\frac34-\alpha\right)R^2.
\]
Taking the positive part only increases the left side, which proves
\eqref{eq:v23-sector-tail-positive-area}.  Its integrand is at most \(E_*\),
so division by \(E_*\) proves \eqref{eq:v23-sector-positive-width}.
\end{proof}

% =============================================================================
\section{One productive helical triad has an exact sector-overlap profile}
\label{sec:v23-triad-overlap}
% =============================================================================

For one helical triad define its compensated sector profile by
\begin{equation}
 \widetilde\Phi_\rho^{s,\triangle}
 :=-\sum_{\{i:s_i=s,\ \lambda_i>\rho\}}
       T_i^{\triangle,s}.
 \label{eq:v23-triad-sector-profile}
\end{equation}
Its two profiles sum to the ordinary complete-triad profile
\eqref{eq:v18-triad-flux-profile}.

\begin{theorem}[Positive-profile table by minority-helicity location]
\label{thm:v23-positive-sector-profile-table}
Assume
\begin{equation}
 0<\lambda_1<\lambda_2<\lambda_3
 \label{eq:v23-distinct-triad-radii}
\end{equation}
and let a mixed-helicity triad satisfy
\(\mathcal G_{1/2}^{\triangle,s}>0\).  Reverse all helicity signs and absorb
the resulting common sign into \(\Xi\), if needed, so that the minority sign
is negative.  Let \(m\) be the minority leg and \(a:=|\Xi|>0\).  Denote the
minority and majority compensated profiles by
\(\widetilde\Phi^{\rm min}\) and
\(\widetilde\Phi^{\rm maj}\).  Up to cutoff endpoints of measure zero, their
positive parts are exactly
\begin{equation}
\begin{array}{c|c|c}
 m
 &(\widetilde\Phi_\rho^{\rm min})_+
 &(\widetilde\Phi_\rho^{\rm maj})_+\rule{0pt}{3ex}\\
\hline
1
 &(\lambda_3-\lambda_2)a\,
     \mathbf1_{[0,\lambda_1)}(\rho)
 &(\lambda_1+\lambda_2)a\,
     \mathbf1_{[\lambda_2,\lambda_3)}(\rho)\\[0.4em]
2
 &(\lambda_3-\lambda_1)a\,
     \mathbf1_{[0,\lambda_2)}(\rho)
 &(\lambda_1+\lambda_2)a\,
     \mathbf1_{[\lambda_1,\lambda_3)}(\rho)\\[0.4em]
3
 &(\lambda_2-\lambda_1)a\,
     \mathbf1_{[0,\lambda_3)}(\rho)
 &(\lambda_3+\lambda_1)a\,
     \mathbf1_{[\lambda_1,\lambda_2)}(\rho).
\end{array}
\label{eq:v23-positive-sector-profile-table}
\end{equation}
Thus the two positive profiles are disjoint for a low-minority triad.  For
either high-minority triad they overlap throughout
\([\lambda_1,\lambda_2)\).
\end{theorem}

\begin{proof}
Use the three work vectors in the proof of
Theorem~\ref{thm:v22-minority-helicity-table}.  If \(m=1\), positivity of
critical work gives \(\Xi=a\), and
\[
 T_1=-(\lambda_3-\lambda_2)a,
 \quad
 T_2=(\lambda_3+\lambda_1)a,
 \quad
 T_3=-(\lambda_1+\lambda_2)a.
\]
The minority sector consists only of leg one, so its compensated profile is
\(-T_1\) below \(\lambda_1\) and zero above.  The majority tail sum is
negative below \(\lambda_2\) and equals \(-T_3>0\) between
\(\lambda_2\) and \(\lambda_3\).  This proves the first row.

If \(m=2\), positivity gives \(\Xi=-a\), hence
\[
 T_1=(\lambda_2+\lambda_3)a,
 \quad
 T_2=-(\lambda_3-\lambda_1)a,
 \quad
 T_3=-(\lambda_1+\lambda_2)a.
\]
The minority profile equals \(-T_2>0\) below \(\lambda_2\).  The majority
profile is negative below \(\lambda_1\), equals \(-T_3>0\) from
\(\lambda_1\) to \(\lambda_3\), and then vanishes.  This is the second row.

If \(m=3\), positivity gives \(\Xi=a\), and
\[
 T_1=(\lambda_2+\lambda_3)a,
 \quad
 T_2=-(\lambda_3+\lambda_1)a,
 \quad
 T_3=-(\lambda_2-\lambda_1)a.
\]
The minority profile equals \(-T_3>0\) below \(\lambda_3\).  The majority
profile is negative below \(\lambda_1\), equals \(-T_2>0\) between
\(\lambda_1\) and \(\lambda_2\), and vanishes above.  This proves the last
row and the overlap assertions.
\end{proof}

\begin{corollary}[Quantitative common receiver corridor]
\label{cor:v23-common-receiver-corridor}
Define the positive sector-overlap area of one triad by
\begin{equation}
 \mathfrak O^{\triangle}
 :=\int_0^\infty
   \min\left\{
    (\widetilde\Phi_\rho^{+,\triangle})_+,
    (\widetilde\Phi_\rho^{-,\triangle})_+
   \right\}\,\dd\rho.
 \label{eq:v23-triad-overlap-area}
\end{equation}
Let
\(W^{\triangle}:=\mathcal G_{1/2}^{\triangle,s}/2\), the common critical
work of either helicity sector.  Then
\begin{equation}
 \mathfrak O^{\triangle}=
 \begin{cases}
  0,&m=1,\\
  (\lambda_2-\lambda_1)(\lambda_3-\lambda_1)a,
     &m=2,\\
  (\lambda_2-\lambda_1)^2a,&m=3,
 \end{cases}
 \label{eq:v23-exact-overlap-area}
\end{equation}
whereas
\begin{equation}
 W^{\triangle}=
 \begin{cases}
  \lambda_1(\lambda_3-\lambda_2)a,&m=1,\\
  \lambda_2(\lambda_3-\lambda_1)a,&m=2,\\
  \lambda_3(\lambda_2-\lambda_1)a,&m=3.
 \end{cases}
 \label{eq:v23-common-sector-triad-work}
\end{equation}
If \(m\in\{2,3\}\) and
\(\lambda_1\le\delta\lambda_2\) with \(0<\delta<1\), then
\begin{equation}
 \boxed{\quad
 \mathfrak O^{\triangle}
 \ge\frac{1-\delta}{1+\delta}\,W^{\triangle}.
 \quad}
 \label{eq:v23-overlap-fraction}
\end{equation}
In particular, the overlap is at least one third of the common sector work
when \(\lambda_1\le\lambda_2/2\).
\end{corollary}

\begin{proof}
For \(m=1\), the positive supports in the first row of
\eqref{eq:v23-positive-sector-profile-table} are disjoint.  For \(m=2\), the
common interval is \([\lambda_1,\lambda_2)\).  The minority height
\((\lambda_3-\lambda_1)a\) is no larger than the majority height
\((\lambda_1+\lambda_2)a\), because
\(\lambda_3\le\lambda_1+\lambda_2\).  This gives the second row of
\eqref{eq:v23-exact-overlap-area}.  For \(m=3\), the minority height
\((\lambda_2-\lambda_1)a\) is the smaller height on the same common interval,
giving the third row.

Equation~\eqref{eq:v23-common-sector-triad-work} is one half of
\eqref{eq:v22-minority-critical-table}, with the sign of \(\Xi\) chosen for
positive work.  Therefore
\[
 \frac{\mathfrak O^{\triangle}}{W^{\triangle}}
 =\frac{\lambda_2-\lambda_1}{\lambda_2}
 \quad(m=2),
 \qquad
 \frac{\mathfrak O^{\triangle}}{W^{\triangle}}
 =\frac{\lambda_2-\lambda_1}{\lambda_3}
 \quad(m=3).
\]
The first ratio is at least \(1-\delta\).  For the second, use
\(\lambda_3\le\lambda_2+\lambda_1\le(1+\delta)\lambda_2\); it is at least
\((1-\delta)/(1+\delta)\).  The latter lower bound works in both cases and
proves \eqref{eq:v23-overlap-fraction}.
\end{proof}

\begin{assessment}[What the overlap theorem adds to the V22 receiver bound]
\label{ass:v23-overlap-gain}
Corollary~\ref{cor:v22-minority-high-receiver} measured injection into one
minority high leg relative to the common high-band energy injection.
Corollary~\ref{cor:v23-common-receiver-corridor} is different: it constructs
a whole cutoff interval on which the \emph{compensated fluxes of both
helicities are simultaneously positive}.  Its radial integral is comparable
to the critical work of either sector, with a one-third lower fraction in the
half-nonlocal class.  This is the desired triad-level common receiver.  The
next question is whether it survives summation over all triads sharing the
same modes and cutoffs.
\end{assessment}

% =============================================================================
\section{Complete aggregation retains an exact overlap-cancellation residual}
\label{sec:v23-overlap-residual}
% =============================================================================

At a fixed time, decompose the sector profiles into complete helical triads:
\begin{equation}
 \widetilde\Phi_\rho^s
 =\sum_{\triangle,\boldsymbol s}
   \widetilde\Phi_\rho^{s,\triangle,\boldsymbol s}.
 \label{eq:v23-sector-profile-triad-sum}
\end{equation}
Let \(\mathfrak P\) be any selected class, for example the positive
high-minority nonlocal triads, and define
\begin{align}
 U_\rho^s&:=\sum_{(\triangle,\boldsymbol s)\in\mathfrak P}
     \widetilde\Phi_\rho^{s,\triangle,\boldsymbol s},
 \label{eq:v23-selected-sector-profile}\\
 C_\rho^s&:=\sum_{(\triangle,\boldsymbol s)\notin\mathfrak P}
     \widetilde\Phi_\rho^{s,\triangle,\boldsymbol s}.
 \label{eq:v23-complement-sector-profile}
\end{align}

\begin{proposition}[Exact aggregate receiver identity]
\label{prop:v23-aggregate-receiver-identity}
For each helicity and cutoff,
\begin{equation}
 \boxed{\quad
 \widetilde\Phi_\rho^s=U_\rho^s+C_\rho^s.
 \quad}
 \label{eq:v23-aggregate-receiver-identity}
\end{equation}
Consequently,
\begin{equation}
 (\widetilde\Phi_\rho^s)_+
 \ge (U_\rho^s)_+-(C_\rho^s)_-,
 \label{eq:v23-overlap-cancellation-pointwise}
\end{equation}
and for the simultaneous positive receiver,
\begin{align}
 \min\{(\widetilde\Phi_\rho^+)_+,
        (\widetilde\Phi_\rho^-)_+\}
 &\ge
 \min\{(U_\rho^+)_+,(U_\rho^-)_+\}
 \notag\\
 &\quad-(C_\rho^+)_--(C_\rho^-)_-.
 \label{eq:v23-overlap-cancellation-residual}
\end{align}
Thus a lower bound obtained by summing selected triad overlaps becomes a net
sector-overlap bound only after the two negative complementary profiles are
paid.
\end{proposition}

\begin{proof}
Equation~\eqref{eq:v23-aggregate-receiver-identity} is the partition of the
absolutely convergent complete-triad sum
\eqref{eq:v23-sector-profile-triad-sum}.  The scalar inequality
\((x+y)_+\ge x_+-y_-\) proves
\eqref{eq:v23-overlap-cancellation-pointwise}.  Apply it in the two sectors
and use, for nonnegative \(a,b,c,d\),
\[
 \min\{(a-c)_+,(b-d)_+\}
 \ge\min\{a,b\}-c-d.
\]
This yields \eqref{eq:v23-overlap-cancellation-residual}.
\end{proof}

\begin{firewall}[Positive-triad overlap cannot be summed before assembly]
\label{fw:v23-positive-triad-overlap}
The quantity
\(\sum_{\triangle:\mathcal G_{1/2}^{\triangle}>0}
  \mathfrak O^{\triangle}\)
takes a positive part before summing complete triads.  It is therefore a
gross-transfer quantity of the same mathematical type as
\eqref{eq:v22-gross-triad-action}.  The equal-sector identity
\eqref{eq:v23-equal-instant-sector-work} is signed and controls only the
assembled sum.  Dropping the two \(C^s\) terms in
\eqref{eq:v23-overlap-cancellation-residual} would repeat the partial-triad
error isolated in Theorem~\ref{thm:v22-partial-triad-balance}.
\end{firewall}

\begin{counterexample}[Equal sector areas and the full flux cap do not force
overlap]
\label{cth:v23-profile-overlap}
The scalar information proved so far cannot by itself force the positive
sets of \(\widetilde F_I^+\) and \(\widetilde F_I^-\) to intersect.  Indeed,
fix \(K=R^2/E_*\) and define the model profiles
\begin{align}
 f_+(\rho)&:=\frac{E_*}{2}
  \mathbf1_{[K/2,\,2K)}(\rho),
 \label{eq:v23-disjoint-model-plus}\\
 f_-(\rho)&:=\frac{E_*}{2}
  \mathbf1_{[2K,\,7K/2)}(\rho).
 \label{eq:v23-disjoint-model-minus}
\end{align}
They have disjoint positive supports up to one endpoint and satisfy
\begin{equation}
 \int_0^\infty f_+\,\dd\rho
 =\int_0^\infty f_-\,\dd\rho
 =\frac34R^2,
 \qquad
 0\le f_++f_-\le\frac{E_*}{2}.
 \label{eq:v23-disjoint-model-data}
\end{equation}
They therefore obey the minimum equal-sector area, both sector height caps,
and the sharper full-flux height cap.  Smooth approximations can preserve
arbitrarily small overlap while changing the equalities by arbitrarily small
amounts.  This is a profile-level logical counterexample, not a claim that
the displayed step functions are realized by a Navier--Stokes trajectory.
It proves that the balances and width bounds alone do not pay the complement
in \eqref{eq:v23-overlap-cancellation-residual}.
\end{counterexample}

\begin{proof}
Each support has length \(3K/2\), so each integral equals
\((E_*/2)(3K/2)=3R^2/4\).  The supports are disjoint except at an endpoint,
which has measure zero, and their sum has height \(E_*/2\).  First shortening the two intervals by an arbitrarily small amount creates a
positive gap; mollifying inside the resulting disjoint open neighborhoods
proves the approximation assertion.
\end{proof}

\begin{assessment}[Packet and V5 audits]
\label{ass:v23-model-audits}
The V5 homochiral tangent has
\(\mathcal G_+=\mathcal G_-=0\), consistent with
Theorem~\ref{thm:v23-equal-sector-work}; it supplies no false positive sector
area.  Activating the mixed \((+,+,-)\) sector on the same geometry puts the
minority sign on the highest leg, so the third row of
\eqref{eq:v23-positive-sector-profile-table} gives a genuine common receiver
interval.  Reversing its phase reverses all work signs, as required.

The V14/V18 concentrated packet permits flux height of order \(R^2/N\) over
width of order \(N\).  The sector area and width conclusions allow the same
height--width trade and therefore do not contradict that family of distinct
solutions.  A valid next step must control cancellation between sector
profiles across one terminal chronology, not impose a per-card height floor.
\end{assessment}

% =============================================================================
\section{Integrated V23 conclusion and V24 acquisition order}
\label{sec:v23-conclusion}
% =============================================================================

\begin{theorem}[What V23 proves]
\label{thm:v23-conclusion}
Relative to V1--V22, the following are proved.
\begin{enumerate}[label=\textup{(V23.\arabic*)},leftmargin=6.2em]
\item Subtracting the total energy exchange of one helicity sector produces
      a compensated cutoff flux which is exactly the nonlinear injection
      into that sector's high tail.
\item The two compensated profiles add to the ordinary kinetic-flux profile,
      and the radial area of either is its sector critical work.
\item Helicity-resolved annular stocks satisfy exact paired-cutoff balances.
      Every terminal band has an inherited-or-innovated certificate carrying
      one helicity label, at the explicit factors \(1/4\) and \(1/8\).
\item Nonlinear helicity cancellation makes the two sector critical works
      equal pointwise.  On every first record doubling, each integrated work
      is at least \(3R^2/4\).
\item Each helicity sector has positive compensated flux on a cutoff set of
      width at least \((3/4-\alpha)K\) above \(\alpha K\), for every
      \(0\le\alpha<3/4\).
\item A positive high-minority triad has simultaneous positive sector flux
      on \([\lambda_1,\lambda_2)\).  In the half-nonlocal class, its overlap
      area is at least one third of the common sector critical work.
\item A positive low-minority triad has disjoint positive sector profiles.
      After complete aggregation, selected common receivers retain the exact
      negative complementary-circulation residual
      \eqref{eq:v23-overlap-cancellation-residual}.
\item Equal sector areas, the sector height bounds, and even the sharper full
      flux cap do not alone force the two positive sector sets to overlap.
\end{enumerate}
No full Navier--Stokes stability or global-regularity theorem is obtained.
\end{theorem}

\begin{proof}
Items \textup{(V23.1)--(V23.2)} are
Proposition~\ref{prop:v23-compensated-flux-identities}.  Item
\textup{(V23.3)} is Theorem~\ref{thm:v23-sector-paired-balance} and
Corollary~\ref{cor:v23-labelled-band-certificate}.  Items
\textup{(V23.4)--(V23.5)} are
Theorem~\ref{thm:v23-equal-sector-work} and
Corollary~\ref{cor:v23-sector-positive-width}.  Items
\textup{(V23.6)--(V23.7)} are
Theorem~\ref{thm:v23-positive-sector-profile-table},
Corollary~\ref{cor:v23-common-receiver-corridor}, and
Proposition~\ref{prop:v23-aggregate-receiver-identity}.  Item
\textup{(V23.8)} is Counterexample~\ref{cth:v23-profile-overlap}.
\end{proof}

\begin{programme}[V24 acquisition order]
\label{prog:v24}
\begin{enumerate}[label=\textup{(AC\arabic*)},leftmargin=4.8em]
\item Use only the compensated sector profiles
      \(\widetilde F_I^s\).  The uncompensated profile omits the exact
      sector exchange \(Z_s\) and is no longer admissible.
\item Construct positive-area quantiles separately in the two helicities
      using \eqref{eq:v23-sector-tail-positive-area}.  Record their radial
      separation before attempting a common cutoff.
\item Split the chronology into an overlap branch and a separation branch.
      On the overlap branch, preserve the signed paired-sector injection and
      seek a common annular terminal stock.
\item On the separation branch, integrate the difference profile
      \(\widetilde F_I^+-\widetilde F_I^-\), whose total radial area is zero.
      Determine whether moving equal sector area across a radial gap forces a
      signed first-moment or dissipation debit.
\item Resolve the complementary terms \(C_\rho^s\) in
      \eqref{eq:v23-overlap-cancellation-residual} by complete mode assembly.
      Do not sum positive triad overlaps or minority receivers separately.
\item Treat low-minority triads as part of the separation branch; their exact
      positive supports in \eqref{eq:v23-positive-sector-profile-table} show
      that they cannot be charged by a common-receiver assertion.
\item Test every proposed sector-separation debit on the disjoint profile
      counterexample and the concentrated packet.  A useful PDE input must
      exclude them dynamically, not only by scalar area and height data.
\item Prepare the exact open quantities for the required SM/YM
      synchronization at V25.
\end{enumerate}
\end{programme}

\begin{assessment}[Position after V23]
The helical receiver proposed in V22 now has a complete PDE-level flux
representation.  The compensation by \(Z_s\) is forced, the two sector
critical works are exactly equal, and each doubled record produces large
positive radial area in both helicities.  At the triad level, the dominant
nonlocal high-minority class has a quantitative common receiver corridor.

These facts do not yet yield a common assembled cutoff.  Low-minority triads
can separate the two positive profiles, and other incident triads can cancel
the selected high-minority overlap through the exact residual
\eqref{eq:v23-overlap-cancellation-residual}.  The live V24 dichotomy is
therefore radial overlap versus sector separation.  A closure step must pay
one of those two assembled branches without returning to gross triadic
action, the inverse-frequency source ledger, or an assumed flux-height floor.
The full stability/global-regularity objective remains open.
\end{assessment}

% =============================================================================
\section*{Version V23 history}
\addcontentsline{toc}{section}{Version V23 history}
% =============================================================================

Introduced the compensated helical cutoff flux and proved its exact high-tail,
recombination, and layer-cake identities.  Derived helicity-resolved annular
balances and a labelled inherited-or-innovation certificate for every
terminal band.  Proved that the nonlinear critical works of the two helicity
sectors are equal pointwise and that each receives at least \(3R^2/4\) on a
first record doubling, yielding sectorwise cofinal positive-flux width.
Computed the complete positive compensated-flux profile of every productive
mixed-helicity triad.  High-minority nonlocal triads have a quantitative
common receiver corridor, while low-minority triads have disjoint positive
sector supports.  Finally isolated the exact complementary-circulation
residual after complete aggregation and gave a profile-level counterexample
showing that equal area and height bounds alone do not force sector overlap.
The full stability/global-regularity goal remains open.

% =============================================================================
% VERSION V24 APPEND
% =============================================================================

\section{V24: assembled sector overlap versus signed radial separation}
\label{sec:v24-entry}
% =============================================================================

V23 produced two exact compensated helical flux profiles on every first
critical-record doubling.  Their signed radial masses are equal and at least
\(3R^2/4\), while positive high-minority triads possess a common receiver
corridor.  The missing step was to distinguish a genuine overlap of the
\emph{assembled} sector profiles from a cancellation-driven separation of
their positive parts.

V24 makes that dichotomy quantitative.  Large assembled overlap is paid by
either a common terminal chiral-tail stock or the full critical viscous
action.  Small overlap forces order-\(R^2\) positive and negative mass in the
zero-mean difference of the two sector profiles.  The latter difference has
an exact high-tail balance.  Its first radial moment is the signed production
of an \(H^1\)-level helicity, so a direct absolute estimate reaches endpoint
enstrophy and palinstrophy rather than the Leray budget.

The positive-quantile version has one further residual: positive parts and
negative parts must be reassembled before their barycentres can be identified
with the signed PDE moment.  This is the sector-profile analogue of the
partial-triad residual in V22.  The calculation therefore supplies exact
objects for the scheduled V25 companion review without claiming a closure of
either branch.

% =============================================================================
\section{A quantitative assembled overlap--separation dichotomy}
\label{sec:v24-overlap-separation}
% =============================================================================

Fix a first-doubling corridor \(I=[a,b]\) as in
\eqref{eq:v23-doubling-scales}.  Abbreviate
\begin{align}
 F_s(\rho)&:=\widetilde F_I^s(\rho),
 \label{eq:v24-sector-profile-abbreviation}\\
 P_s(\rho)&:=(F_s(\rho))_+,
 &N_s(\rho)&:=(-F_s(\rho))_+,
 \label{eq:v24-sector-Jordan-parts}\\
 W_I&:=\int_0^\infty F_+(\rho)\,\dd\rho
     =\int_0^\infty F_-(\rho)\,\dd\rho.
 \label{eq:v24-common-signed-mass}
\end{align}
The equality follows from
\eqref{eq:v23-equal-instant-sector-work}, and
\begin{equation}
 W_I\ge\frac34R^2.
 \label{eq:v24-common-mass-lower}
\end{equation}
Define the assembled positive-overlap mass and sector difference by
\begin{align}
 \Omega_I&:=\int_0^\infty
   \min\{P_+(\rho),P_-(\rho)\}\,\dd\rho,
 \label{eq:v24-assembled-overlap-mass}\\
 D_I(\rho)&:=F_+(\rho)-F_-(\rho).
 \label{eq:v24-sector-difference-profile}
\end{align}

\begin{theorem}[Assembled overlap--separation dichotomy]
\label{thm:v24-overlap-separation}
For every \(0\le\eta<3/4\), at least one of
\begin{align}
 \boxed{\quad \Omega_I\ge\eta R^2,\quad}
 &\tag{O}\label{eq:v24-overlap-branch}\\
 \boxed{\quad
  \int_0^\infty|D_I(\rho)|\,\dd\rho
  \ge\left(\frac32-2\eta\right)R^2
  \quad}
 &\tag{S}\label{eq:v24-separation-branch}
\end{align}
holds.  On branch \textup{(S)},
\begin{align}
 \int_0^\infty(D_I)_+\,\dd\rho
 =\int_0^\infty(D_I)_-\,\dd\rho
 &\ge\left(\frac34-\eta\right)R^2,
 \label{eq:v24-difference-Jordan-mass}\\
 \left|\{\rho:(D_I(\rho))_+>0\}\right|,
 \quad
 \left|\{\rho:(D_I(\rho))_->0\}\right|
 &\ge\left(\frac12-\frac{2\eta}{3}\right)K.
 \label{eq:v24-difference-sign-widths}
\end{align}
Thus failure of a common positive receiver creates two oppositely signed
sector-difference regions of canonical total width.
\end{theorem}

\begin{proof}
Let \(M_s:=\int P_s\) and \(L_s:=\int N_s\).  The Jordan decomposition and
\eqref{eq:v24-common-signed-mass} give
\begin{equation}
 M_s-L_s=W_I,
 \qquad M_s\ge W_I.
 \label{eq:v24-positive-mass-lower}
\end{equation}
Pointwise,
\(|x-y|\ge|x_+-y_+|\).  Hence
\begin{align*}
 \int|D_I|
 &\ge\int|P_+-P_-|\\
 &=M_++M_--2\int\min\{P_+,P_-\}\\
 &\ge2W_I-2\Omega_I.
\end{align*}
If branch \textup{(O)} fails, use
\eqref{eq:v24-common-mass-lower} to obtain
\eqref{eq:v24-separation-branch}.

The two signed masses of \(D_I\) are equal because
\begin{equation}
 \int_0^\infty D_I(\rho)\,\dd\rho=0.
 \label{eq:v24-difference-zero-mass}
\end{equation}
Each is half of \(\int|D_I|\), proving
\eqref{eq:v24-difference-Jordan-mass}.  Finally
\eqref{eq:v23-sector-height-cap} gives
\begin{equation}
 |D_I(\rho)|\le\frac32E_*.
 \label{eq:v24-difference-height-cap}
\end{equation}
Divide each lower mass in
\eqref{eq:v24-difference-Jordan-mass} by \(3E_*/2\) and use
\(R^2=KE_*\).  This proves \eqref{eq:v24-difference-sign-widths}.
\end{proof}

\begin{remark}[The dichotomy uses assembled profiles]
\label{rem:v24-assembled-not-triadwise}
The positive parts in \eqref{eq:v24-assembled-overlap-mass} are taken only
after summing all complete helical triads into \(F_s\).  Thus
Theorem~\ref{thm:v24-overlap-separation} does not invoke the gross positive-
triad overlap excluded by Firewall~\ref{fw:v23-positive-triad-overlap}.
\end{remark}

% =============================================================================
\section{The overlap branch is terminal common stock or critical dissipation}
\label{sec:v24-overlap-payment}
% =============================================================================

For a cutoff \(\rho\), terminal time \(b\), and sign \(s\), put
\begin{align}
 H_\rho^s(b)&:=\|Q_\rho u^s(b)\|_2^2,
 \label{eq:v24-terminal-sector-tail}\\
 V_\rho^s(I)&:=\int_I\|\Lambda Q_\rho u^s(t)\|_2^2\,\dd t.
 \label{eq:v24-sector-tail-dissipation}
\end{align}
Define the common terminal chiral-tail stock
\begin{equation}
 \mathcal C_I(b)
 :=\int_0^\infty
   \min\{H_\rho^+(b),H_\rho^-(b)\}\,\dd\rho.
 \label{eq:v24-common-terminal-tail-stock}
\end{equation}

\begin{theorem}[Exact payment alternatives for assembled sector overlap]
\label{thm:v24-overlap-payment}
The overlap mass obeys
\begin{equation}
 \boxed{\quad
 \Omega_I
 \le\frac12\mathcal C_I(b)
  +\nu\bigl(\mathcal P_+(I)+\mathcal P_-(I)\bigr).
 \quad}
 \label{eq:v24-overlap-payment}
\end{equation}
Moreover,
\begin{equation}
 0\le\mathcal C_I(b)
 \le\min\{X_+(b),X_-(b)\}
 \le2R^2.
 \label{eq:v24-common-tail-cap}
\end{equation}
Consequently, on branch \textup{(O)}, at least one of
\begin{align}
 \boxed{\quad \mathcal C_I(b)\ge\eta R^2,\quad}
 &\tag{CT}\label{eq:v24-common-terminal-branch}\\
 \boxed{\quad
  \nu\bigl(\mathcal P_+(I)+\mathcal P_-(I)\bigr)
  \ge\frac\eta2R^2
  \quad}
 &\tag{CD}\label{eq:v24-critical-dissipation-branch}
\end{align}
holds.
\end{theorem}

\begin{proof}
Equation~\eqref{eq:v23-integrated-sector-high-balance} gives
\[
 F_s(\rho)
 \le\frac12H_\rho^s(b)+\nu V_\rho^s(I),
\]
whose right side is nonnegative.  Therefore
\[
 P_s(\rho)
 \le\frac12H_\rho^s(b)+\nu V_\rho^s(I).
\]
For nonnegative \(a,b,c,d\),
\(\min\{a+b,c+d\}\le\min\{a,c\}+b+d\).  Apply this inequality to the two
sector bounds and integrate in \(\rho\):
\begin{align*}
 \Omega_I
 &\le\frac12\int_0^\infty
      \min\{H_\rho^+(b),H_\rho^-(b)\}\,\dd\rho\\
 &\quad+\nu\sum_{s\in\{+,-\}}
     \int_0^\infty V_\rho^s(I)\,\dd\rho.
\end{align*}
Spectral Fubini gives
\begin{equation}
 \int_0^\infty V_\rho^s(I)\,\dd\rho
 =\int_I\|\Lambda^{3/2}u^s(t)\|_2^2\,\dd t
 =\mathcal P_s(I),
 \label{eq:v24-tail-dissipation-layer-cake}
\end{equation}
which proves \eqref{eq:v24-overlap-payment}.

Similarly,
\[
 \int_0^\infty H_\rho^s(b)\,\dd\rho=X_s(b).
\]
The integral of a pointwise minimum is bounded by the minimum of the two
integrals.  Since \(X_+(b)+X_-(b)=4R^2\), this proves
\eqref{eq:v24-common-tail-cap}.

If \(\Omega_I\ge\eta R^2\) but
\(\mathcal C_I(b)<\eta R^2\), then
\eqref{eq:v24-overlap-payment} forces
\[
 \nu(\mathcal P_++\mathcal P_-)
 >\eta R^2-\frac12\eta R^2
 =\frac\eta2R^2.
\]
This proves the two alternatives.
\end{proof}

\begin{firewall}[Critical dissipation is not the Leray debit]
\label{fw:v24-critical-dissipation-unpaid}
The action in \eqref{eq:v24-critical-dissipation-branch} is
\(\nu\int_I\|\Lambda^{3/2}u\|_2^2\), whereas the global energy inequality
pays only \(\nu\int_I\|\Lambda u\|_2^2\).  The former is exactly the
dissipative term in the critical-record balance
\eqref{eq:v5-critical-record-flux}.  It can be large on every doubling
corridor because it is replenished by the same nonlinear critical work in
\eqref{eq:v23-sector-work-lower}.  Calling branch \textup{(CD)} globally
paid would assume the critical a priori estimate sought by the programme.
\end{firewall}

\begin{assessment}[What remains in the common-terminal branch]
\label{ass:v24-common-terminal-status}
The stock \(\mathcal C_I(b)\) is positive, simultaneous in the two
helicities, and measured at one terminal time.  These properties are stronger
than the triadwise receiver in V23.  They do not yet make stocks from different
doubling times orthogonal: each is a radial integral of nested high tails, so
one Fourier mode is repeated across every lower cutoff.  A V25 import is
useful only if it gives a same-terminal or determinant-type subtraction for
this nested common stock without replacing it by
\(X_+(b)+X_-(b)\).
\end{assessment}

% =============================================================================
\section{The separation branch has an exact signed high-tail balance}
\label{sec:v24-separation-balance}
% =============================================================================

Define the signed helical high-tail stock and dissipation by
\begin{align}
 \mathcal Y_\rho(t)
 &:=\|Q_\rho u^+(t)\|_2^2-
    \|Q_\rho u^-(t)\|_2^2,
 \label{eq:v24-signed-tail-stock}\\
 \mathcal V_\rho(I)
 &:=\int_I\left(
    \|\Lambda Q_\rho u^+(t)\|_2^2-
    \|\Lambda Q_\rho u^-(t)\|_2^2
   \right)\,\dd t.
 \label{eq:v24-signed-tail-dissipation}
\end{align}

\begin{theorem}[Exact sector-difference balance and its available upper
bound]
\label{thm:v24-sector-difference-balance}
For every cutoff,
\begin{equation}
 \boxed{\quad
 D_I(\rho)
 =\frac12\bigl(\mathcal Y_\rho(b)-\mathcal Y_\rho(a)\bigr)
  +\nu\mathcal V_\rho(I).
 \quad}
 \label{eq:v24-sector-difference-balance}
\end{equation}
Its signed radial mass vanishes as in
\eqref{eq:v24-difference-zero-mass}, while its absolute mass satisfies
\begin{equation}
 \boxed{\quad
 \int_0^\infty|D_I(\rho)|\,\dd\rho
 \le\frac12\bigl(
   \mathcal R(a)^2+\mathcal R(b)^2
  \bigr)
  +\nu\bigl(\mathcal P_+(I)+\mathcal P_-(I)\bigr).
 \quad}
 \label{eq:v24-difference-absolute-upper}
\end{equation}
On a first doubling this becomes
\begin{equation}
 \int_0^\infty|D_I|\,\dd\rho
 \le\frac52R^2
  +\nu(\mathcal P_++\mathcal P_-).
 \label{eq:v24-difference-doubling-upper}
\end{equation}
Thus the separation lower bound
\eqref{eq:v24-separation-branch} is compatible with endpoint critical stock
alone; it does not force a Leray debit.
\end{theorem}

\begin{proof}
Subtract the two identities
\eqref{eq:v23-integrated-sector-high-balance} to obtain
\eqref{eq:v24-sector-difference-balance}.  Its zero radial mass follows from
the equality of the two sector critical works.  For the absolute estimate,
use the triangle inequality and
\begin{align*}
 |\mathcal Y_\rho(t)|
 &\le\|Q_\rho u^+(t)\|_2^2+
       \|Q_\rho u^-(t)\|_2^2,\\
 |\mathcal V_\rho(I)|
 &\le\int_I\left(
       \|\Lambda Q_\rho u^+\|_2^2+
       \|\Lambda Q_\rho u^-\|_2^2
      \right)\,\dd t.
\end{align*}
Radial Fubini gives
\[
 \int_0^\infty\sum_s\|Q_\rho u^s(t)\|_2^2\,\dd\rho
 =\sum_sX_s(t)=\mathcal R(t)^2
\]
and gives \eqref{eq:v24-tail-dissipation-layer-cake} for the dissipative
term.  This proves \eqref{eq:v24-difference-absolute-upper}; insert the two
endpoint values to get \eqref{eq:v24-difference-doubling-upper}.
\end{proof}

% =============================================================================
\section{Signed barycentric separation is an \(H^1\)-helicity balance}
\label{sec:v24-barycentric-balance}
% =============================================================================

Define the sector enstrophy works and stocks
\begin{align}
 \mathcal Q_s(t)
 &:=-\langle\mathsf P^s\BNS(u(t)),Au^s(t)\rangle,
 \label{eq:v24-sector-enstrophy-work}\\
 Y_s(t)&:=\|\Lambda u^s(t)\|_2^2,
 \label{eq:v24-sector-enstrophy-stock}\\
 \mathscr H_1(t)&:=Y_+(t)-Y_-(t)
 =\langle\Lambda^{1/2}u(t),
          \operatorname{curl}\Lambda^{1/2}u(t)\rangle.
 \label{eq:v24-H1-helicity}
\end{align}
The last quantity is signed and is not the conserved Euler helicity in
\eqref{eq:v23-absolute-signed-helicity}.

\begin{theorem}[Exact signed sector-barycentre identity]
\label{thm:v24-signed-barycentre}
The first radial moment of one compensated sector profile is
\begin{equation}
 \int_0^\infty\rho F_s(\rho)\,\dd\rho
 =\frac12\int_I\mathcal Q_s(t)\,\dd t.
 \label{eq:v24-sector-first-moment}
\end{equation}
Since \(W_I>0\), define its signed barycentre
\begin{equation}
 B_s(I):=\frac1{W_I}
  \int_0^\infty\rho F_s(\rho)\,\dd\rho.
 \label{eq:v24-signed-sector-barycentre}
\end{equation}
Then
\begin{align}
 \boxed{\quad
 W_I\bigl(B_+(I)-B_-(I)\bigr)
 &=\frac14\bigl(\mathscr H_1(b)-\mathscr H_1(a)\bigr)
 \notag\\
 &\quad+\frac\nu2\int_I\left(
   \|Au^+(t)\|_2^2-\|Au^-(t)\|_2^2
  \right)\,\dd t.
 \quad}
 \label{eq:v24-signed-barycentre-balance}
\end{align}
Consequently,
\begin{align}
 W_I|B_+-B_-|
 &\le\frac14\bigl(
   \|\Lambda u(a)\|_2^2+
   \|\Lambda u(b)\|_2^2
  \bigr)
 \notag\\
 &\quad+\frac\nu2\int_I\|Au(t)\|_2^2\,\dd t.
 \label{eq:v24-barycentre-strong-upper}
\end{align}
The right side is an endpoint-enstrophy/palinstrophy quantity, not a Leray
energy quantity.
\end{theorem}

\begin{proof}
From \eqref{eq:v23-compensated-sector-flux}, radial Fubini gives
\[
 \int_0^\infty\rho\widetilde\Phi_\rho^s(t)\,\dd\rho
 =-\sum_{k\ne0}T_{k,s}(t)
   \int_0^{|k|}\rho\,\dd\rho
 =\frac12\mathcal Q_s(t).
\]
Integrate in time to prove \eqref{eq:v24-sector-first-moment}.

Pair the \(s\)-projected equation with \(Au^s\):
\begin{equation}
 \frac12Y_s'(t)+\nu\|Au^s(t)\|_2^2=\mathcal Q_s(t).
 \label{eq:v24-sector-enstrophy-balance}
\end{equation}
Subtract the two signs, integrate over \(I\), multiply by \(1/2\), and use
\eqref{eq:v24-sector-first-moment}.  This proves
\eqref{eq:v24-signed-barycentre-balance}.  Taking absolute values and using
\(|\mathscr H_1|\le Y_++Y_-\) proves
\eqref{eq:v24-barycentre-strong-upper}.
\end{proof}

\begin{firewall}[The barycentric separation estimate is the classical strong
boundary]
\label{fw:v24-barycentre-strong-boundary}
Neither endpoint enstrophy nor
\(\nu\int_I\|Au\|_2^2\) is bounded by the Leray energy inequality.  The
latter is the dissipation in the enstrophy balance and is coupled to vortex
stretching.  Therefore taking an absolute value in
\eqref{eq:v24-signed-barycentre-balance} returns the classical strong-
solution continuation boundary.  The zero-mass property of \(D_I\) gives an
exact signed identity, but no positive payment for radial transport.
\end{firewall}

% =============================================================================
\section{Positive quantile separation has an additional Jordan residual}
\label{sec:v24-positive-quantile-residual}
% =============================================================================

The signed barycentres in
\eqref{eq:v24-signed-sector-barycentre} need not lie in the positive supports
of \(F_s\), because those profiles may have negative parts.  We record the
exact correction before using positive-flux quantiles.

Let \(0<m\le\min\{M_+,M_-\}\).  Choose measurable densities
\(p_s\) satisfying
\begin{equation}
 0\le p_s\le P_s,
 \qquad
 \int_0^\infty p_s(\rho)\,\dd\rho=m,
 \label{eq:v24-selected-positive-submass}
\end{equation}
and put
\begin{align}
 b_s^{(m)}&:=\frac1m\int_0^\infty\rho p_s(\rho)\,\dd\rho,
 \label{eq:v24-positive-submass-barycentre}\\
 R_s^{(m)}&:=P_s-p_s.
 \label{eq:v24-positive-residual-density}
\end{align}

\begin{proposition}[Exact positive-quantile barycentre ledger]
\label{prop:v24-positive-quantile-ledger}
The selected positive barycentres satisfy
\begin{align}
 \boxed{\quad
 m\bigl(b_+^{(m)}-b_-^{(m)}\bigr)
 &=W_I(B_+-B_-)
 \notag\\
 &\quad-\int_0^\infty\rho
    \bigl(R_+^{(m)}-R_-^{(m)}\bigr)\,\dd\rho
 \notag\\
 &\quad+\int_0^\infty\rho(N_+-N_-)\,\dd\rho.
 \quad}
 \label{eq:v24-positive-quantile-ledger}
\end{align}
If the support of \(p_+\) lies at least a radial distance \(g>0\) below the
support of \(p_-\), then
\begin{equation}
 m\bigl(b_-^{(m)}-b_+^{(m)}\bigr)\ge mg.
 \label{eq:v24-positive-quantile-gap}
\end{equation}
This separation cannot be charged solely to the signed PDE term in
\eqref{eq:v24-signed-barycentre-balance} unless the residual positive and
negative first moments in
\eqref{eq:v24-positive-quantile-ledger} are also controlled.
\end{proposition}

\begin{proof}
The Jordan decomposition gives
\[
 F_s=p_s+R_s^{(m)}-N_s.
\]
Multiply by \(\rho\), integrate, and subtract the two signs.  The left side
from the selected parts is
\(m(b_+^{(m)}-b_-^{(m)})\), while the full signed difference is
\(W_I(B_+-B_-)\).  Rearrangement proves
\eqref{eq:v24-positive-quantile-ledger}.  The support ordering makes every
point charged to \(p_-\) at least \(g\) above every point charged to
\(p_+\); comparison of their averages proves
\eqref{eq:v24-positive-quantile-gap}.
\end{proof}

\begin{firewall}[Dropping the Jordan residual recreates gross spectral
variation]
\label{fw:v24-Jordan-residual}
The last two lines of \eqref{eq:v24-positive-quantile-ledger} contain radial
first moments of positive and negative parts.  Bounding them separately is
an absolute spectral-variation estimate.  It is not supplied by the signed
zero-area identity \eqref{eq:v24-difference-zero-mass} and is of the same
type as the gross action obstruction in V22.  Therefore separated positive
quantiles do not, by themselves, imply a large signed \(H^1\)-helicity
increment.
\end{firewall}

% =============================================================================
\section{Sharp scalar tests of both V24 branches}
\label{sec:v24-tests}
% =============================================================================

\begin{proposition}[The V23 disjoint profiles saturate the separation scale]
\label{prop:v24-disjoint-profile-test}
For the model profiles \eqref{eq:v23-disjoint-model-plus}--%
\eqref{eq:v23-disjoint-model-minus},
\begin{align}
 \Omega_I&=0,
 \label{eq:v24-model-zero-overlap}\\
 \int_0^\infty|D_I(\rho)|\,\dd\rho
 &=\frac32R^2,
 \label{eq:v24-model-separation-mass}\\
 B_+-B_-&=-\frac32K,
 \label{eq:v24-model-barycentre-gap}\\
 W_I|B_+-B_-|
 &=\frac98R^2K
 =\frac98\frac{R^4}{E_*}.
 \label{eq:v24-model-quartic-moment}
\end{align}
Thus the minimum zero-overlap profile already reaches the quartic record
scale of the V4 finite-head estimate.
\end{proposition}

\begin{proof}
The two supports are disjoint and each profile has mass \(3R^2/4\), proving
the first two identities.  Their interval midpoints are respectively
\(5K/4\) and \(11K/4\).  Since the profiles are nonnegative, these are their
signed barycentres.  Their difference is \(-3K/2\).  Multiply by
\(W_I=3R^2/4\) and use \(K=R^2/E_*\) to prove the last identity.
\end{proof}

\begin{assessment}[Packet and triad compatibility]
\label{ass:v24-packet-triad-test}
The disjoint profile test proves that the V24 separation lower bound and
signed barycentric identity permit the same quartic scale already found in
V4.  The V14/V18 packet likewise permits order-\(R^2\) area to be transported
by height \(R^2/N\) across width \(N\); neither
\eqref{eq:v24-difference-sign-widths} nor
\eqref{eq:v24-overlap-payment} imposes a forbidden height floor.

At the one-triad level, V23's low-minority profile occupies the separation
branch and its high-minority profile occupies the overlap branch.  Both are
therefore represented by the dichotomy.  What remains untested by those
finite examples is a PDE estimate on the assembled common terminal stock or
on the Jordan residual across infinitely many record doublings.
\end{assessment}

% =============================================================================
\section{Integrated V24 conclusion and V25 synchronization targets}
\label{sec:v24-conclusion}
% =============================================================================

\begin{theorem}[What V24 proves]
\label{thm:v24-conclusion}
Relative to V1--V23, the following are proved.
\begin{enumerate}[label=\textup{(V24.\arabic*)},leftmargin=6.2em]
\item Every first record doubling obeys the assembled overlap--separation
      alternative \eqref{eq:v24-overlap-branch}--%
      \eqref{eq:v24-separation-branch} for each
      \(0\le\eta<3/4\).
\item On the separation branch, the sector-difference profile has equal
      positive and negative masses at least
      \((3/4-\eta)R^2\), each occupying canonical cutoff width.
\item Assembled overlap is bounded by common terminal chiral-tail stock plus
      critical dissipation.  Large overlap therefore forces one of the exact
      alternatives \eqref{eq:v24-common-terminal-branch}--%
      \eqref{eq:v24-critical-dissipation-branch}.
\item The sector-difference profile has the exact high-tail balance
      \eqref{eq:v24-sector-difference-balance}.  Its available absolute upper
      bound uses endpoint critical stock and critical dissipation and is
      compatible with the separation lower bound.
\item The first radial moment of sector difference is exactly the signed
      \(H^1\)-helicity balance
      \eqref{eq:v24-signed-barycentre-balance}.  Its direct absolute estimate
      requires endpoint enstrophy and palinstrophy.
\item Positive-quantile separation has the exact residual ledger
      \eqref{eq:v24-positive-quantile-ledger}; discarding its residual
      positive and negative first moments would assume a gross spectral-
      variation bound.
\item The disjoint minimum-area profiles saturate the separation lower bound
      and have signed first moment \(9R^4/(8E_*)\), the already permitted
      quartic record scale.
\end{enumerate}
No full Navier--Stokes stability or global-regularity theorem is obtained.
\end{theorem}

\begin{proof}
Items \textup{(V24.1)--(V24.2)} are
Theorem~\ref{thm:v24-overlap-separation}.  Item \textup{(V24.3)} is
Theorem~\ref{thm:v24-overlap-payment}.  Item \textup{(V24.4)} is
Theorem~\ref{thm:v24-sector-difference-balance}.  Item \textup{(V24.5)} is
Theorem~\ref{thm:v24-signed-barycentre}.  Item \textup{(V24.6)} is
Proposition~\ref{prop:v24-positive-quantile-ledger}.  Item
\textup{(V24.7)} is Proposition~\ref{prop:v24-disjoint-profile-test}.
\end{proof}

\begin{programme}[V25 required companion review and acquisition order]
\label{prog:v25}
\begin{enumerate}[label=\textup{(AC\arabic*)},leftmargin=4.8em]
\item Perform the required fifth-version review of the latest Standard Model
      and Yang--Mills notes in \texttt{latex\_notes}.  Freeze filenames and
      hashes before importing any mechanism.
\item Present the exact NS targets to the companion review: the common
      terminal stock \(\mathcal C_I(b)\), the critical dissipation branch,
      the zero-mass difference \(D_I\), and the Jordan residual in
      \eqref{eq:v24-positive-quantile-ledger}.
\item Import a same-terminal subtraction or determinant/locality mechanism
      only if its hypotheses survive the helical sector types and its upper
      ledger is paid by Leray energy rather than critical dissipation.
\item On the overlap/common-terminal branch, test whether disjoint front
      bands turn \(\mathcal C_I(b)\) into a one-use stock.  Preserve the
      nested-tail layer-cake multiplicity explicitly.
\item On the separation branch, do not replace the signed barycentre by
      positive quantiles until the two Jordan residuals have been assembled.
\item Compare any proposed first-moment payment with
      \eqref{eq:v24-model-quartic-moment}; a quartic upper bound is already
      known and does not close doubled records.
\item If the companion mechanisms require an unproved positive determinant,
      sector overlap, or critical-dissipation budget, record the type mismatch
      and continue upstream from the exact NS balances.
\end{enumerate}
\end{programme}

\begin{assessment}[Position after V24]
The two-helicity route is now reduced to an assembled, quantitative
dichotomy.  Overlap produces a genuinely simultaneous terminal stock, but it
can also be paid by unbudgeted critical dissipation.  Separation produces a
large zero-mean signed profile, but its barycentric moment is an
\(H^1\)-helicity balance whose absolute control is the classical
enstrophy/palinstrophy problem.  Passing from signed profiles to positive
quantiles introduces an additional Jordan residual which cannot be omitted.

V25 must therefore test the companion programmes against named NS objects,
not import their conclusions by analogy.  A successful mechanism must make
the common terminal stock one-use or cancel the Jordan residual using only a
globally paid ledger.  Until one of those acquisitions is proved, the full
stability/global-regularity objective remains open.
\end{assessment}

% =============================================================================
\section*{Version V24 history}
\addcontentsline{toc}{section}{Version V24 history}
% =============================================================================

Derived an assembled overlap--separation dichotomy for the two compensated
helical flux profiles.  Proved that large overlap forces common terminal
chiral-tail stock or critical dissipation, while small overlap forces
canonical-width positive and negative mass in the zero-mean sector-difference
profile.  Derived the exact difference high-tail balance and its absolute
upper estimate.  Identified the signed first radial moment with an exact
\(H^1\)-helicity balance, exposing endpoint enstrophy and palinstrophy as the
separation cost.  Added the complete Jordan residual for separated positive
quantiles and proved that the disjoint minimum-area profiles saturate the
separation mass while reaching quartic record scale.  The full stability/
global-regularity goal remains open.

% =============================================================================
% VERSION V25 APPEND
% =============================================================================

\clearpage
\section{V25: the helical-tail meet balance and the exact adaptive
imbalance obstruction}
\label{sec:v25-entry}
% =============================================================================

V24 left four named synchronization targets: the common terminal chiral-tail
stock, the critical-dissipation alternative, the zero-mass sector-difference
profile, and the Jordan residual which appears when signed profiles are
replaced by positive quantiles.  This is the scheduled fifth-version
checkpoint.  We first audit the current companion notes, with their exact
claim boundaries, and then return to the first of those NS objects.

The valid methodological transfer is narrow but useful.  The Yang--Mills
calculation says that a cancelling expression must be assembled before its
summands are estimated absolutely.  Applied here, that instruction says to
form the pointwise minimum of the two complete helical tails before splitting
the nonlinear flux into triads or signs.  The resulting minimum, called the
helical-tail meet below, has an exact balance law.  It yields an
inherited-stock/innovation alternative and an exact adaptive correction to the
equal sector work.  The correction can in turn be decomposed into the V24
time-integrated sector difference and a tail-order-crossing residual.

This calculation does not close the programme.  The crossing residual reduces
to endpoint tail order plus a viscosity-weighted critical action, not to the
Leray debit.  The gain is that the common-terminal branch now has a precise
dynamical source and the failure of a cofinal one-use argument is localized.

% =============================================================================
\section{Required fifth-version companion audit}
\label{sec:v25-companion-audit}
% =============================================================================

\begin{proposition}[Frozen V25 companion checkpoint]
\label{prop:v25-companion-checkpoint}
The two latest companion files inspected before the V25 analysis were
\begin{align*}
 &\texttt{Renewal\_SM\_Einstein\_Continuum\_Research\_Notes\_V49.tex},\\
 &\hspace{1.5em}\text{SHA--256 }
 \texttt{f93274386c93c686501eb29fb513d98b21eee9eca322949ba986b43f3a038bf0},\\
 &\texttt{Renewal\_Yang\_Mills\_Mass\_Gap\_Research\_Notes\_V29.tex},\\
 &\hspace{1.5em}\text{SHA--256 }
 \texttt{3d9d1d3167a274a395e61fba8a85ee06cac3febd56aaac3c08884ae660984cbb}.
\end{align*}
Their proved statements transfer to the present NS problem only as follows.
\begin{enumerate}[label=\textup{(\roman*)},leftmargin=3.7em]
\item SM V49 constructs an explicit positive, local, rotationally covariant
      two-speed bath whose reduced three-port Schur matrix has a uniform
      inverse floor.  Its own firewall states that the unreduced Einstein
      glancing estimate is not yet proved.
\item The SM argument uses a fixed finite-dimensional linear port, positive
      constant Gram matrices, and an exact Schur factorization.  The V24 stock
      \(\mathcal C_I(b)\) is instead a state-dependent nonlinear minimum of
      two infinite cutoff tails.  No fixed port map or determinant identity
      realizing that minimum is available.  Therefore the SM inverse floor is
      not an NS estimate.
\item YM V29 proves fixed-window cancellation only after the complete
      M\"obius/fusion sum is assembled.  It explicitly forbids taking
      multiplier-by-multiplier absolute values before that assembly.
\item The uniform all-order YM fusion-cumulant gate remains open.  In
      particular, neither its conditional compiler nor its fixed-window
      theorem supplies a uniform NS flux, dissipation, or source bound.
\end{enumerate}
Thus the only imported mechanism used in V25 is the proved assembly order:
form the complete helical-tail minimum first and derive its exact PDE balance
before applying absolute values.  No SM determinant or conditional YM
majorant is imported as an estimate.
\end{proposition}

\begin{proof}
The filenames and digests are the frozen filesystem checkpoint.  Items
\textup{(i)--(ii)} are the hypothesis and scope statements of the SM V49
explicit positive-port theorem and its port-stability firewall.  Items
\textup{(iii)--(iv)} are respectively the YM V29 no-termwise-absolute
assessment, fixed-window fusion proposition, and uniform-gate firewall.  The
NS type comparison follows directly from the definition
\eqref{eq:v24-common-terminal-tail-stock}: a pointwise minimum depends on the
current solution and on every cutoff, whereas a fixed positive port Gram is a
linear finite-dimensional datum.  Hence the asserted transfer list neither
strengthens nor suppresses a companion hypothesis.
\end{proof}

\begin{firewall}[Companion positivity is not an NS a priori estimate]
\label{fw:v25-companion-type}
Passivity of a local auxiliary bath and cancellation in a fixed Fourier fusion
window are structural facts in their own problems.  They do not bound
\(\nu\int\|\Lambda^{3/2}u\|_2^2\), do not control the sign changes of the
helical tail difference, and do not turn a time-dependent minimum into a
fixed determinant.  Any use of those conclusions without these missing type
identifications would assume the desired NS closure.
\end{firewall}

% =============================================================================
\section{The common stock is an exact spectral-tail meet}
\label{sec:v25-tail-meet}
% =============================================================================

For every smooth time define
\begin{align}
 H_\rho^s(t)&:=\|Q_\rho u^s(t)\|_2^2,
 \label{eq:v25-sector-tail}\\
 m_\rho(t)&:=\min\{H_\rho^+(t),H_\rho^-(t)\},
 \label{eq:v25-tail-meet-density}\\
 \mathcal C(t)&:=\int_0^\infty m_\rho(t)\,\dd\rho,
 \label{eq:v25-time-dependent-common-stock}\\
 \mathcal L(t)&:=\int_0^\infty
   |\mathcal Y_\rho(t)|\,\dd\rho,
 \label{eq:v25-tail-imbalance-variation}
\end{align}
where \(\mathcal Y_\rho\) is \eqref{eq:v24-signed-tail-stock}.  Thus
\(\mathcal C_I(b)=\mathcal C(b)\).

For \(z\ge0\), introduce the decreasing-tail quantile length
\begin{equation}
 q_s(z,t):=
 \bigl|\{\rho\ge0:H_\rho^s(t)>z\}\bigr|.
 \label{eq:v25-tail-quantile-length}
\end{equation}
It is zero once \(z\ge H_0^s(t)\).  This definition avoids any convention at
the jumps of the discrete Fourier tail.

\begin{theorem}[Exact Jordan and quantile formulae for common stock]
\label{thm:v25-common-stock-formulae}
At every smooth time,
\begin{align}
 \boxed{\quad
 \mathcal C(t)
 =\frac12\bigl(\mathcal R(t)^2-\mathcal L(t)\bigr).
 \quad}
 \label{eq:v25-common-stock-Jordan}\\
 \boxed{\quad
 \mathcal C(t)
 =\int_0^\infty
   \min\{q_+(z,t),q_-(z,t)\}\,\dd z.
 \quad}
 \label{eq:v25-common-stock-quantile}
\end{align}
Consequently,
\begin{equation}
 0\le\mathcal C(t)
 \le\min\{X_+(t),X_-(t)\},
 \qquad
 |X_+(t)-X_-(t)|\le\mathcal L(t)
 =\mathcal R(t)^2-2\mathcal C(t).
 \label{eq:v25-common-stock-helicity-cap}
\end{equation}
At the terminal point of a first record doubling,
\begin{equation}
 \mathcal C(b)=2R^2-\frac12\mathcal L(b).
 \label{eq:v25-terminal-common-Jordan}
\end{equation}
Thus the V24 common-terminal branch is exactly a deficit in the absolute
radial variation of the signed helical tail; it is not shellwise overlap.
\end{theorem}

\begin{proof}
For nonnegative numbers \(x,y\),
\[
 \min\{x,y\}=\frac12(x+y-|x-y|).
\]
Apply this pointwise with \(x=H_\rho^+\), \(y=H_\rho^-\), integrate in
\(\rho\), and use spectral layer cake:
\[
 \int_0^\infty(H_\rho^++H_\rho^-)\,\dd\rho
 =X_+(t)+X_-(t)=\mathcal R(t)^2.
\]
This proves \eqref{eq:v25-common-stock-Jordan}.

For the second identity, ordinary layer cake gives
\begin{align*}
 m_\rho(t)
 &=\int_0^\infty
   \mathbf1_{\{z<H_\rho^+(t)\}}
   \mathbf1_{\{z<H_\rho^-(t)\}}\,\dd z.
\end{align*}
For fixed \(z\), monotonicity of each Fourier tail in \(\rho\) says that
\(\{\rho:H_\rho^s>z\}\) is, up to endpoints, an initial interval of length
\(q_s(z,t)\).  The intersection of the two intervals has length
\(\min\{q_+,q_-\}\).  Tonelli proves
\eqref{eq:v25-common-stock-quantile}.

The first upper bound in \eqref{eq:v25-common-stock-helicity-cap} follows by
integrating \(m_\rho\le H_\rho^s\).  The remaining inequality is the triangle
inequality applied to \(\int\mathcal Y_\rho=X_+-X_-\), followed by
\eqref{eq:v25-common-stock-Jordan}.  Finally use
\(\mathcal R(b)^2=4R^2\) to obtain
\eqref{eq:v25-terminal-common-Jordan}.
\end{proof}

\begin{proposition}[A two-atom test: tail meet is not shell intersection]
\label{prop:v25-two-atom-test}
Suppose at one time the two sector energy measures are concentrated at
radii \(\lambda_+,\lambda_->0\), with total sector energies
\(e_+,e_->0\).  Then
\begin{equation}
 \boxed{\quad
 \mathcal C
 =\min\{e_+,e_-\}\min\{\lambda_+,\lambda_-\}.
 \quad}
 \label{eq:v25-two-atom-common-stock}
\end{equation}
In particular, \(\mathcal C>0\) even when
\(\lambda_+\ne\lambda_-\) and the two shell supports are disjoint.
\end{proposition}

\begin{proof}
Here
\(H_\rho^s=e_s\mathbf1_{[0,\lambda_s)}(\rho)\), up to the irrelevant
cutoff endpoint.  Their pointwise minimum equals
\(\min\{e_+,e_-\}\) on
\([0,\min\{\lambda_+,\lambda_-\})\) and vanishes afterwards.  Integration
proves \eqref{eq:v25-two-atom-common-stock}.
\end{proof}

\begin{assessment}[What the nested-tail multiplicity means]
\label{ass:v25-nested-tail-meaning}
The factor \(\min\{\lambda_+,\lambda_-\}\) in
\eqref{eq:v25-two-atom-common-stock} is the layer-cake multiplicity: one unit
of shell energy is counted at every lower cutoff.  This is precisely why
\(\mathcal C\) has the dimension and size of critical stock rather than
energy.  It also shows that \emph{common tail} must not be read as
\emph{common Fourier shell}.  The quantile formula pairs equal energy ranks
and charges each pair only up to the smaller radial quantile.
\end{assessment}
% =============================================================================
\section{Same-terminal band additivity and its temporal limitation}
\label{sec:v25-band-additivity}
% =============================================================================

For a measurable cutoff band \(B\subset[0,\infty)\), set
\begin{equation}
 \mathcal C_B(t):=\int_Bm_\rho(t)\,\dd\rho.
 \label{eq:v25-band-common-stock}
\end{equation}

\begin{theorem}[The common stock is one-use across bands at one time]
\label{thm:v25-same-time-one-use}
If \(B_1,\ldots,B_N\) are pairwise disjoint measurable cutoff bands, then
\begin{equation}
 \boxed{\quad
 \sum_{j=1}^N\mathcal C_{B_j}(t)
 =\mathcal C_{\cup_jB_j}(t)
 \le\mathcal C(t)
 \le\min\{X_+(t),X_-(t)\}.
 \quad}
 \label{eq:v25-same-time-band-payment}
\end{equation}
The statement remains true for a countable disjoint family by monotone
convergence.
\end{theorem}

\begin{proof}
The first equality is countable additivity of the nonnegative measure
\(m_\rho(t)\,\dd\rho\); the next inequality is monotonicity under restriction
of its domain.  The last inequality is
\eqref{eq:v25-common-stock-helicity-cap}.  Monotone convergence handles a
countable family.
\end{proof}

\begin{firewall}[Record terminal times are not a common terminal time]
\label{fw:v25-no-cofinal-time-payment}
Equation~\eqref{eq:v25-same-time-band-payment} pays disjoint bands only when
they are evaluated at the same time.  Cofinal record doublings end at
different times \(b_j\).  Neither sector energy is conserved: nonlinear
exchange can deplete and replenish both tails.  Therefore
\[
 \sum_j\mathcal C_{B_j}(b_j)
 \le \sup_t\mathcal C(t)
\]
does not follow from countable additivity.  Moving all claims to a later time
would require a propagation or ancestry estimate for the meet.  The next
section derives its exact evolution and shows the unbudgeted source which such
an estimate must control.
\end{firewall}

% =============================================================================
\section{An exact PDE balance for the smaller helical tail}
\label{sec:v25-meet-balance}
% =============================================================================

Write
\begin{align}
 V_\rho^s(t)&:=\|\Lambda Q_\rho u^s(t)\|_2^2,
 \label{eq:v25-instant-sector-tail-dissipation}\\
 d_\rho(t)&:=\widetilde\Phi_\rho^+(t)
             -\widetilde\Phi_\rho^-(t),
 \label{eq:v25-instant-sector-flux-difference}\\
 \sigma_\rho(t)&:=\operatorname{sgn}\mathcal Y_\rho(t),
 \qquad \operatorname{sgn}0:=0,
 \label{eq:v25-tail-order-sign}\\
 \chi_\rho^+(t)&:=\frac{1-\sigma_\rho(t)}2,
 &\chi_\rho^-(t)&:=\frac{1+\sigma_\rho(t)}2.
 \label{eq:v25-meet-selectors}
\end{align}
Thus away from equality, \(\chi^s=1\) precisely for the sector with the
smaller tail.  At equality both weights are \(1/2\).  Define the assembled
meet dissipation and flux by
\begin{align}
 V_\rho^\wedge(t)
 &:=\chi_\rho^+(t)V_\rho^+(t)
   +\chi_\rho^-(t)V_\rho^-(t),
 \label{eq:v25-meet-dissipation-density}\\
 \widetilde\Phi_\rho^\wedge(t)
 &:=\chi_\rho^+(t)\widetilde\Phi_\rho^+(t)
   +\chi_\rho^-(t)\widetilde\Phi_\rho^-(t).
 \label{eq:v25-meet-flux-density}
\end{align}

\begin{lemma}[Absolutely continuous minimum chain rule]
\label{lem:v25-min-chain-rule}
If \(f,g\) are absolutely continuous on a compact interval, put
\(h=f-g\), \(\sigma=\operatorname{sgn}h\), and
\(\chi_f=(1-\sigma)/2\), \(\chi_g=(1+\sigma)/2\).  Then
\begin{equation}
 \frac{\dd}{\dd t}\min\{f,g\}
 =\chi_f f'+\chi_g g'
 \quad\text{for almost every }t.
 \label{eq:v25-min-chain-rule}
\end{equation}
\end{lemma}

\begin{proof}
The Lipschitz chain rule gives
\((|h|)'=\operatorname{sgn}(h)h'\) almost everywhere.  On the part of the
zero set \(\{h=0\}\) where derivatives exist, \(h'=0\) almost everywhere;
this is the standard zero-set property of absolutely continuous functions.
Differentiate
\(\min\{f,g\}=(f+g-|f-g|)/2\) to obtain
\eqref{eq:v25-min-chain-rule}, including the symmetric selector on the zero
set.
\end{proof}

\begin{theorem}[Exact helical-tail meet balance]
\label{thm:v25-meet-balance}
On every compact smooth NS interval, for almost every \((\rho,t)\),
\begin{equation}
 \boxed{\quad
 \frac12\partial_t m_\rho(t)
 +\nu V_\rho^\wedge(t)
 =\widetilde\Phi_\rho^\wedge(t).
 \quad}
 \label{eq:v25-pointwise-meet-balance}
\end{equation}
Define
\begin{align}
 \mathcal P_\wedge(I)
 &:=\int_I\int_0^\infty
     V_\rho^\wedge(t)\,\dd\rho\,\dd t,
 \label{eq:v25-meet-critical-action}\\
 \mathcal B_\wedge(I)
 &:=\int_I\int_0^\infty
     \widetilde\Phi_\rho^\wedge(t)\,\dd\rho\,\dd t.
 \label{eq:v25-meet-nonlinear-work}
\end{align}
Then
\begin{equation}
 \boxed{\quad
 \mathcal B_\wedge(I)
 =\frac12\bigl(\mathcal C(b)-\mathcal C(a)\bigr)
  +\nu\mathcal P_\wedge(I).
 \quad}
 \label{eq:v25-integrated-meet-balance}
\end{equation}
Moreover,
\begin{equation}
 0\le\mathcal P_\wedge(I)
 \le\mathcal P_+(I)+\mathcal P_-(I).
 \label{eq:v25-meet-action-bound}
\end{equation}
\end{theorem}

\begin{proof}
For fixed \(\rho\), apply Lemma~\ref{lem:v25-min-chain-rule} to
\(H_\rho^+\) and \(H_\rho^-\).  Multiply the two established sector
balances \eqref{eq:v23-sector-high-balance} by their selectors and add.  This
gives \eqref{eq:v25-pointwise-meet-balance}.  Smoothness on the compact
interval supplies the spectral summability required for Fubini; alternatively
one may first truncate to finitely many Fourier modes and pass by dominated
convergence.  Integration in \(\rho\) and \(t\) gives
\eqref{eq:v25-integrated-meet-balance}.

The selectors are nonnegative and sum to one.  Hence
\(0\le V_\rho^\wedge\le V_\rho^++V_\rho^-\).  Integrate and use
\eqref{eq:v24-tail-dissipation-layer-cake} to prove
\eqref{eq:v25-meet-action-bound}.
\end{proof}

\begin{corollary}[Inherited common stock or bottleneck-sector innovation]
\label{cor:v25-inheritance-innovation}
Let \(0\le\theta<1\).  On every smooth interval \(I=[a,b]\), at least one
of
\begin{align}
 \mathcal C(a)&\ge\theta\mathcal C(b),
 \tag{INH}\label{eq:v25-inherited-meet}\\
 \mathcal B_\wedge(I)
 &>\frac{1-\theta}{2}\mathcal C(b)
   +\nu\mathcal P_\wedge(I)
 \tag{NEW}\label{eq:v25-new-meet}
\end{align}
holds.  In particular, if the V24 branch
\(\mathcal C(b)\ge\eta R^2\) holds on a first doubling, then at least one of
\begin{align}
 \mathcal C(a)&\ge\frac\eta2R^2,
 \label{eq:v25-half-inherited-common}\\
 \mathcal B_\wedge(I)&>
    \frac\eta4R^2+\nu\mathcal P_\wedge(I)
 \label{eq:v25-quarter-new-common}
\end{align}
holds.
\end{corollary}

\begin{proof}
If \eqref{eq:v25-inherited-meet} fails, then
\(\mathcal C(b)-\mathcal C(a)>(1-\theta)\mathcal C(b)\).  Substitute this
strict inequality into \eqref{eq:v25-integrated-meet-balance}.  This proves
\eqref{eq:v25-new-meet}.  Set \(\theta=1/2\) and use
\(\mathcal C(b)\ge\eta R^2\) for the last alternative.
\end{proof}

\begin{assessment}[Exact progress on the V24 overlap branch]
\label{ass:v25-overlap-progress}
The old terminal stock is no longer merely an endpoint quantity.  If it is
not already present at half strength at the preceding record time, the
nonlinearity must perform positive work on whichever helical tail is smaller
at each cutoff.  This is stronger typing than the sum of the two sector works:
the selector is determined by the assembled tails.  However,
\(\mathcal B_\wedge\) is not a conserved injection and
\(\mathcal P_\wedge\) is a critical, not Leray, action.  Repeating
\eqref{eq:v25-quarter-new-common} over record intervals is therefore not yet a
global payment.
\end{assessment}
% =============================================================================
\section{Adaptive imbalance transport and the tail-order crossing residual}
\label{sec:v25-adaptive-imbalance}
% =============================================================================

The meet selector admits an exact comparison with the equal sector work.
Define
\begin{equation}
 \mathcal A_I
 :=\int_I\int_0^\infty
   \sigma_\rho(t)d_\rho(t)\,\dd\rho\,\dd t.
 \label{eq:v25-adaptive-imbalance-work}
\end{equation}

\begin{theorem}[Exact adaptive correction to equal helical work]
\label{thm:v25-adaptive-correction}
For every smooth interval,
\begin{equation}
 \boxed{\quad
 \mathcal B_\wedge(I)
 =\int_I\mathcal G_+(t)\,\dd t
  -\frac12\mathcal A_I
 =\int_I\mathcal G_-(t)\,\dd t
  -\frac12\mathcal A_I.
 \quad}
 \label{eq:v25-meet-equal-work-correction}
\end{equation}
On a first record doubling,
\begin{align}
 \mathcal B_\wedge(I)
 &=\frac34R^2
   +\frac\nu2\bigl(\mathcal P_+(I)+\mathcal P_-(I)\bigr)
   -\frac12\mathcal A_I,
 \label{eq:v25-doubling-meet-work}\\
 \mathcal A_I
 &=\frac32R^2+\mathcal C(a)-\mathcal C(b)
   +\nu\bigl(
      \mathcal P_+(I)+\mathcal P_-(I)
      -2\mathcal P_\wedge(I)
     \bigr).
 \label{eq:v25-adaptive-ledger}
\end{align}
\end{theorem}

\begin{proof}
The definitions of the selectors give pointwise
\begin{align*}
 \widetilde\Phi_\rho^\wedge
 &=\frac12\bigl(
    \widetilde\Phi_\rho^++\widetilde\Phi_\rho^-
   \bigr)
   -\frac12\sigma_\rho
    \bigl(
     \widetilde\Phi_\rho^+-\widetilde\Phi_\rho^-
    \bigr).
\end{align*}
Integrate in \(\rho\).  The sector layer-cake identity
\eqref{eq:v23-sector-layer-cake} and equal-work identity
\eqref{eq:v23-equal-instant-sector-work} turn the first term into either
\(\mathcal G_+\) or \(\mathcal G_-\).  Time integration proves
\eqref{eq:v25-meet-equal-work-correction}.  Insert
\eqref{eq:v23-sector-work-lower}, which is an equality with its displayed
dissipation term, to obtain \eqref{eq:v25-doubling-meet-work}.  Finally equate
that formula with \eqref{eq:v25-integrated-meet-balance} and solve for
\(\mathcal A_I\), proving \eqref{eq:v25-adaptive-ledger}.
\end{proof}

The weight \(\sigma_\rho(t)\) in \(\mathcal A_I\) changes when the two
helical tails exchange order.  To compare it with the already assembled V24
profile \(D_I\), freeze the sign at an endpoint.  Put
\begin{align}
 \sigma_\rho^a&:=\sigma_\rho(a),
 &\sigma_\rho^b&:=\sigma_\rho(b),
 \label{eq:v25-endpoint-tail-signs}\\
 \mathcal S_I^b
 &:=\int_I\int_0^\infty
  (\sigma_\rho(t)-\sigma_\rho^b)d_\rho(t)
  \,\dd\rho\,\dd t,
 \label{eq:v25-terminal-crossing-residual}\\
 \mathcal S_I^a
 &:=\int_I\int_0^\infty
  (\sigma_\rho(t)-\sigma_\rho^a)d_\rho(t)
  \,\dd\rho\,\dd t.
 \label{eq:v25-initial-crossing-residual}
\end{align}

\begin{theorem}[Exact endpoint freezing and crossing formulae]
\label{thm:v25-crossing-formulae}
The adaptive correction decomposes as
\begin{align}
 \mathcal A_I
 &=\int_0^\infty\sigma_\rho^bD_I(\rho)\,\dd\rho
   +\mathcal S_I^b,
 \label{eq:v25-terminal-freezing}\\
 &=\int_0^\infty\sigma_\rho^aD_I(\rho)\,\dd\rho
   +\mathcal S_I^a.
 \label{eq:v25-initial-freezing}
\end{align}
Let
\(\Delta V_\rho:=V_\rho^+-V_\rho^-\).  Then the crossing residuals have
the exact endpoint forms
\begin{align}
 \boxed{\quad
 \mathcal S_I^b
 &=\frac12\int_0^\infty
   \bigl(\sigma_\rho^b\mathcal Y_\rho(a)
         -|\mathcal Y_\rho(a)|\bigr)\,\dd\rho
 \notag\\
 &\quad+\nu\int_I\int_0^\infty
  (\sigma_\rho(t)-\sigma_\rho^b)
  \Delta V_\rho(t)\,\dd\rho\,\dd t,
 \quad}
 \label{eq:v25-terminal-crossing-exact}\\
 \boxed{\quad
 \mathcal S_I^a
 &=\frac12\int_0^\infty
   \bigl(|\mathcal Y_\rho(b)|
         -\sigma_\rho^a\mathcal Y_\rho(b)\bigr)\,\dd\rho
 \notag\\
 &\quad+\nu\int_I\int_0^\infty
  (\sigma_\rho(t)-\sigma_\rho^a)
  \Delta V_\rho(t)\,\dd\rho\,\dd t.
 \quad}
 \label{eq:v25-initial-crossing-exact}
\end{align}
In particular,
\begin{align}
 \mathcal S_I^b&\le
  2\nu\bigl(\mathcal P_+(I)+\mathcal P_-(I)\bigr),
 \label{eq:v25-terminal-crossing-upper}\\
 \mathcal S_I^a&\ge
 -2\nu\bigl(\mathcal P_+(I)+\mathcal P_-(I)\bigr),
 \label{eq:v25-initial-crossing-lower}
\end{align}
and hence
\begin{align}
 \int\sigma_\rho^aD_I(\rho)\,\dd\rho
 -2\nu(\mathcal P_++\mathcal P_-)
 \le\mathcal A_I
 \le
 \int\sigma_\rho^bD_I(\rho)\,\dd\rho
 +2\nu(\mathcal P_++\mathcal P_-).
 \label{eq:v25-endpoint-sign-sandwich}
\end{align}
Here and below the interval argument \(I\) on \(\mathcal P_s\) is
suppressed only inside long displays.
\end{theorem}

\begin{proof}
By definition and Fubini,
\[
 \int_I d_\rho(t)\,\dd t
 =F_+(\rho)-F_-(\rho)=D_I(\rho).
\]
Add and subtract either endpoint sign in
\eqref{eq:v25-adaptive-imbalance-work} to prove
\eqref{eq:v25-terminal-freezing}--\eqref{eq:v25-initial-freezing}.

Subtract the two instantaneous sector-tail balances
\eqref{eq:v23-sector-high-balance}:
\begin{equation}
 d_\rho(t)
 =\frac12\partial_t\mathcal Y_\rho(t)
  +\nu\Delta V_\rho(t).
 \label{eq:v25-instant-tail-difference}
\end{equation}
For an absolutely continuous scalar function \(Y\),
\begin{equation}
 \int_a^b\operatorname{sgn}(Y)Y'\,\dd t
 =|Y(b)|-|Y(a)|.
 \label{eq:v25-absolute-chain-rule}
\end{equation}
Use \eqref{eq:v25-instant-tail-difference} in
\(\mathcal S_I^b\).  The derivative contribution is
\begin{align*}
 \frac12\int_a^b(\sigma(t)-\sigma^b)Y'(t)\,\dd t
 &=\frac12\bigl(
   |Y(b)|-|Y(a)|-\sigma^b(Y(b)-Y(a))
  \bigr)\\
 &=\frac12\bigl(\sigma^bY(a)-|Y(a)|\bigr),
\end{align*}
because \(\sigma^bY(b)=|Y(b)|\), including when \(Y(b)=0\).
Integrating in \(\rho\) proves
\eqref{eq:v25-terminal-crossing-exact}.  Freezing at \(a\) gives instead
\begin{align*}
 \frac12\int_a^b(\sigma(t)-\sigma^a)Y'(t)\,\dd t
 =\frac12\bigl(|Y(b)|-\sigma^aY(b)\bigr),
\end{align*}
which proves \eqref{eq:v25-initial-crossing-exact}.

The endpoint integrand in
\eqref{eq:v25-terminal-crossing-exact} is nonpositive, while that in
\eqref{eq:v25-initial-crossing-exact} is nonnegative.  Also
\begin{align*}
 |\sigma_\rho(t)-\sigma_\rho^{a,b}|&\le2,
 &|\Delta V_\rho(t)|&\le V_\rho^+(t)+V_\rho^-(t).
\end{align*}
The sector dissipation layer cake
\eqref{eq:v24-tail-dissipation-layer-cake} now proves
\eqref{eq:v25-terminal-crossing-upper} and
\eqref{eq:v25-initial-crossing-lower}.  Combine these with the two freezing
identities to obtain \eqref{eq:v25-endpoint-sign-sandwich}.
\end{proof}

\begin{corollary}[The assembled correction is not the absolute sector
difference]
\label{cor:v25-adaptive-not-absolute}
For either endpoint \(\tau\in\{a,b\}\),
\begin{equation}
 \left|\int_0^\infty
  \sigma_\rho^\tau D_I(\rho)\,\dd\rho\right|
 \le\int_0^\infty|D_I(\rho)|\,\dd\rho
 \le\int_I\int_0^\infty|d_\rho(t)|\,\dd\rho\,\dd t.
 \label{eq:v25-absolute-hierarchy}
\end{equation}
The first upper quantity is the signed endpoint pairing actually occurring
in \(\mathcal A_I\); the middle quantity is the V24 assembled separation
mass; the last is the forbidden termwise-in-time absolute variation.
\end{corollary}

\begin{proof}
The first inequality uses \(|\sigma^\tau|\le1\).  The second follows from
\(D_I(\rho)=\int_I d_\rho(t)\,\dd t\), the triangle inequality in time, and
Tonelli.
\end{proof}

\begin{assessment}[The upstream obstruction after exact assembly]
\label{ass:v25-upstream-obstruction}
The YM assembly lesson removes a false obstacle: changes in which helical
tail is smaller do not require a gross nonlinear sign-count.  The scalar
absolute-value chain rule collapses their derivative contribution to the
endpoint mismatch in
\eqref{eq:v25-terminal-crossing-exact}--%
\eqref{eq:v25-initial-crossing-exact}.  What remains is sharper and more
classical: viscosity couples the adaptive sign to
\(V_\rho^+-V_\rho^-\), whose radial integral is critical dissipation.

Nor does the large separation mass in V24 force a large endpoint pairing;
\(D_I\) may oscillate relative to either endpoint tail sign.  Replacing that
pairing by \(\int|D_I|\) discards the assembled sign and returns the endpoint
critical-stock/critical-dissipation upper bound
\eqref{eq:v24-difference-absolute-upper}.  Thus the exact remaining
acquisition is a signed source or ancestry estimate for
\(\mathcal A_I\), or an independently globally paid bound for its
viscosity-weighted tail-order term.  Neither is supplied by Leray energy.
\end{assessment}

\begin{firewall}[The meet balance is not a monotonicity formula]
\label{fw:v25-meet-not-monotone}
The dissipation in \eqref{eq:v25-pointwise-meet-balance} is nonnegative, but
the selected nonlinear flux \(\widetilde\Phi_\rho^\wedge\) has no fixed sign.
Consequently \(\mathcal C(t)\) need not be monotone by the identities proved
here.  In the absence of nonlinearity, the formula does make
\(\mathcal C\) nonincreasing.  For Navier--Stokes, however, exactly the
bottleneck-sector work in \eqref{eq:v25-meet-nonlinear-work} can recreate the
common stock.  Treating old meet mass as surviving to every later record time
would omit this source.
\end{firewall}
% =============================================================================
\section{Integrated V25 conclusion and next acquisition order}
\label{sec:v25-conclusion}
% =============================================================================

\begin{theorem}[What V25 proves]
\label{thm:v25-conclusion}
Relative to V1--V24, the following are proved.
\begin{enumerate}[label=\textup{(V25.\arabic*)},leftmargin=6.2em]
\item The scheduled SM V49/YM V29 checkpoint has been frozen by filename and
      digest.  Its only applicable transfer is to assemble the complete
      cancelling NS object before taking absolute values; the companion
      determinant and conditional all-order majorant do not transfer.
\item The common terminal chiral-tail stock has the exact Jordan identity
      \eqref{eq:v25-common-stock-Jordan} and quantile representation
      \eqref{eq:v25-common-stock-quantile}.  It is a stochastic tail meet,
      not a shell intersection.
\item At a fixed terminal time, disjoint cutoff bands charge the meet stock
      only once, as in \eqref{eq:v25-same-time-band-payment}.  Record
      terminal times differ, so this does not provide a cofinal payment.
\item The smaller helical tail obeys the exact local and integrated PDE
      balances \eqref{eq:v25-pointwise-meet-balance} and
      \eqref{eq:v25-integrated-meet-balance}.
\item Large terminal common stock is either inherited at the preceding
      record time or created by positive bottleneck-sector nonlinear work,
      with the quantitative alternative
      \eqref{eq:v25-half-inherited-common}--%
      \eqref{eq:v25-quarter-new-common}.
\item Bottleneck-sector work equals either common sector critical work minus
      one half of the adaptive signed imbalance transport.  On a doubling
      this gives the exact ledgers
      \eqref{eq:v25-doubling-meet-work}--%
      \eqref{eq:v25-adaptive-ledger}.
\item Freezing the adaptive sign at either endpoint decomposes it into the
      V24 sector-difference profile and an exact tail-order-crossing residual.
      The nonlinear derivative part of that residual collapses to an endpoint
      mismatch; the remaining error is viscosity-weighted critical action.
\item The exact hierarchy \eqref{eq:v25-absolute-hierarchy} distinguishes the
      endpoint signed pairing, the V24 assembled absolute separation, and the
      still larger termwise temporal variation.  None is silently substituted
      for another.
\end{enumerate}
These statements do not prove full Navier--Stokes stability or global
regularity.
\end{theorem}

\begin{proof}
Item \textup{(V25.1)} is Proposition~\ref{prop:v25-companion-checkpoint} and
Firewall~\ref{fw:v25-companion-type}.  Item \textup{(V25.2)} is
Theorem~\ref{thm:v25-common-stock-formulae} and
Proposition~\ref{prop:v25-two-atom-test}.  Item \textup{(V25.3)} is
Theorem~\ref{thm:v25-same-time-one-use} and its firewall.  Items
\textup{(V25.4)--(V25.5)} are Theorem~\ref{thm:v25-meet-balance} and
Corollary~\ref{cor:v25-inheritance-innovation}.  Item \textup{(V25.6)} is
Theorem~\ref{thm:v25-adaptive-correction}.  Item \textup{(V25.7)} is
Theorem~\ref{thm:v25-crossing-formulae}.  Item \textup{(V25.8)} is
Corollary~\ref{cor:v25-adaptive-not-absolute}.
\end{proof}

\begin{programme}[V26 acquisition order]
\label{prog:v26}
\begin{enumerate}[label=\textup{(AC\arabic*)},leftmargin=4.8em]
\item Keep the V25 meet balance assembled.  Do not expand
      \(\widetilde\Phi^\wedge\) triad by triad before the adaptive sign has
      been retained.
\item Move upstream to the signed source
      \(\mathcal A_I\) in \eqref{eq:v25-adaptive-imbalance-work}.  Test its
      complete helical-triad symmetrization against the endpoint pairings in
      \eqref{eq:v25-endpoint-sign-sandwich}.
\item Determine whether monotonicity of the tail-order sign in \(\rho\), or a
      bounded number of radial order crossings forced by the signed spectral
      measure, yields cancellation not visible in
      \(\int|D_I|\).
\item Preserve the viscosity term in
      \eqref{eq:v25-terminal-crossing-exact}.  A proposed estimate must reduce
      it to Leray dissipation or produce a new one-use stock; renaming it as
      harmless critical action is not closure.
\item Test all proposed endpoint-sign inequalities on the V23 disjoint
      profiles and on the two-atom tail-meet example
      \eqref{eq:v25-two-atom-common-stock}.
\item In parallel, retain the V24 separation branch.  If no signed estimate
      for \(\mathcal A_I\) is available, return to the complete helical triad
      source before applying positive parts, rather than bounding the Jordan
      residual termwise.
\item A stability conclusion may be recorded only after every cofinal record
      interval is charged to a globally finite Leray-level ledger or is
      otherwise excluded by a proved continuation criterion.
\end{enumerate}
\end{programme}

\begin{assessment}[Position after V25]
The overlap/common-terminal branch now has an exact dynamics.  Same-time
cutoff bands are genuinely one-use, but the stock can be recreated between
different record terminals.  Recreation is measured exactly by the work on
the smaller helical tail.  Equal helical critical work reduces that source to
one adaptive signed imbalance term, and endpoint freezing reduces the
tail-order-crossing part further to endpoint mismatch plus critical viscous
action.

This is a sharper obstruction than V24's undifferentiated common stock, but it
is still the critical boundary: neither the endpoint pairing nor the
viscosity-weighted tail-order difference is controlled by the Leray budget.
The next iteration must attack that signed source upstream, while preserving
the assembled cancellation.  The full stability/global-regularity objective
remains open.
\end{assessment}

% =============================================================================
\section*{Version V25 append-only record}
\addcontentsline{toc}{section}{Version V25 append-only record}
% =============================================================================

The complete V24 mathematical source was retained before this continuation.
Its SHA--256 digest was
\begin{center}\scriptsize\ttfamily
7ddde7639706331776a94f503aafb0ccd44a9c8b986845488deb61405cb7c4df.
\end{center}
V25 performed the required SM/YM checkpoint, derived exact Jordan and
quantile formulae for common chiral-tail stock, proved same-terminal band
additivity, and obtained the exact balance of the smaller helical tail.  It
then expressed bottleneck-sector work as equal helical critical work minus an
adaptive imbalance correction, froze that correction at both endpoints, and
collapsed the nonlinear tail-order-crossing derivative to endpoint mismatch.
The residual viscosity term is critical rather than Leray-level, so the full
stability/global-regularity goal remains open.
% =============================================================================
% VERSION V26 APPEND
% =============================================================================

\clearpage
\section{V26: signed-radius divided differences and a robust productive
triad class}
\label{sec:v26-entry}
% =============================================================================

V25 reduced the common-tail branch to one instantaneous signed source,
\(\int\sigma_\rho d_\rho\,\dd\rho\), where the sign is chosen only after the
two complete helical tails have been assembled.  The next acquisition is to
identify the exact complete-triad algebra of that source without replacing
the adaptive sign by an absolute value.

V26 obtains such an algebra.  The primitive of the tail-order sign defines an
odd, one-Lipschitz function on the signed helical radius.  On each complete
triad the adaptive source is its second divided difference, multiplied by the
energy--helicity Vandermonde.  An equivalent cutoff-kernel table is explicit
for every minority-helicity location.  It proves a new robust statement:
every productive mixed triad whose minority helicity lies on a high leg and
whose low radius is less than one half of the middle radius gives
nonnegative work to the smaller assembled helical tail, irrespective of the
shape of the global adaptive selector.  With a stricter nonlocality ratio the
contribution has an explicit positive fraction.

The result is not summed as a gross positive-triad estimate.  Its exact
complement is retained.  Homochiral crossing triads, low-minority triads,
weakly nonlocal high-minority triads, and nonproductive triads remain in that
complement.  We also prove that the number of radial tail-order crossings is
not bounded by either kinetic energy or critical stock, so a finite-crossing
shortcut cannot close the complement.

% =============================================================================
\section{The adaptive source as an order-one commutator}
\label{sec:v26-adaptive-commutator}
% =============================================================================

At a fixed smooth time put
\begin{align}
 \mathfrak a(t)&:=\int_0^\infty
   \sigma_\rho(t)d_\rho(t)\,\dd\rho,
 \label{eq:v26-instant-adaptive-source}\\
 \Sigma_t(r)&:=\int_0^r\sigma_\rho(t)\,\dd\rho,
 \qquad r\ge0.
 \label{eq:v26-sign-primitive}
\end{align}
Then \(\mathcal A_I=\int_I\mathfrak a(t)\,\dd t\).  On mean-zero
divergence-free fields define the helical-sign involution and its adaptive
multiplier by
\begin{equation}
 \mathsf K_{\rm hel}:=\Lambda^{-1}\operatorname{curl}
 =\mathsf P^+-\mathsf P^- ,
 \qquad
 \mathsf M_t:=\mathsf K_{\rm hel}\Sigma_t(\Lambda).
 \label{eq:v26-adaptive-multiplier}
\end{equation}
The symbol of \(\mathsf M_t\) on helicity \(s\) and radius \(r\) is
\(s\Sigma_t(r)\).

\begin{theorem}[Exact multiplier and commutator form of the adaptive source]
\label{thm:v26-adaptive-commutator}
At every smooth time,
\begin{align}
 \boxed{\quad
 \mathfrak a(t)
 =-\langle\BNS(u),\mathsf M_tu\rangle.
 \quad}
 \label{eq:v26-source-multiplier-form}\\
 \boxed{\quad
 \mathfrak a(t)
 =-\frac12\left\langle
   [\mathsf M_t,u\cdot\nabla]u,u
  \right\rangle.
 \quad}
 \label{eq:v26-source-commutator-form}
\end{align}
Moreover,
\begin{equation}
 |\Sigma_t(r)|\le r,
 \qquad
 |\Sigma_t(r)-\Sigma_t(q)|\le|r-q|.
 \label{eq:v26-primitive-Lipschitz}
\end{equation}
Thus \(\mathsf M_t\) has an order-one symbol, but its radial derivative is
the potentially discontinuous tail-order sign.
\end{theorem}

\begin{proof}
By \eqref{eq:v25-instant-sector-flux-difference} and the compensated sector
profile \eqref{eq:v23-compensated-sector-flux},
\[
 d_\rho(t)
 =-\sum_{k\ne0}\sum_{s\in\{+,-\}}
   s\,\mathbf1_{\{|k|>\rho\}}T_{k,s}(t).
\]
Fubini therefore gives
\begin{align*}
 \mathfrak a(t)
 &=-\sum_{k,s}sT_{k,s}(t)
   \int_0^{|k|}\sigma_\rho(t)\,\dd\rho\\
 &=-\sum_{k,s}s\Sigma_t(|k|)T_{k,s}(t)
 =-\langle\BNS(u),\mathsf M_tu\rangle,
\end{align*}
which proves \eqref{eq:v26-source-multiplier-form}.  The multiplier is real
self-adjoint and commutes with the Leray projector.  If
\(I:=\langle u\cdot\nabla u,\mathsf M_tu\rangle\), incompressibility
and integration by parts give
\[
 \langle\mathsf M_t(u\cdot\nabla u),u\rangle=I,
 \qquad
 \langle u\cdot\nabla(\mathsf M_tu),u\rangle=-I.
\]
The inner product of their difference with \(u\) is \(2I\), proving
\eqref{eq:v26-source-commutator-form}.  Finally
\(|\sigma_\rho|\le1\) and the definition \eqref{eq:v26-sign-primitive}
prove both inequalities in \eqref{eq:v26-primitive-Lipschitz}.
\end{proof}

\begin{firewall}[Order one is not a Leray commutator estimate]
\label{fw:v26-order-one-not-Leray}
The pointwise symbol bound \(|\Sigma_t(r)|\le r\) places
\(\mathsf M_t\) at the critical derivative level.  The sign
\(\sigma_\rho(t)\) can jump in \(\rho\), so no uniform smooth-symbol
commutator theorem follows from \eqref{eq:v26-primitive-Lipschitz}.  A direct
Sobolev estimate of \eqref{eq:v26-source-commutator-form} returns critical or
strong norms.  The gain must instead come from the complete-triad second
difference derived next.
\end{firewall}

% =============================================================================
\section{Complete-triad determinant on signed radii}
\label{sec:v26-triad-determinant}
% =============================================================================

Consider one reality-paired helical triad as in V22, with radii
\(0<\lambda_1<\lambda_2<\lambda_3\), helicities \(s_i\), signed radii
\(\mu_i=s_i\lambda_i\), work vector \(T_i^{\triangle,s}\), and normal-form
scalar \(\Xi\).  Define its contribution to \(\mathfrak a(t)\) by
\begin{equation}
 \mathfrak a^{\triangle}(t)
 :=-\sum_{i=1}^3s_i\Sigma_t(\lambda_i)
      T_i^{\triangle,s}(t).
 \label{eq:v26-triad-adaptive-source}
\end{equation}
Also extend the sign primitive oddly to the signed-radius line:
\begin{equation}
 g_t(x):=
 \begin{cases}
  \operatorname{sgn}(x)\Sigma_t(|x|),&x\ne0,\\
  0,&x=0.
 \end{cases}
 \label{eq:v26-odd-sign-primitive}
\end{equation}

\begin{theorem}[Adaptive triad determinant and divided difference]
\label{thm:v26-triad-determinant}
Every complete helical triad obeys
\begin{equation}
 \boxed{\quad
 \mathfrak a^{\triangle}(t)
 =\Xi\det
 \begin{pmatrix}
  1&1&1\\
  \mu_1&\mu_2&\mu_3\\
  g_t(\mu_1)&g_t(\mu_2)&g_t(\mu_3)
 \end{pmatrix}.
 \quad}
 \label{eq:v26-adaptive-triad-determinant}
\end{equation}
For distinct signed radii this is equivalently
\begin{equation}
 \mathfrak a^{\triangle}(t)
 =\Xi\prod_{1\le i<j\le3}(\mu_j-\mu_i)
   \,g_t[\mu_1,\mu_2,\mu_3],
 \label{eq:v26-adaptive-divided-difference}
\end{equation}
where the last factor is the symmetric second divided difference.  The
function \(g_t\) is odd, absolutely continuous, and one-Lipschitz, with
\begin{equation}
 g_t'(x)=\sigma_{|x|}(t)
 \quad\text{for almost every }x\ne0.
 \label{eq:v26-odd-primitive-derivative}
\end{equation}
In particular, if the tail-order sign is constant on
\([0,\lambda_3]\), then every helical triad contribution
\(\mathfrak a^{\triangle}(t)\) vanishes.
\end{theorem}

\begin{proof}
Put \(w_i=s_i\Sigma_t(\lambda_i)=g_t(\mu_i)\).  The helical transfer normal
form \eqref{eq:v22-helical-transfer-vector} gives
\[
 \mathfrak a^{\triangle}
 =-\Xi\bigl[
  w_1(\mu_2-\mu_3)+w_2(\mu_3-\mu_1)
       +w_3(\mu_1-\mu_2)
 \bigr].
\]
Expansion of the determinant in
\eqref{eq:v26-adaptive-triad-determinant} along its last row gives exactly
the negative bracket.  This proves the determinant identity.  The standard
Newton interpolation determinant is the Vandermonde product times the second
divided difference, proving \eqref{eq:v26-adaptive-divided-difference}.

For \(x>0\), \(g_t'(x)=\Sigma_t'(x)=\sigma_x(t)\) almost everywhere.  For
\(x<0\), differentiation of \(-\Sigma_t(-x)\) gives the same value
\(\sigma_{|x|}(t)\).  Hence \(g_t\) is one-Lipschitz and
\eqref{eq:v26-odd-primitive-derivative} holds.  If \(\sigma_\rho=c\) on the
stated interval, then \(g_t(x)=cx\) at all three signed radii.  The last row
of the determinant is then \(c\) times its second row, so the determinant
vanishes.
\end{proof}

\begin{remark}[What survives energy and helicity cancellation]
\label{rem:v26-second-difference-meaning}
The first two determinant rows are precisely the energy and helicity weights.
Thus constant and signed-linear parts of the adaptive multiplier are removed
before estimation.  Only the non-affinity of the odd primitive across the
three signed radii survives.  This is the complete-triad implementation of
the assembly-first instruction imported in V25.
\end{remark}
% =============================================================================
\section{Exact cutoff kernels for every helical triad type}
\label{sec:v26-kernel-table}
% =============================================================================

Reverse all helicities and absorb the resulting common orientation into
\(\Xi\), when necessary.  For a mixed triad let
\(m\in\{1,2,3\}\) denote its unique negative-helicity leg.  Let
\(m={\rm h}\) denote the all-positive homochiral representative.  Write
\(a=\lambda_1\), \(b=\lambda_2\), and \(c=\lambda_3\).  Fourier closure
implies the triangle inequality
\begin{equation}
 c\le a+b.
 \label{eq:v26-triad-triangle}
\end{equation}

\begin{theorem}[Exact adaptive cutoff-kernel table]
\label{thm:v26-adaptive-kernel-table}
For every one of the four representatives,
\begin{equation}
 \boxed{\quad
 \mathfrak a_m^{\triangle}(t)
 =\Xi\int_0^cK_m(\rho)\sigma_\rho(t)\,\dd\rho.
 \quad}
 \label{eq:v26-kernel-representation}
\end{equation}
Up to cutoff endpoints of measure zero, the kernels and their exact
\(L^1\) masses are
\begin{equation}
\begin{array}{c|c|c|c|c}
 m
 &0<\rho<a
 &a<\rho<b
 &b<\rho<c
 &\displaystyle\int_0^c|K_m|\,\dd\rho
 \\[0.2em]\hline
 {\rm h}
 &0
 &-(c-b)
 &b-a
 &2(b-a)(c-b)
 \\[0.35em]
 1
 &-2(c-b)
 &-(c-b)
 &a+b
 &2(a+b)(c-b)
 \\[0.35em]
 2
 &2(c-a)
 &c-b-2a
 &-(a+b)
 &4a(c-a)
 \\[0.35em]
 3
 &-2(b-a)
 &c-b+2a
 &-(b-a)
 &2(b-a)(c-b+2a).
\end{array}
\label{eq:v26-adaptive-kernel-table}
\end{equation}
Every row has zero signed mass:
\begin{equation}
 \int_0^cK_m(\rho)\,\dd\rho=0.
 \label{eq:v26-kernel-zero-mass}
\end{equation}
Consequently,
\begin{align}
 |\mathfrak a_{\rm h}^{\triangle}|
 &\le2(b-a)(c-b)|\Xi|,
 \label{eq:v26-homochiral-adaptive-bound}\\
 |\mathfrak a_1^{\triangle}|
 &\le2(a+b)(c-b)|\Xi|,
 \label{eq:v26-m1-adaptive-bound}\\
 |\mathfrak a_2^{\triangle}|
 &\le4a(c-a)|\Xi|,
 \label{eq:v26-m2-adaptive-bound}\\
 |\mathfrak a_3^{\triangle}|
 &\le2(b-a)(c-b+2a)|\Xi|.
 \label{eq:v26-m3-adaptive-bound}
\end{align}
These constants are sharp when the optimization is over all measurable
weights \(|\sigma_\rho|\le1\).
\end{theorem}

\begin{proof}
Insert the three explicit transfer vectors from the proof of
Theorem~\ref{thm:v22-minority-helicity-table} into
\eqref{eq:v26-triad-adaptive-source}.  With
\(\Sigma_j:=\Sigma_t(\lambda_j)\), this gives
\begin{align}
 \frac{\mathfrak a_{\rm h}^{\triangle}}{\Xi}
 &=(c-b)\Sigma_1-(c-a)\Sigma_2+(b-a)\Sigma_3,
 \label{eq:v26-homochiral-Sigma-form}\\
 \frac{\mathfrak a_1^{\triangle}}{\Xi}
 &=(b-c)\Sigma_1-(c+a)\Sigma_2+(a+b)\Sigma_3,
 \label{eq:v26-m1-Sigma-form}\\
 \frac{\mathfrak a_2^{\triangle}}{\Xi}
 &=(b+c)\Sigma_1+(c-a)\Sigma_2-(a+b)\Sigma_3,
 \label{eq:v26-m2-Sigma-form}\\
 \frac{\mathfrak a_3^{\triangle}}{\Xi}
 &=-(b+c)\Sigma_1+(c+a)\Sigma_2+(a-b)\Sigma_3.
 \label{eq:v26-m3-Sigma-form}
\end{align}
For example, the coefficient of \(\sigma_\rho\) below \(a\) is the sum of
all three coefficients in the corresponding line; between \(a\) and \(b\)
it is the sum of the last two; between \(b\) and \(c\) it is the last one.
This rule applied to the four displayed lines gives every entry of
\eqref{eq:v26-adaptive-kernel-table} and proves
\eqref{eq:v26-kernel-representation}.

For \(m=2\), the middle entry is strictly negative because
\(c-b-2a\le-a<0\) by \eqref{eq:v26-triad-triangle}.  All other displayed
signs are immediate from \(0<a<b<c\).  In each row the positive and negative
areas agree.  Directly, those common areas are respectively
\begin{align*}
 &(b-a)(c-b),
 &&(a+b)(c-b),\\
 &2a(c-a),
 &&(b-a)(c-b+2a).
\end{align*}
This proves both \eqref{eq:v26-kernel-zero-mass} and the last column of the
table.  Since \(|\sigma_\rho|\le1\), the four inequalities follow.  Equality
in each absolute bound is obtained in the enlarged weight class by taking
\(\sigma_\rho\) to be either sign of \(K_m(\rho)\).
\end{proof}

\begin{corollary}[Exact comparison with productive critical work]
\label{cor:v26-adaptive-vs-critical}
For a productive mixed triad let
\begin{equation}
 \mathcal W^{\triangle}
 :=\frac12\mathcal G_{1/2}^{\triangle,s}>0
 \label{eq:v26-common-sector-work}
\end{equation}
be the common critical work of either helical sector.  Then
\begin{align}
 \frac{|\mathfrak a_1^{\triangle}|}
      {2\mathcal W^{\triangle}}
 &\le\frac{a+b}{a},
 \label{eq:v26-m1-ratio}\\
 \frac{|\mathfrak a_2^{\triangle}|}
      {2\mathcal W^{\triangle}}
 &\le\frac{2a}{b},
 \label{eq:v26-m2-ratio}\\
 \frac{|\mathfrak a_3^{\triangle}|}
      {2\mathcal W^{\triangle}}
 &\le\frac{c-b+2a}{c}.
 \label{eq:v26-m3-ratio}
\end{align}
The low-minority coefficient in \eqref{eq:v26-m1-ratio} is unbounded as
\(a/b\to0\), whereas both high-minority coefficients are at most one when
\(a\le b/2\).
\end{corollary}

\begin{proof}
For positive critical work, the exact V23 sector-work table
\eqref{eq:v23-common-sector-triad-work} says
\begin{equation}
 \mathcal W^{\triangle}=
 \begin{cases}
  a(c-b)|\Xi|,&m=1,\\
  b(c-a)|\Xi|,&m=2,\\
  c(b-a)|\Xi|,&m=3.
 \end{cases}
 \label{eq:v26-sector-work-table-recalled}
\end{equation}
Divide \eqref{eq:v26-m1-adaptive-bound}--%
\eqref{eq:v26-m3-adaptive-bound} by twice these quantities.  If
\(a\le b/2\), then \(2a/b\le1\), and
\(c-b+2a\le c\).  This proves the final statement.
\end{proof}

Define the contribution of one complete triad to the assembled smaller-tail
flux by
\begin{equation}
 \mathfrak b_\wedge^{\triangle}(t)
 :=\mathcal W^{\triangle}(t)
   -\frac12\mathfrak a^{\triangle}(t).
 \label{eq:v26-triad-meet-work}
\end{equation}
This is the triad summand of the instantaneous version of
\eqref{eq:v25-meet-equal-work-correction}; the global selector
\(\sigma_\rho(t)\) is held fixed while complete triads are summed.

\begin{theorem}[Robust positive meet work for strongly nonlocal
high-minority triads]
\label{thm:v26-nonlocal-high-minority-positive}
Let \(0<\delta<1/2\).  Suppose a mixed triad has positive critical work, its
minority helicity is on leg \(m=2\) or \(m=3\), and
\begin{equation}
 a\le\delta b.
 \label{eq:v26-strict-nonlocality}
\end{equation}
Then for the fully assembled tail-order selector, without any regularity or
monotonicity assumption on \(\sigma_\rho\),
\begin{equation}
 \boxed{\quad
 \mathfrak b_\wedge^{\triangle}(t)
 \ge\frac{1-2\delta}{1+\delta}
       \mathcal W^{\triangle}(t)>0.
 \quad}
 \label{eq:v26-positive-meet-fraction}
\end{equation}
At the endpoint \(\delta=1/2\), the same proof gives the nonnegative bound
\(\mathfrak b_\wedge^{\triangle}\ge0\).  In particular, for
\(a\le b/3\), at least one quarter of the common sector critical work
survives in the smaller-tail flux.
\end{theorem}

\begin{proof}
For \(m=2\), \eqref{eq:v26-m2-ratio} gives
\[
 |\mathfrak a_2^{\triangle}|
 \le 4\delta\mathcal W^{\triangle}.
\]
Therefore
\[
 \mathfrak b_\wedge^{\triangle}
 \ge(1-2\delta)\mathcal W^{\triangle}
 \ge\frac{1-2\delta}{1+\delta}\mathcal W^{\triangle}.
\]
For \(m=3\), the exact ratio and the Fourier triangle inequality give
\begin{align*}
 1-\frac{c-b+2a}{c}
 &=\frac{b-2a}{c}
 \ge\frac{(1-2\delta)b}{(1+\delta)b}
 =\frac{1-2\delta}{1+\delta}.
\end{align*}
Thus
\(|\mathfrak a_3^{\triangle}|/2\le
 [1-(1-2\delta)/(1+\delta)]\mathcal W^{\triangle}\), and
\eqref{eq:v26-triad-meet-work} proves
\eqref{eq:v26-positive-meet-fraction}.  The endpoint follows by continuity;
setting \(\delta=1/3\) gives the fraction \(1/4\).
\end{proof}

\begin{assessment}[Relation to the V23 overlap corridor]
\label{ass:v26-relation-v23}
V23 proved positive overlap of the two \emph{triad} sector profiles for
high-minority nonlocal interactions, before aggregation.  Theorem
\ref{thm:v26-nonlocal-high-minority-positive} is different: the selector in
\(\mathfrak b_\wedge^{\triangle}\) is computed from the two fully assembled
PDE tails, and may oscillate arbitrarily in the cutoff.  Complete-triad
energy--helicity cancellation nevertheless leaves a positive fraction for
the strongly nonlocal productive high-minority class.  This survives the
adaptive selection step, but not yet the signed summation with all remaining
triads.
\end{assessment}
% =============================================================================
\section{Aggregate productive class and its exact complement}
\label{sec:v26-aggregate-class}
% =============================================================================

For every mixed triad now let
\(\mathcal W^{\triangle}:=mathcal G_{1/2}^{\triangle,s}/2\) carry its
actual sign, and set \(\mathcal W^{\triangle}=0\) for a homochiral triad.
Extend \eqref{eq:v26-triad-meet-work} with this signed convention.  Smooth
Fourier summability permits the complete reality-paired triad expansion
\begin{equation}
 \mathfrak g_\wedge(t)
 :=\int_0^\infty\widetilde\Phi_\rho^\wedge(t)\,\dd\rho
 =\sum_{\triangle,\boldsymbol s}
   \mathfrak b_\wedge^{\triangle}(t).
 \label{eq:v26-global-meet-triad-sum}
\end{equation}

Fix \(0<\delta<1/2\), and let \(\mathfrak P_\delta(t)\) be the class of
mixed triads which at time \(t\) have positive critical work, high minority
helicity \(m\in\{2,3\}\), and
\(\lambda_1\le\delta\lambda_2\).  Define
\begin{align}
 \mathcal W_{\mathfrak P_\delta}(t)
 &:=\sum_{\triangle\in\mathfrak P_\delta(t)}
      \mathcal W^{\triangle}(t),
 \label{eq:v26-selected-productive-work}\\
 \mathfrak C_\delta(t)
 &:=\sum_{\triangle\notin\mathfrak P_\delta(t)}
      \mathfrak b_\wedge^{\triangle}(t).
 \label{eq:v26-meet-complement}
\end{align}

\begin{theorem}[Exact robust-class decomposition]
\label{thm:v26-robust-class-decomposition}
At every smooth time,
\begin{align}
 \mathfrak g_\wedge(t)
 &=\sum_{\triangle\in\mathfrak P_\delta(t)}
    \mathfrak b_\wedge^{\triangle}(t)
   +\mathfrak C_\delta(t),
 \label{eq:v26-selected-plus-complement}\\
 \boxed{\quad
 \mathfrak g_\wedge(t)
 \ge\frac{1-2\delta}{1+\delta}
       \mathcal W_{\mathfrak P_\delta}(t)
    +\mathfrak C_\delta(t).
 \quad}
 \label{eq:v26-robust-class-lower}
\end{align}
The complement consists exactly of homochiral triads, low-minority triads,
high-minority triads with \(\lambda_1>\delta\lambda_2\), and nonproductive
mixed triads.  No member is discarded or estimated by gross absolute action
in \eqref{eq:v26-robust-class-lower}.
\end{theorem}

\begin{proof}
Equation~\eqref{eq:v26-selected-plus-complement} is the partition of the
absolutely convergent signed sum
\eqref{eq:v26-global-meet-triad-sum}.  Apply
Theorem~\ref{thm:v26-nonlocal-high-minority-positive} to each member of
\(\mathfrak P_\delta\) and sum its nonnegative lower bounds.  The stated list
is the set-theoretic complement of the three defining conditions for
\(\mathfrak P_\delta\), with the homochiral class included separately.
\end{proof}

\begin{firewall}[The complement is not paid by robust positivity]
\label{fw:v26-complement-unpaid}
The first term on the right of \eqref{eq:v26-robust-class-lower} is a sum over
a sign-selected productive class.  It becomes a net lower bound only after
\(\mathfrak C_\delta\) is controlled.  Dropping the complement would repeat
the positive-triad aggregation error in
Firewall~\ref{fw:v23-positive-triad-overlap}.  V26 has improved the selected
summands---their sign survives the global adaptive selector---but has not
proved that the full signed sum is positive.
\end{firewall}

% =============================================================================
\section{The low-minority obstruction is realized by one active triad}
\label{sec:v26-low-minority-obstruction}
% =============================================================================

\begin{proposition}[A productive low-minority triad can give negative meet
work]
\label{prop:v26-low-minority-negative}
Consider an active mixed triad with negative minority helicity on the low leg,
\(m=1\).  Choose its modal energy magnitudes so that
\begin{equation}
 e_1>e_2+e_3>0,
 \label{eq:v26-low-minority-energy-order}
\end{equation}
where \(e_i\) is the energy in leg \(i\), and choose the interaction phase so
that \(\Xi>0\).  Then its critical work is productive and its own tail-order
sign is
\begin{equation}
 \sigma_\rho=
 \begin{cases}
  -1,&0<\rho<a,\\
  +1,&a<\rho<c,
 \end{cases}
 \label{eq:v26-low-minority-tail-sign}
\end{equation}
up to cutoff endpoints.  For this selector,
\begin{align}
 \mathfrak a_1^{\triangle}
 &=4a(c-b)\Xi=4\mathcal W^{\triangle},
 \label{eq:v26-low-minority-adaptive-fourfold}\\
 \boxed{\quad
 \mathfrak b_\wedge^{\triangle}
 =-\mathcal W^{\triangle}<0.
 \quad}
 \label{eq:v26-low-minority-negative-meet}
\end{align}
Thus even one complete productive triad can lie in the negative complement;
the obstruction in \eqref{eq:v26-m1-ratio} is not only an artifact of an
arbitrary sign-weight bound.
\end{proposition}

\begin{proof}
For \(m=1\), the negative sector contains only the low leg.  Below \(a\),
the signed tail is
\(\mathcal Y_\rho=e_2+e_3-e_1<0\).  Above \(a\), the negative-sector tail
vanishes while at least one positive-sector high leg remains until \(c\), so
\(\mathcal Y_\rho>0\).  This proves
\eqref{eq:v26-low-minority-tail-sign}.  Insert that sign into the \(m=1\)
row of \eqref{eq:v26-adaptive-kernel-table}:
\begin{align*}
 \frac{\mathfrak a_1^{\triangle}}{\Xi}
 &=2a(c-b)-(b-a)(c-b)+(a+b)(c-b)\\
 &=4a(c-b).
\end{align*}
For \(m=1\), \eqref{eq:v26-sector-work-table-recalled} gives
\(\mathcal W^{\triangle}=a(c-b)\Xi>0\).  Therefore
\(\mathfrak b_\wedge^{\triangle}
 =\mathcal W^{\triangle}-\mathfrak a_1^{\triangle}/2
 =-\mathcal W^{\triangle}\), proving both claims.  Activity permits a phase
with \(\Xi>0\) by Lemma~\ref{lem:v22-helical-transfer-normal-form}.
\end{proof}

\begin{assessment}[Geometric meaning of the dangerous low minority]
\label{ass:v26-low-minority-meaning}
The sign switch in \eqref{eq:v26-low-minority-tail-sign} occurs when the
energetically dominant low negative-helicity leg drops out of the high tail.
Above that radius only the two positive-helicity legs remain.  This is the
same minority placement for which V23 found disjoint positive sector-flux
supports.  V26 strengthens that diagnosis: after forming the fully adaptive
smaller-tail flux, the productive low-minority triad can contribute with the
wrong sign and the same magnitude as its common sector critical work.
\end{assessment}

% =============================================================================
\section{Critical stock does not bound radial tail-order crossings}
\label{sec:v26-unbounded-crossings}
% =============================================================================

The divided-difference formula suggests asking whether
\(\sigma_\rho\) can change sign only a bounded number of times.  The answer is
no even for smooth finite Fourier polynomials.

\begin{theorem}[Arbitrarily many tail-order crossings at fixed critical
stock]
\label{thm:v26-unbounded-tail-crossings}
For every integer \(N\ge2\), there is a real, mean-zero, smooth,
divergence-free Fourier polynomial on \(\mathbb T^3\) such that
\begin{equation}
 \mathcal R^2=1,
 \qquad
 \|u\|_2^2\le1,
 \label{eq:v26-crossing-normalization}
\end{equation}
while the signed helical tail \(\mathcal Y_\rho\) alternates sign on each of
the \(N\) consecutive intervals
\begin{equation}
 [0,1),[1,2),\ldots,[N-1,N).
 \label{eq:v26-crossing-intervals}
\end{equation}
Hence no bound on the number or total variation of the radial selector
\(\sigma_\rho=\operatorname{sgn}\mathcal Y_\rho\) follows from either the
Leray energy ceiling or a fixed \(\dot H^{1/2}\) stock.
\end{theorem}

\begin{proof}
Let \(r_j=j\) and prescribe nonzero tail values
\begin{equation}
 y_j:=(-1)^j\varepsilon,
 \qquad 1\le j\le N,
 \qquad y_{N+1}:=0.
 \label{eq:v26-prescribed-tail-values}
\end{equation}
Set signed shell weights
\begin{equation}
 \alpha_j:=y_j-y_{j+1}.
 \label{eq:v26-signed-shell-weights}
\end{equation}
Thus \(\alpha_j=2(-1)^j\varepsilon\) for \(j<N\), and
\(\alpha_N=(-1)^N\varepsilon\).  At radius \(r_j\), put energy
\((\alpha_j)_+\) in a positive-helicity reality pair and energy
\((\alpha_j)_-\) in a negative-helicity reality pair.  Such pairs exist at
the lattice wavevectors \(\pm(j,0,0)\); choose the usual conjugate helical
coefficients to make the field real.  The result is a finite, mean-zero,
divergence-free Fourier polynomial.

For \(r_{j-1}<\rho<r_j\), telescoping gives
\begin{equation}
 \mathcal Y_\rho
 =\sum_{\ell=j}^N\alpha_\ell
 =y_j,
 \label{eq:v26-tail-telescoping}
\end{equation}
so the signs alternate on all intervals in
\eqref{eq:v26-crossing-intervals}.  Put
\(D_N:=2\sum_{j=1}^{N-1}j+N\) and choose
\begin{equation}
 \varepsilon:=D_N^{-1}.
 \label{eq:v26-crossing-epsilon}
\end{equation}
Then spectral layer cake gives
\[
 \mathcal R^2=\sum_{j=1}^Nr_j|\alpha_j|=1.
\]
Since every \(r_j\ge1\),
\(\|u\|_2^2=\sum_j|\alpha_j|\le
\sum_jr_j|\alpha_j|=1\).  This proves
\eqref{eq:v26-crossing-normalization} and the claimed absence of a crossing
bound.
\end{proof}

\begin{corollary}[A bounded-crossing closure is unavailable]
\label{cor:v26-no-bounded-crossing-closure}
Neither \eqref{eq:v26-primitive-Lipschitz} nor the global energy and critical
stock bounds can be upgraded to a uniform finite partition on which
\(\sigma_\rho\) is constant.  Any use of the vanishing clause in
Theorem~\ref{thm:v26-triad-determinant} must therefore be weighted by the
actual triad incidence or by additional PDE information; it cannot be summed
by counting radial sign intervals.
\end{corollary}

\begin{proof}
Theorem~\ref{thm:v26-unbounded-tail-crossings} supplies, for every proposed
finite bound, a smooth field satisfying both stock ceilings and having more
sign intervals.  The conclusion follows.
\end{proof}
\begin{corollary}[Record-interval form of the robust-class decomposition]
\label{cor:v26-robust-class-interval}
Define
\begin{align}
 \mathcal W_{\mathfrak P_\delta}(I)
 &:=\int_I\mathcal W_{\mathfrak P_\delta}(t)\,\dd t,
 \label{eq:v26-selected-work-interval}\\
 \mathfrak C_\delta(I)
 &:=\int_I\mathfrak C_\delta(t)\,\dd t.
 \label{eq:v26-complement-interval}
\end{align}
Then every smooth interval obeys
\begin{equation}
 \boxed{\quad
 \mathcal B_\wedge(I)
 \ge\frac{1-2\delta}{1+\delta}
       \mathcal W_{\mathfrak P_\delta}(I)
    +\mathfrak C_\delta(I).
 \quad}
 \label{eq:v26-robust-class-interval-lower}
\end{equation}
On a first record doubling this lower bound must still be reconciled with the
exact meet balance \eqref{eq:v25-integrated-meet-balance}; no global sign for
\(\mathfrak C_\delta(I)\) has been proved.
\end{corollary}

\begin{proof}
Integrate \eqref{eq:v26-robust-class-lower} in time and use the definition
\eqref{eq:v25-meet-nonlinear-work} of \(\mathcal B_\wedge(I)\).
\end{proof}

% =============================================================================
\section{Integrated V26 conclusion and next acquisition order}
\label{sec:v26-conclusion}
% =============================================================================

\begin{theorem}[What V26 proves]
\label{thm:v26-conclusion}
Relative to V1--V25, the following are proved.
\begin{enumerate}[label=\textup{(V26.\arabic*)},leftmargin=6.2em]
\item The adaptive imbalance source is exactly the order-one multiplier
      pairing and transport commutator
      \eqref{eq:v26-source-multiplier-form}--%
      \eqref{eq:v26-source-commutator-form}.
\item On every complete helical triad, that source is the determinant
      \eqref{eq:v26-adaptive-triad-determinant}, equivalently the signed-radius
      second divided difference
      \eqref{eq:v26-adaptive-divided-difference}.  Energy and helicity remove
      the affine part before any estimate is taken.
\item The complete cutoff-kernel table
      \eqref{eq:v26-adaptive-kernel-table} and its sharp \(L^1\) masses hold
      for homochiral and all three minority-helicity locations.
\item The high-minority adaptive coefficients are controlled by productive
      critical work as in \eqref{eq:v26-m2-ratio}--%
      \eqref{eq:v26-m3-ratio}, whereas the low-minority ratio
      \eqref{eq:v26-m1-ratio} degenerates in the nonlocal limit.
\item Every productive high-minority triad with
      \(\lambda_1\le\delta\lambda_2\), \(\delta<1/2\), contributes at least
      \((1-2\delta)/(1+\delta)\) of its common sector work to the smaller
      assembled tail.  This conclusion is uniform over the global adaptive
      selector.
\item The selected robust class and its complete signed complement satisfy
      the exact instantaneous and interval ledgers
      \eqref{eq:v26-robust-class-lower} and
      \eqref{eq:v26-robust-class-interval-lower}.
\item An active productive low-minority triad can have
      \(\mathfrak b_\wedge^{\triangle}=-\mathcal W^{\triangle}\), proving
      that the negative complement is genuine.
\item Smooth Fourier polynomials with fixed kinetic-energy and critical-stock
      ceilings can have arbitrarily many radial tail-order sign changes.
      Hence no bounded-crossing closure follows from those stocks.
\end{enumerate}
The robust productive class is a strict new positive result, but the full
Navier--Stokes stability/global-regularity theorem remains unproved because
its signed complement has no Leray-level lower bound.
\end{theorem}

\begin{proof}
Item \textup{(V26.1)} is Theorem~\ref{thm:v26-adaptive-commutator}.  Item
\textup{(V26.2)} is Theorem~\ref{thm:v26-triad-determinant}.  Item
\textup{(V26.3)} is Theorem~\ref{thm:v26-adaptive-kernel-table}.  Item
\textup{(V26.4)} is Corollary~\ref{cor:v26-adaptive-vs-critical}.  Item
\textup{(V26.5)} is
Theorem~\ref{thm:v26-nonlocal-high-minority-positive}.  Item
\textup{(V26.6)} is Theorem~\ref{thm:v26-robust-class-decomposition} and
Corollary~\ref{cor:v26-robust-class-interval}.  Item \textup{(V26.7)} is
Proposition~\ref{prop:v26-low-minority-negative}.  Item \textup{(V26.8)} is
Theorem~\ref{thm:v26-unbounded-tail-crossings}.
\end{proof}

\begin{programme}[V27 complement acquisition order]
\label{prog:v27}
\begin{enumerate}[label=\textup{(AC\arabic*)},leftmargin=4.8em]
\item Preserve the exact decomposition
      \eqref{eq:v26-selected-plus-complement}.  Do not drop
      \(\mathfrak C_\delta\) or replace it by the absolute sum of its triads.
\item Begin with the realized low-minority mechanism in
      Proposition~\ref{prop:v26-low-minority-negative}.  Relate its negative
      meet work to the V24 separation profile and to the endpoint-sign
      pairings in \eqref{eq:v25-endpoint-sign-sandwich}.
\item Treat homochiral crossing triads separately.  Their kernel vanishes
      below \(\lambda_1\) and has the adjacent-gap mass
      \(2(\lambda_2-\lambda_1)(\lambda_3-\lambda_2)\); test whether complete
      shell incidence pays this second difference without gross action.
\item Split the remaining high-minority class into the robust nonlocal region
      and the near-local region.  Any log-shell or neighbor-incidence payment
      for the latter must be derived from the signed complete sum.
\item Keep nonproductive mixed triads in the complement until an exact pairing
      with productive triads is proved.  Phase reversal prevents an
      amplitude-only sign assertion.
\item Do not attempt to control the selector by counting its sign changes;
      Theorem~\ref{thm:v26-unbounded-tail-crossings} rules out that route at
      fixed Leray and critical stock.
\item Carry the viscosity-weighted tail-order term from
      \eqref{eq:v25-terminal-crossing-exact} through every estimate.  Closure
      requires a Leray-level debit, a one-use source stock, or a proved
      cancellation with the complete complement.
\item Schedule the next compulsory SM/YM companion audit for V30, after V27--
      V29 have attacked this exact complement.
\end{enumerate}
\end{programme}

\begin{assessment}[Position after V26]
The adaptive sign is no longer an opaque coefficient.  Its primitive becomes
an odd one-Lipschitz function of signed helical radius, and complete-triad
energy--helicity cancellation extracts exactly its second divided difference.
This proves robust positivity for a geometrically important class even after
the selector has been assembled from the full solution.

The same exact table also identifies why the result is not yet global.
Low-minority productive triads can drive the smaller-tail work negatively;
homochiral and near-local interactions remain signed; and the selector may
have arbitrarily many radial changes under the available stock bounds.  The
next upstream target is therefore the complete signed complement, beginning
with the low-minority/separation correspondence, not a regularity assumption
on the selector.  The full stability/global-regularity objective remains
open.
\end{assessment}

% =============================================================================
\section*{Version V26 append-only record}
\addcontentsline{toc}{section}{Version V26 append-only record}
% =============================================================================

The complete V25 mathematical source was retained before this continuation.
Its SHA--256 digest was
\begin{center}\scriptsize\ttfamily
0d80d813253256366be9a515b5c7da5647a29f8dc8f2c2c76fe12b92bc653415.
\end{center}
V26 rewrote the V25 adaptive source as an exact commutator and as a
complete-triad determinant on signed helical radii.  It derived the four-row
cutoff-kernel table, proved robust positive smaller-tail work for strongly
nonlocal productive high-minority triads, retained the exact signed
complement, exhibited a productive low-minority triad with negative meet
work, and constructed fixed-stock smooth spectra with arbitrarily many radial
tail-order crossings.  The complement remains unpaid, so the full stability/
global-regularity goal remains open.
% =============================================================================
% VERSION V27 APPEND
% =============================================================================

\clearpage
\section{V27: radial crossing currents and the normalized hinge ledger}
\label{sec:v27-entry}
% =============================================================================

V26 identified an exact robust class of productive high-minority triads and
retained its full signed complement.  It also proved that the raw number of
tail-order sign changes is unbounded at fixed kinetic energy and critical
stock.  Thus V27 does not try to count crossings.  It instead asks which
\emph{assembled current} is sampled when the sign changes.

The answer is a capped-helicity, or hinge, current.  Radial integration of the
sector-difference flux produces a quantity \(\mathcal J_r\) which is exactly
the nonlinear production of the signed stock with multiplier
\(s\min\{r,|k|\}\).  After division by \(r\), both its endpoint stock and its
viscous term are paid by the Leray energy ledger at every fixed cutoff.  A
second radial summation-by-parts identity shows that the full V25 adaptive
source samples these currents only at shells where the assembled helical-tail
order changes.

This exposes the remaining nonuniformity precisely.  Each fixed normalized
hinge current is Leray-paid, and the shell amplitudes which cause crossings
are one-use at a fixed time, but the sign jump divides out that amplitude.
Arbitrarily small shell imbalances can therefore carry a full sign jump.  The
missing theorem is a signed crossing-incidence estimate, not a bound on the
number of crossings and not a termwise triad estimate.

% =============================================================================
\section{The cumulative sector-difference current}
\label{sec:v27-hinge-current}
% =============================================================================

For \(r>0\), define
\begin{align}
 \mathcal J_r(t)
 &:=\int_0^r d_\rho(t)\,\dd\rho,
 \label{eq:v27-cumulative-difference-current}\\
 \mathcal Z_r(t)
 &:=\int_0^r\mathcal Y_\rho(t)\,\dd\rho,
 \label{eq:v27-hinge-stock}\\
 \mathcal D_r(t)
 &:=\int_0^r
   \bigl(V_\rho^+(t)-V_\rho^-(t)\bigr)\,\dd\rho.
 \label{eq:v27-hinge-dissipation}
\end{align}

\begin{theorem}[Exact modal and PDE formulae for the hinge current]
\label{thm:v27-hinge-current}
At every smooth time,
\begin{align}
 \boxed{\quad
 \mathcal J_r(t)
 =-\sum_{k\ne0}\sum_s
    s\min\{r,|k|\}T_{k,s}(t)
 =\sum_{|k|>r}\sum_s
    s(|k|-r)T_{k,s}(t).
 \quad}
 \label{eq:v27-hinge-current-modal}\\
 \mathcal Z_r(t)
 &=\sum_{k\ne0}\sum_s
   s\min\{r,|k|\}|\widehat u^s(k,t)|^2,
 \label{eq:v27-hinge-stock-modal}\\
 \boxed{\quad
 \mathcal J_r(t)
 =\frac12\partial_t\mathcal Z_r(t)+\nu\mathcal D_r(t).
 \quad}
 \label{eq:v27-hinge-balance}
\end{align}
The hinge stock and dissipation satisfy
\begin{align}
 |\mathcal Z_r(t)|
 &\le\min\{r\|u(t)\|_2^2,\mathcal R(t)^2\},
 \label{eq:v27-hinge-stock-cap}\\
 |\mathcal D_r(t)|
 &\le r\|\Lambda u(t)\|_2^2.
 \label{eq:v27-hinge-dissipation-cap}
\end{align}
\end{theorem}

\begin{proof}
Integrate the modal formula for \(d_\rho\) used in the proof of
Theorem~\ref{thm:v26-adaptive-commutator}.  A mode of radius \(|k|\) is
present over a cutoff interval of length \(\min\{r,|k|\}\), giving the first
sum in \eqref{eq:v27-hinge-current-modal}.  Nonlinear helicity cancellation
is
\begin{equation}
 \sum_{k,s}s|k|T_{k,s}(t)=0.
 \label{eq:v27-modal-helicity-cancellation}
\end{equation}
Separate the first sum into \(|k|\le r\) and \(|k|>r\), and use
\eqref{eq:v27-modal-helicity-cancellation} on the low part.  The result is
the second sum in \eqref{eq:v27-hinge-current-modal}.

Integrating the signed tail stock in \(\rho\) gives
\eqref{eq:v27-hinge-stock-modal}.  Integrate the instantaneous tail-difference
balance \eqref{eq:v25-instant-tail-difference} from zero to \(r\) to obtain
\eqref{eq:v27-hinge-balance}.

For each mode,
\(\min\{r,|k|\}\le r\) and
\(\min\{r,|k|\}\le|k|\).  Taking absolute values in
\eqref{eq:v27-hinge-stock-modal} proves
\eqref{eq:v27-hinge-stock-cap}.  Similarly, spectral Fubini gives
\[
 |\mathcal D_r|
 \le\sum_{k,s}\min\{r,|k|\}|k|^2|\widehat u^s(k)|^2
 \le r\|\Lambda u\|_2^2,
\]
which proves \eqref{eq:v27-hinge-dissipation-cap}.
\end{proof}

Normalize the hinge quantities by
\begin{equation}
 j_r:=\frac{\mathcal J_r}{r},
 \qquad z_r:=\frac{\mathcal Z_r}{r},
 \qquad q_r:=\frac{\mathcal D_r}{r}.
 \label{eq:v27-normalized-hinge-quantities}
\end{equation}

\begin{corollary}[Every fixed normalized hinge current is Leray-paid]
\label{cor:v27-fixed-hinge-Leray}
For every fixed \(r>0\) and smooth interval \(I=[a,b]\),
\begin{align}
 \int_Ij_r(t)\,\dd t
 &=\frac12\bigl(z_r(b)-z_r(a)\bigr)
   +\nu\int_Iq_r(t)\,\dd t,
 \label{eq:v27-normalized-hinge-balance}\\
 \boxed{\quad
 \left|\int_Ij_r(t)\,\dd t\right|
 \le\frac12\bigl(\|u(a)\|_2^2+\|u(b)\|_2^2\bigr)
   +\nu\int_I\|\Lambda u(t)\|_2^2\,\dd t.
 \quad}
 \label{eq:v27-fixed-hinge-Leray-bound}
\end{align}
In particular, under the programme's energy ceiling and global Leray
dissipation inequality, the right side is at most \(3E_*/2\).
\end{corollary}

\begin{proof}
Divide \eqref{eq:v27-hinge-balance} by \(r\), integrate in time, and apply
\eqref{eq:v27-hinge-stock-cap}--\eqref{eq:v27-hinge-dissipation-cap}.  The
definition \(E_*\ge\sup_t\|u(t)\|_2^2\) bounds the two endpoint halves by
\(E_*\), while the Leray inequality bounds the last term by \(E_*/2\).
\end{proof}

\begin{remark}[Why the normalization matters]
\label{rem:v27-hinge-normalization}
The unnormalized stock \(\mathcal Z_r\) is critical and grows at most like
\(rE_*\).  Division by the hinge radius removes that multiplicity and turns
both the stock and viscous terms into the ordinary kinetic-energy and
enstrophy-dissipation ledger.  The obstruction is therefore no longer a
single cutoff.  It is the simultaneous sampling of many normalized hinge
currents by the crossing measure derived next.
\end{remark}
% =============================================================================
\section{Radial summation by parts: only crossing shells are sampled}
\label{sec:v27-crossing-summation}
% =============================================================================

We first state the identity for a finite Galerkin field, where every radial
quantity is a step function and no bounded-variation hypothesis is hidden.
Let
\begin{equation}
 0<r_1<\cdots<r_N
 \label{eq:v27-occupied-radii}
\end{equation}
be its occupied radii and put \(r_0=0\).  On
\((r_{j-1},r_j)\), write
\begin{equation}
 Y_j(t):=\mathcal Y_\rho(t),
 \qquad \sigma_j(t):=\operatorname{sgn}Y_j(t),
 \qquad 1\le j\le N,
 \label{eq:v27-discrete-tail-values}
\end{equation}
and set \(Y_{N+1}=\sigma_{N+1}=0\).  Define the signed helical energy on the
shell \(r_j\) and the sign increment across it by
\begin{align}
 \eta_j(t)
 &:=\sum_{|k|=r_j}
   \bigl(|\widehat u^+(k,t)|^2-|\widehat u^-(k,t)|^2\bigr)
 =Y_j(t)-Y_{j+1}(t),
 \label{eq:v27-signed-shell-mass}\\
 \Delta\sigma_j(t)&:=\sigma_{j+1}(t)-\sigma_j(t).
 \label{eq:v27-radial-sign-increment}
\end{align}

\begin{theorem}[Exact crossing-shell representation of the adaptive source]
\label{thm:v27-crossing-shell-representation}
For every finite Galerkin field,
\begin{align}
 \boxed{\quad
 \mathfrak a(t)
 =-\sum_{j=1}^{N}
   \mathcal J_{r_j}(t)\Delta\sigma_j(t)
 =-\sum_{j=1}^{N}
   r_j\,j_{r_j}(t)\Delta\sigma_j(t).
 \quad}
 \label{eq:v27-crossing-shell-source}
\end{align}
The top-radius term vanishes because \(\mathcal J_{r_N}=0\).  Thus shells
across which the assembled helical-tail order does not change make no direct
appearance in the adaptive source.
\end{theorem}

\begin{proof}
Since \(d_\rho\) is constant between occupied radii,
\begin{align*}
 \mathfrak a(t)
 &=\sum_{j=1}^N\sigma_j(t)
   \int_{r_{j-1}}^{r_j}d_\rho(t)\,\dd\rho\\
 &=\sum_{j=1}^N\sigma_j(t)
   \bigl(\mathcal J_{r_j}(t)-\mathcal J_{r_{j-1}}(t)\bigr).
\end{align*}
Discrete summation by parts, with \(\mathcal J_{r_0}=0\), gives
\[
 \mathfrak a(t)
 =-\sum_{j=1}^{N-1}\mathcal J_{r_j}(t)
    (\sigma_{j+1}(t)-\sigma_j(t))
  +\sigma_N(t)\mathcal J_{r_N}(t).
\]
Above \(r_N\), both compensated sector fluxes vanish.  Moreover
\(\int_0^\infty d_\rho\,\dd\rho
 =\mathcal G_+-\mathcal G_-=0\) by
\eqref{eq:v23-equal-instant-sector-work}.  Hence
\(\mathcal J_{r_N}=0\).  Adding its zero term with
\(\sigma_{N+1}=0\) proves both forms in
\eqref{eq:v27-crossing-shell-source}.
\end{proof}

\begin{lemma}[Every sign jump is carried by an oppositely signed shell]
\label{lem:v27-crossing-shell-budget}
If \(\Delta\sigma_j\ne0\), then
\begin{align}
 Y_jY_{j+1}&\le0,
 \label{eq:v27-crossing-opposite-tails}\\
 \Delta\sigma_j\,\eta_j&<0,
 \label{eq:v27-crossing-opposite-shell}\\
 |\eta_j|&=|Y_j|+|Y_{j+1}|.
 \label{eq:v27-crossing-amplitude}
\end{align}
At every fixed time, the crossing-shell amplitudes are one-use:
\begin{align}
 \sum_{j:\Delta\sigma_j\ne0}|\eta_j(t)|
 &\le\|u(t)\|_2^2,
 \label{eq:v27-crossing-energy-budget}\\
 \sum_{j:\Delta\sigma_j\ne0}r_j|\eta_j(t)|
 &\le\mathcal R(t)^2.
 \label{eq:v27-crossing-critical-budget}
\end{align}
\end{lemma}

\begin{proof}
A nonzero increment of the three-valued sign function means that
\(Y_j\) and \(Y_{j+1}\) lie in different closed half-lines and are not both
zero.  This proves \eqref{eq:v27-crossing-opposite-tails}.  Since
\(\eta_j=Y_j-Y_{j+1}\), an upward sign change has \(\eta_j<0\), while a
downward sign change has \(\eta_j>0\).  This proves
\eqref{eq:v27-crossing-opposite-shell}, and subtraction of opposite-signed
numbers proves \eqref{eq:v27-crossing-amplitude}.

Let
\(e_j^\pm:=\sum_{|k|=r_j}|\widehat u^\pm(k)|^2\).  Then
\(|\eta_j|=|e_j^+-e_j^-|\le e_j^++e_j^-\).  Sum over the subset of crossing
shells to obtain \eqref{eq:v27-crossing-energy-budget}; insert the weight
\(r_j\) to obtain \eqref{eq:v27-crossing-critical-budget}.
\end{proof}

For every crossing shell define its inverse-amplitude weight
\begin{equation}
 \omega_j(t):=\frac{|\Delta\sigma_j(t)|}{|\eta_j(t)|}.
 \label{eq:v27-crossing-inverse-amplitude}
\end{equation}
The preceding lemma gives
\begin{equation}
 \Delta\sigma_j=-\omega_j\eta_j,
 \qquad
 \mathfrak a(t)
 =\sum_{j:\Delta\sigma_j\ne0}
   r_j\,j_{r_j}(t)\omega_j(t)\eta_j(t).
 \label{eq:v27-crossing-leverage-form}
\end{equation}

\begin{proposition}[The precise crossing leverage left unpaid]
\label{prop:v27-crossing-leverage}
The energy and critical budgets
\eqref{eq:v27-crossing-energy-budget}--%
\eqref{eq:v27-crossing-critical-budget} control the measures
\(|\eta_j|\) and \(r_j|\eta_j|\), but they do not control
\(\omega_j\).  More precisely, for every \(L>0\) there is a smooth finite
Fourier polynomial satisfying
\(\mathcal R^2=1\), \(\|u\|_2^2\le1\), and having a crossing shell with
\begin{equation}
 \omega_j>L.
 \label{eq:v27-unbounded-crossing-leverage}
\end{equation}
Thus a proof which combines the fixed-cutoff estimate
\eqref{eq:v27-fixed-hinge-Leray-bound} with the shell budget must retain the
signed incidence in \eqref{eq:v27-crossing-leverage-form}; H\"older or
triangle inequality leaves an unbounded inverse-amplitude factor.
\end{proposition}

\begin{proof}
Use the construction in
Theorem~\ref{thm:v26-unbounded-tail-crossings}.  There
\(|\Delta\sigma_j|=2\) at every interior sign reversal, while
\(|\eta_j|=2\varepsilon\).  Hence \(\omega_j=\varepsilon^{-1}\).  The
normalization there has
\(\varepsilon^{-1}=2\sum_{\ell=1}^{N-1}\ell+N\to\infty\), while both stock
bounds remain fixed.  Choose \(N\) large enough that
\(\varepsilon^{-1}>L\).
\end{proof}

\begin{firewall}[One-use crossing amplitude is not one-use sign variation]
\label{fw:v27-amplitude-not-sign}
The bounds \eqref{eq:v27-crossing-energy-budget}--%
\eqref{eq:v27-crossing-critical-budget} are genuine same-time one-use
statements.  Equation~\eqref{eq:v27-crossing-shell-source}, however, samples
the unweighted jump \(\Delta\sigma_j\), not the shell mass \(\eta_j\).
Replacing the former by the latter without the factor \(\omega_j\) would
discard the near-zero tail-order regime and is false by
Proposition~\ref{prop:v27-crossing-leverage}.
\end{firewall}
% =============================================================================
\section{Triad cumulative currents and crossing orientation}
\label{sec:v27-triad-cumulative}
% =============================================================================

For the V26 kernel \(K_m\), define its cumulative current
\begin{equation}
 P_m(r):=\int_0^rK_m(\rho)\,\dd\rho.
 \label{eq:v27-triad-cumulative-kernel}
\end{equation}
Extend it by zero for \(r\ge c\).  The zero-mass identity
\eqref{eq:v26-kernel-zero-mass} makes this extension continuous.

\begin{theorem}[Exact cumulative-kernel table]
\label{thm:v27-cumulative-kernel-table}
With \(0<a<b<c\), the four cumulative currents are
\begin{align}
 P_{\rm h}(r)
 &=\begin{cases}
  0,&0\le r\le a,\\
  -(c-b)(r-a),&a\le r\le b,\\
  -(b-a)(c-r),&b\le r\le c,
 \end{cases}
 \label{eq:v27-Ph}\\
 P_1(r)
 &=\begin{cases}
  -2(c-b)r,&0\le r\le a,\\
  -(c-b)(r+a),&a\le r\le b,\\
  -(a+b)(c-r),&b\le r\le c,
 \end{cases}
 \label{eq:v27-P1}\\
 P_2(r)
 &=\begin{cases}
  2(c-a)r,&0\le r\le a,\\
  2a(c-a)+(c-b-2a)(r-a),&a\le r\le b,\\
  (a+b)(c-r),&b\le r\le c,
 \end{cases}
 \label{eq:v27-P2}\\
 P_3(r)
 &=\begin{cases}
  -2(b-a)r,&0\le r\le a,\\
  -2a(b-a)+(c-b+2a)(r-a),&a\le r\le b,\\
  (b-a)(c-r),&b\le r\le c.
 \end{cases}
 \label{eq:v27-P3}
\end{align}
Their signs are
\begin{equation}
 P_{\rm h}\le0,
 \qquad P_1\le0,
 \qquad P_2\ge0,
 \label{eq:v27-cumulative-fixed-signs}
\end{equation}
whereas \(P_3\) has one interior zero
\begin{equation}
 r_*:=a+\frac{2a(b-a)}{c-b+2a}\in(a,b),
 \label{eq:v27-P3-zero}
\end{equation}
being negative below \(r_*\) and positive above it.
\end{theorem}

\begin{proof}
Integrate each constant piece in
\eqref{eq:v26-adaptive-kernel-table} and match continuously at \(a\) and
\(b\).  This gives the four displayed formulae.  The signs of
\(P_{\rm h}\) and \(P_1\) are immediate.  For \(P_2\), the middle piece is
linear; its endpoint values are
\[
 P_2(a)=2a(c-a)>0,
 \qquad P_2(b)=(a+b)(c-b)>0,
\]
so it remains positive.  The middle piece of \(P_3\) starts at
\(-2a(b-a)<0\), ends at \((b-a)(c-b)>0\), and has positive slope
\(c-b+2a\).  Solving for its zero gives \eqref{eq:v27-P3-zero}.
\end{proof}

\begin{theorem}[Triad Stieltjes formula and orientation of dangerous
crossings]
\label{thm:v27-triad-Stieltjes}
For a finite Galerkin selector, or more generally whenever
\(\sigma_\rho\) has bounded variation on \([0,c]\),
\begin{equation}
 \boxed{\quad
 \mathfrak a_m^{\triangle}
 =-\Xi\int_{[0,c]}P_m(r)\,\dd\sigma_r.
 \quad}
 \label{eq:v27-triad-Stieltjes}
\end{equation}
For a productive mixed triad the normal-form signs are
\begin{equation}
 \Xi>0\quad(m=1),
 \qquad \Xi<0\quad(m=2),
 \qquad \Xi>0\quad(m=3).
 \label{eq:v27-productive-Xi-signs}
\end{equation}
Consequently:
\begin{enumerate}[label=\textup{(\roman*)},leftmargin=3.7em]
\item upward tail-order changes \(\dd\sigma>0\) are the dangerous, positive
      contributions to \(\mathfrak a^{\triangle}\) for both productive
      \(m=1\) and productive \(m=2\) triads;
\item for productive \(m=3\), upward changes below \(r_*\) and downward
      changes above \(r_*\) are dangerous;
\item homochiral crossing contributions have no phase-independent sign.
\end{enumerate}
Here ``dangerous'' means positive adaptive correction, which is subtracted as
\(-\mathfrak a^{\triangle}/2\) from the common sector work in
\eqref{eq:v26-triad-meet-work}.
\end{theorem}

\begin{proof}
The kernel representation \eqref{eq:v26-kernel-representation} is
\(\mathfrak a_m^{\triangle}=\Xi\int K_m\sigma\).  Since
\(P_m'=K_m\) almost everywhere and \(P_m(0)=P_m(c)=0\), Stieltjes
integration by parts gives \eqref{eq:v27-triad-Stieltjes}.  The signs in
\eqref{eq:v27-productive-Xi-signs} follow from the three rows of
\eqref{eq:v22-minority-critical-table}.  Combine them with
\eqref{eq:v27-cumulative-fixed-signs} and the sign change of \(P_3\).  For
example, when \(m=1\), an upward increment has
\(-\Xi P_1\,\dd\sigma\ge0\); when \(m=2\), both \(-\Xi\) and \(P_2\) are
positive.  The other cases follow identically.  Homochiral critical work is
zero and phase reversal changes \(\Xi\), so no fixed sign remains.
\end{proof}

\begin{corollary}[Radially nonincreasing tail order makes low-minority work
favourable]
\label{cor:v27-low-minority-monotone}
If \(\sigma_\rho\) is nonincreasing on the radius span of a productive
low-minority triad, then
\begin{equation}
 \mathfrak a_1^{\triangle}\le0,
 \qquad
 \mathfrak b_\wedge^{\triangle}
 \ge\mathcal W^{\triangle}>0.
 \label{eq:v27-low-minority-favourable}
\end{equation}
Thus its negative V26 example necessarily uses an upward tail-order change.
\end{corollary}

\begin{proof}
For a nonincreasing selector, \(\dd\sigma\le0\).  Since \(P_1\le0\) and
\(\Xi>0\), formula \eqref{eq:v27-triad-Stieltjes} gives
\(\mathfrak a_1^{\triangle}\le0\).  Insert this into
\eqref{eq:v26-triad-meet-work}.
\end{proof}

% =============================================================================
\section{Exact identification of the V24 low-minority separation profile}
\label{sec:v27-low-minority-separation}
% =============================================================================

\begin{proposition}[Low-minority danger is the separated zero-mass profile]
\label{prop:v27-low-minority-separation}
For a productive \(m=1\) triad, \(\Xi>0\), its instantaneous sector
difference is
\begin{equation}
 d_\rho^{\triangle}=\Xi K_1(\rho).
 \label{eq:v27-low-minority-difference-kernel}
\end{equation}
It is negative on \((0,b)\), positive on \((b,c)\), and has equal Jordan
masses
\begin{equation}
 \int_0^\infty(d_\rho^{\triangle})_+\,\dd\rho
 =\int_0^\infty(d_\rho^{\triangle})_-\,\dd\rho
 =(a+b)(c-b)\Xi
 =\frac{a+b}{a}\mathcal W^{\triangle}.
 \label{eq:v27-low-minority-Jordan-mass}
\end{equation}
If its low minority leg dominates the two high majority energies as in
\eqref{eq:v26-low-minority-energy-order}, then the assembled tail sign jumps
upward at \(a\), and
\begin{equation}
 P_1(a)=-2a(c-b)=-\frac{2\mathcal W^{\triangle}}{\Xi}.
 \label{eq:v27-low-minority-crossing-current}
\end{equation}
The crossing-shell formula therefore gives
\begin{equation}
 \mathfrak a_1^{\triangle}
 =-\Xi P_1(a)(+2)=4\mathcal W^{\triangle},
 \label{eq:v27-low-minority-crossing-pairing}
\end{equation}
recovering \(\mathfrak b_\wedge^{\triangle}=-\mathcal W^{\triangle}\).
Thus the negative complement is exactly the upward-crossing pairing of the
V24-type separated, zero-mass sector-difference profile.
\end{proposition}

\begin{proof}
The equality \(d^{\triangle}=\Xi K_1\) follows by subtracting the two
complete sector profiles in the first row of
Theorem~\ref{thm:v23-positive-sector-profile-table}; it is also the
integrand from which the \(m=1\) row of
\eqref{eq:v26-adaptive-kernel-table} was obtained.  That row is negative on
the first two intervals and positive on the third.  Its positive area is
\((a+b)(c-b)\Xi\), and zero signed mass gives the same negative area.  Divide
by \(\mathcal W^{\triangle}=a(c-b)\Xi\) to prove
\eqref{eq:v27-low-minority-Jordan-mass}.

Under \eqref{eq:v26-low-minority-energy-order}, the tail sign changes from
\(-1\) to \(+1\) at \(a\).  Formula \eqref{eq:v27-P1} gives
\eqref{eq:v27-low-minority-crossing-current}; insert the jump
\(\Delta\sigma=2\) into \eqref{eq:v27-triad-Stieltjes}.  This proves
\eqref{eq:v27-low-minority-crossing-pairing} and the final statement.
\end{proof}

\begin{assessment}[What V27 gains on the complement]
\label{ass:v27-complement-gain}
The low-minority obstruction is no longer merely a bad coefficient in an
\(L^1\) estimate.  It is localized to upward changes of the assembled tail
order and is sampled by the cumulative current \(P_1\).  The global version
of that current is \(\mathcal J_r\), whose normalized time integral is
Leray-paid at each fixed radius.  The unresolved step is simultaneous
crossing incidence: many radii may be sampled, and the inverse crossing
amplitudes in \eqref{eq:v27-crossing-leverage-form} prevent summation by the
shell-energy budget alone.
\end{assessment}
% =============================================================================
\section{A simultaneous Leray bound for the normalized radial packet}
\label{sec:v27-radial-BV-packet}
% =============================================================================

For \(r,\lambda>0\), put
\begin{equation}
 w_r(\lambda):=\frac{\min\{r,\lambda\}}{r}
 =\min\left\{1,\frac\lambda r\right\}.
 \label{eq:v27-hinge-weight}
\end{equation}
For fixed \(\lambda\), this is nonincreasing in \(r\), equals one below
\(\lambda\), tends to zero at infinity, and has total variation one.

\begin{theorem}[Radial bounded variation of the normalized hinge packet]
\label{thm:v27-radial-BV-packet}
At every smooth time,
\begin{align}
 z_r(t)
 &=\sum_{k,s}s,w_r(|k|)|\widehat u^s(k,t)|^2,
 \label{eq:v27-z-weight-form}\\
 q_r(t)
 &=\sum_{k,s}s,w_r(|k|)|k|^2
   |\widehat u^s(k,t)|^2.
 \label{eq:v27-q-weight-form}
\end{align}
As functions of \(r\in(0,\infty)\), they obey
\begin{align}
 \operatorname{Var}_r z_\bullet(t)
 &\le\|u(t)\|_2^2,
 \label{eq:v27-z-radial-variation}\\
 \operatorname{Var}_r q_\bullet(t)
 &\le\|\Lambda u(t)\|_2^2.
 \label{eq:v27-q-radial-variation}
\end{align}
For an interval \(I=[a,b]\), define the time-integrated normalized current
\begin{equation}
 \mathscr J_I(r):=\int_Ij_r(t)\,\dd t.
 \label{eq:v27-integrated-normalized-current}
\end{equation}
Then the \emph{whole radial packet}, simultaneously over all cutoffs, has the
Leray bound
\begin{align}
 \boxed{\quad
 \operatorname{Var}_r\mathscr J_I
 \le\frac12\bigl(\|u(a)\|_2^2+\|u(b)\|_2^2\bigr)
  +\nu\int_I\|\Lambda u(t)\|_2^2\,\dd t
 \le\frac32E_*.
 \quad}
 \label{eq:v27-integrated-current-BV}
\end{align}
The same right side bounds \(\sup_{r>0}|\mathscr J_I(r)|\).
\end{theorem}

\begin{proof}
Equations \eqref{eq:v27-z-weight-form}--\eqref{eq:v27-q-weight-form} follow
from the modal expressions in Theorem~\ref{thm:v27-hinge-current} after
division by \(r\).  For any finite partition
\(0<r_0<\cdots<r_M\), the triangle inequality gives
\begin{align*}
 \sum_{\ell=1}^M|z_{r_\ell}-z_{r_{\ell-1}}|
 &\le\sum_{k,s}|\widehat u^s(k)|^2
   \sum_{\ell=1}^M
   |w_{r_\ell}(|k|)-w_{r_{\ell-1}}(|k|)|\\
 &\le\sum_{k,s}|\widehat u^s(k)|^2.
\end{align*}
Take the supremum over partitions to prove
\eqref{eq:v27-z-radial-variation}.  Insert \(|k|^2\) in the same argument
to prove \eqref{eq:v27-q-radial-variation}.  Absolute convergence justifies
the infinite Fourier sum; finite truncation and monotone convergence give an
equivalent proof.

The normalized hinge balance gives
\begin{equation}
 \mathscr J_I(r)
 =\frac12\bigl(z_r(b)-z_r(a)\bigr)
  +\nu\int_Iq_r(t)\,\dd t.
 \label{eq:v27-integrated-current-formula}
\end{equation}
Total variation is subadditive, and variation of a time integral is bounded
by the integral of the variations.  Apply
\eqref{eq:v27-z-radial-variation}--\eqref{eq:v27-q-radial-variation} to
\eqref{eq:v27-integrated-current-formula}; the Leray energy inequality then
gives the last bound in \eqref{eq:v27-integrated-current-BV}.  Finally each
term tends to zero as \(r\to\infty\), by dominated convergence in the two
Fourier sums.  A function tending to zero at infinity has supremum bounded by
its total variation, proving the last assertion.
\end{proof}

\begin{corollary}[Endpoint-frozen sector difference samples the BV packet]
\label{cor:v27-endpoint-frozen-BV}
For a finite Galerkin solution, let \(\sigma_j^\tau\) and
\(\Delta\sigma_j^\tau\) be the radial tail signs and increments at either
endpoint \(\tau\in\{a,b\}\).  Then
\begin{align}
 \int_0^\infty\sigma_\rho^\tau D_I(\rho)\,\dd\rho
 &=-\sum_{j=1}^Nr_j\mathscr J_I(r_j)
       \Delta\sigma_j^\tau.
 \label{eq:v27-frozen-crossing-BV-pairing}
\end{align}
Here \(\mathscr J_I\) satisfies the simultaneous Leray bound
\eqref{eq:v27-integrated-current-BV}.  The remaining factors are the radius
and the unweighted endpoint sign increment.
\end{corollary}

\begin{proof}
Fubini and \eqref{eq:v27-cumulative-difference-current} give
\begin{equation}
 \int_0^rD_I(\rho)\,\dd\rho
 =\int_I\mathcal J_r(t)\,\dd t
 =r\mathscr J_I(r).
 \label{eq:v27-cumulative-integrated-difference}
\end{equation}
Apply the discrete radial summation-by-parts argument from
Theorem~\ref{thm:v27-crossing-shell-representation} with the endpoint sign
held fixed and with \(D_I\) in place of \(d(t)\).  Equation
\eqref{eq:v27-cumulative-integrated-difference} gives
\eqref{eq:v27-frozen-crossing-BV-pairing}.
\end{proof}

\begin{firewall}[Radial BV does not absorb unweighted sign jumps]
\label{fw:v27-BV-not-crossing-closure}
Bounded variation of \(\mathscr J_I\) controls its values and differences
across all radii with one Leray ledger.  It does not control the pairing in
\eqref{eq:v27-frozen-crossing-BV-pairing}, because that formula multiplies
the values by \(r_j\Delta\sigma_j^\tau\), rather than pairing
\(\dd\mathscr J_I\) with a bounded test function.  Radial integration by
parts has already placed the derivative on the sign selector.  Moving it back
recovers \(D_I\) and the V24 absolute-variation obstruction.  Thus
\eqref{eq:v27-integrated-current-BV} is a genuine simultaneous Leray estimate,
but a crossing-incidence cancellation is still required to use it.
\end{firewall}
% =============================================================================
\section{Large-amplitude crossings versus near-balanced tails}
\label{sec:v27-crossing-threshold}
% =============================================================================

\begin{proposition}[A Leray-paid crossing threshold dichotomy]
\label{prop:v27-crossing-threshold}
Fix \(\varepsilon>0\) at one finite Galerkin time and divide the crossing
shells into
\begin{align}
 \mathfrak X_{\ge\varepsilon}
 &:=\{j:\Delta\sigma_j\ne0,\ |\eta_j|\ge\varepsilon\},
 \label{eq:v27-large-crossings}\\
 \mathfrak X_{<\varepsilon}
 &:=\{j:\Delta\sigma_j\ne0,\ |\eta_j|<\varepsilon\}.
 \label{eq:v27-small-crossings}
\end{align}
Then
\begin{align}
 |\mathfrak X_{\ge\varepsilon}|
 &\le\frac{\|u(t)\|_2^2}{\varepsilon},
 \label{eq:v27-large-crossing-count}\\
 \sum_{j\in\mathfrak X_{\ge\varepsilon}}r_j
 &\le\frac{\mathcal R(t)^2}{\varepsilon}.
 \label{eq:v27-large-crossing-radius-sum}
\end{align}
For every \(j\in\mathfrak X_{<\varepsilon}\), both adjacent tail values are
near balance:
\begin{equation}
 |Y_j|<\varepsilon,
 \qquad |Y_{j+1}|<\varepsilon.
 \label{eq:v27-small-crossing-near-balance}
\end{equation}
Consequently, on either adjacent cutoff interval,
\begin{equation}
 m_\rho(t)
 =\frac12\bigl(H_\rho^+(t)+H_\rho^-(t)
               -|\mathcal Y_\rho(t)|\bigr)
 \ge\frac12\bigl(H_\rho^+(t)+H_\rho^-(t)-\varepsilon\bigr).
 \label{eq:v27-small-crossing-common-tail}
\end{equation}
\end{proposition}

\begin{proof}
Each large crossing contributes at least \(\varepsilon\) to the left side of
\eqref{eq:v27-crossing-energy-budget}; this proves
\eqref{eq:v27-large-crossing-count}.  Apply the same argument to
\eqref{eq:v27-crossing-critical-budget} to prove
\eqref{eq:v27-large-crossing-radius-sum}.  For a small crossing,
\eqref{eq:v27-crossing-amplitude} says
\(|Y_j|+|Y_{j+1}|=|\eta_j|<\varepsilon\), proving
\eqref{eq:v27-small-crossing-near-balance}.  Finally use the pointwise Jordan
identity for the minimum, as in
\eqref{eq:v25-common-stock-Jordan}, on either adjacent interval.
\end{proof}

\begin{assessment}[The remaining two-scale choice]
\label{ass:v27-crossing-threshold-status}
Large-amplitude crossings have a one-use count and radius sum, while
small-amplitude crossings occur only where the two helical tails are nearly
equal.  This is the first exact bridge between the V27 crossing-current
representation and the V25 common-tail meet.  It is not yet a closure:
\eqref{eq:v27-small-crossing-common-tail} also needs a lower bound on the
total tail height over a cutoff interval, and the large-crossing estimates
carry \(\varepsilon^{-1}\).  V28 must choose the threshold relative to the
record/front scale and keep both costs in one ledger.
\end{assessment}

% =============================================================================
\section{Integrated V27 conclusion and next acquisition order}
\label{sec:v27-conclusion}
% =============================================================================

\begin{theorem}[What V27 proves]
\label{thm:v27-conclusion}
Relative to V1--V26, the following are proved.
\begin{enumerate}[label=\textup{(V27.\arabic*)},leftmargin=6.2em]
\item Radial integration of the sector-difference flux gives the exact hinge
      current \eqref{eq:v27-hinge-current-modal}, the production of the capped
      signed-helicity stock \eqref{eq:v27-hinge-stock-modal}.
\item After normalization by the cutoff radius, every fixed hinge-current
      time integral is bounded by endpoint kinetic energy and Leray
      dissipation, as in \eqref{eq:v27-fixed-hinge-Leray-bound}.
\item Discrete radial summation by parts gives the exact crossing-shell source
      \eqref{eq:v27-crossing-shell-source}: only shells where the assembled
      tail order changes are sampled.
\item The signed shell carrying a sign jump has the opposite orientation and
      amplitude \(|Y_j|+|Y_{j+1}|\).  Crossing-shell amplitudes are one-use
      under both kinetic and critical stock at a fixed time.
\item The sign jump introduces the exact inverse-amplitude factor
      \(\omega_j\), which is unbounded under the available stock ceilings.
      Hence shell-energy H\"older estimates do not close the crossing source.
\item The cumulative triad-current table
      \eqref{eq:v27-Ph}--\eqref{eq:v27-P3} and Stieltjes formula
      \eqref{eq:v27-triad-Stieltjes} classify the dangerous orientation of
      crossings for every helical triad type.
\item The V26 negative low-minority triad is exactly the upward-crossing
      pairing of a separated zero-mass sector-difference profile, with Jordan
      mass \eqref{eq:v27-low-minority-Jordan-mass}.
\item The whole time-integrated normalized hinge packet has the simultaneous
      radial BV estimate \eqref{eq:v27-integrated-current-BV}, paid by one
      Leray ledger.
\item At any threshold, large-amplitude crossings have the one-use bounds
      \eqref{eq:v27-large-crossing-count}--%
      \eqref{eq:v27-large-crossing-radius-sum}, while small crossings force
      adjacent near-balanced helical tails.
\end{enumerate}
The exact obstruction is now a signed embedding of the Leray-paid radial BV
current packet against unweighted crossing jumps.  The full Navier--Stokes
stability/global-regularity theorem remains unproved.
\end{theorem}

\begin{proof}
Items \textup{(V27.1)--(V27.2)} are
Theorem~\ref{thm:v27-hinge-current} and
Corollary~\ref{cor:v27-fixed-hinge-Leray}.  Items
\textup{(V27.3)--(V27.5)} are
Theorem~\ref{thm:v27-crossing-shell-representation},
Lemma~\ref{lem:v27-crossing-shell-budget}, and
Proposition~\ref{prop:v27-crossing-leverage}.  Item \textup{(V27.6)} is
Theorems~\ref{thm:v27-cumulative-kernel-table} and
\ref{thm:v27-triad-Stieltjes}.  Item \textup{(V27.7)} is
Proposition~\ref{prop:v27-low-minority-separation}.  Item \textup{(V27.8)} is
Theorem~\ref{thm:v27-radial-BV-packet}.  Item \textup{(V27.9)} is
Proposition~\ref{prop:v27-crossing-threshold}.
\end{proof}

\begin{programme}[V28 crossing-incidence acquisition order]
\label{prog:v28}
\begin{enumerate}[label=\textup{(AC\arabic*)},leftmargin=4.8em]
\item Preserve the global crossing identity
      \eqref{eq:v27-crossing-shell-source}; do not estimate its hinge currents
      or sign increments separately before testing their signed incidence.
\item Use the simultaneous BV packet
      \eqref{eq:v27-integrated-current-BV}, not a separate copy of the
      fixed-cutoff bound for every crossing radius.
\item Split crossings by Proposition~\ref{prop:v27-crossing-threshold} with a
      threshold tied to the record scale.  Large crossings must be charged by
      their shell amplitudes; small crossings must be converted into common
      meet stock using an actual total-tail-height interval.
\item For the low-minority class, retain the upward orientation and the exact
      cumulative value \eqref{eq:v27-low-minority-crossing-current}.  Search
      for a pairing with downward crossings or with the robust high-minority
      class before taking positive parts.
\item For homochiral and near-local high-minority triads, apply the cumulative
      kernels \eqref{eq:v27-Ph} and \eqref{eq:v27-P2}--%
      \eqref{eq:v27-P3}; any shell-neighbor payment must be proved after
      complete incidence aggregation.
\item Test whether a clipped sign or convex tail entropy removes the inverse
      factor \(\omega_j\).  Keep the resulting near-zero-tail error explicit
      and compare it with \eqref{eq:v27-small-crossing-common-tail}.
\item Carry the V25 viscosity-weighted crossing term throughout.  The
      normalized hinge balance pays viscosity at one radius by Leray
      dissipation, but its simultaneous crossing embedding is still to be
      proved.
\item Keep the compulsory companion review scheduled for V30; V28 and V29
      should remain focused on the exact NS crossing-incidence gate.
\end{enumerate}
\end{programme}

\begin{assessment}[Position after V27]
The complement has been reduced from a signed list of triads to an assembled
radial incidence problem.  The sampled object is not raw flux: it is a
normalized capped-helicity current whose complete radial time-integrated
packet has bounded variation under the Leray ledger.  The V26 low-minority
counterexample is exactly an upward crossing of this current against the
separated sector-difference profile.

The remaining loss is equally exact.  Tail-order sign jumps are unweighted,
whereas their one-use shell budget is amplitude-weighted; the quotient can be
arbitrarily large near a balanced tail.  A threshold split turns this into a
large-crossing shell cost versus a small-crossing common-tail problem.  V28
must assemble those two branches at record scale.  The full stability/global-
regularity objective remains open.
\end{assessment}

% =============================================================================
\section*{Version V27 append-only record}
\addcontentsline{toc}{section}{Version V27 append-only record}
% =============================================================================

The complete V26 mathematical source was retained before this continuation.
Its SHA--256 digest was
\begin{center}\scriptsize\ttfamily
c3c897ce0a103057072387684d482f6740e0f0b4b68055c722fd408a55db4c86.
\end{center}
V27 derived the exact capped-helicity hinge current, proved fixed-cutoff and
simultaneous radial-BV Leray bounds after normalization, and rewrote the
adaptive source as a signed sum over assembled tail-order crossing shells.  It
proved one-use shell-amplitude budgets, isolated the unbounded inverse-
amplitude leverage, computed all cumulative triad kernels, identified the
low-minority separation mechanism as an upward crossing, and obtained a
large-crossing/near-balanced-tail threshold alternative.  The crossing-
incidence embedding remains open, so the full stability/global-regularity
goal remains open.
% =============================================================================
% VERSION V28 APPEND
% =============================================================================

\clearpage
\section{V28: clipped tail polarization and the near-balance residual}
\label{sec:v28-entry}
% =============================================================================

V27 rewrote the adaptive imbalance source as an assembled pairing between
normalized capped-helicity currents and radial jumps of the helical-tail
order.  The entire time-integrated normalized current packet has a radial BV
bound paid by one Leray ledger.  The remaining loss is that a full sign jump
can be caused by an arbitrarily small signed shell mass, producing the inverse
amplitude \(\omega_j\).

V28 regularizes exactly that discontinuity.  A Huber-clipped tail-order
selector is Lipschitz in the signed tail value, so its shell increments are
bounded by shell energy divided by the clipping threshold.  Combining this
with V27 gives a simultaneous endpoint-pairing estimate with no crossing
count and no inverse shell amplitude.  The clipped source also has an exact
convex-entropy balance.

The hard sign differs from the clipped selector only where the two helical
tails are nearly balanced.  That residual cannot be discarded: an explicit
productive low-minority triad carries an order-one fraction of its critical
work in the residual even as the tail imbalance tends to zero.  In that
example the residual is accompanied by a large common-tail meet.  Thus V28
turns the crossing-incidence gate into a precise alternative---a
Leray--BV-paid clipped pairing versus a common-tail/near-balance residual---
but does not yet make the latter cofinally one-use.

% =============================================================================
\section{Huber clipping and the exact convex tail entropy}
\label{sec:v28-Huber-entropy}
% =============================================================================

Fix an energy threshold \(\varepsilon>0\).  Define the clipped sign and its
convex primitive by
\begin{align}
 \vartheta_\varepsilon(y)
 &:=\begin{cases}
  -1,&y\le-\varepsilon,\\
  y/\varepsilon,&|y|<\varepsilon,\\
  +1,&y\ge\varepsilon,
 \end{cases}
 \label{eq:v28-clipped-sign}\\
 \Psi_\varepsilon(y)
 &:=\int_0^y\vartheta_\varepsilon(z)\,\dd z
 =\begin{cases}
  y^2/(2\varepsilon),&|y|\le\varepsilon,\\
  |y|-\varepsilon/2,&|y|\ge\varepsilon.
 \end{cases}
 \label{eq:v28-Huber-entropy}
\end{align}
Thus \(\Psi_\varepsilon\) is even, convex, and continuously differentiable,
with derivative \(\vartheta_\varepsilon\).  Set
\begin{align}
 \sigma_\rho^\varepsilon(t)
 &:=\vartheta_\varepsilon(\mathcal Y_\rho(t)),
 \label{eq:v28-clipped-tail-selector}\\
 \mathcal L_\varepsilon(t)
 &:=\int_0^\infty
   \Psi_\varepsilon(\mathcal Y_\rho(t))\,\dd\rho,
 \label{eq:v28-Huber-tail-entropy}\\
 \mathfrak a_\varepsilon(t)
 &:=\int_0^\infty
   \sigma_\rho^\varepsilon(t)d_\rho(t)\,\dd\rho,
 \label{eq:v28-clipped-adaptive-source}\\
 \mathcal A_I^\varepsilon
 &:=\int_I\mathfrak a_\varepsilon(t)\,\dd t.
 \label{eq:v28-integrated-clipped-source}
\end{align}

\begin{theorem}[Exact clipped-source entropy balance]
\label{thm:v28-clipped-entropy-balance}
On every compact smooth interval \(I=[a,b]\),
\begin{equation}
 \boxed{\quad
 \mathcal A_I^\varepsilon
 =\frac12\bigl(\mathcal L_\varepsilon(b)
                 -\mathcal L_\varepsilon(a)\bigr)
  +\nu\int_I\int_0^\infty
   \sigma_\rho^\varepsilon(t)
   \bigl(V_\rho^+(t)-V_\rho^-(t)\bigr)
   \,\dd\rho\,\dd t.
 \quad}
 \label{eq:v28-clipped-entropy-balance}
\end{equation}
The entropy is finite and satisfies
\begin{align}
 0\le\mathcal L_\varepsilon(t)
 &\le\mathcal L(t)\le\mathcal R(t)^2,
 \label{eq:v28-Huber-entropy-cap}\\
 0\le\mathcal L(t)-\mathcal L_\varepsilon(t)
 &\le\int_0^\infty
   \min\{|\mathcal Y_\rho(t)|,\varepsilon/2\}\,\dd\rho.
 \label{eq:v28-Huber-defect}
\end{align}
The viscous residual has the exact absolute upper bound
\begin{align}
 &\left|\int_I\int_0^\infty
   \sigma_\rho^\varepsilon
   (V_\rho^+-V_\rho^-)
   \,\dd\rho\,\dd t\right|
 \notag\\
 &\qquad\le
 \int_I\int_0^\infty
  \min\left\{1,\frac{|\mathcal Y_\rho|}{\varepsilon}\right\}
  (V_\rho^++V_\rho^-)
  \,\dd\rho\,\dd t.
 \label{eq:v28-clipped-viscous-bound}
\end{align}
\end{theorem}

\begin{proof}
The instantaneous difference balance \eqref{eq:v25-instant-tail-difference}
is
\[
 d_\rho=\frac12\partial_t\mathcal Y_\rho
          +\nu(V_\rho^+-V_\rho^-).
\]
Multiply by
\(\vartheta_\varepsilon(\mathcal Y_\rho)\).  The ordinary chain rule for the
\(C^1\) function \(\Psi_\varepsilon\) gives
\[
 \vartheta_\varepsilon(\mathcal Y_\rho)
 \partial_t\mathcal Y_\rho
 =\partial_t\Psi_\varepsilon(\mathcal Y_\rho).
\]
Integrate in cutoff and time.  Smooth spectral decay supplies Fubini and
dominated convergence, proving \eqref{eq:v28-clipped-entropy-balance}.

For every scalar \(y\), direct inspection of
\eqref{eq:v28-Huber-entropy} gives
\begin{equation}
 0\le\Psi_\varepsilon(y)\le|y|,
 \qquad
 0\le|y|-\Psi_\varepsilon(y)
 \le\min\{|y|,\varepsilon/2\}.
 \label{eq:v28-pointwise-Huber-bounds}
\end{equation}
Integrate and use \(\mathcal L\le\mathcal R^2\) from
\eqref{eq:v25-common-stock-Jordan}; this proves
\eqref{eq:v28-Huber-entropy-cap}--\eqref{eq:v28-Huber-defect}.  Finally
\[
 |\sigma_\rho^\varepsilon|
 =\min\left\{1,\frac{|\mathcal Y_\rho|}{\varepsilon}\right\},
 \qquad
 |V_\rho^+-V_\rho^-|\le V_\rho^++V_\rho^-,
\]
which proves \eqref{eq:v28-clipped-viscous-bound}.
\end{proof}

\begin{firewall}[The clipped viscous action is still not Leray by itself]
\label{fw:v28-clipped-viscosity}
The weight in \eqref{eq:v28-clipped-viscous-bound} improves the critical
action near a balanced tail, but its radial integral still contains
\(\int(V_\rho^++V_\rho^-)\,\dd\rho
=\|\Lambda^{3/2}u\|_2^2\).  No Leray estimate follows merely from
\(|\sigma^\varepsilon|\le1\).  The Leray gain comes after radial summation by
parts and hinge normalization, as in the next section.
\end{firewall}
% =============================================================================
\section{Clipped crossing increments and the Leray--BV endpoint estimate}
\label{sec:v28-clipped-crossings}
% =============================================================================

For a finite Galerkin field with the V27 occupied radii, define
\begin{equation}
 \Delta\sigma_j^\varepsilon(t)
 :=\vartheta_\varepsilon(Y_{j+1}(t))
   -\vartheta_\varepsilon(Y_j(t)).
 \label{eq:v28-clipped-sign-increment}
\end{equation}

\begin{lemma}[Clipping removes inverse shell amplitude]
\label{lem:v28-clipped-shell-Lipschitz}
For every shell,
\begin{equation}
 \boxed{\quad
 |\Delta\sigma_j^\varepsilon(t)|
 \le\frac{|\eta_j(t)|}{\varepsilon}.
 \quad}
 \label{eq:v28-clipped-shell-increment}
\end{equation}
Consequently,
\begin{align}
 \sum_j|\Delta\sigma_j^\varepsilon|&\,\varepsilon
 \le\|u(t)\|_2^2,
 \label{eq:v28-clipped-increment-energy}\\
 \sum_jr_j|\Delta\sigma_j^\varepsilon|&\,\varepsilon
 \le\mathcal R(t)^2.
 \label{eq:v28-clipped-increment-critical}
\end{align}
Unlike \(\omega_j=|\Delta\sigma_j|/|\eta_j|\), the Lipschitz constant is the
fixed threshold \(\varepsilon^{-1}\), independent of the crossing shell.
\end{lemma}

\begin{proof}
The scalar map \(\vartheta_\varepsilon\) is
\(1/\varepsilon\)-Lipschitz.  Since
\(Y_{j+1}-Y_j=-\eta_j\), this proves
\eqref{eq:v28-clipped-shell-increment}.  Sum it with weights one and \(r_j\),
then use the full shell-energy and critical-stock bounds from the proof of
Lemma~\ref{lem:v27-crossing-shell-budget}.  The sums may include all shells,
because \(|\eta_j|\le e_j^++e_j^-\) independently of whether the hard sign
changes.
\end{proof}

\begin{theorem}[Exact clipped crossing representation]
\label{thm:v28-clipped-crossing-representation}
For every finite Galerkin field,
\begin{equation}
 \boxed{\quad
 \mathfrak a_\varepsilon(t)
 =-\sum_{j=1}^N
   \mathcal J_{r_j}(t)\Delta\sigma_j^\varepsilon(t)
 =-\sum_{j=1}^N
   r_j\,j_{r_j}(t)\Delta\sigma_j^\varepsilon(t).
 \quad}
 \label{eq:v28-clipped-crossing-source}
\end{equation}
In particular,
\begin{equation}
 |\mathfrak a_\varepsilon(t)|
 \le\frac{\mathcal R(t)^2}{\varepsilon}
       \sup_{r>0}|j_r(t)|.
 \label{eq:v28-instant-clipped-source-bound}
\end{equation}
\end{theorem}

\begin{proof}
Repeat the discrete summation-by-parts proof of
Theorem~\ref{thm:v27-crossing-shell-representation} with
\(\vartheta_\varepsilon(Y_j)\) in place of \(\sigma_j\).  The cumulative
current still vanishes at the top radius, yielding
\eqref{eq:v28-clipped-crossing-source}.  Apply
\eqref{eq:v28-clipped-shell-increment}, write
\(\mathcal J_{r_j}=r_j\,j_{r_j}\), and use
\(\sum_jr_j|\eta_j|\le\mathcal R^2\) to prove
\eqref{eq:v28-instant-clipped-source-bound}.
\end{proof}

The instantaneous supremum in
\eqref{eq:v28-instant-clipped-source-bound} is not Leray-controlled.  The
time-integrated V27 packet is.  Put
\begin{equation}
 \mathscr E_I
 :=\frac12\bigl(\|u(a)\|_2^2+\|u(b)\|_2^2\bigr)
  +\nu\int_I\|\Lambda u(t)\|_2^2\,\dd t.
 \label{eq:v28-Leray-packet-budget}
\end{equation}
Then \(\mathscr E_I\le3E_*/2\) and V27 gives
\begin{equation}
 \sup_{r>0}|\mathscr J_I(r)|
 \le\operatorname{Var}_r\mathscr J_I
 \le\mathscr E_I.
 \label{eq:v28-integrated-current-budget}
\end{equation}

For either endpoint \(\tau\in\{a,b\}\), define the clipped frozen pairing
\begin{equation}
 \mathcal Q_{I,\varepsilon}^{\tau}
 :=\int_0^\infty
  \vartheta_\varepsilon(\mathcal Y_\rho(\tau))
  D_I(\rho)\,\dd\rho.
 \label{eq:v28-clipped-frozen-pairing}
\end{equation}

\begin{theorem}[Simultaneous Leray--BV bound for the clipped endpoint
pairing]
\label{thm:v28-clipped-endpoint-bound}
For a finite Galerkin solution,
\begin{align}
 \mathcal Q_{I,\varepsilon}^{\tau}
 &=-\sum_{j=1}^Nr_j\mathscr J_I(r_j)
   \Delta\sigma_j^\varepsilon(\tau),
 \label{eq:v28-clipped-frozen-sum}\\
 \boxed{\quad
 |\mathcal Q_{I,\varepsilon}^{\tau}|
 \le\frac{\mathscr E_I}{\varepsilon}\mathcal R(\tau)^2
 \le\frac{3E_*}{2\varepsilon}\mathcal R(\tau)^2.
 \quad}
 \label{eq:v28-clipped-frozen-bound}
\end{align}
The same result holds for a smooth solution by spectral truncation.
\end{theorem}

\begin{proof}
Apply Corollary~\ref{cor:v27-endpoint-frozen-BV} with the clipped endpoint
selector.  This proves \eqref{eq:v28-clipped-frozen-sum}.  Then
\begin{align*}
 |\mathcal Q_{I,\varepsilon}^{\tau}|
 &\le\sup_r|\mathscr J_I(r)|
      \sum_jr_j|\Delta\sigma_j^\varepsilon(\tau)|\\
 &\le\frac{\mathscr E_I}{\varepsilon}
      \sum_jr_j|\eta_j(\tau)|
 \le\frac{\mathscr E_I}{\varepsilon}\mathcal R(\tau)^2,
\end{align*}
by \eqref{eq:v28-clipped-shell-increment},
\eqref{eq:v28-integrated-current-budget}, and the shell critical-stock
budget.  For a smooth field, Galerkin tails converge uniformly in \(\rho\)
at a fixed time, the clipped map is Lipschitz, and the smooth flux sums are
absolutely convergent.  Pass to the limit by dominated convergence.
\end{proof}

\begin{corollary}[The V26 robust triad class survives clipping]
\label{cor:v28-robust-class-clipped}
Replace the hard selector \(\sigma_\rho\) by
\(\sigma_\rho^\varepsilon\) in the V26 triad source.  Since
\(|\sigma_\rho^\varepsilon|\le1\), every conclusion of
Theorem~\ref{thm:v26-nonlocal-high-minority-positive}, including the fraction
\((1-2\delta)/(1+\delta)\), remains valid.
\end{corollary}

\begin{proof}
The proof of Theorem~\ref{thm:v26-nonlocal-high-minority-positive} uses only
the exact kernel table and the pointwise bound on the selector.  The clipped
selector has the same bound.
\end{proof}

\begin{firewall}[The endpoint estimate is record-scale, not a contradiction]
\label{fw:v28-endpoint-bound-scale}
At a record endpoint, \(\mathcal R(\tau)^2\) is itself the growing critical
stock.  Taking \(\varepsilon\) comparable to \(E_*\) in
\eqref{eq:v28-clipped-frozen-bound} therefore gives an order-
\(\mathcal R(\tau)^2\) upper bound.  This removes the crossing count and the
inverse shell amplitudes, but it does not contradict the order-
\(R^2\) separation lower bound.  Choosing a larger clipping threshold reduces
the displayed upper bound only by moving more of the hard sign into the
near-balance residual treated next.
\end{firewall}
% =============================================================================
\section{Exact hard-sign residual and its common-tail localization}
\label{sec:v28-near-balance-residual}
% =============================================================================

Define the hard-minus-clipped scalar residual
\begin{equation}
 R_\varepsilon(y)
 :=\operatorname{sgn}(y)-\vartheta_\varepsilon(y).
 \label{eq:v28-sign-residual}
\end{equation}
It obeys
\begin{equation}
 R_\varepsilon(y)=0\quad\text{if }y=0
   \text{ or }|y|\ge\varepsilon,
 \qquad
 |R_\varepsilon(y)|=1-\frac{|y|}{\varepsilon}le1
 \quad(0<|y|<\varepsilon).
 \label{eq:v28-sign-residual-support}
\end{equation}

Set
\begin{align}
 \mathcal N_I^\varepsilon
 &:=\int_I\int_0^\infty
   R_\varepsilon(\mathcal Y_\rho(t))d_\rho(t)
   \,\dd\rho\,\dd t,
 \label{eq:v28-dynamic-near-balance-residual}\\
 \mathcal N_{I,\varepsilon}^{\tau}
 &:=\int_0^\infty
   R_\varepsilon(\mathcal Y_\rho(\tau))D_I(\rho)
   \,\dd\rho,
 \qquad \tau\in\{a,b\}.
 \label{eq:v28-endpoint-near-balance-residual}
\end{align}

\begin{theorem}[Exact clipped-versus-hard ledgers]
\label{thm:v28-hard-clipped-ledgers}
The adaptive source and both endpoint-frozen pairings split exactly as
\begin{align}
 \mathcal A_I
 &=\mathcal A_I^\varepsilon+\mathcal N_I^\varepsilon,
 \label{eq:v28-hard-clipped-source-split}\\
 \int_0^\infty\sigma_\rho^\tau D_I(\rho)\,\dd\rho
 &=\mathcal Q_{I,\varepsilon}^{\tau}
   +\mathcal N_{I,\varepsilon}^{\tau}.
 \label{eq:v28-hard-clipped-endpoint-split}
\end{align}
Every residual integrand is supported on the near-balance set
\begin{equation}
 \mathfrak N_\varepsilon(t)
 :=\{\rho:0<|\mathcal Y_\rho(t)|<\varepsilon\}.
 \label{eq:v28-near-balance-set}
\end{equation}
Consequently,
\begin{align}
 |\mathcal N_I^\varepsilon|
 &\le\int_I\int_{\mathfrak N_\varepsilon(t)}
   |d_\rho(t)|\,\dd\rho\,\dd t,
 \label{eq:v28-dynamic-residual-upper}\\
 |\mathcal N_{I,\varepsilon}^{\tau}|
 &\le\int_{\mathfrak N_\varepsilon(\tau)}
   |D_I(\rho)|\,\dd\rho.
 \label{eq:v28-endpoint-residual-upper}
\end{align}
\end{theorem}

\begin{proof}
The scalar identity
\(\operatorname{sgn}y=\vartheta_\varepsilon(y)+R_\varepsilon(y)\), applied
inside the definitions of \(\mathcal A_I\) and the endpoint pairing, proves
the two exact splits.  Equation~\eqref{eq:v28-sign-residual-support} localizes
their supports and bounds the residual multiplier by one, giving the last two
inequalities.
\end{proof}

The near-balance set becomes common-tail stock only where the total tail is
not simultaneously small.  For \(h>\varepsilon\), define
\begin{equation}
 \mathfrak N_{\varepsilon,h}(t)
 :=\left\{\rho:
 0<|\mathcal Y_\rho(t)|<\varepsilon,
 \ H_\rho^+(t)+H_\rho^-(t)\ge h
 \right\}.
 \label{eq:v28-high-total-near-balance-set}
\end{equation}

\begin{proposition}[High-total near balance is paid by common-tail stock]
\label{prop:v28-near-balance-common-stock}
On \(\mathfrak N_{\varepsilon,h}(t)\),
\begin{equation}
 m_\rho(t)\ge\frac{h-\varepsilon}{2}.
 \label{eq:v28-near-balance-meet-density}
\end{equation}
Therefore
\begin{equation}
 |\mathfrak N_{\varepsilon,h}(t)|
 \le\frac{2\mathcal C(t)}{h-\varepsilon}.
 \label{eq:v28-near-balance-width}
\end{equation}
If \(I=[a,b]\) is a first record doubling, then for either endpoint
\(\tau\in\{a,b\}\),
\begin{equation}
 \int_{\mathfrak N_{\varepsilon,h}(\tau)}
  |R_\varepsilon(\mathcal Y_\rho(\tau))D_I(\rho)|\,\dd\rho
 \le\frac{3E_*\mathcal C(\tau)}{h-\varepsilon}.
 \label{eq:v28-high-total-residual-common-bound}
\end{equation}
\end{proposition}

\begin{proof}
The pointwise Jordan identity gives
\[
 m_\rho
 =\frac12\bigl(H_\rho^++H_\rho^--|\mathcal Y_\rho|\bigr).
\]
The two defining inequalities for
\(\mathfrak N_{\varepsilon,h}\) prove
\eqref{eq:v28-near-balance-meet-density}.  Integrate it over the set and use
\(\int m_\rho=\mathcal C\) to prove
\eqref{eq:v28-near-balance-width}.  On a first doubling, the V24 height cap
\eqref{eq:v24-difference-height-cap} gives
\(|D_I(\rho)|\le3E_*/2\).  Multiply that cap by
\eqref{eq:v28-near-balance-width} and use \(|R_\varepsilon|\le1\) to obtain
\eqref{eq:v28-high-total-residual-common-bound}.
\end{proof}

\begin{corollary}[Exact residual trichotomy at an endpoint]
\label{cor:v28-endpoint-residual-trichotomy}
For \(h>\varepsilon\), the hard endpoint pairing consists of:
\begin{enumerate}[label=\textup{(\roman*)},leftmargin=3.7em]
\item the clipped pairing, bounded by
      \eqref{eq:v28-clipped-frozen-bound};
\item a high-total near-balance residual, bounded by
      \eqref{eq:v28-high-total-residual-common-bound}; and
\item the remaining low-total residual
\begin{equation}
 \int_{\substack{0<|\mathcal Y_\rho(\tau)|<\varepsilon\\
 H_\rho^+(\tau)+H_\rho^-(\tau)<h}}
 R_\varepsilon(\mathcal Y_\rho(\tau))D_I(\rho)\,\dd\rho,
 \label{eq:v28-low-total-residual}
\end{equation}
for which no Leray upper bound is asserted.
\end{enumerate}
No part of the hard pairing is omitted.
\end{corollary}

\begin{proof}
Use \eqref{eq:v28-hard-clipped-endpoint-split} and partition its residual set
according to whether the total terminal tail is at least \(h\).  Apply the two
stated estimates to the first two pieces.
\end{proof}

\begin{firewall}[Common stock in the residual is not yet cofinally one-use]
\label{fw:v28-residual-common-not-cofinal}
Equation~\eqref{eq:v28-high-total-residual-common-bound} is a same-terminal
bound.  V25 proved that \(\mathcal C(t)\) can be recreated by bottleneck-sector
nonlinear work, and record endpoints occur at different times.  Summing this
bound over cofinal records would therefore reuse a critical stock unless an
ancestry or dissipation payment is added.  The low-total residual
\eqref{eq:v28-low-total-residual} is separate and must also remain in the
ledger.
\end{firewall}
% =============================================================================
\section{A sharp near-balance test on the productive low-minority triad}
\label{sec:v28-near-balance-test}
% =============================================================================

\begin{proposition}[Clipping leaves an order-one near-balance residual]
\label{prop:v28-clipped-low-minority-test}
Fix \(\varepsilon>0\) and an active \(m=1\) triad with radii
\(0<a<b<c\).  Choose high-leg energies \(e_2,e_3\ge\varepsilon\), put
\begin{equation}
 e_1=e_2+e_3+\delta,
 \qquad 0<\delta<\varepsilon,
 \label{eq:v28-almost-balanced-low-minority}
\end{equation}
and choose the phase so that \(\Xi>0\).  Let
\(\mathcal W^{\triangle}=a(c-b)\Xi\) be its positive common sector work.
Then its hard and clipped selectors on the three radial intervals are
\begin{equation}
 (\sigma_1,\sigma_2,\sigma_3)=(-1,+1,+1),
 \qquad
 (\sigma_1^\varepsilon,\sigma_2^\varepsilon,
  \sigma_3^\varepsilon)=(-\delta/\varepsilon,+1,+1).
 \label{eq:v28-hard-clipped-triad-signs}
\end{equation}
The corresponding adaptive sources are
\begin{align}
 \mathfrak a_1^{\triangle}
 &=4\mathcal W^{\triangle},
 \label{eq:v28-hard-low-minority-source}\\
 \mathfrak a_{1,\varepsilon}^{\triangle}
 &=2\left(1+\frac\delta\varepsilon\right)
    \mathcal W^{\triangle},
 \label{eq:v28-clipped-low-minority-source}\\
 \boxed{\quad
 \mathfrak a_1^{\triangle}
  -\mathfrak a_{1,\varepsilon}^{\triangle}
 =2\left(1-\frac\delta\varepsilon\right)
   \mathcal W^{\triangle}.
 \quad}
 \label{eq:v28-low-minority-residual-size}
\end{align}
Thus, as \(\delta/\varepsilon\to0\), the tail imbalance on \((0,a)\)
tends to zero while the clipped-to-hard residual tends to
\(2\mathcal W^{\triangle}\).  The clipped and hard smaller-tail triad works
are respectively
\begin{align}
 \mathfrak b_{\wedge,\varepsilon}^{\triangle}
 &:=\mathcal W^{\triangle}
 -\frac12\mathfrak a_{1,\varepsilon}^{\triangle}
 =-\frac\delta\varepsilon\mathcal W^{\triangle},
 \label{eq:v28-clipped-triad-meet-work}\\
 \mathfrak b_\wedge^{\triangle}
 &=-\mathcal W^{\triangle}.
 \label{eq:v28-hard-triad-meet-work}
\end{align}
Meanwhile the common-tail stock of this triad is
\begin{equation}
 \mathcal C^{\triangle}=a(e_2+e_3).
 \label{eq:v28-low-minority-common-stock}
\end{equation}
Hence the nonvanishing hard-sign residual occurs exactly on a radial interval
carrying common helical-tail stock.
\end{proposition}

\begin{proof}
Below \(a\), the signed tail is
\(e_2+e_3-e_1=-\delta\).  Between \(a\) and \(b\), it is
\(e_2+e_3\ge2\varepsilon\), and between \(b\) and \(c\), it is
\(e_3\ge\varepsilon\).  This proves
\eqref{eq:v28-hard-clipped-triad-signs}.

The hard identity is Proposition~\ref{prop:v26-low-minority-negative}.  For
the clipped identity, insert the three clipped constants into the \(m=1\)
kernel row of \eqref{eq:v26-adaptive-kernel-table}:
\begin{align*}
 \frac{\mathfrak a_{1,\varepsilon}^{\triangle}}{\Xi}
 &=2a(c-b)\frac\delta\varepsilon
   -(b-a)(c-b)+(a+b)(c-b)\\
 &=2a(c-b)\left(1+\frac\delta\varepsilon\right).
\end{align*}
Since \(\mathcal W^{\triangle}=a(c-b)\Xi\), this proves
\eqref{eq:v28-clipped-low-minority-source}.  Subtraction proves
\eqref{eq:v28-low-minority-residual-size}, and the two smaller-tail work
formulae follow from their definitions.

For \(0<\rho<a\), the positive-sector tail has energy \(e_2+e_3\), while
the negative-sector tail has energy \(e_1>e_2+e_3\).  Their minimum is
\(e_2+e_3\).  Above \(a\), the negative tail is zero.  Integration gives
\eqref{eq:v28-low-minority-common-stock}.
\end{proof}

\begin{corollary}[Near-zero imbalance does not make the source residual
small]
\label{cor:v28-no-amplitude-smallness}
There is no function \(\zeta(x)\to0\) as \(x\downarrow0\) such that every
productive low-minority triad satisfies
\begin{equation}
 \left|\mathfrak a_1^{\triangle}
       -\mathfrak a_{1,\varepsilon}^{\triangle}\right|
 \le\zeta(\delta/\varepsilon)\mathcal W^{\triangle}.
 \label{eq:v28-no-residual-smallness}
\end{equation}
Any valid small-imbalance estimate must also use the common-tail stock, radial
width, temporal ancestry, or another PDE incidence quantity.
\end{corollary}

\begin{proof}
Keep \(e_2,e_3\), the active geometry, and the productive phase fixed while
letting \(\delta/\varepsilon\to0\).  The active interaction scalar remains
nonzero and varies continuously with the nonzero modal amplitudes, whereas
the ratio of \eqref{eq:v28-low-minority-residual-size} to
\(\mathcal W^{\triangle}\) tends to two.
\end{proof}

\begin{assessment}[What clipping proves and what the sharp test forbids]
\label{ass:v28-clipping-assessment}
Clipping succeeds at its designated task: it replaces every shell-dependent
inverse amplitude by the single threshold \(\varepsilon^{-1}\), allowing the
whole endpoint pairing to use V27's one radial BV ledger.  The sharp triad
test proves that the discarded hard-sign part is not perturbative in the tail
imbalance.  It is instead coupled to simultaneous chiral tail mass.  Thus the
next acquisition cannot be a better pointwise approximation of
\(\operatorname{sgn}\); it must give the common-tail residual temporal
ancestry or assemble it with an oppositely oriented crossing class.
\end{assessment}
% =============================================================================
\section{Integrated V28 conclusion and next acquisition order}
\label{sec:v28-conclusion}
% =============================================================================

\begin{theorem}[What V28 proves]
\label{thm:v28-conclusion}
Relative to V1--V27, the following are proved.
\begin{enumerate}[label=\textup{(V28.\arabic*)},leftmargin=6.2em]
\item The Huber-clipped tail selector has the exact convex-entropy balance
      \eqref{eq:v28-clipped-entropy-balance}, with entropy and defect bounds
      \eqref{eq:v28-Huber-entropy-cap}--\eqref{eq:v28-Huber-defect}.
\item Its radial shell increments satisfy the amplitude-sensitive Lipschitz
      estimate \eqref{eq:v28-clipped-shell-increment}.  The V27 inverse
      crossing amplitudes are replaced by one fixed factor
      \(\varepsilon^{-1}\).
\item The clipped adaptive source has the exact crossing representation
      \eqref{eq:v28-clipped-crossing-source}; its endpoint-frozen pairing is
      bounded simultaneously over all radii by the Leray--BV estimate
      \eqref{eq:v28-clipped-frozen-bound}.
\item The V26 robust productive high-minority class remains positive with the
      same quantitative fraction after clipping.
\item The hard adaptive source and endpoint pairings split exactly into their
      clipped parts and residuals supported only where
      \(0<|\mathcal Y_\rho|<\varepsilon\), as in
      \eqref{eq:v28-hard-clipped-source-split}--%
      \eqref{eq:v28-hard-clipped-endpoint-split}.
\item Where the total tail is at least \(h>\varepsilon\), the near-balance
      residual set has the common-stock width bound
      \eqref{eq:v28-near-balance-width}, giving the record-corridor estimate
      \eqref{eq:v28-high-total-residual-common-bound}.
\item The full endpoint residual is the exact trichotomy of
      Corollary~\ref{cor:v28-endpoint-residual-trichotomy}; its low-total part
      is retained without an unproved estimate.
\item A productive low-minority triad with imbalance
      \(\delta/\varepsilon\to0\) has hard-minus-clipped source tending to
      \(2\mathcal W^{\triangle}\).  Near-zero imbalance alone therefore does
      not make the residual perturbative; the same interval carries the
      common stock \eqref{eq:v28-low-minority-common-stock}.
\end{enumerate}
Clipping closes the inverse-amplitude defect for the regularized pairing but
leaves a quantitatively identified common-tail/low-total residual.  The full
Navier--Stokes stability/global-regularity theorem remains unproved.
\end{theorem}

\begin{proof}
Item \textup{(V28.1)} is
Theorem~\ref{thm:v28-clipped-entropy-balance}.  Items
\textup{(V28.2)--(V28.3)} are
Lemma~\ref{lem:v28-clipped-shell-Lipschitz},
Theorem~\ref{thm:v28-clipped-crossing-representation}, and
Theorem~\ref{thm:v28-clipped-endpoint-bound}.  Item \textup{(V28.4)} is
Corollary~\ref{cor:v28-robust-class-clipped}.  Item \textup{(V28.5)} is
Theorem~\ref{thm:v28-hard-clipped-ledgers}.  Item \textup{(V28.6)} is
Proposition~\ref{prop:v28-near-balance-common-stock}.  Item
\textup{(V28.7)} is Corollary~\ref{cor:v28-endpoint-residual-trichotomy}.
Item \textup{(V28.8)} is
Proposition~\ref{prop:v28-clipped-low-minority-test} and
Corollary~\ref{cor:v28-no-amplitude-smallness}.
\end{proof}

\begin{programme}[V29 near-balance ancestry acquisition order]
\label{prog:v29}
\begin{enumerate}[label=\textup{(AC\arabic*)},leftmargin=4.8em]
\item Keep the hard/clipped decompositions
      \eqref{eq:v28-hard-clipped-source-split}--%
      \eqref{eq:v28-hard-clipped-endpoint-split} exact.  Do not discard the
      near-balance residual after bounding the clipped term.
\item Combine the dynamic residual
      \eqref{eq:v28-dynamic-near-balance-residual} with the V25 meet balance.
      On high-total near-balance regions, test whether persistence in time
      forces either common-stock ancestry or bottleneck dissipation.
\item Attack the low-total residual \eqref{eq:v28-low-total-residual} from the
      exact sector high-tail balances.  Determine whether low terminal tail
      height forces the time-integrated difference to come from initial stock
      or viscosity, rather than gross nonlinear variation.
\item Choose \(\varepsilon\) and \(h\) relative to \(E_*\), the record
      \(R\), and the front scale \(K=R^2/E_*\).  Track all constants against
      the V24 separation lower bound; an order-
      \(R^2\) upper bound with a larger constant is not closure.
\item Use the exact low-minority test
      \eqref{eq:v28-low-minority-residual-size} as a mandatory obstruction.
      Any proposed residual estimate must pay its common stock
      \eqref{eq:v28-low-minority-common-stock} without reusing it at later
      record terminals.
\item Preserve the clipped viscous residual
      \eqref{eq:v28-clipped-viscous-bound}.  Seek a radial hinge
      summation which uses \eqref{eq:v27-integrated-current-BV}; do not replace
      it by full critical dissipation.
\item Retain homochiral, near-local high-minority, low-minority, and
      nonproductive contributions until their complete signed incidence is
      controlled.
\item V30 must begin with the compulsory review of the then-current SM and YM
      notes before importing any new mechanism.
\end{enumerate}
\end{programme}

\begin{assessment}[Position after V28]
The discontinuous crossing leverage has been separated cleanly.  The clipped
part is controlled through all radii by a single Leray--BV packet and a fixed
threshold.  Everything lost in clipping is localized to nearly equal helical
tails; where their total height is substantial, the loss is accompanied by
common-tail meet stock.

The sharp low-minority test prevents a false perturbative conclusion: the
near-balance residual can remain of critical-work size as the imbalance tends
to zero.  The next step is therefore temporal, not pointwise.  V29 must prove
ancestry or a Leray debit for the common-tail residual and must separately
resolve the low-total tail region.  The full stability/global-regularity
objective remains open.
\end{assessment}

% =============================================================================
\section*{Version V28 append-only record}
\addcontentsline{toc}{section}{Version V28 append-only record}
% =============================================================================

The complete V27 mathematical source was retained before this continuation.
Its SHA--256 digest was
\begin{center}\scriptsize\ttfamily
1566207dd59521e30fff08443554955f18614e67cb7f0ff3113b060c73eb2ab8.
\end{center}
V28 introduced a Huber-clipped tail-order selector, proved its exact convex-
entropy balance, removed the shellwise inverse-amplitude loss, and bounded the
clipped endpoint pairing by the V27 simultaneous Leray--BV packet.  It retained
the exact hard-sign residual, localized it to near-balanced tails, bounded its
high-total part by common-tail stock, and exhibited a productive low-minority
triad whose residual stays at critical-work size as its tail imbalance tends
to zero.  Common-tail ancestry and the low-total residual remain open, so the
full stability/global-regularity goal remains open.
% =============================================================================
% VERSION V29 APPEND
% =============================================================================

\clearpage
\section{V29: soft tail meet and Leray-paid radius--time occupation}
\label{sec:v29-entry}
% =============================================================================

V28 removed the inverse crossing amplitude from the clipped endpoint pairing,
but its sharp low-minority test showed that the hard-minus-clipped residual is
not small merely because the two helical tails are nearly equal.  The residual
must be attached either to simultaneous tail stock through time or to a recent
source event.  This continuation makes both alternatives exact.

First, the Huber entropy is not only a regularized imbalance: subtracting it
from the total tail produces a smooth ``soft meet.''  Its balance is precisely
the V28 clipped-source balance recombined with the sum of the two sector
balances.  The excess of this soft meet over the genuine minimum is exactly
one half of the clipping defect.  This identifies, rather than hides, every
artificial unit of stock introduced by regularization.

Second, the spectral layer cake with one cutoff-radius weight converts tail
occupation into the ordinary Leray dissipation.  Consequently a high-total,
nearly balanced cutoff band cannot persist longer than a parabolic time unless
it spends a globally one-use Leray debit.  If a terminal band has not persisted,
then a two-threshold argument forces a recent total-tail influx or a recent
polarization action.  An analogous argument applies without a height threshold
to the low-total residual.  These are rigorous ancestry alternatives, but the
recent nonlinear actions still require an assembled cofinal estimate; hence no
global-regularity conclusion is asserted.

% =============================================================================
\section{The clipped selector is the gradient of an exact soft meet}
\label{sec:v29-soft-meet}
% =============================================================================

Write the total helical tail as
\begin{equation}
 H_\rho^\Sigma(t):=H_\rho^+(t)+H_\rho^-(t).
 \label{eq:v29-total-tail}
\end{equation}
For the V28 Huber function define
\begin{align}
 m_\rho^\varepsilon(t)
 &:=\frac12\left(
       H_\rho^\Sigma(t)
       -\Psi_\varepsilon(\mathcal Y_\rho(t))\right),
 \label{eq:v29-soft-meet-density}\\
 \mathcal C_\varepsilon(t)
 &:=\int_0^\infty m_\rho^\varepsilon(t)\,\dd\rho,
 \label{eq:v29-soft-meet-stock}\\
 \chi_\rho^{+,\varepsilon}(t)
 &:=\frac{1-\sigma_\rho^\varepsilon(t)}2,
 &
 \chi_\rho^{-,\varepsilon}(t)
 &:=\frac{1+\sigma_\rho^\varepsilon(t)}2.
 \label{eq:v29-soft-meet-selectors}
\end{align}
The two weights lie in \([0,1]\) and sum to one.  Put
\begin{align}
 V_\rho^{\wedge,\varepsilon}
 &:=\chi_\rho^{+,\varepsilon}V_\rho^+
   +\chi_\rho^{-,\varepsilon}V_\rho^-,
 \label{eq:v29-soft-meet-dissipation}\\
 \widetilde\Phi_\rho^{\wedge,\varepsilon}
 &:=\chi_\rho^{+,\varepsilon}\widetilde\Phi_\rho^+
   +\chi_\rho^{-,\varepsilon}\widetilde\Phi_\rho^-,
 \label{eq:v29-soft-meet-flux}\\
 \mathcal P_{\wedge,\varepsilon}(I)
 &:=\int_I\int_0^\infty
       V_\rho^{\wedge,\varepsilon}(t)\,\dd\rho\,\dd t,
 \label{eq:v29-soft-meet-action}\\
 \mathcal B_{\wedge,\varepsilon}(I)
 &:=\int_I\int_0^\infty
       \widetilde\Phi_\rho^{\wedge,\varepsilon}(t)
       \,\dd\rho\,\dd t.
 \label{eq:v29-soft-meet-work}
\end{align}

\begin{theorem}[Exact soft-meet balance]
\label{thm:v29-soft-meet-balance}
On every compact smooth interval, for almost every \((\rho,t)\),
\begin{equation}
 \boxed{\quad
 \frac12\partial_t m_\rho^\varepsilon
 +\nu V_\rho^{\wedge,\varepsilon}
 =\widetilde\Phi_\rho^{\wedge,\varepsilon}.
 \quad}
 \label{eq:v29-pointwise-soft-meet-balance}
\end{equation}
Consequently, for \(I=[a,b]\),
\begin{equation}
 \boxed{\quad
 \mathcal B_{\wedge,\varepsilon}(I)
 =\frac12\bigl(\mathcal C_\varepsilon(b)
                 -\mathcal C_\varepsilon(a)\bigr)
  +\nu\mathcal P_{\wedge,\varepsilon}(I).
 \quad}
 \label{eq:v29-integrated-soft-meet-balance}
\end{equation}
At each smooth time its nonlinear work density also satisfies
\begin{equation}
 \int_0^\infty\widetilde\Phi_\rho^{\wedge,\varepsilon}
 \,\dd\rho
 =\frac12\left(
    \mathcal G_+(t)+\mathcal G_-(t)
    -\mathfrak a_\varepsilon(t)\right).
 \label{eq:v29-soft-meet-work-source-relation}
\end{equation}
Thus V28's clipped adaptive source is exactly the amount subtracted from the
assembled total critical work when the soft minimum is followed.
\end{theorem}

\begin{proof}
Differentiate \eqref{eq:v29-soft-meet-density}.  Since
\(\Psi_\varepsilon'=\vartheta_\varepsilon\),
\begin{align*}
 \frac12\partial_t m_\rho^\varepsilon
 &=\frac14\left[
    \partial_tH_\rho^++\partial_tH_\rho^-
    -\sigma_\rho^\varepsilon
       (\partial_tH_\rho^+-\partial_tH_\rho^-)
   \right]\\
 &=\frac12\chi_\rho^{+,\varepsilon}\partial_tH_\rho^+
   +\frac12\chi_\rho^{-,\varepsilon}\partial_tH_\rho^-.
\end{align*}
Multiply the two sector balances
\eqref{eq:v23-sector-high-balance} by the corresponding soft selectors and
add.  This proves \eqref{eq:v29-pointwise-soft-meet-balance}.  Spectral
decay on a compact smooth interval permits cutoff--time integration and gives
\eqref{eq:v29-integrated-soft-meet-balance}.

Finally, \eqref{eq:v23-sector-sum-full-flux} and
\(d_\rho=\widetilde\Phi_\rho^+-\widetilde\Phi_\rho^-\) give
\[
 \widetilde\Phi_\rho^{\wedge,\varepsilon}
 =\frac12(\widetilde\Phi_\rho^+
           +\widetilde\Phi_\rho^-)
  -\frac12\sigma_\rho^\varepsilon d_\rho.
\]
Integrate in \(\rho\), use
\eqref{eq:v23-sector-layer-cake} for each sector, and use
\eqref{eq:v28-clipped-adaptive-source}.
\end{proof}

% =============================================================================
\section{The hard--soft difference is exactly the clipping defect}
\label{sec:v29-hard-soft-defect}
% =============================================================================

Define the scalar Huber cushion and its cutoff integral by
\begin{align}
 \Delta_\varepsilon(y)
 &:=|y|-\Psi_\varepsilon(y)
 =\begin{cases}
    |y|-y^2/(2\varepsilon),&|y|\le\varepsilon,\\
    \varepsilon/2,&|y|\ge\varepsilon,
   \end{cases}
 \label{eq:v29-Huber-cushion-density}\\
 \mathcal K_\varepsilon(t)
 &:=\int_0^\infty
      \Delta_\varepsilon(\mathcal Y_\rho(t))\,\dd\rho.
 \label{eq:v29-Huber-cushion-stock}
\end{align}

\begin{theorem}[Exact hard--soft defect ledger]
\label{thm:v29-hard-soft-ledger}
At every smooth time,
\begin{align}
 m_\rho^\varepsilon
 &=m_\rho+\frac12
   \Delta_\varepsilon(\mathcal Y_\rho),
 \label{eq:v29-hard-soft-density}\\
 \boxed{\quad
 \mathcal C_\varepsilon(t)
 =\frac12\bigl(\mathcal R(t)^2-\mathcal L_\varepsilon(t)\bigr)
 =\mathcal C(t)+\frac12\mathcal K_\varepsilon(t).
 \quad}
 \label{eq:v29-hard-soft-stock}\\
 0\le\mathcal K_\varepsilon(t)
 &=\mathcal L(t)-\mathcal L_\varepsilon(t)
 \le\mathcal L(t)\le\mathcal R(t)^2.
 \label{eq:v29-Huber-cushion-cap}
\end{align}
If \(R_\varepsilon=\operatorname{sgn}-\vartheta_\varepsilon\) is the V28
residual selector, then
\begin{align}
 V_\rho^{\wedge,\varepsilon}-V_\rho^\wedge
 &=\frac12R_\varepsilon(\mathcal Y_\rho)
       (V_\rho^+-V_\rho^-),
 \label{eq:v29-hard-soft-dissipation-difference}\\
 \widetilde\Phi_\rho^{\wedge,\varepsilon}
  -\widetilde\Phi_\rho^\wedge
 &=\frac12R_\varepsilon(\mathcal Y_\rho)d_\rho.
 \label{eq:v29-hard-soft-flux-difference}
\end{align}
In particular, the V28 dynamic residual obeys both exact identities
\begin{align}
 \boxed{\quad
 \mathcal N_I^\varepsilon
 =\frac12\bigl(\mathcal K_\varepsilon(b)
                -\mathcal K_\varepsilon(a)\bigr)
  +\nu\int_I\int_0^\infty
    R_\varepsilon(\mathcal Y_\rho)
    (V_\rho^+-V_\rho^-)
    \,\dd\rho\,\dd t,
 \quad}
 \label{eq:v29-residual-cushion-balance}\\
 \boxed{\quad
 \mathcal B_{\wedge,\varepsilon}(I)
  -\mathcal B_\wedge(I)
 =\frac12\mathcal N_I^\varepsilon.
 \quad}
 \label{eq:v29-work-residual-half}
\end{align}
No source, endpoint, or viscous term is lost in passing between the two
forms.
\end{theorem}

\begin{proof}
The Jordan formula
\(m_\rho=(H_\rho^\Sigma-|\mathcal Y_\rho|)/2\), together with
\eqref{eq:v29-soft-meet-density}, proves
\eqref{eq:v29-hard-soft-density}.  Integrating
\eqref{eq:v29-soft-meet-density} and using
\(\int_0^\infty H_\rho^\Sigma\,\dd\rho=\mathcal R(t)^2\) proves the
first equality in \eqref{eq:v29-hard-soft-stock}; integrating the density
identity proves the second.

Since
\[
 \chi_\rho^{+,\varepsilon}-\chi_\rho^+
 =\frac12R_\varepsilon(\mathcal Y_\rho),
 \qquad
 \chi_\rho^{-,\varepsilon}-\chi_\rho^-
 =-\frac12R_\varepsilon(\mathcal Y_\rho),
\]
the two pointwise difference formulae follow immediately.  Integrating
\eqref{eq:v29-hard-soft-flux-difference} proves
\eqref{eq:v29-work-residual-half}.  Alternatively, subtract the hard meet
balance \eqref{eq:v25-integrated-meet-balance} from
\eqref{eq:v29-integrated-soft-meet-balance}, then insert
\eqref{eq:v29-hard-soft-stock} and
\eqref{eq:v29-hard-soft-dissipation-difference}; comparison gives
\eqref{eq:v29-residual-cushion-balance}.  This is also the difference of the
hard and clipped entropy balances from V25 and V28.
\end{proof}

\begin{proposition}[The soft meet contains a precisely measured pure-sector
cushion]
\label{prop:v29-pure-chiral-cushion}
At a cutoff where one helical tail vanishes and the other has height \(x>0\),
the genuine meet vanishes, whereas
\begin{equation}
 m_\rho^\varepsilon
 =\begin{cases}
   \displaystyle\frac{x}{2}-\frac{x^2}{4\varepsilon},
      &0<x\le\varepsilon,\\[4pt]
   \displaystyle\frac{\varepsilon}{4},
      &x\ge\varepsilon.
  \end{cases}
 \label{eq:v29-pure-chiral-cushion}
\end{equation}
This excess is exactly
\(\Delta_\varepsilon(x)/2\), as required by
\eqref{eq:v29-hard-soft-density}.
\end{proposition}

\begin{proof}
Put \((H_\rho^+,H_\rho^-)=(x,0)\), or reverse the sectors, in
\eqref{eq:v29-soft-meet-density} and use
\eqref{eq:v28-Huber-entropy}.  The hard minimum is zero.
\end{proof}

\begin{firewall}[Soft meet is a regularized ledger, not common matter]
\label{fw:v29-soft-meet-not-genuine}
Proposition~\ref{prop:v29-pure-chiral-cushion} forbids interpreting
\(\mathcal C_\varepsilon\) as physical overlap of the two helical sectors.
Even an exactly pure-sector tail contributes \(\varepsilon/4\) per saturated
cutoff.  The genuine common stock is \(\mathcal C\); the entire discrepancy
is the displayed cushion \(\mathcal K_\varepsilon/2\).  Any ancestry argument
using the soft balance must return that cushion to the residual ledger
\eqref{eq:v29-residual-cushion-balance}.
\end{firewall}

% =============================================================================
\section{A Leray-funded radius--time occupation theorem}
\label{sec:v29-radius-time-occupation}
% =============================================================================

\begin{theorem}[Weighted occupation of all spectral tails]
\label{thm:v29-weighted-tail-occupation}
For every compact smooth interval \(I\subset[0,T_*)\),
\begin{align}
 \int_0^\infty\rho H_\rho^\Sigma(t)\,\dd\rho
 &=\frac12\|\Lambda u(t)\|_2^2,
 \label{eq:v29-weighted-tail-layer-cake}\\
 \boxed{\quad
 \nu\int_I\int_0^\infty
       \rho H_\rho^\Sigma(t)\,\dd\rho\,\dd t
 \le\frac{E_*}{4}.
 \quad}
 \label{eq:v29-global-tail-occupation}\\
 \nu\int_I\int_0^\infty\rho m_\rho(t)\,\dd\rho\,\dd t
 &\le\frac{E_*}{8},
 \label{eq:v29-hard-meet-occupation}\\
 \nu\int_I\int_0^\infty\rho m_\rho^\varepsilon(t)
       \,\dd\rho\,\dd t
 &\le\frac{E_*}{8}.
 \label{eq:v29-soft-meet-occupation}
\end{align}
All four bounds are simultaneous in the cutoff radius; no crossing number is
present.
\end{theorem}

\begin{proof}
For a Fourier atom of radius \(\lambda\),
\(H_\rho^\Sigma\) contains its energy for
\(0\le\rho<\lambda\).  Hence Tonelli gives
\[
 \int_0^\infty\rho H_\rho^\Sigma\,\dd\rho
 =\sum_{k\ne0}|\widehat u(k)|^2
   \int_0^{|k|}\rho\,\dd\rho
 =\frac12\sum_{k\ne0}|k|^2|\widehat u(k)|^2.
\]
This proves \eqref{eq:v29-weighted-tail-layer-cake}.  Integrate in time and
use the Leray energy inequality
\(\nu\int_I\|\Lambda u\|_2^2\,\dd t\le E_*/2\) to obtain
\eqref{eq:v29-global-tail-occupation}.  Finally
\[
 0\le m_\rho\le\frac12H_\rho^\Sigma,
 \qquad
 0\le m_\rho^\varepsilon\le\frac12H_\rho^\Sigma,
\]
the second inequality following from
\(0\le\Psi_\varepsilon(\mathcal Y_\rho)\le H_\rho^\Sigma\).
Apply \eqref{eq:v29-global-tail-occupation}.
\end{proof}

For \(h>\varepsilon>0\), introduce the high-total near-balance region
\begin{equation}
 \mathfrak O_{\varepsilon,h}
 :=\{(\rho,t):
       |\mathcal Y_\rho(t)|<\varepsilon,
       \ H_\rho^\Sigma(t)\ge h\}.
 \label{eq:v29-high-balanced-region}
\end{equation}

\begin{corollary}[High-balanced bands have at most parabolic lifetime]
\label{cor:v29-parabolic-band-lifetime}
For every smooth interval \(I\),
\begin{equation}
 \boxed{\quad
 \nu h\iint_{\mathfrak O_{\varepsilon,h}\cap
                   ([0,\infty)\times I)}
       \rho\,\dd\rho\,\dd t
 \le\frac{E_*}{4}.
 \quad}
 \label{eq:v29-high-balanced-occupation}
\end{equation}
If \(B\subset[K,\infty)\) is measurable and
\(B\times J\subset\mathfrak O_{\varepsilon,h}\), then
\begin{equation}
 |J|
 \le\frac{E_*}{4\nu h\int_B\rho\,\dd\rho}
 \le\frac{E_*}{4\nu hK|B|}.
 \label{eq:v29-band-lifetime}
\end{equation}
In particular, if \(h=\theta E_*\) and
\(B=[K,(1+\gamma)K]\), where \(\theta,\gamma>0\), then
\begin{equation}
 \boxed{\quad
 |J|\le
 \frac{1}{4\nu\theta(\gamma+\gamma^2/2)K^2}.
 \quad}
 \label{eq:v29-record-parabolic-lifetime}
\end{equation}
Thus a common high-frequency band at the canonical record scale cannot remain
high and nearly balanced for longer than a constant multiple of
\((\nu K^2)^{-1}\).
\end{corollary}

\begin{proof}
On \(\mathfrak O_{\varepsilon,h}\),
\(H_\rho^\Sigma\ge h\), and the V28 Jordan estimate also gives
\(m_\rho\ge(h-\varepsilon)/2>0\).  Restrict
\eqref{eq:v29-global-tail-occupation} to this set to obtain
\eqref{eq:v29-high-balanced-occupation}.  A product rectangle contained in
the set contributes
\(\nu h|J|\int_B\rho\,\dd\rho\), proving
\eqref{eq:v29-band-lifetime}.  For the stated interval,
\[
 \int_K^{(1+\gamma)K}\rho\,\dd\rho
 =(\gamma+\gamma^2/2)K^2.
\]
Insert \(h=\theta E_*\).
\end{proof}

\begin{corollary}[Cofinal persistence times are globally one-use]
\label{cor:v29-cofinal-persistence-one-use}
Let \(J_j\) be pairwise disjoint time intervals and suppose
\(B_j\times J_j\subset\mathfrak O_{\varepsilon_j,h_j}\), with
\(B_j\subset[K_j,\infty)\).  Then
\begin{equation}
 \boxed{\quad
 \sum_j\nu h_jK_j|B_j|\,|J_j|
 \le\frac{E_*}{4}.
 \quad}
 \label{eq:v29-cofinal-persistence-sum}
\end{equation}
For \(h_j=\theta E_*\) and
\(B_j=[K_j,(1+\gamma)K_j]\),
\begin{equation}
 \sum_j\nu K_j^2|J_j|
 \le\frac{1}{4\theta(\gamma+\gamma^2/2)}.
 \label{eq:v29-cofinal-parabolic-sum}
\end{equation}
\end{corollary}

\begin{proof}
The rectangles are disjoint in time, so their weighted occupation integrals
are disjoint subsets of the single global integral in
\eqref{eq:v29-global-tail-occupation}.  Sum the lower bounds used in the
proof of Corollary~\ref{cor:v29-parabolic-band-lifetime}.  The interval
calculation gives the final statement.
\end{proof}

\begin{assessment}[What the occupation theorem buys]
\label{ass:v29-occupation-meaning}
V25's common stock was one-use only across radial bands at one fixed time.
Equation~\eqref{eq:v29-cofinal-persistence-sum} is a genuine time-dependent
one-use statement: persistence rectangles in disjoint record corridors spend
one global Leray ledger.  It is, however, an upper bound on residence time.
A large terminal common stock need not have resided for a parabolic interval;
it may have been created immediately before the endpoint.  The next section
makes that recent-creation alternative quantitative rather than silently
assuming persistence.
\end{assessment}

% =============================================================================
\section{Backward lifetime and exact recent-entry alternatives}
\label{sec:v29-backward-lifetime}
% =============================================================================

Fix a terminal time \(b<T_*\).  For a cutoff satisfying the terminal
high-balanced condition, define its backward residence length
\begin{equation}
 \ell_{\varepsilon,h}(\rho;b)
 :=\sup\left\{0\le\ell\le b:
   (\rho,t)\in\mathfrak O_{\varepsilon,h}
   \text{ for every }t\in[b-\ell,b]\right\}.
 \label{eq:v29-high-balanced-backward-life}
\end{equation}
Endpoint conventions do not affect any integral below.  Smoothness makes the
tail functions continuous in time and the lifetime measurable in \(\rho\).

\begin{proposition}[Long-lived terminal radii are Leray-small]
\label{prop:v29-long-lived-terminal-radii}
Let \(S\subset[K,\infty)\) be measurable and suppose
\((\rho,b)\in\mathfrak O_{\varepsilon,h}\) for every \(\rho\in S\).
Let \(0<\tau\le b\).  Then
\begin{align}
 \int_{\{\rho\in S:\ell_{\varepsilon,h}(\rho;b)\ge\tau\}}
       \rho\,\dd\rho
 &\le\frac{E_*}{4\nu h\tau},
 \label{eq:v29-long-life-weighted-width}\\
 \left|\{\rho\in S:
       \ell_{\varepsilon,h}(\rho;b)\ge\tau\}\right|
 &\le\frac{E_*}{4\nu hK\tau}.
 \label{eq:v29-long-life-width}
\end{align}
On a first record-doubling corridor \(I=[a,b]\), the corresponding part of
the V28 endpoint residual satisfies
\begin{equation}
 \int_{\substack{\rho\in S:\
          \ell_{\varepsilon,h}(\rho;b)\ge\tau}}
 |R_\varepsilon(\mathcal Y_\rho(b))D_I(\rho)|\,\dd\rho
 \le\frac{3E_*^2}{8\nu hK\tau}.
 \label{eq:v29-long-life-residual-bound}
\end{equation}
\end{proposition}

\begin{proof}
For every radius in the long-lived set, the vertical segment
\([b-\tau,b]\) lies in \(\mathfrak O_{\varepsilon,h}\).  Restrict
\eqref{eq:v29-high-balanced-occupation} to the union of those segments and
divide by \(\nu h\tau\).  Since \(\rho\ge K\), the second estimate follows.
On a first doubling, \eqref{eq:v24-difference-height-cap} gives
\(|D_I(\rho)|\le3E_*/2\), while \(|R_\varepsilon|\le1\).  Multiply by
\eqref{eq:v29-long-life-width}.
\end{proof}

The complement in Proposition~\ref{prop:v29-long-lived-terminal-radii} has
recently left the high-balanced region.  A margin between the terminal and
residence thresholds turns that statement into a positive quantitative debit.

\begin{theorem}[Two-threshold recent-entry alternative]
\label{thm:v29-two-threshold-entry}
Fix \(0<\beta<1\), \(\kappa>0\), and a terminal cutoff \(\rho\) such that
\begin{equation}
 H_\rho^\Sigma(b)\ge(1+\kappa)h,
 \qquad
 |\mathcal Y_\rho(b)|\le\beta\varepsilon.
 \label{eq:v29-inner-terminal-region}
\end{equation}
If \(\ell_{\varepsilon,h}(\rho;b)<\tau\), then there is
\(t_\rho\in(b-\tau,b)\) for which at least one of the following holds:
\begin{align}
 \int_{t_\rho}^b\Phi_\rho(t)\,\dd t
 &\ge\frac{\kappa h}{2},
 \tag{TOT}\label{eq:v29-recent-total-influx}\\
 \left|\int_{t_\rho}^b d_\rho(t)\,\dd t\right|
 +\nu\int_{t_\rho}^b
     (V_\rho^+(t)+V_\rho^-(t))\,\dd t
 &\ge\frac{(1-\beta)\varepsilon}{2}.
 \tag{POL}\label{eq:v29-recent-polarization-action}
\end{align}
Here \(\Phi_\rho=\widetilde\Phi_\rho^+
+\widetilde\Phi_\rho^-\) is the complete kinetic-energy flux through the
cutoff.  Thus every short-lived inner near-balance cutoff has a recent total
influx or a recent sector-polarization action.
\end{theorem}

\begin{proof}
By time continuity and the definition of the maximal backward component,
there is a recent exit time at which either
\(H_\rho^\Sigma(t_\rho)\le h\) or
\(|\mathcal Y_\rho(t_\rho)|\ge\varepsilon\).

In the first case, sum the two sector balances and integrate:
\[
 \int_{t_\rho}^b\Phi_\rho(t)\,\dd t
 =\frac12\left(H_\rho^\Sigma(b)
                    -H_\rho^\Sigma(t_\rho)\right)
  +\nu\int_{t_\rho}^b(V_\rho^++V_\rho^-)\,\dd t
 \ge\frac{\kappa h}{2}.
\]
In the second case,
\[
 |\mathcal Y_\rho(t_\rho)-\mathcal Y_\rho(b)|
 \ge(1-\beta)\varepsilon.
\]
Integrate \eqref{eq:v25-instant-tail-difference} from \(t_\rho\) to \(b\),
take absolute values, and use
\(|V_\rho^+-V_\rho^-|\le V_\rho^++V_\rho^-\).  This gives
\eqref{eq:v29-recent-polarization-action}.
\end{proof}

\begin{firewall}[A recent-entry debit is not yet a cofinal source bound]
\label{fw:v29-entry-not-summed}
The exit time \(t_\rho\) depends on \(\rho\).  The positive total flux in
\eqref{eq:v29-recent-total-influx} and the signed source in
\eqref{eq:v29-recent-polarization-action} therefore live on a moving family
of cutoff--time intervals.  Taking their absolute values and integrating in
\(\rho\) would return the critical nonlinear source.  The theorem proves the
correct alternative for each terminal cutoff, but a radial hinge or Duhamel
assembly must still show that these recent debits are one-use before they can
be summed over cofinal records.
\end{firewall}

% =============================================================================
\section{The low-total residual: persistence or recent replenishment}
\label{sec:v29-low-total-ancestry}
% =============================================================================

The high-total threshold cannot cover the V28 low-total residual.  A second
lifetime, normalized to the terminal total tail itself, removes that threshold.
For \(0<\alpha<1\) and \(H_\rho^\Sigma(b)>0\), set
\begin{equation}
 \ell_\alpha^\Sigma(\rho;b)
 :=\sup\left\{0\le\ell\le b:
 H_\rho^\Sigma(t)\ge
 \alpha H_\rho^\Sigma(b)
 \text{ for every }t\in[b-\ell,b]\right\}.
 \label{eq:v29-relative-tail-lifetime}
\end{equation}

\begin{theorem}[Terminal tail mass has a paid-persistence/recent-influx
alternative]
\label{thm:v29-tail-mass-ancestry}
Let \(S\subset[K,\infty)\) be measurable and suppose
\(H_\rho^\Sigma(b)>0\) on \(S\).  For \(0<\tau\le b\), put
\begin{align*}
 S_{\rm long}&:=\{\rho\in S:
      \ell_\alpha^\Sigma(\rho;b)\ge\tau\},\\
 S_{\rm new}&:=S\setminus S_{\rm long}.
\end{align*}
Then
\begin{equation}
 \boxed{\quad
 \int_{S_{\rm long}}H_\rho^\Sigma(b)\,\dd\rho
 \le\frac{E_*}{4\nu\alpha K\tau}.
 \quad}
 \label{eq:v29-long-lived-terminal-tail-mass}
\end{equation}
For every \(\rho\in S_{\rm new}\), there is
\(t_\rho\in(b-\tau,b)\) such that
\begin{equation}
 \boxed{\quad
 \int_{t_\rho}^b\Phi_\rho(t)\,\dd t
 \ge\frac{1-\alpha}{2}H_\rho^\Sigma(b).
 \quad}
 \label{eq:v29-relative-recent-influx}
\end{equation}
\end{theorem}

\begin{proof}
On \(S_{\rm long}\times[b-\tau,b]\), the total tail is at least
\(\alpha H_\rho^\Sigma(b)\).  Restrict
\eqref{eq:v29-global-tail-occupation}, use \(\rho\ge K\), and divide by
\(\nu\alpha K\tau\).  This proves
\eqref{eq:v29-long-lived-terminal-tail-mass}.

If \(\rho\in S_{\rm new}\), time continuity supplies a recent time with
\(H_\rho^\Sigma(t_\rho)\le\alpha H_\rho^\Sigma(b)\).  The integrated sum of
the two sector balances gives
\[
 \int_{t_\rho}^b\Phi_\rho\,\dd t
 =\frac12(H_\rho^\Sigma(b)-H_\rho^\Sigma(t_\rho))
  +\nu\int_{t_\rho}^b(V_\rho^++V_\rho^-)\,\dd t,
\]
which is at least the right side of
\eqref{eq:v29-relative-recent-influx}.
\end{proof}

We now apply the theorem to the exact residual profile rather than to a new
surrogate.  On a first-doubling corridor \(I=[a,b]\), let
\begin{equation}
 \mathfrak L_{\varepsilon,h,K}(b)
 :=\{\rho\ge K:
   0<|\mathcal Y_\rho(b)|<\varepsilon,
   \ H_\rho^\Sigma(b)<h\}.
 \label{eq:v29-low-total-terminal-set}
\end{equation}
This is the V28 low-total set restricted only by the record-visible lower
cutoff \(K\).  Its complementary low-cutoff part is not omitted.  Indeed,
\eqref{eq:v24-difference-height-cap} gives
\begin{equation}
 \int_{\substack{0\le\rho<K:\\
       0<|\mathcal Y_\rho(b)|<\varepsilon,\\
       H_\rho^\Sigma(b)<h}}
 |R_\varepsilon(\mathcal Y_\rho(b))D_I(\rho)|\,\dd\rho
 \le\frac{3E_*K}{2}.
 \label{eq:v29-low-cutoff-residual-bound}
\end{equation}
At the canonical choice \(K=R^2/E_*\), this is \(3R^2/2\); it is exact
bookkeeping but not a separating estimate.

\begin{proposition}[Exact low-total residual ledger after the ancestry split]
\label{prop:v29-low-total-ledger}
Set
\begin{equation}
 \mathcal V_I^\Delta(\rho)
 :=\int_I(V_\rho^+-V_\rho^-)\,\dd t.
 \label{eq:v29-integrated-viscous-difference}
\end{equation}
For every measurable \(S\subset
\mathfrak L_{\varepsilon,h,K}(b)\),
\begin{align}
 \int_S R_\varepsilon(\mathcal Y_\rho(b))D_I(\rho)\,\dd\rho
 &=\frac12\int_S
      R_\varepsilon(\mathcal Y_\rho(b))\mathcal Y_\rho(b)\,\dd\rho
 \notag\\
 &\quad-\frac12\int_S
      R_\varepsilon(\mathcal Y_\rho(b))\mathcal Y_\rho(a)\,\dd\rho
 \notag\\
 &\quad+\nu\int_S
      R_\varepsilon(\mathcal Y_\rho(b))
      \mathcal V_I^\Delta(\rho)\,\dd\rho.
 \label{eq:v29-exact-low-total-ledger}
\end{align}
Moreover,
\begin{equation}
 0\le R_\varepsilon(\mathcal Y_\rho(b))
          \mathcal Y_\rho(b)
 \le |\mathcal Y_\rho(b)|
 \le H_\rho^\Sigma(b).
 \label{eq:v29-terminal-residual-self-bound}
\end{equation}
If \(S=S_{\rm long}\) is the long-lived subset from
Theorem~\ref{thm:v29-tail-mass-ancestry}, then its terminal self term obeys
\begin{equation}
 0\le\frac12\int_{S_{\rm long}}
 R_\varepsilon(\mathcal Y_\rho(b))
 \mathcal Y_\rho(b)\,\dd\rho
 \le\frac{E_*}{8\nu\alpha K\tau}.
 \label{eq:v29-low-total-long-self-bound}
\end{equation}
Every complementary radius has the recent positive influx
\eqref{eq:v29-relative-recent-influx}.  Thus the terminal part of the
low-total residual is rigorously divided into a Leray-paid persistent part and
a recently replenished part.
\end{proposition}

\begin{proof}
Integrating \eqref{eq:v25-instant-tail-difference} over \(I\) gives
\[
 D_I(\rho)
 =\frac12(\mathcal Y_\rho(b)-\mathcal Y_\rho(a))
  +\nu\mathcal V_I^\Delta(\rho).
\]
Multiply by the frozen terminal residual selector and integrate over \(S\) to
obtain \eqref{eq:v29-exact-low-total-ledger}.  On its support,
\[
 R_\varepsilon(y)y=|y|-\frac{y^2}{\varepsilon},
 \qquad 0<|y|<\varepsilon,
\]
which proves \eqref{eq:v29-terminal-residual-self-bound}.  Apply
\eqref{eq:v29-long-lived-terminal-tail-mass} to obtain
\eqref{eq:v29-low-total-long-self-bound}; the recent-influx assertion is the
other branch of Theorem~\ref{thm:v29-tail-mass-ancestry}.
\end{proof}

\begin{firewall}[What remains in the low-total ledger]
\label{fw:v29-low-total-remainder}
Equation~\eqref{eq:v29-low-total-long-self-bound} controls only the terminal
self term in \eqref{eq:v29-exact-low-total-ledger}.  The initial signed tail
and the frozen viscous-difference term remain exact.  On a first doubling they
have only the immediate bounds
\begin{align*}
 \frac12\int_S|\mathcal Y_\rho(a)|\,\dd\rho
 &\le\frac12\mathcal R(a)^2=\frac12R^2,\\
 \nu\int_S|\mathcal V_I^\Delta(\rho)|\,\dd\rho
 &\le\nu\int_I\int_S(V_\rho^++V_\rho^-)
       \,\dd\rho\,\dd t,
\end{align*}
and the latter is critical tail dissipation, not the Leray action.  Likewise,
the recent flux intervals in \eqref{eq:v29-relative-recent-influx} depend on
\(\rho\).  Closing the hard residual requires an assembled moving-interval
identity that combines these terms before absolute values.  Low terminal
height alone does not provide it.
\end{firewall}

% =============================================================================
\section{Integrated V29 conclusion and V30 acquisition order}
\label{sec:v29-conclusion}
% =============================================================================

\begin{theorem}[What V29 proves]
\label{thm:v29-conclusion}
Relative to V1--V28, the following are proved.
\begin{enumerate}[label=\textup{(V29.\arabic*)},leftmargin=6.2em]
\item The Huber selector is the gradient of the soft meet
      \eqref{eq:v29-soft-meet-density}, whose pointwise and integrated PDE
      balances are \eqref{eq:v29-pointwise-soft-meet-balance} and
      \eqref{eq:v29-integrated-soft-meet-balance}.
\item The soft stock exceeds the genuine common stock by exactly half the
      Huber cushion, and the difference of the soft and hard meet works is
      exactly half the V28 dynamic residual; see
      \eqref{eq:v29-hard-soft-stock}--\eqref{eq:v29-work-residual-half}.
\item A pure-sector tail still carries the explicit soft cushion
      \eqref{eq:v29-pure-chiral-cushion}; therefore no artificial soft stock
      is counted as genuine two-sector overlap.
\item The cutoff-radius-weighted occupation of all tails, the hard meet, and
      the soft meet is paid by ordinary Leray dissipation, as in
      \eqref{eq:v29-global-tail-occupation}--%
      \eqref{eq:v29-soft-meet-occupation}.
\item A high-total near-balanced band has the parabolic maximum lifetime
      \eqref{eq:v29-record-parabolic-lifetime}.  Across disjoint record
      corridors, all such persistence times obey the global one-use sum
      \eqref{eq:v29-cofinal-parabolic-sum}.
\item At a terminal high-balanced cutoff, either the backward lifetime is
      Leray-small in aggregate or a recent exit forces the quantitative total
      influx/polarization alternative
      \eqref{eq:v29-recent-total-influx}--%
      \eqref{eq:v29-recent-polarization-action}.
\item Without any lower height threshold, terminal total-tail mass is divided
      between a paid persistent part
      \eqref{eq:v29-long-lived-terminal-tail-mass} and a recent positive
      flux part \eqref{eq:v29-relative-recent-influx}.
\item Applying that split to V28's low-total residual controls its long-lived
      terminal self term by \eqref{eq:v29-low-total-long-self-bound}, while
      retaining the initial, viscous, and moving recent-source terms exactly;
      the complementary low-cutoff piece has the explicit nonseparating bound
      \eqref{eq:v29-low-cutoff-residual-bound}.
\end{enumerate}
These results supply the promised common-tail persistence debit and a rigorous
recent-creation alternative.  They do not yet assemble the radius-dependent
recent source intervals, so the full Navier--Stokes stability/global-
regularity theorem remains unproved.
\end{theorem}

\begin{proof}
Items \textup{(V29.1)--(V29.3)} are
Theorems~\ref{thm:v29-soft-meet-balance} and
\ref{thm:v29-hard-soft-ledger}, and
Proposition~\ref{prop:v29-pure-chiral-cushion}.  Item \textup{(V29.4)} is
Theorem~\ref{thm:v29-weighted-tail-occupation}.  Item \textup{(V29.5)} is
Corollaries~\ref{cor:v29-parabolic-band-lifetime} and
\ref{cor:v29-cofinal-persistence-one-use}.  Item \textup{(V29.6)} is
Proposition~\ref{prop:v29-long-lived-terminal-radii} and
Theorem~\ref{thm:v29-two-threshold-entry}.  Item \textup{(V29.7)} is
Theorem~\ref{thm:v29-tail-mass-ancestry}.  Item \textup{(V29.8)} is
Proposition~\ref{prop:v29-low-total-ledger} and
Firewall~\ref{fw:v29-low-total-remainder}.
\end{proof}

\begin{programme}[V30 companion audit and moving-entry assembly]
\label{prog:v30}
\begin{enumerate}[label=\textup{(AC\arabic*)},leftmargin=4.8em]
\item Begin V30 with the compulsory filesystem review of the latest
      Standard-Model and Yang--Mills notes.  Record their versions and hashes,
      and import only proved mechanisms with matching analytic type.
\item Keep the hard--soft defect identity
      \eqref{eq:v29-residual-cushion-balance} exact.  Any use of the smooth
      soft balance must return the pure-sector cushion exposed by
      Proposition~\ref{prop:v29-pure-chiral-cushion}.
\item Assemble the radius-dependent entry intervals from
      Theorems~\ref{thm:v29-two-threshold-entry} and
      \ref{thm:v29-tail-mass-ancestry} before taking absolute values.  Test a
      stopping-time version of V27's hinge current and radial BV ledger.
\item Use geometric layers in total height and polarization margin so that the
      inner two-threshold region covers the large-residual core, while the
      outer polarization annuli gain the explicit factor
      \(|R_\varepsilon(y)|=1-|y|/\varepsilon\).
\item Reuse, rather than rederive, the exact shell and parabolic Duhamel
      modules of V7 and V10--V21 when converting recent positive influx into
      source ancestry.  Preserve their endpoint-storage and moving-cutoff
      defects.
\item On the low-total branch, combine
      \eqref{eq:v29-exact-low-total-ledger} with the recent-influx identity.
      Seek cancellation between the initial-tail and moving-source terms
      before estimating the frozen viscous difference.
\item Normalize every candidate estimate at the record scale
      \(K=R^2/E_*\).  The parabolic persistence sum is globally paid, but a
      per-record upper bound of order \(R^2\) with a nonseparating constant is
      not closure.
\item Retain the V28 productive low-minority test and all complementary triad
      orientations.  No local sign class may be promoted to a full nonlinear
      incidence estimate.
\end{enumerate}
\end{programme}

\begin{assessment}[Position after V29]
The common-tail obstruction now has a temporal theorem.  Long residence of a
high-balanced band at scale \(K\) is limited by the canonical parabolic time,
and the sum of those residence times over disjoint record corridors is paid by
one Leray ledger.  A terminal band that evades this debit must be recent: with
quantitative threshold margins it requires either full-tail influx or
helical-polarization action.  Low-total terminal mass obeys the same
persistence-or-replenishment alternative without a fixed height threshold.

The remaining bottleneck is more precise than at V28.  It is not common stock
residence and not inverse crossing amplitude.  It is the simultaneous signed
assembly of radius-dependent recent-entry intervals with the initial and
viscous pieces of the hard residual.  V30 will audit the two companion
programmes first, then attack that moving-entry assembly.  The full stability/
global-regularity objective remains open.
\end{assessment}

% =============================================================================
\section*{Version V29 append-only record}
\addcontentsline{toc}{section}{Version V29 append-only record}
% =============================================================================

The complete V28 mathematical source was retained before this continuation,
apart from correcting the pre-existing control-character typo in
\eqref{eq:v27-crossing-inverse-amplitude} from a malformed fraction command
to \(\backslash\mathrm{frac}\); no mathematical statement was changed.
Its SHA--256 digest was
\begin{center}\scriptsize\ttfamily
87f11401db51366fa425a33999bc04f206113b89fd09792df8acb3c4bb5ed5ca.
\end{center}
V29 identified the Huber selector as the exact gradient of a soft spectral-tail
meet, proved its PDE balance, and showed that its excess over the genuine meet
is precisely the clipping cushion.  It proved a cutoff-radius-weighted
tail-occupation theorem paid by ordinary Leray dissipation, giving a parabolic
maximum lifetime and a cofinally one-use sum for persistent high-balanced
bands.  It also proved quantitative recent-influx/polarization alternatives
for short-lived terminal bands and extended the persistence/replenishment
split to V28's low-total residual.  The radius-dependent recent-entry sources,
initial tail, and frozen critical viscous difference remain to be assembled;
no full regularity claim is made.
% =============================================================================
% VERSION V30 APPEND
% =============================================================================

\clearpage
\section{V30: companion audit and probability-marked stopping intervals}
\label{sec:v30-entry}
% =============================================================================

V29 proved that a terminal high-balanced tail either persists long enough to
spend a Leray radius--time occupation debit or has a cutoff-dependent recent
entry time at which total-tail influx or helical-polarization action occurs.
Its open row was the simultaneous assembly of those different entry
intervals.  V30 first performs the required fifth-version companion audit and
then closes the \emph{normalized} form of that moving-time incidence problem.

The useful cross-program idea is a discipline, not an imported estimate.  The
new Yang--Mills notes avoid support-incidence multiplicity by assigning
normalized marks before summing a rooted interaction tree.  In the present NS
problem a finite positive measure on cutoff radii plays the role of a mark.
For every probability mark, every measurable radius-dependent lower stopping
time, and every bounded frozen sign, the stopped total-tail and sector-
difference fluxes obey one exact balance and one Leray bound.  Thus the number
and geometry of distinct entry times cause no loss after normalization.

There is also a sharp limitation.  A mark that is Leray-admissible on the
whole high-frequency half-line must have finite radial mass, whereas domination
of the unweighted critical-tail integral would require a density bounded below
on that half-line.  The two requirements are incompatible.  V30 therefore
converts the moving-entry bottleneck into a precise unmarking/far-tail gate; it
does not claim that a normalized incidence estimate controls the full critical
residual.

% =============================================================================
\section{Compulsory V30 review of the current companion notes}
\label{sec:v30-companion-audit}
% =============================================================================

Before selecting the V30 construction, the two current companion files in
\texttt{latex\_notes} were inspected.  The frozen read-only checkpoint is
\begin{align}
 &\texttt{Renewal\_SM\_Einstein\_Continuum\_Research\_Notes\_V54.tex},
 \notag\\[-2pt]
 &\qquad\text{SHA--256 }
 \texttt{a964602e50cb1660001cd58d80755701c60f978f9e1cf3f147dd6fe88088cbc9},
 \label{eq:v30-SM-checkpoint}\\
 &\texttt{Renewal\_Yang\_Mills\_Mass\_Gap\_Research\_Notes\_V32.tex},
 \notag\\[-2pt]
 &\qquad\text{SHA--256 }
 \texttt{f8e4d8a509c2a247cd6deef4f48009b0c4b3f3ba67f2e5eabf6429b31a409a42}.
 \label{eq:v30-YM-checkpoint}
\end{align}
The previous NS companion audit at V25 had reached SM V49 and YM V29.
Accordingly the present review concentrated on SM V50--V54 and YM V30--V32,
rather than repeating the older comparison.

The relevant proved SM conclusions are as follows.  V50 computes the direct
five-variable Einstein glancing matrix without a singular free inverse and
proves a uniform frozen metric floor, but also proves divergence of the
unweighted total bath-energy resolvent at continuous spectrum.  V51--V52 prove
a weighted inhomogeneous limiting-absorption theory and a maximal normal Green
relation.  Their own firewalls do not promote those frequency-domain results
to a time-domain generator.  V53--V54 then expose an orientation/passivity
tradeoff, an unbounded traction trace on the natural energy space, and the
failure of a finite positive mediator to retain the principal glancing floor.
These are linear half-space boundary statements, not estimates for NS Fourier
tails.

The relevant proved YM conclusions are different.  V30 retains a single
assembled defect before charging one output denominator.  V31 proves the
binary rooted-support sum by placing normalized Casimir marks on each shared
link; it also gives an abstract example showing a factor equal to support size
when those marks are removed.  Its marked-tree compiler is conditional beyond
the proved binary case.  V32 proves the ternary marked-tree estimate on every
fixed thermal window and forces any remaining obstruction to escape that
window; the uniform high-channel ternary atlas and all higher orders remain
open.  Wilson multipliers, Casimir weights, and support trees are not NS
quantities.

\begin{theorem}[V30 companion-audit decision]
\label{thm:v30-companion-decision}
The audit yields exactly the following transfer decisions.
\begin{enumerate}[label=\textup{(\roman*)},leftmargin=3.7em]
\item No SM V50--V54 resolvent, Green-domain, port, traction, or mediator
      constant is imported into the NS programme.  In particular, a
      frequency-domain inverse or an exact local power identity does not
      supply the time-dependent cofinal source bound missing after V29.
\item No YM V30--V32 Wilson, cumulant, Casimir, or marked-tree estimate is
      imported.  The uniform YM high-channel and all-order gates are still
      open in their own problem.
\item The proved YM assembly discipline is type-correct as a construction
      rule: normalize the incidence resource before summing it, and never
      replace a marked sum by the raw number of incidences.  The NS mark and
      its estimate must be derived directly from the two sector-tail PDE
      balances.
\item The SM topology warning is retained: an estimate in a strengthened
      graph or weighted topology cannot silently be read as an estimate in
      the natural energy topology.  V30 therefore prints both its normalized
      cutoff measure and the exact cost of returning to unweighted radial
      integration.
\end{enumerate}
\end{theorem}

\begin{proof}
Items \textup{(i)} and \textup{(iv)} are the explicit claim boundaries in
the SM V50 limiting-absorption assessment, the V51--V52 generator firewalls,
and the V54 traction/mediator no-go results.  Their variables and norms do not
occur in \eqref{eq:v23-sector-high-balance}.  Items \textup{(ii)} and
\textup{(iii)} are the exact scope of the YM V31 binary marked-incidence
theorem, its unmarked incidence counterexample, and the V32 thermal-window
firewall.  Since none has the NS tail functions, viscosity, or Leray budget as
hypotheses, no numerical inequality transfers.  The normalization order can
be adopted only after it is proved below from the NS balances themselves.
\end{proof}

\begin{firewall}[Companion analogies do not close an NS estimate]
\label{fw:v30-companion-types}
A Casimir mark on a finite support graph and a probability measure on the
continuum of Fourier cutoffs have different analytic types.  Likewise, an
Einstein limiting-absorption weight is not a Navier--Stokes time-occupation
weight.  The companion review determines the assembly order and the required
topology disclosure only.  Every quantitative statement in the rest of V30
is proved directly from the NS sector-tail balances and the Leray inequality.
\end{firewall}

% =============================================================================
\section{Finite cutoff marks and exact stopped balances}
\label{sec:v30-stopped-marks}
% =============================================================================

Let \(I=[a,b]\) be a compact smooth NS interval.  A \emph{cutoff mark} is a
finite positive Borel measure \(\mu\) on \((0,\infty)\).  Write
\begin{equation}
 M_\mu:=\mu((0,\infty))<\infty.
 \label{eq:v30-mark-mass}
\end{equation}
Let \(\tau:(0,\infty)\to[a,b]\) be measurable and let
\(\gamma:(0,\infty)\to[-1,1]\) be a measurable frozen multiplier.  Define
the stopped total and difference fluxes by
\begin{align}
 \mathfrak F_{\Sigma}^{\mu}[\gamma,\tau]
 &:=\int_{(0,\infty)}\gamma(\rho)
     \int_{\tau(\rho)}^b\Phi_\rho(t)\,\dd t\,\dd\mu(\rho),
 \label{eq:v30-stopped-total-flux}\\
 \mathfrak F_{\Delta}^{\mu}[\gamma,\tau]
 &:=\int_{(0,\infty)}\gamma(\rho)
     \int_{\tau(\rho)}^b d_\rho(t)\,\dd t\,\dd\mu(\rho),
 \label{eq:v30-stopped-difference-flux}
\end{align}
and put
\begin{align}
 \mathfrak P_{\Sigma}^{\mu}[\gamma,\tau]
 &:=\int\gamma(\rho)
       \int_{\tau(\rho)}^b(V_\rho^++V_\rho^-)(t)
       \,\dd t\,\dd\mu(\rho),
 \label{eq:v30-stopped-total-dissipation}\\
 \mathfrak P_{\Delta}^{\mu}[\gamma,\tau]
 &:=\int\gamma(\rho)
       \int_{\tau(\rho)}^b(V_\rho^+-V_\rho^-)(t)
       \,\dd t\,\dd\mu(\rho).
 \label{eq:v30-stopped-difference-dissipation}
\end{align}

\begin{theorem}[Arbitrary stopping maps have no normalized incidence loss]
\label{thm:v30-stopped-mark-balance}
For every finite cutoff mark, measurable stopping map, and frozen multiplier
as above,
\begin{align}
 \mathfrak F_{\Sigma}^{\mu}[\gamma,\tau]
 &=\frac12\int\gamma(\rho)
   \bigl(H_\rho^\Sigma(b)-H_\rho^\Sigma(\tau(\rho))\bigr)
   \,\dd\mu(\rho)
   +\nu\mathfrak P_{\Sigma}^{\mu}[\gamma,\tau],
 \label{eq:v30-exact-stopped-total}\\
 \mathfrak F_{\Delta}^{\mu}[\gamma,\tau]
 &=\frac12\int\gamma(\rho)
   \bigl(\mathcal Y_\rho(b)-\mathcal Y_\rho(\tau(\rho))\bigr)
   \,\dd\mu(\rho)
   +\nu\mathfrak P_{\Delta}^{\mu}[\gamma,\tau].
 \label{eq:v30-exact-stopped-difference}
\end{align}
Moreover,
\begin{align}
 \nu\int\int_{\tau(\rho)}^b
       (V_\rho^++V_\rho^-)(t)\,\dd t\,\dd\mu(\rho)
 &\le\frac12M_\mu E_*,
 \label{eq:v30-stopped-Leray-action}\\
 \boxed{\quad
 |\mathfrak F_{\Sigma}^{\mu}[\gamma,\tau]|
 \le\frac32M_\mu E_*,
 \qquad
 |\mathfrak F_{\Delta}^{\mu}[\gamma,\tau]|
 \le\frac32M_\mu E_*.
 \quad}
 \label{eq:v30-stopped-mark-bound}
\end{align}
The constants are independent of the number, variation, or geometry of the
distinct values taken by \(\tau\).  If \(\gamma\ge0\) and
\(H_\rho^\Sigma(b)\ge H_\rho^\Sigma(\tau(\rho))\) on its support, then
\begin{equation}
 0\le\mathfrak F_{\Sigma}^{\mu}[\gamma,\tau]
 \le M_\mu E_*.
 \label{eq:v30-positive-stopped-total}
\end{equation}
\end{theorem}

\begin{proof}
Sum and subtract the two sector balances
\eqref{eq:v23-sector-high-balance}, respectively, and integrate from
\(\tau(\rho)\) to \(b\).  This gives the pointwise stopped identities;
integration against \(\gamma\,\dd\mu\) proves
\eqref{eq:v30-exact-stopped-total}--%
\eqref{eq:v30-exact-stopped-difference}.

For the viscous term, Tonelli and
\(0\le V_\rho^++V_\rho^-\le\|\Lambda u(t)\|_2^2\) give
\begin{align*}
 &\int\int_{\tau(\rho)}^b(V_\rho^++V_\rho^-)(t)\,\dd t\,\dd\mu(\rho)\\
 &\qquad=\int_a^b\int
   \mathbf1_{\{\tau(\rho)\le t\}}(V_\rho^++V_\rho^-)(t)
   \,\dd\mu(\rho)\,\dd t
 \le M_\mu\int_a^b\|\Lambda u(t)\|_2^2\,\dd t.
\end{align*}
The Leray inequality proves \eqref{eq:v30-stopped-Leray-action}.
At every cutoff and time,
\begin{equation}
 0\le H_\rho^\Sigma(t)\le\|u(t)\|_2^2\le E_*,
 \qquad
 |\mathcal Y_\rho(t)|\le H_\rho^\Sigma(t).
 \label{eq:v30-pointwise-tail-caps}
\end{equation}
The two endpoint contributions in either stopped identity therefore sum to
at most \(M_\mu E_*\) in absolute value, and the viscous contribution is at
most \(M_\mu E_*/2\).  This proves
\eqref{eq:v30-stopped-mark-bound}.  Under the final hypotheses of the
theorem, the storage difference and viscosity are nonnegative; discard the
negative initial storage and use the terminal half-bound plus
\eqref{eq:v30-stopped-Leray-action} to obtain
\eqref{eq:v30-positive-stopped-total}.
\end{proof}

\begin{remark}[Why this is a marked assembly]
\label{rem:v30-marked-assembly}
When \(M_\mu=1\), the mark distributes one unit of radial incidence among
all stopping intervals.  The proof sums the stopped PDE identity first and
uses \(\int\dd\mu=1\) only afterward.  Replacing the mark by a count of
entry radii would destroy precisely this normalization.
\end{remark}

% =============================================================================
\section{The canonical half-line mark and the hard residual}
\label{sec:v30-canonical-mark}
% =============================================================================

For \(K>0\), define the inverse-square probability mark
\begin{equation}
 \dd\mu_K(\rho)
 :=\mathbf1_{[K,\infty)}(\rho)\frac{K}{\rho^2}\,\dd\rho,
 \qquad
 \mu_K([K,\infty))=1.
 \label{eq:v30-canonical-mark}
\end{equation}
This particular mark resolves the whole half-line while retaining an exact
spectral formula.

\begin{proposition}[Exact spectral symbol of the canonical mark]
\label{prop:v30-canonical-mark-symbol}
At every smooth time,
\begin{align}
 \int_K^\infty H_\rho^\Sigma(t)\,\dd\mu_K(\rho)
 &=\sum_{|k|>K}\sum_s
   \left(1-\frac{K}{|k|}\right)|\widehat u^s(k,t)|^2,
 \label{eq:v30-marked-tail-energy}\\
 \int_K^\infty(V_\rho^++V_\rho^-)(t)\,\dd\mu_K(\rho)
 &=\sum_{|k|>K}\sum_s
   \bigl(|k|^2-K|k|\bigr)|\widehat u^s(k,t)|^2.
 \label{eq:v30-marked-tail-dissipation}
\end{align}
In particular, the first quantity is bounded by \(E_*\), and the second by
\(\|\Lambda u(t)\|_2^2\).
\end{proposition}

\begin{proof}
A mode of radius \(\lambda>K\) occurs in the high tail precisely for
\(K\le\rho<\lambda\), and
\begin{equation}
 \int_K^\lambda\frac{K}{\rho^2}\,\dd\rho
 =1-\frac{K}{\lambda}.
 \label{eq:v30-mark-symbol-integral}
\end{equation}
Use Tonelli first with the sector energy atoms and then with the atoms
multiplied by \(\lambda^2\).  The two upper bounds follow because the factor
in \eqref{eq:v30-mark-symbol-integral} lies in \([0,1]\).
\end{proof}

\begin{corollary}[The marked hard selector and its residual are Leray-paid]
\label{cor:v30-marked-hard-residual}
On every smooth interval \(I=[a,b]\),
\begin{align}
 \max\Bigg\{
 &\left|\int_K^\infty
  \sigma_\rho(b)D_I(\rho)\frac{K}{\rho^2}\,\dd\rho\right|,\notag\\
 &\left|\int_K^\infty
  R_\varepsilon(\mathcal Y_\rho(b))D_I(\rho)
  \frac{K}{\rho^2}\,\dd\rho\right|\Bigg\}
 \le\frac32E_*.
 \label{eq:v30-marked-hard-residual}
\end{align}
More generally, the same bound holds with \(a\) replaced by an arbitrary
measurable lower stopping map \(\tau(\rho)\) and
\(D_I(\rho)\) replaced by
\(\int_{\tau(\rho)}^bd_\rho(t)\,\dd t\), for either displayed
frozen multiplier.
\end{corollary}

\begin{proof}
Use Theorem~\ref{thm:v30-stopped-mark-balance} with
\(\mu=\mu_K\) and first \(\tau\equiv a\).  Take in turn
\(\gamma(\rho)=\sigma_\rho(b)\) and
\(\gamma(\rho)=R_\varepsilon(\mathcal Y_\rho(b))\).  The identity
\(D_I(\rho)=\int_a^bd_\rho(t)\,\dd t\) gives
\eqref{eq:v30-marked-hard-residual}.  The identical theorem with general
\(\tau\) proves the second assertion.
\end{proof}

\begin{assessment}[The moving-time multiplicity is closed after marking]
\label{ass:v30-moving-time-closed-marked}
V29 left open whether a continuum of cutoff-dependent recent-entry times
would force the absolute critical source.  Corollary
\ref{cor:v30-marked-hard-residual} answers the normalized question: it does
not.  Once one unit of cutoff incidence is distributed by a probability mark,
all stopping intervals are assembled by one endpoint-energy and Leray ledger,
regardless of how irregularly their lower endpoints vary.  What is not yet
closed is removal of the mark.
\end{assessment}

% =============================================================================
\section{Marked capacity of the V29 recent-entry branches}
\label{sec:v30-entry-capacity}
% =============================================================================

Fix \(0<\beta<1\), \(\kappa>0\), \(h>\varepsilon>0\), and \(K>0\).
At terminal time \(b\), define the inner high-balanced set
\begin{equation}
 S_{\rm in}:=\left\{\rho\ge K:
 H_\rho^\Sigma(b)\ge(1+\kappa)h,
 \quad |\mathcal Y_\rho(b)|\le\beta\varepsilon\right\}.
 \label{eq:v30-inner-high-balanced-set}
\end{equation}
For \(0<\tau_0\le b\), split it using V29's backward lifetime:
\begin{align}
 S_{\rm long}
 &:=\{\rho\in S_{\rm in}:
       \ell_{\varepsilon,h}(\rho;b)\ge\tau_0\},
 \label{eq:v30-inner-long-set}\\
 S_{\rm short}&:=S_{\rm in}\setminus S_{\rm long}.
 \label{eq:v30-inner-short-set}
\end{align}
On \(S_{\rm short}\), choose the canonical exit time
\begin{equation}
 t_e(\rho):=b-\ell_{\varepsilon,h}(\rho;b).
 \label{eq:v30-canonical-exit-time}
\end{equation}
The lifetime is measurable by V29, so \(t_e\) is a permissible stopping map
on \(S_{\rm short}\); extend it by \(t_e=b\) off that set.
Classify an exit as total if
\(H_\rho^\Sigma(t_e)\le h\); every remaining exit has
\(|\mathcal Y_\rho(t_e)|\ge\varepsilon\).  Denote these two measurable sets
by \(S_{\rm tot}\) and \(S_{\rm pol}\).

\begin{theorem}[Probability-marked capacity of all recent exits]
\label{thm:v30-marked-entry-capacity}
For the canonical mark \(\mu_K\),
\begin{align}
 \mu_K(S_{\rm long})
 &\le\frac{E_*}{4\nu hK^2\tau_0},
 \label{eq:v30-marked-long-capacity}\\
 \mu_K(S_{\rm tot})
 &\le\frac{2E_*}{\kappa h},
 \label{eq:v30-marked-total-entry-capacity}\\
 \mu_K(S_{\rm pol})
 &\le\frac{14E_*}{(1-\beta)\varepsilon}.
 \label{eq:v30-marked-polarization-capacity}
\end{align}
Consequently
\begin{equation}
 \boxed{\quad
 \mu_K(S_{\rm in})
 \le\min\left\{1,
  \frac{E_*}{4\nu hK^2\tau_0}
  +\frac{2E_*}{\kappa h}
  +\frac{14E_*}{(1-\beta)\varepsilon}
  \right\}.
 \quad}
 \label{eq:v30-full-marked-entry-capacity}
\end{equation}
The right side contains no count or variation of exit times.
\end{theorem}

\begin{proof}
On \([K,\infty)\),
\begin{equation}
 \frac{K}{\rho^2}\le\frac{\rho}{K^2}.
 \label{eq:v30-mark-versus-radius}
\end{equation}
Apply V29's long-life weighted-width bound
\eqref{eq:v29-long-life-weighted-width} and then
\eqref{eq:v30-mark-versus-radius}; this proves
\eqref{eq:v30-marked-long-capacity}.

For \(\rho\in S_{\rm tot}\), V29's exact sum-sector calculation gives
\begin{equation}
 F_\Sigma^e(\rho)
 :=\int_{t_e(\rho)}^b\Phi_\rho(t)\,\dd t
 \ge\frac{\kappa h}{2}.
 \label{eq:v30-total-exit-flux-floor}
\end{equation}
Apply \eqref{eq:v30-positive-stopped-total} with
\(\gamma=\mathbf1_{S_{\rm tot}}\), \(\mu=\mu_K\), and
\(\tau=t_e\).  The marked integral of \(F_\Sigma^e\) is at most \(E_*\).
Comparison with \eqref{eq:v30-total-exit-flux-floor} proves
\eqref{eq:v30-marked-total-entry-capacity}.

For a polarization exit, put
\begin{align}
 q_e(\rho)&:=\int_{t_e(\rho)}^bd_\rho(t)\,\dd t,
 \label{eq:v30-exit-difference-action}\\
 p_e(\rho)&:=\nu\int_{t_e(\rho)}^b
       (V_\rho^++V_\rho^-)(t)\,\dd t,
 \label{eq:v30-exit-viscous-action}\\
 c_{\beta,\varepsilon}&:=\frac{(1-\beta)\varepsilon}{2}.
 \label{eq:v30-polarization-margin}
\end{align}
The V29 polarization alternative is
\begin{equation}
 |q_e(\rho)|+p_e(\rho)\ge c_{\beta,\varepsilon}.
 \label{eq:v30-polarization-exit-debit}
\end{equation}
Partition \(S_{\rm pol}\) into
\begin{align*}
 S_v&:=\{p_e\ge c_{\beta,\varepsilon}/2\},\\
 S_+&:=\{p_e<c_{\beta,\varepsilon}/2,
          \ q_e\ge c_{\beta,\varepsilon}/2\},\\
 S_-&:=\{p_e<c_{\beta,\varepsilon}/2,
          \ q_e\le-c_{\beta,\varepsilon}/2\}.
\end{align*}
These sets cover \(S_{\rm pol}\).  The stopped Leray action
\eqref{eq:v30-stopped-Leray-action} gives
\begin{equation}
 \mu_K(S_v)\le\frac{E_*}{c_{\beta,\varepsilon}}.
 \label{eq:v30-viscous-exit-capacity}
\end{equation}
Apply the stopped difference bound to \(S_+\) and \(S_-\) separately.  The
integral of \(q_e\) has one sign on each set, so
\begin{equation}
 \mu_K(S_+)+\mu_K(S_-)
 \le\frac{6E_*}{c_{\beta,\varepsilon}}.
 \label{eq:v30-signed-exit-capacity}
\end{equation}
Adding and substituting
\(c_{\beta,\varepsilon}=(1-\beta)\varepsilon/2\) proves
\eqref{eq:v30-marked-polarization-capacity}.  The three sets
\(S_{\rm long},S_{\rm tot},S_{\rm pol}\) partition \(S_{\rm in}\), and
\(\mu_K\) is a probability measure, proving
\eqref{eq:v30-full-marked-entry-capacity}.
\end{proof}

\begin{firewall}[The capacity constants do not separate a record]
\label{fw:v30-capacity-not-closure}
At the natural choices \(h\asymp E_*\),
\(\varepsilon\asymp E_*\), and
\(\tau_0\asymp(\nu K^2)^{-1}\), the right side of
\eqref{eq:v30-full-marked-entry-capacity} is merely an order-one constant
and can exceed the trivial probability bound.  The theorem closes the
\emph{multiplicity} of cutoff-dependent exit times, not the quantitative
separation needed for a record contradiction.  Treating the minimum with
one as a small bound would be false.
\end{firewall}

% =============================================================================
\section{The low-total recent flux is also assembled after marking}
\label{sec:v30-low-total-marked}
% =============================================================================

Fix \(0<\alpha<1\), \(0<\tau_0\le b\), and let
\(S\subset[K,\infty)\) be a measurable terminal set on which
\(H_\rho^\Sigma(b)>0\).  With the V29 relative lifetime, define
\begin{equation}
 S_{\rm new}^\alpha
 :=\{\rho\in S:\ell_\alpha^\Sigma(\rho;b)<\tau_0\},
 \qquad
 t_\alpha(\rho):=b-\ell_\alpha^\Sigma(\rho;b).
 \label{eq:v30-relative-new-set}
\end{equation}
Extend \(t_\alpha\) measurably by \(b\) off \(S\).

\begin{proposition}[Marked assembly of the low-total replenishment branch]
\label{prop:v30-low-total-marked-entry}
For every probability cutoff mark \(\mu\) supported on \(S\),
\begin{align}
 0&\le
 \int_{S_{\rm new}^\alpha}
   \int_{t_\alpha(\rho)}^b\Phi_\rho(t)\,\dd t\,\dd\mu(\rho)
 \le E_*,
 \label{eq:v30-low-total-marked-flux}\\
 \frac{1-\alpha}{2}
 \int_{S_{\rm new}^\alpha}H_\rho^\Sigma(b)\,\dd\mu(\rho)
 &\le
 \int_{S_{\rm new}^\alpha}
   \int_{t_\alpha(\rho)}^b\Phi_\rho(t)\,\dd t\,\dd\mu(\rho).
 \label{eq:v30-low-total-marked-floor}
\end{align}
Thus the continuum of recent replenishment intervals incurs no marked
incidence factor.  The resulting upper bound
\begin{equation}
 \int_{S_{\rm new}^\alpha}H_\rho^\Sigma(b)\,\dd\mu(\rho)
 \le\frac{2E_*}{1-\alpha}
 \label{eq:v30-low-total-marked-tail-bound}
\end{equation}
is recorded only as an exact consequence; it is not smaller than the direct
energy cap when \(\mu\) is a probability mark.
\end{proposition}

\begin{proof}
At the canonical exit time,
\(H_\rho^\Sigma(t_\alpha)\le
\alpha H_\rho^\Sigma(b)\).  The integrated sum-sector balance proves
pointwise both positivity and V29's lower bound
\eqref{eq:v29-relative-recent-influx}.  Integrate it against \(\mu\) to get
\eqref{eq:v30-low-total-marked-floor}.  Apply
\eqref{eq:v30-positive-stopped-total} with
\(\gamma=\mathbf1_{S_{\rm new}^\alpha}\) to get the upper bound in
\eqref{eq:v30-low-total-marked-flux}.  Division proves
\eqref{eq:v30-low-total-marked-tail-bound}.  Finally
\(H_\rho^\Sigma(b)\le E_*\) shows why the last displayed estimate is not a
gain in size.
\end{proof}

% =============================================================================
\section{A sharp no-go for globally unmarking the cutoff half-line}
\label{sec:v30-unmarking-no-go}
% =============================================================================

The preceding theorems are useful only if their normalization is kept
visible.  The next result characterizes why a single positive Leray mark
cannot be removed on all high cutoffs.

Let \(w:[K,\infty)\to[0,\infty)\) be locally integrable and set
\begin{equation}
 W_K(\lambda):=\int_K^\lambda w(\rho)\,\dd\rho,
 \qquad \lambda\ge K.
 \label{eq:v30-weight-primitive}
\end{equation}

\begin{theorem}[Leray-admissible radial weights are exactly finite marks]
\label{thm:v30-finite-mark-characterization}
For every smooth Fourier field,
\begin{align}
 \int_K^\infty w(\rho)H_\rho^\Sigma(t)\,\dd\rho
 &=\sum_{|k|>K}\sum_s
    W_K(|k|)|\widehat u^s(k,t)|^2,
 \label{eq:v30-general-weight-tail-symbol}\\
 \int_K^\infty w(\rho)(V_\rho^++V_\rho^-)(t)\,\dd\rho
 &=\sum_{|k|>K}\sum_s
    |k|^2W_K(|k|)|\widehat u^s(k,t)|^2.
 \label{eq:v30-general-weight-dissipation-symbol}
\end{align}
Consequently, a constant \(C_w<\infty\) satisfies
\begin{equation}
 \int_K^\infty w(\rho)(V_\rho^++V_\rho^-)(t)\,\dd\rho
 \le C_w\|\Lambda u(t)\|_2^2
 \quad\hbox{for every Fourier field}
 \label{eq:v30-Leray-admissible-weight}
\end{equation}
if and only if
\begin{equation}
 \sup_{\lambda>K}W_K(\lambda)le C_w.
 \label{eq:v30-bounded-weight-primitive}
\end{equation}
For nonnegative \(w\), this is equivalent to
\(w\in L^1([K,\infty))\) with
\(\int_K^\infty w\le C_w\).
\end{theorem}

\begin{proof}
Tonelli assigns a Fourier atom of radius \(\lambda\) the radial multiplier
\(\int_K^\lambda w=W_K(\lambda)\), proving both symbol identities.  If
\eqref{eq:v30-bounded-weight-primitive} holds, insert it into
\eqref{eq:v30-general-weight-dissipation-symbol} to prove
\eqref{eq:v30-Leray-admissible-weight}.  Conversely, test
\eqref{eq:v30-Leray-admissible-weight} on a single nonzero helical Fourier
atom at radius \(n\), using the torus modes \((n,0,0)\) for every integer
\(n>K\).  After cancelling its positive \(n^2\)-weighted energy, one
obtains \(W_K(n)\le C_w\).  Since \(W_K\) is nondecreasing and the
integers are unbounded, \(W_K(\lambda)\le C_w\) for every real
\(\lambda\ge K\).  This proves necessity.  Monotone convergence identifies
the supremum with \(\int_K^\infty w\).
\end{proof}

\begin{corollary}[No positive Leray mark dominates unweighted critical
incidence]
\label{cor:v30-no-global-unmarking}
There is no nonnegative locally integrable \(w\) and finite constants
\(C_w,C_u\) for which both
\eqref{eq:v30-Leray-admissible-weight} holds and
\begin{equation}
 \int_K^\infty f(\rho)\,\dd\rho
 \le C_u\int_K^\infty w(\rho)f(\rho)\,\dd\rho
 \label{eq:v30-global-unmarking-demand}
\end{equation}
holds for every nonnegative compactly supported measurable \(f\).
\end{corollary}

\begin{proof}
Testing \eqref{eq:v30-global-unmarking-demand} on indicator functions of
arbitrary finite-measure sets implies
\(w(\rho)\ge C_u^{-1}\) almost everywhere on \([K,\infty)\).  Hence
\(\int_K^\infty w=\infty\), contradicting
Theorem~\ref{thm:v30-finite-mark-characterization}.
\end{proof}

\begin{proposition}[The one-atom concentration test saturates the unmarking
loss]
\label{prop:v30-one-atom-unmarking-test}
Let a helical Fourier atom of energy \(e>0\) lie at an admissible torus
radius \(\lambda>K\).
Then its unweighted and canonically marked tail stocks above \(K\) are
\begin{align}
 \int_K^\infty H_\rho^\Sigma\,\dd\rho
 &=e(\lambda-K),
 \label{eq:v30-one-atom-unweighted-stock}\\
 \int_K^\infty H_\rho^\Sigma\,\dd\mu_K(\rho)
 &=e\left(1-\frac K\lambda\right)\le e.
 \label{eq:v30-one-atom-marked-stock}
\end{align}
Their ratio is exactly \(\lambda\).  Thus bounded energy and bounded marked
stock permit arbitrarily large unweighted critical-tail multiplicity as
\(\lambda/K\to\infty\).
\end{proposition}

\begin{proof}
The atom belongs to the high tail for \(K\le\rho<\lambda\).  Integrate its
constant height first against Lebesgue measure and then against
\eqref{eq:v30-canonical-mark}.
\end{proof}

\begin{remark}[This is a proof-strategy obstruction, not a singularity]
\label{rem:v30-no-go-scope}
Proposition~\ref{prop:v30-one-atom-unmarking-test} is a kinematic spectral
test.  It neither constructs a nonlinear transfer nor a singular NS
solution.  It proves that energy, Leray dissipation, and a single normalized
cutoff mark cannot by themselves recover the unweighted half-line
multiplicity.  This is consistent with the concentrated families already
used in V5 and V9--V14; it does not repeat their dynamical statements.
\end{remark}

\begin{proposition}[Exact cost of unmarking one finite annulus]
\label{prop:v30-annular-unmarking}
Let \(\Lambda>1\), \(B=[K,\Lambda K]\), and let \(f\ge0\) be measurable.
Then
\begin{equation}
 \boxed{\quad
 \int_Bf(\rho)\,\dd\rho
 \le\Lambda^2K\int_Bf(\rho)\,\dd\mu_K(\rho).
 \quad}
 \label{eq:v30-annular-unmarking}
\end{equation}
For the nonnegative recent total-entry flux
\(F_\Sigma^e\) in \eqref{eq:v30-total-exit-flux-floor},
\begin{equation}
 \int_{B\cap S_{\rm tot}}F_\Sigma^e(\rho)\,\dd\rho
 \le\Lambda^2KE_*.
 \label{eq:v30-annular-total-entry-bound}
\end{equation}
At the record scale \(K=R^2/E_*\), the right side is
\(\Lambda^2R^2\), so the conversion is critical but not separating.
\end{proposition}

\begin{proof}
On \(B\), the canonical density obeys
\(K/\rho^2\ge1/(\Lambda^2K)\), proving
\eqref{eq:v30-annular-unmarking}.  The flux is nonnegative on
\(S_{\rm tot}\), and its marked integral is at most \(E_*\) by the proof of
\eqref{eq:v30-marked-total-entry-capacity}.  Apply the first inequality and
then substitute the record relation.
\end{proof}

\begin{firewall}[Why marked control cannot be quoted as critical control]
\label{fw:v30-no-silent-unmarking}
Theorems~\ref{thm:v30-stopped-mark-balance} and
\ref{thm:v30-marked-entry-capacity} close a probability-normalized incidence
problem.  Corollary~\ref{cor:v30-no-global-unmarking} proves that no positive
Leray-admissible weight can dominate the unweighted half-line.  On a finite
annulus the conversion costs one factor of its radius, and at
\(K=R^2/E_*\) that factor returns an order-\(R^2\) bound.  Any claim that V30
already bounds the unweighted hard residual would therefore be false.
\end{firewall}

% =============================================================================
\section{Late companion drift during V30 validation}
\label{sec:v30-late-companion-drift}
% =============================================================================

After the V30 stopped-mark mathematics had been assembled, the final
filesystem check detected two concurrently completed companion appendices.
They were read before V30 was promoted a second time.  The later checkpoint is
\begin{align}
 &\texttt{Renewal\_SM\_Einstein\_Continuum\_Research\_Notes\_V55.tex},
 \notag\\[-2pt]
 &\qquad\text{SHA--256 }
 \texttt{bac7765373dedc4cc99a808d30a53b029776d567c7c3130abd92c1fdce328e1e},
 \label{eq:v30-late-SM-checkpoint}\\
 &\texttt{Renewal\_Yang\_Mills\_Mass\_Gap\_Research\_Notes\_V33.tex},
 \notag\\[-2pt]
 &\qquad\text{SHA--256 }
 \texttt{2342c192d5784b91c5ba4ae144f80500f42166f624a027b39bb347139d141532}.
 \label{eq:v30-late-YM-checkpoint}
\end{align}

SM V55 audits NS V29 and YM V32, derives the exact principal three-port card
for its gravitational problem, and proves that every constant positive Ohmic
matrix still has a stable-cone zero or a glancing rank defect.  Its own
firewall leaves frequency-dependent boundary realizations open.  This new
no-go concerns an Einstein boundary determinant and supplies no estimate for
NS cutoff marks, stopping times, or fluxes.

YM V33 places its five ternary terms on one representation carrier, proves
quadratic damping of paired high-to-low Racah leakage, and obtains genuine
marks for product-singlet formation and private-link tails.  It also proves
that unitary recoupling alone cannot manufacture a normalized incidence mark.
The nontrivial cold-intermediate formation-mark inequality remains open, so
neither its uniform ternary atlas nor the Yang--Mills mass gap is proved.

\begin{theorem}[Late-drift transfer decision]
\label{thm:v30-late-drift-decision}
The new companion appendices do not alter any estimate proved in V30 and do
not close the NS unmarking gate.  They sharpen the V31 acquisition rule as
follows: an annular NS probability mark must be produced or lifted by an
actual NS incidence mechanism---for example a heat/Duhamel parent factor, a
signed hinge identity, or a moving-tail telescope.  A unitary change of
Fourier coordinates, a radial reparameterization, or an abstract tree count
cannot by itself turn the unweighted critical measure into normalized marks.
\end{theorem}

\begin{proof}
SM V55 has a different operator, state space, and stable-frequency
determinant, so no hypothesis of Theorem~\ref{thm:v30-stopped-mark-balance}
is supplied by its Ohmic calculation.  YM V33's proved marks arise from Schur
dimension ratios and exact Wilson multipliers, neither of which occurs in the
NS sector-tail balance.  Its unitarity counterexample is type-correct only as
a logical warning: norm-preserving recoupling does not create a small
incidence factor.  Therefore the V30 NS estimates remain direct consequences
of their own PDE identities, while V31 must derive any mark-lifting gain from
an identified NS source-to-tail map.
\end{proof}

\begin{assessment}[Direction after the late drift]
\label{ass:v30-late-drift-direction}
The late companion results reinforce rather than replace the V30 frontier.
Probability marking closes arbitrary stopping-time multiplicity, but the mark
cannot be removed by algebraic relabeling.  The next version must locate a
physical NS formation coefficient connecting a recent nonlinear ancestor to
one annular cutoff mark.  If no such coefficient survives the full signed
assembly, the correct output is a finite far-tail witness, not an assumed
normalization.
\end{assessment}

% =============================================================================
\section{Integrated V30 conclusion and next acquisition order}
\label{sec:v30-conclusion}
% =============================================================================

\begin{theorem}[What V30 proves]
\label{thm:v30-conclusion}
Relative to V1--V29, the following are proved.
\begin{enumerate}[label=\textup{(V30.\arabic*)},leftmargin=6.2em]
\item The compulsory companion audit records the initial SM V54/YM V32
      checkpoint in \eqref{eq:v30-SM-checkpoint}--%
      \eqref{eq:v30-YM-checkpoint} and the concurrent SM V55/YM V33 drift
      in \eqref{eq:v30-late-SM-checkpoint}--%
      \eqref{eq:v30-late-YM-checkpoint}.  It imports no mismatched estimate.
\item Every finite positive cutoff mark, bounded frozen radial multiplier,
      and measurable cutoff-dependent stopping map obeys the exact stopped
      total and difference balances
      \eqref{eq:v30-exact-stopped-total}--%
      \eqref{eq:v30-exact-stopped-difference}.
\item Both stopped nonlinear fluxes are bounded by
      \(3M_\mu E_*/2\), independently of the number or variation of stopping
      times.  This closes the probability-normalized moving-time incidence
      loss left by V29.
\item The canonical half-line probability mark
      \(K\rho^{-2}\,\dd\rho\) has the exact spectral symbols
      \eqref{eq:v30-marked-tail-energy}--%
      \eqref{eq:v30-marked-tail-dissipation}; the entire frozen hard
      selector and its hard-minus-clipped residual both satisfy
      \eqref{eq:v30-marked-hard-residual}.
\item V29's long-lived, recent-total, and recent-polarization branches have
      the marked capacities
      \eqref{eq:v30-marked-long-capacity}--%
      \eqref{eq:v30-marked-polarization-capacity}.  Their constants contain
      no exit-time incidence count but do not give record separation at the
      natural thresholds.
\item The low-total recent-replenishment intervals likewise assemble under
      every probability mark, as in
      \eqref{eq:v30-low-total-marked-flux}--%
      \eqref{eq:v30-low-total-marked-floor}; the resulting tail bound is
      honestly recorded as non-improving.
\item A nonnegative radial weight is uniformly Leray-admissible for all
      Fourier fields exactly when it has finite half-line mass.  Such a weight
      cannot dominate unweighted radial integration; see
      Theorem~\ref{thm:v30-finite-mark-characterization} and
      Corollary~\ref{cor:v30-no-global-unmarking}.
\item The one-atom test makes the missing radius factor exact, while finite-
      annulus unmarking costs \(\Lambda^2K\).  At the canonical record scale
      this returns an order-\(R^2\), nonseparating bound.
\end{enumerate}
The open gate is no longer arbitrary variation of recent-entry times after
normalization.  It is the recovery of unweighted critical radial mass from a
family of locally marked annuli without paying the radius factor independently
at every scale.  The full Navier--Stokes stability/global-regularity theorem
remains unproved.
\end{theorem}

\begin{proof}
Item \textup{(V30.1)} is
Theorems~\ref{thm:v30-companion-decision} and
\ref{thm:v30-late-drift-decision}.  Items
\textup{(V30.2)--(V30.3)} are
Theorem~\ref{thm:v30-stopped-mark-balance}.  Item \textup{(V30.4)} is
Proposition~\ref{prop:v30-canonical-mark-symbol} and
Corollary~\ref{cor:v30-marked-hard-residual}.  Item \textup{(V30.5)} is
Theorem~\ref{thm:v30-marked-entry-capacity} and
Firewall~\ref{fw:v30-capacity-not-closure}.  Item \textup{(V30.6)} is
Proposition~\ref{prop:v30-low-total-marked-entry}.  Item
\textup{(V30.7)} is Theorem~\ref{thm:v30-finite-mark-characterization} and
Corollary~\ref{cor:v30-no-global-unmarking}.  Item \textup{(V30.8)} is
Propositions~\ref{prop:v30-one-atom-unmarking-test} and
\ref{prop:v30-annular-unmarking}.
\end{proof}

\begin{programme}[V31 dyadic marked ancestry and far-tail gate]
\label{prog:v31}
\begin{enumerate}[label=\textup{(AC\arabic*)},leftmargin=4.8em]
\item Replace the single half-line mark by probability marks on geometric
      annuli \([2^nK,2^{n+1}K]\).  Derive the complete stopped balance on all
      annuli first; do not sum the annular radius costs termwise.
\item Test whether V19's exact moving-tail telescope and V27's hinge-current
      BV identity combine the annular endpoint terms before the factors
      \(2^nK\) are paid.  Any leftover swept-band energy must remain explicit.
\item On recent total-entry branches, insert the already proved shell and
      parabolic Duhamel identities from V7 and V10--V21.  Prove an actual
      NS formation-mark inequality in which a heat factor or common parent
      lifts one annular probability mark and compensates the unmarking radius;
      a unitary change of basis or radial relabeling is not such a proof.
\item On recent polarization branches, preserve the sign partition used in
      \eqref{eq:v30-signed-exit-capacity}.  Assemble positive and negative
      stopped difference actions through the radial hinge current before
      taking absolute values across annuli.
\item Use Proposition~\ref{prop:v30-one-atom-unmarking-test} as a mandatory
      far-tail obstruction.  A proposed estimate depending only on endpoint
      energy and a normalized mark is already refuted by that atom.
\item Combine the persistent branch with
      \eqref{eq:v29-cofinal-persistence-sum}; its parabolic durations are
      globally paid and must not be charged again in the recent-entry tree.
\item Normalize at \(K=R^2/E_*\) and compare every surviving constant with
      the V24 overlap--separation lower bound.  An order-\(R^2\) annular upper
      bound is not a contradiction.
\item Retain the hard--soft cushion, the low-total initial and viscous pieces,
      and every unresolved triad orientation.  The next compulsory companion
      review is V35, not V31.
\end{enumerate}
\end{programme}

\begin{assessment}[Position after V30]
The moving stopping-time continuum is no longer an unspecified multiplicity.
A probability mark turns every collection of radius-dependent entry intervals
into one exact endpoint-energy/Leray balance, and it controls the hard residual
without smoothing.  This conclusion is genuinely independent of how often or
how irregularly the entry time changes with cutoff.

The same theorem also reveals its own sharp boundary.  Every globally
Leray-admissible positive radial mark is finite, and no finite mark dominates
Lebesgue measure on the infinite cutoff half-line.  A single far-frequency
atom realizes the missing radius factor.  The next attack must therefore
assemble locally marked annuli with heat ancestry, moving-tail telescoping, or
a signed hinge cancellation before unmarking.  Merely summing the V30 marked
bounds returns critical size.  The full stability/global-regularity objective
remains open.
\end{assessment}

% =============================================================================
\section*{Version V30 append-only record}
\addcontentsline{toc}{section}{Version V30 append-only record}
% =============================================================================

The complete V29 mathematical source was retained before this continuation.
Its SHA--256 digest was
\begin{center}\scriptsize\ttfamily
77e7087398ae9ca94469a91464cee3e0c1780d2f7513cbc4620d9881dd795a05.
\end{center}
V30 performed the compulsory SM V54/YM V32 audit and a late drift review of
SM V55/YM V33, imported only the proved normalization discipline, and derived
the NS estimates independently.  It
proved exact stopped total/difference balances for arbitrary measurable
cutoff-dependent lower times under every finite radial mark, with one
endpoint-energy/Leray bound and no stopping-time incidence loss.  It applied
the result to the hard residual and both V29 recent-entry branches.  Finally it
proved that every globally Leray-admissible positive cutoff weight is a finite
mark and therefore cannot dominate unweighted critical radial integration;
the exact remaining gate is annular unmarking with far-tail ancestry.  No full
regularity claim is made.
% =============================================================================
% VERSION V31 APPEND
% =============================================================================

\clearpage
\section{V31: flat dyadic marks and the exact heat-formation ladder}
\label{sec:v31-entry}
% =============================================================================

V30 proved that a probability mark removes the multiplicity of arbitrary
cutoff-dependent stopping times, and also proved that no finite positive mark
can dominate the unweighted cutoff half-line.  The next question is therefore
not whether one can choose a more ingenious probability density.  It is
whether the Navier--Stokes evolution itself supplies the first radial moment
which is lost when a probability mark is removed.

This version makes that question exact.  Uniform probability measure on each
dyadic annulus reconstructs the unweighted radial integral with no comparison
constant at all.  The reconstruction is a first moment in the dyadic radius.
For source-born velocity, the heat equation then gives a three-rung formation
ladder:
\[
 \begin{array}{c|c}
 \text{response stock}&\text{source action required without a scale loss}\\
 \hline
 K_n^{-1}\|r_n\|_2^2&A^{-3/4}\BNS(u)\\
 \|r_n\|_2^2&A^{-1/2}\BNS(u)\\
 K_n\|r_n\|_2^2&A^{-1/4}\BNS(u).
 \end{array}
\]
The last row is precisely the first moment needed to unmark the critical
radial integral.  It is valid, including for different lower stopping times
in different shells, but its right side is the critical source action already
identified in V7.  Combining it with the genuine high-parent estimate of V14
removes dyadic parent multiplicity and returns exactly
\(\int\|\Lambda u\|_2^4\), the V1 continuation stock.  Thus V31 obtains an
actual NS formation-mark inequality and identifies, without a hidden
frequency factor, why the currently available payment is circular.

No companion-note review is due in V31.  The next compulsory SM/YM
synchronization remains V35.

% =============================================================================
\section{Flat probability marks on the dyadic cutoff partition}
\label{sec:v31-flat-dyadic-marks}
% =============================================================================

Fix \(K>0\) and put
\begin{equation}
 K_n:=2^nK,\qquad
 \mathcal A_n:=[K_n,K_{n+1}),\qquad n\ge0.
 \label{eq:v31-dyadic-cutoff-annuli}
\end{equation}
On \(\mathcal A_n\), define the flat probability mark
\begin{equation}
 \dd\mu_n^{\rm fl}(\rho)
 :=\mathbf1_{\mathcal A_n}(\rho)\frac{\dd\rho}{K_n},
 \qquad
 \mu_n^{\rm fl}(\mathcal A_n)=1.
 \label{eq:v31-flat-probability-mark}
\end{equation}
For an arbitrary square-integrable divergence-free Fourier field \(v\), write
\begin{equation}
 H_\rho[v]:=\|Q_\rho v\|_2^2,\qquad
 \mathcal M_n[v]
 :=\int_{\mathcal A_n}H_\rho[v]\,\dd\mu_n^{\rm fl}(\rho).
 \label{eq:v31-flat-marked-tail}
\end{equation}

\begin{proposition}[The flat annular mark has an exact ramp symbol]
\label{prop:v31-flat-mark-symbol}
For every \(n\ge0\),
\begin{align}
 \mathcal M_n[v]
 &=\sum_{k\ne0}m_n(|k|)|\widehat v(k)|^2,
 \label{eq:v31-flat-symbol}\\
 m_n(\lambda)
 &:=
 \begin{cases}
  0,&0<\lambda\le K_n,\\[2pt]
  (\lambda-K_n)/K_n,&K_n<\lambda<K_{n+1},\\[2pt]
  1,&\lambda\ge K_{n+1}.
 \end{cases}
 \label{eq:v31-flat-ramp}
\end{align}
Moreover,
\begin{equation}
 \boxed{\quad
 K_n\mathcal M_n[v]
 =\int_{K_n}^{K_{n+1}}H_\rho[v]\,\dd\rho .
 \quad}
 \label{eq:v31-exact-annular-unmarking}
\end{equation}
Thus flat annular unmarking costs exactly \(K_n\), not the
\(4K_n\) comparison obtained by restricting the V30 inverse-square mark.
\end{proposition}

\begin{proof}
A Fourier atom of radius \(\lambda\) belongs to \(Q_\rho v\) precisely when
\(\rho<\lambda\).  Its \(\mu_n^{\rm fl}\)-mass is therefore
\[
 \frac1{K_n}
 \left|\mathcal A_n\cap[0,\lambda)\right|,
\]
which is exactly \eqref{eq:v31-flat-ramp}.  Tonelli proves
\eqref{eq:v31-flat-symbol}.  Multiplication by \(K_n\), or directly the
definition \eqref{eq:v31-flat-probability-mark}, proves
\eqref{eq:v31-exact-annular-unmarking}.
\end{proof}

\begin{theorem}[Dyadic marked first moment is exactly critical coarea]
\label{thm:v31-dyadic-coarea}
Define
\begin{equation}
 \mathcal C_K[v]
 :=\int_K^\infty H_\rho[v]\,\dd\rho.
 \label{eq:v31-critical-coarea}
\end{equation}
Then
\begin{align}
 \boxed{\quad
 \mathcal C_K[v]
 =\sum_{n=0}^\infty K_n\mathcal M_n[v]
 =\sum_{|k|>K}(|k|-K)|\widehat v(k)|^2.
 \quad}
 \label{eq:v31-exact-dyadic-first-moment}
\end{align}
Let
\begin{equation}
 \Pi_n:=\mathbf1_{(K_n,K_{n+1}]}(\Lambda).
 \label{eq:v31-dyadic-shell-projector}
\end{equation}
The far part obeys the two-sided shell equivalence
\begin{equation}
 \frac12\sum_{n=1}^\infty K_n\|\Pi_nv\|_2^2
 \le
 \mathcal C_K[Q_{2K}v]
 \le
 2\sum_{n=1}^\infty K_n\|\Pi_nv\|_2^2.
 \label{eq:v31-far-coarea-shell-equivalence}
\end{equation}
\end{theorem}

\begin{proof}
The annuli \(\mathcal A_n\) partition \([K,\infty)\) up to endpoints.
Sum \eqref{eq:v31-exact-annular-unmarking} and use monotone convergence.
The Fourier formula follows by Tonelli:
\[
 \int_K^\infty\mathbf1_{\{\rho<|k|\}}\,\dd\rho
 =(|k|-K)_+.
\]
If \(|k|\in(K_n,K_{n+1}]\) with \(n\ge1\), then
\[
 \frac{K_n}{2}
 \le K_n-K
 <|k|-K
 \le |k|
 \le2K_n.
\]
Sum this pointwise comparison over the orthogonal shells.
\end{proof}

\begin{proposition}[A far atom fills its entire lower dyadic mark chain]
\label{prop:v31-far-atom-mark-chain}
Suppose a Fourier atom of energy \(e>0\) lies at radius
\(\lambda=2^mK\), \(m\ge1\).  Then
\begin{equation}
 \mathcal M_n=e\quad(0\le n<m),\qquad
 \mathcal M_n=0\quad(n\ge m),
 \label{eq:v31-atom-mark-chain}
\end{equation}
with the harmless endpoint convention in
\eqref{eq:v31-flat-ramp}, and hence
\begin{equation}
 \sum_{n=0}^\infty K_n\mathcal M_n
 =e\sum_{n=0}^{m-1}2^nK
 =e(\lambda-K).
 \label{eq:v31-atom-first-moment}
\end{equation}
Thus dyadic marking changes the V30 one-atom obstruction into an exact
ancestry chain: the remote atom has unit marked mass at every lower
generation, and its geometric first moment is its unweighted critical mass.
\end{proposition}

\begin{proof}
For \(n<m\), one has \(\lambda\ge K_{n+1}\), so
\(m_n(\lambda)=1\).  For \(n\ge m\), one has
\(\lambda\le K_n\), so \(m_n(\lambda)=0\).  Sum the geometric series.
\end{proof}

\begin{firewall}[Dyadic probability is not dyadic first-moment control]
\label{fw:v31-probability-not-first-moment}
The V30 theorem gives a uniform estimate on each
\(\mathcal M_n\)-normalized balance.  Equation
\eqref{eq:v31-exact-dyadic-first-moment} requires the unbounded coefficients
\(K_n\).  Proposition~\ref{prop:v31-far-atom-mark-chain} proves that these
coefficients are not an artefact of a rough comparison: one atom saturates
the whole geometric sum.  Consequently neither a bound uniform in \(n\) nor
an unweighted count of the probability marks controls critical coarea.
\end{firewall}

% =============================================================================
\section{Stopped balances on all flat annuli}
\label{sec:v31-flat-stopped-balances}
% =============================================================================

For every \(n\), let
\(\tau_n:\mathcal A_n\to[a,b]\) and
\(\gamma_n:\mathcal A_n\to[-1,1]\) be measurable.  Define
\begin{align}
 \mathfrak f_{\Sigma,n}
 &:=
 \frac1{K_n}\int_{\mathcal A_n}\gamma_n(\rho)
   \int_{\tau_n(\rho)}^b\Phi_\rho(t)\,\dd t\,\dd\rho,
 \label{eq:v31-normalized-annular-total-flux}\\
 \mathfrak f_{\Delta,n}
 &:=
 \frac1{K_n}\int_{\mathcal A_n}\gamma_n(\rho)
   \int_{\tau_n(\rho)}^bd_\rho(t)\,\dd t\,\dd\rho.
 \label{eq:v31-normalized-annular-difference-flux}
\end{align}

\begin{theorem}[The complete dyadic stopped family has an
\(\ell^1\)-dual bound]
\label{thm:v31-dyadic-stopped-family}
For every \(n\ge0\),
\begin{equation}
 \max\{|\mathfrak f_{\Sigma,n}|,|\mathfrak f_{\Delta,n}|\}
 \le\frac32E_*.
 \label{eq:v31-one-annulus-stopped-bound}
\end{equation}
More generally, if \(c=(c_n)_{n\ge0}\in\ell^1\), then
\begin{align}
 \left|\sum_{n=0}^\infty c_n\mathfrak f_{\Sigma,n}\right|
 &\le\frac32E_*\sum_{n=0}^\infty|c_n|,
 \label{eq:v31-l1-total-family}\\
 \left|\sum_{n=0}^\infty c_n\mathfrak f_{\Delta,n}\right|
 &\le\frac32E_*\sum_{n=0}^\infty|c_n|.
 \label{eq:v31-l1-difference-family}
\end{align}
The corresponding raw annular actions are exactly
\begin{equation}
 K_n\mathfrak f_{\Sigma,n},\qquad
 K_n\mathfrak f_{\Delta,n},
 \label{eq:v31-raw-annular-actions}
\end{equation}
and are therefore bounded separately by \(3K_nE_*/2\).
\end{theorem}

\begin{proof}
Apply Theorem~\ref{thm:v30-stopped-mark-balance} with
\(\mu=\mu_n^{\rm fl}\).  Its mass is one, proving the individual estimates.
For a finitely supported sequence \(c\), use the finite positive mark
\[
 \mu:=\sum_n|c_n|\mu_n^{\rm fl}
\]
and, on \(\mathcal A_n\), use the frozen multiplier
\(\operatorname{sgn}(c_n)\gamma_n\).  The mass of \(\mu\) is
\(\sum_n|c_n|\), so the same theorem gives
\eqref{eq:v31-l1-total-family}--%
\eqref{eq:v31-l1-difference-family}.  Approximate a general
\(\ell^1\) sequence by its finite truncations.  The raw identities follow
from \(\dd\mu_n^{\rm fl}=\dd\rho/K_n\).
\end{proof}

\begin{remark}[Positive total influx improves the constant, not the scale]
\label{rem:v31-positive-flat-annulus}
If \(\gamma_n\ge0\) and
\[
 H_\rho^\Sigma(b)\ge
 H_\rho^\Sigma(\tau_n(\rho))
\]
on its support, the positive part of
Theorem~\ref{thm:v30-stopped-mark-balance} gives
\[
 0\le K_n\mathfrak f_{\Sigma,n}\le K_nE_*.
\]
The coefficient \(K_n\) is still the exact unmarking coefficient in
\eqref{eq:v31-exact-annular-unmarking}.
\end{remark}

\begin{corollary}[Lebesgue-width form of the V29 recent-exit capacities]
\label{cor:v31-flat-exit-widths}
Use the sets \(S_{\rm long},S_{\rm tot},S_{\rm pol}\) and parameters of
Theorem~\ref{thm:v30-marked-entry-capacity}.  On every
\(\mathcal A_n\),
\begin{align}
 |S_{\rm long}\cap\mathcal A_n|
 &\le\frac{E_*}{4\nu hK_n\tau_0},
 \label{eq:v31-flat-long-width}\\
 |S_{\rm tot}\cap\mathcal A_n|
 &\le\frac{2K_nE_*}{\kappa h},
 \label{eq:v31-flat-total-width}\\
 |S_{\rm pol}\cap\mathcal A_n|
 &\le\frac{14K_nE_*}{(1-\beta)\varepsilon}.
 \label{eq:v31-flat-polarization-width}
\end{align}
At \(h\asymp\varepsilon\asymp E_*\) and
\(\tau_0\asymp(\nu K_n^2)^{-1}\), every right side is of order \(K_n\);
none is smaller than the annular width.
\end{corollary}

\begin{proof}
The first estimate is V29's width bound
\eqref{eq:v29-long-life-width}, applied with lower cutoff \(K_n\).
For the total and polarization exits, repeat the proof of
Theorem~\ref{thm:v30-marked-entry-capacity} with the probability mark
\(\mu_n^{\rm fl}\), restricted to \(\mathcal A_n\).  The signed
polarization partition \(S_v,S_+,S_-\) is unchanged.  This gives the V30
probability-capacity constants.  Finally
\[
 |S\cap\mathcal A_n|
 =K_n\mu_n^{\rm fl}(S\cap\mathcal A_n).
\]
The last assertion follows by substitution.
\end{proof}

\begin{assessment}[What flat marking settles]
\label{ass:v31-flat-marking-settles}
There is no longer any ambiguity about the local unmarking constant.
Probability marking completely removes the number and variation of stopping
times, while returning to Lebesgue cutoff measure costs exactly the dyadic
radius.  The remaining question is therefore whether source formation,
moving-tail telescoping, or signed sector cancellation pays the first moment
\(\sum_nK_n(\cdot)\) before the annuli are separated.
\end{assessment}

% =============================================================================
\section{The scale-matched heat-formation ladder}
\label{sec:v31-heat-formation}
% =============================================================================

Let \(a<b<T_*\), let \(\tau_n\in[a,b]\), and use the shell projectors
\(\Pi_n\) from \eqref{eq:v31-dyadic-shell-projector}.  Define the
source-born velocity response at the common terminal time \(b\) by
\begin{equation}
 r_n:=
 \int_{\tau_n}^b
 e^{-\nu A(b-t)}\Pi_n\BNS(u(t))\,\dd t,
 \qquad n\ge1,
 \label{eq:v31-shell-born-velocity}
\end{equation}
and, for \(1/4\le\sigma\le3/4\), define
\begin{equation}
 \mathscr S_{\sigma,n}
 :=
 \int_{\tau_n}^b
 \|A^{-\sigma}\Pi_n\BNS(u(t))\|_2^2\,\dd t.
 \label{eq:v31-shell-source-action}
\end{equation}

\begin{theorem}[Exact source regularity required by each formation weight]
\label{thm:v31-heat-formation-ladder}
For every \(1/4\le\sigma\le3/4\),
\begin{align}
 \boxed{\quad
 \sum_{n=1}^\infty
 K_n^{\,2-4\sigma}\|r_n\|_2^2
 \le\frac1{\nu}\sum_{n=1}^\infty\mathscr S_{\sigma,n}
 \le\frac1{\nu}
 \int_a^b\|A^{-\sigma}Q_{2K}\BNS(u(t))\|_2^2\,\dd t.
 \quad}
 \label{eq:v31-scale-matched-formation}\\
 \sum_{n=1}^\infty K_n\|r_n\|_2^2
 &\le
 \frac1{\nu}\sum_{n=1}^\infty
 K_n^{\,4\sigma-1}\mathscr S_{\sigma,n}.
 \label{eq:v31-first-moment-from-sigma}
\end{align}
The weighted source action in
\eqref{eq:v31-first-moment-from-sigma} is uniformly equivalent, shell by
shell, to the critical source action:
\begin{align}
 2^{-(4\sigma-1)}
 \int_{\tau_n}^b
 \|A^{-1/4}\Pi_n\BNS(u)\|_2^2\,\dd t
 &\le
 K_n^{\,4\sigma-1}\mathscr S_{\sigma,n}
 \notag\\
 &\le
 \int_{\tau_n}^b
 \|A^{-1/4}\Pi_n\BNS(u)\|_2^2\,\dd t.
 \label{eq:v31-source-weight-collapse}
\end{align}
Consequently,
\begin{equation}
 \boxed{\quad
 \sum_{n=1}^\infty K_n\|r_n\|_2^2
 \le\frac1{\nu}
 \int_a^b\|A^{-1/4}Q_{2K}\BNS(u(t))\|_2^2\,\dd t.
 \quad}
 \label{eq:v31-critical-formation-mark}
\end{equation}
\end{theorem}

\begin{proof}
Put
\[
 g_{\sigma,n}:=A^{-\sigma}\Pi_n\BNS(u).
\]
On a Fourier mode of radius
\(\lambda\in(K_n,K_{n+1}]\), formula
\eqref{eq:v31-shell-born-velocity} reads
\[
 \widehat r_n(\lambda)
 =
 \int_{\tau_n}^b
 e^{-\nu\lambda^2(b-t)}
 \lambda^{2\sigma}\widehat g_{\sigma,n}(\lambda,t)\,\dd t.
\]
Time Cauchy--Schwarz gives
\begin{align}
 |\widehat r_n(\lambda)|^2
 &\le
 \lambda^{4\sigma}
 \left(\int_{\tau_n}^b
 e^{-2\nu\lambda^2(b-t)}\,\dd t\right)
 \int_{\tau_n}^b
 |\widehat g_{\sigma,n}(\lambda,t)|^2\,\dd t
 \notag\\
 &\le
 \frac{\lambda^{4\sigma-2}}{2\nu}
 \int_{\tau_n}^b
 |\widehat g_{\sigma,n}(\lambda,t)|^2\,\dd t.
 \label{eq:v31-modal-heat-formation}
\end{align}
Since \(\lambda/K_n\in(1,2]\) and
\(4\sigma-2\in[-1,1]\),
\[
 \frac12
 K_n^{\,2-4\sigma}\lambda^{4\sigma-2}
 \le1.
\]
Sum \eqref{eq:v31-modal-heat-formation} over modes and shells to prove the
first inequality in \eqref{eq:v31-scale-matched-formation}.
At any fixed time, the surviving projectors
\(\Pi_n\mathbf1_{\{t\ge\tau_n\}}\) are still mutually orthogonal.
Tonelli therefore proves the second inequality.

Multiplying \eqref{eq:v31-modal-heat-formation} by \(K_n\) and using
\[
 \frac12K_n\lambda^{4\sigma-2}
 \le K_n^{4\sigma-1}
\]
proves \eqref{eq:v31-first-moment-from-sigma}.  Finally, on the same shell,
\begin{align*}
 \frac{
 K_n^{4\sigma-1}
 |\lambda^{-2\sigma}\widehat{\BNS(u)}(\lambda)|^2
 }{
 |\lambda^{-1/2}\widehat{\BNS(u)}(\lambda)|^2
 }
 &=
 \left(\frac{K_n}{\lambda}\right)^{4\sigma-1}.
\end{align*}
This ratio lies in
\([2^{-(4\sigma-1)},1]\), which proves
\eqref{eq:v31-source-weight-collapse}.  Sum its upper bound in \(n\),
retain only the active time intervals, and use orthogonality to obtain
\eqref{eq:v31-critical-formation-mark}.
\end{proof}

\begin{corollary}[The three rungs and the common-terminal critical response]
\label{cor:v31-three-rung-formation}
The endpoints \(\sigma=3/4,1/2,1/4\) in
\eqref{eq:v31-scale-matched-formation} give
\begin{align}
 \sum_{n\ge1}K_n^{-1}\|r_n\|_2^2
 &\le\frac1\nu
 \int_a^b\|A^{-3/4}Q_{2K}\BNS(u)\|_2^2\,\dd t,
 \label{eq:v31-weak-formation-rung}\\
 \sum_{n\ge1}\|r_n\|_2^2
 &\le\frac1\nu
 \int_a^b\|A^{-1/2}Q_{2K}\BNS(u)\|_2^2\,\dd t,
 \label{eq:v31-intermediate-formation-rung}\\
 \sum_{n\ge1}K_n\|r_n\|_2^2
 &\le\frac1\nu
 \int_a^b\|A^{-1/4}Q_{2K}\BNS(u)\|_2^2\,\dd t.
 \label{eq:v31-critical-formation-rung}
\end{align}
If
\[
 r^{\rm born}:=\sum_{n\ge1}r_n,
\]
then
\begin{equation}
 \boxed{\quad
 \mathcal C_K[r^{\rm born}]
 \le\frac2\nu
 \int_a^b\|A^{-1/4}Q_{2K}\BNS(u(t))\|_2^2\,\dd t.
 \quad}
 \label{eq:v31-born-critical-coarea}
\end{equation}
If all \(\tau_n=a\), Duhamel's formula and heat contraction also give
\begin{align}
 \sqrt{\mathcal C_K[Q_{2K}u(b)]}
 &\le
 \sqrt{\mathcal C_K[Q_{2K}u(a)]}
 \notag\\
 &\quad+
 \left(
  \frac2\nu
  \int_a^b\|A^{-1/4}Q_{2K}\BNS(u(t))\|_2^2\,\dd t
 \right)^{1/2}.
 \label{eq:v31-critical-coarea-Duhamel}
\end{align}
\end{corollary}

\begin{proof}
The three displayed rungs are direct substitutions.  The responses have
orthogonal shell supports, so
\eqref{eq:v31-far-coarea-shell-equivalence} and
\eqref{eq:v31-critical-formation-rung} prove
\eqref{eq:v31-born-critical-coarea}.  When all lower times agree,
\[
 Q_{2K}u(b)
 =
 e^{-\nu A(b-a)}Q_{2K}u(a)-r^{\rm born}.
\]
The square root of \(\mathcal C_K\) is a Hilbert norm on
\(\operatorname{Ran}Q_{2K}\).  Its Fourier multiplier commutes with the heat
semigroup, which is contractive in that norm.  Apply the triangle inequality
and \eqref{eq:v31-born-critical-coarea}.
\end{proof}

\begin{assessment}[The requested formation mark exists, at critical cost]
\label{ass:v31-formation-mark-status}
Equation~\eqref{eq:v31-critical-formation-mark} is an actual
Navier--Stokes source-to-tail formation-mark inequality.  It permits different
lower times in different dyadic shells, has no shell-counting factor, and
supplies exactly the dyadic first moment required by
\eqref{eq:v31-exact-dyadic-first-moment}.  Its price is
\(A^{-1/4}\BNS(u)\), not the globally energy-paid
\(A^{-3/4}\BNS(u)\).  Equation
\eqref{eq:v31-source-weight-collapse} proves that applying the weak source
and restoring the missing \(K_n^2\) is algebraically the same critical
source action.  A heat factor alone therefore does not create a subcritical
mark.
\end{assessment}

% =============================================================================
\section{High-parent lifting removes dyadic multiplicity and returns V1}
\label{sec:v31-high-parent-lifting}
% =============================================================================

The critical source in
\eqref{eq:v31-critical-formation-mark} can be rewritten through the genuine
high-parent identity, but this operation must be performed before summing the
shells.  Define the active-shell indicators
\begin{equation}
 \chi_n(t):=\mathbf1_{[\tau_n,b]}(t)
 \label{eq:v31-active-shell-indicator}
\end{equation}
and the restricted parent-formation occupation
\begin{equation}
 \mathscr Q_{\rm form}
 :=
 \int_a^b\|\Lambda u(t)\|_2^2
 \sum_{n\ge1}\chi_n(t)K_n^2
 H_{K_n/2}^\Sigma(t)\,\dd t.
 \label{eq:v31-restricted-parent-formation}
\end{equation}

\begin{lemma}[The dyadic high-parent first moment has no logarithmic loss]
\label{lem:v31-parent-geometric-sum}
At every smooth time,
\begin{equation}
 \boxed{\quad
 \sum_{n\ge1}K_n^2H_{K_n/2}^\Sigma(t)
 \le\frac{16}{3}\|\Lambda u(t)\|_2^2.
 \quad}
 \label{eq:v31-parent-geometric-sum}
\end{equation}
The same bound holds after any factors \(0\le\chi_n(t)\le1\) are inserted.
\end{lemma}

\begin{proof}
For a Fourier atom of radius \(\lambda\), its coefficient on the left is
\[
 \sum_{\substack{n\ge1\\K_n<2\lambda}}K_n^2.
\]
The \(K_n^2\) form a geometric progression of ratio four.  Its last term is
strictly smaller than \(4\lambda^2\), and the sum is at most \(4/3\) times
that last-term bound.  Thus its coefficient is at most
\(16\lambda^2/3\).  Sum over all Fourier atoms.  Removing terms by factors
between zero and one can only decrease the nonnegative sum.
\end{proof}

\begin{theorem}[The nonlinear high-parent formation mark is exactly
continuation-class]
\label{thm:v31-nonlinear-parent-formation}
With the responses \eqref{eq:v31-shell-born-velocity},
\begin{align}
 \sum_{n\ge1}K_n\|r_n\|_2^2
 &\le\frac{C_{\rm hp}^2}{\nu}\mathscr Q_{\rm form},
 \label{eq:v31-parent-formation-sharp}\\
 \mathscr Q_{\rm form}
 &\le\frac{16}{3}
 \int_a^b\|\Lambda u(t)\|_2^4\,\dd t,
 \label{eq:v31-parent-formation-continuation}\\
 \boxed{\quad
 \sum_{n\ge1}K_n\|r_n\|_2^2
 \le\frac{16C_{\rm hp}^2}{3\nu}
 \int_a^b\|\Lambda u(t)\|_2^4\,\dd t.
 \quad}
 \label{eq:v31-parent-formation-final}
\end{align}
Consequently,
\begin{equation}
 \mathcal C_K[r^{\rm born}]
 \le\frac{32C_{\rm hp}^2}{3\nu}
 \int_a^b\|\Lambda u(t)\|_2^4\,\dd t.
 \label{eq:v31-parent-critical-coarea}
\end{equation}
\end{theorem}

\begin{proof}
Because
\(\Pi_n=\Pi_nQ_{K_n}\), the proof of
Theorem~\ref{thm:v14-high-parent-weak} gives
\begin{equation}
 \|A^{-3/4}\Pi_n\BNS(u(t))\|_2^2
 \le
 C_{\rm hp}^2
 H_{K_n/2}^\Sigma(t)\|\Lambda u(t)\|_2^2.
 \label{eq:v31-shell-high-parent-source}
\end{equation}
Use \eqref{eq:v31-first-moment-from-sigma} at \(\sigma=3/4\), insert
\eqref{eq:v31-shell-high-parent-source}, sum first in \(n\), and then
integrate in time.  This proves
\eqref{eq:v31-parent-formation-sharp}.  Apply
Lemma~\ref{lem:v31-parent-geometric-sum} inside
\eqref{eq:v31-restricted-parent-formation} to obtain
\eqref{eq:v31-parent-formation-continuation}.  Combine the two estimates,
and then use
\eqref{eq:v31-far-coarea-shell-equivalence} for the last assertion.
\end{proof}

\begin{firewall}[A one-use parent sum can still be analytically circular]
\label{fw:v31-parent-formation-circular}
Lemma~\ref{lem:v31-parent-geometric-sum} completely removes the logarithmic
parent multiplicity recorded in V14: every parent atom is charged by a
geometric square sum comparable to its \(\lambda^2\)-energy.  Nevertheless
the product with the instantaneous Leray density produces
\(\|\Lambda u\|_2^4\).  Proposition~\ref{prop:v7-critical-source-return}
and Theorem~\ref{thm:enstrophy-continuation} identify finiteness of its time
integral as the continuation-class stock.  Therefore
\eqref{eq:v31-parent-formation-final} is a rigorous formation theorem but
not a global a priori estimate.  Using it as though its right side were
Leray-paid would assume the desired conclusion.
\end{firewall}

\begin{assessment}[The sharper remaining scalar is restricted formation
occupation]
\label{ass:v31-restricted-formation-target}
The first line
\eqref{eq:v31-parent-formation-sharp} is sharper than its continuation-class
corollary.  It asks only for
\(\mathscr Q_{\rm form}\), in which shell \(n\) is present after its own
formation time and is weighted by the actual high-parent tail
\(H_{K_n/2}^\Sigma\).  A future argument need not bound the full
\(\int\|\Lambda u\|_2^4\); it would suffice to prove a record-separating
bound for this restricted, stopping-time-correlated occupation.  V29's
Leray occupation controls a linear tail factor, not its product with the
instantaneous dissipation density, so it does not yet provide that bound.
\end{assessment}

% =============================================================================
\section{Sharp tests of heat lifting and signed hinge lifting}
\label{sec:v31-sharp-tests}
% =============================================================================

\begin{proposition}[The heat-formation exponent is sharp on one forced mode]
\label{prop:v31-forced-heat-sharpness}
Fix \(L>0\) and \(1/4\le\sigma\le3/4\).  For each integer \(N\ge1\), let
\(\phi_N\) be an \(L^2\)-normalized divergence-free Fourier eigenmode with
\(A\phi_N=N^2\phi_N\), put
\[
 K:=\frac{2N}{3},
 \qquad
 T:=\frac{L}{\nu N^2},
\]
and solve the forced Stokes equation
\[
 \partial_tv+\nu Av=f,\qquad v(0)=0,
\]
where
\begin{equation}
 A^{-\sigma}f(t)
 =c\,e^{-\nu N^2(T-t)}\phi_N.
 \label{eq:v31-extremizing-heat-source}
\end{equation}
Then, with
\[
 \mathscr S_\sigma
 :=\int_0^T\|A^{-\sigma}f(t)\|_2^2\,\dd t,
\]
one has the exact ratio
\begin{equation}
 \boxed{\quad
 \frac{\displaystyle
   \int_K^{2K}\|Q_\rho v(T)\|_2^2\,\dd\rho}
 {\mathscr S_\sigma}
 =
 \frac{1-e^{-2L}}{6\nu}
 N^{4\sigma-1}
 =
 c_{\sigma,L}\frac{K^{4\sigma-1}}{\nu},
 \quad}
 \label{eq:v31-forced-heat-ratio}
\end{equation}
where \(c_{\sigma,L}>0\) is independent of \(N\).  Hence a heat-only
estimate of annular critical coarea from an unweighted
\(A^{-\sigma}\)-source action must lose
\(K^{4\sigma-1}\).  This loss is scale-free only at
\(\sigma=1/4\).
\end{proposition}

\begin{proof}
Set
\[
 I_N:=\int_0^T e^{-2\nu N^2(T-t)}\,\dd t
 =\frac{1-e^{-2L}}{2\nu N^2}.
\]
Then
\[
 \mathscr S_\sigma=|c|^2I_N,
 \qquad
 v(T)=cN^{2\sigma}I_N\phi_N.
\]
Since \(N=3K/2\), the mode belongs to \(Q_\rho v(T)\) for
\(K\le\rho<N\), an interval of length \(N-K=N/3\).  Therefore
\[
 \int_K^{2K}\|Q_\rho v(T)\|_2^2\,\dd\rho
 =
 \frac N3|c|^2N^{4\sigma}I_N^2.
\]
Divide by \(|c|^2I_N\) and insert the value of \(I_N\).
\end{proof}

\begin{firewall}[The forced-mode test has the correct scope]
\label{fw:v31-forced-mode-scope}
Proposition~\ref{prop:v31-forced-heat-sharpness} is a sharp statement about
the heat/Duhamel operator.  Its forcing is not asserted to equal
\(\BNS(u)\), so it does not exclude a cancellation special to the complete
Navier--Stokes quadratic source.  It does prove that heat smoothing by itself
cannot improve \eqref{eq:v31-source-weight-collapse}; any improvement must
use the signed nonlinear structure before taking the source norm.
\end{firewall}

\begin{proposition}[The actual NS packet saturates flat annular unmarking]
\label{prop:v31-packet-flat-mark}
Use the smooth packet solutions, record cutoffs \(K_j\), physical
frequencies \(N_j\), and integrated fluxes \(F_j\) from V14 and V18.  On the
dyadic grid
\[
 K_{j,n}:=2^nK_j,\qquad
 \mathcal A_{j,n}:=[K_{j,n},2K_{j,n}),
\]
there are indices \(n_j\to\infty\) and measurable sets
\[
 G_j\subset
 \mathcal A_{j,n_j}\cap[r_-N_j,r_+N_j]
 \cap\{\rho:F_j(\rho)>0\}
\]
such that
\begin{align}
 K_{j,n_j}&\asymp N_j,
 \label{eq:v31-packet-selected-annulus}\\
 \int_{G_j}F_j(\rho)\,\dd\rho
 &\ge c_0R_j^2,
 \label{eq:v31-packet-raw-annular-area}\\
 \frac1{K_{j,n_j}}\int_{G_j}F_j(\rho)\,\dd\rho
 &\asymp\frac{R_j^2}{N_j}.
 \label{eq:v31-packet-marked-annular-height}
\end{align}
Thus the normalized flat-marked action has the small kinetic scale
\(R_j^2/N_j\), while multiplication by its exact unmarking factor
\(K_{j,n_j}\asymp N_j\) returns the order-\(R_j^2\) critical area.
\end{proposition}

\begin{proof}
The compact ratio interval \([r_-,r_+]\) intersects at most
\[
 N_*:=3+\left\lceil\log_2(r_+/r_-)\right\rceil
\]
annuli of any dyadic grid.  By
\eqref{eq:v18-packet-positive-flux-area}, one of those intersections carries
positive flux area at least \(cR_j^2/N_*\); choose its positive-flux part as
\(G_j\).  Every annulus which intersects the fixed physical window has
\(K_{j,n}\asymp N_j\), proving
\eqref{eq:v31-packet-selected-annulus}.  The lower estimate in
\eqref{eq:v31-packet-marked-annular-height} follows by division.  Its upper
estimate follows from the pointwise packet height bound
\eqref{eq:v18-packet-flux-height} and
\(|G_j|\le K_{j,n_j}\).
\end{proof}

\begin{assessment}[Packet audit of the formation route]
\label{ass:v31-packet-formation-audit}
For the packet,
\[
 K_j/N_j\to0,\qquad
 \mathscr A_j^{\rm far}
 =\int_{J_j}\|A^{-1/4}Q_{2K_j}\BNS(u^{(j)})\|_2^2\,\dd t
 \asymp R_j^3
\]
by \eqref{eq:v14-packet-critical-action}.  Thus the critical formation bound
\eqref{eq:v31-critical-formation-mark} supplies an order-\(R_j^3/\nu\)
upper ledger for an order-\(R_j^2\) critical flux event.  Equivalently,
the weak source has size \(R_j^3/N_j^2\), and restoration of the physical
annular factor squared gives \(N_j^2(R_j^3/N_j^2)=R_j^3\).
The packet therefore passes every V31 identity and gives no record
separation.  As in V14 and V18, these are distinct smooth solutions and are
not a concatenated blow-up trajectory.
\end{assessment}

\begin{proposition}[The normalized hinge has an exact balanced-sector
nullspace]
\label{prop:v31-hinge-balanced-nullspace}
For every admissible torus frequency \(\lambda\) and every \(e>0\), there is
a real divergence-free one-wavevector datum \(v\) with total energy \(e\)
and equal positive and negative helical energies.  The corresponding exact
Navier--Stokes solution is
\begin{equation}
 u(t)=e^{-\nu\lambda^2t}v.
 \label{eq:v31-balanced-plane-wave-solution}
\end{equation}
For every cutoff \(\rho\) and hinge radius \(r\),
\begin{align}
 \mathcal Y_\rho(t)&=0,
 &
 z_r(t)&=0,
 &
 q_r(t)&=0,
 &
 j_r(t)&=0,
 \label{eq:v31-hinge-null-quantities}
\end{align}
whereas, for \(0<K<\lambda\),
\begin{equation}
 \mathcal C_K[u(t)]
 =e^{-2\nu\lambda^2t}e(\lambda-K)>0.
 \label{eq:v31-hinge-null-total-coarea}
\end{equation}
Consequently no estimate of total critical tail coarea can be obtained from
the sector-difference hinge packet alone.
\end{proposition}

\begin{proof}
Choose a nonzero lattice vector \(k\), \(|k|=\lambda\), and a real unit
vector \(a\perp k\).  The real plane wave
\[
 v(x)=c\,a\cos(k\cdot x)
\]
has equal-magnitude coefficients in the positive and negative helical basis;
choose \(c\) to give energy \(e\).  Since \(v\cdot k=0\) pointwise,
\((v\cdot\nabla)v=0\), so \eqref{eq:v31-balanced-plane-wave-solution} solves
Navier--Stokes exactly.  Equal helical energies make every sector difference
tail vanish.  The modal formulae
\eqref{eq:v27-z-weight-form}--\eqref{eq:v27-q-weight-form} then give
\(z_r=q_r=0\), and the nonlinear hinge current is zero.  Finally use the
Fourier formula in \eqref{eq:v31-exact-dyadic-first-moment}.
\end{proof}

\begin{firewall}[Hinge cancellation belongs only to the polarization branch]
\label{fw:v31-hinge-only-polarization}
The nullspace in
Proposition~\ref{prop:v31-hinge-balanced-nullspace} is exactly the
high-total, low-polarization geometry that generated the V28--V30 inner
branch.  V27's radial BV theorem remains valid and useful for assembling
signed polarization exits, but it cannot supply the missing positive total
formation moment.  Any V32 use of the hinge must therefore remain inside the
signed \(S_+,S_-\) partition and must be coupled to an independent total-tail
or ancestor estimate.
\end{firewall}

% =============================================================================
\section{Consequences for the V19 moving-tail and V29 recent-entry gates}
\label{sec:v31-consequences}
% =============================================================================

\begin{assessment}[Why the three inherited mechanisms do not yet combine to
closure]
\label{ass:v31-mechanism-combination}
The exact status of the three mechanisms proposed at the end of V30 is now
as follows.
\begin{enumerate}[label=\textup{(\roman*)},leftmargin=3.7em]
\item A flat dyadic probability mark handles every stopping-time family, but
      its removal is exactly the first moment
      \(\sum K_n(\cdot)\).
\item The heat/Duhamel map supplies that first moment with no shell
      multiplicity through
      \eqref{eq:v31-critical-formation-mark}.  The scale-matched source is
      \(A^{-1/4}\BNS(u)\); use of the energy-paid weak source merely hides
      the same critical action in the weights \(K_n^2\).
\item The V14 high-parent identity converts the weighted weak source into the
      restricted formation occupation
      \(\mathscr Q_{\rm form}\).  Dyadic parent incidence sums
      geometrically rather than logarithmically, but the available universal
      bound is the V1 continuation stock.
\item The V19 moving-tail telescope still exposes actual swept kinetic stock
      \(J_n\).  V20's common-root orthogonal ancestry controls its
      inverse-frequency response, which is exactly the weak rung
      \eqref{eq:v31-weak-formation-rung}.  Upgrading it to raw swept energy
      requires the intermediate rung, and upgrading further to critical
      coarea requires the critical rung.
\item The V27 hinge BV ledger sees the signed sector-difference branch but
      has the exact total-sector nullspace
      \eqref{eq:v31-hinge-null-quantities}.
\end{enumerate}
Thus none of the three mechanisms may be quoted alone as a completed
unmarking theorem.  The remaining opportunity is their
\emph{correlation}: on the actual recent-entry set, the active parent tails,
signed triad work, and endpoint exits are not arbitrary independent
functions.
\end{assessment}

\begin{firewall}[The persistent branch is not charged a second time]
\label{fw:v31-persistent-no-double-charge}
Nothing in V31 changes the cofinally one-use persistence payment
\eqref{eq:v29-cofinal-persistence-sum}.  Long-lived high-balanced rectangles
remain paid by the radius--time Leray occupation.  The new formation ledger
is to be used only on recent-entry shells; adding the persistence duration to
\(\mathscr Q_{\rm form}\) as a second debit would double count the same time
interval.
\end{firewall}

% =============================================================================
\section{Integrated V31 conclusion and next acquisition order}
\label{sec:v31-conclusion}
% =============================================================================

\begin{theorem}[What V31 proves]
\label{thm:v31-conclusion}
Relative to V1--V30, the following are proved.
\begin{enumerate}[label=\textup{(V31.\arabic*)},leftmargin=6.2em]
\item Uniform probability measure on every dyadic cutoff annulus has the
      exact ramp symbol \eqref{eq:v31-flat-ramp}, and multiplying its marked
      tail by \(K_n\) gives the unweighted annular tail exactly.
\item The full critical cutoff coarea is precisely the first moment
      \(\sum_nK_n\mathcal M_n\).  A far Fourier atom fills every lower
      dyadic mark and saturates that geometric first moment.
\item Arbitrary stopping maps on all dyadic annuli obey the
      \(\ell^1\)-dual family bounds
      \eqref{eq:v31-l1-total-family}--%
      \eqref{eq:v31-l1-difference-family}.  Returning to raw annular action
      costs exactly \(K_n\), and the V29 exit capacities remain order-annular
      at natural thresholds.
\item Source-born dyadic velocity responses obey the complete heat-formation
      ladder \eqref{eq:v31-scale-matched-formation}.  The weak,
      intermediate, and critical source spaces pay respectively
      inverse-scale kinetic stock, kinetic stock, and critical first-moment
      stock.
\item Every attempt to recover the first moment through a weaker source norm
      restores the exact missing shell weight and collapses to the critical
      source action, as in \eqref{eq:v31-source-weight-collapse}.
\item Combining the formation ladder with the actual V14 high-parent
      identity yields the restricted occupation
      \(\mathscr Q_{\rm form}\), with no dyadic parent multiplicity.
      Its presently available bound is
      \(\int\|\Lambda u\|_2^4\), the V1 continuation stock.
\item A single forced heat mode proves the source/formation scale exponent
      sharp for the Duhamel operator.  The actual V14/V18 smooth NS packet
      saturates the flat mark-to-unmarked annular conversion at the remote
      physical frequency.
\item A balanced-helicity plane wave lies in the exact nullspace of the
      normalized hinge packet while retaining positive total critical
      coarea.  Hinge BV can therefore be used for polarization exits, not as
      a stand-alone total-tail formation estimate.
\end{enumerate}
The full Navier--Stokes stability/global-regularity theorem remains unproved.
The open gate has been sharpened to a stopping-time-correlated estimate for
\(\mathscr Q_{\rm form}\), or an equivalent signed nonlinear cancellation
which gains one record factor over its continuation-class upper bound.
\end{theorem}

\begin{proof}
Items \textup{(V31.1)--(V31.2)} are
Proposition~\ref{prop:v31-flat-mark-symbol},
Theorem~\ref{thm:v31-dyadic-coarea}, and
Proposition~\ref{prop:v31-far-atom-mark-chain}.  Item
\textup{(V31.3)} is
Theorem~\ref{thm:v31-dyadic-stopped-family} and
Corollary~\ref{cor:v31-flat-exit-widths}.  Items
\textup{(V31.4)--(V31.5)} are
Theorem~\ref{thm:v31-heat-formation-ladder} and
Corollary~\ref{cor:v31-three-rung-formation}.  Item
\textup{(V31.6)} is
Lemma~\ref{lem:v31-parent-geometric-sum} and
Theorem~\ref{thm:v31-nonlinear-parent-formation}.  Item
\textup{(V31.7)} is
Propositions~\ref{prop:v31-forced-heat-sharpness} and
\ref{prop:v31-packet-flat-mark}.  Item \textup{(V31.8)} is
Proposition~\ref{prop:v31-hinge-balanced-nullspace}.
\end{proof}

\begin{programme}[V32 signed formation-correlation acquisition order]
\label{prog:v32}
\begin{enumerate}[label=\textup{(AC\arabic*)},leftmargin=4.8em]
\item Retain the exact restricted quantity
      \(\mathscr Q_{\rm form}\) from
      \eqref{eq:v31-restricted-parent-formation}.  Do not replace it at the
      outset by the full enstrophy-square integral.
\item Specialize the active indicators \(\chi_n\) to the actual V29
      total-exit and polarization-exit times.  Determine which pieces of
      \(H_{K_n/2}^\Sigma\) are inherited across the exit and which are born
      on the recent interval.
\item On total exits, insert the complete signed triad layer-cake formula
      \eqref{eq:v18-triad-gap-work} before taking positive parts in separate
      annuli.  Test whether the low leg of a high--high/low triad pays the
      long dyadic lever arm through low-mode time variation.
\item Preserve V21's one-crossing incidence and V20's common-root
      orthogonality.  Use the scale-matched ladder: weak source may pay only
      inverse-scale response, intermediate source raw swept energy, and
      critical source the first moment.
\item On polarization exits, retain the \(S_v,S_+,S_-\) sign partition and
      assemble \(S_+\) against \(S_-\) through the V27 radial hinge current
      before taking absolute values.  Do not apply the hinge to the total
      branch, forbidden by
      Proposition~\ref{prop:v31-hinge-balanced-nullspace}.
\item Keep V19's moving-tail junction stock explicit.  A remote mode fills
      all lower probability marks, so only the one-use swept-band projection,
      not a sum of tail marks, may be treated as orthogonal ancestry.
\item Seek a gain of one critical-record factor relative to
      \(\mathscr Q_{\rm form}\lesssim\int\|\Lambda u\|_2^4\).
      A bound of the same continuation-class size is not progress toward
      closure.
\item Test every proposed correlation estimate on both
      Proposition~\ref{prop:v31-forced-heat-sharpness} and the actual
      V14/V18 packet.  The forced mode tests the heat operator; the packet
      tests genuine short-time NS nonlinearity.
\item Reuse the V29 cofinal persistence sum without charging its intervals
      again.  The next compulsory companion audit is V35.
\end{enumerate}
\end{programme}

\begin{assessment}[Position after V31]
\label{ass:v31-position}
V31 has carried out the formation-mark step requested by V30.  Dyadic flat
marks give an exact critical first moment, and the Duhamel map lifts that
moment with no stopping-time or shell multiplicity.  The calculation also
proves its boundary: heat lifting is scale-free only in the critical source
space, and high-parent lifting of the globally paid weak source returns the
enstrophy-square continuation stock.

The next attack must therefore use information that has not yet been
discarded in the norm estimates: the correlation between actual recent-entry
times, complete signed triad work, and the active high-parent tail.  The
polarization hinge can contribute only on its signed branch, while the total
branch must retain the complete energy-transfer cancellation.  No full
stability or global-regularity claim is made.
\end{assessment}

% =============================================================================
\section*{Version V31 append-only record}
\addcontentsline{toc}{section}{Version V31 append-only record}
% =============================================================================

The complete V30 mathematical source was retained before this continuation.
Its SHA--256 digest was
\begin{center}\scriptsize\ttfamily
2da179be947d136d7830eb9cfeb4282f95140f5633afb1cab82f07ba492aecdd.
\end{center}
V31 introduced flat probability marks on dyadic cutoff annuli and proved that
their radius-weighted first moment is exactly the unweighted critical coarea.
It derived the complete weak/intermediate/critical heat-formation ladder for
shell responses with arbitrary shell-dependent lower times.  The resulting
NS first-moment formation inequality is multiplicity-free but is paid by the
critical source action.  Combining it with the actual high-parent identity
eliminates dyadic parent multiplicity and returns the restricted formation
occupation, whose universal estimate is the V1 enstrophy-square continuation
stock.  Forced-heat, genuine smooth-packet, and balanced-helicity nullspace
tests delimit the remaining route.  No full regularity claim is made.

% =============================================================================
\section{V32: cutoff-free exit ledgers and the type-I enstrophy firewall}
\label{sec:v32-entry}
% =============================================================================

V31 ended with a stopping-time-correlated target, the restricted formation
occupation \(\mathscr Q_{\rm form}\), and asked for a gain of one record
factor over its continuation-class bound.  Before pursuing the correlation
itself, this version goes one step upstream and asks a simpler question: in
which \emph{ledger} is a V29 recent-entry exit paid, if the cutoff is not
allowed to enter the estimate at all?

The answer is that every nonlinear term in the total and sector-difference
tail balances is dominated, uniformly in the cutoff radius, by one density,
\[
 C_4E_*^{1/4}\|\Lambda u(t)\|_2^{5/2},
\]
where \(C_4\) is the periodic Ladyzhenskaya constant.  Consequently all V29
recent-entry alternatives, total, polarization, and low-total, are paid on
the terminal window by a single number, the recent \(L^{5/2}_tH^1\) ledger.
This removes the cutoff multiplicity recorded in
Firewall~\ref{fw:v29-entry-not-summed} for every exit family, but the
ledger is no longer the Leray ledger.  It lies strictly between the Leray
class \(L^2_tH^1\) and the V1 continuation class \(L^4_tH^1\); its exchange
rate to the continuation class is a quarter power.

Two consequences follow with complete proofs.  First, a high-total spectral
tail at any cutoff is bounded by the recent ledger over one parabolic window,
with no record factor, and the Fourier-side critical record
\(\|A^{1/4}u(b)\|_2^2\) is bounded by the recent ledger times the terminal
enstrophy.  Second, the Leray blow-up rate forces divergence of
\(\int\|\Lambda u\|_2^p\,\dd t\) only for \(p\ge4\); a type-I enstrophy
rate therefore has a finite \(L^{5/2}\) ledger, and under that rate every
high-total branch of V24--V31 is impossible, the trajectory is precompact in
\(L^2\), and the critical record grows no faster than
\((T_*-b)^{-2/11}\), improving the interpolation exponent \(1/4\).

No global regularity claim is made.  No companion review is due in V32; the
next compulsory SM/YM synchronization remains V35.

% =============================================================================
\section{Two-stress anatomy of the cutoff flux}
\label{sec:v32-anatomy}
% =============================================================================

Throughout, \(\BNS(a,c):=\PP\operatorname{div}(a\otimes c)
=\PP\bigl((a\cdot\nabla)c\bigr)\) for divergence-free \(a\), so that
\(\BNS(u)=\BNS(u,u)\).  For a cutoff \(\rho>0\) write
\begin{equation}
 v:=P_\rho u,\qquad w:=Q_\rho u,\qquad
 w_1:=\Pi_{(\rho,2\rho]}u,\qquad
 H_\rho^\Sigma=\|w\|_2^2,\qquad
 V_\rho^\Sigma:=\|\Lambda w\|_2^2.
 \label{eq:v32-low-high-split}
\end{equation}
Here \(\Pi_{(\rho,2\rho]}=\mathbf 1_{(\rho,2\rho]}(\Lambda)\) is the
first exterior octave of V12, and \(H_\rho^\Sigma\) agrees with
\eqref{eq:v29-total-tail} by helical orthogonality.  For a matrix field
\(M\) and a vector field \(c\) we write
\(M{:}\nabla c=\sum_{i,j}M_{ij}\partial_ic_j\).

\begin{proposition}[Exact two-stress form of the cutoff flux]
\label{prop:v32-two-stress}
At every smooth time and every cutoff \(\rho>0\),
\begin{align}
 \Phi_\rho
 &=\langle\BNS(u),v\rangle
 =-\langle\BNS(u),w\rangle
 =-\int_{\Torus^3}(u\otimes u){:}\nabla v,
 \label{eq:v32-flux-dual-forms}\\
 \boxed{\quad
 \Phi_\rho
 =\Phi_\rho^{\rm LL}-\Phi_\rho^{\rm HH},
 \qquad
 \Phi_\rho^{\rm LL}:=\int_{\Torus^3}(v\otimes v){:}\nabla w_1,
 \qquad
 \Phi_\rho^{\rm HH}:=\int_{\Torus^3}(w\otimes w){:}\nabla v.
 \quad}
 \label{eq:v32-two-stress}
\end{align}
The low--low stress acts only on the first exterior octave: replacing
\(w_1\) by \(w\) in \(\Phi_\rho^{\rm LL}\) does not change its value.
\end{proposition}

\begin{proof}
By \eqref{eq:v4-cumulative-flux}, \(\Phi_\rho=\langle\BNS(u),P_\rho u\rangle
=\langle\BNS(u),v\rangle\); the Leray projector may be dropped because
\(v\) is divergence-free.  Energy cancellation
\(\langle\BNS(u),u\rangle=0\) gives the second form.  Integration by parts
with \(\operatorname{div}u=0\) gives
\(\int u_i\partial_iu_j\,v_j=-\int u_iu_j\partial_iv_j\), the third form.

Expand \(\BNS(u)=\BNS(v,v)+\BNS(v,w)+\BNS(w,v)+\BNS(w,w)\) in
\(-\langle\BNS(u),w\rangle\).  Since \(v\) and \(w\) are divergence-free,
\(\langle\BNS(w,w),w\rangle=0\) and
\(\langle\BNS(v,w),w\rangle=\int v_i\partial_iw_j\,w_j=0\).  Hence
\[
 \Phi_\rho=-\langle\BNS(v,v),w\rangle-\langle\BNS(w,v),w\rangle
 =\int(v\otimes v){:}\nabla w-\int(w\otimes w){:}\nabla v,
\]
the first term again by parts with \(\operatorname{div}v=0\), the second
by definition.  Finally, \(v\otimes v\) has Fourier support in
\(\{|\xi|\le2\rho\}\), because each of its factors has support in
\(\{|\xi|\le\rho\}\) and the lattice sum of two such vectors has modulus at
most \(2\rho\); by Parseval the pairing with \(\nabla w\) sees only the
modes of \(w\) with \(\rho<|\xi|\le2\rho\), that is, \(\nabla w_1\).
\end{proof}

\begin{corollary}[The two stresses are the two V18 lever-arm slots]
\label{cor:v32-slots}
Decompose the modal transfers into complete triads as in
Theorem~\ref{thm:v18-triad-gap-formula}, with ordered radii
\(\lambda_1\le\lambda_2\le\lambda_3\).  Then, at every cutoff,
\begin{equation}
 \Phi_\rho^{\rm LL}
 =\sum_{\triangle:\ \lambda_2\le\rho<\lambda_3}
   \bigl(-T_3^\triangle\bigr),
 \qquad
 -\Phi_\rho^{\rm HH}
 =\sum_{\triangle:\ \lambda_1\le\rho<\lambda_2}T_1^\triangle .
 \label{eq:v32-slot-identification}
\end{equation}
Thus the low--low stress is exactly the short-gap slot
\([\lambda_2,\lambda_3)\) of every triad, and the high--high stress is
exactly the long lever-arm slot \([\lambda_1,\lambda_2)\) of
Corollary~\ref{cor:v18-high-high-low-gap}.  The octave localization in
Proposition~\ref{prop:v32-two-stress} is the triangle inequality
\(\lambda_3\le\lambda_1+\lambda_2\le2\rho\) on the short-gap slot.
\end{corollary}

\begin{proof}
A triad contributes to \(-\langle\BNS(v,v),w\rangle\) precisely when two of
its legs lie in \(\{|k|\le\rho\}\) and the third in \(\{|k|>\rho\}\), that
is, when \(\lambda_2\le\rho<\lambda_3\); the value of its contribution is
the transfer into the high leg, which by
\eqref{eq:v18-triad-flux-profile} equals \(-T_3^\triangle\).  Likewise a
triad contributes to \(\langle\BNS(w,w),v\rangle=-\Phi_\rho^{\rm HH}\)
exactly when one leg is low and two are high, that is, when
\(\lambda_1\le\rho<\lambda_2\), and its contribution is the transfer
\(T_1^\triangle\) into the low leg.  Triads with three low or three high
legs cancel by \eqref{eq:v18-triad-energy-cancellation}.
\end{proof}

% =============================================================================
\section{The Ladyzhenskaya ledger is cutoff-free}
\label{sec:v32-ladyzhenskaya}
% =============================================================================

\begin{lemma}[Periodic Ladyzhenskaya inequality]
\label{lem:v32-ladyzhenskaya}
There is a constant \(C_4\), depending only on the torus, such that every
mean-zero \(f\in H^1(\Torus^3;\R^3)\) satisfies
\begin{equation}
 \|f\|_4^2\le C_4\|f\|_2^{1/2}\|\Lambda f\|_2^{3/2},
 \qquad
 \|f\|_4^2\le C_4\|\Lambda^{3/4}f\|_2^2 .
 \label{eq:v32-ladyzhenskaya}
\end{equation}
\end{lemma}

\begin{proof}
For mean-zero periodic functions the Sobolev embedding
\(\dot H^{3/4}(\Torus^3)\hookrightarrow L^4(\Torus^3)\) holds, since
\(3/4=3(1/2-1/4)\); this is the second inequality, with
\(C_4\) the squared embedding constant.  H\"older's inequality on the lattice
with exponents \(4\) and \(4/3\) gives
\[
 \|\Lambda^{3/4}f\|_2^2
 =\sum_{k\ne0}|k|^{3/2}|\widehat f(k)|^2
 =\sum_{k\ne0}\bigl(|\widehat f(k)|^2\bigr)^{1/4}
   \bigl(|k|^2|\widehat f(k)|^2\bigr)^{3/4}
 \le\|f\|_2^{1/2}\|\Lambda f\|_2^{3/2}.
\]
\end{proof}

\begin{theorem}[Every tail balance has a cutoff-free nonlinear density]
\label{thm:v32-cutoff-free-density}
At every smooth time \(t\) and every cutoff \(\rho>0\),
\begin{align}
 \boxed{\quad
 |\Phi_\rho(t)|
 \le\|u\|_4^2\,\|\Lambda P_\rho u\|_2
 \le C_4E_*^{1/4}\|\Lambda u(t)\|_2^{5/2},
 \quad}
 \label{eq:v32-flux-cutoff-free}\\
 \boxed{\quad
 |d_\rho(t)|
 \le\|u\|_4^2\,\|\Lambda Q_\rho u\|_2
 \le C_4E_*^{1/4}\|\Lambda u(t)\|_2^{5/2},
 \quad}
 \label{eq:v32-difference-cutoff-free}
\end{align}
where \(d_\rho=\widetilde\Phi_\rho^+-\widetilde\Phi_\rho^-\) is the
V25 sector-difference current.  The two stress slots obey the finer bounds
\begin{align}
 |\Phi_\rho^{\rm LL}|
 &\le\|P_\rho u\|_4^2\,\|\Lambda w_1\|_2
 \le2\rho\,C_4E_*^{1/4}\|\Lambda P_\rho u\|_2^{3/2}
   \bigl(H_\rho^\Sigma-H_{2\rho}^\Sigma\bigr)^{1/2},
 \label{eq:v32-LL-slot-bound}\\
 |\Phi_\rho^{\rm HH}|
 &\le\|Q_\rho u\|_4^2\,\|\Lambda P_\rho u\|_2
 \le C_4\,\rho^{-1/2}\,V_\rho^\Sigma\,\|\Lambda P_\rho u\|_2
 \le C_4\,\rho^{1/2}E_*^{1/2}\,V_\rho^\Sigma .
 \label{eq:v32-HH-slot-bound}
\end{align}
\end{theorem}

\begin{proof}
From \eqref{eq:v32-flux-dual-forms} and Cauchy--Schwarz,
\(|\Phi_\rho|\le\|u\otimes u\|_2\|\nabla v\|_2=\|u\|_4^2\|\Lambda v\|_2\),
because \(\|u\otimes u\|_2^2=\int|u|^4\) and, for divergence-free \(v\),
\(\|\nabla v\|_2=\|\Lambda v\|_2\).  Apply
\eqref{eq:v32-ladyzhenskaya}, \(\|u\|_2^2\le E_*\), and
\(\|\Lambda v\|_2\le\|\Lambda u\|_2\).

For the difference current, \eqref{eq:v23-compensated-sector-flux} and the
self-adjointness of the helical projectors give
\[
 d_\rho
 =-\bigl\langle\BNS(u),Q_\rho(u^+-u^-)\bigr\rangle
 =-\bigl\langle\BNS(u),Q_\rho\widetilde u\bigr\rangle,
 \qquad
 \widetilde u:=\operatorname{curl}\Lambda^{-1}u,
\]
by \eqref{eq:v23-helical-projectors}.  The field
\(\widetilde u\) is divergence-free and mean-zero, and
\(|k\times\widehat u(k)|=|k||\widehat u(k)|\) for
\(\widehat u(k)\perp k\), so
\(\|\Lambda^sQ_\rho\widetilde u\|_2=\|\Lambda^sQ_\rho u\|_2\) for every
\(s\).  Integrating by parts as before,
\(|d_\rho|\le\|u\|_4^2\|\Lambda Q_\rho\widetilde u\|_2
=\|u\|_4^2\|\Lambda Q_\rho u\|_2\), and
\eqref{eq:v32-ladyzhenskaya} finishes.

The first slot bound is Cauchy--Schwarz in \eqref{eq:v32-two-stress},
followed by \eqref{eq:v32-ladyzhenskaya} for \(v\) and
\(\|\Lambda w_1\|_2\le2\rho\|w_1\|_2
=2\rho(H_\rho^\Sigma-H_{2\rho}^\Sigma)^{1/2}\).  For the second, apply the
Sobolev form of \eqref{eq:v32-ladyzhenskaya} to \(w\) and use
\(|k|^{3/2}\le\rho^{-1/2}|k|^2\) on its support \(\{|k|>\rho\}\), so
\(\|w\|_4^2\le C_4\rho^{-1/2}V_\rho^\Sigma\).  Finally
\(\|\Lambda P_\rho u\|_2\le\rho\|u\|_2\le\rho E_*^{1/2}\).
\end{proof}

\begin{proposition}[Corridor form: the same densities at record scale]
\label{prop:v32-corridor-form}
At every smooth time and cutoff,
\begin{equation}
 |\Phi_\rho|+|d_\rho|
 \le2C_4\|A^{1/4}u\|_2\,\|\Lambda u\|_2^2 .
 \label{eq:v32-record-leray-density}
\end{equation}
On a V1 card satisfying \eqref{eq:record-corridor} the right side is at
most \(4C_4\mathcal R_n\|\Lambda u\|_2^2\).
\end{proposition}

\begin{proof}
By interpolation on the lattice,
\(\|\Lambda^{3/4}u\|_2^2\le\|A^{1/4}u\|_2\|\Lambda u\|_2\); insert this in
the Sobolev form of \eqref{eq:v32-ladyzhenskaya} and in the first
inequalities of
\eqref{eq:v32-flux-cutoff-free}--\eqref{eq:v32-difference-cutoff-free}.
\end{proof}

\begin{assessment}[Two ledgers for the same current]
\label{ass:v32-two-ledgers}
Equation \eqref{eq:v32-record-leray-density} is the estimate implicit in
V4--V7: the recent influx is paid by the Leray density times one record
factor, and a doubled-record sequence requires only the summable Leray
share \(\gtrsim E_*/\mathcal R_n\) per card.  Equation
\eqref{eq:v32-flux-cutoff-free} replaces the record factor by the
\(L^{5/2}_tH^1\) density, which has no record and no cutoff.  The two
ledgers are not comparable pointwise; the rest of V32 works out what the
second one buys and what it costs.
\end{assessment}

% =============================================================================
\section{Recent-entry exits pay one cutoff-free window debit}
\label{sec:v32-exit-debit}
% =============================================================================

For a terminal time \(b\) and \(0<\tau\le b\), define the recent ledgers
\begin{equation}
 \mathscr L_{5/2}(b,\tau)
 :=\int_{b-\tau}^b\|\Lambda u(t)\|_2^{5/2}\,\dd t,
 \qquad
 \mathscr L_2(b,\tau)
 :=\int_{b-\tau}^b\|\Lambda u(t)\|_2^2\,\dd t .
 \label{eq:v32-recent-ledgers}
\end{equation}

\begin{theorem}[All three V29 exits are paid by the same window number]
\label{thm:v32-exit-debit}
Fix \(b\) and \(0<\tau\le b\).
\begin{enumerate}[label=\textup{(\roman*)},leftmargin=3.7em]
\item If a cutoff \(\rho\) satisfies the total alternative
      \eqref{eq:v29-recent-total-influx} of
      Theorem~\ref{thm:v29-two-threshold-entry} with exit time
      \(t_\rho\in(b-\tau,b)\), then
      \begin{equation}
       \mathscr L_{5/2}(b,\tau)
       \ge\frac{\kappa h}{2C_4E_*^{1/4}} .
       \label{eq:v32-total-exit-debit}
      \end{equation}
\item If \(\rho\) satisfies the polarization alternative
      \eqref{eq:v29-recent-polarization-action}, then
      \begin{equation}
       C_4E_*^{1/4}\mathscr L_{5/2}(b,\tau)
       +\nu\mathscr L_2(b,\tau)
       \ge\frac{(1-\beta)\varepsilon}{2}.
       \label{eq:v32-polarization-exit-debit}
      \end{equation}
\item If \(\rho\in S_{\rm new}\) in
      Theorem~\ref{thm:v29-tail-mass-ancestry}, then
      \begin{equation}
       H_\rho^\Sigma(b)
       \le\frac{2C_4E_*^{1/4}}{1-\alpha}\,
         \mathscr L_{5/2}(b,\tau).
       \label{eq:v32-low-total-exit-debit}
      \end{equation}
\end{enumerate}
The right sides of \textup{(i)}--\textup{(ii)} and the coefficient in
\textup{(iii)} do not depend on \(\rho\).
\end{theorem}

\begin{proof}
In case (i), \(\kappa h/2\le\int_{t_\rho}^b\Phi_\rho\,\dd t
\le\int_{b-\tau}^b|\Phi_\rho|\,\dd t\), and
\eqref{eq:v32-flux-cutoff-free} bounds the integrand.  In case (ii) apply
\eqref{eq:v32-difference-cutoff-free} to the integral of \(d_\rho\) and
\(V_\rho^++V_\rho^-=V_\rho^\Sigma\le\|\Lambda u\|_2^2\) to the dissipation
term.  In case (iii) the influx bound
\eqref{eq:v29-relative-recent-influx} gives
\(\frac{1-\alpha}2H_\rho^\Sigma(b)\le\int_{t_\rho}^b\Phi_\rho\,\dd t\);
bound the right side as in (i).
\end{proof}

\begin{corollary}[Removal of the moving-family multiplicity]
\label{cor:v32-no-multiplicity}
Let \(S\subset[K,K']\) be measurable, and let every \(\rho\in S\) be a
recent-entry cutoff at time \(b\) in any of the three senses above, with
common parameters.  Then the total, polarization, or low-total debits of
all \(\rho\in S\) are simultaneously certified by the one window number
\(C_4E_*^{1/4}\mathscr L_{5/2}(b,\tau)+\nu\mathscr L_2(b,\tau)\).  In
particular, in the low-total case,
\begin{equation}
 \int_{S}H_\rho^\Sigma(b)\,\dd\rho
 \le\frac{E_*}{4\nu\alpha K\tau}
 +|S|\,\frac{2C_4E_*^{1/4}}{1-\alpha}\,\mathscr L_{5/2}(b,\tau),
 \label{eq:v32-band-coarea-from-ledger}
\end{equation}
for every measurable \(S\subset[K,\infty)\) on which
\(H_\rho^\Sigma(b)>0\).
\end{corollary}

\begin{proof}
The first assertion restates Theorem~\ref{thm:v32-exit-debit}: the
certifying number does not involve \(\rho\), so it certifies all cutoffs at
once.  For \eqref{eq:v32-band-coarea-from-ledger}, split
\(S=S_{\rm long}\cup S_{\rm new}\) as in
Theorem~\ref{thm:v29-tail-mass-ancestry}, use
\eqref{eq:v29-long-lived-terminal-tail-mass} on the first part and
\eqref{eq:v32-low-total-exit-debit} on the second.
\end{proof}

\begin{firewall}[What has and has not been removed]
\label{fw:v32-ledger-shift}
Firewall~\ref{fw:v29-entry-not-summed} objected that the exit time and the
signed source live on a moving family of cutoff--time intervals, so that
absolute values integrated in \(\rho\) return the critical source.
Corollary~\ref{cor:v32-no-multiplicity} shows that the objection is a
statement about the Leray ledger only.  In the \(L^{5/2}_tH^1\) ledger the
family collapses to one number per terminal window.  The price is that
finiteness of \(\int_0^{T_*}\|\Lambda u\|_2^{5/2}\,\dd t\) is not known for
a hypothetical singular solution; Section~\ref{sec:v32-type-I} makes the
status of that ledger exact.
\end{firewall}

\begin{theorem}[Recent-ledger tail bound]
\label{thm:v32-recent-ledger-tail}
For every smooth \(b\), every \(0<\tau\le b\), and every \(\rho>0\),
\begin{equation}
 \boxed{\quad
 H_\rho^\Sigma(b)
 \le e^{-2\nu\rho^2\tau}H_\rho^\Sigma(b-\tau)
   +2C_4E_*^{1/4}
    \int_{b-\tau}^be^{-2\nu\rho^2(b-t)}
    \|\Lambda u(t)\|_2^{5/2}\,\dd t
 \le e^{-2\nu\rho^2\tau}E_*
   +2C_4E_*^{1/4}\mathscr L_{5/2}(b,\tau).
 \quad}
 \label{eq:v32-recent-ledger-tail}
\end{equation}
Consequently, on the parabolic window \(\tau=\theta^{-1}(\nu\rho^2)^{-1}\)
with \(\theta>0\),
\begin{equation}
 H_\rho^\Sigma(b)\ge(1+\kappa)\theta E_*
 \ \Longrightarrow\
 \mathscr L_{5/2}\bigl(b,\theta^{-1}(\nu\rho^2)^{-1}\bigr)
 \ge\frac{(1+\kappa)\theta-e^{-2/\theta}}{2C_4}E_*^{3/4}.
 \label{eq:v32-high-total-window-debit}
\end{equation}
\end{theorem}

\begin{proof}
Sum the two sector balances \eqref{eq:v23-sector-high-balance}:
\(\frac12\partial_tH_\rho^\Sigma+\nu V_\rho^\Sigma=\Phi_\rho\).  On
\(\{|k|>\rho\}\) one has \(V_\rho^\Sigma\ge\rho^2H_\rho^\Sigma\), so by
\eqref{eq:v32-flux-cutoff-free}
\[
 \partial_tH_\rho^\Sigma+2\nu\rho^2H_\rho^\Sigma
 \le2C_4E_*^{1/4}\|\Lambda u\|_2^{5/2}.
\]
Multiply by \(e^{2\nu\rho^2t}\) and integrate over \([b-\tau,b]\).  The
second inequality uses \(H_\rho^\Sigma\le E_*\) and drops the exponential
weight.  For \eqref{eq:v32-high-total-window-debit} insert
\(\tau=\theta^{-1}(\nu\rho^2)^{-1}\) and rearrange.
\end{proof}

\begin{assessment}[The high-total branch has a record-free order-one debit]
\label{ass:v32-high-total-debit}
For \(\theta\) fixed and \(\kappa>0\), the right side of
\eqref{eq:v32-high-total-window-debit} is a positive constant times
\(E_*^{3/4}\), independent of \(\rho\), of any record, and of any card
length.  Every V24--V31 branch that carried tail energy of order \(E_*\)
above a record cutoff, in particular the V28--V30 high-total near-balanced
inner branch, must therefore spend an order-one amount of the recent
\(L^{5/2}\) ledger on the parabolic window preceding the terminal time.
Compare Assessment~\ref{ass:v32-two-ledgers}: in the Leray ledger the same
event costs only the summable share \(E_*/\mathcal R_n\).
\end{assessment}

\begin{corollary}[Enstrophy spikes inside a high-total window]
\label{cor:v32-spike}
Let \(w=[b-\tau,b]\) carry \(\mathscr L_{5/2}(b,\tau)\ge D>0\) and put
\(\epsilon:=\mathscr L_2(b,\tau)\).  Then
\begin{equation}
 \sup_{t\in w}\|\Lambda u(t)\|_2^2\ge\frac{D^4}{\epsilon^4},
 \qquad
 \int_w\|\Lambda u\|_2^4\,\dd t\ge\frac{D^4}{\epsilon^3}.
 \label{eq:v32-spike}
\end{equation}
Along pairwise disjoint high-total windows \(w_j\) with
\(D_j\ge D_0>0\), the Leray shares satisfy
\(\sum_j\epsilon_j\le E_*/(2\nu)\), so \(\epsilon_j\to0\) and both
quantities in \eqref{eq:v32-spike} diverge at least like
\(\epsilon_j^{-3}\).
\end{corollary}

\begin{proof}
\(\mathscr L_{5/2}\le(\sup_w\|\Lambda u\|_2)^{1/2}\mathscr L_2\) and,
by H\"older, \(\mathscr L_{5/2}\le\mathscr L_2^{3/4}
(\int_w\|\Lambda u\|_2^4)^{1/4}\).  Disjointness and the Leray identity
give the summability.
\end{proof}

% =============================================================================
\section{The critical record is paid by the recent ledger}
\label{sec:v32-record-from-ledger}
% =============================================================================

Recall from Theorem~\ref{thm:v31-dyadic-coarea} that
\(\|A^{1/4}u(b)\|_2^2=\int_0^\infty H_\rho^\Sigma(b)\,\dd\rho\) and that the
critical coarea above \(K\) is \(\mathcal C_K[u(b)]=\int_K^\infty
H_\rho^\Sigma(b)\,\dd\rho\).

\begin{theorem}[Record from recent ledger and terminal enstrophy]
\label{thm:v32-record-from-ledger}
For every smooth \(b\), every \(0<\tau\le b\), and every
\(K'\ge K:=(\nu\tau)^{-1/2}\),
\begin{equation}
 \boxed{\quad
 \|A^{1/4}u(b)\|_2^2
 \le\Bigl(1+\frac{e^{-2}}4\Bigr)\frac{E_*}{(\nu\tau)^{1/2}}
 +2C_4E_*^{1/4}\mathscr L_{5/2}(b,\tau)\,K'
 +\frac{\|\Lambda u(b)\|_2^2}{2K'} .
 \quad}
 \label{eq:v32-record-three-terms}
\end{equation}
If moreover
\(\|\Lambda u(b)\|_2^2\ge4C_4E_*^{1/4}\mathscr L_{5/2}(b,\tau)/(\nu\tau)\),
then
\begin{equation}
 \boxed{\quad
 \|A^{1/4}u(b)\|_2^2
 \le\frac{1.04\,E_*}{(\nu\tau)^{1/2}}
 +2\|\Lambda u(b)\|_2
 \bigl(C_4E_*^{1/4}\mathscr L_{5/2}(b,\tau)\bigr)^{1/2}.
 \quad}
 \label{eq:v32-record-optimized}
\end{equation}
\end{theorem}

\begin{proof}
Split the coarea at \(K\) and \(K'\).  Below \(K\), \(H_\rho^\Sigma\le E_*\)
gives at most \(KE_*\).  On \([K,K']\) use
\eqref{eq:v32-recent-ledger-tail}:
\[
 \int_K^{K'}H_\rho^\Sigma(b)\,\dd\rho
 \le E_*\int_K^\infty e^{-2\nu\rho^2\tau}\,\dd\rho
   +2C_4E_*^{1/4}\mathscr L_{5/2}(b,\tau)(K'-K),
\]
and \(\int_K^\infty e^{-2\nu\rho^2\tau}\dd\rho
\le\int_K^\infty(\rho/K)e^{-2\nu\rho^2\tau}\dd\rho
=e^{-2\nu K^2\tau}/(4\nu K\tau)\).  Above \(K'\),
\eqref{eq:v29-weighted-tail-layer-cake} gives
\(\int_{K'}^\infty H_\rho^\Sigma\,\dd\rho\le(K')^{-1}\int_0^\infty\rho
H_\rho^\Sigma\,\dd\rho=\|\Lambda u(b)\|_2^2/(2K')\).  With
\(K=(\nu\tau)^{-1/2}\) one has \(KE_*=E_*(\nu\tau)^{-1/2}\) and
\(e^{-2\nu K^2\tau}/(4\nu K\tau)=e^{-2}(\nu\tau)^{-1/2}/4\), which proves
\eqref{eq:v32-record-three-terms}.  The function
\(K'\mapsto aK'+c/K'\), with \(a=2C_4E_*^{1/4}\mathscr L_{5/2}\) and
\(c=\|\Lambda u(b)\|_2^2/2\), attains its minimum \(2\sqrt{ac}\) at
\(K'_*=\sqrt{c/a}\); the stated hypothesis is exactly \(K'_*\ge K\).
\end{proof}

\begin{assessment}[Comparison with the trivial interpolation and with V31]
\label{ass:v32-record-comparison}
The lattice Cauchy--Schwarz inequality gives
\(\|A^{1/4}u\|_2^2\le E_*^{1/2}\|\Lambda u\|_2\).  Estimate
\eqref{eq:v32-record-optimized} replaces the factor \(E_*^{1/2}\) by
\(2(C_4E_*^{1/4}\mathscr L_{5/2}(b,\tau))^{1/2}\), at the cost of the
persistent term \(1.04E_*(\nu\tau)^{-1/2}\), which is the coarea of an
order-\(E_*\) tail located at the parabolic frequency of the window.  It
improves the trivial bound exactly when the recent \(L^{5/2}\) ledger of
the window is small compared with \(E_*^{3/2}/C_4^2\).  In the language of
V31, the whole dyadic first moment \(\sum_nK_n\mathcal M_n\) has been paid
by one window ledger and one terminal enstrophy, with no critical source
action and no parent tail; the V31 restricted occupation
\(\mathscr Q_{\rm form}\) does not appear.
\end{assessment}

% =============================================================================
\section{The type-I enstrophy firewall}
\label{sec:v32-type-I}
% =============================================================================

\begin{lemma}[Leray's enstrophy rate on the torus]
\label{lem:v32-leray-rate}
There is an absolute periodic constant \(c_L>0\) such that every smooth
solution with maximal time \(T_*<\infty\) satisfies
\begin{equation}
 \|\Lambda u(t)\|_2\ge c_L\nu^{3/4}(T_*-t)^{-1/4}
 \qquad(0\le t<T_*).
 \label{eq:v32-leray-rate}
\end{equation}
\end{lemma}

\begin{proof}
Let \(Y=\|\Lambda u\|_2^2\) and \(Z=\|Au\|_2^2\).  As in the proof of
Theorem~\ref{thm:enstrophy-continuation},
\(\frac12Y'+\nu Z=-\langle\BNS(u),Au\rangle\), and
\[
 |\langle\BNS(u),Au\rangle|
 \le\|u\|_6\|\nabla u\|_3\|Au\|_2
 \le C\,Y^{1/2}\bigl(Y^{1/4}Z^{1/4}\bigr)Z^{1/2}
 \le\frac\nu2Z+C_1\nu^{-3}Y^3,
\]
using \(\|u\|_6\le C\|\nabla u\|_2\), \(\|\nabla u\|_3\le\|\nabla u\|_2^{1/2}
\|\nabla u\|_6^{1/2}\), \(\|\nabla u\|_6\le C\|Au\|_2\), and Young's
inequality.  Hence \(Y'\le2C_1\nu^{-3}Y^3\), so
\((Y^{-2})'\ge-4C_1\nu^{-3}\) wherever \(Y>0\).  If \(Y\) stayed bounded
along a sequence \(s_k\to T_*\) with \(\sup_kY(s_k)<\infty\), the \(H^1\)
local existence time, which depends only on \(\nu\) and \(Y(s_k)\), would
exceed \(T_*-s_k\) for large \(k\), contradicting maximality; thus
\(Y(s)\to\infty\) as \(s\to T_*\).  Integrating \((Y^{-2})'\ge
-4C_1\nu^{-3}\) from \(t\) to \(s\) and letting \(s\to T_*\) gives
\(Y(t)^{-2}\le4C_1\nu^{-3}(T_*-t)\), which is
\eqref{eq:v32-leray-rate} with \(c_L=(4C_1)^{-1/4}\).
\end{proof}

\begin{proposition}[Which time-integrated enstrophy ledgers the rate
forces to diverge]
\label{prop:v32-ledger-exponents}
Let \(T_*<\infty\) be a maximal time.  For \(p\ge4\),
\(\int_{t_0}^{T_*}\|\Lambda u\|_2^p\,\dd t=\infty\) for every
\(t_0<T_*\).  For \(0<p<4\), the rate \eqref{eq:v32-leray-rate} alone does
not force divergence: every function \(Y\ge0\) satisfying
\(c_L^2\nu^{3/2}(T_*-t)^{-1/2}\le Y(t)\le C_I^2\nu^{3/2}(T_*-t)^{-1/2}\)
has \(\int^{T_*}Y^{p/2}\,\dd t<\infty\).  Moreover, for \(p<4\),
\(L^p_tH^1\) embeds only in the supercritical class \(L^p_tL^6_x\), with
Prodi--Serrin index \(2/p+3/6>1\), so finiteness of the \(L^p_tH^1\)
ledger is not a known continuation criterion either.
\end{proposition}

\begin{proof}
For \(p\ge4\), \eqref{eq:v32-leray-rate} gives
\(\int_{t_0}^{T_*}\|\Lambda u\|_2^p\ge c_L^p\nu^{3p/4}\int_{t_0}^{T_*}
(T_*-t)^{-p/4}\,\dd t=\infty\).  For \(p<4\),
\(\int^{T_*}(T_*-t)^{-p/4}\,\dd t<\infty\).  The embedding statement is the
mean-zero periodic Sobolev inequality \(\|u\|_6\le C\|\Lambda u\|_2\) and
the arithmetic \(2/p+1/2>1\) for \(p<4\).
\end{proof}

We say that a singular solution has a \emph{type-I enstrophy rate} on
\([T_*-\delta,T_*)\) if
\begin{equation}
 \|\Lambda u(t)\|_2\le C_I\nu^{3/4}(T_*-t)^{-1/4}
 \qquad(T_*-\delta\le t<T_*).
 \tag{I}\label{eq:v32-type-I}
\end{equation}
By Lemma~\ref{lem:v32-leray-rate} this is the two-sided Leray rate; it
holds, for instance, for any hypothetical self-similar or discretely
self-similar blow-up with finite-energy profile.

\begin{theorem}[Type-I enstrophy rate forces uniformly small spectral tails
and kills every high-total branch]
\label{thm:v32-type-I-firewall}
Assume \eqref{eq:v32-type-I}.  Then for every
\(b\in[T_*-\delta/2,T_*)\) and every \(\rho>0\), with
\(\tau_\rho:=\min\{\delta/2,(\nu\rho^{3/2})^{-1}\}\),
\begin{equation}
 \boxed{\quad
 \|Q_\rho u(b)\|_2^2
 \le E_*e^{-2\nu\rho^2\tau_\rho}
 +\frac{16}{3}C_4C_I^{5/2}E_*^{1/4}\nu^{15/8}\tau_\rho^{3/8}
 \xrightarrow[\rho\to\infty]{}0
 \quad\text{uniformly in }b.
 \quad}
 \label{eq:v32-uniform-tail-smallness}
\end{equation}
Consequently:
\begin{enumerate}[label=\textup{(\roman*)},leftmargin=3.7em]
\item the trajectory \(\{u(t):0\le t<T_*\}\) is precompact in
      \(L^2(\Torus^3)\), and \(u(t)\) converges strongly in \(L^2\) to a
      limit \(u(T_*)\) as \(t\to T_*\), with no energy loss at \(T_*\);
\item for every \(\theta,\kappa>0\) there is \(\rho_0\) such that no
      cutoff \(\rho\ge\rho_0\) and no \(b\in[T_*-\delta/2,T_*)\) satisfy
      \(H_\rho^\Sigma(b)\ge(1+\kappa)\theta E_*\); every cofinal
      high-total branch of V24--V31, including the V28--V30 high-total
      near-balanced inner branch and the V29 persistence rectangles with
      \(h_j=\theta E_*\), is empty;
\item the critical record obeys, for \(b\) close to \(T_*\),
      \begin{equation}
       \boxed{\quad
       \|A^{1/4}u(b)\|_2^2
       \le C_{\rm I}'\,(T_*-b)^{-2/11},
       \quad}
       \label{eq:v32-type-I-record-rate}
      \end{equation}
      with \(C_{\rm I}'\) depending only on
      \(C_4,C_I,E_*,\nu\); this improves the exponent \(1/4\) obtained by
      interpolating \eqref{eq:v32-type-I} with the energy bound.
\end{enumerate}
\end{theorem}

\begin{proof}
Since \(b-\tau_\rho\ge T_*-\delta\), \eqref{eq:v32-type-I} holds on the
window and
\[
 \mathscr L_{5/2}(b,\tau_\rho)
 \le C_I^{5/2}\nu^{15/8}\int_{b-\tau_\rho}^b(T_*-t)^{-5/8}\,\dd t
 =\frac83C_I^{5/2}\nu^{15/8}
  \bigl[(T_*-b+\tau_\rho)^{3/8}-(T_*-b)^{3/8}\bigr]
 \le\frac83C_I^{5/2}\nu^{15/8}\tau_\rho^{3/8},
\]
by subadditivity of \(s\mapsto s^{3/8}\).  Insert this in
\eqref{eq:v32-recent-ledger-tail}.  For large \(\rho\),
\(\tau_\rho=(\nu\rho^{3/2})^{-1}\), so
\(e^{-2\nu\rho^2\tau_\rho}=e^{-2\rho^{1/2}}\) and
\(\tau_\rho^{3/8}=\nu^{-3/8}\rho^{-9/16}\); both tend to zero.  This is
\eqref{eq:v32-uniform-tail-smallness}.

(i) On \([0,T_*-\delta/2]\) the solution is smooth, hence continuous into
\(H^1\), and its tails are uniformly small by compactness of that interval.
Together with \eqref{eq:v32-uniform-tail-smallness} the whole trajectory
has uniformly small tails and bounded energy, so it is precompact in
\(L^2\).  The weak limit \(u(T_*)\) exists because \(u\) is a Leray--Hopf
solution on \([0,T_*)\); uniform tail smallness upgrades weak to strong
convergence and rules out an energy jump.

(ii) Choose \(\rho_0\) so that the right side of
\eqref{eq:v32-uniform-tail-smallness} is below \((1+\kappa)\theta E_*\)
for \(\rho\ge\rho_0\).  Persistence rectangles with \(h_j=\theta E_*\) and
\(B_j\subset[K_j,\infty)\), \(K_j\to\infty\), require
\(H_\rho^\Sigma\ge\theta E_*\) at some \((\rho,t)\) with
\(\rho\ge\rho_0\) for large \(j\), which is excluded.

(iii) Apply \eqref{eq:v32-record-optimized} with
\(\tau=(T_*-b)^{4/11}\), which is at most \(\delta/2\) for \(b\) close to
\(T_*\), and with the ledger bound above:
\[
 \|A^{1/4}u(b)\|_2^2
 \le1.04E_*\nu^{-1/2}(T_*-b)^{-2/11}
 +2C_I\nu^{3/4}(T_*-b)^{-1/4}
  \Bigl(\tfrac83C_4C_I^{5/2}E_*^{1/4}\nu^{15/8}\Bigr)^{1/2}
  (T_*-b)^{3/22}.
\]
Both exponents equal \(-2/11\).  The hypothesis of
\eqref{eq:v32-record-optimized} reads
\(\|\Lambda u(b)\|_2^2\ge4C_4E_*^{1/4}\mathscr L_{5/2}/(\nu\tau)\); its
left side is at least \(c_L^2\nu^{3/2}(T_*-b)^{-1/2}\) by
Lemma~\ref{lem:v32-leray-rate}, while its right side is at most a constant
times \(\tau^{3/8-1}=(T_*-b)^{-5/22}\), and \(-1/2<-5/22\), so the
hypothesis holds for \(b\) close to \(T_*\).
\end{proof}

\begin{firewall}[Scope of the type-I firewall]
\label{fw:v32-type-I-scope}
Theorem~\ref{thm:v32-type-I-firewall} does not exclude type-I blow-up.  It
proves that under a type-I enstrophy rate the singular mechanism cannot be
a high-total spectral cascade: order-\(E_*\) energy never reaches a remote
cutoff, the solution stays \(L^2\)-precompact, and its critical record
grows at most like \((T_*-b)^{-1/11}\) in norm.  Conversely, by
Assessment~\ref{ass:v32-high-total-debit}, a cofinal high-total branch
forces \(\int_0^{T_*}\|\Lambda u\|_2^{5/2}\,\dd t=\infty\), a
type-II enstrophy behaviour in the \(L^{5/2}\) sense which
Proposition~\ref{prop:v32-ledger-exponents} shows is not implied by the
Leray rate.  The two regimes are therefore separated by the finiteness of
one ledger, and neither has yet been contradicted.
\end{firewall}

\begin{corollary}[Cofinal high-total records force the \(L^{5/2}\) ledger
to diverge]
\label{cor:v32-high-total-divergence}
Let \(T_*<\infty\), \(\theta,\kappa>0\), and suppose there are
\(b_j\to T_*\) and \(\rho_j\to\infty\) with
\(H_{\rho_j}^\Sigma(b_j)\ge(1+\kappa)\theta E_*\).  Then
\begin{equation}
 \int_0^{T_*}\|\Lambda u(t)\|_2^{5/2}\,\dd t=\infty .
 \label{eq:v32-high-total-divergence}
\end{equation}
\end{corollary}

\begin{proof}
Let \(\tau_j=\theta^{-1}(\nu\rho_j^2)^{-1}\to0\).  Passing to a
subsequence, the windows \([b_j-\tau_j,b_j]\) are pairwise disjoint and
contained in \([0,T_*)\).  By \eqref{eq:v32-high-total-window-debit} each
carries ledger at least \(((1+\kappa)\theta-e^{-2/\theta})E_*^{3/4}/(2C_4)\),
which is positive because \(e^{-2/\theta}<\theta\) for every
\(\theta>0\).  Sum over \(j\).
\end{proof}

% =============================================================================
\section{Sharp tests and non-duplication}
\label{sec:v32-tests}
% =============================================================================

\begin{proposition}[The V14/V18 packet is low-total and pays a vanishing
window ledger]
\label{prop:v32-packet-test}
For the packet solutions of Theorem~\ref{thm:v14-two-parameter-packet},
\begin{equation}
 \sup_{J_j}\|Q_{K_j}u^{(j)}\|_2^2\le C\frac{R_j^2}{N_j}\to0,
 \qquad
 \int_{J_j}\|\Lambda u^{(j)}\|_2^{5/2}\,\dd t
 \le C\Bigl(\frac{R_j^2}{N_j}\Bigr)^{3/4}\to0 .
 \label{eq:v32-packet-ledgers}
\end{equation}
Thus the packet never satisfies the hypothesis of
\eqref{eq:v32-high-total-window-debit}, and
Theorem~\ref{thm:v32-recent-ledger-tail} is consistent with it without
supplying record separation.  On its interval the persistent term
\(e^{-2\nu\rho^2|J_j|}E_*\) is not small at any \(\rho\le N_j\), because
\(\nu N_j^2|J_j|=\nu s_\#/R_j\to0\): the packet lives on a sub-parabolic
interval at every relevant scale.
\end{proposition}

\begin{proof}
The tail bound is \eqref{eq:v14-packet-high-energy}.  On \(J_j\) the
packet field has \(\|\Lambda u^{(j)}\|_2^2\le C\beta_j^2N_j^2=CR_j^2N_j\)
up to the fixed reservoir, and \eqref{eq:v14-packet-Leray-debit} gives
\(\int_{J_j}\|\Lambda u^{(j)}\|_2^2\le CR_j/N_j\).  Hence
\(\int_{J_j}\|\Lambda u\|_2^{5/2}\le(\sup\|\Lambda u\|_2^2)^{1/4}
\int_{J_j}\|\Lambda u\|_2^2\le C(R_j^2N_j)^{1/4}R_j/N_j
=C(R_j^2/N_j)^{3/4}\), which tends to zero by
\eqref{eq:v14-two-parameter-separation}.  The last claim is the
interval length \eqref{eq:v14-two-parameter-interval}.
\end{proof}

\begin{proposition}[The Ladyzhenskaya density is sharp on one shell]
\label{prop:v32-density-sharpness}
For every lattice radius \(\lambda\) and every divergence-free
trigonometric polynomial \(u\) supported on \(\{|k|=\lambda\}\),
\[
 \|u\|_2^{1/2}\|\Lambda u\|_2^{3/2}
 =\|\Lambda^{3/4}u\|_2^2
 =\lambda^{3/2}\|u\|_2^2 .
\]
Hence both lattice H\"older steps in Lemma~\ref{lem:v32-ladyzhenskaya} are
equalities on one shell, and no power of the cutoff can be gained in
\eqref{eq:v32-flux-cutoff-free} from those steps alone.  The only
cutoff-dependent gain available is the slot structure of
Corollary~\ref{cor:v32-slots}, exhibited in
\eqref{eq:v32-HH-slot-bound}.
\end{proposition}

\begin{proof}
On one shell \(|k|=\lambda\) is constant, so every H\"older step is an
identity.  The remaining assertion is the content of
Theorem~\ref{thm:v32-cutoff-free-density}.
\end{proof}

\begin{firewall}[No double charge against V29 persistence]
\label{fw:v32-no-double-charge}
The persistent term \(e^{-2\nu\rho^2\tau}H_\rho^\Sigma(b-\tau)\) in
\eqref{eq:v32-recent-ledger-tail} is the linear heat decay of the tail
already present at the start of the window; it is paid by nothing.  The
V29 persistence rectangles are paid once by the Leray occupation
\eqref{eq:v29-cofinal-persistence-sum}.  The window ledger
\(\mathscr L_{5/2}(b,\tau)\) charges only the nonlinear influx during the
window.  The three payments are disjoint and are not to be added as a
second debit for the same interval.
\end{firewall}

\begin{firewall}[Non-duplication audit]
\label{fw:v32-nonduplication}
V4 bounds the net time-integrated flux by \(E_*/2\) and the V7--V14
sections bound the source action by the record--Leray product; neither uses
\(L^4\) norms, the density \(\|\Lambda u\|_2^{5/2}\), or a window ledger.
V29--V31 use Leray occupation, probability marks, and heat formation.  The
recent-ledger tail bound, the record formula
\eqref{eq:v32-record-optimized}, and the type-I firewall are new relative
to V1--V31 and to the Grand Tensor NS section, whose relative-form
continuation certificate is conditional on an \(H^s\) relative estimate
that V32 does not supply.
\end{firewall}

% =============================================================================
\section{Integrated V32 conclusion and next acquisition order}
\label{sec:v32-conclusion}
% =============================================================================

\begin{theorem}[What V32 proves]
\label{thm:v32-conclusion}
Relative to V1--V31, the following are proved.
\begin{enumerate}[label=\textup{(V32.\arabic*)},leftmargin=6.2em]
\item The cutoff flux is exactly the difference of a low--low stress acting
      on the first exterior octave and a high--high stress acting on the
      low strain, \eqref{eq:v32-two-stress}; the two stresses are the
      short-gap and long lever-arm slots of every complete V18 triad.
\item Both the total flux and the sector-difference current are bounded at
      every cutoff by the one density
      \(C_4E_*^{1/4}\|\Lambda u\|_2^{5/2}\),
      \eqref{eq:v32-flux-cutoff-free}--%
      \eqref{eq:v32-difference-cutoff-free}; the high--high slot carries
      an additional factor \(\rho^{-1/2}\) against the tail dissipation.
\item All three V29 recent-entry alternatives are certified by one window
      number, Theorem~\ref{thm:v32-exit-debit}; the moving-family
      multiplicity of Firewall~\ref{fw:v29-entry-not-summed} is absent in
      the \(L^{5/2}_tH^1\) ledger.
\item The tail energy at any cutoff is bounded by heat decay plus the recent
      window ledger, \eqref{eq:v32-recent-ledger-tail}; a high-total tail
      costs an order-\(E_*^{3/4}\) window ledger with no record factor.
\item The critical record is bounded by a parabolic persistent term plus the
      terminal enstrophy times the square root of the window ledger,
      \eqref{eq:v32-record-optimized}.
\item Leray's rate holds on the torus, forces divergence of the
      \(L^p_tH^1\) ledger exactly for \(p\ge4\), and leaves the
      \(L^{5/2}\) ledger undetermined.
\item Under a type-I enstrophy rate the spectral tails are uniformly small,
      the trajectory is \(L^2\)-precompact with no energy jump, every
      cofinal high-total branch is empty, and the critical record grows at
      most like \((T_*-b)^{-2/11}\).
\item Conversely a cofinal high-total branch forces
      \(\int_0^{T_*}\|\Lambda u\|_2^{5/2}\,\dd t=\infty\).  The V14/V18
      packet is low-total and pays a vanishing window ledger.
\end{enumerate}
The full Navier--Stokes stability/global-regularity theorem remains
unproved.  The open gate is now a two-regime alternative: either the
\(L^{5/2}_tH^1\) ledger diverges at \(T_*\), in which case the high-total
branches are alive but every one of them costs an order-one window ledger;
or it stays finite, in which case all high-total branches are dead and the
singularity must be carried by uniformly small, radially spread tails.
\end{theorem}

\begin{proof}
Items \textup{(V32.1)}--\textup{(V32.2)} are
Proposition~\ref{prop:v32-two-stress},
Corollary~\ref{cor:v32-slots}, and
Theorem~\ref{thm:v32-cutoff-free-density}.  Item \textup{(V32.3)} is
Theorem~\ref{thm:v32-exit-debit} and
Corollary~\ref{cor:v32-no-multiplicity}.  Items
\textup{(V32.4)}--\textup{(V32.5)} are
Theorems~\ref{thm:v32-recent-ledger-tail} and
\ref{thm:v32-record-from-ledger}.  Item \textup{(V32.6)} is
Lemma~\ref{lem:v32-leray-rate} and
Proposition~\ref{prop:v32-ledger-exponents}.  Items
\textup{(V32.7)}--\textup{(V32.8)} are
Theorem~\ref{thm:v32-type-I-firewall},
Corollary~\ref{cor:v32-high-total-divergence}, and
Proposition~\ref{prop:v32-packet-test}.
\end{proof}

\begin{programme}[V33 low-total spread acquisition order]
\label{prog:v33}
\begin{enumerate}[label=\textup{(AC\arabic*)},leftmargin=4.8em]
\item Work in the low-total regime: assume
      \(\sup_{\rho\ge K}H_\rho^\Sigma(t)\le\eta\) on a terminal interval,
      as forced by Theorem~\ref{thm:v32-type-I-firewall}, and prove the
      exact radial spread that a doubled record
      \(\|A^{1/4}u\|_2^2\ge\mathcal R_n^2\) then requires above
      \(K_n/2\).  Quantify the enstrophy and the
      \(\|A^{3/4}u\|_2^2\) dissipation of the spread configuration.
\item Insert the tail smallness into the two slot bounds
      \eqref{eq:v32-LL-slot-bound}--\eqref{eq:v32-HH-slot-bound}:
      the high--high slot gains \(H_\rho^{1/4}\le\eta^{1/4}\) through
      \(\|Q_\rho u\|_4^2\le C_4H_\rho^{1/4}(V_\rho^\Sigma)^{3/4}\), and
      the low--low slot is confined to the first octave.  Determine whether
      the \(\dot H^{1/2}\) record balance
      \(\frac12\frac{\dd}{\dd t}\|A^{1/4}u\|_2^2
      +\nu\|A^{3/4}u\|_2^2=\int_0^\infty\Phi_\rho\,\dd\rho\)
      then acquires a relative-form bound of the Grand Tensor type on the
      spread regime.
\item Sharpen Theorem~\ref{thm:v32-record-from-ledger} by replacing the
      single window with the V31 dyadic windows
      \(\tau_n\asymp(\nu K_n^2)^{-1}\log\) and the flat marks
      \(\mathcal M_n\); obtain the record as a half-order fractional
      integral of the \(L^{5/2}\) density and test it on the type-I
      profile and on the packet.
\item In the high-total regime, combine
      Corollary~\ref{cor:v32-spike} with the head bound
      \eqref{eq:v4-terminal-quartic-bound}: the spike must lie in the
      tail above \(K_n\), and its tail dissipation is Leray-paid.  Determine
      whether the spike height \(\epsilon_j^{-4}\) is compatible with the
      corridor length and the one-use Leray share.
\item Retain the V29 persistence sum without a second charge; do not use
      the hinge on the total branch.  The next compulsory companion audit
      is V35.
\end{enumerate}
\end{programme}

\begin{assessment}[Position after V32]
\label{ass:v32-position}
V32 went upstream of the V31 correlation target and found that the
cutoff-dependence which had blocked every summation since V19 is an artefact
of insisting on the Leray ledger.  In the intermediate ledger
\(L^{5/2}_tH^1\) the exits, the tails, and the critical record are all
paid by one window number, with no record factor and no parent tail.  What
is bought is a clean regime split: high-total cascades are type-II in the
\(L^{5/2}\) sense, and type-I enstrophy rates are low-total with
\(L^2\)-precompact trajectories and a quantitative record rate.  What is
not bought is a contradiction in either regime.  The next work is on the
low-total spread regime, where the new slot bounds carry an explicit
smallness factor.  No full stability or global-regularity claim is made.
\end{assessment}

% =============================================================================
\section*{Version V32 append-only record}
\addcontentsline{toc}{section}{Version V32 append-only record}
% =============================================================================

The complete V31 mathematical source was retained before this continuation.
Its SHA--256 digest was
\begin{center}\scriptsize\ttfamily
f1121b62238ad5491bb7cf03635ea9934942edf573ed8faaa9ee79e1d38a6696.
\end{center}
V32 proved the two-stress anatomy of the cutoff flux, the cutoff-free
Ladyzhenskaya density for the total and sector-difference tail balances,
the one-number certification of all V29 recent-entry exits, the
recent-ledger tail and record bounds, Leray's rate on the torus with the
exact ledger-exponent threshold \(p=4\), and the type-I enstrophy firewall
with its \((T_*-b)^{-2/11}\) record rate.  A cofinal high-total branch was
shown to force divergence of the \(L^{5/2}_tH^1\) ledger.  No full
regularity claim is made.

% =============================================================================
\section{V33: the critical record as a half-order fractional integral}
\label{sec:v33-entry}
% =============================================================================

V32 bounded the tail at one cutoff by the recent \(L^{5/2}\) ledger on one
window, and paid the critical record with a persistent parabolic term plus a
window ledger.  The persistent term \(1.04E_*(\nu\tau)^{-1/2}\) was the
price of cutting the history at \(b-\tau\).  This version removes the cut.
Keeping the heat weights \(e^{-2\nu\rho^2(b-t)}\) inside the dyadic sum of
V31's flat marks, and summing them over \emph{all} dyadic scales, produces
the kernel \((\nu(b-t))^{-1/2}\) exactly.  The result is a fractional
integral law: for any density \(\varphi(t)\) that dominates twice the
cutoff flux uniformly in the cutoff,
\[
 \|A^{1/4}u(b)\|_2^2
 \le2\|A^{1/4}u(a)\|_2^2
 +2.3\,\nu^{-1/2}\int_a^b\frac{\varphi(t)}{(b-t)^{1/2}}\,\dd t .
\]
Two densities are available from V32: the Ladyzhenskaya density
\(2C_4E_*^{1/4}\|\Lambda u\|_2^{5/2}\) and the record--Leray density
\(4C_4\|A^{1/4}u\|_2\|\Lambda u\|_2^2\).

With the second density the law closes on itself.  The running record is
bounded by the running half-order fractional integral of the Leray
dissipation density, unconditionally.  Under a type-I enstrophy rate the
fractional integral of \((T_*-t)^{-1/2}\) is a logarithm, and a one-line
bootstrap gives
\[
 \|A^{1/4}u(b)\|_2^2\le C\bigl(1+\log^2\tfrac{\delta}{T_*-b}\bigr),
\]
which supersedes the \((T_*-b)^{-2/11}\) rate of V32: under a type-I
enstrophy rate the critical norm can blow up at most logarithmically.  The
V1 record doubling times are then double-exponentially close to \(T_*\).
The low-total spread lemma and the high-total spike dichotomy requested in
Programme~\ref{prog:v33} are also settled, in the negative sense that
neither contradicts the Leray budget.  No global regularity claim is made.
No companion review is due in V33; the next compulsory SM/YM
synchronization remains V35.

% =============================================================================
\section{The dyadic heat sum has the half-order kernel}
\label{sec:v33-dyadic-heat-sum}
% =============================================================================

Fix any \(K>0\) and let \(K_n:=2^nK\) for \emph{all} integers
\(n\in\mathbb Z\), so that the annuli \(\mathcal A_n=[K_n,K_{n+1})\)
partition \((0,\infty)\).

\begin{lemma}[Two-sided dyadic domination of the critical record]
\label{lem:v33-dyadic-domination}
For every square-integrable divergence-free \(v\),
\begin{equation}
 \frac12\sum_{n\in\mathbb Z}K_nH_{K_n}[v]
 \le\|A^{1/4}v\|_2^2
 =\int_0^\infty H_\rho[v]\,\dd\rho
 \le\sum_{n\in\mathbb Z}K_nH_{K_n}[v].
 \label{eq:v33-dyadic-domination}
\end{equation}
\end{lemma}

\begin{proof}
The identity is the layer cake
\(\sum_k|k||\widehat v(k)|^2=\sum_k|\widehat v(k)|^2\int_0^{|k|}\dd\rho\).
Since \(\rho\mapsto H_\rho\) is nonincreasing,
\(\int_{\mathcal A_n}H_\rho\,\dd\rho\le K_nH_{K_n}\) and
\(\int_{[K_n/2,K_n)}H_\rho\,\dd\rho\ge\frac12K_nH_{K_n}\).  Sum over
\(n\); the intervals \([K_n/2,K_n)\) are pairwise disjoint.
\end{proof}

\begin{lemma}[The dyadic heat sum is the half-order kernel]
\label{lem:v33-dyadic-heat-sum}
For every \(s>0\) and every \(K>0\),
\begin{equation}
 \boxed{\quad
 \sum_{n\in\mathbb Z}K_ne^{-2\nu K_n^2s}
 \le\frac{2.3}{(\nu s)^{1/2}} .
 \quad}
 \label{eq:v33-dyadic-heat-sum}
\end{equation}
\end{lemma}

\begin{proof}
Put \(x_n:=K_n(\nu s)^{1/2}=2^nx_0\).  The terms with \(x_n\le1\) form a
geometric sequence with largest term at most \(1\), so their sum is at most
\(2\).  If \(x_{n_1}\) is the smallest term exceeding \(1\), then
\(x_{n_1}\in(1,2]\) and the remaining terms are
\(2^mx_{n_1}e^{-2\cdot4^mx_{n_1}^2}\le2^{m+1}e^{-2\cdot4^m}\),
\(m\ge0\), whose sum is \(2e^{-2}+4e^{-8}+8e^{-32}+\dots<0.28\).  Divide
by \((\nu s)^{1/2}\).
\end{proof}

% =============================================================================
\section{The fractional-integral record law}
\label{sec:v33-record-law}
% =============================================================================

Call a measurable \(\varphi\ge0\) on \([a,b]\) an \emph{admissible flux
density} if
\begin{equation}
 2\,|\Phi_\rho(t)|\le\varphi(t)
 \qquad\text{for every }\rho>0\text{ and a.e. }t\in[a,b].
 \label{eq:v33-admissible-density}
\end{equation}
By Theorem~\ref{thm:v32-cutoff-free-density} and
Proposition~\ref{prop:v32-corridor-form}, both
\begin{equation}
 \varphi_{5/2}(t):=2C_4E_*^{1/4}\|\Lambda u(t)\|_2^{5/2},
 \qquad
 \varphi_{\rm rec}(t):=4C_4\|A^{1/4}u(t)\|_2\|\Lambda u(t)\|_2^2,
 \label{eq:v33-two-densities}
\end{equation}
and their pointwise minimum, are admissible on every smooth interval.

\begin{theorem}[Fractional-integral law for the critical record]
\label{thm:v33-record-law}
Let \(\varphi\) be admissible on \([a,b]\subset[0,T_*)\).  Then
\begin{equation}
 \boxed{\quad
 \|A^{1/4}u(b)\|_2^2
 \le2\|A^{1/4}u(a)\|_2^2
 +\frac{2.3}{\nu^{1/2}}
 \int_a^b\frac{\varphi(t)}{(b-t)^{1/2}}\,\dd t .
 \quad}
 \label{eq:v33-record-law}
\end{equation}
In particular
\begin{align}
 \|A^{1/4}u(b)\|_2^2
 &\le2\|A^{1/4}u(a)\|_2^2
 +\frac{4.6\,C_4E_*^{1/4}}{\nu^{1/2}}
 \int_a^b\frac{\|\Lambda u(t)\|_2^{5/2}}{(b-t)^{1/2}}\,\dd t,
 \label{eq:v33-record-law-52}\\
 \|A^{1/4}u(b)\|_2^2
 &\le2\|A^{1/4}u(a)\|_2^2
 +\frac{9.2\,C_4}{\nu^{1/2}}
 \int_a^b\frac{\|A^{1/4}u(t)\|_2\|\Lambda u(t)\|_2^2}{(b-t)^{1/2}}\,\dd t .
 \label{eq:v33-record-law-rec}
\end{align}
\end{theorem}

\begin{proof}
As in the proof of Theorem~\ref{thm:v32-recent-ledger-tail}, for every
\(\rho>0\),
\(\partial_tH_\rho^\Sigma+2\nu\rho^2H_\rho^\Sigma\le\varphi\), hence
\[
 H_{K_n}^\Sigma(b)
 \le e^{-2\nu K_n^2(b-a)}H_{K_n}^\Sigma(a)
 +\int_a^be^{-2\nu K_n^2(b-t)}\varphi(t)\,\dd t .
\]
Multiply by \(K_n\), sum over \(n\in\mathbb Z\), and use
Lemma~\ref{lem:v33-dyadic-domination} on the left and for the initial
term, where \(e^{-2\nu K_n^2(b-a)}\le1\).  For the source term, Tonelli
and Lemma~\ref{lem:v33-dyadic-heat-sum} with \(s=b-t\) give the kernel.
The two displayed cases insert \eqref{eq:v33-two-densities}.
\end{proof}

\begin{remark}[Scaling and the removed persistent term]
\label{rem:v33-scaling}
Under \(u_\lambda(t,x)=\lambda u(\lambda^2t,\lambda x)\) both sides of
\eqref{eq:v33-record-law-52} and \eqref{eq:v33-record-law-rec} are
invariant: the critical record is scale-free, and
\(\|\Lambda u\|_2^{5/2}E_*^{1/4}\,\dd t\,(b-t)^{-1/2}\) carries weight
\(\lambda^{5/4-1/4-2+1}=\lambda^0\).  The V32 persistent term
\(1.04E_*(\nu\tau)^{-1/2}\) has disappeared because the heat weights were
kept inside the dyadic sum; the only memory of the past is the record at
the starting time, with a factor two.
\end{remark}

\begin{corollary}[The running record is paid by the Leray fractional
integral]
\label{cor:v33-running-record}
For \(a\le b<T_*\) put
\begin{equation}
 F(b):=\sup_{a\le s\le b}\|A^{1/4}u(s)\|_2^2,
 \qquad
 \mathcal I(s):=\int_a^s\frac{\|\Lambda u(t)\|_2^2}{(s-t)^{1/2}}\,\dd t,
 \qquad
 \mathcal I^*(b):=\sup_{a\le s\le b}\mathcal I(s).
 \label{eq:v33-running-quantities}
\end{equation}
Then
\begin{equation}
 \boxed{\quad
 F(b)^{1/2}
 \le\sqrt2\,\|A^{1/4}u(a)\|_2
 +\frac{9.2\,C_4}{\nu^{1/2}}\,\mathcal I^*(b).
 \quad}
 \label{eq:v33-running-record}
\end{equation}
Moreover \(\mathcal I\) is Leray-integrable:
\begin{equation}
 \int_a^b\mathcal I(s)\,\dd s
 \le2(b-a)^{1/2}\int_a^b\|\Lambda u\|_2^2\,\dd t
 \le\frac{(b-a)^{1/2}E_*}{\nu}.
 \label{eq:v33-fractional-integral-L1}
\end{equation}
\end{corollary}

\begin{proof}
Apply \eqref{eq:v33-record-law-rec} at every \(s\in[a,b]\), bound
\(\|A^{1/4}u(t)\|_2\le F(b)^{1/2}\) inside the integral, and take the
supremum: \(F(b)\le2\|A^{1/4}u(a)\|_2^2+9.2C_4\nu^{-1/2}F(b)^{1/2}
\mathcal I^*(b)\).  The quadratic inequality
\(F\le A+BF^{1/2}\) implies \(F^{1/2}\le\sqrt A+B\).  For
\eqref{eq:v33-fractional-integral-L1}, exchange the order of integration:
\(\int_t^b(s-t)^{-1/2}\dd s\le2(b-a)^{1/2}\), then use the Leray
identity.
\end{proof}

\begin{assessment}[What the record law says about the V1 corridor]
\label{ass:v33-record-law-meaning}
The V1 records \(\mathcal R_n=2^n\mathcal R_*\) are reached only at
times \(b_n\) where the running Leray fractional integral satisfies
\(\mathcal I^*(b_n)\ge(\mathcal R_n-\sqrt2\mathcal R_0)\nu^{1/2}/(9.2C_4)\).
Since \(\mathcal I\) has the Leray-paid time average
\eqref{eq:v33-fractional-integral-L1}, the record can be large only on a
time set of small measure, in agreement with the trivial
\(\|A^{1/4}u\|_2^4\le E_*\|\Lambda u\|_2^2\in L^1_t\); the fractional law
adds no new time integrability.  What it adds is pointwise control near the
terminal time in terms of the concentration of the Leray density, which is
what the type-I bootstrap below uses.  Blow-up with unbounded critical
record requires \(\mathcal I^*(b)\to\infty\); by
Lemma~\ref{lem:v32-leray-rate} the Leray rate already gives
\(\mathcal I(b)\ge c_L^2\nu^{3/2}\int(T_*-t)^{-1/2}(b-t)^{-1/2}\,\dd t
\to\infty\) logarithmically, so no contradiction arises here.
\end{assessment}

% =============================================================================
\section{Type-I enstrophy rate allows only logarithmic critical growth}
\label{sec:v33-type-I-log}
% =============================================================================

\begin{lemma}[The half-order fractional integral of the Leray rate is a
logarithm]
\label{lem:v33-log-kernel}
For \(T_*-\delta\le b<T_*\),
\begin{equation}
 \int_{T_*-\delta}^b\frac{\dd t}{(T_*-t)^{1/2}(b-t)^{1/2}}
 =2\operatorname{arcsinh}\Bigl(\frac{b-T_*+\delta}{T_*-b}\Bigr)^{1/2}
 \le\log\frac{9\delta}{T_*-b}=:L_\delta(b).
 \label{eq:v33-log-kernel}
\end{equation}
\end{lemma}

\begin{proof}
With \(s=b-t\), \(c=T_*-b\), and \(A=b-T_*+\delta\), the integral is
\(\int_0^A(c+s)^{-1/2}s^{-1/2}\,\dd s=2\operatorname{arcsinh}\sqrt{A/c}\).
Since \(\operatorname{arcsinh}x\le\log(2x+1)\) for \(x\ge0\) and
\(A\le\delta\), \(2\operatorname{arcsinh}\sqrt{A/c}
\le2\log\bigl(1+2\sqrt{\delta/c}\bigr)\le2\log\bigl(3\sqrt{\delta/c}\bigr)
=\log(9\delta/c)\), using \(\delta\ge c\).
\end{proof}

\begin{theorem}[Type-I enstrophy rate: the critical norm grows at most like
\(\log\)]
\label{thm:v33-type-I-log}
Assume the type-I enstrophy rate \eqref{eq:v32-type-I} on
\([T_*-\delta,T_*)\), and let
\(R_\delta:=\sup_{0\le t\le T_*-\delta}\|A^{1/4}u(t)\|_2\), which is
finite by smoothness.  Then for every \(b\in[T_*-\delta,T_*)\),
\begin{equation}
 \boxed{\quad
 \|A^{1/4}u(b)\|_2^2
 \le2A_\delta+2B^2L_\delta(b)^2,
 \qquad
 A_\delta:=2R_\delta^2+\frac{4.6\,C_4E_*R_\delta}{\nu^{3/2}\delta^{1/2}},
 \qquad
 B:=9.2\,C_4C_I^2\nu .
 \quad}
 \label{eq:v33-type-I-log-bound}
\end{equation}
Consequently \(\|u(b)\|_{\dot H^{1/2}}\le C\log(9\delta/(T_*-b))\) for
\(b\) close to \(T_*\), and the V1 record doubling times satisfy
\begin{equation}
 T_*-b_n\le9\delta\exp\bigl(-c\,2^n\bigr),
 \qquad c:=\frac{\mathcal R_*}{2B}\ (\text{for }n\ge n_0).
 \label{eq:v33-record-times}
\end{equation}
This supersedes the rate \eqref{eq:v32-type-I-record-rate}.
\end{theorem}

\begin{proof}
Apply \eqref{eq:v33-record-law-rec} with \(a=0\) and split the integral
at \(T_*-\delta\).  On \([0,T_*-\delta]\), \(\|A^{1/4}u\|_2\le R_\delta\)
and \((b-t)^{-1/2}\le\delta^{-1/2}\), so that part is at most
\(9.2C_4\nu^{-1/2}R_\delta\delta^{-1/2}E_*/(2\nu)\) by the Leray
identity.  On \([T_*-\delta,b]\), with
\(F(b):=\sup_{T_*-\delta\le s\le b}\|A^{1/4}u(s)\|_2^2\) and
\eqref{eq:v32-type-I},
\[
 \int_{T_*-\delta}^b
 \frac{\|A^{1/4}u\|_2\|\Lambda u\|_2^2}{(b-t)^{1/2}}\,\dd t
 \le F(b)^{1/2}C_I^2\nu^{3/2}L_\delta(b)
\]
by Lemma~\ref{lem:v33-log-kernel}.  Hence every \(s\in[T_*-\delta,b]\)
satisfies \(\|A^{1/4}u(s)\|_2^2\le A_\delta+BF(b)^{1/2}L_\delta(b)\),
because \(F(s)\le F(b)\) and \(L_\delta(s)\le L_\delta(b)\).  Taking the
supremum, \(F\le A_\delta+BF^{1/2}L_\delta\), so
\(F^{1/2}\le\sqrt{A_\delta}+BL_\delta\) and
\(F\le2A_\delta+2B^2L_\delta^2\).  For \eqref{eq:v33-record-times}, the
record \(\mathcal R_n^2=4^n\mathcal R_*^2\) at time \(b_n\) forces
\(2^n\mathcal R_*\le\sqrt{2A_\delta}+\sqrt2BL_\delta(b_n)\), and for
\(n\ge n_0\) the first term is at most half the left side.
\end{proof}

\begin{firewall}[The logarithm is the exact fixed point of this route]
\label{fw:v33-log-fixed-point}
Theorem~\ref{thm:v33-type-I-log} cannot be improved to a bounded critical
norm as long as the admissible density is a bounded multiple of
\(\|\Lambda u\|_2^2\), because under \eqref{eq:v32-type-I} the half-order
fractional integral of \(\|\Lambda u\|_2^2\) is exactly the divergent
logarithm of Lemma~\ref{lem:v33-log-kernel}, and
Lemma~\ref{lem:v32-leray-rate} shows that the divergence is forced, not
assumed.  Any improvement must therefore come from a density that is
smaller than \(\|A^{1/4}u\|_2\|\Lambda u\|_2^2\) on the parabolic scales
that carry the logarithm.  A bounded critical norm under
the type-I enstrophy rate would, together with the endpoint
\(L^\infty_tL^3_x\) theorem of \cite{ESS} which this series has not
invoked, exclude type-I enstrophy blow-up altogether.  V33 does not claim
that.  It claims only that the type-I regime is now confined to at most
logarithmic critical growth with uniformly small spectral tails.
\end{firewall}

% =============================================================================
\section{Low-total spread and the high-total spike, settled}
\label{sec:v33-spread-spike}
% =============================================================================

\begin{lemma}[Low-total records are enstrophy-expensive]
\label{lem:v33-spread}
If \(H_\rho^\Sigma(t)\le\eta\) for every \(\rho\ge K\), then
\begin{equation}
 \boxed{\quad
 \|\Lambda u(t)\|_2
 \ge\frac{\bigl(\|A^{1/4}u(t)\|_2^2-KE_*\bigr)_++\eta K}
        {(2\eta)^{1/2}} .
 \quad}
 \label{eq:v33-spread}
\end{equation}
In particular, at a V1 record time with
\(K=K_n/2=\mathcal R_n^2/(2E_*)\),
\(\|\Lambda u\|_2\ge\mathcal R_n^2/(2(2\eta)^{1/2})\), to be compared
with the trivial \(\|\Lambda u\|_2\ge\mathcal R_n^2/E_*^{1/2}\).
\end{lemma}

\begin{proof}
For every \(\rho_1\ge K\),
\(\|A^{1/4}u\|_2^2\le KE_*+\eta(\rho_1-K)
+\int_{\rho_1}^\infty H_\rho^\Sigma\,\dd\rho\), and the last integral is
at most \(\|\Lambda u\|_2^2/(2\rho_1)\) by
\eqref{eq:v29-weighted-tail-layer-cake}.  Choose
\(\rho_1=\|\Lambda u\|_2(2\eta)^{-1/2}\), which is at least \(K\)
whenever the right side of \eqref{eq:v33-spread} is positive; otherwise
the claim is trivial.
\end{proof}

\begin{assessment}[Spread does not beat the record law]
\label{ass:v33-spread-versus-record}
Combining \eqref{eq:v33-spread} with the uniform tail smallness
\eqref{eq:v32-uniform-tail-smallness} and the type-I rate gives
\(\|A^{1/4}u\|_2^2\le KE_*+CK^{-9/32}(T_*-b)^{-1/4}\), hence at best
\((T_*-b)^{-8/41}\) after optimizing \(K\).  This is weaker than
Theorem~\ref{thm:v33-type-I-log}.  The spread lemma is retained for the
type-II regime, where it shows that any record carried by uniformly small
tails needs enstrophy \(\gtrsim\mathcal R_n^2\eta^{-1/2}\), a factor
\((E_*/\eta)^{1/2}\) above the trivial requirement.
\end{assessment}

\begin{proposition}[The high-total spike is Leray-consistent]
\label{prop:v33-spike-dichotomy}
Let \(w=[b-\tau,b]\subset I_n\) be a V32 high-total window at the record
cutoff \(K_n\) with ledger \(\mathscr L_{5/2}(b,\tau)\ge D\) and Leray
share \(\epsilon:=\mathscr L_2(b,\tau)\).  Then, with
\(V_{K_n}=\|\Lambda Q_{K_n}u\|_2^2\),
\begin{equation}
 \int_wV_{K_n}^2\,\dd t
 \ge\frac{D^4}{2\epsilon^3}-\frac{3\mathcal R_n^4}{\nu},
 \qquad
 \sup_wV_{K_n}
 \ge\frac1\epsilon\Bigl(\frac{D^4}{2\epsilon^3}
 -\frac{3\mathcal R_n^4}{\nu}\Bigr).
 \label{eq:v33-spike-dichotomy}
\end{equation}
Hence either \(\epsilon\ge(\nu D^4/12)^{1/3}\mathcal R_n^{-4/3}\), or the
tail dissipation above \(K_n\) has a spike of height at least
\(D^4/(4\epsilon^4)\) inside the window.  Both alternatives are compatible
with \(\sum_n\epsilon_n\le E_*/(2\nu)\).
\end{proposition}

\begin{proof}
By \eqref{eq:v32-spike}, \(\int_w\|\Lambda u\|_2^4\ge D^4/\epsilon^3\).
Write \(\|\Lambda u\|_2^2=E_{K_n}+V_{K_n}\), so
\(\|\Lambda u\|_2^4\le2E_{K_n}^2+2V_{K_n}^2\), and bound
\(\int_wE_{K_n}^2\le3\mathcal R_n^4/\nu\) by
\eqref{eq:v4-terminal-quartic-bound}.  Then
\(\int_wV_{K_n}^2\le\sup_wV_{K_n}\int_wV_{K_n}\le\epsilon\sup_wV_{K_n}\).
The first alternative is the case \(3\mathcal R_n^4/\nu\ge D^4/(4\epsilon^3)\);
otherwise the bracket exceeds \(D^4/(4\epsilon^3)\), which is the second.  For the last
claim, \(\sum_n\mathcal R_n^{-4/3}<\infty\), and a spike of finite height on
a window of length \(\tau\asymp(\nu K_n^2)^{-1}\) contributes at most
\(\tau\sup_wV_{K_n}\) to the Leray share, which is not constrained from
below.
\end{proof}

% =============================================================================
\section{Tests}
\label{sec:v33-tests}
% =============================================================================

\begin{proposition}[Neither test field saturates the record law]
\label{prop:v33-tests}
\begin{enumerate}[label=\textup{(\roman*)},leftmargin=3.7em]
\item For the V14/V18 packet on \(J_j\), the right side of
      \eqref{eq:v33-record-law-rec} is of order
      \(\nu^{-1/2}R_j^{5/2}\), while the left side is of order
      \(R_j^2\); the ratio \(\nu^{-1/2}R_j^{1/2}\to\infty\).
\item For the exact Leray-rate profile
      \(\|\Lambda u\|_2^2=c\nu^{3/2}(T_*-t)^{-1/2}\) with bounded critical
      norm, the right side of \eqref{eq:v33-record-law-rec} grows like
      \(\log(1/(T_*-b))\) while the left side stays bounded; the law is
      not saturated by the type-I profile either.
\end{enumerate}
\end{proposition}

\begin{proof}
(i) On \(J_j\), \(\|A^{1/4}u\|_2\asymp R_j\),
\(\|\Lambda u\|_2^2\lesssim R_j^2N_j\), and
\(\int_{J_j}(b-t)^{-1/2}\dd t=2|J_j|^{1/2}\asymp R_j^{-1/2}N_j^{-1}\),
so the integral is of order \(R_j^3N_jR_j^{-1/2}N_j^{-1}=R_j^{5/2}\).
(ii) is Lemma~\ref{lem:v33-log-kernel}.
\end{proof}

\begin{firewall}[Which fields would saturate it]
\label{fw:v33-saturation}
Equality in \eqref{eq:v33-record-law} requires simultaneous equality in the
dyadic domination \eqref{eq:v33-dyadic-domination}, in the heat sum
\eqref{eq:v33-dyadic-heat-sum}, and in the flux density, that is, a field
whose tail at every dyadic scale is created at the parabolic time of that
scale by a flux saturating Ladyzhenskaya.  This is the forced-heat geometry
of Proposition~\ref{prop:v31-forced-heat-sharpness} extended across all
scales at once.  Whether an actual Navier--Stokes trajectory can do this
is exactly the open type-II question; V33 does not decide it.
\end{firewall}

% =============================================================================
\section{Integrated V33 conclusion and next acquisition order}
\label{sec:v33-conclusion}
% =============================================================================

\begin{theorem}[What V33 proves]
\label{thm:v33-conclusion}
Relative to V1--V32, the following are proved.
\begin{enumerate}[label=\textup{(V33.\arabic*)},leftmargin=6.2em]
\item The two-sided dyadic domination
      \eqref{eq:v33-dyadic-domination} and the dyadic heat sum
      \eqref{eq:v33-dyadic-heat-sum} with the exact kernel
      \(2.3(\nu s)^{-1/2}\).
\item The fractional-integral record law \eqref{eq:v33-record-law} for
      every admissible flux density, with the two explicit densities
      \eqref{eq:v33-two-densities}; it is scale-invariant and has no
      persistent term.
\item The running record is bounded by the running half-order fractional
      integral of the Leray density, \eqref{eq:v33-running-record}; that
      integral is Leray-integrable in time.
\item The fractional integral of the Leray rate is exactly logarithmic,
      Lemma~\ref{lem:v33-log-kernel}.
\item Under a type-I enstrophy rate the critical norm grows at most
      logarithmically, \eqref{eq:v33-type-I-log-bound}, and the record
      doubling times are double-exponentially close to \(T_*\).
\item The low-total spread lemma \eqref{eq:v33-spread} and the high-total
      spike dichotomy \eqref{eq:v33-spike-dichotomy}; neither contradicts
      the Leray budget.
\item Neither the packet nor the Leray-rate profile saturates the record
      law.
\end{enumerate}
The full Navier--Stokes stability/global-regularity theorem remains
unproved.  The V32 two-regime alternative is sharpened: in the type-I regime
the critical norm is at most logarithmic and the tails are uniformly small;
in the type-II regime the \(L^{5/2}\) ledger diverges and the record is a
half-order fractional integral of a divergent density.
\end{theorem}

\begin{proof}
Items \textup{(V33.1)}--\textup{(V33.2)} are
Lemmas~\ref{lem:v33-dyadic-domination}--\ref{lem:v33-dyadic-heat-sum}
and Theorem~\ref{thm:v33-record-law}.  Item \textup{(V33.3)} is
Corollary~\ref{cor:v33-running-record}.  Items
\textup{(V33.4)}--\textup{(V33.5)} are Lemma~\ref{lem:v33-log-kernel} and
Theorem~\ref{thm:v33-type-I-log}.  Item \textup{(V33.6)} is
Lemma~\ref{lem:v33-spread} and
Proposition~\ref{prop:v33-spike-dichotomy}.  Item \textup{(V33.7)} is
Proposition~\ref{prop:v33-tests}.
\end{proof}

\begin{programme}[V34 acquisition order]
\label{prog:v34}
\begin{enumerate}[label=\textup{(AC\arabic*)},leftmargin=4.8em]
\item Attack the logarithm.  Replace the record--Leray density in
      \eqref{eq:v33-record-law-rec} by a slot-resolved density built from
      \eqref{eq:v32-LL-slot-bound}--\eqref{eq:v32-HH-slot-bound}, using the
      uniform tail smallness \eqref{eq:v32-uniform-tail-smallness} in the
      high--high slot.  Determine whether the flux at cutoffs above the
      parabolic frequency of \(T_*-t\) is small enough that the fractional
      integral converges.  A bounded critical norm under the type-I
      enstrophy rate would be the target; its consequence through the
      endpoint theorem must be stated as an import, not proved here.
\item In the type-II regime, combine
      Corollary~\ref{cor:v32-high-total-divergence} with the record law:
      the record at \(b_n\) is bounded by the fractional integral of
      \(\|\Lambda u\|_2^{5/2}\), whose window pieces are the V32
      high-total debits.  Determine the minimal growth of
      \(\mathscr L_{5/2}\) forced by a doubled-record sequence and compare
      with the spike dichotomy.
\item Examine whether the fractional law holds for the dyadic tails
      individually, \(K_nH_{K_n}(b)\), rather than for their sum, giving a
      scale-resolved record law that could feed the V31 flat marks.
\item Prepare the compulsory V35 companion audit: freeze the current SM and
      YM notes by digest and record what, if anything, they contribute to
      the type-I/type-II split.
\end{enumerate}
\end{programme}

\begin{assessment}[Position after V33]
\label{ass:v33-position}
V33 converted the V32 window bounds into a global fractional-integral law
for the critical record with no persistent term.  The law is unconditional
and scale-invariant.  Its principal consequence is that a type-I enstrophy
rate confines the critical norm to logarithmic growth, so that a type-I
singularity, if one exists, is an almost-critical event with uniformly
small spectral tails.  The type-II regime is untouched except that its
record is now a fractional integral of the divergent \(L^{5/2}\) density.
No full stability or global-regularity claim is made.
\end{assessment}

% =============================================================================
\section*{Version V33 append-only record}
\addcontentsline{toc}{section}{Version V33 append-only record}
% =============================================================================

The complete V32 mathematical source was retained before this continuation.
Its SHA--256 digest was
\begin{center}\scriptsize\ttfamily
5bce84e922a52ec10becbcf77e45ce007ed811babd2a4ab9573da019764b8ff0.
\end{center}
V33 proved the dyadic heat-sum kernel, the fractional-integral law for the
critical record with the Ladyzhenskaya and record--Leray densities, the
running-record bound by the Leray fractional integral, the logarithmic
critical-norm bound under a type-I enstrophy rate with double-exponential
record times, the low-total spread lemma, and the high-total spike
dichotomy.  No full regularity claim is made.

% =============================================================================
\section{V34: scale-resolved record law, Leray concentration at records, and
the type-I death of the imported terminal branch}
\label{sec:v34-entry}
% =============================================================================

V33 proved the fractional-integral record law and its logarithmic
consequence under a type-I enstrophy rate.  This version does three things
with it.  First, the law is resolved scale by scale: under a type-I rate the
dyadic critical masses \(K_nH_{K_n}^\Sigma(b)\) are bounded by one
logarithm at every scale below the parabolic frequency of \(T_*-b\) and
decay geometrically above it, which exhibits the V33 \(\log^2\) as
``\(\log\) scales times \(\log\) mass per scale'' and identifies exactly
which quantity would have to be summable for the critical norm to stay
bounded.  Second, the running form of the law is turned around: at every
record time the half-order Leray fractional integral must reach the record,
so the Leray density must concentrate within a lag of order
\(\mathcal R_n^{-2}\) before the record, unconditionally.  Third, the
double-exponential record times of V33 are tested against the branch
imported in V1 from the earlier terminal programme.  Under a type-I rate the
record-\(n\) corridor is shorter than \(9\delta e^{-c2^n}\), whereas the
imported raw pulse \eqref{eq:imported-raw-pulse} needs a corridor of length
at least a fixed multiple of \(\mathcal R_n^{-6}\).  Hence the imported
terminal branch is finite under a type-I rate: whatever survives of it is
type-II.  No global regularity claim is made.  The compulsory SM/YM
companion audit is due in V35 and is not performed here.

% =============================================================================
\section{The scale-resolved record law}
\label{sec:v34-scale-resolved}
% =============================================================================

\begin{proposition}[Per-scale law]
\label{prop:v34-per-scale-law}
Let \(\varphi\) be admissible on \([a,b]\) in the sense of
\eqref{eq:v33-admissible-density}.  Then for every \(\rho>0\),
\begin{equation}
 \rho H_\rho^\Sigma(b)
 \le\rho e^{-2\nu\rho^2(b-a)}H_\rho^\Sigma(a)
 +\rho\int_a^be^{-2\nu\rho^2(b-t)}\varphi(t)\,\dd t .
 \label{eq:v34-per-scale-law}
\end{equation}
Summed over the dyadic scales \(\rho=K_n\), \(n\in\mathbb Z\), this is
\eqref{eq:v33-record-law}; restricted to one scale it is the V32 tail bound
\eqref{eq:v32-recent-ledger-tail} with the heat weight retained.
\end{proposition}

\begin{proof}
This is the first display in the proof of
Theorem~\ref{thm:v33-record-law}, multiplied by \(\rho\).
\end{proof}

For \(T_*-\delta\le b<T_*\) write
\begin{equation}
 K_{\rm par}(b):=\bigl(\nu(T_*-b)\bigr)^{-1/2},
 \qquad
 M_\delta(b):=\bigl(2A_\delta+2B^2L_\delta(b)^2\bigr)^{1/2},
 \label{eq:v34-parabolic-frequency}
\end{equation}
so that \(\|A^{1/4}u(s)\|_2\le M_\delta(b)\) for all
\(s\in[T_*-\delta,b]\) by Theorem~\ref{thm:v33-type-I-log}, and put
\begin{equation}
 P_\delta(\rho):=\rho\,e^{-\nu\rho^2\delta}E_*
 \Bigl(1+\frac{2C_4R_\delta}{\nu}\Bigr).
 \label{eq:v34-early-term}
\end{equation}

\begin{theorem}[Flat critical-mass profile under a type-I enstrophy rate]
\label{thm:v34-flat-profile}
Assume \eqref{eq:v32-type-I} on \([T_*-\delta,T_*)\) with
\(\delta\le T_*\).  Then for every \(b\in[T_*-\delta/2,T_*)\) and every
\(\rho>0\),
\begin{align}
 \rho H_\rho^\Sigma(b)
 &\le P_\delta(\rho)+8.3\,C_4C_I^2\nu\,M_\delta(b)
 \qquad(\text{all }\rho),
 \label{eq:v34-flat-profile-low}\\
 \rho H_\rho^\Sigma(b)
 &\le P_\delta(\rho)+2\,C_4C_I^2\nu\,M_\delta(b)\,
   \frac{K_{\rm par}(b)}{\rho}
 \qquad(\rho\ge K_{\rm par}(b)).
 \label{eq:v34-flat-profile-high}
\end{align}
Consequently the dyadic critical masses are at most one logarithm at every
scale between \((\nu\delta)^{-1/2}\) and \(K_{\rm par}(b)\), and decay
geometrically above \(K_{\rm par}(b)\); their sum over the
\(\tfrac12\log_2(\delta/(T_*-b))\) scales in between reproduces the
\(\log^2\) of \eqref{eq:v33-type-I-log-bound}.
\end{theorem}

\begin{proof}
Apply \eqref{eq:v34-per-scale-law} with \(a=0\) and
\(\varphi=\varphi_{\rm rec}\).  The initial term is at most
\(\rho e^{-2\nu\rho^2b}E_*\le\rho e^{-\nu\rho^2\delta}E_*\), since
\(b\ge T_*-\delta/2\ge\delta/2\).  The part of the integral over
\([0,T_*-\delta]\) is at most
\(\rho e^{-\nu\rho^2\delta}\cdot4C_4R_\delta E_*/(2\nu)\), because
\(b-t\ge\delta/2\) there and the Leray identity bounds
\(\int_0^{T_*-\delta}\|\Lambda u\|_2^2\le E_*/(2\nu)\).  These two pieces
are \(P_\delta(\rho)\).

On \([T_*-\delta,b]\) insert \(\|A^{1/4}u\|_2\le M_\delta(b)\) and
\eqref{eq:v32-type-I}; it remains to bound
\(J_\rho(b):=\rho\int_{T_*-\delta}^be^{-2\nu\rho^2(b-t)}(T_*-t)^{-1/2}\dd t\).
If \(\rho\ge K_{\rm par}(b)\), then \((T_*-t)^{-1/2}\le(T_*-b)^{-1/2}\)
and \(\int_{-\infty}^be^{-2\nu\rho^2(b-t)}\dd t=(2\nu\rho^2)^{-1}\), so
\(J_\rho\le(T_*-b)^{-1/2}/(2\nu\rho)=K_{\rm par}(b)/(2\nu^{1/2}\rho)\).
If \(\rho<K_{\rm par}(b)\), put \(s_\rho:=(\nu\rho^2)^{-1}>T_*-b\) and
split at lag \(s_\rho\).  The recent part is
\(\rho\int_{b-s_\rho}^b(T_*-t)^{-1/2}\dd t
=2\rho[(T_*-b+s_\rho)^{1/2}-(T_*-b)^{1/2}]\le2\rho s_\rho^{1/2}
=2\nu^{-1/2}\); the remote part is at most
\(\rho\int_{s_\rho}^\infty e^{-2\nu\rho^2r}r^{-1/2}\dd r
\le\rho s_\rho^{-1/2}e^{-2}/(2\nu\rho^2)=e^{-2}\nu^{-1/2}/2\).
Hence \(J_\rho\le2.07\,\nu^{-1/2}\).  Multiplying by
\(4C_4C_I^2\nu^{3/2}M_\delta(b)\) gives the two displays.  The counting
statement is \(\log_2\bigl(K_{\rm par}(b)(\nu\delta)^{1/2}\bigr)
=\tfrac12\log_2(\delta/(T_*-b))\).
\end{proof}

\begin{proposition}[Exact log-free criterion]
\label{prop:v34-log-free-criterion}
Assume \eqref{eq:v32-type-I}.  The critical norm stays bounded on
\([T_*-\delta,T_*)\) if and only if
\begin{equation}
 \sup_{T_*-\delta\le b<T_*}\
 \sum_{n:\ (\nu\delta)^{-1/2}\le K_n\le K_{\rm par}(b)}
 K_nH_{K_n}^\Sigma(b)<\infty .
 \label{eq:v34-log-free-criterion}
\end{equation}
Moreover, for \(K_n\le K_{\rm par}(b)\) the dyadic mass at time \(b\) is
paid by the flux through the cutoff \(K_n\) during the lag window of its
own parabolic time:
\begin{equation}
 K_nH_{K_n}^\Sigma(b)
 \le P_\delta(K_n)+\frac{e^{-2}}{2\nu^{1/2}}\,
 \sup_{t\le b-s_n}\varphi_{\rm rec}(t)\,s_n^{1/2}
 +2K_n\int_{b-s_n}^b|\Phi_{K_n}(t)|\,\dd t,
 \qquad s_n=(\nu K_n^2)^{-1}.
 \label{eq:v34-own-window}
\end{equation}
\end{proposition}

\begin{proof}
Necessity is the lower bound in Lemma~\ref{lem:v33-dyadic-domination}.
For sufficiency, let \(S\) be the supremum in
\eqref{eq:v34-log-free-criterion} and let
\(F(b):=\sup_{T_*-\delta\le s\le b}\|A^{1/4}u(s)\|_2^2\).  The proof of
Theorem~\ref{thm:v34-flat-profile} used the type-I rate and the bound
\(\|A^{1/4}u\|_2\le M_\delta(b)\) on \([T_*-\delta,b]\) only; replacing
the latter by \(F(b)^{1/2}\) gives, for every \(s\le b\) and every scale
\(K_n\ge K_{\rm par}(s)\), \(K_nH_{K_n}^\Sigma(s)\le P_\delta(K_n)
+2C_4C_I^2\nu F(b)^{1/2}K_{\rm par}(s)/K_n\).  By the upper bound in
Lemma~\ref{lem:v33-dyadic-domination}, \(\|A^{1/4}u(s)\|_2^2\) is at most
the sum of three groups of dyadic masses: scales below
\((\nu\delta)^{-1/2}\), at most \(2(\nu\delta)^{-1/2}E_*\); the middle
scales, at most \(S\); and scales above \(K_{\rm par}(s)\), at most
\(\sum_nP_\delta(K_n)+4C_4C_I^2\nu F(b)^{1/2}\) by the geometric sum.
Taking the supremum over \(s\le b\) gives \(F(b)\le C_0+S+C_1F(b)^{1/2}\)
with \(C_0,C_1\) independent of \(b\), hence \(F\) is bounded.  The own-window formula is
\eqref{eq:v34-per-scale-law} with the integral split at lag \(s_n\), the
remote part bounded as in the proof of Theorem~\ref{thm:v34-flat-profile},
and on the recent part the admissible density replaced by the actual flux
\(2|\Phi_{K_n}|\) with \(e^{-2\nu K_n^2(b-t)}\le1\).
\end{proof}

\begin{assessment}[What the logarithm is made of]
\label{ass:v34-log-anatomy}
Each dyadic scale below the parabolic frequency of \(T_*-b\) has critical
mass at most \(O(M_\delta)\), and there are
\(\tfrac12\log_2(\delta/(T_*-b))\) such scales.  Equation
\eqref{eq:v34-own-window} says that the mass at scale \(K_n\) is created,
up to a geometrically decaying memory, by the flux through \(K_n\) in the
lag window \([b-s_n,b]\), that is, during the times at which \(K_n\) was
itself the parabolic frequency.  A bounded critical norm under type-I
therefore requires the own-window fluxes
\(K_n\int_{b-s_n}^b|\Phi_{K_n}|\) to be summable over the scales, whereas
the record--Leray density only bounds each by \(O(M_\delta)\).  This is the
exact remaining gap in the type-I regime.  It is a statement about the
actual signed flux at the parabolic scale, which is the object the V18 triad
slots and the V32 two-stress anatomy describe; it is not a norm estimate.
\end{assessment}

% =============================================================================
\section{Leray concentration at record times}
\label{sec:v34-leray-concentration}
% =============================================================================

\begin{theorem}[The Leray density concentrates before every record]
\label{thm:v34-leray-concentration}
Let \(a<b<T_*\) and suppose \(b\) is a running record time on
\([a,b]\), that is, \(\|A^{1/4}u(b)\|_2^2=F(b)=\sup_{[a,b]}\|A^{1/4}u\|_2^2
=:\mathcal R^2\), with \(\mathcal R^2\ge4\|A^{1/4}u(a)\|_2^2\).  Then
the half-order Leray fractional integral satisfies
\begin{equation}
 \boxed{\quad
 \mathcal I(b)=\int_a^b\frac{\|\Lambda u(t)\|_2^2}{(b-t)^{1/2}}\,\dd t
 \ge\frac{\nu^{1/2}\mathcal R}{18.4\,C_4},
 \quad}
 \label{eq:v34-record-forces-concentration}
\end{equation}
and for every lag \(0<s\le b-a\),
\begin{equation}
 \boxed{\quad
 \int_{b-s}^b\frac{\|\Lambda u(t)\|_2^2}{(b-t)^{1/2}}\,\dd t
 \ge\frac{\nu^{1/2}\mathcal R}{18.4\,C_4}-\frac{E_*}{2\nu s^{1/2}} .
 \quad}
 \label{eq:v34-lag-concentration}
\end{equation}
In particular, with \(s_{\mathcal R}:=(18.4\,C_4E_*)^2/(\nu^3\mathcal R^2)\),
the weighted Leray share within lag \(s_{\mathcal R}\) of the record is at
least \(\nu^{1/2}\mathcal R/(36.8\,C_4)\).
\end{theorem}

\begin{proof}
By \eqref{eq:v33-record-law-rec} at \(s=b\), with
\(\|A^{1/4}u(t)\|_2\le\mathcal R\) on \([a,b]\),
\(\mathcal R^2\le2\|A^{1/4}u(a)\|_2^2+9.2C_4\nu^{-1/2}\mathcal R\,
\mathcal I(b)\le\tfrac12\mathcal R^2+9.2C_4\nu^{-1/2}\mathcal R\,
\mathcal I(b)\), which is \eqref{eq:v34-record-forces-concentration}.
For \eqref{eq:v34-lag-concentration}, the part of \(\mathcal I(b)\) with
lag greater than \(s\) is at most \(s^{-1/2}\int_a^b\|\Lambda u\|_2^2
\le s^{-1/2}E_*/(2\nu)\).  The last claim is the choice of \(s\) making
the subtracted term half the main term.
\end{proof}

\begin{assessment}[Record scale versus concentration lag]
\label{ass:v34-concentration-scales}
At the V1 record \(\mathcal R_n\) with cutoff \(K_n=\mathcal R_n^2/E_*\),
the parabolic time of the record cutoff is
\((\nu K_n^2)^{-1}=E_*^2/(\nu\mathcal R_n^4)\), while the concentration lag
of Theorem~\ref{thm:v34-leray-concentration} is
\(s_{\mathcal R_n}\asymp E_*^2/(\nu^3\mathcal R_n^2)\), larger by the factor
\(\mathcal R_n^2/\nu^2\).  The Leray density must therefore concentrate
within a window that is long compared with the record's own parabolic time
but short compared with any fixed time, and its weighted share there must
grow linearly with the record.  This is consistent with the type-I regime,
where \(\mathcal I(b)\) is logarithmic and \(\mathcal R\) is logarithmic,
and with any type-II regime.  It is a necessary condition for record
growth in the Leray ledger itself, with no cutoff and no source action; it
does not yet contradict the Leray budget because the
\((b-t)^{-1/2}\) weight is unbounded at zero lag.
\end{assessment}

\begin{remark}[The record density dominates the Ladyzhenskaya density]
\label{rem:v34-density-comparison}
The lattice Cauchy--Schwarz inequality
\(\|A^{1/4}u\|_2^2\le E_*^{1/2}\|\Lambda u\|_2\) gives
\(\|\Lambda u\|_2^2\ge\|A^{1/4}u\|_2^4/E_*\), hence
\(\varphi_{\rm rec}=4C_4\|A^{1/4}u\|_2\|\Lambda u\|_2^2
\le4C_4E_*^{1/4}\|\Lambda u\|_2^{5/2}=2\varphi_{5/2}\) always.  The
\(5/2\)-density is therefore never sharper than twice the record density;
its use in V32 was to obtain statements free of the record, not to obtain
smaller numbers.
\end{remark}

% =============================================================================
\section{Type-I death of the imported terminal branch}
\label{sec:v34-imported-branch}
% =============================================================================

Recall the V1 imports: on a record-\(n\) card \(I_n=[a_n,b_n]\) of length
\(L_n\) the corridor bound \eqref{eq:record-corridor} holds, the head
enstrophy is capped by \eqref{eq:head-cap}, and the imported raw pulse
\eqref{eq:imported-raw-pulse} asserts
\(\int_{J_n}\mathcal E_n^2\ge c_{\rm raw}a_-\nu\mathcal R_n^2\) on some
\(J_n\subset I_n\).  We make explicit the one property of a ``record-\(n\)
corridor'' that is used:
\begin{equation}
 a_n\ \ge\ t_n^{\rm hit}
 :=\inf\{t:\ \|A^{1/4}u(t)\|_2\ge\mathcal R_n\}.
 \tag{C\(_n\)}\label{eq:v34-corridor-after-hit}
\end{equation}

\begin{theorem}[Under a type-I enstrophy rate the imported raw pulse occurs
for finitely many records]
\label{thm:v34-raw-pulse-finite}
Assume \eqref{eq:v32-type-I} and \eqref{eq:v34-corridor-after-hit}.  Then
there is \(n_1\) such that for every \(n\ge n_1\) no card in the
record-\(n\) corridor satisfies \eqref{eq:imported-raw-pulse}.  More
precisely, for \(n\ge n_0\),
\begin{equation}
 L_n\le T_*-t_n^{\rm hit}\le9\delta\exp(-c\,2^n),
 \qquad
 \int_{J_n}\mathcal E_n^2\,\dd t
 \le\frac{16\mathcal R_n^8}{E_*^2}\,9\delta\exp(-c\,2^n),
 \label{eq:v34-raw-pulse-upper}
\end{equation}
and the right side is below \(c_{\rm raw}a_-\nu\mathcal R_n^2\) as soon as
\(\mathcal R_n^6\exp(-c2^n)<c_{\rm raw}a_-\nu E_*^2/(144\delta)\).
\end{theorem}

\begin{proof}
By \eqref{eq:v33-record-times}, \(T_*-t_n^{\rm hit}\le9\delta e^{-c2^n}\)
for \(n\ge n_0\); with \eqref{eq:v34-corridor-after-hit},
\(L_n\le b_n-a_n\le T_*-t_n^{\rm hit}\).  On the card,
\eqref{eq:head-cap} gives \(\mathcal E_n\le4\mathcal R_n^4/E_*\), so
\(\int_{J_n}\mathcal E_n^2\le L_n(4\mathcal R_n^4/E_*)^2\).  Since
\(\mathcal R_n=2^n\mathcal R_*\) grows geometrically while
\(e^{-c2^n}\) decays double-exponentially, the stated inequality holds for
all large \(n\).
\end{proof}

\begin{corollary}[The surviving imported branch is type-II]
\label{cor:v34-imported-branch-type-II}
If the imported raw pulse \eqref{eq:imported-raw-pulse} occurs for
infinitely many records, in corridors satisfying
\eqref{eq:v34-corridor-after-hit}, then the type-I enstrophy rate
\eqref{eq:v32-type-I} fails on every \([T_*-\delta,T_*)\): the solution has
a type-II enstrophy rate, \(\limsup_{t\to T_*}(T_*-t)^{1/4}\|\Lambda u(t)\|_2
=\infty\).
\end{corollary}

\begin{proof}
Contrapositive of Theorem~\ref{thm:v34-raw-pulse-finite}.
\end{proof}

\begin{firewall}[What has been killed and what has not]
\label{fw:v34-imported-branch-scope}
Theorem~\ref{thm:v34-raw-pulse-finite} does not exclude a singularity.  It
shows that the branch on which the whole V1--V31 development was
conditioned, the imported raw pulse at cofinally many doubled records, is
incompatible with a type-I enstrophy rate, purely because the corridor is
then too short to hold the pulse.  The V32 firewall
\ref{fw:v32-type-I-scope} already placed the high-total branches in the
type-II regime; the present theorem places the imported raw branch there
too.  The type-I regime is thus left with: uniformly small tails, at most
logarithmic critical growth, no high-total exits, and no imported raw
pulses.  What remains open in that regime is the own-window flux criterion
of Proposition~\ref{prop:v34-log-free-criterion}.
\end{firewall}

% =============================================================================
\section{Tests}
\label{sec:v34-tests}
% =============================================================================

\begin{proposition}[Consistency checks]
\label{prop:v34-tests}
\begin{enumerate}[label=\textup{(\roman*)},leftmargin=3.7em]
\item Under the exact Leray-rate profile
      \(\|\Lambda u\|_2^2=c\nu^{3/2}(T_*-t)^{-1/2}\), the concentration
      bound \eqref{eq:v34-lag-concentration} at lag \(s=T_*-b\) reads
      \(\mathcal R\lesssim\nu^{-1/2}\bigl(\log\tfrac{\delta}{T_*-b}+1\bigr)\),
      in agreement with Theorem~\ref{thm:v33-type-I-log}; the record law
      and the concentration theorem are mutually consistent on the type-I
      profile.
\item For the V14/V18 packet on \(J_j\), \eqref{eq:v34-raw-pulse-upper}
      does not apply, because the packet's record is attained by its
      initial datum and the packet is not a singular trajectory; the
      packet interval satisfies \(|J_j|\asymp R_j^{-1}N_j^{-2}\), which is
      polynomially, not double-exponentially, small in \(R_j\).
\end{enumerate}
\end{proposition}

\begin{proof}
(i) \(\int_{b-(T_*-b)}^b(T_*-t)^{-1/2}(b-t)^{-1/2}\dd t
=2\operatorname{arcsinh}(1)\) is bounded, and the full
\(\mathcal I(b)\) is the logarithm of Lemma~\ref{lem:v33-log-kernel}.
(ii) is \eqref{eq:v14-two-parameter-interval}.
\end{proof}

% =============================================================================
\section{Integrated V34 conclusion and next acquisition order}
\label{sec:v34-conclusion}
% =============================================================================

\begin{theorem}[What V34 proves]
\label{thm:v34-conclusion}
Relative to V1--V33, the following are proved.
\begin{enumerate}[label=\textup{(V34.\arabic*)},leftmargin=6.2em]
\item The per-scale record law \eqref{eq:v34-per-scale-law}.
\item Under a type-I enstrophy rate the dyadic critical masses obey the
      flat profile
      \eqref{eq:v34-flat-profile-low}--\eqref{eq:v34-flat-profile-high}:
      one logarithm per scale below the parabolic frequency, geometric decay
      above it.
\item The exact log-free criterion \eqref{eq:v34-log-free-criterion} and
      the own-window formula \eqref{eq:v34-own-window}: bounded critical
      norm under type-I is equivalent to summability of the own-window
      fluxes over the scales.
\item Unconditionally, every running record forces the half-order Leray
      fractional integral to reach \(\nu^{1/2}\mathcal R/(18.4C_4)\), with
      the concentration lag \(s_{\mathcal R}\asymp E_*^2\nu^{-3}
      \mathcal R^{-2}\), \eqref{eq:v34-record-forces-concentration}--%
      \eqref{eq:v34-lag-concentration}.
\item The record density dominates the Ladyzhenskaya density,
      Remark~\ref{rem:v34-density-comparison}.
\item Under a type-I enstrophy rate and \eqref{eq:v34-corridor-after-hit},
      the imported raw pulse \eqref{eq:imported-raw-pulse} occurs for
      finitely many records; cofinal raw pulses force a type-II enstrophy
      rate.
\end{enumerate}
The full Navier--Stokes stability/global-regularity theorem remains
unproved.  The type-I regime is reduced to the own-window flux criterion;
the type-II regime carries every surviving branch of the earlier programme.
\end{theorem}

\begin{proof}
Items \textup{(V34.1)}--\textup{(V34.3)} are
Proposition~\ref{prop:v34-per-scale-law},
Theorem~\ref{thm:v34-flat-profile}, and
Proposition~\ref{prop:v34-log-free-criterion}.  Items
\textup{(V34.4)}--\textup{(V34.5)} are
Theorem~\ref{thm:v34-leray-concentration} and
Remark~\ref{rem:v34-density-comparison}.  Item \textup{(V34.6)} is
Theorem~\ref{thm:v34-raw-pulse-finite} and
Corollary~\ref{cor:v34-imported-branch-type-II}.
\end{proof}

\begin{programme}[V35 acquisition order, with compulsory companion audit]
\label{prog:v35}
\begin{enumerate}[label=\textup{(AC\arabic*)},leftmargin=4.8em]
\item Perform the compulsory fifth-version review of the current SM and YM
      notes in \texttt{latex\_notes}: freeze them by digest, and decide
      whether anything in them bears on the own-window flux criterion or
      on the type-II ledger; import nothing that is not proved from the NS
      balances.
\item Own-window flux.  For \(K_n\le K_{\rm par}(b)\), write the
      own-window flux through the V32 two-stress anatomy: the low--low slot
      is confined to the octave \((K_n,2K_n]\) and the high--high slot
      carries \(K_n^{-1/2}\) against the tail dissipation.  Determine
      whether, under type-I, the tail dissipation above \(K_n\) during the
      lag window \([b-s_n,b]\) is Leray-summable over the scales; if so the
      high--high slot is log-free and only the low--low octave slot remains.
\item Type-II regime.  Combine Theorem~\ref{thm:v34-leray-concentration}
      with Corollary~\ref{cor:v32-high-total-divergence}: at a high-total
      record the Leray density concentrates within lag
      \(s_{\mathcal R}\) and the \(L^{5/2}\) ledger diverges.  Determine
      whether the two concentrations can be placed on disjoint windows for
      cofinally many records without exceeding the Leray budget, using the
      spike dichotomy \eqref{eq:v33-spike-dichotomy}.
\item Keep the V29 persistence sum one-use; keep the hinge off the total
      branch; do not re-import the raw pulse in the type-I regime.
\end{enumerate}
\end{programme}

\begin{assessment}[Position after V34]
\label{ass:v34-position}
The type-I regime is now sharply characterized: uniformly small tails, a
flat logarithmic critical-mass profile below the parabolic frequency, no
high-total exits, no imported raw pulses, and one explicit remaining gap,
summability of the own-window fluxes.  The type-II regime inherits every
earlier branch and must additionally exhibit both a divergent
\(L^{5/2}\) ledger and Leray concentration before each record.  No full
stability or global-regularity claim is made.
\end{assessment}

% =============================================================================
\section*{Version V34 append-only record}
\addcontentsline{toc}{section}{Version V34 append-only record}
% =============================================================================

The complete V33 mathematical source was retained before this continuation.
Its SHA--256 digest was
\begin{center}\scriptsize\ttfamily
5c855f87f849839562e0c9937a5f3ed7f719a7c854e41e9b9288c96bf2723656.
\end{center}
V34 proved the per-scale record law, the flat critical-mass profile under a
type-I enstrophy rate with its exact log-free criterion and own-window
formula, the unconditional Leray concentration theorem at record times, the
domination of the Ladyzhenskaya density by the record density, and the
finiteness of the imported raw-pulse branch under a type-I rate.  No full
regularity claim is made.

% =============================================================================
\section{V35: companion audit, the slot no-go for the type-I logarithm, and
the record--share inequality for created tail energy}
\label{sec:v35-entry}
% =============================================================================

This is a fifth-version continuation, so the compulsory review of the two
companion programmes is performed first.  The mathematical content then
settles two items of Programme~\ref{prog:v35}.  For the own-window flux
criterion of V34 it is proved that neither stress slot of the V32 anatomy,
bounded by norms, contributes less than a constant per dyadic scale under a
type-I enstrophy rate; the logarithm is therefore not removable by any
slot-wise norm bound and requires the signed own-window flux.  For the
type-II regime the record density turns the V32 tail bound into a
record--share inequality: the tail energy created at any cutoff during any
window is at most \(4C_4\) times the window record times the window Leray
share.  On a cofinal high-total sequence this forces the reciprocals of the
in-window records to be summable.  No global regularity claim is made.

% =============================================================================
\section{Compulsory V35 review of the current companion notes}
\label{sec:v35-companion-audit}
% =============================================================================

The two current companion files in \texttt{latex\_notes} were inspected
read-only.  The frozen checkpoints are
\begin{align}
 &\texttt{Renewal\_SM\_Einstein\_Continuum\_Research\_Notes\_V51.tex},
 \notag\\[-2pt]
 &\qquad\text{SHA--256 }
 \texttt{1ed8350f3d5a05582a7e60baf70f77fc6be3fb529b568c7d02098cde421d9101},
 \label{eq:v35-SM-checkpoint}\\
 &\texttt{Renewal\_Yang\_Mills\_Mass\_Gap\_Research\_Notes\_V37.tex},
 \notag\\[-2pt]
 &\qquad\text{SHA--256 }
 \texttt{9d5fcb54b3593d914777ef97e444d4173984b109e5075cccf3702c1af361a0e2}.
 \label{eq:v35-YM-checkpoint}
\end{align}
Both companion programmes have restarted their numbering since the V30
audit; the present SM file is the sixth series and the YM file the fifth.
Both audited NS V31 (digest \texttt{f1121b62\ldots}) and imported nothing
from it beyond proof discipline.  The review concentrated on SM V46--V51
and YM V33--V37.

The relevant proved SM conclusions are as follows.  V46--V47 integrate an
abelian plaquette star exactly and close an independent-star capacity by a
latent-link quotient.  V48--V49 disintegrate abelian cochains, restrict to
the Bianchi cycle space, and prove that the compact all-flux cube law has no
order-\(\beta\) Poincar\'e gap because its nonzero lift sectors contain
strict metastable wells.  V50 selects a noncompact hypercharge branch and
proves a finite-cutoff theorem and a free continuum endpoint on it; its
audit assessment insists that every potential-coordinate estimate keep the
least positive singular value of the lattice curl, and that
volume-independence be claimed only in curvature coordinates where the
inverse derivative cancels algebraically.  V51 derives the exact score and
Hessian of the reduced chiral determinant, proves that raw uniform bounds
lose one and two inverse powers of the least nonzero singular value and that
this loss is sharp, and computes that Standard Model hypercharge anomaly
cancellation leaves a square sum \(10/3\), so anomaly cancellation does not
supply the needed negative-Hessian bound; the remaining input is named DSSC.
None of these are Navier--Stokes quantities.

The relevant proved YM conclusions are as follows.  V33--V35 build a
partition-local Casimir square, an exact heat spectral tangent score, a
one-root quantum forest with its root payment, and a resummed root
certificate; V35 proves that coefficientwise infinite jets do not close and
that the complete jet must be summed before norms are taken.  V36 proves a
bounded marked-moment row for every tensor block, the exact binary
heat-normalized activity, exact binary WPRAC, and identifies a
superexponential loss \(e^{B_A^2}\) in a separately normalized high-block
heat reference.  V37 writes the spectral score as a differentiated heat
kernel, proves uniform three-block multiplier expansions, and closes ternary
WPRAC.  All-arity WPRAC and everything downstream of it remain open in
that problem.

\begin{theorem}[V35 companion-audit decision]
\label{thm:v35-companion-decision}
The audit yields exactly the following transfer decisions.
\begin{enumerate}[label=\textup{(\roman*)},leftmargin=3.7em]
\item No SM V46--V51 capacity, quotient, determinant, or anomaly constant
      is imported.  The SM discipline of keeping the inverse singular value
      visible is already the form of the V33 record law, whose kernel
      \((\nu(b-t))^{-1/2}\) is the inverse half-power that the dyadic
      heat sum produces; V33 did not hide it in a ``heat gain.''
\item No YM V33--V37 Casimir, activity, or WPRAC estimate is imported.
      The YM discipline that the complete jet be summed before norms are
      taken is the form of the V33 proof, which kept every heat weight
      inside the dyadic sum; the YM identification of a superexponential
      loss in a separately normalized reference is the exact analogue of
      the V32 persistent term \(E_*(\nu\tau)^{-1/2}\) removed in V33.
\item Neither companion supplies the own-window signed flux of
      Proposition~\ref{prop:v34-log-free-criterion} or any estimate of the
      \(L^{5/2}_tH^1\) ledger; the type-I/type-II split of V32--V34 is
      untouched by the audit.
\end{enumerate}
\end{theorem}

\begin{proof}
Items (i) and (ii) are the explicit claim boundaries in the SM V50 audit
assessment and V51 status theorem and in the YM V36--V37 claim-boundary
firewalls; their hypotheses are cochain complexes, determinant charts,
thermal blocks, and Casimir marks, none of which is a Navier--Stokes tail
balance.  Item (iii) follows because the only NS-relevant statements in
either file are the digests and one-paragraph summaries of NS V31.
\end{proof}

\begin{firewall}[Companion analogies remain analogies]
\label{fw:v35-companion-types}
The two disciplines cited in Theorem~\ref{thm:v35-companion-decision} were
already followed in V33; the audit confirms them and adds no inequality.
Every quantitative statement below is proved from the NS balances and the
Leray identity.
\end{firewall}

% =============================================================================
\section{Slot-wise norm bounds cannot remove the type-I logarithm}
\label{sec:v35-slot-no-go}
% =============================================================================

Fix a type-I enstrophy rate \eqref{eq:v32-type-I}, a time
\(b\in[T_*-\delta/2,T_*)\), a dyadic scale \(K_n\le K_{\rm par}(b)\), its
own window \(w_n:=[b-s_n,b]\), \(s_n=(\nu K_n^2)^{-1}\), and the bound
\(\|A^{1/4}u\|_2\le M_\delta(b)\) on \([T_*-\delta,b]\) from
Theorem~\ref{thm:v33-type-I-log}.

\begin{proposition}[Each stress slot contributes at most a constant per
scale]
\label{prop:v35-slot-no-go}
Under the above hypotheses,
\begin{align}
 2K_n\int_{w_n}|\Phi_{K_n}^{\rm HH}(t)|\,\dd t
 &\le8\,C_4C_I^3\nu^2 ,
 \label{eq:v35-HH-own-window}\\
 2K_n\int_{w_n}|\Phi_{K_n}^{\rm LL}(t)|\,\dd t
 &\le4\,C_4C_I^2\nu\,M_\delta(b),
 \label{eq:v35-LL-own-window}
\end{align}
and, in octave-resolved form,
\begin{equation}
 2K_n\int_{w_n}|\Phi_{K_n}^{\rm LL}|\,\dd t
 \le2C_4M_\delta(b)\,K_n
 \Bigl(\int_{w_n}\|\Lambda u\|_2^2\Bigr)^{1/2}
 \Bigl(\int_{w_n}\|\Lambda\Pi_{(K_n,2K_n]}u\|_2^2\Bigr)^{1/2}.
 \label{eq:v35-LL-octave-form}
\end{equation}
Neither right side in
\eqref{eq:v35-HH-own-window}--\eqref{eq:v35-LL-own-window} decreases
with \(n\).  Consequently, over the
\(\tfrac12\log_2(\delta/(T_*-b))\) scales between \((\nu\delta)^{-1/2}\)
and \(K_{\rm par}(b)\), the slot-wise own-window bounds sum to at least a
logarithm times \(M_\delta(b)\), and the bootstrap of
Theorem~\ref{thm:v33-type-I-log} returns \(\log^2\); no improvement of
the exponent is available from
\eqref{eq:v32-LL-slot-bound}--\eqref{eq:v32-HH-slot-bound} taken with
absolute values.
\end{proposition}

\begin{proof}
For the high--high slot use \eqref{eq:v32-HH-slot-bound} in the form
\(|\Phi^{\rm HH}_{K_n}|\le C_4K_n^{-1/2}V_{K_n}^\Sigma\|\Lambda P_{K_n}u\|_2
\le C_4K_n^{-1/2}\|\Lambda u\|_2^3\) and the type-I rate:
\[
 2K_n\int_{w_n}|\Phi_{K_n}^{\rm HH}|
 \le2C_4K_n^{1/2}C_I^3\nu^{9/4}\int_{b-s_n}^b(T_*-t)^{-3/4}\,\dd t
 \le2C_4K_n^{1/2}C_I^3\nu^{9/4}\cdot4s_n^{1/4}
 =8C_4C_I^3\nu^2,
\]
since \(\int_{b-s}^b(T_*-t)^{-3/4}\dd t\le4s^{1/4}\) by subadditivity of
\(s\mapsto s^{1/4}\), and \(s_n^{1/4}K_n^{1/2}=\nu^{-1/4}\).  For the
low--low slot use \eqref{eq:v32-LL-slot-bound} in the form
\(|\Phi^{\rm LL}_{K_n}|\le\|P_{K_n}u\|_4^2\|\Lambda\Pi_{(K_n,2K_n]}u\|_2
\le C_4\|A^{1/4}u\|_2\|\Lambda u\|_2\|\Lambda\Pi_{(K_n,2K_n]}u\|_2\),
by the interpolation of Proposition~\ref{prop:v32-corridor-form}; then
\(\|\Lambda\Pi_{(K_n,2K_n]}u\|_2\le\|\Lambda u\|_2\), the bound
\(M_\delta(b)\), the type-I rate, and
\(\int_{w_n}(T_*-t)^{-1/2}\dd t\le2s_n^{1/2}=2\nu^{-1/2}K_n^{-1}\) give
\eqref{eq:v35-LL-own-window}.  The octave-resolved form is
Cauchy--Schwarz in time before the last two substitutions.  The counting
statement is that of Theorem~\ref{thm:v34-flat-profile}.
\end{proof}

\begin{proposition}[Octave orthogonality does not help either]
\label{prop:v35-octave-no-go}
Summing the octave-resolved low--low bound over the scales and applying
Cauchy--Schwarz in \(n\),
\begin{equation}
 \sum_{n:\ K_n\le K_{\rm par}(b)}
 2K_n\int_{w_n}|\Phi_{K_n}^{\rm LL}|
 \le2C_4M_\delta(b)\,
 \Bigl(\sum_n K_n\int_{w_n}\|\Lambda u\|_2^2\Bigr)^{1/2}
 \Bigl(\sum_n K_n\int_{w_n}\|\Lambda\Pi_{(K_n,2K_n]}u\|_2^2\Bigr)^{1/2},
 \label{eq:v35-octave-sum}
\end{equation}
and both sums are at most \(2.3\,\nu^{-1/2}\mathcal I(b)\), the Leray
fractional integral of Corollary~\ref{cor:v33-running-record}, which is
logarithmic under type-I by Lemma~\ref{lem:v33-log-kernel}.  Thus
orthogonality of the octaves converts ``\(\log\) scales times \(O(1)\)''
into ``\(\sqrt{\log}\times\sqrt{\log}\)'' and nothing more.
\end{proposition}

\begin{proof}
Exchange sum and integral: \(\sum_nK_n\mathbf 1_{\{s_n\ge b-t\}}
\le\sum_{K_n\le(\nu(b-t))^{-1/2}}K_n\le2(\nu(b-t))^{-1/2}\), and for the
octave sum use \(\sum_n\|\Lambda\Pi_{(K_n,2K_n]}u\|_2^2=\|\Lambda u\|_2^2\)
after bounding each weight by the same \(2(\nu(b-t))^{-1/2}\).
\end{proof}

\begin{assessment}[The logarithm is a signed-flux question]
\label{ass:v35-signed-flux}
Propositions~\ref{prop:v35-slot-no-go}--\ref{prop:v35-octave-no-go} close
item (AC2) of Programme~\ref{prog:v35} negatively: under a type-I rate,
every slot-wise or octave-wise norm bound of the own-window flux is
\(O(M_\delta)\) per scale, exactly the flat profile of
Theorem~\ref{thm:v34-flat-profile}.  Removing the logarithm therefore
requires the signed own-window flux \(\int_{w_n}\Phi_{K_n}\), for which
the only unconditional signed estimate is the V4 net-flux bound
\(|\int\Phi_\rho|\le E_*/2\), too weak by the factor
\(K_nE_*/M_\delta\).  Between the norm bound and the net bound lies the
signed triad structure of V18 and the two-stress anatomy of V32; V36 must
attack it there or accept the logarithm as the type-I endpoint of this
route.
\end{assessment}

% =============================================================================
\section{The record--share inequality for created tail energy}
\label{sec:v35-record-share}
% =============================================================================

For a window \(w=[b-\tau,b]\) write
\(\mathcal R_w:=\sup_{t\in w}\|A^{1/4}u(t)\|_2\) and
\(\epsilon_w:=\int_w\|\Lambda u\|_2^2\,\dd t\).

\begin{theorem}[Created tail energy is at most record times Leray share]
\label{thm:v35-record-share}
For every cutoff \(\rho>0\) and every window \(w=[b-\tau,b]\),
\begin{equation}
 \boxed{\quad
 H_\rho^\Sigma(b)-e^{-2\nu\rho^2\tau}H_\rho^\Sigma(b-\tau)
 \le4C_4\,\mathcal R_w\,\epsilon_w .
 \quad}
 \label{eq:v35-record-share}
\end{equation}
Consequently, if \(H_\rho^\Sigma(b)\ge(1+\kappa)\theta E_*\) and
\(\tau=\theta^{-1}(\nu\rho^2)^{-1}\), then
\begin{equation}
 \boxed{\quad
 \mathcal R_w\,\epsilon_w
 \ge c_\theta E_*,
 \qquad
 c_\theta:=\frac{(1+\kappa)\theta-e^{-2/\theta}}{4C_4}>0 .
 \quad}
 \label{eq:v35-high-total-uncertainty}
\end{equation}
\end{theorem}

\begin{proof}
Apply \eqref{eq:v34-per-scale-law} with \(\varphi=\varphi_{\rm rec}\),
divide by \(\rho\), bound \(\|A^{1/4}u\|_2\le\mathcal R_w\) and the heat
weight by one.  The second display is the case
\(H_\rho^\Sigma(b-\tau)\le E_*\) with \(e^{-2/\theta}<\theta\).
\end{proof}

\begin{corollary}[Cofinal high-total windows have summable reciprocal
records]
\label{cor:v35-summable-reciprocals}
Let \(w_j=[b_j-\tau_j,b_j]\) be pairwise disjoint windows with
\(H_{\rho_j}^\Sigma(b_j)\ge(1+\kappa)\theta E_*\) and
\(\tau_j=\theta^{-1}(\nu\rho_j^2)^{-1}\).  Then
\begin{equation}
 \boxed{\quad
 \sum_j\frac1{\mathcal R_{w_j}}
 \le\frac{1}{2\nu c_\theta}.
 \quad}
 \label{eq:v35-reciprocal-sum}
\end{equation}
In particular at most \(M/(2\nu c_\theta)\) such windows have in-window
record at most \(M\), and along a cofinal high-total sequence
\(\mathcal R_{w_j}\to\infty\) fast enough for
\eqref{eq:v35-reciprocal-sum}.
\end{corollary}

\begin{proof}
\(\mathcal R_{w_j}^{-1}\le\epsilon_{w_j}/(c_\theta E_*)\) by
\eqref{eq:v35-high-total-uncertainty}; sum, and use disjointness with the
Leray identity \(\sum_j\epsilon_{w_j}\le E_*/(2\nu)\).
\end{proof}

\begin{assessment}[Comparison with the V1 corridor and with V32]
\label{ass:v35-record-share-meaning}
On the V1 doubled-record corridor the in-window record is at most
\(2\mathcal R_n\), so \eqref{eq:v35-high-total-uncertainty} requires the
Leray share \(\epsilon_n\ge c_\theta E_*/(2\mathcal R_n)\), the summable
doubled-record debit already found in Proposition~\ref{prop:raw-energy-price}
for the raw pulse; the record--share inequality is therefore the
high-total analogue of that debit, obtained with no source action and no
card structure.  Conversely, V32's Ladyzhenskaya form of the same event
charged the window ledger \(\mathscr L_{5/2}\) an amount independent of
the record.  The two forms are consistent: by
Remark~\ref{rem:v34-density-comparison} the record density is the smaller
one, and \eqref{eq:v35-reciprocal-sum} is the sharpest unconditional
growth condition this series has for a high-total cascade.  It is not a
contradiction, since doubled records satisfy it.
\end{assessment}

\begin{proposition}[Type-II synthesis]
\label{prop:v35-type-II-synthesis}
Let \(u\) have a cofinal high-total sequence in the sense of
Corollary~\ref{cor:v32-high-total-divergence}, with disjoint windows
\(w_j\).  Then simultaneously:
\begin{enumerate}[label=\textup{(\roman*)},leftmargin=3.7em]
\item \(\int_0^{T_*}\|\Lambda u\|_2^{5/2}\,\dd t=\infty\), each window
      carrying ledger at least \(2c_\theta E_*^{3/4}\)
      \textup{(V32)};
\item \(\sum_j\mathcal R_{w_j}^{-1}\le(2\nu c_\theta)^{-1}\)
      \textup{(this version)};
\item at every running record time \(b\) with record \(\mathcal R\), the
      Leray fractional integral obeys
      \(\mathcal I(b)\ge\nu^{1/2}\mathcal R/(18.4C_4)\)
      \textup{(V34)};
\item the enstrophy rate is type-II:
      \(\limsup_{t\to T_*}(T_*-t)^{1/4}\|\Lambda u(t)\|_2=\infty\)
      \textup{(V32)}.
\end{enumerate}
No subset of these conditions is contradictory by arithmetic alone.  The
following scalar spike schedule satisfies all four with a finite Leray
budget: consecutive windows \(w_j\) of length \(\ell_j=L_016^{-j}\), each
containing one spike on which \(\|\Lambda u\|_2^2=A\,16^j\) for a duration
\(\sigma_j=a\,32^{-j}\), with \(\|\Lambda u\|_2^2\) negligible elsewhere,
records \(\mathcal R_{w_j}=2^j\), and tails of order \(E_*\) at
\(\rho_j\asymp4^j\).
\end{proposition}

\begin{proof}
(i)--(iv) are the cited results.  For the schedule: the Leray share is
\(\epsilon_j=Aa\,2^{-j}\), summable; the window ledger is
\(\int_{w_j}\|\Lambda u\|_2^{5/2}=A^{5/4}a\,32^{j}32^{-j}=A^{5/4}a\),
constant in \(j\), so (i) holds and the total ledger diverges; the
record--share product is \(\mathcal R_{w_j}\epsilon_j=Aa\), so (ii) holds
once \(Aa\ge c_\theta E_*\), and \(\sum_j\mathcal R_{w_j}^{-1}<\infty\);
the record is compatible with the spike height by
\(\mathcal R_{w_j}^4=16^j\le E_*A\,16^j\) for \(A\ge E_*^{-1}\); the
fractional integral at the end of the \(j\)-th spike is at least
\(A16^j\cdot2\sigma_j^{1/2}=2Aa^{1/2}(16/\sqrt{32})^j\ge c\,2^j\), so
(iii) holds; finally \(T_*-t\asymp16^{-j}\) on \(w_j\), so
\((T_*-t)^{1/4}\|\Lambda u\|_2\asymp2^{-j}\cdot A^{1/2}4^j\to\infty\),
which is (iv).  The high-total window length
\(\theta^{-1}(\nu\rho_j^2)^{-1}\asymp16^{-j}\) matches \(\ell_j\), and the
spike fits inside it since \(32^{-j}<16^{-j}\).
\end{proof}

\begin{firewall}[The synthesis is a consistency check, not a construction]
\label{fw:v35-synthesis-scope}
The profile in Proposition~\ref{prop:v35-type-II-synthesis} is a scalar
schedule of enstrophy heights and window lengths, not a Navier--Stokes
solution.  It shows that the four necessary conditions accumulated in
V32--V35 do not yet exclude the type-II high-total regime by arithmetic
alone.  Any exclusion must use the signed flux structure or a new ledger.
\end{firewall}

% =============================================================================
\section{Integrated V35 conclusion and next acquisition order}
\label{sec:v35-conclusion}
% =============================================================================

\begin{theorem}[What V35 proves]
\label{thm:v35-conclusion}
Relative to V1--V34, the following are proved.
\begin{enumerate}[label=\textup{(V35.\arabic*)},leftmargin=6.2em]
\item The compulsory companion audit
      (Theorem~\ref{thm:v35-companion-decision}): no SM or YM inequality is
      imported; the resummation and visible-inverse-power disciplines were
      already followed in V33.
\item Under a type-I rate, each stress slot and each octave contributes
      at most a constant times \(M_\delta\) per dyadic scale to the
      own-window flux,
      \eqref{eq:v35-HH-own-window}--\eqref{eq:v35-octave-sum}; the
      logarithm cannot be removed by slot-wise or octave-wise norm bounds.
\item The record--share inequality \eqref{eq:v35-record-share}: created
      tail energy is at most \(4C_4\) times the window record times the
      window Leray share.
\item Cofinal high-total windows have summable reciprocal in-window
      records, \eqref{eq:v35-reciprocal-sum}.
\item The four necessary conditions of the type-II regime are mutually
      consistent with a finite Leray budget
      (Proposition~\ref{prop:v35-type-II-synthesis}).
\end{enumerate}
The full Navier--Stokes stability/global-regularity theorem remains
unproved.
\end{theorem}

\begin{proof}
The items are Theorem~\ref{thm:v35-companion-decision},
Propositions~\ref{prop:v35-slot-no-go}--\ref{prop:v35-octave-no-go},
Theorem~\ref{thm:v35-record-share},
Corollary~\ref{cor:v35-summable-reciprocals}, and
Proposition~\ref{prop:v35-type-II-synthesis}.
\end{proof}

\begin{programme}[V36 acquisition order]
\label{prog:v36}
\begin{enumerate}[label=\textup{(AC\arabic*)},leftmargin=4.8em]
\item Signed own-window flux under type-I.  Write
      \(\int_{w_n}\Phi_{K_n}\,\dd t\) through the V18 triad layer cake and
      the V32 slots, and determine whether the low--low octave slot, which
      is the only slot whose energy lands in \((K_n,2K_n]\), has a signed
      time integral controlled by the octave energy itself (a
      ``self-limiting'' octave inequality).  Test on the balanced plane wave
      of Proposition~\ref{prop:v31-hinge-balanced-nullspace} and on the
      V12 witness.
\item Type-II: replace the scalar schedule of
      Proposition~\ref{prop:v35-type-II-synthesis} by the actual
      requirement that the tail at \(\rho_j\) be formed by the first-octave
      low--low slot from parents below \(\rho_j\), and determine whether
      the parent energy required at \(\rho_j/2\) contradicts the
      record--share inequality one scale down.
\item Keep the V29 persistence sum one-use; keep the hinge off the total
      branch.  The next compulsory companion audit is V40.
\end{enumerate}
\end{programme}

\begin{assessment}[Position after V35]
\label{ass:v35-position}
The type-I regime has one gate, the signed own-window flux, and V35
proved that no norm bound reaches it.  The type-II regime has four
necessary conditions that are jointly satisfiable by a scalar schedule, so
it too needs the signed structure.  Both regimes now point at the same
object: the sign of the first-octave low--low flux over one parabolic
window.  No full stability or global-regularity claim is made.
\end{assessment}

% =============================================================================
\section*{Version V35 append-only record}
\addcontentsline{toc}{section}{Version V35 append-only record}
% =============================================================================

The complete V34 mathematical source was retained before this continuation.
Its SHA--256 digest was
\begin{center}\scriptsize\ttfamily
0f1360cd7fc4d06ebf24bc1bd3dee106e18584af8a1e071a3460dce5d1fd8374.
\end{center}
V35 performed the compulsory companion audit, proved the slot-wise and
octave-wise no-go for removing the type-I logarithm, proved the
record--share inequality for created tail energy with its summable
reciprocal-record consequence, and checked the joint consistency of the
type-II necessary conditions.  No full regularity claim is made.

% =============================================================================
\section{V36: subcritical record laws, the sign of the octave source, and the
logarithmic gap to local regularity}
\label{sec:v36-entry}
% =============================================================================

V35 left both regimes pointing at one object, the signed first-octave
low--low flux over one parabolic window.  This version does three things.
It extends the V33 record law to every subcritical order
\(\dot H^s\), \(1/2\le s<1\), and shows that the kernel becomes
non-integrable exactly at \(s=1\), where no a priori bound exists.  It
decomposes the signed integrated low--low flux exactly: the part generated
by the octave's own low--low source is a nonnegative quadratic form, equal
to the V31 formed octave energy plus its dissipation, so no sign gain is
available from self-interaction; the only indefinite term is the
interaction of the low--low source with the mixed and high--high sources.
It then measures the type-I regime against local regularity: the
scale-invariant local energies at every point and scale are bounded by
powers of the same logarithm, so the gap between the type-I regime and the
small-local-energy hypothesis of partial-regularity theory is exactly
logarithmic.  A two-slot formation alternative for high-total events is
added for the type-II regime.  No global regularity claim is made.  The
next compulsory companion audit is V40.

% =============================================================================
\section{Record laws at every subcritical order}
\label{sec:v36-subcritical-laws}
% =============================================================================

\begin{theorem}[\(\dot H^s\) fractional-integral laws]
\label{thm:v36-Hs-law}
Let \(1/2\le s<1\), let \(\varphi\) be admissible on \([a,b]\) in the
sense of \eqref{eq:v33-admissible-density}, and put
\(C_s:=(1-4^{-s})^{-1}+0.14\cdot4^s\le2.3\).  Then
\begin{equation}
 \boxed{\quad
 \|\Lambda^su(b)\|_2^2
 \le4^s\|\Lambda^su(a)\|_2^2
 +(4^s-1)\,C_s\,\nu^{-s}
 \int_a^b\frac{\varphi(t)}{(b-t)^{s}}\,\dd t .
 \quad}
 \label{eq:v36-Hs-law}
\end{equation}
At \(s=1/2\) this is \eqref{eq:v33-record-law}.  At \(s=1\) the dyadic
heat sum is \(\sum_nK_n^2e^{-2\nu K_n^2r}\asymp(\nu r)^{-1}\), the kernel
is not integrable at \(r=0\), and the same argument gives no bound; this
is consistent with the absence of any a priori \(H^1\) bound, and with
Lemma~\ref{lem:v32-leray-rate}, since
\(\int^b(T_*-t)^{-1/2}(b-t)^{-1}\,\dd t=\infty\).
\end{theorem}

\begin{proof}
The layer cake \(|k|^{2s}=\int_0^{|k|}2s\rho^{2s-1}\dd\rho\) gives
\(\|\Lambda^sv\|_2^2=\int_0^\infty2s\rho^{2s-1}H_\rho[v]\dd\rho\).  Since
\(H_\rho\) is nonincreasing, the annulus \(\mathcal A_n\) contributes at
most \((4^s-1)K_n^{2s}H_{K_n}\), and the interval \([K_n/2,K_n)\) at
least \((1-4^{-s})K_n^{2s}H_{K_n}\); hence
\[
 (1-4^{-s})\sum_{n\in\mathbb Z}K_n^{2s}H_{K_n}[v]
 \le\|\Lambda^sv\|_2^2
 \le(4^s-1)\sum_{n\in\mathbb Z}K_n^{2s}H_{K_n}[v].
\]
Multiply \eqref{eq:v34-per-scale-law} (divided by \(\rho\)) by
\(K_n^{2s}\) and sum.  The initial term is at most
\((4^s-1)\sum_nK_n^{2s}H_{K_n}(a)\le\frac{4^s-1}{1-4^{-s}}
\|\Lambda^su(a)\|_2^2=4^s\|\Lambda^su(a)\|_2^2\).  For the source term,
with \(x_n=K_n(\nu r)^{1/2}\), the terms \(x_n\le1\) of
\(\sum_nx_n^{2s}e^{-2x_n^2}\) sum to at most \((1-4^{-s})^{-1}\), and the
terms \(x_n>1\), starting at \(x_{n_1}\in(1,2]\), to at most
\(\sum_{m\ge0}4^{s(m+1)}e^{-2\cdot4^m}\le4^s(e^{-2}+4e^{-8}+\dots)
\le0.14\cdot4^s\).  Hence
\(\sum_nK_n^{2s}e^{-2\nu K_n^2r}\le C_s(\nu r)^{-s}\), and Tonelli
finishes.  The \(s=1\) statement is the same computation with
\(\sum_nx_n^2e^{-2x_n^2}\) bounded above and below by positive constants.
\end{proof}

\begin{corollary}[Type-I subcritical rates are the interpolation rates]
\label{cor:v36-type-I-Hs}
Under \eqref{eq:v32-type-I} and for \(1/2<s<1\),
\begin{equation}
 \|\Lambda^su(b)\|_2^2
 \le C\Bigl(1+\log\frac{9\delta}{T_*-b}\Bigr)(T_*-b)^{-(s-1/2)},
 \label{eq:v36-type-I-Hs-rate}
\end{equation}
which agrees, up to the logarithm, with the interpolation of
\eqref{eq:v33-type-I-log-bound} and \eqref{eq:v32-type-I}; the
fractional-integral laws at orders above \(1/2\) therefore add nothing to
the type-I picture, and the critical order \(s=1/2\) is the only one at
which the law improves on interpolation.
\end{corollary}

\begin{proof}
Insert \(\varphi_{\rm rec}\le4C_4M_\delta(b)C_I^2\nu^{3/2}(T_*-t)^{-1/2}\)
in \eqref{eq:v36-Hs-law} and use
\(\int_{T_*-\delta}^b(T_*-t)^{-1/2}(b-t)^{-s}\dd t
\le C(T_*-b)^{1/2-s}\) for \(s>1/2\).  The interpolation
\(\|\Lambda^su\|_2^2\le\|A^{1/4}u\|_2^{2(2-2s)}\|\Lambda u\|_2^{2(2s-1)}\)
gives the same power with \(\log^{2(2-2s)}\).
\end{proof}

% =============================================================================
\section{The exact sign structure of the integrated octave source}
\label{sec:v36-octave-sign}
% =============================================================================

Fix a dyadic scale \(K=K_n\), the head \(v=P_Ku\), the octave field
\(w_1=\Pi_{(K,2K]}u\), and on \([a,b]\) the two octave sources
\begin{equation}
 g:=\Pi_{(K,2K]}\BNS(v,v),\qquad
 h:=\Pi_{(K,2K]}\bigl[\BNS(u)-\BNS(v,v)\bigr],
 \label{eq:v36-octave-sources}
\end{equation}
so that \(\Pi_{(K,2K]}\BNS(u)=g+h\) and, by
Proposition~\ref{prop:v32-two-stress},
\(\Phi_K^{\rm LL}=-\langle g,w_1\rangle\).  Let
\begin{equation}
 r_g(t):=\int_a^te^{-\nu A(t-s)}g(s)\,\dd s,\qquad
 r_h(t):=\int_a^te^{-\nu A(t-s)}h(s)\,\dd s .
 \label{eq:v36-octave-responses}
\end{equation}

\begin{theorem}[Exact decomposition of the integrated low--low flux]
\label{thm:v36-LL-sign}
\begin{equation}
 \boxed{\quad
 \int_a^b\Phi_K^{\rm LL}\,\dd t
 =-\int_a^b\bigl\langle g(t),e^{-\nu A(t-a)}w_1(a)\bigr\rangle\dd t
 +\Bigl(\frac12\|r_g(b)\|_2^2+\nu\int_a^b\|\Lambda r_g\|_2^2\,\dd t\Bigr)
 +\int_a^b\langle g(t),r_h(t)\rangle\,\dd t .
 \quad}
 \label{eq:v36-LL-decomposition}
\end{equation}
The middle term is nonnegative and satisfies
\begin{equation}
 0\le\frac12\|r_g(b)\|_2^2+\nu\int_a^b\|\Lambda r_g\|_2^2
 \le\frac1{\nu}\int_a^b\|A^{-1/2}g\|_2^2\,\dd t
 \le\frac{1}{\nu}\int_a^b\|P_Ku\|_4^4\,\dd t ,
 \label{eq:v36-self-formation-bound}
\end{equation}
and it is exactly the V31 intermediate-rung formed octave energy of the
head source plus its own dissipation.
\end{theorem}

\begin{proof}
Duhamel on the octave gives
\(w_1(t)=e^{-\nu A(t-a)}w_1(a)-r_g(t)-r_h(t)\).  Insert in
\(-\langle g,w_1\rangle\) and integrate.  For the middle term,
\(\partial_tr_g+\nu Ar_g=g\) with \(r_g(a)=0\), so
\(\langle g,r_g\rangle=\frac12\partial_t\|r_g\|_2^2+\nu\|\Lambda r_g\|_2^2\);
integrate.  For the bound, \(\langle g,r_g\rangle\le\|A^{-1/2}g\|_2
\|\Lambda r_g\|_2\le\frac\nu2\|\Lambda r_g\|_2^2+\frac1{2\nu}\|A^{-1/2}g\|_2^2\),
hence \(\frac12\|r_g(b)\|_2^2+\frac\nu2\int\|\Lambda r_g\|_2^2\le
\frac1{2\nu}\int\|A^{-1/2}g\|_2^2\).  The middle term of
\eqref{eq:v36-LL-decomposition} is at most twice the left side of this
inequality, which gives the constant \(1/\nu\).  Finally
\(\|A^{-1/2}\PP\operatorname{div}(v\otimes v)\|_2
\le\|v\otimes v\|_2=\|v\|_4^2\).
\end{proof}

\begin{corollary}[No sign gain from self-interaction; the indefinite term
is the cross term]
\label{cor:v36-no-sign-gain}
Over any window, the integrated low--low flux is at least
\(-\int|\langle g,e^{-\nu A(t-a)}w_1(a)\rangle|+\int\langle g,r_h\rangle\).
A window on which the low--low flux is negative, that is, on which the head
loses no energy to the octave, must have either a memory term or a cross
term \(\int\langle g,r_h\rangle\) at least as large in absolute value as
the formed energy \eqref{eq:v36-self-formation-bound}.  The cross term
pairs the head source with the octave response to the mixed and
high--high sources; it is the only place in the octave balance where the
V18 triad cancellation \eqref{eq:v18-triad-energy-cancellation} between
different slot types can act.
\end{corollary}

\begin{proof}
Drop the nonnegative middle term in \eqref{eq:v36-LL-decomposition} for
the lower bound; the rest restates the decomposition.
\end{proof}

\begin{assessment}[AC1 of Programme~\ref{prog:v36} is closed negatively]
\label{ass:v36-AC1}
A self-limiting octave inequality, in the sense of a signed bound on
\(\int\Phi_K^{\rm LL}\) by the octave energy alone, does not exist: the
self-formation part is exactly the formed energy, which is what one is
trying to bound.  The type-I logarithm and the type-II tail formation both
depend on the cross term \(\int\langle g,r_h\rangle\), whose sign is
governed by completed triads with one low--low slot and one mixed slot.
\end{assessment}

% =============================================================================
\section{The logarithmic gap between type-I and local regularity}
\label{sec:v36-local-gap}
% =============================================================================

Let \(C_3\) be the periodic constant of \(\dot H^{1/2}\hookrightarrow L^3\)
for mean-zero fields.  Write \(B_r(x_0)\) for a ball and
\(Q_r(x_0,t):=B_r(x_0)\times[t-r^2,t]\) for a parabolic cylinder.

\begin{proposition}[Type-I bounds on all scale-invariant local energies]
\label{prop:v36-local-energies}
Assume \eqref{eq:v32-type-I}, and let \(M_\delta(b)\) be as in
\eqref{eq:v34-parabolic-frequency}.  Then for every
\(t\in[T_*-\delta/2,T_*)\), every \(x_0\), and every
\(0<r\le\min\{(t-T_*+\delta)^{1/2},\pi\}\),
\begin{align}
 \frac1r\int_{B_r(x_0)}|u(t)|^2\,\dd x
 &\le\Bigl(\frac{4\pi}3\Bigr)^{1/3}C_3^2M_\delta(t)^2,
 \label{eq:v36-local-L2}\\
 \frac1{r^2}\iint_{Q_r(x_0,t)}|u|^3\,\dd x\,\dd s
 &\le C_3^3M_\delta(t)^3,
 \label{eq:v36-local-L3}\\
 \frac1r\iint_{Q_r(x_0,t)}|\nabla u|^2\,\dd x\,\dd s
 &\le\frac{C_I^2\nu^{3/2}}{r}\int_{t-r^2}^t(T_*-s)^{-1/2}\dd s
 \le\frac{2C_I^2\nu^{3/2}(T_*-t+r^2)^{1/2}}{r}
 \le2\sqrt2\,C_I^2\nu^{3/2}\quad(r^2\ge T_*-t).
 \label{eq:v36-local-dissipation}
\end{align}
All three are bounded by fixed powers of
\(M_\delta(t)\lesssim\log(9\delta/(T_*-t))\); the dissipation quantity
is bounded outright on cylinders of parabolic size at least \(T_*-t\).
\end{proposition}

\begin{proof}
By H\"older and the embedding,
\(\int_{B_r}|u|^2\le\|u\|_3^2|B_r|^{1/3}\le C_3^2\|A^{1/4}u\|_2^2
(4\pi r^3/3)^{1/3}\), and \(\|A^{1/4}u(t)\|_2\le M_\delta(t)\).  For the
cubic quantity integrate \(\|u(s)\|_3^3\le C_3^3M_\delta(t)^3\) over a
time interval of length \(r^2\).  For the dissipation use
\eqref{eq:v32-type-I} and \(\int_{t-r^2}^t(T_*-s)^{-1/2}\dd s
\le2(T_*-t+r^2)^{1/2}\le2\sqrt2\,r\) when \(r^2\ge T_*-t\).
\end{proof}

\begin{firewall}[Partial regularity is an external theorem; the gap is a
logarithm]
\label{fw:v36-CKN-gap}
The Caffarelli--Kohn--Nirenberg \(\varepsilon\)-regularity theorem, which
this series does not use and does not re-prove, states that a suitable weak
solution is regular at \((x_0,t)\) if
\(r^{-2}\iint_{Q_r}(|u|^3+|u||p-\bar p|)\le\varepsilon_0\) for one \(r\),
with \(\varepsilon_0\) absolute.  Proposition~\ref{prop:v36-local-energies}
shows that, under a type-I enstrophy rate, the velocity part of that
quantity is at most \(C_3^3M_\delta(t)^3\lesssim\log^3(9\delta/(T_*-t))\)
at every point and every parabolic scale.  Thus the type-I regime misses
the small-energy hypothesis by exactly a cube of the V33 logarithm, and by
whatever the pressure term contributes; a bounded critical norm would
remove the velocity part of that gap entirely.  For orientation only: the
known unconditional lower bound on the critical norm at a singularity, due
to Tao, is a triple logarithm, far below the single logarithm that the
type-I regime permits here; there is no contradiction between the two.
\end{firewall}

% =============================================================================
\section{A two-slot formation alternative for high-total events}
\label{sec:v36-two-slot}
% =============================================================================

\begin{theorem}[High-total tails are formed by one of three payments]
\label{thm:v36-two-slot-alternative}
Let \(w=[b-\tau,b]\), \(\tau=\theta^{-1}(\nu\rho^2)^{-1}\), and suppose
\(H_\rho^\Sigma(b)\ge(1+\kappa)\theta E_*\).  Put
\(c_\theta':=(1+\kappa)\theta-e^{-2/\theta}>0\),
\(\mathcal R_w=\sup_w\|A^{1/4}u\|_2\), \(\epsilon_w=\int_w\|\Lambda u\|_2^2\),
\(\epsilon_w^{\rm tail}=\int_wV_\rho^\Sigma\), and let
\(e_1^{\max}:=\sup_w\bigl(H_\rho^\Sigma-H_{2\rho}^\Sigma\bigr)\) be the
largest first-octave energy on the window.  Then at least one of the
following holds:
\begin{align}
 &\text{\textup{(O)}}\quad e_1^{\max}\ge\eta E_*,
 \label{eq:v36-alt-octave}\\
 &\text{\textup{(LL)}}\quad
 \mathcal R_w^2\,\epsilon_w\ge
 \frac{c_\theta'^2\,\theta\nu E_*}{64\,C_4^2\,\eta},
 \label{eq:v36-alt-LL}\\
 &\text{\textup{(HH)}}\quad
 \epsilon_w^{\rm tail}\ge\frac{c_\theta'E_*^{1/2}}{4C_4\,\rho^{1/2}},
 \label{eq:v36-alt-HH}
\end{align}
for every \(\eta\in(0,1]\).
\end{theorem}

\begin{proof}
From the tail balance and \eqref{eq:v32-two-stress},
\(\partial_tH_\rho^\Sigma+2\nu\rho^2H_\rho^\Sigma\le2|\Phi_\rho^{\rm LL}|
+2|\Phi_\rho^{\rm HH}|\), so
\(c_\theta'E_*\le2\int_w|\Phi_\rho^{\rm LL}|+2\int_w|\Phi_\rho^{\rm HH}|\).
If (O) fails, \eqref{eq:v32-LL-slot-bound} with
\(\|P_\rho u\|_4^2\le C_4\|A^{1/4}u\|_2\|\Lambda P_\rho u\|_2\) gives
\(|\Phi_\rho^{\rm LL}|\le2C_4\mathcal R_w\rho(\eta E_*)^{1/2}
\|\Lambda P_\rho u\|_2\), hence by Cauchy--Schwarz in time
\(2\int_w|\Phi_\rho^{\rm LL}|\le4C_4\mathcal R_w\rho\eta^{1/2}E_*^{1/2}
\tau^{1/2}\epsilon_w^{1/2}=4C_4\mathcal R_w\eta^{1/2}E_*^{1/2}
(\theta\nu)^{-1/2}\epsilon_w^{1/2}\).  By \eqref{eq:v32-HH-slot-bound},
\(2\int_w|\Phi_\rho^{\rm HH}|\le2C_4\rho^{1/2}E_*^{1/2}\epsilon_w^{\rm tail}\).
One of the two terms is at least \(c_\theta'E_*/2\); solving gives (LL) or
(HH).
\end{proof}

\begin{assessment}[What the alternative adds to the record--share
inequality]
\label{ass:v36-two-slot-meaning}
Option (LL) is weaker than the V35 record--share inequality
\(\mathcal R_w\epsilon_w\ge c_\theta E_*\), because it charges the
first-octave energy instead of the full tail.  Its content is the
coupling with (O): a high-total tail that is not carried by its first
octave, and is not formed by the high--high slot, needs
\(\mathcal R_w^2\epsilon_w\gtrsim\eta^{-1}\), which for small \(\eta\)
exceeds the record--share requirement.  Option (HH) is the only one with a
scale gain, \(\rho^{-1/2}\), but the tail dissipation share it requires is
summable over dyadic \(\rho\), so it does not contradict the Leray budget
either.  A type-II high-total cascade can therefore run entirely on
(HH) payments as far as these estimates can tell: order-\(E_*\) tails
formed by high--high stress against low strain, each costing tail
dissipation \(\asymp\rho_j^{-1/2}\).  This identifies the high--high slot,
not the low--low octave slot, as the cheapest formation route in the
type-II regime, which is the opposite of the type-I regime where the
octave slot carries the logarithm.
\end{assessment}

% =============================================================================
\section{Integrated V36 conclusion and next acquisition order}
\label{sec:v36-conclusion}
% =============================================================================

\begin{theorem}[What V36 proves]
\label{thm:v36-conclusion}
Relative to V1--V35, the following are proved.
\begin{enumerate}[label=\textup{(V36.\arabic*)},leftmargin=6.2em]
\item The \(\dot H^s\) fractional-integral laws \eqref{eq:v36-Hs-law} for
      \(1/2\le s<1\), with non-integrable kernel at \(s=1\); under type-I
      they reproduce the interpolation rates, so \(s=1/2\) is the only
      order at which the law gains.
\item The exact decomposition \eqref{eq:v36-LL-decomposition} of the
      integrated low--low flux: memory, a nonnegative self-formation form
      equal to the V31 formed energy plus dissipation, and the cross term
      with the mixed and high--high sources; no sign gain from
      self-interaction.
\item Under type-I, all scale-invariant local energies are bounded by
      powers of the V33 logarithm at every point and scale
      (Proposition~\ref{prop:v36-local-energies}); the gap to the
      \(\varepsilon\)-regularity hypothesis is exactly logarithmic in the
      velocity part.
\item The three-way formation alternative
      \eqref{eq:v36-alt-octave}--\eqref{eq:v36-alt-HH} for high-total
      tails; the high--high slot is the cheapest type-II formation route.
\end{enumerate}
The full Navier--Stokes stability/global-regularity theorem remains
unproved.
\end{theorem}

\begin{proof}
Theorem~\ref{thm:v36-Hs-law} and Corollary~\ref{cor:v36-type-I-Hs};
Theorem~\ref{thm:v36-LL-sign} and Corollary~\ref{cor:v36-no-sign-gain};
Proposition~\ref{prop:v36-local-energies};
Theorem~\ref{thm:v36-two-slot-alternative}.
\end{proof}

\begin{programme}[V37 acquisition order]
\label{prog:v37}
\begin{enumerate}[label=\textup{(AC\arabic*)},leftmargin=4.8em]
\item Cross term.  Expand \(\int\langle g,r_h\rangle\) with
      \(h=\Pi_{(K,2K]}[\BNS(v,w)+\BNS(w,v)+\BNS(w,w)]\) and identify, triad
      by triad, which completed V18 triads contribute; determine whether
      the part of \(h\) built from the octave field itself
      (\(w=w_1\)) produces a term of definite sign against \(g\).
\item High--high route in type-II.  The tail is high for a summable
      fraction \(f_j\) of its window (V29 persistence), so the (HH)
      formation must occur while the tail is below the threshold.  Combine
      the (HH) slot bound with the exit structure of
      Theorem~\ref{thm:v29-two-threshold-entry}: between the exit time
      and \(b_j\) the tail rises from \(h\) to \((1+\kappa)h\) with
      \(V_{\rho_j}\ge\rho_j^2h\) throughout, so the rise itself costs tail
      share at least \(\rho_j^2h\,(b_j-t_{\rho_j})\).  Determine whether
      the (HH) share \(\rho_j^{-1/2}\) and the rise duration are
      compatible with the record--share inequality one scale down, or
      whether the rise must be faster than the high--high stress permits.
\item Pressure.  For the local-regularity comparison, bound the pressure
      part \(r^{-2}\iint_{Q_r}|u||p-\bar p|\) under type-I by the same
      logarithm, using the periodic Calder\'on--Zygmund bound
      \(\|p\|_{3/2}\le C\|u\|_3^2\); record it as an external comparison
      only.
\item Keep the V29 persistence sum one-use; keep the hinge off the total
      branch.  The next compulsory companion audit is V40.
\end{enumerate}
\end{programme}

\begin{assessment}[Position after V36]
\label{ass:v36-position}
The norm-based route is now complete in both regimes: subcritical orders
add nothing, the octave self-interaction has the wrong sign to help, and
the type-I regime sits a logarithm away from the local-regularity
hypothesis.  Item (AC2) is the one place where a purely quantitative
contradiction may still be available.  A high tail dissipates at rate at
least \(\rho^2\theta E_*\), so if it is high for a fraction \(f\) of its
parabolic window \(\theta^{-1}(\nu\rho^2)^{-1}\) it spends tail Leray
share at least \(fE_*/\nu\); by the Leray budget the fractions \(f_j\) of
cofinally many windows are summable, which is the V29 one-use persistence
statement in another form.  The (HH) formation must therefore take place
while the tail is low, during the complementary part of the window, and
its share \(\asymp\rho_j^{-1/2}\) must be produced there.  No full
stability or global-regularity claim is made.
\end{assessment}

% =============================================================================
\section*{Version V36 append-only record}
\addcontentsline{toc}{section}{Version V36 append-only record}
% =============================================================================

The complete V35 mathematical source was retained before this continuation.
Its SHA--256 digest was
\begin{center}\scriptsize\ttfamily
aea7449895cd39b2d866869048033b96a47140d037dd9e3738ebb26f3fa55c01.
\end{center}
V36 proved the subcritical \(\dot H^s\) record laws, the exact sign
decomposition of the integrated low--low octave flux, the logarithmic bounds
on all scale-invariant local energies under a type-I rate, and the
three-way formation alternative for high-total tails.  No full regularity
claim is made.

% =============================================================================
\section{V37: the backward Leray cone, spike anatomy, and a correction of the
type-II schedule}
\label{sec:v37-entry}
% =============================================================================

Two of the three items of Programme~\ref{prog:v37} are settled here, and
one error in V35 is corrected.  The pressure part of the local-regularity
comparison is bounded by the same logarithm as the velocity part, so under
a type-I enstrophy rate the complete Caffarelli--Kohn--Nirenberg quantity
is at most a cube of the V33 logarithm at every point and scale.  The
Leray differential inequality of Lemma~\ref{lem:v32-leray-rate} is then
read backward from any time of large enstrophy: every enstrophy spike
casts a backward cone, \(\|\Lambda u(t)\|_2^2\gtrsim\nu^{3/2}(t_2-t)^{-1/2}\)
for \(t<t_2\), which forces the enstrophy floor between spikes, makes the
square roots of the inter-spike gaps summable, and bounds the rise time of
every spike from below by its base.  The V35 scalar schedule, which placed
negligible enstrophy between spikes, violates this cone; the corrected
schedule is written down and still satisfies all four type-II necessary
conditions.  The cross-term item (AC1) remains open; its two natural test
fields are degenerate for it.  No global regularity claim is made.  The
next compulsory companion audit is V40.

% =============================================================================
\section{The pressure part of the local-regularity comparison}
\label{sec:v37-pressure}
% =============================================================================

On the torus the pressure of a smooth divergence-free solution is the
mean-zero field \(p=-\Delta^{-1}\partial_i\partial_j(u_iu_j)\), and the
Riesz transforms are bounded on \(L^{3/2}(\Torus^3)\); let \(C_{\rm CZ}\)
be the corresponding constant, so that
\begin{equation}
 \|p\|_{3/2}\le C_{\rm CZ}\|u\otimes u\|_{3/2}=C_{\rm CZ}\|u\|_3^2 .
 \label{eq:v37-pressure-CZ}
\end{equation}

\begin{proposition}[Type-I bound on the complete local cubic quantity]
\label{prop:v37-pressure-local}
Under the hypotheses of Proposition~\ref{prop:v36-local-energies}, with
\(\bar p\) the mean of \(p\) over \(B_r(x_0)\),
\begin{equation}
 \frac1{r^2}\iint_{Q_r(x_0,t)}\bigl(|u|^3+|u||p-\bar p|\bigr)\,\dd x\,\dd s
 \le(1+2C_{\rm CZ})\,C_3^3\,M_\delta(t)^3 .
 \label{eq:v37-local-CKN-quantity}
\end{equation}
\end{proposition}

\begin{proof}
The cubic term is \eqref{eq:v36-local-L3}.  For the pressure term, H\"older
on \(B_r\) gives \(\int_{B_r}|u||p-\bar p|\le\|u\|_{L^3(B_r)}
\|p-\bar p\|_{L^{3/2}(B_r)}\), and
\(\|\bar p\|_{L^{3/2}(B_r)}=|\bar p|\,|B_r|^{2/3}
\le|B_r|^{-1/3}\|p\|_{L^1(B_r)}\le\|p\|_{L^{3/2}(B_r)}\), so
\(\|p-\bar p\|_{L^{3/2}(B_r)}\le2\|p\|_{3/2}\le2C_{\rm CZ}\|u\|_3^2\).
Hence the integrand in time is at most \(2C_{\rm CZ}\|u(s)\|_3^3
\le2C_{\rm CZ}C_3^3M_\delta(t)^3\); integrate over a time interval of
length \(r^2\).
\end{proof}

\begin{firewall}[The gap to \(\varepsilon\)-regularity is exactly
\(\log^3\), pressure included]
\label{fw:v37-CKN-gap}
Firewall~\ref{fw:v36-CKN-gap} is now complete: under a type-I enstrophy
rate, the full scale-invariant quantity of the external
\(\varepsilon\)-regularity theorem is at most
\((1+2C_{\rm CZ})C_3^3M_\delta(t)^3\lesssim\log^3(9\delta/(T_*-t))\) at
every point and every parabolic scale not exceeding
\((t-T_*+\delta)^{1/2}\).  Whether that theorem is applied is an import
decision outside this series; what is proved here is that its hypothesis
fails, if at all, by no more than a cube of the V33 logarithm.
\end{firewall}

% =============================================================================
\section{The backward Leray cone}
\label{sec:v37-backward-cone}
% =============================================================================

Write \(Y(t):=\|\Lambda u(t)\|_2^2\) and let \(C_1\) be the constant of
Lemma~\ref{lem:v32-leray-rate}, so that \(Y'\le2C_1\nu^{-3}Y^3\) on every
smooth interval.

\begin{lemma}[Backward cone of an enstrophy value]
\label{lem:v37-backward-cone}
For every \(0\le t<t_2<T_*\) with \(Y(t_2)>0\),
\begin{equation}
 \boxed{\quad
 Y(t)\ \ge\ \bigl[Y(t_2)^{-2}+4C_1\nu^{-3}(t_2-t)\bigr]^{-1/2}
 \ \ge\ \Bigl(\frac{\nu^3}{8C_1(t_2-t)}\Bigr)^{1/2}
 \quad\text{whenever }t_2-t\ge\frac{\nu^3}{4C_1Y(t_2)^2}.
 \quad}
 \label{eq:v37-backward-cone}
\end{equation}
Equivalently, a spike of height \(Y(t_2)\) cannot have risen from a base
\(Y(t_1)\) in less time than
\begin{equation}
 t_2-t_1\ \ge\ \frac{\nu^3}{4C_1}\bigl(Y(t_1)^{-2}-Y(t_2)^{-2}\bigr).
 \label{eq:v37-rise-time}
\end{equation}
\end{lemma}

\begin{proof}
\((Y^{-2})'\ge-4C_1\nu^{-3}\) integrated over \([t,t_2]\) gives
\(Y(t)^{-2}\le Y(t_2)^{-2}+4C_1\nu^{-3}(t_2-t)\).  If
\(t_2-t\ge\nu^3/(4C_1Y(t_2)^2)\), the first term is at most the second,
which gives the second inequality.  The rise-time bound is the same
integrated inequality with \(t=t_1\), solved for \(t_2-t_1\).
\end{proof}

\begin{theorem}[Leray share cast by a spike, and summable inter-spike gaps]
\label{thm:v37-cone-share}
Let \(t_2<T_*\) and \(S\le t_2\).  Then
\begin{equation}
 \boxed{\quad
 \int_{t_2-S}^{t_2}Y(t)\,\dd t
 \ \ge\ \Bigl(\frac{\nu^3}{2C_1}\Bigr)^{1/2}
 \Bigl(S^{1/2}-\bigl(\tfrac{\nu^3}{4C_1}\bigr)^{1/2}Y(t_2)^{-1}\Bigr)_+ .
 \quad}
 \label{eq:v37-cone-share}
\end{equation}
Consequently, if \(t_1<t_2<\dots\) are times with
\(Y(t_j)\ge Y_{\min}\) and gaps \(\Delta_j:=t_{j+1}-t_j\), then
\begin{equation}
 \boxed{\quad
 \sum_j\Bigl(\Delta_j^{1/2}-\bigl(\tfrac{\nu^3}{4C_1}\bigr)^{1/2}Y_{\min}^{-1}\Bigr)_+
 \ \le\ \Bigl(\frac{2C_1}{\nu^3}\Bigr)^{1/2}\frac{E_*}{2\nu}.
 \quad}
 \label{eq:v37-gap-summability}
\end{equation}
In particular the square roots of the gaps between cofinally many spikes of
unbounded height are summable, and a solution singular at \(T_*\)
satisfies \(T_*\le\bigl(E_*(2C_1)^{1/2}/(2\nu^{5/2})\bigr)^2\).
\end{theorem}

\begin{proof}
Integrate \eqref{eq:v37-backward-cone} over
\(t\in[t_2-S,t_2-s_0]\), \(s_0:=\nu^3/(4C_1Y(t_2)^2)\):
\(\int(\nu^3/(8C_1(t_2-t)))^{1/2}\dd t=(\nu^3/(2C_1))^{1/2}
(S^{1/2}-s_0^{1/2})\).  For the gaps apply this on each
\([t_j,t_{j+1}]\) with \(S=\Delta_j\) and sum; the intervals are disjoint
and the Leray identity bounds the sum of the shares by \(E_*/(2\nu)\).  The
last statement is the case of a single time \(t_2\to T_*\) with
\(S=t_2\), using Lemma~\ref{lem:v32-leray-rate} to send
\(Y(t_2)\to\infty\).
\end{proof}

\begin{proposition}[Spike anatomy]
\label{prop:v37-spike-anatomy}
Let \(w=[b-\tau,b]\) carry \(\mathscr L_{5/2}(b,\tau)\ge D\) with Leray
share \(\epsilon=\mathscr L_2(b,\tau)\), so that by
Corollary~\ref{cor:v32-spike} there is \(t_2\in w\) with
\(Y(t_2)=:Y_{\max}\ge D^4/\epsilon^4\).  Then:
\begin{enumerate}[label=\textup{(\roman*)},leftmargin=3.7em]
\item the set \(\{t\in w:Y(t)\ge Y_{\max}/2\}\) has measure at most
      \(2\epsilon/Y_{\max}\le2\epsilon^5/D^4\);
\item the enstrophy at the start of the window satisfies
      \(Y(b-\tau)\ge\bigl[Y_{\max}^{-2}+4C_1\nu^{-3}\tau\bigr]^{-1/2}\),
      and if \(\tau\ge\nu^3/(4C_1Y_{\max}^2)\) then
      \(Y(b-\tau)\ge(\nu^3/(8C_1\tau))^{1/2}\);
\item the share spent before the spike obeys
      \(\int_{b-\tau}^{t_2}Y\ge(\nu^3/(2C_1))^{1/2}
      \bigl((t_2-b+\tau)^{1/2}-(\nu^3/(4C_1))^{1/2}Y_{\max}^{-1}\bigr)_+\),
      so a spike sitting at lag \(\lambda\) after the start of its window
      costs share at least \(\asymp\nu^{3/2}\lambda^{1/2}\), independently
      of its height.
\end{enumerate}
\end{proposition}

\begin{proof}
(i) is Chebyshev.  (ii) is \eqref{eq:v37-backward-cone} with
\(t=b-\tau\).  (iii) is \eqref{eq:v37-cone-share} with \(S=t_2-b+\tau\).
\end{proof}

% =============================================================================
\section{Correction of the V35 type-II schedule}
\label{sec:v37-schedule-correction}
% =============================================================================

\begin{counterexample}[The V35 schedule violates the backward cone]
\label{cex:v37-V35-schedule}
The scalar schedule of Proposition~\ref{prop:v35-type-II-synthesis}
prescribes \(Y=A\,16^j\) on a spike of duration \(a\,32^{-j}\) inside a
window of length \(L_016^{-j}\), and \(Y\) negligible elsewhere.  By
\eqref{eq:v37-backward-cone}, at lag \(s\) before the \(j\)-th spike one
must have \(Y\ge(\nu^3/(8C_1s))^{1/2}\) as soon as
\(s\ge\nu^3/(4C_1A^2 16^{2j})\); at lag \(s=L_016^{-j}/2\) this floor is
\(\asymp\nu^{3/2}L_0^{-1/2}4^j\), which is not negligible.  The schedule
is therefore not the enstrophy of any smooth Navier--Stokes solution.
\end{counterexample}

\begin{proposition}[Corrected type-II schedule]
\label{prop:v37-corrected-schedule}
Replace the negligible background of the V35 schedule by the cone floor
\begin{equation}
 Y(t):=\Bigl(\frac{\nu^3}{8C_1(t_{j}-t)}\Bigr)^{1/2}
 \quad\text{for }t_{j-1}<t<t_j-\sigma_j,
 \qquad
 Y:=A\,16^j\ \text{on }[t_j-\sigma_j,t_j],
 \qquad \sigma_j=a\,32^{-j},
 \label{eq:v37-corrected-schedule}
\end{equation}
with spike times \(t_j\) spaced by \(\Delta_j=L_016^{-j}\).  Then the
schedule is compatible with \eqref{eq:v37-backward-cone} and
\eqref{eq:v37-rise-time} for \(A\) large, and it still satisfies the four
conditions of Proposition~\ref{prop:v35-type-II-synthesis}: the Leray share
of window \(j\) is
\(\asymp\nu^{3/2}\Delta_j^{1/2}+Aa\,2^{-j}\asymp2^{-j}\), summable; the
window ledger is at least \(A^{5/4}a\), so the \(L^{5/2}\) ledger
diverges; the record--share and reciprocal-record conditions hold with the
same constants; and the enstrophy rate is type-II.
\end{proposition}

\begin{proof}
The cone floor is exactly the equality case of
\eqref{eq:v37-backward-cone} relative to the next spike, and the rise from
the floor value at lag \(\sigma_j\), namely
\((\nu^3/(8C_1\sigma_j))^{1/2}\asymp\nu^{3/2}a^{-1/2}(4\sqrt2)^{j}\), to
\(A16^j\) needs, by \eqref{eq:v37-rise-time}, time at least
\(\nu^3/(4C_1)\cdot8C_1\sigma_j/\nu^3=2\sigma_j\), so the spike plateau
must be preceded by a rise of duration \(\asymp\sigma_j\), which is
absorbed by doubling \(a\).  The floor's contribution to the Leray share is
\((\nu^3/(2C_1))^{1/2}\Delta_j^{1/2}\asymp4^{-j}\), and to the
\(L^{5/2}\) ledger \(\int(\nu^3/(8C_1s))^{5/8}\dd s\asymp\Delta_j^{3/8}\),
both summable; the spike contributions are those computed in the proof of
Proposition~\ref{prop:v35-type-II-synthesis}.  Records
\(\mathcal R_{w_j}=2^j\) remain compatible with the spike height, and
\((T_*-t)^{1/4}Y^{1/2}\) on the spike is \(\asymp16^{-j/4}A^{1/2}4^j\to\infty\).
\end{proof}

\begin{assessment}[What the cone changes]
\label{ass:v37-cone-meaning}
The backward cone does not exclude the type-II high-total regime; it
raises its price.  Between two spikes the enstrophy cannot fall below the
Leray rate relative to the later spike, so a cofinal spike sequence is
accompanied by a cofinal sequence of Leray-rate floors, each paid once, with
summable square-root gaps.  Combined with V34's concentration theorem, the
Leray density before each record is now bounded below both by the
fractional-integral requirement \(\mathcal I(b)\gtrsim\nu^{1/2}\mathcal R\)
and by the cone floor of the next spike.  Neither bound contradicts the
budget for dyadic records, and the corrected schedule shows why: the cone
floor costs \(\asymp\Delta_j^{1/2}\), and dyadic windows have
\(\sum_j\Delta_j^{1/2}<\infty\).
\end{assessment}

% =============================================================================
\section{Status of the cross term}
\label{sec:v37-cross-term}
% =============================================================================

\begin{proposition}[Both standard test fields are degenerate for the cross
term]
\label{prop:v37-cross-term-tests}
For the balanced plane wave of
Proposition~\ref{prop:v31-hinge-balanced-nullspace}, \(\BNS(u)=0\), so
\(g=h=0\) and every term of \eqref{eq:v36-LL-decomposition} vanishes.  For
the V12 witness of Proposition~\ref{prop:v12-annular-independence} at
cutoff \(K=N\), \(P_Ku=0\), so \(g=0\) and the octave is fed only by
\(h\); the cross term vanishes identically while the octave energy is
created.  Neither field tests the sign of \(\int\langle g,r_h\rangle\).
\end{proposition}

\begin{proof}
Both statements are the cited propositions read against
\eqref{eq:v36-octave-sources}.
\end{proof}

\begin{firewall}[AC1 of Programme~\ref{prog:v37} remains open]
\label{fw:v37-cross-term-open}
No signed estimate of the cross term is proved in V37.  A test field with
\(g\ne0\), \(h\ne0\), and a controlled sign would need at least three
active shells straddling one octave; constructing one and computing the
cross term explicitly is left to V38.
\end{firewall}

% =============================================================================
\section{Integrated V37 conclusion and next acquisition order}
\label{sec:v37-conclusion}
% =============================================================================

\begin{theorem}[What V37 proves]
\label{thm:v37-conclusion}
Relative to V1--V36, the following are proved.
\begin{enumerate}[label=\textup{(V37.\arabic*)},leftmargin=6.2em]
\item Under a type-I enstrophy rate the complete local cubic quantity,
      pressure included, is at most \((1+2C_{\rm CZ})C_3^3M_\delta^3\)
      at every point and parabolic scale, \eqref{eq:v37-local-CKN-quantity}.
\item The backward Leray cone \eqref{eq:v37-backward-cone} and the rise-time
      bound \eqref{eq:v37-rise-time}.
\item The Leray share cast by any enstrophy value,
      \eqref{eq:v37-cone-share}, the summability of square-root gaps between
      cofinal spikes, \eqref{eq:v37-gap-summability}, and the resulting
      bound on the singular time.
\item Spike anatomy: plateau measure at most \(2\epsilon^5/D^4\), the
      enstrophy floor at the window start, and the share cast before the
      spike (Proposition~\ref{prop:v37-spike-anatomy}).
\item The V35 scalar schedule violates the cone; the corrected schedule
      \eqref{eq:v37-corrected-schedule} satisfies all four type-II
      conditions.
\item The two standard test fields are degenerate for the cross term.
\end{enumerate}
The full Navier--Stokes stability/global-regularity theorem remains
unproved.
\end{theorem}

\begin{proof}
Proposition~\ref{prop:v37-pressure-local};
Lemma~\ref{lem:v37-backward-cone};
Theorem~\ref{thm:v37-cone-share};
Proposition~\ref{prop:v37-spike-anatomy};
Counterexample~\ref{cex:v37-V35-schedule} and
Proposition~\ref{prop:v37-corrected-schedule};
Proposition~\ref{prop:v37-cross-term-tests}.
\end{proof}

\begin{programme}[V38 acquisition order]
\label{prog:v38}
\begin{enumerate}[label=\textup{(AC\arabic*)},leftmargin=4.8em]
\item Cross-term witness.  Build a three-shell divergence-free
      trigonometric polynomial with one shell below \(K\), one in
      \((K,2K]\), and one above \(2K\), such that \(g\ne0\) and
      \(h\ne0\), and compute \(\langle g,r_h\rangle\) to leading order in
      time from the datum; determine its sign as a function of the shell
      phases.
\item Cone plus record law.  Insert the cone floor
      \(Y\ge(\nu^3/(8C_1(t_2-t)))^{1/2}\) into the fractional-integral
      law at a record time \(b=t_2\): the memory integral then carries a
      logarithm from the floor alone, exactly as in type-I.  Determine
      whether this forces \(\mathcal R(b)^2\gtrsim\log\) at every
      enstrophy maximum in every regime, and whether the reverse inequality
      \(\mathcal I(b)\gtrsim\nu^{1/2}\mathcal R\) of V34 then closes to a
      two-sided log law at spikes.
\item Type-II budget.  Combine \eqref{eq:v37-gap-summability} with
      \eqref{eq:v35-reciprocal-sum}: gaps have summable square roots and
      records have summable reciprocals; test whether a cofinal high-total
      sequence with both properties can have windows of parabolic length
      \(\theta^{-1}(\nu\rho_j^2)^{-1}\) and records
      \(\mathcal R_{w_j}\lesssim\rho_j^{1/2}\) (the minimal record
      compatible with \(H_{\rho_j}\ge\theta E_*\)).
\item Keep the V29 persistence sum one-use; keep the hinge off the total
      branch.  The next compulsory companion audit is V40.
\end{enumerate}
\end{programme}

\begin{assessment}[Position after V37]
\label{ass:v37-position}
The type-II regime now carries five necessary conditions: divergent
\(L^{5/2}\) ledger, summable reciprocal records, Leray concentration at
records, type-II enstrophy rate, and the backward-cone floors with summable
square-root gaps.  A corrected scalar schedule satisfies all five.  The
type-I regime is a cube of one logarithm away from the small-energy
hypothesis of partial regularity, pressure included.  No full stability or
global-regularity claim is made.
\end{assessment}

% =============================================================================
\section*{Version V37 append-only record}
\addcontentsline{toc}{section}{Version V37 append-only record}
% =============================================================================

The complete V36 mathematical source was retained before this continuation.
Its SHA--256 digest was
\begin{center}\scriptsize\ttfamily
413a4fa749b8bda476927411b8742389ff8fa7187a57b5b190d933c684a055ac.
\end{center}
V37 completed the local-regularity comparison with the pressure term, proved
the backward Leray cone with its share, gap-summability, and rise-time
consequences, described the anatomy of enstrophy spikes, corrected the V35
type-II schedule, and recorded that the standard test fields are degenerate
for the octave cross term.  No full regularity claim is made.

% =============================================================================
\section{V38: an explicit cross-term witness, the cone inside the fractional
integral, and summable reciprocal cutoffs}
\label{sec:v38-entry}
% =============================================================================

The three items of Programme~\ref{prog:v38} are settled.  An explicit
three-shell divergence-free datum is constructed for which the octave
cross term of Theorem~\ref{thm:v36-LL-sign} is nonzero at leading order in
time and changes sign under a phase shift of the high shell; the cross term
therefore has no definite sign, and the type-I logarithm cannot be removed
by a sign argument on the octave balance either.  The backward cone of V37
is inserted into the Leray fractional integral: at every time the
fractional integral is at least a constant times the logarithm of the
enstrophy at that time, so every spike is seen logarithmically, while the
record law remains one-sided.  Finally, the cone gives every high-total
window a Leray share of order \(\nu/\rho_j\), so the cutoffs of a cofinal
high-total sequence must have summable reciprocals.  No global regularity
claim is made.  The next compulsory companion audit is V40.

% =============================================================================
\section{A three-shell witness for the octave cross term}
\label{sec:v38-witness}
% =============================================================================

Let \(N>M\ge1\) be integers with \(4M^2>3N^2\), for instance
\((N,M)=(10,9)\).  Put \(K:=N\), so that the octave is
\((N,2N]\), and define the real divergence-free trigonometric datum
\begin{equation}
 u_0(x)
 =\alpha\,\mathbf e_2\cos(Nx_1)
 +\beta\,\mathbf e_1\cos(Mx_2)
 +\gamma\,c\,\cos(k_2\cdot x+\phi),
 \qquad
 k_2:=(N,2M,0),\quad c:=(2M,-N,0),
 \label{eq:v38-witness-datum}
\end{equation}
with real \(\alpha,\beta,\gamma\) and phase \(\phi\).  The three
wavevectors \(k_1=(N,0,0)\), \(k_1'=(0,M,0)\), \(k_2\) satisfy
\(|k_1|=N\le K\), \(|k_1'|=M<K\), and \(|k_2|^2=N^2+4M^2>4N^2\), so the
head is \(v=P_Ku_0=\alpha\mathbf e_2\cos(Nx_1)+\beta\mathbf e_1\cos(Mx_2)\),
the octave field \(w_1=\Pi_{(N,2N]}u_0\) vanishes, and
\(w=Q_Nu_0=\gamma c\cos(k_2\cdot x+\phi)\) is a single plane wave.

\begin{proposition}[The leading-order cross term has either sign]
\label{prop:v38-witness}
Let \(u\) be the smooth solution with datum \eqref{eq:v38-witness-datum}
and let \(g,h,r_h\) be as in \eqref{eq:v36-octave-sources}--%
\eqref{eq:v36-octave-responses} with \(a=0\) and \(K=N\).  Then
\begin{align}
 g(0)&=-\frac{\alpha\beta(N^2-M^2)}{2(N^2+M^2)}
 \Bigl[\sin(m_+\!\cdot x)\,(-M,N,0)+\sin(m_-\!\cdot x)\,(M,N,0)\Bigr],
 \qquad m_\pm:=(N,\pm M,0),
 \label{eq:v38-g}\\
 \langle g(0),h(0)\rangle
 &=-\frac{\alpha\beta^2\gamma\,N(N^2-M^2)(N^2+3M^2)}{8(N^2+M^2)}\cos\phi ,
 \label{eq:v38-cross-pairing}
\end{align}
and consequently
\begin{equation}
 \boxed{\quad
 \int_0^b\langle g(t),r_h(t)\rangle\,\dd t
 =\langle g(0),h(0)\rangle\frac{b^2}{2}+O(b^3)
 =-\frac{\alpha\beta^2\gamma\,N(N^2-M^2)(N^2+3M^2)\cos\phi}{16(N^2+M^2)}\,b^2
 +O(b^3).
 \quad}
 \label{eq:v38-cross-leading}
\end{equation}
The sign of the leading term is reversed by \(\phi\mapsto\phi+\pi\), by
\(\alpha\mapsto-\alpha\), or by \(\gamma\mapsto-\gamma\).
\end{proposition}

\begin{proof}
\emph{Head source.}  With \(v_1=\beta\cos(Mx_2)\), \(v_2=\alpha\cos(Nx_1)\),
\[
 (v\cdot\nabla)v
 =-\alpha\beta\bigl[N\sin(Nx_1)\cos(Mx_2)\,\mathbf e_2
 +M\cos(Nx_1)\sin(Mx_2)\,\mathbf e_1\bigr]
 =-\frac{\alpha\beta}2\bigl[\sin(m_+\!\cdot x)(M,N,0)
 +\sin(m_-\!\cdot x)(-M,N,0)\bigr],
\]
by the product-to-sum formulas.  Projecting \((M,N,0)\) orthogonally to
\(m_+\) gives \(\frac{N^2-M^2}{N^2+M^2}(-M,N,0)\), and projecting
\((-M,N,0)\) orthogonally to \(m_-\) gives
\(\frac{N^2-M^2}{N^2+M^2}(M,N,0)\); both modes have
\(|m_\pm|=(N^2+M^2)^{1/2}\in(N,2N]\), so \(g(0)=\BNS(v,v)\) is
\eqref{eq:v38-g}, and it is nonzero because \(N\ne M\).

\emph{Mixed source.}  Since \(w\) is a single plane wave,
\(\BNS(w,w)=0\), so \(h(0)=\Pi_{(N,2N]}\bigl[\BNS(v,w)+\BNS(w,v)\bigr]\).
Its wavevectors are \(k_2\pm k_1=(2N,2M,0),(0,2M,0)\) and
\(k_2\pm k_1'=(N,3M,0),m_+\); of these only \((0,2M,0)\) and \(m_+\) can
lie in the octave, and \((0,2M,0)\) is orthogonal to \(g(0)\).  Only the
interaction of \(v':=\beta\mathbf e_1\cos(Mx_2)\) with \(w\) produces
\(m_+\).  With \(\theta:=k_2\cdot x+\phi\),
\[
 (v'\cdot\nabla)w=-\beta\gamma N\,c\cos(Mx_2)\sin\theta,
 \qquad
 (w\cdot\nabla)v'=-\beta\gamma M\,c_2\,\mathbf e_1\cos\theta\sin(Mx_2),
\]
and \(\theta-Mx_2=m_+\!\cdot x+\phi\); the \(m_+\) parts are
\(-\frac{\beta\gamma N}2c\sin(m_+\!\cdot x+\phi)\) and
\(+\frac{\beta\gamma Mc_2}2\mathbf e_1\sin(m_+\!\cdot x+\phi)\).  Hence the
\(m_+\) component of \(h(0)\) is the Leray projection of
\(\frac{\beta\gamma}2\sin(m_+\!\cdot x+\phi)\bigl(-Nc+Mc_2\mathbf e_1\bigr)\).

\emph{Pairing.}  Since the \(m_+\) component of \(g(0)\) is already
orthogonal to \(m_+\), the Leray projection may be dropped in the pairing,
which reduces to the mode \(m_+\):
\[
 \langle g(0),h(0)\rangle
 =-\frac{\alpha\beta(N^2-M^2)}{2(N^2+M^2)}\cdot\frac{\beta\gamma}2
 \cdot\bigl[(-M,N,0)\cdot(-Nc+Mc_2\mathbf e_1)\bigr]
 \cdot\int_{\Torus^3}\sin(m_+\!\cdot x)\sin(m_+\!\cdot x+\phi)\,\dd x .
\]
With \(c=(2M,-N,0)\): \((-M,N,0)\cdot(-Nc)=NMc_1-N^2c_2=2NM^2+N^3\) and
\((-M,N,0)\cdot(Mc_2\mathbf e_1)=-M^2c_2=NM^2\), so the bracket is
\(N(N^2+3M^2)\); the integral is \(\frac12\cos\phi\) for normalized Haar
measure.  This is \eqref{eq:v38-cross-pairing}.

\emph{Leading order.}  For a smooth solution,
\(r_h(t)=\int_0^te^{-\nu A(t-s)}h(s)\dd s=t\,h(0)+O(t^2)\) and
\(g(t)=g(0)+O(t)\) in \(L^2\), so
\(\int_0^b\langle g,r_h\rangle=\langle g(0),h(0)\rangle b^2/2+O(b^3)\).
\end{proof}

\begin{corollary}[AC1 of Programme~\ref{prog:v38} is closed negatively]
\label{cor:v38-no-sign}
The octave cross term \(\int\langle g,r_h\rangle\) has no definite sign on
smooth Navier--Stokes solutions: for the datum
\eqref{eq:v38-witness-datum} it is positive for \(\alpha\beta^2\gamma\cos\phi<0\)
and negative for \(\alpha\beta^2\gamma\cos\phi>0\), for all sufficiently
small \(b\).  Together with Corollary~\ref{cor:v36-no-sign-gain}, no term
of the exact decomposition \eqref{eq:v36-LL-decomposition} carries a sign
usable against the type-I logarithm.
\end{corollary}

\begin{proof}
Equation \eqref{eq:v38-cross-leading}.
\end{proof}

\begin{assessment}[The type-I logarithm is the endpoint of the octave route]
\label{ass:v38-type-I-endpoint}
Norm bounds (V35), self-interaction (V36), and the cross term (this
version) have all been examined for a gain over one logarithm per scale,
and none provides it.  Within the tools of this series the type-I regime is
therefore characterized as: uniformly small spectral tails,
\(L^2\)-precompact trajectory, critical norm at most logarithmic, no
high-total exits, no imported raw pulses, and a local cubic quantity at
most a cube of the logarithm.  Removing the logarithm would require a
genuinely new mechanism; V39--V40 turn to the type-II regime and to the
compulsory audit.
\end{assessment}

% =============================================================================
\section{The cone inside the Leray fractional integral}
\label{sec:v38-cone-fractional}
% =============================================================================

\begin{proposition}[The fractional integral logs the enstrophy]
\label{prop:v38-log-enstrophy}
For every \(a<b<T_*\) with \(4C_1Y(b)^2(b-a)\ge\nu^3\),
\begin{equation}
 \boxed{\quad
 \mathcal I(b)=\int_a^b\frac{Y(t)}{(b-t)^{1/2}}\,\dd t
 \ \ge\ \Bigl(\frac{\nu^3}{8C_1}\Bigr)^{1/2}
 \log\frac{4C_1Y(b)^2(b-a)}{\nu^3}.
 \quad}
 \label{eq:v38-log-enstrophy}
\end{equation}
At a V32 spike of height \(Y(b)\ge D^4/\epsilon^4\) this is at least
\((\nu^3/(8C_1))^{1/2}\log(4C_1D^8(b-a)/(\nu^3\epsilon^8))\); at a V1
record time, where \(Y(b)\ge\mathcal R^4/E_*\), it is at least
\((\nu^3/(8C_1))^{1/2}\log(4C_1\mathcal R^8(b-a)/(\nu^3E_*^2))\).
\end{proposition}

\begin{proof}
Insert \eqref{eq:v37-backward-cone} for \(b-t\ge s_0=\nu^3/(4C_1Y(b)^2)\):
\(\int_{s_0}^{b-a}(\nu^3/(8C_1s))^{1/2}s^{-1/2}\dd s
=(\nu^3/(8C_1))^{1/2}\log((b-a)/s_0)\).
\end{proof}

\begin{assessment}[AC2 of Programme~\ref{prog:v38}: no two-sided law]
\label{ass:v38-no-two-sided}
The cone gives \(\mathcal I(b)\gtrsim\nu^{3/2}\log\mathcal R\) at a
record, whereas Theorem~\ref{thm:v34-leray-concentration} requires
\(\mathcal I(b)\gtrsim\nu^{1/2}\mathcal R\).  The cone therefore supplies
only a logarithm of what a record needs, and the record law is an upper
bound on \(\mathcal R\), not a lower one; no two-sided law at spikes
follows.  In the type-I regime the cone reproduces the V33 logarithm from
below: the memory integral at any time carries at least
\((\nu^3/(8C_1))^{1/2}\log(Y(b)^2(b-a))\), which under
\eqref{eq:v32-type-I} is again \(\asymp\log(1/(T_*-b))\).
\end{assessment}

% =============================================================================
\section{High-total cutoffs have summable reciprocals}
\label{sec:v38-summable-cutoffs}
% =============================================================================

\begin{theorem}[Leray share cast by a high-total tail]
\label{thm:v38-high-total-share}
Let \(H_{\rho}^\Sigma(b)\ge\theta E_*\) and
\(\rho^2\ge\nu^4/(C_1\theta E_*^2)\).  Then the parabolic window
\(w=[b-\tau,b]\), \(\tau=\theta^{-1}(\nu\rho^2)^{-1}\), carries Leray
share
\begin{equation}
 \boxed{\quad
 \epsilon_w=\int_wY\,\dd t\ \ge\ \frac{\nu}{2(2C_1\theta)^{1/2}}\,\frac1\rho .
 \quad}
 \label{eq:v38-high-total-share}
\end{equation}
Consequently, if \(b_j\to T_*\), \(\rho_j\to\infty\),
\(H_{\rho_j}^\Sigma(b_j)\ge\theta E_*\), and the windows \(w_j\) are
pairwise disjoint, then
\begin{equation}
 \boxed{\quad
 \sum_j\frac1{\rho_j}\ \le\ \frac{(2C_1\theta)^{1/2}E_*}{\nu^2}.
 \quad}
 \label{eq:v38-summable-reciprocals}
\end{equation}
\end{theorem}

\begin{proof}
A tail of energy \(\theta E_*\) above \(\rho\) has enstrophy
\(Y(b)\ge\rho^2\theta E_*\).  The cone \eqref{eq:v37-backward-cone} applies
for \(b-t\ge s_0=\nu^3/(4C_1Y(b)^2)\le\nu^3/(4C_1\rho^4\theta^2E_*^2)\), and
the hypothesis on \(\rho\) makes \(s_0\le\tau/4\).  Then
\eqref{eq:v37-cone-share} with \(S=\tau\) gives
\(\epsilon_w\ge(\nu^3/(2C_1))^{1/2}(\tau^{1/2}-s_0^{1/2})
\ge\frac12(\nu^3/(2C_1))^{1/2}\tau^{1/2}\), and
\(\tau^{1/2}=(\theta\nu)^{-1/2}\rho^{-1}\).  Sum over disjoint windows
against the Leray identity.
\end{proof}

\begin{assessment}[AC3 of Programme~\ref{prog:v38}: the budget in three
ledgers]
\label{ass:v38-budget}
A cofinal high-total sequence now satisfies three summability conditions
in the Leray ledger alone: \(\sum_j\rho_j^{-1}<\infty\) (this version),
\(\sum_j\mathcal R_{w_j}^{-1}<\infty\) (V35), and
\(\sum_j\Delta_j^{1/2}<\infty\) for the gaps between its spikes (V37).
Doubled cutoffs and doubled records satisfy all three, so the Leray
ledger alone does not exclude the regime.  The one ledger in which each
window costs a fixed amount is \(L^{5/2}_tH^1\) (V32), whose finiteness is
not known.  The corrected schedule of Proposition~\ref{prop:v37-corrected-schedule}
satisfies all conditions with \(\rho_j\asymp4^j\).
\end{assessment}

% =============================================================================
\section{Integrated V38 conclusion and next acquisition order}
\label{sec:v38-conclusion}
% =============================================================================

\begin{theorem}[What V38 proves]
\label{thm:v38-conclusion}
Relative to V1--V37, the following are proved.
\begin{enumerate}[label=\textup{(V38.\arabic*)},leftmargin=6.2em]
\item The explicit three-shell datum \eqref{eq:v38-witness-datum} has
      nonzero head and mixed octave sources, and its leading-order cross
      term \eqref{eq:v38-cross-leading} takes either sign; the octave
      cross term has no definite sign.
\item The Leray fractional integral at any time is at least
      \((\nu^3/(8C_1))^{1/2}\) times the logarithm of the enstrophy at
      that time, \eqref{eq:v38-log-enstrophy}.
\item Every high-total window carries Leray share at least
      \(\nu/(2(2C_1\theta)^{1/2}\rho)\), so the cutoffs of a cofinal
      high-total sequence have summable reciprocals,
      \eqref{eq:v38-summable-reciprocals}.
\end{enumerate}
The full Navier--Stokes stability/global-regularity theorem remains
unproved.
\end{theorem}

\begin{proof}
Proposition~\ref{prop:v38-witness} and Corollary~\ref{cor:v38-no-sign};
Proposition~\ref{prop:v38-log-enstrophy};
Theorem~\ref{thm:v38-high-total-share}.
\end{proof}

\begin{programme}[V39 acquisition order]
\label{prog:v39}
\begin{enumerate}[label=\textup{(AC\arabic*)},leftmargin=4.8em]
\item Type-II parents.  A high-total tail at \(\rho_j\) is formed, by
      Theorem~\ref{thm:v36-two-slot-alternative}, either through the
      first octave or through the high--high slot; in both cases the
      forming interactions have at least one leg at or below \(\rho_j\)
      and one at or above \(\rho_j\).  Determine, using the V14 high-parent
      identity \eqref{eq:v14-high-parent-identity} at \(M=\rho_j\), the
      minimal energy that must sit in \((\rho_j/2,\rho_j]\) during the
      formation window, and whether that energy is itself a high-total
      event one scale down, so that the sequence \(\rho_j\) must be
      dyadically dense rather than sparse.
\item Combine (AC1) with \eqref{eq:v38-summable-reciprocals}: if a
      high-total event at \(\rho\) forces one at \(\rho/2\) within a
      bounded number of parabolic times, then the cutoff sequence contains
      all dyadic scales between consecutive \(\rho_j\), and
      \(\sum\rho^{-1}\) over that dense set must be compared with the
      budget.
\item Prepare the V40 audit.
\end{enumerate}
\end{programme}

\begin{assessment}[Position after V38]
\label{ass:v38-position}
The type-I regime is closed within this series at one logarithm.  The
type-II regime has a Leray-ledger price sheet with three summable series
and one order-one cost per window in the intermediate ledger.  The only
remaining lever inside the Leray ledger is density of the cutoff sequence:
if formation at one scale forces formation at the scale below, the
summable series \(\sum\rho_j^{-1}\) runs over a dyadically dense set and
its terms no longer decay geometrically.  No full stability or
global-regularity claim is made.
\end{assessment}

% =============================================================================
\section*{Version V38 append-only record}
\addcontentsline{toc}{section}{Version V38 append-only record}
% =============================================================================

The complete V37 mathematical source was retained before this continuation.
Its SHA--256 digest was
\begin{center}\scriptsize\ttfamily
946ce9bde7ad0b43aecf8dcb90b8473f75e739bfa31d90e6f1de129d3256db2a.
\end{center}
V38 constructed an explicit three-shell witness proving that the octave
cross term has no definite sign, inserted the backward cone into the Leray
fractional integral, and proved that the cutoffs of a cofinal high-total
sequence have summable reciprocals.  No full regularity claim is made.

% =============================================================================
\section{V39: tail-weighted flux densities and the scale-invariance no-go}
\label{sec:v39-entry}
% =============================================================================

Programme~\ref{prog:v39} asked whether a high-total tail at \(\rho\) forces
a high-total parent at \(\rho/2\), so that the cutoff sequence of a
type-II cascade would have to be dyadically dense.  The answer is no, and
the reason is a general one.  This version first proves two sharper flux
densities in which the tail energy at the cutoff itself appears as a
weight, \(H_\rho^{1/8}\), together with the parent-weighted form obtained
through the V14 high-parent identity.  It then proves a scale-invariance
theorem: for every density that is cubic in the amplitude and consistent
with the Navier--Stokes scaling, the own-window contribution of a dyadic
scale to the critical record is, on a type-I profile, proportional to
\(K_n^{a}\) where \(a\ge0\) is the exponent of the energy in the density.
Every density in V32--V38 has \(a\in\{0,1/4\}\), and the tail weights do
not change \(a\).  Hence no density of this kind removes the type-I
logarithm, which is the number of dyadic scales between the outer scale and
the parabolic frequency; and no parent factor forces dyadic density of a
type-II cascade, because the mixed source is controlled by the parent tail
only through a factor that is trivially of order one for a high tail.
No global regularity claim is made.  The compulsory SM/YM companion audit
is due in V40.

% =============================================================================
\section{Tail-weighted and parent-weighted flux densities}
\label{sec:v39-weighted-densities}
% =============================================================================

\begin{theorem}[Self-weighted densities]
\label{thm:v39-self-weighted}
At every smooth time and every cutoff \(\rho\), with \(w=Q_\rho u\),
\begin{align}
 |\Phi_\rho|
 &\le\|u\|_4\,\|\Lambda u\|_2\,\|w\|_4
 \le C_4\,E_*^{1/8}\,H_\rho^\Sigma{}^{1/8}\,\|\Lambda u\|_2^{5/2},
 \label{eq:v39-self-weighted-52}\\
 |\Phi_\rho|
 &\le C_4\,\|A^{1/4}u\|_2^{1/2}\,H_\rho^\Sigma{}^{1/8}\,\|\Lambda u\|_2^{9/4}.
 \label{eq:v39-self-weighted-record}
\end{align}
The same bounds hold for the sector-difference current \(d_\rho\).
\end{theorem}

\begin{proof}
\(\Phi_\rho=-\langle\BNS(u),w\rangle=-\int(u\cdot\nabla)u\cdot w\), so
H\"older with exponents \(4,2,4\) gives the first inequality.  By
Lemma~\ref{lem:v32-ladyzhenskaya}, \(\|u\|_4\le C_4^{1/2}\|u\|_2^{1/4}
\|\Lambda u\|_2^{3/4}\) and, by the lattice H\"older inequality on the
support \(\{|k|>\rho\}\), \(\|\Lambda^{3/4}w\|_2^2\le\|w\|_2^{1/2}
\|\Lambda w\|_2^{3/2}\), so \(\|w\|_4\le C_4^{1/2}H_\rho^{\Sigma\,1/8}
\|\Lambda u\|_2^{3/4}\); this is \eqref{eq:v39-self-weighted-52}.  For
\eqref{eq:v39-self-weighted-record} use instead
\(\|u\|_4\le C_4^{1/2}(\|A^{1/4}u\|_2\|\Lambda u\|_2)^{1/2}\) from
Proposition~\ref{prop:v32-corridor-form}.  For \(d_\rho\) replace \(w\) by
\(Q_\rho\widetilde u\) as in the proof of
Theorem~\ref{thm:v32-cutoff-free-density}; the tail norms are unchanged.
\end{proof}

\begin{proposition}[Parent-weighted density]
\label{prop:v39-parent-weighted}
With \(v'=P_{\rho/2}u\), \(w'=Q_{\rho/2}u\),
\begin{equation}
 |\Phi_\rho|
 \le C_4\Bigl(H_{\rho/2}^\Sigma{}^{1/8}H_\rho^\Sigma{}^{1/8}
 +E_*^{1/8}H_{\rho/2}^\Sigma{}^{1/8}\Bigr)\|\Lambda u\|_2^{5/2}
 \le2C_4E_*^{1/8}H_{\rho/2}^\Sigma{}^{1/8}\|\Lambda u\|_2^{5/2}.
 \label{eq:v39-parent-weighted}
\end{equation}
Since \(H_{\rho/2}^\Sigma\ge H_\rho^\Sigma\), this is never sharper than
\eqref{eq:v39-self-weighted-52}; its content is that the flux into the
tail above \(\rho\) is controlled by the parent tail above \(\rho/2\) with
the same eighth-power weight.
\end{proposition}

\begin{proof}
By \eqref{eq:v14-high-parent-identity},
\(\Phi_\rho=-\langle\BNS(w',u),w\rangle-\langle\BNS(v',w'),w\rangle\).
The first pairing is at most \(\|w'\|_4\|\Lambda u\|_2\|w\|_4\); the
second, after integrating by parts with \(\operatorname{div}v'=0\), equals
\(\langle\BNS(v',w),w'\rangle\) and is at most
\(\|v'\|_4\|\Lambda w\|_2\|w'\|_4\).  Apply the tail form of the
Ladyzhenskaya bound to \(w'\) and \(w\), and the energy form to \(v'\).
\end{proof}

\begin{corollary}[No dyadic-density forcing in the type-II regime]
\label{cor:v39-no-density-forcing}
For a high-total tail \(H_\rho^\Sigma\ge\theta E_*\) the weights in
\eqref{eq:v39-self-weighted-52} and \eqref{eq:v39-parent-weighted} are at
least \(\theta^{1/8}E_*^{1/8}\) throughout its rise, so neither density
forces the parent tail to exceed a fixed fraction of \(E_*\) beyond the
trivial \(H_{\rho/2}^\Sigma\ge H_\rho^\Sigma\).  Items (AC1)--(AC2) of
Programme~\ref{prog:v39} are closed negatively: a cofinal high-total
sequence need not be dyadically dense, and the sparse dyadic schedule of
Proposition~\ref{prop:v37-corrected-schedule} remains admissible.
\end{corollary}

\begin{proof}
Immediate from the two densities and \(H_{\rho/2}^\Sigma\ge H_\rho^\Sigma\).
\end{proof}

% =============================================================================
\section{The scale-invariance no-go}
\label{sec:v39-scale-invariance}
% =============================================================================

Call a function
\begin{equation}
 \varphi_\rho(t)
 =C\,E_*^{a}\,\|A^{1/4}u(t)\|_2^{b}\,\|\Lambda u(t)\|_2^{c}\,
 H_\rho^\Sigma(t)^{d}\,\rho^{e}
 \label{eq:v39-monomial-density}
\end{equation}
a \emph{monomial flux density} if it bounds \(2|\Phi_\rho|\) at every
time and cutoff, with \(a,b,c,d\ge0\), \(e\in\R\).  Amplitude homogeneity
of the cubic flux forces \(2a+b+c+2d=3\), and invariance under the
Navier--Stokes scaling \(u_\lambda(t,x)=\lambda u(\lambda^2t,\lambda x)\),
under which \(E_*\mapsto\lambda^{-1}E_*\), \(\|A^{1/4}u\|_2\mapsto
\|A^{1/4}u\|_2\), \(\|\Lambda u\|_2\mapsto\lambda^{1/2}\|\Lambda u\|_2\),
\(H_\rho\mapsto\lambda^{-1}H_{\rho/\lambda}\), \(\rho\mapsto\lambda\rho\),
and \(\Phi_\rho\mapsto\lambda\Phi_{\rho/\lambda}\), forces
\begin{equation}
 -a+\frac c2-d+e=1 .
 \label{eq:v39-scaling-identity}
\end{equation}
The densities of V32--V39 are monomial flux densities:
\(\varphi_{5/2}\) has \((a,b,c,d,e)=(\tfrac14,0,\tfrac52,0,0)\),
\(\varphi_{\rm rec}\) has \((0,1,2,0,0)\),
\eqref{eq:v39-self-weighted-52} has \((\tfrac18,0,\tfrac52,\tfrac18,0)\),
\eqref{eq:v39-self-weighted-record} has \((0,\tfrac12,\tfrac94,\tfrac18,0)\),
and the high--high slot bound \eqref{eq:v32-HH-slot-bound} in the form
\(C_4\rho^{-1/2}\|\Lambda u\|_2^3\) has \((0,0,3,0,-\tfrac12)\).

\begin{theorem}[Own-window contributions scale like \(K_n^{a}\) on a type-I
profile]
\label{thm:v39-scale-no-go}
Assume the type-I rate \eqref{eq:v32-type-I}, the critical bound
\(\|A^{1/4}u\|_2\le M\) and the flat profile
\(H_{K_n}^\Sigma\le Q/K_n\) of Theorem~\ref{thm:v34-flat-profile} on the
relevant times, and let \(\varphi\) be a monomial flux density with
\(c<4\).  Then for every dyadic scale \(K_n\le K_{\rm par}(b)\) with own
window \(w_n=[b-s_n,b]\), \(s_n=(\nu K_n^2)^{-1}\),
\begin{equation}
 \boxed{\quad
 K_n\int_{w_n}\varphi_{K_n}(t)\,\dd t
 \ \le\ C\,\frac{2^{1-c/4}}{1-c/4}\,E_*^{a}M^{b}Q^{d}
 \bigl(C_I^2\nu^{3/2}\bigr)^{c/2}\nu^{-(1-c/4)}\;K_n^{\,a}.
 \quad}
 \label{eq:v39-own-window-scaling}
\end{equation}
In particular the own-window contribution is bounded independently of the
scale when \(a=0\) and grows when \(a>0\); it never decays.  Consequently
no monomial flux density yields, through the dyadic record law, a bound on
the critical norm better than one logarithm, and the type-I logarithm of
Theorem~\ref{thm:v33-type-I-log} is the endpoint of the entire
monomial-density route.
\end{theorem}

\begin{proof}
On \(w_n\) one has \(T_*-t\le T_*-b+s_n\le2s_n\) because
\(K_n\le K_{\rm par}(b)\) means \(s_n\ge T_*-b\).  Insert the profile
bounds:
\[
 K_n\int_{w_n}\varphi_{K_n}
 \le CE_*^aM^bQ^dK_n^{1+e-d}(C_I^2\nu^{3/2})^{c/2}
 \int_{b-s_n}^b(T_*-t)^{-c/4}\dd t
 \le CE_*^aM^bQ^dK_n^{1+e-d}(C_I^2\nu^{3/2})^{c/2}
 \frac{(2s_n)^{1-c/4}}{1-c/4}.
\]
With \(s_n=(\nu K_n^2)^{-1}\) the power of \(K_n\) is
\(1+e-d-2+c/2=a\) by \eqref{eq:v39-scaling-identity}.  For the last
statement: the dyadic record law sums own-window contributions over the
\(\tfrac12\log_2(\delta/(T_*-b))\) scales between \((\nu\delta)^{-1/2}\)
and \(K_{\rm par}(b)\) (the remote parts of the windows and the scales
above \(K_{\rm par}(b)\) are handled exactly as in
Theorem~\ref{thm:v34-flat-profile} and contribute the same orders), so the
resulting bound on \(\|A^{1/4}u(b)\|_2^2\) is at least a constant times the
number of scales whenever \(a\ge0\).
\end{proof}

\begin{remark}[Why the energy exponent decides]
\label{rem:v39-energy-exponent}
The Navier--Stokes scaling leaves the critical norm invariant and multiplies
the energy by \(\lambda^{-1}\); on a type-I profile every other factor in a
monomial density is fixed by the profile, so the only trace of the scaling
left in an own-window contribution is \(E_*^aK_n^a\).  A density with
\(a<0\) would bound the cubic flux by a negative power of the energy; no
H\"older--Sobolev bound of the kind used in this series has that form,
and the amplitude constraint \(2a+b+c+2d=3\) with \(a<0\) would require
\(b+c+2d>3\), that is, more than three amplitude factors from the critical,
enstrophy, and tail norms alone.
\end{remark}

\begin{assessment}[Unified reading of V35, V36, V38 and V39]
\label{ass:v39-unified}
The slot no-go (V35), the sign decomposition (V36), the cross-term witness
(V38), and the tail weights of this version are all instances of
Theorem~\ref{thm:v39-scale-no-go}: each candidate improvement is a
monomial density with \(a\ge0\).  The tail weights \(H_\rho^{1/8}\) gain
\(K_n^{-1/8}\) per scale but cost exactly \(K_n^{1/8}\) in the
accompanying interpolation between \(\|\Lambda u\|_2^{5/2}\) and
\(\|\Lambda u\|_2^{2}\), so their \(a\) is unchanged.  Removing the
logarithm requires information that is not scale-invariant: an actual
smallness of the signed own-window flux relative to its scale-invariant
bound at cofinally many scales.  That is a dynamical statement about the
type-I profile, not a norm inequality.
\end{assessment}

% =============================================================================
\section{Integrated V39 conclusion and next acquisition order}
\label{sec:v39-conclusion}
% =============================================================================

\begin{theorem}[What V39 proves]
\label{thm:v39-conclusion}
Relative to V1--V38, the following are proved.
\begin{enumerate}[label=\textup{(V39.\arabic*)},leftmargin=6.2em]
\item The self-weighted densities
      \eqref{eq:v39-self-weighted-52}--\eqref{eq:v39-self-weighted-record}
      and the parent-weighted density \eqref{eq:v39-parent-weighted}, for
      the flux and for the sector-difference current.
\item A high-total tail does not force a high-total parent beyond the
      trivial monotonicity; the type-II cutoff sequence need not be dense.
\item The scale-invariance theorem \eqref{eq:v39-own-window-scaling}: on a
      type-I profile the own-window contribution of a dyadic scale under
      any monomial flux density is proportional to \(K_n^{a}\), \(a\ge0\)
      the energy exponent; the type-I logarithm is the endpoint of the
      monomial-density route.
\end{enumerate}
The full Navier--Stokes stability/global-regularity theorem remains
unproved.
\end{theorem}

\begin{proof}
Theorem~\ref{thm:v39-self-weighted},
Proposition~\ref{prop:v39-parent-weighted},
Corollary~\ref{cor:v39-no-density-forcing}, and
Theorem~\ref{thm:v39-scale-no-go}.
\end{proof}

\begin{programme}[V40 acquisition order, with compulsory companion audit]
\label{prog:v40}
\begin{enumerate}[label=\textup{(AC\arabic*)},leftmargin=4.8em]
\item Perform the compulsory review of the current SM and YM notes and of
      the Grand Tensor Navier--Stokes section in its current revision;
      record digests; import nothing numerical.
\item Consolidate.  State, in one theorem with pointers, the complete
      type-I/type-II classification proved in V32--V39: the necessary
      conditions of each regime, the counterexamples and schedules that
      show what the Leray ledger cannot exclude, and the two exact open
      gates (signed own-window flux in type I; finiteness of the
      \(L^{5/2}_tH^1\) ledger in type II).
\item Non-monomial information.  Identify which quantities in the series
      are not scale-invariant and are not the energy: the V29 backward
      lifetimes, the V31 stopping times, and the exit times of
      Theorem~\ref{thm:v29-two-threshold-entry}.  Determine whether an
      own-window flux bound conditioned on the exit time, rather than on
      norms, escapes Theorem~\ref{thm:v39-scale-no-go}.
\item Keep the V29 persistence sum one-use; keep the hinge off the total
      branch.
\end{enumerate}
\end{programme}

\begin{assessment}[Position after V39]
\label{ass:v39-position}
The series now has a theorem explaining its own limits: scale-invariant
densities cannot beat one logarithm per dyadic scale on a type-I profile,
and the type-II regime is not forced to be dense.  The next version must
audit the companions and consolidate, and then look for information that
is genuinely not scale-invariant, of which the stopping-time structure of
V29--V31 is the only candidate already present in the series.  No full
stability or global-regularity claim is made.
\end{assessment}

% =============================================================================
\section*{Version V39 append-only record}
\addcontentsline{toc}{section}{Version V39 append-only record}
% =============================================================================

The complete V38 mathematical source was retained before this continuation.
Its SHA--256 digest was
\begin{center}\scriptsize\ttfamily
b2d6db9a9ba7a54da42f2419f8f1bfc8e0d1ec9baf6fc86f49caed2f195b5707.
\end{center}
V39 proved the self- and parent-weighted flux densities, closed the
dyadic-density question negatively, and proved the scale-invariance
theorem showing that every monomial flux density yields at best one
logarithm per dyadic scale on a type-I profile.  No full regularity claim
is made.

% =============================================================================
\section{V40: compulsory companion audit and the consolidated two-regime
classification}
\label{sec:v40-entry}
% =============================================================================

This is a tenth-version continuation.  The compulsory review of the two
companion programmes is performed first, together with a read-only review
of the current Grand Tensor Navier--Stokes section.  The mathematical work
of the version is consolidation: V32--V39 are collected into one
classification theorem with pointers, so that the two exact open gates and
the schedules that show what the Leray ledger cannot exclude are stated in
one place.  A short scale-covariance lemma then settles item (AC3) of
Programme~\ref{prog:v40}: the V29--V31 stopping-time structure transforms
covariantly under the Navier--Stokes scaling with the energy-type
thresholds carrying the same exponent as the energy, so conditioning on
exit times does not escape the scale-invariance theorem of V39.  No global
regularity claim is made.

% =============================================================================
\section{Compulsory V40 review of the companion notes and the Grand Tensor}
\label{sec:v40-companion-audit}
% =============================================================================

The files inspected read-only, with their SHA--256 digests, were
\begin{align}
 &\texttt{Renewal\_SM\_Einstein\_Continuum\_Research\_Notes\_V59.tex},
 \notag\\[-2pt]
 &\qquad
 \texttt{79b7cea05ef2097067e8421d6913ac7ca660e4e1f10c0ddc596fbceceb31c468},
 \label{eq:v40-SM-checkpoint}\\
 &\texttt{Renewal\_Yang\_Mills\_Mass\_Gap\_Research\_Notes\_V42.tex},
 \notag\\[-2pt]
 &\qquad
 \texttt{980514090d33a8c6799bf9aef6a011234e1ba62415f5b58de5e0434f4e5af7dc},
 \label{eq:v40-YM-checkpoint}\\
 &\texttt{Renewal\_Grand\_Tensor\_Monograph\_V073.tex},
 \notag\\[-2pt]
 &\qquad
 \texttt{5e1ca07ab0094a79152c3fa25f87038a375aaa8c0a201f599a76e822be6c0d0c}.
 \label{eq:v40-GT-checkpoint}
\end{align}
The review concentrated on SM V55--V59, YM V38--V41, and the theorems of
the Grand Tensor section ``Navier--Stokes: the exact relative-form
continuation certificate'' added since the V067 revision audited in V1.

\emph{SM V55--V59.}  V55 shows that the polarization logarithm of the
noncompact hypercharge branch must be treated as a complete-loop local
counterterm and that the running coefficient loses every fixed margin at a
finite horizon.  V56--V57 replace the frozen hypercharge coefficient by a
heavy-field Schur complement, prove its relative stability, extend it to
trace-class operator perturbations with the exact Fredholm normalization,
and record that heavy zero modes must be fixed or retained as boundary
sectors.  V58 proves that a positive gauge block cannot repair a negative
metric principal block and that the Maxwell action is conformally
invariant in four dimensions.  V59 reduces the linearized graviton to two
transverse--traceless oscillators with positive transfer and a reflection
positive continuum limit, and states that all nonlinear, conformal, and
constraint inputs remain open.  These are linear or Gaussian
constructions in field theory; none has a Navier--Stokes tail balance as
hypothesis.

\emph{YM V38--V41.}  V38 sums an infinite block-local root and proves its
all-arity score card for a calibrated even-entire spectral model.  V39
derives exact Gaussian--Laguerre transforms and proves that the separated
Haar series fails the even-entire criterion while its paired images cancel
the antipodal jet.  V40 locates a nontrivial imaginary zero of the exact
image pair, proves collapse of the physical analytic radius, and rules out
a \(\tau\)-uniform zero-free-strip proof.  V41 proves that the exact
heat-normalized logarithmic score grows faster than every exponential
along a spectral subsequence, disproving the previous geometric-bridge
gates, replaces the bridge by an exact additive one, and isolates a merge
cocycle as the corrected gate.  The YM discipline relevant here is that a
conjectured sufficient gate was disproved by an exact computation rather
than weakened; V38 of this series did the same for the octave cross term.

\emph{Grand Tensor V073, NS section.}  Relative to V067 the section
contains, among others, theorems titled ``A finite maximal time forces a
scale-critical dyadic shell,'' ``Nested-cylinder compactness and
dissipation-defect alternative,'' ``Terminal analytic closure and
critical-shell radius collapse,'' ``Scale-vanishing entropy and
palinstrophy escalation,'' and ``Two-centre source/pressure rectangle and
high--high backscatter,'' with firewalls stating that the scale boundary
must remain open and that no global regularity follows from the critical
dipole.  The first is the Grand Tensor form of the V1 record corridor; the
second is the compactness route whose type-I instance V32 made exact
(uniform tail smallness and \(L^2\)-precompactness); the last is the
high--high slot of the V32 two-stress anatomy.  None of them bounds the
\(L^{5/2}_tH^1\) ledger or the signed own-window flux.

\begin{theorem}[V40 companion-audit decision]
\label{thm:v40-companion-decision}
\begin{enumerate}[label=\textup{(\roman*)},leftmargin=3.7em]
\item No SM V55--V59 Schur, Fredholm, conformal, or transverse--traceless
      constant is imported; the SM observation that a positive block cannot
      repair a negative principal block is type-correct for the V36
      decomposition, where the nonnegative self-formation form cannot
      repair the indefinite cross term, and was proved there directly.
\item No YM V38--V41 root, Laguerre, image-pair, or bridge estimate is
      imported.  The YM practice of disproving a candidate gate by an exact
      witness is the practice of V38 here; nothing further transfers.
\item The Grand Tensor V073 NS theorems are consistent with, and in three
      cases anticipated by, V32--V39; none supplies either open gate.  The
      relative-form certificate remains conditional on an \(H^s\) relative
      estimate that this series has not proved.
\end{enumerate}
\end{theorem}

\begin{proof}
The hypotheses of every cited companion or Grand Tensor result are
cochain, determinant, Gaussian, thermal-block, or abstract-tensor objects,
or, in the Grand Tensor NS section, statements whose conclusions are
explicitly marked conditional or open by their own firewalls.  None has the
V23 sector-tail balances, the Ladyzhenskaya density, or the Leray identity
on the torus as hypothesis, so no inequality transfers.
\end{proof}

% =============================================================================
\section{The consolidated classification}
\label{sec:v40-consolidation}
% =============================================================================

\begin{theorem}[Two-regime classification of a hypothetical finite-time
singularity, V32--V39]
\label{thm:v40-classification}
Let \(u\) be a smooth mean-zero solution on \([0,T_*)\times\Torus^3\) with
\(T_*<\infty\) maximal, \(E_*=\sup\|u\|_2^2\).  Exactly one of the
following holds.

\smallskip
\noindent\textbf{Regime I (finite intermediate ledger).}\quad
\(\int_0^{T_*}\|\Lambda u\|_2^{5/2}\,\dd t<\infty\).  Then:
\begin{enumerate}[label=\textup{(I.\arabic*)},leftmargin=4.4em]
\item no cofinal high-total sequence exists
      (Corollary~\ref{cor:v32-high-total-divergence});
\item every high-total window that does occur costs a fixed
      \(L^{5/2}\) ledger \eqref{eq:v32-high-total-window-debit}, so only
      finitely many occur.
\end{enumerate}
If moreover the enstrophy rate is type-I, \eqref{eq:v32-type-I}, then in
addition:
\begin{enumerate}[label=\textup{(I.\arabic*)},leftmargin=4.4em,resume]
\item the spectral tails are uniformly small, the trajectory is
      \(L^2\)-precompact, and \(u(t)\to u(T_*)\) strongly in \(L^2\)
      (Theorem~\ref{thm:v32-type-I-firewall});
\item \(\|A^{1/4}u(b)\|_2^2\le2A_\delta+2B^2\log^2(9\delta/(T_*-b))\)
      (Theorem~\ref{thm:v33-type-I-log}), record doubling times are
      double-exponentially close to \(T_*\), and the dyadic critical masses
      have the flat profile of Theorem~\ref{thm:v34-flat-profile};
\item the imported raw pulse \eqref{eq:imported-raw-pulse} occurs for
      finitely many records (Theorem~\ref{thm:v34-raw-pulse-finite});
\item every scale-invariant local energy, pressure included, is bounded by
      a cube of the same logarithm at every point and parabolic scale
      (Propositions~\ref{prop:v36-local-energies},
      \ref{prop:v37-pressure-local});
\item the critical norm is bounded if and only if the own-window fluxes
      are summable over the dyadic scales
      (Proposition~\ref{prop:v34-log-free-criterion}), and no monomial
      flux density, slot bound, self-interaction sign, or cross-term sign
      achieves this
      (Theorems~\ref{thm:v39-scale-no-go}, \ref{thm:v36-LL-sign},
      Propositions~\ref{prop:v35-slot-no-go}, \ref{prop:v38-witness}).
\end{enumerate}

\smallskip
\noindent\textbf{Regime II (divergent intermediate ledger).}\quad
\(\int_0^{T_*}\|\Lambda u\|_2^{5/2}\,\dd t=\infty\).  Then the enstrophy
rate is type-II, and if a cofinal high-total sequence with windows
\(w_j\) and cutoffs \(\rho_j\) exists, the following hold simultaneously:
\begin{enumerate}[label=\textup{(II.\arabic*)},leftmargin=4.8em]
\item each window carries \(L^{5/2}\) ledger at least
      \(2c_\theta E_*^{3/4}\) and contains an enstrophy spike of height at
      least \(D^4/\epsilon_{w_j}^4\) with plateau measure at most
      \(2\epsilon_{w_j}^5/D^4\)
      (Corollary~\ref{cor:v32-spike}, Proposition~\ref{prop:v37-spike-anatomy});
\item \(\mathcal R_{w_j}\epsilon_{w_j}\ge c_\theta E_*\) and
      \(\sum_j\mathcal R_{w_j}^{-1}<\infty\)
      (Theorem~\ref{thm:v35-record-share},
      Corollary~\ref{cor:v35-summable-reciprocals});
\item \(\sum_j\rho_j^{-1}<\infty\)
      (Theorem~\ref{thm:v38-high-total-share});
\item the square roots of the gaps between consecutive spikes are
      summable, and the enstrophy between spikes lies above the backward
      cone of the next spike (Theorem~\ref{thm:v37-cone-share});
\item at every running record the Leray fractional integral reaches
      \(\nu^{1/2}\mathcal R/(18.4C_4)\), concentrated within lag
      \(\asymp E_*^2\nu^{-3}\mathcal R^{-2}\)
      (Theorem~\ref{thm:v34-leray-concentration});
\item each high-total tail is formed through the first octave, the
      low--low slot with \(\mathcal R_w^2\epsilon_w\gtrsim\eta^{-1}\), or
      the high--high slot with tail share \(\gtrsim\rho_j^{-1/2}\)
      (Theorem~\ref{thm:v36-two-slot-alternative}), and no parent
      high-total event is forced
      (Corollary~\ref{cor:v39-no-density-forcing}).
\end{enumerate}
The corrected scalar schedule \eqref{eq:v37-corrected-schedule} satisfies
(II.1)--(II.6) with a finite Leray budget; hence Regime II is not excluded
by these conditions.  In both regimes the critical record obeys the
fractional-integral law \eqref{eq:v33-record-law} and its running form
\eqref{eq:v33-running-record}.
\end{theorem}

\begin{proof}
The dichotomy is the finiteness or not of one integral.  Each item is the
cited result; the hypotheses of the cited results are those of the
respective regime, and the schedule statement is
Proposition~\ref{prop:v37-corrected-schedule} together with
Theorem~\ref{thm:v38-high-total-share}, whose condition
\(\sum4^{-j}<\infty\) it satisfies.
\end{proof}

\begin{firewall}[What the classification does not say]
\label{fw:v40-classification-scope}
Theorem~\ref{thm:v40-classification} does not assert that either regime
occurs, nor that either is empty.  Regime I with a type-I rate is confined
to logarithmic critical growth but not excluded; Regime I without a type-I
rate is constrained only by (I.1)--(I.2); Regime II is constrained by six
necessary conditions that a scalar schedule satisfies.  The two exact open
gates are the summability of the signed own-window fluxes in Regime I and
a mechanism beyond the Leray ledger in Regime II.
\end{firewall}

% =============================================================================
\section{Scale covariance of the stopping-time structure}
\label{sec:v40-scale-covariance}
% =============================================================================

\begin{lemma}[The V29 exit structure is scale-covariant]
\label{lem:v40-exit-covariance}
Under \(u_\lambda(t,x)=\lambda u(\lambda^2t,\lambda x)\) the tail
functions transform as \(H_\rho^\Sigma[u_\lambda](t)
=\lambda^{-1}H_{\rho/\lambda}^\Sigma[u](\lambda^2t)\), and likewise for
\(\mathcal Y_\rho\), \(V_\rho^\Sigma\), and \(\Phi_\rho\) up to the
factors \(\lambda^{-1},\lambda,\lambda\).  Consequently the high-balanced
region, the backward lifetimes, the exit times, and the relative lifetimes
of V29 satisfy
\begin{equation}
 \mathfrak O_{\varepsilon,h}[u_\lambda]
 =\{(\rho,t):(\rho/\lambda,\lambda^2t)\in
 \mathfrak O_{\lambda\varepsilon,\lambda h}[u]\},
 \qquad
 \ell_{\varepsilon,h}(\rho;b)[u_\lambda]
 =\lambda^{-2}\ell_{\lambda\varepsilon,\lambda h}(\rho/\lambda;\lambda^2b)[u],
 \label{eq:v40-exit-covariance}
\end{equation}
and the relative lifetime \(\ell_\alpha^\Sigma\) transforms with
\(\alpha\) unchanged.  Thresholds proportional to the energy,
\(h=\theta E_*\), \(\varepsilon=\theta'E_*\), are invariant as
proportions.
\end{lemma}

\begin{proof}
\(\widehat{u_\lambda}(k,t)=\lambda\cdot\lambda^{-3}\widehat u(k/\lambda,
\lambda^2t)\) on the rescaled lattice, so a tail sum over \(|k|>\rho\)
becomes \(\lambda^{-1}\) times the tail sum over \(|k'|>\rho/\lambda\) at
time \(\lambda^2t\), by \(\|u_\lambda(t)\|_2^2=\lambda^{-1}
\|u(\lambda^2t)\|_2^2\).  The other transformations follow from the
balances \eqref{eq:v23-sector-high-balance}.  Substituting in the
definitions of \(\mathfrak O_{\varepsilon,h}\), of the lifetimes, and of
the exit times gives \eqref{eq:v40-exit-covariance}.
\end{proof}

\begin{proposition}[Exit-conditioned bounds do not escape the no-go]
\label{prop:v40-exit-no-escape}
Let \(\psi_\rho(t)\) be any bound on \(2|\Phi_\rho(t)|\) of the form
\eqref{eq:v39-monomial-density} multiplied by a function of the V29
lifetimes and exit times with energy-proportional thresholds, which is
homogeneous of degree zero under \eqref{eq:v40-exit-covariance}.  Then on a
type-I profile the own-window contribution of scale \(K_n\) under
\(\psi\) is again proportional to \(K_n^{a}\), \(a\ge0\).
\end{proposition}

\begin{proof}
On the type-I profile the lifetimes at scale \(K_n\), measured in units
of the parabolic time \((\nu K_n^2)^{-1}\), are bounded independently of
\(n\) by Corollary~\ref{cor:v29-parabolic-band-lifetime} with
\(h=\theta E_*\), and the relative lifetimes are dimensionless.  A
degree-zero function of them is therefore bounded uniformly in \(n\), and
the computation in the proof of Theorem~\ref{thm:v39-scale-no-go} applies
unchanged.
\end{proof}

\begin{assessment}[AC3 of Programme~\ref{prog:v40} is closed]
\label{ass:v40-AC3}
The stopping-time structure carries no information that is not scale
covariant with the energy as the only dimensionful scale.  The remaining
non-invariant inputs are the outer scale of the torus, the initial datum,
and the crossover region between the outer scale and the type-I profile;
none of them acts at cofinally many dyadic scales, so none can supply the
summability required by Proposition~\ref{prop:v34-log-free-criterion}.
Within this series the type-I logarithm is final.
\end{assessment}

% =============================================================================
\section{Integrated V40 conclusion and next acquisition order}
\label{sec:v40-conclusion}
% =============================================================================

\begin{theorem}[What V40 proves]
\label{thm:v40-conclusion}
Relative to V1--V39: the companion and Grand Tensor audit decision
(Theorem~\ref{thm:v40-companion-decision}); the consolidated two-regime
classification with its six Regime-II necessary conditions and the
schedule that satisfies them (Theorem~\ref{thm:v40-classification}); the
scale covariance of the V29 exit structure
(Lemma~\ref{lem:v40-exit-covariance}) and the resulting closure of the
exit-conditioned route (Proposition~\ref{prop:v40-exit-no-escape}).  The
full Navier--Stokes stability/global-regularity theorem remains unproved.
\end{theorem}

\begin{programme}[V41 acquisition order: new mechanisms only]
\label{prog:v41}
\begin{enumerate}[label=\textup{(AC\arabic*)},leftmargin=4.8em]
\item Type-I compactness.  Rescale by \(\lambda_k=(T_*-t_k)^{1/2}\) along
      \(t_k\to T_*\); under \eqref{eq:v32-type-I} the rescaled fields
      have uniformly bounded \(\dot H^1\) norm on the rescaled tori and
      critical norm at most \(\log(1/\lambda_k)\).  Determine what the
      V32 \(L^2\)-precompactness and the V34 flat profile say about local
      limits, and whether an ancient limit with bounded \(\dot H^1\) is
      forced to be trivial; state any Liouville input as an import.
\item Type-II beyond the Leray ledger.  The six conditions of
      Theorem~\ref{thm:v40-classification} all live in the Leray or
      \(L^{5/2}\) ledgers.  Test whether the helicity balance of
      V22--V28, which is not a norm of \(u\), gives a seventh condition of
      a different type on the corrected schedule.
\item Physical-space localization.  Determine whether the fractional
      record law has a localized form for \(\|u\|_{L^3(B_r)}\) with a
      local density; if so, the V36--V37 local-energy bounds could be
      sharpened at the singular points rather than globally.
\item The next compulsory companion audit is V45.
\end{enumerate}
\end{programme}

\begin{assessment}[Position after V40]
\label{ass:v40-position}
V32--V40 replaced the cutoff-multiplicity bottleneck of V19--V31 by a
ledger classification with explicit necessary conditions in each regime,
one logarithmic endpoint in the type-I regime, and a proof that the
density route cannot go further.  The next versions must bring a
mechanism that is not a scale-invariant norm inequality: compactness of
rescaled type-I profiles, the helicity balance, or physical-space
localization.  No full stability or global-regularity claim is made.
\end{assessment}

% =============================================================================
\section*{Version V40 append-only record}
\addcontentsline{toc}{section}{Version V40 append-only record}
% =============================================================================

The complete V39 mathematical source was retained before this continuation.
Its SHA--256 digest was
\begin{center}\scriptsize\ttfamily
691ff4014f9a822d58e076eeefb504230f61fed7dd0cf3f7fb4640b4e9c03737.
\end{center}
V40 performed the compulsory companion and Grand Tensor audit, consolidated
V32--V39 into the two-regime classification theorem, and proved the scale
covariance of the stopping-time structure, closing the exit-conditioned
route.  No full regularity claim is made.

% =============================================================================
\section{V41: rescaled type-I profiles, the ancient limit, and the
finite-energy Liouville theorem}
\label{sec:v41-entry}
% =============================================================================

V40 closed the density route and asked for mechanisms that are not
scale-invariant norm inequalities.  This version implements the first of
them, item (AC1) of Programme~\ref{prog:v41}.  A hypothetical singularity
with a type-I enstrophy rate is rescaled at a point of the singular set by
its parabolic scale.  Three external theorems are imported and marked as
such: the dissipation form of the Caffarelli--Kohn--Nirenberg
\(\varepsilon\)-regularity criterion, which supplies a concentration point;
the Aubin--Lions compactness lemma; and the Fujita--Kato small-data theorem
with its uniqueness class.  With these, the rescaled solutions converge on
compact sets to a nontrivial smooth ancient solution \(U\) of the
Navier--Stokes equations on \((-\infty,0)\times\R^3\) which inherits the
type-I enstrophy bound \(\|\nabla U(s)\|_2^2\le C_I^2\nu^{3/2}(-s)^{-1/2}\).
The new theorem is a Liouville statement for this class: an ancient
solution with uniformly bounded enstrophy on \((-\infty,-1]\) and finite
energy is zero.  Consequently the ancient limit of a type-I singularity has
infinite energy, and the singular point does not admit a finite-energy
blow-up profile.  The remaining gate in the type-I regime is therefore a
Liouville theorem for infinite-energy ancient solutions with bounded
enstrophy, and it is stated exactly.  No global regularity claim is made.
The next compulsory companion audit is V45.

% =============================================================================
\section{Imports and standing conventions}
\label{sec:v41-imports}
% =============================================================================

\begin{firewall}[External theorems used in V41]
\label{fw:v41-imports}
The following are used as black boxes and are not re-proved.
\begin{enumerate}[label=\textup{(E\arabic*)},leftmargin=3.6em]
\item \emph{CKN dissipation criterion.}  There is an absolute
      \(\varepsilon_1>0\) such that if a suitable weak solution satisfies
      \(\limsup_{r\to0}r^{-1}\iint_{Q_r(z)}|\nabla u|^2\le\varepsilon_1\)
      at a point \(z=(x_0,t_0)\), then \(u\) is bounded in a neighbourhood
      of \(z\).  Smooth solutions on \([0,T_*)\times\Torus^3\) are suitable.
\item \emph{Aubin--Lions.}  A sequence bounded in
      \(L^2(J;H^2(B))\) with time derivatives bounded in \(L^2(J;L^2(B))\)
      is precompact in \(L^2(J;H^1(B))\), for a bounded ball \(B\) and a
      compact interval \(J\).
\item \emph{Fujita--Kato.}  There is \(\varepsilon_{\rm FK}>0\) such that
      for divergence-free \(U_0\in\dot H^{1/2}(\R^3)\) with
      \(\|U_0\|_{\dot H^{1/2}}\le\varepsilon_{\rm FK}\) the Navier--Stokes
      equations have a unique global solution in
      \(C([0,\infty);\dot H^{1/2})\cap L^2_{\rm loc}([0,\infty);\dot H^{3/2})\),
      with \(\sup_t\|U(t)\|_{\dot H^{1/2}}\le2\varepsilon_{\rm FK}\).
\item \emph{Serrin regularity.}  A distributional solution on an open
      space-time set with \(U\in L^\infty_{\rm loc}L^6\) and
      \(\nabla U\in L^2_{\rm loc}\) is smooth there.
\end{enumerate}
The Sobolev, Agmon, and Riesz-transform inequalities used below hold on
every torus \(\Torus^3_L=(\R/2\pi L^{-1}\mathbb Z)^3\) for mean-zero
fields with constants independent of \(L\), because both sides scale
identically under dilation.
\end{firewall}

Throughout, \(u\) is a smooth mean-zero solution on
\([0,T_*)\times\Torus^3\), \(T_*<\infty\) is maximal, and the type-I
enstrophy rate \eqref{eq:v32-type-I} holds on \([T_*-\delta,T_*)\).  Write
\(Y=\|\Lambda u\|_2^2\), \(Z=\|Au\|_2^2\).

% =============================================================================
\section{A concentration point and the rescaled sequence}
\label{sec:v41-rescaling}
% =============================================================================

\begin{lemma}[A singular time has a dissipation-concentration point]
\label{lem:v41-concentration-point}
There exist \(x_0\in\Torus^3\) and \(r_k\downarrow0\) with
\begin{equation}
 r_k^{-1}\iint_{Q_{r_k}(x_0,T_*)}|\nabla u|^2\,\dd x\,\dd t>\varepsilon_1
 \qquad(k\ge1),
 \qquad Q_r(x_0,T_*)=B_r(x_0)\times(T_*-r^2,T_*).
 \label{eq:v41-concentration}
\end{equation}
\end{lemma}

\begin{proof}
If every point \((x,T_*)\) had
\(\limsup_{r\to0}r^{-1}\iint_{Q_r}|\nabla u|^2\le\varepsilon_1\), then by
(E1) \(u\) would be bounded in a neighbourhood of each such point; by
compactness of \(\Torus^3\), \(u\) would be bounded on
\([T_*-\eta,T_*)\times\Torus^3\) for some \(\eta>0\), and a bounded
solution extends smoothly past \(T_*\) (Serrin class
\(L^\infty_tL^\infty_x\)), contradicting maximality.  Hence some \(x_0\)
has limit superior exceeding \(\varepsilon_1\); choose \(r_k\) realizing
it.
\end{proof}

Define the rescaled fields on the tori \(\Torus^3_k:=\Torus^3_{r_k}\)
\begin{equation}
 U_k(s,y):=r_k\,u\bigl(T_*+r_k^2s,\ x_0+r_ky\bigr),
 \qquad
 P_k(s,y):=r_k^2\,p\bigl(T_*+r_k^2s,\ x_0+r_ky\bigr),
 \qquad
 s\in[-T_*r_k^{-2},0).
 \label{eq:v41-rescaled}
\end{equation}
Each \((U_k,P_k)\) is a smooth mean-zero divergence-free solution of the
Navier--Stokes equations on its torus with the same viscosity \(\nu\).

\begin{proposition}[Uniform bounds on the rescaled sequence]
\label{prop:v41-uniform-bounds}
Put \(Y_k(s):=\|\nabla U_k(s)\|_{L^2(\Torus^3_k)}^2\),
\(Z_k(s):=\|\Delta U_k(s)\|_{L^2(\Torus^3_k)}^2\).  For all \(k\) with
\(r_k^2\le\delta\) and all \(s\in[-\delta r_k^{-2},0)\):
\begin{align}
 Y_k(s)&\le C_I^2\nu^{3/2}(-s)^{-1/2},
 \label{eq:v41-rescaled-type-I}\\
 \|U_k(s)\|_{L^6(\Torus^3_k)}&\le C_6\,C_I\nu^{3/4}(-s)^{-1/4},
 \qquad
 \|P_k(s)\|_{L^3(\Torus^3_k)}\le C_{\rm CZ}C_6^2C_I^2\nu^{3/2}(-s)^{-1/2},
 \label{eq:v41-rescaled-L6}\\
 \nu\int_{-S}^{-\eta}Z_k(s)\,\dd s
 &\le C_I^2\nu^{3/2}S^{-1/2}
   +2C_1\nu^{-3}C_I^6\nu^{9/2}\int_{-S}^{-\eta}(-s)^{-3/2}\dd s
 \le C(S,\eta),
 \label{eq:v41-rescaled-H2}\\
 \int_{-S}^{-\eta}\|\partial_sU_k\|_{L^2(\Torus^3_k)}^2\,\dd s
 &\le C'(S,\eta),
 \label{eq:v41-rescaled-time-derivative}
\end{align}
for every \(0<\eta<S\), with constants independent of \(k\).  Moreover
\begin{equation}
 \iint_{B_1(0)\times(-1,0)}|\nabla U_k|^2\,\dd y\,\dd s>\varepsilon_1 .
 \label{eq:v41-rescaled-nontrivial}
\end{equation}
\end{proposition}

\begin{proof}
\(\|\nabla U_k(s)\|_{L^2(\Torus^3_k)}^2=r_k\|\Lambda u(T_*+r_k^2s)\|_2^2\),
and \eqref{eq:v32-type-I} gives \eqref{eq:v41-rescaled-type-I}.  The
mean-zero Sobolev inequality on \(\Torus^3_k\) gives the \(L^6\) bound,
and \(P_k=-\Delta^{-1}\partial_i\partial_j(U_{k,i}U_{k,j})\) with the
Riesz bound on \(L^3\) gives the pressure bound.  The \(H^1\) balance
\(\frac12Y_k'+\nu Z_k\le\frac\nu2Z_k+C_1\nu^{-3}Y_k^3\) of
Lemma~\ref{lem:v32-leray-rate}, valid on every torus with the same
constant, integrated over \([-S,-\eta]\), gives
\eqref{eq:v41-rescaled-H2} after inserting
\eqref{eq:v41-rescaled-type-I}.  For the time derivative,
\(\partial_sU_k=\nu\Delta U_k-\PP(U_k\cdot\nabla U_k)\), and
\(\|U_k\cdot\nabla U_k\|_2\le\|U_k\|_6\|\nabla U_k\|_3
\le C\,Y_k^{1/2}\,Y_k^{1/4}Z_k^{1/4}\), so
\(\|\partial_sU_k\|_2^2\le2\nu^2Z_k+CY_k^{3/2}Z_k^{1/2}\), whose integral
over \([-S,-\eta]\) is bounded by \eqref{eq:v41-rescaled-H2} and
Cauchy--Schwarz.  Finally the change of variables
\(\iint_{B_1\times(-1,0)}|\nabla U_k|^2=r_k^{-1}\iint_{Q_{r_k}}|\nabla u|^2\)
and \eqref{eq:v41-concentration} give \eqref{eq:v41-rescaled-nontrivial}.
\end{proof}

% =============================================================================
\section{The ancient limit}
\label{sec:v41-ancient-limit}
% =============================================================================

\begin{theorem}[A type-I singularity produces a nontrivial ancient solution
with the type-I enstrophy bound]
\label{thm:v41-ancient-limit}
There is a subsequence, not relabelled, and a pair \((U,P)\) on
\((-\infty,0)\times\R^3\) such that, for every \(0<\eta<S\) and
\(R>0\),
\[
 U_k\to U\ \text{in }L^2\bigl((-S,-\eta);H^1(B_R)\bigr),
 \qquad
 P_k\rightharpoonup P\ \text{weakly in }L^{3}\bigl((-S,-\eta)\times B_R\bigr).
\]
The limit is a smooth divergence-free solution of the Navier--Stokes
equations on \((-\infty,0)\times\R^3\) with viscosity \(\nu\), and
\begin{align}
 \|\nabla U(s)\|_{L^2(\R^3)}^2&\le C_I^2\nu^{3/2}(-s)^{-1/2},
 \qquad
 \|U(s)\|_{L^6(\R^3)}\le C_6C_I\nu^{3/4}(-s)^{-1/4}
 \qquad(s<0),
 \label{eq:v41-limit-type-I}\\
 \iint_{B_1(0)\times(-1,-\eta_0)}|\nabla U|^2\,\dd y\,\dd s
 &\ge\frac{\varepsilon_1}2,
 \qquad
 \eta_0:=\Bigl(\frac{\varepsilon_1}{4C_I^2\nu^{3/2}}\Bigr)^{2}.
 \label{eq:v41-limit-nontrivial}
\end{align}
In particular \(U\not\equiv0\), and \(\nabla U\in L^\infty((-\infty,-1];
L^2(\R^3))\) with \(\|\nabla U(s)\|_2^2\le C_I^2\nu^{3/2}\) for \(s\le-1\).
\end{theorem}

\begin{proof}
Fix \(\eta,S,R\).  On \(B_R\), Proposition~\ref{prop:v41-uniform-bounds}
bounds \(U_k\) in \(L^2((-S,-\eta);H^2(B_R))\) and \(\partial_sU_k\) in
\(L^2((-S,-\eta);L^2(B_R))\) uniformly in \(k\) (for \(k\) large the torus
contains \(B_R\)), so (E2) gives strong convergence in
\(L^2((-S,-\eta);H^1(B_R))\) along a subsequence; a diagonal argument over
\(\eta\downarrow0\), \(S\uparrow\infty\), \(R\uparrow\infty\) gives a
single subsequence and a limit \(U\).  The pressure is bounded in
\(L^\infty((-S,-\eta);L^3)\), so it converges weakly along a further
subsequence.  Passing to the limit in the weak formulation, in which the
nonlinear term converges because \(U_k\to U\) strongly in
\(L^2_{\rm loc}\) and is bounded in \(L^\infty L^6\), shows that
\((U,P)\) is a distributional solution.  The bounds
\eqref{eq:v41-limit-type-I} follow from
\eqref{eq:v41-rescaled-type-I}--\eqref{eq:v41-rescaled-L6} by weak lower
semicontinuity on every ball, letting \(R\to\infty\).  By (E4) the limit
is smooth.

For \eqref{eq:v41-limit-nontrivial}, split
\eqref{eq:v41-rescaled-nontrivial} at \(s=-\eta_0\):
\(\int_{-\eta_0}^0\|\nabla U_k(s)\|_{L^2(\Torus^3_k)}^2\dd s
\le2C_I^2\nu^{3/2}\eta_0^{1/2}=\varepsilon_1/2\), so
\(\iint_{B_1\times(-1,-\eta_0)}|\nabla U_k|^2\ge\varepsilon_1/2\), and the
strong \(H^1\) convergence on \(B_1\times(-1,-\eta_0)\) passes this to
\(U\).
\end{proof}

\begin{remark}[What the limit inherits]
\label{rem:v41-inheritance}
Every scale-invariant statement proved in V32--V39 for \(u\) holds for
\(U\) with the same constants, since the rescaling is the Navier--Stokes
symmetry and the bounds pass to local limits: the fractional-integral
record law on every ball in the sense of local weak limits, the backward
cone \eqref{eq:v37-backward-cone} for the global enstrophy of \(U\), and
the two-stress anatomy of the flux.  What is \emph{not} inherited is any
bound involving the energy \(E_*\), because
\(\|U_k(s)\|_{L^2(\Torus^3_k)}^2=r_k^{-1}\|u\|_2^2\to\infty\); in
particular the critical norm bound of Theorem~\ref{thm:v33-type-I-log}
becomes \(\|U_k(s)\|_{\dot H^{1/2}}\lesssim\log(1/(r_k^2|s|))\) and does
not pass to the limit.
\end{remark}

% =============================================================================
\section{The finite-energy Liouville theorem}
\label{sec:v41-liouville}
% =============================================================================

\begin{theorem}[Ancient solutions with bounded enstrophy and finite energy
vanish]
\label{thm:v41-liouville}
Let \(V\) be a smooth divergence-free solution of the Navier--Stokes
equations on \((-\infty,s_1)\times\R^3\), \(s_1\le0\), with
\begin{equation}
 \sup_{s\le s_1-1}\|V(s)\|_{L^2(\R^3)}^2=:E_V<\infty,
 \qquad
 \sup_{s\le s_1-1}\|\nabla V(s)\|_{L^2(\R^3)}^2<\infty .
 \label{eq:v41-liouville-hypotheses}
\end{equation}
Then \(V\equiv0\) on \((-\infty,s_1)\times\R^3\).
\end{theorem}

\begin{proof}
On \((-\infty,s_1-1]\), \(V\) is a smooth solution with finite energy and
finite enstrophy, hence a strong solution satisfying the energy equality
\(\frac12\|V(s)\|_2^2+\nu\int_{s'}^s\|\nabla V\|_2^2=\frac12\|V(s')\|_2^2\)
for \(s'<s\le s_1-1\).  Therefore
\(\nu\int_{-\infty}^{s_1-1}\|\nabla V\|_2^2\,\dd s\le E_V/2\), and there is
a sequence \(s_0\to-\infty\) with \(\|\nabla V(s_0)\|_2\to0\).  By the
lattice-free interpolation \(\|V\|_{\dot H^{1/2}}^2\le\|V\|_2\|\nabla V\|_2\)
on \(\R^3\), \(\|V(s_0)\|_{\dot H^{1/2}}\le E_V^{1/4}\|\nabla V(s_0)\|_2^{1/2}
\to0\).  Fix \(s_0\) with \(\|V(s_0)\|_{\dot H^{1/2}}\le\varepsilon_{\rm FK}\)
and with \(2C_2\varepsilon_{\rm FK}\le\nu/2\), where \(C_2\) is the constant
in \(|\langle(V\cdot\nabla)V,\Delta V\rangle|\le C_2\|V\|_{\dot H^{1/2}}
\|\Delta V\|_2^2\), which follows from H\"older with exponents \(3,6,2\),
\(\|V\|_3\le C\|V\|_{\dot H^{1/2}}\), and \(\|\nabla V\|_6\le C\|\Delta V\|_2\).

On \([s_0,s_1)\), \(V\) is a smooth solution with
\(V\in C([s_0,s_1');\dot H^{1/2})\cap L^2(s_0,s_1';\dot H^{3/2})\) for every
\(s_1'<s_1\), because \(\|V\|_{\dot H^{3/2}}^2\le\|\nabla V\|_2\|\Delta V\|_2\)
and \(\Delta V\in L^2_{\rm loc}L^2\) by the \(H^1\) balance.  By the
uniqueness in (E3), \(V\) coincides on \([s_0,s_1)\) with the Fujita--Kato
solution from \(V(s_0)\), so \(\|V(s)\|_{\dot H^{1/2}}\le2\varepsilon_{\rm FK}\)
there.  The \(H^1\) balance then reads
\(\frac12Y'+\nu Z\le C_2\cdot2\varepsilon_{\rm FK}Z\le\frac\nu2Z\), with
\(Y=\|\nabla V\|_2^2\), \(Z=\|\Delta V\|_2^2\), so \(Y'\le-\nu Z\).  Since
\(Y\le\|V\|_2\|\Delta V\|_2\le E_V^{1/2}Z^{1/2}\), one has
\(Z\ge Y^2/E_V\) and hence \(Y'\le-\nu Y^2/E_V\), which integrates to
\begin{equation}
 \|\nabla V(s)\|_2^2\le\frac{E_V}{\nu(s-s_0)}\qquad(s_0<s<s_1).
 \label{eq:v41-small-data-decay}
\end{equation}
Letting \(s_0\to-\infty\) along the sequence shows
\(\nabla V(s)=0\) for every \(s<s_1\); a field with zero gradient is
constant, and a constant in \(L^2(\R^3)\) is zero (the energy on
\([s_0,s_1)\) is finite by the energy equality from \(s_0\)).
\end{proof}

\begin{corollary}[The ancient limit of a type-I singularity has infinite
energy]
\label{cor:v41-infinite-energy}
The limit \(U\) of Theorem~\ref{thm:v41-ancient-limit} satisfies
\begin{equation}
 \sup_{s\le-1}\|U(s)\|_{L^2(\R^3)}=\infty .
 \label{eq:v41-infinite-energy}
\end{equation}
Equivalently, the singular point \(x_0\) has no finite-energy blow-up
profile: there is no constant \(A\) such that
\(r_k^{-1}\int_{B_{Rr_k}(x_0)}|u(T_*+r_k^2s)|^2\,\dd x\le A\) for all
\(R>0\), all \(s\le-1\), and all large \(k\).
\end{corollary}

\begin{proof}
\(U\) has \(\|\nabla U(s)\|_2^2\le C_I^2\nu^{3/2}\) for \(s\le-1\) and is
nontrivial on \(B_1\times(-1,-\eta_0)\), so
Theorem~\ref{thm:v41-liouville} with \(s_1=0\) excludes
\(\sup_{s\le-1}\|U(s)\|_2<\infty\).  If the displayed uniform bound held,
then for every \(s\le-1\) and \(R\), weak lower semicontinuity on
\(B_R\) would give \(\|U(s)\|_{L^2(B_R)}^2\le A\), hence
\(\|U(s)\|_{L^2(\R^3)}^2\le A\), contradicting
\eqref{eq:v41-infinite-energy}.
\end{proof}

\begin{assessment}[Comparison with V36 and with the global picture]
\label{ass:v41-comparison}
Proposition~\ref{prop:v36-local-energies} bounds the scaled local energy
\(r^{-1}\int_{B_r}|u|^2\) at every point by a square of the V33 logarithm.
Corollary~\ref{cor:v41-infinite-energy} shows that at a singular point,
under a type-I rate, this scaled local energy cannot stay bounded
uniformly in the rescaled radius: energy must be present at all rescaled
distances from the singular point, with total divergent in the limit.  The
type-I singular profile is therefore squeezed between two logarithmic
facts, unbounded scaled local energy from below and \(\log^2\) from
above.  The mechanism used, small-data decay of an ancient finite-energy
solution, is not scale-invariant in the sense of
Theorem~\ref{thm:v39-scale-no-go}: it uses the energy as an absolute
quantity in the far past, which is exactly the non-invariant input
Assessment~\ref{ass:v40-AC3} identified.
\end{assessment}

% =============================================================================
\section{The exact remaining type-I gate}
\label{sec:v41-gate}
% =============================================================================

\begin{definition}[Bounded-enstrophy ancient class]
\label{def:v41-class}
Let \(\mathcal A_\nu(C)\) be the set of smooth divergence-free solutions
\(U\) of the Navier--Stokes equations on \((-\infty,0)\times\R^3\) with
\(\|\nabla U(s)\|_{L^2(\R^3)}^2\le C(-s)^{-1/2}\) for all \(s<0\).
\end{definition}

\begin{theorem}[Type-I exclusion is equivalent to a Liouville theorem on
\(\mathcal A_\nu\)]
\label{thm:v41-gate}
If every element of \(\mathcal A_\nu(C_I^2\nu^{3/2})\) vanishes
identically, then no smooth solution on \([0,T_*)\times\Torus^3\) with
maximal time \(T_*<\infty\) satisfies the type-I enstrophy rate
\eqref{eq:v32-type-I}.  Conversely, every nonzero element of
\(\mathcal A_\nu(C)\) is itself a solution with a type-I enstrophy rate
on \((-\infty,0)\) whose maximal time is at most \(0\) or which continues
as a global type-I-bounded solution; the finite-energy elements of
\(\mathcal A_\nu(C)\) are zero.
\end{theorem}

\begin{proof}
The first statement is Theorem~\ref{thm:v41-ancient-limit}.  The second is
the definition together with Theorem~\ref{thm:v41-liouville}.
\end{proof}

\begin{firewall}[Status of the gate]
\label{fw:v41-gate-status}
Theorem~\ref{thm:v41-gate} does not prove the Liouville statement for
infinite-energy elements of \(\mathcal A_\nu\).  Nonzero ancient solutions
with bounded enstrophy on \((-\infty,-1]\) but infinite energy are not
excluded by anything in this series; the class contains, for instance, any
nonzero stationary solution of the Navier--Stokes equations on \(\R^3\)
with finite Dirichlet integral and infinite energy, should one exist with
the decay implied by \(L^6\) (whether such solutions exist is itself an
open question closely related to the Liouville problem for the stationary
equations).  The type-I regime is therefore reduced, within this series,
to that Liouville problem; the reduction is rigorous and its imports are
listed in Firewall~\ref{fw:v41-imports}.
\end{firewall}

% =============================================================================
\section{Integrated V41 conclusion and next acquisition order}
\label{sec:v41-conclusion}
% =============================================================================

\begin{theorem}[What V41 proves]
\label{thm:v41-conclusion}
Relative to V1--V40, and modulo the imports (E1)--(E4):
\begin{enumerate}[label=\textup{(V41.\arabic*)},leftmargin=6.2em]
\item a singular time has a dissipation-concentration point
      (Lemma~\ref{lem:v41-concentration-point});
\item the rescaled sequence at that point has the uniform bounds of
      Proposition~\ref{prop:v41-uniform-bounds} and converges to a
      nontrivial smooth ancient solution with the type-I enstrophy bound
      (Theorem~\ref{thm:v41-ancient-limit});
\item ancient solutions with bounded enstrophy and finite energy vanish
      (Theorem~\ref{thm:v41-liouville}), so the ancient limit has infinite
      energy and the singular point has no finite-energy profile
      (Corollary~\ref{cor:v41-infinite-energy});
\item type-I exclusion is equivalent to a Liouville theorem on the class
      \(\mathcal A_\nu\) (Theorem~\ref{thm:v41-gate}).
\end{enumerate}
The full Navier--Stokes stability/global-regularity theorem remains
unproved.
\end{theorem}

\begin{proof}
The cited results.
\end{proof}

\begin{programme}[V42 acquisition order]
\label{prog:v42}
\begin{enumerate}[label=\textup{(AC\arabic*)},leftmargin=4.8em]
\item Infinite-energy ancient solutions.  For \(U\in\mathcal A_\nu(C)\)
      derive the local energy inequality on balls \(B_R\) with the
      pressure bounded in \(L^3\), and test whether the local energy
      \(\int_{B_R}|U(s)|^2\) can grow in \(R\) slower than \(R\) (which
      \(L^6\) permits up to \(R^{2}\)) while the enstrophy stays bounded;
      determine the exact growth class in which the Fujita--Kato argument
      of Theorem~\ref{thm:v41-liouville} still runs after localization.
\item Backward cone for \(U\).  The cone \eqref{eq:v37-backward-cone}
      applies to \(Y_U(s)=\|\nabla U(s)\|_2^2\) and gives, for \(s\le-1\),
      the two-sided bound \(c(-s)^{-1/2}\le Y_U(s)\le C(-s)^{-1/2}\) only
      if \(Y_U\) is unbounded as \(s\uparrow0\); determine whether
      \(U\) is regular at \(s=0\) (then \(U\) is a global solution with
      uniformly bounded enstrophy on \((-\infty,\infty)\)) or singular
      (then \(U\) is a self-similar-type type-I singularity of its own).
\item Type-II.  Item (AC2) of Programme~\ref{prog:v41} remains.
\item The next compulsory companion audit is V45.
\end{enumerate}
\end{programme}

\begin{assessment}[Position after V41]
\label{ass:v41-position}
The type-I regime has been reduced from a logarithm to a Liouville
problem, and the finite-energy half of that problem is proved.  The
reduction uses the energy as an absolute quantity in the far past of the
rescaled solution, which is the first non-scale-invariant input in the
series.  No full stability or global-regularity claim is made.
\end{assessment}

% =============================================================================
\section*{Version V41 append-only record}
\addcontentsline{toc}{section}{Version V41 append-only record}
% =============================================================================

The complete V40 mathematical source was retained before this continuation.
Its SHA--256 digest was
\begin{center}\scriptsize\ttfamily
f5db6ca77c1dfaabb890c2ec3a229cb39fa304c96fab2a9477429850e92952ca.
\end{center}
V41 rescaled a type-I singularity at a concentration point to a nontrivial
ancient solution with the type-I enstrophy bound, proved that ancient
solutions with bounded enstrophy and finite energy vanish, concluded that
the ancient limit has infinite energy, and reduced type-I exclusion to a
Liouville theorem on the bounded-enstrophy ancient class.  No full
regularity claim is made.

% =============================================================================
\section{V42: structure of the bounded-enstrophy ancient class}
\label{sec:v42-entry}
% =============================================================================

V41 reduced the type-I regime to a Liouville problem on the class
\(\mathcal A_\nu(C)\) of ancient solutions on \(\R^3\) with
\(\|\nabla U(s)\|_2^2\le C(-s)^{-1/2}\), and proved the finite-energy
case.  This version describes the class.  Every nonzero element is either
eternal, extending to a smooth solution on all of \(\R\times\R^3\) with
enstrophy bounded on every half-line, or singular at \(s=0\) with the
two-sided rate \(c(-s)^{-1/2}\le\|\nabla U(s)\|_2^2\le C(-s)^{-1/2}\)
forced by the backward cone.  Uniform palinstrophy averages, a localized
energy inequality with explicit growth rates in the ball radius, and a
normalization of the ancient limit at a point of maximal unit-scale
dissipation are proved.  The stationary solutions with finite Dirichlet
integral are identified as a subclass of the eternal case, so the type-I
gate of V41 contains Leray's stationary Liouville problem; the known
partial answer to that problem, which requires \(L^{9/2}\) decay, is
recorded as an external fact and not used.  No global regularity claim is
made.  The next compulsory companion audit is V45.

% =============================================================================
\section{The trichotomy}
\label{sec:v42-trichotomy}
% =============================================================================

Throughout, \(U\in\mathcal A_\nu(C)\) as in Definition~\ref{def:v41-class},
\(Y_U(s):=\|\nabla U(s)\|_{L^2(\R^3)}^2\), \(Z_U(s):=\|\Delta U(s)\|_2^2\).
All inequalities of Lemma~\ref{lem:v32-leray-rate} and
Lemma~\ref{lem:v37-backward-cone} hold on \(\R^3\) with the same
constants, because they use only Sobolev and interpolation inequalities
that are valid on \(\R^3\) for fields in \(L^6\cap\dot H^1\cap\dot H^2\).

\begin{theorem}[Trichotomy for the ancient class]
\label{thm:v42-trichotomy}
Let \(U\in\mathcal A_\nu(C)\) be nonzero.  Then
\(\sup_{s\le-1}\|U(s)\|_{L^2(\R^3)}=\infty\), and exactly one of the
following holds.
\begin{enumerate}[label=\textup{(\Alph*)},leftmargin=3.4em]
\item \textbf{Eternal.}  \(\limsup_{s\uparrow0}Y_U(s)<\infty\).  Then
      \(U\) extends to a smooth solution on \(\R\times\R^3\), and for every
      \(S\in\R\), \(\sup_{s\le S}Y_U(s)<\infty\).
\item \textbf{Two-sided type-I singular.}  \(Y_U(s)\to\infty\) as
      \(s\uparrow0\).  Then
      \begin{equation}
       \Bigl(\frac{\nu^3}{8C_1(-s)}\Bigr)^{1/2}
       \le Y_U(s)\le\frac{C}{(-s)^{1/2}}
       \qquad(s<0),
       \label{eq:v42-two-sided}
      \end{equation}
      and \(U\) cannot be continued as a strong solution past \(s=0\).
\end{enumerate}
\end{theorem}

\begin{proof}
Infinite energy is Theorem~\ref{thm:v41-liouville}.  If
\(\limsup_{s\uparrow0}Y_U<\infty\), then \(Y_U\) is bounded on
\([-1,0)\) along a sequence, and since \(Y_U'\le2C_1\nu^{-3}Y_U^3\) on
\(\R^3\) as in Lemma~\ref{lem:v32-leray-rate}, \(Y_U\) is bounded on all of
\([-1,0)\); the \(H^1\) local existence theorem on \(\R^3\) for data in
\(L^6\cap\dot H^1\) (whose existence time depends only on \(\nu\) and
\(\|\nabla U\|_2\)) continues \(U\) past \(0\), and repeating the argument
gives a smooth solution on \(\R\times\R^3\) with \(Y_U\) bounded on every
half-line, which is (A).  Otherwise \(Y_U(s_j)\to\infty\) along some
\(s_j\uparrow0\); the backward cone \eqref{eq:v37-backward-cone} with
\(t_2=s_j\), \(t=s\), and \(Y_U(s_j)\to\infty\) gives
\(Y_U(s)^{-2}\le4C_1\nu^{-3}(-s)\), the lower bound in
\eqref{eq:v42-two-sided}; the upper bound is the definition of
\(\mathcal A_\nu(C)\).  Finally \(Y_U(s)\to\infty\) excludes strong
continuation.
\end{proof}

\begin{proposition}[Uniform palinstrophy averages]
\label{prop:v42-palinstrophy}
For \(U\in\mathcal A_\nu(C)\) and every \(s\le-1\),
\begin{equation}
 \nu\int_{s-1}^{s}Z_U(\sigma)\,\dd\sigma
 \le Y_U(s-1)+2C_1\nu^{-3}\int_{s-1}^sY_U^3
 \le C+2C_1\nu^{-3}C^3=:C_Z ,
 \label{eq:v42-palinstrophy}
\end{equation}
and hence, by Agmon's inequality
\(\|U\|_\infty^2\le C_A\|\nabla U\|_2\|\Delta U\|_2\),
\begin{equation}
 \int_{s-1}^s\|U(\sigma)\|_{L^\infty(\R^3)}^4\,\dd\sigma
 \le C_A^2\,C\,\frac{C_Z}{\nu}.
 \label{eq:v42-Linfty-average}
\end{equation}
Thus every element of the class lies in \(L^4_{\rm loc}L^\infty\) on
\((-\infty,-1]\) with a uniform bound per unit time, and in particular in
the Serrin class \(L^4_tL^\infty_x\) there.
\end{proposition}

\begin{proof}
Integrate \(\frac12Y_U'+\frac\nu2Z_U\le C_1\nu^{-3}Y_U^3\) over
\([s-1,s]\) and drop \(Y_U(s)\ge0\).  Agmon on \(\R^3\) for
\(L^6\cap\dot H^2\) fields and Cauchy--Schwarz give the second display.
\end{proof}

% =============================================================================
\section{Localized energy and the growth of local energy}
\label{sec:v42-local-energy}
% =============================================================================

Let \(\phi_R(y)=\phi(y/R)\) with \(\phi\in C_c^\infty(B_2)\),
\(0\le\phi\le1\), \(\phi=1\) on \(B_1\), and write
\(e_R(s):=\frac12\int|U(s)|^2\phi_R\).  Let \(P\) be the pressure of
\(U\), which by the \(L^6\) bound and the Riesz transforms satisfies
\(\|P(s)\|_{L^3(\R^3)}\le C_{\rm CZ}\|U(s)\|_6^2\).

\begin{proposition}[Localized energy inequality on the ancient class]
\label{prop:v42-local-energy}
For \(U\in\mathcal A_\nu(C)\), \(s'<s\le-1\), and \(R\ge1\),
\begin{align}
 e_R(s)+\nu\int_{s'}^{s}\int|\nabla U|^2\phi_R
 &\le e_R(s')+(s-s')\bigl(C_\Delta+C_{\rm fl}R^{1/2}\bigr),
 \label{eq:v42-local-energy-inequality}\\
 e_R(s)&\le\tfrac12|B_{2R}|^{2/3}\|U(s)\|_6^2\le C_{\rm loc}R^2 ,
 \label{eq:v42-local-energy-size}
\end{align}
with \(C_\Delta,C_{\rm fl},C_{\rm loc}\) depending only on
\(\nu,C,\phi\).  Consequently the time-averaged local enstrophy obeys
\begin{equation}
 \frac1L\int_{s-L}^s\int_{B_R}|\nabla U|^2
 \le\frac{C_{\rm loc}R^2}{\nu L}+\frac{C_\Delta+C_{\rm fl}R^{1/2}}{\nu}
 \qquad(L>0),
 \label{eq:v42-averaged-local-enstrophy}
\end{equation}
which for \(L\ge R^{3/2}\) is at most a constant times \(1+R^{1/2}\).
\end{proposition}

\begin{proof}
Multiply the equation by \(U\phi_R\) and integrate:
\(e_R'+\nu\int|\nabla U|^2\phi_R=\frac\nu2\int|U|^2\Delta\phi_R
+\frac12\int\bigl(|U|^2+2P\bigr)U\cdot\nabla\phi_R\).  With
\(|\Delta\phi_R|\le CR^{-2}\mathbf 1_{B_{2R}}\) and
\(|\nabla\phi_R|\le CR^{-1}\mathbf 1_{B_{2R}}\), H\"older gives
\(\int|U|^2|\Delta\phi_R|\le CR^{-2}|B_{2R}|^{2/3}\|U\|_6^2\le C'\),
\(\int|U|^3|\nabla\phi_R|\le CR^{-1}|B_{2R}|^{1/2}\|U\|_6^3\le C'R^{1/2}\),
and \(\int|P||U||\nabla\phi_R|\le CR^{-1}\|P\|_3\|U\|_{L^{3/2}(B_{2R})}
\le CR^{-1}\|P\|_3|B_{2R}|^{1/2}\|U\|_6\le C'R^{1/2}\); the \(L^6\) norm is
bounded by \(C_6C^{1/2}\) for \(s\le-1\).  Integrate in time.  The size
bound is H\"older on \(B_{2R}\).  Divide the inequality with
\(s'=s-L\) by \(L\) and use \eqref{eq:v42-local-energy-size} at \(s-L\).
\end{proof}

\begin{assessment}[Where the finite-energy argument breaks and what
replaces it]
\label{ass:v42-local-growth}
The proof of Theorem~\ref{thm:v41-liouville} used the global energy
equality to find far-past times of small enstrophy.  On the ancient class
the local energy may grow like \(R^2\) in the ball radius, and the flux
through the boundary of \(B_R\) is of order \(R^{1/2}\) per unit time; the
localized inequality \eqref{eq:v42-local-energy-inequality} therefore
controls only time averages, \eqref{eq:v42-averaged-local-enstrophy}, and
those are no better than the global bound \(Y_U\le C\) for
\(R\lesssim1\).  A Liouville theorem on \(\mathcal A_\nu\) cannot be
obtained by localizing V41; it needs either a decay hypothesis at spatial
infinity or a different mechanism.
\end{assessment}

% =============================================================================
\section{Normalization of the ancient limit}
\label{sec:v42-normalization}
% =============================================================================

\begin{proposition}[The ancient limit can be normalized at a point of
maximal unit-scale dissipation]
\label{prop:v42-normalization}
In Theorem~\ref{thm:v41-ancient-limit} the centres may be chosen as
\(x_k\in\Torus^3\) maximizing
\(x\mapsto r_k^{-1}\iint_{Q_{r_k}(x,T_*)}|\nabla u|^2\), with the same
radii \(r_k\).  The resulting limit \(U\in\mathcal A_\nu(C_I^2\nu^{3/2})\)
then satisfies, in addition to \eqref{eq:v41-limit-nontrivial},
\begin{equation}
 \sup_{y\in\R^3}\iint_{B_1(y)\times(-1,0)}|\nabla U|^2
 \le2C_I^2\nu^{3/2},
 \qquad
 \iint_{B_1(0)\times(-1,-\eta_0)}|\nabla U|^2\ge\frac{\varepsilon_1}{2}.
 \label{eq:v42-normalized}
\end{equation}
\end{proposition}

\begin{proof}
The map \(x\mapsto r_k^{-1}\iint_{Q_{r_k}(x,T_*)}|\nabla u|^2\) is
continuous on the compact torus, so a maximizer \(x_k\) exists, and its
value is at least the value at \(x_0\), which exceeds \(\varepsilon_1\).
For every \(x\), \(r_k^{-1}\iint_{Q_{r_k}(x,T_*)}|\nabla u|^2
\le r_k^{-1}\int_{T_*-r_k^2}^{T_*}Y\,\dd t\le r_k^{-1}\cdot2C_I^2\nu^{3/2}r_k\)
by \eqref{eq:v32-type-I}.  After rescaling at \(x_k\), the unit-scale
dissipation of \(U_k\) is bounded by \(2C_I^2\nu^{3/2}\) at every point
and exceeds \(\varepsilon_1\) at the origin; the proof of
Theorem~\ref{thm:v41-ancient-limit} applies verbatim (the centres do not
enter the uniform bounds), and weak lower semicontinuity on each
\(B_1(y)\times(-1,-\eta)\), followed by \(\eta\downarrow0\) with the
uniform tail bound \(\int_{-\eta}^0Y_k\le2C_I^2\nu^{3/2}\eta^{1/2}\), gives
the first inequality of \eqref{eq:v42-normalized}.
\end{proof}

% =============================================================================
\section{The stationary subclass}
\label{sec:v42-stationary}
% =============================================================================

\begin{proposition}[Finite-Dirichlet stationary solutions lie in the
eternal class]
\label{prop:v42-stationary}
Let \(V\) be a smooth divergence-free stationary solution of the
Navier--Stokes equations on \(\R^3\) with \(\int|\nabla V|^2<\infty\) and
\(V\in L^6(\R^3)\).  Then \(U(s,y):=V(y)\) belongs to
\(\mathcal A_\nu(C)\) for \(C=\|\nabla V\|_2^2\) and is of type (A) in
Theorem~\ref{thm:v42-trichotomy}.  Consequently, if the Liouville statement
of Theorem~\ref{thm:v41-gate} holds, then every such \(V\) is zero.
\end{proposition}

\begin{proof}
\(Y_U\equiv\|\nabla V\|_2^2\le C\le C(-s)^{-1/2}\) for \(s\le-1\), and
for \(-1<s<0\) the bound holds because \((-s)^{-1/2}\ge1\).  The solution
is eternal and \(Y_U\) is bounded on \(\R\).
\end{proof}

\begin{firewall}[The gate contains Leray's stationary Liouville problem]
\label{fw:v42-Leray-problem}
Whether a stationary solution with finite Dirichlet integral and
\(V\to0\) at infinity must vanish is a classical open problem of Leray.
It is known, as an external fact not used here, that the answer is yes
under the additional integrability \(V\in L^{9/2}(\R^3)\), and \(L^6\)
alone is not known to suffice.  Since Proposition~\ref{prop:v42-stationary}
places these solutions inside \(\mathcal A_\nu\), any proof of the type-I
Liouville statement of Theorem~\ref{thm:v41-gate} would in particular
settle Leray's problem for solutions in \(L^6\cap\dot H^1\).  This series
does not claim to do so.  The honest position is therefore: the type-I
regime of a hypothetical Navier--Stokes singularity is at least as hard as
Leray's stationary Liouville problem, and V41--V42 prove that it is
exactly a Liouville problem on the slightly larger class
\(\mathcal A_\nu\), in which finite-energy elements have been removed.
\end{firewall}

% =============================================================================
\section{Integrated V42 conclusion and next acquisition order}
\label{sec:v42-conclusion}
% =============================================================================

\begin{theorem}[What V42 proves]
\label{thm:v42-conclusion}
Relative to V1--V41:
\begin{enumerate}[label=\textup{(V42.\arabic*)},leftmargin=6.2em]
\item the trichotomy for \(\mathcal A_\nu(C)\): zero if finite energy,
      eternal with enstrophy bounded on every half-line, or two-sided
      type-I singular with \eqref{eq:v42-two-sided};
\item uniform palinstrophy averages and the \(L^4_{\rm loc}L^\infty\)
      bound \eqref{eq:v42-Linfty-average} on the class;
\item the localized energy inequality with local energy growth at most
      \(R^2\) and boundary flux at most \(R^{1/2}\) per unit time, and the
      time-averaged local enstrophy bound
      \eqref{eq:v42-averaged-local-enstrophy};
\item the normalization \eqref{eq:v42-normalized} of the ancient limit at
      a point of maximal unit-scale dissipation;
\item finite-Dirichlet stationary solutions in \(L^6\) belong to the
      eternal class, so the type-I gate contains Leray's stationary
      Liouville problem.
\end{enumerate}
The full Navier--Stokes stability/global-regularity theorem remains
unproved.
\end{theorem}

\begin{proof}
Theorem~\ref{thm:v42-trichotomy},
Propositions~\ref{prop:v42-palinstrophy}, \ref{prop:v42-local-energy},
\ref{prop:v42-normalization}, and \ref{prop:v42-stationary}.
\end{proof}

\begin{programme}[V43 acquisition order]
\label{prog:v43}
\begin{enumerate}[label=\textup{(AC\arabic*)},leftmargin=4.8em]
\item Decay at infinity for the ancient limit.  Determine whether the
      construction of Theorem~\ref{thm:v41-ancient-limit} can be
      supplemented, using the uniform tail smallness
      \eqref{eq:v32-uniform-tail-smallness} of the original solution at
      frequencies below the parabolic scale, by a bound on
      \(\int_{B_R}|U(s)|^2\) growing slower than \(R^2\); the V34 flat
      profile is the candidate input, since it bounds the energy of \(u\)
      at scales \(\rho\ll K_{\rm par}\) by \(Q/\rho\).
\item The singular case (B).  Apply the V32--V34 machinery to \(U\) on
      \(\R^3\) in the form of local Fourier cutoffs on balls, and test
      whether the two-sided rate \eqref{eq:v42-two-sided} together with
      the backward cone forces the enstrophy of \(U\) to concentrate at a
      single point as \(s\uparrow0\).
\item Type-II: item (AC2) of Programme~\ref{prog:v41} remains.
\item The next compulsory companion audit is V45.
\end{enumerate}
\end{programme}

\begin{assessment}[Position after V42]
\label{ass:v42-position}
The type-I gate is now a precisely described Liouville problem with three
subcases, one of them closed.  The stationary subcase is a classical open
problem, which fixes the level of difficulty honestly.  The remaining
opportunity specific to this series is that the ancient limit is not an
arbitrary element of the class: it comes with the flat critical-mass
profile and uniform tail smallness of the original solution, which may
translate into decay at spatial infinity for \(U\).  No full stability or
global-regularity claim is made.
\end{assessment}

% =============================================================================
\section*{Version V42 append-only record}
\addcontentsline{toc}{section}{Version V42 append-only record}
% =============================================================================

The complete V41 mathematical source was retained before this continuation.
Its SHA--256 digest was
\begin{center}\scriptsize\ttfamily
7fa4c4a1a2db94209ddd66281993749e56d2509de4d6eb062f33c873824c8fa6.
\end{center}
V42 proved the trichotomy for the bounded-enstrophy ancient class, uniform
palinstrophy and \(L^4L^\infty\) bounds, the localized energy inequality
with its growth rates, the normalization of the ancient limit, and the
inclusion of Leray's stationary Liouville problem in the type-I gate.  No
full regularity claim is made.

% =============================================================================
\section{V43: dissipation concentration and the dimension of the type-I
singular set}
\label{sec:v43-entry}
% =============================================================================

Item (AC1) of Programme~\ref{prog:v43}, decay at infinity for the ancient
limit, does not follow from the flat profile: the local energy of the
rescaled fields on a ball of rescaled radius \(R\) is bounded only by
\(R\log^2(1/r_k^2)\), and the logarithm diverges.  This is recorded and the
version turns to what the concentration structure of V41 gives in physical
space.  Under a type-I enstrophy rate the total dissipation in the slab
\((T_*-r^2,T_*)\) is two-sided, of order \(r\) exactly, by the Leray rate
and the type-I bound.  At a concentration scale the parabolic cylinder of
the singular point captures a fixed fraction of it.  Two consequences
follow with short proofs: the type-I constant is bounded below by the
Caffarelli--Kohn--Nirenberg threshold, and at any scale at most
\(N_*=4C_I^2\nu^{3/2}/\varepsilon_1\) well-separated points can be active,
so that the singular set at the singular time has Hausdorff dimension zero.
No global regularity claim is made.  The next compulsory companion audit is
V45.

% =============================================================================
\section{Why the flat profile does not give decay of the ancient limit}
\label{sec:v43-no-decay}
% =============================================================================

\begin{proposition}[Local energy of the rescaled fields]
\label{prop:v43-local-energy-rescaled}
Under \eqref{eq:v32-type-I}, for the rescaled fields
\eqref{eq:v41-rescaled}, every \(R>0\) and \(s\le-1\) with
\(T_*+r_k^2s\ge T_*-\delta/2\),
\begin{equation}
 \int_{B_R}|U_k(s)|^2\,\dd y
 =\frac1{r_k}\int_{B_{Rr_k}(x_0)}|u(T_*+r_k^2s)|^2\,\dd x
 \le\Bigl(\frac{4\pi}3\Bigr)^{1/3}C_3^2\,R\,M_\delta(T_*+r_k^2s)^2,
 \label{eq:v43-rescaled-local-energy}
\end{equation}
with \(M_\delta\lesssim\log(9\delta/(r_k^2|s|))\to\infty\).  The bound is
linear in \(R\) but not uniform in \(k\); no uniform bound of the form
\(\int_{B_R}|U_k(s)|^2\le AR^\gamma\) with \(\gamma<2\) follows from
V32--V34, and \(\gamma=2\) is what \(L^6\) already gives.
\end{proposition}

\begin{proof}
Change variables and apply \eqref{eq:v36-local-L2} with \(r=Rr_k\).
\end{proof}

\begin{firewall}[AC1 of Programme~\ref{prog:v43} is not achieved]
\label{fw:v43-AC1}
The flat profile bounds the Fourier-side energy of \(u\) at frequency
\(\rho\) by \(Q/\rho\) with \(Q\) logarithmic; in rescaled variables this
is \(Q_k/\kappa\) with \(Q_k\to\infty\), and Bernstein-type localization
of low frequencies to balls loses further powers of the radius.  Decay at
spatial infinity for the ancient limit is therefore not available from the
present bounds; Assessment~\ref{ass:v42-local-growth} stands.
\end{firewall}

% =============================================================================
\section{Two-sided slab dissipation and the concentration fraction}
\label{sec:v43-slab}
% =============================================================================

\begin{lemma}[The slab dissipation is two-sided of order \(r\)]
\label{lem:v43-slab}
Under \eqref{eq:v32-type-I}, for \(0<r\le\delta^{1/2}\),
\begin{equation}
 2c_L^2\nu^{3/2}\,r
 \le\int_{T_*-r^2}^{T_*}\|\Lambda u(t)\|_2^2\,\dd t
 \le2C_I^2\nu^{3/2}\,r .
 \label{eq:v43-slab}
\end{equation}
\end{lemma}

\begin{proof}
Integrate the Leray lower bound \eqref{eq:v32-leray-rate} and the type-I
upper bound over \((T_*-r^2,T_*)\).
\end{proof}

For \(x\in\Torus^3\) and \(r>0\) put
\begin{equation}
 f_x(r):=\frac1r\iint_{Q_r(x,T_*)}|\nabla u|^2\,\dd x'\,\dd t,
 \qquad
 S_r:=\{x\in\Torus^3:\ f_x(r)\ge\varepsilon_1/2\},
 \label{eq:v43-active-set}
\end{equation}
with \(\varepsilon_1\) the CKN threshold of Firewall~\ref{fw:v41-imports}.

\begin{proposition}[Concentration fraction and the type-I constant]
\label{prop:v43-fraction}
Under \eqref{eq:v32-type-I}, every \(x\) and \(r\le\delta^{1/2}\) satisfy
\(f_x(r)\le2C_I^2\nu^{3/2}\).  At a concentration point \(x_0\) with
scales \(r_k\) as in Lemma~\ref{lem:v41-concentration-point}, the cylinder
\(Q_{r_k}(x_0,T_*)\) carries a fraction at least
\(c_*:=\varepsilon_1/(2C_I^2\nu^{3/2})\) of the total dissipation of its
slab.  Consequently
\begin{equation}
 \boxed{\quad
 C_I^2\ \ge\ \frac{\varepsilon_1}{2\nu^{3/2}}
 \quad}
 \label{eq:v43-type-I-lower}
\end{equation}
for every singular solution with a type-I enstrophy rate: type-I constants
below the CKN threshold are impossible.
\end{proposition}

\begin{proof}
\(f_x(r)\le r^{-1}\int_{T_*-r^2}^{T_*}\|\Lambda u\|_2^2\le2C_I^2\nu^{3/2}\)
by \eqref{eq:v43-slab}.  At \(x_0\), \(f_{x_0}(r_k)>\varepsilon_1\) by
\eqref{eq:v41-concentration}, so the fraction is at least
\(\varepsilon_1r_k/(2C_I^2\nu^{3/2}r_k)=c_*\), and \(c_*\le1\) is
\eqref{eq:v43-type-I-lower}.
\end{proof}

\begin{lemma}[Active points at one scale are few]
\label{lem:v43-few-active}
Under \eqref{eq:v32-type-I}, for \(r\le\delta^{1/2}\), any subset of
\(S_r\) whose points are pairwise at distance at least \(2r\) has at most
\begin{equation}
 N_*:=\frac{4C_I^2\nu^{3/2}}{\varepsilon_1}
 \label{eq:v43-N-star}
\end{equation}
elements.  Hence \(S_r\) is covered by at most \(N_*\) balls of radius
\(2r\).  Moreover, each singular point \(x\) (in the sense of
Firewall~\ref{fw:v41-imports}(E1)) belongs to \(S_{2^{-m}}\) for
infinitely many integers \(m\).
\end{lemma}

\begin{proof}
For \(2r\)-separated points the cylinders \(Q_r(x_i,T_*)\) are pairwise
disjoint and lie in the slab \((T_*-r^2,T_*)\), so
\(\sum_i(\varepsilon_1/2)r\le\sum_i\iint_{Q_r(x_i)}|\nabla u|^2
\le2C_I^2\nu^{3/2}r\), which bounds the count by \(N_*\).  A maximal
\(2r\)-separated subset of \(S_r\) has at most \(N_*\) points, and the
balls of radius \(2r\) about them cover \(S_r\) by maximality.  For a
singular point \(x\), \(\limsup_{r\to0}f_x(r)>\varepsilon_1\) gives
radii \(r_k\downarrow0\) with \(f_x(r_k)>\varepsilon_1\); for
\(r\in[r_k,2r_k]\), \(Q_{r_k}(x)\subset Q_r(x)\) gives
\(f_x(r)\ge(r_k/r)f_x(r_k)>\varepsilon_1/2\), and each interval
\([r_k,2r_k]\) contains a dyadic \(2^{-m}\).
\end{proof}

% =============================================================================
\section{The type-I singular set has dimension zero}
\label{sec:v43-dimension}
% =============================================================================

\begin{theorem}[Dimension of the singular set under a type-I enstrophy rate]
\label{thm:v43-dimension-zero}
Let \(u\) be a smooth solution on \([0,T_*)\times\Torus^3\) with maximal
time \(T_*<\infty\) satisfying \eqref{eq:v32-type-I}, and let
\(\Sigma\subset\Torus^3\) be the set of \(x\) such that \(u\) is unbounded
in every neighbourhood of \((x,T_*)\).  Then \(\Sigma\ne\emptyset\), and
for every \(\alpha>0\)
\begin{equation}
 \boxed{\quad
 \mathcal H^\alpha(\Sigma)=0 ,
 \quad}
 \label{eq:v43-dimension-zero}
\end{equation}
so \(\Sigma\) has Hausdorff dimension zero.  Moreover, for every
\(m_0\), \(\Sigma\subset\bigcup_{m\ge m_0}S_{2^{-m}}\), where each
\(S_{2^{-m}}\) is covered by \(N_*\) balls of radius \(2^{-m+1}\).
\end{theorem}

\begin{proof}
\(\Sigma\ne\emptyset\) is Lemma~\ref{lem:v41-concentration-point}.  By
(E1), every \(x\in\Sigma\) has \(\limsup_{r\to0}f_x(r)>\varepsilon_1\),
hence by Lemma~\ref{lem:v43-few-active} lies in \(S_{2^{-m}}\) for
infinitely many \(m\), which gives the inclusion.  For \(\alpha>0\) and
\(m_0\) with \(2^{-m_0}\le\delta^{1/2}\), the covering of each
\(S_{2^{-m}}\) by \(N_*\) balls of radius \(2^{-m+1}\) gives
\[
 \mathcal H^\alpha_{2^{-m_0+2}}(\Sigma)
 \le\sum_{m\ge m_0}N_*\bigl(2^{-m+2}\bigr)^\alpha
 =N_*4^\alpha\frac{2^{-m_0\alpha}}{1-2^{-\alpha}}
 \xrightarrow[m_0\to\infty]{}0 .
\]
\end{proof}

\begin{assessment}[Comparison with the unconditional partial-regularity
bound]
\label{ass:v43-comparison}
Without a rate hypothesis, the external CKN theorem gives that the
parabolic one-dimensional Hausdorff measure of the space-time singular set
vanishes.  Under the type-I enstrophy rate the time slice of the singular
set has dimension zero, because the two-sided slab dissipation of
Lemma~\ref{lem:v43-slab} fixes the total mass at each scale while
\(\varepsilon\)-regularity fixes the minimal mass a singular point must
capture.  The bound \(N_*\) on active points is not a bound on the number
of singular points, because distinct singular points may be active at
disjoint sets of scales; finiteness of \(\Sigma\) is not proved here.
\end{assessment}

\begin{corollary}[The ancient limit sees a fraction of the slab dissipation]
\label{cor:v43-limit-fraction}
For the normalized ancient limit of Proposition~\ref{prop:v42-normalization},
\begin{equation}
 \iint_{B_1(0)\times(-1,0)}|\nabla U|^2\ \ge\ \frac{\varepsilon_1}{2}
 \ \ge\ \frac{c_*}{2}\,\int_{-1}^0\|\nabla U(s)\|_{L^2(\R^3)}^2\,\dd s ,
 \label{eq:v43-limit-fraction}
\end{equation}
that is, the unit ball at the origin carries at least the fraction
\(c_*/2\) of the dissipation of \(U\) in the unit slab.
\end{corollary}

\begin{proof}
The first inequality is \eqref{eq:v41-limit-nontrivial} with
\(\eta_0\) replaced by \(0\), which only enlarges the domain.  The second
uses \(\int_{-1}^0\|\nabla U(s)\|_2^2\le2C_I^2\nu^{3/2}\) from
\eqref{eq:v41-limit-type-I} and the definition of \(c_*\).
\end{proof}

% =============================================================================
\section{Integrated V43 conclusion and next acquisition order}
\label{sec:v43-conclusion}
% =============================================================================

\begin{theorem}[What V43 proves]
\label{thm:v43-conclusion}
Relative to V1--V42, and modulo the import (E1):
\begin{enumerate}[label=\textup{(V43.\arabic*)},leftmargin=6.2em]
\item the flat profile does not give decay of the ancient limit
      (Proposition~\ref{prop:v43-local-energy-rescaled});
\item the slab dissipation is two-sided of order \(r\)
      (Lemma~\ref{lem:v43-slab}), a concentration cylinder captures the
      fraction \(c_*\), and the type-I constant obeys
      \eqref{eq:v43-type-I-lower};
\item at most \(N_*\) well-separated points are active at any scale
      (Lemma~\ref{lem:v43-few-active});
\item the singular set at the singular time has Hausdorff dimension zero
      under a type-I enstrophy rate (Theorem~\ref{thm:v43-dimension-zero});
\item the normalized ancient limit carries a fixed fraction of its slab
      dissipation in the unit ball (Corollary~\ref{cor:v43-limit-fraction}).
\end{enumerate}
The full Navier--Stokes stability/global-regularity theorem remains
unproved.
\end{theorem}

\begin{proof}
The cited results.
\end{proof}

\begin{programme}[V44 acquisition order]
\label{prog:v44}
\begin{enumerate}[label=\textup{(AC\arabic*)},leftmargin=4.8em]
\item Finiteness of \(\Sigma\).  Determine whether the two-sided slab
      dissipation forces every singular point to be active at a positive
      proportion of dyadic scales; if so, pigeonholing over scales bounds
      \(|\Sigma|\) by a multiple of \(N_*\).  The candidate tool is the
      backward cone applied to the local dissipation
      \(\iint_{Q_r(x)}|\nabla u|^2\) rather than to the global enstrophy.
\item Local backward cone.  Prove or disprove a localized form of
      Lemma~\ref{lem:v37-backward-cone} for the local enstrophy
      \(\int_{B_r(x)}|\nabla u(t)|^2\), using the localized energy
      inequality of Proposition~\ref{prop:v42-local-energy} at the level of
      vorticity.
\item Type-II: item (AC2) of Programme~\ref{prog:v41} remains.
\item Prepare the compulsory V45 companion audit.
\end{enumerate}
\end{programme}

\begin{assessment}[Position after V43]
\label{ass:v43-position}
The type-I regime now has a geometric description: dimension-zero singular
set, at most \(N_*\) active points per scale, an ancient limit with a
fixed fraction of its slab dissipation in the unit ball, and the CKN
threshold as a lower bound on the type-I constant.  The Liouville gate of
V41 is unchanged.  No full stability or global-regularity claim is made.
\end{assessment}

% =============================================================================
\section*{Version V43 append-only record}
\addcontentsline{toc}{section}{Version V43 append-only record}
% =============================================================================

The complete V42 mathematical source was retained before this continuation.
Its SHA--256 digest was
\begin{center}\scriptsize\ttfamily
cfe70f62c7d20cf8d0f3b0d0144899f8deb082bf2689f15a46ea09647e4fc1b5.
\end{center}
V43 recorded that the flat profile does not yield decay of the ancient
limit, proved the two-sided slab dissipation, the concentration fraction,
the CKN lower bound on the type-I constant, the bound \(N_*\) on active
points per scale, and the dimension-zero theorem for the type-I singular
set.  No full regularity claim is made.

% =============================================================================
\section{V44: localization of the type-I singularity to a dimension-zero
set}
\label{sec:v44-entry}
% =============================================================================

Items (AC1)--(AC2) of Programme~\ref{prog:v44} are not achieved: no
argument was found forcing a singular point to be active at a positive
proportion of scales, and no localized backward cone is available because
the vortex-stretching term has no local bound in terms of local enstrophy.
Both are recorded.  The version instead assembles what V32, V41, and V43
give jointly in physical space.  Under a type-I enstrophy rate the singular
set \(\Sigma\) at the singular time is closed and of Hausdorff dimension
zero; the solution extends smoothly up to and including \(T_*\) away from
\(\Sigma\); the strong \(L^2\) limit \(u(T_*)\) of V32 coincides with that
extension off \(\Sigma\); the energy in any shrinking neighbourhood of
\(\Sigma\) tends to zero uniformly in time, while the enstrophy in every
fixed neighbourhood of \(\Sigma\) carries the full Leray blow-up rate and
every concentration cylinder carries a fixed fraction of the slab
dissipation.  A type-I singularity is therefore an energy-free,
dissipation-carrying event on a dimension-zero set.  No global regularity
claim is made.  The compulsory companion audit is due in V45.

% =============================================================================
\section{Two negative records}
\label{sec:v44-negative}
% =============================================================================

\begin{firewall}[Sparse activity is not excluded]
\label{fw:v44-sparse-activity}
For a singular point \(x\) with concentration scales \(r_k\), the
inclusion \(Q_{r_k}(x)\subset Q_r(x)\) gives \(f_x(r)\ge(r_k/r)
f_x(r_k)\) for \(r\ge r_k\), so activity at \(r_k\) implies activity only
up to \(2r_k\), and nothing constrains \(f_x\) at scales below \(r_k\)
except the trivial \(f_x\le2C_I^2\nu^{3/2}\).  A singular point active
only on a lacunary sequence of scales is consistent with every bound in
this series, so the pigeonhole argument of (AC1) does not start.
\end{firewall}

\begin{firewall}[No local backward cone]
\label{fw:v44-no-local-cone}
The local enstrophy \(y_\phi(t):=\int|\omega|^2\phi\), for a cutoff
\(\phi\), satisfies
\(\frac12y_\phi'+\nu\int|\nabla\omega|^2\phi
=\int(\omega\cdot\nabla u)\cdot\omega\,\phi
+\frac\nu2\int|\omega|^2\Delta\phi+\frac12\int|\omega|^2u\cdot\nabla\phi\).
The stretching term is bounded by
\(\|\nabla u\|_{L^\infty(\operatorname{supp}\phi)}y_\phi\), not by a power
of \(y_\phi\) alone, and the transport term needs \(u\) on the support.  A
differential inequality \(y_\phi'\le Cy_\phi^3\), which is what a
backward cone requires, does not follow.  Item (AC2) is closed
negatively within the present tools.
\end{firewall}

% =============================================================================
\section{Smooth extension off the singular set}
\label{sec:v44-extension}
% =============================================================================

Let \(\Sigma\) be as in Theorem~\ref{thm:v43-dimension-zero}.  For
\(\varepsilon>0\) write \(N_\varepsilon:=\{x:\operatorname{dist}(x,\Sigma)<\varepsilon\}\).

\begin{theorem}[Localization of a type-I singularity]
\label{thm:v44-localization}
Assume \eqref{eq:v32-type-I}.  Then:
\begin{enumerate}[label=\textup{(\roman*)},leftmargin=3.7em]
\item \(\Sigma\) is closed, nonempty, and \(\mathcal H^\alpha(\Sigma)=0\)
      for every \(\alpha>0\);
\item \(u\) extends to a smooth function on
      \([0,T_*]\times(\Torus^3\setminus\Sigma)\), and for every
      \(\varepsilon>0\) and every \(k\), \(\sup_{t<T_*}
      \|u(t)\|_{C^k(\Torus^3\setminus N_\varepsilon)}<\infty\);
\item the strong \(L^2\) limit \(u(T_*)\) of
      Theorem~\ref{thm:v32-type-I-firewall}(i) coincides with the smooth
      extension on \(\Torus^3\setminus\Sigma\), belongs to \(L^2(\Torus^3)\)
      with \(\|u(T_*)\|_2=\lim_{t\to T_*}\|u(t)\|_2\), and is divergence-free;
\item energy does not concentrate on \(\Sigma\):
      \begin{equation}
       \lim_{\varepsilon\to0}\ \sup_{t<T_*}\int_{N_\varepsilon}|u(t)|^2\,\dd x=0 ;
       \label{eq:v44-no-energy-concentration}
      \end{equation}
\item enstrophy does concentrate on \(\Sigma\): for every
      \(\varepsilon>0\) there is \(C_\varepsilon\) with
      \begin{equation}
       \int_{N_\varepsilon}|\nabla u(t)|^2\,\dd x
       \ \ge\ c_L^2\nu^{3/2}(T_*-t)^{-1/2}-C_\varepsilon
       \qquad(t<T_*),
       \label{eq:v44-enstrophy-concentration}
      \end{equation}
      and every concentration cylinder \(Q_{r_k}(x_0)\), \(x_0\in\Sigma\),
      carries at least the fraction \(c_*\) of the slab dissipation.
\end{enumerate}
\end{theorem}

\begin{proof}
(i) Nonemptiness and dimension are Theorem~\ref{thm:v43-dimension-zero}.
The set of points near which \(u\) is bounded is open by definition, so
\(\Sigma\) is closed.

(ii) Each \(x\notin\Sigma\) has a neighbourhood \(B_{r}(x)\times(T_*-r^2,T_*)\)
on which \(u\) is bounded; a bounded solution is smooth on the cylinder with
all derivatives bounded up to \(T_*\) by the local Serrin regularity theory
(import (E4) in its local form, applied to \(u\) on the cylinder, where
\(L^\infty\) is a Serrin class).  Compact subsets of
\(\Torus^3\setminus\Sigma\) are covered by finitely many such cylinders,
which gives the uniform \(C^k\) bounds off \(N_\varepsilon\) and the
smooth extension.

(iii) By V32, \(u(t)\to u(T_*)\) strongly in \(L^2\).  On
\(\Torus^3\setminus N_\varepsilon\), the uniform \(C^1\) bound and
Arzel\`a--Ascoli give uniform convergence of \(u(t)\) to the smooth
extension as \(t\to T_*\); the \(L^2\) limit agrees with it almost
everywhere there, hence everywhere by continuity, and \(\varepsilon\) is
arbitrary.  Strong convergence gives the norm identity, and the
divergence-free condition passes to weak limits.

(iv) \(\int_{N_\varepsilon}|u(t)|^2\le2\int_{N_\varepsilon}|u(T_*)|^2
+2\|u(t)-u(T_*)\|_2^2\).  Given \(\eta>0\), strong convergence gives
\(t_\eta\) with the second term below \(\eta\) for \(t\ge t_\eta\); since
\(\Sigma\) has Lebesgue measure zero and \(u(T_*)\in L^2\), the first term
tends to zero as \(\varepsilon\to0\); for \(t\le t_\eta\) the solution is
smooth and \(\int_{N_\varepsilon}|u(t)|^2\le|N_\varepsilon|
\sup_{t\le t_\eta}\|u(t)\|_\infty^2\to0\).

(v) By (ii), \(\int_{\Torus^3\setminus N_\varepsilon}|\nabla u(t)|^2\le
C_\varepsilon\) uniformly in \(t\), and the global enstrophy is at least
\(c_L^2\nu^{3/2}(T_*-t)^{-1/2}\) by Lemma~\ref{lem:v32-leray-rate}.  The
fraction statement is Proposition~\ref{prop:v43-fraction}.
\end{proof}

\begin{corollary}[The terminal datum of a type-I singularity]
\label{cor:v44-terminal-datum}
\(u(T_*)\in L^2(\Torus^3)\) is divergence-free, smooth off a closed set of
Hausdorff dimension zero, and satisfies
\(\int_{\Torus^3\setminus N_\varepsilon}|\nabla u(T_*)|^2<\infty\) for every
\(\varepsilon>0\).  Whether \(\nabla u(T_*)\in L^2(\Torus^3)\) is not
determined: if it were, the \(H^1\) local theory would restart the solution
from \(u(T_*)\) as a strong solution, and \(T_*\) would be a singular time
of the trajectory only in the sense that the enstrophy diverges as
\(t\uparrow T_*\) and is finite again at \(T_*\); this is not excluded by
anything in this series.
\end{corollary}

\begin{proof}
Immediate from Theorem~\ref{thm:v44-localization}(ii)--(iii) and weak
lower semicontinuity on \(\Torus^3\setminus N_\varepsilon\).
\end{proof}

\begin{assessment}[The shape of a type-I singularity after V44]
\label{ass:v44-shape}
Under a type-I enstrophy rate a singularity is now described as follows.
Its trace on \(\Torus^3\) is a closed set of dimension zero.  Away from it
the flow is smooth through the singular time.  It carries no energy in the
limit, and at most \(C\,r\log^2\) energy in a ball of radius \(r\) at its
points.  It carries the entire divergent part of the enstrophy, with a
fixed fraction \(c_*\) of each slab's dissipation in one parabolic
cylinder at concentration scales.  Rescaled at such a cylinder it produces
a nontrivial infinite-energy ancient solution with bounded enstrophy in the
past.  The critical norm of the solution grows at most logarithmically.
Each of these statements is proved; none of them, singly or jointly,
excludes the singularity, and Firewall~\ref{fw:v41-gate-status} names the
Liouville problem that would.
\end{assessment}

% =============================================================================
\section{Local energy at a singular point}
\label{sec:v44-local-energy}
% =============================================================================

\begin{proposition}[Two-sided information at a singular point]
\label{prop:v44-local-two-sided}
Let \(x_0\in\Sigma\) with concentration scales \(r_k\).  Then for \(r_k\)
small,
\begin{align}
 \int_{B_{r_k}(x_0)}|u(T_*-r_k^2)|^2\,\dd x
 &\le\Bigl(\frac{4\pi}3\Bigr)^{1/3}C_3^2\,r_k\,M_\delta(T_*-r_k^2)^2,
 \label{eq:v44-local-energy-upper}\\
 \frac1{r_k}\iint_{Q_{r_k}(x_0)}|\nabla u|^2
 &\ \ge\ \varepsilon_1,
 \qquad
 \frac1{r_k}\iint_{Q_{r_k}(x_0)}|u|^3
 \ \le\ C\,r_k\Bigl(1+\log^3\frac{9\delta}{r_k^2}\Bigr),
 \label{eq:v44-local-dissipation-cubic}
\end{align}
Thus the scale-invariant dissipation at a concentration cylinder is
bounded below by an absolute constant while the scale-invariant local
energy and cubic quantities are at most logarithmic in
\(1/r_k\); no bound in this series makes the latter small.
\end{proposition}

\begin{proof}
The first display is \eqref{eq:v36-local-L2} with \(r=r_k\) and
\(t=T_*-r_k^2\).  For the third, \(\|u(t)\|_3^3\le C_3^3\|A^{1/4}u(t)\|_2^3
\le C_3^3(2A_\delta+2B^2L_\delta(t)^2)^{3/2}\) by
Theorem~\ref{thm:v33-type-I-log}, and
\(\int_{T_*-r^2}^{T_*}L_\delta(t)^3\,\dd t=\int_0^{r^2}\log^3(9\delta/s)\,\dd s
\le Cr^2(1+\log^3(9\delta/r^2))\); integrate over the slab and divide by
\(r_k\).  The second display is \eqref{eq:v41-concentration}.
\end{proof}

% =============================================================================
\section{Integrated V44 conclusion and next acquisition order}
\label{sec:v44-conclusion}
% =============================================================================

\begin{theorem}[What V44 proves]
\label{thm:v44-conclusion}
Relative to V1--V43, and modulo the imports (E1) and (E4): the
localization theorem \ref{thm:v44-localization} with its five parts, the
description of the terminal datum (Corollary~\ref{cor:v44-terminal-datum}),
and the two-sided local information at a singular point
(Proposition~\ref{prop:v44-local-two-sided}).  Items (AC1)--(AC2) of
Programme~\ref{prog:v44} are closed negatively within the present tools.
The full Navier--Stokes stability/global-regularity theorem remains
unproved.
\end{theorem}

\begin{proof}
The cited results and firewalls.
\end{proof}

\begin{programme}[V45 acquisition order, with compulsory companion audit]
\label{prog:v45}
\begin{enumerate}[label=\textup{(AC\arabic*)},leftmargin=4.8em]
\item Perform the compulsory review of the current SM and YM notes;
      record digests; import nothing numerical.
\item Terminal datum.  Determine whether \(\nabla u(T_*)\in L^2\) can be
      decided from the type-I structure: the enstrophy diverges as
      \(t\uparrow T_*\) at the rate \((T_*-t)^{-1/2}\), and by weak lower
      semicontinuity along any sequence, \(\|\nabla u(T_*)\|_2\le
      \liminf\|\nabla u(t)\|_2=\infty\) is no information; test instead
      whether the local enstrophy off \(N_\varepsilon\) converges and
      whether \(\int_{N_\varepsilon}|\nabla u(T_*)|^2\) is finite.
\item Restart.  If \(u(T_*)\in H^1\) fails, the trajectory continues past
      \(T_*\) only as a Leray--Hopf weak solution; determine what the
      strong \(L^2\) continuity of V32 implies for the energy inequality of
      that continuation at \(T_*\).
\item Type-II: item (AC2) of Programme~\ref{prog:v41} remains.
\end{enumerate}
\end{programme}

\begin{assessment}[Position after V44]
\label{ass:v44-position}
The type-I regime is fully localized: a dimension-zero set carrying
dissipation but no energy, with a logarithmic critical norm and an
infinite-energy ancient profile.  The next version audits the companions
and examines the terminal datum.  No full stability or global-regularity
claim is made.
\end{assessment}

% =============================================================================
\section*{Version V44 append-only record}
\addcontentsline{toc}{section}{Version V44 append-only record}
% =============================================================================

The complete V43 mathematical source was retained before this continuation.
Its SHA--256 digest was
\begin{center}\scriptsize\ttfamily
bb0294bddf584735f0640d1c97b681d16ed898d39741a3a58d953b6d349b44a7.
\end{center}
V44 localized a type-I singularity to its dimension-zero trace: smooth
extension off the trace, identification of the strong \(L^2\) limit,
vanishing energy and full enstrophy concentration on the trace, and the
two-sided local information at a concentration cylinder.  No full
regularity claim is made.

% =============================================================================
\section{V45: compulsory companion audit and the terminal datum of a type-I
singularity}
\label{sec:v45-entry}
% =============================================================================

This is a fifteenth-version continuation.  The compulsory review of the
two companion programmes is performed first.  The mathematical content
then settles items (AC2)--(AC3) of Programme~\ref{prog:v45} as far as the
present tools allow: under a type-I enstrophy rate the energy equality
holds up to and including the singular time, the trajectory glues to a
Leray--Hopf weak solution on \([0,\infty)\) that is strongly
\(L^2\)-continuous at \(T_*\) and satisfies the strong energy inequality
there, and the question whether the terminal datum lies in \(H^1\) is
shown to be undecided by every bound in the series, with the exact
consequence of each answer recorded.  No global regularity claim is made.
The next compulsory companion audit is V50.

% =============================================================================
\section{Compulsory V45 review of the current companion notes}
\label{sec:v45-companion-audit}
% =============================================================================

The two current companion files in \texttt{latex\_notes} were inspected
read-only.  The frozen checkpoints are
\begin{align}
 &\texttt{Renewal\_SM\_Einstein\_Continuum\_Research\_Notes\_V65.tex},
 \notag\\[-2pt]
 &\qquad
 \texttt{e1740d57875819e9a152e48e878063522b96cb7d1276c36a5d1b98e9eb9c91e6},
 \label{eq:v45-SM-checkpoint}\\
 &\texttt{Renewal\_Yang\_Mills\_Mass\_Gap\_Research\_Notes\_V45.tex},
 \notag\\[-2pt]
 &\qquad
 \texttt{3eda55900b68bbe85a1dc43eb0f51433374bf3559047967899aa7967772d706f}.
 \label{eq:v45-YM-checkpoint}
\end{align}
The review concentrated on SM V60--V65 and YM V43--V45.

\emph{SM V60--V65.}  V60 performs its own audit and, in its words,
``forbids removal of scale logarithms by an unproved sign''; it proves a
positive additive bridge between finite-cutoff endpoint transfers and
isolates a merge cocycle.  V61 reduces the linearized vacuum ADM
constraint surface at nonzero momentum to two transverse--traceless
canonical pairs.  V62--V64 study the negative Newton form induced by
eliminating a scalar constraint carrier: a volume-dependent stability
margin on the torus, form-boundedness on \(\R^3\) at fixed charge, an exact
minimum-flux characterization, dyadic cell decompositions with a scale
loss, and an exact spectral-shell Newton ledger in which an order-one
contribution on every active shell produces an unavoidable logarithm.  V65
computes the Wick-square spectrum of a free massive scalar and its
logarithmic Newton stock, one positive payment per ultraviolet shell.  The
SM V64 statement that order-one payments on every active shell force a
logarithm is the same arithmetic as Theorem~\ref{thm:v34-flat-profile}
here; it was proved independently in each programme.

\emph{YM V43--V45.}  V43 computes a shifted Bessel Wilson source, its
zero set, and a uniform zero-free disk, and proves scalar all-arity
cumulant cards.  V44 proves a Banach-logarithm tensor-cumulant identity,
uniform weighted-simplex invertibility of the operator source, and the
binary additive matrix cards, and shows that the direct source argument
loses exactly the higher-arity Gaussian cancellation.  V45 performs its
audit, derives exact additive covariance ledgers, and closes the ternary
unmarked and marked matrix cards.  None of these involves a fluid
quantity.

\begin{theorem}[V45 companion-audit decision]
\label{thm:v45-companion-decision}
\begin{enumerate}[label=\textup{(\roman*)},leftmargin=3.7em]
\item No SM V60--V65 constraint, Newton, Wick, or shell constant is
      imported.  The SM discipline against removing a logarithm by an
      unproved sign coincides with the outcome of V36--V39 here, where the
      sign was tested and found indefinite.
\item No YM V43--V45 Bessel, cumulant, or matrix-card estimate is
      imported.
\item Neither companion bears on the type-I Liouville gate of V41 or on
      the type-II ledger; the audit changes nothing in V41--V44.
\end{enumerate}
\end{theorem}

\begin{proof}
The hypotheses of every cited result are lattice-gauge, constraint, or
Gaussian-field objects; none is a Navier--Stokes balance.
\end{proof}

% =============================================================================
\section{The energy equality up to the singular time}
\label{sec:v45-energy-equality}
% =============================================================================

\begin{proposition}[Energy equality at \(T_*\) under a type-I rate]
\label{prop:v45-energy-equality}
Assume \eqref{eq:v32-type-I}, and let \(u(T_*)\) be the strong \(L^2\)
limit of Theorem~\ref{thm:v32-type-I-firewall}(i).  Then
\begin{equation}
 \boxed{\quad
 \frac12\|u(T_*)\|_2^2+\nu\int_0^{T_*}\|\nabla u(t)\|_2^2\,\dd t
 =\frac12\|u_0\|_2^2 .
 \quad}
 \label{eq:v45-energy-equality}
\end{equation}
In particular the dissipation up to the singular time is finite and equals
the energy lost, with no energy carried by the singularity.
\end{proposition}

\begin{proof}
The smooth solution satisfies the equality with \(T_*\) replaced by
\(t<T_*\).  As \(t\uparrow T_*\), \(\|u(t)\|_2\to\|u(T_*)\|_2\) by strong
convergence and the dissipation integral converges monotonically.
\end{proof}

\begin{theorem}[Glued Leray--Hopf continuation, strongly continuous at
\(T_*\)]
\label{thm:v45-glued-continuation}
Assume \eqref{eq:v32-type-I}.  Let \(v\) be a Leray--Hopf weak solution on
\([T_*,\infty)\) with initial datum \(u(T_*)\), which exists by Leray's
theorem (import), and may be chosen suitable in the sense of
Firewall~\ref{fw:v41-imports}(E1).  Define \(\bar u:=u\) on \([0,T_*)\)
and \(\bar u:=v\) on \([T_*,\infty)\).  Then \(\bar u\) is a Leray--Hopf
weak solution on \([0,\infty)\) with datum \(u_0\), it is strongly
continuous in \(L^2\) at \(T_*\), it satisfies the energy equality on
\([0,T_*]\) and the energy inequality
\begin{equation}
 \frac12\|\bar u(t)\|_2^2+\nu\int_{s}^{t}\|\nabla\bar u\|_2^2
 \le\frac12\|\bar u(s)\|_2^2
 \qquad\text{for }s=T_*\text{ and a.e. }s\ge T_*,\ t\ge s,
 \label{eq:v45-strong-energy-inequality}
\end{equation}
and it is smooth on \([0,T_*)\times\Torus^3\) and on
\([0,T_*]\times(\Torus^3\setminus\Sigma)\).
\end{theorem}

\begin{proof}
\(\bar u\in L^\infty(0,\infty;L^2)\cap L^2_{\rm loc}(0,\infty;H^1)\) by
\eqref{eq:v45-energy-equality} and the Leray--Hopf properties of \(v\).
The weak formulation holds on \([0,T_*)\) and on \([T_*,\infty)\); at
\(T_*\) the two pieces match because \(u(t)\to u(T_*)=v(T_*)\) strongly in
\(L^2\), so the distributional time derivative of \(\bar u\) has no
singular part at \(T_*\), and the weak formulation on \([0,\infty)\)
follows by splitting a test function at \(T_*\).  Strong continuity at
\(T_*\) from the left is V32 and from the right is the strong continuity
of Leray--Hopf solutions at their initial time.  The energy equality on
\([0,T_*]\) is Proposition~\ref{prop:v45-energy-equality}; the inequality
from \(s=T_*\) is the energy inequality of \(v\) from its initial time,
and from a.e. \(s\ge T_*\) is Leray's strong energy inequality for \(v\).
Smoothness off \(\Sigma\) up to \(T_*\) is
Theorem~\ref{thm:v44-localization}(ii).
\end{proof}

\begin{assessment}[What is and is not gained]
\label{ass:v45-glued}
For a general Leray--Hopf solution, the energy may drop at a singular time
and the solution need not be strongly continuous there.  Under a type-I
enstrophy rate neither happens: the type-I singular time is a point of
strong \(L^2\) continuity and of energy equality, exactly as at a regular
time.  What is not gained is any statement about \(v\) after \(T_*\):
\(v\) is a weak solution about which nothing more is proved here, and
uniqueness of the continuation is not asserted.
\end{assessment}

% =============================================================================
\section{Is the terminal datum in \(H^1\)?}
\label{sec:v45-terminal-H1}
% =============================================================================

\begin{proposition}[Both answers are consistent with every bound in the
series]
\label{prop:v45-H1-undecided}
Assume \eqref{eq:v32-type-I}.
\begin{enumerate}[label=\textup{(\roman*)},leftmargin=3.7em]
\item If \(\nabla u(T_*)\in L^2(\Torus^3)\), then the continuation \(v\)
      of Theorem~\ref{thm:v45-glued-continuation} coincides on
      \([T_*,T_*+\tau)\), for some \(\tau>0\), with the unique strong
      solution from \(u(T_*)\), the function \(t\mapsto\|\nabla\bar u(t)\|_2^2\)
      is finite on \([T_*,T_*+\tau)\) and tends to \(+\infty\) as
      \(t\uparrow T_*\), and no bound of V32--V44 is violated.
\item If \(\nabla u(T_*)\notin L^2\), then \(u(T_*)\in L^2\setminus H^1\)
      is smooth off \(\Sigma\) with \(\int_{\Torus^3\setminus N_\varepsilon}
      |\nabla u(T_*)|^2<\infty\) for every \(\varepsilon\), so the infinite
      Dirichlet integral is concentrated on every neighbourhood of a
      dimension-zero set, and again no bound of V32--V44 is violated.
\end{enumerate}
Weak lower semicontinuity gives only
\(\|\nabla u(T_*)\|_2\le\liminf_{t\to T_*}\|\nabla u(t)\|_2=\infty\), which
is no information.
\end{proposition}

\begin{proof}
(i) The \(H^1\) local existence theorem and weak--strong uniqueness of
Leray--Hopf solutions against strong solutions (imports) identify \(v\)
with the strong solution on its existence interval.  The left limit is
Lemma~\ref{lem:v32-leray-rate}.  The bounds of V32--V44 concern
\(t<T_*\) only, except Theorem~\ref{thm:v44-localization}(iii)--(iv),
which are consistent with either alternative.  (ii) is
Corollary~\ref{cor:v44-terminal-datum}.
\end{proof}

\begin{firewall}[A Fourier-side remark on alternative (i)]
\label{fw:v45-alternative-i}
Alternative (i) requires \(\sum_k|k|^2|\widehat u(k,T_*)|^2<\infty\) while
\(\sum_k|k|^2|\widehat u(k,t)|^2\asymp(T_*-t)^{-1/2}\) and each
\(\widehat u(k,t)\to\widehat u(k,T_*)\).  This is the Fourier-side
statement that the divergent enstrophy lives, at each \(t\), at
frequencies that tend to infinity as \(t\to T_*\), carrying vanishing
energy, which is precisely the uniform tail smallness of V32.  Thus
alternative (i) is not merely consistent with the type-I picture; it is
its most natural reading.  It is nevertheless not proved, and in
alternative (i) the singular time would be one at which the enstrophy
blows up from the left and is finite from the right, a behaviour that no
theorem in this series excludes.
\end{firewall}

% =============================================================================
\section{Integrated V45 conclusion and next acquisition order}
\label{sec:v45-conclusion}
% =============================================================================

\begin{theorem}[What V45 proves]
\label{thm:v45-conclusion}
Relative to V1--V44, modulo the stated imports: the companion-audit
decision (Theorem~\ref{thm:v45-companion-decision}); the energy equality up
to and including the singular time under a type-I rate
(Proposition~\ref{prop:v45-energy-equality}); the glued Leray--Hopf
continuation with strong \(L^2\) continuity and strong energy inequality at
\(T_*\) (Theorem~\ref{thm:v45-glued-continuation}); and the exact status of
the \(H^1\) question for the terminal datum
(Proposition~\ref{prop:v45-H1-undecided}).  The full Navier--Stokes
stability/global-regularity theorem remains unproved.
\end{theorem}

\begin{proof}
The cited results.
\end{proof}

\begin{programme}[V46 acquisition order]
\label{prog:v46}
\begin{enumerate}[label=\textup{(AC\arabic*)},leftmargin=4.8em]
\item Alternative (i) of Proposition~\ref{prop:v45-H1-undecided}.  Test
      whether a strong solution on \([T_*,T_*+\tau)\) with datum
      \(u(T_*)\), together with the type-I structure from the left, yields
      a contradiction through the H\"older continuity in time of
      \(\|\nabla\bar u\|_2^2\) on the right: the right derivative of the
      enstrophy at \(T_*\) is finite by the \(H^1\) balance, while the
      enstrophy from the left is infinite; determine whether the weak
      formulation across \(T_*\) transmits any information from left to
      right at the level of \(H^1\).
\item Alternative (ii).  Determine whether an \(L^2\) datum smooth off a
      dimension-zero set with locally finite Dirichlet integral off that
      set but infinite total Dirichlet integral can be the strong \(L^2\)
      limit of a family with uniformly small Fourier tails; this is a
      question about \(u(T_*)\) alone.
\item Type-II: item (AC2) of Programme~\ref{prog:v41} remains.
\item The next compulsory companion audit is V50.
\end{enumerate}
\end{programme}

\begin{assessment}[Position after V45]
\label{ass:v45-position}
A type-I singularity is now known to be an energy-conserving, strongly
\(L^2\)-continuous event whose terminal datum is either an \(H^1\) field
(in which case the trajectory restarts as a strong solution) or a field
whose infinite Dirichlet integral sits on a dimension-zero set.  Deciding
between the two is the next concrete question in the type-I regime.  No
full stability or global-regularity claim is made.
\end{assessment}

% =============================================================================
\section*{Version V45 append-only record}
\addcontentsline{toc}{section}{Version V45 append-only record}
% =============================================================================

The complete V44 mathematical source was retained before this continuation.
Its SHA--256 digest was
\begin{center}\scriptsize\ttfamily
52597298bbb739d1b522fe8a3c32eff68ecadddd772885ab9052b232b9376492.
\end{center}
V45 performed the compulsory companion audit, proved the energy equality
up to the singular time and the strongly continuous glued Leray--Hopf
continuation under a type-I rate, and recorded the exact status of the
\(H^1\) question for the terminal datum.  No full regularity claim is
made.

% =============================================================================
\section{V46: the \(L^\infty\) upgrade of the type-I rate and the
non-\(L^3\) terminal datum}
\label{sec:v46-entry}
% =============================================================================

Item (AC1) of Programme~\ref{prog:v46} is settled, and more.  Two facts
are proved, each modulo an explicitly imported external theorem.  First,
\(\varepsilon\)-regularity applied on cylinders of radius a small multiple
of the parabolic distance to the singular time upgrades the type-I
enstrophy rate to the classical type-I velocity bound
\(|u(x,t)|\le C_\infty(T_*-t)^{-1/2}\).  Second, if the terminal datum
\(u(T_*)\) were cubically integrable near a concentration point, the
ancient limit of V41 would vanish at its final time in \(L^2_{\rm loc}\),
would be bounded near that time outside a ball, and would therefore vanish
identically by the backward-uniqueness and unique-continuation package of
Escauriaza, Seregin, and \v Sver\'ak, contradicting its nontriviality.
Hence the terminal datum of a type-I singularity is not in
\(L^3\) on any neighbourhood of a concentration point, in particular not in
\(H^1\), and alternative (i) of Proposition~\ref{prop:v45-H1-undecided}
is excluded.  No global regularity claim is made.  The next compulsory
companion audit is V50.

% =============================================================================
\section{Additional imports}
\label{sec:v46-imports}
% =============================================================================

\begin{firewall}[External theorems added in V46]
\label{fw:v46-imports}
\begin{enumerate}[label=\textup{(E\arabic*)},leftmargin=3.6em,start=5]
\item \emph{CKN, cubic form.}  There are absolute \(\varepsilon_2>0\) and
      \(C_{\rm CKN}\) such that if a suitable weak solution \((u,p)\)
      satisfies \(r^{-2}\iint_{Q_r(z)}\bigl(|u|^3+|u||p-\bar p|\bigr)
      \le\varepsilon_2\), with \(\bar p\) the spatial mean of \(p\) over the
      ball, then \(|u|\le C_{\rm CKN}/r\) on \(Q_{r/2}(z)\), and \(u\) is
      smooth there with \(|\nabla u|\le C_{\rm CKN}'/r^2\) on
      \(Q_{r/4}(z)\).  Smooth solutions with \(u\in L^\infty_{\rm loc}L^6\)
      and \(p\in L^\infty_{\rm loc}L^3\) are suitable.
\item \emph{ESS backward uniqueness and unique continuation.}  Let
      \(\omega\) be smooth and bounded, with bounded gradient, on
      \(\{|y|>R\}\times(a,0)\), satisfying
      \(|\partial_s\omega-\nu\Delta\omega|\le M(|\omega|+|\nabla\omega|)\)
      there and \(\omega(\cdot,0)=0\) on \(\{|y|>R\}\); then
      \(\omega\equiv0\) on \(\{|y|>R\}\times(a,0)\).  If moreover the
      inequality holds on \(\R^3\times(a,b)\) with \(\omega\) bounded
      there and \(\omega(\cdot,s)=0\) outside a ball for every
      \(s\in(a,b)\), then \(\omega\equiv0\) on \(\R^3\times(a,b)\).
\end{enumerate}
\end{firewall}

% =============================================================================
\section{From the enstrophy rate to the velocity rate}
\label{sec:v46-Linfty}
% =============================================================================

\begin{theorem}[Type-I in enstrophy implies type-I in \(L^\infty\)]
\label{thm:v46-Linfty-upgrade}
Assume \eqref{eq:v32-type-I} on \([T_*-\delta,T_*)\).  There is
\(c_0\in(0,1)\), depending only on \(C_I,\nu,\varepsilon_2,C_6,C_{\rm CZ}\),
such that for every \(x\in\Torus^3\) and every
\(t\in[T_*-\delta/4,T_*)\),
\begin{equation}
 \boxed{\quad
 |u(x,t)|\le\frac{2C_{\rm CKN}}{c_0}\,(T_*-t)^{-1/2},
 \qquad
 |\nabla u(x,t)|\le\frac{4C_{\rm CKN}'}{c_0^2}\,(T_*-t)^{-1}.
 \quad}
 \label{eq:v46-Linfty-type-I}
\end{equation}
\end{theorem}

\begin{proof}
Fix \((x,t)\), put \(\theta:=T_*-t\), \(r:=c(\theta/2)^{1/2}\) with
\(c\le1\) to be chosen, and \(z:=(x,t+r^2/2)\), so that \((x,t)\) lies in
the interior of \(Q_{r/2}(z)\) and \(Q_r(z)\subset\Torus^3\times
(t-r^2/2,t+r^2/2)\subset\Torus^3\times[T_*-\delta,T_*)\).  On
\(Q_r(z)\), \(T_*-\tau\ge\theta-r^2/2\ge\theta/2\), so
\(\|u(\tau)\|_6\le C_6C_I\nu^{3/4}(\theta/2)^{-1/4}\) and
\(\|p(\tau)\|_3\le C_{\rm CZ}C_6^2C_I^2\nu^{3/2}(\theta/2)^{-1/2}\) by
\eqref{eq:v37-pressure-CZ} with exponent \(3\) in place of \(3/2\), which
holds by the same Riesz bound.  H\"older on \(B_r\) gives
\(\int_{B_r}|u|^3\le|B_r|^{1/2}\|u\|_6^3\) and
\(\int_{B_r}|u||p-\bar p|\le\|u\|_{L^3(B_r)}\cdot2\|p\|_{L^{3/2}(B_r)}
\le|B_r|^{1/6}\|u\|_6\cdot2|B_r|^{1/3}\|p\|_3\), hence
\[
 \frac1{r^2}\iint_{Q_r(z)}\bigl(|u|^3+|u||p-\bar p|\bigr)
 \le C_\sharp\,r^{3/2}(\theta/2)^{-3/4}
 =C_\sharp\,c^{3/2},
\]
with \(C_\sharp\) depending only on \(C_I,\nu,C_6,C_{\rm CZ}\).  Choose
\(c=c_0\) with \(C_\sharp c_0^{3/2}\le\varepsilon_2\).  By (E5),
\(|u|\le C_{\rm CKN}/r\) on \(Q_{r/2}(z)\) and
\(|\nabla u|\le C'_{\rm CKN}/r^2\) on \(Q_{r/4}(z)\); insert
\(r=c_0(\theta/2)^{1/2}\).
\end{proof}

\begin{assessment}[The two type-I notions coincide]
\label{ass:v46-two-notions}
The converse direction, from the velocity rate to the enstrophy rate, does
not follow by interpolation alone, because \(L^\infty\) and \(\dot H^1\) are
not comparable on the torus without an energy factor: \(\|\nabla u\|_2^2
\le\|u\|_2\|\Delta u\|_2\) needs \(H^2\).  Theorem~\ref{thm:v46-Linfty-upgrade}
therefore says that the enstrophy-rate hypothesis used throughout
V32--V45 is at least as strong as the classical type-I hypothesis, and all
results of this series in the type-I regime hold under the enstrophy rate
alone.  Under the velocity rate one has, in addition, the classical
rescaling theory of type-I singularities; V46 does not need it.
\end{assessment}

% =============================================================================
\section{Vanishing of the ancient limit at its final time}
\label{sec:v46-vanishing}
% =============================================================================

Let \(x_0\in\Sigma\) be a concentration point with scales \(r_k\), and
let \(U_k\), \(U\), \(P\) be as in Theorem~\ref{thm:v41-ancient-limit}.
Write \(e_R^{(k)}(s):=\frac12\int|U_k(s)|^2\phi_R\) and
\(e_R(s):=\frac12\int|U(s)|^2\phi_R\) with \(\phi_R\) as in
Section~\ref{sec:v42-local-energy}.

\begin{lemma}[Local energy of the rescaled fields near the final time]
\label{lem:v46-local-energy-near-zero}
There is \(C_R\), independent of \(k\), such that for \(-1\le s<0\),
\begin{equation}
 e_R^{(k)}(s)\le e_R^{(k)}(0^-)+C_R(-s)^{1/4},
 \qquad
 e_R^{(k)}(0^-):=\lim_{\sigma\uparrow0}e_R^{(k)}(\sigma)
 =\frac1{2r_k}\int_{\Torus^3}|u(T_*,x)|^2\phi_R\Bigl(\frac{x-x_0}{r_k}\Bigr)\dd x .
 \label{eq:v46-local-energy-near-zero}
\end{equation}
\end{lemma}

\begin{proof}
For \(\sigma<0\) the smooth solution \(U_k\) satisfies the local energy
identity
\[
 e_R^{(k)}(0^-)-e_R^{(k)}(s)
 =-\nu\int_s^0\!\!\int|\nabla U_k|^2\phi_R
 +\frac\nu2\int_s^0\!\!\int|U_k|^2\Delta\phi_R
 +\frac12\int_s^0\!\!\int\bigl(|U_k|^2+2P_k\bigr)U_k\cdot\nabla\phi_R,
\]
where the left side exists because \(u\) is strongly \(L^2\)-continuous at
\(T_*\) (Theorem~\ref{thm:v32-type-I-firewall}(i)) and rescaling preserves
this, and the time integrals converge absolutely: by
\eqref{eq:v41-rescaled-L6}, the second term is bounded by
\(C\int_s^0(-\sigma)^{-1/2}\dd\sigma\) and the third by
\(C_R\int_s^0(-\sigma)^{-3/4}\dd\sigma\), as in the proof of
Proposition~\ref{prop:v42-local-energy}.  Solving for \(e_R^{(k)}(s)\)
and bounding the dissipation term by \(\nu\int_s^0Y_k\) gives
\(e_R^{(k)}(s)\le e_R^{(k)}(0^-)+\nu\int_s^0Y_k+C(-s)^{1/2}+C_R(-s)^{1/4}\),
and \(\nu\int_s^0Y_k\le2\nu C_I^2\nu^{3/2}(-s)^{1/2}\).  The formula for
\(e_R^{(k)}(0^-)\) is the change of variables.
\end{proof}

\begin{proposition}[Local cubic integrability of the terminal datum forces
vanishing of the ancient limit]
\label{prop:v46-vanishing}
If \(u(T_*)\in L^3(B_{\rho_0}(x_0))\) for some \(\rho_0>0\), then for
every \(R>0\)
\begin{equation}
 e_R(s)\le C_R(-s)^{1/4}\qquad(-1\le s<0),
 \qquad\text{hence}\qquad
 U(s)\to0\ \text{in }L^2_{\rm loc}(\R^3)\ \text{as }s\uparrow0 .
 \label{eq:v46-vanishing}
\end{equation}
\end{proposition}

\begin{proof}
By H\"older, \(\int_{B_{2Rr_k}(x_0)}|u(T_*)|^2\le
\bigl(\int_{B_{2Rr_k}(x_0)}|u(T_*)|^3\bigr)^{2/3}|B_{2Rr_k}|^{1/3}\), and
the first factor tends to zero as \(k\to\infty\) by absolute continuity of
the integral, while the second is \(CR\,r_k\).  Hence
\(e_R^{(k)}(0^-)\le\frac1{2r_k}\int_{B_{2Rr_k}(x_0)}|u(T_*)|^2\to0\).
For fixed \(s<0\), \(U_k(s)\to U(s)\) in \(L^2(B_{2R})\) by
Theorem~\ref{thm:v41-ancient-limit} (strong convergence in
\(L^2((-S,-\eta);H^1(B_{2R}))\) and equicontinuity in time from the
uniform bound on \(\partial_sU_k\) give convergence at every fixed time),
so \(e_R(s)=\lim_ke_R^{(k)}(s)\le C_R(-s)^{1/4}\) by
Lemma~\ref{lem:v46-local-energy-near-zero}.
\end{proof}

% =============================================================================
\section{Boundedness of the ancient limit away from the origin near the
final time}
\label{sec:v46-exterior-bounded}
% =============================================================================

\begin{lemma}[Exterior boundedness near \(s=0\), and interior boundedness
away from \(s=0\)]
\label{lem:v46-bounded}
For the ancient limit \(U\) of Theorem~\ref{thm:v41-ancient-limit}:
\begin{enumerate}[label=\textup{(\roman*)},leftmargin=3.7em]
\item there is \(R_0\) such that \(U\), \(\nabla U\), and \(\nabla^2U\)
      are bounded on \(\{|y|\ge R_0\}\times(-1/16,0)\), and \(U\) extends
      continuously to \(s=0\) there;
\item for every \(\eta>0\), \(U\) and \(\nabla U\) are bounded on
      \(\R^3\times(-\infty,-\eta]\), with bounds depending on \(\eta\).
\end{enumerate}
\end{lemma}

\begin{proof}
(i) For \(y\in\R^3\) consider the unit cylinder \(Q_1(y,0)\).  By
\eqref{eq:v41-limit-type-I} and the Riesz bound,
\(\|U(s)\|_6\le C(-s)^{-1/4}\) and \(\|P(s)\|_3\le C(-s)^{-1/2}\), so the
integrands of \(F(y):=\iint_{Q_1(y,0)}(|U|^3+|U||P-\bar P|)\) are
dominated, after H\"older on \(B_1(y)\), by
\(C(-s)^{-3/4}\), which is integrable on \((-1,0)\), and for each fixed
\(s\) the local norms \(\|U(s)\|_{L^6(B_1(y))}\), \(\|P(s)\|_{L^3(B_1(y))}\)
tend to zero as \(|y|\to\infty\) because \(U(s)\in L^6(\R^3)\),
\(P(s)\in L^3(\R^3)\).  By dominated convergence \(F(y)\to0\) as
\(|y|\to\infty\), and \(F\) is continuous, so there is \(R_0\) with
\(F(y)\le\varepsilon_2\) for \(|y|\ge R_0-1\).  Since \(U\) is smooth with
the stated integrability it is suitable, and (E5) gives boundedness of
\(U\) on \(Q_{1/2}(y,0)\) and of \(\nabla U,\nabla^2U\) on
\(Q_{1/4}(y,0)\) for all such \(y\); these cylinders cover
\(\{|y|\ge R_0\}\times(-1/16,0)\).  A bounded smooth solution on a
cylinder is uniformly continuous up to its top, giving the extension.

(ii) For \((y,s)\) with \(s\le-\eta\), apply (E5) on the cylinder
\(Q_r(y,s+r^2/2)\) with \(r=c_0'(-s)^{1/2}\), exactly as in the proof of
Theorem~\ref{thm:v46-Linfty-upgrade} with \(\R^3\) in place of the torus;
the cubic quantity is at most \(C_\sharp'c_0'^{3/2}\le\varepsilon_2\) for
\(c_0'\) small, and (E5) gives \(|U(y,s)|\le2C_{\rm CKN}/(c_0'(-s)^{1/2})\)
and the corresponding gradient bound.
\end{proof}

% =============================================================================
\section{The terminal datum is not locally cubically integrable}
\label{sec:v46-main}
% =============================================================================

\begin{theorem}[Non-\(L^3\) terminal datum]
\label{thm:v46-non-L3}
Assume the type-I enstrophy rate \eqref{eq:v32-type-I}, and let
\(x_0\in\Sigma\) be a concentration point.  Then, modulo the imports
(E1)--(E6),
\begin{equation}
 \boxed{\quad
 u(T_*)\notin L^3\bigl(B_\rho(x_0)\bigr)\quad\text{for every }\rho>0 .
 \quad}
 \label{eq:v46-non-L3}
\end{equation}
In particular \(u(T_*)\notin L^3(\Torus^3)\), \(u(T_*)\notin H^1(\Torus^3)\),
and alternative (ii) of Proposition~\ref{prop:v45-H1-undecided} holds:
the trajectory cannot be restarted from \(u(T_*)\) as a strong solution.
\end{theorem}

\begin{proof}
Suppose \(u(T_*)\in L^3(B_{\rho_0}(x_0))\).  By
Proposition~\ref{prop:v46-vanishing}, \(U(s)\to0\) in \(L^2_{\rm loc}\) as
\(s\uparrow0\).  By Lemma~\ref{lem:v46-bounded}(i), \(U\) is bounded and
continuous up to \(s=0\) on \(\{|y|\ge R_0\}\times(-1/16,0]\), so
\(U(\cdot,0)=0\) there, and the vorticity \(\omega=\operatorname{curl}U\)
satisfies \(\omega(\cdot,0)=0\) on \(\{|y|\ge R_0\}\).  On that exterior
region \(\omega\) is smooth and bounded with bounded gradient, and it
satisfies the vorticity equation
\(\partial_s\omega-\nu\Delta\omega=(\omega\cdot\nabla)U-(U\cdot\nabla)\omega\),
whose right side is bounded by \(M(|\omega|+|\nabla\omega|)\) with
\(M=\sup(|U|+|\nabla U|)\) over the region.  By (E6),
\(\omega\equiv0\) on \(\{|y|\ge R_0\}\times(-1/16,0)\).

Fix \(\eta\in(0,1/16)\).  On \(\R^3\times(-1/16,-\eta)\), \(U\) and
\(\nabla U\) are bounded by Lemma~\ref{lem:v46-bounded}(ii), the
vorticity inequality holds with bounded coefficients, and
\(\omega(\cdot,s)\) vanishes outside \(B_{R_0}\) for every \(s\) in the
interval.  By the second part of (E6), \(\omega\equiv0\) on
\(\R^3\times(-1/16,-\eta)\).  Then \(\operatorname{curl}U=0\) and
\(\operatorname{div}U=0\) make each component of \(U(s)\) harmonic on
\(\R^3\); a harmonic function in \(L^6(\R^3)\) is zero.  Hence
\(U\equiv0\) on \(\R^3\times(-1/16,0)\), as \(\eta\) is arbitrary.

Now argue backward.  On \(\R^3\times(-1/8,-1/16)\), \(U\) is bounded with
bounded gradient (Lemma~\ref{lem:v46-bounded}(ii) with \(\eta=1/16\)),
and \(\omega(\cdot,-1/16)=0\).  Applying (E6) on the exterior of any ball
and then its second part as above gives \(U\equiv0\) on
\(\R^3\times(-1/8,-1/16)\).  Repeating on \((-1/4,-1/8)\),
\((-1/2,-1/4)\), \((-1,-1/2)\) gives \(U\equiv0\) on
\(\R^3\times(-1,0)\), which contradicts
\eqref{eq:v41-limit-nontrivial}.  Therefore
\(u(T_*)\notin L^3(B_{\rho_0}(x_0))\).  Since \(H^1(\Torus^3)\subset
L^6\subset L^3\), the remaining assertions follow, and the restart
statement is Proposition~\ref{prop:v45-H1-undecided}(i) read
contrapositively.
\end{proof}

\begin{corollary}[Scale-critical local energy of the terminal datum]
\label{cor:v46-critical-local-energy}
Under the hypotheses of Theorem~\ref{thm:v46-non-L3},
\begin{equation}
 \limsup_{\rho\to0}\ \frac1\rho\int_{B_\rho(x_0)}|u(T_*,x)|^2\,\dd x\ >\ 0 ,
 \label{eq:v46-critical-local-energy}
\end{equation}
while \(\int_{B_\rho(x_0)}|u(T_*)|^2\to0\) as \(\rho\to0\) because
\(u(T_*)\in L^2\).  Thus the terminal datum carries energy of exactly
scale-critical size \(\asymp\rho\) in balls around each concentration
point, on a sequence of radii.
\end{corollary}

\begin{proof}
The proof of Proposition~\ref{prop:v46-vanishing} used only
\(e_R^{(k)}(0^-)\to0\), which follows from
\(\rho^{-1}\int_{B_\rho(x_0)}|u(T_*)|^2\to0\) along \(\rho=2Rr_k\).  If
\eqref{eq:v46-critical-local-energy} failed, this would hold for every
\(R\), and the proof of Theorem~\ref{thm:v46-non-L3} would go through
verbatim.
\end{proof}

\begin{assessment}[The terminal datum of a type-I singularity]
\label{ass:v46-terminal-datum}
Combining V32, V44, V45, and V46: the terminal datum \(u(T_*)\) is a
divergence-free field in \(L^2(\Torus^3)\), attained strongly with energy
equality, smooth off a closed set \(\Sigma\) of Hausdorff dimension zero,
with locally finite Dirichlet integral off \(\Sigma\), and at every
concentration point of \(\Sigma\) it fails to be cubically integrable and
carries scale-critical local energy on a sequence of scales.  This is the
profile of a \(-1\)-homogeneous field such as \(|x-x_0|^{-1}\), which is
the terminal datum of a self-similar singularity; V46 proves that every
type-I singularity, self-similar or not, has a terminal datum of this
critical type at its concentration points.  The result does not exclude
the singularity.  It does exclude the possibility, left open in V45, that
a type-I singular time is a mere left-sided enstrophy blow-up followed by
a strong restart.
\end{assessment}

% =============================================================================
\section{Integrated V46 conclusion and next acquisition order}
\label{sec:v46-conclusion}
% =============================================================================

\begin{theorem}[What V46 proves]
\label{thm:v46-conclusion}
Relative to V1--V45, modulo (E1)--(E6):
\begin{enumerate}[label=\textup{(V46.\arabic*)},leftmargin=6.2em]
\item the type-I enstrophy rate implies the type-I velocity and gradient
      rates \eqref{eq:v46-Linfty-type-I};
\item local cubic integrability of the terminal datum at a concentration
      point forces the ancient limit to vanish at its final time in
      \(L^2_{\rm loc}\) (Proposition~\ref{prop:v46-vanishing});
\item the ancient limit is bounded, with bounded derivatives, outside a
      ball up to its final time and everywhere away from it
      (Lemma~\ref{lem:v46-bounded});
\item the terminal datum of a type-I singularity is not in \(L^3\) near
      any concentration point, hence not in \(H^1\), and carries
      scale-critical local energy there
      (Theorem~\ref{thm:v46-non-L3}, Corollary~\ref{cor:v46-critical-local-energy}).
\end{enumerate}
The full Navier--Stokes stability/global-regularity theorem remains
unproved.
\end{theorem}

\begin{proof}
The cited results.
\end{proof}

\begin{programme}[V47 acquisition order]
\label{prog:v47}
\begin{enumerate}[label=\textup{(AC\arabic*)},leftmargin=4.8em]
\item Quantify the critical local energy.  Combine
      \eqref{eq:v46-critical-local-energy} with the upper bound
      \(\rho^{-1}\int_{B_\rho}|u(t)|^2\le C\log^2\) of V36 at \(t<T_*\) and
      the strong convergence \(u(t)\to u(T_*)\); determine whether the
      terminal datum's scaled local energy is bounded, so that
      \(u(T_*)\in L^{3,\infty}\) locally, or can diverge logarithmically.
\item Use the velocity rate.  With \eqref{eq:v46-Linfty-type-I} the
      rescaled fields are bounded by \(C(-s)^{-1/2}\); revisit the
      Liouville gate of V41 in the class of ancient solutions with this
      bound and bounded enstrophy, and record what the backward-uniqueness
      package gives when the final-time trace is not zero.
\item Type-II: item (AC2) of Programme~\ref{prog:v41} remains.
\item The next compulsory companion audit is V50.
\end{enumerate}
\end{programme}

\begin{assessment}[Position after V46]
\label{ass:v46-position}
The type-I regime is now equivalent to the classical type-I regime, and
its terminal datum is proved to be of scale-critical, non-\(L^3\) type at
the singular points.  The Liouville gate of V41 remains the obstruction to
excluding the regime.  No full stability or global-regularity claim is
made.
\end{assessment}

% =============================================================================
\section*{Version V46 append-only record}
\addcontentsline{toc}{section}{Version V46 append-only record}
% =============================================================================

The complete V45 mathematical source was retained before this continuation.
Its SHA--256 digest was
\begin{center}\scriptsize\ttfamily
b51d63bad08f5046fdeb27499edda7d3bc7ab3f0f1a46baf04ff60b3e2e7a634.
\end{center}
V46 proved the \(L^\infty\) upgrade of the type-I enstrophy rate, the
vanishing of the ancient limit at its final time under local cubic
integrability of the terminal datum, the exterior and interior boundedness
of the ancient limit, and, through the imported backward-uniqueness
package, that the terminal datum of a type-I singularity is not locally
\(L^3\) and not \(H^1\).  No full regularity claim is made.

% =============================================================================
\section{V47: the uniform critical Morrey bound}
\label{sec:v47-entry}
% =============================================================================

Item (AC1) of Programme~\ref{prog:v47} is settled in the strong form.
Under a type-I enstrophy rate the local energy in every ball of radius
\(\rho\), at every time up to and including the singular time, is at most
a constant times \(\rho\).  The proof is the local energy identity on a
parabolic cylinder ending at the singular time, in which every term is of
scale-critical size by the \(L^6\) and pressure bounds already in the
series and the strong \(L^2\) continuity of V32.  This replaces the
logarithmic local-energy bound of V36 by a constant, shows that the
terminal datum lies in the critical Morrey space, and, with V46, gives a
two-sided scale-critical description of the terminal datum at every
concentration point.  The final-time trace of the ancient limit is
identified as a weak limit in the same Morrey class.  No global regularity
claim is made.  The next compulsory companion audit is V50.

% =============================================================================
\section{Local energy at the parabolic scale and beyond}
\label{sec:v47-morrey}
% =============================================================================

Throughout, \eqref{eq:v32-type-I} holds on \([T_*-\delta,T_*)\), and
\(\phi_\rho(x')=\phi((x'-x)/\rho)\) is the cutoff of
Section~\ref{sec:v42-local-energy} centred at \(x\).  Write
\(C_6\), \(C_{\rm CZ}\) for the mean-zero Sobolev and Riesz constants on
\(\Torus^3\), so that for \(\tau<T_*\)
\begin{equation}
 \|u(\tau)\|_6\le C_6C_I\nu^{3/4}(T_*-\tau)^{-1/4},
 \qquad
 \|p(\tau)\|_3\le C_{\rm CZ}C_6^2C_I^2\nu^{3/2}(T_*-\tau)^{-1/2}.
 \label{eq:v47-L6-L3}
\end{equation}

\begin{theorem}[Uniform critical Morrey bound up to the singular time]
\label{thm:v47-morrey}
There is \(C_M\), depending only on \(C_I,\nu,C_6,C_{\rm CZ},\phi\), such
that for every \(x\in\Torus^3\), every \(0<\rho\le(\delta/2)^{1/2}\), and
every \(t\in[T_*-\delta/2,T_*]\),
\begin{equation}
 \boxed{\quad
 \int_{B_\rho(x)}|u(t,x')|^2\,\dd x'\ \le\ C_M\,\rho ,
 \quad}
 \label{eq:v47-morrey}
\end{equation}
where \(u(T_*)\) is the strong \(L^2\) limit of V32.
\end{theorem}

\begin{proof}
Fix \(x\), \(\rho\), \(t\).  \emph{Case \(\rho^2\le T_*-t\).}  H\"older
and \eqref{eq:v47-L6-L3} give
\(\int_{B_\rho}|u(t)|^2\le|B_\rho|^{2/3}\|u(t)\|_6^2
\le C\rho^2(T_*-t)^{-1/2}\le C\rho\).

\emph{Case \(\rho^2>T_*-t\), including \(t=T_*\).}  Put
\(t_0:=t-\rho^2\ge T_*-\delta\).  For \(t'<T_*\) the smooth solution
satisfies the local energy identity
\[
 \frac12\int|u(t')|^2\phi_\rho-\frac12\int|u(t_0)|^2\phi_\rho
 =-\nu\int_{t_0}^{t'}\!\!\int|\nabla u|^2\phi_\rho
 +\frac\nu2\int_{t_0}^{t'}\!\!\int|u|^2\Delta\phi_\rho
 +\frac12\int_{t_0}^{t'}\!\!\int\bigl(|u|^2+2(p-\bar p)\bigr)u\cdot\nabla\phi_\rho,
\]
with \(\bar p\) any constant, since \(\int\bar p\,u\cdot\nabla\phi_\rho
=-\bar p\int\phi_\rho\operatorname{div}u=0\).  If \(t=T_*\), let
\(t'\uparrow T_*\): the left side converges by strong \(L^2\) continuity,
the dissipation term converges monotonically, and the remaining terms
converge absolutely by the bounds below.  Dropping the dissipation term,
\[
 \int_{B_\rho}|u(t)|^2
 \le\int_{B_{2\rho}}|u(t_0)|^2
 +\nu\int_{t_0}^{t}\!\!\int|u|^2|\Delta\phi_\rho|
 +\int_{t_0}^{t}\!\!\int\bigl(|u|^3+2|u||p-\bar p|\bigr)|\nabla\phi_\rho| .
\]
First term: \(|B_{2\rho}|^{2/3}\|u(t_0)\|_6^2\le C\rho^2(T_*-t_0)^{-1/2}
\le C\rho^2\rho^{-1}=C\rho\), since \(T_*-t_0\ge\rho^2\).  Second term:
\(|\Delta\phi_\rho|\le C\rho^{-2}\mathbf 1_{B_{2\rho}}\) and
\(\int_{B_{2\rho}}|u(\tau)|^2\le C\rho^2(T_*-\tau)^{-1/2}\), so it is at most
\(C\nu\int_{t_0}^{T_*}(T_*-\tau)^{-1/2}\dd\tau\le C\nu\cdot2(T_*-t_0)^{1/2}
\le C\nu\cdot2\sqrt2\,\rho\), using \(T_*-t_0\le2\rho^2\) because
\(t\ge T_*-\rho^2\).  Third term: \(|\nabla\phi_\rho|\le C\rho^{-1}
\mathbf 1_{B_{2\rho}}\), and on \(B_{2\rho}\), with \(\bar p\) the mean of
\(p(\tau)\) over \(B_{2\rho}\),
\(\int_{B_{2\rho}}|u|^3\le|B_{2\rho}|^{1/2}\|u\|_6^3\le C\rho^{3/2}(T_*-\tau)^{-3/4}\)
and \(\int_{B_{2\rho}}|u||p-\bar p|\le|B_{2\rho}|^{1/6}\|u\|_6\cdot
2|B_{2\rho}|^{1/3}\|p\|_3\le C\rho^{3/2}(T_*-\tau)^{-3/4}\); integrating,
\(\int_{t_0}^{T_*}(T_*-\tau)^{-3/4}\dd\tau=4(T_*-t_0)^{1/4}\le C\rho^{1/2}\),
so the third term is at most \(C\rho^{-1}\rho^{3/2}\rho^{1/2}=C\rho\).
\end{proof}

\begin{corollary}[The terminal datum lies in the critical Morrey space]
\label{cor:v47-terminal-morrey}
\(u(T_*)\) satisfies \(\sup_{x,\,\rho\le(\delta/2)^{1/2}}\rho^{-1}
\int_{B_\rho(x)}|u(T_*)|^2\le C_M\).  At every concentration point
\(x_0\in\Sigma\), by Corollary~\ref{cor:v46-critical-local-energy},
\begin{equation}
 0<\limsup_{\rho\to0}\frac1\rho\int_{B_\rho(x_0)}|u(T_*)|^2\le C_M ,
 \label{eq:v47-two-sided}
\end{equation}
while \(u(T_*)\notin L^3(B_\rho(x_0))\) for every \(\rho>0\)
(Theorem~\ref{thm:v46-non-L3}).
\end{corollary}

\begin{proof}
Theorem~\ref{thm:v47-morrey} at \(t=T_*\), and the cited results.
\end{proof}

\begin{assessment}[Local versus global critical size]
\label{ass:v47-local-vs-global}
Proposition~\ref{prop:v36-local-energies} bounded
\(\rho^{-1}\int_{B_\rho}|u(t)|^2\) by a square of the V33 logarithm,
through the global critical norm.  Theorem~\ref{thm:v47-morrey} shows the
logarithm is absent locally: every ball at every scale carries at most
\(C_M\rho\) of energy.  The global critical norm can still be logarithmic,
because \(\|u\|_{\dot H^{1/2}}^2\) is not a sum of local energies; the
V34 flat profile says that \(\log\)-many dyadic scales each carry
order-one critical mass globally, and that is compatible with order-one
local energy at every scale near every point.  Thus the local
\(\varepsilon\)-regularity quantities of V36--V37 are all bounded by
constants, not logarithms; they are simply not small.
\end{assessment}

% =============================================================================
\section{The final-time trace of the ancient limit}
\label{sec:v47-trace}
% =============================================================================

\begin{proposition}[Final-time trace in the critical Morrey class]
\label{prop:v47-trace}
Let \(x_0\in\Sigma\) with scales \(r_k\), and let
\(T_k(y):=r_ku(T_*,x_0+r_ky)\) be the rescaled terminal data.  Then
\(\int_{B_R}|T_k|^2\le C_MR\) for all \(k\) and \(R\), so along a
subsequence \(T_k\rightharpoonup U_0\) weakly in \(L^2_{\rm loc}(\R^3)\),
with \(\int_{B_R(y)}|U_0|^2\le C_MR\) for all \(y\), \(R\), and
\(\operatorname{div}U_0=0\).  Moreover the ancient limit \(U\) of
Theorem~\ref{thm:v41-ancient-limit} satisfies
\begin{equation}
 \int_{B_R(y)}|U(s)|^2\le C_M(R+(-s)^{1/2})\qquad(s<0),
 \label{eq:v47-limit-morrey}
\end{equation}
and \(e_R(s)\) converges as \(s\uparrow0\) for every \(R\); the
\(L^2_{\rm loc}\)-weak limit \(U(0^-)\) of \(U(s)\) exists and satisfies
the same Morrey bound.
\end{proposition}

\begin{proof}
\(\int_{B_R}|T_k|^2=r_k^{-1}\int_{B_{Rr_k}(x_0)}|u(T_*)|^2\le C_MR\) by
Corollary~\ref{cor:v47-terminal-morrey}.  Weak compactness and lower
semicontinuity give \(U_0\) with the same bound on every ball (translate
\(x_0\) to \(x_0+r_ky\)); the divergence condition is preserved.  For
\eqref{eq:v47-limit-morrey}, apply Theorem~\ref{thm:v47-morrey} to
\(u\) at time \(T_*+r_k^2s\) on the ball of radius
\(\rho=r_k\max\{R,(-s)^{1/2}\}\), rescale, and pass to the limit by
lower semicontinuity.  Convergence of \(e_R(s)\) as \(s\uparrow0\) follows
from the local energy identity for \(U\) on \((s_1,s_2)\), whose right
side is bounded by \(C_R\int_{s_1}^{s_2}(-\sigma)^{-3/4}\dd\sigma\) plus
the dissipation, so \(e_R\) is of bounded variation near \(0\); weak
compactness in \(L^2(B_R)\) then gives the trace.
\end{proof}

\begin{firewall}[The trace and the rescaled terminal data need not agree]
\label{fw:v47-trace-identification}
\(U(0^-)\) is the trace of the ancient limit; \(U_0\) is the weak limit of
the rescaled terminal data.  For each \(k\), \(U_k(s)\to T_k\) strongly in
\(L^2_{\rm loc}\) as \(s\uparrow0\), but the two limits \(k\to\infty\) and
\(s\uparrow0\) are exchanged only if the convergence is uniform in \(k\),
which the bound \eqref{eq:v46-local-energy-near-zero} gives for the
energies but not for the fields.  Whether \(U(0^-)=U_0\) is not decided.
Neither is proved nonzero: Corollary~\ref{cor:v46-critical-local-energy}
shows only that the local energies of \(T_k\) do not all tend to zero,
which weak limits may lose.
\end{firewall}

% =============================================================================
\section{Consequences for the earlier local bounds}
\label{sec:v47-consequences}
% =============================================================================

\begin{corollary}[All local scale-invariant quantities are bounded by
constants]
\label{cor:v47-local-constants}
Under \eqref{eq:v32-type-I}, for every \(x\), every
\(\rho\le(\delta/4)^{1/2}\), and every \(t\in[T_*-\delta/4,T_*]\),
\begin{equation}
 \frac1\rho\sup_{t-\rho^2\le\tau\le t}\int_{B_\rho(x)}|u(\tau)|^2
 +\frac1\rho\iint_{Q_\rho(x,t)}|\nabla u|^2
 +\frac1{\rho^2}\iint_{Q_\rho(x,t)}\bigl(|u|^3+|u||p-\bar p|\bigr)
 \ \le\ C_M+2\sqrt2\,C_I^2\nu^{3/2}+C_\sharp' ,
 \label{eq:v47-all-local}
\end{equation}
with \(C_\sharp'\) from the cubic computation in the proof of
Theorem~\ref{thm:v46-Linfty-upgrade} with \(c=1\).
\end{corollary}

\begin{proof}
The first term is Theorem~\ref{thm:v47-morrey}; the second is
\eqref{eq:v36-local-dissipation} extended to \(t=T_*\) by monotone
convergence; the third is the computation in the proof of
Theorem~\ref{thm:v46-Linfty-upgrade} at scale \(\rho\) with
\(\rho^2\le T_*-t+\rho^2\) in place of \(\theta/2\), which gives a bound
independent of \(t\) and \(\rho\).
\end{proof}

% =============================================================================
\section{Integrated V47 conclusion and next acquisition order}
\label{sec:v47-conclusion}
% =============================================================================

\begin{theorem}[What V47 proves]
\label{thm:v47-conclusion}
Relative to V1--V46: the uniform critical Morrey bound
\eqref{eq:v47-morrey} on the trajectory up to and including the singular
time; the critical Morrey membership and the two-sided scale-critical
local energy \eqref{eq:v47-two-sided} of the terminal datum at
concentration points; the Morrey bounds for the ancient limit and its
final-time trace (Proposition~\ref{prop:v47-trace}); and the constant
bounds \eqref{eq:v47-all-local} on all local scale-invariant quantities.
The full Navier--Stokes stability/global-regularity theorem remains
unproved.
\end{theorem}

\begin{proof}
The cited results.
\end{proof}

\begin{programme}[V48 acquisition order]
\label{prog:v48}
\begin{enumerate}[label=\textup{(AC\arabic*)},leftmargin=4.8em]
\item Identify the trace.  Determine whether \(U(0^-)=U_0\) by proving
      uniform-in-\(k\) convergence of \(U_k(s)\to T_k\) in
      \(L^2_{\rm loc}\) as \(s\uparrow0\), for instance through the
      \(H^{-1}_{\rm loc}\) time-derivative bound of the rescaled equation on
      \((-1,0)\), which the Morrey bounds and the \(L^3\) pressure bound
      make uniform.
\item Nontriviality of the trace.  Combine the Morrey bound (no point
      concentration of \(T_k\)) with
      Corollary~\ref{cor:v46-critical-local-energy}; determine whether the
      failure of strong convergence must be oscillatory and whether the
      backward-uniqueness package can be run with a nonzero trace of
      critical Morrey type by first subtracting \(U_0\).
\item Type-II: item (AC2) of Programme~\ref{prog:v41} remains.
\item The next compulsory companion audit is V50.
\end{enumerate}
\end{programme}

\begin{assessment}[Position after V47]
\label{ass:v47-position}
The type-I regime has, locally, no logarithms at all: every scale-critical
local quantity is bounded by a constant, the terminal datum is of critical
Morrey type and non-\(L^3\) at the singular points, and the singular set
has dimension zero.  The logarithm survives only in the global critical
norm.  The Liouville gate of V41 remains.  No full stability or
global-regularity claim is made.
\end{assessment}

% =============================================================================
\section*{Version V47 append-only record}
\addcontentsline{toc}{section}{Version V47 append-only record}
% =============================================================================

The complete V46 mathematical source was retained before this continuation.
Its SHA--256 digest was
\begin{center}\scriptsize\ttfamily
0c6329efac2a2cea6d6d48c7c21fee4ae72b52a5cc1bc88f35d9bebcbef9be42.
\end{center}
V47 proved the uniform critical Morrey bound up to the singular time, the
two-sided scale-critical local energy of the terminal datum at
concentration points, the Morrey bounds for the ancient limit and its
trace, and constant bounds on all local scale-invariant quantities.  No
full regularity claim is made.

% =============================================================================
\section{V48: identification and nontriviality of the final-time trace}
\label{sec:v48-entry}
% =============================================================================

Items (AC1)--(AC2) of Programme~\ref{prog:v48} are settled.  The rescaled
equation, with the Morrey and pressure bounds of V47, gives a uniform
H\"older bound in \(H^{-1}_{\rm loc}\) for the rescaled fields near the
final time, which identifies the trace \(U(0^-)\) of the ancient limit
with the weak limit \(U_0\) of the rescaled terminal data.  The
backward-uniqueness argument of V46, which used only the vanishing of the
trace on an exterior region, then shows that the trace does not vanish on
any exterior region: its vorticity is nonzero arbitrarily far from the
origin.  The final-time trace of the ancient limit of a type-I singularity
is therefore a nontrivial divergence-free field of critical Morrey type
that does not decay to zero in the weak sense at spatial infinity.  No
global regularity claim is made.  The next compulsory companion audit is
V50.

% =============================================================================
\section{Uniform H\"older continuity in \(H^{-1}_{\rm loc}\) near the final
time}
\label{sec:v48-Hminus1}
% =============================================================================

Let \(x_0\in\Sigma\), \(r_k\), \(U_k\), \(P_k\), \(U\), \(P\), and
\(T_k=U_k(0^-)\) be as in V41 and V47.  For a ball \(B_R\subset\R^3\)
write \(\|\cdot\|_{H^{-1}(B_R)}\) for the dual norm of \(H^1_0(B_R)\).

\begin{lemma}[Uniform \(H^{-1}_{\rm loc}\) bound on the time increment]
\label{lem:v48-Hminus1}
There is \(C_R\), independent of \(k\), such that for \(-1\le s<0\),
\begin{equation}
 \|U_k(s)-T_k\|_{H^{-1}(B_R)}\le C_R(-s)^{1/2}.
 \label{eq:v48-Hminus1}
\end{equation}
\end{lemma}

\begin{proof}
For \(\psi\in H^1_0(B_R;\R^3)\) and \(s<s'<0\), the equation gives
\[
 \langle U_k(s')-U_k(s),\psi\rangle
 =\int_s^{s'}\Bigl(-\nu\langle\nabla U_k,\nabla\psi\rangle
 +\langle U_k\otimes U_k,\nabla\psi\rangle
 +\langle P_k,\operatorname{div}\psi\rangle\Bigr)\dd\sigma ,
\]
because \(\PP\operatorname{div}(U\otimes U)=\operatorname{div}(U\otimes U)
+\nabla P\) and \(\psi\) is compactly supported.  Bound the three integrands
by \(\nu\|\nabla U_k(\sigma)\|_{L^2(B_R)}\|\nabla\psi\|_2\),
\(\|U_k(\sigma)\|_{L^4(B_R)}^2\|\nabla\psi\|_2\le|B_R|^{1/6}
\|U_k(\sigma)\|_6^2\|\nabla\psi\|_2\), and
\(|B_R|^{1/6}\|P_k(\sigma)\|_3\|\nabla\psi\|_2\).  By
\eqref{eq:v41-rescaled-type-I}--\eqref{eq:v41-rescaled-L6} these are at
most \(C(-\sigma)^{-1/4}\), \(C_R(-\sigma)^{-1/2}\), and
\(C_R(-\sigma)^{-1/2}\) times \(\|\nabla\psi\|_2\), all integrable on
\((s,0)\) with integral at most \(C_R(-s)^{1/2}\).  Let \(s'\uparrow0\),
using the strong \(L^2\) convergence \(U_k(s')\to T_k\), which holds
because \(u(t)\to u(T_*)\) strongly in \(L^2(\Torus^3)\).
\end{proof}

\begin{theorem}[The trace of the ancient limit is the weak limit of the
rescaled terminal data]
\label{thm:v48-trace-identification}
Along the subsequence of Proposition~\ref{prop:v47-trace},
\begin{equation}
 U(s)\longrightarrow U_0\quad\text{in }H^{-1}_{\rm loc}(\R^3)
 \text{ and weakly in }L^2_{\rm loc}(\R^3)\quad\text{as }s\uparrow0,
 \qquad
 \|U(s)-U_0\|_{H^{-1}(B_R)}\le C_R(-s)^{1/2},
 \label{eq:v48-trace-identification}
\end{equation}
so \(U(0^-)=U_0\).  Moreover \(U_0=U(\cdot,0)\) almost everywhere on the
exterior region \(\{|y|\ge R_0\}\) of Lemma~\ref{lem:v46-bounded}, where
\(U\) is bounded and continuous up to \(s=0\).
\end{theorem}

\begin{proof}
For fixed \(s<0\), \(U_k(s)\to U(s)\) strongly in \(L^2(B_R)\), hence in
\(H^{-1}(B_R)\).  Since \(T_k\rightharpoonup U_0\) weakly in \(L^2(B_R)\)
and \(L^2(B_R)\hookrightarrow H^{-1}(B_R)\) compactly,
\(T_k\to U_0\) strongly in \(H^{-1}(B_R)\).  Then
\[
 \|U(s)-U_0\|_{H^{-1}(B_R)}
 \le\|U(s)-U_k(s)\|_{H^{-1}(B_R)}+\|U_k(s)-T_k\|_{H^{-1}(B_R)}
 +\|T_k-U_0\|_{H^{-1}(B_R)},
\]
and letting \(k\to\infty\) with Lemma~\ref{lem:v48-Hminus1} gives the
bound.  Since \(U(s)\) is bounded in \(L^2(B_R)\) by
\eqref{eq:v47-limit-morrey}, every weak \(L^2(B_R)\) limit point as
\(s\uparrow0\) equals the \(H^{-1}\) limit \(U_0\), so
\(U(s)\rightharpoonup U_0\).  On the exterior region \(U(s)\) converges
locally uniformly to the continuous extension \(U(\cdot,0)\) by
Lemma~\ref{lem:v46-bounded}(i), and uniform limits agree with weak limits.
\end{proof}

% =============================================================================
\section{The trace does not vanish on any exterior region}
\label{sec:v48-nontrivial}
% =============================================================================

\begin{theorem}[Nontriviality of the trace at spatial infinity]
\label{thm:v48-trace-nontrivial}
Modulo the imports (E1)--(E6), for every \(R\ge R_0\),
\begin{equation}
 \boxed{\quad
 \operatorname{curl}U_0\not\equiv0\ \text{ on }\{|y|\ge R\},
 \qquad\text{in particular}\qquad
 U_0\not\equiv0\ \text{ on }\{|y|\ge R\}.
 \quad}
 \label{eq:v48-trace-nontrivial}
\end{equation}
\end{theorem}

\begin{proof}
Fix \(R\ge R_0\).  By Lemma~\ref{lem:v46-bounded}(i), \(U\), \(\nabla U\),
\(\nabla^2U\) are bounded on \(\{|y|\ge R\}\times(-1/16,0)\) and \(U\)
extends continuously to \(s=0\) there, with \(U(\cdot,0)=U_0\) by
Theorem~\ref{thm:v48-trace-identification}.  Suppose
\(\operatorname{curl}U_0=0\) on \(\{|y|\ge R\}\).  Then
\(\omega=\operatorname{curl}U\) is smooth and bounded with bounded gradient
on \(\{|y|\ge R\}\times(-1/16,0)\), satisfies the vorticity inequality with
bounded coefficients there, and vanishes at \(s=0\).  The proof of
Theorem~\ref{thm:v46-non-L3} from this point on used exactly these facts
with \(R_0\) in place of \(R\): (E6) gives \(\omega\equiv0\) on the
exterior region for \(s\in(-1/16,0)\); Lemma~\ref{lem:v46-bounded}(ii)
and the second part of (E6) give \(\omega\equiv0\) on
\(\R^3\times(-1/16,-\eta)\) for every \(\eta\), hence \(U\equiv0\) on
\(\R^3\times(-1/16,0)\) since harmonic fields in \(L^6\) vanish; and the
backward iteration gives \(U\equiv0\) on \(\R^3\times(-1,0)\),
contradicting \eqref{eq:v41-limit-nontrivial}.  Hence
\(\operatorname{curl}U_0\not\equiv0\) on \(\{|y|\ge R\}\), and a
curl-free assertion fails a fortiori for \(U_0\) itself.
\end{proof}

\begin{corollary}[The rescaled terminal data do not converge weakly to zero
on any exterior region]
\label{cor:v48-terminal-data}
For every \(R\ge R_0\) and every subsequence along which
\(T_k\rightharpoonup U_0\), the restriction of \(U_0\) to
\(\{|y|\ge R\}\) is nonzero.  Consequently, for every \(R\ge R_0\) there
are \(R'>R\) and \(c>0\) with
\begin{equation}
 \liminf_{k\to\infty}\ \frac1{r_k}\int_{Rr_k\le|x-x_0|\le R'r_k}|u(T_*,x)|^2\,\dd x
 \ \ge\ c ,
 \label{eq:v48-annular-energy}
\end{equation}
so the terminal datum carries energy of scale-critical size in every
sufficiently wide dyadic-type annulus around \(x_0\), at all small scales
along the concentration sequence.
\end{corollary}

\begin{proof}
The first statement is Theorem~\ref{thm:v48-trace-nontrivial} applied
along the subsequence (the proof did not depend on which subsequence
realizes \(U_0\)).  For the second, \(U_0\ne0\) on \(\{|y|\ge R\}\) gives
\(R'\) and \(c\) with \(\int_{R\le|y|\le R'}|U_0|^2\ge2c\), and weak lower
semicontinuity gives \(\liminf_k\int_{R\le|y|\le R'}|T_k|^2\ge2c\), which is
\eqref{eq:v48-annular-energy} after the change of variables.
\end{proof}

\begin{assessment}[The terminal datum at a singular point]
\label{ass:v48-terminal-datum}
V46 showed that \(u(T_*)\) is not cubically integrable near \(x_0\); V47
showed that its local energy is at most \(C_M\rho\) in every ball; V48
shows that along the concentration scales it has energy at least
\(c\,r_k\) in annuli of every fixed rescaled width around \(x_0\), and that
its rescaled weak limit has vorticity at arbitrarily large rescaled
distance.  In the scale-invariant picture this is the statement that the
terminal datum looks, at the concentration scales, like a field of size
\(|x-x_0|^{-1}\) with nonzero vorticity in every annulus: it is neither
smoother nor more concentrated than a homogeneous \(-1\) profile.  The
picture is fully consistent with a self-similar singularity and does not
exclude one.
\end{assessment}

\begin{firewall}[What remains open about the trace]
\label{fw:v48-open}
It is not proved that \(T_k\to U_0\) strongly in \(L^2_{\rm loc}\), that
\(U_0\) is \(-1\)-homogeneous, that \(U_0\) is smooth away from the
origin, or that \(U_0\) determines the ancient limit.  The nontriviality
proved here is at spatial infinity of the rescaled picture, that is, at
distances from \(x_0\) large compared with \(r_k\) but small compared with
\(1\); it says nothing about \(U_0\) near the origin.
\end{firewall}

% =============================================================================
\section{Integrated V48 conclusion and next acquisition order}
\label{sec:v48-conclusion}
% =============================================================================

\begin{theorem}[What V48 proves]
\label{thm:v48-conclusion}
Relative to V1--V47, modulo (E1)--(E6): the uniform
\(H^{-1}_{\rm loc}\) H\"older bound \eqref{eq:v48-Hminus1}; the
identification \(U(0^-)=U_0\) of the trace with the weak limit of the
rescaled terminal data (Theorem~\ref{thm:v48-trace-identification}); the
nontriviality of the trace and of its vorticity on every exterior region
(Theorem~\ref{thm:v48-trace-nontrivial}); and the scale-critical annular
energy of the terminal datum around every concentration point
(Corollary~\ref{cor:v48-terminal-data}).  The full Navier--Stokes
stability/global-regularity theorem remains unproved.
\end{theorem}

\begin{proof}
The cited results.
\end{proof}

\begin{programme}[V49 acquisition order]
\label{prog:v49}
\begin{enumerate}[label=\textup{(AC\arabic*)},leftmargin=4.8em]
\item Regularity of the trace away from the origin.  On the exterior
      region \(U_0=U(\cdot,0)\) is the continuous trace of a bounded
      solution; determine whether \(U_0\) is smooth on \(\{|y|\ge R_0\}\)
      and whether the exterior boundedness radius \(R_0\) can be made
      uniform along a further rescaling of \(U\) itself.
\item Homogeneity.  Rescale the ancient limit \(U\) by
      \(U^\lambda(s,y)=\lambda U(\lambda^2s,\lambda y)\); the class
      \(\mathcal A_\nu(C)\) and all V47--V48 bounds are invariant.
      Determine whether a limit as \(\lambda\to\infty\) or \(\lambda\to0\)
      exists and is a discretely or continuously self-similar element of
      the class; a nonzero such limit would be a self-similar ancient
      solution with critical Morrey trace.
\item Type-II: item (AC2) of Programme~\ref{prog:v41} remains.
\item The next compulsory companion audit is V50.
\end{enumerate}
\end{programme}

\begin{assessment}[Position after V48]
\label{ass:v48-position}
The ancient limit of a type-I singularity is now known to have a
nontrivial final-time trace of critical Morrey type with vorticity out to
spatial infinity, and the terminal datum of the original solution is
correspondingly of scale-critical size in every annulus around each
concentration point.  The Liouville gate of V41 remains.  No full
stability or global-regularity claim is made.
\end{assessment}

% =============================================================================
\section*{Version V48 append-only record}
\addcontentsline{toc}{section}{Version V48 append-only record}
% =============================================================================

The complete V47 mathematical source was retained before this continuation.
Its SHA--256 digest was
\begin{center}\scriptsize\ttfamily
b283fded689659a6922c993e1e94e70fcf8bf66b4c924ce437c7c99fc3403f78.
\end{center}
V48 identified the final-time trace of the ancient limit with the weak
limit of the rescaled terminal data and proved that this trace, and its
vorticity, are nonzero on every exterior region, with the corresponding
annular energy lower bound for the terminal datum.  No full regularity
claim is made.

% =============================================================================
\section{V49: regularity and decay of the trace, and the scaling
dichotomies of the ancient limit}
\label{sec:v49-entry}
% =============================================================================

Item (AC1) of Programme~\ref{prog:v49} is settled: on the exterior region
the final-time trace of the ancient limit is smooth and tends to zero at
spatial infinity, by the quantitative form of \(\varepsilon\)-regularity.
Item (AC2) is analysed but not settled.  The ancient class and every bound
of V47--V48 are invariant under the Navier--Stokes scaling; rescaling the
ancient limit itself gives two dichotomies, at large scales and at small
scales, each of which either produces a further nontrivial element of the
class or forces a dispersion statement about the dissipation of \(U\).
Neither branch is excluded.  No global regularity claim is made.  The
compulsory companion audit is due in V50.

% =============================================================================
\section{Regularity and decay of the trace on the exterior}
\label{sec:v49-trace-regularity}
% =============================================================================

\begin{firewall}[Quantitative form of the import (E5)]
\label{fw:v49-import}
\begin{enumerate}[label=\textup{(E\arabic*)},leftmargin=3.6em,start=7]
\item The \(\varepsilon\)-regularity theorem (E5) holds in the
      quantitative form: if
      \(\Phi_r(z):=r^{-2}\iint_{Q_r(z)}(|u|^3+|u||p-\bar p|)\le\varepsilon_2\),
      then \(\sup_{Q_{r/2}(z)}|u|\le C_{\rm CKN}\,\Phi_r(z)^{1/3}r^{-1}\),
      and all derivatives of \(u\) are bounded on \(Q_{r/4}(z)\) by
      constants depending on \(\Phi_r(z)\), \(r\), \(\nu\), and the order.
\end{enumerate}
This is the form in which the theorem is proved (the iteration bounds the
cube of the velocity by the scale-invariant quantity), and it is used
only for the qualitative decay below.
\end{firewall}

\begin{proposition}[The trace is smooth and decays on the exterior region]
\label{prop:v49-trace-smooth-decay}
Let \(R_0\) be as in Lemma~\ref{lem:v46-bounded}.  Then
\(U_0=U(\cdot,0)\) is smooth on \(\{|y|>R_0\}\), and
\begin{equation}
 \lim_{|y|\to\infty}\Bigl(|U_0(y)|+|\nabla U_0(y)|\Bigr)=0 .
 \label{eq:v49-trace-decay}
\end{equation}
Consequently \(U_0\) is a smooth divergence-free field on the exterior
region, decaying at infinity together with its gradient, with nonzero
vorticity on every exterior region by Theorem~\ref{thm:v48-trace-nontrivial}.
\end{proposition}

\begin{proof}
For \(|y|\ge R_0\), \(U\) is a bounded solution on \(Q_{1/2}(y,0)\) by
Lemma~\ref{lem:v46-bounded}(i), hence smooth up to \(s=0\) on
\(Q_{1/4}(y,0)\) with all derivatives bounded (E7), so \(U_0\) is smooth on
\(\{|y|>R_0\}\).  For the decay, the proof of
Lemma~\ref{lem:v46-bounded}(i) showed \(F(y)=\Phi_1((y,0))\to0\) as
\(|y|\to\infty\); (E7) then gives \(\sup_{Q_{1/2}(y,0)}|U|\le
C_{\rm CKN}F(y)^{1/3}\to0\), and the derivative bounds on
\(Q_{1/4}(y,0)\), which depend continuously on \(F(y)\) and vanish with it
by the same iteration applied to the equation for \(\nabla U\), give the
gradient decay.  Continuity up to \(s=0\) transfers both to \(U_0\).
\end{proof}

\begin{assessment}[The trace at infinity]
\label{ass:v49-trace-at-infinity}
\(U_0\) decays pointwise at infinity but is nonzero on every exterior
region and has critical Morrey size \(\int_{B_R(y)}|U_0|^2\le C_MR\).  A
field with all three properties is, for instance, any smooth field of size
\(|y|^{-1}\) with nonvanishing curl in every annulus.  The pointwise decay
rate is not proved; the Morrey bound permits \(|U_0(y)|\) to decay as
slowly as \(|y|^{-1}\) in an averaged sense and forbids nothing faster.
\end{assessment}

% =============================================================================
\section{Scaling invariance and the two dichotomies}
\label{sec:v49-dichotomies}
% =============================================================================

For \(\lambda>0\) put \(U^\lambda(s,y):=\lambda U(\lambda^2s,\lambda y)\).

\begin{lemma}[Invariance of the class and of the V47--V48 bounds]
\label{lem:v49-invariance}
If \(U\in\mathcal A_\nu(C)\), then \(U^\lambda\in\mathcal A_\nu(C)\) for
every \(\lambda>0\); the critical Morrey bound \eqref{eq:v47-limit-morrey}
and the trace bounds of Proposition~\ref{prop:v47-trace} hold for
\(U^\lambda\) with the same constants; and the trace of \(U^\lambda\) is
\(\lambda U_0(\lambda\cdot)\).  The unit-scale dissipation
\(\iint_{B_1\times(-1,0)}|\nabla U^\lambda|^2\) equals
\(\lambda^{-1}\iint_{B_\lambda\times(-\lambda^2,0)}|\nabla U|^2\).
\end{lemma}

\begin{proof}
Direct computation, as in Remark~\ref{rem:v33-scaling}.
\end{proof}

Define the scale-\(\lambda\) dissipation concentration of \(U\) at the
origin and its supremum over centres:
\begin{equation}
 d_U(\lambda):=\frac1\lambda\iint_{B_\lambda(0)\times(-\lambda^2,0)}|\nabla U|^2,
 \qquad
 d_U^*(\lambda):=\sup_{y\in\R^3}\frac1\lambda\iint_{B_\lambda(y)\times(-\lambda^2,0)}|\nabla U|^2
 \le2C .
 \label{eq:v49-concentration-function}
\end{equation}

\begin{theorem}[Large-scale dichotomy]
\label{thm:v49-large-scale}
Let \(U\in\mathcal A_\nu(C)\) be the ancient limit of
Theorem~\ref{thm:v41-ancient-limit}.  Exactly one of the following holds.
\begin{enumerate}[label=\textup{(L\arabic*)},leftmargin=3.8em]
\item \(\limsup_{\lambda\to\infty}d_U^*(\lambda)>0\).  Then there are
      \(\lambda_j\to\infty\) and centres \(y_j\) such that the translated
      rescalings \(\lambda_jU(\lambda_j^2s,\lambda_j y+y_j)\) converge, in
      the sense of Theorem~\ref{thm:v41-ancient-limit}, to a nontrivial
      element \(U^{(\infty)}\in\mathcal A_\nu(C)\) which satisfies all
      V47--V48 bounds and has unit-scale dissipation at least
      \(\tfrac12\limsup_\lambda d_U^*(\lambda)\).
\item \(d_U^*(\lambda)\to0\) as \(\lambda\to\infty\): the dissipation of
      \(U\) in every parabolic cylinder of size \(\lambda\) is
      \(o(\lambda)\) uniformly in the centre, although in case (B) of
      Theorem~\ref{thm:v42-trichotomy} the total slab dissipation
      \(\int_{-\lambda^2}^0Y_U\) is at least \(2(\nu^3/(8C_1))^{1/2}\lambda\).
\end{enumerate}
\end{theorem}

\begin{proof}
If (L1), choose \(\lambda_j\) and near-maximizing centres \(y_j\).  The
translated rescalings are in \(\mathcal A_\nu(C)\) with the same uniform
bounds as the sequence \(U_k\) of Proposition~\ref{prop:v41-uniform-bounds}
(those bounds used only the enstrophy rate), so the compactness argument
of Theorem~\ref{thm:v41-ancient-limit} applies verbatim and gives a limit
in the class; nontriviality follows as in \eqref{eq:v41-limit-nontrivial}
from the unit-scale dissipation of the rescalings, which is
\(d_U^*(\lambda_j)(1-o(1))\), after removing the part of the cylinder near
\(s=0\), which carries at most \(2C\eta^{1/2}\).  The Morrey and trace
bounds pass to the limit by lower semicontinuity, and the nontriviality of
the trace on exterior regions follows by repeating V46--V48 for the new
limit.  Otherwise (L2) holds, and the slab bound is the backward cone
\eqref{eq:v42-two-sided} integrated.
\end{proof}

\begin{theorem}[Small-scale dichotomy]
\label{thm:v49-small-scale}
Let \(U\) be as above and suppose \(U\) is in case (B).  Exactly one of the
following holds.
\begin{enumerate}[label=\textup{(S\arabic*)},leftmargin=3.8em]
\item \(\limsup_{\lambda\to0}d_U^*(\lambda)>0\).  Then, as in (L1) with
      \(\lambda_j\to0\), a nontrivial element \(U^{(0)}\in\mathcal A_\nu(C)\)
      is obtained, and \(U\) has a singular point in the sense that its
      velocity is unbounded near \((y_j,0)\) along a bounded sequence
      \(y_j\) or the concentration escapes to infinity.
\item \(d_U^*(\lambda)\to0\) as \(\lambda\to0\).  Then every point
      \((y,0)\) is regular for \(U\) by the dissipation criterion (E1), so
      \(U\) is smooth in a neighbourhood of every point of
      \(\R^3\times\{0\}\), and its enstrophy divergence
      \(Y_U(s)\to\infty\) is carried entirely by spatial infinity: for every
      \(R\), \(\int_{B_R}|\nabla U(s)|^2\) stays bounded as \(s\uparrow0\).
\end{enumerate}
\end{theorem}

\begin{proof}
(S1) is the small-scale version of (L1), with the same compactness.
Under (S2), for every \(y\), \(\limsup_{\lambda\to0}\lambda^{-1}
\iint_{Q_\lambda(y,0)}|\nabla U|^2\le\lim d_U^*(\lambda)=0<\varepsilon_1\),
so (E1) makes \((y,0)\) a regular point; local boundedness near each point
and compactness of \(\overline{B_R}\) give bounded enstrophy on \(B_R\)
up to \(s=0\), while \(Y_U(s)\to\infty\) in case (B).
\end{proof}

\begin{assessment}[What the dichotomies leave open]
\label{ass:v49-dichotomies}
Item (AC2) asked for a self-similar limit.  The scaling limits of (L1) and
(S1) exist and lie in the class, but nothing proved here makes them
self-similar: the rescalings \(\lambda_j\) are a sequence, not a
continuum, and the class is not known to be compact under the full scaling
group.  The alternative branches (L2) and (S2) are dispersion statements,
in which the dissipation of \(U\) at large or small scales spreads out
rather than concentrating; (S2) additionally forces the enstrophy blow-up
of \(U\) to escape to spatial infinity.  None of the four branches is
excluded, and none contradicts the Liouville gate; they refine its
statement into cases.
\end{assessment}

% =============================================================================
\section{Integrated V49 conclusion and next acquisition order}
\label{sec:v49-conclusion}
% =============================================================================

\begin{theorem}[What V49 proves]
\label{thm:v49-conclusion}
Relative to V1--V48, modulo (E1)--(E7): the trace of the ancient limit is
smooth and decays with its gradient on the exterior region
(Proposition~\ref{prop:v49-trace-smooth-decay}); the class and all V47--V48
bounds are scaling invariant (Lemma~\ref{lem:v49-invariance}); and the
large-scale and small-scale dichotomies
(Theorems~\ref{thm:v49-large-scale}--\ref{thm:v49-small-scale}).  The full
Navier--Stokes stability/global-regularity theorem remains unproved.
\end{theorem}

\begin{proof}
The cited results.
\end{proof}

\begin{programme}[V50 acquisition order, with compulsory companion audit]
\label{prog:v50}
\begin{enumerate}[label=\textup{(AC\arabic*)},leftmargin=4.8em]
\item Perform the compulsory review of the current SM and YM notes;
      record digests; import nothing numerical.
\item Consolidate V41--V49 into one theorem describing the ancient limit
      of a type-I singularity: class, trichotomy, Morrey bounds, trace
      identification, nontriviality and decay of the trace, and the four
      scaling branches; state the Liouville gate in its final form.
\item Branch (S2).  Determine whether an element of the class in case (B)
      can have its enstrophy blow-up at spatial infinity while its trace is
      of critical Morrey type with vorticity on every exterior region;
      test against the backward cone applied to the local enstrophy on
      exterior regions, where \(U\) is bounded.
\item Type-II: item (AC2) of Programme~\ref{prog:v41} remains.
\end{enumerate}
\end{programme}

\begin{assessment}[Position after V49]
\label{ass:v49-position}
The ancient limit of a type-I singularity has a trace that is smooth,
decaying, and vortical at infinity, and its behaviour at large and small
scales splits into concentration branches, which reproduce the class, and
dispersion branches, which are constrained but not excluded.  No full
stability or global-regularity claim is made.
\end{assessment}

% =============================================================================
\section*{Version V49 append-only record}
\addcontentsline{toc}{section}{Version V49 append-only record}
% =============================================================================

The complete V48 mathematical source was retained before this continuation.
Its SHA--256 digest was
\begin{center}\scriptsize\ttfamily
0f206b06f6120ed12ecfb24699b80c8874d7e39746d020832bbda71039d6bb08.
\end{center}
V49 proved the smoothness and decay of the final-time trace on the exterior
region, the scaling invariance of the ancient class with its V47--V48
bounds, and the large-scale and small-scale dichotomies for the ancient
limit.  No full regularity claim is made.

% =============================================================================
\section{V50: compulsory companion audit and the consolidated description of
the type-I ancient limit}
\label{sec:v50-entry}
% =============================================================================

This is a twentieth-version continuation.  The compulsory review of the
two companion programmes is performed first.  The version then
consolidates V41--V49 into one theorem describing the ancient limit of a
type-I singularity and the final form of the Liouville gate, and analyses
the dispersion branch (S2) of V49, in which the enstrophy blow-up of the
ancient limit escapes to spatial infinity: in that branch the limit is
globally bounded with bounded derivatives up to its final time, and its
divergent enstrophy is carried by a region of volume at least a constant
times \((-s)^{-1/2}\).  Neither this nor any other branch is excluded.
No global regularity claim is made.  The next compulsory companion audit
is V55.

% =============================================================================
\section{Compulsory V50 review of the current companion notes}
\label{sec:v50-companion-audit}
% =============================================================================

The two current companion files in \texttt{latex\_notes} were inspected
read-only.  The frozen checkpoints are
\begin{align}
 &\texttt{Renewal\_SM\_Einstein\_Continuum\_Research\_Notes\_V67.tex},
 \notag\\[-2pt]
 &\qquad
 \texttt{147f3d120b0d03e506676d5f9dcbfbacbb0029925a3fe29d91da8508b0297cdb},
 \label{eq:v50-SM-checkpoint}\\
 &\texttt{Renewal\_Yang\_Mills\_Mass\_Gap\_Research\_Notes\_V48.tex},
 \notag\\[-2pt]
 &\qquad
 \texttt{180af15cee5cc06df2646b7f6e007f0f89db24210974f99aa2573c4303d5f38c}.
 \label{eq:v50-YM-checkpoint}
\end{align}
The review concentrated on SM V66--V67 and YM V46--V48.

\emph{SM V66--V67.}  V66 proves that the improved free-scalar
two-particle stress vertex is transverse with a strictly positive
scheme-independent spectral submeasure at high momentum, so that
scalar-constraint projection yields a Newton witness growing at least like
the sixth power of the cutoff, and concludes that a continuum source card
must be distributional with spatial test order at least \(H^{3/2}\).  V67
factors the improved scalar energy numerator so that the Newton pole
cancels before norms, obtains a finite positive spacetime-smeared source
with an explicit Osterwalder--Schrader form on separated tests, and closes
the free separated distributional Newton card, leaving the sharp boundary
trace and the dynamical-gravity transfer open.

\emph{YM V46--V48.}  These versions extend calibrated Gaussian moment
comparisons through degree six, decompose matrix increments into odd and
even parts with uniform thermal rows, prove vanishing of a Gaussian
connected tensor and the \(s\)-order of a calibrated residual, and close
the quintic and six-input unmarked and marked matrix cards, isolating the
repeated density-root degree census as the next step.

\begin{theorem}[V50 companion-audit decision]
\label{thm:v50-companion-decision}
No SM V66--V67 stress, Newton, or Osterwalder--Schrader constant and no YM
V46--V48 moment or matrix-card estimate is imported.  The SM practice of
cancelling a pole before taking norms is the practice of V33 here, where
heat weights were kept inside the dyadic sum; nothing further transfers.
Neither companion bears on the Liouville gate of V41 or on the type-II
ledger.
\end{theorem}

\begin{proof}
The hypotheses of every cited result are free-field, lattice-gauge, or
Gaussian objects; none is a Navier--Stokes balance.
\end{proof}

% =============================================================================
\section{The ancient limit of a type-I singularity: consolidated statement}
\label{sec:v50-consolidation}
% =============================================================================

\begin{theorem}[The ancient limit of a type-I singularity, V41--V49]
\label{thm:v50-ancient-limit}
Let \(u\) be a smooth mean-zero solution on \([0,T_*)\times\Torus^3\) with
maximal time \(T_*<\infty\) and a type-I enstrophy rate
\eqref{eq:v32-type-I}.  Modulo the imports (E1)--(E7), the following
hold.
\begin{enumerate}[label=\textup{(\arabic*)},leftmargin=3.2em]
\item \emph{Velocity rate and singular set.}  \(|u(x,t)|\le C_\infty(T_*-t)^{-1/2}\)
      and \(|\nabla u|\le C_\infty'(T_*-t)^{-1}\); the singular set
      \(\Sigma\subset\Torus^3\) is closed, nonempty, of Hausdorff dimension
      zero, and at most \(N_*\) points of it are active at any one scale.
\item \emph{Terminal datum.}  \(u(t)\to u(T_*)\) strongly in \(L^2\) with
      energy equality; \(u(T_*)\) is smooth off \(\Sigma\), satisfies the
      critical Morrey bound \(\int_{B_\rho(x)}|u(T_*)|^2\le C_M\rho\)
      everywhere, is not in \(L^3\) near any concentration point, hence not
      in \(H^1\), and carries scale-critical energy in annuli around every
      concentration point along the concentration scales.  The trajectory
      glues to a Leray--Hopf solution strongly continuous at \(T_*\) which
      cannot restart as a strong solution.
\item \emph{Ancient limit.}  Rescaling at a concentration point produces
      a nontrivial smooth ancient solution \(U\in\mathcal A_\nu(C_I^2\nu^{3/2})\)
      on \((-\infty,0)\times\R^3\) with infinite energy, uniform
      palinstrophy averages, the critical Morrey bound
      \(\int_{B_R(y)}|U(s)|^2\le C_M(R+(-s)^{1/2})\), and at least the
      fraction \(c_*/2\) of its unit-slab dissipation in the unit ball;
      it is either eternal with enstrophy bounded on every half-line, or
      two-sided type-I singular at \(s=0\).
\item \emph{Trace.}  \(U(s)\to U_0\) in \(H^{-1}_{\rm loc}\) and weakly in
      \(L^2_{\rm loc}\) as \(s\uparrow0\), where \(U_0\) is the weak limit
      of the rescaled terminal data; \(U_0\) is of critical Morrey type,
      smooth and decaying with its gradient on an exterior region, and its
      vorticity is nonzero on every exterior region.
\item \emph{Scaling branches.}  At large scales \(U\) either concentrates
      dissipation, producing a further nontrivial element of the class, or
      disperses it, \(d_U^*(\lambda)\to0\); at small scales, in the
      singular case, it either concentrates, producing a further element,
      or disperses, in which case every point of \(\R^3\times\{0\}\) is
      regular and the enstrophy blow-up escapes to spatial infinity.
\end{enumerate}
The Liouville gate is: \emph{every nonzero element of
\(\mathcal A_\nu(C)\) with the properties in (3)--(5) is excluded}.  Its
finite-energy case is proved (Theorem~\ref{thm:v41-liouville}); its
stationary case is Leray's problem for \(L^6\cap\dot H^1\) solutions
(Firewall~\ref{fw:v42-Leray-problem}).  If the gate holds, no smooth
solution on the torus has a singularity with a type-I enstrophy rate.
\end{theorem}

\begin{proof}
(1) is Theorems~\ref{thm:v46-Linfty-upgrade} and
\ref{thm:v43-dimension-zero} with Lemma~\ref{lem:v43-few-active}.  (2)
is Theorem~\ref{thm:v32-type-I-firewall}(i),
Proposition~\ref{prop:v45-energy-equality},
Theorem~\ref{thm:v44-localization}, Theorem~\ref{thm:v47-morrey},
Theorem~\ref{thm:v46-non-L3}, Corollary~\ref{cor:v48-terminal-data}, and
Theorem~\ref{thm:v45-glued-continuation}.  (3) is
Theorem~\ref{thm:v41-ancient-limit}, Corollary~\ref{cor:v41-infinite-energy},
Propositions~\ref{prop:v42-palinstrophy} and \ref{prop:v47-trace},
Corollary~\ref{cor:v43-limit-fraction}, and
Theorem~\ref{thm:v42-trichotomy}.  (4) is
Theorems~\ref{thm:v48-trace-identification} and
\ref{thm:v48-trace-nontrivial} with
Proposition~\ref{prop:v49-trace-smooth-decay}.  (5) is
Theorems~\ref{thm:v49-large-scale}--\ref{thm:v49-small-scale}.  The gate
statement is Theorem~\ref{thm:v41-gate}.
\end{proof}

% =============================================================================
\section{The dispersion branch (S2)}
\label{sec:v50-S2}
% =============================================================================

\begin{proposition}[Structure of the ancient limit in branch (S2)]
\label{prop:v50-S2}
Suppose \(U\) is in case (B) and branch (S2) of
Theorem~\ref{thm:v49-small-scale}.  Then there is \(s_0<0\) such that
\(U\), \(\nabla U\), and \(\nabla^2U\) are bounded on
\(\R^3\times[s_0,0]\), \(U\) is smooth up to \(s=0\) on all of \(\R^3\),
\(U_0=U(\cdot,0)\) is a smooth divergence-free field on \(\R^3\) of
critical Morrey type with vorticity on every exterior region, and for
\(s\in[s_0,0)\)
\begin{equation}
 \bigl|\{y:\ |\nabla U(y,s)|\ge\tfrac12\}\bigr|
 \ \ge\ \frac{1}{\sup|\nabla U|^2}\Bigl(\Bigl(\frac{\nu^3}{8C_1(-s)}\Bigr)^{1/2}
 -\tfrac14\,\bigl|\{|\nabla U(\cdot,s)|<\tfrac12\}\cap\operatorname{supp}\nabla U(\cdot,s)\bigr|\Bigr)_+ ,
 \label{eq:v50-S2-volume}
\end{equation}
so that, in particular, the set where \(|\nabla U(\cdot,s)|\ge1/2\) has
infinite measure in the limit \(s\uparrow0\) whenever the low-gradient part
of the support has finite measure, and in any case
\(\int|\nabla U(s)|^2\ge(\nu^3/(8C_1(-s)))^{1/2}\) is carried by a region
whose measure tends to infinity.
\end{proposition}

\begin{proof}
In branch (S2) every \((y,0)\) is a regular point, so \(U\) is bounded
near each such point; together with Lemma~\ref{lem:v46-bounded}(i) on the
exterior region and compactness of \(\overline{B_{R_0}}\), \(U\) is
bounded on \(\R^3\times[s_0,0]\) for some \(s_0\), and (E5), (E7) give the
derivative bounds and smoothness up to \(s=0\); the trace properties are
Theorems~\ref{thm:v48-trace-nontrivial} and
Proposition~\ref{prop:v49-trace-smooth-decay}.  For the volume estimate,
split \(Y_U(s)=\int|\nabla U|^2\) over \(\{|\nabla U|\ge1/2\}\), where the
integrand is at most \(\sup|\nabla U|^2\), and its complement within the
support, where it is at most \(1/4\); use the lower bound
\eqref{eq:v42-two-sided}.  Since \(Y_U(s)\to\infty\) while
\(\sup|\nabla U|<\infty\), the measure of \(\{|\nabla U(\cdot,s)|\ge c\}\)
tends to infinity for every \(c\) smaller than the limit superior of
\(\sup_y|\nabla U(y,s)|\), and the support statement follows.
\end{proof}

\begin{assessment}[Branch (S2) is a bounded ancient solution with escaping
enstrophy]
\label{ass:v50-S2}
In branch (S2) the ancient limit is a bounded smooth solution on a slab
ending at \(s=0\), with bounded derivatives, whose enstrophy nevertheless
diverges like \((-s)^{-1/2}\) because its gradient of order one occupies
an ever larger region.  Its trace is a smooth field on all of \(\R^3\)
that is neither in \(L^2\) nor decaying faster than the critical Morrey
rate.  The Liouville question in this branch is therefore about bounded
ancient solutions on a slab whose spatial support of order-one gradient
grows without bound, a setting in which the backward-uniqueness package
does not apply directly because the trace is nonzero.  Nothing in this
series excludes the branch.
\end{assessment}

% =============================================================================
\section{Integrated V50 conclusion and next acquisition order}
\label{sec:v50-conclusion}
% =============================================================================

\begin{theorem}[What V50 proves]
\label{thm:v50-conclusion}
Relative to V1--V49: the companion-audit decision
(Theorem~\ref{thm:v50-companion-decision}); the consolidated description
of the ancient limit of a type-I singularity and the final form of the
Liouville gate (Theorem~\ref{thm:v50-ancient-limit}); and the structure of
the dispersion branch (Proposition~\ref{prop:v50-S2}).  The full
Navier--Stokes stability/global-regularity theorem remains unproved.
\end{theorem}

\begin{proof}
The cited results.
\end{proof}

\begin{programme}[V51 acquisition order]
\label{prog:v51}
\begin{enumerate}[label=\textup{(AC\arabic*)},leftmargin=4.8em]
\item Concentration branches.  In (L1) and (S1), iterate the rescaling
      and determine whether the sequence of limits stabilizes: a fixed
      point of the construction would be an element of the class equal to
      one of its own rescaling limits, hence discretely self-similar.
      Record precisely what compactness is missing for a fixed point to
      exist.
\item Dispersion branches.  In (L2) and (S2), test the local energy
      identity on cylinders of size \(\lambda\) against the Morrey bound
      and the backward cone; determine whether dissipation of order
      \(\lambda\) in the slab can be spread over volume \(\gg\lambda^3\)
      while every unit ball carries energy at most \(C_M\).
\item Type-II: item (AC2) of Programme~\ref{prog:v41} remains.
\item The next compulsory companion audit is V55.
\end{enumerate}
\end{programme}

\begin{assessment}[Position after V50]
\label{ass:v50-position}
Twenty versions have converted the cutoff-multiplicity bottleneck of the
first thirty into a two-regime classification, closed the type-I regime
up to a precisely described Liouville problem on an explicit class of
ancient solutions with proved trace properties, and left the type-II
regime with a price sheet of necessary conditions in three ledgers.  No
full stability or global-regularity claim is made.
\end{assessment}

% =============================================================================
\section*{Version V50 append-only record}
\addcontentsline{toc}{section}{Version V50 append-only record}
% =============================================================================

The complete V49 mathematical source was retained before this continuation.
Its SHA--256 digest was
\begin{center}\scriptsize\ttfamily
3e4a757e13487edabf3afcd8eba71baeef663442da98fbeaa8b843637abd4175.
\end{center}
V50 performed the compulsory companion audit, consolidated V41--V49 into
one description of the ancient limit of a type-I singularity with the
final form of the Liouville gate, and described the dispersion branch.  No
full regularity claim is made.

% =============================================================================
\section{V51: the singular set as the critical-density set of the terminal
datum}
\label{sec:v51-entry}
% =============================================================================

Two exact characterizations of the type-I singular set follow from
V44--V48 with short proofs, and are recorded here together with a local
blow-up statement.  The singular set \(\Sigma\) is exactly the set of
points at which the terminal datum has positive upper critical density,
\(\limsup_{\rho\to0}\rho^{-1}\int_{B_\rho(x)}|u(T_*)|^2>0\), and exactly
the set of points near which the terminal datum is not cubically
integrable.  At every singular point the local \(L^3\) norm of \(u(t)\) on
every ball tends to infinity as \(t\uparrow T_*\), as a limit and not
only a limit superior, and at most logarithmically.  The dispersion
branches of V49 are then checked against the Morrey bound and found
consistent, and the fixed-point question of Programme~\ref{prog:v51} is
recorded as open for lack of compactness under the scaling group.  No
global regularity claim is made.  The next compulsory companion audit is
V55.

% =============================================================================
\section{Two characterizations of the singular set}
\label{sec:v51-characterizations}
% =============================================================================

Throughout, \eqref{eq:v32-type-I} holds, \(\Sigma\) is the singular set of
Theorem~\ref{thm:v44-localization}, and \(u(T_*)\) is the terminal datum.

\begin{theorem}[The singular set is the positive-density set and the
\(L^3\)-singular support of the terminal datum]
\label{thm:v51-characterization}
Modulo the imports (E1)--(E6), for \(x\in\Torus^3\) the following are
equivalent:
\begin{enumerate}[label=\textup{(\roman*)},leftmargin=3.7em]
\item \(x\in\Sigma\);
\item \(\displaystyle\limsup_{\rho\to0}\frac1\rho\int_{B_\rho(x)}|u(T_*)|^2>0\);
\item \(u(T_*)\notin L^3(B_\rho(x))\) for every \(\rho>0\);
\item \(\displaystyle\lim_{t\uparrow T_*}\|u(t)\|_{L^3(B_\rho(x))}=\infty\)
      for every \(\rho>0\).
\end{enumerate}
Moreover in case (iv) the divergence is at most logarithmic:
\(\|u(t)\|_{L^3(B_\rho(x))}\le C_3(2A_\delta+2B^2L_\delta(t)^2)^{1/2}\).
\end{theorem}

\begin{proof}
(i)\(\Rightarrow\)(ii): every singular point is a concentration point by
(E1), so Corollary~\ref{cor:v46-critical-local-energy} applies.
(ii)\(\Rightarrow\)(iii): if \(u(T_*)\in L^3(B_{\rho_0}(x))\) then by
H\"older \(\rho^{-1}\int_{B_\rho(x)}|u(T_*)|^2\le\rho^{-1}
\|u(T_*)\|_{L^3(B_\rho)}^2|B_\rho|^{1/3}=C\|u(T_*)\|_{L^3(B_\rho(x))}^2\to0\)
by absolute continuity of the integral.
(iii)\(\Rightarrow\)(i): if \(x\notin\Sigma\), \(u\) is bounded on a
neighbourhood of \((x,T_*)\), so \(u(T_*)\), which is the uniform limit of
\(u(t)\) there by Theorem~\ref{thm:v44-localization}(ii)--(iii), is
bounded near \(x\) and hence in \(L^3(B_\rho(x))\) for small \(\rho\).
(iii)\(\Rightarrow\)(iv): \(u(t)\to u(T_*)\) strongly in \(L^2\), so every
sequence \(t_j\uparrow T_*\) has a subsequence along which
\(u(t_j)\to u(T_*)\) almost everywhere; Fatou's lemma gives
\(\liminf_j\int_{B_\rho}|u(t_j)|^3\ge\int_{B_\rho}|u(T_*)|^3=\infty\).
Since every sequence has such a subsequence, the limit is \(+\infty\).
(iv)\(\Rightarrow\)(i): if \(x\notin\Sigma\), \(u\) is bounded near
\((x,T_*)\), so the local \(L^3\) norm on a small ball stays bounded.
The logarithmic bound is \(\|u(t)\|_3\le C_3\|A^{1/4}u(t)\|_2\) with
Theorem~\ref{thm:v33-type-I-log}.
\end{proof}

\begin{corollary}[Local blow-up of the critical norm at every singular
point]
\label{cor:v51-local-blowup}
Under a type-I enstrophy rate, at every \(x\in\Sigma\) and every
\(\rho>0\),
\begin{equation}
 \lim_{t\uparrow T_*}\|u(t)\|_{L^3(B_\rho(x))}=\infty,
 \qquad
 \|u(t)\|_{L^3(B_\rho(x))}\le C\log\frac{9\delta}{T_*-t}
 \quad(t\text{ close to }T_*).
 \label{eq:v51-local-blowup}
\end{equation}
\end{corollary}

\begin{proof}
Theorem~\ref{thm:v51-characterization}.
\end{proof}

\begin{assessment}[Relation to the external local blow-up theorem]
\label{ass:v51-Seregin}
It is known externally, without any rate hypothesis, that the global
\(L^3\) norm of a singular solution tends to infinity at the singular
time; this series has not used that theorem.  Under the type-I enstrophy
rate, Corollary~\ref{cor:v51-local-blowup} gives the local version on
every ball around every singular point, with the logarithmic upper rate,
and Theorem~\ref{thm:v51-characterization} identifies the singular set
with the \(L^3\)-singular support of the terminal datum.  The route is
entirely through the terminal datum: strong \(L^2\) convergence (V32),
the non-\(L^3\) property at concentration points (V46), and Fatou.
\end{assessment}

% =============================================================================
\section{Dispersion branches against the Morrey bound}
\label{sec:v51-dispersion}
% =============================================================================

\begin{proposition}[Local dissipation of the ancient limit at scale
\(\lambda\)]
\label{prop:v51-local-dissipation}
For \(U\in\mathcal A_\nu(C)\) with the Morrey bound
\eqref{eq:v47-limit-morrey}, every \(y\in\R^3\), and every \(\lambda>0\),
\begin{equation}
 \nu\iint_{B_\lambda(y)\times(-\lambda^2,0)}|\nabla U|^2
 \le C_{\rm diss}\,\lambda ,
 \label{eq:v51-local-dissipation}
\end{equation}
with \(C_{\rm diss}\) depending only on \(C,C_M,\nu,\phi\).  Hence in
branch (L2) of Theorem~\ref{thm:v49-large-scale}, and in case (B), the
slab dissipation \(\int_{-\lambda^2}^0Y_U\ge2(\nu^3/(8C_1))^{1/2}\lambda\)
is distributed over at least
\(2(\nu^3/(8C_1))^{1/2}/\bigl(\nu\,d_U^*(\lambda)\bigr)\to\infty\)
disjoint cylinders of size \(\lambda\), and no contradiction with the
Morrey bound arises.
\end{proposition}

\begin{proof}
Apply the local energy inequality on \(B_{2\lambda}(y)\) over
\((-\lambda^2,0)\) as in the proof of Theorem~\ref{thm:v47-morrey}, with
\(\lambda\) in place of \(\rho\), the Morrey bound for the initial local
energy, and the \(L^6\) and \(L^3\) bounds \eqref{eq:v41-limit-type-I}
and the Riesz bound for the remaining terms; every term is at most a
constant times \(\lambda\).  The counting statement divides the slab
lower bound of \eqref{eq:v42-two-sided} by the per-cylinder bound
\(\lambda d_U^*(\lambda)\).
\end{proof}

\begin{firewall}[Fixed points of the rescaling are not available]
\label{fw:v51-fixed-point}
Item (AC1) of Programme~\ref{prog:v51} asks whether iterating the
concentration rescalings stabilizes at an element equal to its own
rescaling limit.  The construction gives, from any element of the class
with \(\limsup_\lambda d^*(\lambda)>0\), a new element with unit-scale
dissipation bounded below, but the sequence of rescaling factors is chosen
by the previous element and there is no compactness of the class under the
continuous scaling group and translations that would make the iteration
converge.  A discretely self-similar element would follow only from an
additional monotonicity, which no quantity in this series provides
(the natural candidate, the dissipation concentration \(d^*(\lambda)\), is
bounded but not monotone in \(\lambda\)).  The item is recorded as open.
\end{firewall}

% =============================================================================
\section{Integrated V51 conclusion and next acquisition order}
\label{sec:v51-conclusion}
% =============================================================================

\begin{theorem}[What V51 proves]
\label{thm:v51-conclusion}
Relative to V1--V50, modulo (E1)--(E7): the four-way characterization of
the singular set (Theorem~\ref{thm:v51-characterization}); the local
\(L^3\) blow-up as a limit at every singular point, with logarithmic upper
rate (Corollary~\ref{cor:v51-local-blowup}); and the local dissipation
bound \eqref{eq:v51-local-dissipation} with the consistency of the
dispersion branches.  The fixed-point question is open.  The full
Navier--Stokes stability/global-regularity theorem remains unproved.
\end{theorem}

\begin{proof}
The cited results.
\end{proof}

\begin{programme}[V52 acquisition order]
\label{prog:v52}
\begin{enumerate}[label=\textup{(AC\arabic*)},leftmargin=4.8em]
\item Density of the terminal datum.  Determine whether the upper
      critical density \(\limsup_\rho\rho^{-1}\int_{B_\rho(x)}|u(T_*)|^2\)
      at singular points is bounded below by an absolute constant times
      \(\varepsilon_1\) or \(c_*\); a uniform lower bound would make
      \(\Sigma\) finite by the Morrey bound and a covering argument.
\item Lower density.  Determine whether the lower density
      \(\liminf_\rho\rho^{-1}\int_{B_\rho(x)}|u(T_*)|^2\) is positive at
      singular points; positivity at all small scales would upgrade the
      annular statement of Corollary~\ref{cor:v48-terminal-data} to all
      scales and bear on finiteness of \(\Sigma\).
\item Type-II: item (AC2) of Programme~\ref{prog:v41} remains.
\item The next compulsory companion audit is V55.
\end{enumerate}
\end{programme}

\begin{assessment}[Position after V51]
\label{ass:v51-position}
The singular set of a type-I singularity is now characterized three ways
through the terminal datum alone.  A uniform lower bound on its critical
density at singular points is the natural next target, since it would
force finiteness of the singular set through the Morrey bound.  No full
stability or global-regularity claim is made.
\end{assessment}

% =============================================================================
\section*{Version V51 append-only record}
\addcontentsline{toc}{section}{Version V51 append-only record}
% =============================================================================

The complete V50 mathematical source was retained before this continuation.
Its SHA--256 digest was
\begin{center}\scriptsize\ttfamily
a6ff0ce7d0b350ec67a27b2c686b5c8126084712b751282036adc2d0e268091e.
\end{center}
V51 proved the characterization of the type-I singular set as the
positive-density set and \(L^3\)-singular support of the terminal datum,
the local \(L^3\) blow-up as a limit with logarithmic rate, and the
consistency of the dispersion branches with the Morrey bound.  No full
regularity claim is made.

% =============================================================================
\section{V52: a uniform lower bound on the critical density of the terminal
datum over the singular set}
\label{sec:v52-entry}
% =============================================================================

Item (AC1) of Programme~\ref{prog:v52} is settled in the form that the
present tools permit.  The upper critical density of the terminal datum,
which V51 showed to be positive at every singular point, is bounded below
by a positive constant uniformly over the singular set.  The proof is a
diagonal compactness argument: a sequence of singular points with
densities tending to zero would produce, by rescaling the same solution at
those points, an ancient limit with vanishing trace, which the
backward-uniqueness argument of V46 forbids.  The constant depends on the
solution; a bound depending only on the type-I constants is not obtained.
It is also shown, by an explicit superposition, that no density bound on
the terminal datum, upper or lower, can count singular points, so
finiteness of the singular set is a common-scale activity question and not
a density question.  Item (AC2), positivity of the lower density, remains
open.  No global regularity claim is made.  The next compulsory companion
audit is V55.

% =============================================================================
\section{The uniform lower bound}
\label{sec:v52-lower-bound}
% =============================================================================

Throughout, \eqref{eq:v32-type-I} holds and the imports (E1)--(E7) are in
force.  For \(x\in\Torus^3\) put
\begin{equation}
 \overline d(x):=\limsup_{\rho\to0}\frac1\rho\int_{B_\rho(x)}|u(T_*,x')|^2\,\dd x' ,
 \label{eq:v52-upper-density}
\end{equation}
so that, by Theorem~\ref{thm:v51-characterization},
\(\Sigma=\{\overline d>0\}\), and by Theorem~\ref{thm:v47-morrey},
\(\overline d\le C_M\) everywhere.

\begin{theorem}[Uniform positive lower bound over the singular set]
\label{thm:v52-uniform-lower-bound}
\begin{equation}
 \boxed{\quad
 c_0(u):=\inf_{x\in\Sigma}\overline d(x)\ >\ 0 .
 \quad}
 \label{eq:v52-uniform-lower-bound}
\end{equation}
\end{theorem}

\begin{proof}
Suppose \(c_0(u)=0\); then there are \(x_j\in\Sigma\) with
\(\overline d(x_j)\le2^{-j}\).  By the definition of the limit superior
there is \(\rho_j>0\) such that
\(\rho^{-1}\int_{B_\rho(x_j)}|u(T_*)|^2\le2\cdot2^{-j}\) for all
\(\rho\le\rho_j\).  Each \(x_j\) is a concentration point by (E1); choose
one of its concentration scales \(r_j\) with \(r_j\le\rho_j/2^{j+1}\), so
that for every \(R\le2^j\),
\begin{equation}
 \frac1{r_j}\int_{B_{2Rr_j}(x_j)}|u(T_*)|^2
 \le2R\cdot2\cdot2^{-j}=4R\,2^{-j}.
 \label{eq:v52-small-terminal-energy}
\end{equation}
Let \(U_j(s,y):=r_ju(T_*+r_j^2s,x_j+r_jy)\).  These are rescalings of the
one solution \(u\) with \(r_j\to0\), and the uniform bounds of
Proposition~\ref{prop:v41-uniform-bounds} hold for them, since those
bounds used only the enstrophy rate and not the centre; also
\(\iint_{B_1\times(-1,0)}|\nabla U_j|^2>\varepsilon_1\) by the choice of
\(r_j\) as a concentration scale of \(x_j\).  The proof of
Theorem~\ref{thm:v41-ancient-limit} therefore applies verbatim and gives,
along a subsequence, a nontrivial ancient limit \(U\in\mathcal A_\nu(C_I^2\nu^{3/2})\)
with \eqref{eq:v41-limit-nontrivial}.  For fixed \(R\) and
\(-1\le s<0\), Lemma~\ref{lem:v46-local-energy-near-zero} applied to
\(U_j\) gives \(e_R^{(j)}(s)\le e_R^{(j)}(0^-)+C_R(-s)^{1/4}\), and
\(e_R^{(j)}(0^-)\le\frac1{2r_j}\int_{B_{2Rr_j}(x_j)}|u(T_*)|^2
\le2R\,2^{-j}\to0\) by \eqref{eq:v52-small-terminal-energy} once
\(2^j\ge R\).  Passing to the limit as in
Proposition~\ref{prop:v46-vanishing}, \(e_R(s)\le C_R(-s)^{1/4}\), so
\(U(s)\to0\) in \(L^2_{\rm loc}\) as \(s\uparrow0\).  The proof of
Theorem~\ref{thm:v46-non-L3} from this point on, which used only this
vanishing, Lemma~\ref{lem:v46-bounded} (valid for any element of the class
with the Morrey bound, which \(U\) has by Proposition~\ref{prop:v47-trace}),
and (E6), yields \(U\equiv0\) on \(\R^3\times(-1,0)\), contradicting
\eqref{eq:v41-limit-nontrivial}.
\end{proof}

\begin{corollary}[The singular set is a level set of the density]
\label{cor:v52-level-set}
\(\Sigma=\{x:\overline d(x)\ge c_0(u)\}\), and \(c_0(u)\le C_M\).  In
particular the density function \(\overline d\) takes no values in
\((0,c_0(u))\).
\end{corollary}

\begin{proof}
\(\Sigma=\{\overline d>0\}\) by Theorem~\ref{thm:v51-characterization},
and \(\overline d\ge c_0(u)\) on \(\Sigma\) by
Theorem~\ref{thm:v52-uniform-lower-bound}; the upper bound is the Morrey
bound at \(t=T_*\).
\end{proof}

% =============================================================================
\section{Why finiteness does not yet follow}
\label{sec:v52-finiteness}
% =============================================================================

\begin{firewall}[Density bounds do not count singular points]
\label{fw:v52-no-count}
Theorem~\ref{thm:v52-uniform-lower-bound} says that around every singular
point \(x\) there are arbitrarily small radii \(\rho\) with
\(\int_{B_\rho(x)}|u(T_*)|^2\ge\frac12c_0(u)\rho\).  For \(m\) distinct
singular points with pairwise separation \(\sigma\), balls of a common
radius \(\rho<\sigma/2\) are disjoint, and if the density were bounded
below by \(c\) at that common scale the total energy in them would be at
least \(mc\rho\).  But the Morrey bound on a ball containing all of them
gives only \(C_M(\operatorname{diam}+\rho)\), and the finite energy of
\(u(T_*)\) gives only a constant, so the resulting inequality
\(m\le C_M(\operatorname{diam}+\rho)/(c\rho)\) does not bound \(m\)
independently of \(\rho\).  Concretely, a superposition of profiles of
size \(a_j|x-x_j|^{-1}\) at points \(x_j\to x_\infty\) with
\(a_j\ge c^{1/2}\) and rapidly decreasing spacings has uniformly positive
lower density at every \(x_j\), critical Morrey bound, and finite energy.
Neither the upper nor the lower density therefore counts singular points;
a count needs an argument at a common scale with a scale-critical budget,
such as the slab dissipation of Lemma~\ref{lem:v43-slab}, and V43 showed
that the slab dissipation counts only points that are active at a common
scale.
\end{firewall}

\begin{assessment}[The remaining question about the singular set]
\label{ass:v52-lower-density}
After V52 the singular set of a type-I singularity is a closed set of
dimension zero on which the terminal datum has upper critical density in
\([c_0(u),C_M]\).  Finiteness is not a consequence of any density bound
on the terminal datum.  It would follow from a common-scale activity
statement: if every singular point were active, in the sense of
\eqref{eq:v43-active-set}, at every sufficiently small dyadic scale, then
Lemma~\ref{lem:v43-few-active} would bound \(|\Sigma|\) by \(N_*\).
Whether the enstrophy of a type-I singularity can concentrate at a point
only along a lacunary sequence of scales is therefore the precise question
behind finiteness; Firewall~\ref{fw:v44-sparse-activity} recorded that
sparse activity is not excluded by dissipation bounds, and V52 shows that
the terminal datum does not exclude it either.
\end{assessment}

% =============================================================================
\section{Integrated V52 conclusion and next acquisition order}
\label{sec:v52-conclusion}
% =============================================================================

\begin{theorem}[What V52 proves]
\label{thm:v52-conclusion}
Relative to V1--V51, modulo (E1)--(E7): the uniform positive lower bound
\eqref{eq:v52-uniform-lower-bound} on the upper critical density of the
terminal datum over the singular set; the level-set description of
\(\Sigma\); and the record that neither density bound counts singular
points (Firewall~\ref{fw:v52-no-count}).  The full Navier--Stokes
stability/global-regularity theorem remains unproved.
\end{theorem}

\begin{proof}
The cited results.
\end{proof}

\begin{programme}[V53 acquisition order]
\label{prog:v53}
\begin{enumerate}[label=\textup{(AC\arabic*)},leftmargin=4.8em]
\item Lower density.  For a singular point \(x\) and a scale \(\rho\) at
      which the density \(\rho^{-1}\int_{B_\rho(x)}|u(T_*)|^2\) is small,
      rescale at \((x,\rho)\); the limit exists in the class but may be
      trivial.  Determine whether the backward cone or the two-sided slab
      dissipation forces a nontrivial limit along every sequence of scales
      at a singular point, or construct the obstruction.
\item Uniform constants.  Determine whether \(c_0(u)\) can be replaced by
      a constant depending only on \(C_I\), \(\nu\), and the absolute
      constants, by running the diagonal argument over a family of
      solutions with common type-I constants; record the compactness
      needed.
\item Type-II: item (AC2) of Programme~\ref{prog:v41} remains.
\item The next compulsory companion audit is V55.
\end{enumerate}
\end{programme}

\begin{assessment}[Position after V52]
\label{ass:v52-position}
The type-I singular set carries a uniformly positive upper critical
density of the terminal datum; its finiteness is a common-scale activity
question that no density bound decides.  The type-I regime is otherwise unchanged: the
Liouville gate of V41 remains.  No full stability or global-regularity
claim is made.
\end{assessment}

% =============================================================================
\section*{Version V52 append-only record}
\addcontentsline{toc}{section}{Version V52 append-only record}
% =============================================================================

The complete V51 mathematical source was retained before this continuation.
Its SHA--256 digest was
\begin{center}\scriptsize\ttfamily
22b6c1f0ad3674ee83b8b6834bf085daf9963c2f8ff0272adb8d0fc366873a22.
\end{center}
V52 proved a uniform positive lower bound on the upper critical density of
the terminal datum over the singular set by a diagonal compactness
argument, described the singular set as a density level set, and showed that no
density bound on the terminal datum counts singular points.  No full
regularity claim is made.

% =============================================================================
\section{V53: dissipation activity forces terminal density, scale by scale}
\label{sec:v53-entry}
% =============================================================================

Item (AC1) of Programme~\ref{prog:v53} is settled in the form of a
per-scale implication.  At a singular point, along any sequence of scales
tending to zero, either the rescaled local dissipation vanishes in the
limit, or the terminal datum has positive critical density at comparable
scales; the two cannot fail together.  The proof is the V46 argument run
along an arbitrary sequence of scales instead of a concentration sequence.
By compactness the implication is quantitative: small terminal density at
all comparable scales forces small rescaled dissipation.  The version also
records that, after the \(L^\infty\) upgrade of V46, the Liouville gate of
V41 is a Liouville problem for bounded ancient solutions with finite
Dirichlet integral, a subclass of the class in which the Liouville
conjecture for bounded ancient solutions is open.  No global regularity
claim is made.  The next compulsory companion audit is V55.

% =============================================================================
\section{The per-scale implication}
\label{sec:v53-per-scale}
% =============================================================================

Throughout, \eqref{eq:v32-type-I} holds and (E1)--(E7) are in force.
For \(x\in\Torus^3\), \(\rho>0\), \(R\ge1\), and \(0<\eta<S\), define the
rescaled local dissipation and the terminal density at scale \(\rho\):
\begin{equation}
 \mathcal D_x(\rho;R,S,\eta)
 :=\frac1\rho\iint_{B_{R\rho}(x)\times(T_*-S\rho^2,\,T_*-\eta\rho^2)}|\nabla u|^2,
 \qquad
 \mathcal E_x(\rho;R):=\frac1\rho\int_{B_{R\rho}(x)}|u(T_*)|^2 .
 \label{eq:v53-quantities}
\end{equation}
Both are bounded: \(\mathcal D_x\le2C_I^2\nu^{3/2}S^{1/2}\) by
Lemma~\ref{lem:v43-slab}, and \(\mathcal E_x\le C_MR\) by
Theorem~\ref{thm:v47-morrey}.

\begin{theorem}[Nonvanishing dissipation at a scale sequence forces
terminal density at comparable scales]
\label{thm:v53-per-scale}
Let \(x\in\Torus^3\) and \(\rho_j\downarrow0\).  If for some
\(R,S,\eta\)
\begin{equation}
 \liminf_{j\to\infty}\mathcal D_x(\rho_j;R,S,\eta)>0,
 \label{eq:v53-active}
\end{equation}
then there is \(R'\) with
\begin{equation}
 \liminf_{j\to\infty}\mathcal E_x(\rho_j;R')>0 .
 \label{eq:v53-density}
\end{equation}
Equivalently: if \(\mathcal E_x(\rho_j;R')\to0\) for every \(R'\), then
\(\mathcal D_x(\rho_j;R,S,\eta)\to0\) for every \(R,S,\eta\).
\end{theorem}

\begin{proof}
Let \(U_j(s,y):=\rho_ju(T_*+\rho_j^2s,x+\rho_jy)\).  The uniform bounds of
Proposition~\ref{prop:v41-uniform-bounds} hold, and the compactness of
Theorem~\ref{thm:v41-ancient-limit} gives, along a subsequence, a limit
\(U\in\mathcal A_\nu(C_I^2\nu^{3/2})\) with strong convergence in
\(L^2((-S,-\eta);H^1(B_R))\), so
\(\iint_{B_R\times(-S,-\eta)}|\nabla U|^2=\lim_j\mathcal D_x(\rho_j;R,S,\eta)>0\):
the limit is nontrivial.  Suppose \eqref{eq:v53-density} fails for every
\(R'\); passing to a further subsequence (diagonal in \(R'\)),
\(\mathcal E_x(\rho_j;R')\to0\) for every \(R'\).  Then, exactly as in
Proposition~\ref{prop:v46-vanishing}, \(e_{R'}^{(j)}(0^-)\le\frac12
\mathcal E_x(\rho_j;2R')\to0\), Lemma~\ref{lem:v46-local-energy-near-zero}
gives \(e_{R'}(s)\le C_{R'}(-s)^{1/4}\) for the limit, and the proof of
Theorem~\ref{thm:v46-non-L3} gives \(U\equiv0\), a contradiction.  The
second formulation is the contrapositive with the subsequence argument
reversed: if the dissipation did not tend to zero along the full
sequence, a subsequence with \eqref{eq:v53-active} would exist, and the
first part would give \eqref{eq:v53-density} along it, contradicting
\(\mathcal E_x\to0\).
\end{proof}

\begin{corollary}[Quantitative form]
\label{cor:v53-quantitative}
For every \(\varepsilon>0\) and every \(R,S,\eta\) there are
\(\kappa>0\) and \(R'\ge R\), depending only on
\(\varepsilon,R,S,\eta,C_I,\nu\) and the absolute constants, such that
for every \(x\) and every \(\rho\le\rho_\varepsilon\),
\begin{equation}
 \mathcal E_x(\rho;R'')\le\kappa\ \text{ for all }R''\le R'
 \quad\Longrightarrow\quad
 \mathcal D_x(\rho;R,S,\eta)\le\varepsilon .
 \label{eq:v53-quantitative}
\end{equation}
\end{corollary}

\begin{proof}
If not, there are \(\varepsilon,R,S,\eta\) and, for each \(n\), points
\(x_n\) and scales \(\rho_n\to0\) with
\(\mathcal E_{x_n}(\rho_n;R'')\le1/n\) for all \(R''\le n\) but
\(\mathcal D_{x_n}(\rho_n;R,S,\eta)>\varepsilon\).  Rescaling the one
solution \(u\) at \((x_n,\rho_n)\) and repeating the proof of
Theorem~\ref{thm:v53-per-scale}, whose compactness did not depend on the
centre, gives a nontrivial limit with vanishing trace and the same
contradiction.  The threshold \(\rho_\varepsilon\) is needed only so that
the cylinders lie in \([T_*-\delta,T_*)\).
\end{proof}

\begin{assessment}[Terminal datum as a lower control on the dynamics]
\label{ass:v53-terminal-controls}
V47 bounded the terminal datum from above by the dynamics through the
critical Morrey bound.  Theorem~\ref{thm:v53-per-scale} is the reverse
direction: the terminal datum bounds the recent local dissipation from
above at every scale, in the sense that a terminal datum that is small at
all scales comparable to \(\rho\) around \(x\) forces the dissipation in
the parabolic cylinders of size \(\rho\) at \(x\) to be small.  At a
singular point the concentration scales are active, and V52 showed the
terminal density is positive at some comparable scale; the new statement
says that this holds scale by scale, not only in the limit superior.
Finiteness of \(\Sigma\), which by Assessment~\ref{ass:v52-lower-density}
needs activity at a common scale, is therefore equivalent to a common-scale
positivity of the terminal density, and V52's superposition shows that
such a common-scale positivity is not forced by any of the proved
properties of the terminal datum.
\end{assessment}

% =============================================================================
\section{The gate as a bounded-ancient-solution problem}
\label{sec:v53-gate}
% =============================================================================

\begin{proposition}[Elements of the class are bounded ancient solutions on
every half-line]
\label{prop:v53-bounded-ancient}
Let \(U\in\mathcal A_\nu(C)\).  Then for every \(\eta>0\),
\(\sup_{s\le-\eta}\|U(s)\|_{L^\infty(\R^3)}\le C_{\rm b}(\eta)<\infty\),
with \(C_{\rm b}(\eta)\) depending only on \(C,\nu,\eta\) and the absolute
constants; \(\nabla U\) and \(\nabla^2U\) are likewise bounded there.
Consequently \(U|_{(-\infty,-1]}\) is a bounded ancient solution of the
Navier--Stokes equations on \(\R^3\) with uniformly bounded enstrophy
\(\|\nabla U(s)\|_2^2\le C\) and \(U(s)\in L^6(\R^3)\).
\end{proposition}

\begin{proof}
Lemma~\ref{lem:v46-bounded}(ii) and (E7).
\end{proof}

\begin{firewall}[Position of the gate among Liouville problems]
\label{fw:v53-gate-position}
The Liouville conjecture for bounded ancient solutions of the
three-dimensional Navier--Stokes equations, that every such solution is
constant, is open; it is known to hold in two dimensions and in
axisymmetric settings, and in three dimensions under additional
assumptions.  Proposition~\ref{prop:v53-bounded-ancient} places the
type-I gate of V41 inside this conjecture with two extra hypotheses,
finite Dirichlet integral at every time and membership in \(L^6\), and
with the additional structure of V47--V49: a critical Morrey trace at the
final time with vorticity on every exterior region.  Whether these extra
hypotheses make the conjecture tractable is not decided here; what is
decided is that a proof of the bounded-ancient Liouville conjecture in the
finite-Dirichlet subclass would exclude every singularity with a type-I
enstrophy rate.
\end{firewall}

% =============================================================================
\section{Integrated V53 conclusion and next acquisition order}
\label{sec:v53-conclusion}
% =============================================================================

\begin{theorem}[What V53 proves]
\label{thm:v53-conclusion}
Relative to V1--V52, modulo (E1)--(E7): the per-scale implication
(Theorem~\ref{thm:v53-per-scale}) and its quantitative form
(Corollary~\ref{cor:v53-quantitative}); and the identification of the
gate as a Liouville problem for bounded ancient solutions with finite
Dirichlet integral (Proposition~\ref{prop:v53-bounded-ancient}).  The full
Navier--Stokes stability/global-regularity theorem remains unproved.
\end{theorem}

\begin{proof}
The cited results.
\end{proof}

\begin{programme}[V54 acquisition order]
\label{prog:v54}
\begin{enumerate}[label=\textup{(AC\arabic*)},leftmargin=4.8em]
\item Bounded ancient solutions with finite Dirichlet integral.  Test the
      known two-dimensional and axisymmetric Liouville arguments for
      bounded ancient solutions against the extra hypotheses of
      Proposition~\ref{prop:v53-bounded-ancient}; determine whether the
      finite Dirichlet integral supplies the decay at infinity those
      arguments use in place of symmetry.
\item Common-scale activity.  Determine whether the backward cone applied
      on exterior regions, where \(U\) is bounded up to \(s=0\), forces
      activity of a singular point at every sufficiently small scale, or
      construct a lacunary schedule consistent with all proved bounds.
\item Type-II: item (AC2) of Programme~\ref{prog:v41} remains.
\item The next compulsory companion audit is V55.
\end{enumerate}
\end{programme}

\begin{assessment}[Position after V53]
\label{ass:v53-position}
The type-I regime is reduced to a Liouville problem for bounded ancient
solutions with finite Dirichlet integral and a vortical critical-Morrey
trace, and the terminal datum is proved to control the dynamics from above
and to be controlled by it, scale by scale.  No full stability or
global-regularity claim is made.
\end{assessment}

% =============================================================================
\section*{Version V53 append-only record}
\addcontentsline{toc}{section}{Version V53 append-only record}
% =============================================================================

The complete V52 mathematical source was retained before this continuation.
Its SHA--256 digest was
\begin{center}\scriptsize\ttfamily
99f6290540d0219c9cb23a7225d320664ba4269a4556d90e1bb28a354407c3db.
\end{center}
V53 proved that nonvanishing rescaled dissipation at a singular point
forces positive terminal density at comparable scales, gave the
quantitative form by compactness, and placed the type-I Liouville gate
inside the bounded-ancient-solution Liouville problem with finite
Dirichlet integral.  No full regularity claim is made.

% =============================================================================
\section{V54: the dyadic Leray-share dichotomy between the regimes}
\label{sec:v54-entry}
% =============================================================================

Item (AC2) of Programme~\ref{prog:v54} is closed negatively: a lacunary
activity schedule at a singular point is consistent with every proved
bound, because the V53 quantities exclude the final fraction of each
parabolic window, and activity at a much smaller scale lives exactly there.
Item (AC1) is not advanced.  The version instead returns to the type-II
regime, item (AC2) of Programme~\ref{prog:v41}, and proves a dichotomy in
the Leray ledger alone.  Writing \(\epsilon_k^{(\theta)}\) for the dissipation in the
\(k\)-th geometric slab of ratio \(\theta\) before the singular time, the
Leray rate gives \(\epsilon_k^{(\theta)}\gtrsim\theta^{k/2}\) always, and
the solution has a type-I enstrophy bound along a sequence of times if and
only if \(\epsilon_k^{(\theta)}\lesssim\theta^{k/2}\) along a subsequence,
for one and then for every ratio close enough to one.  Thus type II, in
the strong sense that no sequence of times is type-I, is exactly the
statement that every geometric slab spends a Leray share large compared
with its parabolic size, by a factor tending to infinity, while the total
remains finite.  At a type-I time the rescaled enstrophy is bounded on a forward
parabolic window, which gives a local compactness statement without
ancientness.  No global regularity claim is made.  The next compulsory
companion audit is V55.

% =============================================================================
\section{Dyadic slabs and their Leray shares}
\label{sec:v54-slabs}
% =============================================================================

Let \(T_*<\infty\) be the maximal time, \(Y=\|\Lambda u\|_2^2\), and for
\(k\ge k_0\) let
\begin{equation}
 I_k:=[T_*-2^{-k},\,T_*-2^{-k-1}),
 \qquad
 \epsilon_k:=\int_{I_k}Y(t)\,\dd t,
 \qquad
 \sum_{k\ge k_0}\epsilon_k\le\frac{E_*}{2\nu}.
 \label{eq:v54-slabs}
\end{equation}

\begin{lemma}[Leray lower bound on every slab]
\label{lem:v54-lower}
\(\epsilon_k\ge2(1-2^{-1/2})\,c_L^2\nu^{3/2}\,2^{-k/2}\) for every
\(k\ge k_0\).
\end{lemma}

\begin{proof}
Integrate \eqref{eq:v32-leray-rate}:
\(\int_{I_k}(T_*-t)^{-1/2}\dd t=2(2^{-k/2}-2^{-(k+1)/2})\).
\end{proof}

For a ratio \(\theta\in(0,1)\) let also
\(I_k^{(\theta)}:=[T_*-\theta^k,T_*-\theta^{k+1})\) and
\(\epsilon_k^{(\theta)}:=\int_{I_k^{(\theta)}}Y\); the dyadic case is
\(\theta=1/2\), and Lemma~\ref{lem:v54-lower} holds for every ratio with
\(\epsilon_k^{(\theta)}\ge2(1-\theta^{1/2})c_L^2\nu^{3/2}\theta^{k/2}\).

\begin{theorem}[Geometric Leray-share dichotomy]
\label{thm:v54-dichotomy}
\begin{enumerate}[label=\textup{(\roman*)},leftmargin=3.7em]
\item For every ratio \(\theta\) and every \(k\) there is
      \(t_k\in I_k^{(\theta)}\) with
      \((T_*-t_k)^{1/2}Y(t_k)\le(1-\theta)^{-1}\theta^{-k/2}\epsilon_k^{(\theta)}\).
\item The following are equivalent:
      \begin{enumerate}[label=\textup{(\alph*)},leftmargin=2.6em]
      \item there are \(t_j\uparrow T_*\) and \(M<\infty\) with
            \((T_*-t_j)^{1/2}Y(t_j)\le M\) (type-I along a sequence);
      \item for some ratio \(\theta\),
            \(\liminf_{k\to\infty}\theta^{-k/2}\epsilon_k^{(\theta)}<\infty\);
      \item for every ratio \(\theta\ge(1-\theta_M)^{1/2}\), where
            \(\theta_M:=\min\{1/2,\,3\nu^3/(16C_1M^2)\}\) with the \(M\) of
            (a), \(\liminf_k\theta^{-k/2}\epsilon_k^{(\theta)}\le2M\).
      \end{enumerate}
\item Consequently \((T_*-t)^{1/2}Y(t)\to\infty\) as \(t\uparrow T_*\)
      (strong type-II) if and only if
      \(\theta^{-k/2}\epsilon_k^{(\theta)}\to\infty\) for every ratio
      \(\theta\): every geometric slab spends a Leray share large compared
      with its parabolic size \(\theta^{k/2}\) by a factor tending to
      infinity, while \(\sum_k\epsilon_k^{(\theta)}<\infty\).
\end{enumerate}
\end{theorem}

\begin{proof}
(i) The mean of \(Y\) over \(I_k^{(\theta)}\), an interval of length
\((1-\theta)\theta^k\), is \(\epsilon_k^{(\theta)}/((1-\theta)\theta^k)\),
so some \(t_k\) has \(Y(t_k)\) at most this, and
\((T_*-t_k)^{1/2}\le\theta^{k/2}\).  (b)\(\Rightarrow\)(a) is (i) along a
subsequence, and (c)\(\Rightarrow\)(b) is trivial.  (a)\(\Rightarrow\)(c):
the differential inequality \(Y'\le2C_1\nu^{-3}Y^3\) of
Lemma~\ref{lem:v32-leray-rate} gives \(Y(t)\le2Y(t_j)\) on
\([t_j,t_j+\tau_j]\) with \(\tau_j=3\nu^3/(16C_1Y(t_j)^2)\ge\theta_M(T_*-t_j)\),
so \(Y\le2M(T_*-t_j)^{-1/2}\) on the window
\(W_j:=\{t:\ (1-\theta_M)(T_*-t_j)\le T_*-t\le T_*-t_j\}\).  Choose
\(k\) with \(\theta^{k}\le T_*-t_j<\theta^{k-1}\); then
\(\theta^{k+1}\ge\theta^2(T_*-t_j)\ge(1-\theta_M)(T_*-t_j)\), so
\(I_k^{(\theta)}\subset W_j\), and
\(\epsilon_k^{(\theta)}\le(1-\theta)\theta^k\cdot2M(T_*-t_j)^{-1/2}
\le2M\theta^k\theta^{-k/2}=2M\theta^{k/2}\).  As \(j\to\infty\) the
corresponding \(k\to\infty\), which gives (c).  (iii) is the negation of
(ii)(a) and (ii)(b) for every ratio.
\end{proof}

\begin{assessment}[The two regimes in the Leray ledger]
\label{ass:v54-regimes}
V32 separated the regimes by the finiteness of the \(L^{5/2}_tH^1\)
ledger.  Theorem~\ref{thm:v54-dichotomy} gives a second separation, purely
in the Leray ledger, between ``type-I along a sequence of times'' and
``strongly type-II'': the former holds exactly when infinitely many dyadic
slabs spend a Leray share of the parabolic size \(2^{-k/2}\), the latter
exactly when every slab spends more by a divergent factor.  The two
separations are not the same: the type-I enstrophy \emph{rate} of V32
requires the bound at all times near \(T_*\), and strong type-II is its
negation along every sequence; between them lie solutions that are type-I
along some sequences and not others, to which V41--V53 do not apply as
stated.  A finite-Leray-budget schedule with \(2^{k/2}\epsilon_k\to\infty\)
exists: \(\epsilon_k=k\,2^{-k/2}\) is summable and has
\(2^{k/2}\epsilon_k=k\to\infty\), so strong type II is consistent with
the Leray budget.
\end{assessment}

% =============================================================================
\section{Local compactness at a type-I time without ancientness}
\label{sec:v54-forward-window}
% =============================================================================

\begin{proposition}[Forward parabolic window at a type-I time]
\label{prop:v54-forward-window}
Let \(t_0<T_*\) satisfy \((T_*-t_0)^{1/2}Y(t_0)\le M\), and put
\(\lambda:=(T_*-t_0)^{1/2}\), \(\theta_M:=3\nu^3/(16C_1M^2)\).  Then the
rescaled field \(V(s,y):=\lambda u(t_0+\lambda^2s,x+\lambda y)\), for any
centre \(x\), satisfies on \(s\in[0,\theta_M]\)
\begin{equation}
 \|\nabla V(s)\|_{L^2(\Torus^3_\lambda)}^2\le2M,
 \qquad
 \nu\int_0^{\theta_M}\|\Delta V\|_{L^2(\Torus^3_\lambda)}^2\,\dd s
 \le M+2C_1\nu^{-3}\theta_M(2M)^3,
 \label{eq:v54-forward-bounds}
\end{equation}
and \(\|V(s)\|_{L^6}\), \(\|P_V(s)\|_{L^3}\), \(\|\partial_sV\|_{L^2(0,\theta_M;L^2)}\)
are bounded by constants depending only on \(M,\nu\) and the absolute
constants.  Along any sequence of such times \(t_0^{(j)}\uparrow T_*\) with
common \(M\), the rescaled fields converge on compact subsets of
\((0,\theta_M)\times\R^3\), in the sense of Theorem~\ref{thm:v41-ancient-limit},
to a smooth solution \(V\) on \((0,\theta_M)\times\R^3\) with
\(\|\nabla V(s)\|_2^2\le2M\).
\end{proposition}

\begin{proof}
The enstrophy bound is the differential inequality as in the proof of
Theorem~\ref{thm:v54-dichotomy}; the rest is
Proposition~\ref{prop:v41-uniform-bounds} and
Theorem~\ref{thm:v41-ancient-limit} with the interval \([0,\theta_M]\) in
place of \([-S,-\eta]\).
\end{proof}

\begin{firewall}[What the forward window lacks]
\label{fw:v54-forward-lacks}
The limit \(V\) of Proposition~\ref{prop:v54-forward-window} is a
bounded-enstrophy solution on a time interval of fixed length, not an
ancient solution, and nothing in the construction makes it nontrivial: the
type-I time \(t_0\) need not be aligned with a concentration cylinder.
The type-I \emph{rate} of V32, which holds at all times, was what allowed
V41 to combine concentration and compactness.  Along a type-I sequence one
obtains compactness at the chosen times and concentration at the
concentration scales, and the two need not meet; this is the precise
reason V41--V53 assume the rate rather than the sequence condition.
\end{firewall}

% =============================================================================
\section{Integrated V54 conclusion and next acquisition order}
\label{sec:v54-conclusion}
% =============================================================================

\begin{theorem}[What V54 proves]
\label{thm:v54-conclusion}
Relative to V1--V53: the Leray lower bound on every dyadic slab
(Lemma~\ref{lem:v54-lower}); the dyadic Leray-share dichotomy between
type-I along a sequence and strong type-II
(Theorem~\ref{thm:v54-dichotomy}); and the forward-window compactness at a
type-I time (Proposition~\ref{prop:v54-forward-window}).  Items (AC1) and
(AC2) of Programme~\ref{prog:v54} are not advanced, for the reasons
recorded.  The full Navier--Stokes stability/global-regularity theorem
remains unproved.
\end{theorem}

\begin{proof}
The cited results.
\end{proof}

\begin{programme}[V55 acquisition order, with compulsory companion audit]
\label{prog:v55}
\begin{enumerate}[label=\textup{(AC\arabic*)},leftmargin=4.8em]
\item Perform the compulsory review of the current SM and YM notes;
      record digests; import nothing numerical.
\item Intermediate solutions.  For solutions that are type-I along some
      sequence but not at all times, determine which of V41--V53 survive
      when the type-I rate is replaced by the sequence condition together
      with the concentration structure; in particular whether a
      concentration cylinder can be chosen inside a forward window of a
      type-I time.
\item Strong type-II.  Combine \(2^{k/2}\epsilon_k\to\infty\) with the
      record--share and cone conditions of V35--V38 to see whether the
      divergent factor is constrained.
\item The next compulsory companion audit after V55 is V60.
\end{enumerate}
\end{programme}

\begin{assessment}[Position after V54]
\label{ass:v54-position}
The regimes are now separated twice, by the intermediate ledger and by the
dyadic Leray shares, and the gap between ``type-I at all times'' and
``type-I along a sequence'' is identified as the reason the ancient-limit
machinery needs the rate.  No full stability or global-regularity claim
is made.
\end{assessment}

% =============================================================================
\section*{Version V54 append-only record}
\addcontentsline{toc}{section}{Version V54 append-only record}
% =============================================================================

The complete V53 mathematical source was retained before this continuation.
Its SHA--256 digest was
\begin{center}\scriptsize\ttfamily
9cfcdbfbeaa08e677b3237d5dfa720d851054b34ff868d60884c1d60d286997d.
\end{center}
V54 proved the dyadic Leray-share dichotomy between type-I along a
sequence and strong type-II, and the forward-window compactness at a
type-I time, and recorded why the ancient-limit machinery requires the
type-I rate.  No full regularity claim is made.

% =============================================================================
\section{V55: compulsory companion audit, the alignment obstruction, and
strong type II against the price sheet}
\label{sec:v55-entry}
% =============================================================================

This is a twenty-fifth-version continuation.  The compulsory review of
the two companion programmes is performed first; notably, the SM V70
audit identifies in all three programmes the structural gap that fixed
local information need not possess a common all-scale budget, which is
exactly the finding of V39 and V52 here.  Items (AC2)--(AC3) of
Programme~\ref{prog:v55} are then settled negatively with precise
reasons: without the type-I rate, the concentration cylinders supplied by
\(\varepsilon\)-regularity and the compact forward windows supplied by
type-I times need not align, so none of V41--V53 survives on the sequence
condition alone; and strong type II is consistent with every condition of
the type-II price sheet, because the divergent slab factor is compatible
with the record--share, cone, and reciprocal-cutoff bounds.  No global
regularity claim is made.  The next compulsory companion audit is V60.

% =============================================================================
\section{Compulsory V55 review of the current companion notes}
\label{sec:v55-companion-audit}
% =============================================================================

The two current companion files in \texttt{latex\_notes} were inspected
read-only.  The frozen checkpoints are
\begin{align}
 &\texttt{Renewal\_SM\_Einstein\_Continuum\_Research\_Notes\_V70.tex},
 \notag\\[-2pt]
 &\qquad
 \texttt{784fd6f1e103bef7ad432146a76517fdfc99044b4fafcdd92b84ca192877106a},
 \label{eq:v55-SM-checkpoint}\\
 &\texttt{Renewal\_Yang\_Mills\_Mass\_Gap\_Research\_Notes\_V50.tex},
 \notag\\[-2pt]
 &\qquad
 \texttt{226f4204f409a5aeae615738fda88bb644f9d465c6418a5d8d2fc0e7370a328f}.
 \label{eq:v55-YM-checkpoint}
\end{align}
The review concentrated on SM V68--V70 and YM V49--V50.

\emph{SM V68--V70.}  V68 proves that a Newton-smeared improved scalar
stress insertion has a norm diverging like the inverse distance to the
reflection plane, a local boundary gradient contact, whose
contact-subtracted finite part exists as a tempered distribution but is
not automatically positive.  V69 closes the free reduced Einstein--matter
source certificate: factorized positive reflected TT graviton form, an
instantaneous Newton scalar, and a finite nonnegative Osterwalder--Schrader
norm for a joint graviton and two-scalar insertion.  V70 performs its
audit and proves, on an exact two-Gaussian model, that every fixed
interaction moment may exist while the exponential diverges for every
nonzero coupling, with an exact quartic stabilization threshold; it
records that ``fixed local information need not possess a common
all-arity or all-scale budget'' as the shared gap of the three programmes.

\emph{YM V49--V50.}  V49 proves a connected Wick leg inequality and a
grade/mark condition, computes eighth-order cumulants, and disproves the
additive strong cards at arity eight.  V50 performs its audit, proves
degree bounds for heat and Wilson grades, the connected rooted-Wick
coefficient formula, thermal homogeneity cards at every fixed arity, and
isolates a uniform-arity gate.

\begin{theorem}[V55 companion-audit decision]
\label{thm:v55-companion-decision}
\begin{enumerate}[label=\textup{(\roman*)},leftmargin=3.7em]
\item No SM V68--V70 contact, source, or threshold constant and no YM
      V49--V50 cumulant or grade estimate is imported.
\item The SM V70 diagnosis, that fixed local information need not have a
      common all-scale budget, is the content of
      Theorem~\ref{thm:v39-scale-no-go} (per-scale bounds do not sum to a
      bound) and of Firewall~\ref{fw:v52-no-count} (per-point density
      bounds do not count points); it was reached here independently and
      is confirmed, not imported.
\item Neither companion bears on the Liouville gate of V41 or on the
      type-II ledgers.
\end{enumerate}
\end{theorem}

\begin{proof}
The hypotheses of every cited result are free-field, Gaussian-model, or
lattice-gauge objects; none is a Navier--Stokes balance.
\end{proof}

% =============================================================================
\section{The alignment obstruction without the type-I rate}
\label{sec:v55-alignment}
% =============================================================================

\begin{proposition}[Concentration cylinders and type-I windows need not
meet]
\label{prop:v55-alignment}
Let \(T_*<\infty\) be a maximal time and suppose only that
\((T_*-t_j)^{1/2}Y(t_j)\le M\) along some \(t_j\uparrow T_*\).  Let
\(x_0\in\Sigma\) with concentration scales \(r_k\), and let
\(W_j=[t_j,t_j+\theta_M(T_*-t_j)]\) be the forward windows of
Proposition~\ref{prop:v54-forward-window}.  Then:
\begin{enumerate}[label=\textup{(\roman*)},leftmargin=3.7em]
\item the cylinder \(Q_{r_k}(x_0,T_*)\) meets \(W_j\) in a set of times
      of relative measure at least a fixed fraction only if
      \(r_k^2\asymp T_*-t_j\), and the sequences \(r_k\) and
      \((T_*-t_j)^{1/2}\) are chosen independently by the singular point
      and by the enstrophy;
\item if no subsequence satisfies \(r_{k}^2\asymp T_*-t_{j}\), the
      rescaled solutions at the concentration scales have no uniform
      enstrophy bound on any fixed window \((-S,-\eta)\), and the
      compactness of Theorem~\ref{thm:v41-ancient-limit} fails; the
      rescaled solutions at the type-I times have compactness on
      \((0,\theta_M)\) but no lower bound on their dissipation there.
\end{enumerate}
Hence none of V41--V53 is available on the sequence condition alone.
\end{proposition}

\begin{proof}
(i) is the geometry of the two families of intervals: \(W_j\) has length
\(\theta_M(T_*-t_j)\) and ends at \(t_j+\theta_M(T_*-t_j)<T_*\), while
\(Q_{r_k}\) spans \((T_*-r_k^2,T_*)\); their intersection has length
comparable to both only if \(r_k^2\asymp T_*-t_j\).  (ii) restates what
each construction uses: Proposition~\ref{prop:v41-uniform-bounds} needs
\(Y\le C(T_*-t)^{-1/2}\) on \((T_*-Sr_k^2,T_*-\eta r_k^2)\), which the
sequence condition gives only inside the windows \(W_j\); and
\eqref{eq:v41-rescaled-nontrivial} needs a concentration cylinder, which
the type-I times do not supply.
\end{proof}

\begin{assessment}[Why the rate is the right hypothesis]
\label{ass:v55-rate}
The type-I enstrophy rate at all times near \(T_*\) is precisely the
statement that every window is a type-I window, so that every
concentration cylinder is automatically aligned.  Under the weaker
sequence condition the theory of V41--V53 would need an additional
alignment hypothesis, that concentration occurs at the parabolic scales
of the type-I times, which is not implied by anything proved.  The
intermediate class between the rate and strong type II is therefore
governed by the same open alignment question, and V55 records it rather
than assuming it.
\end{assessment}

% =============================================================================
\section{Strong type II against the price sheet}
\label{sec:v55-strong-type-II}
% =============================================================================

\begin{proposition}[Consistency of the divergent slab factor with the
type-II conditions]
\label{prop:v55-consistency}
Let the solution be strongly type II, so that
\(m_k:=2^{k/2}\epsilon_k\to\infty\) with \(\sum_k\epsilon_k<\infty\).
Then the following are all compatible for a cofinal high-total sequence
with cutoffs \(\rho_j\asymp2^{k_j/2}\) and windows in the slabs \(I_{k_j}\):
the high-total share \(\epsilon_{k_j}\ge c\nu\rho_j^{-1}\asymp2^{-k_j/2}\)
of Theorem~\ref{thm:v38-high-total-share}, which is implied by
\(m_{k_j}\to\infty\); the record--share inequality
\(\mathcal R_{w_j}\epsilon_{k_j}\ge c_\theta E_*\) of
Theorem~\ref{thm:v35-record-share}, which bounds the in-window record
below by \(c_\theta E_*2^{k_j/2}/m_{k_j}\) and is compatible with the
lower bound \(\mathcal R_{w_j}\ge(\theta E_*\rho_j)^{1/2}\asymp2^{k_j/4}\)
forced by the high tail; the summable reciprocal records of
Corollary~\ref{cor:v35-summable-reciprocals}; the cone floors and
summable square-root gaps of Theorem~\ref{thm:v37-cone-share}; and the
divergent \(L^{5/2}\) ledger of Corollary~\ref{cor:v32-high-total-divergence}.
An explicit schedule is \(\epsilon_k=k\,2^{-k/2}\), \(\rho_j=2^{j}\),
\(k_j=2j\), \(\mathcal R_{w_j}=2^{j}\), with the corrected spike structure
of Proposition~\ref{prop:v37-corrected-schedule} inside each window.
\end{proposition}

\begin{proof}
Each listed condition is a lower bound of order \(2^{-k_j/2}\) on the
slab share, or a growth condition on records or cutoffs satisfied by
geometric sequences; \(m_{k_j}\to\infty\) only enlarges the shares while
\(\sum_k k2^{-k/2}<\infty\) keeps the budget finite.  The record bound
\(c_\theta E_*2^{k_j/2}/m_{k_j}=c_\theta E_*2^{j}/(2j)\) is below the
schedule's record \(2^j\) for large \(j\), and above the high-tail
requirement \(2^{j/2}\).  The remaining conditions were verified for this
type of schedule in Proposition~\ref{prop:v37-corrected-schedule} and
Assessment~\ref{ass:v38-budget}.
\end{proof}

\begin{firewall}[The type-II regime after V55]
\label{fw:v55-type-II}
Strong type II is the negation of every type-I hypothesis this series
can use, and it is consistent with every necessary condition this series
has proved.  The conditions are all in the Leray and \(L^{5/2}\) ledgers
and all concern global scalar quantities; excluding the regime would
require a mechanism that sees the spatial structure of an enstrophy spike,
which no estimate here provides.  This is recorded as the exact boundary
of the type-II analysis.
\end{firewall}

% =============================================================================
\section{Integrated V55 conclusion and next acquisition order}
\label{sec:v55-conclusion}
% =============================================================================

\begin{theorem}[What V55 proves]
\label{thm:v55-conclusion}
Relative to V1--V54: the companion-audit decision
(Theorem~\ref{thm:v55-companion-decision}); the alignment obstruction
(Proposition~\ref{prop:v55-alignment}); and the consistency of strong
type II with the type-II price sheet
(Proposition~\ref{prop:v55-consistency}).  The full Navier--Stokes
stability/global-regularity theorem remains unproved.
\end{theorem}

\begin{proof}
The cited results.
\end{proof}

\begin{programme}[V56 acquisition order]
\label{prog:v56}
\begin{enumerate}[label=\textup{(AC\arabic*)},leftmargin=4.8em]
\item Spatial structure of spikes.  At a strongly type-II singular
      point, an enstrophy spike of height \(Y\) at time \(t\) with
      \((T_*-t)^{1/2}Y\to\infty\) has, by the Leray rate applied
      backward from \(t\), a floor before it; determine whether
      \(\varepsilon\)-regularity at the spike's own parabolic scale
      \(Y^{-1}\nu^{3}\)-type, rather than at \((T_*-t)^{1/2}\), yields a
      local structure that the global ledgers miss.
\item Alignment.  Determine whether the concentration scales of a
      singular point are constrained by the enstrophy through the
      two-sided slab dissipation of Lemma~\ref{lem:v43-slab} in the
      type-I case, and what replaces that lemma under the sequence
      condition.
\item The next compulsory companion audit is V60.
\end{enumerate}
\end{programme}

\begin{assessment}[Position after V55]
\label{ass:v55-position}
Twenty-five versions after V31, the series has: a two-regime
classification with proved necessary conditions; a complete reduction of
the type-I regime to a Liouville problem for bounded ancient solutions
with finite Dirichlet integral and vortical critical-Morrey trace, with a
detailed description of the singular set and the terminal datum; and a
proof that the type-II regime is consistent with everything proved.  No
full stability or global-regularity claim is made.
\end{assessment}

% =============================================================================
\section*{Version V55 append-only record}
\addcontentsline{toc}{section}{Version V55 append-only record}
% =============================================================================

The complete V54 mathematical source was retained before this continuation.
Its SHA--256 digest was
\begin{center}\scriptsize\ttfamily
20134c1a5188b1c3b9f12fe8105157b8ccbcb0378d2ca67be7f645e583134ed7.
\end{center}
V55 performed the compulsory companion audit, proved the alignment
obstruction that prevents the ancient-limit theory from working under a
sequence-only type-I condition, and verified that strong type II is
consistent with the full type-II price sheet.  No full regularity claim
is made.

% =============================================================================
\section{V56: the time-uniform Morrey bound for the ancient limit and the
spike scale in type II}
\label{sec:v56-entry}
% =============================================================================

Two small improvements and one negative record.  First, the critical
Morrey bound of V47 for the ancient limit, which was stated with an
additional term \((-s)^{1/2}\), holds in the time-uniform form
\(\int_{B_R(y)}|U(s)|^2\le C_MR\) for all \(s<0\), all centres, and all
radii; the ancient limit therefore lies in
\(L^\infty_s\bigl(M^{2,1}(\R^3)\bigr)\), the critical Morrey space, in
addition to \(L^\infty_sL^6\) and, on every half-line, \(L^\infty_s\dot H^1
\cap L^\infty_{s,y}\).  Second, in the type-II regime,
\(\varepsilon\)-regularity at the parabolic scale of an enstrophy value
gives the classical forward smoothing bound and nothing beyond it, which
records why the spatial structure of spikes is invisible to the present
tools.  The alignment item of Programme~\ref{prog:v56} is not advanced.
No global regularity claim is made.  The next compulsory companion audit
is V60.

% =============================================================================
\section{Time-uniform Morrey bound for the ancient limit}
\label{sec:v56-morrey}
% =============================================================================

\begin{theorem}[The ancient limit is bounded in the critical Morrey space
uniformly in time]
\label{thm:v56-uniform-morrey}
For the ancient limit \(U\) of Theorem~\ref{thm:v41-ancient-limit}, and
for every \(y\in\R^3\), \(R>0\), \(s<0\),
\begin{equation}
 \boxed{\quad
 \int_{B_R(y)}|U(s,y')|^2\,\dd y'\ \le\ C_M\,R .
 \quad}
 \label{eq:v56-uniform-morrey}
\end{equation}
The same holds for the rescaling limits of Theorem~\ref{thm:v49-large-scale}
and Theorem~\ref{thm:v49-small-scale}, and for every element of the class
obtained from a type-I singularity by the constructions of V41--V49.
\end{theorem}

\begin{proof}
For the rescaled fields, \(\int_{B_R(y)}|U_k(s)|^2=r_k^{-1}
\int_{B_{Rr_k}(x_0+r_ky)}|u(T_*+r_k^2s)|^2\).  Apply
Theorem~\ref{thm:v47-morrey} with \(\rho=Rr_k\) and \(t=T_*+r_k^2s\),
for \(k\) large enough that \(\rho\le(\delta/2)^{1/2}\) and
\(t\ge T_*-\delta/2\).  In the case \(\rho^2\le T_*-t\), the proof of that
theorem gives the sharper bound \(C\rho^2(T_*-t)^{-1/2}
=CR^2r_k^2\,r_k^{-1}(-s)^{-1/2}=CR^2r_k(-s)^{-1/2}\), and
\(\rho^2\le T_*-t\) means \(R^2\le-s\), so this is at most \(CRr_k\); in
the other case the bound is \(C_M\rho=C_MRr_k\).  Dividing by \(r_k\)
gives \(\int_{B_R(y)}|U_k(s)|^2\le\max\{C,C_M\}R\), and weak lower
semicontinuity passes it to \(U\).  Rescaling limits of \(U\) inherit the
bound by Lemma~\ref{lem:v49-invariance} and lower semicontinuity.
\end{proof}

\begin{corollary}[Function-space description of the ancient limit]
\label{cor:v56-class}
The ancient limit of a type-I singularity satisfies
\begin{equation}
 U\in L^\infty\bigl((-\infty,0);M^{2,1}(\R^3)\bigr),
 \qquad
 \|U(s)\|_{L^6}\le C(-s)^{-1/4},
 \qquad
 \|\nabla U(s)\|_{L^2}^2\le C(-s)^{-1/2},
 \label{eq:v56-class-bounds}
\end{equation}
and, on every half-line \((-\infty,-\eta]\), \(U\), \(\nabla U\),
\(\nabla^2U\) are bounded.  Here \(M^{2,1}\) is the space of locally
square-integrable fields with \(\sup_{y,R}R^{-1}\int_{B_R(y)}|f|^2<\infty\),
which contains \(L^3\) and the weak space \(L^{3,\infty}\) and is
invariant under the Navier--Stokes scaling.
\end{corollary}

\begin{proof}
Theorem~\ref{thm:v56-uniform-morrey}, \eqref{eq:v41-limit-type-I}, and
Proposition~\ref{prop:v53-bounded-ancient}.  The inclusions are H\"older's
inequality and the definition of \(L^{3,\infty}\); the invariance is the
computation \(\int_{B_R}|\lambda f(\lambda\cdot)|^2=\lambda^{-1}
\int_{B_{\lambda R}}|f|^2\le\lambda^{-1}C\lambda R\).
\end{proof}

\begin{assessment}[The gate in function-space terms]
\label{ass:v56-gate}
The Liouville gate of V41 now reads: a nonzero ancient solution on
\(\R^3\) that is bounded on every half-line, has bounded enstrophy there,
lies in \(L^\infty_sM^{2,1}\), and has a critical-Morrey trace at its
final time with vorticity on every exterior region, is impossible.  The
critical Morrey class is the largest scale-invariant class in which local
energy solutions of the Navier--Stokes equations are known to be
constructed; the gate asks for a Liouville theorem in that class with the
additional bounded-enstrophy and boundedness hypotheses.  No such theorem
is proved here.
\end{assessment}

% =============================================================================
\section{The spike scale in type II}
\label{sec:v56-spike-scale}
% =============================================================================

\begin{proposition}[Forward smoothing at the parabolic scale of an
enstrophy value]
\label{prop:v56-forward-smoothing}
Let \(t<T_*\) and \(Y(t)=\|\Lambda u(t)\|_2^2\), and put
\(r_t:=c_1\nu^{3/2}Y(t)^{-1}\) with \(c_1\) small enough that
\(r_t^2\le\nu^3/(8C_1Y(t)^2)\).  Then for every \(x\),
\begin{equation}
 \sup_{B_{r_t/2}(x)\times[t+r_t^2/2,\,t+r_t^2]}|u|
 \le\frac{C_{\rm CKN}}{r_t}
 =\frac{C_{\rm CKN}}{c_1\nu^{3/2}}\,Y(t),
 \label{eq:v56-forward-smoothing}
\end{equation}
provided \(t+r_t^2<T_*\).  This is the classical statement that an
\(H^1\) bound at one time gives an \(L^\infty\) bound after a parabolic
time, and it holds in every regime.
\end{proposition}

\begin{proof}
On \([t,t+r_t^2]\), \(Y\le2Y(t)\) by the differential inequality of
Lemma~\ref{lem:v32-leray-rate}, so \(\|u\|_6\le C_6(2Y(t))^{1/2}\) and
\(\|p\|_3\le C_{\rm CZ}C_6^2\,2Y(t)\) there.  On the cylinder
\(Q_{r_t}(x,t+r_t^2)=B_{r_t}(x)\times(t,t+r_t^2)\) the scale-invariant
cubic quantity is at most \(Cr_t^{3/2}Y(t)^{3/2}=Cc_1^{3/2}\nu^{9/4}\),
which is below \(\varepsilon_2\) for \(c_1\) small; (E5) gives the bound
on \(Q_{r_t/2}\).
\end{proof}

\begin{firewall}[Why the spatial structure of spikes is invisible]
\label{fw:v56-spikes-invisible}
At a strongly type-II time, \((T_*-t)^{1/2}Y(t)\to\infty\), so the
parabolic scale \(r_t\) of the enstrophy value is much smaller than the
parabolic distance \((T_*-t)^{1/2}\) to the singular time:
\(r_t/(T_*-t)^{1/2}=c_1\nu^{3/2}/((T_*-t)^{1/2}Y(t))\to0\).  The forward
smoothing bound \eqref{eq:v56-forward-smoothing} therefore controls the
solution on cylinders far smaller than the singular scale and only after
a delay \(r_t^2\), during which, by the backward cone, the enstrophy may
have changed by a bounded factor but its spatial distribution is
unconstrained.  Every estimate of this series concerns global scalar
quantities of \(u(t)\) or local quantities on cylinders of the singular
scale, and none relates the location of the gradient at a spike to the
location at the next spike.  Item (AC1) of Programme~\ref{prog:v56} is
closed negatively: the present tools do not see spike geometry.
\end{firewall}

% =============================================================================
\section{Integrated V56 conclusion and next acquisition order}
\label{sec:v56-conclusion}
% =============================================================================

\begin{theorem}[What V56 proves]
\label{thm:v56-conclusion}
Relative to V1--V55, modulo (E1)--(E7): the time-uniform critical Morrey
bound \eqref{eq:v56-uniform-morrey} for the ancient limit and its
function-space description \eqref{eq:v56-class-bounds}; and the forward
smoothing bound \eqref{eq:v56-forward-smoothing} with the record that
spike geometry is invisible to the present tools.  Item (AC2) of
Programme~\ref{prog:v56} is not advanced.  The full Navier--Stokes
stability/global-regularity theorem remains unproved.
\end{theorem}

\begin{proof}
The cited results.
\end{proof}

\begin{programme}[V57 acquisition order]
\label{prog:v57}
\begin{enumerate}[label=\textup{(AC\arabic*)},leftmargin=4.8em]
\item Liouville in the critical Morrey class.  Determine what is known,
      as external facts to be recorded and not used, about ancient
      solutions bounded in \(M^{2,1}\) or in \(L^{3,\infty}\), and
      whether any such result applies with the extra hypotheses of
      Corollary~\ref{cor:v56-class}; state exactly which hypothesis is
      missing in each case.
\item Alignment under the sequence condition (item (AC2) of
      Programme~\ref{prog:v56}) remains.
\item The next compulsory companion audit is V60.
\end{enumerate}
\end{programme}

\begin{assessment}[Position after V56]
\label{ass:v56-position}
The ancient limit of a type-I singularity is now placed in the critical
Morrey space uniformly in time, which is the natural scale-invariant
setting for the Liouville question that the type-I regime has become.
The type-II regime is unchanged.  No full stability or global-regularity
claim is made.
\end{assessment}

% =============================================================================
\section*{Version V56 append-only record}
\addcontentsline{toc}{section}{Version V56 append-only record}
% =============================================================================

The complete V55 mathematical source was retained before this continuation.
Its SHA--256 digest was
\begin{center}\scriptsize\ttfamily
4dd9ba8e6ef33608febc79c0492b755f3aea3d8a2ebec3849ddac6f3ae506e64.
\end{center}
V56 proved the time-uniform critical Morrey bound for the ancient limit
of a type-I singularity, described its function-space class, and recorded
through the forward smoothing bound why spike geometry in the type-II
regime is invisible to the present tools.  No full regularity claim is
made.

% =============================================================================
\section{V57: the singular points of the ancient limit and the \(L^3\)
dividing line}
\label{sec:v57-entry}
% =============================================================================

The V51 characterization of the singular set through the terminal datum is
transferred to the ancient limit itself: a point \((y_0,0)\) is a singular
point of \(U\) if and only if the trace \(U_0\) fails to be cubically
integrable near \(y_0\).  The mechanism is a dilation lemma, that
\(\lambda f(\lambda\,\cdot)\rightharpoonup0\) weakly in \(L^3\) for
\(f\in L^3\), which is the exact reason \(L^3\) is the dividing line in
V46 and here.  The version then records, as the final form of the type-I
gate, the list of statements any one of which would close the regime,
each with the hypothesis that is missing.  No global regularity claim is
made.  The next compulsory companion audit is V60.

% =============================================================================
\section{The dilation lemma}
\label{sec:v57-dilation}
% =============================================================================

\begin{lemma}[Critical dilations of an \(L^3\) field vanish weakly]
\label{lem:v57-dilation}
Let \(f\in L^3(\R^3)\) or \(f\in L^3(B_\rho(y_0))\), and
\(f_\lambda(y):=\lambda f(y_0+\lambda y)\).  Then for every \(R>0\),
\(\int_{B_R}|f_\lambda|^2\,\dd y\to0\) as \(\lambda\to0\).  Conversely, if
\(\limsup_{\lambda\to0}\int_{B_R}|f_\lambda|^2>0\) for some \(R\), then
\(f\notin L^3(B_\rho(y_0))\) for every \(\rho>0\).
\end{lemma}

\begin{proof}
\(\int_{B_R}|f_\lambda|^2=\lambda^{-1}\int_{B_{\lambda R}(y_0)}|f|^2
\le\lambda^{-1}\bigl(\int_{B_{\lambda R}(y_0)}|f|^3\bigr)^{2/3}|B_{\lambda R}|^{1/3}
=CR\bigl(\int_{B_{\lambda R}(y_0)}|f|^3\bigr)^{2/3}\to0\) by absolute
continuity of the integral.  The converse is the contrapositive.
\end{proof}

% =============================================================================
\section{Singular points of the ancient limit}
\label{sec:v57-singular-points-U}
% =============================================================================

Let \(U\) be the ancient limit of a type-I singularity, with trace
\(U_0\) (Theorem~\ref{thm:v48-trace-identification}).  Call
\((y_0,0)\) a singular point of \(U\) if \(U\) is unbounded on every
neighbourhood of \((y_0,0)\) in \(\R^3\times(-\infty,0)\).

\begin{theorem}[Singular points of the ancient limit are the
\(L^3\)-singular points of its trace]
\label{thm:v57-singular-points-U}
Modulo (E1)--(E7), for \(y_0\in\R^3\) the following are equivalent:
\begin{enumerate}[label=\textup{(\roman*)},leftmargin=3.7em]
\item \((y_0,0)\) is a singular point of \(U\);
\item \(U_0\notin L^3(B_\rho(y_0))\) for every \(\rho>0\);
\item \(\limsup_{\lambda\to0}\lambda^{-1}\int_{B_{\lambda R}(y_0)}|U_0|^2>0\)
      for some \(R\).
\end{enumerate}
In particular every singular point of \(U\) lies in the ball
\(\overline{B_{R_0}}\) of Lemma~\ref{lem:v46-bounded}, and if \(U_0\in
L^3_{\rm loc}(\R^3)\) then \(U\) has no singular point at \(s=0\) and is
in branch (S2) of Theorem~\ref{thm:v49-small-scale} whenever it is in
case (B).
\end{theorem}

\begin{proof}
(i)\(\Rightarrow\)(ii): at a singular point \((y_0,0)\), (E1) gives
concentration scales \(\lambda_k\to0\) with
\(\lambda_k^{-1}\iint_{Q_{\lambda_k}(y_0,0)}|\nabla U|^2>\varepsilon_1\).
The rescalings \(U^{(k)}(s,y):=\lambda_kU(\lambda_k^2s,y_0+\lambda_ky)\)
are in the class with the bounds of Corollary~\ref{cor:v56-class}, the
compactness of Theorem~\ref{thm:v41-ancient-limit} applies (its proof
used only the enstrophy rate, the \(L^6\) and pressure bounds, and the
concentration inequality), and gives a nontrivial limit \(U'\).  If
\(U_0\in L^3(B_\rho(y_0))\), Lemma~\ref{lem:v57-dilation} gives
\(\int_{B_R}|\lambda_kU_0(y_0+\lambda_k\cdot)|^2\to0\) for every \(R\),
that is, the rescaled traces vanish in \(L^2_{\rm loc}\); the argument of
Lemma~\ref{lem:v46-local-energy-near-zero} and
Proposition~\ref{prop:v46-vanishing}, applied to \(U\) in place of \(u\)
(it used only the \(L^6\), pressure, and enstrophy bounds and the strong
\(L^2_{\rm loc}\) continuity of the trace, which for \(U\) is the weak
continuity of Theorem~\ref{thm:v48-trace-identification} together with the
local energy identity), shows that \(U'\) has vanishing trace, and the
proof of Theorem~\ref{thm:v46-non-L3} gives \(U'\equiv0\), a
contradiction.  (ii)\(\Leftrightarrow\)(iii) is Lemma~\ref{lem:v57-dilation}
together with H\"older as in the proof of
Theorem~\ref{thm:v51-characterization}.  (ii)\(\Rightarrow\)(i): if
\((y_0,0)\) is regular, \(U\) is bounded near it, so \(U_0\) is bounded
near \(y_0\).  The localization to \(\overline{B_{R_0}}\) is
Lemma~\ref{lem:v46-bounded}(i), and the last assertion is
Theorem~\ref{thm:v49-small-scale}.
\end{proof}

\begin{remark}[A point of care in the transfer]
\label{rem:v57-care}
For the original solution the strong \(L^2\) continuity at \(T_*\) was
used to define \(e_R^{(k)}(0^-)\); for \(U\) the trace is only a weak
\(L^2_{\rm loc}\) limit, but Lemma~\ref{lem:v46-local-energy-near-zero}
needs only that the local energy \(e_R(s)\) converges as \(s\uparrow0\)
and that its limit is bounded by the local energy of the trace, which
holds because the local energy identity for \(U\) has absolutely
integrable right side (Proposition~\ref{prop:v47-trace}) and weak limits
do not increase norms.  The bound used is therefore
\(\lim_{s\uparrow0}e_R(s)\le\frac12\int|U_0|^2\phi_R+\) (defect), where
the defect is nonnegative; the argument needs the reverse inequality,
\(e_R(s)\le\lim e_R+C_R(-s)^{1/4}\), which is what
Lemma~\ref{lem:v46-local-energy-near-zero} provides once \(e_R(0^-)\) is
read as \(\lim_{s\uparrow0}e_R(s)\), and then the vanishing of the
rescaled traces in \(L^2_{\rm loc}\) must be replaced by vanishing of the
rescaled local energies \(\lim_{s\uparrow0}e_R^{(k)}(s)\).  These are
\(\lambda_k^{-1}\lim_{s\uparrow0}\frac12\int|U(s)|^2\phi_{\lambda_kR}(\cdot-y_0)\),
and by the local energy identity for \(U\) on \((s,0)\) with the
integrable bounds, they equal \(\lambda_k^{-1}\frac12\int|U_0|^2\phi_{\lambda_kR}\)
plus a defect that vanishes as \(\lambda_k\to0\) by the same
absolute-continuity argument; this closes the transfer.
\end{remark}

% =============================================================================
\section{The type-I gate: what would suffice}
\label{sec:v57-what-would-suffice}
% =============================================================================

\begin{theorem}[Sufficient closures of the type-I regime]
\label{thm:v57-sufficient}
Modulo (E1)--(E7), each of the following statements, if proved, implies
that no smooth solution on \(\Torus^3\) has a singularity with a type-I
enstrophy rate.
\begin{enumerate}[label=\textup{(S\arabic*)},leftmargin=3.8em]
\item Every bounded ancient solution on \(\R^3\) with uniformly bounded
      enstrophy and \(L^6\) values is zero.
\item Every element of \(\mathcal A_\nu(C)\) whose trace is in
      \(L^3_{\rm loc}(\R^3)\) is zero, together with: every element of
      \(\mathcal A_\nu(C)\) has trace in \(L^3_{\rm loc}(\R^3)\).
\item Every element of \(\mathcal A_\nu(C)\) with a critical-Morrey trace
      that has vorticity on every exterior region is zero.
\item The dispersion branches (L2) and (S2) are impossible, and the
      concentration branches (L1) and (S1) admit a fixed point that is
      self-similar and zero.
\end{enumerate}
The missing hypotheses are, respectively: in (S1) nothing, this is the
bounded-ancient Liouville conjecture in the finite-Dirichlet subclass; in
(S2) the second half, since Theorem~\ref{thm:v57-singular-points-U} makes
the first half a consequence of branch (S2) being impossible; in (S3)
nothing beyond what V48--V49 proved about the trace; in (S4) both halves.
\end{theorem}

\begin{proof}
(S1) contains the ancient limit by Proposition~\ref{prop:v53-bounded-ancient}.
(S2): by Theorem~\ref{thm:v57-singular-points-U}, an \(L^3_{\rm loc}\)
trace puts \(U\) in branch (S2) or in case (A); the first half then
excludes both.  (S3) restates the gate with the V48--V49 trace properties.
(S4) restates Theorems~\ref{thm:v49-large-scale}--\ref{thm:v49-small-scale}.
\end{proof}

% =============================================================================
\section{Integrated V57 conclusion and next acquisition order}
\label{sec:v57-conclusion}
% =============================================================================

\begin{theorem}[What V57 proves]
\label{thm:v57-conclusion}
Relative to V1--V56, modulo (E1)--(E7): the dilation lemma
(Lemma~\ref{lem:v57-dilation}); the characterization of the singular
points of the ancient limit by its trace
(Theorem~\ref{thm:v57-singular-points-U}); and the list of sufficient
closures (Theorem~\ref{thm:v57-sufficient}).  The full Navier--Stokes
stability/global-regularity theorem remains unproved.
\end{theorem}

\begin{proof}
The cited results.
\end{proof}

\begin{programme}[V58 acquisition order]
\label{prog:v58}
\begin{enumerate}[label=\textup{(AC\arabic*)},leftmargin=4.8em]
\item Trace integrability.  Determine whether the trace \(U_0\) is in
      \(L^3_{\rm loc}\) away from the origin using the exterior smoothness
      and decay of Proposition~\ref{prop:v49-trace-smooth-decay}, and
      whether the singular points of \(U\) at \(s=0\) can be reduced to
      the origin by the normalization of Proposition~\ref{prop:v42-normalization}.
\item Branch (S2) revisited with Theorem~\ref{thm:v57-singular-points-U}:
      in (S2) the trace is \(L^3_{\rm loc}\) everywhere and smooth; test
      whether a bounded ancient solution on a slab with smooth trace and
      enstrophy escaping to infinity contradicts the time-uniform Morrey
      bound through the local energy identity on large balls.
\item The next compulsory companion audit is V60.
\end{enumerate}
\end{programme}

\begin{assessment}[Position after V57]
\label{ass:v57-position}
The ancient limit now mirrors the original solution: its singular points
at the final time are exactly the \(L^3\)-singular points of its trace,
and all lie in a fixed ball.  The gate is stated as four sufficient
closures with their missing hypotheses.  No full stability or
global-regularity claim is made.
\end{assessment}

% =============================================================================
\section*{Version V57 append-only record}
\addcontentsline{toc}{section}{Version V57 append-only record}
% =============================================================================

The complete V56 mathematical source was retained before this continuation.
Its SHA--256 digest was
\begin{center}\scriptsize\ttfamily
e462f9aba202065ad773c089ea6a0860783b35768dc235a903ebc189206b015b.
\end{center}
V57 proved the dilation lemma, characterized the singular points of the
ancient limit by the \(L^3\)-singular support of its trace, and listed
the sufficient closures of the type-I regime with their missing
hypotheses.  No full regularity claim is made.

% =============================================================================
\section{V58: the terminal-datum theorems under a window type-I hypothesis}
\label{sec:v58-entry}
% =============================================================================

The arguments of V46--V48 used the type-I enstrophy rate only on the time
spans of the concentration cylinders of the singular point under study,
together with local-in-time compactness.  This version isolates that
hypothesis, a type-I bound on the windows \((T_*-r_k^2,T_*)\) of the
concentration scales \(r_k\) of one singular point, and proves under it
that the left local energy of the solution around that point is of
scale-critical size along the concentration scales, and that the terminal
datum, when it is a strong local limit, is not cubically integrable there.
This is the alignment condition sought in V55, made explicit: it is not
implied by type-I along a sequence of times, but it is exactly what the
terminal-datum theorems need.  Items (AC1)--(AC2) of
Programme~\ref{prog:v58} are recorded as not advanced.  No global
regularity claim is made.  The next compulsory companion audit is V60.

% =============================================================================
\section{The window hypothesis}
\label{sec:v58-hypothesis}
% =============================================================================

Let \(T_*<\infty\) be a maximal time, let \(x_0\in\Sigma\) be a singular
point, and let \(r_k\downarrow0\) be concentration scales,
\(r_k^{-1}\iint_{Q_{r_k}(x_0,T_*)}|\nabla u|^2>\varepsilon_1\).  The
\emph{window type-I hypothesis} at \((x_0,(r_k))\) is
\begin{equation}
 \sup_k\ \sup_{T_*-r_k^2<t<T_*}(T_*-t)^{1/2}\,\|\Lambda u(t)\|_2^2
 \ \le\ M<\infty .
 \tag{W}\label{eq:v58-window-type-I}
\end{equation}
It is implied by the global rate \eqref{eq:v32-type-I} with
\(M=C_I^2\nu^{3/2}\), and it implies type-I along the sequence
\(t_k=T_*-r_k^2\), but not conversely.  Define the left local energies
\begin{equation}
 \ell_\rho(x_0):=\lim_{t\uparrow T_*}\int_{B_\rho(x_0)}|u(t)|^2\phi\Bigl(\frac{\cdot-x_0}{\rho}\Bigr),
 \label{eq:v58-left-local-energy}
\end{equation}
whenever the limit exists.

\begin{lemma}[Under (W) the left local energies exist at the concentration
scales]
\label{lem:v58-left-energy-exists}
Under \eqref{eq:v58-window-type-I}, for every \(k\) and every \(R\) with
\(\rho=Rr_k\), the limit \eqref{eq:v58-left-local-energy} exists, and for
\(T_*-r_k^2\le t<T_*\)
\begin{equation}
 \int|u(t)|^2\phi_\rho\ \le\ \ell_\rho(x_0)+C_{R,M}\,r_k\Bigl(\frac{T_*-t}{r_k^2}\Bigr)^{1/4}.
 \label{eq:v58-left-energy-bound}
\end{equation}
\end{lemma}

\begin{proof}
On the window, \eqref{eq:v58-window-type-I} gives
\(\|u(\tau)\|_6\le C_6M^{1/2}(T_*-\tau)^{-1/4}\) and
\(\|p(\tau)\|_3\le C_{\rm CZ}C_6^2M(T_*-\tau)^{-1/2}\), so the right side
of the local energy identity on \(B_{2\rho}\) is absolutely integrable up
to \(T_*\), as in the proof of Lemma~\ref{lem:v46-local-energy-near-zero}
after rescaling; hence \(\int|u(t)|^2\phi_\rho\) is Cauchy as
\(t\uparrow T_*\), and the same computation gives
\eqref{eq:v58-left-energy-bound}.
\end{proof}

% =============================================================================
\section{Scale-critical left local energy}
\label{sec:v58-main}
% =============================================================================

\begin{theorem}[Window type-I forces scale-critical left local energy]
\label{thm:v58-critical-left-energy}
Modulo (E1)--(E7), under \eqref{eq:v58-window-type-I} there is \(R\ge1\)
with
\begin{equation}
 \boxed{\quad
 \liminf_{k\to\infty}\ \frac1{r_k}\,\ell_{Rr_k}(x_0)\ >\ 0 .
 \quad}
 \label{eq:v58-critical-left-energy}
\end{equation}
If in addition \(u(t)\to u(T_*)\) strongly in \(L^2(B_{\rho_0}(x_0))\) for
some \(\rho_0>0\), then \(u(T_*)\notin L^3(B_\rho(x_0))\) for every
\(\rho>0\).
\end{theorem}

\begin{proof}
Rescale: \(U_k(s,y)=r_ku(T_*+r_k^2s,x_0+r_ky)\) on \(s\in(-1,0)\).  By
\eqref{eq:v58-window-type-I}, \(Y_k(s)\le M(-s)^{-1/2}\) on \((-1,0)\),
and the bounds of Proposition~\ref{prop:v41-uniform-bounds} hold on every
\((-1,-\eta)\); the compactness of Theorem~\ref{thm:v41-ancient-limit}
gives a limit \(U\) on \((-1,0)\times\R^3\), smooth, with
\(\|\nabla U(s)\|_2^2\le M(-s)^{-1/2}\), \(\|U(s)\|_6\le C(-s)^{-1/4}\),
\(\|P(s)\|_3\le C(-s)^{-1/2}\), and nontrivial by
\eqref{eq:v41-limit-nontrivial} (whose proof used only the window).
Suppose \eqref{eq:v58-critical-left-energy} fails for every \(R\); passing
to a diagonal subsequence, \(r_k^{-1}\ell_{2Rr_k}(x_0)\to0\) for every
\(R\).  Then \(e_R^{(k)}(0^-)=\frac1{2r_k}\ell_{2Rr_k}(x_0)\to0\) by the
change of variables, and \eqref{eq:v58-left-energy-bound} rescaled is
\(e_R^{(k)}(s)\le e_R^{(k)}(0^-)+C_{R,M}(-s)^{1/4}\); the limit satisfies
\(e_R(s)\le C_{R,M}(-s)^{1/4}\), so \(U(s)\to0\) in \(L^2_{\rm loc}\) as
\(s\uparrow0\).  Lemma~\ref{lem:v46-bounded}(i)--(ii) hold for \(U\) on
\((-1,0)\), since their proofs used only the \(L^6\), pressure, and
enstrophy bounds on \((-1,0)\) and (E5).  The proof of
Theorem~\ref{thm:v46-non-L3} then gives \(U\equiv0\) on
\(\R^3\times(-1/16,0)\) and, by the backward iteration through
\((-1/8,-1/16)\), \((-1/4,-1/8)\), \((-1/2,-1/4)\), \((-1,-1/2)\), on
\(\R^3\times(-1,0)\), contradicting nontriviality.  For the second
statement, strong local convergence identifies \(\ell_\rho(x_0)\) with
\(\int|u(T_*)|^2\phi_\rho\) for small \(\rho\), and
Lemma~\ref{lem:v57-dilation} shows that \(u(T_*)\in L^3(B_\rho(x_0))\)
would force \(r_k^{-1}\ell_{2Rr_k}(x_0)\to0\) for every \(R\).
\end{proof}

\begin{corollary}[Alignment is exactly the window hypothesis]
\label{cor:v58-alignment}
Let the solution be type-I along a sequence \(t_j\uparrow T_*\) with
constant \(M_0\), let \(\theta_{M_0}\) be as in
Proposition~\ref{prop:v54-forward-window}, and let \(x_0\in\Sigma\).  If
infinitely many concentration scales satisfy
\(r_k^2\le\theta_{M_0}(T_*-t_{j(k)})\) with
\(T_*-r_k^2\ge t_{j(k)}\), then \eqref{eq:v58-window-type-I} holds along
those scales with \(M=2M_0(1-\theta_{M_0})^{-1/2}\), and
Theorem~\ref{thm:v58-critical-left-energy} applies at \(x_0\).
\end{corollary}

\begin{proof}
On the forward window \(W_j=[t_j,t_j+\theta_{M_0}(T_*-t_j)]\) one has
\(Y\le2M_0(T_*-t_j)^{-1/2}\), and for \(t\in W_j\),
\((T_*-t)\ge(1-\theta_{M_0})(T_*-t_j)\), so
\((T_*-t)^{1/2}Y(t)\le2M_0(1-\theta_{M_0})^{-1/2}\); the concentration
window \((T_*-r_k^2,T_*)\) lies in \(W_j\) exactly under the stated
inequalities, except for its final part, on which the same bound holds by
the differential inequality extended forward from any point of \(W_j\)
with the same constant, since the extension time
\(3\nu^3/(16C_1Y^2)\) at a point of \(W_j\) is at least
\(\theta_{M_0}(T_*-t_j)/4\), and one repeats the extension finitely often
across \((T_*-r_k^2,T_*)\) with a constant that doubles at each step over
at most a bounded number of steps because the remaining time shrinks
geometrically; the resulting bound is uniform in \(k\).
\end{proof}

\begin{firewall}[The final part of the window]
\label{fw:v58-final-part}
The last step of Corollary~\ref{cor:v58-alignment} extends a type-I bound
from a forward window to the closing part of the concentration window by
iterating the forward smoothing of the differential inequality; each step
doubles the constant and covers a fixed fraction of the remaining time, so
a bound of the form \((T_*-t)^{1/2}Y(t)\le M\) with \(M\) independent of
\(k\) follows only if the number of steps is bounded, which requires the
concentration window to end a fixed fraction of \(T_*-t_j\) before
\(T_*\), that is, \(r_k^2\ge c(T_*-t_j)\).  Together with
\(r_k^2\le\theta_{M_0}(T_*-t_j)\) this is the alignment
\(r_k^2\asymp T_*-t_{j(k)}\) of Proposition~\ref{prop:v55-alignment}, now
stated as a sufficient condition.  Without alignment the corollary does
not apply, and the theorem requires (W) directly.
\end{firewall}

% =============================================================================
\section{Integrated V58 conclusion and next acquisition order}
\label{sec:v58-conclusion}
% =============================================================================

\begin{theorem}[What V58 proves]
\label{thm:v58-conclusion}
Relative to V1--V57, modulo (E1)--(E7): existence of the left local
energies under the window hypothesis (Lemma~\ref{lem:v58-left-energy-exists});
scale-critical left local energy along the concentration scales, and
non-\(L^3\) terminal datum when the terminal limit is strong
(Theorem~\ref{thm:v58-critical-left-energy}); and the identification of the
alignment condition with the window hypothesis
(Corollary~\ref{cor:v58-alignment}, Firewall~\ref{fw:v58-final-part}).
Items (AC1)--(AC2) of Programme~\ref{prog:v58} are not advanced.  The
full Navier--Stokes stability/global-regularity theorem remains unproved.
\end{theorem}

\begin{proof}
The cited results.
\end{proof}

\begin{programme}[V59 acquisition order]
\label{prog:v59}
\begin{enumerate}[label=\textup{(AC\arabic*)},leftmargin=4.8em]
\item Window Morrey bound.  Determine whether the critical Morrey bound of
      V47 holds at the concentration scales under (W) alone, giving a
      two-sided scale-critical left local energy along \(r_k\).
\item Strong local limit.  Determine whether (W) implies strong
      \(L^2\) convergence of \(u(t)\) on \(B_{\rho}(x_0)\), through the
      rescaled tail smallness restricted to the window, so that the second
      statement of Theorem~\ref{thm:v58-critical-left-energy} becomes
      unconditional.
\item Prepare the compulsory V60 companion audit.
\end{enumerate}
\end{programme}

\begin{assessment}[Position after V58]
\label{ass:v58-position}
The terminal-datum theorems are now known to depend only on a type-I
bound over the concentration windows of one singular point, which is the
precise alignment condition between type-I times and concentration scales.
Under the global rate the condition is automatic; under type-I along a
sequence it is an additional hypothesis.  No full stability or
global-regularity claim is made.
\end{assessment}

% =============================================================================
\section*{Version V58 append-only record}
\addcontentsline{toc}{section}{Version V58 append-only record}
% =============================================================================

The complete V57 mathematical source was retained before this continuation.
Its SHA--256 digest was
\begin{center}\scriptsize\ttfamily
2dd0b5036de732e1171cd9b601e69d0c006313579dfb324349a10a4e85d965b5.
\end{center}
V58 isolated the window type-I hypothesis, proved under it that the left
local energy around a singular point is scale-critical along the
concentration scales and that a strong terminal limit is not locally
cubically integrable, and identified alignment with the window hypothesis.
No full regularity claim is made.

% =============================================================================
\section{V59: retraction of the alignment corollary and the window Morrey
bound}
\label{sec:v59-entry}
% =============================================================================

Corollary~\ref{cor:v58-alignment} of V58 is withdrawn.  Its proof
extended a type-I bound from a forward window to the closing part of a
concentration window by iterating the differential inequality
\(Y'\le2C_1\nu^{-3}Y^3\), claiming that each step covers a fixed fraction
of the remaining time.  That claim is false: after \(m\) doublings of the
enstrophy bound the step length is a fixed fraction of \(T_*-t_j\)
divided by \(4^m\), so the iteration covers only a bounded fraction of
\(T_*-t_j\) in total and never reaches \(T_*\).  Consequently a type-I
bound at one time, or along a sequence of times, never implies the window
hypothesis (W) of V58, which requires the bound up to the singular time
itself; the window hypothesis is strictly stronger than alignment.
Theorem~\ref{thm:v58-critical-left-energy} is unaffected, since it assumes
(W) directly.  The version then proves the window Morrey bound at scales
at most the concentration scale, which gives a two-sided statement only
when the radius in Theorem~\ref{thm:v58-critical-left-energy} can be taken
at most one, and records item (AC2) of Programme~\ref{prog:v59} as open.
No global regularity claim is made.  The compulsory companion audit is
due in V60.

% =============================================================================
\section{The retraction}
\label{sec:v59-retraction}
% =============================================================================

\begin{counterexample}[The forward extension does not reach the singular
time]
\label{cex:v59-forward-extension}
Suppose \(Y(t)\le A\) at a time \(t\).  The differential inequality gives
\(Y\le2^{m+1}A\) on \([t,t+\sigma_m]\) with
\(\sigma_m=\sum_{i=0}^{m}\frac{3\nu^3}{16C_1(2^iA)^2}
=\frac{3\nu^3}{16C_1A^2}\sum_{i=0}^m4^{-i}<\frac{\nu^3}{4C_1A^2}\).
Thus no number of iterations extends a bound of the type
\(Y\le2^{m+1}A\) beyond time \(t+\nu^3/(4C_1A^2)\), which is exactly the
Leray existence-time bound.  With \(A=2M_0(T_*-t_j)^{-1/2}\) at the end of
a forward window, the total extension is at most
\(\nu^3(T_*-t_j)/(16C_1M_0^2)\), a fixed fraction of \(T_*-t_j\), and
the bound \((T_*-t)^{1/2}Y(t)\le M\) is obtained only on a set of times
whose distance to \(T_*\) is bounded below by a fixed fraction of
\(T_*-t_j\).  The proof of Corollary~\ref{cor:v58-alignment} asserted the
contrary and is therefore invalid.
\end{counterexample}

\begin{theorem}[Corrected statement]
\label{thm:v59-corrected}
Let the solution be type-I along a sequence \(t_j\uparrow T_*\) with
constant \(M_0\).  Then for every \(j\) there is a window
\(W_j'=[t_j,\,T_*-\gamma(T_*-t_j)]\), with
\(\gamma=\gamma(M_0,\nu)\in(0,1)\), on which
\((T_*-t)^{1/2}Y(t)\le M\) with \(M=M(M_0,\nu)\); no bound near
\(T_*\) follows.  In particular the window hypothesis (W), which requires
the bound on \((T_*-r_k^2,T_*)\), is not implied by type-I along a
sequence, whether or not the concentration scales are aligned with the
sequence.
\end{theorem}

\begin{proof}
The first statement is Proposition~\ref{prop:v54-forward-window} with the
extension of Counterexample~\ref{cex:v59-forward-extension}.  The second
is the definition of (W) together with the counterexample.
\end{proof}

\begin{firewall}[Status of V58 after the retraction]
\label{fw:v59-status-of-V58}
Lemma~\ref{lem:v58-left-energy-exists} and
Theorem~\ref{thm:v58-critical-left-energy} assume (W) and are unchanged.
Corollary~\ref{cor:v58-alignment} and Firewall~\ref{fw:v58-final-part}
are withdrawn; the sentence in Assessment~\ref{ass:v58-position} stating
that (W) ``is the precise alignment condition'' is corrected to: (W) is a
stronger condition than alignment, and it is the hypothesis that the
terminal-datum theorems need.  Assessment~\ref{ass:v55-rate} stands: the
global rate is the right hypothesis because it supplies (W) at every
singular point.
\end{firewall}

% =============================================================================
\section{The window Morrey bound}
\label{sec:v59-window-morrey}
% =============================================================================

\begin{proposition}[Critical Morrey bound at scales up to the
concentration scale]
\label{prop:v59-window-morrey}
Under (W) at \((x_0,(r_k))\), for every \(k\), every \(x\in\Torus^3\),
every \(0<\rho\le r_k\), and every \(t\in[T_*-\rho^2,T_*)\),
\begin{equation}
 \int_{B_\rho(x)}|u(t)|^2\le C_M'\,\rho ,
 \qquad
 \ell_\rho(x_0)\le C_M'\,\rho ,
 \label{eq:v59-window-morrey}
\end{equation}
with \(C_M'\) depending only on \(M,\nu,C_6,C_{\rm CZ},\phi\).
Consequently, if the radius \(R\) in
Theorem~\ref{thm:v58-critical-left-energy} can be taken at most one, the
left local energy at \(x_0\) is two-sided along the concentration scales:
\(c\,r_k\le\ell_{Rr_k}(x_0)\le C_M'Rr_k\) for \(k\) large.
\end{proposition}

\begin{proof}
The proof of Theorem~\ref{thm:v47-morrey} uses the rate only on
\([t-\rho^2,T_*)\subset(T_*-r_k^2,T_*)\), where (W) gives
\(\|u(\tau)\|_6\le C_6M^{1/2}(T_*-\tau)^{-1/4}\) and the pressure bound;
repeat it verbatim, and let \(t\uparrow T_*\) using
Lemma~\ref{lem:v58-left-energy-exists}.
\end{proof}

\begin{firewall}[Why the radius may exceed one]
\label{fw:v59-radius}
The proof of Theorem~\ref{thm:v58-critical-left-energy} obtains the
vanishing of the trace of the limit on every ball, hence on an exterior
region, from vanishing of the rescaled left local energies for all
\(R\); its conclusion therefore asserts positivity for some \(R\), which
may be large.  A large \(R\) means that the scale-critical energy sits at
distances up to \(Rr_k\) from \(x_0\), where (W) gives no Morrey bound
unless \(Rr_k\le r_{k'}\) for an earlier concentration scale; in that
case the bound holds with the window of \(r_{k'}\).  Under the global
rate every radius is covered, which is why V47 is unconditional.
\end{firewall}

\begin{firewall}[Item (AC2) is open]
\label{fw:v59-AC2}
Strong \(L^2\) convergence of \(u(t)\) on a ball around \(x_0\) under
(W) alone is not proved.  The V32 argument used the uniform Fourier tail
smallness, which is global and needs the rate on a full interval; a local
version would need a localized tail estimate, which the present tools do
not provide.
\end{firewall}

% =============================================================================
\section{Integrated V59 conclusion and next acquisition order}
\label{sec:v59-conclusion}
% =============================================================================

\begin{theorem}[What V59 establishes]
\label{thm:v59-conclusion}
Relative to V1--V58: the retraction of Corollary~\ref{cor:v58-alignment}
with the counterexample \ref{cex:v59-forward-extension} and the corrected
Theorem~\ref{thm:v59-corrected}; and the window Morrey bound
\eqref{eq:v59-window-morrey}.  The full Navier--Stokes
stability/global-regularity theorem remains unproved.
\end{theorem}

\begin{proof}
The cited results.
\end{proof}

\begin{programme}[V60 acquisition order, with compulsory companion audit]
\label{prog:v60}
\begin{enumerate}[label=\textup{(AC\arabic*)},leftmargin=4.8em]
\item Perform the compulsory review of the current SM and YM notes;
      record digests; import nothing numerical.
\item Audit V41--V59 for any other use of a forward extension of a
      type-I bound; confirm that every result assuming the global rate
      is unaffected.
\item Consolidate the hypotheses actually used by each theorem of
      V41--V58 into a table: global rate, window hypothesis, or none.
\end{enumerate}
\end{programme}

\begin{assessment}[Position after V59]
\label{ass:v59-position}
An error in V58 was found and withdrawn within the series: type-I along a
sequence does not extend to the singular time, so the window hypothesis
is genuinely stronger than alignment.  Everything proved under the global
rate stands.  No full stability or global-regularity claim is made.
\end{assessment}

% =============================================================================
\section*{Version V59 append-only record}
\addcontentsline{toc}{section}{Version V59 append-only record}
% =============================================================================

The complete V58 mathematical source was retained before this continuation.
Its SHA--256 digest was
\begin{center}\scriptsize\ttfamily
609a19e99e8609407db2db903d2b4d873c0f0fdfd770f87dcd82821979c45437.
\end{center}
V59 withdrew the alignment corollary of V58 with an explicit
counterexample to its forward-extension step, stated the corrected
theorem, and proved the window Morrey bound at scales up to the
concentration scale.  No full regularity claim is made.

% =============================================================================
\section{V60: compulsory companion audit, internal audit, and the hypothesis
table}
\label{sec:v60-entry}
% =============================================================================

This is a thirtieth-version continuation.  The compulsory review of the
two companion programmes is performed first.  The version then carries
out the internal audit requested by Programme~\ref{prog:v60}: every place
in V41--V59 where a type-I bound is propagated in time is checked, and no
other use of the invalid forward extension of V58 is found.  Finally the
hypotheses actually used by each theorem of V32--V59 are collected in one
table, with three entries: none beyond smoothness and the Leray identity,
the global type-I enstrophy rate, or the window hypothesis (W).  No
global regularity claim is made.  The next compulsory companion audit is
V65.

% =============================================================================
\section{Compulsory V60 review of the current companion notes}
\label{sec:v60-companion-audit}
% =============================================================================

The current companion files in \texttt{latex\_notes} were inspected
read-only.  The frozen checkpoints are
\begin{align}
 &\texttt{Renewal\_SM\_Einstein\_Continuum\_Research\_Notes\_V71.tex},
 \notag\\[-2pt]
 &\qquad
 \texttt{fbbeafff3167299cb8882df4550126ddd5de12b027af6912fcd2f52da1ee1410},
 \label{eq:v60-SM-checkpoint}\\
 &\texttt{Renewal\_Yang\_Mills\_Mass\_Gap\_Research\_Notes\_V51.tex},
 \notag\\[-2pt]
 &\qquad
 \texttt{226f4204f409a5aeae615738fda88bb644f9d465c6418a5d8d2fc0e7370a328f}.
 \label{eq:v60-YM-checkpoint}
\end{align}
At the time of reading, the YM V51 file had the same digest as the YM V50
file audited in V55, that is, it was a renamed copy not yet extended; the
YM review therefore adds nothing to V55.

\emph{SM V71.}  For an explicit two-mode packet the transverse--traceless
source form and the quadratic Newton source are computed exactly and
shown to cross finite frequency horizons at fixed couplings; including the
Higgs potential in the linearized Newton source produces a negative
amplitude-eight term and destroys global normalizability.  The
linearized global stability gate of SM V70 is thereby disproved by an
explicit computation, and the programme is directed to a nonlinear or
effective-field-theory fork.

\begin{theorem}[V60 companion-audit decision]
\label{thm:v60-companion-decision}
No SM V71 constant is imported.  The SM practice of disproving a
conjectured gate by an explicit packet, and V59's withdrawal of a
corollary by an explicit counterexample, are the same discipline;
nothing further transfers.  Neither companion bears on the type-I
Liouville gate or the type-II ledgers.
\end{theorem}

\begin{proof}
The SM V71 hypotheses are free-field packets and quartic forms; none is a
Navier--Stokes balance.
\end{proof}

% =============================================================================
\section{Internal audit of time propagation of type-I bounds}
\label{sec:v60-internal-audit}
% =============================================================================

\begin{proposition}[No other forward extension is used]
\label{prop:v60-internal-audit}
In V41--V59 the differential inequality \(Y'\le2C_1\nu^{-3}Y^3\) is used
to propagate an enstrophy bound forward in time in exactly the following
places, each of which is valid.
\begin{enumerate}[label=\textup{(\roman*)},leftmargin=3.7em]
\item Proposition~\ref{prop:v54-forward-window}: a bound on a forward
      window of length a fixed fraction of \(T_*-t_0\), which is the
      correct single-step statement.
\item Theorem~\ref{thm:v54-dichotomy}: the same window, used only to
      locate a geometric slab inside it; no extension toward \(T_*\).
\item Proposition~\ref{prop:v56-forward-smoothing}: a single step at the
      parabolic scale of the enstrophy value.
\item Corollary~\ref{cor:v58-alignment}: withdrawn in V59.
\end{enumerate}
All other results of V41--V59 assume the global rate \eqref{eq:v32-type-I}
or the window hypothesis (W) and use the backward cone
\eqref{eq:v37-backward-cone}, which is a backward propagation and is
valid.
\end{proposition}

\begin{proof}
Inspection of the proofs; each use is cited at its occurrence.
\end{proof}

% =============================================================================
\section{Hypothesis table for V32--V59}
\label{sec:v60-table}
% =============================================================================

\begin{theorem}[Hypotheses used]
\label{thm:v60-table}
Beyond smoothness on \([0,T_*)\times\Torus^3\), the Leray identity, and
the listed external imports, the results of V32--V59 use the following
hypotheses.
\begin{center}
\small
\begin{tabular}{p{0.30\textwidth}p{0.62\textwidth}}
\toprule
\textbf{Hypothesis} & \textbf{Results}\\
\midrule
None &
V32 cutoff-free density, exit debits, recent-ledger tail bound, record
bounds, Leray rate, high-total divergence; V33 fractional record law,
running-record bound; V34 per-scale law, Leray concentration; V35
record--share, reciprocal records; V36 \(\dot H^s\) laws, sign
decomposition, two-slot alternative; V37 backward cone, gap summability,
schedule; V38 witness, cone in the fractional integral, reciprocal
cutoffs; V39 tail-weighted densities, scale-invariance no-go; V40 exit
covariance; V54 geometric Leray-share dichotomy and forward window; V55
alignment obstruction and type-II consistency; V56 forward smoothing.\\[4pt]
Global type-I enstrophy rate \eqref{eq:v32-type-I} &
V32 uniform tails and \(L^2\)-precompactness; V33 \(\log^2\) bound; V34
flat profile, log-free criterion, raw-pulse finiteness; V35 slot no-go;
V36--V37 local energies and pressure; V39 own-window scaling; V41--V53
ancient limit, Liouville reduction, dimension zero, localization,
terminal datum, \(L^\infty\) upgrade, Morrey bound, trace, dichotomies,
density bounds, per-scale implication; V56 uniform Morrey; V57 singular
points of the limit; V59 corrected statement (sequence condition only).\\[4pt]
Window hypothesis (W) at one singular point &
V58 left local energies and scale-critical left energy (with the non-\(L^3\)
conclusion also needing a strong local terminal limit); V59 window Morrey
bound.\\
\bottomrule
\end{tabular}
\end{center}
Every result in the second and third rows becomes unconditional under
the global rate; no result in the first row needs any rate.
\end{theorem}

\begin{proof}
Each entry is the hypothesis stated in the cited result.
\end{proof}

\begin{assessment}[Reading the table]
\label{ass:v60-table}
The unconditional row contains the two-regime classification and the
type-II price sheet; the global-rate row contains the entire type-I
theory; the window row contains its partial extension.  The table makes
explicit that the type-II regime has no analogue of the compactness
theory, because every compactness statement needs a rate, and that the
type-I theory is complete up to one Liouville problem.
\end{assessment}

% =============================================================================
\section{Integrated V60 conclusion and next acquisition order}
\label{sec:v60-conclusion}
% =============================================================================

\begin{theorem}[What V60 establishes]
\label{thm:v60-conclusion}
Relative to V1--V59: the companion-audit decision
(Theorem~\ref{thm:v60-companion-decision}); the internal audit
(Proposition~\ref{prop:v60-internal-audit}); and the hypothesis table
(Theorem~\ref{thm:v60-table}).  The full Navier--Stokes
stability/global-regularity theorem remains unproved.
\end{theorem}

\begin{proof}
The cited results.
\end{proof}

\begin{programme}[V61 acquisition order]
\label{prog:v61}
\begin{enumerate}[label=\textup{(AC\arabic*)},leftmargin=4.8em]
\item A rate-free compactness statement.  Determine whether the geometric
      Leray-share dichotomy of V54 can be combined with the concentration
      structure to produce, in the regime of type-I along a sequence, a
      limit on a forward window that is nontrivial; the obstruction is
      alignment, so the target is a pigeonhole over scales showing that
      some concentration scale lies inside some forward window.
\item Type-II: record any rate-free consequence of the backward cone for
      the spatial distribution of dissipation, through the local energy
      identity with the Leray-rate floor as the only input.
\item The next compulsory companion audit is V65.
\end{enumerate}
\end{programme}

\begin{assessment}[Position after V60]
\label{ass:v60-position}
Thirty versions after V30, the series is internally audited, its
hypotheses are tabulated, and its one withdrawn corollary is documented
with a counterexample.  No full stability or global-regularity claim is
made.
\end{assessment}

% =============================================================================
\section*{Version V60 append-only record}
\addcontentsline{toc}{section}{Version V60 append-only record}
% =============================================================================

The complete V59 mathematical source was retained before this continuation.
Its SHA--256 digest was
\begin{center}\scriptsize\ttfamily
9c9c3a0f810aced3728362ba19a1eac53da1abb732f7318c47dd3176ca4cc6bd.
\end{center}
V60 performed the compulsory companion audit, the internal audit of
time-propagation arguments, and tabulated the hypotheses of V32--V59.
No full regularity claim is made.

% =============================================================================
\section{V61: the trace at infinity and the large-scale branch}
\label{sec:v61-entry}
% =============================================================================

The V46 argument, applied to the large-scale rescalings of the ancient
limit, gives a dichotomy at infinity that mirrors V57 at the origin:
either the trace \(U_0\) carries scale-critical energy at large scales,
\(\limsup_{R\to\infty}R^{-1}\int_{B_R}|U_0|^2>0\), or the ancient limit
disperses at large scales in the centred sense, its centred dissipation
concentration tending to zero.  Items (AC1)--(AC2) of
Programme~\ref{prog:v61} are recorded as not advanced: no pigeonhole over
scales produces alignment, because the forward windows never reach the
singular time (V59), and the backward cone is a global statement with no
spatial content.  No global regularity claim is made.  The next
compulsory companion audit is V65.

% =============================================================================
\section{The large-scale dichotomy}
\label{sec:v61-large-scale}
% =============================================================================

Let \(U\) be the ancient limit of a type-I singularity with trace
\(U_0\), and for \(\lambda>0\) let \(U^\lambda(s,y)=\lambda U(\lambda^2s,\lambda y)\),
whose trace is \(U_0^\lambda:=\lambda U_0(\lambda\,\cdot)\)
(Lemma~\ref{lem:v49-invariance}).

\begin{lemma}[Rescaled traces at large scales]
\label{lem:v61-rescaled-trace}
For every \(R>0\), \(\int_{B_R}|U_0^\lambda|^2=\lambda^{-1}\int_{B_{\lambda R}}|U_0|^2\).
Hence \(\int_{B_R}|U_0^\lambda|^2\to0\) as \(\lambda\to\infty\) for every
\(R\) if and only if \(\rho^{-1}\int_{B_\rho}|U_0|^2\to0\) as
\(\rho\to\infty\).
\end{lemma}

\begin{proof}
Change of variables.
\end{proof}

\begin{theorem}[Critical energy at infinity or centred dispersion]
\label{thm:v61-dichotomy}
Modulo (E1)--(E7), with \(d_U(\lambda)\) the centred dissipation
concentration of \eqref{eq:v49-concentration-function}, the following
implication holds:
\begin{equation}
 \lim_{\rho\to\infty}\frac1\rho\int_{B_\rho}|U_0|^2=0
 \quad\Longrightarrow\quad
 \lim_{\lambda\to\infty}d_U(\lambda)=0,
 \label{eq:v61-implication}
\end{equation}
and every centred large-scale limit of \(U\) is then zero.  Equivalently,
if some centred large-scale rescaling of \(U\) has a nontrivial limit,
then \(\limsup_{\rho\to\infty}\rho^{-1}\int_{B_\rho}|U_0|^2>0\): the
trace carries scale-critical energy at infinity.  In all cases
\(\rho^{-1}\int_{B_\rho}|U_0|^2\le C_M\).
\end{theorem}

\begin{proof}
Suppose the left side of \eqref{eq:v61-implication} holds and
\(\limsup_\lambda d_U(\lambda)>0\).  As in Theorem~\ref{thm:v49-large-scale}
with centres at the origin, there are \(\lambda_j\to\infty\) and a
nontrivial limit \(U^{(\infty)}\) of \(U^{\lambda_j}\), satisfying the
bounds of Corollary~\ref{cor:v56-class}.  Its trace is the weak
\(L^2_{\rm loc}\) limit of \(U_0^{\lambda_j}\), and
Lemma~\ref{lem:v61-rescaled-trace} with the hypothesis gives
\(\int_{B_R}|U_0^{\lambda_j}|^2\to0\) for every \(R\).  The local energy
lemma for \(U^{\lambda_j}\) near \(s=0\), Lemma~\ref{lem:v46-local-energy-near-zero}
transferred to the class as in Remark~\ref{rem:v57-care}, then gives
\(e_R(s)\le C_R(-s)^{1/4}\) for \(U^{(\infty)}\), so
\(U^{(\infty)}(s)\to0\) in \(L^2_{\rm loc}\) as \(s\uparrow0\), and the
proof of Theorem~\ref{thm:v46-non-L3} yields \(U^{(\infty)}\equiv0\), a
contradiction.  Hence \(d_U(\lambda)\to0\), and any centred large-scale
limit, being a limit of fields whose dissipation on every fixed cylinder
tends to zero, is zero.  The equivalent formulation is the contrapositive.
The upper bound is Theorem~\ref{thm:v56-uniform-morrey} at the trace.
\end{proof}

\begin{remark}[Centred versus translated rescalings]
\label{rem:v61-centred}
The implication concerns centred rescalings.  Translated rescalings with
centres \(y_j\) see the trace near \(y_j\); if \(|y_j|/\lambda_j\to\infty\)
the relevant quantity is \(\lambda_j^{-1}\int_{B_{\lambda_jR}(y_j)}|U_0|^2\),
which the hypothesis of \eqref{eq:v61-implication} does not control.  The
branch (L2) of Theorem~\ref{thm:v49-large-scale} refers to
\(d_U^*\), the supremum over centres; the present theorem refines only its
centred part.
\end{remark}

\begin{corollary}[Two-sided picture of the trace]
\label{cor:v61-picture}
The trace \(U_0\) of the ancient limit of a type-I singularity is a
divergence-free field in \(M^{2,1}(\R^3)\), smooth with pointwise decay
on an exterior region, with vorticity on every exterior region, with
\(L^3\)-singular points exactly at the singular points of \(U\) at
\(s=0\), all inside \(\overline{B_{R_0}}\); and it has scale-critical
energy at infinity unless the ancient limit disperses at large scales.
\end{corollary}

\begin{proof}
V47--V49, V57, and Theorem~\ref{thm:v61-dichotomy}.
\end{proof}

% =============================================================================
\section{Negative records}
\label{sec:v61-negative}
% =============================================================================

\begin{firewall}[No pigeonhole produces alignment]
\label{fw:v61-no-pigeonhole}
By Theorem~\ref{thm:v59-corrected}, a type-I bound along a sequence lives
on windows that end a fixed fraction of \(T_*-t_j\) before \(T_*\), while
concentration cylinders end at \(T_*\).  No choice of subsequences can
place a cylinder ending at \(T_*\) inside a window ending before \(T_*\).
Item (AC1) of Programme~\ref{prog:v61} is closed negatively.
\end{firewall}

\begin{firewall}[The cone has no spatial content]
\label{fw:v61-cone-spatial}
The backward cone \eqref{eq:v37-backward-cone} bounds the global
enstrophy from below; the local energy identity on a ball involves the
local dissipation, which the cone does not bound below.  Item (AC2) is
closed negatively.
\end{firewall}

% =============================================================================
\section{Integrated V61 conclusion and next acquisition order}
\label{sec:v61-conclusion}
% =============================================================================

\begin{theorem}[What V61 proves]
\label{thm:v61-conclusion}
Relative to V1--V60, modulo (E1)--(E7): the large-scale dichotomy for the
trace, in the centred form (Theorem~\ref{thm:v61-dichotomy} with
Remark~\ref{rem:v61-centred}), and the two-sided picture of the trace
(Corollary~\ref{cor:v61-picture}).  The full Navier--Stokes
stability/global-regularity theorem remains unproved.
\end{theorem}

\begin{proof}
The cited results.
\end{proof}

\begin{programme}[V62 acquisition order]
\label{prog:v62}
\begin{enumerate}[label=\textup{(AC\arabic*)},leftmargin=4.8em]
\item Branch (L2) with critical energy at infinity.  If the trace has
      critical energy at infinity but the centred large-scale limits
      vanish, the energy at infinity must be at centres escaping faster
      than the scale; determine whether the exterior decay of
      Proposition~\ref{prop:v49-trace-smooth-decay} can be quantified to
      exclude this.
\item Exhaust the trace: combine Theorem~\ref{thm:v57-singular-points-U}
      and Theorem~\ref{thm:v61-dichotomy} into a single statement of the
      possible shapes of \(U_0\), and test each against the requirement
      that \(U_0\) be the final-time trace of an element of the class.
\item The next compulsory companion audit is V65.
\end{enumerate}
\end{programme}

\begin{assessment}[Position after V61]
\label{ass:v61-position}
The trace of the ancient limit is now described at the origin, on
exterior regions, and at infinity.  No full stability or
global-regularity claim is made.
\end{assessment}

% =============================================================================
\section*{Version V61 append-only record}
\addcontentsline{toc}{section}{Version V61 append-only record}
% =============================================================================

The complete V60 mathematical source was retained before this continuation.
Its SHA--256 digest was
\begin{center}\scriptsize\ttfamily
5594e8b3c07dedfe329b64df28bc7a366df1a8d2e5cfcfa4b80433c52d9cd9ff.
\end{center}
V61 proved the large-scale dichotomy between critical energy of the trace
at infinity and dispersion of the ancient limit at large scales, in the
centred form, and closed the two remaining alignment and cone items
negatively.  No full regularity claim is made.

% =============================================================================
\section{V62: the shape of the trace}
\label{sec:v62-entry}
% =============================================================================

Item (AC2) of Programme~\ref{prog:v62} is settled by a consolidated
theorem describing every proved property of the final-time trace of the
ancient limit, with one new lemma: a cubically integrable field on
\(\R^3\) has vanishing scale-critical energy on large balls, so a
nontrivial centred large-scale limit forces the trace out of
\(L^3(\R^3)\), mirroring at infinity the \(L^3\) dividing line of V57 at
the origin.  Item (AC1), a quantitative exterior decay, is recorded as
open.  No global regularity claim is made.  The next compulsory companion
audit is V65.

% =============================================================================
\section{Cubic integrability at infinity}
\label{sec:v62-L3-infinity}
% =============================================================================

\begin{lemma}[\(L^3\) fields have no critical energy at infinity]
\label{lem:v62-L3-infinity}
If \(f\in L^3(\R^3)\), then \(\rho^{-1}\int_{B_\rho}|f|^2\to0\) as
\(\rho\to\infty\).  Consequently, if
\(\limsup_{\rho\to\infty}\rho^{-1}\int_{B_\rho}|f|^2>0\), then
\(f\notin L^3(\R^3)\), and in fact \(f\notin L^3(\{|y|\ge R\})\) for every
\(R\).
\end{lemma}

\begin{proof}
For \(\rho_0<\rho\), H\"older gives
\(\int_{B_\rho\setminus B_{\rho_0}}|f|^2\le\|f\|_{L^3(\{|y|\ge\rho_0\})}^2
|B_\rho|^{1/3}=C\rho\,\|f\|_{L^3(\{|y|\ge\rho_0\})}^2\), so
\(\rho^{-1}\int_{B_\rho}|f|^2\le\rho^{-1}\int_{B_{\rho_0}}|f|^2
+C\|f\|_{L^3(\{|y|\ge\rho_0\})}^2\); let \(\rho\to\infty\) and then
\(\rho_0\to\infty\).  The second statement applies the first to
\(f\mathbf 1_{\{|y|\ge R\}}\), whose large-ball energies differ from those
of \(f\) by at most \(\rho^{-1}\int_{B_R}|f|^2\to0\).
\end{proof}

\begin{corollary}[A nontrivial centred large-scale limit forces a
non-\(L^3\) trace at infinity]
\label{cor:v62-trace-not-L3}
If some centred large-scale rescaling of the ancient limit has a
nontrivial limit, then \(U_0\notin L^3(\{|y|\ge R\})\) for every \(R\),
although \(U_0\) is smooth and tends to zero pointwise there.
\end{corollary}

\begin{proof}
Theorem~\ref{thm:v61-dichotomy} and Lemma~\ref{lem:v62-L3-infinity}, with
Proposition~\ref{prop:v49-trace-smooth-decay}.
\end{proof}

% =============================================================================
\section{The shape theorem}
\label{sec:v62-shape}
% =============================================================================

\begin{theorem}[Shape of the final-time trace of a type-I ancient limit]
\label{thm:v62-shape}
Let \(u\) be a smooth solution on \([0,T_*)\times\Torus^3\) with maximal
time \(T_*<\infty\) and a type-I enstrophy rate, let \(U\) be the ancient
limit at a concentration point, and let \(U_0\) be its final-time trace.
Modulo (E1)--(E7):
\begin{enumerate}[label=\textup{(\arabic*)},leftmargin=3.2em]
\item \(U_0\) is a divergence-free field on \(\R^3\) with
      \(\int_{B_R(y)}|U_0|^2\le C_MR\) for all \(y,R\), and it is the
      weak \(L^2_{\rm loc}\) limit both of the rescaled terminal data and
      of \(U(s)\) as \(s\uparrow0\), with
      \(\|U(s)-U_0\|_{H^{-1}(B_R)}\le C_R(-s)^{1/2}\).
\item Its \(L^3\)-singular points are exactly the singular points of
      \(U\) at \(s=0\); they lie in \(\overline{B_{R_0}}\); and
      \(U_0\in L^3_{\rm loc}\) away from them.
\item On \(\{|y|>R_0\}\) it is smooth, tends to zero with its gradient
      at infinity, and its vorticity does not vanish identically on any
      exterior region.
\item Either \(\limsup_{\rho\to\infty}\rho^{-1}\int_{B_\rho}|U_0|^2>0\),
      in which case \(U_0\notin L^3\) on every exterior region, or every
      centred large-scale limit of \(U\) is zero.
\item If \(U_0\in L^3_{\rm loc}(\R^3)\), then \(U\) has no singular point
      at \(s=0\), so \(U\) is either eternal or in the dispersion branch
      (S2) with enstrophy escaping to infinity.
\end{enumerate}
No proved property excludes \(U_0\) from being, on exterior regions, a
smooth field of size comparable to \(|y|^{-1}\) with nonvanishing curl in
every annulus, nor from being singular at finitely or infinitely many
points of \(\overline{B_{R_0}}\).
\end{theorem}

\begin{proof}
(1) Theorem~\ref{thm:v56-uniform-morrey}, Proposition~\ref{prop:v47-trace},
Theorem~\ref{thm:v48-trace-identification}.  (2)
Theorem~\ref{thm:v57-singular-points-U} and Lemma~\ref{lem:v46-bounded}(i).
(3) Proposition~\ref{prop:v49-trace-smooth-decay} and
Theorem~\ref{thm:v48-trace-nontrivial}.  (4)
Theorem~\ref{thm:v61-dichotomy} and Corollary~\ref{cor:v62-trace-not-L3}.
(5) Theorem~\ref{thm:v57-singular-points-U} with
Theorem~\ref{thm:v42-trichotomy}.  The last sentence collects the
firewalls of V48--V52.
\end{proof}

\begin{firewall}[Exterior decay is qualitative only]
\label{fw:v62-decay}
Item (AC1) of Programme~\ref{prog:v62} asked for a rate in the decay of
\(U_0\) at infinity.  The decay of Proposition~\ref{prop:v49-trace-smooth-decay}
comes from dominated convergence of the local \(L^6\) and \(L^3\) norms of
\(U(s)\) and \(P(s)\) at infinity, which gives no rate; the Morrey bound
gives no pointwise information.  Consequently critical energy at infinity
around escaping centres, the case left open in Remark~\ref{rem:v61-centred},
is not excluded.  The item is recorded as open.
\end{firewall}

% =============================================================================
\section{Integrated V62 conclusion and next acquisition order}
\label{sec:v62-conclusion}
% =============================================================================

\begin{theorem}[What V62 proves]
\label{thm:v62-conclusion}
Relative to V1--V61, modulo (E1)--(E7): the \(L^3\)-at-infinity lemma
(Lemma~\ref{lem:v62-L3-infinity}), its consequence for the trace
(Corollary~\ref{cor:v62-trace-not-L3}), and the shape theorem
(Theorem~\ref{thm:v62-shape}).  The full Navier--Stokes
stability/global-regularity theorem remains unproved.
\end{theorem}

\begin{proof}
The cited results.
\end{proof}

\begin{programme}[V63 acquisition order]
\label{prog:v63}
\begin{enumerate}[label=\textup{(AC\arabic*)},leftmargin=4.8em]
\item Vorticity of the trace.  On the exterior region the trace is smooth
      with nonvanishing curl on every annulus; determine whether the
      vorticity equation at \(s=0\), which holds classically there,
      constrains \(\omega_0=\operatorname{curl}U_0\) through the value of
      \(\partial_s\omega\) at \(s=0\), for instance whether
      \(\omega_0\) can be harmonic-like or must have nonzero Laplacian on
      every exterior region.
\item The remaining items of Programmes~\ref{prog:v61}--\ref{prog:v62}
      are open.
\item The next compulsory companion audit is V65.
\end{enumerate}
\end{programme}

\begin{assessment}[Position after V62]
\label{ass:v62-position}
The trace of the type-I ancient limit is described in one theorem; its
only undetermined features are the number of its singular points and its
decay rate at infinity.  No full stability or global-regularity claim is
made.
\end{assessment}

% =============================================================================
\section*{Version V62 append-only record}
\addcontentsline{toc}{section}{Version V62 append-only record}
% =============================================================================

The complete V61 mathematical source was retained before this continuation.
Its SHA--256 digest was
\begin{center}\scriptsize\ttfamily
765b81a6f693ca6af024a468c338ee0b177bc87eeca9176024e7d4e4b893885c.
\end{center}
V62 proved that cubically integrable fields have no critical energy at
infinity, concluded that a nontrivial centred large-scale limit forces the
trace out of \(L^3\) on exterior regions, and consolidated the shape of
the trace into one theorem.  No full regularity claim is made.

% =============================================================================
\section{V63: vorticity at infinity at every time}
\label{sec:v63-entry}
% =============================================================================

Item (AC1) of Programme~\ref{prog:v63} is answered in a stronger form
than asked.  The backward-uniqueness argument of V46, which was run from
the final time, runs from any time \(s_1<0\), where the ancient limit is
bounded everywhere; combined with forward uniqueness in the \(L^6\)
class, imported and marked, it shows that the vorticity of the ancient
limit is nonzero on every exterior region at every time.  The ancient
limit of a type-I singularity therefore never has compactly supported
vorticity, never is irrotational outside a ball, and this holds at all
times, not only at the trace.  No global regularity claim is made.  The
next compulsory companion audit is V65.

% =============================================================================
\section{An additional import}
\label{sec:v63-import}
% =============================================================================

\begin{firewall}[Forward uniqueness in \(L^6\)]
\label{fw:v63-import}
\begin{enumerate}[label=\textup{(E\arabic*)},leftmargin=3.6em,start=8]
\item Mild solutions of the Navier--Stokes equations on \(\R^3\) in the
      class \(C([a,b];L^6(\R^3))\) are unique; in particular a mild
      solution in that class with zero datum at time \(a\) vanishes on
      \([a,b]\).  Smooth solutions with \(U\in C_s L^6\), \(\nabla U\in
      L^\infty_sL^2\), and pressure given by the Riesz transforms are mild.
\end{enumerate}
\end{firewall}

% =============================================================================
\section{The theorem}
\label{sec:v63-theorem}
% =============================================================================

\begin{theorem}[The ancient limit has vorticity at infinity at every time]
\label{thm:v63-vorticity-every-time}
Let \(U\) be the ancient limit of a type-I singularity, or any element of
the class obtained from it by the constructions of V41--V62.  Modulo
(E1)--(E8), for every \(s_1<0\) and every \(R>0\),
\begin{equation}
 \boxed{\quad
 \operatorname{curl}U(\cdot,s_1)\not\equiv0\ \text{ on }\{|y|\ge R\}.
 \quad}
 \label{eq:v63-vorticity-every-time}
\end{equation}
\end{theorem}

\begin{proof}
Fix \(s_1<0\) and \(R\), and suppose \(\operatorname{curl}U(\cdot,s_1)=0\)
on \(\{|y|\ge R\}\).  By Lemma~\ref{lem:v46-bounded}(ii), \(U\), \(\nabla U\),
\(\nabla^2U\) are bounded on \(\R^3\times(-\infty,s_1]\), so the vorticity
\(\omega\) satisfies the inequality
\(|\partial_s\omega-\nu\Delta\omega|\le M(|\omega|+|\nabla\omega|)\) with
bounded coefficients on \(\R^3\times(s_1-1,s_1)\), is bounded there with
bounded gradient, and vanishes at \(s_1\) on the exterior region.  The
first part of (E6) gives \(\omega\equiv0\) on \(\{|y|\ge R\}\times(s_1-1,s_1)\);
the second part gives \(\omega\equiv0\) on \(\R^3\times(s_1-1,s_1)\), hence
\(U\equiv0\) there since harmonic fields in \(L^6\) vanish.  Iterating
backward over \((s_1-2,s_1-1)\), and so on, gives \(U\equiv0\) on
\((-\infty,s_1)\).  Forward from \(s_1\), \(U\) is a smooth solution in
\(C([s_1,s_1']; L^6)\) for every \(s_1'<0\), with datum
\(U(s_1)=0\); by (E8), \(U\equiv0\) on \([s_1,0)\).  Thus \(U\equiv0\) on
\((-\infty,0)\), contradicting \eqref{eq:v41-limit-nontrivial}.
\end{proof}

\begin{corollary}[Consequences]
\label{cor:v63-consequences}
For every \(s<0\): the vorticity \(\omega(s)\) of the ancient limit is not
compactly supported; \(U(s)\) is not a potential flow outside any ball;
and \(U(s)\) is not identically zero on any exterior region.  In the
eternal case (A) the same holds for every \(s\in\R\) after the extension
of Theorem~\ref{thm:v42-trichotomy}, by applying the argument on
\((-\infty,s_1)\) for any \(s_1\) at which \(U\) is bounded.
\end{corollary}

\begin{proof}
The theorem, and the fact that \(\operatorname{curl}U(s)=0\) on an
exterior region is implied by each of the listed conditions.  In case (A)
the boundedness of \(U\) on \((-\infty,s_1]\) for every real \(s_1\)
follows from bounded enstrophy on every half-line and (E5) as in
Proposition~\ref{prop:v53-bounded-ancient}.
\end{proof}

\begin{assessment}[What this adds to the Liouville gate]
\label{ass:v63-gate}
The gate of V41 is a Liouville statement for a class of ancient
solutions.  Theorem~\ref{thm:v63-vorticity-every-time} shows that every
element of the class produced by a type-I singularity has vorticity
reaching spatial infinity at all times; equivalently, any Liouville
theorem that assumes compactly supported or exterior-vanishing vorticity
at some time already applies, and the gate is a statement about
solutions with vorticity everywhere at infinity.  This narrows the gate
without closing it.
\end{assessment}

% =============================================================================
\section{Integrated V63 conclusion and next acquisition order}
\label{sec:v63-conclusion}
% =============================================================================

\begin{theorem}[What V63 proves]
\label{thm:v63-conclusion}
Relative to V1--V62, modulo (E1)--(E8): the vorticity of the ancient limit
is nonzero on every exterior region at every time
(Theorem~\ref{thm:v63-vorticity-every-time}) with the consequences of
Corollary~\ref{cor:v63-consequences}.  The full Navier--Stokes
stability/global-regularity theorem remains unproved.
\end{theorem}

\begin{proof}
The cited results.
\end{proof}

\begin{programme}[V64 acquisition order]
\label{prog:v64}
\begin{enumerate}[label=\textup{(AC\arabic*)},leftmargin=4.8em]
\item Quantify the vorticity at infinity: determine whether the exterior
      vorticity of \(U(s)\) has a lower bound in a weighted \(L^2\) norm
      uniform in \(s\le-1\), through the bounded enstrophy and the
      backward-uniqueness structure, or whether it can tend to zero in
      every norm as \(|y|\to\infty\) while remaining nonzero.
\item Determine whether the same argument, applied to the original
      solution on the torus with the exterior region replaced by the
      complement of a neighbourhood of the singular set, gives that the
      vorticity of \(u(t)\) is nonzero on the regular region at every
      time \(t<T_*\), which is trivial for smooth solutions, and whether a
      quantitative version exists.
\item Prepare the compulsory V65 companion audit.
\end{enumerate}
\end{programme}

\begin{assessment}[Position after V63]
\label{ass:v63-position}
The ancient limit of a type-I singularity is vortical at infinity at all
times.  No full stability or global-regularity claim is made.
\end{assessment}

% =============================================================================
\section*{Version V63 append-only record}
\addcontentsline{toc}{section}{Version V63 append-only record}
% =============================================================================

The complete V62 mathematical source was retained before this continuation.
Its SHA--256 digest was
\begin{center}\scriptsize\ttfamily
ded78f9f9a972dcb49c4a5a8f8eb65ee256dae43839fdae6480d6e8abd760a98.
\end{center}
V63 proved that the vorticity of the ancient limit is nonzero on every
exterior region at every time, using the backward-uniqueness package and
an imported forward uniqueness in \(L^6\).  No full regularity claim is
made.

% =============================================================================
\section{V64: limits of the vorticity-at-infinity statement}
\label{sec:v64-entry}
% =============================================================================

Both items of Programme~\ref{prog:v64} are settled negatively with
precise reasons, and one refinement of V63 is added.  The exterior
nontriviality of the vorticity cannot be localized to bounded annuli,
because backward uniqueness for the vorticity inequality holds on
exterior domains and on the whole space but fails on bounded domains
without boundary data; and it cannot be quantified in any norm by the
present tools, because the only ingredient is a uniqueness statement,
which is qualitative.  The refinement is that the exterior vorticity is
nonzero on every exterior region at every time \emph{and} the exterior
region may be taken with respect to any centre, not only the origin.  No
global regularity claim is made.  The compulsory companion audit is due in
V65.

% =============================================================================
\section{Negative records}
\label{sec:v64-negative}
% =============================================================================

\begin{firewall}[No localization to annuli]
\label{fw:v64-annuli}
The proof of Theorem~\ref{thm:v63-vorticity-every-time} needs the
vorticity to vanish at one time on a set on which a backward uniqueness
theorem holds without boundary conditions; the import (E6) provides this
for exterior domains and for \(\R^3\).  On a bounded annulus a solution of
a parabolic inequality can vanish at the final time without vanishing
earlier, since information enters through the boundary; no conclusion
follows from \(\operatorname{curl}U(s_1)=0\) on an annulus.  Item (AC1)
in its localized form is closed negatively.
\end{firewall}

\begin{firewall}[No quantitative version]
\label{fw:v64-quantitative}
A lower bound on \(\|\operatorname{curl}U(s)\|\) in any weighted norm on
an exterior region would need a quantitative backward-uniqueness
estimate, that is, a bound on the solution on the exterior region at
earlier times in terms of its trace, which (E6) does not provide and
which is not available in this series.  Item (AC1) in its quantitative
form is closed negatively.
\end{firewall}

\begin{firewall}[The torus version is trivial]
\label{fw:v64-torus}
For the original smooth solution on the torus, vorticity that vanishes on
the complement of a neighbourhood of the singular set at some time
\(t<T_*\) would, by the same argument on the torus minus that
neighbourhood, force vanishing on the whole torus for earlier times and
hence \(u\equiv0\); but a smooth nonzero solution never has vanishing
vorticity on an open set at every earlier time except by this route, and
the statement carries no information beyond nontriviality.  Item (AC2)
is closed as trivial.
\end{firewall}

% =============================================================================
\section{A refinement}
\label{sec:v64-refinement}
% =============================================================================

\begin{proposition}[Exterior regions about any centre]
\label{prop:v64-any-centre}
Under the hypotheses of Theorem~\ref{thm:v63-vorticity-every-time}, for
every \(s_1<0\), every \(y_0\in\R^3\), and every \(R>0\),
\(\operatorname{curl}U(\cdot,s_1)\not\equiv0\) on \(\{|y-y_0|\ge R\}\).
\end{proposition}

\begin{proof}
The proof of Theorem~\ref{thm:v63-vorticity-every-time} used only that
the set on which the vorticity vanishes at \(s_1\) is an exterior domain
and that \(U\) is bounded with bounded derivatives on
\(\R^3\times(-\infty,s_1]\); both hold for exterior domains about any
centre.
\end{proof}

% =============================================================================
\section{Integrated V64 conclusion and next acquisition order}
\label{sec:v64-conclusion}
% =============================================================================

\begin{theorem}[What V64 establishes]
\label{thm:v64-conclusion}
Relative to V1--V63: the three negative records and
Proposition~\ref{prop:v64-any-centre}.  The full Navier--Stokes
stability/global-regularity theorem remains unproved.
\end{theorem}

\begin{proof}
The cited items.
\end{proof}

\begin{programme}[V65 acquisition order, with compulsory companion audit]
\label{prog:v65}
\begin{enumerate}[label=\textup{(AC\arabic*)},leftmargin=4.8em]
\item Perform the compulsory review of the current SM and YM notes;
      record digests; import nothing numerical.
\item Consolidate V56--V64 into the description of the ancient limit
      begun in Theorem~\ref{thm:v50-ancient-limit}, updating it with the
      time-uniform Morrey bound, the trace shape, and the vorticity at
      infinity at every time.
\item Record the complete list of imports (E1)--(E8) in one place with
      the versions in which each is used.
\end{enumerate}
\end{programme}

\begin{assessment}[Position after V64]
\label{ass:v64-position}
The vorticity-at-infinity statement is qualitative and cannot be
localized or quantified with the present tools.  No full stability or
global-regularity claim is made.
\end{assessment}

% =============================================================================
\section*{Version V64 append-only record}
\addcontentsline{toc}{section}{Version V64 append-only record}
% =============================================================================

The complete V63 mathematical source was retained before this continuation.
Its SHA--256 digest was
\begin{center}\scriptsize\ttfamily
4c2fa01c1910de4b0404c072a0cf19d06083b9b74f3badeb093b3c3240a55715.
\end{center}
V64 recorded that the vorticity-at-infinity statement cannot be localized
to annuli or quantified, that its torus version is trivial, and extended
it to exterior regions about any centre.  No full regularity claim is
made.

% =============================================================================
\section{V65: compulsory companion audit, the updated description of the
ancient limit, and the import ledger}
\label{sec:v65-entry}
% =============================================================================

This is a thirty-fifth-version continuation.  The compulsory review of
the two companion programmes is performed first.  The description of the
ancient limit of a type-I singularity given in V50 is then updated with
V56--V64, and the external theorems used anywhere in V41--V64 are listed
in one place with the versions in which each is used.  No global
regularity claim is made.  The next compulsory companion audit is V70.

% =============================================================================
\section{Compulsory V65 review of the current companion notes}
\label{sec:v65-companion-audit}
% =============================================================================

The current companion files in \texttt{latex\_notes} were inspected
read-only.  The frozen checkpoints are
\begin{align}
 &\texttt{Renewal\_SM\_Einstein\_Continuum\_Research\_Notes\_V72.tex},
 \notag\\[-2pt]
 &\qquad
 \texttt{e4192c52b10810e7d6ef555da4d7f74b64311acbd084cff27fcac99f2bbeb1a4},
 \label{eq:v65-SM-checkpoint}\\
 &\texttt{Renewal\_Yang\_Mills\_Mass\_Gap\_Research\_Notes\_V51.tex},
 \notag\\[-2pt]
 &\qquad
 \texttt{daa6256c65efab59f718c1548018e0e8c7893598e3801b529da7f52a60be1214}.
 \label{eq:v65-YM-checkpoint}
\end{align}

\emph{SM V72.}  On a flat compact slice the nonlinear constant-mean-curvature
Hamiltonian constraint is solved as a Lichnerowicz equation with an exact
integrated identity; a time-symmetric nonnegative source is excluded, a
homogeneous Higgs amplitude is admitted only at a mean curvature growing
with the amplitude squared, and with constant transverse--traceless
momentum solutions exist only below a density ceiling at which the volume
diverges.  The V71 negative ray is thereby replaced by a sector-existence
boundary.

\emph{YM V51.}  A tilted Hartman--Watson time law is constructed with an
exact Bessel transform, exact Wilson multiplier and heat subordination are
proved, together with tensor-block and Bernstein functional calculi and an
amplified mixed Duhamel formula; the independent block-time structure of
Wilson cumulants is recorded.

\begin{theorem}[V65 companion-audit decision]
\label{thm:v65-companion-decision}
No SM V72 constraint or ceiling constant and no YM V51 subordination or
Duhamel formula is imported.  Neither bears on the type-I Liouville gate
or on the type-II ledgers.
\end{theorem}

\begin{proof}
The hypotheses of every cited result are constraint equations on
Riemannian slices or lattice-gauge heat laws; none is a Navier--Stokes
balance.
\end{proof}

% =============================================================================
\section{The ancient limit after V64}
\label{sec:v65-ancient-limit}
% =============================================================================

\begin{theorem}[The ancient limit of a type-I singularity, V41--V64]
\label{thm:v65-ancient-limit}
Under the hypotheses of Theorem~\ref{thm:v50-ancient-limit}, and modulo
(E1)--(E8), items (1)--(5) of that theorem hold, and in addition:
\begin{enumerate}[label=\textup{(\arabic*)},leftmargin=3.2em,start=6]
\item \emph{Uniform Morrey class.}  \(U\in L^\infty((-\infty,0);M^{2,1}(\R^3))\)
      with \(\int_{B_R(y)}|U(s)|^2\le C_MR\) for all \(s<0\), and the
      same for every rescaling limit of \(U\).
\item \emph{Singular points of the limit.}  The singular points of \(U\)
      at \(s=0\) are exactly the \(L^3\)-singular points of the trace
      \(U_0\), and they lie in \(\overline{B_{R_0}}\); if
      \(U_0\in L^3_{\rm loc}\), the limit is eternal or in the dispersion
      branch (S2).
\item \emph{Shape of the trace.}  \(U_0\) has the properties of
      Theorem~\ref{thm:v62-shape}; in particular it carries critical
      energy at infinity, and is not in \(L^3\) on any exterior region,
      whenever some centred large-scale rescaling of \(U\) has a
      nontrivial limit.
\item \emph{Vorticity at infinity at every time.}  For every \(s<0\), every
      centre, and every radius, \(\operatorname{curl}U(s)\) does not
      vanish identically on the exterior region; the same holds at every
      real time in the eternal case.
\item \emph{Sufficient closures.}  Any of the statements (S1)--(S4) of
      Theorem~\ref{thm:v57-sufficient} excludes the regime; (S1) is the
      bounded-ancient Liouville conjecture restricted to finite Dirichlet
      integral, \(L^6\) values, and vorticity at infinity at all times.
\end{enumerate}
\end{theorem}

\begin{proof}
(6) Theorem~\ref{thm:v56-uniform-morrey} and Corollary~\ref{cor:v56-class}.
(7) Theorem~\ref{thm:v57-singular-points-U}.  (8)
Theorem~\ref{thm:v62-shape} and Corollary~\ref{cor:v62-trace-not-L3}.
(9) Theorem~\ref{thm:v63-vorticity-every-time},
Corollary~\ref{cor:v63-consequences}, and
Proposition~\ref{prop:v64-any-centre}.  (10)
Theorem~\ref{thm:v57-sufficient} with (9).
\end{proof}

% =============================================================================
\section{The import ledger}
\label{sec:v65-imports}
% =============================================================================

\begin{theorem}[External theorems used in V41--V64]
\label{thm:v65-import-ledger}
The following external theorems are used, each only in the versions
listed, and none is re-proved.
\begin{center}
\small
\begin{tabular}{p{0.07\textwidth}p{0.50\textwidth}p{0.35\textwidth}}
\toprule
& \textbf{Statement} & \textbf{Used in}\\
\midrule
(E1) & CKN dissipation criterion & V41, V43, V44, V49, V51--V53, V57\\
(E2) & Aubin--Lions compactness & V41, V49, V52, V53, V58, V61\\
(E3) & Fujita--Kato small data and uniqueness & V41, V42\\
(E4) & Serrin regularity for \(L^\infty_{\rm loc}L^6\) solutions & V41, V44\\
(E5) & CKN cubic criterion with \(L^\infty\) and derivative bounds & V46--V50, V53, V56\\
(E6) & ESS backward uniqueness and unique continuation & V46, V48, V52, V53, V57, V58, V61, V63, V64\\
(E7) & Quantitative form of (E5) & V49, V50, V53\\
(E8) & Forward uniqueness of mild solutions in \(C_tL^6\) & V63, V64\\
\bottomrule
\end{tabular}
\end{center}
V32--V40, V54, and V55 use no external theorem beyond the classical
periodic Sobolev, Ladyzhenskaya, Agmon, and Riesz-transform inequalities
and Leray's existence theorem (V45).
\end{theorem}

\begin{proof}
Inspection of the firewalls \ref{fw:v41-imports}, \ref{fw:v46-imports},
\ref{fw:v49-import}, \ref{fw:v63-import}, and of the proofs citing them.
\end{proof}

% =============================================================================
\section{Integrated V65 conclusion and next acquisition order}
\label{sec:v65-conclusion}
% =============================================================================

\begin{theorem}[What V65 establishes]
\label{thm:v65-conclusion}
Relative to V1--V64: the companion-audit decision
(Theorem~\ref{thm:v65-companion-decision}), the updated description of the
ancient limit (Theorem~\ref{thm:v65-ancient-limit}), and the import ledger
(Theorem~\ref{thm:v65-import-ledger}).  The full Navier--Stokes
stability/global-regularity theorem remains unproved.
\end{theorem}

\begin{proof}
The cited results.
\end{proof}

\begin{programme}[V66 acquisition order]
\label{prog:v66}
\begin{enumerate}[label=\textup{(AC\arabic*)},leftmargin=4.8em]
\item Vorticity at infinity as a Liouville input.  Search, as external
      facts to be recorded and not used, for Liouville theorems on
      bounded ancient solutions whose hypotheses exclude vorticity at
      infinity, and determine whether their proofs use that hypothesis
      essentially; if not, whether they extend.
\item The (S2) branch with the updated description: a bounded solution on
      a slab, smooth up to the final time, with vorticity at infinity at
      every time and enstrophy escaping to infinity; determine whether the
      time-uniform Morrey bound and the vorticity equation on exterior
      regions, where all coefficients are bounded, give a maximum
      principle for a weighted vorticity quantity.
\item The next compulsory companion audit is V70.
\end{enumerate}
\end{programme}

\begin{assessment}[Position after V65]
\label{ass:v65-position}
Thirty-five versions after V30, the series has a complete, audited,
import-tabulated description of the type-I regime as a Liouville problem
and of the type-II regime as a set of consistent necessary conditions.
No full stability or global-regularity claim is made.
\end{assessment}

% =============================================================================
\section*{Version V65 append-only record}
\addcontentsline{toc}{section}{Version V65 append-only record}
% =============================================================================

The complete V64 mathematical source was retained before this continuation.
Its SHA--256 digest was
\begin{center}\scriptsize\ttfamily
d8534196494f98d832f71d42bda48a3e044dd2967d30ce437c2c04de873adf69.
\end{center}
V65 performed the compulsory companion audit, updated the description of
the ancient limit with V56--V64, and tabulated the external theorems used
in V41--V64.  No full regularity claim is made.

% =============================================================================
\section{V66: a rate-free local energy bound and the \(\tau^{1/3}\)
threshold for spatial energy concentration}
\label{sec:v66-entry}
% =============================================================================

Item (AC2) of Programme~\ref{prog:v66} is not advanced; item (AC1) is
recorded as not decidable inside the series.  The version instead returns
to the type-II regime with a rate-free version of the V47 Morrey
argument.  The local energy identity, with the initial time chosen at a
low-enstrophy instant of the preceding parabolic interval and with the
cubic and pressure terms bounded through energy and enstrophy only,
bounds the energy in every ball of radius \(\rho\) by the Leray share of
the preceding \(2\rho^2\) time plus \(\rho^{-1/2}\) times the \(3/4\)
power of that share.  Consequently spatial energy concentration at the
singular time is impossible whenever the Leray shares of parabolic slabs
decay faster than \(\tau^{1/3}\), which includes the type-I regime, where
they decay like \(\tau^{1/2}\), and a sub-regime of type II.  Conversely,
energy concentration at a point forces slab shares of order at least
\(\tau^{1/3}\).  No global regularity claim is made.  The next compulsory
companion audit is V70.

% =============================================================================
\section{The rate-free local energy bound}
\label{sec:v66-bound}
% =============================================================================

For \(t\in(0,T_*)\) and \(0<\tau\le t\) write
\begin{equation}
 \epsilon(t;\tau):=\int_{t-\tau}^t\|\Lambda u\|_2^2\,\dd\sigma,
 \qquad
 \bar\epsilon(\tau):=\sup_{\tau\le t<T_*}\epsilon(t;\tau)\le\frac{E_*}{2\nu}.
 \label{eq:v66-shares}
\end{equation}

\begin{theorem}[Local energy from recent Leray shares]
\label{thm:v66-local-energy}
There is a constant \(C_0\), depending only on \(\nu\) and the periodic
Sobolev and Riesz constants, such that for every \(x\in\Torus^3\), every
\(\rho\in(0,\pi]\), and every \(t\in[2\rho^2,T_*)\),
\begin{equation}
 \boxed{\quad
 \int_{B_\rho(x)}|u(t)|^2\,\dd x'
 \le C_0\Bigl(\epsilon(t;2\rho^2)
 +\rho^{-1/2}E_*^{3/4}\,\epsilon(t;2\rho^2)^{3/4}\Bigr).
 \quad}
 \label{eq:v66-local-energy}
\end{equation}
\end{theorem}

\begin{proof}
Let \(Y=\|\Lambda u\|_2^2\).  The mean of \(Y\) over
\([t-2\rho^2,t-\rho^2]\) is at most \(\rho^{-2}\epsilon(t;2\rho^2)\), so
there is \(t_0\) in that interval with \(Y(t_0)\le\rho^{-2}\epsilon(t;2\rho^2)\).
The local energy identity on \([t_0,t]\) with the cutoff \(\phi_\rho\)
centred at \(x\) gives, after dropping the dissipation term,
\[
 \int_{B_\rho}|u(t)|^2
 \le\int_{B_{2\rho}}|u(t_0)|^2
 +\nu\int_{t_0}^t\!\!\int|u|^2|\Delta\phi_\rho|
 +\int_{t_0}^t\!\!\int\bigl(|u|^3+2|u||p-\bar p|\bigr)|\nabla\phi_\rho| .
\]
First term: \(|B_{2\rho}|^{2/3}\|u(t_0)\|_6^2\le C\rho^2C_6^2Y(t_0)
\le C\epsilon(t;2\rho^2)\).  Second term: \(|\Delta\phi_\rho|\le C\rho^{-2}
\mathbf 1_{B_{2\rho}}\) and \(\int_{B_{2\rho}}|u|^2\le C\rho^2\|u\|_6^2\le
C\rho^2Y\), so the term is at most \(C\nu\int_{t_0}^tY\le C\nu\,\epsilon(t;2\rho^2)\).
Third term: \(|\nabla\phi_\rho|\le C\rho^{-1}\mathbf 1_{B_{2\rho}}\);
\(\int|u|^3\le\|u\|_2^{3/2}\|u\|_6^{3/2}\le E_*^{3/4}C_6^{3/2}Y^{3/4}\)
by interpolation, and
\(\int|u||p-\bar p|\le\|u\|_3\cdot2\|p\|_{3/2}\le\|u\|_3\cdot2C_{\rm CZ}
\|u\|_3^2\le CE_*^{3/4}Y^{3/4}\) by \(\|u\|_3^2\le\|u\|_2\|u\|_6\); hence
the term is at most \(C\rho^{-1}E_*^{3/4}\int_{t_0}^tY^{3/4}
\le C\rho^{-1}E_*^{3/4}\bigl(\int_{t_0}^tY\bigr)^{3/4}(t-t_0)^{1/4}
\le C\rho^{-1}E_*^{3/4}\epsilon(t;2\rho^2)^{3/4}(2\rho^2)^{1/4}\),
which is the second term of \eqref{eq:v66-local-energy}.
\end{proof}

\begin{corollary}[The \(\tau^{1/3}\) threshold]
\label{cor:v66-threshold}
\begin{enumerate}[label=\textup{(\roman*)},leftmargin=3.7em]
\item If \(\bar\epsilon(\tau)\le A\tau^{1/3}\eta(\tau)\) with
      \(\eta(\tau)\to0\) as \(\tau\to0\), then
      \begin{equation}
       \sup_{x\in\Torus^3,\ 2\rho^2\le t<T_*}\int_{B_\rho(x)}|u(t)|^2
       \ \le\ C_0\Bigl(A\,2^{1/3}\rho^{2/3}\eta(2\rho^2)
       +2^{1/4}A^{3/4}E_*^{3/4}\eta(2\rho^2)^{3/4}\Bigr)
       \xrightarrow[\rho\to0]{}0 :
       \label{eq:v66-no-concentration}
      \end{equation}
      the energy does not concentrate in space at the singular time,
      uniformly in time and position.
\item Conversely, if there are \(x_0\), \(c>0\), and \(\rho_j\to0\),
      \(t_j\to T_*\) with \(\int_{B_{\rho_j}(x_0)}|u(t_j)|^2\ge c\), then
      \(\epsilon(t_j;2\rho_j^2)\ge c'\rho_j^{2/3}\) with
      \(c'=c'(c,C_0,E_*)\): energy concentration at a point forces Leray
      shares of order at least \(\tau^{1/3}\) on the preceding parabolic
      slabs.
\item In the type-I regime \(\bar\epsilon(\tau)\le2C_I^2\nu^{3/2}\tau^{1/2}\),
      so (i) applies with \(\eta(\tau)=\tau^{1/6}\), recovering the
      absence of spatial concentration proved in V44 by a different
      route, and giving the explicit bound
      \(\int_{B_\rho(x)}|u(t)|^2\le C\rho^{5/6}\) uniformly; the Morrey
      bound \(C_M\rho\) of V47 is sharper and needs the rate.
\end{enumerate}
\end{corollary}

\begin{proof}
(i) Insert the hypothesis in \eqref{eq:v66-local-energy}:
\(\rho^{-1/2}(A(2\rho^2)^{1/3}\eta)^{3/4}=2^{1/4}A^{3/4}\eta^{3/4}\).
(ii) If \(\epsilon(t_j;2\rho_j^2)\le c'\rho_j^{2/3}\) with \(c'\) small,
\eqref{eq:v66-local-energy} gives
\(\int_{B_{\rho_j}}|u(t_j)|^2\le C_0(c'\rho_j^{2/3}+E_*^{3/4}c'^{3/4})<c\)
for \(c'\) small and \(\rho_j\) small.  (iii) Lemma~\ref{lem:v43-slab}
and the substitution.
\end{proof}

\begin{assessment}[A sub-regime of type II without spatial concentration]
\label{ass:v66-subregime}
The geometric Leray-share dichotomy of V54 separated type-I along a
sequence (shares \(\asymp\tau^{1/2}\) along a subsequence) from strong
type II (shares \(\gg\tau^{1/2}\) always).  Corollary~\ref{cor:v66-threshold}
inserts a second threshold: as long as the shares stay below
\(\tau^{1/3}\) by a vanishing factor, the singularity cannot concentrate
energy at a point.  Between \(\tau^{1/2}\) and \(\tau^{1/3}\) lies a
sub-regime of type II which shares with type I the absence of spatial
energy concentration, though not the compactness theory.  Above
\(\tau^{1/3}\) energy concentration is not excluded by this argument;
whether it can actually occur for a Navier--Stokes solution is not
decided here.  The bound is rate-free and holds for every smooth solution
up to its maximal time.
\end{assessment}

\begin{firewall}[Items of Programme~\ref{prog:v66}]
\label{fw:v66-items}
Item (AC1) asks for external Liouville theorems and cannot be settled
inside the series; item (AC2) was attempted and yields only the trivial
maximum-principle bound \(\sup|\omega(s)|\le e^{2M(s-s_0)}\sup|\omega(s_0)|\)
on exterior regions, which adds nothing to boundedness.  Both are
recorded as open.
\end{firewall}

% =============================================================================
\section{Integrated V66 conclusion and next acquisition order}
\label{sec:v66-conclusion}
% =============================================================================

\begin{theorem}[What V66 proves]
\label{thm:v66-conclusion}
Relative to V1--V65, with no external import: the rate-free local energy
bound \eqref{eq:v66-local-energy} and the \(\tau^{1/3}\) threshold
(Corollary~\ref{cor:v66-threshold}).  The full Navier--Stokes
stability/global-regularity theorem remains unproved.
\end{theorem}

\begin{proof}
The cited results.
\end{proof}

\begin{programme}[V67 acquisition order]
\label{prog:v67}
\begin{enumerate}[label=\textup{(AC\arabic*)},leftmargin=4.8em]
\item Sharpen the exponent.  The \(3/4\) power in
      \eqref{eq:v66-local-energy} came from bounding \(\int|u|^3\) by
      \(E_*^{3/4}Y^{3/4}\); test whether using the local energy itself in
      the cubic term (\(\int_{B}|u|^3\le(\int_B|u|^2)^{1/2}\|u\|_6\cdot
      |B|^{\ldots}\)) closes a Gr\"onwall-type inequality for the local
      energy that improves the threshold \(\tau^{1/3}\) toward
      \(\tau^{1/2}\).
\item Combine with the type-II price sheet: high-total windows have
      shares at least \(c\nu/\rho_j\asymp\tau_j^{1/2}\); determine whether
      the record--share inequality forces shares above \(\tau^{1/3}\) on
      some windows, which would make spatial concentration possible only
      in the strongly dissipative part of type II.
\item The next compulsory companion audit is V70.
\end{enumerate}
\end{programme}

\begin{assessment}[Position after V66]
\label{ass:v66-position}
A rate-free local energy bound gives every smooth solution a threshold
below which its singularity, if any, cannot concentrate energy in space.
No full stability or global-regularity claim is made.
\end{assessment}

% =============================================================================
\section*{Version V66 append-only record}
\addcontentsline{toc}{section}{Version V66 append-only record}
% =============================================================================

The complete V65 mathematical source was retained before this continuation.
Its SHA--256 digest was
\begin{center}\scriptsize\ttfamily
a8c468ffa977d807d8ef08d4c81a28247293942163bb9f113e7276ad2f7903e0.
\end{center}
V66 proved a rate-free bound on the local energy in terms of recent Leray
shares and derived the \(\tau^{1/3}\) threshold for spatial energy
concentration at the singular time.  No full regularity claim is made.

% =============================================================================
\section{V67: the concentration threshold against the price sheet}
\label{sec:v67-entry}
% =============================================================================

Both items of Programme~\ref{prog:v67} are examined.  Using the local
energy itself in the cubic term of the V66 bound gives a self-improving
recursion for the concentration function, but the pressure term is
global and keeps the threshold at \(\tau^{1/3}\); a localized pressure
decomposition would be needed to move it, and is not carried out.  The
high-total windows of the type-II price sheet have Leray shares of
parabolic order \(\tau^{1/2}\), below the concentration threshold, so the
price sheet does not by itself force or forbid spatial concentration.
Energy concentration at a point, if it occurs, forces the continuation
stock of the preceding parabolic slab to diverge like \(\rho^{-2/3}\).
No global regularity claim is made.  The next compulsory companion audit
is V70.

% =============================================================================
\section{The recursion for the concentration function}
\label{sec:v67-recursion}
% =============================================================================

Let \(M(\rho):=\sup_{x,\,t\ge2\rho^2}\int_{B_\rho(x)}|u(t)|^2\).

\begin{proposition}[Self-improving recursion]
\label{prop:v67-recursion}
For \(\rho\in(0,\pi/2]\),
\begin{equation}
 M(\rho)\le C_0\Bigl(\bar\epsilon(2\rho^2)
 +\rho^{-1/2}\bar\epsilon(2\rho^2)^{3/4}
 \bigl[E_*^{1/4}M(2\rho)^{1/2}+E_*^{7/12}M(2\rho)^{1/6}\bigr]\Bigr).
 \label{eq:v67-recursion}
\end{equation}
\end{proposition}

\begin{proof}
Repeat the proof of Theorem~\ref{thm:v66-local-energy}, bounding the
cubic term on \(B_{2\rho}\) by
\(\int_{B_{2\rho}}|u|^3\le\bigl(\int_{B_{2\rho}}|u|^2\bigr)^{1/2}\|u\|_4^2
\le M(2\rho)^{1/2}C_4E_*^{1/4}Y^{3/4}\) and the pressure term by
\(\|u\|_{L^3(B_{2\rho})}\cdot2\|p\|_{3/2}\) with
\(\|u\|_{L^3(B_{2\rho})}^3\le M(2\rho)^{1/2}C_4E_*^{1/4}Y^{3/4}\) and
\(\|p\|_{3/2}\le C_{\rm CZ}E_*^{1/2}Y^{1/2}\); then integrate in time
with H\"older as before.
\end{proof}

\begin{firewall}[The threshold does not move]
\label{fw:v67-threshold}
Since \(M\le E_*\), \eqref{eq:v67-recursion} reproduces
\eqref{eq:v66-local-energy}; when \(M(2\rho)\to0\) it improves the
constant but not the power \(\rho^{-1/2}\bar\epsilon^{3/4}\), whose
balance against \(\bar\epsilon\asymp\tau^{1/3}\) sets the threshold.
The obstruction is the pressure, bounded through the global
\(\|u\|_3^2\); a localized pressure estimate of the
Caffarelli--Kohn--Nirenberg type would replace \(E_*^{7/12}M^{1/6}\) by
a quantity of the order of \(M^{1/2}\) plus a harmonic remainder and could
move the threshold toward \(\tau^{1/2}\).  This is not carried out.
Item (AC1) is closed as not done.
\end{firewall}

% =============================================================================
\section{High-total windows and the threshold}
\label{sec:v67-price-sheet}
% =============================================================================

\begin{proposition}[Price-sheet shares are of parabolic order]
\label{prop:v67-price-sheet}
A high-total window \(w_j=[b_j-\tau_j,b_j]\), \(\tau_j=\theta^{-1}(\nu\rho_j^2)^{-1}\),
has Leray share \(\epsilon_{w_j}\ge\nu/(2(2C_1\theta)^{1/2}\rho_j)
=c\,\tau_j^{1/2}\) by Theorem~\ref{thm:v38-high-total-share}, and the
record--share inequality requires \(\epsilon_{w_j}\ge c_\theta E_*/\mathcal R_{w_j}\).
Neither forces \(\epsilon_{w_j}\gtrsim\tau_j^{1/3}\): with doubled
records \(\mathcal R_{w_j}\asymp2^j\) and \(\tau_j\asymp16^{-j}\) both are
satisfied by shares \(\asymp\tau_j^{1/2}\).  Hence a cofinal high-total
sequence is compatible with the absence of spatial energy concentration
at the singular time, and also with its presence; the price sheet does
not decide.
\end{proposition}

\begin{proof}
Direct substitution of the cited bounds.
\end{proof}

\begin{proposition}[Cost of energy concentration in the continuation
stock]
\label{prop:v67-continuation-cost}
If \(\int_{B_{\rho_j}(x_0)}|u(t_j)|^2\ge c\) with \(\rho_j\to0\),
\(t_j\to T_*\), then for the slabs \(S_j=[t_j-2\rho_j^2,t_j]\),
\begin{equation}
 \int_{S_j}\|\Lambda u\|_2^2\ge c'\rho_j^{2/3},
 \qquad
 \int_{S_j}\|\Lambda u\|_2^4\ge\frac{c'^2}{2}\rho_j^{-2/3},
 \qquad
 \sup_{S_j}\|\Lambda u\|_2^2\ge\frac{c'}{2}\rho_j^{-4/3}.
 \label{eq:v67-continuation-cost}
\end{equation}
\end{proposition}

\begin{proof}
The first is Corollary~\ref{cor:v66-threshold}(ii); the second is
Cauchy--Schwarz \(\int_SY\le(\int_SY^2)^{1/2}|S|^{1/2}\); the third is
\(\int_SY\le|S|\sup_SY\).
\end{proof}

\begin{assessment}[Reading]
\label{ass:v67-reading}
Energy concentration at a point of the singular set is a strongly
type-II event: it needs enstrophy of order \(\rho^{-4/3}\) on the
preceding parabolic slab, against the Leray rate \(\rho^{-1}\), and slab
shares \(\rho^{2/3}\) against the parabolic \(\rho\).  The type-I regime,
and the sub-regime of type II below the \(\tau^{1/3}\) threshold, are
free of it.  Whether energy concentration at a point can happen at all
for a Navier--Stokes singularity is not decided.
\end{assessment}

% =============================================================================
\section{Integrated V67 conclusion and next acquisition order}
\label{sec:v67-conclusion}
% =============================================================================

\begin{theorem}[What V67 proves]
\label{thm:v67-conclusion}
Relative to V1--V66, with no external import: the recursion
\eqref{eq:v67-recursion}, the compatibility of the price sheet with both
alternatives (Proposition~\ref{prop:v67-price-sheet}), and the cost of
energy concentration \eqref{eq:v67-continuation-cost}.  The full
Navier--Stokes stability/global-regularity theorem remains unproved.
\end{theorem}

\begin{proof}
The cited results.
\end{proof}

\begin{programme}[V68 acquisition order]
\label{prog:v68}
\begin{enumerate}[label=\textup{(AC\arabic*)},leftmargin=4.8em]
\item Localized pressure.  Carry out the decomposition of the pressure
      on \(B_{2\rho}\) into a local part bounded by the local \(L^3\)
      norm of \(u\) and a harmonic part bounded by averages over larger
      balls, and rerun the recursion; determine the resulting threshold.
\item Energy concentration versus the \(L^{5/2}\) ledger: the shares
      \(\rho^{2/3}\) give ledger at least \(c\rho^{1/3}\) per slab, which
      is summable over dyadic slabs; determine whether concentration at
      infinitely many scales is compatible with a finite intermediate
      ledger, which would place it in Regime I of
      Theorem~\ref{thm:v40-classification}.
\item The next compulsory companion audit is V70.
\end{enumerate}
\end{programme}

\begin{assessment}[Position after V67]
\label{ass:v67-position}
Spatial energy concentration at the singular time is now placed above a
\(\tau^{1/3}\) dissipation threshold and shown to be undecided by the
type-II price sheet.  No full stability or global-regularity claim is
made.
\end{assessment}

% =============================================================================
\section*{Version V67 append-only record}
\addcontentsline{toc}{section}{Version V67 append-only record}
% =============================================================================

The complete V66 mathematical source was retained before this continuation.
Its SHA--256 digest was
\begin{center}\scriptsize\ttfamily
a47bcd717e6b3b4841dd8f9a494bb479ec43d07a929ff752e212026e5cbcb27d.
\end{center}
V67 derived the self-improving recursion for the concentration function,
recorded that the pressure keeps the threshold at \(\tau^{1/3}\), showed
that the type-II price sheet does not decide spatial concentration, and
computed the continuation-stock cost of concentration.  No full
regularity claim is made.

% =============================================================================
\section{V68: exponent bookkeeping for the localized pressure}
\label{sec:v68-entry}
% =============================================================================

Item (AC1) of Programme~\ref{prog:v68} is carried out for the local part
of the pressure and shown not to move the \(\tau^{1/3}\) threshold: with
Leray shares of order \(\tau^\alpha\) and a concentration function of
order \(\rho^\beta\), the local cubic and local pressure terms close
exactly when \(\beta\le\min\{2\alpha,3\alpha-1\}\), so the concentration
function vanishes as \(\rho\to0\) if and only if \(\alpha>1/3\), as in
V66.  The harmonic part of the pressure requires the scale iteration of
the Caffarelli--Kohn--Nirenberg theory and is not carried out; the
threshold is therefore intrinsic to energy-based arguments and would move
only through a genuinely different mechanism.  Item (AC2) is settled:
energy concentration at a point costs a summable amount of the
\(L^{5/2}\) ledger per dyadic slab, so it is compatible with Regime I of
Theorem~\ref{thm:v40-classification} without the type-I rate, and
excluded within the type-I rate by V47.  No global regularity claim is
made.  The next compulsory companion audit is V70.

% =============================================================================
\section{Local part of the pressure}
\label{sec:v68-local-pressure}
% =============================================================================

Fix \(x\) and \(\rho\), let \(\chi\) be a cutoff equal to one on
\(B_{4\rho}(x)\) and supported in \(B_{8\rho}(x)\), and split the
mean-zero pressure as \(p=p_1+p_2\), \(p_1:=R_iR_j(\chi u_iu_j)\),
\(p_2:=R_iR_j((1-\chi)u_iu_j)\), with \(R_i\) the periodic Riesz
transforms.  Then \(\|p_1\|_{3/2}\le C_{\rm CZ}\|u\|_{L^3(B_{8\rho})}^2\),
and \(p_2\) is smooth on \(B_{2\rho}\), where its source vanishes.

\begin{proposition}[Local terms with Leray shares of order \(\tau^\alpha\)]
\label{prop:v68-bookkeeping}
Suppose \(\bar\epsilon(\tau)\le A\tau^\alpha\) for \(\tau\le\tau_0\) and
\(M(r)\le Br^\beta\) for \(\rho\le r\le8\rho\), with \(M\) as in
Proposition~\ref{prop:v67-recursion}.  Then the initial, Laplacian,
cubic, and local-pressure terms in the local energy identity on
\(B_\rho\) over \([t_0,t]\) as in the proof of
Theorem~\ref{thm:v66-local-energy} are bounded, respectively, by
\begin{equation}
 CA\rho^{2\alpha},\qquad CA\rho^{2\alpha},\qquad
 CA^{3/4}B^{1/2}E_*^{1/4}\rho^{-1/2+3\alpha/2+\beta/2},\qquad
 CA^{3/4}B^{1/2}E_*^{1/4}\rho^{-1/2+3\alpha/2+\beta/2},
 \label{eq:v68-terms}
\end{equation}
all of order at most \(\rho^\beta\) if and only if
\(\beta\le\min\{2\alpha,\,3\alpha-1\}\).  In particular a positive
\(\beta\) exists if and only if \(\alpha>1/3\).
\end{proposition}

\begin{proof}
The initial and Laplacian terms are as in Theorem~\ref{thm:v66-local-energy}
with \(\epsilon\le A(2\rho^2)^\alpha\).  For the cubic term,
\(\int_{B_{2\rho}}|u|^3\le M(2\rho)^{1/2}\|u\|_4^2\le(B(2\rho)^\beta)^{1/2}
C_4E_*^{1/4}Y^{3/4}\), integrated over a time interval of length
\(2\rho^2\) with H\"older: \(\int Y^{3/4}\le\epsilon^{3/4}(2\rho^2)^{1/4}\);
multiply by \(\rho^{-1}\) from \(|\nabla\phi_\rho|\).  For the local
pressure, \(\int_{B_{2\rho}}|u||p_1-\bar p_1|\le\|u\|_{L^3(B_{2\rho})}
\cdot2C_{\rm CZ}\|u\|_{L^3(B_{8\rho})}^2\le C\|u\|_{L^3(B_{8\rho})}^3
\le CM(8\rho)^{1/2}C_4E_*^{1/4}Y^{3/4}\), which is the same order as the
cubic term.  The exponent conditions are read off.
\end{proof}

\begin{firewall}[The harmonic part]
\label{fw:v68-harmonic}
On \(B_{2\rho}\), \(\nabla p_2\) is bounded by
\(C\sum_{k\ge0}(2^k\rho)^{-4}\int_{B_{2^{k+3}\rho}}|u|^2\), which with
\(M(r)\le Br^\beta\) on all larger scales gives \(|\nabla p_2|\le
CB\rho^{\beta-4}\) and a flux contribution that does not close under the
same bookkeeping: its exponent condition is \(\beta\ge3-3\alpha\), which
contradicts \(\beta\le2\alpha\) for every \(\alpha\le1\).  The standard
resolution, in the theory of suitable weak solutions, is to use the
harmonicity of \(p_2\) to compare its oscillation on \(B_\rho\) with that
on \(B_{2\rho}\) and to iterate a decay estimate across scales; this
iteration is the substance of the \(\varepsilon\)-regularity theorem
imported as (E5) and is not reproduced here.  Consequently the present
version establishes only that localizing the \emph{local} part of the
pressure does not move the threshold; whether the full CKN iteration,
combined with Leray shares of order \(\tau^\alpha\) for \(1/3<\alpha<1/2\),
gives a Morrey bound is not decided.
\end{firewall}

% =============================================================================
\section{Energy concentration and the ledgers}
\label{sec:v68-ledgers}
% =============================================================================

\begin{proposition}[Concentration is ledger-cheap]
\label{prop:v68-ledger-cheap}
If \(\int_{B_{\rho_j}(x_0)}|u(t_j)|^2\ge c\) with \(\rho_j=2^{-j}\rho_0\)
and \(t_j\to T_*\), the slabs \(S_j=[t_j-2\rho_j^2,t_j]\) carry
\begin{equation}
 \int_{S_j}\|\Lambda u\|_2^{5/2}\ \ge\ \frac{\epsilon(t_j;2\rho_j^2)^{5/4}}{(2\rho_j^2)^{1/4}}
 \ \ge\ c''\rho_j^{1/3},
 \qquad
 \sum_j c''\rho_j^{1/3}<\infty .
 \label{eq:v68-ledger-cheap}
\end{equation}
Hence energy concentration at a point along dyadic scales is compatible
with a finite \(L^{5/2}\) ledger, that is, with Regime I of
Theorem~\ref{thm:v40-classification}; within Regime I it is excluded by
Theorem~\ref{thm:v47-morrey} exactly when the type-I rate holds, and is
otherwise undecided.
\end{proposition}

\begin{proof}
H\"older \(\int_SY\le(\int_SY^{5/4})^{4/5}|S|^{1/5}\) and
Corollary~\ref{cor:v66-threshold}(ii) give the first inequality; the sum
is geometric.  The regime statements are the cited theorems.
\end{proof}

\begin{assessment}[The intermediate class again]
\label{ass:v68-intermediate}
Between the type-I rate and strong type II, and inside Regime I, there
is room for solutions with slab shares between \(\tau^{1/2}\) and
\(\tau^{1/3}\) that concentrate no energy, and for solutions with shares
above \(\tau^{1/3}\) that might.  The series has no tool that decides
what happens in this intermediate class beyond the rate-free bounds of
V66--V68; every finer statement has needed the rate.
\end{assessment}

% =============================================================================
\section{Integrated V68 conclusion and next acquisition order}
\label{sec:v68-conclusion}
% =============================================================================

\begin{theorem}[What V68 proves]
\label{thm:v68-conclusion}
Relative to V1--V67, with no external import: the exponent bookkeeping
\eqref{eq:v68-terms} and its threshold, the record that the harmonic
pressure part needs the CKN iteration, and the ledger cost of
concentration \eqref{eq:v68-ledger-cheap}.  The full Navier--Stokes
stability/global-regularity theorem remains unproved.
\end{theorem}

\begin{proof}
The cited results.
\end{proof}

\begin{programme}[V69 acquisition order]
\label{prog:v69}
\begin{enumerate}[label=\textup{(AC\arabic*)},leftmargin=4.8em]
\item Determine whether the \(\varepsilon\)-regularity import (E5), used
      as a black box on cylinders where the scale-invariant cubic
      quantity is small, gives a Morrey-type bound under slab shares of
      order \(\tau^\alpha\), \(\alpha>1/3\), through the bound
      \(r^{-2}\iint_{Q_r}|u|^3\le Cr^{-2}\int_{t-r^2}^tE_*^{3/4}Y^{3/4}
      \le CE_*^{3/4}r^{-3/2}\epsilon^{3/4}\), which is small at scale \(r\)
      when \(\epsilon(t;r^2)\ll r^2\), that is, never at the singular time;
      record the exact obstruction.
\item Prepare the compulsory V70 companion audit.
\end{enumerate}
\end{programme}

\begin{assessment}[Position after V68]
\label{ass:v68-position}
The \(\tau^{1/3}\) threshold is intrinsic to energy-based local
arguments.  No full stability or global-regularity claim is made.
\end{assessment}

% =============================================================================
\section*{Version V68 append-only record}
\addcontentsline{toc}{section}{Version V68 append-only record}
% =============================================================================

The complete V67 mathematical source was retained before this continuation.
Its SHA--256 digest was
\begin{center}\scriptsize\ttfamily
53446b92a5578a67efc8e3efc8e12623bdfa5111c9c4ae30d14d50cc05e288cb.
\end{center}
V68 carried out the exponent bookkeeping for the localized pressure,
showed the threshold is unchanged by the local part, recorded that the
harmonic part needs the CKN iteration, and showed that energy
concentration is compatible with a finite intermediate ledger.  No full
regularity claim is made.

% =============================================================================
\section{V69: why \(\varepsilon\)-regularity cannot be applied at the singular
time without a rate}
\label{sec:v69-entry}
% =============================================================================

Item (AC1) of Programme~\ref{prog:v69} is settled negatively with the
exact obstruction.  The scale-invariant cubic quantity on a parabolic
cylinder ending at the singular time is bounded above by
\(r^{-1/2}\) times the \(L^3_tH^1\) ledger of the last \(r^2\) time and,
by the Leray rate, that ledger is at least a constant times \(r^{1/2}\);
the two bounds are of the same order, so the quantity is bounded but
never small, and \(\varepsilon\)-regularity gives nothing at the singular
time without a rate hypothesis that makes the constant small.  The same
computation yields, for every singular time, a lower bound on the
\(L^3_tH^1\) ledger of every final slab, which is of the order already
implied by the Leray rate.  No global regularity claim is made.  The
compulsory companion audit is due in V70.

% =============================================================================
\section{The cubic quantity at the singular time}
\label{sec:v69-cubic}
% =============================================================================

\begin{proposition}[Two-sided order of the cubic quantity]
\label{prop:v69-two-sided}
For every \(x\in\Torus^3\) and \(0<r\le\pi\), with
\(Q_r=Q_r(x,T_*)\),
\begin{equation}
 \frac1{r^2}\iint_{Q_r}\bigl(|u|^3+|u||p-\bar p|\bigr)
 \ \le\ C_\flat\,\frac1{r^{1/2}}\int_{T_*-r^2}^{T_*}\|\Lambda u\|_2^3\,\dd t ,
 \label{eq:v69-upper}
\end{equation}
with \(C_\flat\) depending only on the periodic Sobolev and Riesz
constants, and
\begin{equation}
 \frac1{r^{1/2}}\int_{T_*-r^2}^{T_*}\|\Lambda u\|_2^3\,\dd t
 \ \ge\ 4c_L^3\nu^{9/4}
 \label{eq:v69-lower}
\end{equation}
by the Leray rate.  Hence the right side of \eqref{eq:v69-upper} is
bounded below by a positive constant at every singular time, and the
hypothesis of (E5) can be verified at \((x,T_*)\) only if
\(C_\flat r^{-1/2}\int_{T_*-r^2}^{T_*}\|\Lambda u\|_2^3\le\varepsilon_2\),
which a rate-free hypothesis cannot supply because the left side is at
least \(4C_\flat c_L^3\nu^{9/4}\).
\end{proposition}

\begin{proof}
\(\int_{B_r}|u|^3\le|B_r|^{1/2}\|u\|_6^3\le Cr^{3/2}\|\Lambda u\|_2^3\) and
\(\int_{B_r}|u||p-\bar p|\le|B_r|^{1/6}\|u\|_6\cdot2|B_r|^{1/3}\|p\|_3
\le Cr^{3/2}\|\Lambda u\|_2\cdot\|\Lambda u\|_2^2\); integrate in time
and divide by \(r^2\).  The lower bound integrates \eqref{eq:v32-leray-rate}
cubed: \(\int_{T_*-r^2}^{T_*}(T_*-t)^{-3/4}\dd t=4r^{1/2}\).
\end{proof}

\begin{assessment}[The role of the rate]
\label{ass:v69-rate}
Under the type-I rate the same quantity is bounded by
\(C_\sharp c^{3/2}\) at scale \(c(T_*-t)^{1/2}\) for \(t<T_*\), which is
small for small \(c\) (Theorem~\ref{thm:v46-Linfty-upgrade}); the
smallness comes from choosing the cylinder shorter than the parabolic
distance to the singular time, not from any smallness at \(T_*\) itself.
Proposition~\ref{prop:v69-two-sided} shows that at the singular time
the quantity is of critical order for every solution, so
\(\varepsilon\)-regularity can only ever be applied at times before
\(T_*\), on cylinders whose length is controlled by a rate.  This is
the exact reason every compactness statement of the series needed the
rate (Theorem~\ref{thm:v60-table}), and it closes item (AC1).
\end{assessment}

\begin{corollary}[An \(L^3_tH^1\) lower bound at every singular time,
from \(\varepsilon\)-regularity]
\label{cor:v69-L3-lower}
Modulo (E5): if \(T_*\) is singular then for every \(x\) in the singular
set and every small \(r\),
\begin{equation}
 \int_{T_*-r^2}^{T_*}\|\Lambda u\|_2^3\,\dd t\ \ge\ \frac{\varepsilon_2}{C_\flat}\,r^{1/2}.
 \label{eq:v69-L3-lower}
\end{equation}
This is of the same order as the consequence \eqref{eq:v69-lower} of the
Leray rate and adds only a comparison of constants.
\end{corollary}

\begin{proof}
If the left side were smaller for some small \(r\), (E5) would make
\((x,T_*)\) regular.
\end{proof}

% =============================================================================
\section{Integrated V69 conclusion and next acquisition order}
\label{sec:v69-conclusion}
% =============================================================================

\begin{theorem}[What V69 establishes]
\label{thm:v69-conclusion}
Relative to V1--V68: the two-sided order of the cubic quantity at the
singular time (Proposition~\ref{prop:v69-two-sided}), the exact reason
\(\varepsilon\)-regularity needs a rate, and the \(L^3_tH^1\) lower bound
\eqref{eq:v69-L3-lower}.  The full Navier--Stokes stability/global-regularity
theorem remains unproved.
\end{theorem}

\begin{proof}
The cited results.
\end{proof}

\begin{programme}[V70 acquisition order, with compulsory companion audit]
\label{prog:v70}
\begin{enumerate}[label=\textup{(AC\arabic*)},leftmargin=4.8em]
\item Perform the compulsory review of the current SM and YM notes;
      record digests; import nothing numerical.
\item Consolidate the rate-free results of V54, V55, V66--V69 into one
      theorem on the intermediate class between the type-I rate and
      strong type II.
\item State, as the position of the series after forty versions, the
      exact list of open gates in both regimes.
\end{enumerate}
\end{programme}

\begin{assessment}[Position after V69]
\label{ass:v69-position}
The rate is necessary for every local-regularity argument at the
singular time, and the series has now recorded this as a theorem rather
than an observation.  No full stability or global-regularity claim is
made.
\end{assessment}

% =============================================================================
\section*{Version V69 append-only record}
\addcontentsline{toc}{section}{Version V69 append-only record}
% =============================================================================

The complete V68 mathematical source was retained before this continuation.
Its SHA--256 digest was
\begin{center}\scriptsize\ttfamily
3a976bab63455b77440acaa5fdfe9c7d3a77974148e67ca7ff1b51f988936d87.
\end{center}
V69 proved the two-sided critical order of the cubic quantity at the
singular time, recorded the exact reason \(\varepsilon\)-regularity needs
a rate, and derived the \(L^3_tH^1\) lower bound at singular times.  No
full regularity claim is made.

% =============================================================================
\section{V70: compulsory companion audit, the intermediate class, and the
open gates after forty versions}
\label{sec:v70-entry}
% =============================================================================

This is a fortieth-version continuation.  The compulsory review of the
two companion programmes is performed first.  The rate-free results of
V54, V55, and V66--V69 are then consolidated into one theorem on the
class of solutions between the type-I rate and strong type II, and the
open gates of both regimes are listed exactly.  No global regularity
claim is made.  The next compulsory companion audit is V75.

% =============================================================================
\section{Compulsory V70 review of the current companion notes}
\label{sec:v70-companion-audit}
% =============================================================================

The current companion files in \texttt{latex\_notes} were inspected
read-only.  The frozen checkpoints are
\begin{align}
 &\texttt{Renewal\_SM\_Einstein\_Continuum\_Research\_Notes\_V74.tex},
 \notag\\[-2pt]
 &\qquad
 \texttt{a07e6c1dc5d96755cf5739ad3aed99a2deee8b83a5d7679761726f5223ffa3c3},
 \label{eq:v70-SM-checkpoint}\\
 &\texttt{Renewal\_Yang\_Mills\_Mass\_Gap\_Research\_Notes\_V52.tex},
 \notag\\[-2pt]
 &\qquad
 \texttt{674d2758caf95893ee5534a455afba81fe3e1f7ad1cf1a9d5742ce1c1aa21ff4}.
 \label{eq:v70-YM-checkpoint}
\end{align}

\emph{SM V73--V74.}  The minimally coupled scalar constraint is solved by
explicit barriers under a transverse--traceless momentum floor and a
potential ceiling, with pointwise bounds; for a high-frequency packet the
physical volume grows at least like \(N^{12/5}\), converting the
linearized frequency instability into conformal-volume escape, and volume
normalization forces York time to grow like \(N^{12/13}\) or gradient
energy to concentrate where the conformal factor collapses.  The
compactness rows of the corresponding cards remain open.

\emph{YM V52.}  The Bessel order dependence of the tilted time law is
split into parity functions, two-sided bounds for a Bessel ratio are
proved and identified with the scale \(e^{-1/s}\), the \(r!\,r^{-3/2}\)
high-moment law is proved, and a blockwise absolute card is ruled out by
a scalar diagnostic.

\begin{theorem}[V70 companion-audit decision]
\label{thm:v70-companion-decision}
No SM V73--V74 barrier, volume, or York-time constant and no YM V52
Bessel bound is imported.  The SM finding that a normalization moves but
does not remove a compactness obstruction is the same phenomenon as
Theorem~\ref{thm:v59-corrected} here, where a type-I bound along a
sequence moves the window but does not reach the singular time; it was
reached independently.  Neither companion bears on the type-I Liouville
gate or the type-II ledgers.
\end{theorem}

\begin{proof}
The hypotheses of every cited result are constraint equations on slices or
Bessel-function identities; none is a Navier--Stokes balance.
\end{proof}

% =============================================================================
\section{The intermediate class}
\label{sec:v70-intermediate}
% =============================================================================

\begin{theorem}[Rate-free facts about a singularity, V54--V69]
\label{thm:v70-intermediate}
Let \(u\) be smooth on \([0,T_*)\times\Torus^3\) with \(T_*<\infty\)
maximal, and let \(\epsilon_k^{(\theta)}\) be the geometric slab shares
of V54 and \(\bar\epsilon(\tau)\) the maximal share of V66.  Without any
rate hypothesis:
\begin{enumerate}[label=\textup{(\arabic*)},leftmargin=3.2em]
\item every slab has share at least \(2(1-\theta^{1/2})c_L^2\nu^{3/2}\theta^{k/2}\),
      and type-I along a sequence of times holds if and only if
      \(\theta^{-k/2}\epsilon_k^{(\theta)}\) is bounded along a subsequence
      (Theorem~\ref{thm:v54-dichotomy});
\item at a type-I time the rescaled solution is compact on a forward
      window of length a fixed fraction of the time remaining, and no
      further (Proposition~\ref{prop:v54-forward-window},
      Theorem~\ref{thm:v59-corrected});
\item the local energy in every ball of radius \(\rho\) is bounded by
      \(C_0(\bar\epsilon(2\rho^2)+\rho^{-1/2}E_*^{3/4}\bar\epsilon(2\rho^2)^{3/4})\)
      (Theorem~\ref{thm:v66-local-energy}), so energy cannot concentrate
      at a point of the singular time when
      \(\bar\epsilon(\tau)=o(\tau^{1/3})\), and concentration at a point
      forces \(\bar\epsilon(2\rho^2)\ge c'\rho^{2/3}\), enstrophy at
      least \(c'\rho^{-4/3}/2\) on the preceding slab, and continuation
      stock at least \(c'^2\rho^{-2/3}/2\) (Corollary~\ref{cor:v66-threshold},
      Proposition~\ref{prop:v67-continuation-cost});
\item the \(\tau^{1/3}\) threshold is unchanged by localizing the local
      part of the pressure (Proposition~\ref{prop:v68-bookkeeping});
\item energy concentration at dyadic scales is compatible with a finite
      \(L^{5/2}\) ledger (Proposition~\ref{prop:v68-ledger-cheap});
\item at the singular time the scale-invariant cubic quantity is of
      critical order from both sides, so \(\varepsilon\)-regularity
      applies only on cylinders shorter than the parabolic distance to
      \(T_*\), which a rate must supply (Proposition~\ref{prop:v69-two-sided}).
\end{enumerate}
\end{theorem}

\begin{proof}
The cited results.
\end{proof}

% =============================================================================
\section{The open gates after forty versions}
\label{sec:v70-gates}
% =============================================================================

\begin{theorem}[Open gates]
\label{thm:v70-gates}
Modulo the imports (E1)--(E8), a finite-time singularity of a smooth
periodic solution is excluded if each of the following is proved.
\begin{enumerate}[label=\textup{(G\arabic*)},leftmargin=3.8em]
\item \textbf{Type-I Liouville.}  Every element of the class described in
      Theorem~\ref{thm:v65-ancient-limit}, a nonzero ancient solution on
      \(\R^3\) bounded on every half-line with bounded enstrophy, in
      \(L^\infty_sM^{2,1}\cap L^\infty_sL^6\), with vorticity on every
      exterior region at every time and a critical-Morrey trace, is
      impossible.  Proved cases: finite energy.  Contains Leray's
      stationary problem for \(L^6\cap\dot H^1\) solutions.
\item \textbf{Intermediate class.}  Solutions with a type-I bound along a
      sequence but no rate, or with slab shares between the parabolic
      order and \(\tau^{1/3}\), are excluded.  Known: no energy
      concentration below the \(\tau^{1/3}\) threshold; no compactness
      at the singular time.
\item \textbf{Strong type II.}  Solutions with \(\theta^{-k/2}\epsilon_k^{(\theta)}\to\infty\)
      for every ratio, equivalently \((T_*-t)^{1/2}\|\Lambda u\|_2^2\to\infty\),
      are excluded.  Known: six necessary conditions in the Leray and
      \(L^{5/2}\) ledgers, all satisfied by an explicit scalar schedule;
      the spatial structure of enstrophy spikes is invisible to the
      present tools.
\end{enumerate}
Every theorem of V32--V69 is a statement about one of these three
classes; none excludes any of them.
\end{theorem}

\begin{proof}
The three classes are exhaustive by Theorem~\ref{thm:v54-dichotomy}(iii)
and the definition of the type-I rate; the listed known facts are the
cited results.
\end{proof}

\begin{assessment}[Position after V70]
\label{ass:v70-position}
Forty versions after V30, the series has converted a single stalled
estimate into three explicitly described open classes with proved
partial results in each, one withdrawn corollary, and a tabulated import
ledger.  It has not proved global regularity, and its own theorems
(Theorem~\ref{thm:v39-scale-no-go}, Proposition~\ref{prop:v69-two-sided},
Firewall~\ref{fw:v55-type-II}) identify why the tools used cannot close
any of the three gates: scale-invariant norm inequalities cannot beat
one logarithm per scale, local regularity needs a rate at the singular
time, and global ledgers do not see spike geometry.  Progress beyond this
point requires a mechanism of a different kind, and the series records
that honestly.  No full stability or global-regularity claim is made.
\end{assessment}

\begin{programme}[V71 acquisition order]
\label{prog:v71}
\begin{enumerate}[label=\textup{(AC\arabic*)},leftmargin=4.8em]
\item Gate (G1): any argument specific to the extra hypotheses of the
      class, in particular vorticity at infinity at every time and the
      critical-Morrey trace, that is not a scale-invariant norm
      inequality.
\item Gate (G3): any quantity that sees the spatial location of the
      gradient at consecutive spikes.
\item The next compulsory companion audit is V75.
\end{enumerate}
\end{programme}

% =============================================================================
\section*{Version V70 append-only record}
\addcontentsline{toc}{section}{Version V70 append-only record}
% =============================================================================

The complete V69 mathematical source was retained before this continuation.
Its SHA--256 digest was
\begin{center}\scriptsize\ttfamily
d797bc2a7a390480983d9530deaa2a743b1c7a4ac69e721582ce48ed9d37d4c1.
\end{center}
V70 performed the compulsory companion audit, consolidated the rate-free
results into one theorem on the intermediate class, and stated the three
open gates with their proved partial results.  No full regularity claim
is made.

% =============================================================================
\section{V71: the energy of the ancient limit in time}
\label{sec:v71-entry}
% =============================================================================

V41 proved that the ancient limit has infinite energy on the half-line
\((-\infty,-1]\), in the sense that its energy is not uniformly bounded
there.  This version refines the statement in time: the energy of the
ancient limit is either infinite at every time or finite from some time
onward, and in the second case the trace is square integrable, carries no
critical energy at infinity, and the centred large-scale limits vanish.
Items (AC1)--(AC2) of Programme~\ref{prog:v71} are recorded as not
advanced.  No global regularity claim is made.  The next compulsory
companion audit is V75.

% =============================================================================
\section{Finite energy propagates forward}
\label{sec:v71-energy}
% =============================================================================

\begin{proposition}[Energy dichotomy in time]
\label{prop:v71-energy-dichotomy}
Let \(U\) be the ancient limit of a type-I singularity.  Then exactly one
of the following holds.
\begin{enumerate}[label=\textup{(\roman*)},leftmargin=3.7em]
\item \(\|U(s)\|_{L^2(\R^3)}=\infty\) for every \(s<0\).
\item There is \(s_0<0\) with \(\|U(s_0)\|_2<\infty\); then
      \(\|U(s)\|_2\le\|U(s_0)\|_2\) for all \(s\in[s_0,0)\), the energy
      equality \(\frac12\|U(s)\|_2^2+\nu\int_{s_0}^s\|\nabla U\|_2^2
      =\frac12\|U(s_0)\|_2^2\) holds, the trace satisfies
      \(U_0\in L^2(\R^3)\) with \(\|U_0\|_2\le\|U(s_0)\|_2\), the trace
      has no critical energy at infinity,
      \(\rho^{-1}\int_{B_\rho}|U_0|^2\to0\), and every centred
      large-scale limit of \(U\) is zero.  Moreover
      \(\|U(s)\|_2=\infty\) for all \(s\) below
      \(s_*:=\inf\{s:\|U(s)\|_2<\infty\}\in[-\infty,0)\), and if
      \(s_*=-\infty\) then \(\|U(s)\|_2\to\infty\) as \(s\to-\infty\).
\end{enumerate}
\end{proposition}

\begin{proof}
If \(\|U(s_0)\|_2<\infty\), then on \([s_0,-\eta]\) the solution is
smooth with finite energy and bounded enstrophy, hence a strong solution
with the energy equality, which gives the monotonicity and the equality
up to \(-\eta\); let \(\eta\downarrow0\), the dissipation integral being
bounded by \(\frac12\|U(s_0)\|_2^2\).  The trace is the weak
\(L^2_{\rm loc}\) limit of \(U(s)\), which is bounded in \(L^2(\R^3)\),
so \(U_0\in L^2\) by weak lower semicontinuity.  An \(L^2\) field has
\(\rho^{-1}\int_{B_\rho}|U_0|^2\le\rho^{-1}\|U_0\|_2^2\to0\), and
Theorem~\ref{thm:v61-dichotomy} gives the vanishing of centred
large-scale limits.  Finally, below \(s_*\) the energy is infinite by definition of \(s_*\)
and by the forward monotonicity just proved (finite energy at a time
\(s'<s_*\) would propagate to \([s',0)\)).  If \(s_*=-\infty\), the
energy is finite and nonincreasing on \((-\infty,0)\), and
Corollary~\ref{cor:v41-infinite-energy} gives
\(\sup_{s\le-1}\|U(s)\|_2=\infty\); a nonincreasing function unbounded on
a half-line to the past tends to infinity as \(s\to-\infty\).
\end{proof}

\begin{assessment}[Two shapes of the ancient limit]
\label{ass:v71-shapes}
In case (i) the ancient limit has infinite energy at every time and its
trace may carry critical energy at infinity.  In case (ii) it has finite
energy from \(s_*\) onward, where either \(s_*\) is finite and the energy
is infinite before it, or \(s_*=-\infty\) and the energy is finite at
every time but grows without bound into the past; after \(s_*\) it is an
ordinary finite-energy strong solution with a singular or regular final
time; the trace
is then an \(L^2\) field of critical Morrey type, not locally \(L^3\) at
the singular points of \(U\).  Neither case is excluded.  In case (ii)
the finite-energy Liouville theorem of V41 does not apply, since it needs
bounded energy on a half-line to the past.
\end{assessment}

% =============================================================================
\section{Negative records}
\label{sec:v71-negative}
% =============================================================================

\begin{firewall}[Gates (G1) and (G3)]
\label{fw:v71-gates}
No argument specific to vorticity at infinity or to the critical-Morrey
trace was found that is not a scale-invariant norm inequality; and no
quantity seeing the spatial location of the gradient at consecutive
spikes was found.  Items (AC1)--(AC2) of Programme~\ref{prog:v71} are
recorded as open.
\end{firewall}

% =============================================================================
\section{Integrated V71 conclusion and next acquisition order}
\label{sec:v71-conclusion}
% =============================================================================

\begin{theorem}[What V71 proves]
\label{thm:v71-conclusion}
Relative to V1--V70, modulo (E1)--(E8): the energy dichotomy in time for
the ancient limit (Proposition~\ref{prop:v71-energy-dichotomy}).  The
full Navier--Stokes stability/global-regularity theorem remains unproved.
\end{theorem}

\begin{proof}
The cited result.
\end{proof}

\begin{programme}[V72 acquisition order]
\label{prog:v72}
\begin{enumerate}[label=\textup{(AC\arabic*)},leftmargin=4.8em]
\item Case (ii) of Proposition~\ref{prop:v71-energy-dichotomy}: a
      finite-energy strong solution on \([s_*,0)\) with a critical-Morrey
      trace that is not locally \(L^3\) at its singular points is itself a
      finite-energy singularity of the Navier--Stokes equations on
      \(\R^3\) if \(U\) is singular at \(0\); determine whether the V32
      strong \(L^2\) continuity and the V44--V46 structure transfer to it
      from the type-I bound it inherits, so that case (ii) reproduces on
      \(\R^3\) the situation on the torus.
\item Case (i): determine whether infinite energy at every time forces
      critical energy of the trace at infinity, or whether the energy can
      be infinite while the trace decays faster than critically.
\item The next compulsory companion audit is V75.
\end{enumerate}
\end{programme}

\begin{assessment}[Position after V71]
\label{ass:v71-position}
The ancient limit either has infinite energy at all times or has finite
energy from some time onward, possibly from the infinite past with energy
unbounded into the past; the second case reproduces the
original problem on \(\R^3\).  No full stability or global-regularity
claim is made.
\end{assessment}

% =============================================================================
\section*{Version V71 append-only record}
\addcontentsline{toc}{section}{Version V71 append-only record}
% =============================================================================

The complete V70 mathematical source was retained before this continuation.
Its SHA--256 digest was
\begin{center}\scriptsize\ttfamily
6aac28fe849c169a80e38f1a93ce497aabd8780545790a84cec83ec5aba4c2e6.
\end{center}
V71 proved the energy dichotomy in time for the ancient limit and
recorded the gate items as open.  No full regularity claim is made.

% =============================================================================
\section{V72: the finite-energy case of the ancient limit reproduces the
type-I theory on \(\R^3\)}
\label{sec:v72-entry}
% =============================================================================

Item (AC1) of Programme~\ref{prog:v72} is settled.  In case (ii) of
Proposition~\ref{prop:v71-energy-dichotomy} the ancient limit is, from
some time onward, a finite-energy strong solution on \(\R^3\) with the
type-I enstrophy bound \(\|\nabla U(s)\|_2^2\le C(-s)^{-1/2}\); every
step of the torus theory of V32--V47 that uses only the Fourier
truncations, the Ladyzhenskaya inequality, the Leray identity, and the
local energy identity transfers to \(\R^3\) verbatim, and the steps using
mean-zero Sobolev inequalities transfer with the corresponding
\(\R^3\) inequalities.  Consequently, in case (ii) the Fourier tails of \(U(s)\) are uniformly
small, the trace is attained strongly on every ball and weakly on
\(\R^3\), the critical Morrey bound holds up to \(s=0\), and the trace is
not locally cubically integrable at the singular points of \(U\); whether
energy can escape to spatial infinity at the singular time is recorded
as a gap.  Case
(ii) with a singular final time is therefore a finite-energy type-I
singularity on \(\R^3\), to which the whole series applies again.  Item
(AC2) is recorded as open.  No global regularity claim is made.  The
next compulsory companion audit is V75.

% =============================================================================
\section{The transfer}
\label{sec:v72-transfer}
% =============================================================================

\begin{theorem}[Case (ii) on \(\R^3\)]
\label{thm:v72-transfer}
Let \(U\) be the ancient limit of a type-I singularity, and suppose
\(\|U(s_0)\|_{L^2(\R^3)}<\infty\) for some \(s_0<0\).  Then, on
\([s_0,0)\), \(U\) is a smooth finite-energy solution on \(\R^3\) with
\(E_U:=\sup_{s\ge s_0}\|U(s)\|_2^2=\|U(s_0)\|_2^2\) and
\(\|\nabla U(s)\|_2^2\le C(-s)^{-1/2}\), and modulo (E1)--(E8) the
following hold, with the Fourier truncations \(P_\rho,Q_\rho\) of
\(\R^3\), the Ladyzhenskaya constant of \(\R^3\), and \(E_U\) in place of
\(E_*\).
\begin{enumerate}[label=\textup{(\roman*)},leftmargin=3.7em]
\item \emph{Cutoff-free density and tail bound (V32).}
      \(|\Phi_\rho(s)|\le C_4E_U^{1/4}\|\nabla U(s)\|_2^{5/2}\) and
      \(\|Q_\rho U(s)\|_2^2\le e^{-2\nu\rho^2\tau}E_U
      +2C_4E_U^{1/4}\int_{s-\tau}^s\|\nabla U\|_2^{5/2}\).
\item \emph{Uniform tail smallness and the trace.}
      \(\sup_{s\in[s_0,0)}\|Q_\rho U(s)\|_2^2\to0\) as \(\rho\to\infty\);
      \(U(s)\to U_0\) strongly in \(L^2(B_R)\) for every \(R\) and
      weakly in \(L^2(\R^3)\) as \(s\uparrow0\); the energy equality
      holds on \([s_0,s]\) for every \(s<0\), and
      \(\|U_0\|_2\le\lim_{s\uparrow0}\|U(s)\|_2\).
\item \emph{Critical norm.}  \(\|U(s)\|_{\dot H^{1/2}(\R^3)}^2
      \le C(1+\log^2(1/(-s)))\) for \(s\) close to \(0\), by the
      fractional record law on \(\R^3\), which holds with the same dyadic
      heat-sum kernel.
\item \emph{Morrey bound.}  \(\int_{B_\rho(y)}|U(s)|^2\le C_M'\rho\) for
      all \(y\), \(\rho\), \(s\in[s_0+\rho^2,0]\), including at the trace.
\item \emph{Terminal datum.}  At every concentration point \(y_0\) of
      \(U\) at \(s=0\), \(U_0\notin L^3(B_\rho(y_0))\) for every
      \(\rho>0\); the singular set of \(U\) at \(s=0\) is the
      \(L^3\)-singular support of \(U_0\), of Hausdorff dimension zero,
      contained in \(\overline{B_{R_0}}\).
\end{enumerate}
\end{theorem}

\begin{proof}
Finite energy from \(s_0\) is Proposition~\ref{prop:v71-energy-dichotomy}.
(i) The proof of Theorem~\ref{thm:v32-cutoff-free-density} uses
\(\Phi_\rho=-\langle\BNS(U),Q_\rho U\rangle\), H\"older, and the
Ladyzhenskaya inequality, all valid on \(\R^3\); the Gr\"onwall step of
Theorem~\ref{thm:v32-recent-ledger-tail} uses \(V_\rho\ge\rho^2H_\rho\),
which holds for the \(\R^3\) truncation.  (ii) Insert the type-I bound in (i) as in
Theorem~\ref{thm:v32-type-I-firewall}: the window ledger is at most
\(C\tau^{3/8}\), so the Fourier tails are uniformly small.  On a ball
\(B_R\), the low-frequency parts \(P_\rho U(s)\) have uniformly bounded
derivatives of all orders (Bernstein on \(\R^3\) with the energy bound)
and converge in \(H^{-1}(B_R)\) by the argument of
Theorem~\ref{thm:v48-trace-identification}, hence in \(L^2(B_R)\) by
interpolation; the high-frequency parts are uniformly small in
\(L^2(\R^3)\); so \(U(s)\) is Cauchy in \(L^2(B_R)\) as \(s\uparrow0\)
and converges strongly there to \(U_0\).  Weak convergence in
\(L^2(\R^3)\) follows from the uniform energy bound and the local strong
convergence.  The energy equality on \([s_0,s]\) is that of a strong
solution, and \(\|U_0\|_2\le\lim_{s\uparrow0}\|U(s)\|_2\) is weak lower
semicontinuity; equality would require that no energy escape to spatial
infinity as \(s\uparrow0\), which the present bounds do not decide
(Firewall~\ref{fw:v72-gap}).
(iii) The dyadic heat sum and domination of V33 are lattice-free.  (iv)
The proof of Theorem~\ref{thm:v47-morrey} uses the \(L^6\) bound
\(\|U\|_6\le C_6\|\nabla U\|_2\), valid on \(\R^3\), the Riesz bound,
and strong local convergence at the trace, which (ii) provides on balls.
(v) The proofs of Theorems~\ref{thm:v46-non-L3}, \ref{thm:v43-dimension-zero},
and \ref{thm:v51-characterization} are local in space and use only the
bounds above and (E1)--(E6).
\end{proof}

\begin{firewall}[A gap in the transfer, recorded]
\label{fw:v72-gap}
The proof of (ii) shows that strong \(L^2(\R^3)\) convergence of \(U(s)\)
to its trace is not obtained by the torus argument, whose compactness came
from the boundedness of the torus; on \(\R^3\) a defect of energy at
spatial infinity is not excluded by the present bounds.  The statements
(iv)--(v), which are local, do not need it.  Whether the energy of a
finite-energy type-I solution on \(\R^3\) can escape to spatial infinity
at the singular time is left open.
\end{firewall}

\begin{assessment}[Case (ii) is the original problem on \(\R^3\)]
\label{ass:v72-reproduction}
If the ancient limit has finite energy at some time and is singular at
\(s=0\), it is a finite-energy solution on \(\R^3\) with a type-I
enstrophy rate at its singular time, with the local structure of
Theorem~\ref{thm:v72-transfer}.  Rescaling it at its own concentration
points produces again an element of the class.  The construction can
therefore be iterated, and each iteration either enters case (i)
(infinite energy at every time) or reproduces a finite-energy type-I
singularity.  Nothing proved here prevents the iteration from
reproducing the same object indefinitely, which is what a self-similar
singularity would do.
\end{assessment}

% =============================================================================
\section{Integrated V72 conclusion and next acquisition order}
\label{sec:v72-conclusion}
% =============================================================================

\begin{theorem}[What V72 proves]
\label{thm:v72-conclusion}
Relative to V1--V71, modulo (E1)--(E8): the transfer of the type-I local
theory to the finite-energy case of the ancient limit
(Theorem~\ref{thm:v72-transfer}) and the record of the gap at spatial
infinity (Firewall~\ref{fw:v72-gap}).
Item (AC2) of Programme~\ref{prog:v72} is open.  The full Navier--Stokes
stability/global-regularity theorem remains unproved.
\end{theorem}

\begin{proof}
The cited results.
\end{proof}

\begin{programme}[V73 acquisition order]
\label{prog:v73}
\begin{enumerate}[label=\textup{(AC\arabic*)},leftmargin=4.8em]
\item Energy at spatial infinity for finite-energy type-I solutions on
      \(\R^3\): use the uniform Fourier tail smallness of
      Theorem~\ref{thm:v72-transfer}(ii) together with a spatial cutoff to
      decide whether energy can escape to infinity at the singular time,
      that is, whether \(\|U_0\|_2<\lim_{s\uparrow0}\|U(s)\|_2\) is
      possible.
\item Iteration.  Formalize the iteration of Assessment~\ref{ass:v72-reproduction}
      and determine whether an invariant quantity, such as the fraction
      \(c_*\) of slab dissipation in the concentration cylinder, is
      monotone along it.
\item The next compulsory companion audit is V75.
\end{enumerate}
\end{programme}

\begin{assessment}[Position after V72]
\label{ass:v72-position}
The finite-energy case of the ancient limit reproduces the original
problem on \(\R^3\), with one gap at spatial infinity recorded.  No full
stability or global-regularity claim is made.
\end{assessment}

% =============================================================================
\section*{Version V72 append-only record}
\addcontentsline{toc}{section}{Version V72 append-only record}
% =============================================================================

The complete V71 mathematical source was retained before this continuation.
Its SHA--256 digest was
\begin{center}\scriptsize\ttfamily
108229cff23f5f9bb7de8815928650817a196c73397156b423b0c93c6a6520ef.
\end{center}
V72 transferred the type-I local theory to the finite-energy case of the
ancient limit on \(\R^3\), recorded the gap at spatial infinity, and
noted that the construction iterates.  No full regularity claim is made.

% =============================================================================
\section{V73: no energy escapes to spatial infinity at a type-I singular
time on \(\R^3\)}
\label{sec:v73-entry}
% =============================================================================

Item (AC1) of Programme~\ref{prog:v73} is settled, and the gap recorded
in Firewall~\ref{fw:v72-gap} is closed.  For a finite-energy solution on
\(\R^3\) with a type-I enstrophy rate, the band-limited parts
\(P_\rho U(s)\) have a time derivative whose \(L^2(\R^3)\) norm is
integrable up to the singular time, so they converge strongly in
\(L^2(\R^3)\); combined with the uniform smallness of the Fourier tails,
this gives strong convergence of \(U(s)\) to its trace in \(L^2(\R^3)\),
the energy equality up to the singular time, and the impossibility of
energy escaping to spatial infinity.  Item (AC2) is recorded as open.  No
global regularity claim is made.  The next compulsory companion audit is
V75.

% =============================================================================
\section{Strong attainment of the trace on \(\R^3\)}
\label{sec:v73-strong}
% =============================================================================

\begin{theorem}[Strong \(L^2(\R^3)\) continuity at a type-I singular time]
\label{thm:v73-strong}
Let \(U\) be a smooth divergence-free solution on \([s_0,0)\times\R^3\)
with finite energy \(E_U=\sup_s\|U(s)\|_2^2\), pressure given by the Riesz
transforms, and \(\|\nabla U(s)\|_2^2\le C(-s)^{-1/2}\).  Then \(U(s)\)
converges strongly in \(L^2(\R^3)\) as \(s\uparrow0\) to a limit
\(U_0\), and
\begin{equation}
 \boxed{\quad
 \frac12\|U_0\|_2^2+\nu\int_{s_0}^0\|\nabla U(s)\|_2^2\,\dd s
 =\frac12\|U(s_0)\|_2^2 .
 \quad}
 \label{eq:v73-energy-equality}
\end{equation}
In particular no energy escapes to spatial infinity at the singular time,
and Theorem~\ref{thm:v72-transfer}(ii) holds with strong convergence in
\(L^2(\R^3)\).
\end{theorem}

\begin{proof}
Let \(P_\rho\) be the Fourier truncation to \(\{|\xi|\le\rho\}\) on
\(\R^3\).  Applying \(P_\rho\) to the equation,
\(\partial_sP_\rho U=\nu\Delta P_\rho U-P_\rho\PP\operatorname{div}(U\otimes U)\),
and the multiplier bounds on \(\{|\xi|\le\rho\}\) give
\[
 \|\partial_sP_\rho U(s)\|_2
 \le\nu\rho^2\|U(s)\|_2+\rho\|U(s)\otimes U(s)\|_2
 \le\nu\rho^2E_U^{1/2}+\rho\,C_4E_U^{1/4}\|\nabla U(s)\|_2^{3/2}
 \le\nu\rho^2E_U^{1/2}+C\rho(-s)^{-3/8},
\]
which is integrable on \((s_0,0)\).  Hence \(P_\rho U(s)\) is Cauchy in
\(L^2(\R^3)\) as \(s\uparrow0\), for each fixed \(\rho\), and converges
strongly to some \(W_\rho\).  By Theorem~\ref{thm:v72-transfer}(i)--(ii),
\(\eta(\rho):=\sup_{s\in[s_0,0)}\|Q_\rho U(s)\|_2^2\to0\) as
\(\rho\to\infty\).  For \(s,s'<0\),
\[
 \|U(s)-U(s')\|_2\le\|P_\rho U(s)-P_\rho U(s')\|_2+2\eta(\rho)^{1/2},
\]
so \(\limsup_{s,s'\uparrow0}\|U(s)-U(s')\|_2\le2\eta(\rho)^{1/2}\) for
every \(\rho\), hence \(U(s)\) is Cauchy in \(L^2(\R^3)\) and converges
strongly to a limit \(U_0\), which coincides with the weak limit of
Theorem~\ref{thm:v72-transfer}(ii).  The energy equality on
\([s_0,s]\) passes to the limit \(s\uparrow0\) by strong convergence and
monotone convergence of the dissipation integral.
\end{proof}

\begin{corollary}[The gap of V72 is closed]
\label{cor:v73-gap-closed}
Firewall~\ref{fw:v72-gap} is withdrawn: in case (ii) of
Proposition~\ref{prop:v71-energy-dichotomy}, energy cannot escape to
spatial infinity at the singular time, and the ancient limit on
\([s_0,0]\) has exactly the structure of the torus solution described in
Theorem~\ref{thm:v50-ancient-limit}(2), with \(\R^3\) in place of the
torus.
\end{corollary}

\begin{proof}
Theorem~\ref{thm:v73-strong} supplies the strong convergence that
Theorem~\ref{thm:v72-transfer}(ii) lacked; the remaining items of
Theorem~\ref{thm:v72-transfer} were already local.
\end{proof}

\begin{remark}[The same argument on the torus]
\label{rem:v73-torus}
The proof of Theorem~\ref{thm:v73-strong} is the argument behind
Theorem~\ref{thm:v32-type-I-firewall}(i) made explicit: on the torus,
uniform tail smallness alone gives precompactness by the Kolmogorov--Riesz
criterion, and the strong convergence of the band-limited part follows
from the same time-derivative bound.  On \(\R^3\) the time-derivative
bound is essential, because precompactness in \(L^2(\R^3)\) also
requires tightness in space, which is exactly what the strong convergence
of \(P_\rho U(s)\) provides.
\end{remark}

% =============================================================================
\section{Integrated V73 conclusion and next acquisition order}
\label{sec:v73-conclusion}
% =============================================================================

\begin{theorem}[What V73 proves]
\label{thm:v73-conclusion}
Relative to V1--V72, with no additional import: strong \(L^2(\R^3)\)
continuity at a type-I singular time for finite-energy solutions on
\(\R^3\) with the energy equality \eqref{eq:v73-energy-equality}, closing
the gap of V72.  Item (AC2) of Programme~\ref{prog:v73} is open.  The
full Navier--Stokes stability/global-regularity theorem remains unproved.
\end{theorem}

\begin{proof}
The cited results.
\end{proof}

\begin{programme}[V74 acquisition order]
\label{prog:v74}
\begin{enumerate}[label=\textup{(AC\arabic*)},leftmargin=4.8em]
\item Iteration of the ancient-limit construction: with case (ii) now
      fully reproducing the torus structure on \(\R^3\), formalize the
      map from a finite-energy type-I singularity to the class of its
      ancient limits and determine what, if anything, decreases along it.
\item Case (i): infinite energy at every time; determine whether the
      time-derivative bound of Theorem~\ref{thm:v73-strong}, which needs
      only the energy at a single time in the band-limited part, gives a
      local version of strong continuity on balls without the global
      energy.
\item The next compulsory companion audit is V75.
\end{enumerate}
\end{programme}

\begin{assessment}[Position after V73]
\label{ass:v73-position}
The finite-energy case of the ancient limit is now a faithful copy of
the original type-I problem on \(\R^3\), with strong continuity at the
singular time and no energy loss.  No full stability or global-regularity
claim is made.
\end{assessment}

% =============================================================================
\section*{Version V73 append-only record}
\addcontentsline{toc}{section}{Version V73 append-only record}
% =============================================================================

The complete V72 mathematical source was retained before this continuation.
Its SHA--256 digest was
\begin{center}\scriptsize\ttfamily
a54cb89f717149e2b91cdc1dfb406c08bf7f9bd67064d60a8799b030e2efff5e.
\end{center}
V73 proved strong \(L^2(\R^3)\) continuity at a type-I singular time for
finite-energy solutions on \(\R^3\), with the energy equality, closing the
gap recorded in V72.  No full regularity claim is made.

% =============================================================================
\section{V74: the two defects coincide}
\label{sec:v74-entry}
% =============================================================================

The local energy identity near the final time gives a two-sided bound,
uniform in the rescaling index, on the difference between the local
energy of a rescaled field at time \(s\) and at its final time.  Passing
to the limit shows that the local energy of the ancient limit at its
final time equals the limit of the local energies of the rescaled
terminal data.  Consequently the trace \(U_0\) is attained strongly in
\(L^2_{\rm loc}\) by \(U(s)\) if and only if the rescaled terminal data
converge strongly in \(L^2_{\rm loc}\) to \(U_0\): the energy defect of
the trace and the weak-limit defect of the rescaled terminal data are the
same number.  Both items of Programme~\ref{prog:v74} are recorded as not
advanced.  No global regularity claim is made.  The next compulsory
companion audit is V75.

% =============================================================================
\section{Two-sided local energy continuity}
\label{sec:v74-two-sided}
% =============================================================================

Let \(x_0\in\Sigma\), \(r_k\), \(U_k\), \(U\), \(T_k\), \(U_0\) be as in
V41 and V47--V48, with \(e_R^{(k)}(s)=\frac12\int|U_k(s)|^2\phi_R\),
\(e_R(s)=\frac12\int|U(s)|^2\phi_R\), and
\(e_R^{(k)}(0^-)=\frac12\int|T_k|^2\phi_R\).

\begin{lemma}[Two-sided bound]
\label{lem:v74-two-sided}
There is \(C_R\), independent of \(k\), with
\begin{equation}
 \bigl|e_R^{(k)}(s)-e_R^{(k)}(0^-)\bigr|
 \le\nu\int_s^0Y_k+C_R(-s)^{1/4}
 \le C_R'(-s)^{1/4}
 \qquad(-1\le s<0).
 \label{eq:v74-two-sided}
\end{equation}
\end{lemma}

\begin{proof}
In the local energy identity of Lemma~\ref{lem:v46-local-energy-near-zero},
\(e^{(k)}_R(0^-)-e^{(k)}_R(s)=-\nu\iint_s^0|\nabla U_k|^2\phi_R+[\Delta]+[{\rm flux}]\),
the Laplacian and flux terms are bounded in absolute value by
\(C(-s)^{1/2}+C_R(-s)^{1/4}\) uniformly in \(k\), and the dissipation
term by \(\nu\int_s^0Y_k\le2\nu C_I^2\nu^{3/2}(-s)^{1/2}\); this gives
both inequalities.
\end{proof}

\begin{theorem}[The local energy at the trace time is the limit of the
rescaled terminal local energies]
\label{thm:v74-defects}
Along the subsequence realizing \(T_k\rightharpoonup U_0\), the limits
\(\ell_R:=\lim_k\frac12\int|T_k|^2\phi_R\) exist for every \(R\) after a
further subsequence, and
\begin{equation}
 \lim_{s\uparrow0}e_R(s)=\ell_R,
 \qquad
 \lim_{s\uparrow0}e_R(s)-\frac12\int|U_0|^2\phi_R
 =\ell_R-\frac12\int|U_0|^2\phi_R\ \ge0 .
 \label{eq:v74-defects}
\end{equation}
Consequently \(U(s)\to U_0\) strongly in \(L^2_{\rm loc}\) as
\(s\uparrow0\) if and only if \(T_k\to U_0\) strongly in \(L^2_{\rm loc}\)
along the subsequence.
\end{theorem}

\begin{proof}
The rescaled terminal local energies are bounded by \(C_M(2R)\)
(Theorem~\ref{thm:v47-morrey}), so a diagonal subsequence makes all
\(\ell_R\) exist.  For fixed \(s<0\), \(e_R^{(k)}(s)\to e_R(s)\) by
strong local convergence at fixed time.  Letting \(k\to\infty\) in
\eqref{eq:v74-two-sided} gives \(|e_R(s)-\ell_R|\le C_R'(-s)^{1/4}\),
which is the first identity.  Weak convergence \(T_k\rightharpoonup U_0\)
gives \(\ell_R\ge\frac12\int|U_0|^2\phi_R\), and \(U(s)\rightharpoonup U_0\)
gives \(\lim e_R(s)\ge\frac12\int|U_0|^2\phi_R\); the two defects are
equal by the first identity.  Strong convergence in \(L^2(B_R)\) of a
weakly convergent family is equivalent to convergence of norms, that is,
to vanishing of the defect, for every \(R\); the two defects being equal,
the two strong convergences are equivalent.
\end{proof}

\begin{corollary}[Strong attainment in the type-I case is a property of
the terminal datum alone]
\label{cor:v74-terminal}
The trace of the ancient limit is attained strongly in \(L^2_{\rm loc}\)
if and only if the terminal datum \(u(T_*)\) of the original solution
satisfies: the rescalings \(r_ku(T_*,x_0+r_k\cdot)\) converge strongly in
\(L^2_{\rm loc}\) along the concentration scales.  In the finite-energy
case (ii) of Proposition~\ref{prop:v71-energy-dichotomy} the trace is
attained strongly by Theorem~\ref{thm:v73-strong}, so in that case the
rescaled terminal data converge strongly.
\end{corollary}

\begin{proof}
The first statement is Theorem~\ref{thm:v74-defects}; the second combines
it with Theorem~\ref{thm:v73-strong}.
\end{proof}

\begin{firewall}[Programme~\ref{prog:v74}]
\label{fw:v74-items}
Item (AC1), an invariant decreasing along the iteration, was not found;
item (AC2), local strong continuity in case (i), is exactly the
vanishing of the defect of Theorem~\ref{thm:v74-defects}, which is not
decided in case (i).  Both are open.
\end{firewall}

% =============================================================================
\section{Integrated V74 conclusion and next acquisition order}
\label{sec:v74-conclusion}
% =============================================================================

\begin{theorem}[What V74 proves]
\label{thm:v74-conclusion}
Relative to V1--V73, modulo (E1)--(E8): the two-sided local energy
continuity \eqref{eq:v74-two-sided}, the identity of the two defects
\eqref{eq:v74-defects}, and the equivalence of strong attainment of the
trace with strong convergence of the rescaled terminal data.  The full
Navier--Stokes stability/global-regularity theorem remains unproved.
\end{theorem}

\begin{proof}
The cited results.
\end{proof}

\begin{programme}[V75 acquisition order, with compulsory companion audit]
\label{prog:v75}
\begin{enumerate}[label=\textup{(AC\arabic*)},leftmargin=4.8em]
\item Perform the compulsory review of the current SM and YM notes;
      record digests; import nothing numerical.
\item Strong convergence of rescaled terminal data.  Determine whether
      the terminal datum of a type-I singularity on the torus, which is
      in \(L^2\) with critical Morrey bound and not locally \(L^3\) at
      concentration points, has strongly convergent rescalings along the
      concentration scales; a positive answer would place every ancient
      limit in the strongly attained case and make the defect zero.
\item The next compulsory companion audit after V75 is V80.
\end{enumerate}
\end{programme}

\begin{assessment}[Position after V74]
\label{ass:v74-position}
Strong attainment of the trace of the ancient limit is now a question
about the terminal datum alone.  No full stability or global-regularity
claim is made.
\end{assessment}

% =============================================================================
\section*{Version V74 append-only record}
\addcontentsline{toc}{section}{Version V74 append-only record}
% =============================================================================

The complete V73 mathematical source was retained before this continuation.
Its SHA--256 digest was
\begin{center}\scriptsize\ttfamily
8c256edcd36fd79a760bf4fcb628887251fbf39e7031c699d404fa4b14ed0d71.
\end{center}
V74 proved the two-sided local energy continuity near the final time and
the coincidence of the trace defect with the weak-limit defect of the
rescaled terminal data.  No full regularity claim is made.

% =============================================================================
\section{V75: compulsory companion audit and the rescaled terminal data}
\label{sec:v75-entry}
% =============================================================================

This is a forty-fifth-version continuation.  The compulsory review of
the two companion programmes is performed first.  Item (AC2) of
Programme~\ref{prog:v75} is then analysed: strong \(L^2_{\rm loc}\)
convergence of the rescaled terminal data along the concentration scales
is equivalent, by V74, to strong attainment of the trace of the ancient
limit, and it is shown to hold whenever the ancient limit has finite
energy at some time (case (ii) of V71), by V73.  In case (i) the question
reduces to whether the terminal datum can oscillate at the concentration
scales, which the critical Morrey bound does not exclude; a scalar
example shows that a field with all the proved properties of the terminal
datum can have non-strongly-convergent rescalings, so the answer cannot
come from those properties alone.  No global regularity claim is made.
The next compulsory companion audit is V80.

% =============================================================================
\section{Compulsory V75 review of the current companion notes}
\label{sec:v75-companion-audit}
% =============================================================================

The current companion files in \texttt{latex\_notes} were inspected
read-only.  The frozen checkpoints are
\begin{align}
 &\texttt{Renewal\_SM\_Einstein\_Continuum\_Research\_Notes\_V75.tex},
 \notag\\[-2pt]
 &\qquad
 \texttt{6678a93c8e0910e32898c18e97dfd3e8d58e91a07132ef2d3654062c3431a0d2},
 \label{eq:v75-SM-checkpoint}\\
 &\texttt{Renewal\_Yang\_Mills\_Mass\_Gap\_Research\_Notes\_V54.tex},
 \notag\\[-2pt]
 &\qquad
 \texttt{e34c14c00621502eb74c6b1ed2397dea9c4ca42cfe67fcde6cddebd62178eeaf}.
 \label{eq:v75-YM-checkpoint}
\end{align}
At the time of reading the YM V54 file had the same digest as the YM V53
file, that is, it was a renamed copy not yet extended.

\emph{SM V75.}  A negative DeWitt trace coordinate has no positive real
Gaussian law and its real Euclidean integral diverges for every source;
rotation to an imaginary contour gives an exact convergent factor but no
positive law for trace observables, and the programme records a
contour-or-constraint fork.

\emph{YM V53.}  An exact compact/wrapped split of a Schl\"afli-type
integral is given, with a uniform complex maximum on the compact sector,
coefficient bounds \((C/\alpha)^r\), a strict-disk wrapped maximum, a
two-sector moment majorant, localization of every all-order thermal-radius
loss to the wrapped hyperbolic integral, and an impossibility result for
post-differentiation nonperturbative absorption.

\begin{theorem}[V75 companion-audit decision]
\label{thm:v75-companion-decision}
No SM V75 contour or DeWitt constant and no YM V53 Schl\"afli bound is
imported.  Neither bears on the type-I Liouville gate, the intermediate
class, or the type-II ledgers.
\end{theorem}

\begin{proof}
The hypotheses of every cited result are Gaussian integrals over metric
coordinates or Bessel-type integrals; none is a Navier--Stokes balance.
\end{proof}

% =============================================================================
\section{Strong convergence of the rescaled terminal data}
\label{sec:v75-rescaled-terminal}
% =============================================================================

\begin{proposition}[Case (ii) settles the question positively]
\label{prop:v75-case-ii}
If the ancient limit at \(x_0\) has finite energy at some time, then the
rescaled terminal data \(T_k=r_ku(T_*,x_0+r_k\cdot)\) converge strongly
in \(L^2_{\rm loc}(\R^3)\) to \(U_0\) along the subsequence of
Theorem~\ref{thm:v74-defects}, and the defect vanishes.
\end{proposition}

\begin{proof}
Theorem~\ref{thm:v73-strong} gives strong attainment of the trace in
\(L^2(\R^3)\), hence in \(L^2_{\rm loc}\); Theorem~\ref{thm:v74-defects}
transfers it to the rescaled terminal data.
\end{proof}

\begin{counterexample}[The proved properties of the terminal datum do not
force strong convergence]
\label{cex:v75-oscillation}
Let \(f(x)=|x|^{-1}\bigl(1+\tfrac12\sin(\log|x|)\bigr)\,\mathbf 1_{|x|\le1}\)
on \(\R^3\), or its divergence-free analogue obtained by multiplying a
fixed unit vector field tangent to spheres.  Then \(f\in L^2\), \(f\) has
the critical Morrey bound \(\int_{B_\rho}|f|^2\le C\rho\), \(f\notin
L^3(B_\rho(0))\) for every \(\rho\), and the upper critical density at
\(0\) is positive.  Along \(r_k=e^{-2\pi k}\) the rescalings
\(r_kf(r_k\cdot)\) are all equal to \(f\) on \(B_{e^{2\pi k}}\), so they
converge strongly; but along \(r_k=e^{-2\pi k-\pi}\) they equal
\(|y|^{-1}(1-\tfrac12\sin\log|y|)\), a different field, so the family of
all rescalings has two distinct limit points and does not converge along
the full sequence \(r=e^{-\pi k}\).  Adding an oscillation at scales
faster than any concentration sequence, such as
\(|x|^{-1}\sin(|x|^{-1})\) suitably cut off, keeps every listed property
and produces rescalings that converge weakly but not strongly.
\end{counterexample}

\begin{assessment}[What decides strong convergence]
\label{ass:v75-what-decides}
Counterexample~\ref{cex:v75-oscillation} shows that membership in
\(L^2\), the critical Morrey bound, non-\(L^3\) at the point, and
positive upper density do not force strong convergence of rescalings; the
dynamics must supply it.  In case (ii) it does, through the energy.  In
case (i) the only dynamical input available is the two-sided local
energy continuity of V74, which converts the question into the vanishing
of a single defect number per ball and does not decide it.  Item (AC2)
is therefore settled in case (ii) and open in case (i).
\end{assessment}

% =============================================================================
\section{Integrated V75 conclusion and next acquisition order}
\label{sec:v75-conclusion}
% =============================================================================

\begin{theorem}[What V75 establishes]
\label{thm:v75-conclusion}
Relative to V1--V74: the companion-audit decision
(Theorem~\ref{thm:v75-companion-decision}); strong convergence of the
rescaled terminal data in the finite-energy case
(Proposition~\ref{prop:v75-case-ii}); and the counterexample showing
that the proved properties of the terminal datum do not force it
(Counterexample~\ref{cex:v75-oscillation}).  The full Navier--Stokes
stability/global-regularity theorem remains unproved.
\end{theorem}

\begin{proof}
The cited items.
\end{proof}

\begin{programme}[V76 acquisition order]
\label{prog:v76}
\begin{enumerate}[label=\textup{(AC\arabic*)},leftmargin=4.8em]
\item Case (i) versus case (ii).  Determine whether the ancient limit of
      a type-I singularity on the torus can be in case (i) at all: the
      rescaled fields have energy \(r_k^{-1}E_*\) on tori of size
      \(r_k^{-1}\), and the local limit has infinite energy only if the
      energy does not stay within bounded rescaled distance of the
      concentration point; test this against the Morrey bound, which
      allows energy of order \(R\) in \(B_R\), and against the
      dimension-zero singular set, which forces energy away from \(x_0\)
      to sit near regular points where \(u(T_*)\) is bounded.
\item The next compulsory companion audit is V80.
\end{enumerate}
\end{programme}

\begin{assessment}[Position after V75]
\label{ass:v75-position}
Forty-five versions after V30, the ancient limit of a type-I singularity
is described in detail, with one remaining dichotomy, finite or infinite
energy at the trace time, on which the strong attainment of the trace
depends.  No full stability or global-regularity claim is made.
\end{assessment}

% =============================================================================
\section*{Version V75 append-only record}
\addcontentsline{toc}{section}{Version V75 append-only record}
% =============================================================================

The complete V74 mathematical source was retained before this continuation.
Its SHA--256 digest was
\begin{center}\scriptsize\ttfamily
fe01054a54011ab3bb7a797697de0586fbe0df8024afb2113e5a41ed805973b6.
\end{center}
V75 performed the compulsory companion audit, settled strong convergence
of the rescaled terminal data in the finite-energy case, and gave a
counterexample showing that the proved properties of the terminal datum
do not force it in general.  No full regularity claim is made.

% =============================================================================
\section{V76: the energy cases and the terminal datum}
\label{sec:v76-entry}
% =============================================================================

Item (AC1) of Programme~\ref{prog:v76} is analysed and found undecidable
from the proved properties.  The rescaled fields at a concentration point
have energy \(r_k^{-1}E_*\), which tends to infinity, so the ancient limit
has finite energy at some time only if almost all of that energy escapes
to rescaled spatial infinity in the limit; nothing proved forbids this,
and nothing forces it.  A terminal datum with a homogeneous profile of
degree \(-1\) at the singular point, the self-similar shape, gives an
ancient limit in case (i); a terminal datum whose critical local energy
at the singular point is carried by scales comparable to the concentration
scales only, with subcritical energy at larger rescaled radii, is
compatible with case (ii).  Both shapes satisfy every property proved in
V44--V75.  No global regularity claim is made.  The next compulsory
companion audit is V80.

% =============================================================================
\section{Energy of the rescaled fields and of the limit}
\label{sec:v76-energy}
% =============================================================================

\begin{proposition}[Where the energy goes]
\label{prop:v76-energy}
Let \(U\) be the ancient limit at \(x_0\) with scales \(r_k\), and
\(s<0\).  Then for every \(R\),
\begin{equation}
 \int_{B_R}|U(s)|^2\le\liminf_{k\to\infty}\frac1{r_k}\int_{B_{Rr_k}(x_0)}|u(T_*+r_k^2s)|^2
 \le C_MR ,
 \label{eq:v76-local}
\end{equation}
while \(\|U_k(s)\|_{L^2(\Torus^3_k)}^2=r_k^{-1}\|u(T_*+r_k^2s)\|_2^2\to\infty\).
Consequently:
\begin{enumerate}[label=\textup{(\roman*)},leftmargin=3.7em]
\item in case (ii) of Proposition~\ref{prop:v71-energy-dichotomy}, for
      every \(\varepsilon>0\) there is \(R\) such that
      \(\liminf_kr_k^{-1}\int_{B_{Rr_k}(x_0)}|u(T_*+r_k^2s)|^2
      \ge\|U(s)\|_2^2-\varepsilon\), and the remaining energy
      \(r_k^{-1}\int_{\Torus^3\setminus B_{Rr_k}(x_0)}|u|^2\to\infty\)
      sits at rescaled distance at least \(R\); the fraction of the
      rescaled energy within any fixed rescaled radius tends to zero;
\item in case (i), for every \(A\) there is \(R\) with
      \(\int_{B_R}|U(s)|^2\ge A\), so the local energy of \(u\) near
      \(x_0\) at scale \(Rr_k\) is at least \(Ar_k\) for large \(k\): the
      terminal energy profile is of critical size \(\asymp\rho\) on a
      range of scales \(\rho\in[r_k,Rr_k]\) growing without bound in
      rescaled units.
\end{enumerate}
\end{proposition}

\begin{proof}
\eqref{eq:v76-local} is weak lower semicontinuity on \(B_R\) and the
Morrey bound; the energy identity of the rescaling gives the divergence.
(i) and (ii) read off the definitions of the two cases.
\end{proof}

\begin{assessment}[Two admissible shapes]
\label{ass:v76-shapes}
A terminal datum behaving like \(|x-x_0|^{-1}\) on \(B_{\rho_0}(x_0)\)
has \(\rho^{-1}\int_{B_\rho}|u(T_*)|^2\asymp1\) for all \(\rho\le\rho_0\),
its rescalings converge to \(|y|^{-1}\) times a profile, the trace has
infinite energy, and the ancient limit is in case (i); this is the
self-similar shape, consistent with every result of the series.  A
terminal datum with critical energy only at scales near each
concentration scale, for instance a sum over \(k\) of bumps of
\(L^2\) mass \(\asymp r_k\) supported on the annuli
\(r_k\le|x-x_0|\le2r_k\), with lacunary \(r_k\),
has rescalings converging to a single bump, a trace of finite energy, and
is compatible with case (ii).  The dynamics may exclude one of these, but
no theorem here does; in particular the second shape is compatible with
the dimension-zero singular set, the Morrey bound, and the non-\(L^3\)
property.
\end{assessment}

% =============================================================================
\section{Integrated V76 conclusion and next acquisition order}
\label{sec:v76-conclusion}
% =============================================================================

\begin{theorem}[What V76 establishes]
\label{thm:v76-conclusion}
Relative to V1--V75: Proposition~\ref{prop:v76-energy} and the record
that both energy cases are compatible with every proved property of the
terminal datum.  The full Navier--Stokes stability/global-regularity
theorem remains unproved.
\end{theorem}

\begin{proof}
The cited items.
\end{proof}

\begin{programme}[V77 acquisition order]
\label{prog:v77}
\begin{enumerate}[label=\textup{(AC\arabic*)},leftmargin=4.8em]
\item Determine whether the lacunary shape of Assessment~\ref{ass:v76-shapes}
      is compatible with the two-sided slab dissipation of
      Lemma~\ref{lem:v43-slab} and the per-scale implication of
      Theorem~\ref{thm:v53-per-scale}: between concentration scales the
      terminal density is subcritical, so the rescaled dissipation at
      intermediate scales must vanish, while the slab dissipation is of
      order \(\rho\) at every scale; test whether this forces the
      dissipation at intermediate scales to sit away from \(x_0\), at
      other singular points or spread out.
\item The next compulsory companion audit is V80.
\end{enumerate}
\end{programme}

\begin{assessment}[Position after V76]
\label{ass:v76-position}
The finite/infinite energy dichotomy of the ancient limit corresponds to
two admissible shapes of the terminal datum at the singular point, and
the series does not decide between them.  No full stability or
global-regularity claim is made.
\end{assessment}

% =============================================================================
\section*{Version V76 append-only record}
\addcontentsline{toc}{section}{Version V76 append-only record}
% =============================================================================

The complete V75 mathematical source was retained before this continuation.
Its SHA--256 digest was
\begin{center}\scriptsize\ttfamily
ed584d59a6163c2b59d866f31a6ef776ea46fd2060d79be6603753a60d51df0c.
\end{center}
V76 related the energy cases of the ancient limit to the terminal energy
profile at the singular point and recorded that both cases are admissible.
No full regularity claim is made.

% =============================================================================
\section{V77: lacunary singular points need satellites}
\label{sec:v77-entry}
% =============================================================================

Item (AC1) of Programme~\ref{prog:v77} is settled by a satellite theorem.
Along scales at which the terminal datum has subcritical energy around a
singular point at all fixed rescaled radii, the per-scale implication of
V53 makes the rescaled dissipation near the point vanish, while the
two-sided slab dissipation of V43 keeps the total of order the scale;
since regular regions contribute only the square of the scale, the
dissipation must sit in every neighbourhood of the singular set but
outside every fixed rescaled ball around the point.  Under the
\(L^\infty\) type-I bound it occupies volume of order the cube of the
scale there.  Thus the lacunary shape of V76 forces, at each intermediate
scale, a satellite structure at rescaled distance tending to infinity
from the point and converging to the singular set.  For a singular set
reduced to one point the satellites converge to that point.  No global
regularity claim is made.  The next compulsory companion audit is V80.

% =============================================================================
\section{The satellite theorem}
\label{sec:v77-satellite}
% =============================================================================

Throughout, \eqref{eq:v32-type-I} holds, \(x_0\in\Sigma\), and
\(\rho_j\downarrow0\) is a sequence of scales along which
\begin{equation}
 \frac1{\rho_j}\int_{B_{R\rho_j}(x_0)}|u(T_*)|^2\longrightarrow0
 \qquad\text{for every fixed }R>0
 \label{eq:v77-lacunary}
\end{equation}
(the intermediate scales of the lacunary shape of
Assessment~\ref{ass:v76-shapes}).  Fix \(0<\eta<S\) and the slabs
\(I_j:=(T_*-S\rho_j^2,\,T_*-\eta\rho_j^2)\).

\begin{theorem}[Satellites]
\label{thm:v77-satellite}
Modulo (E1)--(E7), for every \(R>0\) and \(\varepsilon>0\),
\begin{align}
 \int_{I_j}\int_{B_{R\rho_j}(x_0)}|\nabla u|^2
 &=o(\rho_j),
 \label{eq:v77-near}\\
 \int_{I_j}\int_{\{\operatorname{dist}(\cdot,\Sigma)\ge\varepsilon\}}|\nabla u|^2
 &\le C_\varepsilon\,\rho_j^2,
 \label{eq:v77-regular}\\
 \int_{I_j}\int_{N_\varepsilon(\Sigma)\setminus B_{R\rho_j}(x_0)}|\nabla u|^2
 &\ge2c_L^2\nu^{3/2}(S^{1/2}-\eta^{1/2})\,\rho_j-o(\rho_j).
 \label{eq:v77-satellite}
\end{align}
If moreover the gradient bound \eqref{eq:v46-Linfty-type-I} is used,
then with \(c_1:=c_L^2\nu^{3/2}(S^{1/2}-\eta^{1/2})\) and the satellite
region \(\Omega_j(t):=\{x\in N_\varepsilon(\Sigma)\setminus B_{R\rho_j}(x_0):
|\nabla u(x,t)|^2\ge c_1/(4(S-\eta)|\Torus^3|\rho_j)\}\),
\begin{equation}
 \frac1{|I_j|}\int_{I_j}|\Omega_j(t)|\,\dd t
 \ \ge\ \frac{3c_1\eta^2}{4(S-\eta)(C_\infty')^2}\,\rho_j^3-o(\rho_j^3):
 \label{eq:v77-volume}
\end{equation}
the satellites have time-averaged volume of order \(\rho_j^3\) on which
the gradient is at least of order \(\rho_j^{-1/2}\), while it is at most
of order \(\rho_j^{-2}\) everywhere.
\end{theorem}

\begin{proof}
\eqref{eq:v77-near} is Theorem~\ref{thm:v53-per-scale} in its second
formulation: \(\mathcal E_{x_0}(\rho_j;R')\to0\) for every \(R'\) by
\eqref{eq:v77-lacunary}, hence \(\mathcal D_{x_0}(\rho_j;R,S,\eta)\to0\).
\eqref{eq:v77-regular}: by Theorem~\ref{thm:v44-localization}(ii),
\(\nabla u\) is bounded on \(\{\operatorname{dist}(\cdot,\Sigma)\ge\varepsilon\}
\times[T_*-\delta/2,T_*)\), so the integral is at most a constant times
\(|I_j|=(S-\eta)\rho_j^2\).  \eqref{eq:v77-satellite}: the slab
dissipation is at least \(\int_{I_j}c_L^2\nu^{3/2}(T_*-t)^{-1/2}\dd t
=2c_L^2\nu^{3/2}(S^{1/2}-\eta^{1/2})\rho_j\) by Lemma~\ref{lem:v32-leray-rate};
subtract \eqref{eq:v77-near} and \eqref{eq:v77-regular}.  For \eqref{eq:v77-volume}, split the satellite dissipation rate at time
\(t\in I_j\) into the parts over \(\Omega_j(t)\) and over its complement
in the satellite region.  On the complement, \(|\nabla u|^2\) is below the
threshold \(c_1/(4(S-\eta)|\Torus^3|\rho_j)\), so that part of the rate is
at most \(c_1/(4(S-\eta)\rho_j)\), and its integral over \(I_j\) is at
most \(c_1\rho_j/4\).  On \(\Omega_j(t)\), \eqref{eq:v46-Linfty-type-I}
gives \(|\nabla u|^2\le(C_\infty')^2(T_*-t)^{-2}\le(C_\infty')^2\eta^{-2}\rho_j^{-4}\),
so that part of the rate is at most \((C_\infty')^2\eta^{-2}\rho_j^{-4}|\Omega_j(t)|\).
By \eqref{eq:v77-satellite} the total over \(I_j\) is at least
\(2c_1\rho_j-o(\rho_j)\), hence
\((C_\infty')^2\eta^{-2}\rho_j^{-4}\int_{I_j}|\Omega_j(t)|\dd t
\ge\tfrac74c_1\rho_j-o(\rho_j)\ge\tfrac34c_1\rho_j-o(\rho_j)\); divide by
\(|I_j|=(S-\eta)\rho_j^2\).
\end{proof}

\begin{corollary}[One-point singular set]
\label{cor:v77-one-point}
If \(\Sigma=\{x_0\}\) and the terminal datum is lacunary at \(x_0\) in
the sense of \eqref{eq:v77-lacunary} along some \(\rho_j\to0\), then for
every \(R,\varepsilon\) and large \(j\) there is a satellite region in
\(B_\varepsilon(x_0)\setminus B_{R\rho_j}(x_0)\) carrying dissipation of
order \(\rho_j\) during \(I_j\); the satellites converge to \(x_0\) while
their rescaled distance \(|x-x_0|/\rho_j\) tends to infinity.
\end{corollary}

\begin{proof}
Theorem~\ref{thm:v77-satellite} with \(N_\varepsilon(\Sigma)=B_\varepsilon(x_0)\).
\end{proof}

\begin{assessment}[Satellites and the concentration scales]
\label{ass:v77-satellites}
The satellite regions at scale \(\rho_j\) sit at distance \(d_j\) with
\(d_j/\rho_j\to\infty\) and \(d_j\to0\).  Their own parabolic scale is
\(\rho_j\), their natural time is \(\rho_j^2\), and by the time \(T_*\)
they must have regularized, since their centres are regular points; a
structure that dissipates at the maximal type-I rate for a parabolic time
and then regularizes is not excluded by anything proved.  Whether the
satellites are the dissipation of a previous concentration scale
\(r_{k}\) with \(r_k\gg\rho_j\), seen at rescaled distance
\(d_j/\rho_j\), or new structures, is not decided; the lacunary shape
therefore remains admissible, but it is now a picture with concentric
generations of satellites rather than an isolated point.
\end{assessment}

% =============================================================================
\section{Integrated V77 conclusion and next acquisition order}
\label{sec:v77-conclusion}
% =============================================================================

\begin{theorem}[What V77 proves]
\label{thm:v77-conclusion}
Relative to V1--V76, modulo (E1)--(E7): the satellite theorem
\eqref{eq:v77-near}--\eqref{eq:v77-satellite} with its volume statement,
and the one-point corollary.  The full Navier--Stokes
stability/global-regularity theorem remains unproved.
\end{theorem}

\begin{proof}
The cited results.
\end{proof}

\begin{programme}[V78 acquisition order]
\label{prog:v78}
\begin{enumerate}[label=\textup{(AC\arabic*)},leftmargin=4.8em]
\item Satellites as previous generations.  Determine whether the
      dissipation at intermediate scale \(\rho_j\), located at distance
      \(d_j\gg\rho_j\), can be the tail of the concentration event at the
      previous scale \(r_k\gg\rho_j\); the concentration cylinder
      \(Q_{r_k}(x_0)\) contains the satellite region if \(d_j\le r_k\),
      and its dissipation is of order \(r_k\), so the slab \(I_j\) of
      length \(\rho_j^2\ll r_k^2\) would receive a fraction
      \(\asymp\rho_j/r_k\) of it if the dissipation were spread in time,
      which is consistent.  Record the exact compatibility condition.
\item The next compulsory companion audit is V80.
\end{enumerate}
\end{programme}

\begin{assessment}[Position after V77]
\label{ass:v77-position}
A lacunary singular point is forced to have satellites at every
intermediate scale.  No full stability or global-regularity claim is
made.
\end{assessment}

% =============================================================================
\section*{Version V77 append-only record}
\addcontentsline{toc}{section}{Version V77 append-only record}
% =============================================================================

The complete V76 mathematical source was retained before this continuation.
Its SHA--256 digest was
\begin{center}\scriptsize\ttfamily
d234321ad3e8d14fc6423b30828a14c4b72235d537b4a5b879fcc8ceee275609.
\end{center}
V77 proved that a lacunary terminal datum at a singular point forces
satellite dissipation regions at every intermediate scale, converging to
the singular set at rescaled distance tending to infinity.  No full
regularity claim is made.

% =============================================================================
\section{V78: satellites are active at their own scale}
\label{sec:v78-entry}
% =============================================================================

Item (AC1) of Programme~\ref{prog:v78} asked whether the satellites of
V77 could be the remnants of the concentration event at the preceding
concentration scale.  The question is not well posed: under a type-I
rate the dissipation rate at time \(T_*-\rho^2\) is at most a constant
times \(\rho^{-1}\) whatever structure carries it, and the two-sided
slab bound already fixes the satellite dissipation in the slab at order
\(\rho_j\); there is no separate accounting of ``old'' and ``new''
dissipation.  What is rigorous is that the satellite region at scale
\(\rho_j\) dissipates at the maximal type-I rate, up to constants, and
has spatial extent at least a constant times \(\rho_j\), so it is a
structure of the scale \(\rho_j\) itself, located at rescaled distance
tending to infinity from the singular point and regularizing before the
singular time.  No global regularity claim is made.  The next compulsory
companion audit is V80.

% =============================================================================
\section{Extent and rate of the satellites}
\label{sec:v78-extent}
% =============================================================================

\begin{proposition}[Satellites have the extent and the rate of their scale]
\label{prop:v78-extent-rate}
Under the hypotheses of Theorem~\ref{thm:v77-satellite}, with
\(\Omega_j(t)\) as in \eqref{eq:v77-volume}:
\begin{enumerate}[label=\textup{(\roman*)},leftmargin=3.7em]
\item for a set of times \(t\in I_j\) of measure at least
      \(\kappa|I_j|\), \(|\Omega_j(t)|\ge c\rho_j^3\), with \(c,\kappa>0\)
      depending only on \(S,\eta,C_\infty',C_I,c_L,\nu\); hence
      \(\operatorname{diam}\Omega_j(t)\ge c'\rho_j\) for those times;
\item the satellite dissipation rate
      \(\int_{N_\varepsilon(\Sigma)\setminus B_{R\rho_j}(x_0)}|\nabla u(t)|^2\)
      averaged over \(I_j\) is at least
      \(\bigl(2c_1-o(1)\bigr)/((S-\eta)\rho_j)\) and at most
      \(C_I^2\nu^{3/2}(\eta\rho_j^2)^{-1/2}\), that is, comparable to the
      global type-I rate \(\rho_j^{-1}\) from both sides;
\item every point of \(\Omega_j(t)\), for those times, lies at distance
      at least \(R\rho_j\) from \(x_0\) and within \(\varepsilon\) of
      \(\Sigma\); if \(\Sigma=\{x_0\}\), the satellites lie in
      \(B_\varepsilon(x_0)\setminus B_{R\rho_j}(x_0)\) and every point of
      them is a regular point of \(u\).
\end{enumerate}
\end{proposition}

\begin{proof}
(i) Let \(G:=\{t\in I_j:|\Omega_j(t)|\ge c\rho_j^3\}\).  For
\(t\in I_j\setminus G\), the satellite dissipation rate is at most
\((C_\infty')^2\eta^{-2}\rho_j^{-4}\cdot c\rho_j^3\) on \(\Omega_j(t)\) plus
\(c_1/(4(S-\eta)\rho_j)\) outside it, so its integral over
\(I_j\setminus G\) is at most \(\bigl((C_\infty')^2\eta^{-2}(S-\eta)c
+c_1/4\bigr)\rho_j\), which is below \(c_1\rho_j\) once \(c\) is small.
For \(t\in G\) the rate is at most \(C_G\rho_j^{-1}\) with
\(C_G:=C_I^2\nu^{3/2}\eta^{-1/2}\), by (ii).  Since the total is at least
\(2c_1\rho_j-o(\rho_j)\) by \eqref{eq:v77-satellite},
\(C_G\rho_j^{-1}|G|\ge c_1\rho_j-o(\rho_j)\), so
\(|G|\ge(c_1/(2C_G))\rho_j^2=\kappa|I_j|\) with
\(\kappa=c_1/(2C_G(S-\eta))\) for large \(j\).  The diameter bound is the
isodiametric inequality.  (ii) The lower bound is \eqref{eq:v77-satellite} divided by
\(|I_j|\); the upper bound is the type-I rate at \(t\le T_*-\eta\rho_j^2\).
(iii) is the definition of the satellite region and, for
\(\Sigma=\{x_0\}\), the fact that all points other than \(x_0\) are
regular.
\end{proof}

\begin{assessment}[The satellite picture]
\label{ass:v78-picture}
At each intermediate scale of a lacunary singular point there is a
structure of that scale, dissipating at the type-I rate, located at
rescaled distance tending to infinity from the point and made of regular
points which must regularize within their own parabolic time.  This is
consistent with every proved bound; excluding it would require a
statement that a structure dissipating at the maximal type-I rate for a
parabolic time cannot regularize, which is false in general (small
perturbations of the trivial solution can dissipate at any rate briefly).
The lacunary shape therefore remains admissible.
\end{assessment}

% =============================================================================
\section{Integrated V78 conclusion and next acquisition order}
\label{sec:v78-conclusion}
% =============================================================================

\begin{theorem}[What V78 proves]
\label{thm:v78-conclusion}
Relative to V1--V77, modulo (E1)--(E7): Proposition~\ref{prop:v78-extent-rate}.
The full Navier--Stokes stability/global-regularity theorem remains
unproved.
\end{theorem}

\begin{proof}
The cited result.
\end{proof}

\begin{programme}[V79 acquisition order]
\label{prog:v79}
\begin{enumerate}[label=\textup{(AC\arabic*)},leftmargin=4.8em]
\item Determine whether the satellite regions at consecutive
      intermediate scales can be nested, that is, whether the satellite at
      scale \(\rho_{j+1}\) can sit inside the satellite at scale
      \(\rho_j\), which would make the satellites themselves lacunary
      concentration structures at regular points, and whether the
      per-scale implication of V53 applied at a satellite centre
      constrains this.
\item Prepare the compulsory V80 companion audit.
\end{enumerate}
\end{programme}

\begin{assessment}[Position after V78]
\label{ass:v78-position}
Satellites are structures of their own scale.  No full stability or
global-regularity claim is made.
\end{assessment}

% =============================================================================
\section*{Version V78 append-only record}
\addcontentsline{toc}{section}{Version V78 append-only record}
% =============================================================================

The complete V77 mathematical source was retained before this continuation.
Its SHA--256 digest was
\begin{center}\scriptsize\ttfamily
a4098fb9f7f92a4da00a7be68b6491dd511a11289fdc8d49a04789726caa4a73.
\end{center}
V78 proved that the satellites have the extent and the dissipation rate
of their own scale and recorded that the lacunary shape remains
admissible.  No full regularity claim is made.

% =============================================================================
\section{V79: compact satellites force terminal bumps}
\label{sec:v79-entry}
% =============================================================================

Item (AC1) of Programme~\ref{prog:v79} is answered in a stronger form.
The per-scale implication of V53, which was stated at a singular point,
holds with moving centres at arbitrary points, because its proof used
only the global type-I bounds, compactness, nontriviality from a
dissipation lower bound, and backward uniqueness.  Applied at the centres
of the satellites of V77, it shows that a satellite carried by boundedly
many balls of its own scale forces the terminal datum to have positive
critical density around the satellite centre at that scale, hence values
of order the inverse scale there; since the satellite sits at a distance
much larger than its scale from the singular point, these values exceed
the critical profile at that distance.  Satellites spread over regions
much larger than their scale escape this conclusion.  No global
regularity claim is made.  The next compulsory companion audit is V80.

% =============================================================================
\section{The per-scale implication with moving centres}
\label{sec:v79-moving-centres}
% =============================================================================

\begin{lemma}[Moving centres]
\label{lem:v79-moving-centres}
Assume \eqref{eq:v32-type-I} and (E1)--(E7).  Let \(x_j\in\Torus^3\) and
\(\rho_j\downarrow0\) satisfy, for some \(R,S,\eta\),
\begin{equation}
 \liminf_{j\to\infty}\frac1{\rho_j}\iint_{B_{R\rho_j}(x_j)\times(T_*-S\rho_j^2,\,T_*-\eta\rho_j^2)}|\nabla u|^2>0 .
 \label{eq:v79-active-moving}
\end{equation}
Then there is \(R'\) with
\(\liminf_j\rho_j^{-1}\int_{B_{R'\rho_j}(x_j)}|u(T_*)|^2>0\).
\end{lemma}

\begin{proof}
Identical to Theorem~\ref{thm:v53-per-scale} with \(x\) replaced by
\(x_j\): the uniform bounds of Proposition~\ref{prop:v41-uniform-bounds}
and the compactness of Theorem~\ref{thm:v41-ancient-limit} do not involve
the centre, nontriviality of the limit is \eqref{eq:v79-active-moving},
and the vanishing-trace contradiction is
Proposition~\ref{prop:v46-vanishing} with Theorem~\ref{thm:v46-non-L3},
whose proofs are centre-independent.
\end{proof}

% =============================================================================
\section{Compact satellites}
\label{sec:v79-compact}
% =============================================================================

Under the hypotheses of Theorem~\ref{thm:v77-satellite}, call the
satellites \emph{compact} if there are \(N\) and \(R_1\), independent of
\(j\), and centres \(x_j^{(1)},\dots,x_j^{(N)}\) such that a fixed
fraction of the satellite dissipation of the slab \(I_j\) lies in
\(\bigcup_{i}B_{R_1\rho_j}(x_j^{(i)})\times I_j\).

\begin{theorem}[Compact satellites force terminal bumps]
\label{thm:v79-bumps}
If the satellites are compact, then for some \(i\), some \(R'\), and
\(c>0\), along a subsequence,
\begin{equation}
 \frac1{\rho_j}\int_{B_{R'\rho_j}(x_j^{(i)})}|u(T_*)|^2\ \ge\ c,
 \qquad\text{hence}\qquad
 \sup_{B_{R'\rho_j}(x_j^{(i)})}|u(T_*)|\ \ge\ \frac{c'}{\rho_j},
 \label{eq:v79-bumps}
\end{equation}
while \(|x_j^{(i)}-x_0|\ge R\rho_j\) with \(R\) arbitrary, so that
\(|x_j^{(i)}-x_0|/\rho_j\to\infty\) along the full construction.  In
particular the terminal datum has, near each such satellite centre,
values larger by the unbounded factor \(|x_j^{(i)}-x_0|/\rho_j\) than the
critical profile \(|x-x_0|^{-1}\) at that distance.
\end{theorem}

\begin{proof}
A fixed fraction of at least \(c_1\rho_j-o(\rho_j)\) lies in \(N\) balls,
so one ball \(B_{R_1\rho_j}(x_j^{(i)})\) carries at least
\((c_1/N)\rho_j-o(\rho_j)\) along a subsequence, which is
\eqref{eq:v79-active-moving} with \(R=R_1\).  Lemma~\ref{lem:v79-moving-centres}
gives the density bound; \(\int_{B_{R'\rho_j}}|u(T_*)|^2\le|B_{R'\rho_j}|
\sup|u(T_*)|^2\) gives the supremum bound, \(u(T_*)\) being bounded near
the regular point \(x_j^{(i)}\) by Theorem~\ref{thm:v44-localization}(ii).
The distance statement is the definition of the satellite region.
\end{proof}

\begin{assessment}[Compact versus spread satellites]
\label{ass:v79-compact-spread}
Compact satellites make the terminal datum spiky: at distance \(d\) from
the singular point it must reach size \(\rho^{-1}\gg d^{-1}\) on balls of
radius \(\rho\), with \(\rho\) an intermediate scale.  This is compatible
with the Morrey bound, since a bump of height \(\rho^{-1}\) on a ball of
radius \(\rho\) has energy \(\rho\), and with smoothness off the singular
point, since the bumps sit at regular points at positive distance from
it.  Spread satellites, whose dissipation is distributed over regions of
size much larger than \(\rho_j\), for instance a shell of radius
\(d_j\) and thickness \(\rho_j\), are not constrained by
Lemma~\ref{lem:v79-moving-centres}, because no ball of radius
\(R_1\rho_j\) carries a fixed fraction of their dissipation.  Neither kind
is excluded.
\end{assessment}

% =============================================================================
\section{Integrated V79 conclusion and next acquisition order}
\label{sec:v79-conclusion}
% =============================================================================

\begin{theorem}[What V79 proves]
\label{thm:v79-conclusion}
Relative to V1--V78, modulo (E1)--(E7): the per-scale implication with
moving centres (Lemma~\ref{lem:v79-moving-centres}) and the terminal
bumps forced by compact satellites (Theorem~\ref{thm:v79-bumps}).  The
full Navier--Stokes stability/global-regularity theorem remains unproved.
\end{theorem}

\begin{proof}
The cited results.
\end{proof}

\begin{programme}[V80 acquisition order, with compulsory companion audit]
\label{prog:v80}
\begin{enumerate}[label=\textup{(AC\arabic*)},leftmargin=4.8em]
\item Perform the compulsory review of the current SM and YM notes;
      record digests; import nothing numerical.
\item Spread satellites.  Determine whether the Morrey bound and the
      \(L^\infty\) type-I bound constrain a shell-shaped satellite of
      radius \(d_j\) and thickness \(\rho_j\): its volume is
      \(\asymp d_j^2\rho_j\), and dissipation \(\asymp\rho_j\) over the
      slab of length \(\rho_j^2\) requires a mean gradient square
      \(\asymp d_j^{-2}\rho_j^{-2}\) on the shell, which is below the
      type-I ceiling \(\rho_j^{-4}\) by the factor \((\rho_j/d_j)^2\);
      record the exact scaling and test it against the Morrey bound on
      \(B_{2d_j}(x_0)\).
\item Consolidate V76--V79 into a description of the lacunary shape.
\end{enumerate}
\end{programme}

\begin{assessment}[Position after V79]
\label{ass:v79-position}
Compact satellites force spiky terminal data; spread satellites are
unconstrained.  No full stability or global-regularity claim is made.
\end{assessment}

% =============================================================================
\section*{Version V79 append-only record}
\addcontentsline{toc}{section}{Version V79 append-only record}
% =============================================================================

The complete V78 mathematical source was retained before this continuation.
Its SHA--256 digest was
\begin{center}\scriptsize\ttfamily
e7534ab35eb4848adb0c4f8c00a9ab651fdfbcd3cd973a0153deb39262402255.
\end{center}
V79 extended the per-scale implication to moving centres and proved that
compact satellites force terminal bumps of size the inverse scale at the
satellite centres.  No full regularity claim is made.

% =============================================================================
\section{V80: compulsory companion audit, spread satellites, and the lacunary
shape}
\label{sec:v80-entry}
% =============================================================================

This is a fiftieth-version continuation.  The compulsory review of the
two companion programmes is performed first.  Item (AC2) of
Programme~\ref{prog:v80} is then settled: a shell-shaped satellite of
radius \(d_j\) and thickness \(\rho_j\) around the singular point is
compatible with the Morrey bound, the \(L^\infty\) type-I bound, and the
per-scale implication, so spread satellites are admissible.  Finally
V76--V79 are consolidated into a description of the lacunary shape of a
type-I singularity: concentric generations of satellites at every
intermediate scale, either compact and accompanied by terminal bumps, or
spread over shells.  No global regularity claim is made.  The next
compulsory companion audit is V85.

% =============================================================================
\section{Compulsory V80 review of the current companion notes}
\label{sec:v80-companion-audit}
% =============================================================================

The current companion files in \texttt{latex\_notes} were inspected
read-only.  The frozen checkpoints are
\begin{align}
 &\texttt{Renewal\_SM\_Einstein\_Continuum\_Research\_Notes\_V76.tex},
 \notag\\[-2pt]
 &\qquad
 \texttt{e346d07ce63a7d8275c7f693a0171d73159399611c2c8688a91ae425a2696325},
 \label{eq:v80-SM-checkpoint}\\
 &\texttt{Renewal\_Yang\_Mills\_Mass\_Gap\_Research\_Notes\_V54.tex},
 \notag\\[-2pt]
 &\qquad
 \texttt{cc98f8c79e51e5e90c2a31ebaa734f0c2c852fc4fff433b1d3dbac6265229da6}.
 \label{eq:v80-YM-checkpoint}
\end{align}

\emph{SM V76.}  Real-lapse group averaging gives a positive
delta-constraint inner product, identified with the coarea measure for
regular finite-dimensional constraints; the homogeneous CMC--Higgs shell
has an explicit branch weight; a quartic Higgs tail gives a radial measure
finite only in one amplitude dimension; simple York branch meetings are
integrable while critical meetings need a singular-stratum analysis.

\emph{YM V54.}  The V51 branch is located relative to the physical
total-spin spectrum, small-\(s\) bounds for the expectation and the
multiplier are proved, a strong two-leg singlet bound and a rare-mixture
logarithm lemma are established, and Bernoulli source zeros are located at
distance of order \(1/s\).

\begin{theorem}[V80 companion-audit decision]
\label{thm:v80-companion-decision}
No SM V76 coarea or branch constant and no YM V54 spectral bound is
imported.  Neither bears on the three gates of Theorem~\ref{thm:v70-gates}.
\end{theorem}

\begin{proof}
The hypotheses of every cited result are constraint measures or
lattice-gauge spectral objects; none is a Navier--Stokes balance.
\end{proof}

% =============================================================================
\section{Spread satellites}
\label{sec:v80-spread}
% =============================================================================

\begin{proposition}[A shell satellite is admissible by scaling]
\label{prop:v80-shell}
Let a satellite at intermediate scale \(\rho_j\) occupy the shell
\(A_j:=\{d_j-\rho_j\le|x-x_0|\le d_j+\rho_j\}\) with \(d_j/\rho_j\to\infty\),
\(d_j\to0\), during the slab \(I_j\).  The requirements are: dissipation
\(\int_{I_j}\int_{A_j}|\nabla u|^2\asymp\rho_j\); the type-I ceiling
\(|\nabla u|\le C_\infty'(\eta\rho_j^2)^{-1}\); the Morrey bound
\(\int_{B_{2d_j}(x_0)}|u(t)|^2\le C_Md_j\); and, if the shell is to carry
the dissipation with the gradient concentrated on it, the per-scale
implication at any centre on the shell with radius \(R\rho_j\), which
sees only the fraction \(\asymp(\rho_j/d_j)^2\) of the shell.  Then the
mean gradient square on the shell over the slab is
\(\asymp\rho_j/(\rho_j^2\cdot d_j^2\rho_j)=d_j^{-2}\rho_j^{-2}\), below the
ceiling by the factor \((\rho_j/d_j)^2\); the corresponding velocity of
size \(\asymp d_j^{-1}\) on the shell has energy
\(\asymp d_j^{-2}\cdot d_j^2\rho_j=\rho_j\ll C_Md_j\); and no ball of
radius \(R\rho_j\) carries a fixed fraction of the dissipation, so
Lemma~\ref{lem:v79-moving-centres} does not apply.  All three
requirements are satisfied, and the shell satellite is admissible.
\end{proposition}

\begin{proof}
Direct computation of the stated orders of magnitude.
\end{proof}

\begin{assessment}[The two satellite types]
\label{ass:v80-types}
A compact satellite is a ball-sized structure of scale \(\rho_j\) with
gradient of order \(\rho_j^{-2}\) and velocity \(\rho_j^{-1}\), forced to
leave a bump of size \(\rho_j^{-1}\) in the terminal datum at its centre.
A shell satellite is a thin shell of radius \(d_j\) with gradient
\(d_j^{-1}\rho_j^{-1}\) and velocity \(d_j^{-1}\), which is exactly the
size of the self-similar profile at distance \(d_j\), and it leaves no
bump.  The second is the natural shape if the singular structure at the
larger scale \(d_j\) has thin internal layers; the first is a genuinely
separate small-scale event.  Nothing proved distinguishes them.
\end{assessment}

% =============================================================================
\section{The lacunary shape, consolidated}
\label{sec:v80-lacunary}
% =============================================================================

\begin{theorem}[Shape of a lacunary type-I singular point, V76--V80]
\label{thm:v80-lacunary}
Let \(x_0\in\Sigma\) under the type-I rate, with concentration scales
\(r_k\) and intermediate scales \(\rho_j\) satisfying
\eqref{eq:v77-lacunary}.  Modulo (E1)--(E8):
\begin{enumerate}[label=\textup{(\arabic*)},leftmargin=3.2em]
\item the terminal datum has positive upper critical density at \(x_0\)
      along \(r_k\), bounded below uniformly over \(\Sigma\), and
      subcritical density along \(\rho_j\) at every fixed rescaled radius;
\item along \(\rho_j\) the rescaled dissipation near \(x_0\) vanishes,
      and the slab dissipation of order \(\rho_j\) sits in satellites
      within every neighbourhood of \(\Sigma\), at rescaled distance
      tending to infinity, with extent at least of order \(\rho_j\) and
      rate of order \(\rho_j^{-1}\);
\item compact satellites force terminal bumps of size \(\rho_j^{-1}\) at
      their centres; shell satellites are admissible without bumps;
\item the ancient limit at \(x_0\) along \(r_k\) has the properties of
      Theorem~\ref{thm:v65-ancient-limit}, and along \(\rho_j\) the
      rescaled limit is zero;
\item nothing proved excludes the lacunary shape, and nothing proved
      excludes the homogeneous shape in which every scale is a
      concentration scale.
\end{enumerate}
\end{theorem}

\begin{proof}
(1) Corollary~\ref{cor:v46-critical-local-energy},
Theorem~\ref{thm:v52-uniform-lower-bound}, and the definition of
intermediate scales.  (2) Theorem~\ref{thm:v77-satellite} and
Proposition~\ref{prop:v78-extent-rate}.  (3) Theorem~\ref{thm:v79-bumps}
and Proposition~\ref{prop:v80-shell}.  (4) Theorem~\ref{thm:v65-ancient-limit}
and Theorem~\ref{thm:v53-per-scale}.  (5) collects the firewalls of
V76--V80.
\end{proof}

% =============================================================================
\section{Integrated V80 conclusion and next acquisition order}
\label{sec:v80-conclusion}
% =============================================================================

\begin{theorem}[What V80 establishes]
\label{thm:v80-conclusion}
Relative to V1--V79: the companion-audit decision
(Theorem~\ref{thm:v80-companion-decision}), the admissibility of shell
satellites (Proposition~\ref{prop:v80-shell}), and the consolidated
lacunary shape (Theorem~\ref{thm:v80-lacunary}).  The full Navier--Stokes
stability/global-regularity theorem remains unproved.
\end{theorem}

\begin{proof}
The cited results.
\end{proof}

\begin{programme}[V81 acquisition order]
\label{prog:v81}
\begin{enumerate}[label=\textup{(AC\arabic*)},leftmargin=4.8em]
\item Homogeneous shape.  If every small scale is a concentration scale
      at \(x_0\), the terminal density is positive at every scale and the
      rescaled limits along every sequence are nontrivial; determine
      whether the family of all rescaled limits is then compact under the
      scaling group, which would give a self-similar element by a
      standard argument, and record exactly which compactness is missing.
\item Distinguishing the shapes.  Determine whether the two-sided slab
      dissipation forces, for a homogeneous shape, a lower bound on the
      dissipation fraction inside \(B_{R\rho}(x_0)\) uniform in \(\rho\),
      which would be the per-scale converse of V53.
\item The next compulsory companion audit is V85.
\end{enumerate}
\end{programme}

\begin{assessment}[Position after V80]
\label{ass:v80-position}
Fifty versions after V30, the type-I singular point has two admissible
shapes, homogeneous and lacunary, each described in detail, and neither
excluded.  No full stability or global-regularity claim is made.
\end{assessment}

% =============================================================================
\section*{Version V80 append-only record}
\addcontentsline{toc}{section}{Version V80 append-only record}
% =============================================================================

The complete V79 mathematical source was retained before this continuation.
Its SHA--256 digest was
\begin{center}\scriptsize\ttfamily
1a6210f24b58249c12f23c8b9df2dabbaffab6847cf4f69fed97db2ca91510d9.
\end{center}
V80 performed the compulsory companion audit, showed shell satellites are
admissible by scaling, and consolidated the lacunary shape of a type-I
singular point.  No full regularity claim is made.

% =============================================================================
\section{V81: the per-scale equivalence and the count of uniformly
homogeneous singular points}
\label{sec:v81-entry}
% =============================================================================

Item (AC2) of Programme~\ref{prog:v81}, the converse of the per-scale
implication of V53, is proved: at any point and along any sequence of
scales, the terminal datum has positive critical density if and only if
the rescaled dissipation does not vanish.  The new direction follows from
the two-sided local-energy continuity of V74, which transports positive
terminal energy at the final time to positive energy at a fixed earlier
rescaled time, where strong local convergence makes the limit
nontrivial.  Both directions are quantitative by compactness.
Consequently the homogeneous shape of a singular point, positive terminal
density at every scale, is exactly dissipation activity at every scale,
and the V43 count applies: singular points that are active at every
scale with a common lower bound \(c\) number at most
\(2C_I^2\nu^{3/2}/c\).  Item (AC1), compactness under the scaling group,
is recorded as open.  No global regularity claim is made.  The next
compulsory companion audit is V85.

% =============================================================================
\section{The equivalence}
\label{sec:v81-equivalence}
% =============================================================================

Throughout, \eqref{eq:v32-type-I} and (E1)--(E7) hold, and
\(\mathcal D_x\), \(\mathcal E_x\) are as in \eqref{eq:v53-quantities}.

\begin{theorem}[Per-scale equivalence]
\label{thm:v81-equivalence}
For every \(x\in\Torus^3\) and every \(\rho_j\downarrow0\), the following
are equivalent:
\begin{enumerate}[label=\textup{(\roman*)},leftmargin=3.7em]
\item \(\liminf_j\mathcal D_x(\rho_j;R,S,\eta)>0\) for some
      \(R,S,\eta\);
\item \(\liminf_j\mathcal E_x(\rho_j;R')>0\) for some \(R'\).
\end{enumerate}
The same holds with moving centres \(x_j\) in place of \(x\).
\end{theorem}

\begin{proof}
(i)\(\Rightarrow\)(ii) is Theorem~\ref{thm:v53-per-scale} and
Lemma~\ref{lem:v79-moving-centres}.  For (ii)\(\Rightarrow\)(i), let
\(U_j\) be the rescalings at \((x,\rho_j)\), with local energies
\(e^{(j)}_{R'}\).  By (ii), \(e^{(j)}_{R'}(0^-)=\frac12\int|T_j|^2\phi_{R'}
\ge\frac12\mathcal E_x(\rho_j;R')\ge c/2\) for large \(j\).  The two-sided
bound of Lemma~\ref{lem:v74-two-sided}, whose proof used only the global
type-I bounds and holds for any centre and scale, gives
\(e^{(j)}_{R'}(s)\ge c/2-C_{R'}(-s)^{1/4}\); fix \(\eta_1>0\) with
\(C_{R'}\eta_1^{1/4}\le c/4\), so that \(e^{(j)}_{R'}(-\eta_1)\ge c/4\) for
large \(j\).  By the compactness of Theorem~\ref{thm:v41-ancient-limit},
along a subsequence \(U_j\to U\) strongly in \(L^2((-S,-\eta);H^1(B_{2R'}))\)
for every \(S>\eta>0\), and at the fixed time \(-\eta_1\) strongly in
\(L^2(B_{2R'})\) by the equicontinuity in time; hence
\(\frac12\int|U(-\eta_1)|^2\phi_{R'}\ge c/4>0\) and \(U\not\equiv0\).  If
(i) failed for all \(R,S,\eta\), then along a further subsequence
\(\iint_{B_R\times(-S,-\eta)}|\nabla U_j|^2\to0\) for all \(R,S,\eta\),
so \(\nabla U=0\) almost everywhere on \((-\infty,0)\times\R^3\), and
\(U(s)\) would be constant in space and in \(L^6(\R^3)\), hence zero, a
contradiction.
\end{proof}

\begin{corollary}[Quantitative forms]
\label{cor:v81-quantitative}
For every \(c>0\) and \(R'\) there are \(\varepsilon>0\) and
\(R,S,\eta\), depending only on \(c,R',C_I,\nu\) and the absolute
constants, such that for all \(x\) and all small \(\rho\),
\begin{equation}
 \mathcal E_x(\rho;R')\ge c\ \Longrightarrow\ \mathcal D_x(\rho;R,S,\eta)\ge\varepsilon ,
 \label{eq:v81-quantitative}
\end{equation}
and conversely, as in Corollary~\ref{cor:v53-quantitative}, small
terminal density at all radii up to \(R'\) forces small dissipation.
\end{corollary}

\begin{proof}
If \eqref{eq:v81-quantitative} failed, there would be \(x_n\),
\(\rho_n\to0\) with \(\mathcal E_{x_n}(\rho_n;R')\ge c\) and
\(\mathcal D_{x_n}(\rho_n;n,n,1/n)\to0\); the proof of
Theorem~\ref{thm:v81-equivalence} with moving centres gives a
contradiction.
\end{proof}

% =============================================================================
\section{Homogeneous points and their count}
\label{sec:v81-count}
% =============================================================================

Call \(x\in\Sigma\) \emph{homogeneous with constant \(c\)} if
\(\mathcal E_x(\rho;1)\ge c\) for all sufficiently small \(\rho\), and
\emph{homogeneous} if it is homogeneous with some constant.

\begin{theorem}[Count of uniformly homogeneous singular points]
\label{thm:v81-count}
Let \(c>0\).  The set \(\Sigma_c\) of singular points homogeneous with
constant \(c\) satisfies
\begin{equation}
 \boxed{\quad
 |\Sigma_c|\ \le\ \frac{2C_I^2\nu^{3/2}}{\varepsilon(c)}\,,
 \quad}
 \label{eq:v81-count}
\end{equation}
where \(\varepsilon(c)\) is the constant of \eqref{eq:v81-quantitative}
with \(R'=1\) and the parameters \(R,S,\eta\) it provides, after
adjusting the cylinder to \(Q_{r}(x,T_*)\) with \(r=S^{1/2}\rho\).
Consequently the set of all homogeneous singular points is countable,
and it is finite if the constants of homogeneity are bounded below.
\end{theorem}

\begin{proof}
For \(x\in\Sigma_c\) and small \(\rho\), \eqref{eq:v81-quantitative}
gives dissipation at least \(\varepsilon\rho\) in
\(B_{R\rho}(x)\times(T_*-S\rho^2,T_*-\eta\rho^2)\subset Q_{r}(x,T_*)\)
with \(r:=\max\{R,S^{1/2}\}\rho\), so \(f_x(r)\ge\varepsilon\rho/r
=\varepsilon/\max\{R,S^{1/2}\}=:\varepsilon'\) for all small \(r\).  For
\(2r\)-separated points of \(\Sigma_c\) the cylinders are disjoint, and
Lemma~\ref{lem:v43-few-active} with \(\varepsilon'\) in place of
\(\varepsilon_1/2\) bounds their number by \(2C_I^2\nu^{3/2}/\varepsilon'\);
since this holds for all small \(r\), it bounds the number of points of
any finite subset of \(\Sigma_c\), hence \(|\Sigma_c|\).  Countability of
\(\bigcup_{c>0}\Sigma_c=\bigcup_n\Sigma_{1/n}\) follows.
\end{proof}

\begin{assessment}[Homogeneous versus lacunary, after V81]
\label{ass:v81-shapes}
A singular point is either homogeneous, in which case it is active at
every scale and belongs to a countable set, finite for each constant of
homogeneity, or lacunary, in which case it has satellites at every
intermediate scale (V77--V80).  The singular set is therefore the union
of a countable homogeneous part and a lacunary part; the lacunary part
is not counted by anything proved.  An uncountable singular set, if one
exists, consists of lacunary points up to a countable set.
\end{assessment}

\begin{firewall}[Item (AC1)]
\label{fw:v81-AC1}
Compactness of the family of rescaled limits under the continuous
scaling group was not obtained; the equivalence proved here makes every
rescaled limit at a homogeneous point nontrivial but does not relate
limits at different scales.  Recorded as open.
\end{firewall}

% =============================================================================
\section{Integrated V81 conclusion and next acquisition order}
\label{sec:v81-conclusion}
% =============================================================================

\begin{theorem}[What V81 proves]
\label{thm:v81-conclusion}
Relative to V1--V80, modulo (E1)--(E7): the per-scale equivalence
(Theorem~\ref{thm:v81-equivalence}), its quantitative forms
(Corollary~\ref{cor:v81-quantitative}), and the count of uniformly
homogeneous singular points (Theorem~\ref{thm:v81-count}).  The full
Navier--Stokes stability/global-regularity theorem remains unproved.
\end{theorem}

\begin{proof}
The cited results.
\end{proof}

\begin{programme}[V82 acquisition order]
\label{prog:v82}
\begin{enumerate}[label=\textup{(AC\arabic*)},leftmargin=4.8em]
\item Lacunary points.  Determine whether an uncountable set of lacunary
      singular points is compatible with the satellite theorem: each
      lacunary point needs satellites at its intermediate scales, located
      within every neighbourhood of \(\Sigma\); test whether the total
      dissipation budget of a slab, of order the scale, can serve the
      satellites of uncountably many points at once.
\item Homogeneous points and self-similarity.  At a homogeneous point
      every rescaled limit is nontrivial; determine whether the
      concentration function \(d_U(\lambda)\) of the ancient limit is
      then bounded below at all scales, which would place the ancient
      limit in the concentration branches (L1) and (S1) simultaneously.
\item The next compulsory companion audit is V85.
\end{enumerate}
\end{programme}

\begin{assessment}[Position after V81]
\label{ass:v81-position}
Terminal density and dissipation activity are equivalent scale by scale,
and uniformly homogeneous singular points are finite in number.  No full
stability or global-regularity claim is made.
\end{assessment}

% =============================================================================
\section*{Version V81 append-only record}
\addcontentsline{toc}{section}{Version V81 append-only record}
% =============================================================================

The complete V80 mathematical source was retained before this continuation.
Its SHA--256 digest was
\begin{center}\scriptsize\ttfamily
b08c6f09693eddf4b003bbb07967df3d5903bd81e263aa4b8c1024ba3f889853.
\end{center}
V81 proved the per-scale equivalence between terminal density and
dissipation activity, its quantitative forms, and the finiteness of the
set of singular points homogeneous with a given constant.  No full
regularity claim is made.

% =============================================================================
\section{V82: homogeneous points give uniformly concentrated ancient limits}
\label{sec:v82-entry}
% =============================================================================

Item (AC2) of Programme~\ref{prog:v82} is settled: at a homogeneous
singular point the concentration function of the ancient limit is
bounded below at every scale, so the ancient limit lies in both
concentration branches of V49, is singular at the origin at its final
time, and all of its centred rescalings, at every scale, are nontrivial
with a common lower bound.  Item (AC1) is settled negatively: the
satellites of distinct lacunary points at a common scale are the same
dissipation, so an uncountable lacunary set is compatible with the slab
budget.  No global regularity claim is made.  The next compulsory
companion audit is V85.

% =============================================================================
\section{Uniform concentration}
\label{sec:v82-uniform}
% =============================================================================

\begin{theorem}[The ancient limit at a homogeneous point is uniformly
concentrated]
\label{thm:v82-uniform}
Let \(x_0\in\Sigma\) be homogeneous with constant \(c\), with
\(\varepsilon'>0\) the dissipation constant given by
Corollary~\ref{cor:v81-quantitative} and Theorem~\ref{thm:v81-count}, so
that \(f_{x_0}(r)\ge\varepsilon'\) for all small \(r\).  Let \(U\) be the
ancient limit at \(x_0\) along any sequence \(r_k\to0\).  Then, modulo
(E1)--(E7),
\begin{equation}
 \boxed{\quad
 d_U(\lambda)=\frac1\lambda\iint_{B_\lambda(0)\times(-\lambda^2,0)}|\nabla U|^2
 \ \ge\ \frac{\varepsilon'}2
 \qquad\text{for every }\lambda\in(0,\infty).
 \quad}
 \label{eq:v82-uniform}
\end{equation}
Consequently: \(U\) is in branch (L1) and in branch (S1) of
Theorems~\ref{thm:v49-large-scale}--\ref{thm:v49-small-scale}; \((0,0)\)
is a singular point of \(U\), so \(U\) is in case (B) and
\(U_0\notin L^3(B_\rho(0))\) for every \(\rho\); and every centred
rescaling \(U^\lambda\) has unit-scale dissipation at least
\(\varepsilon'/2\), so every centred rescaling limit of \(U\), at any
sequence of scales, is nontrivial.
\end{theorem}

\begin{proof}
For the rescaled fields \(U_k\) at \((x_0,r_k)\), the change of variables
gives \(d_{U_k}(\lambda)=f_{x_0}(\lambda r_k)\ge\varepsilon'\) for all
\(\lambda\) and all large \(k\).  Split the cylinder at \(s=-\eta\lambda^2\):
the part with \(s\in(-\eta\lambda^2,0)\) carries at most
\(\lambda^{-1}\int_{-\eta\lambda^2}^0Y_k\le2C_I^2\nu^{3/2}\eta^{1/2}\),
which is at most \(\varepsilon'/2\) for \(\eta\) small; on the part
\(B_\lambda\times(-\lambda^2,-\eta\lambda^2)\) the convergence
\(U_k\to U\) is strong in \(L^2_sH^1_y\), so
\(\lambda^{-1}\iint_{B_\lambda\times(-\lambda^2,-\eta\lambda^2)}|\nabla U|^2
=\lim_k(\cdot)\ge\varepsilon'/2\), which gives \eqref{eq:v82-uniform}.
Branches (L1) and (S1) are the statements \(\limsup_{\lambda\to\infty}
d_U^*\ge\limsup d_U>0\) and \(\limsup_{\lambda\to0}d_U^*>0\).  The origin
is singular: if \(U\) were bounded near \((0,0)\), then by (E5) its
gradient would be bounded on a cylinder \(Q_{\lambda_0}(0,0)\), and
\(\lambda^{-1}\iint_{Q_\lambda(0,0)}|\nabla U|^2\le C\lambda^{-1}\lambda^5\to0\)
as \(\lambda\to0\), contradicting \eqref{eq:v82-uniform}; the
non-\(L^3\) statement is Theorem~\ref{thm:v57-singular-points-U}.  The last assertion is
Lemma~\ref{lem:v49-invariance}: the unit-scale dissipation of \(U^\lambda\)
is \(d_U(\lambda)\), and the compactness argument of
Theorem~\ref{thm:v49-large-scale} applied to any sequence of scales gives
a limit with dissipation at least \(\varepsilon'/2\) on the unit cylinder
minus its final part.
\end{proof}

\begin{assessment}[The homogeneous ancient limit]
\label{ass:v82-homogeneous-limit}
At a homogeneous point the ancient limit reproduces the homogeneity: it
is singular at the origin, concentrated at every scale, and each of its
rescalings is again such an element.  This is the structure one would
expect from a self-similar singularity, and it is exactly what the
missing compactness under the scaling group (Firewall~\ref{fw:v81-AC1})
would turn into a self-similar element.  With the uniform lower bound
\eqref{eq:v82-uniform} the obstruction is no longer possible triviality of
limits but only the absence of a fixed point.
\end{assessment}

% =============================================================================
\section{Uncountable lacunary sets and the slab budget}
\label{sec:v82-lacunary}
% =============================================================================

\begin{proposition}[Satellites are shared]
\label{prop:v82-shared}
Let \(L\subset\Sigma\) be any set of lacunary points and \(\rho\) a
common intermediate scale for all of them, in the sense of
\eqref{eq:v77-lacunary} along a common sequence.  Then the satellite
theorem applied at each \(x\in L\) locates the slab dissipation, up to
\(o(\rho)\), outside \(\bigcup_{x\in L}B_{R\rho}(x)\) and inside
\(N_\varepsilon(\Sigma)\); it is one and the same dissipation of order
\(\rho\) for all \(x\in L\).  In particular no upper bound on the
cardinality of \(L\) follows from the slab budget.
\end{proposition}

\begin{proof}
Theorem~\ref{thm:v77-satellite} gives, for each \(x\in L\), that the
dissipation in \(B_{R\rho}(x)\) is \(o(\rho)\); if \(L\) is finite the
union of the balls contributes \(|L|\,o(\rho)\), and for infinite \(L\)
the statement is applied to finite subsets with the same \(o(\rho)\)
uniform in \(x\in L\) only if the density smallness in
\eqref{eq:v77-lacunary} is uniform on \(L\), which is assumed by the
common-sequence hypothesis; the remaining dissipation is the satellite
dissipation, common to all points.
\end{proof}

\begin{firewall}[Where a count could still come from]
\label{fw:v82-count}
A bound on the number of lacunary points would need, at each scale, a
lower bound on the dissipation that each point individually forces
somewhere, disjointly from the others.  The satellite theorem forces
dissipation away from the point, not attributed to it.  Only the
concentration scales of a point force dissipation at the point itself,
and lacunary points have concentration scales that may be disjoint from
those of other points; a count would need common scales, which is the
same obstruction as in Firewall~\ref{fw:v44-sparse-activity}.
\end{firewall}

% =============================================================================
\section{Integrated V82 conclusion and next acquisition order}
\label{sec:v82-conclusion}
% =============================================================================

\begin{theorem}[What V82 proves]
\label{thm:v82-conclusion}
Relative to V1--V81, modulo (E1)--(E7): the uniform concentration of the
ancient limit at a homogeneous point (Theorem~\ref{thm:v82-uniform}) and
the sharing of satellites (Proposition~\ref{prop:v82-shared}).  The full
Navier--Stokes stability/global-regularity theorem remains unproved.
\end{theorem}

\begin{proof}
The cited results.
\end{proof}

\begin{programme}[V83 acquisition order]
\label{prog:v83}
\begin{enumerate}[label=\textup{(AC\arabic*)},leftmargin=4.8em]
\item Fixed point at a homogeneous point.  With \(d_U(\lambda)\ge\varepsilon'/2\)
      at all scales, the set of centred rescaling limits of \(U\) is a
      nonempty family closed under the constructions; determine whether
      a diagonal argument over \(\lambda\to0\) and \(\lambda\to\infty\)
      produces an element whose small-scale and large-scale limits are
      itself, and record exactly which property of the family, such as
      sequential closedness in the topology of local convergence, is
      needed.
\item The next compulsory companion audit is V85.
\end{enumerate}
\end{programme}

\begin{assessment}[Position after V82]
\label{ass:v82-position}
Homogeneous singular points produce uniformly concentrated ancient
limits; lacunary points share their satellites and are not counted.  No
full stability or global-regularity claim is made.
\end{assessment}

% =============================================================================
\section*{Version V82 append-only record}
\addcontentsline{toc}{section}{Version V82 append-only record}
% =============================================================================

The complete V81 mathematical source was retained before this continuation.
Its SHA--256 digest was
\begin{center}\scriptsize\ttfamily
a3f99a1ea2809456b478403663324076af6a22f1f9d92ce6b1c6e3e738aa5833.
\end{center}
V82 proved that the ancient limit at a homogeneous singular point is
uniformly concentrated at all scales and singular at the origin, and that
satellites of lacunary points are shared.  No full regularity claim is
made.

% =============================================================================
\section{V83: the family of rescaling limits at a homogeneous point}
\label{sec:v83-entry}
% =============================================================================

Item (AC1) of Programme~\ref{prog:v83} is analysed to its exact
obstruction.  At a homogeneous singular point the centred rescaling
limits of the ancient limit form a family that is invariant under the
scaling group and closed under rescaling limits along sequences, whose
members are all uniformly concentrated, singular at the origin, and of
the class of V65.  A self-similar element would be a fixed point of the
scaling action on this family; the standard route to such a fixed point
is a quantity monotone along the scaling, such as the monotonicity
formulas of harmonic-map or mean-curvature flows, and no such quantity
is known for the Navier--Stokes equations or provided by this series.
The item is therefore recorded as reduced to a monotonicity formula.  No
global regularity claim is made.  The next compulsory companion audit is
V85.

% =============================================================================
\section{The family}
\label{sec:v83-family}
% =============================================================================

Let \(x_0\) be homogeneous with constant \(c\), \(U\) its ancient limit,
and let \(\mathcal F(U)\) be the set of all \(W\) obtained as local
limits, in the sense of Theorem~\ref{thm:v41-ancient-limit}, of centred
rescalings \(U^{\lambda_j}\) along sequences \(\lambda_j\in(0,\infty)\)
with a limit in \([0,\infty]\).

\begin{proposition}[Structure of \(\mathcal F(U)\)]
\label{prop:v83-family}
Modulo (E1)--(E8):
\begin{enumerate}[label=\textup{(\roman*)},leftmargin=3.7em]
\item \(U\in\mathcal F(U)\), and \(W\in\mathcal F(U)\) implies
      \(W^\lambda\in\mathcal F(U)\) for every \(\lambda>0\);
\item \(\mathcal F(U)\) is closed under rescaling limits: if
      \(W\in\mathcal F(U)\) and \(W^{\mu_j}\to W'\) locally, then
      \(W'\in\mathcal F(U)\);
\item every \(W\in\mathcal F(U)\) belongs to
      \(\mathcal A_\nu(C_I^2\nu^{3/2})\), satisfies the bounds of
      Corollary~\ref{cor:v56-class}, has \(d_W(\lambda)\ge\varepsilon'/2\)
      for all \(\lambda\), is singular at \((0,0)\), has a trace not in
      \(L^3\) near the origin, and has vorticity on every exterior region
      at every time;
\item \(\mathcal F(U)\) is nonempty and sequentially compact in the
      topology of local convergence, and every sequence in it has a
      subsequence converging to an element of \(\mathcal F(U)\).
\end{enumerate}
\end{proposition}

\begin{proof}
(i) \(U\) is the limit of the constant sequence \(\lambda_j=1\); and
\(W^\lambda\) is the limit of \(U^{\lambda\lambda_j}\).  (ii) A diagonal
argument: \(W\) is a local limit of \(U^{\lambda_j}\), \(W^{\mu_i}\) is
the local limit of \(U^{\mu_i\lambda_j}\) as \(j\to\infty\), and
\(W'\) is the limit of \(W^{\mu_i}\); choosing \(j(i)\) so that
\(U^{\mu_i\lambda_{j(i)}}\) is within \(1/i\) of \(W^{\mu_i}\) on the
\(i\)-th compact set gives \(U^{\mu_i\lambda_{j(i)}}\to W'\).  (iii)
Each property is preserved under local limits of rescalings: the class
bounds and the Morrey bound by lower semicontinuity
(Theorem~\ref{thm:v56-uniform-morrey}); the concentration bound by the
proof of Theorem~\ref{thm:v82-uniform} applied to \(W\), whose
rescalings have \(d_{W^\lambda}(\mu)=d_W(\lambda\mu)\ge\varepsilon'/2\);
singularity at the origin and the trace property as in
Theorem~\ref{thm:v82-uniform}; vorticity at infinity by
Theorem~\ref{thm:v63-vorticity-every-time}, whose proof applies to every
element of the class produced by these constructions.  (iv) The uniform
bounds of Proposition~\ref{prop:v41-uniform-bounds} hold for all
rescalings, so the compactness of Theorem~\ref{thm:v41-ancient-limit}
applies to any sequence in \(\mathcal F(U)\), and (ii) places the limit
in \(\mathcal F(U)\).
\end{proof}

\begin{firewall}[What a fixed point needs]
\label{fw:v83-fixed-point}
A discretely self-similar element is \(W\in\mathcal F(U)\) with
\(W^\lambda=W\) for some \(\lambda\ne1\); a continuously self-similar one
has \(W^\lambda=W\) for all \(\lambda\).  Proposition~\ref{prop:v83-family}
gives a compact family invariant under the scaling action, but a compact
invariant family need not contain a fixed point.  The classical route to
a fixed point, in the theory of tangent cones and tangent flows, is a
quantity \(\Theta(\lambda)\) monotone in \(\lambda\) along the rescalings
whose constancy characterizes self-similarity; its limits as
\(\lambda\to0\) and \(\lambda\to\infty\) then force every limit point to
be self-similar.  No monotonicity formula of this kind is known for the
three-dimensional Navier--Stokes equations, and none of the quantities
of this series is monotone along the scaling: the concentration
\(d_U(\lambda)\) is bounded below but not monotone, the Morrey ratios are
bounded but not monotone, and the dissipation fractions are not
monotone.  The fixed-point question is therefore reduced exactly to the
existence of a monotonicity formula, which is outside the reach of the
present tools.
\end{firewall}

% =============================================================================
\section{Integrated V83 conclusion and next acquisition order}
\label{sec:v83-conclusion}
% =============================================================================

\begin{theorem}[What V83 proves]
\label{thm:v83-conclusion}
Relative to V1--V82, modulo (E1)--(E8): the structure of the family of
rescaling limits at a homogeneous point (Proposition~\ref{prop:v83-family})
and the reduction of the fixed-point question to a monotonicity formula
(Firewall~\ref{fw:v83-fixed-point}).  The full Navier--Stokes
stability/global-regularity theorem remains unproved.
\end{theorem}

\begin{proof}
The cited items.
\end{proof}

\begin{programme}[V84 acquisition order]
\label{prog:v84}
\begin{enumerate}[label=\textup{(AC\arabic*)},leftmargin=4.8em]
\item Candidate monotone quantities.  Test, for the elements of
      \(\mathcal F(U)\), the scaled local energy
      \(\lambda^{-1}\int_{B_\lambda}|W(-\lambda^2)|^2\) and the scaled
      dissipation \(d_W(\lambda)\) for monotonicity in \(\lambda\) using
      the local energy identity; record the sign of the derivative and
      the terms that obstruct monotonicity.
\item Prepare the compulsory V85 companion audit.
\end{enumerate}
\end{programme}

\begin{assessment}[Position after V83]
\label{ass:v83-position}
The self-similar question at a homogeneous point is reduced to a
monotonicity formula.  No full stability or global-regularity claim is
made.
\end{assessment}

% =============================================================================
\section*{Version V83 append-only record}
\addcontentsline{toc}{section}{Version V83 append-only record}
% =============================================================================

The complete V82 mathematical source was retained before this continuation.
Its SHA--256 digest was
\begin{center}\scriptsize\ttfamily
ff9be76826f273dcd2a9e77aa9320e120a6406dc289446095296074954c5e7b8.
\end{center}
V83 described the family of rescaling limits at a homogeneous point as a
compact scale-invariant family of uniformly concentrated elements and
reduced the existence of a self-similar element to a monotonicity
formula.  No full regularity claim is made.

% =============================================================================
\section{V84: the scaled local energy is not monotone along the scaling}
\label{sec:v84-entry}
% =============================================================================

Item (AC1) of Programme~\ref{prog:v84} is carried out for the scaled
local energy at the parabolic time.  Its derivative in the scale is
computed exactly from the local energy identity: it contains a
nonnegative dissipation term, a negative term equal to the quantity
itself divided by the scale, and boundary terms of both signs coming from
the cutoff.  Neither sign dominates in general, so the scaled local
energy is not monotone, and no modification by the available quantities
makes it so; the fixed-point question of V83 therefore remains reduced to
a monotonicity formula that this series does not have.  No global
regularity claim is made.  The compulsory companion audit is due in V85.

% =============================================================================
\section{The derivative formula}
\label{sec:v84-derivative}
% =============================================================================

Let \(W\in\mathcal A_\nu(C)\) with pressure \(P\), let \(\phi\) be the
cutoff of Section~\ref{sec:v42-local-energy}, \(\phi_\lambda=\phi(\cdot/\lambda)\),
and define the scaled local energy at the parabolic time
\begin{equation}
 \mathcal E_W(\lambda):=\frac1\lambda\int_{\R^3}|W(-\lambda^2,y)|^2\phi_\lambda(y)\,\dd y .
 \label{eq:v84-scaled-energy}
\end{equation}
By Theorem~\ref{thm:v56-uniform-morrey}, \(0\le\mathcal E_W(\lambda)\le2C_M\)
for all \(\lambda\), and \(\mathcal E_{W^\mu}(\lambda)=\mathcal E_W(\mu\lambda)\).

\begin{proposition}[Exact derivative]
\label{prop:v84-derivative}
\(\mathcal E_W\) is differentiable on \((0,\infty)\) and, with all fields
evaluated at time \(-\lambda^2\),
\begin{equation}
 \boxed{\quad
 \frac{\dd}{\dd\lambda}\mathcal E_W(\lambda)
 =4\nu\int|\nabla W|^2\phi_\lambda
 -\frac1\lambda\mathcal E_W(\lambda)
 -\frac1{\lambda^3}\int|W|^2\,(y\cdot\nabla\phi)\Bigl(\frac y\lambda\Bigr)
 -2\nu\int|W|^2\Delta\phi_\lambda
 -2\int\bigl(|W|^2+2P\bigr)W\cdot\nabla\phi_\lambda .
 \quad}
 \label{eq:v84-derivative}
\end{equation}
The first term is nonnegative and, by \eqref{eq:v82-uniform} at a
homogeneous point, of order \(\lambda^{-1}\) on average over dyadic
scales; the second is nonpositive and of order \(\lambda^{-1}\); the
remaining three are supported on the annulus \(\lambda\le|y|\le2\lambda\)
and are bounded by \(C\lambda^{-1}\) (using the Morrey bound and the
\(L^6\), pressure bounds), with no sign.
\end{proposition}

\begin{proof}
Differentiate \eqref{eq:v84-scaled-energy}: the factor \(\lambda^{-1}\)
gives the second term; the time argument \(-\lambda^2\) gives
\(-2\lambda\cdot\lambda^{-1}\int\partial_s|W|^2\phi_\lambda
=-2\int\partial_s|W|^2\phi_\lambda\), and the local energy identity
\(\frac12\int\partial_s|W|^2\phi_\lambda=-\nu\int|\nabla W|^2\phi_\lambda
+\frac\nu2\int|W|^2\Delta\phi_\lambda+\frac12\int(|W|^2+2P)W\cdot\nabla\phi_\lambda\)
gives the first, fourth, and fifth terms; the cutoff argument gives the
third term via \(\partial_\lambda\phi(y/\lambda)=-\lambda^{-2}(y\cdot\nabla\phi)(y/\lambda)\).
The size statements: \(\int|\nabla W(-\lambda^2)|^2\phi_\lambda\) integrated
over \(\lambda\in[\Lambda,2\Lambda]\) against \(\dd\lambda\) is comparable
to \(\Lambda^{-1}\iint_{B_{2\Lambda}\times(-4\Lambda^2,-\Lambda^2)}|\nabla W|^2\),
which is at least \(c\,\varepsilon'\) at a homogeneous point by the proof
of Theorem~\ref{thm:v82-uniform} with the window
\((-4\Lambda^2,-\Lambda^2)\); the second term is \(\mathcal E_W/\lambda\le2C_M/\lambda\);
the annulus terms are bounded as in the proof of
Proposition~\ref{prop:v42-local-energy} with \(R=\lambda\) and the
Morrey bound, each by \(C\lambda^{-1}\).
\end{proof}

\begin{firewall}[No monotonicity]
\label{fw:v84-no-monotonicity}
The dissipation term and the self-similar term in
\eqref{eq:v84-derivative} are of the same order \(\lambda^{-1}\) and of
opposite signs, and the annulus terms are of the same order with
indefinite sign.  A self-similar \(W\) has \(\mathcal E_W\) constant, so
the five terms cancel exactly; for a general element none of the three
groups controls the others, and \(\mathcal E_W\) can increase or decrease
with \(\lambda\).  The same computation for the scaled dissipation
\(d_W(\lambda)\) produces a time derivative of the dissipation, that is,
the palinstrophy and the vortex-stretching term, again without sign.
A monotone quantity would have to combine energy, dissipation, and
pressure so that the annulus terms become a total derivative and the
stretching term acquires a sign; no such combination is known.  Item
(AC1) is closed negatively.
\end{firewall}

% =============================================================================
\section{Integrated V84 conclusion and next acquisition order}
\label{sec:v84-conclusion}
% =============================================================================

\begin{theorem}[What V84 proves]
\label{thm:v84-conclusion}
Relative to V1--V83: the exact derivative formula \eqref{eq:v84-derivative}
for the scaled local energy along the scaling, and the record that it is
not monotone.  The full Navier--Stokes stability/global-regularity
theorem remains unproved.
\end{theorem}

\begin{proof}
Proposition~\ref{prop:v84-derivative}.
\end{proof}

\begin{programme}[V85 acquisition order, with compulsory companion audit]
\label{prog:v85}
\begin{enumerate}[label=\textup{(AC\arabic*)},leftmargin=4.8em]
\item Perform the compulsory review of the current SM and YM notes;
      record digests; import nothing numerical.
\item Consolidate V81--V84 into a statement on homogeneous singular
      points: uniform concentration, the family of rescaling limits, and
      the exact reduction of self-similarity to a monotonicity formula.
\item The next compulsory companion audit after V85 is V90.
\end{enumerate}
\end{programme}

\begin{assessment}[Position after V84]
\label{ass:v84-position}
The scaling derivative of the natural local quantity has been computed
and has no sign.  No full stability or global-regularity claim is made.
\end{assessment}

% =============================================================================
\section*{Version V84 append-only record}
\addcontentsline{toc}{section}{Version V84 append-only record}
% =============================================================================

The complete V83 mathematical source was retained before this continuation.
Its SHA--256 digest was
\begin{center}\scriptsize\ttfamily
79ec8c3e76ce2b5e97503bf7fa00cf142877b6e10810800a6074f5f83a0ff1f5.
\end{center}
V84 computed the exact derivative of the scaled local energy along the
scaling and recorded that it has no sign.  No full regularity claim is
made.

% =============================================================================
\section{V85: compulsory companion audit and the consolidated statement on
homogeneous singular points}
\label{sec:v85-entry}
% =============================================================================

This is a fifty-fifth-version continuation.  The compulsory review of the
two companion programmes is performed first.  V81--V84 are then
consolidated into one statement on homogeneous singular points: the
per-scale equivalence, the finite count for each constant of
homogeneity, the uniform concentration of the ancient limit, the compact
scale-invariant family of its rescaling limits, and the exact reduction
of self-similarity to a monotonicity formula, together with the explicit
non-monotonicity of the natural candidate.  No global regularity claim is
made.  The next compulsory companion audit is V90.

% =============================================================================
\section{Compulsory V85 review of the current companion notes}
\label{sec:v85-companion-audit}
% =============================================================================

The current companion files in \texttt{latex\_notes} were inspected
read-only.  The frozen checkpoints are
\begin{align}
 &\texttt{Renewal\_SM\_Einstein\_Continuum\_Research\_Notes\_V78.tex},
 \notag\\[-2pt]
 &\qquad
 \texttt{559de44ca6e6b8ff09d89f8fbf9638ffd5efa9a4fb6eb7480f87357f38ed1396},
 \label{eq:v85-SM-checkpoint}\\
 &\texttt{Renewal\_Yang\_Mills\_Mass\_Gap\_Research\_Notes\_V55.tex},
 \notag\\[-2pt]
 &\qquad
 \texttt{9f990713e4c4e256d150a839514ffc0b4149cf4189ed0d63c6688dbdd733ecb7}.
 \label{eq:v85-YM-checkpoint}
\end{align}

\emph{SM V77--V78.}  Real-lapse and shift group averaging of each
nonzero linearized ADM mode is positive, removes the gauge pairs, and
recovers the Newton multiplier and the reduced vacuum space from one
rigging construction; for affine sourced constraints the surface
Jacobian cancels exactly, a uniform relative form bound on the reduced
potential suffices for a lower-bounded Hamiltonian, and the earlier
packet still produces a negative polynomial ray beyond its frequency
horizon, so the global linearized stability card is not repaired by
group averaging.

\emph{YM V55.}  The YM programme completed its own SM/NS audit, proved a
uniform antipodal heat-kernel lower bound, computed the normalized Wilson
antipodal density, proved a positive tilted-time tail theorem and an exact
positive conditional Wilson split, a truncated-time Duhamel card at every
arity, and a support-diameter zero-free strip, closing two rows of its
current card and isolating a logarithmic source gate.

\begin{theorem}[V85 companion-audit decision]
\label{thm:v85-companion-decision}
No SM V77--V78 averaging, Jacobian, or form-bound constant and no YM V55
heat-kernel or Duhamel estimate is imported.  Neither bears on the three
gates of Theorem~\ref{thm:v70-gates}.
\end{theorem}

\begin{proof}
The hypotheses of every cited result are constraint-averaging measures or
lattice-gauge heat kernels; none is a Navier--Stokes balance.
\end{proof}

% =============================================================================
\section{Homogeneous singular points, consolidated}
\label{sec:v85-homogeneous}
% =============================================================================

\begin{theorem}[Homogeneous singular points, V81--V84]
\label{thm:v85-homogeneous}
Under the type-I rate and modulo (E1)--(E8), let \(x_0\in\Sigma\) be
homogeneous with constant \(c\), and let \(\varepsilon'=\varepsilon'(c)\)
be the dissipation constant of Theorem~\ref{thm:v81-count}.  Then:
\begin{enumerate}[label=\textup{(\arabic*)},leftmargin=3.2em]
\item \emph{Equivalence.}  Positive terminal density and nonvanishing
      rescaled dissipation are equivalent at every point along every
      sequence of scales, with quantitative constants; \(x_0\) is active
      at every small scale with \(f_{x_0}(r)\ge\varepsilon'\).
\item \emph{Count.}  The singular points homogeneous with constant \(c\)
      number at most \(2C_I^2\nu^{3/2}/\varepsilon'\); all homogeneous
      points form a countable set.
\item \emph{Uniform concentration.}  The ancient limit \(U\) at \(x_0\),
      along any sequence of scales, has \(d_U(\lambda)\ge\varepsilon'/2\)
      for every \(\lambda>0\), is singular at \((0,0)\), is in case (B)
      with a non-\(L^3\) trace at the origin, and lies in both
      concentration branches.
\item \emph{The family.}  The centred rescaling limits of \(U\) form a
      nonempty, scale-invariant, sequentially compact family closed under
      rescaling limits, all of whose members have the properties in (3),
      the class properties of Theorem~\ref{thm:v65-ancient-limit}, and
      vorticity on every exterior region at every time.
\item \emph{Self-similarity.}  A self-similar element of that family
      exists if a quantity monotone along the scaling exists on it; the
      scaled local energy has the exact derivative
      \eqref{eq:v84-derivative}, which has no sign, and no monotone
      quantity is provided.
\end{enumerate}
\end{theorem}

\begin{proof}
(1) Theorem~\ref{thm:v81-equivalence} and Corollary~\ref{cor:v81-quantitative}.
(2) Theorem~\ref{thm:v81-count}.  (3) Theorem~\ref{thm:v82-uniform}.
(4) Proposition~\ref{prop:v83-family}.  (5) Firewall~\ref{fw:v83-fixed-point}
and Proposition~\ref{prop:v84-derivative}.
\end{proof}

\begin{assessment}[The two shapes after fifty-five versions]
\label{ass:v85-shapes}
A type-I singular point is homogeneous or lacunary.  Homogeneous points
are countable, finitely many for each constant, and carry uniformly
concentrated ancient limits whose rescaling family is compact and
scale-invariant; their self-similarity is reduced to a monotonicity
formula.  Lacunary points have satellites at every intermediate scale,
shared among all such points, and are not counted.  The type-I Liouville
gate of V41 stands over both.  No full stability or global-regularity
claim is made.
\end{assessment}

% =============================================================================
\section{Integrated V85 conclusion and next acquisition order}
\label{sec:v85-conclusion}
% =============================================================================

\begin{theorem}[What V85 establishes]
\label{thm:v85-conclusion}
Relative to V1--V84: the companion-audit decision
(Theorem~\ref{thm:v85-companion-decision}) and the consolidated statement
on homogeneous singular points (Theorem~\ref{thm:v85-homogeneous}).  The
full Navier--Stokes stability/global-regularity theorem remains unproved.
\end{theorem}

\begin{proof}
The cited results.
\end{proof}

\begin{programme}[V86 acquisition order]
\label{prog:v86}
\begin{enumerate}[label=\textup{(AC\arabic*)},leftmargin=4.8em]
\item Weighted monotone candidates.  Test the Gaussian-weighted local
      energy \(\lambda^{-1}\int|W(-\lambda^2)|^2e^{-|y|^2/(4\nu\lambda^2)}\),
      whose weight is the backward heat kernel, for monotonicity: the
      Laplacian and cutoff terms of \eqref{eq:v84-derivative} then combine
      through the heat equation, and only the nonlinear flux term
      remains; record its sign structure.
\item The next compulsory companion audit is V90.
\end{enumerate}
\end{programme}

\begin{assessment}[Position after V85]
\label{ass:v85-position}
The series is consolidated on homogeneous singular points.  No full
stability or global-regularity claim is made.
\end{assessment}

% =============================================================================
\section*{Version V85 append-only record}
\addcontentsline{toc}{section}{Version V85 append-only record}
% =============================================================================

The complete V84 mathematical source was retained before this continuation.
Its SHA--256 digest was
\begin{center}\scriptsize\ttfamily
166367e94a14f5b4354ecb4229ef4267768457b35f66b92e245059231e62b9d0.
\end{center}
V85 performed the compulsory companion audit and consolidated V81--V84
into one statement on homogeneous singular points.  No full regularity
claim is made.

% =============================================================================
\section{V86: the Gaussian-weighted scaled energy and the inward radial flux}
\label{sec:v86-entry}
% =============================================================================

Item (AC1) of Programme~\ref{prog:v86} is carried out and yields a
positive result.  With the backward Gaussian weight in place of a cutoff,
the scaled local energy of an element of the ancient class has an exact
scaling derivative consisting of two nonnegative terms, the weighted
dissipation and twice the quantity itself divided by the scale, and one
indefinite term, the Gaussian-weighted radial momentum flux.  The
quantity is bounded uniformly in the scale by the Morrey bound.  At a
homogeneous singular point the weighted dissipation contributes a fixed
positive amount on every scale interval of fixed ratio.  Consequently the
radial flux term must be negative on average over scales, with an
integral over scales that diverges logarithmically: the flow of the
ancient limit is, in the Gaussian-weighted sense, inward at every scale.
This is the first signed statement about the ancient limit in the series.
It is not a monotonicity formula, but it identifies exactly the quantity
whose sign would give one.  No global regularity claim is made.  The
next compulsory companion audit is V90.

% =============================================================================
\section{The weighted energy and its derivative}
\label{sec:v86-identity}
% =============================================================================

Let \(W\in\mathcal A_\nu(C)\) with pressure \(P\), and put
\begin{equation}
 \psi(y,s):=\exp\Bigl(-\frac{|y|^2}{4\nu(-s)}\Bigr),
 \qquad
 \Phi_W(\lambda):=\frac1\lambda\int_{\R^3}|W(-\lambda^2,y)|^2\,\psi(y,-\lambda^2)\,\dd y
 \qquad(\lambda>0).
 \label{eq:v86-Phi}
\end{equation}

\begin{lemma}[Properties of the weight]
\label{lem:v86-weight}
With \(T=-s\): \(\nabla\psi=-\frac{y}{2\nu T}\psi\),
\(\nu\Delta\psi=\partial_T\psi-\frac{3}{2T}\psi\), hence
\(\partial_s\psi=-\nu\Delta\psi-\frac{3}{2T}\psi\).  Moreover
\(\Phi_W(\lambda)\le C_\Phi:=C_M\sum_{k\ge0}2^{k}e^{-4^{k-1}/(4\nu)}<\infty\)
for all \(\lambda\), and \(\Phi_{W^\mu}(\lambda)=\Phi_W(\mu\lambda)\).
\end{lemma}

\begin{proof}
Direct differentiation.  For the bound, split \(\R^3\) into
\(B_\lambda\) and the shells \(2^{k-1}\lambda\le|y|<2^k\lambda\), use
\(\psi\le e^{-4^{k-1}/(4\nu)}\) on the \(k\)-th shell, and the
time-uniform Morrey bound \(\int_{B_{2^k\lambda}}|W|^2\le C_M2^k\lambda\).
Scaling is the change of variables.
\end{proof}

\begin{theorem}[Exact scaling derivative of the weighted energy]
\label{thm:v86-derivative}
\(\Phi_W\) is differentiable and, with all fields at time \(s=-\lambda^2\),
\begin{equation}
 \boxed{\quad
 \Phi_W'(\lambda)
 =4\nu\int|\nabla W|^2\psi
 +\frac2\lambda\Phi_W(\lambda)
 +\frac1{\nu\lambda^2}\int\bigl(|W|^2+2P\bigr)\,(y\cdot W)\,\psi .
 \quad}
 \label{eq:v86-derivative}
\end{equation}
\end{theorem}

\begin{proof}
Let \(F(s):=\int|W(s)|^2\psi(\cdot,s)\).  From the equation,
\(\partial_s|W|^2=\nu(\Delta|W|^2-2|\nabla W|^2)-(W\cdot\nabla)|W|^2-2W\cdot\nabla P\);
multiply by \(\psi\), integrate, and integrate by parts using
\(\operatorname{div}W=0\):
\(\int\partial_s|W|^2\psi=\nu\int|W|^2\Delta\psi-2\nu\int|\nabla W|^2\psi
+\int(|W|^2+2P)W\cdot\nabla\psi\).  Adding \(\int|W|^2\partial_s\psi\) with
Lemma~\ref{lem:v86-weight} cancels the Laplacian terms:
\[
 F'(s)=-2\nu\int|\nabla W|^2\psi-\frac{3}{2T}F(s)
 -\frac1{2\nu T}\int(|W|^2+2P)(y\cdot W)\psi .
\]
Since \(\Phi_W(\lambda)=\lambda^{-1}F(-\lambda^2)\),
\(\Phi_W'(\lambda)=-\lambda^{-2}F-2F'(-\lambda^2)\), and inserting
\(F'\) with \(T=\lambda^2\) and \(F=\lambda\Phi_W\) gives
\eqref{eq:v86-derivative}.  All integrals converge by the \(L^6\),
enstrophy, pressure, and Morrey bounds and the Gaussian decay.
\end{proof}

% =============================================================================
\section{The inward flux at a homogeneous point}
\label{sec:v86-inward}
% =============================================================================

\begin{theorem}[The weighted radial flux is inward on average over scales]
\label{thm:v86-inward}
Let \(W\) be the ancient limit at a homogeneous singular point, or any
element of its rescaling family \(\mathcal F(U)\), with the window
parameters \(R,S,\eta\) and the dissipation constant \(\varepsilon\) of
Corollary~\ref{cor:v81-quantitative}.  Then there is \(c_2>0\), depending
only on \(R,S,\eta,\varepsilon,\nu\), such that for all \(0<a<b\),
\begin{equation}
 \boxed{\quad
 \int_a^b\frac1{\nu\lambda^2}\int\bigl(|W|^2+2P\bigr)(y\cdot W)\,\psi\,\dd y\,\dd\lambda
 \ \le\ 2C_\Phi-c_2\Bigl(\frac{\log(b/a)}{\log(S/\eta)}-1\Bigr),
 \quad}
 \label{eq:v86-inward}
\end{equation}
all fields evaluated at \(s=-\lambda^2\).  In particular the integral
tends to \(-\infty\) as \(b/a\to\infty\): the Gaussian-weighted radial
momentum flux of the ancient limit is negative on average at every scale.
\end{theorem}

\begin{proof}
Integrate \eqref{eq:v86-derivative} over \([a,b]\):
\(\Phi_W(b)-\Phi_W(a)=\int_a^b4\nu\int|\nabla W|^2\psi+\int_a^b\frac2\lambda\Phi_W
+\int_a^b(\text{flux})\); the left side is at most \(2C_\Phi\) and the
second term is nonnegative, so \(\int_a^b(\text{flux})\le2C_\Phi
-\int_a^b4\nu\int|\nabla W|^2\psi\).  For the dissipation term, fix
\(\Lambda\) and consider \(\lambda\in[\eta^{1/2}\Lambda,S^{1/2}\Lambda]\);
on \(B_{R'\lambda}\) with \(R':=R\eta^{-1/2}\), \(\psi(\cdot,-\lambda^2)\ge e^{-R'^2/(4\nu)}\),
so with the substitution \(s=-\lambda^2\), \(\dd\lambda=\dd s/(2\lambda)
\ge\dd s/(2S^{1/2}\Lambda)\),
\[
 \int_{\eta^{1/2}\Lambda}^{S^{1/2}\Lambda}4\nu\int|\nabla W(-\lambda^2)|^2\psi\,\dd\lambda
 \ge\frac{4\nu e^{-R'^2/(4\nu)}}{2S^{1/2}\Lambda}
 \iint_{B_{R\Lambda}\times(-S\Lambda^2,-\eta\Lambda^2)}|\nabla W|^2
 \ge\frac{2\nu e^{-R'^2/(4\nu)}}{S^{1/2}}\,\varepsilon=:c_2,
\]
because \(B_{R'\lambda}\supset B_{R\Lambda}\) for \(\lambda\ge\eta^{1/2}\Lambda\),
and the cylinder dissipation is at least \(\varepsilon\Lambda\) by the
homogeneity of \(x_0\) transferred to \(W\) as in the proof of
Theorem~\ref{thm:v82-uniform} (and to every element of \(\mathcal F(U)\)
by Proposition~\ref{prop:v83-family}(iii)).  The interval \([a,b]\)
contains at least \(\lfloor\log(b/a)/\log(S/\eta)\rfloor\) disjoint
intervals of the form \([\eta^{1/2}\Lambda,S^{1/2}\Lambda]\), which gives
\eqref{eq:v86-inward}.
\end{proof}

\begin{assessment}[What the inward flux says]
\label{ass:v86-inward}
The quantity \(\int(|W|^2+2P)(y\cdot W)\psi\) is the Gaussian-weighted
flux of energy and pressure through the radial direction; its negativity
means that, weighted by the backward heat kernel at the parabolic scale,
the ancient limit transports energy and pressure work inward toward the
singular point, at every scale, by an amount whose integral over scales
diverges like the logarithm of the scale ratio.  For a self-similar
element the three terms of \eqref{eq:v86-derivative} cancel exactly at
every scale, so the inward flux equals minus the weighted dissipation
minus \(2\Phi/\lambda\) identically; the theorem shows that the same
holds on average for every element of the family.  A monotonicity
formula would follow if the flux term could be shown to be bounded below
by \(-\,\)(the positive terms) pointwise in \(\lambda\), which is not
proved; the theorem gives the averaged version only.
\end{assessment}

% =============================================================================
\section{Integrated V86 conclusion and next acquisition order}
\label{sec:v86-conclusion}
% =============================================================================

\begin{theorem}[What V86 proves]
\label{thm:v86-conclusion}
Relative to V1--V85, modulo (E1)--(E8): the exact scaling derivative
\eqref{eq:v86-derivative} of the Gaussian-weighted scaled energy, its
uniform bound, and the averaged inward radial flux \eqref{eq:v86-inward}
at homogeneous points.  The full Navier--Stokes stability/global-regularity
theorem remains unproved.
\end{theorem}

\begin{proof}
Lemma~\ref{lem:v86-weight}, Theorem~\ref{thm:v86-derivative},
Theorem~\ref{thm:v86-inward}.
\end{proof}

\begin{programme}[V87 acquisition order]
\label{prog:v87}
\begin{enumerate}[label=\textup{(AC\arabic*)},leftmargin=4.8em]
\item Pointwise sign.  Determine whether the flux term can be bounded
      below by \(-\theta\) times the positive terms with \(\theta<1\)
      pointwise in \(\lambda\), using the pressure representation
      \(P=R_iR_j(W_iW_j)\) and the Gaussian weight; a pointwise bound
      would make \(\Phi_W\) monotone and, with V83, give a self-similar
      element at every homogeneous point.
\item Lacunary points.  Determine what \eqref{eq:v86-derivative} gives
      along intermediate scales, where the dissipation term near the
      point vanishes: the flux term is then bounded below by
      \(-2\Phi/\lambda\) on average, and the averaged inward flux is
      weaker; record the exact statement.
\item The next compulsory companion audit is V90.
\end{enumerate}
\end{programme}

\begin{assessment}[Position after V86]
\label{ass:v86-position}
The series has its first signed statement about the ancient limit: its
Gaussian-weighted radial flux is inward on average over scales at every
homogeneous singular point, with a logarithmically divergent integral.
No full stability or global-regularity claim is made.
\end{assessment}

% =============================================================================
\section*{Version V86 append-only record}
\addcontentsline{toc}{section}{Version V86 append-only record}
% =============================================================================

The complete V85 mathematical source was retained before this continuation.
Its SHA--256 digest was
\begin{center}\scriptsize\ttfamily
daff450239c3d056112278b81839a566005ed06d8a410890dad0747014520f1a.
\end{center}
V86 proved the exact scaling derivative of the Gaussian-weighted scaled
energy of the ancient limit and derived the averaged inward radial flux
at homogeneous singular points.  No full regularity claim is made.

% =============================================================================
\section{V87: the inward flux of the original solution at a homogeneous
singular point}
\label{sec:v87-entry}
% =============================================================================

The Gaussian-weighted identity of V86 is transferred from the ancient
limit to the original solution on the torus, using the periodized
backward heat kernel centred at the singular point.  At a homogeneous
singular point the same argument shows that the weighted radial flux of
the solution itself, integrated over scales down to zero, diverges to
minus infinity: the solution transports energy and pressure work inward
toward the point, in the Gaussian-weighted sense, at every small scale.
Item (AC1) of Programme~\ref{prog:v87}, a pointwise sign, is recorded as
open; item (AC2) is recorded with the exact statement available at
lacunary scales.  No global regularity claim is made.  The next
compulsory companion audit is V90.

% =============================================================================
\section{The periodized weight}
\label{sec:v87-periodized}
% =============================================================================

For \(x_0\in\Torus^3\), \(T>0\), and \(x\in\Torus^3\) identified with
\([-\pi,\pi)^3\), let
\begin{equation}
 \Psi(x,T):=\sum_{k\in\mathbb Z^3}\exp\Bigl(-\frac{|x-x_0+2\pi k|^2}{4\nu T}\Bigr),
 \qquad
 \Phi_u(\lambda):=\frac1\lambda\int_{\Torus^3}|u(T_*-\lambda^2,x)|^2\,\Psi(x,\lambda^2)\,\dd x .
 \label{eq:v87-Phi-u}
\end{equation}
\(\Psi\) is smooth and periodic, satisfies
\(\partial_T\Psi=\nu\Delta\Psi+\frac{3}{2T}\Psi\), and
\(\nabla\Psi=-\sum_k\frac{x-x_0+2\pi k}{2\nu T}\exp(\cdots)\).

\begin{theorem}[Scaling derivative for the original solution]
\label{thm:v87-derivative-u}
For \(0<\lambda<T_*^{1/2}\), with all fields at time \(t=T_*-\lambda^2\),
\begin{equation}
 \boxed{\quad
 \Phi_u'(\lambda)
 =4\nu\int|\nabla u|^2\Psi
 +\frac2\lambda\Phi_u(\lambda)
 +\frac1{\nu\lambda^2}\sum_{k\in\mathbb Z^3}\int\bigl(|u|^2+2p\bigr)\,
 \bigl((x-x_0+2\pi k)\cdot u\bigr)\,e^{-|x-x_0+2\pi k|^2/(4\nu\lambda^2)}\,\dd x .
 \quad}
 \label{eq:v87-derivative-u}
\end{equation}
Under the type-I rate, \(\Phi_u(\lambda)\le C_\Phi'\) for
\(\lambda\le(\delta/2)^{1/2}\), with \(C_\Phi'\) depending on \(C_M\).
\end{theorem}

\begin{proof}
Identical to Theorem~\ref{thm:v86-derivative}, with the periodic
integrations by parts on the torus, \(\Psi\) in place of \(\psi\), and the
time \(T_*-\lambda^2\) in place of \(-\lambda^2\); the periodization does
not affect the linear identity \(\partial_T\Psi=\nu\Delta\Psi+\frac3{2T}\Psi\),
which holds term by term.  The bound is Lemma~\ref{lem:v86-weight}'s
argument with Theorem~\ref{thm:v47-morrey}, the images with
\(|k|\ge1\) contributing \(e^{-c/\lambda^2}\).
\end{proof}

% =============================================================================
\section{Divergence of the inward flux at homogeneous points}
\label{sec:v87-divergence}
% =============================================================================

\begin{theorem}[Inward flux of the solution at a homogeneous singular point]
\label{thm:v87-inward-u}
Let \(x_0\in\Sigma\) be homogeneous with constant \(c\), with window
parameters \(R,S,\eta\) and dissipation constant \(\varepsilon\) from
Corollary~\ref{cor:v81-quantitative}.  Then, with \(c_2\) as in
Theorem~\ref{thm:v86-inward}, for all \(0<a<b\le(\delta/2)^{1/2}\) small,
\begin{equation}
 \boxed{\quad
 \int_a^b\frac1{\nu\lambda^2}\sum_k\int\bigl(|u|^2+2p\bigr)
 \bigl((x-x_0+2\pi k)\cdot u\bigr)e^{-|x-x_0+2\pi k|^2/(4\nu\lambda^2)}\,\dd x\,\dd\lambda
 \ \le\ 2C_\Phi'-c_2\Bigl(\frac{\log(b/a)}{\log(S/\eta)}-1\Bigr),
 \quad}
 \label{eq:v87-inward-u}
\end{equation}
so the integral over \((0,b]\) diverges to \(-\infty\).
\end{theorem}

\begin{proof}
As in Theorem~\ref{thm:v86-inward}: integrate \eqref{eq:v87-derivative-u},
bound \(\Phi_u(b)-\Phi_u(a)\le2C_\Phi'\), drop the nonnegative
\(2\Phi_u/\lambda\) term, and bound the dissipation term from below on
each scale interval \([\eta^{1/2}\Lambda,S^{1/2}\Lambda]\) by the cylinder
dissipation of homogeneity, using \(\Psi\ge e^{-R'^2/(4\nu)}\) on
\(B_{R'\lambda}(x_0)\) (the \(k=0\) term alone).
\end{proof}

\begin{assessment}[Meaning for the singularity itself]
\label{ass:v87-meaning}
Theorem~\ref{thm:v87-inward-u} is a statement about the actual solution
near its singular time, not about a limit: at a homogeneous singular
point, the Gaussian-weighted radial flux at parabolic scale \(\lambda\),
averaged over dyadic ranges of \(\lambda\), is inward by a fixed amount
per range.  Since the total energy is finite and the singular point
carries no energy in the limit (Theorem~\ref{thm:v44-localization}(iv)),
the inward transport is balanced by dissipation, exactly as
\eqref{eq:v87-derivative-u} says: the positive dissipation term and the
negative flux term cancel on average, with \(\Phi_u\) staying bounded.
This is the energy-budget picture of a homogeneous singularity, now
proved with the correct weights.
\end{assessment}

% =============================================================================
\section{Records}
\label{sec:v87-records}
% =============================================================================

\begin{firewall}[Pointwise sign]
\label{fw:v87-pointwise}
A pointwise lower bound on the flux term by \(-\theta\) times the
positive terms, \(\theta<1\), is not proved; for a general divergence-free
field at a fixed time the flux can have either sign, and the Navier--Stokes
evolution is used in \eqref{eq:v86-derivative} only through the energy
identity, which does not constrain the sign of the flux at a single
scale.  Item (AC1) is open.
\end{firewall}

\begin{firewall}[Lacunary scales]
\label{fw:v87-lacunary}
Along intermediate scales of a lacunary point the \(k=0\) dissipation
term near the point is \(o(\lambda^{-1})\) by Theorem~\ref{thm:v77-satellite},
while satellites at distance \(d\gg\lambda\) are weighted by
\(e^{-d^2/(4\nu\lambda^2)}\) and contribute negligibly; hence on such
scale intervals the flux term satisfies only
\(\int(\text{flux})\ge-2C_\Phi'-\int2\Phi_u/\lambda-o(1)\), which does not
force divergence.  The inward-flux divergence is therefore a signature of
homogeneous scales.  Item (AC2) is settled with this statement.
\end{firewall}

% =============================================================================
\section{Integrated V87 conclusion and next acquisition order}
\label{sec:v87-conclusion}
% =============================================================================

\begin{theorem}[What V87 proves]
\label{thm:v87-conclusion}
Relative to V1--V86, modulo (E1)--(E8): the scaling derivative of the
Gaussian-weighted scaled energy for the original solution
\eqref{eq:v87-derivative-u} and the divergence of the inward flux at
homogeneous singular points \eqref{eq:v87-inward-u}.  The full
Navier--Stokes stability/global-regularity theorem remains unproved.
\end{theorem}

\begin{proof}
Theorems~\ref{thm:v87-derivative-u}--\ref{thm:v87-inward-u}.
\end{proof}

\begin{programme}[V88 acquisition order]
\label{prog:v88}
\begin{enumerate}[label=\textup{(AC\arabic*)},leftmargin=4.8em]
\item Decompose the flux term into its energy part
      \(\int|u|^2(x\cdot u)\Psi\) and its pressure part
      \(2\int p\,(x\cdot u)\Psi\), and determine which carries the inward
      sign on average; the pressure part can be rewritten through the
      Riesz representation as a quadratic form in \(u\otimes u\).
\item Test the identity on the explicit V38 three-shell datum and on the
      balanced plane wave to see the sign of each part at one time.
\item The next compulsory companion audit is V90.
\end{enumerate}
\end{programme}

\begin{assessment}[Position after V87]
\label{ass:v87-position}
The inward-flux statement now holds for the solution itself at every
homogeneous singular point.  No full stability or global-regularity
claim is made.
\end{assessment}

% =============================================================================
\section*{Version V87 append-only record}
\addcontentsline{toc}{section}{Version V87 append-only record}
% =============================================================================

The complete V86 mathematical source was retained before this continuation.
Its SHA--256 digest was
\begin{center}\scriptsize\ttfamily
25f782af58edadde46bf3d15038f949bc84969f2f08a9789eb9b7e656210c97b.
\end{center}
V87 transferred the Gaussian-weighted identity to the original solution
on the torus and proved the divergence of the inward radial flux over
scales at homogeneous singular points.  No full regularity claim is
made.

% =============================================================================
\section{V88: the two parts of the radial flux and the impossibility of an
instantaneous sign}
\label{sec:v88-entry}
% =============================================================================

Items (AC1) and (AC2) of Programme~\ref{prog:v88} are carried out.  The
Gaussian-weighted radial flux of V86--V87 is split into its energy part
and its pressure part; the pressure part is rewritten, through the Riesz
representation of the pressure, as the integral of the strain--vorticity
balance \(|S|^2-\frac12|\omega|^2\) against a Newtonian potential of the
weighted radial momentum, and each part is shown to be of the same order
\(\lambda^{-1}\) as the positive terms of the identity.  Both parts are
odd under \(W\mapsto-W\), so neither can have a sign among instantaneous
divergence-free states; the parity symmetry of the Navier--Stokes class
fixes the flux, so no symmetry argument excludes a sign along the
evolution either.  Two exact solutions with identically vanishing flux
are recorded, the decaying plane wave on the torus and the columnar
swirl on \(\R^3\), for which the weighted scaled energy is strictly
increasing in the scale.  No global regularity claim is made.  The next
compulsory companion audit is V90.

% =============================================================================
\section{The decomposition}
\label{sec:v88-decomposition}
% =============================================================================

Throughout, \(W\in\mathcal A_\nu(C)\) with pressure
\(P=\partial_i\partial_j(-\Delta)^{-1}(W_iW_j)\), \(\psi=\psi(\cdot,-\lambda^2)\)
as in \eqref{eq:v86-Phi}, all fields at time \(s=-\lambda^2\), and
\begin{equation}
 \mathcal J(\lambda):=\frac1{\nu\lambda^2}\int(|W|^2+2P)(y\cdot W)\psi,
 \qquad
 \mathcal J_E:=\frac1{\nu\lambda^2}\int|W|^2(y\cdot W)\psi,
 \qquad
 \mathcal J_P:=\frac2{\nu\lambda^2}\int P\,(y\cdot W)\psi .
 \label{eq:v88-parts}
\end{equation}

\begin{proposition}[Advective and potential forms of the two parts]
\label{prop:v88-forms}
Let \(q\) be the Newtonian potential of the weighted radial momentum,
\(-\Delta q=(y\cdot W)\psi\), \(q=\frac1{4\pi|\cdot|}*\bigl((y\cdot W)\psi\bigr)\).
Then
\begin{align}
 \mathcal J_E&=-2\int|W|^2\,W\cdot\nabla\psi=2\int\psi\,W\cdot\nabla|W|^2 ,
 \label{eq:v88-JE}\\
 \mathcal J_P&=\frac2{\nu\lambda^2}\int\Bigl(|S|^2-\frac12|\omega|^2\Bigr)q
 =-\frac2{\nu\lambda^2}\int\bigl((W\cdot\nabla)W\bigr)\cdot\nabla q ,
 \label{eq:v88-JP}
\end{align}
where \(S=\frac12(\nabla W+\nabla W^{\mathsf T})\) and \(\omega=\operatorname{curl}W\).
All integrals converge absolutely.
\end{proposition}

\begin{proof}
Since \(\nabla\psi=-\frac{y}{2\nu\lambda^2}\psi\), one has
\((y\cdot W)\psi=-2\nu\lambda^2\,W\cdot\nabla\psi=-2\nu\lambda^2\operatorname{div}(\psi W)\),
which gives the first form of \eqref{eq:v88-JE}; the second follows by
integrating by parts with \(\operatorname{div}W=0\).  For the pressure
part, \(-\Delta P=\partial_iW_j\partial_jW_i=|S|^2-\frac12|\omega|^2\)
pointwise (the identity \(\partial_iW_j\partial_jW_i=|S|^2-|A|^2\) for
the symmetric and antisymmetric parts of \(\nabla W\), with
\(|A|^2=\frac12|\omega|^2\)), so
\(\int P\,(y\cdot W)\psi=\int P(-\Delta q)=\int(-\Delta P)\,q\), the
integrations by parts being justified because \((y\cdot W)\psi\in L^1\cap L^{3}\)
(Morrey bound and \(L^6\) bound with the Gaussian factor), hence
\(q\in L^3\) with \(\nabla q\in L^{3/2}\cap L^\infty\) and
\(\nabla^2q\in L^{3}\), while \(P\in L^3\) and \(\nabla W\in L^2\); the
boundary terms at infinity vanish by the decay of \(q\), \(\nabla q\) and the
integrability of \(P\), \(|\nabla W|^2\).  The second form of \eqref{eq:v88-JP}
follows from \(-\Delta P=\partial_i\partial_j(W_iW_j)\):
\(\int(-\Delta P)q=\int W_iW_j\partial_i\partial_jq=-\int W_i\partial_iW_j\,\partial_jq\),
using \(\operatorname{div}W=0\) once more.  Absolute convergence: the
Gaussian factor makes \(|y||W|\psi\) integrable against \(|W|^2\) and
against \(|P|\) by the Morrey and \(L^3\) bounds.
\end{proof}

\begin{lemma}[Both parts are of order \(\lambda^{-1}\)]
\label{lem:v88-order}
There is \(C_J\) depending only on \(C\), \(C_M\), \(\nu\) such that
\(|\mathcal J_E(\lambda)|+|\mathcal J_P(\lambda)|\le C_J\lambda^{-1}\)
for every \(\lambda>0\).  Consequently the averaged inward flux
\eqref{eq:v86-inward} at a homogeneous point satisfies
\begin{equation}
 -\,C_J\log\frac ba\ \le\ \int_a^b\mathcal J(\lambda)\,\dd\lambda
 \ \le\ 2C_\Phi-c_2\Bigl(\frac{\log(b/a)}{\log(S/\eta)}-1\Bigr),
 \label{eq:v88-two-sided}
\end{equation}
so the inward flux is logarithmic in the scale ratio from both sides.
\end{lemma}

\begin{proof}
At \(s=-\lambda^2\) the class bound gives \(\|\nabla W\|_2\le C^{1/2}\lambda^{-1/2}\),
hence \(\|W\|_6\le C_6C^{1/2}\lambda^{-1/2}\) and, by the Calder\'on--Zygmund
bound for the Riesz representation, \(\|P\|_3\le C_{\rm CZ}\|W\|_6^2\le C'\lambda^{-1}\).
On the \(k\)-th shell \(2^{k-1}\lambda\le|y|<2^k\lambda\) (and on \(B_\lambda\)),
\(|y|\psi\le2^k\lambda e^{-4^{k-1}/(4\nu)}\), so by H\"older
(exponents \(6,2,3\)) and the Morrey bound
\[
 \int|W|^3|y|\psi\le\sum_k2^k\lambda e^{-4^{k-1}/(4\nu)}\|W\|_{L^6}\Bigl(\int_{B_{2^k\lambda}}|W|^2\Bigr)^{1/2}\bigl|B_{2^k\lambda}\bigr|^{1/3}
 \le C''\sum_k2^ke^{-4^{k-1}/(4\nu)}\,\lambda\cdot\lambda^{-1/2}(2^k\lambda)^{1/2}(2^k\lambda)
 \le C'''\lambda^{2},
\]
which after division by \(\nu\lambda^2\) bounds \(\mathcal J_E\) by a
constant only.  The bound of order \(\lambda^{-1}\) uses the \(L^\infty\)
form of the type-I rate, valid in the class by
Theorem~\ref{thm:v46-Linfty-upgrade} transferred to the limit:
\(|W(-\lambda^2)|\le C_\infty\lambda^{-1}\).  Then
\(\int|W|^3|y|\psi\le C_\infty\lambda^{-1}\sum_k2^k\lambda e^{-4^{k-1}/(4\nu)}C_M2^k\lambda\le C\lambda\),
so \(|\mathcal J_E|\le C/(\nu\lambda)\).  For the pressure part,
\(|\mathcal J_P|\le\frac2{\nu\lambda^2}\|P\|_3\,\bigl\||y|W\psi\bigr\|_{3/2}\)
and \(\||y|W\psi\|_{3/2}\le C\lambda\bigl(\int|W|^{3/2}e^{-|y|^2/(8\nu\lambda^2)}\bigr)^{2/3}
\le C\lambda\bigl((C_M\lambda)^{3/4}(C\lambda^3)^{1/4}\bigr)^{2/3}=C\lambda^2\)
by H\"older against the Gaussian and the Morrey bound, so
\(|\mathcal J_P|\le\frac2{\nu\lambda^2}C'\lambda^{-1}C\lambda^2=C/(\nu\lambda)\).
The two-sided bound \eqref{eq:v88-two-sided} follows by integrating
\(|\mathcal J|\le C_J\lambda^{-1}\) and Theorem~\ref{thm:v86-inward}.
\end{proof}

\begin{remark}
\label{rem:v88-order}
The first count in the proof, which uses only \(L^6\), gives
\(|\mathcal J_E|\le C\) rather than \(C\lambda^{-1}\), and is recorded to
show that the \(L^\infty\) form of the rate is what makes the energy part
scale like the other terms of \eqref{eq:v86-derivative}; in the class
\(\mathcal A_\nu(C)\) without the \(L^\infty\) bound only the constant bound
is available.
\end{remark}

% =============================================================================
\section{No instantaneous sign}
\label{sec:v88-sign}
% =============================================================================

\begin{proposition}[Oddness and parity]
\label{prop:v88-oddness}
\begin{enumerate}[label=\textup{(\roman*)},leftmargin=3.2em]
\item For every divergence-free \(W\) at a fixed time for which the
      integrals converge, \(\mathcal J_E(-W)=-\mathcal J_E(W)\) and
      \(\mathcal J_P(-W)=-\mathcal J_P(W)\).  Hence no bound of the form
      \(\mathcal J\ge-\theta\bigl(4\nu\int|\nabla W|^2\psi+\frac2\lambda\Phi_W\bigr)\)
      with \(\theta<1\) can hold for all instantaneous states: it fails
      for \(W\) or for \(-W\) whenever \(\mathcal J(W)\ne0\).
\item The parity map \(W\mapsto\widetilde W\), \(\widetilde W(s,y):=-W(s,-y)\),
      \(\widetilde P(s,y):=P(s,-y)\), preserves \(\mathcal A_\nu(C)\),
      \(\Phi_W\), the dissipation term, and \(\mathcal J_E\), \(\mathcal J_P\)
      separately.  Hence no symmetry of the class reverses the sign of
      the flux, and a dynamical sign statement is not excluded by
      Item~(i).
\end{enumerate}
\end{proposition}

\begin{proof}
(i) \(\mathcal J_E\) is cubic and \(\mathcal J_P\) is cubic (\(P\) is
quadratic in \(W\)), so both are odd.  If \(\mathcal J(W)\ne0\) then one of
\(\mathcal J(\pm W)\) is negative and equals minus the other, while the
positive terms are even; the inequality with \(\theta<1\) for both
signs would force \(|\mathcal J|\le\theta(\cdots)\), which is a statement
that can fail for a suitable instantaneous state, e.g.\ by taking
\(W=\mu V\) with \(V\) fixed with \(\mathcal J(V)\ne0\) and \(\mu\to\infty\):
the flux scales like \(\mu^3\) and the positive terms like \(\mu^2\).
(ii) If \(W\) solves the Navier--Stokes equations so does \(\widetilde W\)
with pressure \(\widetilde P\); the Dirichlet integral is preserved, so
the class is preserved; \(\psi\) is even; \(|\widetilde W|^2(y)=|W|^2(-y)\),
\(y\cdot\widetilde W(y)=(-y)\cdot W(-y)\), \(\widetilde P(y)=P(-y)\), so
each integrand is transported by \(y\mapsto-y\).
\end{proof}

\begin{firewall}[Item (AC1): the sign is dynamical or absent]
\label{fw:v88-AC1}
By Proposition~\ref{prop:v88-oddness}(i) a pointwise lower bound on the
flux by \(-\theta\) times the positive terms with \(\theta<1\) cannot be a
statement about the instantaneous state alone, and by (ii) no symmetry of
the class contradicts it for solutions.  Whether it holds along the
evolution for elements of \(\mathcal F(U)\) at homogeneous points is not
decided.  The series records that the only signed statement available is
the averaged one, \eqref{eq:v88-two-sided}.
\end{firewall}

% =============================================================================
\section{Exact solutions with vanishing flux}
\label{sec:v88-examples}
% =============================================================================

\begin{proposition}[Two exact solutions with \(\mathcal J\equiv0\)]
\label{prop:v88-examples}
\begin{enumerate}[label=\textup{(\roman*)},leftmargin=3.2em]
\item \emph{Plane wave on the torus.}  Let \(k\in\mathbb Z^3\setminus\{0\}\),
      \(e\perp k\), \(u(t,x)=e\cos(k\cdot x)e^{-\nu|k|^2t}\), \(p\equiv0\).
      This solves the Navier--Stokes equations on \(\Torus^3\) for all
      time, and for every centre \(x_0\) the flux term of
      \eqref{eq:v87-derivative-u} vanishes identically in \(\lambda\), so
      \(\Phi_u'(\lambda)=4\nu\int|\nabla u|^2\Psi+\frac2\lambda\Phi_u>0\).
\item \emph{Columnar swirl on \(\R^3\).}  Let \(W(s,y)=w(s,r)\,e_\theta\) in
      cylindrical coordinates about the \(y_3\)-axis, with \(w\) solving
      \(\partial_sw=\nu(\partial_r^2w+r^{-1}\partial_rw-r^{-2}w)\) and
      \(P=\int_0^r w^2(\rho)\rho^{-1}\dd\rho\).  This solves the
      Navier--Stokes equations on \(\R^3\), and \(y\cdot W\equiv0\), so
      \(\mathcal J\equiv0\) and \(\Phi_W'(\lambda)>0\) for all \(\lambda\)
      at which \(\Phi_W\) is finite.  Such \(W\) has infinite energy and
      infinite Dirichlet integral on \(\R^3\), hence is not in
      \(\mathcal A_\nu(C)\); the Gaussian weight makes \(\Phi_W\) finite
      whenever \(w(s,\cdot)\) is bounded.
\end{enumerate}
\end{proposition}

\begin{proof}
(i) \(\operatorname{div}u=-(e\cdot k)\sin(k\cdot x)e^{-\nu|k|^2t}=0\),
\((u\cdot\nabla)u=(e\cdot\nabla\cos(k\cdot x))\,e\,e^{-2\nu|k|^2t}
=-(e\cdot k)\sin(k\cdot x)\,e\,e^{-2\nu|k|^2t}=0\), and
\(\partial_tu=\nu\Delta u\); the pressure is zero.  With \(z=x-x_0+2\pi k'\) one has
\((z\cdot u)e^{-|z|^2/(4\nu\lambda^2)}=-2\nu\lambda^2\,u\cdot\nabla_xe^{-|z|^2/(4\nu\lambda^2)}\),
so the flux term of \eqref{eq:v87-derivative-u} equals
\(-2\int|u|^2u\cdot\nabla\Psi=2\int\Psi\,u\cdot\nabla|u|^2\) after periodic
integration by parts with \(\operatorname{div}u=0\), and
\(u\cdot\nabla|u|^2=-2\cos(k\cdot x)\sin(k\cdot x)(e\cdot k)e^{-2\nu|k|^2t}=0\).
(ii) Standard: with \(u_r=u_z=0\) the momentum equations reduce to
\(\partial_sw=\nu(\Delta w-r^{-2}w)\) in the azimuthal direction and
\(\partial_rP=w^2/r\) in the radial direction, with the axial equation
trivial; the given \(P\) satisfies the radial equation since \(w\) is
independent of \(y_3\).  \(y\cdot e_\theta=0\).  The energy statements are
clear since \(W\) is independent of \(y_3\).
\end{proof}

\begin{assessment}[What the examples show]
\label{ass:v88-examples}
Both exact solutions have vanishing radial flux because their transport
term \(u\cdot\nabla|u|^2\) vanishes identically: the plane wave is a
Stokes flow and the columnar swirl advects nothing.  For such flows the
weighted scaled energy is monotone increasing in the scale, as for the
heat equation.  The averaged inward flux at a homogeneous singular point
is therefore a property of the genuinely nonlinear transport, and the
dissipation-forced inward flux of Theorem~\ref{thm:v87-inward-u} has no
counterpart in any flow with vanishing transport term.  Item (AC2) is
settled with these two computations; the three-shell datum of V38 is not
needed, since Proposition~\ref{prop:v88-oddness}(i) already shows that
the two parts take both signs among instantaneous states.
\end{assessment}

% =============================================================================
\section{Integrated V88 conclusion and next acquisition order}
\label{sec:v88-conclusion}
% =============================================================================

\begin{theorem}[What V88 proves]
\label{thm:v88-conclusion}
Relative to V1--V87, modulo (E1)--(E8): the advective and potential forms
of the two parts of the radial flux \eqref{eq:v88-JE}--\eqref{eq:v88-JP},
the order bound and the two-sided logarithmic bound
\eqref{eq:v88-two-sided}, the oddness and parity statements of
Proposition~\ref{prop:v88-oddness}, and the two exact solutions with
vanishing flux.  The full Navier--Stokes stability/global-regularity
theorem remains unproved.
\end{theorem}

\begin{proof}
Propositions~\ref{prop:v88-forms}, \ref{prop:v88-oddness},
\ref{prop:v88-examples} and Lemma~\ref{lem:v88-order}.
\end{proof}

\begin{programme}[V89 acquisition order]
\label{prog:v89}
\begin{enumerate}[label=\textup{(AC\arabic*)},leftmargin=4.8em]
\item Windowed flux.  Determine the window ratio \(\Theta_0\) such that
      the flux integral over every scale window of ratio \(\Theta_0\) is
      at most \(-c_2\) for every element of \(\mathcal F(U)\) at a
      homogeneous point, and whether the windowed flux functional is
      continuous on the family in the topology of
      Proposition~\ref{prop:v83-family}, so that an extremal element
      exists.
\item Strain versus vorticity.  Using \eqref{eq:v88-JP}, record what the
      averaged inward flux says about the Gaussian-weighted excess of
      strain over vorticity against the potential \(q\), and whether the
      sign of \(q\) can be controlled at homogeneous points.
\item The next compulsory companion audit is V90.
\end{enumerate}
\end{programme}

\begin{assessment}[Position after V88]
\label{ass:v88-position}
The radial flux has been split and both parts have been shown to scale
like the positive terms and to take both signs instantaneously; the
averaged inward flux is logarithmic from both sides.  No full stability
or global-regularity claim is made.
\end{assessment}

% =============================================================================
\section*{Version V88 append-only record}
\addcontentsline{toc}{section}{Version V88 append-only record}
% =============================================================================

The complete V87 mathematical source was retained before this continuation.
Its SHA--256 digest was
\begin{center}\scriptsize\ttfamily
27615589d9e2ccf997823c1043b1eaec27bfa15c49df4d45cc45adf3909604d5.
\end{center}
V88 decomposed the Gaussian-weighted radial flux into its energy and
pressure parts, proved the order and two-sided logarithmic bounds, the
oddness and parity statements, and recorded two exact solutions with
vanishing flux.  No full regularity claim is made.

% =============================================================================
\section{V89: the windowed flux on the rescaling family and the potential
of the weighted radial momentum}
\label{sec:v89-entry}
% =============================================================================

Items (AC1) and (AC2) of Programme~\ref{prog:v89} are carried out.  A
window ratio \(\Theta_0\) is computed such that, for every element of
the rescaling family at a homogeneous singular point and every position
of the window, the flux integral over a scale window of ratio \(\Theta_0\)
is at most \(-c_2\).  The windowed flux is shown to be a continuous
functional on the family in the topology of local convergence, so that
its extreme values over the family are attained; the maximum is at most
\(-c_2\) and the minimum is at least \(-C_J\log\Theta_0\).  For the
pressure part, the weighted radial momentum is shown to have zero total
mass and a dipole moment equal to the Gaussian-weighted momentum, and
the strain--vorticity balance is shown to have zero total, so the
pressure part is a correlation of two mean-zero quantities; the sign of
the potential is recorded as uncontrolled.  No global regularity claim is
made.  The next compulsory companion audit is V90.

% =============================================================================
\section{The windowed flux}
\label{sec:v89-window}
% =============================================================================

For \(W\in\mathcal A_\nu(C)\) and \(\Theta>1\) put
\begin{equation}
 \mathcal J_\Theta(W):=\int_1^\Theta\mathcal J_W(\lambda)\,\dd\lambda
 =\int_{-\Theta^2}^{-1}\frac1{2\nu(-s)^{3/2}}\int_{\R^3}\bigl(|W|^2+2P\bigr)(y\cdot W)\,\psi(y,s)\,\dd y\,\dd s ,
 \label{eq:v89-windowed}
\end{equation}
where \(\mathcal J_W\) is the flux term of \eqref{eq:v86-derivative} and
the second form uses \(s=-\lambda^2\).  By the scaling
\(\Phi_{W^\mu}(\lambda)=\Phi_W(\mu\lambda)\), one has
\(\mathcal J_{W^\mu}(\lambda)=\mu\,\mathcal J_W(\mu\lambda)\) and
\(\mathcal J_\Theta(W^\mu)=\int_\mu^{\mu\Theta}\mathcal J_W\), so the flux
over the window \([\mu,\mu\Theta]\) of \(W\) is the flux over \([1,\Theta]\)
of \(W^\mu\).

\begin{proposition}[Uniformly negative windows]
\label{prop:v89-negative-window}
Let \(x_0\) be a homogeneous singular point, \(\mathcal F(U)\) its
rescaling family, and \(c_2\), \(C_\Phi\), \(S/\eta\) as in
Theorem~\ref{thm:v86-inward}.  Put
\begin{equation}
 N_0:=\Bigl\lceil2+\frac{2C_\Phi}{c_2}\Bigr\rceil,
 \qquad
 \Theta_0:=(S/\eta)^{N_0}.
 \label{eq:v89-Theta0}
\end{equation}
Then for every \(W\in\mathcal F(U)\), every \(\Theta\ge\Theta_0\), and every
\(\mu>0\),
\begin{equation}
 \boxed{\quad
 -\,C_J\log\Theta\ \le\ \int_\mu^{\mu\Theta}\mathcal J_W(\lambda)\,\dd\lambda\ \le\ -\,c_2 .
 \quad}
 \label{eq:v89-window-bound}
\end{equation}
\end{proposition}

\begin{proof}
Theorem~\ref{thm:v86-inward} with \(a=\mu\), \(b=\mu\Theta\) gives the
upper bound \(2C_\Phi-c_2(\log\Theta/\log(S/\eta)-1)\le2C_\Phi-c_2(N_0-1)\le-c_2\)
by the choice of \(N_0\).  The lower bound is Lemma~\ref{lem:v88-order}
integrated.
\end{proof}

\begin{proposition}[Continuity of the windowed flux on the family]
\label{prop:v89-continuity}
Let \(W_k,W\in\mathcal F(U)\) with \(W_k\to W\) locally, in the sense of
Theorem~\ref{thm:v41-ancient-limit}: \(W_k\to W\) in
\(L^2((-S,-\eta);H^1(B_R))\) and \(P_k\rightharpoonup P\) weakly in
\(L^3((-S,-\eta)\times B_R)\) for all \(0<\eta<S\), \(R>0\).  Then
\(\mathcal J_\Theta(W_k)\to\mathcal J_\Theta(W)\) for every \(\Theta>1\).
Consequently \(\mathcal J_\Theta\) attains its supremum and its infimum on
\(\mathcal F(U)\); for \(\Theta\ge\Theta_0\),
\begin{equation}
 -\,C_J\log\Theta\ \le\ \min_{\mathcal F(U)}\mathcal J_\Theta
 \ \le\ \max_{\mathcal F(U)}\mathcal J_\Theta\ \le\ -\,c_2 .
 \label{eq:v89-extremes}
\end{equation}
\end{proposition}

\begin{proof}
Fix \(\Theta\) and write the second form of \eqref{eq:v89-windowed} over
the time interval \(I=(-\Theta^2,-1)\), on which the weight
\((-s)^{-3/2}\) is bounded.  Every element of \(\mathcal F(U)\) satisfies,
uniformly, the bounds of Corollary~\ref{cor:v56-class} and the pointwise
bounds \(|W(s)|\le C_\infty(-s)^{-1/2}\), \(|\nabla W(s)|\le C_\infty'(-s)^{-1}\):
the rescalings \(U^{\lambda}\) of the original solution satisfy them by
Theorem~\ref{thm:v46-Linfty-upgrade}, and they pass to local limits by
lower semicontinuity of the supremum norm under almost-everywhere
convergence.  Split the space integral into \(B_R\) and its complement.

\emph{Tails.}  On \(|y|>R\), by the shell decomposition of the proof of
Lemma~\ref{lem:v88-order}, the cubic part is bounded by
\(C\sum_{2^j\ge R}2^{2j}e^{-4^{j-1}/(4\nu)}\) uniformly in \(k\) and \(s\in I\),
and the pressure part by \(\|P_k(s)\|_3\||y|W_k(s)\psi\|_{L^{3/2}(|y|>R)}\),
where \(\|P_k(s)\|_3\le C\) on \(I\) and the second factor is bounded by
\(C\bigl(\sum_{2^j\ge R}2^{j}e^{-4^{j-1}/(8\nu)}\cdot2^{j}\bigr)^{2/3}\) by the same
H\"older-against-the-Gaussian argument with the Morrey bound.  Both tails
tend to zero as \(R\to\infty\) uniformly in \(k\).

\emph{Cubic part on \(B_R\).}  With
\(|W_k|^2W_k-|W|^2W\) bounded pointwise by \(C(|W_k|^2+|W|^2)|W_k-W|\) and
the uniform \(L^\infty(I\times\R^3)\) bounds,
\(\int_I\int_{B_R}\bigl||W_k|^2W_k-|W|^2W\bigr||y|\psi\le CR\int_I\int_{B_R}|W_k-W|\to0\),
since \(W_k\to W\) in \(L^2(I;L^2(B_R))\).

\emph{Pressure part on \(B_R\).}  \((y\cdot W_k)\psi\to(y\cdot W)\psi\) in
\(L^{3/2}(I\times B_R)\) by the same convergence and the boundedness of the
weight, and \(P_k\rightharpoonup P\) weakly in \(L^3(I\times B_R)\); the
pairing of a weakly convergent sequence with a strongly convergent one
converges.

Combining, \(\limsup_k|\mathcal J_\Theta(W_k)-\mathcal J_\Theta(W)|\) is
bounded by the tails at radius \(R\), which are arbitrarily small.  The
extreme values are attained by the sequential compactness of
Proposition~\ref{prop:v83-family}(iv), and \eqref{eq:v89-extremes}
restates \eqref{eq:v89-window-bound}.
\end{proof}

\begin{firewall}[What an extremal element does not give]
\label{fw:v89-extremal}
The minimizer \(W_*\) of \(\mathcal J_{\Theta}\) on \(\mathcal F(U)\)
satisfies \(\mathcal J_\Theta(W_*^\mu)\ge\mathcal J_\Theta(W_*)\) for every
\(\mu>0\), that is, the flux of \(W_*\) over \([\mu,\mu\Theta]\) is
minimal at \(\mu=1\); differentiation in \(\mu\) gives
\(\mathcal J_{W_*}(\Theta)\Theta=\mathcal J_{W_*}(1)\) whenever the
one-sided derivatives exist, an equality of the flux density at the two
ends of the window, but no self-similarity.  The family is not a
manifold and the functional is not known to be differentiable along any
other direction, so no Euler--Lagrange equation is available.
\end{firewall}

% =============================================================================
\section{The potential of the weighted radial momentum}
\label{sec:v89-potential}
% =============================================================================

\begin{proposition}[Zero mass, dipole moment, and the strain--vorticity total]
\label{prop:v89-moments}
Let \(W\in\mathcal A_\nu(C)\) at time \(s=-\lambda^2\), and let \(q\) be
the Newtonian potential of \(m:=(y\cdot W)\psi\) as in
Proposition~\ref{prop:v88-forms}.  Then
\begin{equation}
 \int m=0,\qquad
 \int y\,m=2\nu\lambda^2\int\psi\,W,\qquad
 \int\Bigl(|S|^2-\frac12|\omega|^2\Bigr)=0 .
 \label{eq:v89-moments}
\end{equation}
Consequently \(q(y)=\frac{1}{4\pi}\frac{y\cdot\int\psi W}{|y|^3}\,2\nu\lambda^2+O(|y|^{-3})\)
as \(|y|\to\infty\), and for every constant \(\bar q\),
\begin{equation}
 \mathcal J_P=\frac2{\nu\lambda^2}\int\Bigl(|S|^2-\frac12|\omega|^2\Bigr)(q-\bar q).
 \label{eq:v89-JP-centred}
\end{equation}
\end{proposition}

\begin{proof}
\(m=-2\nu\lambda^2\operatorname{div}(\psi W)\) with \(\psi W\) integrable
and decaying, so \(\int m=0\); and
\(\int y_im=-2\nu\lambda^2\int y_i\operatorname{div}(\psi W)=2\nu\lambda^2\int\psi W_i\).
The strain--vorticity total: \(|S|^2-\frac12|\omega|^2=\partial_iW_j\partial_jW_i=\partial_i(W_j\partial_jW_i)\)
using \(\operatorname{div}W=0\), and \(W_j\partial_jW_i\in L^{3/2}(\R^3)\)
by \(W\in L^6\), \(\nabla W\in L^2\); for a divergence of an
\(L^{3/2}\) field the integral against a cutoff \(\chi(y/R)\) equals
\(-\int W_j\partial_jW_i\,\partial_i\chi(y/R)\), bounded by
\(CR^{-1}\|W\cdot\nabla W\|_{L^{3/2}(B_{2R}\setminus B_R)}|B_{2R}|^{1/3}\le C\|W\cdot\nabla W\|_{L^{3/2}(B_{2R}\setminus B_R)}\to0\).
The expansion of \(q\) is the multipole expansion of the Newtonian
potential of a compactly decaying, mean-zero source with the computed
dipole moment (the Gaussian factor gives all moments).  The centred
form \eqref{eq:v89-JP-centred} follows from the third identity.
\end{proof}

\begin{firewall}[Item (AC2): sign of the potential]
\label{fw:v89-potential-sign}
By \eqref{eq:v89-JP-centred} the pressure part is the covariance of the
strain excess \(|S|^2-\frac12|\omega|^2\) with the potential \(q\).  The
averaged inward flux says that \(\mathcal J_E+\mathcal J_P\) is negative on
average over scales at a homogeneous point; it does not separate the two
parts.  The potential \(q\) has no monopole, its dipole is the weighted
momentum, and its sign is not controlled: the source \(m\) is the
weighted radial momentum, which is positive where the flow is outward
and negative where it is inward, and \(q\) is its inverse Laplacian.
No statement about the excess of strain over vorticity at homogeneous
points follows.
\end{firewall}

% =============================================================================
\section{Integrated V89 conclusion and next acquisition order}
\label{sec:v89-conclusion}
% =============================================================================

\begin{theorem}[What V89 proves]
\label{thm:v89-conclusion}
Relative to V1--V88, modulo (E1)--(E8): the uniformly negative windows
\eqref{eq:v89-window-bound} with the explicit ratio \(\Theta_0\), the
continuity of the windowed flux on the rescaling family and the
attainment of its extreme values \eqref{eq:v89-extremes}, and the
moment identities \eqref{eq:v89-moments} with the centred form
\eqref{eq:v89-JP-centred} of the pressure part.  The full Navier--Stokes
stability/global-regularity theorem remains unproved.
\end{theorem}

\begin{proof}
Propositions~\ref{prop:v89-negative-window}, \ref{prop:v89-continuity},
\ref{prop:v89-moments}.
\end{proof}

\begin{programme}[V90 acquisition order, with compulsory companion audit]
\label{prog:v90}
\begin{enumerate}[label=\textup{(AC\arabic*)},leftmargin=4.8em]
\item Perform the compulsory review of the current SM and YM notes;
      record digests; import nothing numerical.
\item Consolidate V86--V89 into one statement on the Gaussian-weighted
      scaled energy at homogeneous singular points: the exact
      derivative, the averaged inward flux for the limit and for the
      solution, the two parts and their bounds, the negative windows,
      and the moment identities.
\item The next compulsory companion audit after V90 is V95.
\end{enumerate}
\end{programme}

\begin{assessment}[Position after V89]
\label{ass:v89-position}
The windowed flux is a continuous functional on the compact
scale-invariant family and is uniformly negative on windows of a
computable ratio; the pressure part is a covariance of two mean-zero
quantities.  No full stability or global-regularity claim is made.
\end{assessment}

% =============================================================================
\section*{Version V89 append-only record}
\addcontentsline{toc}{section}{Version V89 append-only record}
% =============================================================================

The complete V88 mathematical source was retained before this continuation.
Its SHA--256 digest was
\begin{center}\scriptsize\ttfamily
efbf5e682f9f9ecb81c7f13aab702abd5ce76e4b43ce9937d150cf21fac9e823.
\end{center}
V89 proved the uniformly negative windows, the continuity and extreme
values of the windowed flux on the rescaling family, and the moment
identities of the weighted radial momentum.  No full regularity claim is
made.

% =============================================================================
\section{V90: compulsory companion audit and the consolidated statement on
the Gaussian-weighted energy at homogeneous singular points}
\label{sec:v90-entry}
% =============================================================================

This is a sixtieth-version continuation.  The compulsory review of the
two companion programmes is performed first.  V86--V89 are then
consolidated into one statement on the Gaussian-weighted scaled energy
at homogeneous singular points: the exact scaling derivative for the
ancient limit and for the original solution, the averaged inward radial
flux with its two-sided logarithmic bounds, the uniformly negative
windows on the compact scale-invariant rescaling family, the two parts
of the flux with their advective and potential forms, the impossibility
of an instantaneous sign, and the moment identities of the weighted
radial momentum.  No global regularity claim is made.  The next
compulsory companion audit is V95.

% =============================================================================
\section{Compulsory V90 review of the current companion notes}
\label{sec:v90-companion-audit}
% =============================================================================

The current companion files in \texttt{latex\_notes} were inspected
read-only.  The frozen checkpoints are
\begin{align}
 &\texttt{Renewal\_SM\_Einstein\_Continuum\_Research\_Notes\_V79.tex},
 \notag\\[-2pt]
 &\qquad
 \texttt{a29083ac3021046f27d49f388ec52c1583228d921a1e1cee71b8c1dec66da213},
 \label{eq:v90-SM-checkpoint}\\
 &\texttt{Renewal\_Yang\_Mills\_Mass\_Gap\_Research\_Notes\_V57.tex},
 \notag\\[-2pt]
 &\qquad
 \texttt{4b480ba75edec24b06d3e0ee70354bb003b6c7b4e70cf94546b8e1c03cd86ab2}.
 \label{eq:v90-YM-checkpoint}
\end{align}
At the time of inspection a file named \(\ldots\)\texttt{V56.tex} with
the same digest as \eqref{eq:v90-YM-checkpoint} was also present; the
two are byte-identical, and the V57 name is recorded.

\emph{SM V79.}  At a spectral cutoff \(\Lambda\), the worst transverse
traceless and quadratic Newton source constants grow at exactly fourth
order; the lower bound is realized by two nearly opposite high-frequency
scalar modes whose product has fixed output momentum while their local
Higgs quartic stays constant, so fixed local quartic and gravitational
couplings cannot satisfy the global linearized relative form bound as
the cutoff is removed; finite spectral and amplitude boxes have
lower-bounded Friedrichs Hamiltonians with explicit cutoff losses.

\emph{YM V56--V57.}  The YM programme verified that its earlier
Banach-algebra invertibility of the complete squarefree operator source
remains uniform for the positive local conditional law, computed exact
singleton and conjugate-singlet zeros, and reduced its remaining
all-arity tangent Gaussian card to a clock cumulant hierarchy
(cumulant bounds \(r!K_0^rs^{2r-1}\)), shown equivalent to a thermal
zero-free scale \(O(s^{-2})\) for the centred local-clock logarithm and
implied by a moderate-deviation card; it records that bounded support
alone gives only the \(O(s^{-1})\) scale.

\begin{theorem}[V90 companion-audit decision]
\label{thm:v90-companion-decision}
No SM V79 cutoff constant or beat-frequency source bound and no YM
V56--V57 invertibility, zero, or cumulant estimate is imported.  Neither
bears on the three gates of Theorem~\ref{thm:v70-gates} or on the
monotonicity question of Firewall~\ref{fw:v88-AC1}.
\end{theorem}

\begin{proof}
The SM results concern finite-cutoff relative form bounds for a
gravitational quartic and the YM results concern cumulants of a lattice
clock variable; none is a Navier--Stokes balance, weight, or flux.
\end{proof}

\begin{assessment}[Direction after the audit]
\label{ass:v90-direction}
One structural parallel is noted and not imported: the YM programme's
reduction of an all-order card to a single hierarchy of cumulant bounds
mirrors this series' reduction of self-similarity at homogeneous points
to a single monotonicity formula (Firewall~\ref{fw:v83-fixed-point}),
and in both programmes the reduced statement is recorded as open rather
than assumed.  The direction of this series is unchanged.
\end{assessment}

% =============================================================================
\section{The Gaussian-weighted energy at homogeneous points, consolidated}
\label{sec:v90-consolidated}
% =============================================================================

\begin{theorem}[Gaussian-weighted energy at homogeneous singular points, V86--V89]
\label{thm:v90-consolidated}
Under the type-I rate and modulo (E1)--(E8), let \(x_0\in\Sigma\) be
homogeneous with constant \(c\), let \(U\) be its ancient limit and
\(\mathcal F(U)\) its rescaling family, and let \(\psi\), \(\Psi\),
\(\Phi_W\), \(\Phi_u\), \(\mathcal J_W\), \(\mathcal J_E\), \(\mathcal J_P\)
be as in V86--V88.  Then:
\begin{enumerate}[label=\textup{(\arabic*)},leftmargin=3.2em]
\item \emph{Exact derivatives.}  For every \(W\in\mathcal A_\nu(C)\),
      \(\Phi_W'(\lambda)=4\nu\int|\nabla W|^2\psi+\frac2\lambda\Phi_W+\mathcal J_W(\lambda)\)
      with \(\Phi_W\le C_\Phi\); the same identity holds for the original
      solution on the torus with the periodized weight, with
      \(\Phi_u\le C_\Phi'\) on small scales.
\item \emph{Averaged inward flux.}  For every \(W\in\mathcal F(U)\) and
      \(0<a<b\),
      \(-C_J\log(b/a)\le\int_a^b\mathcal J_W\le2C_\Phi-c_2(\log(b/a)/\log(S/\eta)-1)\);
      the same upper bound holds for the original solution at \(x_0\) on
      small scales.  In particular the flux integral over \((0,b]\)
      diverges to \(-\infty\), logarithmically.
\item \emph{Negative windows and extremal elements.}  With
      \(\Theta_0=(S/\eta)^{N_0}\), \(N_0=\lceil2+2C_\Phi/c_2\rceil\),
      every window of ratio \(\Theta\ge\Theta_0\) at every position carries
      flux at most \(-c_2\) for every element of \(\mathcal F(U)\); the
      windowed flux is continuous on \(\mathcal F(U)\) in the local
      topology and attains its extreme values, which lie in
      \([-C_J\log\Theta,-c_2]\).
\item \emph{The two parts.}  \(\mathcal J_E=2\int\psi\,W\cdot\nabla|W|^2\)
      and \(\mathcal J_P=\frac2{\nu\lambda^2}\int(|S|^2-\frac12|\omega|^2)(q-\bar q)\)
      for the Newtonian potential \(q\) of the weighted radial momentum
      \(m=(y\cdot W)\psi\), which has zero mass and dipole moment
      \(2\nu\lambda^2\int\psi W\); both parts are bounded by
      \(C_J\lambda^{-1}\) in the class with the \(L^\infty\) rate; both are
      odd under \(W\mapsto-W\) and invariant under the parity symmetry of
      the class.
\item \emph{No instantaneous sign; exact zero-flux flows.}  No pointwise
      lower bound of the flux by \(-\theta\) times the positive terms,
      \(\theta<1\), holds for all instantaneous divergence-free states;
      the decaying plane wave on the torus and the columnar swirl on
      \(\R^3\) have identically vanishing flux and strictly increasing
      weighted scaled energy.
\end{enumerate}
\end{theorem}

\begin{proof}
(1) Theorems~\ref{thm:v86-derivative} and \ref{thm:v87-derivative-u},
Lemma~\ref{lem:v86-weight}.  (2) Theorems~\ref{thm:v86-inward},
\ref{thm:v87-inward-u}, and Lemma~\ref{lem:v88-order}.
(3) Propositions~\ref{prop:v89-negative-window} and \ref{prop:v89-continuity}.
(4) Propositions~\ref{prop:v88-forms}, \ref{prop:v88-oddness}(ii),
\ref{prop:v89-moments}, and Lemma~\ref{lem:v88-order}.
(5) Propositions~\ref{prop:v88-oddness}(i) and \ref{prop:v88-examples}.
\end{proof}

\begin{assessment}[The state of the monotonicity question after sixty versions]
\label{ass:v90-monotonicity}
The natural Gaussian-weighted quantity is not monotone, its scaling
derivative is an exact three-term identity, and the indefinite term is
inward on average with logarithmic divergence at every homogeneous
singular point, for the ancient limit, for every element of its
rescaling family, and for the solution itself.  A pointwise sign of the
indefinite term would give self-similarity at every homogeneous point
(Firewall~\ref{fw:v83-fixed-point}); such a sign cannot be a property of
the instantaneous state and is not excluded by any symmetry.  The three
gates of Theorem~\ref{thm:v70-gates} stand.  No full stability or
global-regularity claim is made.
\end{assessment}

% =============================================================================
\section{Integrated V90 conclusion and next acquisition order}
\label{sec:v90-conclusion}
% =============================================================================

\begin{theorem}[What V90 establishes]
\label{thm:v90-conclusion}
Relative to V1--V89: the companion-audit decision
(Theorem~\ref{thm:v90-companion-decision}) and the consolidated statement
on the Gaussian-weighted energy at homogeneous singular points
(Theorem~\ref{thm:v90-consolidated}).  The full Navier--Stokes
stability/global-regularity theorem remains unproved.
\end{theorem}

\begin{proof}
The cited results.
\end{proof}

\begin{programme}[V91 acquisition order]
\label{prog:v91}
\begin{enumerate}[label=\textup{(AC\arabic*)},leftmargin=4.8em]
\item Lower bound on the weighted energy.  Determine whether
      \(\Phi_W(\lambda)\ge c_\Phi>0\) for all \(\lambda\) and all
      \(W\in\mathcal F(U)\) at a homogeneous point, using the uniform
      concentration \eqref{eq:v82-uniform} and the per-scale
      equivalence of V81 (dissipation activity forces terminal density,
      hence energy at the parabolic scale); record the constants.
\item Limit sets.  Determine the sets of limit points of \(\Phi_W(\lambda)\)
      as \(\lambda\to0\) and \(\lambda\to\infty\) on the family, and whether
      a self-similar element would be forced if \(\Phi_W\) had a limit at
      one end.
\item The next compulsory companion audit is V95.
\end{enumerate}
\end{programme}

\begin{assessment}[Position after V90]
\label{ass:v90-position}
The series is consolidated on the Gaussian-weighted energy at
homogeneous singular points.  No full stability or global-regularity
claim is made.
\end{assessment}

% =============================================================================
\section*{Version V90 append-only record}
\addcontentsline{toc}{section}{Version V90 append-only record}
% =============================================================================

The complete V89 mathematical source was retained before this continuation.
Its SHA--256 digest was
\begin{center}\scriptsize\ttfamily
069d78f300f3b9ea7dca067640f81222e0c85126fd58f5a01d7d09aa60614798.
\end{center}
V90 performed the compulsory companion audit and consolidated V86--V89
into one statement on the Gaussian-weighted energy at homogeneous
singular points.  No full regularity claim is made.

% =============================================================================
\section{V91: the weighted energy is pinched on the rescaling family}
\label{sec:v91-entry}
% =============================================================================

Items (AC1) and (AC2) of Programme~\ref{prog:v91} are carried out.  At a
homogeneous singular point the Gaussian-weighted scaled energy is
bounded below by a positive constant uniformly over the rescaling family
and over all scales: it is pinched between two positive constants.  The
proof is by compactness on the family, the equicontinuity in time of the
weighted energy, and the forward uniqueness import; the lower constant is
therefore not explicit.  The weighted energy at a fixed scale is shown to
be a continuous functional on the family.  For the limit sets it is shown
that if the weighted energy of some element has a limit at scale zero or
at scale infinity, then the family contains an element whose weighted
energy is constant in the scale, for which the flux equals minus the
positive terms at every scale; this is not self-similarity, and the
firewall is recorded.  No global regularity claim is made.  The next
compulsory companion audit is V95.

% =============================================================================
\section{Continuity and the lower bound}
\label{sec:v91-lower}
% =============================================================================

\begin{lemma}[Equicontinuity in time of the weighted energy]
\label{lem:v91-equicontinuity}
Let \(T>0\) and \(G_W(s):=\int|W(s)|^2\psi(\cdot,-T)\) for
\(W\in\mathcal A_\nu(C)\), where \(\psi(y,-T)=e^{-|y|^2/(4\nu T)}\).  For
every compact interval \(I\subset(-\infty,0)\) there is \(L\), depending
only on \(I\), \(T\), \(C\), \(C_M\), \(C_\infty\), \(\nu\), such that
\(|G_W'(s)|\le L\) on \(I\) for every \(W\) in the class with the Morrey
and \(L^\infty\) bounds, in particular for every element of a rescaling
family \(\mathcal F(U)\).
\end{lemma}

\begin{proof}
As in the proof of Theorem~\ref{thm:v86-derivative} with the weight
frozen at \(T\),
\(G_W'(s)=-2\nu\int|\nabla W|^2\psi+\nu\int|W|^2\Delta\psi+\int(|W|^2+2P)\,W\cdot\nabla\psi\).
On \(I\) the first term is bounded by \(2\nu C(-s)^{-1/2}\), the second by
the Morrey bound and the Gaussian decay of \(\Delta\psi\), and the third by
the argument of Lemma~\ref{lem:v88-order} with \(T\) in place of
\(\lambda^2\), using the \(L^\infty\) bound \(|W(s)|\le C_\infty(-s)^{-1/2}\)
and \(\|P(s)\|_3\le C(-s)^{-1/2}\).
\end{proof}

\begin{proposition}[Continuity of the weighted energy on the family]
\label{prop:v91-continuity}
Let \(W_k\to W\) locally in \(\mathcal F(U)\) as in
Proposition~\ref{prop:v89-continuity}.  Then
\(\Phi_{W_k}(\lambda)\to\Phi_W(\lambda)\) for every \(\lambda>0\), uniformly
on compact sets of \(\lambda\).
\end{proposition}

\begin{proof}
Fix a compact interval of \(\lambda\), and let \(I\) be the corresponding
compact set of times \(s=-\lambda^2\).  With the weight frozen at
\(T=\lambda^2\), the functions \(G_{W_k}\) are uniformly Lipschitz on a
neighbourhood of \(I\) by Lemma~\ref{lem:v91-equicontinuity}, and
\(G_{W_k}\to G_W\) in \(L^1(I)\): on \(B_R\) by the strong
\(L^2(I;L^2(B_R))\) convergence, and on \(|y|>R\) uniformly small by the
Morrey bound against the Gaussian tail.  A uniformly Lipschitz sequence
converging in \(L^1(I)\) converges uniformly on \(I\).  Since the weight
scale \(T=\lambda^2\) varies with \(\lambda\), cover the compact interval
of \(\lambda\) by finitely many small intervals and note that
\(\psi(\cdot,-T)\) depends continuously on \(T\) in the norm that
controls \(\int|W|^2|\psi(\cdot,-T)-\psi(\cdot,-T')|\) uniformly on the
family, again by the Morrey bound.
\end{proof}

\begin{theorem}[The weighted energy is pinched]
\label{thm:v91-pinched}
Let \(x_0\) be a homogeneous singular point and \(\mathcal F(U)\) its
rescaling family.  Then, modulo (E1)--(E8), there is \(c_\Phi>0\) such
that
\begin{equation}
 \boxed{\quad
 c_\Phi\ \le\ \Phi_W(\lambda)\ \le\ C_\Phi
 \qquad\text{for every }W\in\mathcal F(U)\text{ and every }\lambda>0 .
 \quad}
 \label{eq:v91-pinched}
\end{equation}
\end{theorem}

\begin{proof}
The upper bound is Lemma~\ref{lem:v86-weight}.  By the scaling
\(\Phi_{W^\mu}(\lambda)=\Phi_W(\mu\lambda)\) and
Proposition~\ref{prop:v83-family}(i), it suffices to prove
\(\inf_{W\in\mathcal F(U)}\Phi_W(1)>0\).  If not, there are
\(W_k\in\mathcal F(U)\) with \(\Phi_{W_k}(1)\to0\); by
Proposition~\ref{prop:v83-family}(iv) a subsequence converges locally to
some \(W\in\mathcal F(U)\), and by Proposition~\ref{prop:v91-continuity}
\(\Phi_W(1)=0\), that is, \(W(-1)\equiv0\) on \(\R^3\).  By the forward
uniqueness import (E8) in the class \(C([-1,0);L^6)\) of
Corollary~\ref{cor:v56-class}, \(W\equiv0\) on \([-1,0)\), so
\(d_W(1)=0\), contradicting \(d_W(1)\ge\varepsilon'/2\) from
Proposition~\ref{prop:v83-family}(iii).
\end{proof}

\begin{firewall}[The lower constant is not explicit]
\label{fw:v91-constant}
\(c_\Phi\) is produced by compactness and depends on the family
\(\mathcal F(U)\), hence on the singular point; no bound in terms of
\(c\), \(C_I\), \(\nu\) alone is given.  An explicit bound would
require a quantitative form of forward uniqueness, namely a lower bound
on the energy at time \(-1\) in terms of the dissipation on
\(B_1\times(-1,0)\), which the local energy inequality does not provide
because of its indefinite cutoff terms.
\end{firewall}

% =============================================================================
\section{Limit sets of the weighted energy}
\label{sec:v91-limits}
% =============================================================================

\begin{proposition}[Limits at the ends of the scale produce constant elements]
\label{prop:v91-limit-sets}
Let \(W\in\mathcal F(U)\).  The sets of limit points of \(\Phi_W(\lambda)\)
as \(\lambda\to0\) and as \(\lambda\to\infty\) are nonempty compact subsets
of \([c_\Phi,C_\Phi]\).  If \(\Phi_W(\lambda)\to\ell\) as \(\lambda\to0\)
(respectively \(\lambda\to\infty\)), then every rescaling limit \(W'\) of
\(W\) along scales \(\mu_j\to0\) (respectively \(\mu_j\to\infty\)) lies in
\(\mathcal F(U)\) and has \(\Phi_{W'}\equiv\ell\), hence
\begin{equation}
 \mathcal J_{W'}(\lambda)=-4\nu\int|\nabla W'(-\lambda^2)|^2\psi-\frac{2\ell}\lambda
 \qquad\text{for every }\lambda>0 .
 \label{eq:v91-constant-element}
\end{equation}
\end{proposition}

\begin{proof}
The first statement is \eqref{eq:v91-pinched}.  For the second,
\(\Phi_{W^{\mu_j}}(\lambda)=\Phi_W(\mu_j\lambda)\to\ell\) for each fixed
\(\lambda\), and \(W^{\mu_j}\to W'\) along a subsequence with
\(W'\in\mathcal F(U)\) by Proposition~\ref{prop:v83-family}(ii),(iv); by
Proposition~\ref{prop:v91-continuity} \(\Phi_{W'}(\lambda)=\ell\) for every
\(\lambda\).  Then \(\Phi_{W'}'\equiv0\) and \eqref{eq:v86-derivative}
gives \eqref{eq:v91-constant-element}.
\end{proof}

\begin{firewall}[Constant weighted energy is not self-similarity]
\label{fw:v91-constant-not-self-similar}
An element with \(\Phi_{W'}\) constant satisfies the pointwise identity
\eqref{eq:v91-constant-element}, in which the flux equals minus the
positive terms at every scale; this is exactly the equality case of the
pointwise sign of Firewall~\ref{fw:v88-AC1} with \(\theta=1\).  It does
not force \(W'\) to be self-similar: the identity constrains one
scalar function of \(\lambda\), while self-similarity is the statement
\(W'^{\mu}=W'\) for all \(\mu\).  Conversely a self-similar element has
constant \(\Phi\).  The reduction of Firewall~\ref{fw:v83-fixed-point}
therefore stands: a monotone quantity is needed, and a quantity that is
merely constant on some element is not it.  Whether \(\Phi_W\) has a
limit at either end for some element of the family is not decided.
\end{firewall}

% =============================================================================
\section{Integrated V91 conclusion and next acquisition order}
\label{sec:v91-conclusion}
% =============================================================================

\begin{theorem}[What V91 proves]
\label{thm:v91-conclusion}
Relative to V1--V90, modulo (E1)--(E8): the equicontinuity and
continuity statements (Lemma~\ref{lem:v91-equicontinuity},
Proposition~\ref{prop:v91-continuity}), the pinching
\eqref{eq:v91-pinched} of the weighted energy on the rescaling family at
homogeneous singular points, and the constant-element statement
\eqref{eq:v91-constant-element}.  The full Navier--Stokes
stability/global-regularity theorem remains unproved.
\end{theorem}

\begin{proof}
Theorem~\ref{thm:v91-pinched} and Proposition~\ref{prop:v91-limit-sets}.
\end{proof}

\begin{programme}[V92 acquisition order]
\label{prog:v92}
\begin{enumerate}[label=\textup{(AC\arabic*)},leftmargin=4.8em]
\item Oscillation.  Using the pinching and the two-sided flux bounds,
      determine whether \(\Phi_W\) must oscillate: over a window of ratio
      \(\Theta_0\) the flux contributes at most \(-c_2\) while the
      positive terms contribute at least the dissipation share; record
      the resulting bound on the total variation of \(\Phi_W\) per window
      and whether it forces \(\Phi_W\) to have no limit at scale zero.
\item Dissipation share.  Compute the windowed dissipation term
      \(\int_\mu^{\mu\Theta}4\nu\int|\nabla W|^2\psi\) from above, using
      the class bound, to close the window budget from both sides.
\item The next compulsory companion audit is V95.
\end{enumerate}
\end{programme}

\begin{assessment}[Position after V91]
\label{ass:v91-position}
The weighted energy is pinched between two positive constants on the
rescaling family at every homogeneous singular point, and a limit at
either end of the scale would produce an element with constant weighted
energy, not a self-similar one.  No full stability or global-regularity
claim is made.
\end{assessment}

% =============================================================================
\section*{Version V91 append-only record}
\addcontentsline{toc}{section}{Version V91 append-only record}
% =============================================================================

The complete V90 mathematical source was retained before this continuation.
Its SHA--256 digest was
\begin{center}\scriptsize\ttfamily
bc95aa4e94c2666c600555b72804ca7fda7a9e1edab75104adec0f6e48647443.
\end{center}
V91 proved the pinching of the Gaussian-weighted scaled energy on the
rescaling family at homogeneous singular points and the constant-element
statement for limits at the ends of the scale.  No full regularity claim
is made.

% =============================================================================
\section{V92: the weighted energy as an integral of the excess inward flux
over larger scales}
\label{sec:v92-entry}
% =============================================================================

Items (AC1) and (AC2) of Programme~\ref{prog:v92} are carried out.  The
scaling identity of V86 is a linear first-order equation in the scale
for the weighted energy with forcing equal to the dissipation plus the
flux; since the weighted energy is bounded, it can be solved from scale
infinity downward, which gives an exact representation of the weighted
energy at any scale as the \(\lambda^2\)-weighted integral over all
larger scales of the excess inward flux, the negative of flux plus
dissipation.  With the pinching of V91 this shows that the excess inward
flux is positive on average against the weight \(\lambda^{-2}\) on every
half-line of scales, with an explicit rate, and it sharpens the averaged
inward flux: the flux integral over a scale window is at most minus
twice the lower constant per e-fold, minus the windowed dissipation.  The
windowed dissipation is bounded above by the class constant per e-fold,
which closes the window budget from both sides.  Oscillation of the
weighted energy is not forced by these bounds.  No global regularity
claim is made.  The next compulsory companion audit is V95.

% =============================================================================
\section{The representation}
\label{sec:v92-representation}
% =============================================================================

For \(W\in\mathcal A_\nu(C)\) write, with all fields at \(s=-\lambda^2\),
\begin{equation}
 D_W(\lambda):=4\nu\int|\nabla W|^2\psi\ \ge0,
 \qquad
 \mathcal X_W(\lambda):=-\mathcal J_W(\lambda)-D_W(\lambda),
 \label{eq:v92-excess}
\end{equation}
the \emph{excess inward flux}: the inward radial flux beyond what the
weighted dissipation absorbs.  The identity \eqref{eq:v86-derivative}
reads \(\Phi_W'-\frac2\lambda\Phi_W=D_W+\mathcal J_W=-\mathcal X_W\).

\begin{theorem}[Representation of the weighted energy]
\label{thm:v92-representation}
For every \(W\in\mathcal A_\nu(C)\) with the Morrey bound and every
\(\lambda>0\),
\begin{equation}
 \boxed{\quad
 \Phi_W(\lambda)=\lambda^2\int_\lambda^\infty\frac{\mathcal X_W(\lambda')}{\lambda'^2}\,\dd\lambda' ,
 \quad}
 \label{eq:v92-representation}
\end{equation}
the integral converging absolutely.  For the original solution at a
singular point \(x_0\) under the type-I rate, with \(\Phi_u\), \(D_u\),
\(\mathcal J_u\) the periodized quantities of V87 and any
\(0<\lambda<\lambda_1\le(\delta/2)^{1/2}\),
\begin{equation}
 \Phi_u(\lambda)=\lambda^2\Bigl[\frac{\Phi_u(\lambda_1)}{\lambda_1^2}
 +\int_\lambda^{\lambda_1}\frac{\mathcal X_u(\lambda')}{\lambda'^2}\,\dd\lambda'\Bigr].
 \label{eq:v92-representation-u}
\end{equation}
\end{theorem}

\begin{proof}
From \eqref{eq:v86-derivative},
\(\frac{\dd}{\dd\lambda}\bigl(\lambda^{-2}\Phi_W\bigr)=\lambda^{-2}(\Phi_W'-2\Phi_W/\lambda)=-\lambda^{-2}\mathcal X_W\).
Integrate over \([\lambda,b]\):
\(b^{-2}\Phi_W(b)-\lambda^{-2}\Phi_W(\lambda)=-\int_\lambda^b\lambda'^{-2}\mathcal X_W\).
By Lemma~\ref{lem:v86-weight} \(\Phi_W(b)\le C_\Phi\), so \(b^{-2}\Phi_W(b)\to0\)
as \(b\to\infty\); and \(|\mathcal X_W(\lambda')|\le|\mathcal J_W|+D_W\le
(C_J+4\nu C)\lambda'^{-1}\) by Lemma~\ref{lem:v88-order} and the class
bound \(\|\nabla W(-\lambda'^2)\|_2^2\le C\lambda'^{-1}\), so
\(\lambda'^{-2}\mathcal X_W\) is integrable at infinity.  Letting
\(b\to\infty\) gives \eqref{eq:v92-representation}.  The same
computation with \eqref{eq:v87-derivative-u} on \([\lambda,\lambda_1]\)
gives \eqref{eq:v92-representation-u}.
\end{proof}

\begin{corollary}[The excess inward flux is positive on average at every scale]
\label{cor:v92-excess-positive}
Let \(x_0\) be a homogeneous singular point and \(W\in\mathcal F(U)\).
Then for every \(\lambda>0\),
\begin{equation}
 \frac{c_\Phi}{\lambda^2}\ \le\ \int_\lambda^\infty\frac{\mathcal X_W(\lambda')}{\lambda'^2}\,\dd\lambda'
 \ \le\ \frac{C_\Phi}{\lambda^2},
 \label{eq:v92-excess-average}
\end{equation}
and for every \(0<a<b\),
\begin{equation}
 \int_a^b\mathcal J_W(\lambda)\,\dd\lambda\ \le\ (C_\Phi-c_\Phi)-2c_\Phi\log\frac ba-\int_a^bD_W(\lambda)\,\dd\lambda .
 \label{eq:v92-sharpened-inward}
\end{equation}
In particular, in the scale-invariant variables \(\tau=\log\lambda\),
\(j_W(\tau):=\lambda\mathcal J_W(\lambda)\), the average of \(j_W\) over any
interval of length \(\ell\) is at most \(-2c_\Phi+(C_\Phi-c_\Phi)/\ell\).
\end{corollary}

\begin{proof}
\eqref{eq:v92-excess-average} is \eqref{eq:v92-representation} with
\eqref{eq:v91-pinched}.  For \eqref{eq:v92-sharpened-inward}, integrate
\eqref{eq:v86-derivative} over \([a,b]\):
\(\int_a^b\mathcal J_W=\Phi_W(b)-\Phi_W(a)-\int_a^bD_W-\int_a^b2\Phi_W/\lambda\),
and use \(\Phi_W(b)-\Phi_W(a)\le C_\Phi-c_\Phi\) and
\(\int_a^b2\Phi_W/\lambda\ge2c_\Phi\log(b/a)\).  The last statement is
the change of variables \(\dd\lambda/\lambda=\dd\tau\) with \(D_W\ge0\).
\end{proof}

\begin{assessment}[What the representation says]
\label{ass:v92-representation}
The weighted energy at scale \(\lambda\) is not a free quantity: it is
determined by the excess inward flux at all larger scales, weighted by
\((\lambda/\lambda')^2\).  At a homogeneous point the excess inward flux
must therefore be positive on average on every half-line of scales, at
the rate \(c_\Phi\) in the \(\lambda^{-2}\)-average, and the flux itself
is inward at the rate \(2c_\Phi\) per e-fold of scale, beyond the
dissipation.  Compared with V86, whose inward rate \(c_2/\log(S/\eta)\)
came from the dissipation term alone, the rate \(2c_\Phi\) comes from the
term \(2\Phi_W/\lambda\), which is the scaling weight of the energy
itself; the dissipation adds to it.  Both rates are per e-fold and
neither is explicit in \(c\), \(C_I\), \(\nu\): \(c_2\) is explicit but
\(c_\Phi\) is not.
\end{assessment}

% =============================================================================
\section{The window budget from both sides}
\label{sec:v92-budget}
% =============================================================================

\begin{proposition}[Two-sided window budget]
\label{prop:v92-budget}
Let \(x_0\) be a homogeneous singular point, \(W\in\mathcal F(U)\),
\(\mu>0\), \(\Theta>1\).  Then
\begin{align}
 c_2\Bigl\lfloor\frac{\log\Theta}{\log(S/\eta)}\Bigr\rfloor
 &\le\int_\mu^{\mu\Theta}D_W\le4\nu C\log\Theta ,
 \label{eq:v92-D-budget}\\
 2c_\Phi\log\Theta&\le\int_\mu^{\mu\Theta}\frac{2\Phi_W}\lambda\le2C_\Phi\log\Theta ,
 \label{eq:v92-Phi-budget}\\
 -(C_J+4\nu C)\log\Theta\ \le\ -\int_\mu^{\mu\Theta}D_W-2C_\Phi\log\Theta-(C_\Phi-c_\Phi)
 &\le\int_\mu^{\mu\Theta}\mathcal J_W\le(C_\Phi-c_\Phi)-2c_\Phi\log\Theta-\int_\mu^{\mu\Theta}D_W ,
 \label{eq:v92-J-budget}
\end{align}
and the three windowed terms sum to \(\Phi_W(\mu\Theta)-\Phi_W(\mu)\in[c_\Phi-C_\Phi,C_\Phi-c_\Phi]\).
\end{proposition}

\begin{proof}
The lower bound in \eqref{eq:v92-D-budget} is the count in the proof of
Theorem~\ref{thm:v86-inward}; the upper bound is
\(D_W(\lambda)\le4\nu\|\nabla W(-\lambda^2)\|_2^2\le4\nu C\lambda^{-1}\)
integrated.  \eqref{eq:v92-Phi-budget} is \eqref{eq:v91-pinched}
integrated.  In \eqref{eq:v92-J-budget} the upper bound is
\eqref{eq:v92-sharpened-inward}; the middle lower bound is the same
integration of \eqref{eq:v86-derivative} with
\(\Phi_W(\mu\Theta)-\Phi_W(\mu)\ge c_\Phi-C_\Phi\) and
\(\int2\Phi_W/\lambda\le2C_\Phi\log\Theta\); the leftmost bound is
Lemma~\ref{lem:v88-order} and is weaker than the middle one only when
\(C_J\) is small.
\end{proof}

\begin{firewall}[Oscillation is not forced]
\label{fw:v92-oscillation}
The budget \eqref{eq:v92-D-budget}--\eqref{eq:v92-J-budget} is consistent
with \(\Phi_W\) constant: then \(\mathcal J_W=-D_W-2\Phi_W/\lambda\)
pointwise, all three windowed terms are of order \(\log\Theta\), and
their sum is zero.  The bounds also permit oscillation of \(\Phi_W\)
between \(c_\Phi\) and \(C_\Phi\) on every window.  Hence neither a limit
of \(\Phi_W\) at scale zero nor its absence follows, and the alternative
of Proposition~\ref{prop:v91-limit-sets} remains open.  Item (AC1) is
closed with this record.
\end{firewall}

% =============================================================================
\section{Integrated V92 conclusion and next acquisition order}
\label{sec:v92-conclusion}
% =============================================================================

\begin{theorem}[What V92 proves]
\label{thm:v92-conclusion}
Relative to V1--V91, modulo (E1)--(E8): the representation
\eqref{eq:v92-representation} of the weighted energy as the
\(\lambda^{-2}\)-weighted integral of the excess inward flux over larger
scales, its version \eqref{eq:v92-representation-u} for the solution,
the positivity on average \eqref{eq:v92-excess-average}, the sharpened
inward flux \eqref{eq:v92-sharpened-inward} at rate \(2c_\Phi\) per
e-fold, and the two-sided window budget of
Proposition~\ref{prop:v92-budget}.  The full Navier--Stokes
stability/global-regularity theorem remains unproved.
\end{theorem}

\begin{proof}
Theorem~\ref{thm:v92-representation}, Corollary~\ref{cor:v92-excess-positive},
Proposition~\ref{prop:v92-budget}.
\end{proof}

\begin{programme}[V93 acquisition order]
\label{prog:v93}
\begin{enumerate}[label=\textup{(AC\arabic*)},leftmargin=4.8em]
\item Representation of the two parts.  Apply the representation to the
      energy part and the pressure part separately: determine whether
      \(\lambda^2\int_\lambda^\infty\lambda'^{-2}(-\mathcal J_E)\) or
      \(\lambda^2\int_\lambda^\infty\lambda'^{-2}(-\mathcal J_P)\) can be
      identified with a weighted energy-type quantity of its own, using
      the advective form \(\mathcal J_E=2\int\psi W\cdot\nabla|W|^2\) and
      the transport structure.
\item Scale-averaged Gaussian weights.  Determine whether the
      \(\lambda^{-2}\)-average over scales of the Gaussian weight,
      \(\int_\lambda^\infty\lambda'^{-2}\psi(\cdot,-\lambda'^2)\,\dd\lambda'\),
      gives a new weight (a power-law kernel) for which the
      representation becomes a space-time identity; compute the kernel.
\item The next compulsory companion audit is V95.
\end{enumerate}
\end{programme}

\begin{assessment}[Position after V92]
\label{ass:v92-position}
The weighted energy at every scale is an integral over larger scales of
the excess inward flux, which is positive on average at every homogeneous
singular point; the flux is inward at an explicit rate per e-fold beyond
dissipation.  No full stability or global-regularity claim is made.
\end{assessment}

% =============================================================================
\section*{Version V92 append-only record}
\addcontentsline{toc}{section}{Version V92 append-only record}
% =============================================================================

The complete V91 mathematical source was retained before this continuation.
Its SHA--256 digest was
\begin{center}\scriptsize\ttfamily
81647a2b6aa9be49938febb13587ab7e7ee208d6144c4b6a86978cefdc1848cf.
\end{center}
V92 proved the representation of the weighted energy as an integral of
the excess inward flux over larger scales, the sharpened inward flux, and
the two-sided window budget.  No full regularity claim is made.

% =============================================================================
\section{V93: correction of the stationary-subclass claim and the enstrophy
decay of the ancient class at time minus infinity}
\label{sec:v93-entry}
% =============================================================================

This version goes upstream to correct an error.  Proposition~\ref{prop:v42-stationary}
asserted that every stationary solution with finite Dirichlet integral
and \(L^6\) values belongs to the class \(\mathcal A_\nu(C)\) of
Definition~\ref{def:v41-class}; its proof used the inequality
\(C\le C(-s)^{-1/2}\) for \(s\le-1\), which is false since
\((-s)^{-1/2}<1\) for \(s<-1\).  The claim is wrong: the class
\(\mathcal A_\nu(C)\) requires \(\|\nabla W(s)\|_2^2\le C(-s)^{-1/2}\) for
all \(s<0\), which forces the enstrophy of every element to tend to zero
as \(s\to-\infty\), so the only stationary element is zero.
Consequently the statements that the type-I Liouville gate ``contains
Leray's stationary Liouville problem'' (Firewall~\ref{fw:v42-Leray-problem},
the corresponding items of Theorem~\ref{thm:v42-conclusion},
Assessment~\ref{ass:v42-position}, the sentence in
Theorem~\ref{thm:v50-ancient-limit}, and the clause in gate (G1) of
Theorem~\ref{thm:v70-gates}) are withdrawn.  What survives is stated
precisely: the sufficient closure (S1) of Theorem~\ref{thm:v57-sufficient}
is a statement about a larger class and does contain Leray's problem,
but it is sufficient, not necessary; the gate proper is a Liouville
problem on ancient solutions whose enstrophy decays like \((-s)^{-1/2}\)
at time minus infinity, a class whose self-similar members have exactly
this decay.  Programme~\ref{prog:v93} is deferred to V94.  No global
regularity claim is made.  The next compulsory companion audit is V95.

% =============================================================================
\section{Enstrophy decay at time minus infinity}
\label{sec:v93-decay}
% =============================================================================

\begin{proposition}[Every element of the class has vanishing enstrophy at minus infinity]
\label{prop:v93-decay}
Let \(W\in\mathcal A_\nu(C)\).  Then
\begin{equation}
 \|\nabla W(s)\|_{L^2(\R^3)}^2\le C(-s)^{-1/2}\longrightarrow0
 \quad(s\to-\infty),
 \qquad
 \|W(s)\|_{L^6(\R^3)}\le C_6C^{1/2}(-s)^{-1/4}\longrightarrow0 .
 \label{eq:v93-decay}
\end{equation}
Consequently:
\begin{enumerate}[label=\textup{(\roman*)},leftmargin=3.2em]
\item the only stationary element of \(\mathcal A_\nu(C)\) is \(W\equiv0\);
\item no element is time-periodic, and no nonzero element has
      \(\liminf_{s\to-\infty}\|\nabla W(s)\|_2>0\);
\item a backward self-similar element, \(W(s,y)=(-s)^{-1/2}V(y/(-s)^{1/2})\)
      with \(V\in\dot H^1(\R^3)\), satisfies
      \(\|\nabla W(s)\|_2^2=(-s)^{-1/2}\|\nabla V\|_2^2\) exactly and lies in
      \(\mathcal A_\nu(\|\nabla V\|_2^2)\).
\end{enumerate}
\end{proposition}

\begin{proof}
The first bound is the definition of the class; the \(L^6\) bound is the
Sobolev inequality on \(\R^3\).  (i) If \(W(s,y)=V(y)\) then
\(\|\nabla V\|_2^2\le C(-s)^{-1/2}\) for all \(s<0\), so \(\nabla V=0\),
and \(V\in L^6\) forces \(V=0\).  (ii) Time-periodicity with period
\(\tau\) gives \(\|\nabla W(s)\|_2=\|\nabla W(s-k\tau)\|_2\to0\) along
\(k\to\infty\), so \(\nabla W(s)=0\) for all \(s\), and \(W=0\) as in (i);
the liminf statement is \eqref{eq:v93-decay}.  (iii) The change of
variables \(y=(-s)^{1/2}z\) gives
\(\int|\nabla W(s)|^2\dd y=(-s)^{-1}(-s)^{-1}(-s)^{3/2}\int|\nabla V|^2\dd z=(-s)^{-1/2}\|\nabla V\|_2^2\).
\end{proof}

\begin{theorem}[Retraction of the stationary-subclass claim]
\label{thm:v93-retraction}
Proposition~\ref{prop:v42-stationary} is false as stated, and its proof
is invalid at the step \(C\le C(-s)^{-1/2}\) for \(s\le-1\).  The
following consequences drawn from it are withdrawn:
\begin{enumerate}[label=\textup{(\alph*)},leftmargin=3.2em]
\item Firewall~\ref{fw:v42-Leray-problem}: the type-I gate is not shown
      to contain Leray's stationary Liouville problem, and the assertion
      that the type-I regime ``is at least as hard as'' that problem is
      withdrawn;
\item the third item of Theorem~\ref{thm:v42-conclusion} and the
      sentence ``the stationary subcase is a classical open problem'' of
      Assessment~\ref{ass:v42-position};
\item the sentence of V41 preceding Theorem~\ref{thm:v41-gate} asserting
      that the class contains nonzero stationary solutions with finite
      Dirichlet integral;
\item the sentence ``its stationary case is Leray's problem for
      \(L^6\cap\dot H^1\) solutions'' in Theorem~\ref{thm:v50-ancient-limit};
\item the clause ``Contains Leray's stationary problem for
      \(L^6\cap\dot H^1\) solutions'' in gate (G1) of
      Theorem~\ref{thm:v70-gates}, and the description of the class in
      (G1) as having ``bounded enstrophy'': the class has enstrophy
      bounded by \(C(-s)^{-1/2}\), which is bounded on every half-line
      \((-\infty,-\eta]\) and tends to zero at \(-\infty\).
\end{enumerate}
No theorem of V41--V92 other than Proposition~\ref{prop:v42-stationary}
used the stationary inclusion in a proof: the inclusion appeared only in
assessments of difficulty and in the descriptions of the gate.  In
particular Theorem~\ref{thm:v41-gate}, Theorem~\ref{thm:v41-liouville},
Theorem~\ref{thm:v65-ancient-limit}, and every statement of V43--V92
about the ancient limit and its rescaling family are unaffected.
\end{theorem}

\begin{proof}
The falsity is Proposition~\ref{prop:v93-decay}(i).  For the last
assertion, the dependency check: Proposition~\ref{prop:v42-stationary}
is cited in V42 (conclusion, assessment, record), V50 (the sentence
listed in (d)), and V70 (the clause listed in (e)); each citation is
descriptive.  The sufficient closures (S1)--(S4) of
Theorem~\ref{thm:v57-sufficient} cite Proposition~\ref{prop:v53-bounded-ancient},
not the stationary inclusion.
\end{proof}

% =============================================================================
\section{The corrected description of the type-I gate}
\label{sec:v93-gate}
% =============================================================================

\begin{theorem}[Gate (G1), corrected]
\label{thm:v93-G1}
Modulo (E1)--(E8), the type-I Liouville gate is the statement:
\emph{every element \(W\) of \(\mathcal A_\nu(C_I^2\nu^{3/2})\) with the
properties of Theorem~\ref{thm:v65-ancient-limit}, hence in particular
with \(\|\nabla W(s)\|_2^2\le C(-s)^{-1/2}\to0\) as \(s\to-\infty\),
bounded on every half-line \((-\infty,-\eta]\), in
\(L^\infty_sM^{2,1}\cap L^\infty_sL^6\), with vorticity on every exterior
region at every time and a critical-Morrey trace, is zero.}
Its subcases stand as follows.
\begin{enumerate}[label=\textup{(\roman*)},leftmargin=3.2em]
\item Finite energy: proved zero (Theorem~\ref{thm:v41-liouville}).
\item Stationary and time-periodic: zero, trivially
      (Proposition~\ref{prop:v93-decay}).
\item Backward self-similar: inside the class with the exact decay
      \((-s)^{-1/2}\); excluded by external results on self-similar
      solutions with locally bounded energy (Ne\v{c}as--R\r{u}\v{z}i\v{c}ka--\v{S}ver\'{a}k
      for \(L^3\) profiles, Tsai for locally energy-bounded ones), which
      are recorded as external facts and are not imported.
\item Discretely self-similar and general elements: open within this
      series; the sufficient closures (S1)--(S4) of
      Theorem~\ref{thm:v57-sufficient} remain sufficient, and (S1) is a
      statement about the larger class of bounded ancient solutions with
      bounded enstrophy, which does contain Leray's stationary problem,
      so (S1) is stronger than the gate.
\end{enumerate}
\end{theorem}

\begin{proof}
The class description is Theorem~\ref{thm:v65-ancient-limit} with the
definition of \(\mathcal A_\nu\); (i)--(iii) are the cited results and
Proposition~\ref{prop:v93-decay}; (iv) is the record of what is proved
and what is not.
\end{proof}

\begin{assessment}[Consequence for the difficulty assessment]
\label{ass:v93-difficulty}
The series had recorded, from V42 onward, that the type-I gate is at
least as hard as Leray's stationary Liouville problem.  That was wrong,
and the corrected position is weaker in both directions: the gate is not
known to contain any classical open problem, and it is not known to be
easier than one.  The enstrophy decay at time minus infinity is a
genuine restriction that a Liouville argument may use: every element is
asymptotically small in \(\dot H^1\) and \(L^6\) at early times, and
its large-scale rescalings \(W^\mu\), \(\mu\to\infty\), which the family
\(\mathcal F(U)\) contains at homogeneous points, carry the same
restriction.  The three gates of Theorem~\ref{thm:v70-gates} otherwise
stand.  No full stability or global-regularity claim is made.
\end{assessment}

% =============================================================================
\section{Integrated V93 conclusion and next acquisition order}
\label{sec:v93-conclusion}
% =============================================================================

\begin{theorem}[What V93 establishes]
\label{thm:v93-conclusion}
Relative to V1--V92: the enstrophy decay \eqref{eq:v93-decay} of every
element of the ancient class, the retraction of
Proposition~\ref{prop:v42-stationary} and its descriptive consequences
(Theorem~\ref{thm:v93-retraction}), and the corrected gate
(Theorem~\ref{thm:v93-G1}).  The full Navier--Stokes
stability/global-regularity theorem remains unproved.
\end{theorem}

\begin{proof}
Proposition~\ref{prop:v93-decay} and the cited retractions.
\end{proof}

\begin{programme}[V94 acquisition order]
\label{prog:v94}
\begin{enumerate}[label=\textup{(AC\arabic*)},leftmargin=4.8em]
\item Carry out items (AC1)--(AC2) of Programme~\ref{prog:v93}: the
      representation of the two parts of the flux, and the
      scale-averaged Gaussian weight; record the kernel
      \(\int_\lambda^\infty\lambda'^{-2}\psi(\cdot,-\lambda'^2)\,\dd\lambda'\)
      in closed form.
\item Early-time behaviour.  Using \eqref{eq:v93-decay}, determine what
      the representation \eqref{eq:v92-representation} says as
      \(\lambda\to\infty\): the weighted dissipation \(D_W(\lambda)\) is
      \(O(\lambda^{-1})\) with the explicit constant \(4\nu C\), and the
      pinching \(\Phi_W\ge c_\Phi\) forces the flux to supply the whole
      weighted energy at large scales; record the resulting lower bound
      on the inward flux at large scales in terms of \(c_\Phi\) alone.
\item The next compulsory companion audit is V95.
\end{enumerate}
\end{programme}

\begin{assessment}[Position after V93]
\label{ass:v93-position}
An error of V42 is corrected; the type-I gate is restated on the class
with enstrophy decay at minus infinity, without the claimed inclusion of
Leray's problem.  No full stability or global-regularity claim is made.
\end{assessment}

% =============================================================================
\section*{Version V93 append-only record}
\addcontentsline{toc}{section}{Version V93 append-only record}
% =============================================================================

The complete V92 mathematical source was retained before this continuation.
Its SHA--256 digest was
\begin{center}\scriptsize\ttfamily
42964da15a062af05d11c6fc8dcf9151d33cfce4d89b6bd7fa0e35fb9acabfa8.
\end{center}
V93 corrected the stationary-subclass claim of V42, proved the enstrophy
decay of the ancient class at time minus infinity, and restated the
type-I gate.  No full regularity claim is made.

% =============================================================================
\section{V94: the space-time form of the representation and the density of
strongly inward scales}
\label{sec:v94-entry}
% =============================================================================

Items (AC1) and (AC2) of Programme~\ref{prog:v94} are carried out.  The
representation of V92 is written as a space-time integral over the
backward region \(s\le-\lambda^2\) with explicit power-Gaussian weights;
the scale-averaged Gaussian kernels are computed in closed form, and it
is recorded that they apply to the weights only, since the fields depend
on time and the class has no stationary elements.  The two parts of the
flux are not identified with energy-type quantities of their own.  From
the representation and the pinching a density statement follows: beyond
every scale, and inside every window of a fixed ratio, the scales at
which the inward flux is at least \(c_\Phi\) in scale-invariant units
have a positive weighted measure with an explicit lower bound.  The
large-scale dissipation share of the representation is scale-invariant,
so the enstrophy decay at early times of V93 gives nothing further in
the representation.  No global regularity claim is made.  The next
compulsory companion audit is V95.

% =============================================================================
\section{The space-time form and the kernels}
\label{sec:v94-spacetime}
% =============================================================================

\begin{theorem}[Space-time form of the representation]
\label{thm:v94-spacetime}
For \(W\in\mathcal A_\nu(C)\) with the Morrey bound and every \(\lambda>0\),
\begin{equation}
 \boxed{\quad
 \Phi_W(\lambda)=\frac{\lambda^2}2\int_{-\infty}^{-\lambda^2}\!\!\int_{\R^3}
 \Bigl[-\frac{(|W|^2+2P)\,(y\cdot W)}{\nu(-s)^{5/2}}-\frac{4\nu|\nabla W|^2}{(-s)^{3/2}}\Bigr]
 e^{-|y|^2/(4\nu(-s))}\,\dd y\,\dd s .
 \quad}
 \label{eq:v94-spacetime}
\end{equation}
The energy and pressure parts of the flux contribute
\(-\frac{\lambda^2}{2\nu}\iint|W|^2(y\cdot W)(-s)^{-5/2}\psi\) and
\(-\frac{\lambda^2}{\nu}\iint P\,(y\cdot W)(-s)^{-5/2}\psi\) respectively.
\end{theorem}

\begin{proof}
In \eqref{eq:v92-representation} substitute \(s=-\lambda'^2\),
\(\dd\lambda'=\dd s/(2(-s)^{1/2})\), \(\lambda'^{-2}=(-s)^{-1}\), and the
definitions \eqref{eq:v92-excess}, \eqref{eq:v88-parts} of
\(\mathcal X_W\), \(\mathcal J_W\), \(D_W\) at time \(s\).
\end{proof}

\begin{lemma}[Scale-averaged Gaussian kernels]
\label{lem:v94-kernels}
For \(r=|y|>0\), \(\lambda>0\), and \(a:=r/(2\sqrt\nu\,\lambda)\),
\begin{align}
 \int_\lambda^\infty\frac{e^{-r^2/(4\nu t^2)}}{t^2}\,\dd t
 &=\frac{\sqrt{\pi\nu}}{r}\operatorname{erf}(a),
 \label{eq:v94-K1}\\
 \int_\lambda^\infty\frac{e^{-r^2/(4\nu t^2)}}{t^4}\,\dd t
 &=\frac{8\nu^{3/2}}{r^3}\Bigl[\frac{\sqrt\pi}4\operatorname{erf}(a)-\frac a2e^{-a^2}\Bigr].
 \label{eq:v94-K2}
\end{align}
Both kernels are positive; the first is \(\lambda^{-1}(1+O(a^2))\) as
\(r\to0\) and \(\sqrt{\pi\nu}/r\) as \(r\to\infty\); the second is
\(\frac13\lambda^{-3}(1+O(a^2))\) as \(r\to0\) and \(2\sqrt\pi\,\nu^{3/2}r^{-3}\)
as \(r\to\infty\).
\end{lemma}

\begin{proof}
Substitute \(\sigma=r/(2\sqrt\nu\,t)\), \(\dd t=-r\,\dd\sigma/(2\sqrt\nu\,\sigma^2)\):
the first integral becomes \(\frac{2\sqrt\nu}r\int_0^ae^{-\sigma^2}\dd\sigma\)
and the second \(\frac{8\nu^{3/2}}{r^3}\int_0^a\sigma^2e^{-\sigma^2}\dd\sigma\),
with \(\int_0^a\sigma^2e^{-\sigma^2}=\frac{\sqrt\pi}4\operatorname{erf}(a)-\frac a2e^{-a^2}\).
The asymptotics are those of \(\operatorname{erf}\).
\end{proof}

\begin{firewall}[The kernels apply to the weights only; the two parts are not identified]
\label{fw:v94-kernels}
If the fields in \eqref{eq:v94-spacetime} were independent of \(s\), the
time integration would produce the kernels of Lemma~\ref{lem:v94-kernels}
and \eqref{eq:v94-spacetime} would become a spatial identity of
Pohozaev type; but by Proposition~\ref{prop:v93-decay} the class has no
nonzero stationary element, so no element realizes this.  For
time-dependent fields the representation is the space-time identity
\eqref{eq:v94-spacetime} and nothing simpler.  The energy part
\(-\frac{\lambda^2}{2\nu}\iint|W|^2(y\cdot W)(-s)^{-5/2}\psi\) is a cubic
space-time moment and is not the value at scale \(\lambda\) of any
quantity built from \(|W|^2\) alone; the pressure part is the same
moment of \(P\).  Item (AC1) is closed with this record.
\end{firewall}

% =============================================================================
\section{Density of strongly inward scales}
\label{sec:v94-density}
% =============================================================================

\begin{theorem}[Strongly inward scales have positive weighted measure]
\label{thm:v94-density}
Let \(x_0\) be a homogeneous singular point, \(W\in\mathcal F(U)\), and
\(c_\Phi\), \(C_\Phi\), \(C_J\) as in V88, V91.  Let
\(S_W:=\{\lambda'>0:\ -\lambda'\mathcal J_W(\lambda')\ge c_\Phi\}\) be the set
of strongly inward scales.  Then for every \(\lambda>0\),
\begin{equation}
 \lambda^2\int_{S_W\cap[\lambda,\infty)}\frac{\dd\lambda'}{\lambda'^3}\ \ge\ \frac{c_\Phi}{2C_J},
 \label{eq:v94-density-large}
\end{equation}
and for every \(a>0\) and \(\lambda\ge a\,(4C_\Phi/c_\Phi)^{1/2}\),
\begin{equation}
 a^2\int_{S_W\cap[a,\lambda]}\frac{\dd\lambda'}{\lambda'^3}\ \ge\ \frac{c_\Phi}{4C_J}.
 \label{eq:v94-density-small}
\end{equation}
In particular \(\limsup_{\lambda\to\infty}\lambda(-\mathcal J_W(\lambda))\ge2c_\Phi\)
and \(\limsup_{\lambda\to0}\lambda(-\mathcal J_W(\lambda))\ge2c_\Phi\).
\end{theorem}

\begin{proof}
From \eqref{eq:v92-representation} with \(D_W\ge0\) and \(\Phi_W\ge c_\Phi\),
\(c_\Phi\le\lambda^2\int_\lambda^\infty(-\mathcal J_W)\lambda'^{-2}\dd\lambda'\).
Split the integral over \(S_W\) and its complement: on the complement
\(-\mathcal J_W<c_\Phi/\lambda'\), on \(S_W\) \(-\mathcal J_W\le C_J/\lambda'\)
by Lemma~\ref{lem:v88-order}.  Hence
\(c_\Phi\le C_J\lambda^2\int_{S_W\cap[\lambda,\infty)}\lambda'^{-3}+c_\Phi\lambda^2\int_\lambda^\infty\lambda'^{-3}
=C_J\lambda^2\int_{S_W\cap[\lambda,\infty)}\lambda'^{-3}+\frac{c_\Phi}2\),
which is \eqref{eq:v94-density-large}.  For \eqref{eq:v94-density-small},
integrate \(\frac{\dd}{\dd\lambda}(\lambda^{-2}\Phi_W)=-\lambda^{-2}\mathcal X_W\)
over \([a,\lambda]\): \(\int_a^\lambda\mathcal X_W\lambda'^{-2}=a^{-2}\Phi_W(a)-\lambda^{-2}\Phi_W(\lambda)
\ge c_\Phi a^{-2}-C_\Phi\lambda^{-2}\ge\frac{3c_\Phi}{4a^2}\) for
\(\lambda\ge a(4C_\Phi/c_\Phi)^{1/2}\); then as before, with
\(\mathcal X_W\le-\mathcal J_W\),
\(\frac{3c_\Phi}{4a^2}\le C_J\int_{S_W\cap[a,\lambda]}\lambda'^{-3}+\frac{c_\Phi}{2a^2}\).
The limsup statements: if \(\lambda'(-\mathcal J_W)\le m\) on \([\lambda,\infty)\)
then \(c_\Phi\le\lambda^2\int_\lambda^\infty m\lambda'^{-3}=m/2\); and if
\(\lambda'(-\mathcal J_W)\le m\) on \((0,\lambda]\) then for every \(a<\lambda\)
\(c_\Phi a^{-2}-C_\Phi\lambda^{-2}\le\int_a^\lambda\mathcal X_W\lambda'^{-2}\le\frac m2(a^{-2}-\lambda^{-2})\);
multiplying by \(a^2\) and letting \(a\to0\) gives \(m\ge2c_\Phi\).
\end{proof}

\begin{assessment}[What the density says]
\label{ass:v94-density}
In the scale-invariant variable \(\tau'=\log(\lambda'/\lambda)\) the
measure \(\lambda^2\lambda'^{-3}\dd\lambda'\) is \(e^{-2\tau'}\dd\tau'\), of
total mass \(\frac12\) on \([0,\infty)\).  The theorem says that, for
every base scale, the strongly inward scales carry at least the fraction
\(c_\Phi/C_J\) of this mass, and the same inside every window of ratio
\((4C_\Phi/c_\Phi)^{1/2}\) counted from its lower end.  Together with
the two-sided budget of V92 this is the complete quantitative picture
of the flux available to the series: inward on log-average at rate
\(2c_\Phi\), strongly inward on a set of scales of positive weighted
density, bounded by \(C_J\) in scale-invariant units, and of no sign at
any single scale.
\end{assessment}

\begin{firewall}[Item (AC2): early times give nothing further]
\label{fw:v94-early}
In \eqref{eq:v92-representation} the dissipation contributes
\(\lambda^2\int_\lambda^\infty D_W\lambda'^{-2}\le\lambda^2\int_\lambda^\infty4\nu C\lambda'^{-3}=2\nu C\)
at every scale: the share is scale-invariant, and the decay of the
enstrophy at \(s\to-\infty\) (Proposition~\ref{prop:v93-decay}) is the
statement that this share is finite, not that it is small at large
scales.  Hence
\(\lambda^2\int_\lambda^\infty(-\mathcal J_W)\lambda'^{-2}\in[c_\Phi,C_\Phi+2\nu C]\)
for all \(\lambda\), and no additional large-scale information follows.
\end{firewall}

% =============================================================================
\section{Integrated V94 conclusion and next acquisition order}
\label{sec:v94-conclusion}
% =============================================================================

\begin{theorem}[What V94 proves]
\label{thm:v94-conclusion}
Relative to V1--V93, modulo (E1)--(E8): the space-time form
\eqref{eq:v94-spacetime} of the representation, the kernels
\eqref{eq:v94-K1}--\eqref{eq:v94-K2}, and the density of strongly inward
scales \eqref{eq:v94-density-large}--\eqref{eq:v94-density-small} with
the limsup bounds.  The full Navier--Stokes stability/global-regularity
theorem remains unproved.
\end{theorem}

\begin{proof}
Theorem~\ref{thm:v94-spacetime}, Lemma~\ref{lem:v94-kernels},
Theorem~\ref{thm:v94-density}.
\end{proof}

\begin{programme}[V95 acquisition order, with compulsory companion audit]
\label{prog:v95}
\begin{enumerate}[label=\textup{(AC\arabic*)},leftmargin=4.8em]
\item Perform the compulsory review of the current SM and YM notes;
      record digests; import nothing numerical.
\item Consolidate V91--V94 into one statement: the pinching, the
      representation and its space-time form, the sharpened inward flux,
      the two-sided budget, the density of strongly inward scales, and
      the corrected gate (G1) of V93.
\item Prepare the V100 restart: list the results of V1--V99 that the new
      V1 summary must carry, in the order of the two-regime
      classification, the ancient-limit theory, the lacunary and
      homogeneous shapes, and the Gaussian-weighted energy.
\item The next compulsory companion audit after V95 is V100, which is
      also the restart.
\end{enumerate}
\end{programme}

\begin{assessment}[Position after V94]
\label{ass:v94-position}
The representation is a space-time identity with explicit weights, and
the strongly inward scales have positive weighted density beyond every
scale and in every window.  No full stability or global-regularity
claim is made.
\end{assessment}

% =============================================================================
\section*{Version V94 append-only record}
\addcontentsline{toc}{section}{Version V94 append-only record}
% =============================================================================

The complete V93 mathematical source was retained before this continuation.
Its SHA--256 digest was
\begin{center}\scriptsize\ttfamily
0bb428f461a6b914692f9ab9ff9adff9cad9cbc7124bb1a7ee561399c8731902.
\end{center}
V94 wrote the representation in space-time form, computed the
scale-averaged kernels, and proved the density of strongly inward
scales.  No full regularity claim is made.

% =============================================================================
\section{V95: compulsory companion audit, the consolidated statement on the
representation of the weighted energy, and the restart manifest}
\label{sec:v95-entry}
% =============================================================================

This is a sixty-fifth-version continuation.  The compulsory review of
the two companion programmes is performed first.  V91--V94 are then
consolidated into one statement: the pinching of the weighted energy on
the rescaling family, its representation as an integral of the excess
inward flux over larger scales in scale and in space-time form, the
sharpened inward flux and the two-sided window budget, the density of
strongly inward scales, and the corrected gate (G1) of V93.  Finally the
manifest for the V100 restart is recorded: the list of results of
V1--V99 that the new first version must carry.  No global regularity
claim is made.  The next compulsory companion audit is V100, which is
also the restart.

% =============================================================================
\section{Compulsory V95 review of the current companion notes}
\label{sec:v95-companion-audit}
% =============================================================================

The current companion files in \texttt{latex\_notes} were inspected
read-only.  The frozen checkpoints are
\begin{align}
 &\texttt{Renewal\_SM\_Einstein\_Continuum\_Research\_Notes\_V80.tex},
 \notag\\[-2pt]
 &\qquad
 \texttt{e72143a6605bbf296e9236607f66e9026886004ad0b0cf0cad2ffc56fab149fb},
 \label{eq:v95-SM-checkpoint}\\
 &\texttt{Renewal\_Yang\_Mills\_Mass\_Gap\_Research\_Notes\_V59.tex},
 \notag\\[-2pt]
 &\qquad
 \texttt{9e07073e181be13c113684a84faff7bc97f6fc8ebb6657fe75aab729b854deda}.
 \label{eq:v95-YM-checkpoint}
\end{align}
At the time of inspection a file named \(\ldots\)\texttt{V58.tex} with
the same digest as \eqref{eq:v95-YM-checkpoint} was also present; the
two are byte-identical, and the V59 name is recorded.

\emph{SM V80.}  The SM programme performed its own compulsory audit of
the YM and NS notes, reading NS V87; it recorded, correctly, that the
signed statement of V86--V87 appears only after the weighted
dissipation, the scale derivative, and the nonlinear flux are kept in
one identity and integrated, that no pointwise-in-scale sign is proved,
and that no global regularity conclusion is claimed.  It then proved a
general high--high to low merger theorem giving, for every
translation-invariant quadratic source, a necessary two-frequency bound
on the cutoff constant; specialized it to the Einstein transverse
traceless and Newton symbols, recovering their fourth-order regulator
loss; showed that the displayed linearized reduced energy tends to
\(-\infty\) on an amplitude-bounded spectral sequence for every fixed
attractive gravitational coefficient and derivative-free local scalar
potential; and identified a nonlinear beat-response condition as its
next gate.

\emph{YM V58--V59.}  The YM programme derived an exact Riccati equation
and a triangular coefficient hierarchy for the Hartman--Watson clock
logarithm, identified the decaying homogeneous solution with the wrapped
sector, proved a logarithmic high-order branch law, proved that its
full-clock cumulant hierarchy is false, and showed that positive time
truncation cancels the complete branch, isolating a residual wrapped
cancellation condition as its gate.

\begin{theorem}[V95 companion-audit decision]
\label{thm:v95-companion-decision}
No SM V80 merger-symbol or cutoff constant and no YM V58--V59 Riccati
coefficient or branch rate is imported.  Neither bears on the gates of
Theorem~\ref{thm:v70-gates} as corrected by Theorem~\ref{thm:v93-G1}.
The SM V80 reading of NS V87 is accurate and requires no correction on
this side.
\end{theorem}

\begin{proof}
The SM results are bounds on quadratic sources of a linearized
gravitational model and the YM results are properties of a Bessel-type
clock transform; none is a Navier--Stokes identity.  The accuracy of the
SM reading is checked against Theorems~\ref{thm:v87-inward-u} and
Firewalls~\ref{fw:v87-pointwise}--\ref{fw:v87-lacunary}.
\end{proof}

% =============================================================================
\section{The representation of the weighted energy, consolidated}
\label{sec:v95-consolidated}
% =============================================================================

\begin{theorem}[Representation and density at homogeneous points, V91--V94]
\label{thm:v95-consolidated}
Under the type-I rate and modulo (E1)--(E8), let \(x_0\in\Sigma\) be
homogeneous, \(\mathcal F(U)\) its rescaling family, and
\(\Phi_W,D_W,\mathcal J_W,\mathcal X_W=-\mathcal J_W-D_W\) as in V86--V92.
Then:
\begin{enumerate}[label=\textup{(\arabic*)},leftmargin=3.2em]
\item \emph{Pinching.}  There is \(c_\Phi>0\), depending on the family,
      with \(c_\Phi\le\Phi_W(\lambda)\le C_\Phi\) for all
      \(W\in\mathcal F(U)\), \(\lambda>0\); \(\Phi_W(\lambda)\) is continuous
      on the family in the local topology, uniformly on compact sets of
      \(\lambda\).
\item \emph{Representation.}  \(\Phi_W(\lambda)=\lambda^2\int_\lambda^\infty\mathcal X_W(\lambda')\lambda'^{-2}\dd\lambda'\)
      for every element of the class with the Morrey bound, with the
      space-time form \eqref{eq:v94-spacetime}; for the solution,
      \eqref{eq:v92-representation-u}.
\item \emph{Sharpened inward flux and budget.}
      \(\int_a^b\mathcal J_W\le(C_\Phi-c_\Phi)-2c_\Phi\log(b/a)-\int_a^bD_W\);
      the windowed dissipation lies in
      \([c_2\lfloor\log\Theta/\log(S/\eta)\rfloor,\,4\nu C\log\Theta]\), the
      windowed scaling term in \([2c_\Phi\log\Theta,2C_\Phi\log\Theta]\),
      and the three windowed terms sum to a number of modulus at most
      \(C_\Phi-c_\Phi\).
\item \emph{Density of strongly inward scales.}  The scales with
      \(-\lambda'\mathcal J_W(\lambda')\ge c_\Phi\) carry weighted measure
      at least \(c_\Phi/(2C_J)\) beyond every scale and at least
      \(c_\Phi/(4C_J)\) in every window of ratio \((4C_\Phi/c_\Phi)^{1/2}\),
      and \(\limsup\lambda(-\mathcal J_W)\ge2c_\Phi\) at both ends.
\item \emph{Constant elements.}  If \(\Phi_W\) has a limit at scale zero
      or infinity for some element, the family contains an element with
      constant weighted energy, for which the flux equals minus the
      positive terms at every scale; this is not self-similarity.
\item \emph{The class at early times.}  Every element has
      \(\|\nabla W(s)\|_2^2\le C(-s)^{-1/2}\to0\) and
      \(\|W(s)\|_6\to0\) as \(s\to-\infty\); the class has no nonzero
      stationary or time-periodic element; the type-I gate is the
      Liouville statement of Theorem~\ref{thm:v93-G1}, and the earlier
      claim that it contains Leray's stationary problem is withdrawn.
\end{enumerate}
\end{theorem}

\begin{proof}
(1) Theorem~\ref{thm:v91-pinched}, Proposition~\ref{prop:v91-continuity}.
(2) Theorems~\ref{thm:v92-representation}, \ref{thm:v94-spacetime}.
(3) Corollary~\ref{cor:v92-excess-positive}, Proposition~\ref{prop:v92-budget}.
(4) Theorem~\ref{thm:v94-density}.  (5) Proposition~\ref{prop:v91-limit-sets}.
(6) Proposition~\ref{prop:v93-decay}, Theorems~\ref{thm:v93-retraction},
\ref{thm:v93-G1}.
\end{proof}

% =============================================================================
\section{Restart manifest for V100}
\label{sec:v95-manifest}
% =============================================================================

\begin{programme}[Manifest: what the new first version must carry]
\label{prog:v95-manifest}
The summary opening the next cycle records, in this order and with the
present labels as its references:
\begin{enumerate}[label=\textup{(M\arabic*)},leftmargin=3.8em]
\item \emph{Setting and imports.}  Smooth solutions on \(\Torus^3\) up to
      a hypothetical first singular time; the type-I enstrophy rate
      \eqref{eq:v32-type-I}; the imports (E1)--(E8) of V41, V46, V49,
      V63 in one table.
\item \emph{Ledgers and the two-regime classification} (V32--V40): the
      cutoff-free Ladyzhenskaya flux density, the fractional-integral
      record law, the \(\log^2\) bound on the critical norm under type
      I, the backward Leray cone, the scale-invariance no-go for
      monomial densities, and the classification theorem of V40.
\item \emph{The ancient limit} (V41--V50, V56--V57, V61--V65, V71--V75):
      the class \(\mathcal A_\nu(C)\), the type-I gate as a Liouville
      problem, finite-energy Liouville, dimension zero of the singular
      set, the \(L^\infty\) upgrade, the non-\(L^3\) terminal datum, the
      uniform critical Morrey bound, the trace and its shape, vorticity
      at infinity at every time, the energy dichotomy in time, strong
      continuity at a type-I time on \(\R^3\).
\item \emph{Rate-free results and the three gates} (V51--V55, V66--V70):
      the singular set as \(L^3\)-singular support, the per-scale
      implication, the geometric Leray-share dichotomy, the rate-free
      local energy bound and the \(\tau^{1/3}\) threshold, the necessity
      of a rate for \(\varepsilon\)-regularity, the three gates (G1)--(G3).
\item \emph{Lacunary and homogeneous shapes} (V76--V85): satellites at
      every intermediate scale of a lacunary point, their extent and
      rate, terminal bumps; per-scale equivalence and the count of
      homogeneous points; uniform concentration, the compact
      scale-invariant rescaling family, self-similarity reduced to a
      monotonicity formula, the non-monotone scaled local energy.
\item \emph{The Gaussian-weighted energy} (V86--V94): the exact scaling
      derivative, the averaged inward flux for the limit and the
      solution, the two parts and their bounds, negative windows and
      extremal elements, moment identities, pinching, the representation
      in scale and space-time form, the sharpened inward flux, the
      budget, the density of strongly inward scales.
\item \emph{Corrections}: V37 (type-II schedule), V59 (alignment
      corollary), V71 (energy dichotomy), V73 (strong convergence), V93
      (stationary subclass).
\item \emph{Status}: the full Navier--Stokes stability/global-regularity
      theorem is unproved; the gates (G1) as corrected, (G2), (G3) are
      the open statements; no external result beyond (E1)--(E8) is used.
\end{enumerate}
\end{programme}

% =============================================================================
\section{Integrated V95 conclusion and next acquisition order}
\label{sec:v95-conclusion}
% =============================================================================

\begin{theorem}[What V95 establishes]
\label{thm:v95-conclusion}
Relative to V1--V94: the companion-audit decision
(Theorem~\ref{thm:v95-companion-decision}), the consolidated statement
(Theorem~\ref{thm:v95-consolidated}), and the restart manifest.  The
full Navier--Stokes stability/global-regularity theorem remains
unproved.
\end{theorem}

\begin{proof}
The cited results.
\end{proof}

\begin{programme}[V96 acquisition order]
\label{prog:v96}
\begin{enumerate}[label=\textup{(AC\arabic*)},leftmargin=4.8em]
\item Pinching for the solution.  Transfer the lower bound to the
      original solution: prove \(\Phi_u(\lambda)\ge c_\Phi'\) for all
      small \(\lambda\) at a homogeneous singular point by rescaling and
      the compactness of the family, and derive the representation,
      the sharpened inward flux, and the density of strongly inward
      scales for the solution itself.
\item Rate-free identity.  Record that the scaling identity
      \eqref{eq:v87-derivative-u} and the finite-window representation
      \eqref{eq:v92-representation-u} hold for every smooth solution
      without any rate, and determine what they say at a strong type-II
      singular point where \(\Phi_u(\lambda)\) is unbounded along a
      sequence of scales: a necessary condition on the excess inward
      flux in the Gaussian ledger.
\item The next compulsory companion audit is V100, with the restart.
\end{enumerate}
\end{programme}

\begin{assessment}[Position after V95]
\label{ass:v95-position}
The series is consolidated on the representation of the weighted energy
at homogeneous singular points, the gate (G1) is corrected, and the
restart manifest is fixed.  No full stability or global-regularity claim
is made.
\end{assessment}

% =============================================================================
\section*{Version V95 append-only record}
\addcontentsline{toc}{section}{Version V95 append-only record}
% =============================================================================

The complete V94 mathematical source was retained before this continuation.
Its SHA--256 digest was
\begin{center}\scriptsize\ttfamily
34ac98a2cf3a07a2b02835b06dfecdb10cbd4d7d9a4a92410dc207286484c578.
\end{center}
V95 performed the compulsory companion audit, consolidated V91--V94, and
fixed the restart manifest for V100.  No full regularity claim is made.

% =============================================================================
\section{V96: pinching and representation for the solution itself, and the
rate-free Gaussian ledger}
\label{sec:v96-entry}
% =============================================================================

Items (AC1) and (AC2) of Programme~\ref{prog:v96} are carried out.  The
lower bound on the Gaussian-weighted scaled energy is transferred from
the rescaling family to the original solution at a homogeneous singular
point, by rescaling along any sequence of scales on which the weighted
energy would degenerate and using the compactness of the rescalings and
the forward uniqueness import; the representation, the sharpened inward
flux, and the density of strongly inward scales then hold for the
solution on small scales.  The scaling identity and the finite-window
representation are then recorded as rate-free identities valid for every
smooth solution, and their content at a point where the weighted energy
is unbounded along a sequence of scales is stated as a necessary
condition on the excess inward flux; boundedness of the weighted energy
at a point is shown to be equivalent to a critical Morrey bound at the
parabolic time.  No global regularity claim is made.  The next
compulsory companion audit is V100, with the restart.

% =============================================================================
\section{Pinching for the solution}
\label{sec:v96-pinching}
% =============================================================================

Throughout this section \(u\) satisfies the type-I rate, \(x_0\in\Sigma\),
and \(\Phi_u\), \(D_u\), \(\mathcal J_u\), \(\mathcal X_u=-\mathcal J_u-D_u\)
are the periodized quantities of V87 and V92 centred at \(x_0\).

\begin{lemma}[Order bound for the solution]
\label{lem:v96-order}
There are \(\lambda_0>0\) and \(C_J'\), depending on \(C_I\), \(\nu\),
\(C_M\), \(C_\infty\) and the periodic Riesz constant, such that
\(|\mathcal J_u(\lambda)|\le C_J'\lambda^{-1}\) and \(D_u(\lambda)\le4\nu C_I^2\nu^{3/2}\lambda^{-1}\)
for \(0<\lambda\le\lambda_0\).
\end{lemma}

\begin{proof}
As in Lemma~\ref{lem:v88-order}, with the \(L^\infty\) rate
\eqref{eq:v46-Linfty-type-I}, the Morrey bound of
Theorem~\ref{thm:v47-morrey}, the periodic Calder\'on--Zygmund bound
\(\|p(t)\|_{L^3(\Torus^3)}\le C_{\rm CZ}\|u(t)\|_{L^6}^2\le C(T_*-t)^{-1/2}\),
and the images \(|k|\ge1\) of the periodized weight contributing
\(e^{-c/\lambda^2}\).  The dissipation bound is the type-I rate at
\(t=T_*-\lambda^2\).
\end{proof}

\begin{theorem}[The weighted energy of the solution is pinched at a homogeneous point]
\label{thm:v96-pinched-u}
Let \(x_0\) be homogeneous.  Then, modulo (E1)--(E8), there are
\(\lambda_1\in(0,\lambda_0]\) and \(c_\Phi'>0\) with
\begin{equation}
 \boxed{\quad
 c_\Phi'\ \le\ \Phi_u(\lambda)\ \le\ C_\Phi'
 \qquad\text{for all }0<\lambda\le\lambda_1 .
 \quad}
 \label{eq:v96-pinched-u}
\end{equation}
\end{theorem}

\begin{proof}
The upper bound is Theorem~\ref{thm:v87-derivative-u}.  If the lower
bound failed there would be \(\lambda_k\to0\) with \(\Phi_u(\lambda_k)\to0\).
Let \(U_k(s,y):=\lambda_ku(x_0+\lambda_ky,\,T_*+\lambda_k^2s)\) on the
torus \(\Torus^3_k\) of side \(2\pi/\lambda_k\).  The weighted energy of
\(U_k\) at scale \(1\) with the \(\R^3\) Gaussian (the \(k'=0\) image of
the periodized weight) satisfies
\(\Phi_{U_k}(1)=\lambda_k^{-1}\int_{\Torus^3}|u(T_*-\lambda_k^2,x)|^2e^{-|x-x_0|^2/(4\nu\lambda_k^2)}\dd x\le\Phi_u(\lambda_k)\to0\),
because \(\Psi\) dominates its \(k'=0\) term.  By
Theorem~\ref{thm:v41-ancient-limit} a subsequence of \(U_k\) converges
locally, in \(L^2((-S,-\eta);H^1(B_R))\) with weak-\(L^3\) pressures, to an
ancient solution \(U'\); by the proof of Theorem~\ref{thm:v82-uniform},
homogeneity of \(x_0\) gives \(d_{U'}(1)\ge\varepsilon'/2\), the
dissipation converging by the strong \(H^1\) convergence.  The proof of
Proposition~\ref{prop:v91-continuity} applies verbatim to the sequence
\(U_k\), whose uniform bounds are the type-I rate, the \(L^\infty\) rate,
and the Morrey bound on \(\Torus^3_k\): hence \(\Phi_{U'}(1)=0\), so
\(U'(-1)\equiv0\), and by (E8) \(U'\equiv0\) on \([-1,0)\), contradicting
\(d_{U'}(1)\ge\varepsilon'/2\).
\end{proof}

\begin{corollary}[Representation, sharpened inward flux, and density for the solution]
\label{cor:v96-solution}
Let \(x_0\) be homogeneous, \(\lambda_1\), \(c_\Phi'\), \(C_\Phi'\), \(C_J'\)
as above.  Then for \(0<a<b\le\lambda_1\):
\begin{align}
 \frac{c_\Phi'}{a^2}-\frac{C_\Phi'}{b^2}
 &\le\int_a^b\frac{\mathcal X_u(\lambda')}{\lambda'^2}\,\dd\lambda'
 \le\frac{C_\Phi'}{a^2}-\frac{c_\Phi'}{b^2},
 \label{eq:v96-representation-window}\\
 \int_a^b\mathcal J_u(\lambda)\,\dd\lambda
 &\le(C_\Phi'-c_\Phi')-2c_\Phi'\log\frac ba-\int_a^bD_u(\lambda)\,\dd\lambda ,
 \label{eq:v96-sharpened-u}
\end{align}
and, with \(S_u:=\{\lambda'\le\lambda_1:\ -\lambda'\mathcal J_u(\lambda')\ge c_\Phi'\}\),
for every \(a>0\) and \(\lambda\in[a(4C_\Phi'/c_\Phi')^{1/2},\lambda_1]\),
\begin{equation}
 a^2\int_{S_u\cap[a,\lambda]}\frac{\dd\lambda'}{\lambda'^3}\ \ge\ \frac{c_\Phi'}{4C_J'} ,
 \qquad
 \limsup_{\lambda\to0}\lambda\bigl(-\mathcal J_u(\lambda)\bigr)\ \ge\ 2c_\Phi' .
 \label{eq:v96-density-u}
\end{equation}
\end{corollary}

\begin{proof}
\eqref{eq:v96-representation-window} is \eqref{eq:v92-representation-u}
with \eqref{eq:v96-pinched-u}; \eqref{eq:v96-sharpened-u} and
\eqref{eq:v96-density-u} are the proofs of
Corollary~\ref{cor:v92-excess-positive} and Theorem~\ref{thm:v94-density}
(small-scale case) with the primed constants and
Lemma~\ref{lem:v96-order}.
\end{proof}

\begin{assessment}[The solution near a homogeneous singular point]
\label{ass:v96-solution}
At a homogeneous singular point the actual solution has, at every small
parabolic scale, a Gaussian-weighted energy comparable to a fixed
positive constant, a radial flux that is inward at the rate \(2c_\Phi'\)
per e-fold beyond its dissipation, and strongly inward scales of positive
weighted density in every window of a computable ratio.  The constants
\(c_\Phi'\) and \(\lambda_1\) depend on the point.
\end{assessment}

% =============================================================================
\section{The rate-free Gaussian ledger}
\label{sec:v96-rate-free}
% =============================================================================

\begin{proposition}[Rate-free identities]
\label{prop:v96-rate-free}
Let \(u\) be any smooth mean-zero solution on \([0,T_*)\times\Torus^3\),
\(x_0\in\Torus^3\), and \(\Phi_u\), \(D_u\), \(\mathcal J_u\), \(\mathcal X_u\)
the periodized quantities centred at \(x_0\).  Then for
\(0<\lambda<T_*^{1/2}\) the identity \eqref{eq:v87-derivative-u} holds,
and for \(0<a<b<T_*^{1/2}\),
\begin{equation}
 \frac{\Phi_u(a)}{a^2}-\frac{\Phi_u(b)}{b^2}=\int_a^b\frac{\mathcal X_u(\lambda')}{\lambda'^2}\,\dd\lambda' .
 \label{eq:v96-ledger}
\end{equation}
Moreover \(\sup_{0<\lambda\le\lambda_1}\Phi_u(\lambda)<\infty\) if and
only if there is \(C\) with
\(\int_{B_{R\lambda}(x_0)}|u(T_*-\lambda^2)|^2\le CR\lambda\) for all
\(R\ge1\) and \(\lambda\le\lambda_1\) (a critical Morrey bound at the
parabolic time); trivially \(\Phi_u(\lambda)\le\lambda^{-1}E_*\sup\Psi\le C\lambda^{-1}E_*\).
\end{proposition}

\begin{proof}
The proof of Theorem~\ref{thm:v87-derivative-u} uses only smoothness on
\([0,T_*)\), the periodic integrations by parts, and the heat identity
of the weight; the type-I rate entered only the bound on \(\Phi_u\).
\eqref{eq:v96-ledger} is the integration of
\(\frac{\dd}{\dd\lambda}(\lambda^{-2}\Phi_u)=-\lambda^{-2}\mathcal X_u\).
For the equivalence: \(\Psi\ge e^{-R^2/(4\nu)}\) on \(B_{R\lambda}(x_0)\)
gives the Morrey bound from \(\Phi_u\le C\) with the constant
\(Ce^{R^2/(4\nu)}\lambda\), which for \(R\) fixed is of the required form,
and a Morrey bound at the parabolic time gives \(\Phi_u\le C\) by the
shell decomposition of Lemma~\ref{lem:v86-weight}.  The trivial bound is \(\Psi\le C\) uniformly for \(\lambda\le\lambda_1\) and
\(\int|u|^2\le E_*\).
\end{proof}

\begin{corollary}[Necessary condition in the Gaussian ledger without a rate]
\label{cor:v96-necessary}
Let \(u\) be any smooth solution with maximal time \(T_*<\infty\) and
\(x_0\in\Torus^3\).  If \(\Phi_u(a_k)\to\infty\) along \(a_k\to0\), then
for every fixed \(b\),
\begin{equation}
 \int_{a_k}^b\frac{\mathcal X_u(\lambda')}{\lambda'^2}\,\dd\lambda'\ \ge\ \frac{\Phi_u(a_k)}{a_k^2}-\frac{\Phi_u(b)}{b^2}\longrightarrow\infty,
 \label{eq:v96-necessary}
\end{equation}
at a rate exceeding \(a_k^{-2}\): the excess inward flux, integrated
against \(\lambda'^{-2}\), must be supercritical.  Under the type-I rate
\(\Phi_u\) is bounded, so this can occur only outside type I.
\end{corollary}

\begin{proof}
\eqref{eq:v96-ledger} with \(\Phi_u(b)\) fixed, and Theorem~\ref{thm:v87-derivative-u}.
\end{proof}

\begin{firewall}[What the rate-free ledger does not give]
\label{fw:v96-rate-free}
Without a rate, \(\mathcal J_u\) has no bound of order \(\lambda^{-1}\),
so the density and limsup statements of Corollary~\ref{cor:v96-solution}
have no rate-free version, and \eqref{eq:v96-necessary} is a necessary
condition only: it does not exclude any of the type-II schedules of V37,
V55, V70, which constrain the enstrophy ledger and say nothing about
\(\Phi_u\).  Whether \(\Phi_u\) is unbounded at a strong type-II
singular point is not decided: strong type II is a statement about
\((T_*-t)^{1/2}\|\Lambda u\|_2^2\), while \(\Phi_u\) is a statement about
the energy at the parabolic time in the parabolic Gaussian, and the two
are related only through the local energy inequality with its
indefinite terms.  Item (AC2) is closed with this record.
\end{firewall}

% =============================================================================
\section{Integrated V96 conclusion and next acquisition order}
\label{sec:v96-conclusion}
% =============================================================================

\begin{theorem}[What V96 proves]
\label{thm:v96-conclusion}
Relative to V1--V95, modulo (E1)--(E8): the pinching
\eqref{eq:v96-pinched-u} of the weighted energy of the solution at a
homogeneous singular point, the representation, sharpened inward flux and
density statements \eqref{eq:v96-representation-window}--\eqref{eq:v96-density-u}
for the solution, the rate-free ledger identity \eqref{eq:v96-ledger}
with the Morrey equivalence, and the necessary condition
\eqref{eq:v96-necessary}.  The full Navier--Stokes stability/global-regularity
theorem remains unproved.
\end{theorem}

\begin{proof}
Theorem~\ref{thm:v96-pinched-u}, Corollary~\ref{cor:v96-solution},
Proposition~\ref{prop:v96-rate-free}, Corollary~\ref{cor:v96-necessary}.
\end{proof}

\begin{programme}[V97 acquisition order]
\label{prog:v97}
\begin{enumerate}[label=\textup{(AC\arabic*)},leftmargin=4.8em]
\item Lacunary points in the Gaussian ledger.  Prove that at every
      singular point \(\limsup_{\lambda\to0}\Phi_u(\lambda)>0\), and that
      \(x_0\) is homogeneous if and only if \(\liminf_{\lambda\to0}\Phi_u(\lambda)>0\):
      along intermediate scales of a lacunary point the rescalings
      converge locally to zero (the satellites escape), so \(\Phi_u\to0\)
      along those scales by the continuity of the weighted energy.
\item Outward excess between an intermediate and an active scale.
      Derive from \eqref{eq:v96-ledger} that between an intermediate
      scale and the next larger active scale the excess flux is net
      outward in the \(\lambda^{-2}\)-average, and quantify it.
\item The next compulsory companion audit is V100, with the restart.
\end{enumerate}
\end{programme}

\begin{assessment}[Position after V96]
\label{ass:v96-position}
The Gaussian-weighted statements now hold for the solution itself at
homogeneous points, and the scaling identity is recorded as a rate-free
ledger.  No full stability or global-regularity claim is made.
\end{assessment}

% =============================================================================
\section*{Version V96 append-only record}
\addcontentsline{toc}{section}{Version V96 append-only record}
% =============================================================================

The complete V95 mathematical source was retained before this continuation.
Its SHA--256 digest was
\begin{center}\scriptsize\ttfamily
dee9cfce57d44e5103dffba507ac77e849a447c233254c4d79f23696f66f77e5.
\end{center}
V96 transferred the pinching, representation, sharpened inward flux and
density statements to the solution at homogeneous points and recorded
the rate-free Gaussian ledger.  No full regularity claim is made.

% =============================================================================
\section{V97: the Gaussian ledger of lacunary points and the minimal
log-period of the active--intermediate alternation}
\label{sec:v97-entry}
% =============================================================================

Items (AC1) and (AC2) of Programme~\ref{prog:v97} are carried out.  At
every singular point the Gaussian-weighted scaled energy stays bounded
below along the active scales, so its upper limit at scale zero is
positive; along the intermediate scales of a lacunary point it tends to
zero, because the rescalings along such scales converge locally to zero.
Hence the weighted energy of a lacunary point oscillates between values
near zero and values above a fixed constant.  Each such oscillation
forces the radial flux to deviate from the equality budget by a fixed
amount in the corresponding direction, and, since the scaling derivative
of the weighted energy is bounded by a constant over the scale, each
oscillation occupies a log-interval of scales of length bounded below:
the active--intermediate alternation of a lacunary point has a minimal
log-period.  No global regularity claim is made.  The next compulsory
companion audit is V100, with the restart.

% =============================================================================
\section{Active and intermediate scales in the Gaussian ledger}
\label{sec:v97-ledger}
% =============================================================================

Throughout, \(u\) satisfies the type-I rate, \(x_0\in\Sigma\), and the
notation of V81 is used: \(\mathcal E_{x_0}(\rho;R')=\rho^{-1}\int_{B_{R'\rho}(x_0)}|u(T_*)|^2\)
and \(\mathcal D_{x_0}(\rho;R,S,\eta)=\rho^{-1}\iint_{B_{R\rho}(x_0)\times(T_*-S\rho^2,T_*-\eta\rho^2)}|\nabla u|^2\).
A sequence \(\rho_j\downarrow0\) is \emph{active} (with parameters
\(R,S,\eta,\varepsilon\)) if \(\mathcal D_{x_0}(\rho_j;R,S,\eta)\ge\varepsilon\)
for all \(j\), and \emph{intermediate} if \eqref{eq:v77-lacunary} holds,
that is, \(\mathcal E_{x_0}(\rho_j;R')\to0\) for every \(R'\).

\begin{theorem}[The weighted energy along active and intermediate scales]
\label{thm:v97-ledger}
Modulo (E1)--(E8):
\begin{enumerate}[label=\textup{(\roman*)},leftmargin=3.2em]
\item If \(\rho_j\) is active with parameters \(R,S,\eta,\varepsilon\), then
      \(\liminf_j\Phi_u(S^{1/2}\rho_j)>0\).  Every singular point has
      active sequences, hence \(\limsup_{\lambda\to0}\Phi_u(\lambda)>0\)
      at every \(x_0\in\Sigma\).
\item If \(\rho_j\) is intermediate, then \(\Phi_u(\mu\rho_j)\to0\) for
      every fixed \(\mu>0\), uniformly for \(\mu\) in compact subsets of
      \((0,\infty)\).
\item If \(x_0\) is homogeneous then \(\liminf_{\lambda\to0}\Phi_u(\lambda)>0\)
      and \(x_0\) has no intermediate sequence.  Conversely, if
      \(\liminf_{\lambda\to0}\Phi_u(\lambda)>0\), then for every sequence
      \(\rho_j\downarrow0\) there are \(R'\) and a subsequence along which
      \(\mathcal E_{x_0}(\rho_j;R')\) stays bounded below: \(x_0\) has no
      intermediate sequence.
\end{enumerate}
\end{theorem}

\begin{proof}
(i) Suppose \(\Phi_u(S^{1/2}\rho_j)\to0\) along a subsequence.  Rescale,
\(U_j(s,y)=\rho_ju(x_0+\rho_jy,T_*+\rho_j^2s)\); by
Theorem~\ref{thm:v41-ancient-limit} a further subsequence converges
locally to \(U'\), with \(\iint_{B_R\times(-S,-\eta)}|\nabla U'|^2=\lim_j\mathcal D_{x_0}(\rho_j;R,S,\eta)\ge\varepsilon\)
by the strong \(H^1\) convergence.  As in the proof of
Theorem~\ref{thm:v96-pinched-u}, \(\Phi_{U_j}(S^{1/2})\le\Phi_u(S^{1/2}\rho_j)\to0\)
and the continuity argument of Proposition~\ref{prop:v91-continuity}
gives \(\Phi_{U'}(S^{1/2})=0\), so \(U'(-S)\equiv0\) and by (E8)
\(U'\equiv0\) on \([-S,0)\), contradicting the dissipation on
\(B_R\times(-S,-\eta)\).  Every singular point has positive upper
terminal density (Theorem~\ref{thm:v51-characterization}), hence an active
sequence by Corollary~\ref{cor:v81-quantitative}.
(ii) By Theorem~\ref{thm:v53-per-scale}, \(\mathcal D_{x_0}(\rho_j;R,S,\eta)\to0\)
for every \(R,S,\eta\).  Let \(U'\) be any local limit of a subsequence
of the rescalings \(U_j\); then \(\iint_{B_R\times(-S,-\eta)}|\nabla U'|^2=0\)
for all \(R,S,\eta\), so \(\nabla U'\equiv0\) on \(\R^3\times(-\infty,0)\),
\(U'(s)\) is constant in space and in \(L^6\), hence \(U'\equiv0\).  By
the continuity argument, \(\Phi_{U_j}(\mu)\to\Phi_{U'}(\mu)=0\) uniformly
for \(\mu\) in compact sets, and
\(\Phi_u(\mu\rho_j)=\Phi_{U_j}(\mu)+O(e^{-c/\rho_j^2}E_*/\rho_j)\to0\); since
every subsequence has a further subsequence along which this holds, the
full sequence converges.
(iii) The first statement is Theorem~\ref{thm:v96-pinched-u} and
Theorem~\ref{thm:v53-per-scale} (an intermediate sequence would have
vanishing dissipation, contradicting \(f_{x_0}(r)\ge\varepsilon'\)).  For
the converse, let \(\rho_j\downarrow0\) and \(U'\) a local limit of the
rescalings along a subsequence; \(\Phi_{U'}(1)=\lim\Phi_{U_j}(1)\ge\liminf\Phi_u(\lambda)>0\)
up to the \(O(e^{-c/\rho_j^2}E_*/\rho_j)\) error, so \(U'(-1)\not\equiv0\)
and \(U'\) is nonzero; then \(\nabla U'\not\equiv0\) (a spatially
constant \(L^6\) field is zero), so some cylinder
\(B_R\times(-S,-\eta)\) carries dissipation of \(U'\), which gives
\(\mathcal D_{x_0}(\rho_j;R,S,\eta)\ge\varepsilon\) along the subsequence,
and Theorem~\ref{thm:v53-per-scale} gives \(R'\) with
\(\liminf\mathcal E_{x_0}(\rho_j;R')>0\).
\end{proof}

\begin{firewall}[Gaussian homogeneity versus homogeneity]
\label{fw:v97-gap}
Theorem~\ref{thm:v97-ledger}(iii) shows
\(\text{homogeneous}\Rightarrow\liminf_{\lambda\to0}\Phi_u>0\Rightarrow\text{no intermediate sequence}\).
Neither converse is proved: homogeneity requires positive terminal
density at the fixed radius \(R'=1\) at every small scale, while the
absence of intermediate sequences gives a radius \(R'\) depending on
the sequence.  The three notions may or may not coincide; the series
records the implications only.
\end{firewall}

% =============================================================================
\section{Oscillation, budget deviation, and the minimal log-period}
\label{sec:v97-oscillation}
% =============================================================================

\begin{theorem}[Budget deviation across an oscillation]
\label{thm:v97-deviation}
Let \(x_0\in\Sigma\) have an intermediate sequence \(\rho_j\) and an active
sequence \(\rho_i'\) with parameters \(R,S,\eta,\varepsilon\), and let
\(c''>0\) be a lower bound for \(\Phi_u(S^{1/2}\rho_i')\) for large \(i\)
(Theorem~\ref{thm:v97-ledger}(i)).  For \(j\) large and any \(i\) with
\(\Phi_u(\rho_j)\le c''/2\), put \(a=\rho_j\), \(b=S^{1/2}\rho_i'\).
\begin{enumerate}[label=\textup{(\roman*)},leftmargin=3.2em]
\item If \(a<b\) (intermediate below active),
      \(\int_a^b\bigl(-\mathcal J_u\bigr)\le\int_a^b\Bigl(D_u+\frac{2\Phi_u}\lambda\Bigr)-\frac{c''}2\):
      the flux is less inward than the equality budget by \(c''/2\).
\item If \(b<a\) (active below intermediate),
      \(\int_b^a\bigl(-\mathcal J_u\bigr)\ge\int_b^a\Bigl(D_u+\frac{2\Phi_u}\lambda\Bigr)+\frac{c''}2\):
      the flux is more inward than the equality budget by \(c''/2\).
\end{enumerate}
\end{theorem}

\begin{proof}
Integrate \eqref{eq:v87-derivative-u}:
\(\Phi_u(b)-\Phi_u(a)=\int_a^b(D_u+2\Phi_u/\lambda+\mathcal J_u)\), and use
\(\Phi_u(b)\ge c''\), \(\Phi_u(a)\le c''/2\) in (i), and the same with the
roles exchanged in (ii).
\end{proof}

\begin{theorem}[Minimal log-period of the alternation]
\label{thm:v97-log-period}
Let \(K:=4\nu C_I^2\nu^{3/2}+2C_\Phi'+C_J'\), with \(C_\Phi'\), \(C_J'\)
from Theorem~\ref{thm:v87-derivative-u} and Lemma~\ref{lem:v96-order}.
Then \(|\Phi_u'(\lambda)|\le K/\lambda\) for \(\lambda\le\lambda_0\), and
consequently, for \(a\), \(b\) as in Theorem~\ref{thm:v97-deviation},
\begin{equation}
 \boxed{\quad
 \Bigl|\log\frac ba\Bigr|\ \ge\ \frac{c''}{2K}.
 \quad}
 \label{eq:v97-log-period}
\end{equation}
At a lacunary point with an intermediate sequence, every scale at which
\(\Phi_u\le c''/2\) is separated from every scale at which
\(\Phi_u\ge c''\) by a log-distance at least \(c''/(2K)\); the set
\(\{\lambda:\Phi_u(\lambda)\le c''/2\}\) and the set
\(\{\lambda:\Phi_u(\lambda)\ge c''\}\), both of which accumulate at
\(0\), alternate with gaps of log-length at least \(c''/(2K)\).
\end{theorem}

\begin{proof}
From \eqref{eq:v87-derivative-u}, \(0\le D_u\le4\nu C_I^2\nu^{3/2}/\lambda\),
\(0\le2\Phi_u/\lambda\le2C_\Phi'/\lambda\), \(|\mathcal J_u|\le C_J'/\lambda\),
so \(|\Phi_u'|\le K/\lambda\).  Then
\(c''/2\le|\Phi_u(b)-\Phi_u(a)|\le K|\log(b/a)|\).  Both sets accumulate
at \(0\) by Theorem~\ref{thm:v97-ledger}(i),(ii).
\end{proof}

\begin{assessment}[The lacunary shape in the Gaussian ledger]
\label{ass:v97-shape}
A lacunary singular point, seen through the Gaussian-weighted energy at
the parabolic scale, is a sequence of active scales at which the energy
near the point is comparable to a fixed constant, separated by
intermediate scales at which it is nearly absent, with the satellites of
V77 carrying the dissipation at rescaled distance tending to infinity
and therefore invisible to the centred Gaussian.  The transitions have
minimal log-length \(c''/(2K)\), and across each transition the radial
flux deviates from the equality budget by \(c''/2\) in the direction that
fills or empties the parabolic Gaussian.  The number of transitions in a
window of ratio \(\Theta\) is therefore at most \(2K\log\Theta/c''\).
No full stability or global-regularity claim is made.
\end{assessment}

% =============================================================================
\section{Integrated V97 conclusion and next acquisition order}
\label{sec:v97-conclusion}
% =============================================================================

\begin{theorem}[What V97 proves]
\label{thm:v97-conclusion}
Relative to V1--V96, modulo (E1)--(E8): the behaviour of the weighted
energy along active and intermediate scales (Theorem~\ref{thm:v97-ledger}),
the budget deviation across an oscillation
(Theorem~\ref{thm:v97-deviation}), and the minimal log-period
\eqref{eq:v97-log-period}.  The full Navier--Stokes
stability/global-regularity theorem remains unproved.
\end{theorem}

\begin{proof}
The three theorems.
\end{proof}

\begin{programme}[V98 acquisition order]
\label{prog:v98}
\begin{enumerate}[label=\textup{(AC\arabic*)},leftmargin=4.8em]
\item Satellite Gaussians.  At an intermediate scale of a lacunary
      point, centre the Gaussian weight at a satellite region of V77--V79
      instead of at \(x_0\): prove that the weighted energy centred at a
      point of maximal satellite dissipation at the parabolic scale is
      bounded below, by the moving-centre compactness of V79, so that the
      energy missing from the centred Gaussian is present in a satellite
      Gaussian at the same scale.
\item Count of satellite Gaussians.  Determine whether the number of
      disjoint satellite Gaussians at one scale with weighted energy at
      least a fixed constant is bounded, using the count \(N_*\) of
      points active at one scale.
\item The next compulsory companion audit is V100, with the restart.
\end{enumerate}
\end{programme}

\begin{assessment}[Position after V97]
\label{ass:v97-position}
The Gaussian ledger now distinguishes the active and intermediate scales
of a lacunary point, with a minimal log-period for their alternation and
a fixed budget deviation across each transition.  No full stability or
global-regularity claim is made.
\end{assessment}

% =============================================================================
\section*{Version V97 append-only record}
\addcontentsline{toc}{section}{Version V97 append-only record}
% =============================================================================

The complete V96 mathematical source was retained before this continuation.
Its SHA--256 digest was
\begin{center}\scriptsize\ttfamily
3c5458efb72a294600e2dac8dbd078a20bb59be82375603f6da2ecb3f69bf4e7.
\end{center}
V97 proved the behaviour of the weighted energy along active and
intermediate scales, the budget deviation across an oscillation, and the
minimal log-period of the alternation at lacunary points.  No full
regularity claim is made.

% =============================================================================
\section{V98: satellite Gaussians and the count of Gaussian-active centres
at one scale}
\label{sec:v98-entry}
% =============================================================================

Items (AC1) and (AC2) of Programme~\ref{prog:v98} are carried out.  The
Gaussian-weighted scaled energy centred at a moving point is bounded
below along any sequence of scales and centres carrying a fixed
dissipation share, by the moving-centre compactness of V79; at an
intermediate scale of a lacunary point with compact satellites this
places the energy missing from the centred Gaussian in a Gaussian
centred at a satellite.  Conversely, a fixed Gaussian energy at a scale
forces a fixed dissipation share at comparable scales, with constants
depending only on the energy bound; hence the number of well-separated
centres at which the Gaussian energy at one scale exceeds a fixed
constant is bounded by the type-I ledger, uniformly in the scale.  For
non-compact satellites no lower bound is available, and this is
recorded.  No global regularity claim is made.  The next compulsory
companion audit is V100, with the restart.

% =============================================================================
\section{Gaussians with moving centres}
\label{sec:v98-moving}
% =============================================================================

For \(x\in\Torus^3\) write \(\Phi_{u,x}(\lambda)\), \(D_{u,x}\),
\(\mathcal J_{u,x}\) for the periodized quantities of V87 centred at
\(x\), and \(\mathcal D_x(\rho;R,S,\eta)\) as in V97.  Throughout, the
type-I rate holds and (E1)--(E8) are assumed.

\begin{proposition}[Lower bound with moving centres]
\label{prop:v98-moving-lower}
Let \(x_j\in\Torus^3\), \(\rho_j\downarrow0\), and suppose
\(\mathcal D_{x_j}(\rho_j;R,S,\eta)\ge\varepsilon\) for all \(j\).  Then
\(\liminf_j\Phi_{u,x_j}(S^{1/2}\rho_j)>0\).
\end{proposition}

\begin{proof}
The proof of Theorem~\ref{thm:v97-ledger}(i) with the rescalings
\(U_j(s,y)=\rho_ju(x_j+\rho_jy,T_*+\rho_j^2s)\): the uniform bounds and
the compactness of Theorem~\ref{thm:v41-ancient-limit} are
centre-independent (Lemma~\ref{lem:v79-moving-centres}), the limit
\(U'\) has dissipation at least \(\varepsilon\) on \(B_R\times(-S,-\eta)\),
and \(\Phi_{U'}(S^{1/2})=0\) would force \(U'\equiv0\) on \([-S,0)\) by
(E8).
\end{proof}

\begin{corollary}[Satellite Gaussians at intermediate scales]
\label{cor:v98-satellite-gaussian}
Let \(x_0\) be lacunary with intermediate sequence \(\rho_j\) and compact
satellites in the sense of Theorem~\ref{thm:v79-bumps}: a fixed fraction
of the satellite dissipation \eqref{eq:v77-satellite} lies in \(N\) balls
\(B_{R_1\rho_j}(x_j^{(i)})\), \(i=1,\dots,N\).  Then for some \(i\) and
\(c>0\), along a subsequence,
\begin{equation}
 \Phi_{u,x_0}(\mu\rho_j)\to0\ \ (\mu\text{ fixed}),
 \qquad
 \Phi_{u,x_j^{(i)}}(S^{1/2}\rho_j)\ \ge\ c,
 \qquad
 \frac{|x_j^{(i)}-x_0|}{\rho_j}\to\infty .
 \label{eq:v98-satellite-gaussian}
\end{equation}
\end{corollary}

\begin{proof}
The first limit is Theorem~\ref{thm:v97-ledger}(ii).  One of the \(N\)
balls carries dissipation at least \((c_1/N)\rho_j-o(\rho_j)\) over
\(I_j=(T_*-S\rho_j^2,T_*-\eta\rho_j^2)\) along a subsequence, which is the
hypothesis of Proposition~\ref{prop:v98-moving-lower} with \(R=R_1\); the
distance statement is the definition of the satellite region.
\end{proof}

\begin{firewall}[Non-compact satellites]
\label{fw:v98-noncompact}
If the satellite dissipation at scale \(\rho_j\) is spread over a number
of balls of radius \(R\rho_j\) tending to infinity, each carrying
\(o(\rho_j)\), then Proposition~\ref{prop:v98-moving-lower} does not
apply, and no Gaussian at scale \(\rho_j\) is shown to carry energy
comparable to a constant.  The series does not exclude this
configuration; Proposition~\ref{prop:v78-extent-rate} bounds the extent
and the rate of the satellites but not their number.
\end{firewall}

% =============================================================================
\section{From Gaussian energy to dissipation, quantitatively}
\label{sec:v98-converse}
% =============================================================================

\begin{proposition}[Gaussian energy forces dissipation at comparable scales]
\label{prop:v98-converse}
For every \(c>0\) there are \(R,S,\eta,\varepsilon>0\) and \(\lambda_c>0\),
depending only on \(c\), \(C_I\), \(\nu\), \(E_*\) and the absolute
constants, such that for every \(x\in\Torus^3\) and every
\(0<\lambda\le\lambda_c\),
\begin{equation}
 \Phi_{u,x}(\lambda)\ge c\ \Longrightarrow\ \mathcal D_x(\lambda;R,S,\eta)\ge\varepsilon .
 \label{eq:v98-converse}
\end{equation}
\end{proposition}

\begin{proof}
If not, there are \(x_n\), \(\lambda_n\to0\) with \(\Phi_{u,x_n}(\lambda_n)\ge c\)
and \(\mathcal D_{x_n}(\lambda_n;n,n,1/n)\to0\).  Rescale at \((x_n,\lambda_n)\)
and extract a local limit \(U'\); for every fixed \(R,S,\eta\),
\(\iint_{B_R\times(-S,-\eta)}|\nabla U'|^2=\lim_n\mathcal D_{x_n}(\lambda_n;R,S,\eta)\le\lim_n\mathcal D_{x_n}(\lambda_n;n,n,1/n)=0\)
once \(n\ge\max(R,S,1/\eta)\), so \(\nabla U'\equiv0\) and \(U'\equiv0\)
as in Theorem~\ref{thm:v97-ledger}(ii).  By the continuity argument
\(\Phi_{U_n}(1)\to\Phi_{U'}(1)=0\), while
\(\Phi_{U_n}(1)=\Phi_{u,x_n}(\lambda_n)-O(e^{-c'/\lambda_n^2}E_*/\lambda_n)\ge c/2\)
for large \(n\), a contradiction.
\end{proof}

\begin{theorem}[Count of Gaussian-active centres at one scale]
\label{thm:v98-count}
Let \(c>0\) and \(R,S,\eta,\varepsilon,\lambda_c\) as in
Proposition~\ref{prop:v98-converse}.  For \(0<\lambda\le\lambda_c\), any set
of points \(x_1,\dots,x_m\in\Torus^3\), pairwise at distance at least
\(2R\lambda\), with \(\Phi_{u,x_i}(\lambda)\ge c\) for all \(i\), has
\begin{equation}
 \boxed{\quad
 m\ \le\ N_c:=\frac{2C_I^2\nu^{3/2}\,(S^{1/2}-\eta^{1/2})}{\varepsilon}.
 \quad}
 \label{eq:v98-count}
\end{equation}
\end{theorem}

\begin{proof}
The cylinders \(B_{R\lambda}(x_i)\times(T_*-S\lambda^2,T_*-\eta\lambda^2)\)
are pairwise disjoint, each carries dissipation at least
\(\varepsilon\lambda\) by \eqref{eq:v98-converse}, and the total
dissipation over the slab is at most
\(\int_{T_*-S\lambda^2}^{T_*-\eta\lambda^2}C_I^2\nu^{3/2}(T_*-t)^{-1/2}\dd t
=2C_I^2\nu^{3/2}(S^{1/2}-\eta^{1/2})\lambda\) by the type-I rate.
\end{proof}

\begin{assessment}[The Gaussian ledger at one scale]
\label{ass:v98-one-scale}
At any small parabolic scale, the centres at which the Gaussian-weighted
energy exceeds a fixed constant are few: at most \(N_c\) of them are
\(2R\lambda\)-separated, with \(N_c\) independent of the scale.  At a
homogeneous singular point the point itself is such a centre at every
scale (V96); at an intermediate scale of a lacunary point with compact
satellites the point is not, and a satellite centre is
(Corollary~\ref{cor:v98-satellite-gaussian}); the total number of
Gaussian-active centres, satellites included, is bounded by \(N_c\).
This is the Gaussian counterpart of Lemma~\ref{lem:v43-few-active}.
No full stability or global-regularity claim is made.
\end{assessment}

% =============================================================================
\section{Integrated V98 conclusion and next acquisition order}
\label{sec:v98-conclusion}
% =============================================================================

\begin{theorem}[What V98 proves]
\label{thm:v98-conclusion}
Relative to V1--V97, modulo (E1)--(E8): the moving-centre lower bound
(Proposition~\ref{prop:v98-moving-lower}), the satellite Gaussians
\eqref{eq:v98-satellite-gaussian} under compact satellites, the
quantitative converse \eqref{eq:v98-converse}, and the count
\eqref{eq:v98-count} of Gaussian-active centres at one scale.  The full
Navier--Stokes stability/global-regularity theorem remains unproved.
\end{theorem}

\begin{proof}
The cited results.
\end{proof}

\begin{programme}[V99 acquisition order]
\label{prog:v99}
\begin{enumerate}[label=\textup{(AC\arabic*)},leftmargin=4.8em]
\item Consolidate V96--V98 into one statement on the Gaussian ledger of
      the solution: pinching at homogeneous points, the rate-free
      identities, active and intermediate scales, the minimal
      log-period, satellite Gaussians, and the count at one scale.
\item Record the final position of the first cycle: what the Gaussian
      route has and has not given relative to the three gates, and the
      exact statement that would close the type-I gate through it.
\item Prepare V100: the compulsory audit, the closing theorem of the
      cycle, the renaming of the legacy file, and the new first version
      according to the manifest of Programme~\ref{prog:v95-manifest}.
\end{enumerate}
\end{programme}

\begin{assessment}[Position after V98]
\label{ass:v98-position}
Gaussian energy and dissipation activity are now equivalent at
comparable scales with moving centres, and the Gaussian-active centres
at one scale are counted.  No full stability or global-regularity claim
is made.
\end{assessment}

% =============================================================================
\section*{Version V98 append-only record}
\addcontentsline{toc}{section}{Version V98 append-only record}
% =============================================================================

The complete V97 mathematical source was retained before this continuation.
Its SHA--256 digest was
\begin{center}\scriptsize\ttfamily
effff4860f052df6ef5ef55546e20f28486314c414ea851e7148db95a99cb789.
\end{center}
V98 proved the moving-centre lower bound, the satellite Gaussians under
compact satellites, the quantitative converse from Gaussian energy to
dissipation, and the count of Gaussian-active centres at one scale.  No
full regularity claim is made.

% =============================================================================
\section{V99: the Gaussian ledger of the solution, consolidated, and the
final position of the first cycle}
\label{sec:v99-entry}
% =============================================================================

Items (AC1)--(AC3) of Programme~\ref{prog:v99} are carried out.  V96--V98
are consolidated into one statement on the Gaussian ledger of the
solution.  The final position of the first cycle is then recorded: what
the Gaussian-weighted route has given relative to the three gates, the
exact statement that would close the type-I gate through it, and the
reasons the series does not have that statement.  The instructions for
V100 are fixed.  No global regularity claim is made.  The next version
is the compulsory companion audit and the restart.

% =============================================================================
\section{The Gaussian ledger of the solution, consolidated}
\label{sec:v99-consolidated}
% =============================================================================

\begin{theorem}[Gaussian ledger of the solution, V96--V98]
\label{thm:v99-consolidated}
Let \(u\) satisfy the type-I rate, and let \(\Phi_{u,x}\), \(D_{u,x}\),
\(\mathcal J_{u,x}\), \(\mathcal X_{u,x}\) be the periodized Gaussian
quantities of V87 and V92 centred at \(x\).  Modulo (E1)--(E8):
\begin{enumerate}[label=\textup{(\arabic*)},leftmargin=3.2em]
\item \emph{Rate-free identities.}  For every smooth solution, every
      centre, and \(0<a<b<T_*^{1/2}\):
      \(\Phi_{u,x}'=D_{u,x}+2\Phi_{u,x}/\lambda+\mathcal J_{u,x}\) and
      \(a^{-2}\Phi_{u,x}(a)-b^{-2}\Phi_{u,x}(b)=\int_a^b\mathcal X_{u,x}\lambda'^{-2}\dd\lambda'\);
      \(\sup_{\lambda\le\lambda_1}\Phi_{u,x}<\infty\) if and only if a
      critical Morrey bound holds at the parabolic time; if
      \(\Phi_{u,x}(a_k)\to\infty\) the excess inward flux is
      supercritical in the sense of \eqref{eq:v96-necessary}.
\item \emph{Homogeneous points.}  At a homogeneous singular point
      \(c_\Phi'\le\Phi_{u,x_0}\le C_\Phi'\) on \((0,\lambda_1]\); the flux
      is inward at rate \(2c_\Phi'\) per e-fold beyond dissipation; the
      strongly inward scales have weighted density at least
      \(c_\Phi'/(4C_J')\) in every window of ratio \((4C_\Phi'/c_\Phi')^{1/2}\);
      \(\limsup_{\lambda\to0}\lambda(-\mathcal J_{u,x_0})\ge2c_\Phi'\).
\item \emph{Active and intermediate scales.}  At every singular point
      \(\Phi_{u,x_0}\) is bounded below along active scales, so
      \(\limsup_{\lambda\to0}\Phi_{u,x_0}>0\); along intermediate scales
      \(\Phi_{u,x_0}\to0\); homogeneous \(\Rightarrow\)
      \(\liminf_{\lambda\to0}\Phi_{u,x_0}>0\) \(\Rightarrow\) no intermediate
      sequence.
\item \emph{Minimal log-period and budget deviation.}
      \(|\Phi_{u,x}'|\le K/\lambda\); at a lacunary point the scales with
      \(\Phi\le c''/2\) and those with \(\Phi\ge c''\) alternate with
      log-gaps at least \(c''/(2K)\), and across each transition the flux
      deviates from the equality budget by \(c''/2\).
\item \emph{Moving centres and the count.}  A fixed dissipation share at
      \((x_j,\rho_j)\) forces \(\liminf\Phi_{u,x_j}(S^{1/2}\rho_j)>0\); a
      Gaussian energy at least \(c\) at scale \(\lambda\) forces a
      dissipation share \(\varepsilon(c)\) at comparable scales; at most
      \(N_c\) well-separated centres have Gaussian energy at least \(c\)
      at any one scale; under compact satellites, an intermediate scale
      of a lacunary point has a satellite Gaussian carrying a fixed
      energy.
\end{enumerate}
\end{theorem}

\begin{proof}
(1) Proposition~\ref{prop:v96-rate-free}, Corollary~\ref{cor:v96-necessary}.
(2) Theorem~\ref{thm:v96-pinched-u}, Corollary~\ref{cor:v96-solution}.
(3) Theorem~\ref{thm:v97-ledger}.  (4) Theorems~\ref{thm:v97-deviation},
\ref{thm:v97-log-period}.  (5) Propositions~\ref{prop:v98-moving-lower},
\ref{prop:v98-converse}, Theorem~\ref{thm:v98-count},
Corollary~\ref{cor:v98-satellite-gaussian}.
\end{proof}

% =============================================================================
\section{Final position of the first cycle}
\label{sec:v99-position}
% =============================================================================

\begin{theorem}[What would close the type-I gate through the Gaussian route]
\label{thm:v99-closing-statement}
Modulo (E1)--(E8), each of the following statements implies that no
smooth solution on \(\Torus^3\) has a singularity with a type-I
enstrophy rate at a homogeneous singular point; together with an
exclusion of lacunary points, each would close gate (G1) of
Theorem~\ref{thm:v93-G1}.
\begin{enumerate}[label=\textup{(C\arabic*)},leftmargin=3.8em]
\item \emph{Pointwise sign.}  For every element \(W\) of a rescaling
      family \(\mathcal F(U)\) at a homogeneous point and every \(\lambda\),
      \(\mathcal J_W(\lambda)\ge-\theta\bigl(D_W(\lambda)+2\Phi_W(\lambda)/\lambda\bigr)\)
      with a fixed \(\theta<1\); or the weaker: \(\Phi_W\) is monotone in
      \(\lambda\) on \(\mathcal F(U)\).
\item \emph{Self-similar Liouville.}  Every backward self-similar or
      discretely self-similar element of \(\mathcal A_\nu(C)\) with the
      properties of Theorem~\ref{thm:v65-ancient-limit} is zero, together
      with the existence of such an element in \(\mathcal F(U)\).
\item \emph{Explicit pinching constant with a contradiction.}  A lower
      bound \(c_\Phi\ge\underline c(C_I,\nu)\) together with an upper
      bound on the windowed inward flux strictly below \(2c_\Phi\log\Theta\)
      for some window, contradicting \eqref{eq:v92-sharpened-inward}.
\end{enumerate}
\end{theorem}

\begin{proof}
(C1) gives monotonicity of \(\Phi_W\), hence by Firewall~\ref{fw:v83-fixed-point}
and Proposition~\ref{prop:v91-limit-sets} a self-similar element of the
family; (C2) then excludes it, contradicting
Proposition~\ref{prop:v83-family}(iv) (the family is nonempty).  (C3) is
a direct contradiction with Corollary~\ref{cor:v92-excess-positive}.
Lacunary points are not reached by any of the three, since their
rescaling limits along intermediate scales vanish
(Theorem~\ref{thm:v97-ledger}(ii)) and the Gaussian route sees only the
satellites, whose compactness is not proved
(Firewall~\ref{fw:v98-noncompact}).
\end{proof}

\begin{assessment}[Final position of the first cycle]
\label{ass:v99-final}
After ninety-nine versions the series has: a two-regime classification
of a hypothetical singularity by the enstrophy ledger; a complete
reduction of the type-I regime to a Liouville problem on the class
\(\mathcal A_\nu(C)\), with the finite-energy, stationary and
time-periodic cases settled and the stationary inclusion of V42
corrected in V93; the structure of the singular set, the terminal datum,
and the ancient limit and its trace under type I; rate-free local
energy bounds and the necessity of a rate for \(\varepsilon\)-regularity;
the lacunary and homogeneous shapes of singular points, satellites, the
count of homogeneous points, and the compact scale-invariant rescaling
family; and, on the Gaussian-weighted scaled energy, the exact scaling
identity, the averaged inward flux, the pinching, the representation as
an integral of the excess inward flux, the density of strongly inward
scales, the ledger of active and intermediate scales with its minimal
log-period, and the count of Gaussian-active centres.  It does not have
a monotone quantity along the scaling, a Liouville theorem for the
class, a bound on the number of satellites, or any rate-free exclusion
of the intermediate and strong type-II classes.  The three gates of
Theorem~\ref{thm:v70-gates}, with (G1) as corrected in
Theorem~\ref{thm:v93-G1}, are the open statements, and
Theorem~\ref{thm:v99-closing-statement} names what the Gaussian route
would need.  The full Navier--Stokes stability/global-regularity theorem
is not proved.
\end{assessment}

% =============================================================================
\section{Instructions for V100}
\label{sec:v99-instructions}
% =============================================================================

\begin{programme}[V100 acquisition order, with compulsory companion audit and restart]
\label{prog:v100}
\begin{enumerate}[label=\textup{(AC\arabic*)},leftmargin=4.8em]
\item Perform the compulsory review of the current SM and YM notes;
      record digests; import nothing numerical.
\item State the closing theorem of the first cycle: the complete list of
      proved statements by block, with the imports, corrections, and
      gates, in the form the manifest of Programme~\ref{prog:v95-manifest}
      prescribes.
\item Close the file as the legacy record of the first cycle, to be
      renamed with the suffix \texttt{\_f1}, and open the second cycle
      with a new first version containing the context, the summary of
      the major results with their statements, and the first new
      acquisition of the second cycle.
\end{enumerate}
\end{programme}

% =============================================================================
\section{Integrated V99 conclusion}
\label{sec:v99-conclusion}
% =============================================================================

\begin{theorem}[What V99 establishes]
\label{thm:v99-conclusion}
Relative to V1--V98: the consolidated Gaussian ledger of the solution
(Theorem~\ref{thm:v99-consolidated}) and the closing statements
(Theorem~\ref{thm:v99-closing-statement}).  The full Navier--Stokes
stability/global-regularity theorem remains unproved.
\end{theorem}

\begin{proof}
The cited results.
\end{proof}

\begin{assessment}[Position after V99]
\label{ass:v99-position}
The first cycle is consolidated and its final position recorded.  No
full stability or global-regularity claim is made.
\end{assessment}

% =============================================================================
\section*{Version V99 append-only record}
\addcontentsline{toc}{section}{Version V99 append-only record}
% =============================================================================

The complete V98 mathematical source was retained before this continuation.
Its SHA--256 digest was
\begin{center}\scriptsize\ttfamily
4e5fafb142a09a51b610ebb45570cef4faf2d3a733ff7e6f1cc85379222b13ee.
\end{center}
V99 consolidated V96--V98 and recorded the final position of the first
cycle with the instructions for V100.  No full regularity claim is made.

% =============================================================================
\section{V100: compulsory companion audit, the closing theorem of the first
cycle, and the restart}
\label{sec:v100-entry}
% =============================================================================

This is the hundredth version and the last of the first cycle.  The
compulsory review of the two companion programmes is performed first.
The closing theorem of the cycle then lists, block by block, what
V1--V99 proved, with the imports, the corrections, and the open gates.
The file is then closed as the legacy record of the first cycle, to be
renamed with the suffix \texttt{\_f1}; the second cycle opens in a new
first version whose opening summary follows the manifest of
Programme~\ref{prog:v95-manifest}.  No global regularity claim is made.

% =============================================================================
\section{Compulsory V100 review of the current companion notes}
\label{sec:v100-companion-audit}
% =============================================================================

The current companion files in \texttt{latex\_notes} were inspected
read-only.  The frozen checkpoints are
\begin{align}
 &\texttt{Renewal\_SM\_Einstein\_Continuum\_Research\_Notes\_V82.tex},
 \notag\\[-2pt]
 &\qquad
 \texttt{e5a5c55215dc56e177de8a9a8ccac03608573764cde119d682008f469f54f054},
 \label{eq:v100-SM-checkpoint}\\
 &\texttt{Renewal\_Yang\_Mills\_Mass\_Gap\_Research\_Notes\_V59.tex},
 \notag\\[-2pt]
 &\qquad
 \texttt{90b5fb5e59ba9440df71f08d599a01ff4cf33fcab9b7ea1f8d79a2c14089699a}.
 \label{eq:v100-YM-checkpoint}
\end{align}
The YM V59 file differs from the byte-identical V58/V59 pair recorded
in \eqref{eq:v95-YM-checkpoint}: it now contains the V59 chapter.

\emph{SM V81--V82.}  The matched beat packet of SM V80 was inserted into
the minimally coupled constant-mean-curvature Lichnerowicz equation; its
gradient density has an exact limit, the conformal solutions and volumes
are uniform in frequency and converge subsequentially to a nonlinear
homogenized equation, and the two-index ancestry lattice is the exact
source algebra (V81).  On the closed torus the unreduced ADM Hamiltonian
vanishes on the constraint surface; with York time as clock the physical
volume is the reduced Hamiltonian, the matched beat increases it at first
order at fixed York time, and at fixed volume the required response is
computed explicitly (V82).

\emph{YM V59.}  The YM programme derived the exact Stieltjes measure of
its Wilson clock exponent, proved exponential phase suppression below a
fixed fraction of the turning scale, decomposed the clock into a
small-jump part and an exponentially rare compound-Poisson wrapped part,
and proved its cumulant hierarchy for the small clock, closing its
tangent Gaussian all-arity card while recording that the full mass gap
is not proved.

\begin{theorem}[V100 companion-audit decision]
\label{thm:v100-companion-decision}
No SM V81--V82 conformal, volume, or York-response constant and no YM
V59 Stieltjes, phase-suppression, or cumulant estimate is imported.
Neither bears on the gates of Theorem~\ref{thm:v70-gates} as corrected
by Theorem~\ref{thm:v93-G1}, nor on the closing statements of
Theorem~\ref{thm:v99-closing-statement}.
\end{theorem}

\begin{proof}
The SM results concern conformal factors and Hamiltonian reduction of a
gravitational model on a torus; the YM results concern a Bernstein
exponent of a lattice clock.  None is a Navier--Stokes identity, weight,
or flux.
\end{proof}

% =============================================================================
\section{Closing theorem of the first cycle}
\label{sec:v100-closing}
% =============================================================================

\begin{theorem}[What the first cycle proved]
\label{thm:v100-closing}
Let \(u\) be a smooth mean-zero solution of the Navier--Stokes equations
on \([0,T_*)\times\Torus^3\) with maximal time \(T_*<\infty\).  Modulo the
imports (E1)--(E8) of Theorem~\ref{thm:v65-import-ledger}, the first
cycle proved the following, and nothing beyond it.
\begin{enumerate}[label=\textup{(B\arabic*)},leftmargin=3.8em]
\item \emph{Ledgers and classification} (V32--V40).  The cutoff-free
      Ladyzhenskaya flux density and the recent-ledger tail bound; the
      fractional-integral record law and, under the type-I rate
      \eqref{eq:v32-type-I}, the \(\log^2\) bound on the critical norm;
      the backward Leray cone and the corrected type-II schedule (V37);
      the scale-invariance no-go for monomial flux densities (V39); the
      two-regime classification by the \(L^{5/2}_tH^1\) ledger (V40).
\item \emph{The ancient limit under type I} (V41--V50, V56--V57,
      V61--V65, V71--V75).  Rescaling at a concentration point yields a
      nonzero element of \(\mathcal A_\nu(C_I^2\nu^{3/2})\); type-I
      exclusion is equivalent to a Liouville statement on this class
      (Theorem~\ref{thm:v41-gate}); finite-energy elements vanish; the
      singular set has dimension zero; the enstrophy rate upgrades to
      \(L^\infty\) and gradient rates; the terminal datum is not in
      \(L^3\) near any singular point and satisfies a critical Morrey
      bound; the ancient limit is in \(L^\infty_sM^{2,1}\), has a trace
      with critical energy at infinity, vorticity on every exterior
      region at every time, an energy dichotomy in time, and strong
      \(L^2\) continuity at a type-I time on \(\R^3\).
\item \emph{Rate-free results and the gates} (V51--V55, V66--V70).  The
      singular set equals the positive-upper-density set and the
      \(L^3\)-singular support of the terminal datum; per-scale
      dissipation activity forces terminal density; the geometric
      Leray-share dichotomy characterizes type I along a sequence; the
      rate-free local energy bound \eqref{eq:v66-local-energy} and the
      \(\tau^{1/3}\) threshold; \(\varepsilon\)-regularity at \(T_*\)
      requires a rate; the three gates (G1)--(G3).
\item \emph{Shapes} (V76--V85).  Lacunary points have satellites at every
      intermediate scale with the scale's extent and rate; homogeneous
      points are countable, finitely many per constant, with uniformly
      concentrated ancient limits whose rescaling family is compact and
      scale-invariant; self-similarity reduces to a monotonicity
      formula; the scaled local energy is not monotone.
\item \emph{The Gaussian-weighted energy} (V86--V99).  The exact scaling
      identity \eqref{eq:v86-derivative} and its torus form
      \eqref{eq:v87-derivative-u}; the averaged inward flux
      \eqref{eq:v86-inward}, \eqref{eq:v87-inward-u}; the two parts of the
      flux, their advective and potential forms, order bounds, oddness
      and parity, exact zero-flux flows; negative windows, continuity
      and extremal elements, moment identities; pinching
      \eqref{eq:v91-pinched}, \eqref{eq:v96-pinched-u}; the representation
      \eqref{eq:v92-representation} in scale and space-time form; the
      sharpened inward flux and the budget; the density of strongly
      inward scales; the rate-free ledger; active and intermediate
      scales, minimal log-period, budget deviation; moving centres,
      satellite Gaussians, and the count \eqref{eq:v98-count}.
\item \emph{Corrections.}  V37 (V35 schedule), V59 (V58 corollary), V71
      (\(s_*\)), V73 (V72 convergence), V93 (V42 stationary subclass,
      with the corrected gate Theorem~\ref{thm:v93-G1}).
\item \emph{Open.}  Gates (G1) as corrected, (G2), (G3); the closing
      statements (C1)--(C3) of Theorem~\ref{thm:v99-closing-statement}.
      The full Navier--Stokes stability/global-regularity theorem is not
      proved.
\end{enumerate}
\end{theorem}

\begin{proof}
Each block is the conjunction of the integrated conclusions of its
versions, Theorems~\ref{thm:v40-conclusion}--\ref{thm:v99-conclusion},
and the consolidations of V50, V65, V70, V85, V90, V95, V99.
\end{proof}

\begin{assessment}[Closing assessment of the first cycle]
\label{ass:v100-closing}
The cycle moved the problem from a stopping-time bottleneck in the
enstrophy ledger to a Liouville problem on a precisely described class
of ancient solutions, described the singular set and the terminal datum
under type I, and built a Gaussian-weighted ledger whose signed
statements are the first of the series.  The gap between what is proved
and the theorem sought is the absence of a monotone quantity along the
scaling and of a Liouville theorem for the class, together with the
rate-free classes (G2), (G3), for which the tools of the cycle give
necessary conditions only.  The second cycle starts from this position.
\end{assessment}

% =============================================================================
\section{Restart record}
\label{sec:v100-restart}
% =============================================================================

\begin{programme}[Second cycle]
\label{prog:v100-restart}
The present file is the complete record of the first cycle, V1--V100,
and is retained under the name with suffix \texttt{\_f1}.  The second
cycle opens in \texttt{Renewal\_NS\_Stability\_Research\_Notes\_V1.tex}
with: the context of the programme; the summary of the major results of
the first cycle in the order (M1)--(M8) of Programme~\ref{prog:v95-manifest},
with statements and the legacy labels named in words; the digest of
this file; and the first new acquisition of the second cycle, the
scale-invariant form of the Gaussian ledger and its log-average law.
The append-only, single-file, five-version-audit, and hundred-version
restart rules continue unchanged.
\end{programme}

% =============================================================================
\section*{Version V100 append-only record}
\addcontentsline{toc}{section}{Version V100 append-only record}
% =============================================================================

The complete V99 mathematical source was retained before this continuation.
Its SHA--256 digest was
\begin{center}\scriptsize\ttfamily
52a5c703aaa7d351a46f90b2ae237a0e2058bd9b199a9d777f651823149c341d.
\end{center}
V100 performed the compulsory companion audit, stated the closing
theorem of the first cycle, and recorded the restart.  No full
regularity claim is made.

\end{document}
